\documentclass[10pt,oneside,openany]{book}
\usepackage[T1]{fontenc}
\usepackage[utf8]{inputenc}
\usepackage{babel}
\babelprovide[import,main]{spanish}
\usepackage{lmodern,amsmath,amssymb,amsthm,mathrsfs}
\usepackage{iftex}
\ifLuaTeX\usepackage{luatex85}\fi
\usepackage[all]{xy}
\usepackage{tikz-cd}
\usepackage{xstring,xparse,xcolor,etoolbox}
\usepackage[a4paper,left=18mm,right=18mm,top=22mm,bottom=23mm,marginparwidth=14mm,marginparsep=2mm]{geometry}
\usepackage{array,longtable,multicol,booktabs}
\usepackage[shortlabels]{enumitem}
\usepackage{marginnote,microtype,graphicx,adjustbox}
\usepackage[unicode=true,pdfencoding=auto,colorlinks=true,linkcolor=blue,citecolor=blue,urlcolor=blue,hypertexnames=false,hyperfootnotes=false,linktocpage=true]{hyperref}
\usepackage{bookmark}
\renewcommand{\thesection}{\arabic{section}}
\renewcommand{\thesubsection}{\thesection.\arabic{subsection}}
% Keep PDF destinations unique when the cumulative reader resets table counters.
\renewcommand{\theHtable}{\arabic{part}.\arabic{table}}
\setcounter{tocdepth}{2}
\setlength{\emergencystretch}{12em}
\tolerance=3000
\binoppenalty=200
\relpenalty=200
\hfuzz=0.1pt
\vfuzz=0.1pt
\raggedbottom
\allowdisplaybreaks[3]
\newtheoremstyle{egaenv}{}{}{\normalfont}{}{\bfseries}{.}{ }{\thmnote{(#3)}}
\newtheoremstyle{egastatement}{}{}{\itshape}{}{\bfseries}{. ---}{ }{\thmname{#1}\thmnote{ (#3)}}
\newtheoremstyle{egadescription}{}{}{\normalfont}{}{\bfseries}{. ---}{ }{\thmname{#1}\thmnote{ (#3)}}
\theoremstyle{egaenv}\newtheorem*{env}{---}
\theoremstyle{egastatement}
\newtheorem*{theorem}{Teorema}
\newtheorem*{proposition}{Proposición}
\newtheorem*{lemma}{Lema}
\newtheorem*{corollary}{Corolario}
\theoremstyle{egadescription}
\newtheorem*{definition}{Definición}
\newtheorem*{remark}{Observación}
\newtheorem*{remarks}{Observaciones}
\newtheorem*{example}{Ejemplo}
\newtheorem*{examples}{Ejemplos}
\newtheorem*{scholium}{Escolio}
\newtheorem*{notation}{Notación}
\newenvironment{namedtheorem}[2][]{\par\medskip\noindent\setmanualcurrentlabel{#1}\textbf{#2\ifstrempty{#1}{}{ (#1)}. ---}\ifstrequal{#2}{Observación}{\normalfont}{\ifstrequal{#2}{Observaciones}{\normalfont}{\itshape}}}{\par\medskip}
\newenvironment{namedlemma}[2][]{\par\medskip\noindent\setmanualcurrentlabel{#1}\textbf{#2\ifstrempty{#1}{}{ (#1)}. ---}\itshape}{\par\medskip}
\renewcommand{\proofname}{Demostración}
\renewcommand{\contentsname}{Índice general}
\renewcommand{\bibname}{Bibliografía}
\providecommand{\refname}{Bibliografía}
\newcommand{\oldpage}[2][]{\marginnote{\normalfont\scriptsize\textbf{#1}\,|\,#2}\ignorespaces}
\newcommand{\sectionbreak}{\begin{center}***\end{center}}
\newcommand{\erratum}[2][]{#2\marginnote{\scriptsize\textbf{Err}\textsubscript{#1}}}
\providecommand{\shHom}{\mathcal{H}\!\mathit{om}}
\providecommand{\shProj}{\mathcal{P}\!\mathit{roj}}
% Literal macro definition from the French EGA II notation index.
\providecommand{\egaShProj}{\mathcal{P}\!\mathit{roj}}
\providecommand{\shExt}{\mathcal{E}\!\mathit{xt}}
\providecommand{\shTor}{\mathcal{T}\!\mathit{or}}
\providecommand{\shGr}{\mathcal{G}\!\mathit{r}}
\providecommand{\shDer}{\mathcal{D}\!\mathit{er}}
\providecommand{\shDiff}{\mathcal{D}\!\mathit{iff}}
\providecommand{\shHomcont}{\mathcal{H}\!\mathit{om.cont}}
\providecommand{\shAut}{\mathcal{A}\!\mathit{ut}}
\providecommand{\Frac}{\operatorname{Frac}}
\providecommand{\unsure}[1]{#1}
\providecommand{\completelyunsure}{\textbf{??}}
\providecommand{\og}{\guillemotleft{}}
\providecommand{\fg}{\guillemotright{}}
\makeatletter
\newcommand{\setmanualcurrentlabel}[1]{\def\@currentlabel{#1}}
\newcommand{\EGARefNumber}[1]{\begingroup\StrSubstitute{#1}{-fr}{}[\egarefkey]\StrSubstitute{\egarefkey}{-es}{}[\egarefkey]\StrSubstitute{\egarefkey}{.fr}{}[\egarefkey]\StrBehind{\egarefkey}{.}\endgroup}
\newcommand{\EGAMaybeRefLink}[2]{\@ifundefined{r@#1}{#2}{\hyperref[#1]{#2}}}
\newcommand{\EGAReferenceAlias}[2]{\ifcsdef{r@#2}{\csletcs{r@#1}{r@#2}}{}}
\NewDocumentCommand{\sref}{o m o}{\textup{(\IfNoValueF{#1}{\textbf{#1},~}\EGAMaybeRefLink{#2}{\EGARefNumber{#2}}\IfNoValueF{#3}{,~#3})}}
\newcommand{\eref}[1]{\textup{(\EGAMaybeRefLink{#1}{\EGARefNumber{#1}})}}
\newcommand{\EGAReaderTarget}[2]{\phantomsection\def\@currentlabel{#2}\label{#1}}
\newcommand{\EGAReadLink}[2]{\EGAMaybeRefLink{#1}{#2}}
\newcommand{\EGAUnavailableHyperref}[2]{\EGAMaybeRefLink{#1}{#2}}
\newcommand{\sgaStableTarget}[1]{\hypertarget{#1}{}}
\newcommand{\sgaStableOrOriginalLink}[2]{\EGAMaybeRefLink{#1}{#2}}
\newcommand{\EGAIVFourExternalLabel}[1]{}
\newcommand{\EGAIVFourExternalCitation}[2]{\expandafter\def\csname b@#1\endcsname{#2}}
\makeatother
\newcommand{\EGAResetVolumeCounters}{\setcounter{section}{0}\setcounter{subsection}{0}\setcounter{equation}{0}\setcounter{footnote}{0}\setcounter{figure}{0}\setcounter{table}{0}}
\newcommand{\EGAVolumeTitlePage}[3]{\clearpage\EGAResetVolumeCounters\chapter*{#2\\#3}\addcontentsline{toc}{chapter}{#2: #3}}
\newcommand{\MR}[1]{MR #1}
\newcommand{\MRhref}[2]{\href{#1}{#2}}
\providecommand{\bysame}{\leavevmode\hbox to3em{\hrulefill}\,}
\pagestyle{plain}

% BEGIN authoritative French CC0 preamble-base.tex
% Notation of the EGA editions: the abbreviations the transcription uses for alphabets, symbols, operators and
% arrows. Written for these editions and dedicated to the public domain under CC0 1.0.
\makeatletter
% a command that may or may not exist already is (re)defined in one step
\newcommand{\ega@define}[2]{\providecommand{#1}{}\renewcommand{#1}{#2}}
\newcommand{\ega@operator}[2]{\providecommand{#1}{}\renewcommand{#1}{\operatorname{#2}}}

% alphabets: sheaves, bold letters, categories
\ega@define{\sh}{\mathcal}
\ega@define{\bb}{\mathbf}
\ega@define{\cat}{\mathtt}
\ega@define{\C}{\cat{C}}
\ega@define{\Set}{\cat{Set}}
\ega@define{\CHom}{\cat{Hom}}
\ega@define{\op}{^{\cat{op}}}
\ega@define{\GL}{\bb{GL}}
\ega@define{\bbGamma}{\boldsymbol{\Gamma}}

% letters and symbols
\renewcommand{\leq}{\leqslant}
\renewcommand{\geq}{\geqslant}
\ega@define{\setmin}{-}
\ega@define{\emp}{\varnothing}
\ega@define{\vphi}{\varphi}
\ega@define{\kres}{k}
\ega@define{\rad}{\mathfrak{r}}
\ega@define{\nilrad}{\mathfrak{N}}
\ega@define{\HH}{\mathrm{H}}
\ega@define{\CHH}{\check{\mathrm{H}}}
\ega@define{\RR}{\mathrm{R}}
\ega@define{\LL}{\mathrm{L}}
\ega@define{\red}{\mathrm{red}}
\providecommand{\dual}{}\renewcommand{\dual}[1]{{#1}^{\vee}}
% the tilde of the sheaf associated with a module, written after it: M\supertilde
\ega@define{\supertilde}{{\,\widetilde{\,}\,}}

% operators
\ega@operator{\rg}{rg}
\ega@operator{\gr}{gr}
\ega@operator{\Gr}{Gr}
\ega@operator{\Hom}{Hom}
\ega@operator{\Proj}{Proj}
\ega@operator{\Tor}{Tor}
\ega@operator{\Ext}{Ext}
\ega@operator{\Supp}{Supp}
\ega@operator{\Ker}{Ker}
\ega@operator{\Im}{Im}
\ega@operator{\Coker}{Coker}
\ega@operator{\Spec}{Spec}
\ega@operator{\Spm}{Spm}
\ega@operator{\Spf}{Spf}
\ega@operator{\grad}{grad}
\ega@operator{\dimc}{dimc}
\ega@operator{\codim}{codim}
\ega@operator{\id}{id}
\ega@operator{\Der}{Der}
\ega@operator{\Diff}{Diff}
\ega@operator{\Hyp}{Hyp}
\ega@operator{\Tr}{Tr}

% arrows: long in displays, short in text, scripts and subscripts
\let\ega@shortmapsto\mapsto
\renewcommand{\to}{\mathchoice{\longrightarrow}{\rightarrow}{\rightarrow}{\rightarrow}}
\providecommand{\from}{}\renewcommand{\from}{\mathchoice{\longleftarrow}{\leftarrow}{\leftarrow}{\leftarrow}}
\renewcommand{\mapsto}{\mathchoice{\longmapsto}{\ega@shortmapsto}{\ega@shortmapsto}{\ega@shortmapsto}}
\ega@define{\isoto}{\simeq}
\ega@define{\simto}{\xrightarrow{\sim}}

% wide accents
\renewcommand{\hat}[1]{\widehat{#1}}
\renewcommand{\tilde}[1]{\widetilde{#1}}
\makeatother

% END authoritative preamble-base.tex

\hypersetup{pdftitle={EGA I},pdfauthor={Alexander Grothendieck y Jean Dieudonné}}
\begin{document}
\begin{titlepage}\centering\vspace*{20mm}
{\Huge\bfseries Elementos de geometría algebraica\par}\vspace{14mm}
{\LARGE EGA I\par}\vspace{10mm}
{\large Alexander Grothendieck\\con la colaboración de Jean Dieudonné\par}
\vfill{\Large Edición española actual\par}
\vspace{8mm}{\large 25 de noviembre de 2026\par}\end{titlepage}
\chapter*{Créditos y derechos}
\addcontentsline{toc}{chapter}{Créditos y derechos}
\noindent\textbf{Trabajo de esta ejecución.} Traducción y preparación por GPT-6.0 Astra (OpenAI), modo Pro de ChatGPT, noviembre de 2026. Traducción producida por IA; correcciones producidas por IA. Cotejo francés heredado: Claude Opus 5.5, esfuerzo max, noviembre de 2026.

\subsection*{Créditos previos conservados}

Las traducciones están producidas por IA, y también las correcciones, la edición y la composición. Tradujeron
modelos de OpenAI Codex (GPT-5.6 Sol y GPT-6.1 Sol, esfuerzo Ultra); GPT-6 Astra (esfuerzo Ultra) revisó, corrigió
y montó las ediciones. En la
preparación de algunas secciones de EGA se usaron fuentes inglesas que incorporan la traducción humana del proyecto
ryankeleti/ega; para FGA, la traducción inglesa de Tim Hosgood sirvió de referencia (véase «Fuentes y créditos»).

Los lectores de EGA I (traducción de OpenAI Codex, GPT-6.1 Sol, esfuerzo de razonamiento ultra, agosto–octubre
de 2026) se publicaron el 15 de noviembre de 2026; la composición (nada sobresale del margen más de 2 pt) y la
publicación las hizo Claude Opus 5.5 con esfuerzo high.

Las ediciones corregidas y las traducciones diplomáticas, que siguen el texto impreso, son documentos separados; las
notas matemáticas también.

\subsection*{Fuentes y créditos del repositorio}

\begin{itemize}
\item EGA: Alexander Grothendieck, con la colaboración de Jean Dieudonné, \emph{Éléments de géométrie algébrique},
  \emph{Publications mathématiques de l'IHÉS} 4, 8, 11, 17, 20, 24, 28 y 32 (1960–1967), en
  \href{https://www.numdam.org/journals/PMIHES/}{NUMDAM}.
\item FGA: Alexander Grothendieck, \emph{Fondements de la géométrie algébrique}, exposiciones 149, 182, 190, 195, 212, 221,
  232 y 236 del Seminario Bourbaki, en \href{https://www.numdam.org/}{NUMDAM}.
\item EGA: las fuentes inglesas usadas para preparar EGA III-1 (capítulo 0, §§8–10 y parte del §11; capítulo III, §§1 y
  3), EGA IV-1 (capítulo 0, §14) y EGA IV-4 (§16 y el comienzo de los §§17 y 20) incorporan la traducción humana de
  \href{https://github.com/ryankeleti/ega/tree/3d7afe28c99d1887ccdb85acf2ca1afd93fb1971}{ryankeleti/ega} (versión
  3d7afe2), cuyos colaboradores son Mateo Carmona, Tim Holzschuh, Tim Hosgood, Ryan Keleti, Thiago Solovera \& Nery y
  Wayne Yang (\href{https://github.com/ryankeleti/ega/blob/3d7afe28c99d1887ccdb85acf2ca1afd93fb1971/README.md}{créditos por sección}).
\item FGA: la traducción inglesa de \href{https://github.com/thosgood/fga/tree/9976dd5f8e2ddcb388f0a07f3aa334f120d1019f}{Tim Hosgood}
  (versión 9976dd5), que agradece las correcciones de Steve Hnizdur y el apoyo técnico de Jon Sterling.
\item La traducción española no se atribuye a estos traductores. \href{../../tree/repo/ATTRIBUTION.md}{ATTRIBUTION.md} da las versiones exactas
  de las fuentes inglesas.
\end{itemize}

\subsection*{Fuentes y autoría}

\subsection*{FGA: traducción inglesa de referencia}

La obra original, \emph{Fondements de la géométrie algébrique}, es de Alexandre Grothendieck. La traducción inglesa de \href{https://github.com/thosgood/fga/tree/9976dd5f8e2ddcb388f0a07f3aa334f120d1019f}{Tim Hosgood} es una fuente de referencia de esta edición; su proyecto agradece las correcciones de Steve Hnizdur y el apoyo técnico de Jon Sterling. La versión 9976dd5f8e2ddcb388f0a07f3aa334f120d1019f identifica el texto inglés usado en la comparación; no establece la versión exacta incorporada al principio. La traducción española y las correcciones y la composición posteriores están producidas por IA y no se atribuyen a Tim Hosgood. Esta atribución no constituye permiso ni aval.

\subsection*{EGA: fuentes inglesas}

Las fuentes inglesas utilizadas para preparar las traducciones españolas de EGA incorporan la traducción humana de \href{https://github.com/ryankeleti/ega}{ryankeleti/ega}, versión \href{https://github.com/ryankeleti/ega/tree/3d7afe28c99d1887ccdb85acf2ca1afd93fb1971}{3d7afe28c99d1887ccdb85acf2ca1afd93fb1971}. El proyecto nombra como colaboradores a Mateo Carmona, Tim Holzschuh, Tim Hosgood, Ryan Keleti, Thiago Solovera \& Nery y Wayne Yang; sus \href{https://github.com/ryankeleti/ega/blob/3d7afe28c99d1887ccdb85acf2ca1afd93fb1971/README.md}{créditos por sección} y su \href{https://github.com/ryankeleti/ega/commits/3d7afe28c99d1887ccdb85acf2ca1afd93fb1971}{historial de contribuciones} precisan la autoría. Las secciones afectadas son:

\begin{itemize}
\item EGA III-1: los §§ 8–10 y parte del § 11 del capítulo 0, y los §§ 1 y 3 del capítulo III.
\item EGA IV-1, capítulo 0: el § 14.
\item EGA IV-4: el § 16 y las partes iniciales de los §§ 17 y 20 del capítulo IV.
\end{itemize}

La traducción inglesa y la estructura TeX heredadas conservan esa procedencia. La traducción española y las adiciones y correcciones producidas por IA no se atribuyen a esos colaboradores. La atribución no implica autorización ni aval.

\subsection*{Traducción española}

La traducción al español, las correcciones, la edición y la composición están producidas por IA. Tradujeron modelos de OpenAI Codex (GPT-5.6 Sol y GPT-6.1 Sol, esfuerzo Ultra); GPT-6 Astra (esfuerzo Ultra) revisó, corrigió y montó las ediciones. No se afirma revisión humana. La autoría de las obras originales y de las contribuciones humanas citadas se conserva.

\begin{itemize}
\item \href{../../tree/repo/editions/II\_CURRENT/FUENTES.md}{EGA II: alcance e identidad de las fuentes}.
\item \href{../../tree/repo/editions/III1\_CURRENT/FUENTES.md}{EGA III-1: alcance e identidad de las fuentes}.
\item \href{../../tree/repo/editions/IV1\_CURRENT/FUENTES.md}{EGA IV-1, capítulo 0: alcance e identidad de las fuentes}.
\item \href{../../tree/repo/editions/IV4/FUENTES.md}{EGA IV-4: alcance e identidad de las fuentes}.
\end{itemize}
\subsection*{Declaración de derechos de esta edición}
Los derechos sobre las obras originales pertenecen a sus autores o a los representantes o herederos de estos. Nuestras traducciones, correcciones y composición tipográfica fueron producidas por IA. No tienen autor humano y, según nuestro entender, no están protegidas por derechos de autor; si aun así existiera algún derecho nuestro, se dedica al dominio público bajo CC0 1.0. Todo este trabajo reproduce o traduce las obras originales, por lo que nada de lo que aquí se publica concede licencia ni autoriza reproducción, traducción u otro uso de ellas. Estas ediciones se publican con fines educativos y sin ánimo de lucro, y no reclamamos ningún derecho que no poseamos. Si usted es titular de derechos sobre alguna de estas obras y se opone a su publicación aquí, abra una incidencia en GitHub y la retiraremos.
\clearpage\tableofcontents\clearpage
\EGAReferenceAlias{0.1}{section:0.1-fr}
\EGAReferenceAlias{0.1-fr}{section:0.1-fr}
\EGAReferenceAlias{0.1-es}{section:0.1-fr}
\EGAReferenceAlias{0.1.fr}{section:0.1-fr}
\EGAReferenceAlias{0.1.es}{section:0.1-fr}
\EGAReferenceAlias{0.1.eq}{section:0.1-fr}
\EGAReferenceAlias{eq:0.1}{section:0.1-fr}
\EGAReferenceAlias{section:0.1}{section:0.1-fr}
\EGAReferenceAlias{section:0.1-es}{section:0.1-fr}
\EGAReferenceAlias{section:0.1.fr}{section:0.1-fr}
\EGAReferenceAlias{section:0.1.es}{section:0.1-fr}
\EGAReferenceAlias{section:0.1.eq}{section:0.1-fr}
\EGAReferenceAlias{eq:section:0.1}{section:0.1-fr}
\EGAReferenceAlias{0.1.0}{subsection:0.1.0-fr}
\EGAReferenceAlias{0.1.0-fr}{subsection:0.1.0-fr}
\EGAReferenceAlias{0.1.0-es}{subsection:0.1.0-fr}
\EGAReferenceAlias{0.1.0.fr}{subsection:0.1.0-fr}
\EGAReferenceAlias{0.1.0.es}{subsection:0.1.0-fr}
\EGAReferenceAlias{0.1.0.eq}{subsection:0.1.0-fr}
\EGAReferenceAlias{eq:0.1.0}{subsection:0.1.0-fr}
\EGAReferenceAlias{subsection:0.1.0}{subsection:0.1.0-fr}
\EGAReferenceAlias{subsection:0.1.0-es}{subsection:0.1.0-fr}
\EGAReferenceAlias{subsection:0.1.0.fr}{subsection:0.1.0-fr}
\EGAReferenceAlias{subsection:0.1.0.es}{subsection:0.1.0-fr}
\EGAReferenceAlias{subsection:0.1.0.eq}{subsection:0.1.0-fr}
\EGAReferenceAlias{eq:subsection:0.1.0}{subsection:0.1.0-fr}
\EGAReferenceAlias{0.1.0.1}{0.1.0.1-fr}
\EGAReferenceAlias{0.1.0.1-es}{0.1.0.1-fr}
\EGAReferenceAlias{0.1.0.1.fr}{0.1.0.1-fr}
\EGAReferenceAlias{0.1.0.1.es}{0.1.0.1-fr}
\EGAReferenceAlias{0.1.0.1.eq}{0.1.0.1-fr}
\EGAReferenceAlias{eq:0.1.0.1}{0.1.0.1-fr}
\EGAReferenceAlias{0.1.0.2}{0.1.0.2-fr}
\EGAReferenceAlias{0.1.0.2-es}{0.1.0.2-fr}
\EGAReferenceAlias{0.1.0.2.fr}{0.1.0.2-fr}
\EGAReferenceAlias{0.1.0.2.es}{0.1.0.2-fr}
\EGAReferenceAlias{0.1.0.2.eq}{0.1.0.2-fr}
\EGAReferenceAlias{eq:0.1.0.2}{0.1.0.2-fr}
\EGAReferenceAlias{0.1.0.3}{0.1.0.3-fr}
\EGAReferenceAlias{0.1.0.3-es}{0.1.0.3-fr}
\EGAReferenceAlias{0.1.0.3.fr}{0.1.0.3-fr}
\EGAReferenceAlias{0.1.0.3.es}{0.1.0.3-fr}
\EGAReferenceAlias{0.1.0.3.eq}{0.1.0.3-fr}
\EGAReferenceAlias{eq:0.1.0.3}{0.1.0.3-fr}
\EGAReferenceAlias{0.1.0.4}{0.1.0.4-fr}
\EGAReferenceAlias{0.1.0.4-es}{0.1.0.4-fr}
\EGAReferenceAlias{0.1.0.4.fr}{0.1.0.4-fr}
\EGAReferenceAlias{0.1.0.4.es}{0.1.0.4-fr}
\EGAReferenceAlias{0.1.0.4.eq}{0.1.0.4-fr}
\EGAReferenceAlias{eq:0.1.0.4}{0.1.0.4-fr}
\EGAReferenceAlias{0.1.0.5}{0.1.0.5-fr}
\EGAReferenceAlias{0.1.0.5-es}{0.1.0.5-fr}
\EGAReferenceAlias{0.1.0.5.fr}{0.1.0.5-fr}
\EGAReferenceAlias{0.1.0.5.es}{0.1.0.5-fr}
\EGAReferenceAlias{0.1.0.5.eq}{0.1.0.5-fr}
\EGAReferenceAlias{eq:0.1.0.5}{0.1.0.5-fr}
\EGAReferenceAlias{0.1.0.6}{0.1.0.6-fr}
\EGAReferenceAlias{0.1.0.6-es}{0.1.0.6-fr}
\EGAReferenceAlias{0.1.0.6.fr}{0.1.0.6-fr}
\EGAReferenceAlias{0.1.0.6.es}{0.1.0.6-fr}
\EGAReferenceAlias{0.1.0.6.eq}{0.1.0.6-fr}
\EGAReferenceAlias{eq:0.1.0.6}{0.1.0.6-fr}
\EGAReferenceAlias{0.1.0.7}{0.1.0.7-fr}
\EGAReferenceAlias{0.1.0.7-es}{0.1.0.7-fr}
\EGAReferenceAlias{0.1.0.7.fr}{0.1.0.7-fr}
\EGAReferenceAlias{0.1.0.7.es}{0.1.0.7-fr}
\EGAReferenceAlias{0.1.0.7.eq}{0.1.0.7-fr}
\EGAReferenceAlias{eq:0.1.0.7}{0.1.0.7-fr}
\EGAReferenceAlias{0.1.1}{subsection:0.1.1-fr}
\EGAReferenceAlias{0.1.1-fr}{subsection:0.1.1-fr}
\EGAReferenceAlias{0.1.1-es}{subsection:0.1.1-fr}
\EGAReferenceAlias{0.1.1.fr}{subsection:0.1.1-fr}
\EGAReferenceAlias{0.1.1.es}{subsection:0.1.1-fr}
\EGAReferenceAlias{0.1.1.eq}{subsection:0.1.1-fr}
\EGAReferenceAlias{eq:0.1.1}{subsection:0.1.1-fr}
\EGAReferenceAlias{subsection:0.1.1}{subsection:0.1.1-fr}
\EGAReferenceAlias{subsection:0.1.1-es}{subsection:0.1.1-fr}
\EGAReferenceAlias{subsection:0.1.1.fr}{subsection:0.1.1-fr}
\EGAReferenceAlias{subsection:0.1.1.es}{subsection:0.1.1-fr}
\EGAReferenceAlias{subsection:0.1.1.eq}{subsection:0.1.1-fr}
\EGAReferenceAlias{eq:subsection:0.1.1}{subsection:0.1.1-fr}
\EGAReferenceAlias{0.1.1.1}{0.1.1.1-fr}
\EGAReferenceAlias{0.1.1.1-es}{0.1.1.1-fr}
\EGAReferenceAlias{0.1.1.1.fr}{0.1.1.1-fr}
\EGAReferenceAlias{0.1.1.1.es}{0.1.1.1-fr}
\EGAReferenceAlias{0.1.1.1.eq}{0.1.1.1-fr}
\EGAReferenceAlias{eq:0.1.1.1}{0.1.1.1-fr}
\EGAReferenceAlias{0.1.1.2}{0.1.1.2-fr}
\EGAReferenceAlias{0.1.1.2-es}{0.1.1.2-fr}
\EGAReferenceAlias{0.1.1.2.fr}{0.1.1.2-fr}
\EGAReferenceAlias{0.1.1.2.es}{0.1.1.2-fr}
\EGAReferenceAlias{0.1.1.2.eq}{0.1.1.2-fr}
\EGAReferenceAlias{eq:0.1.1.2}{0.1.1.2-fr}
\EGAReferenceAlias{0.1.2}{subsection:0.1.2-fr}
\EGAReferenceAlias{0.1.2-fr}{subsection:0.1.2-fr}
\EGAReferenceAlias{0.1.2-es}{subsection:0.1.2-fr}
\EGAReferenceAlias{0.1.2.fr}{subsection:0.1.2-fr}
\EGAReferenceAlias{0.1.2.es}{subsection:0.1.2-fr}
\EGAReferenceAlias{0.1.2.eq}{subsection:0.1.2-fr}
\EGAReferenceAlias{eq:0.1.2}{subsection:0.1.2-fr}
\EGAReferenceAlias{subsection:0.1.2}{subsection:0.1.2-fr}
\EGAReferenceAlias{subsection:0.1.2-es}{subsection:0.1.2-fr}
\EGAReferenceAlias{subsection:0.1.2.fr}{subsection:0.1.2-fr}
\EGAReferenceAlias{subsection:0.1.2.es}{subsection:0.1.2-fr}
\EGAReferenceAlias{subsection:0.1.2.eq}{subsection:0.1.2-fr}
\EGAReferenceAlias{eq:subsection:0.1.2}{subsection:0.1.2-fr}
\EGAReferenceAlias{0.1.2.1}{0.1.2.1-fr}
\EGAReferenceAlias{0.1.2.1-es}{0.1.2.1-fr}
\EGAReferenceAlias{0.1.2.1.fr}{0.1.2.1-fr}
\EGAReferenceAlias{0.1.2.1.es}{0.1.2.1-fr}
\EGAReferenceAlias{0.1.2.1.eq}{0.1.2.1-fr}
\EGAReferenceAlias{eq:0.1.2.1}{0.1.2.1-fr}
\EGAReferenceAlias{0.1.2.2}{0.1.2.2-fr}
\EGAReferenceAlias{0.1.2.2-es}{0.1.2.2-fr}
\EGAReferenceAlias{0.1.2.2.fr}{0.1.2.2-fr}
\EGAReferenceAlias{0.1.2.2.es}{0.1.2.2-fr}
\EGAReferenceAlias{0.1.2.2.eq}{0.1.2.2-fr}
\EGAReferenceAlias{eq:0.1.2.2}{0.1.2.2-fr}
\EGAReferenceAlias{0.1.2.3}{0.1.2.3-fr}
\EGAReferenceAlias{0.1.2.3-es}{0.1.2.3-fr}
\EGAReferenceAlias{0.1.2.3.fr}{0.1.2.3-fr}
\EGAReferenceAlias{0.1.2.3.es}{0.1.2.3-fr}
\EGAReferenceAlias{0.1.2.3.eq}{0.1.2.3-fr}
\EGAReferenceAlias{eq:0.1.2.3}{0.1.2.3-fr}
\EGAReferenceAlias{0.1.2.4}{0.1.2.4-fr}
\EGAReferenceAlias{0.1.2.4-es}{0.1.2.4-fr}
\EGAReferenceAlias{0.1.2.4.fr}{0.1.2.4-fr}
\EGAReferenceAlias{0.1.2.4.es}{0.1.2.4-fr}
\EGAReferenceAlias{0.1.2.4.eq}{0.1.2.4-fr}
\EGAReferenceAlias{eq:0.1.2.4}{0.1.2.4-fr}
\EGAReferenceAlias{0.1.2.5}{0.1.2.5-fr}
\EGAReferenceAlias{0.1.2.5-es}{0.1.2.5-fr}
\EGAReferenceAlias{0.1.2.5.fr}{0.1.2.5-fr}
\EGAReferenceAlias{0.1.2.5.es}{0.1.2.5-fr}
\EGAReferenceAlias{0.1.2.5.eq}{0.1.2.5-fr}
\EGAReferenceAlias{eq:0.1.2.5}{0.1.2.5-fr}
\EGAReferenceAlias{0.1.2.6}{0.1.2.6-fr}
\EGAReferenceAlias{0.1.2.6-es}{0.1.2.6-fr}
\EGAReferenceAlias{0.1.2.6.fr}{0.1.2.6-fr}
\EGAReferenceAlias{0.1.2.6.es}{0.1.2.6-fr}
\EGAReferenceAlias{0.1.2.6.eq}{0.1.2.6-fr}
\EGAReferenceAlias{eq:0.1.2.6}{0.1.2.6-fr}
\EGAReferenceAlias{0.1.2.7}{0.1.2.7-fr}
\EGAReferenceAlias{0.1.2.7-es}{0.1.2.7-fr}
\EGAReferenceAlias{0.1.2.7.fr}{0.1.2.7-fr}
\EGAReferenceAlias{0.1.2.7.es}{0.1.2.7-fr}
\EGAReferenceAlias{0.1.2.7.eq}{0.1.2.7-fr}
\EGAReferenceAlias{eq:0.1.2.7}{0.1.2.7-fr}
\EGAReferenceAlias{0.1.2.8}{0.1.2.8-fr}
\EGAReferenceAlias{0.1.2.8-es}{0.1.2.8-fr}
\EGAReferenceAlias{0.1.2.8.fr}{0.1.2.8-fr}
\EGAReferenceAlias{0.1.2.8.es}{0.1.2.8-fr}
\EGAReferenceAlias{0.1.2.8.eq}{0.1.2.8-fr}
\EGAReferenceAlias{eq:0.1.2.8}{0.1.2.8-fr}
\EGAReferenceAlias{0.1.3}{subsection:0.1.3-fr}
\EGAReferenceAlias{0.1.3-fr}{subsection:0.1.3-fr}
\EGAReferenceAlias{0.1.3-es}{subsection:0.1.3-fr}
\EGAReferenceAlias{0.1.3.fr}{subsection:0.1.3-fr}
\EGAReferenceAlias{0.1.3.es}{subsection:0.1.3-fr}
\EGAReferenceAlias{0.1.3.eq}{subsection:0.1.3-fr}
\EGAReferenceAlias{eq:0.1.3}{subsection:0.1.3-fr}
\EGAReferenceAlias{subsection:0.1.3}{subsection:0.1.3-fr}
\EGAReferenceAlias{subsection:0.1.3-es}{subsection:0.1.3-fr}
\EGAReferenceAlias{subsection:0.1.3.fr}{subsection:0.1.3-fr}
\EGAReferenceAlias{subsection:0.1.3.es}{subsection:0.1.3-fr}
\EGAReferenceAlias{subsection:0.1.3.eq}{subsection:0.1.3-fr}
\EGAReferenceAlias{eq:subsection:0.1.3}{subsection:0.1.3-fr}
\EGAReferenceAlias{0.1.3.1}{0.1.3.1-fr}
\EGAReferenceAlias{0.1.3.1-es}{0.1.3.1-fr}
\EGAReferenceAlias{0.1.3.1.fr}{0.1.3.1-fr}
\EGAReferenceAlias{0.1.3.1.es}{0.1.3.1-fr}
\EGAReferenceAlias{0.1.3.1.eq}{0.1.3.1-fr}
\EGAReferenceAlias{eq:0.1.3.1}{0.1.3.1-fr}
\EGAReferenceAlias{0.1.3.2}{0.1.3.2-fr}
\EGAReferenceAlias{0.1.3.2-es}{0.1.3.2-fr}
\EGAReferenceAlias{0.1.3.2.fr}{0.1.3.2-fr}
\EGAReferenceAlias{0.1.3.2.es}{0.1.3.2-fr}
\EGAReferenceAlias{0.1.3.2.eq}{0.1.3.2-fr}
\EGAReferenceAlias{eq:0.1.3.2}{0.1.3.2-fr}
\EGAReferenceAlias{0.1.3.3}{0.1.3.3-fr}
\EGAReferenceAlias{0.1.3.3-es}{0.1.3.3-fr}
\EGAReferenceAlias{0.1.3.3.fr}{0.1.3.3-fr}
\EGAReferenceAlias{0.1.3.3.es}{0.1.3.3-fr}
\EGAReferenceAlias{0.1.3.3.eq}{0.1.3.3-fr}
\EGAReferenceAlias{eq:0.1.3.3}{0.1.3.3-fr}
\EGAReferenceAlias{0.1.3.4}{0.1.3.4-fr}
\EGAReferenceAlias{0.1.3.4-es}{0.1.3.4-fr}
\EGAReferenceAlias{0.1.3.4.fr}{0.1.3.4-fr}
\EGAReferenceAlias{0.1.3.4.es}{0.1.3.4-fr}
\EGAReferenceAlias{0.1.3.4.eq}{0.1.3.4-fr}
\EGAReferenceAlias{eq:0.1.3.4}{0.1.3.4-fr}
\EGAReferenceAlias{0.1.3.5}{0.1.3.5-fr}
\EGAReferenceAlias{0.1.3.5-es}{0.1.3.5-fr}
\EGAReferenceAlias{0.1.3.5.fr}{0.1.3.5-fr}
\EGAReferenceAlias{0.1.3.5.es}{0.1.3.5-fr}
\EGAReferenceAlias{0.1.3.5.eq}{0.1.3.5-fr}
\EGAReferenceAlias{eq:0.1.3.5}{0.1.3.5-fr}
\EGAReferenceAlias{0.1.4}{subsection:0.1.4-fr}
\EGAReferenceAlias{0.1.4-fr}{subsection:0.1.4-fr}
\EGAReferenceAlias{0.1.4-es}{subsection:0.1.4-fr}
\EGAReferenceAlias{0.1.4.fr}{subsection:0.1.4-fr}
\EGAReferenceAlias{0.1.4.es}{subsection:0.1.4-fr}
\EGAReferenceAlias{0.1.4.eq}{subsection:0.1.4-fr}
\EGAReferenceAlias{eq:0.1.4}{subsection:0.1.4-fr}
\EGAReferenceAlias{subsection:0.1.4}{subsection:0.1.4-fr}
\EGAReferenceAlias{subsection:0.1.4-es}{subsection:0.1.4-fr}
\EGAReferenceAlias{subsection:0.1.4.fr}{subsection:0.1.4-fr}
\EGAReferenceAlias{subsection:0.1.4.es}{subsection:0.1.4-fr}
\EGAReferenceAlias{subsection:0.1.4.eq}{subsection:0.1.4-fr}
\EGAReferenceAlias{eq:subsection:0.1.4}{subsection:0.1.4-fr}
\EGAReferenceAlias{0.1.4.1}{0.1.4.1-fr}
\EGAReferenceAlias{0.1.4.1-es}{0.1.4.1-fr}
\EGAReferenceAlias{0.1.4.1.fr}{0.1.4.1-fr}
\EGAReferenceAlias{0.1.4.1.es}{0.1.4.1-fr}
\EGAReferenceAlias{0.1.4.1.eq}{0.1.4.1-fr}
\EGAReferenceAlias{eq:0.1.4.1}{0.1.4.1-fr}
\EGAReferenceAlias{0.1.4.2}{0.1.4.2-fr}
\EGAReferenceAlias{0.1.4.2-es}{0.1.4.2-fr}
\EGAReferenceAlias{0.1.4.2.fr}{0.1.4.2-fr}
\EGAReferenceAlias{0.1.4.2.es}{0.1.4.2-fr}
\EGAReferenceAlias{0.1.4.2.eq}{0.1.4.2-fr}
\EGAReferenceAlias{eq:0.1.4.2}{0.1.4.2-fr}
\EGAReferenceAlias{0.1.4.3}{0.1.4.3-fr}
\EGAReferenceAlias{0.1.4.3-es}{0.1.4.3-fr}
\EGAReferenceAlias{0.1.4.3.fr}{0.1.4.3-fr}
\EGAReferenceAlias{0.1.4.3.es}{0.1.4.3-fr}
\EGAReferenceAlias{0.1.4.3.eq}{0.1.4.3-fr}
\EGAReferenceAlias{eq:0.1.4.3}{0.1.4.3-fr}
\EGAReferenceAlias{0.1.4.4}{0.1.4.4-fr}
\EGAReferenceAlias{0.1.4.4-es}{0.1.4.4-fr}
\EGAReferenceAlias{0.1.4.4.fr}{0.1.4.4-fr}
\EGAReferenceAlias{0.1.4.4.es}{0.1.4.4-fr}
\EGAReferenceAlias{0.1.4.4.eq}{0.1.4.4-fr}
\EGAReferenceAlias{eq:0.1.4.4}{0.1.4.4-fr}
\EGAReferenceAlias{0.1.4.5}{0.1.4.5-fr}
\EGAReferenceAlias{0.1.4.5-es}{0.1.4.5-fr}
\EGAReferenceAlias{0.1.4.5.fr}{0.1.4.5-fr}
\EGAReferenceAlias{0.1.4.5.es}{0.1.4.5-fr}
\EGAReferenceAlias{0.1.4.5.eq}{0.1.4.5-fr}
\EGAReferenceAlias{eq:0.1.4.5}{0.1.4.5-fr}
\EGAReferenceAlias{0.1.4.6}{0.1.4.6-fr}
\EGAReferenceAlias{0.1.4.6-es}{0.1.4.6-fr}
\EGAReferenceAlias{0.1.4.6.fr}{0.1.4.6-fr}
\EGAReferenceAlias{0.1.4.6.es}{0.1.4.6-fr}
\EGAReferenceAlias{0.1.4.6.eq}{0.1.4.6-fr}
\EGAReferenceAlias{eq:0.1.4.6}{0.1.4.6-fr}
\EGAReferenceAlias{0.1.5}{subsection:0.1.5-fr}
\EGAReferenceAlias{0.1.5-fr}{subsection:0.1.5-fr}
\EGAReferenceAlias{0.1.5-es}{subsection:0.1.5-fr}
\EGAReferenceAlias{0.1.5.fr}{subsection:0.1.5-fr}
\EGAReferenceAlias{0.1.5.es}{subsection:0.1.5-fr}
\EGAReferenceAlias{0.1.5.eq}{subsection:0.1.5-fr}
\EGAReferenceAlias{eq:0.1.5}{subsection:0.1.5-fr}
\EGAReferenceAlias{subsection:0.1.5}{subsection:0.1.5-fr}
\EGAReferenceAlias{subsection:0.1.5-es}{subsection:0.1.5-fr}
\EGAReferenceAlias{subsection:0.1.5.fr}{subsection:0.1.5-fr}
\EGAReferenceAlias{subsection:0.1.5.es}{subsection:0.1.5-fr}
\EGAReferenceAlias{subsection:0.1.5.eq}{subsection:0.1.5-fr}
\EGAReferenceAlias{eq:subsection:0.1.5}{subsection:0.1.5-fr}
\EGAReferenceAlias{0.1.5.1}{0.1.5.1-fr}
\EGAReferenceAlias{0.1.5.1-es}{0.1.5.1-fr}
\EGAReferenceAlias{0.1.5.1.fr}{0.1.5.1-fr}
\EGAReferenceAlias{0.1.5.1.es}{0.1.5.1-fr}
\EGAReferenceAlias{0.1.5.1.eq}{0.1.5.1-fr}
\EGAReferenceAlias{eq:0.1.5.1}{0.1.5.1-fr}
\EGAReferenceAlias{0.1.5.2}{0.1.5.2-fr}
\EGAReferenceAlias{0.1.5.2-es}{0.1.5.2-fr}
\EGAReferenceAlias{0.1.5.2.fr}{0.1.5.2-fr}
\EGAReferenceAlias{0.1.5.2.es}{0.1.5.2-fr}
\EGAReferenceAlias{0.1.5.2.eq}{0.1.5.2-fr}
\EGAReferenceAlias{eq:0.1.5.2}{0.1.5.2-fr}
\EGAReferenceAlias{0.1.5.3}{0.1.5.3-fr}
\EGAReferenceAlias{0.1.5.3-es}{0.1.5.3-fr}
\EGAReferenceAlias{0.1.5.3.fr}{0.1.5.3-fr}
\EGAReferenceAlias{0.1.5.3.es}{0.1.5.3-fr}
\EGAReferenceAlias{0.1.5.3.eq}{0.1.5.3-fr}
\EGAReferenceAlias{eq:0.1.5.3}{0.1.5.3-fr}
\EGAReferenceAlias{0.1.5.4}{0.1.5.4-fr}
\EGAReferenceAlias{0.1.5.4-es}{0.1.5.4-fr}
\EGAReferenceAlias{0.1.5.4.fr}{0.1.5.4-fr}
\EGAReferenceAlias{0.1.5.4.es}{0.1.5.4-fr}
\EGAReferenceAlias{0.1.5.4.eq}{0.1.5.4-fr}
\EGAReferenceAlias{eq:0.1.5.4}{0.1.5.4-fr}
\EGAReferenceAlias{0.1.5.5}{0.1.5.5-fr}
\EGAReferenceAlias{0.1.5.5-es}{0.1.5.5-fr}
\EGAReferenceAlias{0.1.5.5.fr}{0.1.5.5-fr}
\EGAReferenceAlias{0.1.5.5.es}{0.1.5.5-fr}
\EGAReferenceAlias{0.1.5.5.eq}{0.1.5.5-fr}
\EGAReferenceAlias{eq:0.1.5.5}{0.1.5.5-fr}
\EGAReferenceAlias{0.1.5.6}{0.1.5.6-fr}
\EGAReferenceAlias{0.1.5.6-es}{0.1.5.6-fr}
\EGAReferenceAlias{0.1.5.6.fr}{0.1.5.6-fr}
\EGAReferenceAlias{0.1.5.6.es}{0.1.5.6-fr}
\EGAReferenceAlias{0.1.5.6.eq}{0.1.5.6-fr}
\EGAReferenceAlias{eq:0.1.5.6}{0.1.5.6-fr}
\EGAReferenceAlias{0.1.5.7}{0.1.5.7-fr}
\EGAReferenceAlias{0.1.5.7-es}{0.1.5.7-fr}
\EGAReferenceAlias{0.1.5.7.fr}{0.1.5.7-fr}
\EGAReferenceAlias{0.1.5.7.es}{0.1.5.7-fr}
\EGAReferenceAlias{0.1.5.7.eq}{0.1.5.7-fr}
\EGAReferenceAlias{eq:0.1.5.7}{0.1.5.7-fr}
\EGAReferenceAlias{0.1.5.8}{0.1.5.8-fr}
\EGAReferenceAlias{0.1.5.8-es}{0.1.5.8-fr}
\EGAReferenceAlias{0.1.5.8.fr}{0.1.5.8-fr}
\EGAReferenceAlias{0.1.5.8.es}{0.1.5.8-fr}
\EGAReferenceAlias{0.1.5.8.eq}{0.1.5.8-fr}
\EGAReferenceAlias{eq:0.1.5.8}{0.1.5.8-fr}
\EGAReferenceAlias{0.1.6}{subsection:0.1.6-fr}
\EGAReferenceAlias{0.1.6-fr}{subsection:0.1.6-fr}
\EGAReferenceAlias{0.1.6-es}{subsection:0.1.6-fr}
\EGAReferenceAlias{0.1.6.fr}{subsection:0.1.6-fr}
\EGAReferenceAlias{0.1.6.es}{subsection:0.1.6-fr}
\EGAReferenceAlias{0.1.6.eq}{subsection:0.1.6-fr}
\EGAReferenceAlias{eq:0.1.6}{subsection:0.1.6-fr}
\EGAReferenceAlias{subsection:0.1.6}{subsection:0.1.6-fr}
\EGAReferenceAlias{subsection:0.1.6-es}{subsection:0.1.6-fr}
\EGAReferenceAlias{subsection:0.1.6.fr}{subsection:0.1.6-fr}
\EGAReferenceAlias{subsection:0.1.6.es}{subsection:0.1.6-fr}
\EGAReferenceAlias{subsection:0.1.6.eq}{subsection:0.1.6-fr}
\EGAReferenceAlias{eq:subsection:0.1.6}{subsection:0.1.6-fr}
\EGAReferenceAlias{0.1.6.1}{0.1.6.1-fr}
\EGAReferenceAlias{0.1.6.1-es}{0.1.6.1-fr}
\EGAReferenceAlias{0.1.6.1.fr}{0.1.6.1-fr}
\EGAReferenceAlias{0.1.6.1.es}{0.1.6.1-fr}
\EGAReferenceAlias{0.1.6.1.eq}{0.1.6.1-fr}
\EGAReferenceAlias{eq:0.1.6.1}{0.1.6.1-fr}
\EGAReferenceAlias{0.1.6.2}{0.1.6.2-fr}
\EGAReferenceAlias{0.1.6.2-es}{0.1.6.2-fr}
\EGAReferenceAlias{0.1.6.2.fr}{0.1.6.2-fr}
\EGAReferenceAlias{0.1.6.2.es}{0.1.6.2-fr}
\EGAReferenceAlias{0.1.6.2.eq}{0.1.6.2-fr}
\EGAReferenceAlias{eq:0.1.6.2}{0.1.6.2-fr}
\EGAReferenceAlias{0.1.7}{subsection:0.1.7-fr}
\EGAReferenceAlias{0.1.7-fr}{subsection:0.1.7-fr}
\EGAReferenceAlias{0.1.7-es}{subsection:0.1.7-fr}
\EGAReferenceAlias{0.1.7.fr}{subsection:0.1.7-fr}
\EGAReferenceAlias{0.1.7.es}{subsection:0.1.7-fr}
\EGAReferenceAlias{0.1.7.eq}{subsection:0.1.7-fr}
\EGAReferenceAlias{eq:0.1.7}{subsection:0.1.7-fr}
\EGAReferenceAlias{subsection:0.1.7}{subsection:0.1.7-fr}
\EGAReferenceAlias{subsection:0.1.7-es}{subsection:0.1.7-fr}
\EGAReferenceAlias{subsection:0.1.7.fr}{subsection:0.1.7-fr}
\EGAReferenceAlias{subsection:0.1.7.es}{subsection:0.1.7-fr}
\EGAReferenceAlias{subsection:0.1.7.eq}{subsection:0.1.7-fr}
\EGAReferenceAlias{eq:subsection:0.1.7}{subsection:0.1.7-fr}
\EGAReferenceAlias{0.1.7.1}{0.1.7.1-fr}
\EGAReferenceAlias{0.1.7.1-es}{0.1.7.1-fr}
\EGAReferenceAlias{0.1.7.1.fr}{0.1.7.1-fr}
\EGAReferenceAlias{0.1.7.1.es}{0.1.7.1-fr}
\EGAReferenceAlias{0.1.7.1.eq}{0.1.7.1-fr}
\EGAReferenceAlias{eq:0.1.7.1}{0.1.7.1-fr}
\EGAReferenceAlias{0.1.7.2}{0.1.7.2-fr}
\EGAReferenceAlias{0.1.7.2-es}{0.1.7.2-fr}
\EGAReferenceAlias{0.1.7.2.fr}{0.1.7.2-fr}
\EGAReferenceAlias{0.1.7.2.es}{0.1.7.2-fr}
\EGAReferenceAlias{0.1.7.2.eq}{0.1.7.2-fr}
\EGAReferenceAlias{eq:0.1.7.2}{0.1.7.2-fr}
\EGAReferenceAlias{0.1.7.3}{0.1.7.3-fr}
\EGAReferenceAlias{0.1.7.3-es}{0.1.7.3-fr}
\EGAReferenceAlias{0.1.7.3.fr}{0.1.7.3-fr}
\EGAReferenceAlias{0.1.7.3.es}{0.1.7.3-fr}
\EGAReferenceAlias{0.1.7.3.eq}{0.1.7.3-fr}
\EGAReferenceAlias{eq:0.1.7.3}{0.1.7.3-fr}
\EGAReferenceAlias{0.1.7.4}{0.1.7.4-fr}
\EGAReferenceAlias{0.1.7.4-es}{0.1.7.4-fr}
\EGAReferenceAlias{0.1.7.4.fr}{0.1.7.4-fr}
\EGAReferenceAlias{0.1.7.4.es}{0.1.7.4-fr}
\EGAReferenceAlias{0.1.7.4.eq}{0.1.7.4-fr}
\EGAReferenceAlias{eq:0.1.7.4}{0.1.7.4-fr}
\EGAReferenceAlias{0.1.7.5}{0.1.7.5-fr}
\EGAReferenceAlias{0.1.7.5-es}{0.1.7.5-fr}
\EGAReferenceAlias{0.1.7.5.fr}{0.1.7.5-fr}
\EGAReferenceAlias{0.1.7.5.es}{0.1.7.5-fr}
\EGAReferenceAlias{0.1.7.5.eq}{0.1.7.5-fr}
\EGAReferenceAlias{eq:0.1.7.5}{0.1.7.5-fr}
\EGAReferenceAlias{0.2}{section:0.2-fr}
\EGAReferenceAlias{0.2-fr}{section:0.2-fr}
\EGAReferenceAlias{0.2-es}{section:0.2-fr}
\EGAReferenceAlias{0.2.fr}{section:0.2-fr}
\EGAReferenceAlias{0.2.es}{section:0.2-fr}
\EGAReferenceAlias{0.2.eq}{section:0.2-fr}
\EGAReferenceAlias{eq:0.2}{section:0.2-fr}
\EGAReferenceAlias{section:0.2}{section:0.2-fr}
\EGAReferenceAlias{section:0.2-es}{section:0.2-fr}
\EGAReferenceAlias{section:0.2.fr}{section:0.2-fr}
\EGAReferenceAlias{section:0.2.es}{section:0.2-fr}
\EGAReferenceAlias{section:0.2.eq}{section:0.2-fr}
\EGAReferenceAlias{eq:section:0.2}{section:0.2-fr}
\EGAReferenceAlias{0.2.1}{subsection:0.2.1-fr}
\EGAReferenceAlias{0.2.1-fr}{subsection:0.2.1-fr}
\EGAReferenceAlias{0.2.1-es}{subsection:0.2.1-fr}
\EGAReferenceAlias{0.2.1.fr}{subsection:0.2.1-fr}
\EGAReferenceAlias{0.2.1.es}{subsection:0.2.1-fr}
\EGAReferenceAlias{0.2.1.eq}{subsection:0.2.1-fr}
\EGAReferenceAlias{eq:0.2.1}{subsection:0.2.1-fr}
\EGAReferenceAlias{subsection:0.2.1}{subsection:0.2.1-fr}
\EGAReferenceAlias{subsection:0.2.1-es}{subsection:0.2.1-fr}
\EGAReferenceAlias{subsection:0.2.1.fr}{subsection:0.2.1-fr}
\EGAReferenceAlias{subsection:0.2.1.es}{subsection:0.2.1-fr}
\EGAReferenceAlias{subsection:0.2.1.eq}{subsection:0.2.1-fr}
\EGAReferenceAlias{eq:subsection:0.2.1}{subsection:0.2.1-fr}
\EGAReferenceAlias{0.2.1.1}{0.2.1.1-fr}
\EGAReferenceAlias{0.2.1.1-es}{0.2.1.1-fr}
\EGAReferenceAlias{0.2.1.1.fr}{0.2.1.1-fr}
\EGAReferenceAlias{0.2.1.1.es}{0.2.1.1-fr}
\EGAReferenceAlias{0.2.1.1.eq}{0.2.1.1-fr}
\EGAReferenceAlias{eq:0.2.1.1}{0.2.1.1-fr}
\EGAReferenceAlias{0.2.1.2}{0.2.1.2-fr}
\EGAReferenceAlias{0.2.1.2-es}{0.2.1.2-fr}
\EGAReferenceAlias{0.2.1.2.fr}{0.2.1.2-fr}
\EGAReferenceAlias{0.2.1.2.es}{0.2.1.2-fr}
\EGAReferenceAlias{0.2.1.2.eq}{0.2.1.2-fr}
\EGAReferenceAlias{eq:0.2.1.2}{0.2.1.2-fr}
\EGAReferenceAlias{0.2.1.3}{0.2.1.3-fr}
\EGAReferenceAlias{0.2.1.3-es}{0.2.1.3-fr}
\EGAReferenceAlias{0.2.1.3.fr}{0.2.1.3-fr}
\EGAReferenceAlias{0.2.1.3.es}{0.2.1.3-fr}
\EGAReferenceAlias{0.2.1.3.eq}{0.2.1.3-fr}
\EGAReferenceAlias{eq:0.2.1.3}{0.2.1.3-fr}
\EGAReferenceAlias{0.2.1.4}{0.2.1.4-fr}
\EGAReferenceAlias{0.2.1.4-es}{0.2.1.4-fr}
\EGAReferenceAlias{0.2.1.4.fr}{0.2.1.4-fr}
\EGAReferenceAlias{0.2.1.4.es}{0.2.1.4-fr}
\EGAReferenceAlias{0.2.1.4.eq}{0.2.1.4-fr}
\EGAReferenceAlias{eq:0.2.1.4}{0.2.1.4-fr}
\EGAReferenceAlias{0.2.1.5}{0.2.1.5-fr}
\EGAReferenceAlias{0.2.1.5-es}{0.2.1.5-fr}
\EGAReferenceAlias{0.2.1.5.fr}{0.2.1.5-fr}
\EGAReferenceAlias{0.2.1.5.es}{0.2.1.5-fr}
\EGAReferenceAlias{0.2.1.5.eq}{0.2.1.5-fr}
\EGAReferenceAlias{eq:0.2.1.5}{0.2.1.5-fr}
\EGAReferenceAlias{0.2.1.6}{0.2.1.6-fr}
\EGAReferenceAlias{0.2.1.6-es}{0.2.1.6-fr}
\EGAReferenceAlias{0.2.1.6.fr}{0.2.1.6-fr}
\EGAReferenceAlias{0.2.1.6.es}{0.2.1.6-fr}
\EGAReferenceAlias{0.2.1.6.eq}{0.2.1.6-fr}
\EGAReferenceAlias{eq:0.2.1.6}{0.2.1.6-fr}
\EGAReferenceAlias{0.2.1.7}{0.2.1.7-fr}
\EGAReferenceAlias{0.2.1.7-es}{0.2.1.7-fr}
\EGAReferenceAlias{0.2.1.7.fr}{0.2.1.7-fr}
\EGAReferenceAlias{0.2.1.7.es}{0.2.1.7-fr}
\EGAReferenceAlias{0.2.1.7.eq}{0.2.1.7-fr}
\EGAReferenceAlias{eq:0.2.1.7}{0.2.1.7-fr}
\EGAReferenceAlias{0.2.1.8}{0.2.1.8-fr}
\EGAReferenceAlias{0.2.1.8-es}{0.2.1.8-fr}
\EGAReferenceAlias{0.2.1.8.fr}{0.2.1.8-fr}
\EGAReferenceAlias{0.2.1.8.es}{0.2.1.8-fr}
\EGAReferenceAlias{0.2.1.8.eq}{0.2.1.8-fr}
\EGAReferenceAlias{eq:0.2.1.8}{0.2.1.8-fr}
\EGAReferenceAlias{0.2.2}{subsection:0.2.2-fr}
\EGAReferenceAlias{0.2.2-fr}{subsection:0.2.2-fr}
\EGAReferenceAlias{0.2.2-es}{subsection:0.2.2-fr}
\EGAReferenceAlias{0.2.2.fr}{subsection:0.2.2-fr}
\EGAReferenceAlias{0.2.2.es}{subsection:0.2.2-fr}
\EGAReferenceAlias{0.2.2.eq}{subsection:0.2.2-fr}
\EGAReferenceAlias{eq:0.2.2}{subsection:0.2.2-fr}
\EGAReferenceAlias{subsection:0.2.2}{subsection:0.2.2-fr}
\EGAReferenceAlias{subsection:0.2.2-es}{subsection:0.2.2-fr}
\EGAReferenceAlias{subsection:0.2.2.fr}{subsection:0.2.2-fr}
\EGAReferenceAlias{subsection:0.2.2.es}{subsection:0.2.2-fr}
\EGAReferenceAlias{subsection:0.2.2.eq}{subsection:0.2.2-fr}
\EGAReferenceAlias{eq:subsection:0.2.2}{subsection:0.2.2-fr}
\EGAReferenceAlias{0.2.2.1}{0.2.2.1-fr}
\EGAReferenceAlias{0.2.2.1-es}{0.2.2.1-fr}
\EGAReferenceAlias{0.2.2.1.fr}{0.2.2.1-fr}
\EGAReferenceAlias{0.2.2.1.es}{0.2.2.1-fr}
\EGAReferenceAlias{0.2.2.1.eq}{0.2.2.1-fr}
\EGAReferenceAlias{eq:0.2.2.1}{0.2.2.1-fr}
\EGAReferenceAlias{0.2.2.2}{0.2.2.2-fr}
\EGAReferenceAlias{0.2.2.2-es}{0.2.2.2-fr}
\EGAReferenceAlias{0.2.2.2.fr}{0.2.2.2-fr}
\EGAReferenceAlias{0.2.2.2.es}{0.2.2.2-fr}
\EGAReferenceAlias{0.2.2.2.eq}{0.2.2.2-fr}
\EGAReferenceAlias{eq:0.2.2.2}{0.2.2.2-fr}
\EGAReferenceAlias{0.2.2.3}{0.2.2.3-fr}
\EGAReferenceAlias{0.2.2.3-es}{0.2.2.3-fr}
\EGAReferenceAlias{0.2.2.3.fr}{0.2.2.3-fr}
\EGAReferenceAlias{0.2.2.3.es}{0.2.2.3-fr}
\EGAReferenceAlias{0.2.2.3.eq}{0.2.2.3-fr}
\EGAReferenceAlias{eq:0.2.2.3}{0.2.2.3-fr}
\EGAReferenceAlias{0.2.2.4}{0.2.2.4-fr}
\EGAReferenceAlias{0.2.2.4-es}{0.2.2.4-fr}
\EGAReferenceAlias{0.2.2.4.fr}{0.2.2.4-fr}
\EGAReferenceAlias{0.2.2.4.es}{0.2.2.4-fr}
\EGAReferenceAlias{0.2.2.4.eq}{0.2.2.4-fr}
\EGAReferenceAlias{eq:0.2.2.4}{0.2.2.4-fr}
\EGAReferenceAlias{0.2.2.5}{0.2.2.5-fr}
\EGAReferenceAlias{0.2.2.5-es}{0.2.2.5-fr}
\EGAReferenceAlias{0.2.2.5.fr}{0.2.2.5-fr}
\EGAReferenceAlias{0.2.2.5.es}{0.2.2.5-fr}
\EGAReferenceAlias{0.2.2.5.eq}{0.2.2.5-fr}
\EGAReferenceAlias{eq:0.2.2.5}{0.2.2.5-fr}
\EGAReferenceAlias{0.3}{section:0.3-fr}
\EGAReferenceAlias{0.3-fr}{section:0.3-fr}
\EGAReferenceAlias{0.3-es}{section:0.3-fr}
\EGAReferenceAlias{0.3.fr}{section:0.3-fr}
\EGAReferenceAlias{0.3.es}{section:0.3-fr}
\EGAReferenceAlias{0.3.eq}{section:0.3-fr}
\EGAReferenceAlias{eq:0.3}{section:0.3-fr}
\EGAReferenceAlias{section:0.3}{section:0.3-fr}
\EGAReferenceAlias{section:0.3-es}{section:0.3-fr}
\EGAReferenceAlias{section:0.3.fr}{section:0.3-fr}
\EGAReferenceAlias{section:0.3.es}{section:0.3-fr}
\EGAReferenceAlias{section:0.3.eq}{section:0.3-fr}
\EGAReferenceAlias{eq:section:0.3}{section:0.3-fr}
\EGAReferenceAlias{0.3.1}{subsection:0.3.1-fr}
\EGAReferenceAlias{0.3.1-fr}{subsection:0.3.1-fr}
\EGAReferenceAlias{0.3.1-es}{subsection:0.3.1-fr}
\EGAReferenceAlias{0.3.1.fr}{subsection:0.3.1-fr}
\EGAReferenceAlias{0.3.1.es}{subsection:0.3.1-fr}
\EGAReferenceAlias{0.3.1.eq}{subsection:0.3.1-fr}
\EGAReferenceAlias{eq:0.3.1}{subsection:0.3.1-fr}
\EGAReferenceAlias{subsection:0.3.1}{subsection:0.3.1-fr}
\EGAReferenceAlias{subsection:0.3.1-es}{subsection:0.3.1-fr}
\EGAReferenceAlias{subsection:0.3.1.fr}{subsection:0.3.1-fr}
\EGAReferenceAlias{subsection:0.3.1.es}{subsection:0.3.1-fr}
\EGAReferenceAlias{subsection:0.3.1.eq}{subsection:0.3.1-fr}
\EGAReferenceAlias{eq:subsection:0.3.1}{subsection:0.3.1-fr}
\EGAReferenceAlias{0.3.1.1}{0.3.1.1-fr}
\EGAReferenceAlias{0.3.1.1-es}{0.3.1.1-fr}
\EGAReferenceAlias{0.3.1.1.fr}{0.3.1.1-fr}
\EGAReferenceAlias{0.3.1.1.es}{0.3.1.1-fr}
\EGAReferenceAlias{0.3.1.1.eq}{0.3.1.1-fr}
\EGAReferenceAlias{eq:0.3.1.1}{0.3.1.1-fr}
\EGAReferenceAlias{0.3.1.2}{0.3.1.2-fr}
\EGAReferenceAlias{0.3.1.2-es}{0.3.1.2-fr}
\EGAReferenceAlias{0.3.1.2.fr}{0.3.1.2-fr}
\EGAReferenceAlias{0.3.1.2.es}{0.3.1.2-fr}
\EGAReferenceAlias{0.3.1.2.eq}{0.3.1.2-fr}
\EGAReferenceAlias{eq:0.3.1.2}{0.3.1.2-fr}
\EGAReferenceAlias{0.3.1.3}{0.3.1.3-fr}
\EGAReferenceAlias{0.3.1.3-es}{0.3.1.3-fr}
\EGAReferenceAlias{0.3.1.3.fr}{0.3.1.3-fr}
\EGAReferenceAlias{0.3.1.3.es}{0.3.1.3-fr}
\EGAReferenceAlias{0.3.1.3.eq}{0.3.1.3-fr}
\EGAReferenceAlias{eq:0.3.1.3}{0.3.1.3-fr}
\EGAReferenceAlias{0.3.1.4}{0.3.1.4-fr}
\EGAReferenceAlias{0.3.1.4-es}{0.3.1.4-fr}
\EGAReferenceAlias{0.3.1.4.fr}{0.3.1.4-fr}
\EGAReferenceAlias{0.3.1.4.es}{0.3.1.4-fr}
\EGAReferenceAlias{0.3.1.4.eq}{0.3.1.4-fr}
\EGAReferenceAlias{eq:0.3.1.4}{0.3.1.4-fr}
\EGAReferenceAlias{0.3.1.5}{0.3.1.5-fr}
\EGAReferenceAlias{0.3.1.5-es}{0.3.1.5-fr}
\EGAReferenceAlias{0.3.1.5.fr}{0.3.1.5-fr}
\EGAReferenceAlias{0.3.1.5.es}{0.3.1.5-fr}
\EGAReferenceAlias{0.3.1.5.eq}{0.3.1.5-fr}
\EGAReferenceAlias{eq:0.3.1.5}{0.3.1.5-fr}
\EGAReferenceAlias{0.3.1.6}{0.3.1.6-fr}
\EGAReferenceAlias{0.3.1.6-es}{0.3.1.6-fr}
\EGAReferenceAlias{0.3.1.6.fr}{0.3.1.6-fr}
\EGAReferenceAlias{0.3.1.6.es}{0.3.1.6-fr}
\EGAReferenceAlias{0.3.1.6.eq}{0.3.1.6-fr}
\EGAReferenceAlias{eq:0.3.1.6}{0.3.1.6-fr}
\EGAReferenceAlias{0.3.2}{subsection:0.3.2-fr}
\EGAReferenceAlias{0.3.2-fr}{subsection:0.3.2-fr}
\EGAReferenceAlias{0.3.2-es}{subsection:0.3.2-fr}
\EGAReferenceAlias{0.3.2.fr}{subsection:0.3.2-fr}
\EGAReferenceAlias{0.3.2.es}{subsection:0.3.2-fr}
\EGAReferenceAlias{0.3.2.eq}{subsection:0.3.2-fr}
\EGAReferenceAlias{eq:0.3.2}{subsection:0.3.2-fr}
\EGAReferenceAlias{subsection:0.3.2}{subsection:0.3.2-fr}
\EGAReferenceAlias{subsection:0.3.2-es}{subsection:0.3.2-fr}
\EGAReferenceAlias{subsection:0.3.2.fr}{subsection:0.3.2-fr}
\EGAReferenceAlias{subsection:0.3.2.es}{subsection:0.3.2-fr}
\EGAReferenceAlias{subsection:0.3.2.eq}{subsection:0.3.2-fr}
\EGAReferenceAlias{eq:subsection:0.3.2}{subsection:0.3.2-fr}
\EGAReferenceAlias{0.3.2.1}{0.3.2.1-fr}
\EGAReferenceAlias{0.3.2.1-es}{0.3.2.1-fr}
\EGAReferenceAlias{0.3.2.1.fr}{0.3.2.1-fr}
\EGAReferenceAlias{0.3.2.1.es}{0.3.2.1-fr}
\EGAReferenceAlias{0.3.2.1.eq}{0.3.2.1-fr}
\EGAReferenceAlias{eq:0.3.2.1}{0.3.2.1-fr}
\EGAReferenceAlias{0.3.2.2}{0.3.2.2-fr}
\EGAReferenceAlias{0.3.2.2-es}{0.3.2.2-fr}
\EGAReferenceAlias{0.3.2.2.fr}{0.3.2.2-fr}
\EGAReferenceAlias{0.3.2.2.es}{0.3.2.2-fr}
\EGAReferenceAlias{0.3.2.2.eq}{0.3.2.2-fr}
\EGAReferenceAlias{eq:0.3.2.2}{0.3.2.2-fr}
\EGAReferenceAlias{0.3.2.3}{0.3.2.3-fr}
\EGAReferenceAlias{0.3.2.3-es}{0.3.2.3-fr}
\EGAReferenceAlias{0.3.2.3.fr}{0.3.2.3-fr}
\EGAReferenceAlias{0.3.2.3.es}{0.3.2.3-fr}
\EGAReferenceAlias{0.3.2.3.eq}{0.3.2.3-fr}
\EGAReferenceAlias{eq:0.3.2.3}{0.3.2.3-fr}
\EGAReferenceAlias{0.3.2.4}{0.3.2.4-fr}
\EGAReferenceAlias{0.3.2.4-es}{0.3.2.4-fr}
\EGAReferenceAlias{0.3.2.4.fr}{0.3.2.4-fr}
\EGAReferenceAlias{0.3.2.4.es}{0.3.2.4-fr}
\EGAReferenceAlias{0.3.2.4.eq}{0.3.2.4-fr}
\EGAReferenceAlias{eq:0.3.2.4}{0.3.2.4-fr}
\EGAReferenceAlias{0.3.2.5}{0.3.2.5-fr}
\EGAReferenceAlias{0.3.2.5-es}{0.3.2.5-fr}
\EGAReferenceAlias{0.3.2.5.fr}{0.3.2.5-fr}
\EGAReferenceAlias{0.3.2.5.es}{0.3.2.5-fr}
\EGAReferenceAlias{0.3.2.5.eq}{0.3.2.5-fr}
\EGAReferenceAlias{eq:0.3.2.5}{0.3.2.5-fr}
\EGAReferenceAlias{0.3.2.6}{0.3.2.6-fr}
\EGAReferenceAlias{0.3.2.6-es}{0.3.2.6-fr}
\EGAReferenceAlias{0.3.2.6.fr}{0.3.2.6-fr}
\EGAReferenceAlias{0.3.2.6.es}{0.3.2.6-fr}
\EGAReferenceAlias{0.3.2.6.eq}{0.3.2.6-fr}
\EGAReferenceAlias{eq:0.3.2.6}{0.3.2.6-fr}
\EGAReferenceAlias{0.3.3}{subsection:0.3.3-fr}
\EGAReferenceAlias{0.3.3-fr}{subsection:0.3.3-fr}
\EGAReferenceAlias{0.3.3-es}{subsection:0.3.3-fr}
\EGAReferenceAlias{0.3.3.fr}{subsection:0.3.3-fr}
\EGAReferenceAlias{0.3.3.es}{subsection:0.3.3-fr}
\EGAReferenceAlias{0.3.3.eq}{subsection:0.3.3-fr}
\EGAReferenceAlias{eq:0.3.3}{subsection:0.3.3-fr}
\EGAReferenceAlias{subsection:0.3.3}{subsection:0.3.3-fr}
\EGAReferenceAlias{subsection:0.3.3-es}{subsection:0.3.3-fr}
\EGAReferenceAlias{subsection:0.3.3.fr}{subsection:0.3.3-fr}
\EGAReferenceAlias{subsection:0.3.3.es}{subsection:0.3.3-fr}
\EGAReferenceAlias{subsection:0.3.3.eq}{subsection:0.3.3-fr}
\EGAReferenceAlias{eq:subsection:0.3.3}{subsection:0.3.3-fr}
\EGAReferenceAlias{0.3.3.1}{0.3.3.1-fr}
\EGAReferenceAlias{0.3.3.1-es}{0.3.3.1-fr}
\EGAReferenceAlias{0.3.3.1.fr}{0.3.3.1-fr}
\EGAReferenceAlias{0.3.3.1.es}{0.3.3.1-fr}
\EGAReferenceAlias{0.3.3.1.eq}{0.3.3.1-fr}
\EGAReferenceAlias{eq:0.3.3.1}{0.3.3.1-fr}
\EGAReferenceAlias{0.3.3.2}{0.3.3.2-fr}
\EGAReferenceAlias{0.3.3.2-es}{0.3.3.2-fr}
\EGAReferenceAlias{0.3.3.2.fr}{0.3.3.2-fr}
\EGAReferenceAlias{0.3.3.2.es}{0.3.3.2-fr}
\EGAReferenceAlias{0.3.3.2.eq}{0.3.3.2-fr}
\EGAReferenceAlias{eq:0.3.3.2}{0.3.3.2-fr}
\EGAReferenceAlias{0.3.3.2.1}{0.3.3.2.1-fr}
\EGAReferenceAlias{0.3.3.2.1-es}{0.3.3.2.1-fr}
\EGAReferenceAlias{0.3.3.2.1.fr}{0.3.3.2.1-fr}
\EGAReferenceAlias{0.3.3.2.1.es}{0.3.3.2.1-fr}
\EGAReferenceAlias{0.3.3.2.1.eq}{0.3.3.2.1-fr}
\EGAReferenceAlias{eq:0.3.3.2.1}{0.3.3.2.1-fr}
\EGAReferenceAlias{0.3.3.3}{0.3.3.3-fr}
\EGAReferenceAlias{0.3.3.3-es}{0.3.3.3-fr}
\EGAReferenceAlias{0.3.3.3.fr}{0.3.3.3-fr}
\EGAReferenceAlias{0.3.3.3.es}{0.3.3.3-fr}
\EGAReferenceAlias{0.3.3.3.eq}{0.3.3.3-fr}
\EGAReferenceAlias{eq:0.3.3.3}{0.3.3.3-fr}
\EGAReferenceAlias{0.3.4}{subsection:0.3.4-fr}
\EGAReferenceAlias{0.3.4-fr}{subsection:0.3.4-fr}
\EGAReferenceAlias{0.3.4-es}{subsection:0.3.4-fr}
\EGAReferenceAlias{0.3.4.fr}{subsection:0.3.4-fr}
\EGAReferenceAlias{0.3.4.es}{subsection:0.3.4-fr}
\EGAReferenceAlias{0.3.4.eq}{subsection:0.3.4-fr}
\EGAReferenceAlias{eq:0.3.4}{subsection:0.3.4-fr}
\EGAReferenceAlias{subsection:0.3.4}{subsection:0.3.4-fr}
\EGAReferenceAlias{subsection:0.3.4-es}{subsection:0.3.4-fr}
\EGAReferenceAlias{subsection:0.3.4.fr}{subsection:0.3.4-fr}
\EGAReferenceAlias{subsection:0.3.4.es}{subsection:0.3.4-fr}
\EGAReferenceAlias{subsection:0.3.4.eq}{subsection:0.3.4-fr}
\EGAReferenceAlias{eq:subsection:0.3.4}{subsection:0.3.4-fr}
\EGAReferenceAlias{0.3.4.1}{0.3.4.1-fr}
\EGAReferenceAlias{0.3.4.1-es}{0.3.4.1-fr}
\EGAReferenceAlias{0.3.4.1.fr}{0.3.4.1-fr}
\EGAReferenceAlias{0.3.4.1.es}{0.3.4.1-fr}
\EGAReferenceAlias{0.3.4.1.eq}{0.3.4.1-fr}
\EGAReferenceAlias{eq:0.3.4.1}{0.3.4.1-fr}
\EGAReferenceAlias{0.3.4.2}{0.3.4.2-fr}
\EGAReferenceAlias{0.3.4.2-es}{0.3.4.2-fr}
\EGAReferenceAlias{0.3.4.2.fr}{0.3.4.2-fr}
\EGAReferenceAlias{0.3.4.2.es}{0.3.4.2-fr}
\EGAReferenceAlias{0.3.4.2.eq}{0.3.4.2-fr}
\EGAReferenceAlias{eq:0.3.4.2}{0.3.4.2-fr}
\EGAReferenceAlias{0.3.4.3}{0.3.4.3-fr}
\EGAReferenceAlias{0.3.4.3-es}{0.3.4.3-fr}
\EGAReferenceAlias{0.3.4.3.fr}{0.3.4.3-fr}
\EGAReferenceAlias{0.3.4.3.es}{0.3.4.3-fr}
\EGAReferenceAlias{0.3.4.3.eq}{0.3.4.3-fr}
\EGAReferenceAlias{eq:0.3.4.3}{0.3.4.3-fr}
\EGAReferenceAlias{0.3.4.4}{0.3.4.4-fr}
\EGAReferenceAlias{0.3.4.4-es}{0.3.4.4-fr}
\EGAReferenceAlias{0.3.4.4.fr}{0.3.4.4-fr}
\EGAReferenceAlias{0.3.4.4.es}{0.3.4.4-fr}
\EGAReferenceAlias{0.3.4.4.eq}{0.3.4.4-fr}
\EGAReferenceAlias{eq:0.3.4.4}{0.3.4.4-fr}
\EGAReferenceAlias{0.3.4.5}{0.3.4.5-fr}
\EGAReferenceAlias{0.3.4.5-es}{0.3.4.5-fr}
\EGAReferenceAlias{0.3.4.5.fr}{0.3.4.5-fr}
\EGAReferenceAlias{0.3.4.5.es}{0.3.4.5-fr}
\EGAReferenceAlias{0.3.4.5.eq}{0.3.4.5-fr}
\EGAReferenceAlias{eq:0.3.4.5}{0.3.4.5-fr}
\EGAReferenceAlias{0.3.4.6}{0.3.4.6-fr}
\EGAReferenceAlias{0.3.4.6-es}{0.3.4.6-fr}
\EGAReferenceAlias{0.3.4.6.fr}{0.3.4.6-fr}
\EGAReferenceAlias{0.3.4.6.es}{0.3.4.6-fr}
\EGAReferenceAlias{0.3.4.6.eq}{0.3.4.6-fr}
\EGAReferenceAlias{eq:0.3.4.6}{0.3.4.6-fr}
\EGAReferenceAlias{0.3.5.1}{0.3.5.1-fr}
\EGAReferenceAlias{0.3.5.1-es}{0.3.5.1-fr}
\EGAReferenceAlias{0.3.5.1.fr}{0.3.5.1-fr}
\EGAReferenceAlias{0.3.5.1.es}{0.3.5.1-fr}
\EGAReferenceAlias{0.3.5.1.eq}{0.3.5.1-fr}
\EGAReferenceAlias{eq:0.3.5.1}{0.3.5.1-fr}
\EGAReferenceAlias{0.3.5.1.1}{0.3.5.1.1-fr}
\EGAReferenceAlias{0.3.5.1.1-es}{0.3.5.1.1-fr}
\EGAReferenceAlias{0.3.5.1.1.fr}{0.3.5.1.1-fr}
\EGAReferenceAlias{0.3.5.1.1.es}{0.3.5.1.1-fr}
\EGAReferenceAlias{0.3.5.1.1.eq}{0.3.5.1.1-fr}
\EGAReferenceAlias{eq:0.3.5.1.1}{0.3.5.1.1-fr}
\EGAReferenceAlias{0.3.5.2}{0.3.5.2-fr}
\EGAReferenceAlias{0.3.5.2-es}{0.3.5.2-fr}
\EGAReferenceAlias{0.3.5.2.fr}{0.3.5.2-fr}
\EGAReferenceAlias{0.3.5.2.es}{0.3.5.2-fr}
\EGAReferenceAlias{0.3.5.2.eq}{0.3.5.2-fr}
\EGAReferenceAlias{eq:0.3.5.2}{0.3.5.2-fr}
\EGAReferenceAlias{0.3.5.3}{0.3.5.3-fr}
\EGAReferenceAlias{0.3.5.3-es}{0.3.5.3-fr}
\EGAReferenceAlias{0.3.5.3.fr}{0.3.5.3-fr}
\EGAReferenceAlias{0.3.5.3.es}{0.3.5.3-fr}
\EGAReferenceAlias{0.3.5.3.eq}{0.3.5.3-fr}
\EGAReferenceAlias{eq:0.3.5.3}{0.3.5.3-fr}
\EGAReferenceAlias{0.3.5.3.1}{0.3.5.3.1-fr}
\EGAReferenceAlias{0.3.5.3.1-es}{0.3.5.3.1-fr}
\EGAReferenceAlias{0.3.5.3.1.fr}{0.3.5.3.1-fr}
\EGAReferenceAlias{0.3.5.3.1.es}{0.3.5.3.1-fr}
\EGAReferenceAlias{0.3.5.3.1.eq}{0.3.5.3.1-fr}
\EGAReferenceAlias{eq:0.3.5.3.1}{0.3.5.3.1-fr}
\EGAReferenceAlias{0.3.5.3.2}{0.3.5.3.2-fr}
\EGAReferenceAlias{0.3.5.3.2-es}{0.3.5.3.2-fr}
\EGAReferenceAlias{0.3.5.3.2.fr}{0.3.5.3.2-fr}
\EGAReferenceAlias{0.3.5.3.2.es}{0.3.5.3.2-fr}
\EGAReferenceAlias{0.3.5.3.2.eq}{0.3.5.3.2-fr}
\EGAReferenceAlias{eq:0.3.5.3.2}{0.3.5.3.2-fr}
\EGAReferenceAlias{0.3.5.3.3}{0.3.5.3.3-fr}
\EGAReferenceAlias{0.3.5.3.3-es}{0.3.5.3.3-fr}
\EGAReferenceAlias{0.3.5.3.3.fr}{0.3.5.3.3-fr}
\EGAReferenceAlias{0.3.5.3.3.es}{0.3.5.3.3-fr}
\EGAReferenceAlias{0.3.5.3.3.eq}{0.3.5.3.3-fr}
\EGAReferenceAlias{eq:0.3.5.3.3}{0.3.5.3.3-fr}
\EGAReferenceAlias{0.3.5.4}{0.3.5.4-fr}
\EGAReferenceAlias{0.3.5.4-es}{0.3.5.4-fr}
\EGAReferenceAlias{0.3.5.4.fr}{0.3.5.4-fr}
\EGAReferenceAlias{0.3.5.4.es}{0.3.5.4-fr}
\EGAReferenceAlias{0.3.5.4.eq}{0.3.5.4-fr}
\EGAReferenceAlias{eq:0.3.5.4}{0.3.5.4-fr}
\EGAReferenceAlias{0.3.5.4.1}{0.3.5.4.1-fr}
\EGAReferenceAlias{0.3.5.4.1-es}{0.3.5.4.1-fr}
\EGAReferenceAlias{0.3.5.4.1.fr}{0.3.5.4.1-fr}
\EGAReferenceAlias{0.3.5.4.1.es}{0.3.5.4.1-fr}
\EGAReferenceAlias{0.3.5.4.1.eq}{0.3.5.4.1-fr}
\EGAReferenceAlias{eq:0.3.5.4.1}{0.3.5.4.1-fr}
\EGAReferenceAlias{0.3.5.4.2}{0.3.5.4.2-fr}
\EGAReferenceAlias{0.3.5.4.2-es}{0.3.5.4.2-fr}
\EGAReferenceAlias{0.3.5.4.2.fr}{0.3.5.4.2-fr}
\EGAReferenceAlias{0.3.5.4.2.es}{0.3.5.4.2-fr}
\EGAReferenceAlias{0.3.5.4.2.eq}{0.3.5.4.2-fr}
\EGAReferenceAlias{eq:0.3.5.4.2}{0.3.5.4.2-fr}
\EGAReferenceAlias{0.3.5.4.3}{0.3.5.4.3-fr}
\EGAReferenceAlias{0.3.5.4.3-es}{0.3.5.4.3-fr}
\EGAReferenceAlias{0.3.5.4.3.fr}{0.3.5.4.3-fr}
\EGAReferenceAlias{0.3.5.4.3.es}{0.3.5.4.3-fr}
\EGAReferenceAlias{0.3.5.4.3.eq}{0.3.5.4.3-fr}
\EGAReferenceAlias{eq:0.3.5.4.3}{0.3.5.4.3-fr}
\EGAReferenceAlias{0.3.5.4.4}{0.3.5.4.4-fr}
\EGAReferenceAlias{0.3.5.4.4-es}{0.3.5.4.4-fr}
\EGAReferenceAlias{0.3.5.4.4.fr}{0.3.5.4.4-fr}
\EGAReferenceAlias{0.3.5.4.4.es}{0.3.5.4.4-fr}
\EGAReferenceAlias{0.3.5.4.4.eq}{0.3.5.4.4-fr}
\EGAReferenceAlias{eq:0.3.5.4.4}{0.3.5.4.4-fr}
\EGAReferenceAlias{0.3.5.5}{0.3.5.5-fr}
\EGAReferenceAlias{0.3.5.5-es}{0.3.5.5-fr}
\EGAReferenceAlias{0.3.5.5.fr}{0.3.5.5-fr}
\EGAReferenceAlias{0.3.5.5.es}{0.3.5.5-fr}
\EGAReferenceAlias{0.3.5.5.eq}{0.3.5.5-fr}
\EGAReferenceAlias{eq:0.3.5.5}{0.3.5.5-fr}
\EGAReferenceAlias{0.3.5.5.1}{0.3.5.5.1-fr}
\EGAReferenceAlias{0.3.5.5.1-es}{0.3.5.5.1-fr}
\EGAReferenceAlias{0.3.5.5.1.fr}{0.3.5.5.1-fr}
\EGAReferenceAlias{0.3.5.5.1.es}{0.3.5.5.1-fr}
\EGAReferenceAlias{0.3.5.5.1.eq}{0.3.5.5.1-fr}
\EGAReferenceAlias{eq:0.3.5.5.1}{0.3.5.5.1-fr}
\EGAReferenceAlias{0.3.5.6}{0.3.5.6-fr}
\EGAReferenceAlias{0.3.5.6-es}{0.3.5.6-fr}
\EGAReferenceAlias{0.3.5.6.fr}{0.3.5.6-fr}
\EGAReferenceAlias{0.3.5.6.es}{0.3.5.6-fr}
\EGAReferenceAlias{0.3.5.6.eq}{0.3.5.6-fr}
\EGAReferenceAlias{eq:0.3.5.6}{0.3.5.6-fr}
\EGAReferenceAlias{0.3.6.1}{0.3.6.1-fr}
\EGAReferenceAlias{0.3.6.1-es}{0.3.6.1-fr}
\EGAReferenceAlias{0.3.6.1.fr}{0.3.6.1-fr}
\EGAReferenceAlias{0.3.6.1.es}{0.3.6.1-fr}
\EGAReferenceAlias{0.3.6.1.eq}{0.3.6.1-fr}
\EGAReferenceAlias{eq:0.3.6.1}{0.3.6.1-fr}
\EGAReferenceAlias{0.3.6.2}{0.3.6.2-fr}
\EGAReferenceAlias{0.3.6.2-es}{0.3.6.2-fr}
\EGAReferenceAlias{0.3.6.2.fr}{0.3.6.2-fr}
\EGAReferenceAlias{0.3.6.2.es}{0.3.6.2-fr}
\EGAReferenceAlias{0.3.6.2.eq}{0.3.6.2-fr}
\EGAReferenceAlias{eq:0.3.6.2}{0.3.6.2-fr}
\EGAReferenceAlias{0.3.7.1}{0.3.7.1-fr}
\EGAReferenceAlias{0.3.7.1-es}{0.3.7.1-fr}
\EGAReferenceAlias{0.3.7.1.fr}{0.3.7.1-fr}
\EGAReferenceAlias{0.3.7.1.es}{0.3.7.1-fr}
\EGAReferenceAlias{0.3.7.1.eq}{0.3.7.1-fr}
\EGAReferenceAlias{eq:0.3.7.1}{0.3.7.1-fr}
\EGAReferenceAlias{0.3.7.2}{0.3.7.2-fr}
\EGAReferenceAlias{0.3.7.2-es}{0.3.7.2-fr}
\EGAReferenceAlias{0.3.7.2.fr}{0.3.7.2-fr}
\EGAReferenceAlias{0.3.7.2.es}{0.3.7.2-fr}
\EGAReferenceAlias{0.3.7.2.eq}{0.3.7.2-fr}
\EGAReferenceAlias{eq:0.3.7.2}{0.3.7.2-fr}
\EGAReferenceAlias{0.3.8.1}{0.3.8.1-fr}
\EGAReferenceAlias{0.3.8.1-es}{0.3.8.1-fr}
\EGAReferenceAlias{0.3.8.1.fr}{0.3.8.1-fr}
\EGAReferenceAlias{0.3.8.1.es}{0.3.8.1-fr}
\EGAReferenceAlias{0.3.8.1.eq}{0.3.8.1-fr}
\EGAReferenceAlias{eq:0.3.8.1}{0.3.8.1-fr}
\EGAReferenceAlias{0.3.8.2}{0.3.8.2-fr}
\EGAReferenceAlias{0.3.8.2-es}{0.3.8.2-fr}
\EGAReferenceAlias{0.3.8.2.fr}{0.3.8.2-fr}
\EGAReferenceAlias{0.3.8.2.es}{0.3.8.2-fr}
\EGAReferenceAlias{0.3.8.2.eq}{0.3.8.2-fr}
\EGAReferenceAlias{eq:0.3.8.2}{0.3.8.2-fr}
\EGAReferenceAlias{0.3.8.3}{0.3.8.3-fr}
\EGAReferenceAlias{0.3.8.3-es}{0.3.8.3-fr}
\EGAReferenceAlias{0.3.8.3.fr}{0.3.8.3-fr}
\EGAReferenceAlias{0.3.8.3.es}{0.3.8.3-fr}
\EGAReferenceAlias{0.3.8.3.eq}{0.3.8.3-fr}
\EGAReferenceAlias{eq:0.3.8.3}{0.3.8.3-fr}
\EGAReferenceAlias{0.4}{section:0.4-fr}
\EGAReferenceAlias{0.4-fr}{section:0.4-fr}
\EGAReferenceAlias{0.4-es}{section:0.4-fr}
\EGAReferenceAlias{0.4.fr}{section:0.4-fr}
\EGAReferenceAlias{0.4.es}{section:0.4-fr}
\EGAReferenceAlias{0.4.eq}{section:0.4-fr}
\EGAReferenceAlias{eq:0.4}{section:0.4-fr}
\EGAReferenceAlias{section:0.4}{section:0.4-fr}
\EGAReferenceAlias{section:0.4-es}{section:0.4-fr}
\EGAReferenceAlias{section:0.4.fr}{section:0.4-fr}
\EGAReferenceAlias{section:0.4.es}{section:0.4-fr}
\EGAReferenceAlias{section:0.4.eq}{section:0.4-fr}
\EGAReferenceAlias{eq:section:0.4}{section:0.4-fr}
\EGAReferenceAlias{0.4.1}{subsection:0.4.1-fr}
\EGAReferenceAlias{0.4.1-fr}{subsection:0.4.1-fr}
\EGAReferenceAlias{0.4.1-es}{subsection:0.4.1-fr}
\EGAReferenceAlias{0.4.1.fr}{subsection:0.4.1-fr}
\EGAReferenceAlias{0.4.1.es}{subsection:0.4.1-fr}
\EGAReferenceAlias{0.4.1.eq}{subsection:0.4.1-fr}
\EGAReferenceAlias{eq:0.4.1}{subsection:0.4.1-fr}
\EGAReferenceAlias{subsection:0.4.1}{subsection:0.4.1-fr}
\EGAReferenceAlias{subsection:0.4.1-es}{subsection:0.4.1-fr}
\EGAReferenceAlias{subsection:0.4.1.fr}{subsection:0.4.1-fr}
\EGAReferenceAlias{subsection:0.4.1.es}{subsection:0.4.1-fr}
\EGAReferenceAlias{subsection:0.4.1.eq}{subsection:0.4.1-fr}
\EGAReferenceAlias{eq:subsection:0.4.1}{subsection:0.4.1-fr}
\EGAReferenceAlias{0.4.1.1}{0.4.1.1-fr}
\EGAReferenceAlias{0.4.1.1-es}{0.4.1.1-fr}
\EGAReferenceAlias{0.4.1.1.fr}{0.4.1.1-fr}
\EGAReferenceAlias{0.4.1.1.es}{0.4.1.1-fr}
\EGAReferenceAlias{0.4.1.1.eq}{0.4.1.1-fr}
\EGAReferenceAlias{eq:0.4.1.1}{0.4.1.1-fr}
\EGAReferenceAlias{0.4.1.2}{0.4.1.2-fr}
\EGAReferenceAlias{0.4.1.2-es}{0.4.1.2-fr}
\EGAReferenceAlias{0.4.1.2.fr}{0.4.1.2-fr}
\EGAReferenceAlias{0.4.1.2.es}{0.4.1.2-fr}
\EGAReferenceAlias{0.4.1.2.eq}{0.4.1.2-fr}
\EGAReferenceAlias{eq:0.4.1.2}{0.4.1.2-fr}
\EGAReferenceAlias{0.4.1.3}{0.4.1.3-fr}
\EGAReferenceAlias{0.4.1.3-es}{0.4.1.3-fr}
\EGAReferenceAlias{0.4.1.3.fr}{0.4.1.3-fr}
\EGAReferenceAlias{0.4.1.3.es}{0.4.1.3-fr}
\EGAReferenceAlias{0.4.1.3.eq}{0.4.1.3-fr}
\EGAReferenceAlias{eq:0.4.1.3}{0.4.1.3-fr}
\EGAReferenceAlias{0.4.1.4}{0.4.1.4-fr}
\EGAReferenceAlias{0.4.1.4-es}{0.4.1.4-fr}
\EGAReferenceAlias{0.4.1.4.fr}{0.4.1.4-fr}
\EGAReferenceAlias{0.4.1.4.es}{0.4.1.4-fr}
\EGAReferenceAlias{0.4.1.4.eq}{0.4.1.4-fr}
\EGAReferenceAlias{eq:0.4.1.4}{0.4.1.4-fr}
\EGAReferenceAlias{0.4.1.5}{0.4.1.5-fr}
\EGAReferenceAlias{0.4.1.5-es}{0.4.1.5-fr}
\EGAReferenceAlias{0.4.1.5.fr}{0.4.1.5-fr}
\EGAReferenceAlias{0.4.1.5.es}{0.4.1.5-fr}
\EGAReferenceAlias{0.4.1.5.eq}{0.4.1.5-fr}
\EGAReferenceAlias{eq:0.4.1.5}{0.4.1.5-fr}
\EGAReferenceAlias{0.4.1.6}{0.4.1.6-fr}
\EGAReferenceAlias{0.4.1.6-es}{0.4.1.6-fr}
\EGAReferenceAlias{0.4.1.6.fr}{0.4.1.6-fr}
\EGAReferenceAlias{0.4.1.6.es}{0.4.1.6-fr}
\EGAReferenceAlias{0.4.1.6.eq}{0.4.1.6-fr}
\EGAReferenceAlias{eq:0.4.1.6}{0.4.1.6-fr}
\EGAReferenceAlias{0.4.1.7}{0.4.1.7-fr}
\EGAReferenceAlias{0.4.1.7-es}{0.4.1.7-fr}
\EGAReferenceAlias{0.4.1.7.fr}{0.4.1.7-fr}
\EGAReferenceAlias{0.4.1.7.es}{0.4.1.7-fr}
\EGAReferenceAlias{0.4.1.7.eq}{0.4.1.7-fr}
\EGAReferenceAlias{eq:0.4.1.7}{0.4.1.7-fr}
\EGAReferenceAlias{0.4.2}{subsection:0.4.2-fr}
\EGAReferenceAlias{0.4.2-fr}{subsection:0.4.2-fr}
\EGAReferenceAlias{0.4.2-es}{subsection:0.4.2-fr}
\EGAReferenceAlias{0.4.2.fr}{subsection:0.4.2-fr}
\EGAReferenceAlias{0.4.2.es}{subsection:0.4.2-fr}
\EGAReferenceAlias{0.4.2.eq}{subsection:0.4.2-fr}
\EGAReferenceAlias{eq:0.4.2}{subsection:0.4.2-fr}
\EGAReferenceAlias{subsection:0.4.2}{subsection:0.4.2-fr}
\EGAReferenceAlias{subsection:0.4.2-es}{subsection:0.4.2-fr}
\EGAReferenceAlias{subsection:0.4.2.fr}{subsection:0.4.2-fr}
\EGAReferenceAlias{subsection:0.4.2.es}{subsection:0.4.2-fr}
\EGAReferenceAlias{subsection:0.4.2.eq}{subsection:0.4.2-fr}
\EGAReferenceAlias{eq:subsection:0.4.2}{subsection:0.4.2-fr}
\EGAReferenceAlias{0.4.2.1}{0.4.2.1-fr}
\EGAReferenceAlias{0.4.2.1-es}{0.4.2.1-fr}
\EGAReferenceAlias{0.4.2.1.fr}{0.4.2.1-fr}
\EGAReferenceAlias{0.4.2.1.es}{0.4.2.1-fr}
\EGAReferenceAlias{0.4.2.1.eq}{0.4.2.1-fr}
\EGAReferenceAlias{eq:0.4.2.1}{0.4.2.1-fr}
\EGAReferenceAlias{0.4.2.2}{0.4.2.2-fr}
\EGAReferenceAlias{0.4.2.2-es}{0.4.2.2-fr}
\EGAReferenceAlias{0.4.2.2.fr}{0.4.2.2-fr}
\EGAReferenceAlias{0.4.2.2.es}{0.4.2.2-fr}
\EGAReferenceAlias{0.4.2.2.eq}{0.4.2.2-fr}
\EGAReferenceAlias{eq:0.4.2.2}{0.4.2.2-fr}
\EGAReferenceAlias{0.4.2.2.1}{0.4.2.2.1-fr}
\EGAReferenceAlias{0.4.2.2.1-es}{0.4.2.2.1-fr}
\EGAReferenceAlias{0.4.2.2.1.fr}{0.4.2.2.1-fr}
\EGAReferenceAlias{0.4.2.2.1.es}{0.4.2.2.1-fr}
\EGAReferenceAlias{0.4.2.2.1.eq}{0.4.2.2.1-fr}
\EGAReferenceAlias{eq:0.4.2.2.1}{0.4.2.2.1-fr}
\EGAReferenceAlias{0.4.2.2.2}{0.4.2.2.2-fr}
\EGAReferenceAlias{0.4.2.2.2-es}{0.4.2.2.2-fr}
\EGAReferenceAlias{0.4.2.2.2.fr}{0.4.2.2.2-fr}
\EGAReferenceAlias{0.4.2.2.2.es}{0.4.2.2.2-fr}
\EGAReferenceAlias{0.4.2.2.2.eq}{0.4.2.2.2-fr}
\EGAReferenceAlias{eq:0.4.2.2.2}{0.4.2.2.2-fr}
\EGAReferenceAlias{0.4.2.3}{0.4.2.3-fr}
\EGAReferenceAlias{0.4.2.3-es}{0.4.2.3-fr}
\EGAReferenceAlias{0.4.2.3.fr}{0.4.2.3-fr}
\EGAReferenceAlias{0.4.2.3.es}{0.4.2.3-fr}
\EGAReferenceAlias{0.4.2.3.eq}{0.4.2.3-fr}
\EGAReferenceAlias{eq:0.4.2.3}{0.4.2.3-fr}
\EGAReferenceAlias{0.4.2.3.1}{0.4.2.3.1-fr}
\EGAReferenceAlias{0.4.2.3.1-es}{0.4.2.3.1-fr}
\EGAReferenceAlias{0.4.2.3.1.fr}{0.4.2.3.1-fr}
\EGAReferenceAlias{0.4.2.3.1.es}{0.4.2.3.1-fr}
\EGAReferenceAlias{0.4.2.3.1.eq}{0.4.2.3.1-fr}
\EGAReferenceAlias{eq:0.4.2.3.1}{0.4.2.3.1-fr}
\EGAReferenceAlias{0.4.2.4}{0.4.2.4-fr}
\EGAReferenceAlias{0.4.2.4-es}{0.4.2.4-fr}
\EGAReferenceAlias{0.4.2.4.fr}{0.4.2.4-fr}
\EGAReferenceAlias{0.4.2.4.es}{0.4.2.4-fr}
\EGAReferenceAlias{0.4.2.4.eq}{0.4.2.4-fr}
\EGAReferenceAlias{eq:0.4.2.4}{0.4.2.4-fr}
\EGAReferenceAlias{0.4.2.5}{0.4.2.5-fr}
\EGAReferenceAlias{0.4.2.5-es}{0.4.2.5-fr}
\EGAReferenceAlias{0.4.2.5.fr}{0.4.2.5-fr}
\EGAReferenceAlias{0.4.2.5.es}{0.4.2.5-fr}
\EGAReferenceAlias{0.4.2.5.eq}{0.4.2.5-fr}
\EGAReferenceAlias{eq:0.4.2.5}{0.4.2.5-fr}
\EGAReferenceAlias{0.4.2.6}{0.4.2.6-fr}
\EGAReferenceAlias{0.4.2.6-es}{0.4.2.6-fr}
\EGAReferenceAlias{0.4.2.6.fr}{0.4.2.6-fr}
\EGAReferenceAlias{0.4.2.6.es}{0.4.2.6-fr}
\EGAReferenceAlias{0.4.2.6.eq}{0.4.2.6-fr}
\EGAReferenceAlias{eq:0.4.2.6}{0.4.2.6-fr}
\EGAReferenceAlias{0.4.3}{subsection:0.4.3-fr}
\EGAReferenceAlias{0.4.3-fr}{subsection:0.4.3-fr}
\EGAReferenceAlias{0.4.3-es}{subsection:0.4.3-fr}
\EGAReferenceAlias{0.4.3.fr}{subsection:0.4.3-fr}
\EGAReferenceAlias{0.4.3.es}{subsection:0.4.3-fr}
\EGAReferenceAlias{0.4.3.eq}{subsection:0.4.3-fr}
\EGAReferenceAlias{eq:0.4.3}{subsection:0.4.3-fr}
\EGAReferenceAlias{subsection:0.4.3}{subsection:0.4.3-fr}
\EGAReferenceAlias{subsection:0.4.3-es}{subsection:0.4.3-fr}
\EGAReferenceAlias{subsection:0.4.3.fr}{subsection:0.4.3-fr}
\EGAReferenceAlias{subsection:0.4.3.es}{subsection:0.4.3-fr}
\EGAReferenceAlias{subsection:0.4.3.eq}{subsection:0.4.3-fr}
\EGAReferenceAlias{eq:subsection:0.4.3}{subsection:0.4.3-fr}
\EGAReferenceAlias{0.4.3.1}{0.4.3.1-fr}
\EGAReferenceAlias{0.4.3.1-es}{0.4.3.1-fr}
\EGAReferenceAlias{0.4.3.1.fr}{0.4.3.1-fr}
\EGAReferenceAlias{0.4.3.1.es}{0.4.3.1-fr}
\EGAReferenceAlias{0.4.3.1.eq}{0.4.3.1-fr}
\EGAReferenceAlias{eq:0.4.3.1}{0.4.3.1-fr}
\EGAReferenceAlias{0.4.3.2}{0.4.3.2-fr}
\EGAReferenceAlias{0.4.3.2-es}{0.4.3.2-fr}
\EGAReferenceAlias{0.4.3.2.fr}{0.4.3.2-fr}
\EGAReferenceAlias{0.4.3.2.es}{0.4.3.2-fr}
\EGAReferenceAlias{0.4.3.2.eq}{0.4.3.2-fr}
\EGAReferenceAlias{eq:0.4.3.2}{0.4.3.2-fr}
\EGAReferenceAlias{0.4.3.2.1}{0.4.3.2.1-fr}
\EGAReferenceAlias{0.4.3.2.1-es}{0.4.3.2.1-fr}
\EGAReferenceAlias{0.4.3.2.1.fr}{0.4.3.2.1-fr}
\EGAReferenceAlias{0.4.3.2.1.es}{0.4.3.2.1-fr}
\EGAReferenceAlias{0.4.3.2.1.eq}{0.4.3.2.1-fr}
\EGAReferenceAlias{eq:0.4.3.2.1}{0.4.3.2.1-fr}
\EGAReferenceAlias{0.4.3.3.1}{0.4.3.3.1-fr}
\EGAReferenceAlias{0.4.3.3.1-es}{0.4.3.3.1-fr}
\EGAReferenceAlias{0.4.3.3.1.fr}{0.4.3.3.1-fr}
\EGAReferenceAlias{0.4.3.3.1.es}{0.4.3.3.1-fr}
\EGAReferenceAlias{0.4.3.3.1.eq}{0.4.3.3.1-fr}
\EGAReferenceAlias{eq:0.4.3.3.1}{0.4.3.3.1-fr}
\EGAReferenceAlias{0.4.3.4}{0.4.3.4-fr}
\EGAReferenceAlias{0.4.3.4-es}{0.4.3.4-fr}
\EGAReferenceAlias{0.4.3.4.fr}{0.4.3.4-fr}
\EGAReferenceAlias{0.4.3.4.es}{0.4.3.4-fr}
\EGAReferenceAlias{0.4.3.4.eq}{0.4.3.4-fr}
\EGAReferenceAlias{eq:0.4.3.4}{0.4.3.4-fr}
\EGAReferenceAlias{0.4.3.5}{0.4.3.5-fr}
\EGAReferenceAlias{0.4.3.5-es}{0.4.3.5-fr}
\EGAReferenceAlias{0.4.3.5.fr}{0.4.3.5-fr}
\EGAReferenceAlias{0.4.3.5.es}{0.4.3.5-fr}
\EGAReferenceAlias{0.4.3.5.eq}{0.4.3.5-fr}
\EGAReferenceAlias{eq:0.4.3.5}{0.4.3.5-fr}
\EGAReferenceAlias{0.4.3.6}{0.4.3.6-fr}
\EGAReferenceAlias{0.4.3.6-es}{0.4.3.6-fr}
\EGAReferenceAlias{0.4.3.6.fr}{0.4.3.6-fr}
\EGAReferenceAlias{0.4.3.6.es}{0.4.3.6-fr}
\EGAReferenceAlias{0.4.3.6.eq}{0.4.3.6-fr}
\EGAReferenceAlias{eq:0.4.3.6}{0.4.3.6-fr}
\EGAReferenceAlias{0.4.4}{subsection:0.4.4-fr}
\EGAReferenceAlias{0.4.4-fr}{subsection:0.4.4-fr}
\EGAReferenceAlias{0.4.4-es}{subsection:0.4.4-fr}
\EGAReferenceAlias{0.4.4.fr}{subsection:0.4.4-fr}
\EGAReferenceAlias{0.4.4.es}{subsection:0.4.4-fr}
\EGAReferenceAlias{0.4.4.eq}{subsection:0.4.4-fr}
\EGAReferenceAlias{eq:0.4.4}{subsection:0.4.4-fr}
\EGAReferenceAlias{subsection:0.4.4}{subsection:0.4.4-fr}
\EGAReferenceAlias{subsection:0.4.4-es}{subsection:0.4.4-fr}
\EGAReferenceAlias{subsection:0.4.4.fr}{subsection:0.4.4-fr}
\EGAReferenceAlias{subsection:0.4.4.es}{subsection:0.4.4-fr}
\EGAReferenceAlias{subsection:0.4.4.eq}{subsection:0.4.4-fr}
\EGAReferenceAlias{eq:subsection:0.4.4}{subsection:0.4.4-fr}
\EGAReferenceAlias{0.4.4.1}{0.4.4.1-fr}
\EGAReferenceAlias{0.4.4.1-es}{0.4.4.1-fr}
\EGAReferenceAlias{0.4.4.1.fr}{0.4.4.1-fr}
\EGAReferenceAlias{0.4.4.1.es}{0.4.4.1-fr}
\EGAReferenceAlias{0.4.4.1.eq}{0.4.4.1-fr}
\EGAReferenceAlias{eq:0.4.4.1}{0.4.4.1-fr}
\EGAReferenceAlias{0.4.4.2}{0.4.4.2-fr}
\EGAReferenceAlias{0.4.4.2-es}{0.4.4.2-fr}
\EGAReferenceAlias{0.4.4.2.fr}{0.4.4.2-fr}
\EGAReferenceAlias{0.4.4.2.es}{0.4.4.2-fr}
\EGAReferenceAlias{0.4.4.2.eq}{0.4.4.2-fr}
\EGAReferenceAlias{eq:0.4.4.2}{0.4.4.2-fr}
\EGAReferenceAlias{0.4.4.3}{0.4.4.3-fr}
\EGAReferenceAlias{0.4.4.3-es}{0.4.4.3-fr}
\EGAReferenceAlias{0.4.4.3.fr}{0.4.4.3-fr}
\EGAReferenceAlias{0.4.4.3.es}{0.4.4.3-fr}
\EGAReferenceAlias{0.4.4.3.eq}{0.4.4.3-fr}
\EGAReferenceAlias{eq:0.4.4.3}{0.4.4.3-fr}
\EGAReferenceAlias{0.4.4.3.1}{0.4.4.3.1-fr}
\EGAReferenceAlias{0.4.4.3.1-es}{0.4.4.3.1-fr}
\EGAReferenceAlias{0.4.4.3.1.fr}{0.4.4.3.1-fr}
\EGAReferenceAlias{0.4.4.3.1.es}{0.4.4.3.1-fr}
\EGAReferenceAlias{0.4.4.3.1.eq}{0.4.4.3.1-fr}
\EGAReferenceAlias{eq:0.4.4.3.1}{0.4.4.3.1-fr}
\EGAReferenceAlias{0.4.4.3.2}{0.4.4.3.2-fr}
\EGAReferenceAlias{0.4.4.3.2-es}{0.4.4.3.2-fr}
\EGAReferenceAlias{0.4.4.3.2.fr}{0.4.4.3.2-fr}
\EGAReferenceAlias{0.4.4.3.2.es}{0.4.4.3.2-fr}
\EGAReferenceAlias{0.4.4.3.2.eq}{0.4.4.3.2-fr}
\EGAReferenceAlias{eq:0.4.4.3.2}{0.4.4.3.2-fr}
\EGAReferenceAlias{0.4.4.3.3}{0.4.4.3.3-fr}
\EGAReferenceAlias{0.4.4.3.3-es}{0.4.4.3.3-fr}
\EGAReferenceAlias{0.4.4.3.3.fr}{0.4.4.3.3-fr}
\EGAReferenceAlias{0.4.4.3.3.es}{0.4.4.3.3-fr}
\EGAReferenceAlias{0.4.4.3.3.eq}{0.4.4.3.3-fr}
\EGAReferenceAlias{eq:0.4.4.3.3}{0.4.4.3.3-fr}
\EGAReferenceAlias{0.4.4.4}{0.4.4.4-fr}
\EGAReferenceAlias{0.4.4.4-es}{0.4.4.4-fr}
\EGAReferenceAlias{0.4.4.4.fr}{0.4.4.4-fr}
\EGAReferenceAlias{0.4.4.4.es}{0.4.4.4-fr}
\EGAReferenceAlias{0.4.4.4.eq}{0.4.4.4-fr}
\EGAReferenceAlias{eq:0.4.4.4}{0.4.4.4-fr}
\EGAReferenceAlias{0.4.4.5}{0.4.4.5-fr}
\EGAReferenceAlias{0.4.4.5-es}{0.4.4.5-fr}
\EGAReferenceAlias{0.4.4.5.fr}{0.4.4.5-fr}
\EGAReferenceAlias{0.4.4.5.es}{0.4.4.5-fr}
\EGAReferenceAlias{0.4.4.5.eq}{0.4.4.5-fr}
\EGAReferenceAlias{eq:0.4.4.5}{0.4.4.5-fr}
\EGAReferenceAlias{0.4.4.6}{0.4.4.6-fr}
\EGAReferenceAlias{0.4.4.6-es}{0.4.4.6-fr}
\EGAReferenceAlias{0.4.4.6.fr}{0.4.4.6-fr}
\EGAReferenceAlias{0.4.4.6.es}{0.4.4.6-fr}
\EGAReferenceAlias{0.4.4.6.eq}{0.4.4.6-fr}
\EGAReferenceAlias{eq:0.4.4.6}{0.4.4.6-fr}
\EGAReferenceAlias{0.4.4.7}{0.4.4.7-fr}
\EGAReferenceAlias{0.4.4.7-es}{0.4.4.7-fr}
\EGAReferenceAlias{0.4.4.7.fr}{0.4.4.7-fr}
\EGAReferenceAlias{0.4.4.7.es}{0.4.4.7-fr}
\EGAReferenceAlias{0.4.4.7.eq}{0.4.4.7-fr}
\EGAReferenceAlias{eq:0.4.4.7}{0.4.4.7-fr}
\EGAReferenceAlias{0.4.4.8}{0.4.4.8-fr}
\EGAReferenceAlias{0.4.4.8-es}{0.4.4.8-fr}
\EGAReferenceAlias{0.4.4.8.fr}{0.4.4.8-fr}
\EGAReferenceAlias{0.4.4.8.es}{0.4.4.8-fr}
\EGAReferenceAlias{0.4.4.8.eq}{0.4.4.8-fr}
\EGAReferenceAlias{eq:0.4.4.8}{0.4.4.8-fr}
\EGAReferenceAlias{0.5}{section:0.5-fr}
\EGAReferenceAlias{0.5-fr}{section:0.5-fr}
\EGAReferenceAlias{0.5-es}{section:0.5-fr}
\EGAReferenceAlias{0.5.fr}{section:0.5-fr}
\EGAReferenceAlias{0.5.es}{section:0.5-fr}
\EGAReferenceAlias{0.5.eq}{section:0.5-fr}
\EGAReferenceAlias{eq:0.5}{section:0.5-fr}
\EGAReferenceAlias{section:0.5}{section:0.5-fr}
\EGAReferenceAlias{section:0.5-es}{section:0.5-fr}
\EGAReferenceAlias{section:0.5.fr}{section:0.5-fr}
\EGAReferenceAlias{section:0.5.es}{section:0.5-fr}
\EGAReferenceAlias{section:0.5.eq}{section:0.5-fr}
\EGAReferenceAlias{eq:section:0.5}{section:0.5-fr}
\EGAReferenceAlias{0.5.1}{subsection:0.5.1-fr}
\EGAReferenceAlias{0.5.1-fr}{subsection:0.5.1-fr}
\EGAReferenceAlias{0.5.1-es}{subsection:0.5.1-fr}
\EGAReferenceAlias{0.5.1.fr}{subsection:0.5.1-fr}
\EGAReferenceAlias{0.5.1.es}{subsection:0.5.1-fr}
\EGAReferenceAlias{0.5.1.eq}{subsection:0.5.1-fr}
\EGAReferenceAlias{eq:0.5.1}{subsection:0.5.1-fr}
\EGAReferenceAlias{subsection:0.5.1}{subsection:0.5.1-fr}
\EGAReferenceAlias{subsection:0.5.1-es}{subsection:0.5.1-fr}
\EGAReferenceAlias{subsection:0.5.1.fr}{subsection:0.5.1-fr}
\EGAReferenceAlias{subsection:0.5.1.es}{subsection:0.5.1-fr}
\EGAReferenceAlias{subsection:0.5.1.eq}{subsection:0.5.1-fr}
\EGAReferenceAlias{eq:subsection:0.5.1}{subsection:0.5.1-fr}
\EGAReferenceAlias{0.5.1.1}{0.5.1.1-fr}
\EGAReferenceAlias{0.5.1.1-es}{0.5.1.1-fr}
\EGAReferenceAlias{0.5.1.1.fr}{0.5.1.1-fr}
\EGAReferenceAlias{0.5.1.1.es}{0.5.1.1-fr}
\EGAReferenceAlias{0.5.1.1.eq}{0.5.1.1-fr}
\EGAReferenceAlias{eq:0.5.1.1}{0.5.1.1-fr}
\EGAReferenceAlias{0.5.1.2}{0.5.1.2-fr}
\EGAReferenceAlias{0.5.1.2-es}{0.5.1.2-fr}
\EGAReferenceAlias{0.5.1.2.fr}{0.5.1.2-fr}
\EGAReferenceAlias{0.5.1.2.es}{0.5.1.2-fr}
\EGAReferenceAlias{0.5.1.2.eq}{0.5.1.2-fr}
\EGAReferenceAlias{eq:0.5.1.2}{0.5.1.2-fr}
\EGAReferenceAlias{0.5.1.3}{0.5.1.3-fr}
\EGAReferenceAlias{0.5.1.3-es}{0.5.1.3-fr}
\EGAReferenceAlias{0.5.1.3.fr}{0.5.1.3-fr}
\EGAReferenceAlias{0.5.1.3.es}{0.5.1.3-fr}
\EGAReferenceAlias{0.5.1.3.eq}{0.5.1.3-fr}
\EGAReferenceAlias{eq:0.5.1.3}{0.5.1.3-fr}
\EGAReferenceAlias{0.5.1.4}{0.5.1.4-fr}
\EGAReferenceAlias{0.5.1.4-es}{0.5.1.4-fr}
\EGAReferenceAlias{0.5.1.4.fr}{0.5.1.4-fr}
\EGAReferenceAlias{0.5.1.4.es}{0.5.1.4-fr}
\EGAReferenceAlias{0.5.1.4.eq}{0.5.1.4-fr}
\EGAReferenceAlias{eq:0.5.1.4}{0.5.1.4-fr}
\EGAReferenceAlias{0.5.2}{subsection:0.5.2-fr}
\EGAReferenceAlias{0.5.2-fr}{subsection:0.5.2-fr}
\EGAReferenceAlias{0.5.2-es}{subsection:0.5.2-fr}
\EGAReferenceAlias{0.5.2.fr}{subsection:0.5.2-fr}
\EGAReferenceAlias{0.5.2.es}{subsection:0.5.2-fr}
\EGAReferenceAlias{0.5.2.eq}{subsection:0.5.2-fr}
\EGAReferenceAlias{eq:0.5.2}{subsection:0.5.2-fr}
\EGAReferenceAlias{subsection:0.5.2}{subsection:0.5.2-fr}
\EGAReferenceAlias{subsection:0.5.2-es}{subsection:0.5.2-fr}
\EGAReferenceAlias{subsection:0.5.2.fr}{subsection:0.5.2-fr}
\EGAReferenceAlias{subsection:0.5.2.es}{subsection:0.5.2-fr}
\EGAReferenceAlias{subsection:0.5.2.eq}{subsection:0.5.2-fr}
\EGAReferenceAlias{eq:subsection:0.5.2}{subsection:0.5.2-fr}
\EGAReferenceAlias{0.5.2.1}{0.5.2.1-fr}
\EGAReferenceAlias{0.5.2.1-es}{0.5.2.1-fr}
\EGAReferenceAlias{0.5.2.1.fr}{0.5.2.1-fr}
\EGAReferenceAlias{0.5.2.1.es}{0.5.2.1-fr}
\EGAReferenceAlias{0.5.2.1.eq}{0.5.2.1-fr}
\EGAReferenceAlias{eq:0.5.2.1}{0.5.2.1-fr}
\EGAReferenceAlias{0.5.2.2}{0.5.2.2-fr}
\EGAReferenceAlias{0.5.2.2-es}{0.5.2.2-fr}
\EGAReferenceAlias{0.5.2.2.fr}{0.5.2.2-fr}
\EGAReferenceAlias{0.5.2.2.es}{0.5.2.2-fr}
\EGAReferenceAlias{0.5.2.2.eq}{0.5.2.2-fr}
\EGAReferenceAlias{eq:0.5.2.2}{0.5.2.2-fr}
\EGAReferenceAlias{0.5.2.3}{0.5.2.3-fr}
\EGAReferenceAlias{0.5.2.3-es}{0.5.2.3-fr}
\EGAReferenceAlias{0.5.2.3.fr}{0.5.2.3-fr}
\EGAReferenceAlias{0.5.2.3.es}{0.5.2.3-fr}
\EGAReferenceAlias{0.5.2.3.eq}{0.5.2.3-fr}
\EGAReferenceAlias{eq:0.5.2.3}{0.5.2.3-fr}
\EGAReferenceAlias{0.5.2.4}{0.5.2.4-fr}
\EGAReferenceAlias{0.5.2.4-es}{0.5.2.4-fr}
\EGAReferenceAlias{0.5.2.4.fr}{0.5.2.4-fr}
\EGAReferenceAlias{0.5.2.4.es}{0.5.2.4-fr}
\EGAReferenceAlias{0.5.2.4.eq}{0.5.2.4-fr}
\EGAReferenceAlias{eq:0.5.2.4}{0.5.2.4-fr}
\EGAReferenceAlias{0.5.2.5}{0.5.2.5-fr}
\EGAReferenceAlias{0.5.2.5-es}{0.5.2.5-fr}
\EGAReferenceAlias{0.5.2.5.fr}{0.5.2.5-fr}
\EGAReferenceAlias{0.5.2.5.es}{0.5.2.5-fr}
\EGAReferenceAlias{0.5.2.5.eq}{0.5.2.5-fr}
\EGAReferenceAlias{eq:0.5.2.5}{0.5.2.5-fr}
\EGAReferenceAlias{0.5.2.6}{0.5.2.6-fr}
\EGAReferenceAlias{0.5.2.6-es}{0.5.2.6-fr}
\EGAReferenceAlias{0.5.2.6.fr}{0.5.2.6-fr}
\EGAReferenceAlias{0.5.2.6.es}{0.5.2.6-fr}
\EGAReferenceAlias{0.5.2.6.eq}{0.5.2.6-fr}
\EGAReferenceAlias{eq:0.5.2.6}{0.5.2.6-fr}
\EGAReferenceAlias{0.5.2.7}{0.5.2.7-fr}
\EGAReferenceAlias{0.5.2.7-es}{0.5.2.7-fr}
\EGAReferenceAlias{0.5.2.7.fr}{0.5.2.7-fr}
\EGAReferenceAlias{0.5.2.7.es}{0.5.2.7-fr}
\EGAReferenceAlias{0.5.2.7.eq}{0.5.2.7-fr}
\EGAReferenceAlias{eq:0.5.2.7}{0.5.2.7-fr}
\EGAReferenceAlias{0.5.3}{subsection:0.5.3-fr}
\EGAReferenceAlias{0.5.3-fr}{subsection:0.5.3-fr}
\EGAReferenceAlias{0.5.3-es}{subsection:0.5.3-fr}
\EGAReferenceAlias{0.5.3.fr}{subsection:0.5.3-fr}
\EGAReferenceAlias{0.5.3.es}{subsection:0.5.3-fr}
\EGAReferenceAlias{0.5.3.eq}{subsection:0.5.3-fr}
\EGAReferenceAlias{eq:0.5.3}{subsection:0.5.3-fr}
\EGAReferenceAlias{subsection:0.5.3}{subsection:0.5.3-fr}
\EGAReferenceAlias{subsection:0.5.3-es}{subsection:0.5.3-fr}
\EGAReferenceAlias{subsection:0.5.3.fr}{subsection:0.5.3-fr}
\EGAReferenceAlias{subsection:0.5.3.es}{subsection:0.5.3-fr}
\EGAReferenceAlias{subsection:0.5.3.eq}{subsection:0.5.3-fr}
\EGAReferenceAlias{eq:subsection:0.5.3}{subsection:0.5.3-fr}
\EGAReferenceAlias{0.5.3.1}{0.5.3.1-fr}
\EGAReferenceAlias{0.5.3.1-es}{0.5.3.1-fr}
\EGAReferenceAlias{0.5.3.1.fr}{0.5.3.1-fr}
\EGAReferenceAlias{0.5.3.1.es}{0.5.3.1-fr}
\EGAReferenceAlias{0.5.3.1.eq}{0.5.3.1-fr}
\EGAReferenceAlias{eq:0.5.3.1}{0.5.3.1-fr}
\EGAReferenceAlias{0.5.3.2}{0.5.3.2-fr}
\EGAReferenceAlias{0.5.3.2-es}{0.5.3.2-fr}
\EGAReferenceAlias{0.5.3.2.fr}{0.5.3.2-fr}
\EGAReferenceAlias{0.5.3.2.es}{0.5.3.2-fr}
\EGAReferenceAlias{0.5.3.2.eq}{0.5.3.2-fr}
\EGAReferenceAlias{eq:0.5.3.2}{0.5.3.2-fr}
\EGAReferenceAlias{0.5.3.3}{0.5.3.3-fr}
\EGAReferenceAlias{0.5.3.3-es}{0.5.3.3-fr}
\EGAReferenceAlias{0.5.3.3.fr}{0.5.3.3-fr}
\EGAReferenceAlias{0.5.3.3.es}{0.5.3.3-fr}
\EGAReferenceAlias{0.5.3.3.eq}{0.5.3.3-fr}
\EGAReferenceAlias{eq:0.5.3.3}{0.5.3.3-fr}
\EGAReferenceAlias{0.5.3.4}{0.5.3.4-fr}
\EGAReferenceAlias{0.5.3.4-es}{0.5.3.4-fr}
\EGAReferenceAlias{0.5.3.4.fr}{0.5.3.4-fr}
\EGAReferenceAlias{0.5.3.4.es}{0.5.3.4-fr}
\EGAReferenceAlias{0.5.3.4.eq}{0.5.3.4-fr}
\EGAReferenceAlias{eq:0.5.3.4}{0.5.3.4-fr}
\EGAReferenceAlias{0.5.3.5}{0.5.3.5-fr}
\EGAReferenceAlias{0.5.3.5-es}{0.5.3.5-fr}
\EGAReferenceAlias{0.5.3.5.fr}{0.5.3.5-fr}
\EGAReferenceAlias{0.5.3.5.es}{0.5.3.5-fr}
\EGAReferenceAlias{0.5.3.5.eq}{0.5.3.5-fr}
\EGAReferenceAlias{eq:0.5.3.5}{0.5.3.5-fr}
\EGAReferenceAlias{0.5.3.6}{0.5.3.6-fr}
\EGAReferenceAlias{0.5.3.6-es}{0.5.3.6-fr}
\EGAReferenceAlias{0.5.3.6.fr}{0.5.3.6-fr}
\EGAReferenceAlias{0.5.3.6.es}{0.5.3.6-fr}
\EGAReferenceAlias{0.5.3.6.eq}{0.5.3.6-fr}
\EGAReferenceAlias{eq:0.5.3.6}{0.5.3.6-fr}
\EGAReferenceAlias{0.5.3.7}{0.5.3.7-fr}
\EGAReferenceAlias{0.5.3.7-es}{0.5.3.7-fr}
\EGAReferenceAlias{0.5.3.7.fr}{0.5.3.7-fr}
\EGAReferenceAlias{0.5.3.7.es}{0.5.3.7-fr}
\EGAReferenceAlias{0.5.3.7.eq}{0.5.3.7-fr}
\EGAReferenceAlias{eq:0.5.3.7}{0.5.3.7-fr}
\EGAReferenceAlias{0.5.3.8}{0.5.3.8-fr}
\EGAReferenceAlias{0.5.3.8-es}{0.5.3.8-fr}
\EGAReferenceAlias{0.5.3.8.fr}{0.5.3.8-fr}
\EGAReferenceAlias{0.5.3.8.es}{0.5.3.8-fr}
\EGAReferenceAlias{0.5.3.8.eq}{0.5.3.8-fr}
\EGAReferenceAlias{eq:0.5.3.8}{0.5.3.8-fr}
\EGAReferenceAlias{0.5.3.9}{0.5.3.9-fr}
\EGAReferenceAlias{0.5.3.9-es}{0.5.3.9-fr}
\EGAReferenceAlias{0.5.3.9.fr}{0.5.3.9-fr}
\EGAReferenceAlias{0.5.3.9.es}{0.5.3.9-fr}
\EGAReferenceAlias{0.5.3.9.eq}{0.5.3.9-fr}
\EGAReferenceAlias{eq:0.5.3.9}{0.5.3.9-fr}
\EGAReferenceAlias{0.5.3.10}{0.5.3.10-fr}
\EGAReferenceAlias{0.5.3.10-es}{0.5.3.10-fr}
\EGAReferenceAlias{0.5.3.10.fr}{0.5.3.10-fr}
\EGAReferenceAlias{0.5.3.10.es}{0.5.3.10-fr}
\EGAReferenceAlias{0.5.3.10.eq}{0.5.3.10-fr}
\EGAReferenceAlias{eq:0.5.3.10}{0.5.3.10-fr}
\EGAReferenceAlias{0.5.3.11}{0.5.3.11-fr}
\EGAReferenceAlias{0.5.3.11-es}{0.5.3.11-fr}
\EGAReferenceAlias{0.5.3.11.fr}{0.5.3.11-fr}
\EGAReferenceAlias{0.5.3.11.es}{0.5.3.11-fr}
\EGAReferenceAlias{0.5.3.11.eq}{0.5.3.11-fr}
\EGAReferenceAlias{eq:0.5.3.11}{0.5.3.11-fr}
\EGAReferenceAlias{0.5.3.12}{0.5.3.12-fr}
\EGAReferenceAlias{0.5.3.12-es}{0.5.3.12-fr}
\EGAReferenceAlias{0.5.3.12.fr}{0.5.3.12-fr}
\EGAReferenceAlias{0.5.3.12.es}{0.5.3.12-fr}
\EGAReferenceAlias{0.5.3.12.eq}{0.5.3.12-fr}
\EGAReferenceAlias{eq:0.5.3.12}{0.5.3.12-fr}
\EGAReferenceAlias{0.5.4}{subsection:0.5.4-fr}
\EGAReferenceAlias{0.5.4-fr}{subsection:0.5.4-fr}
\EGAReferenceAlias{0.5.4-es}{subsection:0.5.4-fr}
\EGAReferenceAlias{0.5.4.fr}{subsection:0.5.4-fr}
\EGAReferenceAlias{0.5.4.es}{subsection:0.5.4-fr}
\EGAReferenceAlias{0.5.4.eq}{subsection:0.5.4-fr}
\EGAReferenceAlias{eq:0.5.4}{subsection:0.5.4-fr}
\EGAReferenceAlias{subsection:0.5.4}{subsection:0.5.4-fr}
\EGAReferenceAlias{subsection:0.5.4-es}{subsection:0.5.4-fr}
\EGAReferenceAlias{subsection:0.5.4.fr}{subsection:0.5.4-fr}
\EGAReferenceAlias{subsection:0.5.4.es}{subsection:0.5.4-fr}
\EGAReferenceAlias{subsection:0.5.4.eq}{subsection:0.5.4-fr}
\EGAReferenceAlias{eq:subsection:0.5.4}{subsection:0.5.4-fr}
\EGAReferenceAlias{0.5.4.1}{0.5.4.1-fr}
\EGAReferenceAlias{0.5.4.1-es}{0.5.4.1-fr}
\EGAReferenceAlias{0.5.4.1.fr}{0.5.4.1-fr}
\EGAReferenceAlias{0.5.4.1.es}{0.5.4.1-fr}
\EGAReferenceAlias{0.5.4.1.eq}{0.5.4.1-fr}
\EGAReferenceAlias{eq:0.5.4.1}{0.5.4.1-fr}
\EGAReferenceAlias{0.5.4.2}{0.5.4.2-fr}
\EGAReferenceAlias{0.5.4.2-es}{0.5.4.2-fr}
\EGAReferenceAlias{0.5.4.2.fr}{0.5.4.2-fr}
\EGAReferenceAlias{0.5.4.2.es}{0.5.4.2-fr}
\EGAReferenceAlias{0.5.4.2.eq}{0.5.4.2-fr}
\EGAReferenceAlias{eq:0.5.4.2}{0.5.4.2-fr}
\EGAReferenceAlias{0.5.4.2.1}{0.5.4.2.1-fr}
\EGAReferenceAlias{0.5.4.2.1-es}{0.5.4.2.1-fr}
\EGAReferenceAlias{0.5.4.2.1.fr}{0.5.4.2.1-fr}
\EGAReferenceAlias{0.5.4.2.1.es}{0.5.4.2.1-fr}
\EGAReferenceAlias{0.5.4.2.1.eq}{0.5.4.2.1-fr}
\EGAReferenceAlias{eq:0.5.4.2.1}{0.5.4.2.1-fr}
\EGAReferenceAlias{0.5.4.3.1}{0.5.4.3.1-fr}
\EGAReferenceAlias{0.5.4.3.1-es}{0.5.4.3.1-fr}
\EGAReferenceAlias{0.5.4.3.1.fr}{0.5.4.3.1-fr}
\EGAReferenceAlias{0.5.4.3.1.es}{0.5.4.3.1-fr}
\EGAReferenceAlias{0.5.4.3.1.eq}{0.5.4.3.1-fr}
\EGAReferenceAlias{eq:0.5.4.3.1}{0.5.4.3.1-fr}
\EGAReferenceAlias{0.5.4.4}{0.5.4.4-fr}
\EGAReferenceAlias{0.5.4.4-es}{0.5.4.4-fr}
\EGAReferenceAlias{0.5.4.4.fr}{0.5.4.4-fr}
\EGAReferenceAlias{0.5.4.4.es}{0.5.4.4-fr}
\EGAReferenceAlias{0.5.4.4.eq}{0.5.4.4-fr}
\EGAReferenceAlias{eq:0.5.4.4}{0.5.4.4-fr}
\EGAReferenceAlias{0.5.4.4.1}{0.5.4.4.1-fr}
\EGAReferenceAlias{0.5.4.4.1-es}{0.5.4.4.1-fr}
\EGAReferenceAlias{0.5.4.4.1.fr}{0.5.4.4.1-fr}
\EGAReferenceAlias{0.5.4.4.1.es}{0.5.4.4.1-fr}
\EGAReferenceAlias{0.5.4.4.1.eq}{0.5.4.4.1-fr}
\EGAReferenceAlias{eq:0.5.4.4.1}{0.5.4.4.1-fr}
\EGAReferenceAlias{0.5.4.6}{0.5.4.6-fr}
\EGAReferenceAlias{0.5.4.6-es}{0.5.4.6-fr}
\EGAReferenceAlias{0.5.4.6.fr}{0.5.4.6-fr}
\EGAReferenceAlias{0.5.4.6.es}{0.5.4.6-fr}
\EGAReferenceAlias{0.5.4.6.eq}{0.5.4.6-fr}
\EGAReferenceAlias{eq:0.5.4.6}{0.5.4.6-fr}
\EGAReferenceAlias{0.5.4.8}{0.5.4.8-fr}
\EGAReferenceAlias{0.5.4.8-es}{0.5.4.8-fr}
\EGAReferenceAlias{0.5.4.8.fr}{0.5.4.8-fr}
\EGAReferenceAlias{0.5.4.8.es}{0.5.4.8-fr}
\EGAReferenceAlias{0.5.4.8.eq}{0.5.4.8-fr}
\EGAReferenceAlias{eq:0.5.4.8}{0.5.4.8-fr}
\EGAReferenceAlias{0.5.4.8.1}{0.5.4.8.1-fr}
\EGAReferenceAlias{0.5.4.8.1-es}{0.5.4.8.1-fr}
\EGAReferenceAlias{0.5.4.8.1.fr}{0.5.4.8.1-fr}
\EGAReferenceAlias{0.5.4.8.1.es}{0.5.4.8.1-fr}
\EGAReferenceAlias{0.5.4.8.1.eq}{0.5.4.8.1-fr}
\EGAReferenceAlias{eq:0.5.4.8.1}{0.5.4.8.1-fr}
\EGAReferenceAlias{0.5.4.9}{0.5.4.9-fr}
\EGAReferenceAlias{0.5.4.9-es}{0.5.4.9-fr}
\EGAReferenceAlias{0.5.4.9.fr}{0.5.4.9-fr}
\EGAReferenceAlias{0.5.4.9.es}{0.5.4.9-fr}
\EGAReferenceAlias{0.5.4.9.eq}{0.5.4.9-fr}
\EGAReferenceAlias{eq:0.5.4.9}{0.5.4.9-fr}
\EGAReferenceAlias{0.5.4.10}{0.5.4.10-fr}
\EGAReferenceAlias{0.5.4.10-es}{0.5.4.10-fr}
\EGAReferenceAlias{0.5.4.10.fr}{0.5.4.10-fr}
\EGAReferenceAlias{0.5.4.10.es}{0.5.4.10-fr}
\EGAReferenceAlias{0.5.4.10.eq}{0.5.4.10-fr}
\EGAReferenceAlias{eq:0.5.4.10}{0.5.4.10-fr}
\EGAReferenceAlias{0.5.4.10.1}{0.5.4.10.1-fr}
\EGAReferenceAlias{0.5.4.10.1-es}{0.5.4.10.1-fr}
\EGAReferenceAlias{0.5.4.10.1.fr}{0.5.4.10.1-fr}
\EGAReferenceAlias{0.5.4.10.1.es}{0.5.4.10.1-fr}
\EGAReferenceAlias{0.5.4.10.1.eq}{0.5.4.10.1-fr}
\EGAReferenceAlias{eq:0.5.4.10.1}{0.5.4.10.1-fr}
\EGAReferenceAlias{0.5.5}{subsection:0.5.5-fr}
\EGAReferenceAlias{0.5.5-fr}{subsection:0.5.5-fr}
\EGAReferenceAlias{0.5.5-es}{subsection:0.5.5-fr}
\EGAReferenceAlias{0.5.5.fr}{subsection:0.5.5-fr}
\EGAReferenceAlias{0.5.5.es}{subsection:0.5.5-fr}
\EGAReferenceAlias{0.5.5.eq}{subsection:0.5.5-fr}
\EGAReferenceAlias{eq:0.5.5}{subsection:0.5.5-fr}
\EGAReferenceAlias{subsection:0.5.5}{subsection:0.5.5-fr}
\EGAReferenceAlias{subsection:0.5.5-es}{subsection:0.5.5-fr}
\EGAReferenceAlias{subsection:0.5.5.fr}{subsection:0.5.5-fr}
\EGAReferenceAlias{subsection:0.5.5.es}{subsection:0.5.5-fr}
\EGAReferenceAlias{subsection:0.5.5.eq}{subsection:0.5.5-fr}
\EGAReferenceAlias{eq:subsection:0.5.5}{subsection:0.5.5-fr}
\EGAReferenceAlias{0.5.5.1}{0.5.5.1-fr}
\EGAReferenceAlias{0.5.5.1-es}{0.5.5.1-fr}
\EGAReferenceAlias{0.5.5.1.fr}{0.5.5.1-fr}
\EGAReferenceAlias{0.5.5.1.es}{0.5.5.1-fr}
\EGAReferenceAlias{0.5.5.1.eq}{0.5.5.1-fr}
\EGAReferenceAlias{eq:0.5.5.1}{0.5.5.1-fr}
\EGAReferenceAlias{0.5.5.2}{0.5.5.2-fr}
\EGAReferenceAlias{0.5.5.2-es}{0.5.5.2-fr}
\EGAReferenceAlias{0.5.5.2.fr}{0.5.5.2-fr}
\EGAReferenceAlias{0.5.5.2.es}{0.5.5.2-fr}
\EGAReferenceAlias{0.5.5.2.eq}{0.5.5.2-fr}
\EGAReferenceAlias{eq:0.5.5.2}{0.5.5.2-fr}
\EGAReferenceAlias{0.5.5.3}{0.5.5.3-fr}
\EGAReferenceAlias{0.5.5.3-es}{0.5.5.3-fr}
\EGAReferenceAlias{0.5.5.3.fr}{0.5.5.3-fr}
\EGAReferenceAlias{0.5.5.3.es}{0.5.5.3-fr}
\EGAReferenceAlias{0.5.5.3.eq}{0.5.5.3-fr}
\EGAReferenceAlias{eq:0.5.5.3}{0.5.5.3-fr}
\EGAReferenceAlias{0.5.5.4}{0.5.5.4-fr}
\EGAReferenceAlias{0.5.5.4-es}{0.5.5.4-fr}
\EGAReferenceAlias{0.5.5.4.fr}{0.5.5.4-fr}
\EGAReferenceAlias{0.5.5.4.es}{0.5.5.4-fr}
\EGAReferenceAlias{0.5.5.4.eq}{0.5.5.4-fr}
\EGAReferenceAlias{eq:0.5.5.4}{0.5.5.4-fr}
\EGAReferenceAlias{0.5.5.5}{0.5.5.5-fr}
\EGAReferenceAlias{0.5.5.5-es}{0.5.5.5-fr}
\EGAReferenceAlias{0.5.5.5.fr}{0.5.5.5-fr}
\EGAReferenceAlias{0.5.5.5.es}{0.5.5.5-fr}
\EGAReferenceAlias{0.5.5.5.eq}{0.5.5.5-fr}
\EGAReferenceAlias{eq:0.5.5.5}{0.5.5.5-fr}
\EGAReferenceAlias{0.6}{section:0.6-fr}
\EGAReferenceAlias{0.6-fr}{section:0.6-fr}
\EGAReferenceAlias{0.6-es}{section:0.6-fr}
\EGAReferenceAlias{0.6.fr}{section:0.6-fr}
\EGAReferenceAlias{0.6.es}{section:0.6-fr}
\EGAReferenceAlias{0.6.eq}{section:0.6-fr}
\EGAReferenceAlias{eq:0.6}{section:0.6-fr}
\EGAReferenceAlias{section:0.6}{section:0.6-fr}
\EGAReferenceAlias{section:0.6-es}{section:0.6-fr}
\EGAReferenceAlias{section:0.6.fr}{section:0.6-fr}
\EGAReferenceAlias{section:0.6.es}{section:0.6-fr}
\EGAReferenceAlias{section:0.6.eq}{section:0.6-fr}
\EGAReferenceAlias{eq:section:0.6}{section:0.6-fr}
\EGAReferenceAlias{0.6.0.0}{0.6.0.0-fr}
\EGAReferenceAlias{0.6.0.0-es}{0.6.0.0-fr}
\EGAReferenceAlias{0.6.0.0.fr}{0.6.0.0-fr}
\EGAReferenceAlias{0.6.0.0.es}{0.6.0.0-fr}
\EGAReferenceAlias{0.6.0.0.eq}{0.6.0.0-fr}
\EGAReferenceAlias{eq:0.6.0.0}{0.6.0.0-fr}
\EGAReferenceAlias{0.6.1}{subsection:0.6.1-fr}
\EGAReferenceAlias{0.6.1-fr}{subsection:0.6.1-fr}
\EGAReferenceAlias{0.6.1-es}{subsection:0.6.1-fr}
\EGAReferenceAlias{0.6.1.fr}{subsection:0.6.1-fr}
\EGAReferenceAlias{0.6.1.es}{subsection:0.6.1-fr}
\EGAReferenceAlias{0.6.1.eq}{subsection:0.6.1-fr}
\EGAReferenceAlias{eq:0.6.1}{subsection:0.6.1-fr}
\EGAReferenceAlias{subsection:0.6.1}{subsection:0.6.1-fr}
\EGAReferenceAlias{subsection:0.6.1-es}{subsection:0.6.1-fr}
\EGAReferenceAlias{subsection:0.6.1.fr}{subsection:0.6.1-fr}
\EGAReferenceAlias{subsection:0.6.1.es}{subsection:0.6.1-fr}
\EGAReferenceAlias{subsection:0.6.1.eq}{subsection:0.6.1-fr}
\EGAReferenceAlias{eq:subsection:0.6.1}{subsection:0.6.1-fr}
\EGAReferenceAlias{0.6.1.1}{0.6.1.1-fr}
\EGAReferenceAlias{0.6.1.1-es}{0.6.1.1-fr}
\EGAReferenceAlias{0.6.1.1.fr}{0.6.1.1-fr}
\EGAReferenceAlias{0.6.1.1.es}{0.6.1.1-fr}
\EGAReferenceAlias{0.6.1.1.eq}{0.6.1.1-fr}
\EGAReferenceAlias{eq:0.6.1.1}{0.6.1.1-fr}
\EGAReferenceAlias{0.6.1.3}{0.6.1.3-fr}
\EGAReferenceAlias{0.6.1.3-es}{0.6.1.3-fr}
\EGAReferenceAlias{0.6.1.3.fr}{0.6.1.3-fr}
\EGAReferenceAlias{0.6.1.3.es}{0.6.1.3-fr}
\EGAReferenceAlias{0.6.1.3.eq}{0.6.1.3-fr}
\EGAReferenceAlias{eq:0.6.1.3}{0.6.1.3-fr}
\EGAReferenceAlias{0.6.1.4}{0.6.1.4-fr}
\EGAReferenceAlias{0.6.1.4-es}{0.6.1.4-fr}
\EGAReferenceAlias{0.6.1.4.fr}{0.6.1.4-fr}
\EGAReferenceAlias{0.6.1.4.es}{0.6.1.4-fr}
\EGAReferenceAlias{0.6.1.4.eq}{0.6.1.4-fr}
\EGAReferenceAlias{eq:0.6.1.4}{0.6.1.4-fr}
\EGAReferenceAlias{0.6.2}{subsection:0.6.2-fr}
\EGAReferenceAlias{0.6.2-fr}{subsection:0.6.2-fr}
\EGAReferenceAlias{0.6.2-es}{subsection:0.6.2-fr}
\EGAReferenceAlias{0.6.2.fr}{subsection:0.6.2-fr}
\EGAReferenceAlias{0.6.2.es}{subsection:0.6.2-fr}
\EGAReferenceAlias{0.6.2.eq}{subsection:0.6.2-fr}
\EGAReferenceAlias{eq:0.6.2}{subsection:0.6.2-fr}
\EGAReferenceAlias{subsection:0.6.2}{subsection:0.6.2-fr}
\EGAReferenceAlias{subsection:0.6.2-es}{subsection:0.6.2-fr}
\EGAReferenceAlias{subsection:0.6.2.fr}{subsection:0.6.2-fr}
\EGAReferenceAlias{subsection:0.6.2.es}{subsection:0.6.2-fr}
\EGAReferenceAlias{subsection:0.6.2.eq}{subsection:0.6.2-fr}
\EGAReferenceAlias{eq:subsection:0.6.2}{subsection:0.6.2-fr}
\EGAReferenceAlias{0.6.2.1}{0.6.2.1-fr}
\EGAReferenceAlias{0.6.2.1-es}{0.6.2.1-fr}
\EGAReferenceAlias{0.6.2.1.fr}{0.6.2.1-fr}
\EGAReferenceAlias{0.6.2.1.es}{0.6.2.1-fr}
\EGAReferenceAlias{0.6.2.1.eq}{0.6.2.1-fr}
\EGAReferenceAlias{eq:0.6.2.1}{0.6.2.1-fr}
\EGAReferenceAlias{0.6.2.2}{0.6.2.2-fr}
\EGAReferenceAlias{0.6.2.2-es}{0.6.2.2-fr}
\EGAReferenceAlias{0.6.2.2.fr}{0.6.2.2-fr}
\EGAReferenceAlias{0.6.2.2.es}{0.6.2.2-fr}
\EGAReferenceAlias{0.6.2.2.eq}{0.6.2.2-fr}
\EGAReferenceAlias{eq:0.6.2.2}{0.6.2.2-fr}
\EGAReferenceAlias{0.6.2.3}{0.6.2.3-fr}
\EGAReferenceAlias{0.6.2.3-es}{0.6.2.3-fr}
\EGAReferenceAlias{0.6.2.3.fr}{0.6.2.3-fr}
\EGAReferenceAlias{0.6.2.3.es}{0.6.2.3-fr}
\EGAReferenceAlias{0.6.2.3.eq}{0.6.2.3-fr}
\EGAReferenceAlias{eq:0.6.2.3}{0.6.2.3-fr}
\EGAReferenceAlias{0.6.3}{subsection:0.6.3-fr}
\EGAReferenceAlias{0.6.3-fr}{subsection:0.6.3-fr}
\EGAReferenceAlias{0.6.3-es}{subsection:0.6.3-fr}
\EGAReferenceAlias{0.6.3.fr}{subsection:0.6.3-fr}
\EGAReferenceAlias{0.6.3.es}{subsection:0.6.3-fr}
\EGAReferenceAlias{0.6.3.eq}{subsection:0.6.3-fr}
\EGAReferenceAlias{eq:0.6.3}{subsection:0.6.3-fr}
\EGAReferenceAlias{subsection:0.6.3}{subsection:0.6.3-fr}
\EGAReferenceAlias{subsection:0.6.3-es}{subsection:0.6.3-fr}
\EGAReferenceAlias{subsection:0.6.3.fr}{subsection:0.6.3-fr}
\EGAReferenceAlias{subsection:0.6.3.es}{subsection:0.6.3-fr}
\EGAReferenceAlias{subsection:0.6.3.eq}{subsection:0.6.3-fr}
\EGAReferenceAlias{eq:subsection:0.6.3}{subsection:0.6.3-fr}
\EGAReferenceAlias{0.6.3.1}{0.6.3.1-fr}
\EGAReferenceAlias{0.6.3.1-es}{0.6.3.1-fr}
\EGAReferenceAlias{0.6.3.1.fr}{0.6.3.1-fr}
\EGAReferenceAlias{0.6.3.1.es}{0.6.3.1-fr}
\EGAReferenceAlias{0.6.3.1.eq}{0.6.3.1-fr}
\EGAReferenceAlias{eq:0.6.3.1}{0.6.3.1-fr}
\EGAReferenceAlias{0.6.3.2}{0.6.3.2-fr}
\EGAReferenceAlias{0.6.3.2-es}{0.6.3.2-fr}
\EGAReferenceAlias{0.6.3.2.fr}{0.6.3.2-fr}
\EGAReferenceAlias{0.6.3.2.es}{0.6.3.2-fr}
\EGAReferenceAlias{0.6.3.2.eq}{0.6.3.2-fr}
\EGAReferenceAlias{eq:0.6.3.2}{0.6.3.2-fr}
\EGAReferenceAlias{0.6.3.3}{0.6.3.3-fr}
\EGAReferenceAlias{0.6.3.3-es}{0.6.3.3-fr}
\EGAReferenceAlias{0.6.3.3.fr}{0.6.3.3-fr}
\EGAReferenceAlias{0.6.3.3.es}{0.6.3.3-fr}
\EGAReferenceAlias{0.6.3.3.eq}{0.6.3.3-fr}
\EGAReferenceAlias{eq:0.6.3.3}{0.6.3.3-fr}
\EGAReferenceAlias{0.6.3.4}{0.6.3.4-fr}
\EGAReferenceAlias{0.6.3.4-es}{0.6.3.4-fr}
\EGAReferenceAlias{0.6.3.4.fr}{0.6.3.4-fr}
\EGAReferenceAlias{0.6.3.4.es}{0.6.3.4-fr}
\EGAReferenceAlias{0.6.3.4.eq}{0.6.3.4-fr}
\EGAReferenceAlias{eq:0.6.3.4}{0.6.3.4-fr}
\EGAReferenceAlias{0.6.4}{subsection:0.6.4-fr}
\EGAReferenceAlias{0.6.4-fr}{subsection:0.6.4-fr}
\EGAReferenceAlias{0.6.4-es}{subsection:0.6.4-fr}
\EGAReferenceAlias{0.6.4.fr}{subsection:0.6.4-fr}
\EGAReferenceAlias{0.6.4.es}{subsection:0.6.4-fr}
\EGAReferenceAlias{0.6.4.eq}{subsection:0.6.4-fr}
\EGAReferenceAlias{eq:0.6.4}{subsection:0.6.4-fr}
\EGAReferenceAlias{subsection:0.6.4}{subsection:0.6.4-fr}
\EGAReferenceAlias{subsection:0.6.4-es}{subsection:0.6.4-fr}
\EGAReferenceAlias{subsection:0.6.4.fr}{subsection:0.6.4-fr}
\EGAReferenceAlias{subsection:0.6.4.es}{subsection:0.6.4-fr}
\EGAReferenceAlias{subsection:0.6.4.eq}{subsection:0.6.4-fr}
\EGAReferenceAlias{eq:subsection:0.6.4}{subsection:0.6.4-fr}
\EGAReferenceAlias{0.6.4.1}{0.6.4.1-fr}
\EGAReferenceAlias{0.6.4.1-es}{0.6.4.1-fr}
\EGAReferenceAlias{0.6.4.1.fr}{0.6.4.1-fr}
\EGAReferenceAlias{0.6.4.1.es}{0.6.4.1-fr}
\EGAReferenceAlias{0.6.4.1.eq}{0.6.4.1-fr}
\EGAReferenceAlias{eq:0.6.4.1}{0.6.4.1-fr}
\EGAReferenceAlias{0.6.4.2}{0.6.4.2-fr}
\EGAReferenceAlias{0.6.4.2-es}{0.6.4.2-fr}
\EGAReferenceAlias{0.6.4.2.fr}{0.6.4.2-fr}
\EGAReferenceAlias{0.6.4.2.es}{0.6.4.2-fr}
\EGAReferenceAlias{0.6.4.2.eq}{0.6.4.2-fr}
\EGAReferenceAlias{eq:0.6.4.2}{0.6.4.2-fr}
\EGAReferenceAlias{0.6.4.3}{0.6.4.3-fr}
\EGAReferenceAlias{0.6.4.3-es}{0.6.4.3-fr}
\EGAReferenceAlias{0.6.4.3.fr}{0.6.4.3-fr}
\EGAReferenceAlias{0.6.4.3.es}{0.6.4.3-fr}
\EGAReferenceAlias{0.6.4.3.eq}{0.6.4.3-fr}
\EGAReferenceAlias{eq:0.6.4.3}{0.6.4.3-fr}
\EGAReferenceAlias{0.6.4.4}{0.6.4.4-fr}
\EGAReferenceAlias{0.6.4.4-es}{0.6.4.4-fr}
\EGAReferenceAlias{0.6.4.4.fr}{0.6.4.4-fr}
\EGAReferenceAlias{0.6.4.4.es}{0.6.4.4-fr}
\EGAReferenceAlias{0.6.4.4.eq}{0.6.4.4-fr}
\EGAReferenceAlias{eq:0.6.4.4}{0.6.4.4-fr}
\EGAReferenceAlias{0.6.4.5}{0.6.4.5-fr}
\EGAReferenceAlias{0.6.4.5-es}{0.6.4.5-fr}
\EGAReferenceAlias{0.6.4.5.fr}{0.6.4.5-fr}
\EGAReferenceAlias{0.6.4.5.es}{0.6.4.5-fr}
\EGAReferenceAlias{0.6.4.5.eq}{0.6.4.5-fr}
\EGAReferenceAlias{eq:0.6.4.5}{0.6.4.5-fr}
\EGAReferenceAlias{0.6.5}{subsection:0.6.5-fr}
\EGAReferenceAlias{0.6.5-fr}{subsection:0.6.5-fr}
\EGAReferenceAlias{0.6.5-es}{subsection:0.6.5-fr}
\EGAReferenceAlias{0.6.5.fr}{subsection:0.6.5-fr}
\EGAReferenceAlias{0.6.5.es}{subsection:0.6.5-fr}
\EGAReferenceAlias{0.6.5.eq}{subsection:0.6.5-fr}
\EGAReferenceAlias{eq:0.6.5}{subsection:0.6.5-fr}
\EGAReferenceAlias{subsection:0.6.5}{subsection:0.6.5-fr}
\EGAReferenceAlias{subsection:0.6.5-es}{subsection:0.6.5-fr}
\EGAReferenceAlias{subsection:0.6.5.fr}{subsection:0.6.5-fr}
\EGAReferenceAlias{subsection:0.6.5.es}{subsection:0.6.5-fr}
\EGAReferenceAlias{subsection:0.6.5.eq}{subsection:0.6.5-fr}
\EGAReferenceAlias{eq:subsection:0.6.5}{subsection:0.6.5-fr}
\EGAReferenceAlias{0.6.5.1}{0.6.5.1-fr}
\EGAReferenceAlias{0.6.5.1-es}{0.6.5.1-fr}
\EGAReferenceAlias{0.6.5.1.fr}{0.6.5.1-fr}
\EGAReferenceAlias{0.6.5.1.es}{0.6.5.1-fr}
\EGAReferenceAlias{0.6.5.1.eq}{0.6.5.1-fr}
\EGAReferenceAlias{eq:0.6.5.1}{0.6.5.1-fr}
\EGAReferenceAlias{0.6.5.2}{0.6.5.2-fr}
\EGAReferenceAlias{0.6.5.2-es}{0.6.5.2-fr}
\EGAReferenceAlias{0.6.5.2.fr}{0.6.5.2-fr}
\EGAReferenceAlias{0.6.5.2.es}{0.6.5.2-fr}
\EGAReferenceAlias{0.6.5.2.eq}{0.6.5.2-fr}
\EGAReferenceAlias{eq:0.6.5.2}{0.6.5.2-fr}
\EGAReferenceAlias{0.6.6}{subsection:0.6.6-fr}
\EGAReferenceAlias{0.6.6-fr}{subsection:0.6.6-fr}
\EGAReferenceAlias{0.6.6-es}{subsection:0.6.6-fr}
\EGAReferenceAlias{0.6.6.fr}{subsection:0.6.6-fr}
\EGAReferenceAlias{0.6.6.es}{subsection:0.6.6-fr}
\EGAReferenceAlias{0.6.6.eq}{subsection:0.6.6-fr}
\EGAReferenceAlias{eq:0.6.6}{subsection:0.6.6-fr}
\EGAReferenceAlias{subsection:0.6.6}{subsection:0.6.6-fr}
\EGAReferenceAlias{subsection:0.6.6-es}{subsection:0.6.6-fr}
\EGAReferenceAlias{subsection:0.6.6.fr}{subsection:0.6.6-fr}
\EGAReferenceAlias{subsection:0.6.6.es}{subsection:0.6.6-fr}
\EGAReferenceAlias{subsection:0.6.6.eq}{subsection:0.6.6-fr}
\EGAReferenceAlias{eq:subsection:0.6.6}{subsection:0.6.6-fr}
\EGAReferenceAlias{0.6.6.1}{0.6.6.1-fr}
\EGAReferenceAlias{0.6.6.1-es}{0.6.6.1-fr}
\EGAReferenceAlias{0.6.6.1.fr}{0.6.6.1-fr}
\EGAReferenceAlias{0.6.6.1.es}{0.6.6.1-fr}
\EGAReferenceAlias{0.6.6.1.eq}{0.6.6.1-fr}
\EGAReferenceAlias{eq:0.6.6.1}{0.6.6.1-fr}
\EGAReferenceAlias{0.6.6.2}{0.6.6.2-fr}
\EGAReferenceAlias{0.6.6.2-es}{0.6.6.2-fr}
\EGAReferenceAlias{0.6.6.2.fr}{0.6.6.2-fr}
\EGAReferenceAlias{0.6.6.2.es}{0.6.6.2-fr}
\EGAReferenceAlias{0.6.6.2.eq}{0.6.6.2-fr}
\EGAReferenceAlias{eq:0.6.6.2}{0.6.6.2-fr}
\EGAReferenceAlias{0.6.6.3}{0.6.6.3-fr}
\EGAReferenceAlias{0.6.6.3-es}{0.6.6.3-fr}
\EGAReferenceAlias{0.6.6.3.fr}{0.6.6.3-fr}
\EGAReferenceAlias{0.6.6.3.es}{0.6.6.3-fr}
\EGAReferenceAlias{0.6.6.3.eq}{0.6.6.3-fr}
\EGAReferenceAlias{eq:0.6.6.3}{0.6.6.3-fr}
\EGAReferenceAlias{0.6.6.4}{0.6.6.4-fr}
\EGAReferenceAlias{0.6.6.4-es}{0.6.6.4-fr}
\EGAReferenceAlias{0.6.6.4.fr}{0.6.6.4-fr}
\EGAReferenceAlias{0.6.6.4.es}{0.6.6.4-fr}
\EGAReferenceAlias{0.6.6.4.eq}{0.6.6.4-fr}
\EGAReferenceAlias{eq:0.6.6.4}{0.6.6.4-fr}
\EGAReferenceAlias{0.6.7}{subsection:0.6.7-fr}
\EGAReferenceAlias{0.6.7-fr}{subsection:0.6.7-fr}
\EGAReferenceAlias{0.6.7-es}{subsection:0.6.7-fr}
\EGAReferenceAlias{0.6.7.fr}{subsection:0.6.7-fr}
\EGAReferenceAlias{0.6.7.es}{subsection:0.6.7-fr}
\EGAReferenceAlias{0.6.7.eq}{subsection:0.6.7-fr}
\EGAReferenceAlias{eq:0.6.7}{subsection:0.6.7-fr}
\EGAReferenceAlias{subsection:0.6.7}{subsection:0.6.7-fr}
\EGAReferenceAlias{subsection:0.6.7-es}{subsection:0.6.7-fr}
\EGAReferenceAlias{subsection:0.6.7.fr}{subsection:0.6.7-fr}
\EGAReferenceAlias{subsection:0.6.7.es}{subsection:0.6.7-fr}
\EGAReferenceAlias{subsection:0.6.7.eq}{subsection:0.6.7-fr}
\EGAReferenceAlias{eq:subsection:0.6.7}{subsection:0.6.7-fr}
\EGAReferenceAlias{0.6.7.1}{0.6.7.1-fr}
\EGAReferenceAlias{0.6.7.1-es}{0.6.7.1-fr}
\EGAReferenceAlias{0.6.7.1.fr}{0.6.7.1-fr}
\EGAReferenceAlias{0.6.7.1.es}{0.6.7.1-fr}
\EGAReferenceAlias{0.6.7.1.eq}{0.6.7.1-fr}
\EGAReferenceAlias{eq:0.6.7.1}{0.6.7.1-fr}
\EGAReferenceAlias{0.6.7.2}{0.6.7.2-fr}
\EGAReferenceAlias{0.6.7.2-es}{0.6.7.2-fr}
\EGAReferenceAlias{0.6.7.2.fr}{0.6.7.2-fr}
\EGAReferenceAlias{0.6.7.2.es}{0.6.7.2-fr}
\EGAReferenceAlias{0.6.7.2.eq}{0.6.7.2-fr}
\EGAReferenceAlias{eq:0.6.7.2}{0.6.7.2-fr}
\EGAReferenceAlias{0.6.7.3}{0.6.7.3-fr}
\EGAReferenceAlias{0.6.7.3-es}{0.6.7.3-fr}
\EGAReferenceAlias{0.6.7.3.fr}{0.6.7.3-fr}
\EGAReferenceAlias{0.6.7.3.es}{0.6.7.3-fr}
\EGAReferenceAlias{0.6.7.3.eq}{0.6.7.3-fr}
\EGAReferenceAlias{eq:0.6.7.3}{0.6.7.3-fr}
\EGAReferenceAlias{0.6.7.4}{0.6.7.4-fr}
\EGAReferenceAlias{0.6.7.4-es}{0.6.7.4-fr}
\EGAReferenceAlias{0.6.7.4.fr}{0.6.7.4-fr}
\EGAReferenceAlias{0.6.7.4.es}{0.6.7.4-fr}
\EGAReferenceAlias{0.6.7.4.eq}{0.6.7.4-fr}
\EGAReferenceAlias{eq:0.6.7.4}{0.6.7.4-fr}
\EGAReferenceAlias{0.6.7.5}{0.6.7.5-fr}
\EGAReferenceAlias{0.6.7.5-es}{0.6.7.5-fr}
\EGAReferenceAlias{0.6.7.5.fr}{0.6.7.5-fr}
\EGAReferenceAlias{0.6.7.5.es}{0.6.7.5-fr}
\EGAReferenceAlias{0.6.7.5.eq}{0.6.7.5-fr}
\EGAReferenceAlias{eq:0.6.7.5}{0.6.7.5-fr}
\EGAReferenceAlias{0.6.7.6}{0.6.7.6-fr}
\EGAReferenceAlias{0.6.7.6-es}{0.6.7.6-fr}
\EGAReferenceAlias{0.6.7.6.fr}{0.6.7.6-fr}
\EGAReferenceAlias{0.6.7.6.es}{0.6.7.6-fr}
\EGAReferenceAlias{0.6.7.6.eq}{0.6.7.6-fr}
\EGAReferenceAlias{eq:0.6.7.6}{0.6.7.6-fr}
\EGAReferenceAlias{0.6.7.6.1}{0.6.7.6.1-fr}
\EGAReferenceAlias{0.6.7.6.1-es}{0.6.7.6.1-fr}
\EGAReferenceAlias{0.6.7.6.1.fr}{0.6.7.6.1-fr}
\EGAReferenceAlias{0.6.7.6.1.es}{0.6.7.6.1-fr}
\EGAReferenceAlias{0.6.7.6.1.eq}{0.6.7.6.1-fr}
\EGAReferenceAlias{eq:0.6.7.6.1}{0.6.7.6.1-fr}
\EGAReferenceAlias{0.6.7.8}{0.6.7.8-fr}
\EGAReferenceAlias{0.6.7.8-es}{0.6.7.8-fr}
\EGAReferenceAlias{0.6.7.8.fr}{0.6.7.8-fr}
\EGAReferenceAlias{0.6.7.8.es}{0.6.7.8-fr}
\EGAReferenceAlias{0.6.7.8.eq}{0.6.7.8-fr}
\EGAReferenceAlias{eq:0.6.7.8}{0.6.7.8-fr}
\EGAReferenceAlias{0.7}{section:0.7-fr}
\EGAReferenceAlias{0.7-fr}{section:0.7-fr}
\EGAReferenceAlias{0.7-es}{section:0.7-fr}
\EGAReferenceAlias{0.7.fr}{section:0.7-fr}
\EGAReferenceAlias{0.7.es}{section:0.7-fr}
\EGAReferenceAlias{0.7.eq}{section:0.7-fr}
\EGAReferenceAlias{eq:0.7}{section:0.7-fr}
\EGAReferenceAlias{section:0.7}{section:0.7-fr}
\EGAReferenceAlias{section:0.7-es}{section:0.7-fr}
\EGAReferenceAlias{section:0.7.fr}{section:0.7-fr}
\EGAReferenceAlias{section:0.7.es}{section:0.7-fr}
\EGAReferenceAlias{section:0.7.eq}{section:0.7-fr}
\EGAReferenceAlias{eq:section:0.7}{section:0.7-fr}
\EGAReferenceAlias{0.7.1}{subsection:0.7.1-fr}
\EGAReferenceAlias{0.7.1-fr}{subsection:0.7.1-fr}
\EGAReferenceAlias{0.7.1-es}{subsection:0.7.1-fr}
\EGAReferenceAlias{0.7.1.fr}{subsection:0.7.1-fr}
\EGAReferenceAlias{0.7.1.es}{subsection:0.7.1-fr}
\EGAReferenceAlias{0.7.1.eq}{subsection:0.7.1-fr}
\EGAReferenceAlias{eq:0.7.1}{subsection:0.7.1-fr}
\EGAReferenceAlias{subsection:0.7.1}{subsection:0.7.1-fr}
\EGAReferenceAlias{subsection:0.7.1-es}{subsection:0.7.1-fr}
\EGAReferenceAlias{subsection:0.7.1.fr}{subsection:0.7.1-fr}
\EGAReferenceAlias{subsection:0.7.1.es}{subsection:0.7.1-fr}
\EGAReferenceAlias{subsection:0.7.1.eq}{subsection:0.7.1-fr}
\EGAReferenceAlias{eq:subsection:0.7.1}{subsection:0.7.1-fr}
\EGAReferenceAlias{0.7.1.1}{0.7.1.1-fr}
\EGAReferenceAlias{0.7.1.1-es}{0.7.1.1-fr}
\EGAReferenceAlias{0.7.1.1.fr}{0.7.1.1-fr}
\EGAReferenceAlias{0.7.1.1.es}{0.7.1.1-fr}
\EGAReferenceAlias{0.7.1.1.eq}{0.7.1.1-fr}
\EGAReferenceAlias{eq:0.7.1.1}{0.7.1.1-fr}
\EGAReferenceAlias{0.7.1.2}{0.7.1.2-fr}
\EGAReferenceAlias{0.7.1.2-es}{0.7.1.2-fr}
\EGAReferenceAlias{0.7.1.2.fr}{0.7.1.2-fr}
\EGAReferenceAlias{0.7.1.2.es}{0.7.1.2-fr}
\EGAReferenceAlias{0.7.1.2.eq}{0.7.1.2-fr}
\EGAReferenceAlias{eq:0.7.1.2}{0.7.1.2-fr}
\EGAReferenceAlias{0.7.1.3}{0.7.1.3-fr}
\EGAReferenceAlias{0.7.1.3-es}{0.7.1.3-fr}
\EGAReferenceAlias{0.7.1.3.fr}{0.7.1.3-fr}
\EGAReferenceAlias{0.7.1.3.es}{0.7.1.3-fr}
\EGAReferenceAlias{0.7.1.3.eq}{0.7.1.3-fr}
\EGAReferenceAlias{eq:0.7.1.3}{0.7.1.3-fr}
\EGAReferenceAlias{0.7.1.4}{0.7.1.4-fr}
\EGAReferenceAlias{0.7.1.4-es}{0.7.1.4-fr}
\EGAReferenceAlias{0.7.1.4.fr}{0.7.1.4-fr}
\EGAReferenceAlias{0.7.1.4.es}{0.7.1.4-fr}
\EGAReferenceAlias{0.7.1.4.eq}{0.7.1.4-fr}
\EGAReferenceAlias{eq:0.7.1.4}{0.7.1.4-fr}
\EGAReferenceAlias{0.7.1.5}{0.7.1.5-fr}
\EGAReferenceAlias{0.7.1.5-es}{0.7.1.5-fr}
\EGAReferenceAlias{0.7.1.5.fr}{0.7.1.5-fr}
\EGAReferenceAlias{0.7.1.5.es}{0.7.1.5-fr}
\EGAReferenceAlias{0.7.1.5.eq}{0.7.1.5-fr}
\EGAReferenceAlias{eq:0.7.1.5}{0.7.1.5-fr}
\EGAReferenceAlias{0.7.1.6}{0.7.1.6-fr}
\EGAReferenceAlias{0.7.1.6-es}{0.7.1.6-fr}
\EGAReferenceAlias{0.7.1.6.fr}{0.7.1.6-fr}
\EGAReferenceAlias{0.7.1.6.es}{0.7.1.6-fr}
\EGAReferenceAlias{0.7.1.6.eq}{0.7.1.6-fr}
\EGAReferenceAlias{eq:0.7.1.6}{0.7.1.6-fr}
\EGAReferenceAlias{0.7.1.7}{0.7.1.7-fr}
\EGAReferenceAlias{0.7.1.7-es}{0.7.1.7-fr}
\EGAReferenceAlias{0.7.1.7.fr}{0.7.1.7-fr}
\EGAReferenceAlias{0.7.1.7.es}{0.7.1.7-fr}
\EGAReferenceAlias{0.7.1.7.eq}{0.7.1.7-fr}
\EGAReferenceAlias{eq:0.7.1.7}{0.7.1.7-fr}
\EGAReferenceAlias{0.7.1.8}{0.7.1.8-fr}
\EGAReferenceAlias{0.7.1.8-es}{0.7.1.8-fr}
\EGAReferenceAlias{0.7.1.8.fr}{0.7.1.8-fr}
\EGAReferenceAlias{0.7.1.8.es}{0.7.1.8-fr}
\EGAReferenceAlias{0.7.1.8.eq}{0.7.1.8-fr}
\EGAReferenceAlias{eq:0.7.1.8}{0.7.1.8-fr}
\EGAReferenceAlias{0.7.1.9}{0.7.1.9-fr}
\EGAReferenceAlias{0.7.1.9-es}{0.7.1.9-fr}
\EGAReferenceAlias{0.7.1.9.fr}{0.7.1.9-fr}
\EGAReferenceAlias{0.7.1.9.es}{0.7.1.9-fr}
\EGAReferenceAlias{0.7.1.9.eq}{0.7.1.9-fr}
\EGAReferenceAlias{eq:0.7.1.9}{0.7.1.9-fr}
\EGAReferenceAlias{0.7.1.10}{0.7.1.10-fr}
\EGAReferenceAlias{0.7.1.10-es}{0.7.1.10-fr}
\EGAReferenceAlias{0.7.1.10.fr}{0.7.1.10-fr}
\EGAReferenceAlias{0.7.1.10.es}{0.7.1.10-fr}
\EGAReferenceAlias{0.7.1.10.eq}{0.7.1.10-fr}
\EGAReferenceAlias{eq:0.7.1.10}{0.7.1.10-fr}
\EGAReferenceAlias{0.7.1.11}{0.7.1.11-fr}
\EGAReferenceAlias{0.7.1.11-es}{0.7.1.11-fr}
\EGAReferenceAlias{0.7.1.11.fr}{0.7.1.11-fr}
\EGAReferenceAlias{0.7.1.11.es}{0.7.1.11-fr}
\EGAReferenceAlias{0.7.1.11.eq}{0.7.1.11-fr}
\EGAReferenceAlias{eq:0.7.1.11}{0.7.1.11-fr}
\EGAReferenceAlias{0.7.1.12}{0.7.1.12-fr}
\EGAReferenceAlias{0.7.1.12-es}{0.7.1.12-fr}
\EGAReferenceAlias{0.7.1.12.fr}{0.7.1.12-fr}
\EGAReferenceAlias{0.7.1.12.es}{0.7.1.12-fr}
\EGAReferenceAlias{0.7.1.12.eq}{0.7.1.12-fr}
\EGAReferenceAlias{eq:0.7.1.12}{0.7.1.12-fr}
\EGAReferenceAlias{0.7.1.13}{0.7.1.13-fr}
\EGAReferenceAlias{0.7.1.13-es}{0.7.1.13-fr}
\EGAReferenceAlias{0.7.1.13.fr}{0.7.1.13-fr}
\EGAReferenceAlias{0.7.1.13.es}{0.7.1.13-fr}
\EGAReferenceAlias{0.7.1.13.eq}{0.7.1.13-fr}
\EGAReferenceAlias{eq:0.7.1.13}{0.7.1.13-fr}
\EGAReferenceAlias{0.7.1.14}{0.7.1.14-fr}
\EGAReferenceAlias{0.7.1.14-es}{0.7.1.14-fr}
\EGAReferenceAlias{0.7.1.14.fr}{0.7.1.14-fr}
\EGAReferenceAlias{0.7.1.14.es}{0.7.1.14-fr}
\EGAReferenceAlias{0.7.1.14.eq}{0.7.1.14-fr}
\EGAReferenceAlias{eq:0.7.1.14}{0.7.1.14-fr}
\EGAReferenceAlias{0.7.2}{subsection:0.7.2-fr}
\EGAReferenceAlias{0.7.2-fr}{subsection:0.7.2-fr}
\EGAReferenceAlias{0.7.2-es}{subsection:0.7.2-fr}
\EGAReferenceAlias{0.7.2.fr}{subsection:0.7.2-fr}
\EGAReferenceAlias{0.7.2.es}{subsection:0.7.2-fr}
\EGAReferenceAlias{0.7.2.eq}{subsection:0.7.2-fr}
\EGAReferenceAlias{eq:0.7.2}{subsection:0.7.2-fr}
\EGAReferenceAlias{subsection:0.7.2}{subsection:0.7.2-fr}
\EGAReferenceAlias{subsection:0.7.2-es}{subsection:0.7.2-fr}
\EGAReferenceAlias{subsection:0.7.2.fr}{subsection:0.7.2-fr}
\EGAReferenceAlias{subsection:0.7.2.es}{subsection:0.7.2-fr}
\EGAReferenceAlias{subsection:0.7.2.eq}{subsection:0.7.2-fr}
\EGAReferenceAlias{eq:subsection:0.7.2}{subsection:0.7.2-fr}
\EGAReferenceAlias{0.7.2.1}{0.7.2.1-fr}
\EGAReferenceAlias{0.7.2.1-es}{0.7.2.1-fr}
\EGAReferenceAlias{0.7.2.1.fr}{0.7.2.1-fr}
\EGAReferenceAlias{0.7.2.1.es}{0.7.2.1-fr}
\EGAReferenceAlias{0.7.2.1.eq}{0.7.2.1-fr}
\EGAReferenceAlias{eq:0.7.2.1}{0.7.2.1-fr}
\EGAReferenceAlias{0.7.2.2}{0.7.2.2-fr}
\EGAReferenceAlias{0.7.2.2-es}{0.7.2.2-fr}
\EGAReferenceAlias{0.7.2.2.fr}{0.7.2.2-fr}
\EGAReferenceAlias{0.7.2.2.es}{0.7.2.2-fr}
\EGAReferenceAlias{0.7.2.2.eq}{0.7.2.2-fr}
\EGAReferenceAlias{eq:0.7.2.2}{0.7.2.2-fr}
\EGAReferenceAlias{0.7.2.3}{0.7.2.3-fr}
\EGAReferenceAlias{0.7.2.3-es}{0.7.2.3-fr}
\EGAReferenceAlias{0.7.2.3.fr}{0.7.2.3-fr}
\EGAReferenceAlias{0.7.2.3.es}{0.7.2.3-fr}
\EGAReferenceAlias{0.7.2.3.eq}{0.7.2.3-fr}
\EGAReferenceAlias{eq:0.7.2.3}{0.7.2.3-fr}
\EGAReferenceAlias{0.7.2.4}{0.7.2.4-fr}
\EGAReferenceAlias{0.7.2.4-es}{0.7.2.4-fr}
\EGAReferenceAlias{0.7.2.4.fr}{0.7.2.4-fr}
\EGAReferenceAlias{0.7.2.4.es}{0.7.2.4-fr}
\EGAReferenceAlias{0.7.2.4.eq}{0.7.2.4-fr}
\EGAReferenceAlias{eq:0.7.2.4}{0.7.2.4-fr}
\EGAReferenceAlias{0.7.2.5}{0.7.2.5-fr}
\EGAReferenceAlias{0.7.2.5-es}{0.7.2.5-fr}
\EGAReferenceAlias{0.7.2.5.fr}{0.7.2.5-fr}
\EGAReferenceAlias{0.7.2.5.es}{0.7.2.5-fr}
\EGAReferenceAlias{0.7.2.5.eq}{0.7.2.5-fr}
\EGAReferenceAlias{eq:0.7.2.5}{0.7.2.5-fr}
\EGAReferenceAlias{0.7.2.6}{0.7.2.6-fr}
\EGAReferenceAlias{0.7.2.6-es}{0.7.2.6-fr}
\EGAReferenceAlias{0.7.2.6.fr}{0.7.2.6-fr}
\EGAReferenceAlias{0.7.2.6.es}{0.7.2.6-fr}
\EGAReferenceAlias{0.7.2.6.eq}{0.7.2.6-fr}
\EGAReferenceAlias{eq:0.7.2.6}{0.7.2.6-fr}
\EGAReferenceAlias{0.7.2.7}{0.7.2.7-fr}
\EGAReferenceAlias{0.7.2.7-es}{0.7.2.7-fr}
\EGAReferenceAlias{0.7.2.7.fr}{0.7.2.7-fr}
\EGAReferenceAlias{0.7.2.7.es}{0.7.2.7-fr}
\EGAReferenceAlias{0.7.2.7.eq}{0.7.2.7-fr}
\EGAReferenceAlias{eq:0.7.2.7}{0.7.2.7-fr}
\EGAReferenceAlias{0.7.2.8}{0.7.2.8-fr}
\EGAReferenceAlias{0.7.2.8-es}{0.7.2.8-fr}
\EGAReferenceAlias{0.7.2.8.fr}{0.7.2.8-fr}
\EGAReferenceAlias{0.7.2.8.es}{0.7.2.8-fr}
\EGAReferenceAlias{0.7.2.8.eq}{0.7.2.8-fr}
\EGAReferenceAlias{eq:0.7.2.8}{0.7.2.8-fr}
\EGAReferenceAlias{0.7.2.9}{0.7.2.9-fr}
\EGAReferenceAlias{0.7.2.9-es}{0.7.2.9-fr}
\EGAReferenceAlias{0.7.2.9.fr}{0.7.2.9-fr}
\EGAReferenceAlias{0.7.2.9.es}{0.7.2.9-fr}
\EGAReferenceAlias{0.7.2.9.eq}{0.7.2.9-fr}
\EGAReferenceAlias{eq:0.7.2.9}{0.7.2.9-fr}
\EGAReferenceAlias{0.7.2.10}{0.7.2.10-fr}
\EGAReferenceAlias{0.7.2.10-es}{0.7.2.10-fr}
\EGAReferenceAlias{0.7.2.10.fr}{0.7.2.10-fr}
\EGAReferenceAlias{0.7.2.10.es}{0.7.2.10-fr}
\EGAReferenceAlias{0.7.2.10.eq}{0.7.2.10-fr}
\EGAReferenceAlias{eq:0.7.2.10}{0.7.2.10-fr}
\EGAReferenceAlias{0.7.2.11}{0.7.2.11-fr}
\EGAReferenceAlias{0.7.2.11-es}{0.7.2.11-fr}
\EGAReferenceAlias{0.7.2.11.fr}{0.7.2.11-fr}
\EGAReferenceAlias{0.7.2.11.es}{0.7.2.11-fr}
\EGAReferenceAlias{0.7.2.11.eq}{0.7.2.11-fr}
\EGAReferenceAlias{eq:0.7.2.11}{0.7.2.11-fr}
\EGAReferenceAlias{0.7.2.12}{0.7.2.12-fr}
\EGAReferenceAlias{0.7.2.12-es}{0.7.2.12-fr}
\EGAReferenceAlias{0.7.2.12.fr}{0.7.2.12-fr}
\EGAReferenceAlias{0.7.2.12.es}{0.7.2.12-fr}
\EGAReferenceAlias{0.7.2.12.eq}{0.7.2.12-fr}
\EGAReferenceAlias{eq:0.7.2.12}{0.7.2.12-fr}
\EGAReferenceAlias{0.7.3}{subsection:0.7.3-fr}
\EGAReferenceAlias{0.7.3-fr}{subsection:0.7.3-fr}
\EGAReferenceAlias{0.7.3-es}{subsection:0.7.3-fr}
\EGAReferenceAlias{0.7.3.fr}{subsection:0.7.3-fr}
\EGAReferenceAlias{0.7.3.es}{subsection:0.7.3-fr}
\EGAReferenceAlias{0.7.3.eq}{subsection:0.7.3-fr}
\EGAReferenceAlias{eq:0.7.3}{subsection:0.7.3-fr}
\EGAReferenceAlias{subsection:0.7.3}{subsection:0.7.3-fr}
\EGAReferenceAlias{subsection:0.7.3-es}{subsection:0.7.3-fr}
\EGAReferenceAlias{subsection:0.7.3.fr}{subsection:0.7.3-fr}
\EGAReferenceAlias{subsection:0.7.3.es}{subsection:0.7.3-fr}
\EGAReferenceAlias{subsection:0.7.3.eq}{subsection:0.7.3-fr}
\EGAReferenceAlias{eq:subsection:0.7.3}{subsection:0.7.3-fr}
\EGAReferenceAlias{0.7.3.1}{0.7.3.1-fr}
\EGAReferenceAlias{0.7.3.1-es}{0.7.3.1-fr}
\EGAReferenceAlias{0.7.3.1.fr}{0.7.3.1-fr}
\EGAReferenceAlias{0.7.3.1.es}{0.7.3.1-fr}
\EGAReferenceAlias{0.7.3.1.eq}{0.7.3.1-fr}
\EGAReferenceAlias{eq:0.7.3.1}{0.7.3.1-fr}
\EGAReferenceAlias{0.7.3.2}{0.7.3.2-fr}
\EGAReferenceAlias{0.7.3.2-es}{0.7.3.2-fr}
\EGAReferenceAlias{0.7.3.2.fr}{0.7.3.2-fr}
\EGAReferenceAlias{0.7.3.2.es}{0.7.3.2-fr}
\EGAReferenceAlias{0.7.3.2.eq}{0.7.3.2-fr}
\EGAReferenceAlias{eq:0.7.3.2}{0.7.3.2-fr}
\EGAReferenceAlias{0.7.3.2.1}{0.7.3.2.1-fr}
\EGAReferenceAlias{0.7.3.2.1-es}{0.7.3.2.1-fr}
\EGAReferenceAlias{0.7.3.2.1.fr}{0.7.3.2.1-fr}
\EGAReferenceAlias{0.7.3.2.1.es}{0.7.3.2.1-fr}
\EGAReferenceAlias{0.7.3.2.1.eq}{0.7.3.2.1-fr}
\EGAReferenceAlias{eq:0.7.3.2.1}{0.7.3.2.1-fr}
\EGAReferenceAlias{0.7.3.3}{0.7.3.3-fr}
\EGAReferenceAlias{0.7.3.3-es}{0.7.3.3-fr}
\EGAReferenceAlias{0.7.3.3.fr}{0.7.3.3-fr}
\EGAReferenceAlias{0.7.3.3.es}{0.7.3.3-fr}
\EGAReferenceAlias{0.7.3.3.eq}{0.7.3.3-fr}
\EGAReferenceAlias{eq:0.7.3.3}{0.7.3.3-fr}
\EGAReferenceAlias{0.7.3.4}{0.7.3.4-fr}
\EGAReferenceAlias{0.7.3.4-es}{0.7.3.4-fr}
\EGAReferenceAlias{0.7.3.4.fr}{0.7.3.4-fr}
\EGAReferenceAlias{0.7.3.4.es}{0.7.3.4-fr}
\EGAReferenceAlias{0.7.3.4.eq}{0.7.3.4-fr}
\EGAReferenceAlias{eq:0.7.3.4}{0.7.3.4-fr}
\EGAReferenceAlias{0.7.3.5}{0.7.3.5-fr}
\EGAReferenceAlias{0.7.3.5-es}{0.7.3.5-fr}
\EGAReferenceAlias{0.7.3.5.fr}{0.7.3.5-fr}
\EGAReferenceAlias{0.7.3.5.es}{0.7.3.5-fr}
\EGAReferenceAlias{0.7.3.5.eq}{0.7.3.5-fr}
\EGAReferenceAlias{eq:0.7.3.5}{0.7.3.5-fr}
\EGAReferenceAlias{0.7.3.6}{0.7.3.6-fr}
\EGAReferenceAlias{0.7.3.6-es}{0.7.3.6-fr}
\EGAReferenceAlias{0.7.3.6.fr}{0.7.3.6-fr}
\EGAReferenceAlias{0.7.3.6.es}{0.7.3.6-fr}
\EGAReferenceAlias{0.7.3.6.eq}{0.7.3.6-fr}
\EGAReferenceAlias{eq:0.7.3.6}{0.7.3.6-fr}
\EGAReferenceAlias{0.7.3.7}{0.7.3.7-fr}
\EGAReferenceAlias{0.7.3.7-es}{0.7.3.7-fr}
\EGAReferenceAlias{0.7.3.7.fr}{0.7.3.7-fr}
\EGAReferenceAlias{0.7.3.7.es}{0.7.3.7-fr}
\EGAReferenceAlias{0.7.3.7.eq}{0.7.3.7-fr}
\EGAReferenceAlias{eq:0.7.3.7}{0.7.3.7-fr}
\EGAReferenceAlias{0.7.4}{subsection:0.7.4-fr}
\EGAReferenceAlias{0.7.4-fr}{subsection:0.7.4-fr}
\EGAReferenceAlias{0.7.4-es}{subsection:0.7.4-fr}
\EGAReferenceAlias{0.7.4.fr}{subsection:0.7.4-fr}
\EGAReferenceAlias{0.7.4.es}{subsection:0.7.4-fr}
\EGAReferenceAlias{0.7.4.eq}{subsection:0.7.4-fr}
\EGAReferenceAlias{eq:0.7.4}{subsection:0.7.4-fr}
\EGAReferenceAlias{subsection:0.7.4}{subsection:0.7.4-fr}
\EGAReferenceAlias{subsection:0.7.4-es}{subsection:0.7.4-fr}
\EGAReferenceAlias{subsection:0.7.4.fr}{subsection:0.7.4-fr}
\EGAReferenceAlias{subsection:0.7.4.es}{subsection:0.7.4-fr}
\EGAReferenceAlias{subsection:0.7.4.eq}{subsection:0.7.4-fr}
\EGAReferenceAlias{eq:subsection:0.7.4}{subsection:0.7.4-fr}
\EGAReferenceAlias{0.7.4.1}{0.7.4.1-fr}
\EGAReferenceAlias{0.7.4.1-es}{0.7.4.1-fr}
\EGAReferenceAlias{0.7.4.1.fr}{0.7.4.1-fr}
\EGAReferenceAlias{0.7.4.1.es}{0.7.4.1-fr}
\EGAReferenceAlias{0.7.4.1.eq}{0.7.4.1-fr}
\EGAReferenceAlias{eq:0.7.4.1}{0.7.4.1-fr}
\EGAReferenceAlias{0.7.4.2}{0.7.4.2-fr}
\EGAReferenceAlias{0.7.4.2-es}{0.7.4.2-fr}
\EGAReferenceAlias{0.7.4.2.fr}{0.7.4.2-fr}
\EGAReferenceAlias{0.7.4.2.es}{0.7.4.2-fr}
\EGAReferenceAlias{0.7.4.2.eq}{0.7.4.2-fr}
\EGAReferenceAlias{eq:0.7.4.2}{0.7.4.2-fr}
\EGAReferenceAlias{0.7.4.3}{0.7.4.3-fr}
\EGAReferenceAlias{0.7.4.3-es}{0.7.4.3-fr}
\EGAReferenceAlias{0.7.4.3.fr}{0.7.4.3-fr}
\EGAReferenceAlias{0.7.4.3.es}{0.7.4.3-fr}
\EGAReferenceAlias{0.7.4.3.eq}{0.7.4.3-fr}
\EGAReferenceAlias{eq:0.7.4.3}{0.7.4.3-fr}
\EGAReferenceAlias{0.7.4.4}{0.7.4.4-fr}
\EGAReferenceAlias{0.7.4.4-es}{0.7.4.4-fr}
\EGAReferenceAlias{0.7.4.4.fr}{0.7.4.4-fr}
\EGAReferenceAlias{0.7.4.4.es}{0.7.4.4-fr}
\EGAReferenceAlias{0.7.4.4.eq}{0.7.4.4-fr}
\EGAReferenceAlias{eq:0.7.4.4}{0.7.4.4-fr}
\EGAReferenceAlias{0.7.4.5}{0.7.4.5-fr}
\EGAReferenceAlias{0.7.4.5-es}{0.7.4.5-fr}
\EGAReferenceAlias{0.7.4.5.fr}{0.7.4.5-fr}
\EGAReferenceAlias{0.7.4.5.es}{0.7.4.5-fr}
\EGAReferenceAlias{0.7.4.5.eq}{0.7.4.5-fr}
\EGAReferenceAlias{eq:0.7.4.5}{0.7.4.5-fr}
\EGAReferenceAlias{0.7.5}{subsection:0.7.5-fr}
\EGAReferenceAlias{0.7.5-fr}{subsection:0.7.5-fr}
\EGAReferenceAlias{0.7.5-es}{subsection:0.7.5-fr}
\EGAReferenceAlias{0.7.5.fr}{subsection:0.7.5-fr}
\EGAReferenceAlias{0.7.5.es}{subsection:0.7.5-fr}
\EGAReferenceAlias{0.7.5.eq}{subsection:0.7.5-fr}
\EGAReferenceAlias{eq:0.7.5}{subsection:0.7.5-fr}
\EGAReferenceAlias{subsection:0.7.5}{subsection:0.7.5-fr}
\EGAReferenceAlias{subsection:0.7.5-es}{subsection:0.7.5-fr}
\EGAReferenceAlias{subsection:0.7.5.fr}{subsection:0.7.5-fr}
\EGAReferenceAlias{subsection:0.7.5.es}{subsection:0.7.5-fr}
\EGAReferenceAlias{subsection:0.7.5.eq}{subsection:0.7.5-fr}
\EGAReferenceAlias{eq:subsection:0.7.5}{subsection:0.7.5-fr}
\EGAReferenceAlias{0.7.5.1}{0.7.5.1-fr}
\EGAReferenceAlias{0.7.5.1-es}{0.7.5.1-fr}
\EGAReferenceAlias{0.7.5.1.fr}{0.7.5.1-fr}
\EGAReferenceAlias{0.7.5.1.es}{0.7.5.1-fr}
\EGAReferenceAlias{0.7.5.1.eq}{0.7.5.1-fr}
\EGAReferenceAlias{eq:0.7.5.1}{0.7.5.1-fr}
\EGAReferenceAlias{0.7.5.2}{0.7.5.2-fr}
\EGAReferenceAlias{0.7.5.2-es}{0.7.5.2-fr}
\EGAReferenceAlias{0.7.5.2.fr}{0.7.5.2-fr}
\EGAReferenceAlias{0.7.5.2.es}{0.7.5.2-fr}
\EGAReferenceAlias{0.7.5.2.eq}{0.7.5.2-fr}
\EGAReferenceAlias{eq:0.7.5.2}{0.7.5.2-fr}
\EGAReferenceAlias{0.7.5.3}{0.7.5.3-fr}
\EGAReferenceAlias{0.7.5.3-es}{0.7.5.3-fr}
\EGAReferenceAlias{0.7.5.3.fr}{0.7.5.3-fr}
\EGAReferenceAlias{0.7.5.3.es}{0.7.5.3-fr}
\EGAReferenceAlias{0.7.5.3.eq}{0.7.5.3-fr}
\EGAReferenceAlias{eq:0.7.5.3}{0.7.5.3-fr}
\EGAReferenceAlias{0.7.5.4}{0.7.5.4-fr}
\EGAReferenceAlias{0.7.5.4-es}{0.7.5.4-fr}
\EGAReferenceAlias{0.7.5.4.fr}{0.7.5.4-fr}
\EGAReferenceAlias{0.7.5.4.es}{0.7.5.4-fr}
\EGAReferenceAlias{0.7.5.4.eq}{0.7.5.4-fr}
\EGAReferenceAlias{eq:0.7.5.4}{0.7.5.4-fr}
\EGAReferenceAlias{0.7.5.5}{0.7.5.5-fr}
\EGAReferenceAlias{0.7.5.5-es}{0.7.5.5-fr}
\EGAReferenceAlias{0.7.5.5.fr}{0.7.5.5-fr}
\EGAReferenceAlias{0.7.5.5.es}{0.7.5.5-fr}
\EGAReferenceAlias{0.7.5.5.eq}{0.7.5.5-fr}
\EGAReferenceAlias{eq:0.7.5.5}{0.7.5.5-fr}
\EGAReferenceAlias{0.7.6}{subsection:0.7.6-fr}
\EGAReferenceAlias{0.7.6-fr}{subsection:0.7.6-fr}
\EGAReferenceAlias{0.7.6-es}{subsection:0.7.6-fr}
\EGAReferenceAlias{0.7.6.fr}{subsection:0.7.6-fr}
\EGAReferenceAlias{0.7.6.es}{subsection:0.7.6-fr}
\EGAReferenceAlias{0.7.6.eq}{subsection:0.7.6-fr}
\EGAReferenceAlias{eq:0.7.6}{subsection:0.7.6-fr}
\EGAReferenceAlias{subsection:0.7.6}{subsection:0.7.6-fr}
\EGAReferenceAlias{subsection:0.7.6-es}{subsection:0.7.6-fr}
\EGAReferenceAlias{subsection:0.7.6.fr}{subsection:0.7.6-fr}
\EGAReferenceAlias{subsection:0.7.6.es}{subsection:0.7.6-fr}
\EGAReferenceAlias{subsection:0.7.6.eq}{subsection:0.7.6-fr}
\EGAReferenceAlias{eq:subsection:0.7.6}{subsection:0.7.6-fr}
\EGAReferenceAlias{0.7.6.1}{0.7.6.1-fr}
\EGAReferenceAlias{0.7.6.1-es}{0.7.6.1-fr}
\EGAReferenceAlias{0.7.6.1.fr}{0.7.6.1-fr}
\EGAReferenceAlias{0.7.6.1.es}{0.7.6.1-fr}
\EGAReferenceAlias{0.7.6.1.eq}{0.7.6.1-fr}
\EGAReferenceAlias{eq:0.7.6.1}{0.7.6.1-fr}
\EGAReferenceAlias{0.7.6.2}{0.7.6.2-fr}
\EGAReferenceAlias{0.7.6.2-es}{0.7.6.2-fr}
\EGAReferenceAlias{0.7.6.2.fr}{0.7.6.2-fr}
\EGAReferenceAlias{0.7.6.2.es}{0.7.6.2-fr}
\EGAReferenceAlias{0.7.6.2.eq}{0.7.6.2-fr}
\EGAReferenceAlias{eq:0.7.6.2}{0.7.6.2-fr}
\EGAReferenceAlias{0.7.6.3}{0.7.6.3-fr}
\EGAReferenceAlias{0.7.6.3-es}{0.7.6.3-fr}
\EGAReferenceAlias{0.7.6.3.fr}{0.7.6.3-fr}
\EGAReferenceAlias{0.7.6.3.es}{0.7.6.3-fr}
\EGAReferenceAlias{0.7.6.3.eq}{0.7.6.3-fr}
\EGAReferenceAlias{eq:0.7.6.3}{0.7.6.3-fr}
\EGAReferenceAlias{0.7.6.4}{0.7.6.4-fr}
\EGAReferenceAlias{0.7.6.4-es}{0.7.6.4-fr}
\EGAReferenceAlias{0.7.6.4.fr}{0.7.6.4-fr}
\EGAReferenceAlias{0.7.6.4.es}{0.7.6.4-fr}
\EGAReferenceAlias{0.7.6.4.eq}{0.7.6.4-fr}
\EGAReferenceAlias{eq:0.7.6.4}{0.7.6.4-fr}
\EGAReferenceAlias{0.7.6.5}{0.7.6.5-fr}
\EGAReferenceAlias{0.7.6.5-es}{0.7.6.5-fr}
\EGAReferenceAlias{0.7.6.5.fr}{0.7.6.5-fr}
\EGAReferenceAlias{0.7.6.5.es}{0.7.6.5-fr}
\EGAReferenceAlias{0.7.6.5.eq}{0.7.6.5-fr}
\EGAReferenceAlias{eq:0.7.6.5}{0.7.6.5-fr}
\EGAReferenceAlias{0.7.6.6}{0.7.6.6-fr}
\EGAReferenceAlias{0.7.6.6-es}{0.7.6.6-fr}
\EGAReferenceAlias{0.7.6.6.fr}{0.7.6.6-fr}
\EGAReferenceAlias{0.7.6.6.es}{0.7.6.6-fr}
\EGAReferenceAlias{0.7.6.6.eq}{0.7.6.6-fr}
\EGAReferenceAlias{eq:0.7.6.6}{0.7.6.6-fr}
\EGAReferenceAlias{0.7.6.7}{0.7.6.7-fr}
\EGAReferenceAlias{0.7.6.7-es}{0.7.6.7-fr}
\EGAReferenceAlias{0.7.6.7.fr}{0.7.6.7-fr}
\EGAReferenceAlias{0.7.6.7.es}{0.7.6.7-fr}
\EGAReferenceAlias{0.7.6.7.eq}{0.7.6.7-fr}
\EGAReferenceAlias{eq:0.7.6.7}{0.7.6.7-fr}
\EGAReferenceAlias{0.7.6.8}{0.7.6.8-fr}
\EGAReferenceAlias{0.7.6.8-es}{0.7.6.8-fr}
\EGAReferenceAlias{0.7.6.8.fr}{0.7.6.8-fr}
\EGAReferenceAlias{0.7.6.8.es}{0.7.6.8-fr}
\EGAReferenceAlias{0.7.6.8.eq}{0.7.6.8-fr}
\EGAReferenceAlias{eq:0.7.6.8}{0.7.6.8-fr}
\EGAReferenceAlias{0.7.6.9}{0.7.6.9-fr}
\EGAReferenceAlias{0.7.6.9-es}{0.7.6.9-fr}
\EGAReferenceAlias{0.7.6.9.fr}{0.7.6.9-fr}
\EGAReferenceAlias{0.7.6.9.es}{0.7.6.9-fr}
\EGAReferenceAlias{0.7.6.9.eq}{0.7.6.9-fr}
\EGAReferenceAlias{eq:0.7.6.9}{0.7.6.9-fr}
\EGAReferenceAlias{0.7.6.10}{0.7.6.10-fr}
\EGAReferenceAlias{0.7.6.10-es}{0.7.6.10-fr}
\EGAReferenceAlias{0.7.6.10.fr}{0.7.6.10-fr}
\EGAReferenceAlias{0.7.6.10.es}{0.7.6.10-fr}
\EGAReferenceAlias{0.7.6.10.eq}{0.7.6.10-fr}
\EGAReferenceAlias{eq:0.7.6.10}{0.7.6.10-fr}
\EGAReferenceAlias{0.7.6.11}{0.7.6.11-fr}
\EGAReferenceAlias{0.7.6.11-es}{0.7.6.11-fr}
\EGAReferenceAlias{0.7.6.11.fr}{0.7.6.11-fr}
\EGAReferenceAlias{0.7.6.11.es}{0.7.6.11-fr}
\EGAReferenceAlias{0.7.6.11.eq}{0.7.6.11-fr}
\EGAReferenceAlias{eq:0.7.6.11}{0.7.6.11-fr}
\EGAReferenceAlias{0.7.6.12}{0.7.6.12-fr}
\EGAReferenceAlias{0.7.6.12-es}{0.7.6.12-fr}
\EGAReferenceAlias{0.7.6.12.fr}{0.7.6.12-fr}
\EGAReferenceAlias{0.7.6.12.es}{0.7.6.12-fr}
\EGAReferenceAlias{0.7.6.12.eq}{0.7.6.12-fr}
\EGAReferenceAlias{eq:0.7.6.12}{0.7.6.12-fr}
\EGAReferenceAlias{0.7.6.13}{0.7.6.13-fr}
\EGAReferenceAlias{0.7.6.13-es}{0.7.6.13-fr}
\EGAReferenceAlias{0.7.6.13.fr}{0.7.6.13-fr}
\EGAReferenceAlias{0.7.6.13.es}{0.7.6.13-fr}
\EGAReferenceAlias{0.7.6.13.eq}{0.7.6.13-fr}
\EGAReferenceAlias{eq:0.7.6.13}{0.7.6.13-fr}
\EGAReferenceAlias{0.7.6.14}{0.7.6.14-fr}
\EGAReferenceAlias{0.7.6.14-es}{0.7.6.14-fr}
\EGAReferenceAlias{0.7.6.14.fr}{0.7.6.14-fr}
\EGAReferenceAlias{0.7.6.14.es}{0.7.6.14-fr}
\EGAReferenceAlias{0.7.6.14.eq}{0.7.6.14-fr}
\EGAReferenceAlias{eq:0.7.6.14}{0.7.6.14-fr}
\EGAReferenceAlias{0.7.6.15}{0.7.6.15-fr}
\EGAReferenceAlias{0.7.6.15-es}{0.7.6.15-fr}
\EGAReferenceAlias{0.7.6.15.fr}{0.7.6.15-fr}
\EGAReferenceAlias{0.7.6.15.es}{0.7.6.15-fr}
\EGAReferenceAlias{0.7.6.15.eq}{0.7.6.15-fr}
\EGAReferenceAlias{eq:0.7.6.15}{0.7.6.15-fr}
\EGAReferenceAlias{0.7.6.16}{0.7.6.16-fr}
\EGAReferenceAlias{0.7.6.16-es}{0.7.6.16-fr}
\EGAReferenceAlias{0.7.6.16.fr}{0.7.6.16-fr}
\EGAReferenceAlias{0.7.6.16.es}{0.7.6.16-fr}
\EGAReferenceAlias{0.7.6.16.eq}{0.7.6.16-fr}
\EGAReferenceAlias{eq:0.7.6.16}{0.7.6.16-fr}
\EGAReferenceAlias{0.7.6.17}{0.7.6.17-fr}
\EGAReferenceAlias{0.7.6.17-es}{0.7.6.17-fr}
\EGAReferenceAlias{0.7.6.17.fr}{0.7.6.17-fr}
\EGAReferenceAlias{0.7.6.17.es}{0.7.6.17-fr}
\EGAReferenceAlias{0.7.6.17.eq}{0.7.6.17-fr}
\EGAReferenceAlias{eq:0.7.6.17}{0.7.6.17-fr}
\EGAReferenceAlias{0.7.6.18}{0.7.6.18-fr}
\EGAReferenceAlias{0.7.6.18-es}{0.7.6.18-fr}
\EGAReferenceAlias{0.7.6.18.fr}{0.7.6.18-fr}
\EGAReferenceAlias{0.7.6.18.es}{0.7.6.18-fr}
\EGAReferenceAlias{0.7.6.18.eq}{0.7.6.18-fr}
\EGAReferenceAlias{eq:0.7.6.18}{0.7.6.18-fr}
\EGAReferenceAlias{0.7.7}{subsection:0.7.7-fr}
\EGAReferenceAlias{0.7.7-fr}{subsection:0.7.7-fr}
\EGAReferenceAlias{0.7.7-es}{subsection:0.7.7-fr}
\EGAReferenceAlias{0.7.7.fr}{subsection:0.7.7-fr}
\EGAReferenceAlias{0.7.7.es}{subsection:0.7.7-fr}
\EGAReferenceAlias{0.7.7.eq}{subsection:0.7.7-fr}
\EGAReferenceAlias{eq:0.7.7}{subsection:0.7.7-fr}
\EGAReferenceAlias{subsection:0.7.7}{subsection:0.7.7-fr}
\EGAReferenceAlias{subsection:0.7.7-es}{subsection:0.7.7-fr}
\EGAReferenceAlias{subsection:0.7.7.fr}{subsection:0.7.7-fr}
\EGAReferenceAlias{subsection:0.7.7.es}{subsection:0.7.7-fr}
\EGAReferenceAlias{subsection:0.7.7.eq}{subsection:0.7.7-fr}
\EGAReferenceAlias{eq:subsection:0.7.7}{subsection:0.7.7-fr}
\EGAReferenceAlias{0.7.7.1}{0.7.7.1-fr}
\EGAReferenceAlias{0.7.7.1-es}{0.7.7.1-fr}
\EGAReferenceAlias{0.7.7.1.fr}{0.7.7.1-fr}
\EGAReferenceAlias{0.7.7.1.es}{0.7.7.1-fr}
\EGAReferenceAlias{0.7.7.1.eq}{0.7.7.1-fr}
\EGAReferenceAlias{eq:0.7.7.1}{0.7.7.1-fr}
\EGAReferenceAlias{0.7.7.2}{0.7.7.2-fr}
\EGAReferenceAlias{0.7.7.2-es}{0.7.7.2-fr}
\EGAReferenceAlias{0.7.7.2.fr}{0.7.7.2-fr}
\EGAReferenceAlias{0.7.7.2.es}{0.7.7.2-fr}
\EGAReferenceAlias{0.7.7.2.eq}{0.7.7.2-fr}
\EGAReferenceAlias{eq:0.7.7.2}{0.7.7.2-fr}
\EGAReferenceAlias{0.7.7.3}{0.7.7.3-fr}
\EGAReferenceAlias{0.7.7.3-es}{0.7.7.3-fr}
\EGAReferenceAlias{0.7.7.3.fr}{0.7.7.3-fr}
\EGAReferenceAlias{0.7.7.3.es}{0.7.7.3-fr}
\EGAReferenceAlias{0.7.7.3.eq}{0.7.7.3-fr}
\EGAReferenceAlias{eq:0.7.7.3}{0.7.7.3-fr}
\EGAReferenceAlias{0.7.7.4}{0.7.7.4-fr}
\EGAReferenceAlias{0.7.7.4-es}{0.7.7.4-fr}
\EGAReferenceAlias{0.7.7.4.fr}{0.7.7.4-fr}
\EGAReferenceAlias{0.7.7.4.es}{0.7.7.4-fr}
\EGAReferenceAlias{0.7.7.4.eq}{0.7.7.4-fr}
\EGAReferenceAlias{eq:0.7.7.4}{0.7.7.4-fr}
\EGAReferenceAlias{0.7.7.5}{0.7.7.5-fr}
\EGAReferenceAlias{0.7.7.5-es}{0.7.7.5-fr}
\EGAReferenceAlias{0.7.7.5.fr}{0.7.7.5-fr}
\EGAReferenceAlias{0.7.7.5.es}{0.7.7.5-fr}
\EGAReferenceAlias{0.7.7.5.eq}{0.7.7.5-fr}
\EGAReferenceAlias{eq:0.7.7.5}{0.7.7.5-fr}
\EGAReferenceAlias{0.7.7.6}{0.7.7.6-fr}
\EGAReferenceAlias{0.7.7.6-es}{0.7.7.6-fr}
\EGAReferenceAlias{0.7.7.6.fr}{0.7.7.6-fr}
\EGAReferenceAlias{0.7.7.6.es}{0.7.7.6-fr}
\EGAReferenceAlias{0.7.7.6.eq}{0.7.7.6-fr}
\EGAReferenceAlias{eq:0.7.7.6}{0.7.7.6-fr}
\EGAReferenceAlias{0.7.7.7}{0.7.7.7-fr}
\EGAReferenceAlias{0.7.7.7-es}{0.7.7.7-fr}
\EGAReferenceAlias{0.7.7.7.fr}{0.7.7.7-fr}
\EGAReferenceAlias{0.7.7.7.es}{0.7.7.7-fr}
\EGAReferenceAlias{0.7.7.7.eq}{0.7.7.7-fr}
\EGAReferenceAlias{eq:0.7.7.7}{0.7.7.7-fr}
\EGAReferenceAlias{0.7.7.8}{0.7.7.8-fr}
\EGAReferenceAlias{0.7.7.8-es}{0.7.7.8-fr}
\EGAReferenceAlias{0.7.7.8.fr}{0.7.7.8-fr}
\EGAReferenceAlias{0.7.7.8.es}{0.7.7.8-fr}
\EGAReferenceAlias{0.7.7.8.eq}{0.7.7.8-fr}
\EGAReferenceAlias{eq:0.7.7.8}{0.7.7.8-fr}
\EGAReferenceAlias{0.7.8}{subsection:0.7.8-fr}
\EGAReferenceAlias{0.7.8-fr}{subsection:0.7.8-fr}
\EGAReferenceAlias{0.7.8-es}{subsection:0.7.8-fr}
\EGAReferenceAlias{0.7.8.fr}{subsection:0.7.8-fr}
\EGAReferenceAlias{0.7.8.es}{subsection:0.7.8-fr}
\EGAReferenceAlias{0.7.8.eq}{subsection:0.7.8-fr}
\EGAReferenceAlias{eq:0.7.8}{subsection:0.7.8-fr}
\EGAReferenceAlias{subsection:0.7.8}{subsection:0.7.8-fr}
\EGAReferenceAlias{subsection:0.7.8-es}{subsection:0.7.8-fr}
\EGAReferenceAlias{subsection:0.7.8.fr}{subsection:0.7.8-fr}
\EGAReferenceAlias{subsection:0.7.8.es}{subsection:0.7.8-fr}
\EGAReferenceAlias{subsection:0.7.8.eq}{subsection:0.7.8-fr}
\EGAReferenceAlias{eq:subsection:0.7.8}{subsection:0.7.8-fr}
\EGAReferenceAlias{0.7.8.1}{0.7.8.1-fr}
\EGAReferenceAlias{0.7.8.1-es}{0.7.8.1-fr}
\EGAReferenceAlias{0.7.8.1.fr}{0.7.8.1-fr}
\EGAReferenceAlias{0.7.8.1.es}{0.7.8.1-fr}
\EGAReferenceAlias{0.7.8.1.eq}{0.7.8.1-fr}
\EGAReferenceAlias{eq:0.7.8.1}{0.7.8.1-fr}
\EGAReferenceAlias{0.7.8.2}{0.7.8.2-fr}
\EGAReferenceAlias{0.7.8.2-es}{0.7.8.2-fr}
\EGAReferenceAlias{0.7.8.2.fr}{0.7.8.2-fr}
\EGAReferenceAlias{0.7.8.2.es}{0.7.8.2-fr}
\EGAReferenceAlias{0.7.8.2.eq}{0.7.8.2-fr}
\EGAReferenceAlias{eq:0.7.8.2}{0.7.8.2-fr}
\EGAReferenceAlias{0.7.8.3}{0.7.8.3-fr}
\EGAReferenceAlias{0.7.8.3-es}{0.7.8.3-fr}
\EGAReferenceAlias{0.7.8.3.fr}{0.7.8.3-fr}
\EGAReferenceAlias{0.7.8.3.es}{0.7.8.3-fr}
\EGAReferenceAlias{0.7.8.3.eq}{0.7.8.3-fr}
\EGAReferenceAlias{eq:0.7.8.3}{0.7.8.3-fr}
\EGAReferenceAlias{I}{chapter:I-fr}
\EGAReferenceAlias{I-fr}{chapter:I-fr}
\EGAReferenceAlias{I-es}{chapter:I-fr}
\EGAReferenceAlias{I.fr}{chapter:I-fr}
\EGAReferenceAlias{I.es}{chapter:I-fr}
\EGAReferenceAlias{I.eq}{chapter:I-fr}
\EGAReferenceAlias{eq:I}{chapter:I-fr}
\EGAReferenceAlias{chapter:I}{chapter:I-fr}
\EGAReferenceAlias{chapter:I-es}{chapter:I-fr}
\EGAReferenceAlias{chapter:I.fr}{chapter:I-fr}
\EGAReferenceAlias{chapter:I.es}{chapter:I-fr}
\EGAReferenceAlias{chapter:I.eq}{chapter:I-fr}
\EGAReferenceAlias{eq:chapter:I}{chapter:I-fr}
\EGAReferenceAlias{I.1}{section:I.1}
\EGAReferenceAlias{I.1-fr}{section:I.1}
\EGAReferenceAlias{I.1-es}{section:I.1}
\EGAReferenceAlias{I.1.fr}{section:I.1}
\EGAReferenceAlias{I.1.es}{section:I.1}
\EGAReferenceAlias{I.1.eq}{section:I.1}
\EGAReferenceAlias{eq:I.1}{section:I.1}
\EGAReferenceAlias{section:I.1-fr}{section:I.1}
\EGAReferenceAlias{section:I.1-es}{section:I.1}
\EGAReferenceAlias{section:I.1.fr}{section:I.1}
\EGAReferenceAlias{section:I.1.es}{section:I.1}
\EGAReferenceAlias{section:I.1.eq}{section:I.1}
\EGAReferenceAlias{eq:section:I.1}{section:I.1}
\EGAReferenceAlias{I.1.1}{subsection:I.1.1}
\EGAReferenceAlias{I.1.1-fr}{subsection:I.1.1}
\EGAReferenceAlias{I.1.1-es}{subsection:I.1.1}
\EGAReferenceAlias{I.1.1.fr}{subsection:I.1.1}
\EGAReferenceAlias{I.1.1.es}{subsection:I.1.1}
\EGAReferenceAlias{I.1.1.eq}{subsection:I.1.1}
\EGAReferenceAlias{eq:I.1.1}{subsection:I.1.1}
\EGAReferenceAlias{subsection:I.1.1-fr}{subsection:I.1.1}
\EGAReferenceAlias{subsection:I.1.1-es}{subsection:I.1.1}
\EGAReferenceAlias{subsection:I.1.1.fr}{subsection:I.1.1}
\EGAReferenceAlias{subsection:I.1.1.es}{subsection:I.1.1}
\EGAReferenceAlias{subsection:I.1.1.eq}{subsection:I.1.1}
\EGAReferenceAlias{eq:subsection:I.1.1}{subsection:I.1.1}
\EGAReferenceAlias{I.1.1.1-fr}{I.1.1.1}
\EGAReferenceAlias{I.1.1.1-es}{I.1.1.1}
\EGAReferenceAlias{I.1.1.1.fr}{I.1.1.1}
\EGAReferenceAlias{I.1.1.1.es}{I.1.1.1}
\EGAReferenceAlias{I.1.1.1.eq}{I.1.1.1}
\EGAReferenceAlias{eq:I.1.1.1}{I.1.1.1}
\EGAReferenceAlias{I.1.1.1.1}{I.1.1.1.1-fr}
\EGAReferenceAlias{I.1.1.1.1-es}{I.1.1.1.1-fr}
\EGAReferenceAlias{I.1.1.1.1.fr}{I.1.1.1.1-fr}
\EGAReferenceAlias{I.1.1.1.1.es}{I.1.1.1.1-fr}
\EGAReferenceAlias{I.1.1.1.1.eq}{I.1.1.1.1-fr}
\EGAReferenceAlias{eq:I.1.1.1.1}{I.1.1.1.1-fr}
\EGAReferenceAlias{I.1.1.2-fr}{I.1.1.2}
\EGAReferenceAlias{I.1.1.2-es}{I.1.1.2}
\EGAReferenceAlias{I.1.1.2.fr}{I.1.1.2}
\EGAReferenceAlias{I.1.1.2.es}{I.1.1.2}
\EGAReferenceAlias{I.1.1.2.eq}{I.1.1.2}
\EGAReferenceAlias{eq:I.1.1.2}{I.1.1.2}
\EGAReferenceAlias{I.1.1.3-fr}{I.1.1.3}
\EGAReferenceAlias{I.1.1.3-es}{I.1.1.3}
\EGAReferenceAlias{I.1.1.3.fr}{I.1.1.3}
\EGAReferenceAlias{I.1.1.3.es}{I.1.1.3}
\EGAReferenceAlias{I.1.1.3.eq}{I.1.1.3}
\EGAReferenceAlias{eq:I.1.1.3}{I.1.1.3}
\EGAReferenceAlias{I.1.1.3.1}{I.1.1.3.1-fr}
\EGAReferenceAlias{I.1.1.3.1-es}{I.1.1.3.1-fr}
\EGAReferenceAlias{I.1.1.3.1.fr}{I.1.1.3.1-fr}
\EGAReferenceAlias{I.1.1.3.1.es}{I.1.1.3.1-fr}
\EGAReferenceAlias{I.1.1.3.1.eq}{I.1.1.3.1-fr}
\EGAReferenceAlias{eq:I.1.1.3.1}{I.1.1.3.1-fr}
\EGAReferenceAlias{I.1.1.3.2}{I.1.1.3.2-fr}
\EGAReferenceAlias{I.1.1.3.2-es}{I.1.1.3.2-fr}
\EGAReferenceAlias{I.1.1.3.2.fr}{I.1.1.3.2-fr}
\EGAReferenceAlias{I.1.1.3.2.es}{I.1.1.3.2-fr}
\EGAReferenceAlias{I.1.1.3.2.eq}{I.1.1.3.2-fr}
\EGAReferenceAlias{eq:I.1.1.3.2}{I.1.1.3.2-fr}
\EGAReferenceAlias{I.1.1.4-fr}{I.1.1.4}
\EGAReferenceAlias{I.1.1.4-es}{I.1.1.4}
\EGAReferenceAlias{I.1.1.4.fr}{I.1.1.4}
\EGAReferenceAlias{I.1.1.4.es}{I.1.1.4}
\EGAReferenceAlias{I.1.1.4.eq}{I.1.1.4}
\EGAReferenceAlias{eq:I.1.1.4}{I.1.1.4}
\EGAReferenceAlias{I.1.1.5-fr}{I.1.1.5}
\EGAReferenceAlias{I.1.1.5-es}{I.1.1.5}
\EGAReferenceAlias{I.1.1.5.fr}{I.1.1.5}
\EGAReferenceAlias{I.1.1.5.es}{I.1.1.5}
\EGAReferenceAlias{I.1.1.5.eq}{I.1.1.5}
\EGAReferenceAlias{eq:I.1.1.5}{I.1.1.5}
\EGAReferenceAlias{I.1.1.6-fr}{I.1.1.6}
\EGAReferenceAlias{I.1.1.6-es}{I.1.1.6}
\EGAReferenceAlias{I.1.1.6.fr}{I.1.1.6}
\EGAReferenceAlias{I.1.1.6.es}{I.1.1.6}
\EGAReferenceAlias{I.1.1.6.eq}{I.1.1.6}
\EGAReferenceAlias{eq:I.1.1.6}{I.1.1.6}
\EGAReferenceAlias{I.1.1.7-fr}{I.1.1.7}
\EGAReferenceAlias{I.1.1.7-es}{I.1.1.7}
\EGAReferenceAlias{I.1.1.7.fr}{I.1.1.7}
\EGAReferenceAlias{I.1.1.7.es}{I.1.1.7}
\EGAReferenceAlias{I.1.1.7.eq}{I.1.1.7}
\EGAReferenceAlias{eq:I.1.1.7}{I.1.1.7}
\EGAReferenceAlias{I.1.1.8-fr}{I.1.1.8}
\EGAReferenceAlias{I.1.1.8-es}{I.1.1.8}
\EGAReferenceAlias{I.1.1.8.fr}{I.1.1.8}
\EGAReferenceAlias{I.1.1.8.es}{I.1.1.8}
\EGAReferenceAlias{I.1.1.8.eq}{I.1.1.8}
\EGAReferenceAlias{eq:I.1.1.8}{I.1.1.8}
\EGAReferenceAlias{I.1.1.9-fr}{I.1.1.9}
\EGAReferenceAlias{I.1.1.9-es}{I.1.1.9}
\EGAReferenceAlias{I.1.1.9.fr}{I.1.1.9}
\EGAReferenceAlias{I.1.1.9.es}{I.1.1.9}
\EGAReferenceAlias{I.1.1.9.eq}{I.1.1.9}
\EGAReferenceAlias{eq:I.1.1.9}{I.1.1.9}
\EGAReferenceAlias{I.1.1.9.1}{I.1.1.9.1-fr}
\EGAReferenceAlias{I.1.1.9.1-es}{I.1.1.9.1-fr}
\EGAReferenceAlias{I.1.1.9.1.fr}{I.1.1.9.1-fr}
\EGAReferenceAlias{I.1.1.9.1.es}{I.1.1.9.1-fr}
\EGAReferenceAlias{I.1.1.9.1.eq}{I.1.1.9.1-fr}
\EGAReferenceAlias{eq:I.1.1.9.1}{I.1.1.9.1-fr}
\EGAReferenceAlias{I.1.1.10-fr}{I.1.1.10}
\EGAReferenceAlias{I.1.1.10-es}{I.1.1.10}
\EGAReferenceAlias{I.1.1.10.fr}{I.1.1.10}
\EGAReferenceAlias{I.1.1.10.es}{I.1.1.10}
\EGAReferenceAlias{I.1.1.10.eq}{I.1.1.10}
\EGAReferenceAlias{eq:I.1.1.10}{I.1.1.10}
\EGAReferenceAlias{I.1.1.11-fr}{I.1.1.11}
\EGAReferenceAlias{I.1.1.11-es}{I.1.1.11}
\EGAReferenceAlias{I.1.1.11.fr}{I.1.1.11}
\EGAReferenceAlias{I.1.1.11.es}{I.1.1.11}
\EGAReferenceAlias{I.1.1.11.eq}{I.1.1.11}
\EGAReferenceAlias{eq:I.1.1.11}{I.1.1.11}
\EGAReferenceAlias{I.1.1.12-fr}{I.1.1.12}
\EGAReferenceAlias{I.1.1.12-es}{I.1.1.12}
\EGAReferenceAlias{I.1.1.12.fr}{I.1.1.12}
\EGAReferenceAlias{I.1.1.12.es}{I.1.1.12}
\EGAReferenceAlias{I.1.1.12.eq}{I.1.1.12}
\EGAReferenceAlias{eq:I.1.1.12}{I.1.1.12}
\EGAReferenceAlias{I.1.1.13-fr}{I.1.1.13}
\EGAReferenceAlias{I.1.1.13-es}{I.1.1.13}
\EGAReferenceAlias{I.1.1.13.fr}{I.1.1.13}
\EGAReferenceAlias{I.1.1.13.es}{I.1.1.13}
\EGAReferenceAlias{I.1.1.13.eq}{I.1.1.13}
\EGAReferenceAlias{eq:I.1.1.13}{I.1.1.13}
\EGAReferenceAlias{I.1.1.14-fr}{I.1.1.14}
\EGAReferenceAlias{I.1.1.14-es}{I.1.1.14}
\EGAReferenceAlias{I.1.1.14.fr}{I.1.1.14}
\EGAReferenceAlias{I.1.1.14.es}{I.1.1.14}
\EGAReferenceAlias{I.1.1.14.eq}{I.1.1.14}
\EGAReferenceAlias{eq:I.1.1.14}{I.1.1.14}
\EGAReferenceAlias{I.1.1.15-fr}{I.1.1.15}
\EGAReferenceAlias{I.1.1.15-es}{I.1.1.15}
\EGAReferenceAlias{I.1.1.15.fr}{I.1.1.15}
\EGAReferenceAlias{I.1.1.15.es}{I.1.1.15}
\EGAReferenceAlias{I.1.1.15.eq}{I.1.1.15}
\EGAReferenceAlias{eq:I.1.1.15}{I.1.1.15}
\EGAReferenceAlias{I.1.2}{subsection:I.1.2}
\EGAReferenceAlias{I.1.2-fr}{subsection:I.1.2}
\EGAReferenceAlias{I.1.2-es}{subsection:I.1.2}
\EGAReferenceAlias{I.1.2.fr}{subsection:I.1.2}
\EGAReferenceAlias{I.1.2.es}{subsection:I.1.2}
\EGAReferenceAlias{I.1.2.eq}{subsection:I.1.2}
\EGAReferenceAlias{eq:I.1.2}{subsection:I.1.2}
\EGAReferenceAlias{subsection:I.1.2-fr}{subsection:I.1.2}
\EGAReferenceAlias{subsection:I.1.2-es}{subsection:I.1.2}
\EGAReferenceAlias{subsection:I.1.2.fr}{subsection:I.1.2}
\EGAReferenceAlias{subsection:I.1.2.es}{subsection:I.1.2}
\EGAReferenceAlias{subsection:I.1.2.eq}{subsection:I.1.2}
\EGAReferenceAlias{eq:subsection:I.1.2}{subsection:I.1.2}
\EGAReferenceAlias{I.1.2.1-fr}{I.1.2.1}
\EGAReferenceAlias{I.1.2.1-es}{I.1.2.1}
\EGAReferenceAlias{I.1.2.1.fr}{I.1.2.1}
\EGAReferenceAlias{I.1.2.1.es}{I.1.2.1}
\EGAReferenceAlias{I.1.2.1.eq}{I.1.2.1}
\EGAReferenceAlias{eq:I.1.2.1}{I.1.2.1}
\EGAReferenceAlias{I.1.2.1.1}{I.1.2.1.1-fr}
\EGAReferenceAlias{I.1.2.1.1-es}{I.1.2.1.1-fr}
\EGAReferenceAlias{I.1.2.1.1.fr}{I.1.2.1.1-fr}
\EGAReferenceAlias{I.1.2.1.1.es}{I.1.2.1.1-fr}
\EGAReferenceAlias{I.1.2.1.1.eq}{I.1.2.1.1-fr}
\EGAReferenceAlias{eq:I.1.2.1.1}{I.1.2.1.1-fr}
\EGAReferenceAlias{I.1.2.2-fr}{I.1.2.2}
\EGAReferenceAlias{I.1.2.2-es}{I.1.2.2}
\EGAReferenceAlias{I.1.2.2.fr}{I.1.2.2}
\EGAReferenceAlias{I.1.2.2.es}{I.1.2.2}
\EGAReferenceAlias{I.1.2.2.eq}{I.1.2.2}
\EGAReferenceAlias{eq:I.1.2.2}{I.1.2.2}
\EGAReferenceAlias{I.1.2.2.1}{I.1.2.2.1-fr}
\EGAReferenceAlias{I.1.2.2.1-es}{I.1.2.2.1-fr}
\EGAReferenceAlias{I.1.2.2.1.fr}{I.1.2.2.1-fr}
\EGAReferenceAlias{I.1.2.2.1.es}{I.1.2.2.1-fr}
\EGAReferenceAlias{I.1.2.2.1.eq}{I.1.2.2.1-fr}
\EGAReferenceAlias{eq:I.1.2.2.1}{I.1.2.2.1-fr}
\EGAReferenceAlias{I.1.2.2.2}{I.1.2.2.2-fr}
\EGAReferenceAlias{I.1.2.2.2-es}{I.1.2.2.2-fr}
\EGAReferenceAlias{I.1.2.2.2.fr}{I.1.2.2.2-fr}
\EGAReferenceAlias{I.1.2.2.2.es}{I.1.2.2.2-fr}
\EGAReferenceAlias{I.1.2.2.2.eq}{I.1.2.2.2-fr}
\EGAReferenceAlias{eq:I.1.2.2.2}{I.1.2.2.2-fr}
\EGAReferenceAlias{I.1.2.2.3}{I.1.2.2.3-fr}
\EGAReferenceAlias{I.1.2.2.3-es}{I.1.2.2.3-fr}
\EGAReferenceAlias{I.1.2.2.3.fr}{I.1.2.2.3-fr}
\EGAReferenceAlias{I.1.2.2.3.es}{I.1.2.2.3-fr}
\EGAReferenceAlias{I.1.2.2.3.eq}{I.1.2.2.3-fr}
\EGAReferenceAlias{eq:I.1.2.2.3}{I.1.2.2.3-fr}
\EGAReferenceAlias{I.1.2.3-fr}{I.1.2.3}
\EGAReferenceAlias{I.1.2.3-es}{I.1.2.3}
\EGAReferenceAlias{I.1.2.3.fr}{I.1.2.3}
\EGAReferenceAlias{I.1.2.3.es}{I.1.2.3}
\EGAReferenceAlias{I.1.2.3.eq}{I.1.2.3}
\EGAReferenceAlias{eq:I.1.2.3}{I.1.2.3}
\EGAReferenceAlias{I.1.2.4-fr}{I.1.2.4}
\EGAReferenceAlias{I.1.2.4-es}{I.1.2.4}
\EGAReferenceAlias{I.1.2.4.fr}{I.1.2.4}
\EGAReferenceAlias{I.1.2.4.es}{I.1.2.4}
\EGAReferenceAlias{I.1.2.4.eq}{I.1.2.4}
\EGAReferenceAlias{eq:I.1.2.4}{I.1.2.4}
\EGAReferenceAlias{I.1.2.5-fr}{I.1.2.5}
\EGAReferenceAlias{I.1.2.5-es}{I.1.2.5}
\EGAReferenceAlias{I.1.2.5.fr}{I.1.2.5}
\EGAReferenceAlias{I.1.2.5.es}{I.1.2.5}
\EGAReferenceAlias{I.1.2.5.eq}{I.1.2.5}
\EGAReferenceAlias{eq:I.1.2.5}{I.1.2.5}
\EGAReferenceAlias{I.1.2.6-fr}{I.1.2.6}
\EGAReferenceAlias{I.1.2.6-es}{I.1.2.6}
\EGAReferenceAlias{I.1.2.6.fr}{I.1.2.6}
\EGAReferenceAlias{I.1.2.6.es}{I.1.2.6}
\EGAReferenceAlias{I.1.2.6.eq}{I.1.2.6}
\EGAReferenceAlias{eq:I.1.2.6}{I.1.2.6}
\EGAReferenceAlias{I.1.2.7-fr}{I.1.2.7}
\EGAReferenceAlias{I.1.2.7-es}{I.1.2.7}
\EGAReferenceAlias{I.1.2.7.fr}{I.1.2.7}
\EGAReferenceAlias{I.1.2.7.es}{I.1.2.7}
\EGAReferenceAlias{I.1.2.7.eq}{I.1.2.7}
\EGAReferenceAlias{eq:I.1.2.7}{I.1.2.7}
\EGAReferenceAlias{I.1.3-fr}{subsection:I.1.3}
\EGAReferenceAlias{I.1.3-es}{subsection:I.1.3}
\EGAReferenceAlias{I.1.3.fr}{subsection:I.1.3}
\EGAReferenceAlias{I.1.3.es}{subsection:I.1.3}
\EGAReferenceAlias{I.1.3.eq}{subsection:I.1.3}
\EGAReferenceAlias{eq:I.1.3}{subsection:I.1.3}
\EGAReferenceAlias{subsection:I.1.3-fr}{subsection:I.1.3}
\EGAReferenceAlias{subsection:I.1.3-es}{subsection:I.1.3}
\EGAReferenceAlias{subsection:I.1.3.fr}{subsection:I.1.3}
\EGAReferenceAlias{subsection:I.1.3.es}{subsection:I.1.3}
\EGAReferenceAlias{subsection:I.1.3.eq}{subsection:I.1.3}
\EGAReferenceAlias{eq:subsection:I.1.3}{subsection:I.1.3}
\EGAReferenceAlias{I.1.3.1-fr}{I.1.3.1}
\EGAReferenceAlias{I.1.3.1-es}{I.1.3.1}
\EGAReferenceAlias{I.1.3.1.fr}{I.1.3.1}
\EGAReferenceAlias{I.1.3.1.es}{I.1.3.1}
\EGAReferenceAlias{I.1.3.1.eq}{I.1.3.1}
\EGAReferenceAlias{eq:I.1.3.1}{I.1.3.1}
\EGAReferenceAlias{I.1.3.2-fr}{I.1.3.2}
\EGAReferenceAlias{I.1.3.2-es}{I.1.3.2}
\EGAReferenceAlias{I.1.3.2.fr}{I.1.3.2}
\EGAReferenceAlias{I.1.3.2.es}{I.1.3.2}
\EGAReferenceAlias{I.1.3.2.eq}{I.1.3.2}
\EGAReferenceAlias{eq:I.1.3.2}{I.1.3.2}
\EGAReferenceAlias{I.1.3.3-fr}{I.1.3.3}
\EGAReferenceAlias{I.1.3.3-es}{I.1.3.3}
\EGAReferenceAlias{I.1.3.3.fr}{I.1.3.3}
\EGAReferenceAlias{I.1.3.3.es}{I.1.3.3}
\EGAReferenceAlias{I.1.3.3.eq}{I.1.3.3}
\EGAReferenceAlias{eq:I.1.3.3}{I.1.3.3}
\EGAReferenceAlias{I.1.3.3.1}{I.1.3.3.1-fr}
\EGAReferenceAlias{I.1.3.3.1-es}{I.1.3.3.1-fr}
\EGAReferenceAlias{I.1.3.3.1.fr}{I.1.3.3.1-fr}
\EGAReferenceAlias{I.1.3.3.1.es}{I.1.3.3.1-fr}
\EGAReferenceAlias{I.1.3.3.1.eq}{I.1.3.3.1-fr}
\EGAReferenceAlias{eq:I.1.3.3.1}{I.1.3.3.1-fr}
\EGAReferenceAlias{I.1.3.4-fr}{I.1.3.4}
\EGAReferenceAlias{I.1.3.4-es}{I.1.3.4}
\EGAReferenceAlias{I.1.3.4.fr}{I.1.3.4}
\EGAReferenceAlias{I.1.3.4.es}{I.1.3.4}
\EGAReferenceAlias{I.1.3.4.eq}{I.1.3.4}
\EGAReferenceAlias{eq:I.1.3.4}{I.1.3.4}
\EGAReferenceAlias{I.1.3.4.1}{I.1.3.4.1-fr}
\EGAReferenceAlias{I.1.3.4.1-es}{I.1.3.4.1-fr}
\EGAReferenceAlias{I.1.3.4.1.fr}{I.1.3.4.1-fr}
\EGAReferenceAlias{I.1.3.4.1.es}{I.1.3.4.1-fr}
\EGAReferenceAlias{I.1.3.4.1.eq}{I.1.3.4.1-fr}
\EGAReferenceAlias{eq:I.1.3.4.1}{I.1.3.4.1-fr}
\EGAReferenceAlias{I.1.3.5-fr}{I.1.3.5}
\EGAReferenceAlias{I.1.3.5-es}{I.1.3.5}
\EGAReferenceAlias{I.1.3.5.fr}{I.1.3.5}
\EGAReferenceAlias{I.1.3.5.es}{I.1.3.5}
\EGAReferenceAlias{I.1.3.5.eq}{I.1.3.5}
\EGAReferenceAlias{eq:I.1.3.5}{I.1.3.5}
\EGAReferenceAlias{I.1.3.6-fr}{I.1.3.6}
\EGAReferenceAlias{I.1.3.6-es}{I.1.3.6}
\EGAReferenceAlias{I.1.3.6.fr}{I.1.3.6}
\EGAReferenceAlias{I.1.3.6.es}{I.1.3.6}
\EGAReferenceAlias{I.1.3.6.eq}{I.1.3.6}
\EGAReferenceAlias{eq:I.1.3.6}{I.1.3.6}
\EGAReferenceAlias{I.1.3.7-fr}{I.1.3.7}
\EGAReferenceAlias{I.1.3.7-es}{I.1.3.7}
\EGAReferenceAlias{I.1.3.7.fr}{I.1.3.7}
\EGAReferenceAlias{I.1.3.7.es}{I.1.3.7}
\EGAReferenceAlias{I.1.3.7.eq}{I.1.3.7}
\EGAReferenceAlias{eq:I.1.3.7}{I.1.3.7}
\EGAReferenceAlias{I.1.3.7.1}{I.1.3.7.1-fr}
\EGAReferenceAlias{I.1.3.7.1-es}{I.1.3.7.1-fr}
\EGAReferenceAlias{I.1.3.7.1.fr}{I.1.3.7.1-fr}
\EGAReferenceAlias{I.1.3.7.1.es}{I.1.3.7.1-fr}
\EGAReferenceAlias{I.1.3.7.1.eq}{I.1.3.7.1-fr}
\EGAReferenceAlias{eq:I.1.3.7.1}{I.1.3.7.1-fr}
\EGAReferenceAlias{I.1.3.7.2}{I.1.3.7.2-fr}
\EGAReferenceAlias{I.1.3.7.2-es}{I.1.3.7.2-fr}
\EGAReferenceAlias{I.1.3.7.2.fr}{I.1.3.7.2-fr}
\EGAReferenceAlias{I.1.3.7.2.es}{I.1.3.7.2-fr}
\EGAReferenceAlias{I.1.3.7.2.eq}{I.1.3.7.2-fr}
\EGAReferenceAlias{eq:I.1.3.7.2}{I.1.3.7.2-fr}
\EGAReferenceAlias{I.1.3.8-fr}{I.1.3.8}
\EGAReferenceAlias{I.1.3.8-es}{I.1.3.8}
\EGAReferenceAlias{I.1.3.8.fr}{I.1.3.8}
\EGAReferenceAlias{I.1.3.8.es}{I.1.3.8}
\EGAReferenceAlias{I.1.3.8.eq}{I.1.3.8}
\EGAReferenceAlias{eq:I.1.3.8}{I.1.3.8}
\EGAReferenceAlias{I.1.3.9-fr}{I.1.3.9}
\EGAReferenceAlias{I.1.3.9-es}{I.1.3.9}
\EGAReferenceAlias{I.1.3.9.fr}{I.1.3.9}
\EGAReferenceAlias{I.1.3.9.es}{I.1.3.9}
\EGAReferenceAlias{I.1.3.9.eq}{I.1.3.9}
\EGAReferenceAlias{eq:I.1.3.9}{I.1.3.9}
\EGAReferenceAlias{I.1.3.10-fr}{I.1.3.10}
\EGAReferenceAlias{I.1.3.10-es}{I.1.3.10}
\EGAReferenceAlias{I.1.3.10.fr}{I.1.3.10}
\EGAReferenceAlias{I.1.3.10.es}{I.1.3.10}
\EGAReferenceAlias{I.1.3.10.eq}{I.1.3.10}
\EGAReferenceAlias{eq:I.1.3.10}{I.1.3.10}
\EGAReferenceAlias{I.1.3.10.1}{I.1.3.10.1-fr}
\EGAReferenceAlias{I.1.3.10.1-es}{I.1.3.10.1-fr}
\EGAReferenceAlias{I.1.3.10.1.fr}{I.1.3.10.1-fr}
\EGAReferenceAlias{I.1.3.10.1.es}{I.1.3.10.1-fr}
\EGAReferenceAlias{I.1.3.10.1.eq}{I.1.3.10.1-fr}
\EGAReferenceAlias{eq:I.1.3.10.1}{I.1.3.10.1-fr}
\EGAReferenceAlias{I.1.3.10.2}{I.1.3.10.2-fr}
\EGAReferenceAlias{I.1.3.10.2-es}{I.1.3.10.2-fr}
\EGAReferenceAlias{I.1.3.10.2.fr}{I.1.3.10.2-fr}
\EGAReferenceAlias{I.1.3.10.2.es}{I.1.3.10.2-fr}
\EGAReferenceAlias{I.1.3.10.2.eq}{I.1.3.10.2-fr}
\EGAReferenceAlias{eq:I.1.3.10.2}{I.1.3.10.2-fr}
\EGAReferenceAlias{I.1.3.11-fr}{I.1.3.11}
\EGAReferenceAlias{I.1.3.11-es}{I.1.3.11}
\EGAReferenceAlias{I.1.3.11.fr}{I.1.3.11}
\EGAReferenceAlias{I.1.3.11.es}{I.1.3.11}
\EGAReferenceAlias{I.1.3.11.eq}{I.1.3.11}
\EGAReferenceAlias{eq:I.1.3.11}{I.1.3.11}
\EGAReferenceAlias{I.1.3.12-fr}{I.1.3.12}
\EGAReferenceAlias{I.1.3.12-es}{I.1.3.12}
\EGAReferenceAlias{I.1.3.12.fr}{I.1.3.12}
\EGAReferenceAlias{I.1.3.12.es}{I.1.3.12}
\EGAReferenceAlias{I.1.3.12.eq}{I.1.3.12}
\EGAReferenceAlias{eq:I.1.3.12}{I.1.3.12}
\EGAReferenceAlias{I.1.3.13-fr}{I.1.3.13}
\EGAReferenceAlias{I.1.3.13-es}{I.1.3.13}
\EGAReferenceAlias{I.1.3.13.fr}{I.1.3.13}
\EGAReferenceAlias{I.1.3.13.es}{I.1.3.13}
\EGAReferenceAlias{I.1.3.13.eq}{I.1.3.13}
\EGAReferenceAlias{eq:I.1.3.13}{I.1.3.13}
\EGAReferenceAlias{I.1.3.14-fr}{I.1.3.14}
\EGAReferenceAlias{I.1.3.14-es}{I.1.3.14}
\EGAReferenceAlias{I.1.3.14.fr}{I.1.3.14}
\EGAReferenceAlias{I.1.3.14.es}{I.1.3.14}
\EGAReferenceAlias{I.1.3.14.eq}{I.1.3.14}
\EGAReferenceAlias{eq:I.1.3.14}{I.1.3.14}
\EGAReferenceAlias{I.1.4}{subsection:I.1.4}
\EGAReferenceAlias{I.1.4-fr}{subsection:I.1.4}
\EGAReferenceAlias{I.1.4-es}{subsection:I.1.4}
\EGAReferenceAlias{I.1.4.fr}{subsection:I.1.4}
\EGAReferenceAlias{I.1.4.es}{subsection:I.1.4}
\EGAReferenceAlias{I.1.4.eq}{subsection:I.1.4}
\EGAReferenceAlias{eq:I.1.4}{subsection:I.1.4}
\EGAReferenceAlias{subsection:I.1.4-fr}{subsection:I.1.4}
\EGAReferenceAlias{subsection:I.1.4-es}{subsection:I.1.4}
\EGAReferenceAlias{subsection:I.1.4.fr}{subsection:I.1.4}
\EGAReferenceAlias{subsection:I.1.4.es}{subsection:I.1.4}
\EGAReferenceAlias{subsection:I.1.4.eq}{subsection:I.1.4}
\EGAReferenceAlias{eq:subsection:I.1.4}{subsection:I.1.4}
\EGAReferenceAlias{I.1.4.1-fr}{I.1.4.1}
\EGAReferenceAlias{I.1.4.1-es}{I.1.4.1}
\EGAReferenceAlias{I.1.4.1.fr}{I.1.4.1}
\EGAReferenceAlias{I.1.4.1.es}{I.1.4.1}
\EGAReferenceAlias{I.1.4.1.eq}{I.1.4.1}
\EGAReferenceAlias{eq:I.1.4.1}{I.1.4.1}
\EGAReferenceAlias{I.1.4.1.1-fr}{I.1.4.1.1}
\EGAReferenceAlias{I.1.4.1.1-es}{I.1.4.1.1}
\EGAReferenceAlias{I.1.4.1.1.fr}{I.1.4.1.1}
\EGAReferenceAlias{I.1.4.1.1.es}{I.1.4.1.1}
\EGAReferenceAlias{I.1.4.1.1.eq}{I.1.4.1.1}
\EGAReferenceAlias{eq:I.1.4.1.1}{I.1.4.1.1}
\EGAReferenceAlias{I.1.4.2-fr}{I.1.4.2}
\EGAReferenceAlias{I.1.4.2-es}{I.1.4.2}
\EGAReferenceAlias{I.1.4.2.fr}{I.1.4.2}
\EGAReferenceAlias{I.1.4.2.es}{I.1.4.2}
\EGAReferenceAlias{I.1.4.2.eq}{I.1.4.2}
\EGAReferenceAlias{eq:I.1.4.2}{I.1.4.2}
\EGAReferenceAlias{I.1.4.3-fr}{I.1.4.3}
\EGAReferenceAlias{I.1.4.3-es}{I.1.4.3}
\EGAReferenceAlias{I.1.4.3.fr}{I.1.4.3}
\EGAReferenceAlias{I.1.4.3.es}{I.1.4.3}
\EGAReferenceAlias{I.1.4.3.eq}{I.1.4.3}
\EGAReferenceAlias{eq:I.1.4.3}{I.1.4.3}
\EGAReferenceAlias{I.1.5}{subsection:I.1.5}
\EGAReferenceAlias{I.1.5-fr}{subsection:I.1.5}
\EGAReferenceAlias{I.1.5-es}{subsection:I.1.5}
\EGAReferenceAlias{I.1.5.fr}{subsection:I.1.5}
\EGAReferenceAlias{I.1.5.es}{subsection:I.1.5}
\EGAReferenceAlias{I.1.5.eq}{subsection:I.1.5}
\EGAReferenceAlias{eq:I.1.5}{subsection:I.1.5}
\EGAReferenceAlias{subsection:I.1.5-fr}{subsection:I.1.5}
\EGAReferenceAlias{subsection:I.1.5-es}{subsection:I.1.5}
\EGAReferenceAlias{subsection:I.1.5.fr}{subsection:I.1.5}
\EGAReferenceAlias{subsection:I.1.5.es}{subsection:I.1.5}
\EGAReferenceAlias{subsection:I.1.5.eq}{subsection:I.1.5}
\EGAReferenceAlias{eq:subsection:I.1.5}{subsection:I.1.5}
\EGAReferenceAlias{I.1.5.1-fr}{I.1.5.1}
\EGAReferenceAlias{I.1.5.1-es}{I.1.5.1}
\EGAReferenceAlias{I.1.5.1.fr}{I.1.5.1}
\EGAReferenceAlias{I.1.5.1.es}{I.1.5.1}
\EGAReferenceAlias{I.1.5.1.eq}{I.1.5.1}
\EGAReferenceAlias{eq:I.1.5.1}{I.1.5.1}
\EGAReferenceAlias{I.1.5.2-fr}{I.1.5.2}
\EGAReferenceAlias{I.1.5.2-es}{I.1.5.2}
\EGAReferenceAlias{I.1.5.2.fr}{I.1.5.2}
\EGAReferenceAlias{I.1.5.2.es}{I.1.5.2}
\EGAReferenceAlias{I.1.5.2.eq}{I.1.5.2}
\EGAReferenceAlias{eq:I.1.5.2}{I.1.5.2}
\EGAReferenceAlias{I.1.5.3-fr}{I.1.5.3}
\EGAReferenceAlias{I.1.5.3-es}{I.1.5.3}
\EGAReferenceAlias{I.1.5.3.fr}{I.1.5.3}
\EGAReferenceAlias{I.1.5.3.es}{I.1.5.3}
\EGAReferenceAlias{I.1.5.3.eq}{I.1.5.3}
\EGAReferenceAlias{eq:I.1.5.3}{I.1.5.3}
\EGAReferenceAlias{I.1.5.4-fr}{I.1.5.4}
\EGAReferenceAlias{I.1.5.4-es}{I.1.5.4}
\EGAReferenceAlias{I.1.5.4.fr}{I.1.5.4}
\EGAReferenceAlias{I.1.5.4.es}{I.1.5.4}
\EGAReferenceAlias{I.1.5.4.eq}{I.1.5.4}
\EGAReferenceAlias{eq:I.1.5.4}{I.1.5.4}
\EGAReferenceAlias{I.1.6}{subsection:I.1.6}
\EGAReferenceAlias{I.1.6-fr}{subsection:I.1.6}
\EGAReferenceAlias{I.1.6-es}{subsection:I.1.6}
\EGAReferenceAlias{I.1.6.fr}{subsection:I.1.6}
\EGAReferenceAlias{I.1.6.es}{subsection:I.1.6}
\EGAReferenceAlias{I.1.6.eq}{subsection:I.1.6}
\EGAReferenceAlias{eq:I.1.6}{subsection:I.1.6}
\EGAReferenceAlias{subsection:I.1.6-fr}{subsection:I.1.6}
\EGAReferenceAlias{subsection:I.1.6-es}{subsection:I.1.6}
\EGAReferenceAlias{subsection:I.1.6.fr}{subsection:I.1.6}
\EGAReferenceAlias{subsection:I.1.6.es}{subsection:I.1.6}
\EGAReferenceAlias{subsection:I.1.6.eq}{subsection:I.1.6}
\EGAReferenceAlias{eq:subsection:I.1.6}{subsection:I.1.6}
\EGAReferenceAlias{I.1.6.1-fr}{I.1.6.1}
\EGAReferenceAlias{I.1.6.1-es}{I.1.6.1}
\EGAReferenceAlias{I.1.6.1.fr}{I.1.6.1}
\EGAReferenceAlias{I.1.6.1.es}{I.1.6.1}
\EGAReferenceAlias{I.1.6.1.eq}{I.1.6.1}
\EGAReferenceAlias{eq:I.1.6.1}{I.1.6.1}
\EGAReferenceAlias{I.1.6.2-fr}{I.1.6.2}
\EGAReferenceAlias{I.1.6.2-es}{I.1.6.2}
\EGAReferenceAlias{I.1.6.2.fr}{I.1.6.2}
\EGAReferenceAlias{I.1.6.2.es}{I.1.6.2}
\EGAReferenceAlias{I.1.6.2.eq}{I.1.6.2}
\EGAReferenceAlias{eq:I.1.6.2}{I.1.6.2}
\EGAReferenceAlias{I.1.6.3-fr}{I.1.6.3}
\EGAReferenceAlias{I.1.6.3-es}{I.1.6.3}
\EGAReferenceAlias{I.1.6.3.fr}{I.1.6.3}
\EGAReferenceAlias{I.1.6.3.es}{I.1.6.3}
\EGAReferenceAlias{I.1.6.3.eq}{I.1.6.3}
\EGAReferenceAlias{eq:I.1.6.3}{I.1.6.3}
\EGAReferenceAlias{I.1.6.4-fr}{I.1.6.4}
\EGAReferenceAlias{I.1.6.4-es}{I.1.6.4}
\EGAReferenceAlias{I.1.6.4.fr}{I.1.6.4}
\EGAReferenceAlias{I.1.6.4.es}{I.1.6.4}
\EGAReferenceAlias{I.1.6.4.eq}{I.1.6.4}
\EGAReferenceAlias{eq:I.1.6.4}{I.1.6.4}
\EGAReferenceAlias{I.1.6.5-fr}{I.1.6.5}
\EGAReferenceAlias{I.1.6.5-es}{I.1.6.5}
\EGAReferenceAlias{I.1.6.5.fr}{I.1.6.5}
\EGAReferenceAlias{I.1.6.5.es}{I.1.6.5}
\EGAReferenceAlias{I.1.6.5.eq}{I.1.6.5}
\EGAReferenceAlias{eq:I.1.6.5}{I.1.6.5}
\EGAReferenceAlias{I.1.6.5.1}{I.1.6.5.1-fr}
\EGAReferenceAlias{I.1.6.5.1-es}{I.1.6.5.1-fr}
\EGAReferenceAlias{I.1.6.5.1.fr}{I.1.6.5.1-fr}
\EGAReferenceAlias{I.1.6.5.1.es}{I.1.6.5.1-fr}
\EGAReferenceAlias{I.1.6.5.1.eq}{I.1.6.5.1-fr}
\EGAReferenceAlias{eq:I.1.6.5.1}{I.1.6.5.1-fr}
\EGAReferenceAlias{I.1.6.6-fr}{I.1.6.6}
\EGAReferenceAlias{I.1.6.6-es}{I.1.6.6}
\EGAReferenceAlias{I.1.6.6.fr}{I.1.6.6}
\EGAReferenceAlias{I.1.6.6.es}{I.1.6.6}
\EGAReferenceAlias{I.1.6.6.eq}{I.1.6.6}
\EGAReferenceAlias{eq:I.1.6.6}{I.1.6.6}
\EGAReferenceAlias{I.1.6.7-fr}{I.1.6.7}
\EGAReferenceAlias{I.1.6.7-es}{I.1.6.7}
\EGAReferenceAlias{I.1.6.7.fr}{I.1.6.7}
\EGAReferenceAlias{I.1.6.7.es}{I.1.6.7}
\EGAReferenceAlias{I.1.6.7.eq}{I.1.6.7}
\EGAReferenceAlias{eq:I.1.6.7}{I.1.6.7}
\EGAReferenceAlias{I.1.6.8-fr}{I.1.6.8}
\EGAReferenceAlias{I.1.6.8-es}{I.1.6.8}
\EGAReferenceAlias{I.1.6.8.fr}{I.1.6.8}
\EGAReferenceAlias{I.1.6.8.es}{I.1.6.8}
\EGAReferenceAlias{I.1.6.8.eq}{I.1.6.8}
\EGAReferenceAlias{eq:I.1.6.8}{I.1.6.8}
\EGAReferenceAlias{I.1.6.9-fr}{I.1.6.9}
\EGAReferenceAlias{I.1.6.9-es}{I.1.6.9}
\EGAReferenceAlias{I.1.6.9.fr}{I.1.6.9}
\EGAReferenceAlias{I.1.6.9.es}{I.1.6.9}
\EGAReferenceAlias{I.1.6.9.eq}{I.1.6.9}
\EGAReferenceAlias{eq:I.1.6.9}{I.1.6.9}
\EGAReferenceAlias{I.1.6.10-fr}{I.1.6.10}
\EGAReferenceAlias{I.1.6.10-es}{I.1.6.10}
\EGAReferenceAlias{I.1.6.10.fr}{I.1.6.10}
\EGAReferenceAlias{I.1.6.10.es}{I.1.6.10}
\EGAReferenceAlias{I.1.6.10.eq}{I.1.6.10}
\EGAReferenceAlias{eq:I.1.6.10}{I.1.6.10}
\EGAReferenceAlias{I.1.7}{subsection:I.1.7}
\EGAReferenceAlias{I.1.7-fr}{subsection:I.1.7}
\EGAReferenceAlias{I.1.7-es}{subsection:I.1.7}
\EGAReferenceAlias{I.1.7.fr}{subsection:I.1.7}
\EGAReferenceAlias{I.1.7.es}{subsection:I.1.7}
\EGAReferenceAlias{I.1.7.eq}{subsection:I.1.7}
\EGAReferenceAlias{eq:I.1.7}{subsection:I.1.7}
\EGAReferenceAlias{subsection:I.1.7-fr}{subsection:I.1.7}
\EGAReferenceAlias{subsection:I.1.7-es}{subsection:I.1.7}
\EGAReferenceAlias{subsection:I.1.7.fr}{subsection:I.1.7}
\EGAReferenceAlias{subsection:I.1.7.es}{subsection:I.1.7}
\EGAReferenceAlias{subsection:I.1.7.eq}{subsection:I.1.7}
\EGAReferenceAlias{eq:subsection:I.1.7}{subsection:I.1.7}
\EGAReferenceAlias{I.1.7.1-fr}{I.1.7.1}
\EGAReferenceAlias{I.1.7.1-es}{I.1.7.1}
\EGAReferenceAlias{I.1.7.1.fr}{I.1.7.1}
\EGAReferenceAlias{I.1.7.1.es}{I.1.7.1}
\EGAReferenceAlias{I.1.7.1.eq}{I.1.7.1}
\EGAReferenceAlias{eq:I.1.7.1}{I.1.7.1}
\EGAReferenceAlias{I.1.7.2-fr}{I.1.7.2}
\EGAReferenceAlias{I.1.7.2-es}{I.1.7.2}
\EGAReferenceAlias{I.1.7.2.fr}{I.1.7.2}
\EGAReferenceAlias{I.1.7.2.es}{I.1.7.2}
\EGAReferenceAlias{I.1.7.2.eq}{I.1.7.2}
\EGAReferenceAlias{eq:I.1.7.2}{I.1.7.2}
\EGAReferenceAlias{I.1.7.3-fr}{I.1.7.3}
\EGAReferenceAlias{I.1.7.3-es}{I.1.7.3}
\EGAReferenceAlias{I.1.7.3.fr}{I.1.7.3}
\EGAReferenceAlias{I.1.7.3.es}{I.1.7.3}
\EGAReferenceAlias{I.1.7.3.eq}{I.1.7.3}
\EGAReferenceAlias{eq:I.1.7.3}{I.1.7.3}
\EGAReferenceAlias{I.1.7.3.diagram}{I.1.7.3.diagram-fr}
\EGAReferenceAlias{I.1.7.3.diagram-es}{I.1.7.3.diagram-fr}
\EGAReferenceAlias{I.1.7.3.diagram.fr}{I.1.7.3.diagram-fr}
\EGAReferenceAlias{I.1.7.3.diagram.es}{I.1.7.3.diagram-fr}
\EGAReferenceAlias{I.1.7.3.diagram.eq}{I.1.7.3.diagram-fr}
\EGAReferenceAlias{eq:I.1.7.3.diagram}{I.1.7.3.diagram-fr}
\EGAReferenceAlias{I.1.7.4-fr}{I.1.7.4}
\EGAReferenceAlias{I.1.7.4-es}{I.1.7.4}
\EGAReferenceAlias{I.1.7.4.fr}{I.1.7.4}
\EGAReferenceAlias{I.1.7.4.es}{I.1.7.4}
\EGAReferenceAlias{I.1.7.4.eq}{I.1.7.4}
\EGAReferenceAlias{eq:I.1.7.4}{I.1.7.4}
\EGAReferenceAlias{I.1.7.5-fr}{I.1.7.5}
\EGAReferenceAlias{I.1.7.5-es}{I.1.7.5}
\EGAReferenceAlias{I.1.7.5.fr}{I.1.7.5}
\EGAReferenceAlias{I.1.7.5.es}{I.1.7.5}
\EGAReferenceAlias{I.1.7.5.eq}{I.1.7.5}
\EGAReferenceAlias{eq:I.1.7.5}{I.1.7.5}
\EGAReferenceAlias{I.1.8}{subsection:I.1.8}
\EGAReferenceAlias{I.1.8-fr}{subsection:I.1.8}
\EGAReferenceAlias{I.1.8-es}{subsection:I.1.8}
\EGAReferenceAlias{I.1.8.fr}{subsection:I.1.8}
\EGAReferenceAlias{I.1.8.es}{subsection:I.1.8}
\EGAReferenceAlias{I.1.8.eq}{subsection:I.1.8}
\EGAReferenceAlias{eq:I.1.8}{subsection:I.1.8}
\EGAReferenceAlias{subsection:I.1.8-fr}{subsection:I.1.8}
\EGAReferenceAlias{subsection:I.1.8-es}{subsection:I.1.8}
\EGAReferenceAlias{subsection:I.1.8.fr}{subsection:I.1.8}
\EGAReferenceAlias{subsection:I.1.8.es}{subsection:I.1.8}
\EGAReferenceAlias{subsection:I.1.8.eq}{subsection:I.1.8}
\EGAReferenceAlias{eq:subsection:I.1.8}{subsection:I.1.8}
\EGAReferenceAlias{I.1.8.1-fr}{I.1.8.1}
\EGAReferenceAlias{I.1.8.1-es}{I.1.8.1}
\EGAReferenceAlias{I.1.8.1.fr}{I.1.8.1}
\EGAReferenceAlias{I.1.8.1.es}{I.1.8.1}
\EGAReferenceAlias{I.1.8.1.eq}{I.1.8.1}
\EGAReferenceAlias{eq:I.1.8.1}{I.1.8.1}
\EGAReferenceAlias{I.1.8.1.1-fr}{I.1.8.1.1}
\EGAReferenceAlias{I.1.8.1.1-es}{I.1.8.1.1}
\EGAReferenceAlias{I.1.8.1.1.fr}{I.1.8.1.1}
\EGAReferenceAlias{I.1.8.1.1.es}{I.1.8.1.1}
\EGAReferenceAlias{I.1.8.1.1.eq}{I.1.8.1.1}
\EGAReferenceAlias{eq:I.1.8.1.1}{I.1.8.1.1}
\EGAReferenceAlias{I.1.8.2-fr}{I.1.8.2}
\EGAReferenceAlias{I.1.8.2-es}{I.1.8.2}
\EGAReferenceAlias{I.1.8.2.fr}{I.1.8.2}
\EGAReferenceAlias{I.1.8.2.es}{I.1.8.2}
\EGAReferenceAlias{I.1.8.2.eq}{I.1.8.2}
\EGAReferenceAlias{eq:I.1.8.2}{I.1.8.2}
\EGAReferenceAlias{I.1.8.2.1-fr}{I.1.8.2.1}
\EGAReferenceAlias{I.1.8.2.1-es}{I.1.8.2.1}
\EGAReferenceAlias{I.1.8.2.1.fr}{I.1.8.2.1}
\EGAReferenceAlias{I.1.8.2.1.es}{I.1.8.2.1}
\EGAReferenceAlias{I.1.8.2.1.eq}{I.1.8.2.1}
\EGAReferenceAlias{eq:I.1.8.2.1}{I.1.8.2.1}
\EGAReferenceAlias{I.1.8.2.2-fr}{I.1.8.2.2}
\EGAReferenceAlias{I.1.8.2.2-es}{I.1.8.2.2}
\EGAReferenceAlias{I.1.8.2.2.fr}{I.1.8.2.2}
\EGAReferenceAlias{I.1.8.2.2.es}{I.1.8.2.2}
\EGAReferenceAlias{I.1.8.2.2.eq}{I.1.8.2.2}
\EGAReferenceAlias{eq:I.1.8.2.2}{I.1.8.2.2}
\EGAReferenceAlias{I.1.8.3-fr}{I.1.8.3}
\EGAReferenceAlias{I.1.8.3-es}{I.1.8.3}
\EGAReferenceAlias{I.1.8.3.fr}{I.1.8.3}
\EGAReferenceAlias{I.1.8.3.es}{I.1.8.3}
\EGAReferenceAlias{I.1.8.3.eq}{I.1.8.3}
\EGAReferenceAlias{eq:I.1.8.3}{I.1.8.3}
\EGAReferenceAlias{I.1.8.4-fr}{I.1.8.4}
\EGAReferenceAlias{I.1.8.4-es}{I.1.8.4}
\EGAReferenceAlias{I.1.8.4.fr}{I.1.8.4}
\EGAReferenceAlias{I.1.8.4.es}{I.1.8.4}
\EGAReferenceAlias{I.1.8.4.eq}{I.1.8.4}
\EGAReferenceAlias{eq:I.1.8.4}{I.1.8.4}
\EGAReferenceAlias{I.1.8.5-fr}{I.1.8.5}
\EGAReferenceAlias{I.1.8.5-es}{I.1.8.5}
\EGAReferenceAlias{I.1.8.5.fr}{I.1.8.5}
\EGAReferenceAlias{I.1.8.5.es}{I.1.8.5}
\EGAReferenceAlias{I.1.8.5.eq}{I.1.8.5}
\EGAReferenceAlias{eq:I.1.8.5}{I.1.8.5}
\EGAReferenceAlias{I.1.8.5.1-fr}{I.1.8.5.1}
\EGAReferenceAlias{I.1.8.5.1-es}{I.1.8.5.1}
\EGAReferenceAlias{I.1.8.5.1.fr}{I.1.8.5.1}
\EGAReferenceAlias{I.1.8.5.1.es}{I.1.8.5.1}
\EGAReferenceAlias{I.1.8.5.1.eq}{I.1.8.5.1}
\EGAReferenceAlias{eq:I.1.8.5.1}{I.1.8.5.1}
\EGAReferenceAlias{I.1.8.6-fr}{I.1.8.6}
\EGAReferenceAlias{I.1.8.6-es}{I.1.8.6}
\EGAReferenceAlias{I.1.8.6.fr}{I.1.8.6}
\EGAReferenceAlias{I.1.8.6.es}{I.1.8.6}
\EGAReferenceAlias{I.1.8.6.eq}{I.1.8.6}
\EGAReferenceAlias{eq:I.1.8.6}{I.1.8.6}
\EGAReferenceAlias{I.1.8.6.1-fr}{I.1.8.6.1}
\EGAReferenceAlias{I.1.8.6.1-es}{I.1.8.6.1}
\EGAReferenceAlias{I.1.8.6.1.fr}{I.1.8.6.1}
\EGAReferenceAlias{I.1.8.6.1.es}{I.1.8.6.1}
\EGAReferenceAlias{I.1.8.6.1.eq}{I.1.8.6.1}
\EGAReferenceAlias{eq:I.1.8.6.1}{I.1.8.6.1}
\EGAReferenceAlias{I.1.8.7-fr}{I.1.8.7}
\EGAReferenceAlias{I.1.8.7-es}{I.1.8.7}
\EGAReferenceAlias{I.1.8.7.fr}{I.1.8.7}
\EGAReferenceAlias{I.1.8.7.es}{I.1.8.7}
\EGAReferenceAlias{I.1.8.7.eq}{I.1.8.7}
\EGAReferenceAlias{eq:I.1.8.7}{I.1.8.7}
\EGAReferenceAlias{I.1.8.7.1-fr}{I.1.8.7.1}
\EGAReferenceAlias{I.1.8.7.1-es}{I.1.8.7.1}
\EGAReferenceAlias{I.1.8.7.1.fr}{I.1.8.7.1}
\EGAReferenceAlias{I.1.8.7.1.es}{I.1.8.7.1}
\EGAReferenceAlias{I.1.8.7.1.eq}{I.1.8.7.1}
\EGAReferenceAlias{eq:I.1.8.7.1}{I.1.8.7.1}
\EGAReferenceAlias{I.1.8.8-fr}{I.1.8.8}
\EGAReferenceAlias{I.1.8.8-es}{I.1.8.8}
\EGAReferenceAlias{I.1.8.8.fr}{I.1.8.8}
\EGAReferenceAlias{I.1.8.8.es}{I.1.8.8}
\EGAReferenceAlias{I.1.8.8.eq}{I.1.8.8}
\EGAReferenceAlias{eq:I.1.8.8}{I.1.8.8}
\EGAReferenceAlias{I.1.8.8.1-fr}{I.1.8.8.1}
\EGAReferenceAlias{I.1.8.8.1-es}{I.1.8.8.1}
\EGAReferenceAlias{I.1.8.8.1.fr}{I.1.8.8.1}
\EGAReferenceAlias{I.1.8.8.1.es}{I.1.8.8.1}
\EGAReferenceAlias{I.1.8.8.1.eq}{I.1.8.8.1}
\EGAReferenceAlias{eq:I.1.8.8.1}{I.1.8.8.1}
\EGAReferenceAlias{I.1.8.8.2-fr}{I.1.8.8.2}
\EGAReferenceAlias{I.1.8.8.2-es}{I.1.8.8.2}
\EGAReferenceAlias{I.1.8.8.2.fr}{I.1.8.8.2}
\EGAReferenceAlias{I.1.8.8.2.es}{I.1.8.8.2}
\EGAReferenceAlias{I.1.8.8.2.eq}{I.1.8.8.2}
\EGAReferenceAlias{eq:I.1.8.8.2}{I.1.8.8.2}
\EGAReferenceAlias{I.1.8.9-fr}{I.1.8.9}
\EGAReferenceAlias{I.1.8.9-es}{I.1.8.9}
\EGAReferenceAlias{I.1.8.9.fr}{I.1.8.9}
\EGAReferenceAlias{I.1.8.9.es}{I.1.8.9}
\EGAReferenceAlias{I.1.8.9.eq}{I.1.8.9}
\EGAReferenceAlias{eq:I.1.8.9}{I.1.8.9}
\EGAReferenceAlias{I.1.8.9.1-fr}{I.1.8.9.1}
\EGAReferenceAlias{I.1.8.9.1-es}{I.1.8.9.1}
\EGAReferenceAlias{I.1.8.9.1.fr}{I.1.8.9.1}
\EGAReferenceAlias{I.1.8.9.1.es}{I.1.8.9.1}
\EGAReferenceAlias{I.1.8.9.1.eq}{I.1.8.9.1}
\EGAReferenceAlias{eq:I.1.8.9.1}{I.1.8.9.1}
\EGAReferenceAlias{I.1.8.10-fr}{I.1.8.10}
\EGAReferenceAlias{I.1.8.10-es}{I.1.8.10}
\EGAReferenceAlias{I.1.8.10.fr}{I.1.8.10}
\EGAReferenceAlias{I.1.8.10.es}{I.1.8.10}
\EGAReferenceAlias{I.1.8.10.eq}{I.1.8.10}
\EGAReferenceAlias{eq:I.1.8.10}{I.1.8.10}
\EGAReferenceAlias{I.1.8.11-fr}{I.1.8.11}
\EGAReferenceAlias{I.1.8.11-es}{I.1.8.11}
\EGAReferenceAlias{I.1.8.11.fr}{I.1.8.11}
\EGAReferenceAlias{I.1.8.11.es}{I.1.8.11}
\EGAReferenceAlias{I.1.8.11.eq}{I.1.8.11}
\EGAReferenceAlias{eq:I.1.8.11}{I.1.8.11}
\EGAReferenceAlias{I.2}{section:I.2}
\EGAReferenceAlias{I.2-fr}{section:I.2}
\EGAReferenceAlias{I.2-es}{section:I.2}
\EGAReferenceAlias{I.2.fr}{section:I.2}
\EGAReferenceAlias{I.2.es}{section:I.2}
\EGAReferenceAlias{I.2.eq}{section:I.2}
\EGAReferenceAlias{eq:I.2}{section:I.2}
\EGAReferenceAlias{section:I.2-fr}{section:I.2}
\EGAReferenceAlias{section:I.2-es}{section:I.2}
\EGAReferenceAlias{section:I.2.fr}{section:I.2}
\EGAReferenceAlias{section:I.2.es}{section:I.2}
\EGAReferenceAlias{section:I.2.eq}{section:I.2}
\EGAReferenceAlias{eq:section:I.2}{section:I.2}
\EGAReferenceAlias{I.2.1}{subsection:I.2.1}
\EGAReferenceAlias{I.2.1-fr}{subsection:I.2.1}
\EGAReferenceAlias{I.2.1-es}{subsection:I.2.1}
\EGAReferenceAlias{I.2.1.fr}{subsection:I.2.1}
\EGAReferenceAlias{I.2.1.es}{subsection:I.2.1}
\EGAReferenceAlias{I.2.1.eq}{subsection:I.2.1}
\EGAReferenceAlias{eq:I.2.1}{subsection:I.2.1}
\EGAReferenceAlias{subsection:I.2.1-fr}{subsection:I.2.1}
\EGAReferenceAlias{subsection:I.2.1-es}{subsection:I.2.1}
\EGAReferenceAlias{subsection:I.2.1.fr}{subsection:I.2.1}
\EGAReferenceAlias{subsection:I.2.1.es}{subsection:I.2.1}
\EGAReferenceAlias{subsection:I.2.1.eq}{subsection:I.2.1}
\EGAReferenceAlias{eq:subsection:I.2.1}{subsection:I.2.1}
\EGAReferenceAlias{I.2.1.1-fr}{I.2.1.1}
\EGAReferenceAlias{I.2.1.1-es}{I.2.1.1}
\EGAReferenceAlias{I.2.1.1.fr}{I.2.1.1}
\EGAReferenceAlias{I.2.1.1.es}{I.2.1.1}
\EGAReferenceAlias{I.2.1.1.eq}{I.2.1.1}
\EGAReferenceAlias{eq:I.2.1.1}{I.2.1.1}
\EGAReferenceAlias{I.2.1.2-fr}{I.2.1.2}
\EGAReferenceAlias{I.2.1.2-es}{I.2.1.2}
\EGAReferenceAlias{I.2.1.2.fr}{I.2.1.2}
\EGAReferenceAlias{I.2.1.2.es}{I.2.1.2}
\EGAReferenceAlias{I.2.1.2.eq}{I.2.1.2}
\EGAReferenceAlias{eq:I.2.1.2}{I.2.1.2}
\EGAReferenceAlias{I.2.1.3-fr}{I.2.1.3}
\EGAReferenceAlias{I.2.1.3-es}{I.2.1.3}
\EGAReferenceAlias{I.2.1.3.fr}{I.2.1.3}
\EGAReferenceAlias{I.2.1.3.es}{I.2.1.3}
\EGAReferenceAlias{I.2.1.3.eq}{I.2.1.3}
\EGAReferenceAlias{eq:I.2.1.3}{I.2.1.3}
\EGAReferenceAlias{I.2.1.4-fr}{I.2.1.4}
\EGAReferenceAlias{I.2.1.4-es}{I.2.1.4}
\EGAReferenceAlias{I.2.1.4.fr}{I.2.1.4}
\EGAReferenceAlias{I.2.1.4.es}{I.2.1.4}
\EGAReferenceAlias{I.2.1.4.eq}{I.2.1.4}
\EGAReferenceAlias{eq:I.2.1.4}{I.2.1.4}
\EGAReferenceAlias{I.2.1.5-fr}{I.2.1.5}
\EGAReferenceAlias{I.2.1.5-es}{I.2.1.5}
\EGAReferenceAlias{I.2.1.5.fr}{I.2.1.5}
\EGAReferenceAlias{I.2.1.5.es}{I.2.1.5}
\EGAReferenceAlias{I.2.1.5.eq}{I.2.1.5}
\EGAReferenceAlias{eq:I.2.1.5}{I.2.1.5}
\EGAReferenceAlias{I.2.1.6-fr}{I.2.1.6}
\EGAReferenceAlias{I.2.1.6-es}{I.2.1.6}
\EGAReferenceAlias{I.2.1.6.fr}{I.2.1.6}
\EGAReferenceAlias{I.2.1.6.es}{I.2.1.6}
\EGAReferenceAlias{I.2.1.6.eq}{I.2.1.6}
\EGAReferenceAlias{eq:I.2.1.6}{I.2.1.6}
\EGAReferenceAlias{I.2.1.7-fr}{I.2.1.7}
\EGAReferenceAlias{I.2.1.7-es}{I.2.1.7}
\EGAReferenceAlias{I.2.1.7.fr}{I.2.1.7}
\EGAReferenceAlias{I.2.1.7.es}{I.2.1.7}
\EGAReferenceAlias{I.2.1.7.eq}{I.2.1.7}
\EGAReferenceAlias{eq:I.2.1.7}{I.2.1.7}
\EGAReferenceAlias{I.2.1.8-fr}{I.2.1.8}
\EGAReferenceAlias{I.2.1.8-es}{I.2.1.8}
\EGAReferenceAlias{I.2.1.8.fr}{I.2.1.8}
\EGAReferenceAlias{I.2.1.8.es}{I.2.1.8}
\EGAReferenceAlias{I.2.1.8.eq}{I.2.1.8}
\EGAReferenceAlias{eq:I.2.1.8}{I.2.1.8}
\EGAReferenceAlias{I.2.2-fr}{subsection:I.2.2}
\EGAReferenceAlias{I.2.2-es}{subsection:I.2.2}
\EGAReferenceAlias{I.2.2.fr}{subsection:I.2.2}
\EGAReferenceAlias{I.2.2.es}{subsection:I.2.2}
\EGAReferenceAlias{I.2.2.eq}{subsection:I.2.2}
\EGAReferenceAlias{eq:I.2.2}{subsection:I.2.2}
\EGAReferenceAlias{subsection:I.2.2-fr}{subsection:I.2.2}
\EGAReferenceAlias{subsection:I.2.2-es}{subsection:I.2.2}
\EGAReferenceAlias{subsection:I.2.2.fr}{subsection:I.2.2}
\EGAReferenceAlias{subsection:I.2.2.es}{subsection:I.2.2}
\EGAReferenceAlias{subsection:I.2.2.eq}{subsection:I.2.2}
\EGAReferenceAlias{eq:subsection:I.2.2}{subsection:I.2.2}
\EGAReferenceAlias{I.2.2.1-fr}{I.2.2.1}
\EGAReferenceAlias{I.2.2.1-es}{I.2.2.1}
\EGAReferenceAlias{I.2.2.1.fr}{I.2.2.1}
\EGAReferenceAlias{I.2.2.1.es}{I.2.2.1}
\EGAReferenceAlias{I.2.2.1.eq}{I.2.2.1}
\EGAReferenceAlias{eq:I.2.2.1}{I.2.2.1}
\EGAReferenceAlias{I.2.2.2-fr}{I.2.2.2}
\EGAReferenceAlias{I.2.2.2-es}{I.2.2.2}
\EGAReferenceAlias{I.2.2.2.fr}{I.2.2.2}
\EGAReferenceAlias{I.2.2.2.es}{I.2.2.2}
\EGAReferenceAlias{I.2.2.2.eq}{I.2.2.2}
\EGAReferenceAlias{eq:I.2.2.2}{I.2.2.2}
\EGAReferenceAlias{I.2.2.3-fr}{I.2.2.3}
\EGAReferenceAlias{I.2.2.3-es}{I.2.2.3}
\EGAReferenceAlias{I.2.2.3.fr}{I.2.2.3}
\EGAReferenceAlias{I.2.2.3.es}{I.2.2.3}
\EGAReferenceAlias{I.2.2.3.eq}{I.2.2.3}
\EGAReferenceAlias{eq:I.2.2.3}{I.2.2.3}
\EGAReferenceAlias{I.2.2.4-fr}{I.2.2.4}
\EGAReferenceAlias{I.2.2.4-es}{I.2.2.4}
\EGAReferenceAlias{I.2.2.4.fr}{I.2.2.4}
\EGAReferenceAlias{I.2.2.4.es}{I.2.2.4}
\EGAReferenceAlias{I.2.2.4.eq}{I.2.2.4}
\EGAReferenceAlias{eq:I.2.2.4}{I.2.2.4}
\EGAReferenceAlias{I.2.2.4.1}{I.2.2.4.1-fr}
\EGAReferenceAlias{I.2.2.4.1-es}{I.2.2.4.1-fr}
\EGAReferenceAlias{I.2.2.4.1.fr}{I.2.2.4.1-fr}
\EGAReferenceAlias{I.2.2.4.1.es}{I.2.2.4.1-fr}
\EGAReferenceAlias{I.2.2.4.1.eq}{I.2.2.4.1-fr}
\EGAReferenceAlias{eq:I.2.2.4.1}{I.2.2.4.1-fr}
\EGAReferenceAlias{I.2.2.5-fr}{I.2.2.5}
\EGAReferenceAlias{I.2.2.5-es}{I.2.2.5}
\EGAReferenceAlias{I.2.2.5.fr}{I.2.2.5}
\EGAReferenceAlias{I.2.2.5.es}{I.2.2.5}
\EGAReferenceAlias{I.2.2.5.eq}{I.2.2.5}
\EGAReferenceAlias{eq:I.2.2.5}{I.2.2.5}
\EGAReferenceAlias{I.2.2.6-fr}{I.2.2.6}
\EGAReferenceAlias{I.2.2.6-es}{I.2.2.6}
\EGAReferenceAlias{I.2.2.6.fr}{I.2.2.6}
\EGAReferenceAlias{I.2.2.6.es}{I.2.2.6}
\EGAReferenceAlias{I.2.2.6.eq}{I.2.2.6}
\EGAReferenceAlias{eq:I.2.2.6}{I.2.2.6}
\EGAReferenceAlias{I.2.2.7-fr}{I.2.2.7}
\EGAReferenceAlias{I.2.2.7-es}{I.2.2.7}
\EGAReferenceAlias{I.2.2.7.fr}{I.2.2.7}
\EGAReferenceAlias{I.2.2.7.es}{I.2.2.7}
\EGAReferenceAlias{I.2.2.7.eq}{I.2.2.7}
\EGAReferenceAlias{eq:I.2.2.7}{I.2.2.7}
\EGAReferenceAlias{I.2.2.8-fr}{I.2.2.8}
\EGAReferenceAlias{I.2.2.8-es}{I.2.2.8}
\EGAReferenceAlias{I.2.2.8.fr}{I.2.2.8}
\EGAReferenceAlias{I.2.2.8.es}{I.2.2.8}
\EGAReferenceAlias{I.2.2.8.eq}{I.2.2.8}
\EGAReferenceAlias{eq:I.2.2.8}{I.2.2.8}
\EGAReferenceAlias{I.2.2.9-fr}{I.2.2.9}
\EGAReferenceAlias{I.2.2.9-es}{I.2.2.9}
\EGAReferenceAlias{I.2.2.9.fr}{I.2.2.9}
\EGAReferenceAlias{I.2.2.9.es}{I.2.2.9}
\EGAReferenceAlias{I.2.2.9.eq}{I.2.2.9}
\EGAReferenceAlias{eq:I.2.2.9}{I.2.2.9}
\EGAReferenceAlias{I.2.2.10-fr}{I.2.2.10}
\EGAReferenceAlias{I.2.2.10-es}{I.2.2.10}
\EGAReferenceAlias{I.2.2.10.fr}{I.2.2.10}
\EGAReferenceAlias{I.2.2.10.es}{I.2.2.10}
\EGAReferenceAlias{I.2.2.10.eq}{I.2.2.10}
\EGAReferenceAlias{eq:I.2.2.10}{I.2.2.10}
\EGAReferenceAlias{I.2.3}{subsection:I.2.3}
\EGAReferenceAlias{I.2.3-fr}{subsection:I.2.3}
\EGAReferenceAlias{I.2.3-es}{subsection:I.2.3}
\EGAReferenceAlias{I.2.3.fr}{subsection:I.2.3}
\EGAReferenceAlias{I.2.3.es}{subsection:I.2.3}
\EGAReferenceAlias{I.2.3.eq}{subsection:I.2.3}
\EGAReferenceAlias{eq:I.2.3}{subsection:I.2.3}
\EGAReferenceAlias{subsection:I.2.3-fr}{subsection:I.2.3}
\EGAReferenceAlias{subsection:I.2.3-es}{subsection:I.2.3}
\EGAReferenceAlias{subsection:I.2.3.fr}{subsection:I.2.3}
\EGAReferenceAlias{subsection:I.2.3.es}{subsection:I.2.3}
\EGAReferenceAlias{subsection:I.2.3.eq}{subsection:I.2.3}
\EGAReferenceAlias{eq:subsection:I.2.3}{subsection:I.2.3}
\EGAReferenceAlias{I.2.3.1-fr}{I.2.3.1}
\EGAReferenceAlias{I.2.3.1-es}{I.2.3.1}
\EGAReferenceAlias{I.2.3.1.fr}{I.2.3.1}
\EGAReferenceAlias{I.2.3.1.es}{I.2.3.1}
\EGAReferenceAlias{I.2.3.1.eq}{I.2.3.1}
\EGAReferenceAlias{eq:I.2.3.1}{I.2.3.1}
\EGAReferenceAlias{I.2.3.2-fr}{I.2.3.2}
\EGAReferenceAlias{I.2.3.2-es}{I.2.3.2}
\EGAReferenceAlias{I.2.3.2.fr}{I.2.3.2}
\EGAReferenceAlias{I.2.3.2.es}{I.2.3.2}
\EGAReferenceAlias{I.2.3.2.eq}{I.2.3.2}
\EGAReferenceAlias{eq:I.2.3.2}{I.2.3.2}
\EGAReferenceAlias{I.2.4}{subsection:I.2.4}
\EGAReferenceAlias{I.2.4-fr}{subsection:I.2.4}
\EGAReferenceAlias{I.2.4-es}{subsection:I.2.4}
\EGAReferenceAlias{I.2.4.fr}{subsection:I.2.4}
\EGAReferenceAlias{I.2.4.es}{subsection:I.2.4}
\EGAReferenceAlias{I.2.4.eq}{subsection:I.2.4}
\EGAReferenceAlias{eq:I.2.4}{subsection:I.2.4}
\EGAReferenceAlias{subsection:I.2.4-fr}{subsection:I.2.4}
\EGAReferenceAlias{subsection:I.2.4-es}{subsection:I.2.4}
\EGAReferenceAlias{subsection:I.2.4.fr}{subsection:I.2.4}
\EGAReferenceAlias{subsection:I.2.4.es}{subsection:I.2.4}
\EGAReferenceAlias{subsection:I.2.4.eq}{subsection:I.2.4}
\EGAReferenceAlias{eq:subsection:I.2.4}{subsection:I.2.4}
\EGAReferenceAlias{I.2.4.1-fr}{I.2.4.1}
\EGAReferenceAlias{I.2.4.1-es}{I.2.4.1}
\EGAReferenceAlias{I.2.4.1.fr}{I.2.4.1}
\EGAReferenceAlias{I.2.4.1.es}{I.2.4.1}
\EGAReferenceAlias{I.2.4.1.eq}{I.2.4.1}
\EGAReferenceAlias{eq:I.2.4.1}{I.2.4.1}
\EGAReferenceAlias{I.2.4.1.diagram}{I.2.4.1.diagram-fr}
\EGAReferenceAlias{I.2.4.1.diagram-es}{I.2.4.1.diagram-fr}
\EGAReferenceAlias{I.2.4.1.diagram.fr}{I.2.4.1.diagram-fr}
\EGAReferenceAlias{I.2.4.1.diagram.es}{I.2.4.1.diagram-fr}
\EGAReferenceAlias{I.2.4.1.diagram.eq}{I.2.4.1.diagram-fr}
\EGAReferenceAlias{eq:I.2.4.1.diagram}{I.2.4.1.diagram-fr}
\EGAReferenceAlias{I.2.4.2-fr}{I.2.4.2}
\EGAReferenceAlias{I.2.4.2-es}{I.2.4.2}
\EGAReferenceAlias{I.2.4.2.fr}{I.2.4.2}
\EGAReferenceAlias{I.2.4.2.es}{I.2.4.2}
\EGAReferenceAlias{I.2.4.2.eq}{I.2.4.2}
\EGAReferenceAlias{eq:I.2.4.2}{I.2.4.2}
\EGAReferenceAlias{I.2.4.3-fr}{I.2.4.3}
\EGAReferenceAlias{I.2.4.3-es}{I.2.4.3}
\EGAReferenceAlias{I.2.4.3.fr}{I.2.4.3}
\EGAReferenceAlias{I.2.4.3.es}{I.2.4.3}
\EGAReferenceAlias{I.2.4.3.eq}{I.2.4.3}
\EGAReferenceAlias{eq:I.2.4.3}{I.2.4.3}
\EGAReferenceAlias{I.2.4.4-fr}{I.2.4.4}
\EGAReferenceAlias{I.2.4.4-es}{I.2.4.4}
\EGAReferenceAlias{I.2.4.4.fr}{I.2.4.4}
\EGAReferenceAlias{I.2.4.4.es}{I.2.4.4}
\EGAReferenceAlias{I.2.4.4.eq}{I.2.4.4}
\EGAReferenceAlias{eq:I.2.4.4}{I.2.4.4}
\EGAReferenceAlias{I.2.4.5-fr}{I.2.4.5}
\EGAReferenceAlias{I.2.4.5-es}{I.2.4.5}
\EGAReferenceAlias{I.2.4.5.fr}{I.2.4.5}
\EGAReferenceAlias{I.2.4.5.es}{I.2.4.5}
\EGAReferenceAlias{I.2.4.5.eq}{I.2.4.5}
\EGAReferenceAlias{eq:I.2.4.5}{I.2.4.5}
\EGAReferenceAlias{I.2.4.6-fr}{I.2.4.6}
\EGAReferenceAlias{I.2.4.6-es}{I.2.4.6}
\EGAReferenceAlias{I.2.4.6.fr}{I.2.4.6}
\EGAReferenceAlias{I.2.4.6.es}{I.2.4.6}
\EGAReferenceAlias{I.2.4.6.eq}{I.2.4.6}
\EGAReferenceAlias{eq:I.2.4.6}{I.2.4.6}
\EGAReferenceAlias{I.2.4.7-fr}{I.2.4.7}
\EGAReferenceAlias{I.2.4.7-es}{I.2.4.7}
\EGAReferenceAlias{I.2.4.7.fr}{I.2.4.7}
\EGAReferenceAlias{I.2.4.7.es}{I.2.4.7}
\EGAReferenceAlias{I.2.4.7.eq}{I.2.4.7}
\EGAReferenceAlias{eq:I.2.4.7}{I.2.4.7}
\EGAReferenceAlias{I.2.4.8-fr}{I.2.4.8}
\EGAReferenceAlias{I.2.4.8-es}{I.2.4.8}
\EGAReferenceAlias{I.2.4.8.fr}{I.2.4.8}
\EGAReferenceAlias{I.2.4.8.es}{I.2.4.8}
\EGAReferenceAlias{I.2.4.8.eq}{I.2.4.8}
\EGAReferenceAlias{eq:I.2.4.8}{I.2.4.8}
\EGAReferenceAlias{I.2.5}{subsection:I.2.5}
\EGAReferenceAlias{I.2.5-fr}{subsection:I.2.5}
\EGAReferenceAlias{I.2.5-es}{subsection:I.2.5}
\EGAReferenceAlias{I.2.5.fr}{subsection:I.2.5}
\EGAReferenceAlias{I.2.5.es}{subsection:I.2.5}
\EGAReferenceAlias{I.2.5.eq}{subsection:I.2.5}
\EGAReferenceAlias{eq:I.2.5}{subsection:I.2.5}
\EGAReferenceAlias{subsection:I.2.5-fr}{subsection:I.2.5}
\EGAReferenceAlias{subsection:I.2.5-es}{subsection:I.2.5}
\EGAReferenceAlias{subsection:I.2.5.fr}{subsection:I.2.5}
\EGAReferenceAlias{subsection:I.2.5.es}{subsection:I.2.5}
\EGAReferenceAlias{subsection:I.2.5.eq}{subsection:I.2.5}
\EGAReferenceAlias{eq:subsection:I.2.5}{subsection:I.2.5}
\EGAReferenceAlias{I.2.5.1-fr}{I.2.5.1}
\EGAReferenceAlias{I.2.5.1-es}{I.2.5.1}
\EGAReferenceAlias{I.2.5.1.fr}{I.2.5.1}
\EGAReferenceAlias{I.2.5.1.es}{I.2.5.1}
\EGAReferenceAlias{I.2.5.1.eq}{I.2.5.1}
\EGAReferenceAlias{eq:I.2.5.1}{I.2.5.1}
\EGAReferenceAlias{I.2.5.2-fr}{I.2.5.2}
\EGAReferenceAlias{I.2.5.2-es}{I.2.5.2}
\EGAReferenceAlias{I.2.5.2.fr}{I.2.5.2}
\EGAReferenceAlias{I.2.5.2.es}{I.2.5.2}
\EGAReferenceAlias{I.2.5.2.eq}{I.2.5.2}
\EGAReferenceAlias{eq:I.2.5.2}{I.2.5.2}
\EGAReferenceAlias{I.2.5.2.diagram}{I.2.5.2.diagram-fr}
\EGAReferenceAlias{I.2.5.2.diagram-es}{I.2.5.2.diagram-fr}
\EGAReferenceAlias{I.2.5.2.diagram.fr}{I.2.5.2.diagram-fr}
\EGAReferenceAlias{I.2.5.2.diagram.es}{I.2.5.2.diagram-fr}
\EGAReferenceAlias{I.2.5.2.diagram.eq}{I.2.5.2.diagram-fr}
\EGAReferenceAlias{eq:I.2.5.2.diagram}{I.2.5.2.diagram-fr}
\EGAReferenceAlias{I.2.5.3-fr}{I.2.5.3}
\EGAReferenceAlias{I.2.5.3-es}{I.2.5.3}
\EGAReferenceAlias{I.2.5.3.fr}{I.2.5.3}
\EGAReferenceAlias{I.2.5.3.es}{I.2.5.3}
\EGAReferenceAlias{I.2.5.3.eq}{I.2.5.3}
\EGAReferenceAlias{eq:I.2.5.3}{I.2.5.3}
\EGAReferenceAlias{I.2.5.4-fr}{I.2.5.4}
\EGAReferenceAlias{I.2.5.4-es}{I.2.5.4}
\EGAReferenceAlias{I.2.5.4.fr}{I.2.5.4}
\EGAReferenceAlias{I.2.5.4.es}{I.2.5.4}
\EGAReferenceAlias{I.2.5.4.eq}{I.2.5.4}
\EGAReferenceAlias{eq:I.2.5.4}{I.2.5.4}
\EGAReferenceAlias{I.2.5.5-fr}{I.2.5.5}
\EGAReferenceAlias{I.2.5.5-es}{I.2.5.5}
\EGAReferenceAlias{I.2.5.5.fr}{I.2.5.5}
\EGAReferenceAlias{I.2.5.5.es}{I.2.5.5}
\EGAReferenceAlias{I.2.5.5.eq}{I.2.5.5}
\EGAReferenceAlias{eq:I.2.5.5}{I.2.5.5}
\EGAReferenceAlias{I.3}{section:I.3}
\EGAReferenceAlias{I.3-fr}{section:I.3}
\EGAReferenceAlias{I.3-es}{section:I.3}
\EGAReferenceAlias{I.3.fr}{section:I.3}
\EGAReferenceAlias{I.3.es}{section:I.3}
\EGAReferenceAlias{I.3.eq}{section:I.3}
\EGAReferenceAlias{eq:I.3}{section:I.3}
\EGAReferenceAlias{section:I.3-fr}{section:I.3}
\EGAReferenceAlias{section:I.3-es}{section:I.3}
\EGAReferenceAlias{section:I.3.fr}{section:I.3}
\EGAReferenceAlias{section:I.3.es}{section:I.3}
\EGAReferenceAlias{section:I.3.eq}{section:I.3}
\EGAReferenceAlias{eq:section:I.3}{section:I.3}
\EGAReferenceAlias{I.3.1-fr}{subsection:I.3.1}
\EGAReferenceAlias{I.3.1-es}{subsection:I.3.1}
\EGAReferenceAlias{I.3.1.fr}{subsection:I.3.1}
\EGAReferenceAlias{I.3.1.es}{subsection:I.3.1}
\EGAReferenceAlias{I.3.1.eq}{subsection:I.3.1}
\EGAReferenceAlias{eq:I.3.1}{subsection:I.3.1}
\EGAReferenceAlias{subsection:I.3.1-fr}{subsection:I.3.1}
\EGAReferenceAlias{subsection:I.3.1-es}{subsection:I.3.1}
\EGAReferenceAlias{subsection:I.3.1.fr}{subsection:I.3.1}
\EGAReferenceAlias{subsection:I.3.1.es}{subsection:I.3.1}
\EGAReferenceAlias{subsection:I.3.1.eq}{subsection:I.3.1}
\EGAReferenceAlias{eq:subsection:I.3.1}{subsection:I.3.1}
\EGAReferenceAlias{I.3.1.1-fr}{I.3.1.1}
\EGAReferenceAlias{I.3.1.1-es}{I.3.1.1}
\EGAReferenceAlias{I.3.1.1.fr}{I.3.1.1}
\EGAReferenceAlias{I.3.1.1.es}{I.3.1.1}
\EGAReferenceAlias{I.3.1.1.eq}{I.3.1.1}
\EGAReferenceAlias{eq:I.3.1.1}{I.3.1.1}
\EGAReferenceAlias{I.3.2}{subsection:I.3.2}
\EGAReferenceAlias{I.3.2-fr}{subsection:I.3.2}
\EGAReferenceAlias{I.3.2-es}{subsection:I.3.2}
\EGAReferenceAlias{I.3.2.fr}{subsection:I.3.2}
\EGAReferenceAlias{I.3.2.es}{subsection:I.3.2}
\EGAReferenceAlias{I.3.2.eq}{subsection:I.3.2}
\EGAReferenceAlias{eq:I.3.2}{subsection:I.3.2}
\EGAReferenceAlias{subsection:I.3.2-fr}{subsection:I.3.2}
\EGAReferenceAlias{subsection:I.3.2-es}{subsection:I.3.2}
\EGAReferenceAlias{subsection:I.3.2.fr}{subsection:I.3.2}
\EGAReferenceAlias{subsection:I.3.2.es}{subsection:I.3.2}
\EGAReferenceAlias{subsection:I.3.2.eq}{subsection:I.3.2}
\EGAReferenceAlias{eq:subsection:I.3.2}{subsection:I.3.2}
\EGAReferenceAlias{I.3.2.1-fr}{I.3.2.1}
\EGAReferenceAlias{I.3.2.1-es}{I.3.2.1}
\EGAReferenceAlias{I.3.2.1.fr}{I.3.2.1}
\EGAReferenceAlias{I.3.2.1.es}{I.3.2.1}
\EGAReferenceAlias{I.3.2.1.eq}{I.3.2.1}
\EGAReferenceAlias{eq:I.3.2.1}{I.3.2.1}
\EGAReferenceAlias{I.3.2.2-fr}{I.3.2.2}
\EGAReferenceAlias{I.3.2.2-es}{I.3.2.2}
\EGAReferenceAlias{I.3.2.2.fr}{I.3.2.2}
\EGAReferenceAlias{I.3.2.2.es}{I.3.2.2}
\EGAReferenceAlias{I.3.2.2.eq}{I.3.2.2}
\EGAReferenceAlias{eq:I.3.2.2}{I.3.2.2}
\EGAReferenceAlias{I.3.2.3-fr}{I.3.2.3}
\EGAReferenceAlias{I.3.2.3-es}{I.3.2.3}
\EGAReferenceAlias{I.3.2.3.fr}{I.3.2.3}
\EGAReferenceAlias{I.3.2.3.es}{I.3.2.3}
\EGAReferenceAlias{I.3.2.3.eq}{I.3.2.3}
\EGAReferenceAlias{eq:I.3.2.3}{I.3.2.3}
\EGAReferenceAlias{I.3.2.4-fr}{I.3.2.4}
\EGAReferenceAlias{I.3.2.4-es}{I.3.2.4}
\EGAReferenceAlias{I.3.2.4.fr}{I.3.2.4}
\EGAReferenceAlias{I.3.2.4.es}{I.3.2.4}
\EGAReferenceAlias{I.3.2.4.eq}{I.3.2.4}
\EGAReferenceAlias{eq:I.3.2.4}{I.3.2.4}
\EGAReferenceAlias{I.3.2.5-fr}{I.3.2.5}
\EGAReferenceAlias{I.3.2.5-es}{I.3.2.5}
\EGAReferenceAlias{I.3.2.5.fr}{I.3.2.5}
\EGAReferenceAlias{I.3.2.5.es}{I.3.2.5}
\EGAReferenceAlias{I.3.2.5.eq}{I.3.2.5}
\EGAReferenceAlias{eq:I.3.2.5}{I.3.2.5}
\EGAReferenceAlias{I.3.2.6-fr}{I.3.2.6}
\EGAReferenceAlias{I.3.2.6-es}{I.3.2.6}
\EGAReferenceAlias{I.3.2.6.fr}{I.3.2.6}
\EGAReferenceAlias{I.3.2.6.es}{I.3.2.6}
\EGAReferenceAlias{I.3.2.6.eq}{I.3.2.6}
\EGAReferenceAlias{eq:I.3.2.6}{I.3.2.6}
\EGAReferenceAlias{I.3.2.6.1-fr}{I.3.2.6.1}
\EGAReferenceAlias{I.3.2.6.1-es}{I.3.2.6.1}
\EGAReferenceAlias{I.3.2.6.1.fr}{I.3.2.6.1}
\EGAReferenceAlias{I.3.2.6.1.es}{I.3.2.6.1}
\EGAReferenceAlias{I.3.2.6.1.eq}{I.3.2.6.1}
\EGAReferenceAlias{eq:I.3.2.6.1}{I.3.2.6.1}
\EGAReferenceAlias{I.3.2.6.2-fr}{I.3.2.6.2}
\EGAReferenceAlias{I.3.2.6.2-es}{I.3.2.6.2}
\EGAReferenceAlias{I.3.2.6.2.fr}{I.3.2.6.2}
\EGAReferenceAlias{I.3.2.6.2.es}{I.3.2.6.2}
\EGAReferenceAlias{I.3.2.6.2.eq}{I.3.2.6.2}
\EGAReferenceAlias{eq:I.3.2.6.2}{I.3.2.6.2}
\EGAReferenceAlias{I.3.2.6.3-fr}{I.3.2.6.3}
\EGAReferenceAlias{I.3.2.6.3-es}{I.3.2.6.3}
\EGAReferenceAlias{I.3.2.6.3.fr}{I.3.2.6.3}
\EGAReferenceAlias{I.3.2.6.3.es}{I.3.2.6.3}
\EGAReferenceAlias{I.3.2.6.3.eq}{I.3.2.6.3}
\EGAReferenceAlias{eq:I.3.2.6.3}{I.3.2.6.3}
\EGAReferenceAlias{I.3.2.6.4-fr}{I.3.2.6.4}
\EGAReferenceAlias{I.3.2.6.4-es}{I.3.2.6.4}
\EGAReferenceAlias{I.3.2.6.4.fr}{I.3.2.6.4}
\EGAReferenceAlias{I.3.2.6.4.es}{I.3.2.6.4}
\EGAReferenceAlias{I.3.2.6.4.eq}{I.3.2.6.4}
\EGAReferenceAlias{eq:I.3.2.6.4}{I.3.2.6.4}
\EGAReferenceAlias{I.3.2.6.5-fr}{I.3.2.6.5}
\EGAReferenceAlias{I.3.2.6.5-es}{I.3.2.6.5}
\EGAReferenceAlias{I.3.2.6.5.fr}{I.3.2.6.5}
\EGAReferenceAlias{I.3.2.6.5.es}{I.3.2.6.5}
\EGAReferenceAlias{I.3.2.6.5.eq}{I.3.2.6.5}
\EGAReferenceAlias{eq:I.3.2.6.5}{I.3.2.6.5}
\EGAReferenceAlias{I.3.2.7-fr}{I.3.2.7}
\EGAReferenceAlias{I.3.2.7-es}{I.3.2.7}
\EGAReferenceAlias{I.3.2.7.fr}{I.3.2.7}
\EGAReferenceAlias{I.3.2.7.es}{I.3.2.7}
\EGAReferenceAlias{I.3.2.7.eq}{I.3.2.7}
\EGAReferenceAlias{eq:I.3.2.7}{I.3.2.7}
\EGAReferenceAlias{I.3.2.8-fr}{I.3.2.8}
\EGAReferenceAlias{I.3.2.8-es}{I.3.2.8}
\EGAReferenceAlias{I.3.2.8.fr}{I.3.2.8}
\EGAReferenceAlias{I.3.2.8.es}{I.3.2.8}
\EGAReferenceAlias{I.3.2.8.eq}{I.3.2.8}
\EGAReferenceAlias{eq:I.3.2.8}{I.3.2.8}
\EGAReferenceAlias{I.3.2.9-fr}{I.3.2.9}
\EGAReferenceAlias{I.3.2.9-es}{I.3.2.9}
\EGAReferenceAlias{I.3.2.9.fr}{I.3.2.9}
\EGAReferenceAlias{I.3.2.9.es}{I.3.2.9}
\EGAReferenceAlias{I.3.2.9.eq}{I.3.2.9}
\EGAReferenceAlias{eq:I.3.2.9}{I.3.2.9}
\EGAReferenceAlias{I.3.3}{subsection:I.3.3}
\EGAReferenceAlias{I.3.3-fr}{subsection:I.3.3}
\EGAReferenceAlias{I.3.3-es}{subsection:I.3.3}
\EGAReferenceAlias{I.3.3.fr}{subsection:I.3.3}
\EGAReferenceAlias{I.3.3.es}{subsection:I.3.3}
\EGAReferenceAlias{I.3.3.eq}{subsection:I.3.3}
\EGAReferenceAlias{eq:I.3.3}{subsection:I.3.3}
\EGAReferenceAlias{subsection:I.3.3-fr}{subsection:I.3.3}
\EGAReferenceAlias{subsection:I.3.3-es}{subsection:I.3.3}
\EGAReferenceAlias{subsection:I.3.3.fr}{subsection:I.3.3}
\EGAReferenceAlias{subsection:I.3.3.es}{subsection:I.3.3}
\EGAReferenceAlias{subsection:I.3.3.eq}{subsection:I.3.3}
\EGAReferenceAlias{eq:subsection:I.3.3}{subsection:I.3.3}
\EGAReferenceAlias{I.3.3.1-fr}{I.3.3.1}
\EGAReferenceAlias{I.3.3.1-es}{I.3.3.1}
\EGAReferenceAlias{I.3.3.1.fr}{I.3.3.1}
\EGAReferenceAlias{I.3.3.1.es}{I.3.3.1}
\EGAReferenceAlias{I.3.3.1.eq}{I.3.3.1}
\EGAReferenceAlias{eq:I.3.3.1}{I.3.3.1}
\EGAReferenceAlias{I.3.3.2-fr}{I.3.3.2}
\EGAReferenceAlias{I.3.3.2-es}{I.3.3.2}
\EGAReferenceAlias{I.3.3.2.fr}{I.3.3.2}
\EGAReferenceAlias{I.3.3.2.es}{I.3.3.2}
\EGAReferenceAlias{I.3.3.2.eq}{I.3.3.2}
\EGAReferenceAlias{eq:I.3.3.2}{I.3.3.2}
\EGAReferenceAlias{I.3.3.3-fr}{I.3.3.3}
\EGAReferenceAlias{I.3.3.3-es}{I.3.3.3}
\EGAReferenceAlias{I.3.3.3.fr}{I.3.3.3}
\EGAReferenceAlias{I.3.3.3.es}{I.3.3.3}
\EGAReferenceAlias{I.3.3.3.eq}{I.3.3.3}
\EGAReferenceAlias{eq:I.3.3.3}{I.3.3.3}
\EGAReferenceAlias{I.3.3.4-fr}{I.3.3.4}
\EGAReferenceAlias{I.3.3.4-es}{I.3.3.4}
\EGAReferenceAlias{I.3.3.4.fr}{I.3.3.4}
\EGAReferenceAlias{I.3.3.4.es}{I.3.3.4}
\EGAReferenceAlias{I.3.3.4.eq}{I.3.3.4}
\EGAReferenceAlias{eq:I.3.3.4}{I.3.3.4}
\EGAReferenceAlias{I.3.3.5-fr}{I.3.3.5}
\EGAReferenceAlias{I.3.3.5-es}{I.3.3.5}
\EGAReferenceAlias{I.3.3.5.fr}{I.3.3.5}
\EGAReferenceAlias{I.3.3.5.es}{I.3.3.5}
\EGAReferenceAlias{I.3.3.5.eq}{I.3.3.5}
\EGAReferenceAlias{eq:I.3.3.5}{I.3.3.5}
\EGAReferenceAlias{I.3.3.6-fr}{I.3.3.6}
\EGAReferenceAlias{I.3.3.6-es}{I.3.3.6}
\EGAReferenceAlias{I.3.3.6.fr}{I.3.3.6}
\EGAReferenceAlias{I.3.3.6.es}{I.3.3.6}
\EGAReferenceAlias{I.3.3.6.eq}{I.3.3.6}
\EGAReferenceAlias{eq:I.3.3.6}{I.3.3.6}
\EGAReferenceAlias{I.3.3.6.base-change-diagram}{I.3.3.6.base-change-diagram-fr}
\EGAReferenceAlias{I.3.3.6.base-change-diagram-es}{I.3.3.6.base-change-diagram-fr}
\EGAReferenceAlias{I.3.3.6.base-change-diagram.fr}{I.3.3.6.base-change-diagram-fr}
\EGAReferenceAlias{I.3.3.6.base-change-diagram.es}{I.3.3.6.base-change-diagram-fr}
\EGAReferenceAlias{I.3.3.6.base-change-diagram.eq}{I.3.3.6.base-change-diagram-fr}
\EGAReferenceAlias{eq:I.3.3.6.base-change-diagram}{I.3.3.6.base-change-diagram-fr}
\EGAReferenceAlias{I.3.3.7-fr}{I.3.3.7}
\EGAReferenceAlias{I.3.3.7-es}{I.3.3.7}
\EGAReferenceAlias{I.3.3.7.fr}{I.3.3.7}
\EGAReferenceAlias{I.3.3.7.es}{I.3.3.7}
\EGAReferenceAlias{I.3.3.7.eq}{I.3.3.7}
\EGAReferenceAlias{eq:I.3.3.7}{I.3.3.7}
\EGAReferenceAlias{I.3.3.8-fr}{I.3.3.8}
\EGAReferenceAlias{I.3.3.8-es}{I.3.3.8}
\EGAReferenceAlias{I.3.3.8.fr}{I.3.3.8}
\EGAReferenceAlias{I.3.3.8.es}{I.3.3.8}
\EGAReferenceAlias{I.3.3.8.eq}{I.3.3.8}
\EGAReferenceAlias{eq:I.3.3.8}{I.3.3.8}
\EGAReferenceAlias{I.3.3.9-fr}{I.3.3.9}
\EGAReferenceAlias{I.3.3.9-es}{I.3.3.9}
\EGAReferenceAlias{I.3.3.9.fr}{I.3.3.9}
\EGAReferenceAlias{I.3.3.9.es}{I.3.3.9}
\EGAReferenceAlias{I.3.3.9.eq}{I.3.3.9}
\EGAReferenceAlias{eq:I.3.3.9}{I.3.3.9}
\EGAReferenceAlias{I.3.3.9.transitivity-diagram}{I.3.3.9.transitivity-diagram-fr}
\EGAReferenceAlias{I.3.3.9.transitivity-diagram-es}{I.3.3.9.transitivity-diagram-fr}
\EGAReferenceAlias{I.3.3.9.transitivity-diagram.fr}{I.3.3.9.transitivity-diagram-fr}
\EGAReferenceAlias{I.3.3.9.transitivity-diagram.es}{I.3.3.9.transitivity-diagram-fr}
\EGAReferenceAlias{I.3.3.9.transitivity-diagram.eq}{I.3.3.9.transitivity-diagram-fr}
\EGAReferenceAlias{eq:I.3.3.9.transitivity-diagram}{I.3.3.9.transitivity-diagram-fr}
\EGAReferenceAlias{I.3.3.9.1}{I.3.3.9.1-fr}
\EGAReferenceAlias{I.3.3.9.1-es}{I.3.3.9.1-fr}
\EGAReferenceAlias{I.3.3.9.1.fr}{I.3.3.9.1-fr}
\EGAReferenceAlias{I.3.3.9.1.es}{I.3.3.9.1-fr}
\EGAReferenceAlias{I.3.3.9.1.eq}{I.3.3.9.1-fr}
\EGAReferenceAlias{eq:I.3.3.9.1}{I.3.3.9.1-fr}
\EGAReferenceAlias{I.3.3.9.2}{I.3.3.9.2-fr}
\EGAReferenceAlias{I.3.3.9.2-es}{I.3.3.9.2-fr}
\EGAReferenceAlias{I.3.3.9.2.fr}{I.3.3.9.2-fr}
\EGAReferenceAlias{I.3.3.9.2.es}{I.3.3.9.2-fr}
\EGAReferenceAlias{I.3.3.9.2.eq}{I.3.3.9.2-fr}
\EGAReferenceAlias{eq:I.3.3.9.2}{I.3.3.9.2-fr}
\EGAReferenceAlias{I.3.3.10-fr}{I.3.3.10}
\EGAReferenceAlias{I.3.3.10-es}{I.3.3.10}
\EGAReferenceAlias{I.3.3.10.fr}{I.3.3.10}
\EGAReferenceAlias{I.3.3.10.es}{I.3.3.10}
\EGAReferenceAlias{I.3.3.10.eq}{I.3.3.10}
\EGAReferenceAlias{eq:I.3.3.10}{I.3.3.10}
\EGAReferenceAlias{I.3.3.10.1}{I.3.3.10.1-fr}
\EGAReferenceAlias{I.3.3.10.1-es}{I.3.3.10.1-fr}
\EGAReferenceAlias{I.3.3.10.1.fr}{I.3.3.10.1-fr}
\EGAReferenceAlias{I.3.3.10.1.es}{I.3.3.10.1-fr}
\EGAReferenceAlias{I.3.3.10.1.eq}{I.3.3.10.1-fr}
\EGAReferenceAlias{eq:I.3.3.10.1}{I.3.3.10.1-fr}
\EGAReferenceAlias{I.3.3.11-fr}{I.3.3.11}
\EGAReferenceAlias{I.3.3.11-es}{I.3.3.11}
\EGAReferenceAlias{I.3.3.11.fr}{I.3.3.11}
\EGAReferenceAlias{I.3.3.11.es}{I.3.3.11}
\EGAReferenceAlias{I.3.3.11.eq}{I.3.3.11}
\EGAReferenceAlias{eq:I.3.3.11}{I.3.3.11}
\EGAReferenceAlias{I.3.3.11.diagram}{I.3.3.11.diagram-fr}
\EGAReferenceAlias{I.3.3.11.diagram-es}{I.3.3.11.diagram-fr}
\EGAReferenceAlias{I.3.3.11.diagram.fr}{I.3.3.11.diagram-fr}
\EGAReferenceAlias{I.3.3.11.diagram.es}{I.3.3.11.diagram-fr}
\EGAReferenceAlias{I.3.3.11.diagram.eq}{I.3.3.11.diagram-fr}
\EGAReferenceAlias{eq:I.3.3.11.diagram}{I.3.3.11.diagram-fr}
\EGAReferenceAlias{I.3.3.12-fr}{I.3.3.12}
\EGAReferenceAlias{I.3.3.12-es}{I.3.3.12}
\EGAReferenceAlias{I.3.3.12.fr}{I.3.3.12}
\EGAReferenceAlias{I.3.3.12.es}{I.3.3.12}
\EGAReferenceAlias{I.3.3.12.eq}{I.3.3.12}
\EGAReferenceAlias{eq:I.3.3.12}{I.3.3.12}
\EGAReferenceAlias{I.3.3.12.base-change-monomorphism}{I.3.3.12.base-change-monomorphism-fr}
\EGAReferenceAlias{I.3.3.12.base-change-monomorphism-es}{I.3.3.12.base-change-monomorphism-fr}
\EGAReferenceAlias{I.3.3.12.base-change-monomorphism.fr}{I.3.3.12.base-change-monomorphism-fr}
\EGAReferenceAlias{I.3.3.12.base-change-monomorphism.es}{I.3.3.12.base-change-monomorphism-fr}
\EGAReferenceAlias{I.3.3.12.base-change-monomorphism.eq}{I.3.3.12.base-change-monomorphism-fr}
\EGAReferenceAlias{eq:I.3.3.12.base-change-monomorphism}{I.3.3.12.base-change-monomorphism-fr}
\EGAReferenceAlias{I.3.3.13-fr}{I.3.3.13}
\EGAReferenceAlias{I.3.3.13-es}{I.3.3.13}
\EGAReferenceAlias{I.3.3.13.fr}{I.3.3.13}
\EGAReferenceAlias{I.3.3.13.es}{I.3.3.13}
\EGAReferenceAlias{I.3.3.13.eq}{I.3.3.13}
\EGAReferenceAlias{eq:I.3.3.13}{I.3.3.13}
\EGAReferenceAlias{I.3.3.14-fr}{I.3.3.14}
\EGAReferenceAlias{I.3.3.14-es}{I.3.3.14}
\EGAReferenceAlias{I.3.3.14.fr}{I.3.3.14}
\EGAReferenceAlias{I.3.3.14.es}{I.3.3.14}
\EGAReferenceAlias{I.3.3.14.eq}{I.3.3.14}
\EGAReferenceAlias{eq:I.3.3.14}{I.3.3.14}
\EGAReferenceAlias{I.3.3.14.correspondence}{I.3.3.14.correspondence-fr}
\EGAReferenceAlias{I.3.3.14.correspondence-es}{I.3.3.14.correspondence-fr}
\EGAReferenceAlias{I.3.3.14.correspondence.fr}{I.3.3.14.correspondence-fr}
\EGAReferenceAlias{I.3.3.14.correspondence.es}{I.3.3.14.correspondence-fr}
\EGAReferenceAlias{I.3.3.14.correspondence.eq}{I.3.3.14.correspondence-fr}
\EGAReferenceAlias{eq:I.3.3.14.correspondence}{I.3.3.14.correspondence-fr}
\EGAReferenceAlias{I.3.3.15-fr}{I.3.3.15}
\EGAReferenceAlias{I.3.3.15-es}{I.3.3.15}
\EGAReferenceAlias{I.3.3.15.fr}{I.3.3.15}
\EGAReferenceAlias{I.3.3.15.es}{I.3.3.15}
\EGAReferenceAlias{I.3.3.15.eq}{I.3.3.15}
\EGAReferenceAlias{eq:I.3.3.15}{I.3.3.15}
\EGAReferenceAlias{I.3.4}{subsection:I.3.4}
\EGAReferenceAlias{I.3.4-fr}{subsection:I.3.4}
\EGAReferenceAlias{I.3.4-es}{subsection:I.3.4}
\EGAReferenceAlias{I.3.4.fr}{subsection:I.3.4}
\EGAReferenceAlias{I.3.4.es}{subsection:I.3.4}
\EGAReferenceAlias{I.3.4.eq}{subsection:I.3.4}
\EGAReferenceAlias{eq:I.3.4}{subsection:I.3.4}
\EGAReferenceAlias{subsection:I.3.4-fr}{subsection:I.3.4}
\EGAReferenceAlias{subsection:I.3.4-es}{subsection:I.3.4}
\EGAReferenceAlias{subsection:I.3.4.fr}{subsection:I.3.4}
\EGAReferenceAlias{subsection:I.3.4.es}{subsection:I.3.4}
\EGAReferenceAlias{subsection:I.3.4.eq}{subsection:I.3.4}
\EGAReferenceAlias{eq:subsection:I.3.4}{subsection:I.3.4}
\EGAReferenceAlias{I.3.4.1-fr}{I.3.4.1}
\EGAReferenceAlias{I.3.4.1-es}{I.3.4.1}
\EGAReferenceAlias{I.3.4.1.fr}{I.3.4.1}
\EGAReferenceAlias{I.3.4.1.es}{I.3.4.1}
\EGAReferenceAlias{I.3.4.1.eq}{I.3.4.1}
\EGAReferenceAlias{eq:I.3.4.1}{I.3.4.1}
\EGAReferenceAlias{I.3.4.2-fr}{I.3.4.2}
\EGAReferenceAlias{I.3.4.2-es}{I.3.4.2}
\EGAReferenceAlias{I.3.4.2.fr}{I.3.4.2}
\EGAReferenceAlias{I.3.4.2.es}{I.3.4.2}
\EGAReferenceAlias{I.3.4.2.eq}{I.3.4.2}
\EGAReferenceAlias{eq:I.3.4.2}{I.3.4.2}
\EGAReferenceAlias{I.3.4.2.1}{I.3.4.2.1-fr}
\EGAReferenceAlias{I.3.4.2.1-es}{I.3.4.2.1-fr}
\EGAReferenceAlias{I.3.4.2.1.fr}{I.3.4.2.1-fr}
\EGAReferenceAlias{I.3.4.2.1.es}{I.3.4.2.1-fr}
\EGAReferenceAlias{I.3.4.2.1.eq}{I.3.4.2.1-fr}
\EGAReferenceAlias{eq:I.3.4.2.1}{I.3.4.2.1-fr}
\EGAReferenceAlias{I.3.4.3-fr}{I.3.4.3}
\EGAReferenceAlias{I.3.4.3-es}{I.3.4.3}
\EGAReferenceAlias{I.3.4.3.fr}{I.3.4.3}
\EGAReferenceAlias{I.3.4.3.es}{I.3.4.3}
\EGAReferenceAlias{I.3.4.3.eq}{I.3.4.3}
\EGAReferenceAlias{eq:I.3.4.3}{I.3.4.3}
\EGAReferenceAlias{I.3.4.3.1}{I.3.4.3.1-fr}
\EGAReferenceAlias{I.3.4.3.1-es}{I.3.4.3.1-fr}
\EGAReferenceAlias{I.3.4.3.1.fr}{I.3.4.3.1-fr}
\EGAReferenceAlias{I.3.4.3.1.es}{I.3.4.3.1-fr}
\EGAReferenceAlias{I.3.4.3.1.eq}{I.3.4.3.1-fr}
\EGAReferenceAlias{eq:I.3.4.3.1}{I.3.4.3.1-fr}
\EGAReferenceAlias{I.3.4.3.2}{I.3.4.3.2-fr}
\EGAReferenceAlias{I.3.4.3.2-es}{I.3.4.3.2-fr}
\EGAReferenceAlias{I.3.4.3.2.fr}{I.3.4.3.2-fr}
\EGAReferenceAlias{I.3.4.3.2.es}{I.3.4.3.2-fr}
\EGAReferenceAlias{I.3.4.3.2.eq}{I.3.4.3.2-fr}
\EGAReferenceAlias{eq:I.3.4.3.2}{I.3.4.3.2-fr}
\EGAReferenceAlias{I.3.4.3.product-diagram}{I.3.4.3.product-diagram-fr}
\EGAReferenceAlias{I.3.4.3.product-diagram-es}{I.3.4.3.product-diagram-fr}
\EGAReferenceAlias{I.3.4.3.product-diagram.fr}{I.3.4.3.product-diagram-fr}
\EGAReferenceAlias{I.3.4.3.product-diagram.es}{I.3.4.3.product-diagram-fr}
\EGAReferenceAlias{I.3.4.3.product-diagram.eq}{I.3.4.3.product-diagram-fr}
\EGAReferenceAlias{eq:I.3.4.3.product-diagram}{I.3.4.3.product-diagram-fr}
\EGAReferenceAlias{I.3.4.4-fr}{I.3.4.4}
\EGAReferenceAlias{I.3.4.4-es}{I.3.4.4}
\EGAReferenceAlias{I.3.4.4.fr}{I.3.4.4}
\EGAReferenceAlias{I.3.4.4.es}{I.3.4.4}
\EGAReferenceAlias{I.3.4.4.eq}{I.3.4.4}
\EGAReferenceAlias{eq:I.3.4.4}{I.3.4.4}
\EGAReferenceAlias{I.3.4.5-fr}{I.3.4.5}
\EGAReferenceAlias{I.3.4.5-es}{I.3.4.5}
\EGAReferenceAlias{I.3.4.5.fr}{I.3.4.5}
\EGAReferenceAlias{I.3.4.5.es}{I.3.4.5}
\EGAReferenceAlias{I.3.4.5.eq}{I.3.4.5}
\EGAReferenceAlias{eq:I.3.4.5}{I.3.4.5}
\EGAReferenceAlias{I.3.4.6-fr}{I.3.4.6}
\EGAReferenceAlias{I.3.4.6-es}{I.3.4.6}
\EGAReferenceAlias{I.3.4.6.fr}{I.3.4.6}
\EGAReferenceAlias{I.3.4.6.es}{I.3.4.6}
\EGAReferenceAlias{I.3.4.6.eq}{I.3.4.6}
\EGAReferenceAlias{eq:I.3.4.6}{I.3.4.6}
\EGAReferenceAlias{I.3.4.7-fr}{I.3.4.7}
\EGAReferenceAlias{I.3.4.7-es}{I.3.4.7}
\EGAReferenceAlias{I.3.4.7.fr}{I.3.4.7}
\EGAReferenceAlias{I.3.4.7.es}{I.3.4.7}
\EGAReferenceAlias{I.3.4.7.eq}{I.3.4.7}
\EGAReferenceAlias{eq:I.3.4.7}{I.3.4.7}
\EGAReferenceAlias{I.3.4.7.underlying-map}{I.3.4.7.underlying-map-fr}
\EGAReferenceAlias{I.3.4.7.underlying-map-es}{I.3.4.7.underlying-map-fr}
\EGAReferenceAlias{I.3.4.7.underlying-map.fr}{I.3.4.7.underlying-map-fr}
\EGAReferenceAlias{I.3.4.7.underlying-map.es}{I.3.4.7.underlying-map-fr}
\EGAReferenceAlias{I.3.4.7.underlying-map.eq}{I.3.4.7.underlying-map-fr}
\EGAReferenceAlias{eq:I.3.4.7.underlying-map}{I.3.4.7.underlying-map-fr}
\EGAReferenceAlias{I.3.4.8-fr}{I.3.4.8}
\EGAReferenceAlias{I.3.4.8-es}{I.3.4.8}
\EGAReferenceAlias{I.3.4.8.fr}{I.3.4.8}
\EGAReferenceAlias{I.3.4.8.es}{I.3.4.8}
\EGAReferenceAlias{I.3.4.8.eq}{I.3.4.8}
\EGAReferenceAlias{eq:I.3.4.8}{I.3.4.8}
\EGAReferenceAlias{I.3.4.8.base-change-image}{I.3.4.8.base-change-image-fr}
\EGAReferenceAlias{I.3.4.8.base-change-image-es}{I.3.4.8.base-change-image-fr}
\EGAReferenceAlias{I.3.4.8.base-change-image.fr}{I.3.4.8.base-change-image-fr}
\EGAReferenceAlias{I.3.4.8.base-change-image.es}{I.3.4.8.base-change-image-fr}
\EGAReferenceAlias{I.3.4.8.base-change-image.eq}{I.3.4.8.base-change-image-fr}
\EGAReferenceAlias{eq:I.3.4.8.base-change-image}{I.3.4.8.base-change-image-fr}
\EGAReferenceAlias{I.3.4.8.proof-diagram}{I.3.4.8.proof-diagram-fr}
\EGAReferenceAlias{I.3.4.8.proof-diagram-es}{I.3.4.8.proof-diagram-fr}
\EGAReferenceAlias{I.3.4.8.proof-diagram.fr}{I.3.4.8.proof-diagram-fr}
\EGAReferenceAlias{I.3.4.8.proof-diagram.es}{I.3.4.8.proof-diagram-fr}
\EGAReferenceAlias{I.3.4.8.proof-diagram.eq}{I.3.4.8.proof-diagram-fr}
\EGAReferenceAlias{eq:I.3.4.8.proof-diagram}{I.3.4.8.proof-diagram-fr}
\EGAReferenceAlias{I.3.4.9-fr}{I.3.4.9}
\EGAReferenceAlias{I.3.4.9-es}{I.3.4.9}
\EGAReferenceAlias{I.3.4.9.fr}{I.3.4.9}
\EGAReferenceAlias{I.3.4.9.es}{I.3.4.9}
\EGAReferenceAlias{I.3.4.9.eq}{I.3.4.9}
\EGAReferenceAlias{eq:I.3.4.9}{I.3.4.9}
\EGAReferenceAlias{I.3.4.9.tensor-product-fibre}{I.3.4.9.tensor-product-fibre-fr}
\EGAReferenceAlias{I.3.4.9.tensor-product-fibre-es}{I.3.4.9.tensor-product-fibre-fr}
\EGAReferenceAlias{I.3.4.9.tensor-product-fibre.fr}{I.3.4.9.tensor-product-fibre-fr}
\EGAReferenceAlias{I.3.4.9.tensor-product-fibre.es}{I.3.4.9.tensor-product-fibre-fr}
\EGAReferenceAlias{I.3.4.9.tensor-product-fibre.eq}{I.3.4.9.tensor-product-fibre-fr}
\EGAReferenceAlias{eq:I.3.4.9.tensor-product-fibre}{I.3.4.9.tensor-product-fibre-fr}
\EGAReferenceAlias{I.3.5}{subsection:I.3.5}
\EGAReferenceAlias{I.3.5-fr}{subsection:I.3.5}
\EGAReferenceAlias{I.3.5-es}{subsection:I.3.5}
\EGAReferenceAlias{I.3.5.fr}{subsection:I.3.5}
\EGAReferenceAlias{I.3.5.es}{subsection:I.3.5}
\EGAReferenceAlias{I.3.5.eq}{subsection:I.3.5}
\EGAReferenceAlias{eq:I.3.5}{subsection:I.3.5}
\EGAReferenceAlias{subsection:I.3.5-fr}{subsection:I.3.5}
\EGAReferenceAlias{subsection:I.3.5-es}{subsection:I.3.5}
\EGAReferenceAlias{subsection:I.3.5.fr}{subsection:I.3.5}
\EGAReferenceAlias{subsection:I.3.5.es}{subsection:I.3.5}
\EGAReferenceAlias{subsection:I.3.5.eq}{subsection:I.3.5}
\EGAReferenceAlias{eq:subsection:I.3.5}{subsection:I.3.5}
\EGAReferenceAlias{I.3.5.1-fr}{I.3.5.1}
\EGAReferenceAlias{I.3.5.1-es}{I.3.5.1}
\EGAReferenceAlias{I.3.5.1.fr}{I.3.5.1}
\EGAReferenceAlias{I.3.5.1.es}{I.3.5.1}
\EGAReferenceAlias{I.3.5.1.eq}{I.3.5.1}
\EGAReferenceAlias{eq:I.3.5.1}{I.3.5.1}
\EGAReferenceAlias{I.3.5.1.product-composition}{I.3.5.1.product-composition-fr}
\EGAReferenceAlias{I.3.5.1.product-composition-es}{I.3.5.1.product-composition-fr}
\EGAReferenceAlias{I.3.5.1.product-composition.fr}{I.3.5.1.product-composition-fr}
\EGAReferenceAlias{I.3.5.1.product-composition.es}{I.3.5.1.product-composition-fr}
\EGAReferenceAlias{I.3.5.1.product-composition.eq}{I.3.5.1.product-composition-fr}
\EGAReferenceAlias{eq:I.3.5.1.product-composition}{I.3.5.1.product-composition-fr}
\EGAReferenceAlias{I.3.5.2-fr}{I.3.5.2}
\EGAReferenceAlias{I.3.5.2-es}{I.3.5.2}
\EGAReferenceAlias{I.3.5.2.fr}{I.3.5.2}
\EGAReferenceAlias{I.3.5.2.es}{I.3.5.2}
\EGAReferenceAlias{I.3.5.2.eq}{I.3.5.2}
\EGAReferenceAlias{eq:I.3.5.2}{I.3.5.2}
\EGAReferenceAlias{I.3.5.3-fr}{I.3.5.3}
\EGAReferenceAlias{I.3.5.3-es}{I.3.5.3}
\EGAReferenceAlias{I.3.5.3.fr}{I.3.5.3}
\EGAReferenceAlias{I.3.5.3.es}{I.3.5.3}
\EGAReferenceAlias{I.3.5.3.eq}{I.3.5.3}
\EGAReferenceAlias{eq:I.3.5.3}{I.3.5.3}
\EGAReferenceAlias{I.3.5.3.surjectivity-diagram}{I.3.5.3.surjectivity-diagram-fr}
\EGAReferenceAlias{I.3.5.3.surjectivity-diagram-es}{I.3.5.3.surjectivity-diagram-fr}
\EGAReferenceAlias{I.3.5.3.surjectivity-diagram.fr}{I.3.5.3.surjectivity-diagram-fr}
\EGAReferenceAlias{I.3.5.3.surjectivity-diagram.es}{I.3.5.3.surjectivity-diagram-fr}
\EGAReferenceAlias{I.3.5.3.surjectivity-diagram.eq}{I.3.5.3.surjectivity-diagram-fr}
\EGAReferenceAlias{eq:I.3.5.3.surjectivity-diagram}{I.3.5.3.surjectivity-diagram-fr}
\EGAReferenceAlias{I.3.5.4-fr}{I.3.5.4}
\EGAReferenceAlias{I.3.5.4-es}{I.3.5.4}
\EGAReferenceAlias{I.3.5.4.fr}{I.3.5.4}
\EGAReferenceAlias{I.3.5.4.es}{I.3.5.4}
\EGAReferenceAlias{I.3.5.4.eq}{I.3.5.4}
\EGAReferenceAlias{eq:I.3.5.4}{I.3.5.4}
\EGAReferenceAlias{I.3.5.5-fr}{I.3.5.5}
\EGAReferenceAlias{I.3.5.5-es}{I.3.5.5}
\EGAReferenceAlias{I.3.5.5.fr}{I.3.5.5}
\EGAReferenceAlias{I.3.5.5.es}{I.3.5.5}
\EGAReferenceAlias{I.3.5.5.eq}{I.3.5.5}
\EGAReferenceAlias{eq:I.3.5.5}{I.3.5.5}
\EGAReferenceAlias{I.3.5.5.algebraic-closure-diagram}{I.3.5.5.algebraic-closure-diagram-fr}
\EGAReferenceAlias{I.3.5.5.algebraic-closure-diagram-es}{I.3.5.5.algebraic-closure-diagram-fr}
\EGAReferenceAlias{I.3.5.5.algebraic-closure-diagram.fr}{I.3.5.5.algebraic-closure-diagram-fr}
\EGAReferenceAlias{I.3.5.5.algebraic-closure-diagram.es}{I.3.5.5.algebraic-closure-diagram-fr}
\EGAReferenceAlias{I.3.5.5.algebraic-closure-diagram.eq}{I.3.5.5.algebraic-closure-diagram-fr}
\EGAReferenceAlias{eq:I.3.5.5.algebraic-closure-diagram}{I.3.5.5.algebraic-closure-diagram-fr}
\EGAReferenceAlias{I.3.5.6-fr}{I.3.5.6}
\EGAReferenceAlias{I.3.5.6-es}{I.3.5.6}
\EGAReferenceAlias{I.3.5.6.fr}{I.3.5.6}
\EGAReferenceAlias{I.3.5.6.es}{I.3.5.6}
\EGAReferenceAlias{I.3.5.6.eq}{I.3.5.6}
\EGAReferenceAlias{eq:I.3.5.6}{I.3.5.6}
\EGAReferenceAlias{I.3.5.7-fr}{I.3.5.7}
\EGAReferenceAlias{I.3.5.7-es}{I.3.5.7}
\EGAReferenceAlias{I.3.5.7.fr}{I.3.5.7}
\EGAReferenceAlias{I.3.5.7.es}{I.3.5.7}
\EGAReferenceAlias{I.3.5.7.eq}{I.3.5.7}
\EGAReferenceAlias{eq:I.3.5.7}{I.3.5.7}
\EGAReferenceAlias{I.3.5.7.field-valued-product}{I.3.5.7.field-valued-product-fr}
\EGAReferenceAlias{I.3.5.7.field-valued-product-es}{I.3.5.7.field-valued-product-fr}
\EGAReferenceAlias{I.3.5.7.field-valued-product.fr}{I.3.5.7.field-valued-product-fr}
\EGAReferenceAlias{I.3.5.7.field-valued-product.es}{I.3.5.7.field-valued-product-fr}
\EGAReferenceAlias{I.3.5.7.field-valued-product.eq}{I.3.5.7.field-valued-product-fr}
\EGAReferenceAlias{eq:I.3.5.7.field-valued-product}{I.3.5.7.field-valued-product-fr}
\EGAReferenceAlias{I.3.5.8-fr}{I.3.5.8}
\EGAReferenceAlias{I.3.5.8-es}{I.3.5.8}
\EGAReferenceAlias{I.3.5.8.fr}{I.3.5.8}
\EGAReferenceAlias{I.3.5.8.es}{I.3.5.8}
\EGAReferenceAlias{I.3.5.8.eq}{I.3.5.8}
\EGAReferenceAlias{eq:I.3.5.8}{I.3.5.8}
\EGAReferenceAlias{I.3.5.9-fr}{I.3.5.9}
\EGAReferenceAlias{I.3.5.9-es}{I.3.5.9}
\EGAReferenceAlias{I.3.5.9.fr}{I.3.5.9}
\EGAReferenceAlias{I.3.5.9.es}{I.3.5.9}
\EGAReferenceAlias{I.3.5.9.eq}{I.3.5.9}
\EGAReferenceAlias{eq:I.3.5.9}{I.3.5.9}
\EGAReferenceAlias{I.3.5.10-fr}{I.3.5.10}
\EGAReferenceAlias{I.3.5.10-es}{I.3.5.10}
\EGAReferenceAlias{I.3.5.10.fr}{I.3.5.10}
\EGAReferenceAlias{I.3.5.10.es}{I.3.5.10}
\EGAReferenceAlias{I.3.5.10.eq}{I.3.5.10}
\EGAReferenceAlias{eq:I.3.5.10}{I.3.5.10}
\EGAReferenceAlias{I.3.5.10.base-change-points-diagram}{I.3.5.10.base-change-points-diagram-fr}
\EGAReferenceAlias{I.3.5.10.base-change-points-diagram-es}{I.3.5.10.base-change-points-diagram-fr}
\EGAReferenceAlias{I.3.5.10.base-change-points-diagram.fr}{I.3.5.10.base-change-points-diagram-fr}
\EGAReferenceAlias{I.3.5.10.base-change-points-diagram.es}{I.3.5.10.base-change-points-diagram-fr}
\EGAReferenceAlias{I.3.5.10.base-change-points-diagram.eq}{I.3.5.10.base-change-points-diagram-fr}
\EGAReferenceAlias{eq:I.3.5.10.base-change-points-diagram}{I.3.5.10.base-change-points-diagram-fr}
\EGAReferenceAlias{I.3.5.11-fr}{I.3.5.11}
\EGAReferenceAlias{I.3.5.11-es}{I.3.5.11}
\EGAReferenceAlias{I.3.5.11.fr}{I.3.5.11}
\EGAReferenceAlias{I.3.5.11.es}{I.3.5.11}
\EGAReferenceAlias{I.3.5.11.eq}{I.3.5.11}
\EGAReferenceAlias{eq:I.3.5.11}{I.3.5.11}
\EGAReferenceAlias{I.3.6}{subsection:I.3.6}
\EGAReferenceAlias{I.3.6-fr}{subsection:I.3.6}
\EGAReferenceAlias{I.3.6-es}{subsection:I.3.6}
\EGAReferenceAlias{I.3.6.fr}{subsection:I.3.6}
\EGAReferenceAlias{I.3.6.es}{subsection:I.3.6}
\EGAReferenceAlias{I.3.6.eq}{subsection:I.3.6}
\EGAReferenceAlias{eq:I.3.6}{subsection:I.3.6}
\EGAReferenceAlias{subsection:I.3.6-fr}{subsection:I.3.6}
\EGAReferenceAlias{subsection:I.3.6-es}{subsection:I.3.6}
\EGAReferenceAlias{subsection:I.3.6.fr}{subsection:I.3.6}
\EGAReferenceAlias{subsection:I.3.6.es}{subsection:I.3.6}
\EGAReferenceAlias{subsection:I.3.6.eq}{subsection:I.3.6}
\EGAReferenceAlias{eq:subsection:I.3.6}{subsection:I.3.6}
\EGAReferenceAlias{I.3.6.1-fr}{I.3.6.1}
\EGAReferenceAlias{I.3.6.1-es}{I.3.6.1}
\EGAReferenceAlias{I.3.6.1.fr}{I.3.6.1}
\EGAReferenceAlias{I.3.6.1.es}{I.3.6.1}
\EGAReferenceAlias{I.3.6.1.eq}{I.3.6.1}
\EGAReferenceAlias{eq:I.3.6.1}{I.3.6.1}
\EGAReferenceAlias{I.3.6.1.localization-fraction}{I.3.6.1.localization-fraction-fr}
\EGAReferenceAlias{I.3.6.1.localization-fraction-es}{I.3.6.1.localization-fraction-fr}
\EGAReferenceAlias{I.3.6.1.localization-fraction.fr}{I.3.6.1.localization-fraction-fr}
\EGAReferenceAlias{I.3.6.1.localization-fraction.es}{I.3.6.1.localization-fraction-fr}
\EGAReferenceAlias{I.3.6.1.localization-fraction.eq}{I.3.6.1.localization-fraction-fr}
\EGAReferenceAlias{eq:I.3.6.1.localization-fraction}{I.3.6.1.localization-fraction-fr}
\EGAReferenceAlias{I.3.6.2-fr}{I.3.6.2}
\EGAReferenceAlias{I.3.6.2-es}{I.3.6.2}
\EGAReferenceAlias{I.3.6.2.fr}{I.3.6.2}
\EGAReferenceAlias{I.3.6.2.es}{I.3.6.2}
\EGAReferenceAlias{I.3.6.2.eq}{I.3.6.2}
\EGAReferenceAlias{eq:I.3.6.2}{I.3.6.2}
\EGAReferenceAlias{I.3.6.3-fr}{I.3.6.3}
\EGAReferenceAlias{I.3.6.3-es}{I.3.6.3}
\EGAReferenceAlias{I.3.6.3.fr}{I.3.6.3}
\EGAReferenceAlias{I.3.6.3.es}{I.3.6.3}
\EGAReferenceAlias{I.3.6.3.eq}{I.3.6.3}
\EGAReferenceAlias{eq:I.3.6.3}{I.3.6.3}
\EGAReferenceAlias{I.3.6.3.fibre-composition}{I.3.6.3.fibre-composition-fr}
\EGAReferenceAlias{I.3.6.3.fibre-composition-es}{I.3.6.3.fibre-composition-fr}
\EGAReferenceAlias{I.3.6.3.fibre-composition.fr}{I.3.6.3.fibre-composition-fr}
\EGAReferenceAlias{I.3.6.3.fibre-composition.es}{I.3.6.3.fibre-composition-fr}
\EGAReferenceAlias{I.3.6.3.fibre-composition.eq}{I.3.6.3.fibre-composition-fr}
\EGAReferenceAlias{eq:I.3.6.3.fibre-composition}{I.3.6.3.fibre-composition-fr}
\EGAReferenceAlias{I.3.6.4-fr}{I.3.6.4}
\EGAReferenceAlias{I.3.6.4-es}{I.3.6.4}
\EGAReferenceAlias{I.3.6.4.fr}{I.3.6.4}
\EGAReferenceAlias{I.3.6.4.es}{I.3.6.4}
\EGAReferenceAlias{I.3.6.4.eq}{I.3.6.4}
\EGAReferenceAlias{eq:I.3.6.4}{I.3.6.4}
\EGAReferenceAlias{I.3.6.5-fr}{I.3.6.5}
\EGAReferenceAlias{I.3.6.5-es}{I.3.6.5}
\EGAReferenceAlias{I.3.6.5.fr}{I.3.6.5}
\EGAReferenceAlias{I.3.6.5.es}{I.3.6.5}
\EGAReferenceAlias{I.3.6.5.eq}{I.3.6.5}
\EGAReferenceAlias{eq:I.3.6.5}{I.3.6.5}
\EGAReferenceAlias{I.3.7-fr}{subsection:I.3.7}
\EGAReferenceAlias{I.3.7-es}{subsection:I.3.7}
\EGAReferenceAlias{I.3.7.fr}{subsection:I.3.7}
\EGAReferenceAlias{I.3.7.es}{subsection:I.3.7}
\EGAReferenceAlias{I.3.7.eq}{subsection:I.3.7}
\EGAReferenceAlias{eq:I.3.7}{subsection:I.3.7}
\EGAReferenceAlias{subsection:I.3.7-fr}{subsection:I.3.7}
\EGAReferenceAlias{subsection:I.3.7-es}{subsection:I.3.7}
\EGAReferenceAlias{subsection:I.3.7.fr}{subsection:I.3.7}
\EGAReferenceAlias{subsection:I.3.7.es}{subsection:I.3.7}
\EGAReferenceAlias{subsection:I.3.7.eq}{subsection:I.3.7}
\EGAReferenceAlias{eq:subsection:I.3.7}{subsection:I.3.7}
\EGAReferenceAlias{I.3.7.1-fr}{I.3.7.1}
\EGAReferenceAlias{I.3.7.1-es}{I.3.7.1}
\EGAReferenceAlias{I.3.7.1.fr}{I.3.7.1}
\EGAReferenceAlias{I.3.7.1.es}{I.3.7.1}
\EGAReferenceAlias{I.3.7.1.eq}{I.3.7.1}
\EGAReferenceAlias{eq:I.3.7.1}{I.3.7.1}
\EGAReferenceAlias{I.3.7.2-fr}{I.3.7.2}
\EGAReferenceAlias{I.3.7.2-es}{I.3.7.2}
\EGAReferenceAlias{I.3.7.2.fr}{I.3.7.2}
\EGAReferenceAlias{I.3.7.2.es}{I.3.7.2}
\EGAReferenceAlias{I.3.7.2.eq}{I.3.7.2}
\EGAReferenceAlias{eq:I.3.7.2}{I.3.7.2}
\EGAReferenceAlias{I.3.7.3-fr}{I.3.7.3}
\EGAReferenceAlias{I.3.7.3-es}{I.3.7.3}
\EGAReferenceAlias{I.3.7.3.fr}{I.3.7.3}
\EGAReferenceAlias{I.3.7.3.es}{I.3.7.3}
\EGAReferenceAlias{I.3.7.3.eq}{I.3.7.3}
\EGAReferenceAlias{eq:I.3.7.3}{I.3.7.3}
\EGAReferenceAlias{I.4}{section:I.4}
\EGAReferenceAlias{I.4-fr}{section:I.4}
\EGAReferenceAlias{I.4-es}{section:I.4}
\EGAReferenceAlias{I.4.fr}{section:I.4}
\EGAReferenceAlias{I.4.es}{section:I.4}
\EGAReferenceAlias{I.4.eq}{section:I.4}
\EGAReferenceAlias{eq:I.4}{section:I.4}
\EGAReferenceAlias{section:I.4-fr}{section:I.4}
\EGAReferenceAlias{section:I.4-es}{section:I.4}
\EGAReferenceAlias{section:I.4.fr}{section:I.4}
\EGAReferenceAlias{section:I.4.es}{section:I.4}
\EGAReferenceAlias{section:I.4.eq}{section:I.4}
\EGAReferenceAlias{eq:section:I.4}{section:I.4}
\EGAReferenceAlias{I.4.1}{subsection:I.4.1}
\EGAReferenceAlias{I.4.1-fr}{subsection:I.4.1}
\EGAReferenceAlias{I.4.1-es}{subsection:I.4.1}
\EGAReferenceAlias{I.4.1.fr}{subsection:I.4.1}
\EGAReferenceAlias{I.4.1.es}{subsection:I.4.1}
\EGAReferenceAlias{I.4.1.eq}{subsection:I.4.1}
\EGAReferenceAlias{eq:I.4.1}{subsection:I.4.1}
\EGAReferenceAlias{subsection:I.4.1-fr}{subsection:I.4.1}
\EGAReferenceAlias{subsection:I.4.1-es}{subsection:I.4.1}
\EGAReferenceAlias{subsection:I.4.1.fr}{subsection:I.4.1}
\EGAReferenceAlias{subsection:I.4.1.es}{subsection:I.4.1}
\EGAReferenceAlias{subsection:I.4.1.eq}{subsection:I.4.1}
\EGAReferenceAlias{eq:subsection:I.4.1}{subsection:I.4.1}
\EGAReferenceAlias{I.4.1.1-fr}{I.4.1.1}
\EGAReferenceAlias{I.4.1.1-es}{I.4.1.1}
\EGAReferenceAlias{I.4.1.1.fr}{I.4.1.1}
\EGAReferenceAlias{I.4.1.1.es}{I.4.1.1}
\EGAReferenceAlias{I.4.1.1.eq}{I.4.1.1}
\EGAReferenceAlias{eq:I.4.1.1}{I.4.1.1}
\EGAReferenceAlias{I.4.1.2-fr}{I.4.1.2}
\EGAReferenceAlias{I.4.1.2-es}{I.4.1.2}
\EGAReferenceAlias{I.4.1.2.fr}{I.4.1.2}
\EGAReferenceAlias{I.4.1.2.es}{I.4.1.2}
\EGAReferenceAlias{I.4.1.2.eq}{I.4.1.2}
\EGAReferenceAlias{eq:I.4.1.2}{I.4.1.2}
\EGAReferenceAlias{I.4.1.3-fr}{I.4.1.3}
\EGAReferenceAlias{I.4.1.3-es}{I.4.1.3}
\EGAReferenceAlias{I.4.1.3.fr}{I.4.1.3}
\EGAReferenceAlias{I.4.1.3.es}{I.4.1.3}
\EGAReferenceAlias{I.4.1.3.eq}{I.4.1.3}
\EGAReferenceAlias{eq:I.4.1.3}{I.4.1.3}
\EGAReferenceAlias{I.4.1.4-fr}{I.4.1.4}
\EGAReferenceAlias{I.4.1.4-es}{I.4.1.4}
\EGAReferenceAlias{I.4.1.4.fr}{I.4.1.4}
\EGAReferenceAlias{I.4.1.4.es}{I.4.1.4}
\EGAReferenceAlias{I.4.1.4.eq}{I.4.1.4}
\EGAReferenceAlias{eq:I.4.1.4}{I.4.1.4}
\EGAReferenceAlias{I.4.1.5-fr}{I.4.1.5}
\EGAReferenceAlias{I.4.1.5-es}{I.4.1.5}
\EGAReferenceAlias{I.4.1.5.fr}{I.4.1.5}
\EGAReferenceAlias{I.4.1.5.es}{I.4.1.5}
\EGAReferenceAlias{I.4.1.5.eq}{I.4.1.5}
\EGAReferenceAlias{eq:I.4.1.5}{I.4.1.5}
\EGAReferenceAlias{I.4.1.6-fr}{I.4.1.6}
\EGAReferenceAlias{I.4.1.6-es}{I.4.1.6}
\EGAReferenceAlias{I.4.1.6.fr}{I.4.1.6}
\EGAReferenceAlias{I.4.1.6.es}{I.4.1.6}
\EGAReferenceAlias{I.4.1.6.eq}{I.4.1.6}
\EGAReferenceAlias{eq:I.4.1.6}{I.4.1.6}
\EGAReferenceAlias{I.4.1.7-fr}{I.4.1.7}
\EGAReferenceAlias{I.4.1.7-es}{I.4.1.7}
\EGAReferenceAlias{I.4.1.7.fr}{I.4.1.7}
\EGAReferenceAlias{I.4.1.7.es}{I.4.1.7}
\EGAReferenceAlias{I.4.1.7.eq}{I.4.1.7}
\EGAReferenceAlias{eq:I.4.1.7}{I.4.1.7}
\EGAReferenceAlias{I.4.1.8-fr}{I.4.1.8}
\EGAReferenceAlias{I.4.1.8-es}{I.4.1.8}
\EGAReferenceAlias{I.4.1.8.fr}{I.4.1.8}
\EGAReferenceAlias{I.4.1.8.es}{I.4.1.8}
\EGAReferenceAlias{I.4.1.8.eq}{I.4.1.8}
\EGAReferenceAlias{eq:I.4.1.8}{I.4.1.8}
\EGAReferenceAlias{I.4.1.9-fr}{I.4.1.9}
\EGAReferenceAlias{I.4.1.9-es}{I.4.1.9}
\EGAReferenceAlias{I.4.1.9.fr}{I.4.1.9}
\EGAReferenceAlias{I.4.1.9.es}{I.4.1.9}
\EGAReferenceAlias{I.4.1.9.eq}{I.4.1.9}
\EGAReferenceAlias{eq:I.4.1.9}{I.4.1.9}
\EGAReferenceAlias{I.4.1.10-fr}{I.4.1.10}
\EGAReferenceAlias{I.4.1.10-es}{I.4.1.10}
\EGAReferenceAlias{I.4.1.10.fr}{I.4.1.10}
\EGAReferenceAlias{I.4.1.10.es}{I.4.1.10}
\EGAReferenceAlias{I.4.1.10.eq}{I.4.1.10}
\EGAReferenceAlias{eq:I.4.1.10}{I.4.1.10}
\EGAReferenceAlias{I.4.2}{subsection:I.4.2}
\EGAReferenceAlias{I.4.2-fr}{subsection:I.4.2}
\EGAReferenceAlias{I.4.2-es}{subsection:I.4.2}
\EGAReferenceAlias{I.4.2.fr}{subsection:I.4.2}
\EGAReferenceAlias{I.4.2.es}{subsection:I.4.2}
\EGAReferenceAlias{I.4.2.eq}{subsection:I.4.2}
\EGAReferenceAlias{eq:I.4.2}{subsection:I.4.2}
\EGAReferenceAlias{subsection:I.4.2-fr}{subsection:I.4.2}
\EGAReferenceAlias{subsection:I.4.2-es}{subsection:I.4.2}
\EGAReferenceAlias{subsection:I.4.2.fr}{subsection:I.4.2}
\EGAReferenceAlias{subsection:I.4.2.es}{subsection:I.4.2}
\EGAReferenceAlias{subsection:I.4.2.eq}{subsection:I.4.2}
\EGAReferenceAlias{eq:subsection:I.4.2}{subsection:I.4.2}
\EGAReferenceAlias{I.4.2.1-fr}{I.4.2.1}
\EGAReferenceAlias{I.4.2.1-es}{I.4.2.1}
\EGAReferenceAlias{I.4.2.1.fr}{I.4.2.1}
\EGAReferenceAlias{I.4.2.1.es}{I.4.2.1}
\EGAReferenceAlias{I.4.2.1.eq}{I.4.2.1}
\EGAReferenceAlias{eq:I.4.2.1}{I.4.2.1}
\EGAReferenceAlias{I.4.2.2-fr}{I.4.2.2}
\EGAReferenceAlias{I.4.2.2-es}{I.4.2.2}
\EGAReferenceAlias{I.4.2.2.fr}{I.4.2.2}
\EGAReferenceAlias{I.4.2.2.es}{I.4.2.2}
\EGAReferenceAlias{I.4.2.2.eq}{I.4.2.2}
\EGAReferenceAlias{eq:I.4.2.2}{I.4.2.2}
\EGAReferenceAlias{I.4.2.2.theta-diagram}{I.4.2.2.theta-diagram-fr}
\EGAReferenceAlias{I.4.2.2.theta-diagram-es}{I.4.2.2.theta-diagram-fr}
\EGAReferenceAlias{I.4.2.2.theta-diagram.fr}{I.4.2.2.theta-diagram-fr}
\EGAReferenceAlias{I.4.2.2.theta-diagram.es}{I.4.2.2.theta-diagram-fr}
\EGAReferenceAlias{I.4.2.2.theta-diagram.eq}{I.4.2.2.theta-diagram-fr}
\EGAReferenceAlias{eq:I.4.2.2.theta-diagram}{I.4.2.2.theta-diagram-fr}
\EGAReferenceAlias{I.4.2.3-fr}{I.4.2.3}
\EGAReferenceAlias{I.4.2.3-es}{I.4.2.3}
\EGAReferenceAlias{I.4.2.3.fr}{I.4.2.3}
\EGAReferenceAlias{I.4.2.3.es}{I.4.2.3}
\EGAReferenceAlias{I.4.2.3.eq}{I.4.2.3}
\EGAReferenceAlias{eq:I.4.2.3}{I.4.2.3}
\EGAReferenceAlias{I.4.2.4-fr}{I.4.2.4}
\EGAReferenceAlias{I.4.2.4-es}{I.4.2.4}
\EGAReferenceAlias{I.4.2.4.fr}{I.4.2.4}
\EGAReferenceAlias{I.4.2.4.es}{I.4.2.4}
\EGAReferenceAlias{I.4.2.4.eq}{I.4.2.4}
\EGAReferenceAlias{eq:I.4.2.4}{I.4.2.4}
\EGAReferenceAlias{I.4.2.5-fr}{I.4.2.5}
\EGAReferenceAlias{I.4.2.5-es}{I.4.2.5}
\EGAReferenceAlias{I.4.2.5.fr}{I.4.2.5}
\EGAReferenceAlias{I.4.2.5.es}{I.4.2.5}
\EGAReferenceAlias{I.4.2.5.eq}{I.4.2.5}
\EGAReferenceAlias{eq:I.4.2.5}{I.4.2.5}
\EGAReferenceAlias{I.4.3-fr}{subsection:I.4.3}
\EGAReferenceAlias{I.4.3-es}{subsection:I.4.3}
\EGAReferenceAlias{I.4.3.fr}{subsection:I.4.3}
\EGAReferenceAlias{I.4.3.es}{subsection:I.4.3}
\EGAReferenceAlias{I.4.3.eq}{subsection:I.4.3}
\EGAReferenceAlias{eq:I.4.3}{subsection:I.4.3}
\EGAReferenceAlias{subsection:I.4.3-fr}{subsection:I.4.3}
\EGAReferenceAlias{subsection:I.4.3-es}{subsection:I.4.3}
\EGAReferenceAlias{subsection:I.4.3.fr}{subsection:I.4.3}
\EGAReferenceAlias{subsection:I.4.3.es}{subsection:I.4.3}
\EGAReferenceAlias{subsection:I.4.3.eq}{subsection:I.4.3}
\EGAReferenceAlias{eq:subsection:I.4.3}{subsection:I.4.3}
\EGAReferenceAlias{I.4.3.1-fr}{I.4.3.1}
\EGAReferenceAlias{I.4.3.1-es}{I.4.3.1}
\EGAReferenceAlias{I.4.3.1.fr}{I.4.3.1}
\EGAReferenceAlias{I.4.3.1.es}{I.4.3.1}
\EGAReferenceAlias{I.4.3.1.eq}{I.4.3.1}
\EGAReferenceAlias{eq:I.4.3.1}{I.4.3.1}
\EGAReferenceAlias{I.4.3.2-fr}{I.4.3.2}
\EGAReferenceAlias{I.4.3.2-es}{I.4.3.2}
\EGAReferenceAlias{I.4.3.2.fr}{I.4.3.2}
\EGAReferenceAlias{I.4.3.2.es}{I.4.3.2}
\EGAReferenceAlias{I.4.3.2.eq}{I.4.3.2}
\EGAReferenceAlias{eq:I.4.3.2}{I.4.3.2}
\EGAReferenceAlias{I.4.4}{subsection:I.4.4}
\EGAReferenceAlias{I.4.4-fr}{subsection:I.4.4}
\EGAReferenceAlias{I.4.4-es}{subsection:I.4.4}
\EGAReferenceAlias{I.4.4.fr}{subsection:I.4.4}
\EGAReferenceAlias{I.4.4.es}{subsection:I.4.4}
\EGAReferenceAlias{I.4.4.eq}{subsection:I.4.4}
\EGAReferenceAlias{eq:I.4.4}{subsection:I.4.4}
\EGAReferenceAlias{subsection:I.4.4-fr}{subsection:I.4.4}
\EGAReferenceAlias{subsection:I.4.4-es}{subsection:I.4.4}
\EGAReferenceAlias{subsection:I.4.4.fr}{subsection:I.4.4}
\EGAReferenceAlias{subsection:I.4.4.es}{subsection:I.4.4}
\EGAReferenceAlias{subsection:I.4.4.eq}{subsection:I.4.4}
\EGAReferenceAlias{eq:subsection:I.4.4}{subsection:I.4.4}
\EGAReferenceAlias{I.4.4.1-fr}{I.4.4.1}
\EGAReferenceAlias{I.4.4.1-es}{I.4.4.1}
\EGAReferenceAlias{I.4.4.1.fr}{I.4.4.1}
\EGAReferenceAlias{I.4.4.1.es}{I.4.4.1}
\EGAReferenceAlias{I.4.4.1.eq}{I.4.4.1}
\EGAReferenceAlias{eq:I.4.4.1}{I.4.4.1}
\EGAReferenceAlias{I.4.4.2-fr}{I.4.4.2}
\EGAReferenceAlias{I.4.4.2-es}{I.4.4.2}
\EGAReferenceAlias{I.4.4.2.fr}{I.4.4.2}
\EGAReferenceAlias{I.4.4.2.es}{I.4.4.2}
\EGAReferenceAlias{I.4.4.2.eq}{I.4.4.2}
\EGAReferenceAlias{eq:I.4.4.2}{I.4.4.2}
\EGAReferenceAlias{I.4.4.3-fr}{I.4.4.3}
\EGAReferenceAlias{I.4.4.3-es}{I.4.4.3}
\EGAReferenceAlias{I.4.4.3.fr}{I.4.4.3}
\EGAReferenceAlias{I.4.4.3.es}{I.4.4.3}
\EGAReferenceAlias{I.4.4.3.eq}{I.4.4.3}
\EGAReferenceAlias{eq:I.4.4.3}{I.4.4.3}
\EGAReferenceAlias{I.4.4.4-fr}{I.4.4.4}
\EGAReferenceAlias{I.4.4.4-es}{I.4.4.4}
\EGAReferenceAlias{I.4.4.4.fr}{I.4.4.4}
\EGAReferenceAlias{I.4.4.4.es}{I.4.4.4}
\EGAReferenceAlias{I.4.4.4.eq}{I.4.4.4}
\EGAReferenceAlias{eq:I.4.4.4}{I.4.4.4}
\EGAReferenceAlias{I.4.4.5-fr}{I.4.4.5}
\EGAReferenceAlias{I.4.4.5-es}{I.4.4.5}
\EGAReferenceAlias{I.4.4.5.fr}{I.4.4.5}
\EGAReferenceAlias{I.4.4.5.es}{I.4.4.5}
\EGAReferenceAlias{I.4.4.5.eq}{I.4.4.5}
\EGAReferenceAlias{eq:I.4.4.5}{I.4.4.5}
\EGAReferenceAlias{I.4.4.6-fr}{I.4.4.6}
\EGAReferenceAlias{I.4.4.6-es}{I.4.4.6}
\EGAReferenceAlias{I.4.4.6.fr}{I.4.4.6}
\EGAReferenceAlias{I.4.4.6.es}{I.4.4.6}
\EGAReferenceAlias{I.4.4.6.eq}{I.4.4.6}
\EGAReferenceAlias{eq:I.4.4.6}{I.4.4.6}
\EGAReferenceAlias{I.4.5}{subsection:I.4.5}
\EGAReferenceAlias{I.4.5-fr}{subsection:I.4.5}
\EGAReferenceAlias{I.4.5-es}{subsection:I.4.5}
\EGAReferenceAlias{I.4.5.fr}{subsection:I.4.5}
\EGAReferenceAlias{I.4.5.es}{subsection:I.4.5}
\EGAReferenceAlias{I.4.5.eq}{subsection:I.4.5}
\EGAReferenceAlias{eq:I.4.5}{subsection:I.4.5}
\EGAReferenceAlias{subsection:I.4.5-fr}{subsection:I.4.5}
\EGAReferenceAlias{subsection:I.4.5-es}{subsection:I.4.5}
\EGAReferenceAlias{subsection:I.4.5.fr}{subsection:I.4.5}
\EGAReferenceAlias{subsection:I.4.5.es}{subsection:I.4.5}
\EGAReferenceAlias{subsection:I.4.5.eq}{subsection:I.4.5}
\EGAReferenceAlias{eq:subsection:I.4.5}{subsection:I.4.5}
\EGAReferenceAlias{I.4.5.1-fr}{I.4.5.1}
\EGAReferenceAlias{I.4.5.1-es}{I.4.5.1}
\EGAReferenceAlias{I.4.5.1.fr}{I.4.5.1}
\EGAReferenceAlias{I.4.5.1.es}{I.4.5.1}
\EGAReferenceAlias{I.4.5.1.eq}{I.4.5.1}
\EGAReferenceAlias{eq:I.4.5.1}{I.4.5.1}
\EGAReferenceAlias{I.4.5.2-fr}{I.4.5.2}
\EGAReferenceAlias{I.4.5.2-es}{I.4.5.2}
\EGAReferenceAlias{I.4.5.2.fr}{I.4.5.2}
\EGAReferenceAlias{I.4.5.2.es}{I.4.5.2}
\EGAReferenceAlias{I.4.5.2.eq}{I.4.5.2}
\EGAReferenceAlias{eq:I.4.5.2}{I.4.5.2}
\EGAReferenceAlias{I.4.5.3-fr}{I.4.5.3}
\EGAReferenceAlias{I.4.5.3-es}{I.4.5.3}
\EGAReferenceAlias{I.4.5.3.fr}{I.4.5.3}
\EGAReferenceAlias{I.4.5.3.es}{I.4.5.3}
\EGAReferenceAlias{I.4.5.3.eq}{I.4.5.3}
\EGAReferenceAlias{eq:I.4.5.3}{I.4.5.3}
\EGAReferenceAlias{I.4.5.4-fr}{I.4.5.4}
\EGAReferenceAlias{I.4.5.4-es}{I.4.5.4}
\EGAReferenceAlias{I.4.5.4.fr}{I.4.5.4}
\EGAReferenceAlias{I.4.5.4.es}{I.4.5.4}
\EGAReferenceAlias{I.4.5.4.eq}{I.4.5.4}
\EGAReferenceAlias{eq:I.4.5.4}{I.4.5.4}
\EGAReferenceAlias{I.4.5.5-fr}{I.4.5.5}
\EGAReferenceAlias{I.4.5.5-es}{I.4.5.5}
\EGAReferenceAlias{I.4.5.5.fr}{I.4.5.5}
\EGAReferenceAlias{I.4.5.5.es}{I.4.5.5}
\EGAReferenceAlias{I.4.5.5.eq}{I.4.5.5}
\EGAReferenceAlias{eq:I.4.5.5}{I.4.5.5}
\EGAReferenceAlias{I.5}{section:I.5}
\EGAReferenceAlias{I.5-fr}{section:I.5}
\EGAReferenceAlias{I.5-es}{section:I.5}
\EGAReferenceAlias{I.5.fr}{section:I.5}
\EGAReferenceAlias{I.5.es}{section:I.5}
\EGAReferenceAlias{I.5.eq}{section:I.5}
\EGAReferenceAlias{eq:I.5}{section:I.5}
\EGAReferenceAlias{section:I.5-fr}{section:I.5}
\EGAReferenceAlias{section:I.5-es}{section:I.5}
\EGAReferenceAlias{section:I.5.fr}{section:I.5}
\EGAReferenceAlias{section:I.5.es}{section:I.5}
\EGAReferenceAlias{section:I.5.eq}{section:I.5}
\EGAReferenceAlias{eq:section:I.5}{section:I.5}
\EGAReferenceAlias{I.5.1}{subsection:I.5.1}
\EGAReferenceAlias{I.5.1-fr}{subsection:I.5.1}
\EGAReferenceAlias{I.5.1-es}{subsection:I.5.1}
\EGAReferenceAlias{I.5.1.fr}{subsection:I.5.1}
\EGAReferenceAlias{I.5.1.es}{subsection:I.5.1}
\EGAReferenceAlias{I.5.1.eq}{subsection:I.5.1}
\EGAReferenceAlias{eq:I.5.1}{subsection:I.5.1}
\EGAReferenceAlias{subsection:I.5.1-fr}{subsection:I.5.1}
\EGAReferenceAlias{subsection:I.5.1-es}{subsection:I.5.1}
\EGAReferenceAlias{subsection:I.5.1.fr}{subsection:I.5.1}
\EGAReferenceAlias{subsection:I.5.1.es}{subsection:I.5.1}
\EGAReferenceAlias{subsection:I.5.1.eq}{subsection:I.5.1}
\EGAReferenceAlias{eq:subsection:I.5.1}{subsection:I.5.1}
\EGAReferenceAlias{I.5.1.1-fr}{I.5.1.1}
\EGAReferenceAlias{I.5.1.1-es}{I.5.1.1}
\EGAReferenceAlias{I.5.1.1.fr}{I.5.1.1}
\EGAReferenceAlias{I.5.1.1.es}{I.5.1.1}
\EGAReferenceAlias{I.5.1.1.eq}{I.5.1.1}
\EGAReferenceAlias{eq:I.5.1.1}{I.5.1.1}
\EGAReferenceAlias{I.5.1.2-fr}{I.5.1.2}
\EGAReferenceAlias{I.5.1.2-es}{I.5.1.2}
\EGAReferenceAlias{I.5.1.2.fr}{I.5.1.2}
\EGAReferenceAlias{I.5.1.2.es}{I.5.1.2}
\EGAReferenceAlias{I.5.1.2.eq}{I.5.1.2}
\EGAReferenceAlias{eq:I.5.1.2}{I.5.1.2}
\EGAReferenceAlias{I.5.1.3-fr}{I.5.1.3}
\EGAReferenceAlias{I.5.1.3-es}{I.5.1.3}
\EGAReferenceAlias{I.5.1.3.fr}{I.5.1.3}
\EGAReferenceAlias{I.5.1.3.es}{I.5.1.3}
\EGAReferenceAlias{I.5.1.3.eq}{I.5.1.3}
\EGAReferenceAlias{eq:I.5.1.3}{I.5.1.3}
\EGAReferenceAlias{I.5.1.4-fr}{I.5.1.4}
\EGAReferenceAlias{I.5.1.4-es}{I.5.1.4}
\EGAReferenceAlias{I.5.1.4.fr}{I.5.1.4}
\EGAReferenceAlias{I.5.1.4.es}{I.5.1.4}
\EGAReferenceAlias{I.5.1.4.eq}{I.5.1.4}
\EGAReferenceAlias{eq:I.5.1.4}{I.5.1.4}
\EGAReferenceAlias{I.5.1.5-fr}{I.5.1.5}
\EGAReferenceAlias{I.5.1.5-es}{I.5.1.5}
\EGAReferenceAlias{I.5.1.5.fr}{I.5.1.5}
\EGAReferenceAlias{I.5.1.5.es}{I.5.1.5}
\EGAReferenceAlias{I.5.1.5.eq}{I.5.1.5}
\EGAReferenceAlias{eq:I.5.1.5}{I.5.1.5}
\EGAReferenceAlias{I.5.1.6-fr}{I.5.1.6}
\EGAReferenceAlias{I.5.1.6-es}{I.5.1.6}
\EGAReferenceAlias{I.5.1.6.fr}{I.5.1.6}
\EGAReferenceAlias{I.5.1.6.es}{I.5.1.6}
\EGAReferenceAlias{I.5.1.6.eq}{I.5.1.6}
\EGAReferenceAlias{eq:I.5.1.6}{I.5.1.6}
\EGAReferenceAlias{I.5.1.7-fr}{I.5.1.7}
\EGAReferenceAlias{I.5.1.7-es}{I.5.1.7}
\EGAReferenceAlias{I.5.1.7.fr}{I.5.1.7}
\EGAReferenceAlias{I.5.1.7.es}{I.5.1.7}
\EGAReferenceAlias{I.5.1.7.eq}{I.5.1.7}
\EGAReferenceAlias{eq:I.5.1.7}{I.5.1.7}
\EGAReferenceAlias{I.5.1.8-fr}{I.5.1.8}
\EGAReferenceAlias{I.5.1.8-es}{I.5.1.8}
\EGAReferenceAlias{I.5.1.8.fr}{I.5.1.8}
\EGAReferenceAlias{I.5.1.8.es}{I.5.1.8}
\EGAReferenceAlias{I.5.1.8.eq}{I.5.1.8}
\EGAReferenceAlias{eq:I.5.1.8}{I.5.1.8}
\EGAReferenceAlias{I.5.1.9-fr}{I.5.1.9}
\EGAReferenceAlias{I.5.1.9-es}{I.5.1.9}
\EGAReferenceAlias{I.5.1.9.fr}{I.5.1.9}
\EGAReferenceAlias{I.5.1.9.es}{I.5.1.9}
\EGAReferenceAlias{I.5.1.9.eq}{I.5.1.9}
\EGAReferenceAlias{eq:I.5.1.9}{I.5.1.9}
\EGAReferenceAlias{I.5.1.9.1}{I.5.1.9.1-fr}
\EGAReferenceAlias{I.5.1.9.1-es}{I.5.1.9.1-fr}
\EGAReferenceAlias{I.5.1.9.1.fr}{I.5.1.9.1-fr}
\EGAReferenceAlias{I.5.1.9.1.es}{I.5.1.9.1-fr}
\EGAReferenceAlias{I.5.1.9.1.eq}{I.5.1.9.1-fr}
\EGAReferenceAlias{eq:I.5.1.9.1}{I.5.1.9.1-fr}
\EGAReferenceAlias{I.5.1.9.2-fr}{I.5.1.9.2}
\EGAReferenceAlias{I.5.1.9.2-es}{I.5.1.9.2}
\EGAReferenceAlias{I.5.1.9.2.fr}{I.5.1.9.2}
\EGAReferenceAlias{I.5.1.9.2.es}{I.5.1.9.2}
\EGAReferenceAlias{I.5.1.9.2.eq}{I.5.1.9.2}
\EGAReferenceAlias{eq:I.5.1.9.2}{I.5.1.9.2}
\EGAReferenceAlias{I.5.1.10-fr}{I.5.1.10}
\EGAReferenceAlias{I.5.1.10-es}{I.5.1.10}
\EGAReferenceAlias{I.5.1.10.fr}{I.5.1.10}
\EGAReferenceAlias{I.5.1.10.es}{I.5.1.10}
\EGAReferenceAlias{I.5.1.10.eq}{I.5.1.10}
\EGAReferenceAlias{eq:I.5.1.10}{I.5.1.10}
\EGAReferenceAlias{I.5.2}{subsection:I.5.2}
\EGAReferenceAlias{I.5.2-fr}{subsection:I.5.2}
\EGAReferenceAlias{I.5.2-es}{subsection:I.5.2}
\EGAReferenceAlias{I.5.2.fr}{subsection:I.5.2}
\EGAReferenceAlias{I.5.2.es}{subsection:I.5.2}
\EGAReferenceAlias{I.5.2.eq}{subsection:I.5.2}
\EGAReferenceAlias{eq:I.5.2}{subsection:I.5.2}
\EGAReferenceAlias{subsection:I.5.2-fr}{subsection:I.5.2}
\EGAReferenceAlias{subsection:I.5.2-es}{subsection:I.5.2}
\EGAReferenceAlias{subsection:I.5.2.fr}{subsection:I.5.2}
\EGAReferenceAlias{subsection:I.5.2.es}{subsection:I.5.2}
\EGAReferenceAlias{subsection:I.5.2.eq}{subsection:I.5.2}
\EGAReferenceAlias{eq:subsection:I.5.2}{subsection:I.5.2}
\EGAReferenceAlias{I.5.2.1-fr}{I.5.2.1}
\EGAReferenceAlias{I.5.2.1-es}{I.5.2.1}
\EGAReferenceAlias{I.5.2.1.fr}{I.5.2.1}
\EGAReferenceAlias{I.5.2.1.es}{I.5.2.1}
\EGAReferenceAlias{I.5.2.1.eq}{I.5.2.1}
\EGAReferenceAlias{eq:I.5.2.1}{I.5.2.1}
\EGAReferenceAlias{I.5.2.2-fr}{I.5.2.2}
\EGAReferenceAlias{I.5.2.2-es}{I.5.2.2}
\EGAReferenceAlias{I.5.2.2.fr}{I.5.2.2}
\EGAReferenceAlias{I.5.2.2.es}{I.5.2.2}
\EGAReferenceAlias{I.5.2.2.eq}{I.5.2.2}
\EGAReferenceAlias{eq:I.5.2.2}{I.5.2.2}
\EGAReferenceAlias{I.5.2.3-fr}{I.5.2.3}
\EGAReferenceAlias{I.5.2.3-es}{I.5.2.3}
\EGAReferenceAlias{I.5.2.3.fr}{I.5.2.3}
\EGAReferenceAlias{I.5.2.3.es}{I.5.2.3}
\EGAReferenceAlias{I.5.2.3.eq}{I.5.2.3}
\EGAReferenceAlias{eq:I.5.2.3}{I.5.2.3}
\EGAReferenceAlias{I.5.2.4-fr}{I.5.2.4}
\EGAReferenceAlias{I.5.2.4-es}{I.5.2.4}
\EGAReferenceAlias{I.5.2.4.fr}{I.5.2.4}
\EGAReferenceAlias{I.5.2.4.es}{I.5.2.4}
\EGAReferenceAlias{I.5.2.4.eq}{I.5.2.4}
\EGAReferenceAlias{eq:I.5.2.4}{I.5.2.4}
\EGAReferenceAlias{I.5.3-fr}{subsection:I.5.3}
\EGAReferenceAlias{I.5.3-es}{subsection:I.5.3}
\EGAReferenceAlias{I.5.3.fr}{subsection:I.5.3}
\EGAReferenceAlias{I.5.3.es}{subsection:I.5.3}
\EGAReferenceAlias{I.5.3.eq}{subsection:I.5.3}
\EGAReferenceAlias{eq:I.5.3}{subsection:I.5.3}
\EGAReferenceAlias{subsection:I.5.3-fr}{subsection:I.5.3}
\EGAReferenceAlias{subsection:I.5.3-es}{subsection:I.5.3}
\EGAReferenceAlias{subsection:I.5.3.fr}{subsection:I.5.3}
\EGAReferenceAlias{subsection:I.5.3.es}{subsection:I.5.3}
\EGAReferenceAlias{subsection:I.5.3.eq}{subsection:I.5.3}
\EGAReferenceAlias{eq:subsection:I.5.3}{subsection:I.5.3}
\EGAReferenceAlias{I.5.3.1-fr}{I.5.3.1}
\EGAReferenceAlias{I.5.3.1-es}{I.5.3.1}
\EGAReferenceAlias{I.5.3.1.fr}{I.5.3.1}
\EGAReferenceAlias{I.5.3.1.es}{I.5.3.1}
\EGAReferenceAlias{I.5.3.1.eq}{I.5.3.1}
\EGAReferenceAlias{eq:I.5.3.1}{I.5.3.1}
\EGAReferenceAlias{I.5.3.1.1}{I.5.3.1.1-fr}
\EGAReferenceAlias{I.5.3.1.1-es}{I.5.3.1.1-fr}
\EGAReferenceAlias{I.5.3.1.1.fr}{I.5.3.1.1-fr}
\EGAReferenceAlias{I.5.3.1.1.es}{I.5.3.1.1-fr}
\EGAReferenceAlias{I.5.3.1.1.eq}{I.5.3.1.1-fr}
\EGAReferenceAlias{eq:I.5.3.1.1}{I.5.3.1.1-fr}
\EGAReferenceAlias{I.5.3.1.2}{I.5.3.1.2-fr}
\EGAReferenceAlias{I.5.3.1.2-es}{I.5.3.1.2-fr}
\EGAReferenceAlias{I.5.3.1.2.fr}{I.5.3.1.2-fr}
\EGAReferenceAlias{I.5.3.1.2.es}{I.5.3.1.2-fr}
\EGAReferenceAlias{I.5.3.1.2.eq}{I.5.3.1.2-fr}
\EGAReferenceAlias{eq:I.5.3.1.2}{I.5.3.1.2-fr}
\EGAReferenceAlias{I.5.3.2-fr}{I.5.3.2}
\EGAReferenceAlias{I.5.3.2-es}{I.5.3.2}
\EGAReferenceAlias{I.5.3.2.fr}{I.5.3.2}
\EGAReferenceAlias{I.5.3.2.es}{I.5.3.2}
\EGAReferenceAlias{I.5.3.2.eq}{I.5.3.2}
\EGAReferenceAlias{eq:I.5.3.2}{I.5.3.2}
\EGAReferenceAlias{I.5.3.3-fr}{I.5.3.3}
\EGAReferenceAlias{I.5.3.3-es}{I.5.3.3}
\EGAReferenceAlias{I.5.3.3.fr}{I.5.3.3}
\EGAReferenceAlias{I.5.3.3.es}{I.5.3.3}
\EGAReferenceAlias{I.5.3.3.eq}{I.5.3.3}
\EGAReferenceAlias{eq:I.5.3.3}{I.5.3.3}
\EGAReferenceAlias{I.5.3.4-fr}{I.5.3.4}
\EGAReferenceAlias{I.5.3.4-es}{I.5.3.4}
\EGAReferenceAlias{I.5.3.4.fr}{I.5.3.4}
\EGAReferenceAlias{I.5.3.4.es}{I.5.3.4}
\EGAReferenceAlias{I.5.3.4.eq}{I.5.3.4}
\EGAReferenceAlias{eq:I.5.3.4}{I.5.3.4}
\EGAReferenceAlias{I.5.3.5-fr}{I.5.3.5}
\EGAReferenceAlias{I.5.3.5-es}{I.5.3.5}
\EGAReferenceAlias{I.5.3.5.fr}{I.5.3.5}
\EGAReferenceAlias{I.5.3.5.es}{I.5.3.5}
\EGAReferenceAlias{I.5.3.5.eq}{I.5.3.5}
\EGAReferenceAlias{eq:I.5.3.5}{I.5.3.5}
\EGAReferenceAlias{I.5.3.5.1}{I.5.3.5.1-fr}
\EGAReferenceAlias{I.5.3.5.1-es}{I.5.3.5.1-fr}
\EGAReferenceAlias{I.5.3.5.1.fr}{I.5.3.5.1-fr}
\EGAReferenceAlias{I.5.3.5.1.es}{I.5.3.5.1-fr}
\EGAReferenceAlias{I.5.3.5.1.eq}{I.5.3.5.1-fr}
\EGAReferenceAlias{eq:I.5.3.5.1}{I.5.3.5.1-fr}
\EGAReferenceAlias{I.5.3.6-fr}{I.5.3.6}
\EGAReferenceAlias{I.5.3.6-es}{I.5.3.6}
\EGAReferenceAlias{I.5.3.6.fr}{I.5.3.6}
\EGAReferenceAlias{I.5.3.6.es}{I.5.3.6}
\EGAReferenceAlias{I.5.3.6.eq}{I.5.3.6}
\EGAReferenceAlias{eq:I.5.3.6}{I.5.3.6}
\EGAReferenceAlias{I.5.3.7-fr}{I.5.3.7}
\EGAReferenceAlias{I.5.3.7-es}{I.5.3.7}
\EGAReferenceAlias{I.5.3.7.fr}{I.5.3.7}
\EGAReferenceAlias{I.5.3.7.es}{I.5.3.7}
\EGAReferenceAlias{I.5.3.7.eq}{I.5.3.7}
\EGAReferenceAlias{eq:I.5.3.7}{I.5.3.7}
\EGAReferenceAlias{I.5.3.8-fr}{I.5.3.8}
\EGAReferenceAlias{I.5.3.8-es}{I.5.3.8}
\EGAReferenceAlias{I.5.3.8.fr}{I.5.3.8}
\EGAReferenceAlias{I.5.3.8.es}{I.5.3.8}
\EGAReferenceAlias{I.5.3.8.eq}{I.5.3.8}
\EGAReferenceAlias{eq:I.5.3.8}{I.5.3.8}
\EGAReferenceAlias{I.5.3.9-fr}{I.5.3.9}
\EGAReferenceAlias{I.5.3.9-es}{I.5.3.9}
\EGAReferenceAlias{I.5.3.9.fr}{I.5.3.9}
\EGAReferenceAlias{I.5.3.9.es}{I.5.3.9}
\EGAReferenceAlias{I.5.3.9.eq}{I.5.3.9}
\EGAReferenceAlias{eq:I.5.3.9}{I.5.3.9}
\EGAReferenceAlias{I.5.3.10-fr}{I.5.3.10}
\EGAReferenceAlias{I.5.3.10-es}{I.5.3.10}
\EGAReferenceAlias{I.5.3.10.fr}{I.5.3.10}
\EGAReferenceAlias{I.5.3.10.es}{I.5.3.10}
\EGAReferenceAlias{I.5.3.10.eq}{I.5.3.10}
\EGAReferenceAlias{eq:I.5.3.10}{I.5.3.10}
\EGAReferenceAlias{I.5.3.11-fr}{I.5.3.11}
\EGAReferenceAlias{I.5.3.11-es}{I.5.3.11}
\EGAReferenceAlias{I.5.3.11.fr}{I.5.3.11}
\EGAReferenceAlias{I.5.3.11.es}{I.5.3.11}
\EGAReferenceAlias{I.5.3.11.eq}{I.5.3.11}
\EGAReferenceAlias{eq:I.5.3.11}{I.5.3.11}
\EGAReferenceAlias{I.5.3.12-fr}{I.5.3.12}
\EGAReferenceAlias{I.5.3.12-es}{I.5.3.12}
\EGAReferenceAlias{I.5.3.12.fr}{I.5.3.12}
\EGAReferenceAlias{I.5.3.12.es}{I.5.3.12}
\EGAReferenceAlias{I.5.3.12.eq}{I.5.3.12}
\EGAReferenceAlias{eq:I.5.3.12}{I.5.3.12}
\EGAReferenceAlias{I.5.3.13-fr}{I.5.3.13}
\EGAReferenceAlias{I.5.3.13-es}{I.5.3.13}
\EGAReferenceAlias{I.5.3.13.fr}{I.5.3.13}
\EGAReferenceAlias{I.5.3.13.es}{I.5.3.13}
\EGAReferenceAlias{I.5.3.13.eq}{I.5.3.13}
\EGAReferenceAlias{eq:I.5.3.13}{I.5.3.13}
\EGAReferenceAlias{I.5.3.14-fr}{I.5.3.14}
\EGAReferenceAlias{I.5.3.14-es}{I.5.3.14}
\EGAReferenceAlias{I.5.3.14.fr}{I.5.3.14}
\EGAReferenceAlias{I.5.3.14.es}{I.5.3.14}
\EGAReferenceAlias{I.5.3.14.eq}{I.5.3.14}
\EGAReferenceAlias{eq:I.5.3.14}{I.5.3.14}
\EGAReferenceAlias{I.5.3.15-fr}{I.5.3.15}
\EGAReferenceAlias{I.5.3.15-es}{I.5.3.15}
\EGAReferenceAlias{I.5.3.15.fr}{I.5.3.15}
\EGAReferenceAlias{I.5.3.15.es}{I.5.3.15}
\EGAReferenceAlias{I.5.3.15.eq}{I.5.3.15}
\EGAReferenceAlias{eq:I.5.3.15}{I.5.3.15}
\EGAReferenceAlias{I.5.3.15.1}{I.5.3.15.1-fr}
\EGAReferenceAlias{I.5.3.15.1-es}{I.5.3.15.1-fr}
\EGAReferenceAlias{I.5.3.15.1.fr}{I.5.3.15.1-fr}
\EGAReferenceAlias{I.5.3.15.1.es}{I.5.3.15.1-fr}
\EGAReferenceAlias{I.5.3.15.1.eq}{I.5.3.15.1-fr}
\EGAReferenceAlias{eq:I.5.3.15.1}{I.5.3.15.1-fr}
\EGAReferenceAlias{I.5.3.16-fr}{I.5.3.16}
\EGAReferenceAlias{I.5.3.16-es}{I.5.3.16}
\EGAReferenceAlias{I.5.3.16.fr}{I.5.3.16}
\EGAReferenceAlias{I.5.3.16.es}{I.5.3.16}
\EGAReferenceAlias{I.5.3.16.eq}{I.5.3.16}
\EGAReferenceAlias{eq:I.5.3.16}{I.5.3.16}
\EGAReferenceAlias{I.5.3.17-fr}{I.5.3.17}
\EGAReferenceAlias{I.5.3.17-es}{I.5.3.17}
\EGAReferenceAlias{I.5.3.17.fr}{I.5.3.17}
\EGAReferenceAlias{I.5.3.17.es}{I.5.3.17}
\EGAReferenceAlias{I.5.3.17.eq}{I.5.3.17}
\EGAReferenceAlias{eq:I.5.3.17}{I.5.3.17}
\EGAReferenceAlias{I.5.4-fr}{subsection:I.5.4}
\EGAReferenceAlias{I.5.4-es}{subsection:I.5.4}
\EGAReferenceAlias{I.5.4.fr}{subsection:I.5.4}
\EGAReferenceAlias{I.5.4.es}{subsection:I.5.4}
\EGAReferenceAlias{I.5.4.eq}{subsection:I.5.4}
\EGAReferenceAlias{eq:I.5.4}{subsection:I.5.4}
\EGAReferenceAlias{subsection:I.5.4-fr}{subsection:I.5.4}
\EGAReferenceAlias{subsection:I.5.4-es}{subsection:I.5.4}
\EGAReferenceAlias{subsection:I.5.4.fr}{subsection:I.5.4}
\EGAReferenceAlias{subsection:I.5.4.es}{subsection:I.5.4}
\EGAReferenceAlias{subsection:I.5.4.eq}{subsection:I.5.4}
\EGAReferenceAlias{eq:subsection:I.5.4}{subsection:I.5.4}
\EGAReferenceAlias{I.5.4.1-fr}{I.5.4.1}
\EGAReferenceAlias{I.5.4.1-es}{I.5.4.1}
\EGAReferenceAlias{I.5.4.1.fr}{I.5.4.1}
\EGAReferenceAlias{I.5.4.1.es}{I.5.4.1}
\EGAReferenceAlias{I.5.4.1.eq}{I.5.4.1}
\EGAReferenceAlias{eq:I.5.4.1}{I.5.4.1}
\EGAReferenceAlias{I.5.4.2-fr}{I.5.4.2}
\EGAReferenceAlias{I.5.4.2-es}{I.5.4.2}
\EGAReferenceAlias{I.5.4.2.fr}{I.5.4.2}
\EGAReferenceAlias{I.5.4.2.es}{I.5.4.2}
\EGAReferenceAlias{I.5.4.2.eq}{I.5.4.2}
\EGAReferenceAlias{eq:I.5.4.2}{I.5.4.2}
\EGAReferenceAlias{I.5.4.3-fr}{I.5.4.3}
\EGAReferenceAlias{I.5.4.3-es}{I.5.4.3}
\EGAReferenceAlias{I.5.4.3.fr}{I.5.4.3}
\EGAReferenceAlias{I.5.4.3.es}{I.5.4.3}
\EGAReferenceAlias{I.5.4.3.eq}{I.5.4.3}
\EGAReferenceAlias{eq:I.5.4.3}{I.5.4.3}
\EGAReferenceAlias{I.5.4.4-fr}{I.5.4.4}
\EGAReferenceAlias{I.5.4.4-es}{I.5.4.4}
\EGAReferenceAlias{I.5.4.4.fr}{I.5.4.4}
\EGAReferenceAlias{I.5.4.4.es}{I.5.4.4}
\EGAReferenceAlias{I.5.4.4.eq}{I.5.4.4}
\EGAReferenceAlias{eq:I.5.4.4}{I.5.4.4}
\EGAReferenceAlias{I.5.4.5-fr}{I.5.4.5}
\EGAReferenceAlias{I.5.4.5-es}{I.5.4.5}
\EGAReferenceAlias{I.5.4.5.fr}{I.5.4.5}
\EGAReferenceAlias{I.5.4.5.es}{I.5.4.5}
\EGAReferenceAlias{I.5.4.5.eq}{I.5.4.5}
\EGAReferenceAlias{eq:I.5.4.5}{I.5.4.5}
\EGAReferenceAlias{I.5.4.6-fr}{I.5.4.6}
\EGAReferenceAlias{I.5.4.6-es}{I.5.4.6}
\EGAReferenceAlias{I.5.4.6.fr}{I.5.4.6}
\EGAReferenceAlias{I.5.4.6.es}{I.5.4.6}
\EGAReferenceAlias{I.5.4.6.eq}{I.5.4.6}
\EGAReferenceAlias{eq:I.5.4.6}{I.5.4.6}
\EGAReferenceAlias{I.5.4.7-fr}{I.5.4.7}
\EGAReferenceAlias{I.5.4.7-es}{I.5.4.7}
\EGAReferenceAlias{I.5.4.7.fr}{I.5.4.7}
\EGAReferenceAlias{I.5.4.7.es}{I.5.4.7}
\EGAReferenceAlias{I.5.4.7.eq}{I.5.4.7}
\EGAReferenceAlias{eq:I.5.4.7}{I.5.4.7}
\EGAReferenceAlias{I.5.4.8-fr}{I.5.4.8}
\EGAReferenceAlias{I.5.4.8-es}{I.5.4.8}
\EGAReferenceAlias{I.5.4.8.fr}{I.5.4.8}
\EGAReferenceAlias{I.5.4.8.es}{I.5.4.8}
\EGAReferenceAlias{I.5.4.8.eq}{I.5.4.8}
\EGAReferenceAlias{eq:I.5.4.8}{I.5.4.8}
\EGAReferenceAlias{I.5.5}{subsection:I.5.5}
\EGAReferenceAlias{I.5.5-fr}{subsection:I.5.5}
\EGAReferenceAlias{I.5.5-es}{subsection:I.5.5}
\EGAReferenceAlias{I.5.5.fr}{subsection:I.5.5}
\EGAReferenceAlias{I.5.5.es}{subsection:I.5.5}
\EGAReferenceAlias{I.5.5.eq}{subsection:I.5.5}
\EGAReferenceAlias{eq:I.5.5}{subsection:I.5.5}
\EGAReferenceAlias{subsection:I.5.5-fr}{subsection:I.5.5}
\EGAReferenceAlias{subsection:I.5.5-es}{subsection:I.5.5}
\EGAReferenceAlias{subsection:I.5.5.fr}{subsection:I.5.5}
\EGAReferenceAlias{subsection:I.5.5.es}{subsection:I.5.5}
\EGAReferenceAlias{subsection:I.5.5.eq}{subsection:I.5.5}
\EGAReferenceAlias{eq:subsection:I.5.5}{subsection:I.5.5}
\EGAReferenceAlias{I.5.5.1-fr}{I.5.5.1}
\EGAReferenceAlias{I.5.5.1-es}{I.5.5.1}
\EGAReferenceAlias{I.5.5.1.fr}{I.5.5.1}
\EGAReferenceAlias{I.5.5.1.es}{I.5.5.1}
\EGAReferenceAlias{I.5.5.1.eq}{I.5.5.1}
\EGAReferenceAlias{eq:I.5.5.1}{I.5.5.1}
\EGAReferenceAlias{I.5.5.1.1}{I.5.5.1.1-fr}
\EGAReferenceAlias{I.5.5.1.1-es}{I.5.5.1.1-fr}
\EGAReferenceAlias{I.5.5.1.1.fr}{I.5.5.1.1-fr}
\EGAReferenceAlias{I.5.5.1.1.es}{I.5.5.1.1-fr}
\EGAReferenceAlias{I.5.5.1.1.eq}{I.5.5.1.1-fr}
\EGAReferenceAlias{eq:I.5.5.1.1}{I.5.5.1.1-fr}
\EGAReferenceAlias{I.5.5.2-fr}{I.5.5.2}
\EGAReferenceAlias{I.5.5.2-es}{I.5.5.2}
\EGAReferenceAlias{I.5.5.2.fr}{I.5.5.2}
\EGAReferenceAlias{I.5.5.2.es}{I.5.5.2}
\EGAReferenceAlias{I.5.5.2.eq}{I.5.5.2}
\EGAReferenceAlias{eq:I.5.5.2}{I.5.5.2}
\EGAReferenceAlias{I.5.5.3-fr}{I.5.5.3}
\EGAReferenceAlias{I.5.5.3-es}{I.5.5.3}
\EGAReferenceAlias{I.5.5.3.fr}{I.5.5.3}
\EGAReferenceAlias{I.5.5.3.es}{I.5.5.3}
\EGAReferenceAlias{I.5.5.3.eq}{I.5.5.3}
\EGAReferenceAlias{eq:I.5.5.3}{I.5.5.3}
\EGAReferenceAlias{I.5.5.4-fr}{I.5.5.4}
\EGAReferenceAlias{I.5.5.4-es}{I.5.5.4}
\EGAReferenceAlias{I.5.5.4.fr}{I.5.5.4}
\EGAReferenceAlias{I.5.5.4.es}{I.5.5.4}
\EGAReferenceAlias{I.5.5.4.eq}{I.5.5.4}
\EGAReferenceAlias{eq:I.5.5.4}{I.5.5.4}
\EGAReferenceAlias{I.5.5.5-fr}{I.5.5.5}
\EGAReferenceAlias{I.5.5.5-es}{I.5.5.5}
\EGAReferenceAlias{I.5.5.5.fr}{I.5.5.5}
\EGAReferenceAlias{I.5.5.5.es}{I.5.5.5}
\EGAReferenceAlias{I.5.5.5.eq}{I.5.5.5}
\EGAReferenceAlias{eq:I.5.5.5}{I.5.5.5}
\EGAReferenceAlias{I.5.5.6-fr}{I.5.5.6}
\EGAReferenceAlias{I.5.5.6-es}{I.5.5.6}
\EGAReferenceAlias{I.5.5.6.fr}{I.5.5.6}
\EGAReferenceAlias{I.5.5.6.es}{I.5.5.6}
\EGAReferenceAlias{I.5.5.6.eq}{I.5.5.6}
\EGAReferenceAlias{eq:I.5.5.6}{I.5.5.6}
\EGAReferenceAlias{I.5.5.7-fr}{I.5.5.7}
\EGAReferenceAlias{I.5.5.7-es}{I.5.5.7}
\EGAReferenceAlias{I.5.5.7.fr}{I.5.5.7}
\EGAReferenceAlias{I.5.5.7.es}{I.5.5.7}
\EGAReferenceAlias{I.5.5.7.eq}{I.5.5.7}
\EGAReferenceAlias{eq:I.5.5.7}{I.5.5.7}
\EGAReferenceAlias{I.5.5.8-fr}{I.5.5.8}
\EGAReferenceAlias{I.5.5.8-es}{I.5.5.8}
\EGAReferenceAlias{I.5.5.8.fr}{I.5.5.8}
\EGAReferenceAlias{I.5.5.8.es}{I.5.5.8}
\EGAReferenceAlias{I.5.5.8.eq}{I.5.5.8}
\EGAReferenceAlias{eq:I.5.5.8}{I.5.5.8}
\EGAReferenceAlias{I.5.5.9-fr}{I.5.5.9}
\EGAReferenceAlias{I.5.5.9-es}{I.5.5.9}
\EGAReferenceAlias{I.5.5.9.fr}{I.5.5.9}
\EGAReferenceAlias{I.5.5.9.es}{I.5.5.9}
\EGAReferenceAlias{I.5.5.9.eq}{I.5.5.9}
\EGAReferenceAlias{eq:I.5.5.9}{I.5.5.9}
\EGAReferenceAlias{I.5.5.10-fr}{I.5.5.10}
\EGAReferenceAlias{I.5.5.10-es}{I.5.5.10}
\EGAReferenceAlias{I.5.5.10.fr}{I.5.5.10}
\EGAReferenceAlias{I.5.5.10.es}{I.5.5.10}
\EGAReferenceAlias{I.5.5.10.eq}{I.5.5.10}
\EGAReferenceAlias{eq:I.5.5.10}{I.5.5.10}
\EGAReferenceAlias{I.5.5.11-fr}{I.5.5.11}
\EGAReferenceAlias{I.5.5.11-es}{I.5.5.11}
\EGAReferenceAlias{I.5.5.11.fr}{I.5.5.11}
\EGAReferenceAlias{I.5.5.11.es}{I.5.5.11}
\EGAReferenceAlias{I.5.5.11.eq}{I.5.5.11}
\EGAReferenceAlias{eq:I.5.5.11}{I.5.5.11}
\EGAReferenceAlias{I.5.5.12-fr}{I.5.5.12}
\EGAReferenceAlias{I.5.5.12-es}{I.5.5.12}
\EGAReferenceAlias{I.5.5.12.fr}{I.5.5.12}
\EGAReferenceAlias{I.5.5.12.es}{I.5.5.12}
\EGAReferenceAlias{I.5.5.12.eq}{I.5.5.12}
\EGAReferenceAlias{eq:I.5.5.12}{I.5.5.12}
\EGAReferenceAlias{I.5.5.13-fr}{I.5.5.13}
\EGAReferenceAlias{I.5.5.13-es}{I.5.5.13}
\EGAReferenceAlias{I.5.5.13.fr}{I.5.5.13}
\EGAReferenceAlias{I.5.5.13.es}{I.5.5.13}
\EGAReferenceAlias{I.5.5.13.eq}{I.5.5.13}
\EGAReferenceAlias{eq:I.5.5.13}{I.5.5.13}
\EGAReferenceAlias{I.6}{section:I.6}
\EGAReferenceAlias{I.6-fr}{section:I.6}
\EGAReferenceAlias{I.6-es}{section:I.6}
\EGAReferenceAlias{I.6.fr}{section:I.6}
\EGAReferenceAlias{I.6.es}{section:I.6}
\EGAReferenceAlias{I.6.eq}{section:I.6}
\EGAReferenceAlias{eq:I.6}{section:I.6}
\EGAReferenceAlias{section:I.6-fr}{section:I.6}
\EGAReferenceAlias{section:I.6-es}{section:I.6}
\EGAReferenceAlias{section:I.6.fr}{section:I.6}
\EGAReferenceAlias{section:I.6.es}{section:I.6}
\EGAReferenceAlias{section:I.6.eq}{section:I.6}
\EGAReferenceAlias{eq:section:I.6}{section:I.6}
\EGAReferenceAlias{I.6.1}{subsection:I.6.1}
\EGAReferenceAlias{I.6.1-fr}{subsection:I.6.1}
\EGAReferenceAlias{I.6.1-es}{subsection:I.6.1}
\EGAReferenceAlias{I.6.1.fr}{subsection:I.6.1}
\EGAReferenceAlias{I.6.1.es}{subsection:I.6.1}
\EGAReferenceAlias{I.6.1.eq}{subsection:I.6.1}
\EGAReferenceAlias{eq:I.6.1}{subsection:I.6.1}
\EGAReferenceAlias{subsection:I.6.1-fr}{subsection:I.6.1}
\EGAReferenceAlias{subsection:I.6.1-es}{subsection:I.6.1}
\EGAReferenceAlias{subsection:I.6.1.fr}{subsection:I.6.1}
\EGAReferenceAlias{subsection:I.6.1.es}{subsection:I.6.1}
\EGAReferenceAlias{subsection:I.6.1.eq}{subsection:I.6.1}
\EGAReferenceAlias{eq:subsection:I.6.1}{subsection:I.6.1}
\EGAReferenceAlias{I.6.1.1-fr}{I.6.1.1}
\EGAReferenceAlias{I.6.1.1-es}{I.6.1.1}
\EGAReferenceAlias{I.6.1.1.fr}{I.6.1.1}
\EGAReferenceAlias{I.6.1.1.es}{I.6.1.1}
\EGAReferenceAlias{I.6.1.1.eq}{I.6.1.1}
\EGAReferenceAlias{eq:I.6.1.1}{I.6.1.1}
\EGAReferenceAlias{I.6.1.2-fr}{I.6.1.2}
\EGAReferenceAlias{I.6.1.2-es}{I.6.1.2}
\EGAReferenceAlias{I.6.1.2.fr}{I.6.1.2}
\EGAReferenceAlias{I.6.1.2.es}{I.6.1.2}
\EGAReferenceAlias{I.6.1.2.eq}{I.6.1.2}
\EGAReferenceAlias{eq:I.6.1.2}{I.6.1.2}
\EGAReferenceAlias{I.6.1.3-fr}{I.6.1.3}
\EGAReferenceAlias{I.6.1.3-es}{I.6.1.3}
\EGAReferenceAlias{I.6.1.3.fr}{I.6.1.3}
\EGAReferenceAlias{I.6.1.3.es}{I.6.1.3}
\EGAReferenceAlias{I.6.1.3.eq}{I.6.1.3}
\EGAReferenceAlias{eq:I.6.1.3}{I.6.1.3}
\EGAReferenceAlias{I.6.1.4-fr}{I.6.1.4}
\EGAReferenceAlias{I.6.1.4-es}{I.6.1.4}
\EGAReferenceAlias{I.6.1.4.fr}{I.6.1.4}
\EGAReferenceAlias{I.6.1.4.es}{I.6.1.4}
\EGAReferenceAlias{I.6.1.4.eq}{I.6.1.4}
\EGAReferenceAlias{eq:I.6.1.4}{I.6.1.4}
\EGAReferenceAlias{I.6.1.5-fr}{I.6.1.5}
\EGAReferenceAlias{I.6.1.5-es}{I.6.1.5}
\EGAReferenceAlias{I.6.1.5.fr}{I.6.1.5}
\EGAReferenceAlias{I.6.1.5.es}{I.6.1.5}
\EGAReferenceAlias{I.6.1.5.eq}{I.6.1.5}
\EGAReferenceAlias{eq:I.6.1.5}{I.6.1.5}
\EGAReferenceAlias{I.6.1.6-fr}{I.6.1.6}
\EGAReferenceAlias{I.6.1.6-es}{I.6.1.6}
\EGAReferenceAlias{I.6.1.6.fr}{I.6.1.6}
\EGAReferenceAlias{I.6.1.6.es}{I.6.1.6}
\EGAReferenceAlias{I.6.1.6.eq}{I.6.1.6}
\EGAReferenceAlias{eq:I.6.1.6}{I.6.1.6}
\EGAReferenceAlias{I.6.1.7-fr}{I.6.1.7}
\EGAReferenceAlias{I.6.1.7-es}{I.6.1.7}
\EGAReferenceAlias{I.6.1.7.fr}{I.6.1.7}
\EGAReferenceAlias{I.6.1.7.es}{I.6.1.7}
\EGAReferenceAlias{I.6.1.7.eq}{I.6.1.7}
\EGAReferenceAlias{eq:I.6.1.7}{I.6.1.7}
\EGAReferenceAlias{I.6.1.8-fr}{I.6.1.8}
\EGAReferenceAlias{I.6.1.8-es}{I.6.1.8}
\EGAReferenceAlias{I.6.1.8.fr}{I.6.1.8}
\EGAReferenceAlias{I.6.1.8.es}{I.6.1.8}
\EGAReferenceAlias{I.6.1.8.eq}{I.6.1.8}
\EGAReferenceAlias{eq:I.6.1.8}{I.6.1.8}
\EGAReferenceAlias{I.6.1.9-fr}{I.6.1.9}
\EGAReferenceAlias{I.6.1.9-es}{I.6.1.9}
\EGAReferenceAlias{I.6.1.9.fr}{I.6.1.9}
\EGAReferenceAlias{I.6.1.9.es}{I.6.1.9}
\EGAReferenceAlias{I.6.1.9.eq}{I.6.1.9}
\EGAReferenceAlias{eq:I.6.1.9}{I.6.1.9}
\EGAReferenceAlias{I.6.1.10-fr}{I.6.1.10}
\EGAReferenceAlias{I.6.1.10-es}{I.6.1.10}
\EGAReferenceAlias{I.6.1.10.fr}{I.6.1.10}
\EGAReferenceAlias{I.6.1.10.es}{I.6.1.10}
\EGAReferenceAlias{I.6.1.10.eq}{I.6.1.10}
\EGAReferenceAlias{eq:I.6.1.10}{I.6.1.10}
\EGAReferenceAlias{I.6.1.11-fr}{I.6.1.11}
\EGAReferenceAlias{I.6.1.11-es}{I.6.1.11}
\EGAReferenceAlias{I.6.1.11.fr}{I.6.1.11}
\EGAReferenceAlias{I.6.1.11.es}{I.6.1.11}
\EGAReferenceAlias{I.6.1.11.eq}{I.6.1.11}
\EGAReferenceAlias{eq:I.6.1.11}{I.6.1.11}
\EGAReferenceAlias{I.6.1.12-fr}{I.6.1.12}
\EGAReferenceAlias{I.6.1.12-es}{I.6.1.12}
\EGAReferenceAlias{I.6.1.12.fr}{I.6.1.12}
\EGAReferenceAlias{I.6.1.12.es}{I.6.1.12}
\EGAReferenceAlias{I.6.1.12.eq}{I.6.1.12}
\EGAReferenceAlias{eq:I.6.1.12}{I.6.1.12}
\EGAReferenceAlias{I.6.1.13-fr}{I.6.1.13}
\EGAReferenceAlias{I.6.1.13-es}{I.6.1.13}
\EGAReferenceAlias{I.6.1.13.fr}{I.6.1.13}
\EGAReferenceAlias{I.6.1.13.es}{I.6.1.13}
\EGAReferenceAlias{I.6.1.13.eq}{I.6.1.13}
\EGAReferenceAlias{eq:I.6.1.13}{I.6.1.13}
\EGAReferenceAlias{I.6.2-fr}{subsection:I.6.2}
\EGAReferenceAlias{I.6.2-es}{subsection:I.6.2}
\EGAReferenceAlias{I.6.2.fr}{subsection:I.6.2}
\EGAReferenceAlias{I.6.2.es}{subsection:I.6.2}
\EGAReferenceAlias{I.6.2.eq}{subsection:I.6.2}
\EGAReferenceAlias{eq:I.6.2}{subsection:I.6.2}
\EGAReferenceAlias{subsection:I.6.2-fr}{subsection:I.6.2}
\EGAReferenceAlias{subsection:I.6.2-es}{subsection:I.6.2}
\EGAReferenceAlias{subsection:I.6.2.fr}{subsection:I.6.2}
\EGAReferenceAlias{subsection:I.6.2.es}{subsection:I.6.2}
\EGAReferenceAlias{subsection:I.6.2.eq}{subsection:I.6.2}
\EGAReferenceAlias{eq:subsection:I.6.2}{subsection:I.6.2}
\EGAReferenceAlias{I.6.2.1-fr}{I.6.2.1}
\EGAReferenceAlias{I.6.2.1-es}{I.6.2.1}
\EGAReferenceAlias{I.6.2.1.fr}{I.6.2.1}
\EGAReferenceAlias{I.6.2.1.es}{I.6.2.1}
\EGAReferenceAlias{I.6.2.1.eq}{I.6.2.1}
\EGAReferenceAlias{eq:I.6.2.1}{I.6.2.1}
\EGAReferenceAlias{I.6.2.2-fr}{I.6.2.2}
\EGAReferenceAlias{I.6.2.2-es}{I.6.2.2}
\EGAReferenceAlias{I.6.2.2.fr}{I.6.2.2}
\EGAReferenceAlias{I.6.2.2.es}{I.6.2.2}
\EGAReferenceAlias{I.6.2.2.eq}{I.6.2.2}
\EGAReferenceAlias{eq:I.6.2.2}{I.6.2.2}
\EGAReferenceAlias{I.6.3}{subsection:I.6.3}
\EGAReferenceAlias{I.6.3-fr}{subsection:I.6.3}
\EGAReferenceAlias{I.6.3-es}{subsection:I.6.3}
\EGAReferenceAlias{I.6.3.fr}{subsection:I.6.3}
\EGAReferenceAlias{I.6.3.es}{subsection:I.6.3}
\EGAReferenceAlias{I.6.3.eq}{subsection:I.6.3}
\EGAReferenceAlias{eq:I.6.3}{subsection:I.6.3}
\EGAReferenceAlias{subsection:I.6.3-fr}{subsection:I.6.3}
\EGAReferenceAlias{subsection:I.6.3-es}{subsection:I.6.3}
\EGAReferenceAlias{subsection:I.6.3.fr}{subsection:I.6.3}
\EGAReferenceAlias{subsection:I.6.3.es}{subsection:I.6.3}
\EGAReferenceAlias{subsection:I.6.3.eq}{subsection:I.6.3}
\EGAReferenceAlias{eq:subsection:I.6.3}{subsection:I.6.3}
\EGAReferenceAlias{I.6.3.1-fr}{I.6.3.1}
\EGAReferenceAlias{I.6.3.1-es}{I.6.3.1}
\EGAReferenceAlias{I.6.3.1.fr}{I.6.3.1}
\EGAReferenceAlias{I.6.3.1.es}{I.6.3.1}
\EGAReferenceAlias{I.6.3.1.eq}{I.6.3.1}
\EGAReferenceAlias{eq:I.6.3.1}{I.6.3.1}
\EGAReferenceAlias{I.6.3.2-fr}{I.6.3.2}
\EGAReferenceAlias{I.6.3.2-es}{I.6.3.2}
\EGAReferenceAlias{I.6.3.2.fr}{I.6.3.2}
\EGAReferenceAlias{I.6.3.2.es}{I.6.3.2}
\EGAReferenceAlias{I.6.3.2.eq}{I.6.3.2}
\EGAReferenceAlias{eq:I.6.3.2}{I.6.3.2}
\EGAReferenceAlias{I.6.3.2.1-fr}{I.6.3.2.1}
\EGAReferenceAlias{I.6.3.2.1-es}{I.6.3.2.1}
\EGAReferenceAlias{I.6.3.2.1.fr}{I.6.3.2.1}
\EGAReferenceAlias{I.6.3.2.1.es}{I.6.3.2.1}
\EGAReferenceAlias{I.6.3.2.1.eq}{I.6.3.2.1}
\EGAReferenceAlias{eq:I.6.3.2.1}{I.6.3.2.1}
\EGAReferenceAlias{I.6.3.3-fr}{I.6.3.3}
\EGAReferenceAlias{I.6.3.3-es}{I.6.3.3}
\EGAReferenceAlias{I.6.3.3.fr}{I.6.3.3}
\EGAReferenceAlias{I.6.3.3.es}{I.6.3.3}
\EGAReferenceAlias{I.6.3.3.eq}{I.6.3.3}
\EGAReferenceAlias{eq:I.6.3.3}{I.6.3.3}
\EGAReferenceAlias{I.6.3.4-fr}{I.6.3.4}
\EGAReferenceAlias{I.6.3.4-es}{I.6.3.4}
\EGAReferenceAlias{I.6.3.4.fr}{I.6.3.4}
\EGAReferenceAlias{I.6.3.4.es}{I.6.3.4}
\EGAReferenceAlias{I.6.3.4.eq}{I.6.3.4}
\EGAReferenceAlias{eq:I.6.3.4}{I.6.3.4}
\EGAReferenceAlias{I.6.3.5-fr}{I.6.3.5}
\EGAReferenceAlias{I.6.3.5-es}{I.6.3.5}
\EGAReferenceAlias{I.6.3.5.fr}{I.6.3.5}
\EGAReferenceAlias{I.6.3.5.es}{I.6.3.5}
\EGAReferenceAlias{I.6.3.5.eq}{I.6.3.5}
\EGAReferenceAlias{eq:I.6.3.5}{I.6.3.5}
\EGAReferenceAlias{I.6.3.6-fr}{I.6.3.6}
\EGAReferenceAlias{I.6.3.6-es}{I.6.3.6}
\EGAReferenceAlias{I.6.3.6.fr}{I.6.3.6}
\EGAReferenceAlias{I.6.3.6.es}{I.6.3.6}
\EGAReferenceAlias{I.6.3.6.eq}{I.6.3.6}
\EGAReferenceAlias{eq:I.6.3.6}{I.6.3.6}
\EGAReferenceAlias{I.6.3.7-fr}{I.6.3.7}
\EGAReferenceAlias{I.6.3.7-es}{I.6.3.7}
\EGAReferenceAlias{I.6.3.7.fr}{I.6.3.7}
\EGAReferenceAlias{I.6.3.7.es}{I.6.3.7}
\EGAReferenceAlias{I.6.3.7.eq}{I.6.3.7}
\EGAReferenceAlias{eq:I.6.3.7}{I.6.3.7}
\EGAReferenceAlias{I.6.3.8-fr}{I.6.3.8}
\EGAReferenceAlias{I.6.3.8-es}{I.6.3.8}
\EGAReferenceAlias{I.6.3.8.fr}{I.6.3.8}
\EGAReferenceAlias{I.6.3.8.es}{I.6.3.8}
\EGAReferenceAlias{I.6.3.8.eq}{I.6.3.8}
\EGAReferenceAlias{eq:I.6.3.8}{I.6.3.8}
\EGAReferenceAlias{I.6.3.9-fr}{I.6.3.9}
\EGAReferenceAlias{I.6.3.9-es}{I.6.3.9}
\EGAReferenceAlias{I.6.3.9.fr}{I.6.3.9}
\EGAReferenceAlias{I.6.3.9.es}{I.6.3.9}
\EGAReferenceAlias{I.6.3.9.eq}{I.6.3.9}
\EGAReferenceAlias{eq:I.6.3.9}{I.6.3.9}
\EGAReferenceAlias{I.6.3.10-fr}{I.6.3.10}
\EGAReferenceAlias{I.6.3.10-es}{I.6.3.10}
\EGAReferenceAlias{I.6.3.10.fr}{I.6.3.10}
\EGAReferenceAlias{I.6.3.10.es}{I.6.3.10}
\EGAReferenceAlias{I.6.3.10.eq}{I.6.3.10}
\EGAReferenceAlias{eq:I.6.3.10}{I.6.3.10}
\EGAReferenceAlias{I.6.4-fr}{subsection:I.6.4}
\EGAReferenceAlias{I.6.4-es}{subsection:I.6.4}
\EGAReferenceAlias{I.6.4.fr}{subsection:I.6.4}
\EGAReferenceAlias{I.6.4.es}{subsection:I.6.4}
\EGAReferenceAlias{I.6.4.eq}{subsection:I.6.4}
\EGAReferenceAlias{eq:I.6.4}{subsection:I.6.4}
\EGAReferenceAlias{subsection:I.6.4-fr}{subsection:I.6.4}
\EGAReferenceAlias{subsection:I.6.4-es}{subsection:I.6.4}
\EGAReferenceAlias{subsection:I.6.4.fr}{subsection:I.6.4}
\EGAReferenceAlias{subsection:I.6.4.es}{subsection:I.6.4}
\EGAReferenceAlias{subsection:I.6.4.eq}{subsection:I.6.4}
\EGAReferenceAlias{eq:subsection:I.6.4}{subsection:I.6.4}
\EGAReferenceAlias{I.6.4.1-fr}{I.6.4.1}
\EGAReferenceAlias{I.6.4.1-es}{I.6.4.1}
\EGAReferenceAlias{I.6.4.1.fr}{I.6.4.1}
\EGAReferenceAlias{I.6.4.1.es}{I.6.4.1}
\EGAReferenceAlias{I.6.4.1.eq}{I.6.4.1}
\EGAReferenceAlias{eq:I.6.4.1}{I.6.4.1}
\EGAReferenceAlias{I.6.4.2-fr}{I.6.4.2}
\EGAReferenceAlias{I.6.4.2-es}{I.6.4.2}
\EGAReferenceAlias{I.6.4.2.fr}{I.6.4.2}
\EGAReferenceAlias{I.6.4.2.es}{I.6.4.2}
\EGAReferenceAlias{I.6.4.2.eq}{I.6.4.2}
\EGAReferenceAlias{eq:I.6.4.2}{I.6.4.2}
\EGAReferenceAlias{I.6.4.3-fr}{I.6.4.3}
\EGAReferenceAlias{I.6.4.3-es}{I.6.4.3}
\EGAReferenceAlias{I.6.4.3.fr}{I.6.4.3}
\EGAReferenceAlias{I.6.4.3.es}{I.6.4.3}
\EGAReferenceAlias{I.6.4.3.eq}{I.6.4.3}
\EGAReferenceAlias{eq:I.6.4.3}{I.6.4.3}
\EGAReferenceAlias{I.6.4.4-fr}{I.6.4.4}
\EGAReferenceAlias{I.6.4.4-es}{I.6.4.4}
\EGAReferenceAlias{I.6.4.4.fr}{I.6.4.4}
\EGAReferenceAlias{I.6.4.4.es}{I.6.4.4}
\EGAReferenceAlias{I.6.4.4.eq}{I.6.4.4}
\EGAReferenceAlias{eq:I.6.4.4}{I.6.4.4}
\EGAReferenceAlias{I.6.4.5-fr}{I.6.4.5}
\EGAReferenceAlias{I.6.4.5-es}{I.6.4.5}
\EGAReferenceAlias{I.6.4.5.fr}{I.6.4.5}
\EGAReferenceAlias{I.6.4.5.es}{I.6.4.5}
\EGAReferenceAlias{I.6.4.5.eq}{I.6.4.5}
\EGAReferenceAlias{eq:I.6.4.5}{I.6.4.5}
\EGAReferenceAlias{I.6.4.5.1}{I.6.4.5.1-fr}
\EGAReferenceAlias{I.6.4.5.1-es}{I.6.4.5.1-fr}
\EGAReferenceAlias{I.6.4.5.1.fr}{I.6.4.5.1-fr}
\EGAReferenceAlias{I.6.4.5.1.es}{I.6.4.5.1-fr}
\EGAReferenceAlias{I.6.4.5.1.eq}{I.6.4.5.1-fr}
\EGAReferenceAlias{eq:I.6.4.5.1}{I.6.4.5.1-fr}
\EGAReferenceAlias{I.6.4.5.2}{I.6.4.5.2-fr}
\EGAReferenceAlias{I.6.4.5.2-es}{I.6.4.5.2-fr}
\EGAReferenceAlias{I.6.4.5.2.fr}{I.6.4.5.2-fr}
\EGAReferenceAlias{I.6.4.5.2.es}{I.6.4.5.2-fr}
\EGAReferenceAlias{I.6.4.5.2.eq}{I.6.4.5.2-fr}
\EGAReferenceAlias{eq:I.6.4.5.2}{I.6.4.5.2-fr}
\EGAReferenceAlias{I.6.4.6-fr}{I.6.4.6}
\EGAReferenceAlias{I.6.4.6-es}{I.6.4.6}
\EGAReferenceAlias{I.6.4.6.fr}{I.6.4.6}
\EGAReferenceAlias{I.6.4.6.es}{I.6.4.6}
\EGAReferenceAlias{I.6.4.6.eq}{I.6.4.6}
\EGAReferenceAlias{eq:I.6.4.6}{I.6.4.6}
\EGAReferenceAlias{I.6.4.7-fr}{I.6.4.7}
\EGAReferenceAlias{I.6.4.7-es}{I.6.4.7}
\EGAReferenceAlias{I.6.4.7.fr}{I.6.4.7}
\EGAReferenceAlias{I.6.4.7.es}{I.6.4.7}
\EGAReferenceAlias{I.6.4.7.eq}{I.6.4.7}
\EGAReferenceAlias{eq:I.6.4.7}{I.6.4.7}
\EGAReferenceAlias{I.6.4.8-fr}{I.6.4.8}
\EGAReferenceAlias{I.6.4.8-es}{I.6.4.8}
\EGAReferenceAlias{I.6.4.8.fr}{I.6.4.8}
\EGAReferenceAlias{I.6.4.8.es}{I.6.4.8}
\EGAReferenceAlias{I.6.4.8.eq}{I.6.4.8}
\EGAReferenceAlias{eq:I.6.4.8}{I.6.4.8}
\EGAReferenceAlias{I.6.4.8.1}{I.6.4.8.1-fr}
\EGAReferenceAlias{I.6.4.8.1-es}{I.6.4.8.1-fr}
\EGAReferenceAlias{I.6.4.8.1.fr}{I.6.4.8.1-fr}
\EGAReferenceAlias{I.6.4.8.1.es}{I.6.4.8.1-fr}
\EGAReferenceAlias{I.6.4.8.1.eq}{I.6.4.8.1-fr}
\EGAReferenceAlias{eq:I.6.4.8.1}{I.6.4.8.1-fr}
\EGAReferenceAlias{I.6.4.8.2}{I.6.4.8.2-fr}
\EGAReferenceAlias{I.6.4.8.2-es}{I.6.4.8.2-fr}
\EGAReferenceAlias{I.6.4.8.2.fr}{I.6.4.8.2-fr}
\EGAReferenceAlias{I.6.4.8.2.es}{I.6.4.8.2-fr}
\EGAReferenceAlias{I.6.4.8.2.eq}{I.6.4.8.2-fr}
\EGAReferenceAlias{eq:I.6.4.8.2}{I.6.4.8.2-fr}
\EGAReferenceAlias{I.6.4.9-fr}{I.6.4.9}
\EGAReferenceAlias{I.6.4.9-es}{I.6.4.9}
\EGAReferenceAlias{I.6.4.9.fr}{I.6.4.9}
\EGAReferenceAlias{I.6.4.9.es}{I.6.4.9}
\EGAReferenceAlias{I.6.4.9.eq}{I.6.4.9}
\EGAReferenceAlias{eq:I.6.4.9}{I.6.4.9}
\EGAReferenceAlias{I.6.4.10-fr}{I.6.4.10}
\EGAReferenceAlias{I.6.4.10-es}{I.6.4.10}
\EGAReferenceAlias{I.6.4.10.fr}{I.6.4.10}
\EGAReferenceAlias{I.6.4.10.es}{I.6.4.10}
\EGAReferenceAlias{I.6.4.10.eq}{I.6.4.10}
\EGAReferenceAlias{eq:I.6.4.10}{I.6.4.10}
\EGAReferenceAlias{I.6.4.11-fr}{I.6.4.11}
\EGAReferenceAlias{I.6.4.11-es}{I.6.4.11}
\EGAReferenceAlias{I.6.4.11.fr}{I.6.4.11}
\EGAReferenceAlias{I.6.4.11.es}{I.6.4.11}
\EGAReferenceAlias{I.6.4.11.eq}{I.6.4.11}
\EGAReferenceAlias{eq:I.6.4.11}{I.6.4.11}
\EGAReferenceAlias{I.6.4.12-fr}{I.6.4.12}
\EGAReferenceAlias{I.6.4.12-es}{I.6.4.12}
\EGAReferenceAlias{I.6.4.12.fr}{I.6.4.12}
\EGAReferenceAlias{I.6.4.12.es}{I.6.4.12}
\EGAReferenceAlias{I.6.4.12.eq}{I.6.4.12}
\EGAReferenceAlias{eq:I.6.4.12}{I.6.4.12}
\EGAReferenceAlias{I.6.4.13-fr}{I.6.4.13}
\EGAReferenceAlias{I.6.4.13-es}{I.6.4.13}
\EGAReferenceAlias{I.6.4.13.fr}{I.6.4.13}
\EGAReferenceAlias{I.6.4.13.es}{I.6.4.13}
\EGAReferenceAlias{I.6.4.13.eq}{I.6.4.13}
\EGAReferenceAlias{eq:I.6.4.13}{I.6.4.13}
\EGAReferenceAlias{I.6.5-fr}{subsection:I.6.5}
\EGAReferenceAlias{I.6.5-es}{subsection:I.6.5}
\EGAReferenceAlias{I.6.5.fr}{subsection:I.6.5}
\EGAReferenceAlias{I.6.5.es}{subsection:I.6.5}
\EGAReferenceAlias{I.6.5.eq}{subsection:I.6.5}
\EGAReferenceAlias{eq:I.6.5}{subsection:I.6.5}
\EGAReferenceAlias{subsection:I.6.5-fr}{subsection:I.6.5}
\EGAReferenceAlias{subsection:I.6.5-es}{subsection:I.6.5}
\EGAReferenceAlias{subsection:I.6.5.fr}{subsection:I.6.5}
\EGAReferenceAlias{subsection:I.6.5.es}{subsection:I.6.5}
\EGAReferenceAlias{subsection:I.6.5.eq}{subsection:I.6.5}
\EGAReferenceAlias{eq:subsection:I.6.5}{subsection:I.6.5}
\EGAReferenceAlias{I.6.5.1-fr}{I.6.5.1}
\EGAReferenceAlias{I.6.5.1-es}{I.6.5.1}
\EGAReferenceAlias{I.6.5.1.fr}{I.6.5.1}
\EGAReferenceAlias{I.6.5.1.es}{I.6.5.1}
\EGAReferenceAlias{I.6.5.1.eq}{I.6.5.1}
\EGAReferenceAlias{eq:I.6.5.1}{I.6.5.1}
\EGAReferenceAlias{I.6.5.1.1}{I.6.5.1.1-fr}
\EGAReferenceAlias{I.6.5.1.1-es}{I.6.5.1.1-fr}
\EGAReferenceAlias{I.6.5.1.1.fr}{I.6.5.1.1-fr}
\EGAReferenceAlias{I.6.5.1.1.es}{I.6.5.1.1-fr}
\EGAReferenceAlias{I.6.5.1.1.eq}{I.6.5.1.1-fr}
\EGAReferenceAlias{eq:I.6.5.1.1}{I.6.5.1.1-fr}
\EGAReferenceAlias{I.6.5.2-fr}{I.6.5.2}
\EGAReferenceAlias{I.6.5.2-es}{I.6.5.2}
\EGAReferenceAlias{I.6.5.2.fr}{I.6.5.2}
\EGAReferenceAlias{I.6.5.2.es}{I.6.5.2}
\EGAReferenceAlias{I.6.5.2.eq}{I.6.5.2}
\EGAReferenceAlias{eq:I.6.5.2}{I.6.5.2}
\EGAReferenceAlias{I.6.5.3-fr}{I.6.5.3}
\EGAReferenceAlias{I.6.5.3-es}{I.6.5.3}
\EGAReferenceAlias{I.6.5.3.fr}{I.6.5.3}
\EGAReferenceAlias{I.6.5.3.es}{I.6.5.3}
\EGAReferenceAlias{I.6.5.3.eq}{I.6.5.3}
\EGAReferenceAlias{eq:I.6.5.3}{I.6.5.3}
\EGAReferenceAlias{I.6.5.4-fr}{I.6.5.4}
\EGAReferenceAlias{I.6.5.4-es}{I.6.5.4}
\EGAReferenceAlias{I.6.5.4.fr}{I.6.5.4}
\EGAReferenceAlias{I.6.5.4.es}{I.6.5.4}
\EGAReferenceAlias{I.6.5.4.eq}{I.6.5.4}
\EGAReferenceAlias{eq:I.6.5.4}{I.6.5.4}
\EGAReferenceAlias{I.6.5.5-fr}{I.6.5.5}
\EGAReferenceAlias{I.6.5.5-es}{I.6.5.5}
\EGAReferenceAlias{I.6.5.5.fr}{I.6.5.5}
\EGAReferenceAlias{I.6.5.5.es}{I.6.5.5}
\EGAReferenceAlias{I.6.5.5.eq}{I.6.5.5}
\EGAReferenceAlias{eq:I.6.5.5}{I.6.5.5}
\EGAReferenceAlias{I.6.6}{subsection:I.6.6}
\EGAReferenceAlias{I.6.6-fr}{subsection:I.6.6}
\EGAReferenceAlias{I.6.6-es}{subsection:I.6.6}
\EGAReferenceAlias{I.6.6.fr}{subsection:I.6.6}
\EGAReferenceAlias{I.6.6.es}{subsection:I.6.6}
\EGAReferenceAlias{I.6.6.eq}{subsection:I.6.6}
\EGAReferenceAlias{eq:I.6.6}{subsection:I.6.6}
\EGAReferenceAlias{subsection:I.6.6-fr}{subsection:I.6.6}
\EGAReferenceAlias{subsection:I.6.6-es}{subsection:I.6.6}
\EGAReferenceAlias{subsection:I.6.6.fr}{subsection:I.6.6}
\EGAReferenceAlias{subsection:I.6.6.es}{subsection:I.6.6}
\EGAReferenceAlias{subsection:I.6.6.eq}{subsection:I.6.6}
\EGAReferenceAlias{eq:subsection:I.6.6}{subsection:I.6.6}
\EGAReferenceAlias{I.6.6.1-fr}{I.6.6.1}
\EGAReferenceAlias{I.6.6.1-es}{I.6.6.1}
\EGAReferenceAlias{I.6.6.1.fr}{I.6.6.1}
\EGAReferenceAlias{I.6.6.1.es}{I.6.6.1}
\EGAReferenceAlias{I.6.6.1.eq}{I.6.6.1}
\EGAReferenceAlias{eq:I.6.6.1}{I.6.6.1}
\EGAReferenceAlias{I.6.6.2-fr}{I.6.6.2}
\EGAReferenceAlias{I.6.6.2-es}{I.6.6.2}
\EGAReferenceAlias{I.6.6.2.fr}{I.6.6.2}
\EGAReferenceAlias{I.6.6.2.es}{I.6.6.2}
\EGAReferenceAlias{I.6.6.2.eq}{I.6.6.2}
\EGAReferenceAlias{eq:I.6.6.2}{I.6.6.2}
\EGAReferenceAlias{I.6.6.3-fr}{I.6.6.3}
\EGAReferenceAlias{I.6.6.3-es}{I.6.6.3}
\EGAReferenceAlias{I.6.6.3.fr}{I.6.6.3}
\EGAReferenceAlias{I.6.6.3.es}{I.6.6.3}
\EGAReferenceAlias{I.6.6.3.eq}{I.6.6.3}
\EGAReferenceAlias{eq:I.6.6.3}{I.6.6.3}
\EGAReferenceAlias{I.6.6.4-fr}{I.6.6.4}
\EGAReferenceAlias{I.6.6.4-es}{I.6.6.4}
\EGAReferenceAlias{I.6.6.4.fr}{I.6.6.4}
\EGAReferenceAlias{I.6.6.4.es}{I.6.6.4}
\EGAReferenceAlias{I.6.6.4.eq}{I.6.6.4}
\EGAReferenceAlias{eq:I.6.6.4}{I.6.6.4}
\EGAReferenceAlias{I.6.6.5-fr}{I.6.6.5}
\EGAReferenceAlias{I.6.6.5-es}{I.6.6.5}
\EGAReferenceAlias{I.6.6.5.fr}{I.6.6.5}
\EGAReferenceAlias{I.6.6.5.es}{I.6.6.5}
\EGAReferenceAlias{I.6.6.5.eq}{I.6.6.5}
\EGAReferenceAlias{eq:I.6.6.5}{I.6.6.5}
\EGAReferenceAlias{I.6.6.6-fr}{I.6.6.6}
\EGAReferenceAlias{I.6.6.6-es}{I.6.6.6}
\EGAReferenceAlias{I.6.6.6.fr}{I.6.6.6}
\EGAReferenceAlias{I.6.6.6.es}{I.6.6.6}
\EGAReferenceAlias{I.6.6.6.eq}{I.6.6.6}
\EGAReferenceAlias{eq:I.6.6.6}{I.6.6.6}
\EGAReferenceAlias{I.6.6.7-fr}{I.6.6.7}
\EGAReferenceAlias{I.6.6.7-es}{I.6.6.7}
\EGAReferenceAlias{I.6.6.7.fr}{I.6.6.7}
\EGAReferenceAlias{I.6.6.7.es}{I.6.6.7}
\EGAReferenceAlias{I.6.6.7.eq}{I.6.6.7}
\EGAReferenceAlias{eq:I.6.6.7}{I.6.6.7}
\EGAReferenceAlias{I.6.6.8-fr}{I.6.6.8}
\EGAReferenceAlias{I.6.6.8-es}{I.6.6.8}
\EGAReferenceAlias{I.6.6.8.fr}{I.6.6.8}
\EGAReferenceAlias{I.6.6.8.es}{I.6.6.8}
\EGAReferenceAlias{I.6.6.8.eq}{I.6.6.8}
\EGAReferenceAlias{eq:I.6.6.8}{I.6.6.8}
\EGAReferenceAlias{I.7}{section:I.7}
\EGAReferenceAlias{I.7-fr}{section:I.7}
\EGAReferenceAlias{I.7-es}{section:I.7}
\EGAReferenceAlias{I.7.fr}{section:I.7}
\EGAReferenceAlias{I.7.es}{section:I.7}
\EGAReferenceAlias{I.7.eq}{section:I.7}
\EGAReferenceAlias{eq:I.7}{section:I.7}
\EGAReferenceAlias{section:I.7-fr}{section:I.7}
\EGAReferenceAlias{section:I.7-es}{section:I.7}
\EGAReferenceAlias{section:I.7.fr}{section:I.7}
\EGAReferenceAlias{section:I.7.es}{section:I.7}
\EGAReferenceAlias{section:I.7.eq}{section:I.7}
\EGAReferenceAlias{eq:section:I.7}{section:I.7}
\EGAReferenceAlias{I.7.1}{subsection:I.7.1}
\EGAReferenceAlias{I.7.1-fr}{subsection:I.7.1}
\EGAReferenceAlias{I.7.1-es}{subsection:I.7.1}
\EGAReferenceAlias{I.7.1.fr}{subsection:I.7.1}
\EGAReferenceAlias{I.7.1.es}{subsection:I.7.1}
\EGAReferenceAlias{I.7.1.eq}{subsection:I.7.1}
\EGAReferenceAlias{eq:I.7.1}{subsection:I.7.1}
\EGAReferenceAlias{subsection:I.7.1-fr}{subsection:I.7.1}
\EGAReferenceAlias{subsection:I.7.1-es}{subsection:I.7.1}
\EGAReferenceAlias{subsection:I.7.1.fr}{subsection:I.7.1}
\EGAReferenceAlias{subsection:I.7.1.es}{subsection:I.7.1}
\EGAReferenceAlias{subsection:I.7.1.eq}{subsection:I.7.1}
\EGAReferenceAlias{eq:subsection:I.7.1}{subsection:I.7.1}
\EGAReferenceAlias{I.7.1.1-fr}{I.7.1.1}
\EGAReferenceAlias{I.7.1.1-es}{I.7.1.1}
\EGAReferenceAlias{I.7.1.1.fr}{I.7.1.1}
\EGAReferenceAlias{I.7.1.1.es}{I.7.1.1}
\EGAReferenceAlias{I.7.1.1.eq}{I.7.1.1}
\EGAReferenceAlias{eq:I.7.1.1}{I.7.1.1}
\EGAReferenceAlias{I.7.1.2-fr}{I.7.1.2}
\EGAReferenceAlias{I.7.1.2-es}{I.7.1.2}
\EGAReferenceAlias{I.7.1.2.fr}{I.7.1.2}
\EGAReferenceAlias{I.7.1.2.es}{I.7.1.2}
\EGAReferenceAlias{I.7.1.2.eq}{I.7.1.2}
\EGAReferenceAlias{eq:I.7.1.2}{I.7.1.2}
\EGAReferenceAlias{I.7.1.3-fr}{I.7.1.3}
\EGAReferenceAlias{I.7.1.3-es}{I.7.1.3}
\EGAReferenceAlias{I.7.1.3.fr}{I.7.1.3}
\EGAReferenceAlias{I.7.1.3.es}{I.7.1.3}
\EGAReferenceAlias{I.7.1.3.eq}{I.7.1.3}
\EGAReferenceAlias{eq:I.7.1.3}{I.7.1.3}
\EGAReferenceAlias{I.7.1.4-fr}{I.7.1.4}
\EGAReferenceAlias{I.7.1.4-es}{I.7.1.4}
\EGAReferenceAlias{I.7.1.4.fr}{I.7.1.4}
\EGAReferenceAlias{I.7.1.4.es}{I.7.1.4}
\EGAReferenceAlias{I.7.1.4.eq}{I.7.1.4}
\EGAReferenceAlias{eq:I.7.1.4}{I.7.1.4}
\EGAReferenceAlias{I.7.1.5-fr}{I.7.1.5}
\EGAReferenceAlias{I.7.1.5-es}{I.7.1.5}
\EGAReferenceAlias{I.7.1.5.fr}{I.7.1.5}
\EGAReferenceAlias{I.7.1.5.es}{I.7.1.5}
\EGAReferenceAlias{I.7.1.5.eq}{I.7.1.5}
\EGAReferenceAlias{eq:I.7.1.5}{I.7.1.5}
\EGAReferenceAlias{I.7.1.6-fr}{I.7.1.6}
\EGAReferenceAlias{I.7.1.6-es}{I.7.1.6}
\EGAReferenceAlias{I.7.1.6.fr}{I.7.1.6}
\EGAReferenceAlias{I.7.1.6.es}{I.7.1.6}
\EGAReferenceAlias{I.7.1.6.eq}{I.7.1.6}
\EGAReferenceAlias{eq:I.7.1.6}{I.7.1.6}
\EGAReferenceAlias{I.7.1.7-fr}{I.7.1.7}
\EGAReferenceAlias{I.7.1.7-es}{I.7.1.7}
\EGAReferenceAlias{I.7.1.7.fr}{I.7.1.7}
\EGAReferenceAlias{I.7.1.7.es}{I.7.1.7}
\EGAReferenceAlias{I.7.1.7.eq}{I.7.1.7}
\EGAReferenceAlias{eq:I.7.1.7}{I.7.1.7}
\EGAReferenceAlias{I.7.1.8-fr}{I.7.1.8}
\EGAReferenceAlias{I.7.1.8-es}{I.7.1.8}
\EGAReferenceAlias{I.7.1.8.fr}{I.7.1.8}
\EGAReferenceAlias{I.7.1.8.es}{I.7.1.8}
\EGAReferenceAlias{I.7.1.8.eq}{I.7.1.8}
\EGAReferenceAlias{eq:I.7.1.8}{I.7.1.8}
\EGAReferenceAlias{I.7.1.9-fr}{I.7.1.9}
\EGAReferenceAlias{I.7.1.9-es}{I.7.1.9}
\EGAReferenceAlias{I.7.1.9.fr}{I.7.1.9}
\EGAReferenceAlias{I.7.1.9.es}{I.7.1.9}
\EGAReferenceAlias{I.7.1.9.eq}{I.7.1.9}
\EGAReferenceAlias{eq:I.7.1.9}{I.7.1.9}
\EGAReferenceAlias{I.7.1.9.1-fr}{I.7.1.9.1}
\EGAReferenceAlias{I.7.1.9.1-es}{I.7.1.9.1}
\EGAReferenceAlias{I.7.1.9.1.fr}{I.7.1.9.1}
\EGAReferenceAlias{I.7.1.9.1.es}{I.7.1.9.1}
\EGAReferenceAlias{I.7.1.9.1.eq}{I.7.1.9.1}
\EGAReferenceAlias{eq:I.7.1.9.1}{I.7.1.9.1}
\EGAReferenceAlias{I.7.1.10-fr}{I.7.1.10}
\EGAReferenceAlias{I.7.1.10-es}{I.7.1.10}
\EGAReferenceAlias{I.7.1.10.fr}{I.7.1.10}
\EGAReferenceAlias{I.7.1.10.es}{I.7.1.10}
\EGAReferenceAlias{I.7.1.10.eq}{I.7.1.10}
\EGAReferenceAlias{eq:I.7.1.10}{I.7.1.10}
\EGAReferenceAlias{I.7.1.11-fr}{I.7.1.11}
\EGAReferenceAlias{I.7.1.11-es}{I.7.1.11}
\EGAReferenceAlias{I.7.1.11.fr}{I.7.1.11}
\EGAReferenceAlias{I.7.1.11.es}{I.7.1.11}
\EGAReferenceAlias{I.7.1.11.eq}{I.7.1.11}
\EGAReferenceAlias{eq:I.7.1.11}{I.7.1.11}
\EGAReferenceAlias{I.7.1.12-fr}{I.7.1.12}
\EGAReferenceAlias{I.7.1.12-es}{I.7.1.12}
\EGAReferenceAlias{I.7.1.12.fr}{I.7.1.12}
\EGAReferenceAlias{I.7.1.12.es}{I.7.1.12}
\EGAReferenceAlias{I.7.1.12.eq}{I.7.1.12}
\EGAReferenceAlias{eq:I.7.1.12}{I.7.1.12}
\EGAReferenceAlias{I.7.1.13-fr}{I.7.1.13}
\EGAReferenceAlias{I.7.1.13-es}{I.7.1.13}
\EGAReferenceAlias{I.7.1.13.fr}{I.7.1.13}
\EGAReferenceAlias{I.7.1.13.es}{I.7.1.13}
\EGAReferenceAlias{I.7.1.13.eq}{I.7.1.13}
\EGAReferenceAlias{eq:I.7.1.13}{I.7.1.13}
\EGAReferenceAlias{I.7.1.14-fr}{I.7.1.14}
\EGAReferenceAlias{I.7.1.14-es}{I.7.1.14}
\EGAReferenceAlias{I.7.1.14.fr}{I.7.1.14}
\EGAReferenceAlias{I.7.1.14.es}{I.7.1.14}
\EGAReferenceAlias{I.7.1.14.eq}{I.7.1.14}
\EGAReferenceAlias{eq:I.7.1.14}{I.7.1.14}
\EGAReferenceAlias{I.7.1.15-fr}{I.7.1.15}
\EGAReferenceAlias{I.7.1.15-es}{I.7.1.15}
\EGAReferenceAlias{I.7.1.15.fr}{I.7.1.15}
\EGAReferenceAlias{I.7.1.15.es}{I.7.1.15}
\EGAReferenceAlias{I.7.1.15.eq}{I.7.1.15}
\EGAReferenceAlias{eq:I.7.1.15}{I.7.1.15}
\EGAReferenceAlias{I.7.1.16-fr}{I.7.1.16}
\EGAReferenceAlias{I.7.1.16-es}{I.7.1.16}
\EGAReferenceAlias{I.7.1.16.fr}{I.7.1.16}
\EGAReferenceAlias{I.7.1.16.es}{I.7.1.16}
\EGAReferenceAlias{I.7.1.16.eq}{I.7.1.16}
\EGAReferenceAlias{eq:I.7.1.16}{I.7.1.16}
\EGAReferenceAlias{I.7.2-fr}{subsection:I.7.2}
\EGAReferenceAlias{I.7.2-es}{subsection:I.7.2}
\EGAReferenceAlias{I.7.2.fr}{subsection:I.7.2}
\EGAReferenceAlias{I.7.2.es}{subsection:I.7.2}
\EGAReferenceAlias{I.7.2.eq}{subsection:I.7.2}
\EGAReferenceAlias{eq:I.7.2}{subsection:I.7.2}
\EGAReferenceAlias{subsection:I.7.2-fr}{subsection:I.7.2}
\EGAReferenceAlias{subsection:I.7.2-es}{subsection:I.7.2}
\EGAReferenceAlias{subsection:I.7.2.fr}{subsection:I.7.2}
\EGAReferenceAlias{subsection:I.7.2.es}{subsection:I.7.2}
\EGAReferenceAlias{subsection:I.7.2.eq}{subsection:I.7.2}
\EGAReferenceAlias{eq:subsection:I.7.2}{subsection:I.7.2}
\EGAReferenceAlias{I.7.2.1-fr}{I.7.2.1}
\EGAReferenceAlias{I.7.2.1-es}{I.7.2.1}
\EGAReferenceAlias{I.7.2.1.fr}{I.7.2.1}
\EGAReferenceAlias{I.7.2.1.es}{I.7.2.1}
\EGAReferenceAlias{I.7.2.1.eq}{I.7.2.1}
\EGAReferenceAlias{eq:I.7.2.1}{I.7.2.1}
\EGAReferenceAlias{I.7.2.2-fr}{I.7.2.2}
\EGAReferenceAlias{I.7.2.2-es}{I.7.2.2}
\EGAReferenceAlias{I.7.2.2.fr}{I.7.2.2}
\EGAReferenceAlias{I.7.2.2.es}{I.7.2.2}
\EGAReferenceAlias{I.7.2.2.eq}{I.7.2.2}
\EGAReferenceAlias{eq:I.7.2.2}{I.7.2.2}
\EGAReferenceAlias{I.7.2.2.1-fr}{I.7.2.2.1}
\EGAReferenceAlias{I.7.2.2.1-es}{I.7.2.2.1}
\EGAReferenceAlias{I.7.2.2.1.fr}{I.7.2.2.1}
\EGAReferenceAlias{I.7.2.2.1.es}{I.7.2.2.1}
\EGAReferenceAlias{I.7.2.2.1.eq}{I.7.2.2.1}
\EGAReferenceAlias{eq:I.7.2.2.1}{I.7.2.2.1}
\EGAReferenceAlias{I.7.2.3-fr}{I.7.2.3}
\EGAReferenceAlias{I.7.2.3-es}{I.7.2.3}
\EGAReferenceAlias{I.7.2.3.fr}{I.7.2.3}
\EGAReferenceAlias{I.7.2.3.es}{I.7.2.3}
\EGAReferenceAlias{I.7.2.3.eq}{I.7.2.3}
\EGAReferenceAlias{eq:I.7.2.3}{I.7.2.3}
\EGAReferenceAlias{I.7.2.4-fr}{I.7.2.4}
\EGAReferenceAlias{I.7.2.4-es}{I.7.2.4}
\EGAReferenceAlias{I.7.2.4.fr}{I.7.2.4}
\EGAReferenceAlias{I.7.2.4.es}{I.7.2.4}
\EGAReferenceAlias{I.7.2.4.eq}{I.7.2.4}
\EGAReferenceAlias{eq:I.7.2.4}{I.7.2.4}
\EGAReferenceAlias{I.7.2.5-fr}{I.7.2.5}
\EGAReferenceAlias{I.7.2.5-es}{I.7.2.5}
\EGAReferenceAlias{I.7.2.5.fr}{I.7.2.5}
\EGAReferenceAlias{I.7.2.5.es}{I.7.2.5}
\EGAReferenceAlias{I.7.2.5.eq}{I.7.2.5}
\EGAReferenceAlias{eq:I.7.2.5}{I.7.2.5}
\EGAReferenceAlias{I.7.2.6-fr}{I.7.2.6}
\EGAReferenceAlias{I.7.2.6-es}{I.7.2.6}
\EGAReferenceAlias{I.7.2.6.fr}{I.7.2.6}
\EGAReferenceAlias{I.7.2.6.es}{I.7.2.6}
\EGAReferenceAlias{I.7.2.6.eq}{I.7.2.6}
\EGAReferenceAlias{eq:I.7.2.6}{I.7.2.6}
\EGAReferenceAlias{I.7.2.7-fr}{I.7.2.7}
\EGAReferenceAlias{I.7.2.7-es}{I.7.2.7}
\EGAReferenceAlias{I.7.2.7.fr}{I.7.2.7}
\EGAReferenceAlias{I.7.2.7.es}{I.7.2.7}
\EGAReferenceAlias{I.7.2.7.eq}{I.7.2.7}
\EGAReferenceAlias{eq:I.7.2.7}{I.7.2.7}
\EGAReferenceAlias{I.7.2.8-fr}{I.7.2.8}
\EGAReferenceAlias{I.7.2.8-es}{I.7.2.8}
\EGAReferenceAlias{I.7.2.8.fr}{I.7.2.8}
\EGAReferenceAlias{I.7.2.8.es}{I.7.2.8}
\EGAReferenceAlias{I.7.2.8.eq}{I.7.2.8}
\EGAReferenceAlias{eq:I.7.2.8}{I.7.2.8}
\EGAReferenceAlias{I.7.2.8.1-fr}{I.7.2.8.1}
\EGAReferenceAlias{I.7.2.8.1-es}{I.7.2.8.1}
\EGAReferenceAlias{I.7.2.8.1.fr}{I.7.2.8.1}
\EGAReferenceAlias{I.7.2.8.1.es}{I.7.2.8.1}
\EGAReferenceAlias{I.7.2.8.1.eq}{I.7.2.8.1}
\EGAReferenceAlias{eq:I.7.2.8.1}{I.7.2.8.1}
\EGAReferenceAlias{I.7.2.9-fr}{I.7.2.9}
\EGAReferenceAlias{I.7.2.9-es}{I.7.2.9}
\EGAReferenceAlias{I.7.2.9.fr}{I.7.2.9}
\EGAReferenceAlias{I.7.2.9.es}{I.7.2.9}
\EGAReferenceAlias{I.7.2.9.eq}{I.7.2.9}
\EGAReferenceAlias{eq:I.7.2.9}{I.7.2.9}
\EGAReferenceAlias{I.7.3}{subsection:I.7.3}
\EGAReferenceAlias{I.7.3-fr}{subsection:I.7.3}
\EGAReferenceAlias{I.7.3-es}{subsection:I.7.3}
\EGAReferenceAlias{I.7.3.fr}{subsection:I.7.3}
\EGAReferenceAlias{I.7.3.es}{subsection:I.7.3}
\EGAReferenceAlias{I.7.3.eq}{subsection:I.7.3}
\EGAReferenceAlias{eq:I.7.3}{subsection:I.7.3}
\EGAReferenceAlias{subsection:I.7.3-fr}{subsection:I.7.3}
\EGAReferenceAlias{subsection:I.7.3-es}{subsection:I.7.3}
\EGAReferenceAlias{subsection:I.7.3.fr}{subsection:I.7.3}
\EGAReferenceAlias{subsection:I.7.3.es}{subsection:I.7.3}
\EGAReferenceAlias{subsection:I.7.3.eq}{subsection:I.7.3}
\EGAReferenceAlias{eq:subsection:I.7.3}{subsection:I.7.3}
\EGAReferenceAlias{I.7.3.1-fr}{I.7.3.1}
\EGAReferenceAlias{I.7.3.1-es}{I.7.3.1}
\EGAReferenceAlias{I.7.3.1.fr}{I.7.3.1}
\EGAReferenceAlias{I.7.3.1.es}{I.7.3.1}
\EGAReferenceAlias{I.7.3.1.eq}{I.7.3.1}
\EGAReferenceAlias{eq:I.7.3.1}{I.7.3.1}
\EGAReferenceAlias{I.7.3.2-fr}{I.7.3.2}
\EGAReferenceAlias{I.7.3.2-es}{I.7.3.2}
\EGAReferenceAlias{I.7.3.2.fr}{I.7.3.2}
\EGAReferenceAlias{I.7.3.2.es}{I.7.3.2}
\EGAReferenceAlias{I.7.3.2.eq}{I.7.3.2}
\EGAReferenceAlias{eq:I.7.3.2}{I.7.3.2}
\EGAReferenceAlias{I.7.3.3-fr}{I.7.3.3}
\EGAReferenceAlias{I.7.3.3-es}{I.7.3.3}
\EGAReferenceAlias{I.7.3.3.fr}{I.7.3.3}
\EGAReferenceAlias{I.7.3.3.es}{I.7.3.3}
\EGAReferenceAlias{I.7.3.3.eq}{I.7.3.3}
\EGAReferenceAlias{eq:I.7.3.3}{I.7.3.3}
\EGAReferenceAlias{I.7.3.4-fr}{I.7.3.4}
\EGAReferenceAlias{I.7.3.4-es}{I.7.3.4}
\EGAReferenceAlias{I.7.3.4.fr}{I.7.3.4}
\EGAReferenceAlias{I.7.3.4.es}{I.7.3.4}
\EGAReferenceAlias{I.7.3.4.eq}{I.7.3.4}
\EGAReferenceAlias{eq:I.7.3.4}{I.7.3.4}
\EGAReferenceAlias{I.7.3.5-fr}{I.7.3.5}
\EGAReferenceAlias{I.7.3.5-es}{I.7.3.5}
\EGAReferenceAlias{I.7.3.5.fr}{I.7.3.5}
\EGAReferenceAlias{I.7.3.5.es}{I.7.3.5}
\EGAReferenceAlias{I.7.3.5.eq}{I.7.3.5}
\EGAReferenceAlias{eq:I.7.3.5}{I.7.3.5}
\EGAReferenceAlias{I.7.3.6-fr}{I.7.3.6}
\EGAReferenceAlias{I.7.3.6-es}{I.7.3.6}
\EGAReferenceAlias{I.7.3.6.fr}{I.7.3.6}
\EGAReferenceAlias{I.7.3.6.es}{I.7.3.6}
\EGAReferenceAlias{I.7.3.6.eq}{I.7.3.6}
\EGAReferenceAlias{eq:I.7.3.6}{I.7.3.6}
\EGAReferenceAlias{I.7.3.7-fr}{I.7.3.7}
\EGAReferenceAlias{I.7.3.7-es}{I.7.3.7}
\EGAReferenceAlias{I.7.3.7.fr}{I.7.3.7}
\EGAReferenceAlias{I.7.3.7.es}{I.7.3.7}
\EGAReferenceAlias{I.7.3.7.eq}{I.7.3.7}
\EGAReferenceAlias{eq:I.7.3.7}{I.7.3.7}
\EGAReferenceAlias{I.7.3.8-fr}{I.7.3.8}
\EGAReferenceAlias{I.7.3.8-es}{I.7.3.8}
\EGAReferenceAlias{I.7.3.8.fr}{I.7.3.8}
\EGAReferenceAlias{I.7.3.8.es}{I.7.3.8}
\EGAReferenceAlias{I.7.3.8.eq}{I.7.3.8}
\EGAReferenceAlias{eq:I.7.3.8}{I.7.3.8}
\EGAReferenceAlias{I.7.3.8.1}{I.7.3.8.1-fr}
\EGAReferenceAlias{I.7.3.8.1-es}{I.7.3.8.1-fr}
\EGAReferenceAlias{I.7.3.8.1.fr}{I.7.3.8.1-fr}
\EGAReferenceAlias{I.7.3.8.1.es}{I.7.3.8.1-fr}
\EGAReferenceAlias{I.7.3.8.1.eq}{I.7.3.8.1-fr}
\EGAReferenceAlias{eq:I.7.3.8.1}{I.7.3.8.1-fr}
\EGAReferenceAlias{I.7.4-fr}{subsection:I.7.4}
\EGAReferenceAlias{I.7.4-es}{subsection:I.7.4}
\EGAReferenceAlias{I.7.4.fr}{subsection:I.7.4}
\EGAReferenceAlias{I.7.4.es}{subsection:I.7.4}
\EGAReferenceAlias{I.7.4.eq}{subsection:I.7.4}
\EGAReferenceAlias{eq:I.7.4}{subsection:I.7.4}
\EGAReferenceAlias{subsection:I.7.4-fr}{subsection:I.7.4}
\EGAReferenceAlias{subsection:I.7.4-es}{subsection:I.7.4}
\EGAReferenceAlias{subsection:I.7.4.fr}{subsection:I.7.4}
\EGAReferenceAlias{subsection:I.7.4.es}{subsection:I.7.4}
\EGAReferenceAlias{subsection:I.7.4.eq}{subsection:I.7.4}
\EGAReferenceAlias{eq:subsection:I.7.4}{subsection:I.7.4}
\EGAReferenceAlias{I.7.4.1-fr}{I.7.4.1}
\EGAReferenceAlias{I.7.4.1-es}{I.7.4.1}
\EGAReferenceAlias{I.7.4.1.fr}{I.7.4.1}
\EGAReferenceAlias{I.7.4.1.es}{I.7.4.1}
\EGAReferenceAlias{I.7.4.1.eq}{I.7.4.1}
\EGAReferenceAlias{eq:I.7.4.1}{I.7.4.1}
\EGAReferenceAlias{I.7.4.2-fr}{I.7.4.2}
\EGAReferenceAlias{I.7.4.2-es}{I.7.4.2}
\EGAReferenceAlias{I.7.4.2.fr}{I.7.4.2}
\EGAReferenceAlias{I.7.4.2.es}{I.7.4.2}
\EGAReferenceAlias{I.7.4.2.eq}{I.7.4.2}
\EGAReferenceAlias{eq:I.7.4.2}{I.7.4.2}
\EGAReferenceAlias{I.7.4.3-fr}{I.7.4.3}
\EGAReferenceAlias{I.7.4.3-es}{I.7.4.3}
\EGAReferenceAlias{I.7.4.3.fr}{I.7.4.3}
\EGAReferenceAlias{I.7.4.3.es}{I.7.4.3}
\EGAReferenceAlias{I.7.4.3.eq}{I.7.4.3}
\EGAReferenceAlias{eq:I.7.4.3}{I.7.4.3}
\EGAReferenceAlias{I.7.4.4-fr}{I.7.4.4}
\EGAReferenceAlias{I.7.4.4-es}{I.7.4.4}
\EGAReferenceAlias{I.7.4.4.fr}{I.7.4.4}
\EGAReferenceAlias{I.7.4.4.es}{I.7.4.4}
\EGAReferenceAlias{I.7.4.4.eq}{I.7.4.4}
\EGAReferenceAlias{eq:I.7.4.4}{I.7.4.4}
\EGAReferenceAlias{I.7.4.5-fr}{I.7.4.5}
\EGAReferenceAlias{I.7.4.5-es}{I.7.4.5}
\EGAReferenceAlias{I.7.4.5.fr}{I.7.4.5}
\EGAReferenceAlias{I.7.4.5.es}{I.7.4.5}
\EGAReferenceAlias{I.7.4.5.eq}{I.7.4.5}
\EGAReferenceAlias{eq:I.7.4.5}{I.7.4.5}
\EGAReferenceAlias{I.7.4.6-fr}{I.7.4.6}
\EGAReferenceAlias{I.7.4.6-es}{I.7.4.6}
\EGAReferenceAlias{I.7.4.6.fr}{I.7.4.6}
\EGAReferenceAlias{I.7.4.6.es}{I.7.4.6}
\EGAReferenceAlias{I.7.4.6.eq}{I.7.4.6}
\EGAReferenceAlias{eq:I.7.4.6}{I.7.4.6}
\EGAReferenceAlias{I.7.4.7-fr}{I.7.4.7}
\EGAReferenceAlias{I.7.4.7-es}{I.7.4.7}
\EGAReferenceAlias{I.7.4.7.fr}{I.7.4.7}
\EGAReferenceAlias{I.7.4.7.es}{I.7.4.7}
\EGAReferenceAlias{I.7.4.7.eq}{I.7.4.7}
\EGAReferenceAlias{eq:I.7.4.7}{I.7.4.7}
\EGAReferenceAlias{I.8}{section:I.8}
\EGAReferenceAlias{I.8-fr}{section:I.8}
\EGAReferenceAlias{I.8-es}{section:I.8}
\EGAReferenceAlias{I.8.fr}{section:I.8}
\EGAReferenceAlias{I.8.es}{section:I.8}
\EGAReferenceAlias{I.8.eq}{section:I.8}
\EGAReferenceAlias{eq:I.8}{section:I.8}
\EGAReferenceAlias{section:I.8-fr}{section:I.8}
\EGAReferenceAlias{section:I.8-es}{section:I.8}
\EGAReferenceAlias{section:I.8.fr}{section:I.8}
\EGAReferenceAlias{section:I.8.es}{section:I.8}
\EGAReferenceAlias{section:I.8.eq}{section:I.8}
\EGAReferenceAlias{eq:section:I.8}{section:I.8}
\EGAReferenceAlias{I.8.1}{subsection:I.8.1}
\EGAReferenceAlias{I.8.1-fr}{subsection:I.8.1}
\EGAReferenceAlias{I.8.1-es}{subsection:I.8.1}
\EGAReferenceAlias{I.8.1.fr}{subsection:I.8.1}
\EGAReferenceAlias{I.8.1.es}{subsection:I.8.1}
\EGAReferenceAlias{I.8.1.eq}{subsection:I.8.1}
\EGAReferenceAlias{eq:I.8.1}{subsection:I.8.1}
\EGAReferenceAlias{subsection:I.8.1-fr}{subsection:I.8.1}
\EGAReferenceAlias{subsection:I.8.1-es}{subsection:I.8.1}
\EGAReferenceAlias{subsection:I.8.1.fr}{subsection:I.8.1}
\EGAReferenceAlias{subsection:I.8.1.es}{subsection:I.8.1}
\EGAReferenceAlias{subsection:I.8.1.eq}{subsection:I.8.1}
\EGAReferenceAlias{eq:subsection:I.8.1}{subsection:I.8.1}
\EGAReferenceAlias{I.8.1.1-fr}{I.8.1.1}
\EGAReferenceAlias{I.8.1.1-es}{I.8.1.1}
\EGAReferenceAlias{I.8.1.1.fr}{I.8.1.1}
\EGAReferenceAlias{I.8.1.1.es}{I.8.1.1}
\EGAReferenceAlias{I.8.1.1.eq}{I.8.1.1}
\EGAReferenceAlias{eq:I.8.1.1}{I.8.1.1}
\EGAReferenceAlias{I.8.1.2-fr}{I.8.1.2}
\EGAReferenceAlias{I.8.1.2-es}{I.8.1.2}
\EGAReferenceAlias{I.8.1.2.fr}{I.8.1.2}
\EGAReferenceAlias{I.8.1.2.es}{I.8.1.2}
\EGAReferenceAlias{I.8.1.2.eq}{I.8.1.2}
\EGAReferenceAlias{eq:I.8.1.2}{I.8.1.2}
\EGAReferenceAlias{I.8.1.3-fr}{I.8.1.3}
\EGAReferenceAlias{I.8.1.3-es}{I.8.1.3}
\EGAReferenceAlias{I.8.1.3.fr}{I.8.1.3}
\EGAReferenceAlias{I.8.1.3.es}{I.8.1.3}
\EGAReferenceAlias{I.8.1.3.eq}{I.8.1.3}
\EGAReferenceAlias{eq:I.8.1.3}{I.8.1.3}
\EGAReferenceAlias{I.8.1.4-fr}{I.8.1.4}
\EGAReferenceAlias{I.8.1.4-es}{I.8.1.4}
\EGAReferenceAlias{I.8.1.4.fr}{I.8.1.4}
\EGAReferenceAlias{I.8.1.4.es}{I.8.1.4}
\EGAReferenceAlias{I.8.1.4.eq}{I.8.1.4}
\EGAReferenceAlias{eq:I.8.1.4}{I.8.1.4}
\EGAReferenceAlias{I.8.1.5-fr}{I.8.1.5}
\EGAReferenceAlias{I.8.1.5-es}{I.8.1.5}
\EGAReferenceAlias{I.8.1.5.fr}{I.8.1.5}
\EGAReferenceAlias{I.8.1.5.es}{I.8.1.5}
\EGAReferenceAlias{I.8.1.5.eq}{I.8.1.5}
\EGAReferenceAlias{eq:I.8.1.5}{I.8.1.5}
\EGAReferenceAlias{I.8.2}{subsection:I.8.2}
\EGAReferenceAlias{I.8.2-fr}{subsection:I.8.2}
\EGAReferenceAlias{I.8.2-es}{subsection:I.8.2}
\EGAReferenceAlias{I.8.2.fr}{subsection:I.8.2}
\EGAReferenceAlias{I.8.2.es}{subsection:I.8.2}
\EGAReferenceAlias{I.8.2.eq}{subsection:I.8.2}
\EGAReferenceAlias{eq:I.8.2}{subsection:I.8.2}
\EGAReferenceAlias{subsection:I.8.2-fr}{subsection:I.8.2}
\EGAReferenceAlias{subsection:I.8.2-es}{subsection:I.8.2}
\EGAReferenceAlias{subsection:I.8.2.fr}{subsection:I.8.2}
\EGAReferenceAlias{subsection:I.8.2.es}{subsection:I.8.2}
\EGAReferenceAlias{subsection:I.8.2.eq}{subsection:I.8.2}
\EGAReferenceAlias{eq:subsection:I.8.2}{subsection:I.8.2}
\EGAReferenceAlias{I.8.2.1-fr}{I.8.2.1}
\EGAReferenceAlias{I.8.2.1-es}{I.8.2.1}
\EGAReferenceAlias{I.8.2.1.fr}{I.8.2.1}
\EGAReferenceAlias{I.8.2.1.es}{I.8.2.1}
\EGAReferenceAlias{I.8.2.1.eq}{I.8.2.1}
\EGAReferenceAlias{eq:I.8.2.1}{I.8.2.1}
\EGAReferenceAlias{I.8.2.1.1}{I.8.2.1.1-fr}
\EGAReferenceAlias{I.8.2.1.1-es}{I.8.2.1.1-fr}
\EGAReferenceAlias{I.8.2.1.1.fr}{I.8.2.1.1-fr}
\EGAReferenceAlias{I.8.2.1.1.es}{I.8.2.1.1-fr}
\EGAReferenceAlias{I.8.2.1.1.eq}{I.8.2.1.1-fr}
\EGAReferenceAlias{eq:I.8.2.1.1}{I.8.2.1.1-fr}
\EGAReferenceAlias{I.8.2.2-fr}{I.8.2.2}
\EGAReferenceAlias{I.8.2.2-es}{I.8.2.2}
\EGAReferenceAlias{I.8.2.2.fr}{I.8.2.2}
\EGAReferenceAlias{I.8.2.2.es}{I.8.2.2}
\EGAReferenceAlias{I.8.2.2.eq}{I.8.2.2}
\EGAReferenceAlias{eq:I.8.2.2}{I.8.2.2}
\EGAReferenceAlias{I.8.2.3-fr}{I.8.2.3}
\EGAReferenceAlias{I.8.2.3-es}{I.8.2.3}
\EGAReferenceAlias{I.8.2.3.fr}{I.8.2.3}
\EGAReferenceAlias{I.8.2.3.es}{I.8.2.3}
\EGAReferenceAlias{I.8.2.3.eq}{I.8.2.3}
\EGAReferenceAlias{eq:I.8.2.3}{I.8.2.3}
\EGAReferenceAlias{I.8.2.4-fr}{I.8.2.4}
\EGAReferenceAlias{I.8.2.4-es}{I.8.2.4}
\EGAReferenceAlias{I.8.2.4.fr}{I.8.2.4}
\EGAReferenceAlias{I.8.2.4.es}{I.8.2.4}
\EGAReferenceAlias{I.8.2.4.eq}{I.8.2.4}
\EGAReferenceAlias{eq:I.8.2.4}{I.8.2.4}
\EGAReferenceAlias{I.8.2.5-fr}{I.8.2.5}
\EGAReferenceAlias{I.8.2.5-es}{I.8.2.5}
\EGAReferenceAlias{I.8.2.5.fr}{I.8.2.5}
\EGAReferenceAlias{I.8.2.5.es}{I.8.2.5}
\EGAReferenceAlias{I.8.2.5.eq}{I.8.2.5}
\EGAReferenceAlias{eq:I.8.2.5}{I.8.2.5}
\EGAReferenceAlias{I.8.2.6-fr}{I.8.2.6}
\EGAReferenceAlias{I.8.2.6-es}{I.8.2.6}
\EGAReferenceAlias{I.8.2.6.fr}{I.8.2.6}
\EGAReferenceAlias{I.8.2.6.es}{I.8.2.6}
\EGAReferenceAlias{I.8.2.6.eq}{I.8.2.6}
\EGAReferenceAlias{eq:I.8.2.6}{I.8.2.6}
\EGAReferenceAlias{I.8.2.7-fr}{I.8.2.7}
\EGAReferenceAlias{I.8.2.7-es}{I.8.2.7}
\EGAReferenceAlias{I.8.2.7.fr}{I.8.2.7}
\EGAReferenceAlias{I.8.2.7.es}{I.8.2.7}
\EGAReferenceAlias{I.8.2.7.eq}{I.8.2.7}
\EGAReferenceAlias{eq:I.8.2.7}{I.8.2.7}
\EGAReferenceAlias{I.8.2.8-fr}{I.8.2.8}
\EGAReferenceAlias{I.8.2.8-es}{I.8.2.8}
\EGAReferenceAlias{I.8.2.8.fr}{I.8.2.8}
\EGAReferenceAlias{I.8.2.8.es}{I.8.2.8}
\EGAReferenceAlias{I.8.2.8.eq}{I.8.2.8}
\EGAReferenceAlias{eq:I.8.2.8}{I.8.2.8}
\EGAReferenceAlias{I.8.3}{subsection:I.8.3}
\EGAReferenceAlias{I.8.3-fr}{subsection:I.8.3}
\EGAReferenceAlias{I.8.3-es}{subsection:I.8.3}
\EGAReferenceAlias{I.8.3.fr}{subsection:I.8.3}
\EGAReferenceAlias{I.8.3.es}{subsection:I.8.3}
\EGAReferenceAlias{I.8.3.eq}{subsection:I.8.3}
\EGAReferenceAlias{eq:I.8.3}{subsection:I.8.3}
\EGAReferenceAlias{subsection:I.8.3-fr}{subsection:I.8.3}
\EGAReferenceAlias{subsection:I.8.3-es}{subsection:I.8.3}
\EGAReferenceAlias{subsection:I.8.3.fr}{subsection:I.8.3}
\EGAReferenceAlias{subsection:I.8.3.es}{subsection:I.8.3}
\EGAReferenceAlias{subsection:I.8.3.eq}{subsection:I.8.3}
\EGAReferenceAlias{eq:subsection:I.8.3}{subsection:I.8.3}
\EGAReferenceAlias{I.8.3.1-fr}{I.8.3.1}
\EGAReferenceAlias{I.8.3.1-es}{I.8.3.1}
\EGAReferenceAlias{I.8.3.1.fr}{I.8.3.1}
\EGAReferenceAlias{I.8.3.1.es}{I.8.3.1}
\EGAReferenceAlias{I.8.3.1.eq}{I.8.3.1}
\EGAReferenceAlias{eq:I.8.3.1}{I.8.3.1}
\EGAReferenceAlias{I.8.3.2-fr}{I.8.3.2}
\EGAReferenceAlias{I.8.3.2-es}{I.8.3.2}
\EGAReferenceAlias{I.8.3.2.fr}{I.8.3.2}
\EGAReferenceAlias{I.8.3.2.es}{I.8.3.2}
\EGAReferenceAlias{I.8.3.2.eq}{I.8.3.2}
\EGAReferenceAlias{eq:I.8.3.2}{I.8.3.2}
\EGAReferenceAlias{I.9}{section:I.9}
\EGAReferenceAlias{I.9-fr}{section:I.9}
\EGAReferenceAlias{I.9-es}{section:I.9}
\EGAReferenceAlias{I.9.fr}{section:I.9}
\EGAReferenceAlias{I.9.es}{section:I.9}
\EGAReferenceAlias{I.9.eq}{section:I.9}
\EGAReferenceAlias{eq:I.9}{section:I.9}
\EGAReferenceAlias{section:I.9-fr}{section:I.9}
\EGAReferenceAlias{section:I.9-es}{section:I.9}
\EGAReferenceAlias{section:I.9.fr}{section:I.9}
\EGAReferenceAlias{section:I.9.es}{section:I.9}
\EGAReferenceAlias{section:I.9.eq}{section:I.9}
\EGAReferenceAlias{eq:section:I.9}{section:I.9}
\EGAReferenceAlias{I.9.1}{subsection:I.9.1}
\EGAReferenceAlias{I.9.1-fr}{subsection:I.9.1}
\EGAReferenceAlias{I.9.1-es}{subsection:I.9.1}
\EGAReferenceAlias{I.9.1.fr}{subsection:I.9.1}
\EGAReferenceAlias{I.9.1.es}{subsection:I.9.1}
\EGAReferenceAlias{I.9.1.eq}{subsection:I.9.1}
\EGAReferenceAlias{eq:I.9.1}{subsection:I.9.1}
\EGAReferenceAlias{subsection:I.9.1-fr}{subsection:I.9.1}
\EGAReferenceAlias{subsection:I.9.1-es}{subsection:I.9.1}
\EGAReferenceAlias{subsection:I.9.1.fr}{subsection:I.9.1}
\EGAReferenceAlias{subsection:I.9.1.es}{subsection:I.9.1}
\EGAReferenceAlias{subsection:I.9.1.eq}{subsection:I.9.1}
\EGAReferenceAlias{eq:subsection:I.9.1}{subsection:I.9.1}
\EGAReferenceAlias{I.9.1.1-fr}{I.9.1.1}
\EGAReferenceAlias{I.9.1.1-es}{I.9.1.1}
\EGAReferenceAlias{I.9.1.1.fr}{I.9.1.1}
\EGAReferenceAlias{I.9.1.1.es}{I.9.1.1}
\EGAReferenceAlias{I.9.1.1.eq}{I.9.1.1}
\EGAReferenceAlias{eq:I.9.1.1}{I.9.1.1}
\EGAReferenceAlias{I.9.1.2-fr}{I.9.1.2}
\EGAReferenceAlias{I.9.1.2-es}{I.9.1.2}
\EGAReferenceAlias{I.9.1.2.fr}{I.9.1.2}
\EGAReferenceAlias{I.9.1.2.es}{I.9.1.2}
\EGAReferenceAlias{I.9.1.2.eq}{I.9.1.2}
\EGAReferenceAlias{eq:I.9.1.2}{I.9.1.2}
\EGAReferenceAlias{I.9.1.3-fr}{I.9.1.3}
\EGAReferenceAlias{I.9.1.3-es}{I.9.1.3}
\EGAReferenceAlias{I.9.1.3.fr}{I.9.1.3}
\EGAReferenceAlias{I.9.1.3.es}{I.9.1.3}
\EGAReferenceAlias{I.9.1.3.eq}{I.9.1.3}
\EGAReferenceAlias{eq:I.9.1.3}{I.9.1.3}
\EGAReferenceAlias{I.9.1.4-fr}{I.9.1.4}
\EGAReferenceAlias{I.9.1.4-es}{I.9.1.4}
\EGAReferenceAlias{I.9.1.4.fr}{I.9.1.4}
\EGAReferenceAlias{I.9.1.4.es}{I.9.1.4}
\EGAReferenceAlias{I.9.1.4.eq}{I.9.1.4}
\EGAReferenceAlias{eq:I.9.1.4}{I.9.1.4}
\EGAReferenceAlias{I.9.1.5-fr}{I.9.1.5}
\EGAReferenceAlias{I.9.1.5-es}{I.9.1.5}
\EGAReferenceAlias{I.9.1.5.fr}{I.9.1.5}
\EGAReferenceAlias{I.9.1.5.es}{I.9.1.5}
\EGAReferenceAlias{I.9.1.5.eq}{I.9.1.5}
\EGAReferenceAlias{eq:I.9.1.5}{I.9.1.5}
\EGAReferenceAlias{I.9.1.6-fr}{I.9.1.6}
\EGAReferenceAlias{I.9.1.6-es}{I.9.1.6}
\EGAReferenceAlias{I.9.1.6.fr}{I.9.1.6}
\EGAReferenceAlias{I.9.1.6.es}{I.9.1.6}
\EGAReferenceAlias{I.9.1.6.eq}{I.9.1.6}
\EGAReferenceAlias{eq:I.9.1.6}{I.9.1.6}
\EGAReferenceAlias{I.9.1.7-fr}{I.9.1.7}
\EGAReferenceAlias{I.9.1.7-es}{I.9.1.7}
\EGAReferenceAlias{I.9.1.7.fr}{I.9.1.7}
\EGAReferenceAlias{I.9.1.7.es}{I.9.1.7}
\EGAReferenceAlias{I.9.1.7.eq}{I.9.1.7}
\EGAReferenceAlias{eq:I.9.1.7}{I.9.1.7}
\EGAReferenceAlias{I.9.1.8-fr}{I.9.1.8}
\EGAReferenceAlias{I.9.1.8-es}{I.9.1.8}
\EGAReferenceAlias{I.9.1.8.fr}{I.9.1.8}
\EGAReferenceAlias{I.9.1.8.es}{I.9.1.8}
\EGAReferenceAlias{I.9.1.8.eq}{I.9.1.8}
\EGAReferenceAlias{eq:I.9.1.8}{I.9.1.8}
\EGAReferenceAlias{I.9.1.9-fr}{I.9.1.9}
\EGAReferenceAlias{I.9.1.9-es}{I.9.1.9}
\EGAReferenceAlias{I.9.1.9.fr}{I.9.1.9}
\EGAReferenceAlias{I.9.1.9.es}{I.9.1.9}
\EGAReferenceAlias{I.9.1.9.eq}{I.9.1.9}
\EGAReferenceAlias{eq:I.9.1.9}{I.9.1.9}
\EGAReferenceAlias{I.9.1.9.1}{I.9.1.9.1-fr}
\EGAReferenceAlias{I.9.1.9.1-es}{I.9.1.9.1-fr}
\EGAReferenceAlias{I.9.1.9.1.fr}{I.9.1.9.1-fr}
\EGAReferenceAlias{I.9.1.9.1.es}{I.9.1.9.1-fr}
\EGAReferenceAlias{I.9.1.9.1.eq}{I.9.1.9.1-fr}
\EGAReferenceAlias{eq:I.9.1.9.1}{I.9.1.9.1-fr}
\EGAReferenceAlias{I.9.1.10-fr}{I.9.1.10}
\EGAReferenceAlias{I.9.1.10-es}{I.9.1.10}
\EGAReferenceAlias{I.9.1.10.fr}{I.9.1.10}
\EGAReferenceAlias{I.9.1.10.es}{I.9.1.10}
\EGAReferenceAlias{I.9.1.10.eq}{I.9.1.10}
\EGAReferenceAlias{eq:I.9.1.10}{I.9.1.10}
\EGAReferenceAlias{I.9.1.11-fr}{I.9.1.11}
\EGAReferenceAlias{I.9.1.11-es}{I.9.1.11}
\EGAReferenceAlias{I.9.1.11.fr}{I.9.1.11}
\EGAReferenceAlias{I.9.1.11.es}{I.9.1.11}
\EGAReferenceAlias{I.9.1.11.eq}{I.9.1.11}
\EGAReferenceAlias{eq:I.9.1.11}{I.9.1.11}
\EGAReferenceAlias{I.9.1.12-fr}{I.9.1.12}
\EGAReferenceAlias{I.9.1.12-es}{I.9.1.12}
\EGAReferenceAlias{I.9.1.12.fr}{I.9.1.12}
\EGAReferenceAlias{I.9.1.12.es}{I.9.1.12}
\EGAReferenceAlias{I.9.1.12.eq}{I.9.1.12}
\EGAReferenceAlias{eq:I.9.1.12}{I.9.1.12}
\EGAReferenceAlias{I.9.1.13-fr}{I.9.1.13}
\EGAReferenceAlias{I.9.1.13-es}{I.9.1.13}
\EGAReferenceAlias{I.9.1.13.fr}{I.9.1.13}
\EGAReferenceAlias{I.9.1.13.es}{I.9.1.13}
\EGAReferenceAlias{I.9.1.13.eq}{I.9.1.13}
\EGAReferenceAlias{eq:I.9.1.13}{I.9.1.13}
\EGAReferenceAlias{I.9.1.13.1}{I.9.1.13.1-fr}
\EGAReferenceAlias{I.9.1.13.1-es}{I.9.1.13.1-fr}
\EGAReferenceAlias{I.9.1.13.1.fr}{I.9.1.13.1-fr}
\EGAReferenceAlias{I.9.1.13.1.es}{I.9.1.13.1-fr}
\EGAReferenceAlias{I.9.1.13.1.eq}{I.9.1.13.1-fr}
\EGAReferenceAlias{eq:I.9.1.13.1}{I.9.1.13.1-fr}
\EGAReferenceAlias{I.9.2}{subsection:I.9.2}
\EGAReferenceAlias{I.9.2-fr}{subsection:I.9.2}
\EGAReferenceAlias{I.9.2-es}{subsection:I.9.2}
\EGAReferenceAlias{I.9.2.fr}{subsection:I.9.2}
\EGAReferenceAlias{I.9.2.es}{subsection:I.9.2}
\EGAReferenceAlias{I.9.2.eq}{subsection:I.9.2}
\EGAReferenceAlias{eq:I.9.2}{subsection:I.9.2}
\EGAReferenceAlias{subsection:I.9.2-fr}{subsection:I.9.2}
\EGAReferenceAlias{subsection:I.9.2-es}{subsection:I.9.2}
\EGAReferenceAlias{subsection:I.9.2.fr}{subsection:I.9.2}
\EGAReferenceAlias{subsection:I.9.2.es}{subsection:I.9.2}
\EGAReferenceAlias{subsection:I.9.2.eq}{subsection:I.9.2}
\EGAReferenceAlias{eq:subsection:I.9.2}{subsection:I.9.2}
\EGAReferenceAlias{I.9.2.1-fr}{I.9.2.1}
\EGAReferenceAlias{I.9.2.1-es}{I.9.2.1}
\EGAReferenceAlias{I.9.2.1.fr}{I.9.2.1}
\EGAReferenceAlias{I.9.2.1.es}{I.9.2.1}
\EGAReferenceAlias{I.9.2.1.eq}{I.9.2.1}
\EGAReferenceAlias{eq:I.9.2.1}{I.9.2.1}
\EGAReferenceAlias{I.9.2.2-fr}{I.9.2.2}
\EGAReferenceAlias{I.9.2.2-es}{I.9.2.2}
\EGAReferenceAlias{I.9.2.2.fr}{I.9.2.2}
\EGAReferenceAlias{I.9.2.2.es}{I.9.2.2}
\EGAReferenceAlias{I.9.2.2.eq}{I.9.2.2}
\EGAReferenceAlias{eq:I.9.2.2}{I.9.2.2}
\EGAReferenceAlias{I.9.3}{subsection:I.9.3}
\EGAReferenceAlias{I.9.3-fr}{subsection:I.9.3}
\EGAReferenceAlias{I.9.3-es}{subsection:I.9.3}
\EGAReferenceAlias{I.9.3.fr}{subsection:I.9.3}
\EGAReferenceAlias{I.9.3.es}{subsection:I.9.3}
\EGAReferenceAlias{I.9.3.eq}{subsection:I.9.3}
\EGAReferenceAlias{eq:I.9.3}{subsection:I.9.3}
\EGAReferenceAlias{subsection:I.9.3-fr}{subsection:I.9.3}
\EGAReferenceAlias{subsection:I.9.3-es}{subsection:I.9.3}
\EGAReferenceAlias{subsection:I.9.3.fr}{subsection:I.9.3}
\EGAReferenceAlias{subsection:I.9.3.es}{subsection:I.9.3}
\EGAReferenceAlias{subsection:I.9.3.eq}{subsection:I.9.3}
\EGAReferenceAlias{eq:subsection:I.9.3}{subsection:I.9.3}
\EGAReferenceAlias{I.9.3.1-fr}{I.9.3.1}
\EGAReferenceAlias{I.9.3.1-es}{I.9.3.1}
\EGAReferenceAlias{I.9.3.1.fr}{I.9.3.1}
\EGAReferenceAlias{I.9.3.1.es}{I.9.3.1}
\EGAReferenceAlias{I.9.3.1.eq}{I.9.3.1}
\EGAReferenceAlias{eq:I.9.3.1}{I.9.3.1}
\EGAReferenceAlias{I.9.3.2-fr}{I.9.3.2}
\EGAReferenceAlias{I.9.3.2-es}{I.9.3.2}
\EGAReferenceAlias{I.9.3.2.fr}{I.9.3.2}
\EGAReferenceAlias{I.9.3.2.es}{I.9.3.2}
\EGAReferenceAlias{I.9.3.2.eq}{I.9.3.2}
\EGAReferenceAlias{eq:I.9.3.2}{I.9.3.2}
\EGAReferenceAlias{I.9.3.3-fr}{I.9.3.3}
\EGAReferenceAlias{I.9.3.3-es}{I.9.3.3}
\EGAReferenceAlias{I.9.3.3.fr}{I.9.3.3}
\EGAReferenceAlias{I.9.3.3.es}{I.9.3.3}
\EGAReferenceAlias{I.9.3.3.eq}{I.9.3.3}
\EGAReferenceAlias{eq:I.9.3.3}{I.9.3.3}
\EGAReferenceAlias{I.9.3.4-fr}{I.9.3.4}
\EGAReferenceAlias{I.9.3.4-es}{I.9.3.4}
\EGAReferenceAlias{I.9.3.4.fr}{I.9.3.4}
\EGAReferenceAlias{I.9.3.4.es}{I.9.3.4}
\EGAReferenceAlias{I.9.3.4.eq}{I.9.3.4}
\EGAReferenceAlias{eq:I.9.3.4}{I.9.3.4}
\EGAReferenceAlias{I.9.3.5-fr}{I.9.3.5}
\EGAReferenceAlias{I.9.3.5-es}{I.9.3.5}
\EGAReferenceAlias{I.9.3.5.fr}{I.9.3.5}
\EGAReferenceAlias{I.9.3.5.es}{I.9.3.5}
\EGAReferenceAlias{I.9.3.5.eq}{I.9.3.5}
\EGAReferenceAlias{eq:I.9.3.5}{I.9.3.5}
\EGAReferenceAlias{I.9.4}{subsection:I.9.4}
\EGAReferenceAlias{I.9.4-fr}{subsection:I.9.4}
\EGAReferenceAlias{I.9.4-es}{subsection:I.9.4}
\EGAReferenceAlias{I.9.4.fr}{subsection:I.9.4}
\EGAReferenceAlias{I.9.4.es}{subsection:I.9.4}
\EGAReferenceAlias{I.9.4.eq}{subsection:I.9.4}
\EGAReferenceAlias{eq:I.9.4}{subsection:I.9.4}
\EGAReferenceAlias{subsection:I.9.4-fr}{subsection:I.9.4}
\EGAReferenceAlias{subsection:I.9.4-es}{subsection:I.9.4}
\EGAReferenceAlias{subsection:I.9.4.fr}{subsection:I.9.4}
\EGAReferenceAlias{subsection:I.9.4.es}{subsection:I.9.4}
\EGAReferenceAlias{subsection:I.9.4.eq}{subsection:I.9.4}
\EGAReferenceAlias{eq:subsection:I.9.4}{subsection:I.9.4}
\EGAReferenceAlias{I.9.4.1-fr}{I.9.4.1}
\EGAReferenceAlias{I.9.4.1-es}{I.9.4.1}
\EGAReferenceAlias{I.9.4.1.fr}{I.9.4.1}
\EGAReferenceAlias{I.9.4.1.es}{I.9.4.1}
\EGAReferenceAlias{I.9.4.1.eq}{I.9.4.1}
\EGAReferenceAlias{eq:I.9.4.1}{I.9.4.1}
\EGAReferenceAlias{I.9.4.2-fr}{I.9.4.2}
\EGAReferenceAlias{I.9.4.2-es}{I.9.4.2}
\EGAReferenceAlias{I.9.4.2.fr}{I.9.4.2}
\EGAReferenceAlias{I.9.4.2.es}{I.9.4.2}
\EGAReferenceAlias{I.9.4.2.eq}{I.9.4.2}
\EGAReferenceAlias{eq:I.9.4.2}{I.9.4.2}
\EGAReferenceAlias{I.9.4.3-fr}{I.9.4.3}
\EGAReferenceAlias{I.9.4.3-es}{I.9.4.3}
\EGAReferenceAlias{I.9.4.3.fr}{I.9.4.3}
\EGAReferenceAlias{I.9.4.3.es}{I.9.4.3}
\EGAReferenceAlias{I.9.4.3.eq}{I.9.4.3}
\EGAReferenceAlias{eq:I.9.4.3}{I.9.4.3}
\EGAReferenceAlias{I.9.4.4-fr}{I.9.4.4}
\EGAReferenceAlias{I.9.4.4-es}{I.9.4.4}
\EGAReferenceAlias{I.9.4.4.fr}{I.9.4.4}
\EGAReferenceAlias{I.9.4.4.es}{I.9.4.4}
\EGAReferenceAlias{I.9.4.4.eq}{I.9.4.4}
\EGAReferenceAlias{eq:I.9.4.4}{I.9.4.4}
\EGAReferenceAlias{I.9.4.5-fr}{I.9.4.5}
\EGAReferenceAlias{I.9.4.5-es}{I.9.4.5}
\EGAReferenceAlias{I.9.4.5.fr}{I.9.4.5}
\EGAReferenceAlias{I.9.4.5.es}{I.9.4.5}
\EGAReferenceAlias{I.9.4.5.eq}{I.9.4.5}
\EGAReferenceAlias{eq:I.9.4.5}{I.9.4.5}
\EGAReferenceAlias{I.9.4.6-fr}{I.9.4.6}
\EGAReferenceAlias{I.9.4.6-es}{I.9.4.6}
\EGAReferenceAlias{I.9.4.6.fr}{I.9.4.6}
\EGAReferenceAlias{I.9.4.6.es}{I.9.4.6}
\EGAReferenceAlias{I.9.4.6.eq}{I.9.4.6}
\EGAReferenceAlias{eq:I.9.4.6}{I.9.4.6}
\EGAReferenceAlias{I.9.4.7-fr}{I.9.4.7}
\EGAReferenceAlias{I.9.4.7-es}{I.9.4.7}
\EGAReferenceAlias{I.9.4.7.fr}{I.9.4.7}
\EGAReferenceAlias{I.9.4.7.es}{I.9.4.7}
\EGAReferenceAlias{I.9.4.7.eq}{I.9.4.7}
\EGAReferenceAlias{eq:I.9.4.7}{I.9.4.7}
\EGAReferenceAlias{I.9.4.8-fr}{I.9.4.8}
\EGAReferenceAlias{I.9.4.8-es}{I.9.4.8}
\EGAReferenceAlias{I.9.4.8.fr}{I.9.4.8}
\EGAReferenceAlias{I.9.4.8.es}{I.9.4.8}
\EGAReferenceAlias{I.9.4.8.eq}{I.9.4.8}
\EGAReferenceAlias{eq:I.9.4.8}{I.9.4.8}
\EGAReferenceAlias{I.9.4.9-fr}{I.9.4.9}
\EGAReferenceAlias{I.9.4.9-es}{I.9.4.9}
\EGAReferenceAlias{I.9.4.9.fr}{I.9.4.9}
\EGAReferenceAlias{I.9.4.9.es}{I.9.4.9}
\EGAReferenceAlias{I.9.4.9.eq}{I.9.4.9}
\EGAReferenceAlias{eq:I.9.4.9}{I.9.4.9}
\EGAReferenceAlias{I.9.4.10-fr}{I.9.4.10}
\EGAReferenceAlias{I.9.4.10-es}{I.9.4.10}
\EGAReferenceAlias{I.9.4.10.fr}{I.9.4.10}
\EGAReferenceAlias{I.9.4.10.es}{I.9.4.10}
\EGAReferenceAlias{I.9.4.10.eq}{I.9.4.10}
\EGAReferenceAlias{eq:I.9.4.10}{I.9.4.10}
\EGAReferenceAlias{I.9.5}{subsection:I.9.5}
\EGAReferenceAlias{I.9.5-fr}{subsection:I.9.5}
\EGAReferenceAlias{I.9.5-es}{subsection:I.9.5}
\EGAReferenceAlias{I.9.5.fr}{subsection:I.9.5}
\EGAReferenceAlias{I.9.5.es}{subsection:I.9.5}
\EGAReferenceAlias{I.9.5.eq}{subsection:I.9.5}
\EGAReferenceAlias{eq:I.9.5}{subsection:I.9.5}
\EGAReferenceAlias{subsection:I.9.5-fr}{subsection:I.9.5}
\EGAReferenceAlias{subsection:I.9.5-es}{subsection:I.9.5}
\EGAReferenceAlias{subsection:I.9.5.fr}{subsection:I.9.5}
\EGAReferenceAlias{subsection:I.9.5.es}{subsection:I.9.5}
\EGAReferenceAlias{subsection:I.9.5.eq}{subsection:I.9.5}
\EGAReferenceAlias{eq:subsection:I.9.5}{subsection:I.9.5}
\EGAReferenceAlias{I.9.5.1-fr}{I.9.5.1}
\EGAReferenceAlias{I.9.5.1-es}{I.9.5.1}
\EGAReferenceAlias{I.9.5.1.fr}{I.9.5.1}
\EGAReferenceAlias{I.9.5.1.es}{I.9.5.1}
\EGAReferenceAlias{I.9.5.1.eq}{I.9.5.1}
\EGAReferenceAlias{eq:I.9.5.1}{I.9.5.1}
\EGAReferenceAlias{I.9.5.2-fr}{I.9.5.2}
\EGAReferenceAlias{I.9.5.2-es}{I.9.5.2}
\EGAReferenceAlias{I.9.5.2.fr}{I.9.5.2}
\EGAReferenceAlias{I.9.5.2.es}{I.9.5.2}
\EGAReferenceAlias{I.9.5.2.eq}{I.9.5.2}
\EGAReferenceAlias{eq:I.9.5.2}{I.9.5.2}
\EGAReferenceAlias{I.9.5.3-fr}{I.9.5.3}
\EGAReferenceAlias{I.9.5.3-es}{I.9.5.3}
\EGAReferenceAlias{I.9.5.3.fr}{I.9.5.3}
\EGAReferenceAlias{I.9.5.3.es}{I.9.5.3}
\EGAReferenceAlias{I.9.5.3.eq}{I.9.5.3}
\EGAReferenceAlias{eq:I.9.5.3}{I.9.5.3}
\EGAReferenceAlias{I.9.5.4-fr}{I.9.5.4}
\EGAReferenceAlias{I.9.5.4-es}{I.9.5.4}
\EGAReferenceAlias{I.9.5.4.fr}{I.9.5.4}
\EGAReferenceAlias{I.9.5.4.es}{I.9.5.4}
\EGAReferenceAlias{I.9.5.4.eq}{I.9.5.4}
\EGAReferenceAlias{eq:I.9.5.4}{I.9.5.4}
\EGAReferenceAlias{I.9.5.5-fr}{I.9.5.5}
\EGAReferenceAlias{I.9.5.5-es}{I.9.5.5}
\EGAReferenceAlias{I.9.5.5.fr}{I.9.5.5}
\EGAReferenceAlias{I.9.5.5.es}{I.9.5.5}
\EGAReferenceAlias{I.9.5.5.eq}{I.9.5.5}
\EGAReferenceAlias{eq:I.9.5.5}{I.9.5.5}
\EGAReferenceAlias{I.9.5.6-fr}{I.9.5.6}
\EGAReferenceAlias{I.9.5.6-es}{I.9.5.6}
\EGAReferenceAlias{I.9.5.6.fr}{I.9.5.6}
\EGAReferenceAlias{I.9.5.6.es}{I.9.5.6}
\EGAReferenceAlias{I.9.5.6.eq}{I.9.5.6}
\EGAReferenceAlias{eq:I.9.5.6}{I.9.5.6}
\EGAReferenceAlias{I.9.5.7-fr}{I.9.5.7}
\EGAReferenceAlias{I.9.5.7-es}{I.9.5.7}
\EGAReferenceAlias{I.9.5.7.fr}{I.9.5.7}
\EGAReferenceAlias{I.9.5.7.es}{I.9.5.7}
\EGAReferenceAlias{I.9.5.7.eq}{I.9.5.7}
\EGAReferenceAlias{eq:I.9.5.7}{I.9.5.7}
\EGAReferenceAlias{I.9.5.8-fr}{I.9.5.8}
\EGAReferenceAlias{I.9.5.8-es}{I.9.5.8}
\EGAReferenceAlias{I.9.5.8.fr}{I.9.5.8}
\EGAReferenceAlias{I.9.5.8.es}{I.9.5.8}
\EGAReferenceAlias{I.9.5.8.eq}{I.9.5.8}
\EGAReferenceAlias{eq:I.9.5.8}{I.9.5.8}
\EGAReferenceAlias{I.9.5.9-fr}{I.9.5.9}
\EGAReferenceAlias{I.9.5.9-es}{I.9.5.9}
\EGAReferenceAlias{I.9.5.9.fr}{I.9.5.9}
\EGAReferenceAlias{I.9.5.9.es}{I.9.5.9}
\EGAReferenceAlias{I.9.5.9.eq}{I.9.5.9}
\EGAReferenceAlias{eq:I.9.5.9}{I.9.5.9}
\EGAReferenceAlias{I.9.5.10-fr}{I.9.5.10}
\EGAReferenceAlias{I.9.5.10-es}{I.9.5.10}
\EGAReferenceAlias{I.9.5.10.fr}{I.9.5.10}
\EGAReferenceAlias{I.9.5.10.es}{I.9.5.10}
\EGAReferenceAlias{I.9.5.10.eq}{I.9.5.10}
\EGAReferenceAlias{eq:I.9.5.10}{I.9.5.10}
\EGAReferenceAlias{I.9.5.11-fr}{I.9.5.11}
\EGAReferenceAlias{I.9.5.11-es}{I.9.5.11}
\EGAReferenceAlias{I.9.5.11.fr}{I.9.5.11}
\EGAReferenceAlias{I.9.5.11.es}{I.9.5.11}
\EGAReferenceAlias{I.9.5.11.eq}{I.9.5.11}
\EGAReferenceAlias{eq:I.9.5.11}{I.9.5.11}
\EGAReferenceAlias{I.9.6-fr}{subsection:I.9.6}
\EGAReferenceAlias{I.9.6-es}{subsection:I.9.6}
\EGAReferenceAlias{I.9.6.fr}{subsection:I.9.6}
\EGAReferenceAlias{I.9.6.es}{subsection:I.9.6}
\EGAReferenceAlias{I.9.6.eq}{subsection:I.9.6}
\EGAReferenceAlias{eq:I.9.6}{subsection:I.9.6}
\EGAReferenceAlias{subsection:I.9.6-fr}{subsection:I.9.6}
\EGAReferenceAlias{subsection:I.9.6-es}{subsection:I.9.6}
\EGAReferenceAlias{subsection:I.9.6.fr}{subsection:I.9.6}
\EGAReferenceAlias{subsection:I.9.6.es}{subsection:I.9.6}
\EGAReferenceAlias{subsection:I.9.6.eq}{subsection:I.9.6}
\EGAReferenceAlias{eq:subsection:I.9.6}{subsection:I.9.6}
\EGAReferenceAlias{I.9.6.1-fr}{I.9.6.1}
\EGAReferenceAlias{I.9.6.1-es}{I.9.6.1}
\EGAReferenceAlias{I.9.6.1.fr}{I.9.6.1}
\EGAReferenceAlias{I.9.6.1.es}{I.9.6.1}
\EGAReferenceAlias{I.9.6.1.eq}{I.9.6.1}
\EGAReferenceAlias{eq:I.9.6.1}{I.9.6.1}
\EGAReferenceAlias{I.9.6.2-fr}{I.9.6.2}
\EGAReferenceAlias{I.9.6.2-es}{I.9.6.2}
\EGAReferenceAlias{I.9.6.2.fr}{I.9.6.2}
\EGAReferenceAlias{I.9.6.2.es}{I.9.6.2}
\EGAReferenceAlias{I.9.6.2.eq}{I.9.6.2}
\EGAReferenceAlias{eq:I.9.6.2}{I.9.6.2}
\EGAReferenceAlias{I.9.6.3-fr}{I.9.6.3}
\EGAReferenceAlias{I.9.6.3-es}{I.9.6.3}
\EGAReferenceAlias{I.9.6.3.fr}{I.9.6.3}
\EGAReferenceAlias{I.9.6.3.es}{I.9.6.3}
\EGAReferenceAlias{I.9.6.3.eq}{I.9.6.3}
\EGAReferenceAlias{eq:I.9.6.3}{I.9.6.3}
\EGAReferenceAlias{I.9.6.4-fr}{I.9.6.4}
\EGAReferenceAlias{I.9.6.4-es}{I.9.6.4}
\EGAReferenceAlias{I.9.6.4.fr}{I.9.6.4}
\EGAReferenceAlias{I.9.6.4.es}{I.9.6.4}
\EGAReferenceAlias{I.9.6.4.eq}{I.9.6.4}
\EGAReferenceAlias{eq:I.9.6.4}{I.9.6.4}
\EGAReferenceAlias{I.9.6.5-fr}{I.9.6.5}
\EGAReferenceAlias{I.9.6.5-es}{I.9.6.5}
\EGAReferenceAlias{I.9.6.5.fr}{I.9.6.5}
\EGAReferenceAlias{I.9.6.5.es}{I.9.6.5}
\EGAReferenceAlias{I.9.6.5.eq}{I.9.6.5}
\EGAReferenceAlias{eq:I.9.6.5}{I.9.6.5}
\EGAReferenceAlias{I.9.6.6-fr}{I.9.6.6}
\EGAReferenceAlias{I.9.6.6-es}{I.9.6.6}
\EGAReferenceAlias{I.9.6.6.fr}{I.9.6.6}
\EGAReferenceAlias{I.9.6.6.es}{I.9.6.6}
\EGAReferenceAlias{I.9.6.6.eq}{I.9.6.6}
\EGAReferenceAlias{eq:I.9.6.6}{I.9.6.6}
\EGAReferenceAlias{I.10}{section:I.10}
\EGAReferenceAlias{I.10-fr}{section:I.10}
\EGAReferenceAlias{I.10-es}{section:I.10}
\EGAReferenceAlias{I.10.fr}{section:I.10}
\EGAReferenceAlias{I.10.es}{section:I.10}
\EGAReferenceAlias{I.10.eq}{section:I.10}
\EGAReferenceAlias{eq:I.10}{section:I.10}
\EGAReferenceAlias{section:I.10-fr}{section:I.10}
\EGAReferenceAlias{section:I.10-es}{section:I.10}
\EGAReferenceAlias{section:I.10.fr}{section:I.10}
\EGAReferenceAlias{section:I.10.es}{section:I.10}
\EGAReferenceAlias{section:I.10.eq}{section:I.10}
\EGAReferenceAlias{eq:section:I.10}{section:I.10}
\EGAReferenceAlias{I.10.1}{subsection:I.10.1}
\EGAReferenceAlias{I.10.1-fr}{subsection:I.10.1}
\EGAReferenceAlias{I.10.1-es}{subsection:I.10.1}
\EGAReferenceAlias{I.10.1.fr}{subsection:I.10.1}
\EGAReferenceAlias{I.10.1.es}{subsection:I.10.1}
\EGAReferenceAlias{I.10.1.eq}{subsection:I.10.1}
\EGAReferenceAlias{eq:I.10.1}{subsection:I.10.1}
\EGAReferenceAlias{subsection:I.10.1-fr}{subsection:I.10.1}
\EGAReferenceAlias{subsection:I.10.1-es}{subsection:I.10.1}
\EGAReferenceAlias{subsection:I.10.1.fr}{subsection:I.10.1}
\EGAReferenceAlias{subsection:I.10.1.es}{subsection:I.10.1}
\EGAReferenceAlias{subsection:I.10.1.eq}{subsection:I.10.1}
\EGAReferenceAlias{eq:subsection:I.10.1}{subsection:I.10.1}
\EGAReferenceAlias{I.10.1.1-fr}{I.10.1.1}
\EGAReferenceAlias{I.10.1.1-es}{I.10.1.1}
\EGAReferenceAlias{I.10.1.1.fr}{I.10.1.1}
\EGAReferenceAlias{I.10.1.1.es}{I.10.1.1}
\EGAReferenceAlias{I.10.1.1.eq}{I.10.1.1}
\EGAReferenceAlias{eq:I.10.1.1}{I.10.1.1}
\EGAReferenceAlias{I.10.1.2-fr}{I.10.1.2}
\EGAReferenceAlias{I.10.1.2-es}{I.10.1.2}
\EGAReferenceAlias{I.10.1.2.fr}{I.10.1.2}
\EGAReferenceAlias{I.10.1.2.es}{I.10.1.2}
\EGAReferenceAlias{I.10.1.2.eq}{I.10.1.2}
\EGAReferenceAlias{eq:I.10.1.2}{I.10.1.2}
\EGAReferenceAlias{I.10.1.3-fr}{I.10.1.3}
\EGAReferenceAlias{I.10.1.3-es}{I.10.1.3}
\EGAReferenceAlias{I.10.1.3.fr}{I.10.1.3}
\EGAReferenceAlias{I.10.1.3.es}{I.10.1.3}
\EGAReferenceAlias{I.10.1.3.eq}{I.10.1.3}
\EGAReferenceAlias{eq:I.10.1.3}{I.10.1.3}
\EGAReferenceAlias{I.10.1.4-fr}{I.10.1.4}
\EGAReferenceAlias{I.10.1.4-es}{I.10.1.4}
\EGAReferenceAlias{I.10.1.4.fr}{I.10.1.4}
\EGAReferenceAlias{I.10.1.4.es}{I.10.1.4}
\EGAReferenceAlias{I.10.1.4.eq}{I.10.1.4}
\EGAReferenceAlias{eq:I.10.1.4}{I.10.1.4}
\EGAReferenceAlias{I.10.1.5-fr}{I.10.1.5}
\EGAReferenceAlias{I.10.1.5-es}{I.10.1.5}
\EGAReferenceAlias{I.10.1.5.fr}{I.10.1.5}
\EGAReferenceAlias{I.10.1.5.es}{I.10.1.5}
\EGAReferenceAlias{I.10.1.5.eq}{I.10.1.5}
\EGAReferenceAlias{eq:I.10.1.5}{I.10.1.5}
\EGAReferenceAlias{I.10.1.6-fr}{I.10.1.6}
\EGAReferenceAlias{I.10.1.6-es}{I.10.1.6}
\EGAReferenceAlias{I.10.1.6.fr}{I.10.1.6}
\EGAReferenceAlias{I.10.1.6.es}{I.10.1.6}
\EGAReferenceAlias{I.10.1.6.eq}{I.10.1.6}
\EGAReferenceAlias{eq:I.10.1.6}{I.10.1.6}
\EGAReferenceAlias{I.10.2}{subsection:I.10.2}
\EGAReferenceAlias{I.10.2-fr}{subsection:I.10.2}
\EGAReferenceAlias{I.10.2-es}{subsection:I.10.2}
\EGAReferenceAlias{I.10.2.fr}{subsection:I.10.2}
\EGAReferenceAlias{I.10.2.es}{subsection:I.10.2}
\EGAReferenceAlias{I.10.2.eq}{subsection:I.10.2}
\EGAReferenceAlias{eq:I.10.2}{subsection:I.10.2}
\EGAReferenceAlias{subsection:I.10.2-fr}{subsection:I.10.2}
\EGAReferenceAlias{subsection:I.10.2-es}{subsection:I.10.2}
\EGAReferenceAlias{subsection:I.10.2.fr}{subsection:I.10.2}
\EGAReferenceAlias{subsection:I.10.2.es}{subsection:I.10.2}
\EGAReferenceAlias{subsection:I.10.2.eq}{subsection:I.10.2}
\EGAReferenceAlias{eq:subsection:I.10.2}{subsection:I.10.2}
\EGAReferenceAlias{I.10.2.1-fr}{I.10.2.1}
\EGAReferenceAlias{I.10.2.1-es}{I.10.2.1}
\EGAReferenceAlias{I.10.2.1.fr}{I.10.2.1}
\EGAReferenceAlias{I.10.2.1.es}{I.10.2.1}
\EGAReferenceAlias{I.10.2.1.eq}{I.10.2.1}
\EGAReferenceAlias{eq:I.10.2.1}{I.10.2.1}
\EGAReferenceAlias{I.10.2.2-fr}{I.10.2.2}
\EGAReferenceAlias{I.10.2.2-es}{I.10.2.2}
\EGAReferenceAlias{I.10.2.2.fr}{I.10.2.2}
\EGAReferenceAlias{I.10.2.2.es}{I.10.2.2}
\EGAReferenceAlias{I.10.2.2.eq}{I.10.2.2}
\EGAReferenceAlias{eq:I.10.2.2}{I.10.2.2}
\EGAReferenceAlias{I.10.2.3-fr}{I.10.2.3}
\EGAReferenceAlias{I.10.2.3-es}{I.10.2.3}
\EGAReferenceAlias{I.10.2.3.fr}{I.10.2.3}
\EGAReferenceAlias{I.10.2.3.es}{I.10.2.3}
\EGAReferenceAlias{I.10.2.3.eq}{I.10.2.3}
\EGAReferenceAlias{eq:I.10.2.3}{I.10.2.3}
\EGAReferenceAlias{I.10.3}{subsection:I.10.3}
\EGAReferenceAlias{I.10.3-fr}{subsection:I.10.3}
\EGAReferenceAlias{I.10.3-es}{subsection:I.10.3}
\EGAReferenceAlias{I.10.3.fr}{subsection:I.10.3}
\EGAReferenceAlias{I.10.3.es}{subsection:I.10.3}
\EGAReferenceAlias{I.10.3.eq}{subsection:I.10.3}
\EGAReferenceAlias{eq:I.10.3}{subsection:I.10.3}
\EGAReferenceAlias{subsection:I.10.3-fr}{subsection:I.10.3}
\EGAReferenceAlias{subsection:I.10.3-es}{subsection:I.10.3}
\EGAReferenceAlias{subsection:I.10.3.fr}{subsection:I.10.3}
\EGAReferenceAlias{subsection:I.10.3.es}{subsection:I.10.3}
\EGAReferenceAlias{subsection:I.10.3.eq}{subsection:I.10.3}
\EGAReferenceAlias{eq:subsection:I.10.3}{subsection:I.10.3}
\EGAReferenceAlias{I.10.3.1-fr}{I.10.3.1}
\EGAReferenceAlias{I.10.3.1-es}{I.10.3.1}
\EGAReferenceAlias{I.10.3.1.fr}{I.10.3.1}
\EGAReferenceAlias{I.10.3.1.es}{I.10.3.1}
\EGAReferenceAlias{I.10.3.1.eq}{I.10.3.1}
\EGAReferenceAlias{eq:I.10.3.1}{I.10.3.1}
\EGAReferenceAlias{I.10.3.1.1}{I.10.3.1.1-fr}
\EGAReferenceAlias{I.10.3.1.1-es}{I.10.3.1.1-fr}
\EGAReferenceAlias{I.10.3.1.1.fr}{I.10.3.1.1-fr}
\EGAReferenceAlias{I.10.3.1.1.es}{I.10.3.1.1-fr}
\EGAReferenceAlias{I.10.3.1.1.eq}{I.10.3.1.1-fr}
\EGAReferenceAlias{eq:I.10.3.1.1}{I.10.3.1.1-fr}
\EGAReferenceAlias{I.10.3.2-fr}{I.10.3.2}
\EGAReferenceAlias{I.10.3.2-es}{I.10.3.2}
\EGAReferenceAlias{I.10.3.2.fr}{I.10.3.2}
\EGAReferenceAlias{I.10.3.2.es}{I.10.3.2}
\EGAReferenceAlias{I.10.3.2.eq}{I.10.3.2}
\EGAReferenceAlias{eq:I.10.3.2}{I.10.3.2}
\EGAReferenceAlias{I.10.3.2.1}{I.10.3.2.1-fr}
\EGAReferenceAlias{I.10.3.2.1-es}{I.10.3.2.1-fr}
\EGAReferenceAlias{I.10.3.2.1.fr}{I.10.3.2.1-fr}
\EGAReferenceAlias{I.10.3.2.1.es}{I.10.3.2.1-fr}
\EGAReferenceAlias{I.10.3.2.1.eq}{I.10.3.2.1-fr}
\EGAReferenceAlias{eq:I.10.3.2.1}{I.10.3.2.1-fr}
\EGAReferenceAlias{I.10.3.3-fr}{I.10.3.3}
\EGAReferenceAlias{I.10.3.3-es}{I.10.3.3}
\EGAReferenceAlias{I.10.3.3.fr}{I.10.3.3}
\EGAReferenceAlias{I.10.3.3.es}{I.10.3.3}
\EGAReferenceAlias{I.10.3.3.eq}{I.10.3.3}
\EGAReferenceAlias{eq:I.10.3.3}{I.10.3.3}
\EGAReferenceAlias{I.10.3.4-fr}{I.10.3.4}
\EGAReferenceAlias{I.10.3.4-es}{I.10.3.4}
\EGAReferenceAlias{I.10.3.4.fr}{I.10.3.4}
\EGAReferenceAlias{I.10.3.4.es}{I.10.3.4}
\EGAReferenceAlias{I.10.3.4.eq}{I.10.3.4}
\EGAReferenceAlias{eq:I.10.3.4}{I.10.3.4}
\EGAReferenceAlias{I.10.3.5-fr}{I.10.3.5}
\EGAReferenceAlias{I.10.3.5-es}{I.10.3.5}
\EGAReferenceAlias{I.10.3.5.fr}{I.10.3.5}
\EGAReferenceAlias{I.10.3.5.es}{I.10.3.5}
\EGAReferenceAlias{I.10.3.5.eq}{I.10.3.5}
\EGAReferenceAlias{eq:I.10.3.5}{I.10.3.5}
\EGAReferenceAlias{I.10.3.6-fr}{I.10.3.6}
\EGAReferenceAlias{I.10.3.6-es}{I.10.3.6}
\EGAReferenceAlias{I.10.3.6.fr}{I.10.3.6}
\EGAReferenceAlias{I.10.3.6.es}{I.10.3.6}
\EGAReferenceAlias{I.10.3.6.eq}{I.10.3.6}
\EGAReferenceAlias{eq:I.10.3.6}{I.10.3.6}
\EGAReferenceAlias{I.10.3.7-fr}{I.10.3.7}
\EGAReferenceAlias{I.10.3.7-es}{I.10.3.7}
\EGAReferenceAlias{I.10.3.7.fr}{I.10.3.7}
\EGAReferenceAlias{I.10.3.7.es}{I.10.3.7}
\EGAReferenceAlias{I.10.3.7.eq}{I.10.3.7}
\EGAReferenceAlias{eq:I.10.3.7}{I.10.3.7}
\EGAReferenceAlias{I.10.4}{subsection:I.10.4}
\EGAReferenceAlias{I.10.4-fr}{subsection:I.10.4}
\EGAReferenceAlias{I.10.4-es}{subsection:I.10.4}
\EGAReferenceAlias{I.10.4.fr}{subsection:I.10.4}
\EGAReferenceAlias{I.10.4.es}{subsection:I.10.4}
\EGAReferenceAlias{I.10.4.eq}{subsection:I.10.4}
\EGAReferenceAlias{eq:I.10.4}{subsection:I.10.4}
\EGAReferenceAlias{subsection:I.10.4-fr}{subsection:I.10.4}
\EGAReferenceAlias{subsection:I.10.4-es}{subsection:I.10.4}
\EGAReferenceAlias{subsection:I.10.4.fr}{subsection:I.10.4}
\EGAReferenceAlias{subsection:I.10.4.es}{subsection:I.10.4}
\EGAReferenceAlias{subsection:I.10.4.eq}{subsection:I.10.4}
\EGAReferenceAlias{eq:subsection:I.10.4}{subsection:I.10.4}
\EGAReferenceAlias{I.10.4.1-fr}{I.10.4.1}
\EGAReferenceAlias{I.10.4.1-es}{I.10.4.1}
\EGAReferenceAlias{I.10.4.1.fr}{I.10.4.1}
\EGAReferenceAlias{I.10.4.1.es}{I.10.4.1}
\EGAReferenceAlias{I.10.4.1.eq}{I.10.4.1}
\EGAReferenceAlias{eq:I.10.4.1}{I.10.4.1}
\EGAReferenceAlias{I.10.4.2-fr}{I.10.4.2}
\EGAReferenceAlias{I.10.4.2-es}{I.10.4.2}
\EGAReferenceAlias{I.10.4.2.fr}{I.10.4.2}
\EGAReferenceAlias{I.10.4.2.es}{I.10.4.2}
\EGAReferenceAlias{I.10.4.2.eq}{I.10.4.2}
\EGAReferenceAlias{eq:I.10.4.2}{I.10.4.2}
\EGAReferenceAlias{I.10.4.3-fr}{I.10.4.3}
\EGAReferenceAlias{I.10.4.3-es}{I.10.4.3}
\EGAReferenceAlias{I.10.4.3.fr}{I.10.4.3}
\EGAReferenceAlias{I.10.4.3.es}{I.10.4.3}
\EGAReferenceAlias{I.10.4.3.eq}{I.10.4.3}
\EGAReferenceAlias{eq:I.10.4.3}{I.10.4.3}
\EGAReferenceAlias{I.10.4.4-fr}{I.10.4.4}
\EGAReferenceAlias{I.10.4.4-es}{I.10.4.4}
\EGAReferenceAlias{I.10.4.4.fr}{I.10.4.4}
\EGAReferenceAlias{I.10.4.4.es}{I.10.4.4}
\EGAReferenceAlias{I.10.4.4.eq}{I.10.4.4}
\EGAReferenceAlias{eq:I.10.4.4}{I.10.4.4}
\EGAReferenceAlias{I.10.4.5-fr}{I.10.4.5}
\EGAReferenceAlias{I.10.4.5-es}{I.10.4.5}
\EGAReferenceAlias{I.10.4.5.fr}{I.10.4.5}
\EGAReferenceAlias{I.10.4.5.es}{I.10.4.5}
\EGAReferenceAlias{I.10.4.5.eq}{I.10.4.5}
\EGAReferenceAlias{eq:I.10.4.5}{I.10.4.5}
\EGAReferenceAlias{I.10.4.6-fr}{I.10.4.6}
\EGAReferenceAlias{I.10.4.6-es}{I.10.4.6}
\EGAReferenceAlias{I.10.4.6.fr}{I.10.4.6}
\EGAReferenceAlias{I.10.4.6.es}{I.10.4.6}
\EGAReferenceAlias{I.10.4.6.eq}{I.10.4.6}
\EGAReferenceAlias{eq:I.10.4.6}{I.10.4.6}
\EGAReferenceAlias{I.10.4.7-fr}{I.10.4.7}
\EGAReferenceAlias{I.10.4.7-es}{I.10.4.7}
\EGAReferenceAlias{I.10.4.7.fr}{I.10.4.7}
\EGAReferenceAlias{I.10.4.7.es}{I.10.4.7}
\EGAReferenceAlias{I.10.4.7.eq}{I.10.4.7}
\EGAReferenceAlias{eq:I.10.4.7}{I.10.4.7}
\EGAReferenceAlias{I.10.4.8-fr}{I.10.4.8}
\EGAReferenceAlias{I.10.4.8-es}{I.10.4.8}
\EGAReferenceAlias{I.10.4.8.fr}{I.10.4.8}
\EGAReferenceAlias{I.10.4.8.es}{I.10.4.8}
\EGAReferenceAlias{I.10.4.8.eq}{I.10.4.8}
\EGAReferenceAlias{eq:I.10.4.8}{I.10.4.8}
\EGAReferenceAlias{I.10.5}{subsection:I.10.5}
\EGAReferenceAlias{I.10.5-fr}{subsection:I.10.5}
\EGAReferenceAlias{I.10.5-es}{subsection:I.10.5}
\EGAReferenceAlias{I.10.5.fr}{subsection:I.10.5}
\EGAReferenceAlias{I.10.5.es}{subsection:I.10.5}
\EGAReferenceAlias{I.10.5.eq}{subsection:I.10.5}
\EGAReferenceAlias{eq:I.10.5}{subsection:I.10.5}
\EGAReferenceAlias{subsection:I.10.5-fr}{subsection:I.10.5}
\EGAReferenceAlias{subsection:I.10.5-es}{subsection:I.10.5}
\EGAReferenceAlias{subsection:I.10.5.fr}{subsection:I.10.5}
\EGAReferenceAlias{subsection:I.10.5.es}{subsection:I.10.5}
\EGAReferenceAlias{subsection:I.10.5.eq}{subsection:I.10.5}
\EGAReferenceAlias{eq:subsection:I.10.5}{subsection:I.10.5}
\EGAReferenceAlias{I.10.5.1-fr}{I.10.5.1}
\EGAReferenceAlias{I.10.5.1-es}{I.10.5.1}
\EGAReferenceAlias{I.10.5.1.fr}{I.10.5.1}
\EGAReferenceAlias{I.10.5.1.es}{I.10.5.1}
\EGAReferenceAlias{I.10.5.1.eq}{I.10.5.1}
\EGAReferenceAlias{eq:I.10.5.1}{I.10.5.1}
\EGAReferenceAlias{I.10.5.2-fr}{I.10.5.2}
\EGAReferenceAlias{I.10.5.2-es}{I.10.5.2}
\EGAReferenceAlias{I.10.5.2.fr}{I.10.5.2}
\EGAReferenceAlias{I.10.5.2.es}{I.10.5.2}
\EGAReferenceAlias{I.10.5.2.eq}{I.10.5.2}
\EGAReferenceAlias{eq:I.10.5.2}{I.10.5.2}
\EGAReferenceAlias{I.10.5.3-fr}{I.10.5.3}
\EGAReferenceAlias{I.10.5.3-es}{I.10.5.3}
\EGAReferenceAlias{I.10.5.3.fr}{I.10.5.3}
\EGAReferenceAlias{I.10.5.3.es}{I.10.5.3}
\EGAReferenceAlias{I.10.5.3.eq}{I.10.5.3}
\EGAReferenceAlias{eq:I.10.5.3}{I.10.5.3}
\EGAReferenceAlias{I.10.5.4-fr}{I.10.5.4}
\EGAReferenceAlias{I.10.5.4-es}{I.10.5.4}
\EGAReferenceAlias{I.10.5.4.fr}{I.10.5.4}
\EGAReferenceAlias{I.10.5.4.es}{I.10.5.4}
\EGAReferenceAlias{I.10.5.4.eq}{I.10.5.4}
\EGAReferenceAlias{eq:I.10.5.4}{I.10.5.4}
\EGAReferenceAlias{I.10.5.5-fr}{I.10.5.5}
\EGAReferenceAlias{I.10.5.5-es}{I.10.5.5}
\EGAReferenceAlias{I.10.5.5.fr}{I.10.5.5}
\EGAReferenceAlias{I.10.5.5.es}{I.10.5.5}
\EGAReferenceAlias{I.10.5.5.eq}{I.10.5.5}
\EGAReferenceAlias{eq:I.10.5.5}{I.10.5.5}
\EGAReferenceAlias{I.10.5.6-fr}{I.10.5.6}
\EGAReferenceAlias{I.10.5.6-es}{I.10.5.6}
\EGAReferenceAlias{I.10.5.6.fr}{I.10.5.6}
\EGAReferenceAlias{I.10.5.6.es}{I.10.5.6}
\EGAReferenceAlias{I.10.5.6.eq}{I.10.5.6}
\EGAReferenceAlias{eq:I.10.5.6}{I.10.5.6}
\EGAReferenceAlias{I.10.5.6.1}{I.10.5.6.1-fr}
\EGAReferenceAlias{I.10.5.6.1-es}{I.10.5.6.1-fr}
\EGAReferenceAlias{I.10.5.6.1.fr}{I.10.5.6.1-fr}
\EGAReferenceAlias{I.10.5.6.1.es}{I.10.5.6.1-fr}
\EGAReferenceAlias{I.10.5.6.1.eq}{I.10.5.6.1-fr}
\EGAReferenceAlias{eq:I.10.5.6.1}{I.10.5.6.1-fr}
\EGAReferenceAlias{I.10.5.6.2}{I.10.5.6.2-fr}
\EGAReferenceAlias{I.10.5.6.2-es}{I.10.5.6.2-fr}
\EGAReferenceAlias{I.10.5.6.2.fr}{I.10.5.6.2-fr}
\EGAReferenceAlias{I.10.5.6.2.es}{I.10.5.6.2-fr}
\EGAReferenceAlias{I.10.5.6.2.eq}{I.10.5.6.2-fr}
\EGAReferenceAlias{eq:I.10.5.6.2}{I.10.5.6.2-fr}
\EGAReferenceAlias{I.10.6}{subsection:I.10.6}
\EGAReferenceAlias{I.10.6-fr}{subsection:I.10.6}
\EGAReferenceAlias{I.10.6-es}{subsection:I.10.6}
\EGAReferenceAlias{I.10.6.fr}{subsection:I.10.6}
\EGAReferenceAlias{I.10.6.es}{subsection:I.10.6}
\EGAReferenceAlias{I.10.6.eq}{subsection:I.10.6}
\EGAReferenceAlias{eq:I.10.6}{subsection:I.10.6}
\EGAReferenceAlias{subsection:I.10.6-fr}{subsection:I.10.6}
\EGAReferenceAlias{subsection:I.10.6-es}{subsection:I.10.6}
\EGAReferenceAlias{subsection:I.10.6.fr}{subsection:I.10.6}
\EGAReferenceAlias{subsection:I.10.6.es}{subsection:I.10.6}
\EGAReferenceAlias{subsection:I.10.6.eq}{subsection:I.10.6}
\EGAReferenceAlias{eq:subsection:I.10.6}{subsection:I.10.6}
\EGAReferenceAlias{I.10.6.1-fr}{I.10.6.1}
\EGAReferenceAlias{I.10.6.1-es}{I.10.6.1}
\EGAReferenceAlias{I.10.6.1.fr}{I.10.6.1}
\EGAReferenceAlias{I.10.6.1.es}{I.10.6.1}
\EGAReferenceAlias{I.10.6.1.eq}{I.10.6.1}
\EGAReferenceAlias{eq:I.10.6.1}{I.10.6.1}
\EGAReferenceAlias{I.10.6.2-fr}{I.10.6.2}
\EGAReferenceAlias{I.10.6.2-es}{I.10.6.2}
\EGAReferenceAlias{I.10.6.2.fr}{I.10.6.2}
\EGAReferenceAlias{I.10.6.2.es}{I.10.6.2}
\EGAReferenceAlias{I.10.6.2.eq}{I.10.6.2}
\EGAReferenceAlias{eq:I.10.6.2}{I.10.6.2}
\EGAReferenceAlias{I.10.6.2.1}{I.10.6.2.1-fr}
\EGAReferenceAlias{I.10.6.2.1-es}{I.10.6.2.1-fr}
\EGAReferenceAlias{I.10.6.2.1.fr}{I.10.6.2.1-fr}
\EGAReferenceAlias{I.10.6.2.1.es}{I.10.6.2.1-fr}
\EGAReferenceAlias{I.10.6.2.1.eq}{I.10.6.2.1-fr}
\EGAReferenceAlias{eq:I.10.6.2.1}{I.10.6.2.1-fr}
\EGAReferenceAlias{I.10.6.3-fr}{I.10.6.3}
\EGAReferenceAlias{I.10.6.3-es}{I.10.6.3}
\EGAReferenceAlias{I.10.6.3.fr}{I.10.6.3}
\EGAReferenceAlias{I.10.6.3.es}{I.10.6.3}
\EGAReferenceAlias{I.10.6.3.eq}{I.10.6.3}
\EGAReferenceAlias{eq:I.10.6.3}{I.10.6.3}
\EGAReferenceAlias{I.10.6.4-fr}{I.10.6.4}
\EGAReferenceAlias{I.10.6.4-es}{I.10.6.4}
\EGAReferenceAlias{I.10.6.4.fr}{I.10.6.4}
\EGAReferenceAlias{I.10.6.4.es}{I.10.6.4}
\EGAReferenceAlias{I.10.6.4.eq}{I.10.6.4}
\EGAReferenceAlias{eq:I.10.6.4}{I.10.6.4}
\EGAReferenceAlias{I.10.6.5-fr}{I.10.6.5}
\EGAReferenceAlias{I.10.6.5-es}{I.10.6.5}
\EGAReferenceAlias{I.10.6.5.fr}{I.10.6.5}
\EGAReferenceAlias{I.10.6.5.es}{I.10.6.5}
\EGAReferenceAlias{I.10.6.5.eq}{I.10.6.5}
\EGAReferenceAlias{eq:I.10.6.5}{I.10.6.5}
\EGAReferenceAlias{I.10.6.6-fr}{I.10.6.6}
\EGAReferenceAlias{I.10.6.6-es}{I.10.6.6}
\EGAReferenceAlias{I.10.6.6.fr}{I.10.6.6}
\EGAReferenceAlias{I.10.6.6.es}{I.10.6.6}
\EGAReferenceAlias{I.10.6.6.eq}{I.10.6.6}
\EGAReferenceAlias{eq:I.10.6.6}{I.10.6.6}
\EGAReferenceAlias{I.10.6.7-fr}{I.10.6.7}
\EGAReferenceAlias{I.10.6.7-es}{I.10.6.7}
\EGAReferenceAlias{I.10.6.7.fr}{I.10.6.7}
\EGAReferenceAlias{I.10.6.7.es}{I.10.6.7}
\EGAReferenceAlias{I.10.6.7.eq}{I.10.6.7}
\EGAReferenceAlias{eq:I.10.6.7}{I.10.6.7}
\EGAReferenceAlias{I.10.6.7.1}{I.10.6.7.1-fr}
\EGAReferenceAlias{I.10.6.7.1-es}{I.10.6.7.1-fr}
\EGAReferenceAlias{I.10.6.7.1.fr}{I.10.6.7.1-fr}
\EGAReferenceAlias{I.10.6.7.1.es}{I.10.6.7.1-fr}
\EGAReferenceAlias{I.10.6.7.1.eq}{I.10.6.7.1-fr}
\EGAReferenceAlias{eq:I.10.6.7.1}{I.10.6.7.1-fr}
\EGAReferenceAlias{I.10.6.8-fr}{I.10.6.8}
\EGAReferenceAlias{I.10.6.8-es}{I.10.6.8}
\EGAReferenceAlias{I.10.6.8.fr}{I.10.6.8}
\EGAReferenceAlias{I.10.6.8.es}{I.10.6.8}
\EGAReferenceAlias{I.10.6.8.eq}{I.10.6.8}
\EGAReferenceAlias{eq:I.10.6.8}{I.10.6.8}
\EGAReferenceAlias{I.10.6.9-fr}{I.10.6.9}
\EGAReferenceAlias{I.10.6.9-es}{I.10.6.9}
\EGAReferenceAlias{I.10.6.9.fr}{I.10.6.9}
\EGAReferenceAlias{I.10.6.9.es}{I.10.6.9}
\EGAReferenceAlias{I.10.6.9.eq}{I.10.6.9}
\EGAReferenceAlias{eq:I.10.6.9}{I.10.6.9}
\EGAReferenceAlias{I.10.6.10-fr}{I.10.6.10}
\EGAReferenceAlias{I.10.6.10-es}{I.10.6.10}
\EGAReferenceAlias{I.10.6.10.fr}{I.10.6.10}
\EGAReferenceAlias{I.10.6.10.es}{I.10.6.10}
\EGAReferenceAlias{I.10.6.10.eq}{I.10.6.10}
\EGAReferenceAlias{eq:I.10.6.10}{I.10.6.10}
\EGAReferenceAlias{I.10.6.11-fr}{I.10.6.11}
\EGAReferenceAlias{I.10.6.11-es}{I.10.6.11}
\EGAReferenceAlias{I.10.6.11.fr}{I.10.6.11}
\EGAReferenceAlias{I.10.6.11.es}{I.10.6.11}
\EGAReferenceAlias{I.10.6.11.eq}{I.10.6.11}
\EGAReferenceAlias{eq:I.10.6.11}{I.10.6.11}
\EGAReferenceAlias{I.10.7}{subsection:I.10.7}
\EGAReferenceAlias{I.10.7-fr}{subsection:I.10.7}
\EGAReferenceAlias{I.10.7-es}{subsection:I.10.7}
\EGAReferenceAlias{I.10.7.fr}{subsection:I.10.7}
\EGAReferenceAlias{I.10.7.es}{subsection:I.10.7}
\EGAReferenceAlias{I.10.7.eq}{subsection:I.10.7}
\EGAReferenceAlias{eq:I.10.7}{subsection:I.10.7}
\EGAReferenceAlias{subsection:I.10.7-fr}{subsection:I.10.7}
\EGAReferenceAlias{subsection:I.10.7-es}{subsection:I.10.7}
\EGAReferenceAlias{subsection:I.10.7.fr}{subsection:I.10.7}
\EGAReferenceAlias{subsection:I.10.7.es}{subsection:I.10.7}
\EGAReferenceAlias{subsection:I.10.7.eq}{subsection:I.10.7}
\EGAReferenceAlias{eq:subsection:I.10.7}{subsection:I.10.7}
\EGAReferenceAlias{I.10.7.1-fr}{I.10.7.1}
\EGAReferenceAlias{I.10.7.1-es}{I.10.7.1}
\EGAReferenceAlias{I.10.7.1.fr}{I.10.7.1}
\EGAReferenceAlias{I.10.7.1.es}{I.10.7.1}
\EGAReferenceAlias{I.10.7.1.eq}{I.10.7.1}
\EGAReferenceAlias{eq:I.10.7.1}{I.10.7.1}
\EGAReferenceAlias{I.10.7.2-fr}{I.10.7.2}
\EGAReferenceAlias{I.10.7.2-es}{I.10.7.2}
\EGAReferenceAlias{I.10.7.2.fr}{I.10.7.2}
\EGAReferenceAlias{I.10.7.2.es}{I.10.7.2}
\EGAReferenceAlias{I.10.7.2.eq}{I.10.7.2}
\EGAReferenceAlias{eq:I.10.7.2}{I.10.7.2}
\EGAReferenceAlias{I.10.7.3-fr}{I.10.7.3}
\EGAReferenceAlias{I.10.7.3-es}{I.10.7.3}
\EGAReferenceAlias{I.10.7.3.fr}{I.10.7.3}
\EGAReferenceAlias{I.10.7.3.es}{I.10.7.3}
\EGAReferenceAlias{I.10.7.3.eq}{I.10.7.3}
\EGAReferenceAlias{eq:I.10.7.3}{I.10.7.3}
\EGAReferenceAlias{I.10.7.4-fr}{I.10.7.4}
\EGAReferenceAlias{I.10.7.4-es}{I.10.7.4}
\EGAReferenceAlias{I.10.7.4.fr}{I.10.7.4}
\EGAReferenceAlias{I.10.7.4.es}{I.10.7.4}
\EGAReferenceAlias{I.10.7.4.eq}{I.10.7.4}
\EGAReferenceAlias{eq:I.10.7.4}{I.10.7.4}
\EGAReferenceAlias{I.10.8}{subsection:I.10.8}
\EGAReferenceAlias{I.10.8-fr}{subsection:I.10.8}
\EGAReferenceAlias{I.10.8-es}{subsection:I.10.8}
\EGAReferenceAlias{I.10.8.fr}{subsection:I.10.8}
\EGAReferenceAlias{I.10.8.es}{subsection:I.10.8}
\EGAReferenceAlias{I.10.8.eq}{subsection:I.10.8}
\EGAReferenceAlias{eq:I.10.8}{subsection:I.10.8}
\EGAReferenceAlias{subsection:I.10.8-fr}{subsection:I.10.8}
\EGAReferenceAlias{subsection:I.10.8-es}{subsection:I.10.8}
\EGAReferenceAlias{subsection:I.10.8.fr}{subsection:I.10.8}
\EGAReferenceAlias{subsection:I.10.8.es}{subsection:I.10.8}
\EGAReferenceAlias{subsection:I.10.8.eq}{subsection:I.10.8}
\EGAReferenceAlias{eq:subsection:I.10.8}{subsection:I.10.8}
\EGAReferenceAlias{I.10.8.1-fr}{I.10.8.1}
\EGAReferenceAlias{I.10.8.1-es}{I.10.8.1}
\EGAReferenceAlias{I.10.8.1.fr}{I.10.8.1}
\EGAReferenceAlias{I.10.8.1.es}{I.10.8.1}
\EGAReferenceAlias{I.10.8.1.eq}{I.10.8.1}
\EGAReferenceAlias{eq:I.10.8.1}{I.10.8.1}
\EGAReferenceAlias{I.10.8.2-fr}{I.10.8.2}
\EGAReferenceAlias{I.10.8.2-es}{I.10.8.2}
\EGAReferenceAlias{I.10.8.2.fr}{I.10.8.2}
\EGAReferenceAlias{I.10.8.2.es}{I.10.8.2}
\EGAReferenceAlias{I.10.8.2.eq}{I.10.8.2}
\EGAReferenceAlias{eq:I.10.8.2}{I.10.8.2}
\EGAReferenceAlias{I.10.8.3-fr}{I.10.8.3}
\EGAReferenceAlias{I.10.8.3-es}{I.10.8.3}
\EGAReferenceAlias{I.10.8.3.fr}{I.10.8.3}
\EGAReferenceAlias{I.10.8.3.es}{I.10.8.3}
\EGAReferenceAlias{I.10.8.3.eq}{I.10.8.3}
\EGAReferenceAlias{eq:I.10.8.3}{I.10.8.3}
\EGAReferenceAlias{I.10.8.4-fr}{I.10.8.4}
\EGAReferenceAlias{I.10.8.4-es}{I.10.8.4}
\EGAReferenceAlias{I.10.8.4.fr}{I.10.8.4}
\EGAReferenceAlias{I.10.8.4.es}{I.10.8.4}
\EGAReferenceAlias{I.10.8.4.eq}{I.10.8.4}
\EGAReferenceAlias{eq:I.10.8.4}{I.10.8.4}
\EGAReferenceAlias{I.10.8.5-fr}{I.10.8.5}
\EGAReferenceAlias{I.10.8.5-es}{I.10.8.5}
\EGAReferenceAlias{I.10.8.5.fr}{I.10.8.5}
\EGAReferenceAlias{I.10.8.5.es}{I.10.8.5}
\EGAReferenceAlias{I.10.8.5.eq}{I.10.8.5}
\EGAReferenceAlias{eq:I.10.8.5}{I.10.8.5}
\EGAReferenceAlias{I.10.8.6-fr}{I.10.8.6}
\EGAReferenceAlias{I.10.8.6-es}{I.10.8.6}
\EGAReferenceAlias{I.10.8.6.fr}{I.10.8.6}
\EGAReferenceAlias{I.10.8.6.es}{I.10.8.6}
\EGAReferenceAlias{I.10.8.6.eq}{I.10.8.6}
\EGAReferenceAlias{eq:I.10.8.6}{I.10.8.6}
\EGAReferenceAlias{I.10.8.7-fr}{I.10.8.7}
\EGAReferenceAlias{I.10.8.7-es}{I.10.8.7}
\EGAReferenceAlias{I.10.8.7.fr}{I.10.8.7}
\EGAReferenceAlias{I.10.8.7.es}{I.10.8.7}
\EGAReferenceAlias{I.10.8.7.eq}{I.10.8.7}
\EGAReferenceAlias{eq:I.10.8.7}{I.10.8.7}
\EGAReferenceAlias{I.10.8.8-fr}{I.10.8.8}
\EGAReferenceAlias{I.10.8.8-es}{I.10.8.8}
\EGAReferenceAlias{I.10.8.8.fr}{I.10.8.8}
\EGAReferenceAlias{I.10.8.8.es}{I.10.8.8}
\EGAReferenceAlias{I.10.8.8.eq}{I.10.8.8}
\EGAReferenceAlias{eq:I.10.8.8}{I.10.8.8}
\EGAReferenceAlias{I.10.8.8.1}{I.10.8.8.1-fr}
\EGAReferenceAlias{I.10.8.8.1-es}{I.10.8.8.1-fr}
\EGAReferenceAlias{I.10.8.8.1.fr}{I.10.8.8.1-fr}
\EGAReferenceAlias{I.10.8.8.1.es}{I.10.8.8.1-fr}
\EGAReferenceAlias{I.10.8.8.1.eq}{I.10.8.8.1-fr}
\EGAReferenceAlias{eq:I.10.8.8.1}{I.10.8.8.1-fr}
\EGAReferenceAlias{I.10.8.9-fr}{I.10.8.9}
\EGAReferenceAlias{I.10.8.9-es}{I.10.8.9}
\EGAReferenceAlias{I.10.8.9.fr}{I.10.8.9}
\EGAReferenceAlias{I.10.8.9.es}{I.10.8.9}
\EGAReferenceAlias{I.10.8.9.eq}{I.10.8.9}
\EGAReferenceAlias{eq:I.10.8.9}{I.10.8.9}
\EGAReferenceAlias{I.10.8.10-fr}{I.10.8.10}
\EGAReferenceAlias{I.10.8.10-es}{I.10.8.10}
\EGAReferenceAlias{I.10.8.10.fr}{I.10.8.10}
\EGAReferenceAlias{I.10.8.10.es}{I.10.8.10}
\EGAReferenceAlias{I.10.8.10.eq}{I.10.8.10}
\EGAReferenceAlias{eq:I.10.8.10}{I.10.8.10}
\EGAReferenceAlias{I.10.8.10.1}{I.10.8.10.1-fr}
\EGAReferenceAlias{I.10.8.10.1-es}{I.10.8.10.1-fr}
\EGAReferenceAlias{I.10.8.10.1.fr}{I.10.8.10.1-fr}
\EGAReferenceAlias{I.10.8.10.1.es}{I.10.8.10.1-fr}
\EGAReferenceAlias{I.10.8.10.1.eq}{I.10.8.10.1-fr}
\EGAReferenceAlias{eq:I.10.8.10.1}{I.10.8.10.1-fr}
\EGAReferenceAlias{I.10.8.10.2}{I.10.8.10.2-fr}
\EGAReferenceAlias{I.10.8.10.2-es}{I.10.8.10.2-fr}
\EGAReferenceAlias{I.10.8.10.2.fr}{I.10.8.10.2-fr}
\EGAReferenceAlias{I.10.8.10.2.es}{I.10.8.10.2-fr}
\EGAReferenceAlias{I.10.8.10.2.eq}{I.10.8.10.2-fr}
\EGAReferenceAlias{eq:I.10.8.10.2}{I.10.8.10.2-fr}
\EGAReferenceAlias{I.10.8.11-fr}{I.10.8.11}
\EGAReferenceAlias{I.10.8.11-es}{I.10.8.11}
\EGAReferenceAlias{I.10.8.11.fr}{I.10.8.11}
\EGAReferenceAlias{I.10.8.11.es}{I.10.8.11}
\EGAReferenceAlias{I.10.8.11.eq}{I.10.8.11}
\EGAReferenceAlias{eq:I.10.8.11}{I.10.8.11}
\EGAReferenceAlias{I.10.8.12-fr}{I.10.8.12}
\EGAReferenceAlias{I.10.8.12-es}{I.10.8.12}
\EGAReferenceAlias{I.10.8.12.fr}{I.10.8.12}
\EGAReferenceAlias{I.10.8.12.es}{I.10.8.12}
\EGAReferenceAlias{I.10.8.12.eq}{I.10.8.12}
\EGAReferenceAlias{eq:I.10.8.12}{I.10.8.12}
\EGAReferenceAlias{I.10.8.13-fr}{I.10.8.13}
\EGAReferenceAlias{I.10.8.13-es}{I.10.8.13}
\EGAReferenceAlias{I.10.8.13.fr}{I.10.8.13}
\EGAReferenceAlias{I.10.8.13.es}{I.10.8.13}
\EGAReferenceAlias{I.10.8.13.eq}{I.10.8.13}
\EGAReferenceAlias{eq:I.10.8.13}{I.10.8.13}
\EGAReferenceAlias{I.10.8.14-fr}{I.10.8.14}
\EGAReferenceAlias{I.10.8.14-es}{I.10.8.14}
\EGAReferenceAlias{I.10.8.14.fr}{I.10.8.14}
\EGAReferenceAlias{I.10.8.14.es}{I.10.8.14}
\EGAReferenceAlias{I.10.8.14.eq}{I.10.8.14}
\EGAReferenceAlias{eq:I.10.8.14}{I.10.8.14}
\EGAReferenceAlias{I.10.9}{subsection:I.10.9}
\EGAReferenceAlias{I.10.9-fr}{subsection:I.10.9}
\EGAReferenceAlias{I.10.9-es}{subsection:I.10.9}
\EGAReferenceAlias{I.10.9.fr}{subsection:I.10.9}
\EGAReferenceAlias{I.10.9.es}{subsection:I.10.9}
\EGAReferenceAlias{I.10.9.eq}{subsection:I.10.9}
\EGAReferenceAlias{eq:I.10.9}{subsection:I.10.9}
\EGAReferenceAlias{subsection:I.10.9-fr}{subsection:I.10.9}
\EGAReferenceAlias{subsection:I.10.9-es}{subsection:I.10.9}
\EGAReferenceAlias{subsection:I.10.9.fr}{subsection:I.10.9}
\EGAReferenceAlias{subsection:I.10.9.es}{subsection:I.10.9}
\EGAReferenceAlias{subsection:I.10.9.eq}{subsection:I.10.9}
\EGAReferenceAlias{eq:subsection:I.10.9}{subsection:I.10.9}
\EGAReferenceAlias{I.10.9.1-fr}{I.10.9.1}
\EGAReferenceAlias{I.10.9.1-es}{I.10.9.1}
\EGAReferenceAlias{I.10.9.1.fr}{I.10.9.1}
\EGAReferenceAlias{I.10.9.1.es}{I.10.9.1}
\EGAReferenceAlias{I.10.9.1.eq}{I.10.9.1}
\EGAReferenceAlias{eq:I.10.9.1}{I.10.9.1}
\EGAReferenceAlias{I.10.9.2-fr}{I.10.9.2}
\EGAReferenceAlias{I.10.9.2-es}{I.10.9.2}
\EGAReferenceAlias{I.10.9.2.fr}{I.10.9.2}
\EGAReferenceAlias{I.10.9.2.es}{I.10.9.2}
\EGAReferenceAlias{I.10.9.2.eq}{I.10.9.2}
\EGAReferenceAlias{eq:I.10.9.2}{I.10.9.2}
\EGAReferenceAlias{I.10.9.2.1}{I.10.9.2.1-fr}
\EGAReferenceAlias{I.10.9.2.1-es}{I.10.9.2.1-fr}
\EGAReferenceAlias{I.10.9.2.1.fr}{I.10.9.2.1-fr}
\EGAReferenceAlias{I.10.9.2.1.es}{I.10.9.2.1-fr}
\EGAReferenceAlias{I.10.9.2.1.eq}{I.10.9.2.1-fr}
\EGAReferenceAlias{eq:I.10.9.2.1}{I.10.9.2.1-fr}
\EGAReferenceAlias{I.10.9.3-fr}{I.10.9.3}
\EGAReferenceAlias{I.10.9.3-es}{I.10.9.3}
\EGAReferenceAlias{I.10.9.3.fr}{I.10.9.3}
\EGAReferenceAlias{I.10.9.3.es}{I.10.9.3}
\EGAReferenceAlias{I.10.9.3.eq}{I.10.9.3}
\EGAReferenceAlias{eq:I.10.9.3}{I.10.9.3}
\EGAReferenceAlias{I.10.9.4-fr}{I.10.9.4}
\EGAReferenceAlias{I.10.9.4-es}{I.10.9.4}
\EGAReferenceAlias{I.10.9.4.fr}{I.10.9.4}
\EGAReferenceAlias{I.10.9.4.es}{I.10.9.4}
\EGAReferenceAlias{I.10.9.4.eq}{I.10.9.4}
\EGAReferenceAlias{eq:I.10.9.4}{I.10.9.4}
\EGAReferenceAlias{I.10.9.5-fr}{I.10.9.5}
\EGAReferenceAlias{I.10.9.5-es}{I.10.9.5}
\EGAReferenceAlias{I.10.9.5.fr}{I.10.9.5}
\EGAReferenceAlias{I.10.9.5.es}{I.10.9.5}
\EGAReferenceAlias{I.10.9.5.eq}{I.10.9.5}
\EGAReferenceAlias{eq:I.10.9.5}{I.10.9.5}
\EGAReferenceAlias{I.10.9.6-fr}{I.10.9.6}
\EGAReferenceAlias{I.10.9.6-es}{I.10.9.6}
\EGAReferenceAlias{I.10.9.6.fr}{I.10.9.6}
\EGAReferenceAlias{I.10.9.6.es}{I.10.9.6}
\EGAReferenceAlias{I.10.9.6.eq}{I.10.9.6}
\EGAReferenceAlias{eq:I.10.9.6}{I.10.9.6}
\EGAReferenceAlias{I.10.9.7-fr}{I.10.9.7}
\EGAReferenceAlias{I.10.9.7-es}{I.10.9.7}
\EGAReferenceAlias{I.10.9.7.fr}{I.10.9.7}
\EGAReferenceAlias{I.10.9.7.es}{I.10.9.7}
\EGAReferenceAlias{I.10.9.7.eq}{I.10.9.7}
\EGAReferenceAlias{eq:I.10.9.7}{I.10.9.7}
\EGAReferenceAlias{I.10.9.8-fr}{I.10.9.8}
\EGAReferenceAlias{I.10.9.8-es}{I.10.9.8}
\EGAReferenceAlias{I.10.9.8.fr}{I.10.9.8}
\EGAReferenceAlias{I.10.9.8.es}{I.10.9.8}
\EGAReferenceAlias{I.10.9.8.eq}{I.10.9.8}
\EGAReferenceAlias{eq:I.10.9.8}{I.10.9.8}
\EGAReferenceAlias{I.10.9.9-fr}{I.10.9.9}
\EGAReferenceAlias{I.10.9.9-es}{I.10.9.9}
\EGAReferenceAlias{I.10.9.9.fr}{I.10.9.9}
\EGAReferenceAlias{I.10.9.9.es}{I.10.9.9}
\EGAReferenceAlias{I.10.9.9.eq}{I.10.9.9}
\EGAReferenceAlias{eq:I.10.9.9}{I.10.9.9}
\EGAReferenceAlias{I.10.9.9.1}{I.10.9.9.1-fr}
\EGAReferenceAlias{I.10.9.9.1-es}{I.10.9.9.1-fr}
\EGAReferenceAlias{I.10.9.9.1.fr}{I.10.9.9.1-fr}
\EGAReferenceAlias{I.10.9.9.1.es}{I.10.9.9.1-fr}
\EGAReferenceAlias{I.10.9.9.1.eq}{I.10.9.9.1-fr}
\EGAReferenceAlias{eq:I.10.9.9.1}{I.10.9.9.1-fr}
\EGAReferenceAlias{I.10.9.10-fr}{I.10.9.10}
\EGAReferenceAlias{I.10.9.10-es}{I.10.9.10}
\EGAReferenceAlias{I.10.9.10.fr}{I.10.9.10}
\EGAReferenceAlias{I.10.9.10.es}{I.10.9.10}
\EGAReferenceAlias{I.10.9.10.eq}{I.10.9.10}
\EGAReferenceAlias{eq:I.10.9.10}{I.10.9.10}
\EGAReferenceAlias{I.10.10}{subsection:I.10.10}
\EGAReferenceAlias{I.10.10-fr}{subsection:I.10.10}
\EGAReferenceAlias{I.10.10-es}{subsection:I.10.10}
\EGAReferenceAlias{I.10.10.fr}{subsection:I.10.10}
\EGAReferenceAlias{I.10.10.es}{subsection:I.10.10}
\EGAReferenceAlias{I.10.10.eq}{subsection:I.10.10}
\EGAReferenceAlias{eq:I.10.10}{subsection:I.10.10}
\EGAReferenceAlias{subsection:I.10.10-fr}{subsection:I.10.10}
\EGAReferenceAlias{subsection:I.10.10-es}{subsection:I.10.10}
\EGAReferenceAlias{subsection:I.10.10.fr}{subsection:I.10.10}
\EGAReferenceAlias{subsection:I.10.10.es}{subsection:I.10.10}
\EGAReferenceAlias{subsection:I.10.10.eq}{subsection:I.10.10}
\EGAReferenceAlias{eq:subsection:I.10.10}{subsection:I.10.10}
\EGAReferenceAlias{I.10.10.1-fr}{I.10.10.1}
\EGAReferenceAlias{I.10.10.1-es}{I.10.10.1}
\EGAReferenceAlias{I.10.10.1.fr}{I.10.10.1}
\EGAReferenceAlias{I.10.10.1.es}{I.10.10.1}
\EGAReferenceAlias{I.10.10.1.eq}{I.10.10.1}
\EGAReferenceAlias{eq:I.10.10.1}{I.10.10.1}
\EGAReferenceAlias{I.10.10.2-fr}{I.10.10.2}
\EGAReferenceAlias{I.10.10.2-es}{I.10.10.2}
\EGAReferenceAlias{I.10.10.2.fr}{I.10.10.2}
\EGAReferenceAlias{I.10.10.2.es}{I.10.10.2}
\EGAReferenceAlias{I.10.10.2.eq}{I.10.10.2}
\EGAReferenceAlias{eq:I.10.10.2}{I.10.10.2}
\EGAReferenceAlias{I.10.10.2.1}{I.10.10.2.1-fr}
\EGAReferenceAlias{I.10.10.2.1-es}{I.10.10.2.1-fr}
\EGAReferenceAlias{I.10.10.2.1.fr}{I.10.10.2.1-fr}
\EGAReferenceAlias{I.10.10.2.1.es}{I.10.10.2.1-fr}
\EGAReferenceAlias{I.10.10.2.1.eq}{I.10.10.2.1-fr}
\EGAReferenceAlias{eq:I.10.10.2.1}{I.10.10.2.1-fr}
\EGAReferenceAlias{I.10.10.2.2}{I.10.10.2.2-fr}
\EGAReferenceAlias{I.10.10.2.2-es}{I.10.10.2.2-fr}
\EGAReferenceAlias{I.10.10.2.2.fr}{I.10.10.2.2-fr}
\EGAReferenceAlias{I.10.10.2.2.es}{I.10.10.2.2-fr}
\EGAReferenceAlias{I.10.10.2.2.eq}{I.10.10.2.2-fr}
\EGAReferenceAlias{eq:I.10.10.2.2}{I.10.10.2.2-fr}
\EGAReferenceAlias{I.10.10.2.3}{I.10.10.2.3-fr}
\EGAReferenceAlias{I.10.10.2.3-es}{I.10.10.2.3-fr}
\EGAReferenceAlias{I.10.10.2.3.fr}{I.10.10.2.3-fr}
\EGAReferenceAlias{I.10.10.2.3.es}{I.10.10.2.3-fr}
\EGAReferenceAlias{I.10.10.2.3.eq}{I.10.10.2.3-fr}
\EGAReferenceAlias{eq:I.10.10.2.3}{I.10.10.2.3-fr}
\EGAReferenceAlias{I.10.10.3-fr}{I.10.10.3}
\EGAReferenceAlias{I.10.10.3-es}{I.10.10.3}
\EGAReferenceAlias{I.10.10.3.fr}{I.10.10.3}
\EGAReferenceAlias{I.10.10.3.es}{I.10.10.3}
\EGAReferenceAlias{I.10.10.3.eq}{I.10.10.3}
\EGAReferenceAlias{eq:I.10.10.3}{I.10.10.3}
\EGAReferenceAlias{I.10.10.4-fr}{I.10.10.4}
\EGAReferenceAlias{I.10.10.4-es}{I.10.10.4}
\EGAReferenceAlias{I.10.10.4.fr}{I.10.10.4}
\EGAReferenceAlias{I.10.10.4.es}{I.10.10.4}
\EGAReferenceAlias{I.10.10.4.eq}{I.10.10.4}
\EGAReferenceAlias{eq:I.10.10.4}{I.10.10.4}
\EGAReferenceAlias{I.10.10.5-fr}{I.10.10.5}
\EGAReferenceAlias{I.10.10.5-es}{I.10.10.5}
\EGAReferenceAlias{I.10.10.5.fr}{I.10.10.5}
\EGAReferenceAlias{I.10.10.5.es}{I.10.10.5}
\EGAReferenceAlias{I.10.10.5.eq}{I.10.10.5}
\EGAReferenceAlias{eq:I.10.10.5}{I.10.10.5}
\EGAReferenceAlias{I.10.10.6-fr}{I.10.10.6}
\EGAReferenceAlias{I.10.10.6-es}{I.10.10.6}
\EGAReferenceAlias{I.10.10.6.fr}{I.10.10.6}
\EGAReferenceAlias{I.10.10.6.es}{I.10.10.6}
\EGAReferenceAlias{I.10.10.6.eq}{I.10.10.6}
\EGAReferenceAlias{eq:I.10.10.6}{I.10.10.6}
\EGAReferenceAlias{I.10.10.7-fr}{I.10.10.7}
\EGAReferenceAlias{I.10.10.7-es}{I.10.10.7}
\EGAReferenceAlias{I.10.10.7.fr}{I.10.10.7}
\EGAReferenceAlias{I.10.10.7.es}{I.10.10.7}
\EGAReferenceAlias{I.10.10.7.eq}{I.10.10.7}
\EGAReferenceAlias{eq:I.10.10.7}{I.10.10.7}
\EGAReferenceAlias{I.10.10.8-fr}{I.10.10.8}
\EGAReferenceAlias{I.10.10.8-es}{I.10.10.8}
\EGAReferenceAlias{I.10.10.8.fr}{I.10.10.8}
\EGAReferenceAlias{I.10.10.8.es}{I.10.10.8}
\EGAReferenceAlias{I.10.10.8.eq}{I.10.10.8}
\EGAReferenceAlias{eq:I.10.10.8}{I.10.10.8}
\EGAReferenceAlias{I.10.10.9-fr}{I.10.10.9}
\EGAReferenceAlias{I.10.10.9-es}{I.10.10.9}
\EGAReferenceAlias{I.10.10.9.fr}{I.10.10.9}
\EGAReferenceAlias{I.10.10.9.es}{I.10.10.9}
\EGAReferenceAlias{I.10.10.9.eq}{I.10.10.9}
\EGAReferenceAlias{eq:I.10.10.9}{I.10.10.9}
\EGAReferenceAlias{I.10.11}{subsection:I.10.11}
\EGAReferenceAlias{I.10.11-fr}{subsection:I.10.11}
\EGAReferenceAlias{I.10.11-es}{subsection:I.10.11}
\EGAReferenceAlias{I.10.11.fr}{subsection:I.10.11}
\EGAReferenceAlias{I.10.11.es}{subsection:I.10.11}
\EGAReferenceAlias{I.10.11.eq}{subsection:I.10.11}
\EGAReferenceAlias{eq:I.10.11}{subsection:I.10.11}
\EGAReferenceAlias{subsection:I.10.11-fr}{subsection:I.10.11}
\EGAReferenceAlias{subsection:I.10.11-es}{subsection:I.10.11}
\EGAReferenceAlias{subsection:I.10.11.fr}{subsection:I.10.11}
\EGAReferenceAlias{subsection:I.10.11.es}{subsection:I.10.11}
\EGAReferenceAlias{subsection:I.10.11.eq}{subsection:I.10.11}
\EGAReferenceAlias{eq:subsection:I.10.11}{subsection:I.10.11}
\EGAReferenceAlias{I.10.11.1-fr}{I.10.11.1}
\EGAReferenceAlias{I.10.11.1-es}{I.10.11.1}
\EGAReferenceAlias{I.10.11.1.fr}{I.10.11.1}
\EGAReferenceAlias{I.10.11.1.es}{I.10.11.1}
\EGAReferenceAlias{I.10.11.1.eq}{I.10.11.1}
\EGAReferenceAlias{eq:I.10.11.1}{I.10.11.1}
\EGAReferenceAlias{I.10.11.2-fr}{I.10.11.2}
\EGAReferenceAlias{I.10.11.2-es}{I.10.11.2}
\EGAReferenceAlias{I.10.11.2.fr}{I.10.11.2}
\EGAReferenceAlias{I.10.11.2.es}{I.10.11.2}
\EGAReferenceAlias{I.10.11.2.eq}{I.10.11.2}
\EGAReferenceAlias{eq:I.10.11.2}{I.10.11.2}
\EGAReferenceAlias{I.10.11.3-fr}{I.10.11.3}
\EGAReferenceAlias{I.10.11.3-es}{I.10.11.3}
\EGAReferenceAlias{I.10.11.3.fr}{I.10.11.3}
\EGAReferenceAlias{I.10.11.3.es}{I.10.11.3}
\EGAReferenceAlias{I.10.11.3.eq}{I.10.11.3}
\EGAReferenceAlias{eq:I.10.11.3}{I.10.11.3}
\EGAReferenceAlias{I.10.11.4-fr}{I.10.11.4}
\EGAReferenceAlias{I.10.11.4-es}{I.10.11.4}
\EGAReferenceAlias{I.10.11.4.fr}{I.10.11.4}
\EGAReferenceAlias{I.10.11.4.es}{I.10.11.4}
\EGAReferenceAlias{I.10.11.4.eq}{I.10.11.4}
\EGAReferenceAlias{eq:I.10.11.4}{I.10.11.4}
\EGAReferenceAlias{I.10.11.4.1}{I.10.11.4.1-fr}
\EGAReferenceAlias{I.10.11.4.1-es}{I.10.11.4.1-fr}
\EGAReferenceAlias{I.10.11.4.1.fr}{I.10.11.4.1-fr}
\EGAReferenceAlias{I.10.11.4.1.es}{I.10.11.4.1-fr}
\EGAReferenceAlias{I.10.11.4.1.eq}{I.10.11.4.1-fr}
\EGAReferenceAlias{eq:I.10.11.4.1}{I.10.11.4.1-fr}
\EGAReferenceAlias{I.10.11.5-fr}{I.10.11.5}
\EGAReferenceAlias{I.10.11.5-es}{I.10.11.5}
\EGAReferenceAlias{I.10.11.5.fr}{I.10.11.5}
\EGAReferenceAlias{I.10.11.5.es}{I.10.11.5}
\EGAReferenceAlias{I.10.11.5.eq}{I.10.11.5}
\EGAReferenceAlias{eq:I.10.11.5}{I.10.11.5}
\EGAReferenceAlias{I.10.11.6-fr}{I.10.11.6}
\EGAReferenceAlias{I.10.11.6-es}{I.10.11.6}
\EGAReferenceAlias{I.10.11.6.fr}{I.10.11.6}
\EGAReferenceAlias{I.10.11.6.es}{I.10.11.6}
\EGAReferenceAlias{I.10.11.6.eq}{I.10.11.6}
\EGAReferenceAlias{eq:I.10.11.6}{I.10.11.6}
\EGAReferenceAlias{I.10.11.7-fr}{I.10.11.7}
\EGAReferenceAlias{I.10.11.7-es}{I.10.11.7}
\EGAReferenceAlias{I.10.11.7.fr}{I.10.11.7}
\EGAReferenceAlias{I.10.11.7.es}{I.10.11.7}
\EGAReferenceAlias{I.10.11.7.eq}{I.10.11.7}
\EGAReferenceAlias{eq:I.10.11.7}{I.10.11.7}
\EGAReferenceAlias{I.10.11.7.1}{I.10.11.7.1-fr}
\EGAReferenceAlias{I.10.11.7.1-es}{I.10.11.7.1-fr}
\EGAReferenceAlias{I.10.11.7.1.fr}{I.10.11.7.1-fr}
\EGAReferenceAlias{I.10.11.7.1.es}{I.10.11.7.1-fr}
\EGAReferenceAlias{I.10.11.7.1.eq}{I.10.11.7.1-fr}
\EGAReferenceAlias{eq:I.10.11.7.1}{I.10.11.7.1-fr}
\EGAReferenceAlias{I.10.11.7.2}{I.10.11.7.2-fr}
\EGAReferenceAlias{I.10.11.7.2-es}{I.10.11.7.2-fr}
\EGAReferenceAlias{I.10.11.7.2.fr}{I.10.11.7.2-fr}
\EGAReferenceAlias{I.10.11.7.2.es}{I.10.11.7.2-fr}
\EGAReferenceAlias{I.10.11.7.2.eq}{I.10.11.7.2-fr}
\EGAReferenceAlias{eq:I.10.11.7.2}{I.10.11.7.2-fr}
\EGAReferenceAlias{I.10.11.8-fr}{I.10.11.8}
\EGAReferenceAlias{I.10.11.8-es}{I.10.11.8}
\EGAReferenceAlias{I.10.11.8.fr}{I.10.11.8}
\EGAReferenceAlias{I.10.11.8.es}{I.10.11.8}
\EGAReferenceAlias{I.10.11.8.eq}{I.10.11.8}
\EGAReferenceAlias{eq:I.10.11.8}{I.10.11.8}
\EGAReferenceAlias{I.10.11.9-fr}{I.10.11.9}
\EGAReferenceAlias{I.10.11.9-es}{I.10.11.9}
\EGAReferenceAlias{I.10.11.9.fr}{I.10.11.9}
\EGAReferenceAlias{I.10.11.9.es}{I.10.11.9}
\EGAReferenceAlias{I.10.11.9.eq}{I.10.11.9}
\EGAReferenceAlias{eq:I.10.11.9}{I.10.11.9}
\EGAReferenceAlias{I.10.12}{subsection:I.10.12}
\EGAReferenceAlias{I.10.12-fr}{subsection:I.10.12}
\EGAReferenceAlias{I.10.12-es}{subsection:I.10.12}
\EGAReferenceAlias{I.10.12.fr}{subsection:I.10.12}
\EGAReferenceAlias{I.10.12.es}{subsection:I.10.12}
\EGAReferenceAlias{I.10.12.eq}{subsection:I.10.12}
\EGAReferenceAlias{eq:I.10.12}{subsection:I.10.12}
\EGAReferenceAlias{subsection:I.10.12-fr}{subsection:I.10.12}
\EGAReferenceAlias{subsection:I.10.12-es}{subsection:I.10.12}
\EGAReferenceAlias{subsection:I.10.12.fr}{subsection:I.10.12}
\EGAReferenceAlias{subsection:I.10.12.es}{subsection:I.10.12}
\EGAReferenceAlias{subsection:I.10.12.eq}{subsection:I.10.12}
\EGAReferenceAlias{eq:subsection:I.10.12}{subsection:I.10.12}
\EGAReferenceAlias{I.10.12.1-fr}{I.10.12.1}
\EGAReferenceAlias{I.10.12.1-es}{I.10.12.1}
\EGAReferenceAlias{I.10.12.1.fr}{I.10.12.1}
\EGAReferenceAlias{I.10.12.1.es}{I.10.12.1}
\EGAReferenceAlias{I.10.12.1.eq}{I.10.12.1}
\EGAReferenceAlias{eq:I.10.12.1}{I.10.12.1}
\EGAReferenceAlias{I.10.12.2-fr}{I.10.12.2}
\EGAReferenceAlias{I.10.12.2-es}{I.10.12.2}
\EGAReferenceAlias{I.10.12.2.fr}{I.10.12.2}
\EGAReferenceAlias{I.10.12.2.es}{I.10.12.2}
\EGAReferenceAlias{I.10.12.2.eq}{I.10.12.2}
\EGAReferenceAlias{eq:I.10.12.2}{I.10.12.2}
\EGAReferenceAlias{I.10.12.2.1}{I.10.12.2.1-fr}
\EGAReferenceAlias{I.10.12.2.1-es}{I.10.12.2.1-fr}
\EGAReferenceAlias{I.10.12.2.1.fr}{I.10.12.2.1-fr}
\EGAReferenceAlias{I.10.12.2.1.es}{I.10.12.2.1-fr}
\EGAReferenceAlias{I.10.12.2.1.eq}{I.10.12.2.1-fr}
\EGAReferenceAlias{eq:I.10.12.2.1}{I.10.12.2.1-fr}
\EGAReferenceAlias{I.10.12.3-fr}{I.10.12.3}
\EGAReferenceAlias{I.10.12.3-es}{I.10.12.3}
\EGAReferenceAlias{I.10.12.3.fr}{I.10.12.3}
\EGAReferenceAlias{I.10.12.3.es}{I.10.12.3}
\EGAReferenceAlias{I.10.12.3.eq}{I.10.12.3}
\EGAReferenceAlias{eq:I.10.12.3}{I.10.12.3}
\EGAReferenceAlias{I.10.12.3.1-fr}{I.10.12.3.1}
\EGAReferenceAlias{I.10.12.3.1-es}{I.10.12.3.1}
\EGAReferenceAlias{I.10.12.3.1.fr}{I.10.12.3.1}
\EGAReferenceAlias{I.10.12.3.1.es}{I.10.12.3.1}
\EGAReferenceAlias{I.10.12.3.1.eq}{I.10.12.3.1}
\EGAReferenceAlias{eq:I.10.12.3.1}{I.10.12.3.1}
\EGAReferenceAlias{I.10.12.3.2-fr}{I.10.12.3.2}
\EGAReferenceAlias{I.10.12.3.2-es}{I.10.12.3.2}
\EGAReferenceAlias{I.10.12.3.2.fr}{I.10.12.3.2}
\EGAReferenceAlias{I.10.12.3.2.es}{I.10.12.3.2}
\EGAReferenceAlias{I.10.12.3.2.eq}{I.10.12.3.2}
\EGAReferenceAlias{eq:I.10.12.3.2}{I.10.12.3.2}
\EGAReferenceAlias{I.10.13}{subsection:I.10.13}
\EGAReferenceAlias{I.10.13-fr}{subsection:I.10.13}
\EGAReferenceAlias{I.10.13-es}{subsection:I.10.13}
\EGAReferenceAlias{I.10.13.fr}{subsection:I.10.13}
\EGAReferenceAlias{I.10.13.es}{subsection:I.10.13}
\EGAReferenceAlias{I.10.13.eq}{subsection:I.10.13}
\EGAReferenceAlias{eq:I.10.13}{subsection:I.10.13}
\EGAReferenceAlias{subsection:I.10.13-fr}{subsection:I.10.13}
\EGAReferenceAlias{subsection:I.10.13-es}{subsection:I.10.13}
\EGAReferenceAlias{subsection:I.10.13.fr}{subsection:I.10.13}
\EGAReferenceAlias{subsection:I.10.13.es}{subsection:I.10.13}
\EGAReferenceAlias{subsection:I.10.13.eq}{subsection:I.10.13}
\EGAReferenceAlias{eq:subsection:I.10.13}{subsection:I.10.13}
\EGAReferenceAlias{I.10.13.1-fr}{I.10.13.1}
\EGAReferenceAlias{I.10.13.1-es}{I.10.13.1}
\EGAReferenceAlias{I.10.13.1.fr}{I.10.13.1}
\EGAReferenceAlias{I.10.13.1.es}{I.10.13.1}
\EGAReferenceAlias{I.10.13.1.eq}{I.10.13.1}
\EGAReferenceAlias{eq:I.10.13.1}{I.10.13.1}
\EGAReferenceAlias{I.10.13.2-fr}{I.10.13.2}
\EGAReferenceAlias{I.10.13.2-es}{I.10.13.2}
\EGAReferenceAlias{I.10.13.2.fr}{I.10.13.2}
\EGAReferenceAlias{I.10.13.2.es}{I.10.13.2}
\EGAReferenceAlias{I.10.13.2.eq}{I.10.13.2}
\EGAReferenceAlias{eq:I.10.13.2}{I.10.13.2}
\EGAReferenceAlias{I.10.13.3-fr}{I.10.13.3}
\EGAReferenceAlias{I.10.13.3-es}{I.10.13.3}
\EGAReferenceAlias{I.10.13.3.fr}{I.10.13.3}
\EGAReferenceAlias{I.10.13.3.es}{I.10.13.3}
\EGAReferenceAlias{I.10.13.3.eq}{I.10.13.3}
\EGAReferenceAlias{eq:I.10.13.3}{I.10.13.3}
\EGAReferenceAlias{I.10.13.4-fr}{I.10.13.4}
\EGAReferenceAlias{I.10.13.4-es}{I.10.13.4}
\EGAReferenceAlias{I.10.13.4.fr}{I.10.13.4}
\EGAReferenceAlias{I.10.13.4.es}{I.10.13.4}
\EGAReferenceAlias{I.10.13.4.eq}{I.10.13.4}
\EGAReferenceAlias{eq:I.10.13.4}{I.10.13.4}
\EGAReferenceAlias{I.10.13.5-fr}{I.10.13.5}
\EGAReferenceAlias{I.10.13.5-es}{I.10.13.5}
\EGAReferenceAlias{I.10.13.5.fr}{I.10.13.5}
\EGAReferenceAlias{I.10.13.5.es}{I.10.13.5}
\EGAReferenceAlias{I.10.13.5.eq}{I.10.13.5}
\EGAReferenceAlias{eq:I.10.13.5}{I.10.13.5}
\EGAReferenceAlias{I.10.13.6-fr}{I.10.13.6}
\EGAReferenceAlias{I.10.13.6-es}{I.10.13.6}
\EGAReferenceAlias{I.10.13.6.fr}{I.10.13.6}
\EGAReferenceAlias{I.10.13.6.es}{I.10.13.6}
\EGAReferenceAlias{I.10.13.6.eq}{I.10.13.6}
\EGAReferenceAlias{eq:I.10.13.6}{I.10.13.6}
\EGAReferenceAlias{I.10.14}{subsection:I.10.14}
\EGAReferenceAlias{I.10.14-fr}{subsection:I.10.14}
\EGAReferenceAlias{I.10.14-es}{subsection:I.10.14}
\EGAReferenceAlias{I.10.14.fr}{subsection:I.10.14}
\EGAReferenceAlias{I.10.14.es}{subsection:I.10.14}
\EGAReferenceAlias{I.10.14.eq}{subsection:I.10.14}
\EGAReferenceAlias{eq:I.10.14}{subsection:I.10.14}
\EGAReferenceAlias{subsection:I.10.14-fr}{subsection:I.10.14}
\EGAReferenceAlias{subsection:I.10.14-es}{subsection:I.10.14}
\EGAReferenceAlias{subsection:I.10.14.fr}{subsection:I.10.14}
\EGAReferenceAlias{subsection:I.10.14.es}{subsection:I.10.14}
\EGAReferenceAlias{subsection:I.10.14.eq}{subsection:I.10.14}
\EGAReferenceAlias{eq:subsection:I.10.14}{subsection:I.10.14}
\EGAReferenceAlias{I.10.14.1-fr}{I.10.14.1}
\EGAReferenceAlias{I.10.14.1-es}{I.10.14.1}
\EGAReferenceAlias{I.10.14.1.fr}{I.10.14.1}
\EGAReferenceAlias{I.10.14.1.es}{I.10.14.1}
\EGAReferenceAlias{I.10.14.1.eq}{I.10.14.1}
\EGAReferenceAlias{eq:I.10.14.1}{I.10.14.1}
\EGAReferenceAlias{I.10.14.2-fr}{I.10.14.2}
\EGAReferenceAlias{I.10.14.2-es}{I.10.14.2}
\EGAReferenceAlias{I.10.14.2.fr}{I.10.14.2}
\EGAReferenceAlias{I.10.14.2.es}{I.10.14.2}
\EGAReferenceAlias{I.10.14.2.eq}{I.10.14.2}
\EGAReferenceAlias{eq:I.10.14.2}{I.10.14.2}
\EGAReferenceAlias{I.10.14.3-fr}{I.10.14.3}
\EGAReferenceAlias{I.10.14.3-es}{I.10.14.3}
\EGAReferenceAlias{I.10.14.3.fr}{I.10.14.3}
\EGAReferenceAlias{I.10.14.3.es}{I.10.14.3}
\EGAReferenceAlias{I.10.14.3.eq}{I.10.14.3}
\EGAReferenceAlias{eq:I.10.14.3}{I.10.14.3}
\EGAReferenceAlias{I.10.14.4-fr}{I.10.14.4}
\EGAReferenceAlias{I.10.14.4-es}{I.10.14.4}
\EGAReferenceAlias{I.10.14.4.fr}{I.10.14.4}
\EGAReferenceAlias{I.10.14.4.es}{I.10.14.4}
\EGAReferenceAlias{I.10.14.4.eq}{I.10.14.4}
\EGAReferenceAlias{eq:I.10.14.4}{I.10.14.4}
\EGAReferenceAlias{I.10.14.5-fr}{I.10.14.5}
\EGAReferenceAlias{I.10.14.5-es}{I.10.14.5}
\EGAReferenceAlias{I.10.14.5.fr}{I.10.14.5}
\EGAReferenceAlias{I.10.14.5.es}{I.10.14.5}
\EGAReferenceAlias{I.10.14.5.eq}{I.10.14.5}
\EGAReferenceAlias{eq:I.10.14.5}{I.10.14.5}
\EGAReferenceAlias{I.10.14.6-fr}{I.10.14.6}
\EGAReferenceAlias{I.10.14.6-es}{I.10.14.6}
\EGAReferenceAlias{I.10.14.6.fr}{I.10.14.6}
\EGAReferenceAlias{I.10.14.6.es}{I.10.14.6}
\EGAReferenceAlias{I.10.14.6.eq}{I.10.14.6}
\EGAReferenceAlias{eq:I.10.14.6}{I.10.14.6}
\EGAReferenceAlias{I.10.14.6.1}{I.10.14.6.1-fr}
\EGAReferenceAlias{I.10.14.6.1-es}{I.10.14.6.1-fr}
\EGAReferenceAlias{I.10.14.6.1.fr}{I.10.14.6.1-fr}
\EGAReferenceAlias{I.10.14.6.1.es}{I.10.14.6.1-fr}
\EGAReferenceAlias{I.10.14.6.1.eq}{I.10.14.6.1-fr}
\EGAReferenceAlias{eq:I.10.14.6.1}{I.10.14.6.1-fr}
\EGAReferenceAlias{I.10.14.7-fr}{I.10.14.7}
\EGAReferenceAlias{I.10.14.7-es}{I.10.14.7}
\EGAReferenceAlias{I.10.14.7.fr}{I.10.14.7}
\EGAReferenceAlias{I.10.14.7.es}{I.10.14.7}
\EGAReferenceAlias{I.10.14.7.eq}{I.10.14.7}
\EGAReferenceAlias{eq:I.10.14.7}{I.10.14.7}
\EGAReferenceAlias{I.10.15}{subsection:I.10.15}
\EGAReferenceAlias{I.10.15-fr}{subsection:I.10.15}
\EGAReferenceAlias{I.10.15-es}{subsection:I.10.15}
\EGAReferenceAlias{I.10.15.fr}{subsection:I.10.15}
\EGAReferenceAlias{I.10.15.es}{subsection:I.10.15}
\EGAReferenceAlias{I.10.15.eq}{subsection:I.10.15}
\EGAReferenceAlias{eq:I.10.15}{subsection:I.10.15}
\EGAReferenceAlias{subsection:I.10.15-fr}{subsection:I.10.15}
\EGAReferenceAlias{subsection:I.10.15-es}{subsection:I.10.15}
\EGAReferenceAlias{subsection:I.10.15.fr}{subsection:I.10.15}
\EGAReferenceAlias{subsection:I.10.15.es}{subsection:I.10.15}
\EGAReferenceAlias{subsection:I.10.15.eq}{subsection:I.10.15}
\EGAReferenceAlias{eq:subsection:I.10.15}{subsection:I.10.15}
\EGAReferenceAlias{I.10.15.1-fr}{I.10.15.1}
\EGAReferenceAlias{I.10.15.1-es}{I.10.15.1}
\EGAReferenceAlias{I.10.15.1.fr}{I.10.15.1}
\EGAReferenceAlias{I.10.15.1.es}{I.10.15.1}
\EGAReferenceAlias{I.10.15.1.eq}{I.10.15.1}
\EGAReferenceAlias{eq:I.10.15.1}{I.10.15.1}
\EGAReferenceAlias{I.10.15.2-fr}{I.10.15.2}
\EGAReferenceAlias{I.10.15.2-es}{I.10.15.2}
\EGAReferenceAlias{I.10.15.2.fr}{I.10.15.2}
\EGAReferenceAlias{I.10.15.2.es}{I.10.15.2}
\EGAReferenceAlias{I.10.15.2.eq}{I.10.15.2}
\EGAReferenceAlias{eq:I.10.15.2}{I.10.15.2}
\EGAReferenceAlias{I.10.15.3-fr}{I.10.15.3}
\EGAReferenceAlias{I.10.15.3-es}{I.10.15.3}
\EGAReferenceAlias{I.10.15.3.fr}{I.10.15.3}
\EGAReferenceAlias{I.10.15.3.es}{I.10.15.3}
\EGAReferenceAlias{I.10.15.3.eq}{I.10.15.3}
\EGAReferenceAlias{eq:I.10.15.3}{I.10.15.3}
\EGAReferenceAlias{I.10.15.4-fr}{I.10.15.4}
\EGAReferenceAlias{I.10.15.4-es}{I.10.15.4}
\EGAReferenceAlias{I.10.15.4.fr}{I.10.15.4}
\EGAReferenceAlias{I.10.15.4.es}{I.10.15.4}
\EGAReferenceAlias{I.10.15.4.eq}{I.10.15.4}
\EGAReferenceAlias{eq:I.10.15.4}{I.10.15.4}
\EGAReferenceAlias{I.10.15.5-fr}{I.10.15.5}
\EGAReferenceAlias{I.10.15.5-es}{I.10.15.5}
\EGAReferenceAlias{I.10.15.5.fr}{I.10.15.5}
\EGAReferenceAlias{I.10.15.5.es}{I.10.15.5}
\EGAReferenceAlias{I.10.15.5.eq}{I.10.15.5}
\EGAReferenceAlias{eq:I.10.15.5}{I.10.15.5}
\EGAReferenceAlias{I.10.15.6-fr}{I.10.15.6}
\EGAReferenceAlias{I.10.15.6-es}{I.10.15.6}
\EGAReferenceAlias{I.10.15.6.fr}{I.10.15.6}
\EGAReferenceAlias{I.10.15.6.es}{I.10.15.6}
\EGAReferenceAlias{I.10.15.6.eq}{I.10.15.6}
\EGAReferenceAlias{eq:I.10.15.6}{I.10.15.6}
\EGAReferenceAlias{I.10.15.7-fr}{I.10.15.7}
\EGAReferenceAlias{I.10.15.7-es}{I.10.15.7}
\EGAReferenceAlias{I.10.15.7.fr}{I.10.15.7}
\EGAReferenceAlias{I.10.15.7.es}{I.10.15.7}
\EGAReferenceAlias{I.10.15.7.eq}{I.10.15.7}
\EGAReferenceAlias{eq:I.10.15.7}{I.10.15.7}
\EGAReferenceAlias{index-notation}{section:index-notation}
\EGAReferenceAlias{index-notation-fr}{section:index-notation}
\EGAReferenceAlias{index-notation-es}{section:index-notation}
\EGAReferenceAlias{index-notation.fr}{section:index-notation}
\EGAReferenceAlias{index-notation.es}{section:index-notation}
\EGAReferenceAlias{index-notation.eq}{section:index-notation}
\EGAReferenceAlias{eq:index-notation}{section:index-notation}
\EGAReferenceAlias{section:index-notation-fr}{section:index-notation}
\EGAReferenceAlias{section:index-notation-es}{section:index-notation}
\EGAReferenceAlias{section:index-notation.fr}{section:index-notation}
\EGAReferenceAlias{section:index-notation.es}{section:index-notation}
\EGAReferenceAlias{section:index-notation.eq}{section:index-notation}
\EGAReferenceAlias{eq:section:index-notation}{section:index-notation}
\EGAReferenceAlias{index-notation-chapter0}{section:index-notation-chapter0}
\EGAReferenceAlias{index-notation-chapter0-fr}{section:index-notation-chapter0}
\EGAReferenceAlias{index-notation-chapter0-es}{section:index-notation-chapter0}
\EGAReferenceAlias{index-notation-chapter0.fr}{section:index-notation-chapter0}
\EGAReferenceAlias{index-notation-chapter0.es}{section:index-notation-chapter0}
\EGAReferenceAlias{index-notation-chapter0.eq}{section:index-notation-chapter0}
\EGAReferenceAlias{eq:index-notation-chapter0}{section:index-notation-chapter0}
\EGAReferenceAlias{section:index-notation-chapter0-fr}{section:index-notation-chapter0}
\EGAReferenceAlias{section:index-notation-chapter0-es}{section:index-notation-chapter0}
\EGAReferenceAlias{section:index-notation-chapter0.fr}{section:index-notation-chapter0}
\EGAReferenceAlias{section:index-notation-chapter0.es}{section:index-notation-chapter0}
\EGAReferenceAlias{section:index-notation-chapter0.eq}{section:index-notation-chapter0}
\EGAReferenceAlias{eq:section:index-notation-chapter0}{section:index-notation-chapter0}
\EGAReferenceAlias{index-notation-chapterI}{section:index-notation-chapterI}
\EGAReferenceAlias{index-notation-chapterI-fr}{section:index-notation-chapterI}
\EGAReferenceAlias{index-notation-chapterI-es}{section:index-notation-chapterI}
\EGAReferenceAlias{index-notation-chapterI.fr}{section:index-notation-chapterI}
\EGAReferenceAlias{index-notation-chapterI.es}{section:index-notation-chapterI}
\EGAReferenceAlias{index-notation-chapterI.eq}{section:index-notation-chapterI}
\EGAReferenceAlias{eq:index-notation-chapterI}{section:index-notation-chapterI}
\EGAReferenceAlias{section:index-notation-chapterI-fr}{section:index-notation-chapterI}
\EGAReferenceAlias{section:index-notation-chapterI-es}{section:index-notation-chapterI}
\EGAReferenceAlias{section:index-notation-chapterI.fr}{section:index-notation-chapterI}
\EGAReferenceAlias{section:index-notation-chapterI.es}{section:index-notation-chapterI}
\EGAReferenceAlias{section:index-notation-chapterI.eq}{section:index-notation-chapterI}
\EGAReferenceAlias{eq:section:index-notation-chapterI}{section:index-notation-chapterI}
\EGAReferenceAlias{index-terminology}{section:index-terminology}
\EGAReferenceAlias{index-terminology-fr}{section:index-terminology}
\EGAReferenceAlias{index-terminology-es}{section:index-terminology}
\EGAReferenceAlias{index-terminology.fr}{section:index-terminology}
\EGAReferenceAlias{index-terminology.es}{section:index-terminology}
\EGAReferenceAlias{index-terminology.eq}{section:index-terminology}
\EGAReferenceAlias{eq:index-terminology}{section:index-terminology}
\EGAReferenceAlias{section:index-terminology-fr}{section:index-terminology}
\EGAReferenceAlias{section:index-terminology-es}{section:index-terminology}
\EGAReferenceAlias{section:index-terminology.fr}{section:index-terminology}
\EGAReferenceAlias{section:index-terminology.es}{section:index-terminology}
\EGAReferenceAlias{section:index-terminology.eq}{section:index-terminology}
\EGAReferenceAlias{eq:section:index-terminology}{section:index-terminology}
\EGAReferenceAlias{ega2:addendum:subsection:I.1.8}{ega2:addendum:subsection:I.1.8-fr}
\EGAReferenceAlias{ega2:addendum:subsection:I.1.8-es}{ega2:addendum:subsection:I.1.8-fr}
\EGAReferenceAlias{ega2:addendum:subsection:I.1.8.fr}{ega2:addendum:subsection:I.1.8-fr}
\EGAReferenceAlias{ega2:addendum:subsection:I.1.8.es}{ega2:addendum:subsection:I.1.8-fr}
\EGAReferenceAlias{ega2:addendum:subsection:I.1.8.eq}{ega2:addendum:subsection:I.1.8-fr}
\EGAReferenceAlias{eq:ega2:addendum:subsection:I.1.8}{ega2:addendum:subsection:I.1.8-fr}
\EGAReferenceAlias{ega2:addendum:I.1.8.1}{ega2:addendum:I.1.8.1-fr}
\EGAReferenceAlias{ega2:addendum:I.1.8.1-es}{ega2:addendum:I.1.8.1-fr}
\EGAReferenceAlias{ega2:addendum:I.1.8.1.fr}{ega2:addendum:I.1.8.1-fr}
\EGAReferenceAlias{ega2:addendum:I.1.8.1.es}{ega2:addendum:I.1.8.1-fr}
\EGAReferenceAlias{ega2:addendum:I.1.8.1.eq}{ega2:addendum:I.1.8.1-fr}
\EGAReferenceAlias{eq:ega2:addendum:I.1.8.1}{ega2:addendum:I.1.8.1-fr}
\EGAReferenceAlias{ega2:addendum:I.1.8.1.1}{ega2:addendum:I.1.8.1.1-fr}
\EGAReferenceAlias{ega2:addendum:I.1.8.1.1-es}{ega2:addendum:I.1.8.1.1-fr}
\EGAReferenceAlias{ega2:addendum:I.1.8.1.1.fr}{ega2:addendum:I.1.8.1.1-fr}
\EGAReferenceAlias{ega2:addendum:I.1.8.1.1.es}{ega2:addendum:I.1.8.1.1-fr}
\EGAReferenceAlias{ega2:addendum:I.1.8.1.1.eq}{ega2:addendum:I.1.8.1.1-fr}
\EGAReferenceAlias{eq:ega2:addendum:I.1.8.1.1}{ega2:addendum:I.1.8.1.1-fr}
\EGAReferenceAlias{ega2:addendum:I.1.8.2}{ega2:addendum:I.1.8.2-fr}
\EGAReferenceAlias{ega2:addendum:I.1.8.2-es}{ega2:addendum:I.1.8.2-fr}
\EGAReferenceAlias{ega2:addendum:I.1.8.2.fr}{ega2:addendum:I.1.8.2-fr}
\EGAReferenceAlias{ega2:addendum:I.1.8.2.es}{ega2:addendum:I.1.8.2-fr}
\EGAReferenceAlias{ega2:addendum:I.1.8.2.eq}{ega2:addendum:I.1.8.2-fr}
\EGAReferenceAlias{eq:ega2:addendum:I.1.8.2}{ega2:addendum:I.1.8.2-fr}
\EGAReferenceAlias{ega2:addendum:I.1.8.2.1}{ega2:addendum:I.1.8.2.1-fr}
\EGAReferenceAlias{ega2:addendum:I.1.8.2.1-es}{ega2:addendum:I.1.8.2.1-fr}
\EGAReferenceAlias{ega2:addendum:I.1.8.2.1.fr}{ega2:addendum:I.1.8.2.1-fr}
\EGAReferenceAlias{ega2:addendum:I.1.8.2.1.es}{ega2:addendum:I.1.8.2.1-fr}
\EGAReferenceAlias{ega2:addendum:I.1.8.2.1.eq}{ega2:addendum:I.1.8.2.1-fr}
\EGAReferenceAlias{eq:ega2:addendum:I.1.8.2.1}{ega2:addendum:I.1.8.2.1-fr}
\EGAReferenceAlias{ega2:addendum:I.1.8.2.2}{ega2:addendum:I.1.8.2.2-fr}
\EGAReferenceAlias{ega2:addendum:I.1.8.2.2-es}{ega2:addendum:I.1.8.2.2-fr}
\EGAReferenceAlias{ega2:addendum:I.1.8.2.2.fr}{ega2:addendum:I.1.8.2.2-fr}
\EGAReferenceAlias{ega2:addendum:I.1.8.2.2.es}{ega2:addendum:I.1.8.2.2-fr}
\EGAReferenceAlias{ega2:addendum:I.1.8.2.2.eq}{ega2:addendum:I.1.8.2.2-fr}
\EGAReferenceAlias{eq:ega2:addendum:I.1.8.2.2}{ega2:addendum:I.1.8.2.2-fr}
\EGAReferenceAlias{ega2:addendum:I.1.8.3}{ega2:addendum:I.1.8.3-fr}
\EGAReferenceAlias{ega2:addendum:I.1.8.3-es}{ega2:addendum:I.1.8.3-fr}
\EGAReferenceAlias{ega2:addendum:I.1.8.3.fr}{ega2:addendum:I.1.8.3-fr}
\EGAReferenceAlias{ega2:addendum:I.1.8.3.es}{ega2:addendum:I.1.8.3-fr}
\EGAReferenceAlias{ega2:addendum:I.1.8.3.eq}{ega2:addendum:I.1.8.3-fr}
\EGAReferenceAlias{eq:ega2:addendum:I.1.8.3}{ega2:addendum:I.1.8.3-fr}
\EGAReferenceAlias{ega2:addendum:I.1.8.4}{ega2:addendum:I.1.8.4-fr}
\EGAReferenceAlias{ega2:addendum:I.1.8.4-es}{ega2:addendum:I.1.8.4-fr}
\EGAReferenceAlias{ega2:addendum:I.1.8.4.fr}{ega2:addendum:I.1.8.4-fr}
\EGAReferenceAlias{ega2:addendum:I.1.8.4.es}{ega2:addendum:I.1.8.4-fr}
\EGAReferenceAlias{ega2:addendum:I.1.8.4.eq}{ega2:addendum:I.1.8.4-fr}
\EGAReferenceAlias{eq:ega2:addendum:I.1.8.4}{ega2:addendum:I.1.8.4-fr}
\EGAReferenceAlias{ega2:addendum:I.1.8.5}{ega2:addendum:I.1.8.5-fr}
\EGAReferenceAlias{ega2:addendum:I.1.8.5-es}{ega2:addendum:I.1.8.5-fr}
\EGAReferenceAlias{ega2:addendum:I.1.8.5.fr}{ega2:addendum:I.1.8.5-fr}
\EGAReferenceAlias{ega2:addendum:I.1.8.5.es}{ega2:addendum:I.1.8.5-fr}
\EGAReferenceAlias{ega2:addendum:I.1.8.5.eq}{ega2:addendum:I.1.8.5-fr}
\EGAReferenceAlias{eq:ega2:addendum:I.1.8.5}{ega2:addendum:I.1.8.5-fr}
\EGAReferenceAlias{ega2:addendum:I.1.8.5.1}{ega2:addendum:I.1.8.5.1-fr}
\EGAReferenceAlias{ega2:addendum:I.1.8.5.1-es}{ega2:addendum:I.1.8.5.1-fr}
\EGAReferenceAlias{ega2:addendum:I.1.8.5.1.fr}{ega2:addendum:I.1.8.5.1-fr}
\EGAReferenceAlias{ega2:addendum:I.1.8.5.1.es}{ega2:addendum:I.1.8.5.1-fr}
\EGAReferenceAlias{ega2:addendum:I.1.8.5.1.eq}{ega2:addendum:I.1.8.5.1-fr}
\EGAReferenceAlias{eq:ega2:addendum:I.1.8.5.1}{ega2:addendum:I.1.8.5.1-fr}
\EGAReferenceAlias{ega2:addendum:I.1.8.6}{ega2:addendum:I.1.8.6-fr}
\EGAReferenceAlias{ega2:addendum:I.1.8.6-es}{ega2:addendum:I.1.8.6-fr}
\EGAReferenceAlias{ega2:addendum:I.1.8.6.fr}{ega2:addendum:I.1.8.6-fr}
\EGAReferenceAlias{ega2:addendum:I.1.8.6.es}{ega2:addendum:I.1.8.6-fr}
\EGAReferenceAlias{ega2:addendum:I.1.8.6.eq}{ega2:addendum:I.1.8.6-fr}
\EGAReferenceAlias{eq:ega2:addendum:I.1.8.6}{ega2:addendum:I.1.8.6-fr}
\EGAReferenceAlias{ega2:addendum:I.1.8.6.1}{ega2:addendum:I.1.8.6.1-fr}
\EGAReferenceAlias{ega2:addendum:I.1.8.6.1-es}{ega2:addendum:I.1.8.6.1-fr}
\EGAReferenceAlias{ega2:addendum:I.1.8.6.1.fr}{ega2:addendum:I.1.8.6.1-fr}
\EGAReferenceAlias{ega2:addendum:I.1.8.6.1.es}{ega2:addendum:I.1.8.6.1-fr}
\EGAReferenceAlias{ega2:addendum:I.1.8.6.1.eq}{ega2:addendum:I.1.8.6.1-fr}
\EGAReferenceAlias{eq:ega2:addendum:I.1.8.6.1}{ega2:addendum:I.1.8.6.1-fr}
\EGAReferenceAlias{ega2:addendum:I.1.8.7}{ega2:addendum:I.1.8.7-fr}
\EGAReferenceAlias{ega2:addendum:I.1.8.7-es}{ega2:addendum:I.1.8.7-fr}
\EGAReferenceAlias{ega2:addendum:I.1.8.7.fr}{ega2:addendum:I.1.8.7-fr}
\EGAReferenceAlias{ega2:addendum:I.1.8.7.es}{ega2:addendum:I.1.8.7-fr}
\EGAReferenceAlias{ega2:addendum:I.1.8.7.eq}{ega2:addendum:I.1.8.7-fr}
\EGAReferenceAlias{eq:ega2:addendum:I.1.8.7}{ega2:addendum:I.1.8.7-fr}
\EGAReferenceAlias{ega2:addendum:I.1.8.7.1}{ega2:addendum:I.1.8.7.1-fr}
\EGAReferenceAlias{ega2:addendum:I.1.8.7.1-es}{ega2:addendum:I.1.8.7.1-fr}
\EGAReferenceAlias{ega2:addendum:I.1.8.7.1.fr}{ega2:addendum:I.1.8.7.1-fr}
\EGAReferenceAlias{ega2:addendum:I.1.8.7.1.es}{ega2:addendum:I.1.8.7.1-fr}
\EGAReferenceAlias{ega2:addendum:I.1.8.7.1.eq}{ega2:addendum:I.1.8.7.1-fr}
\EGAReferenceAlias{eq:ega2:addendum:I.1.8.7.1}{ega2:addendum:I.1.8.7.1-fr}
\EGAReferenceAlias{ega2:addendum:I.1.8.8}{ega2:addendum:I.1.8.8-fr}
\EGAReferenceAlias{ega2:addendum:I.1.8.8-es}{ega2:addendum:I.1.8.8-fr}
\EGAReferenceAlias{ega2:addendum:I.1.8.8.fr}{ega2:addendum:I.1.8.8-fr}
\EGAReferenceAlias{ega2:addendum:I.1.8.8.es}{ega2:addendum:I.1.8.8-fr}
\EGAReferenceAlias{ega2:addendum:I.1.8.8.eq}{ega2:addendum:I.1.8.8-fr}
\EGAReferenceAlias{eq:ega2:addendum:I.1.8.8}{ega2:addendum:I.1.8.8-fr}
\EGAReferenceAlias{ega2:addendum:I.1.8.8.1}{ega2:addendum:I.1.8.8.1-fr}
\EGAReferenceAlias{ega2:addendum:I.1.8.8.1-es}{ega2:addendum:I.1.8.8.1-fr}
\EGAReferenceAlias{ega2:addendum:I.1.8.8.1.fr}{ega2:addendum:I.1.8.8.1-fr}
\EGAReferenceAlias{ega2:addendum:I.1.8.8.1.es}{ega2:addendum:I.1.8.8.1-fr}
\EGAReferenceAlias{ega2:addendum:I.1.8.8.1.eq}{ega2:addendum:I.1.8.8.1-fr}
\EGAReferenceAlias{eq:ega2:addendum:I.1.8.8.1}{ega2:addendum:I.1.8.8.1-fr}
\EGAReferenceAlias{ega2:addendum:I.1.8.8.2}{ega2:addendum:I.1.8.8.2-fr}
\EGAReferenceAlias{ega2:addendum:I.1.8.8.2-es}{ega2:addendum:I.1.8.8.2-fr}
\EGAReferenceAlias{ega2:addendum:I.1.8.8.2.fr}{ega2:addendum:I.1.8.8.2-fr}
\EGAReferenceAlias{ega2:addendum:I.1.8.8.2.es}{ega2:addendum:I.1.8.8.2-fr}
\EGAReferenceAlias{ega2:addendum:I.1.8.8.2.eq}{ega2:addendum:I.1.8.8.2-fr}
\EGAReferenceAlias{eq:ega2:addendum:I.1.8.8.2}{ega2:addendum:I.1.8.8.2-fr}
\EGAReferenceAlias{ega2:addendum:I.1.8.9}{ega2:addendum:I.1.8.9-fr}
\EGAReferenceAlias{ega2:addendum:I.1.8.9-es}{ega2:addendum:I.1.8.9-fr}
\EGAReferenceAlias{ega2:addendum:I.1.8.9.fr}{ega2:addendum:I.1.8.9-fr}
\EGAReferenceAlias{ega2:addendum:I.1.8.9.es}{ega2:addendum:I.1.8.9-fr}
\EGAReferenceAlias{ega2:addendum:I.1.8.9.eq}{ega2:addendum:I.1.8.9-fr}
\EGAReferenceAlias{eq:ega2:addendum:I.1.8.9}{ega2:addendum:I.1.8.9-fr}
\EGAReferenceAlias{ega2:addendum:I.1.8.9.1}{ega2:addendum:I.1.8.9.1-fr}
\EGAReferenceAlias{ega2:addendum:I.1.8.9.1-es}{ega2:addendum:I.1.8.9.1-fr}
\EGAReferenceAlias{ega2:addendum:I.1.8.9.1.fr}{ega2:addendum:I.1.8.9.1-fr}
\EGAReferenceAlias{ega2:addendum:I.1.8.9.1.es}{ega2:addendum:I.1.8.9.1-fr}
\EGAReferenceAlias{ega2:addendum:I.1.8.9.1.eq}{ega2:addendum:I.1.8.9.1-fr}
\EGAReferenceAlias{eq:ega2:addendum:I.1.8.9.1}{ega2:addendum:I.1.8.9.1-fr}
\EGAReferenceAlias{ega2:addendum:I.1.8.10}{ega2:addendum:I.1.8.10-fr}
\EGAReferenceAlias{ega2:addendum:I.1.8.10-es}{ega2:addendum:I.1.8.10-fr}
\EGAReferenceAlias{ega2:addendum:I.1.8.10.fr}{ega2:addendum:I.1.8.10-fr}
\EGAReferenceAlias{ega2:addendum:I.1.8.10.es}{ega2:addendum:I.1.8.10-fr}
\EGAReferenceAlias{ega2:addendum:I.1.8.10.eq}{ega2:addendum:I.1.8.10-fr}
\EGAReferenceAlias{eq:ega2:addendum:I.1.8.10}{ega2:addendum:I.1.8.10-fr}
\EGAReferenceAlias{ega2:addendum:I.1.8.11}{ega2:addendum:I.1.8.11-fr}
\EGAReferenceAlias{ega2:addendum:I.1.8.11-es}{ega2:addendum:I.1.8.11-fr}
\EGAReferenceAlias{ega2:addendum:I.1.8.11.fr}{ega2:addendum:I.1.8.11-fr}
\EGAReferenceAlias{ega2:addendum:I.1.8.11.es}{ega2:addendum:I.1.8.11-fr}
\EGAReferenceAlias{ega2:addendum:I.1.8.11.eq}{ega2:addendum:I.1.8.11-fr}
\EGAReferenceAlias{eq:ega2:addendum:I.1.8.11}{ega2:addendum:I.1.8.11-fr}
\EGAReferenceAlias{ega2:addendum:I.3.2.9}{ega2:addendum:I.3.2.9-fr}
\EGAReferenceAlias{ega2:addendum:I.3.2.9-es}{ega2:addendum:I.3.2.9-fr}
\EGAReferenceAlias{ega2:addendum:I.3.2.9.fr}{ega2:addendum:I.3.2.9-fr}
\EGAReferenceAlias{ega2:addendum:I.3.2.9.es}{ega2:addendum:I.3.2.9-fr}
\EGAReferenceAlias{ega2:addendum:I.3.2.9.eq}{ega2:addendum:I.3.2.9-fr}
\EGAReferenceAlias{eq:ega2:addendum:I.3.2.9}{ega2:addendum:I.3.2.9-fr}
\EGAReferenceAlias{ega2:addendum:I.7.3.8.1}{ega2:addendum:I.7.3.8.1-fr}
\EGAReferenceAlias{ega2:addendum:I.7.3.8.1-es}{ega2:addendum:I.7.3.8.1-fr}
\EGAReferenceAlias{ega2:addendum:I.7.3.8.1.fr}{ega2:addendum:I.7.3.8.1-fr}
\EGAReferenceAlias{ega2:addendum:I.7.3.8.1.es}{ega2:addendum:I.7.3.8.1-fr}
\EGAReferenceAlias{ega2:addendum:I.7.3.8.1.eq}{ega2:addendum:I.7.3.8.1-fr}
\EGAReferenceAlias{eq:ega2:addendum:I.7.3.8.1}{ega2:addendum:I.7.3.8.1-fr}
\EGAReferenceAlias{0.4.3.3}{canon-ref.0.4.3.3-fr}
\EGAReferenceAlias{0.4.3.3-fr}{canon-ref.0.4.3.3-fr}
\EGAReferenceAlias{0.4.3.3-es}{canon-ref.0.4.3.3-fr}
\EGAReferenceAlias{0.4.3.3.fr}{canon-ref.0.4.3.3-fr}
\EGAReferenceAlias{0.4.3.3.es}{canon-ref.0.4.3.3-fr}
\EGAReferenceAlias{0.4.3.3.eq}{canon-ref.0.4.3.3-fr}
\EGAReferenceAlias{eq:0.4.3.3}{canon-ref.0.4.3.3-fr}
\EGAReferenceAlias{0.5.4.3}{canon-ref.0.5.4.3-fr}
\EGAReferenceAlias{0.5.4.3-fr}{canon-ref.0.5.4.3-fr}
\EGAReferenceAlias{0.5.4.3-es}{canon-ref.0.5.4.3-fr}
\EGAReferenceAlias{0.5.4.3.fr}{canon-ref.0.5.4.3-fr}
\EGAReferenceAlias{0.5.4.3.es}{canon-ref.0.5.4.3-fr}
\EGAReferenceAlias{0.5.4.3.eq}{canon-ref.0.5.4.3-fr}
\EGAReferenceAlias{eq:0.5.4.3}{canon-ref.0.5.4.3-fr}
\EGAReferenceAlias{0.5.4.5}{canon-ref.0.5.4.5-fr}
\EGAReferenceAlias{0.5.4.5-fr}{canon-ref.0.5.4.5-fr}
\EGAReferenceAlias{0.5.4.5-es}{canon-ref.0.5.4.5-fr}
\EGAReferenceAlias{0.5.4.5.fr}{canon-ref.0.5.4.5-fr}
\EGAReferenceAlias{0.5.4.5.es}{canon-ref.0.5.4.5-fr}
\EGAReferenceAlias{0.5.4.5.eq}{canon-ref.0.5.4.5-fr}
\EGAReferenceAlias{eq:0.5.4.5}{canon-ref.0.5.4.5-fr}
\EGAReferenceAlias{0.5.4.7}{canon-ref.0.5.4.7-fr}
\EGAReferenceAlias{0.5.4.7-fr}{canon-ref.0.5.4.7-fr}
\EGAReferenceAlias{0.5.4.7-es}{canon-ref.0.5.4.7-fr}
\EGAReferenceAlias{0.5.4.7.fr}{canon-ref.0.5.4.7-fr}
\EGAReferenceAlias{0.5.4.7.es}{canon-ref.0.5.4.7-fr}
\EGAReferenceAlias{0.5.4.7.eq}{canon-ref.0.5.4.7-fr}
\EGAReferenceAlias{eq:0.5.4.7}{canon-ref.0.5.4.7-fr}
\EGAReferenceAlias{0.6.1.2}{canon-ref.0.6.1.2-fr}
\EGAReferenceAlias{0.6.1.2-fr}{canon-ref.0.6.1.2-fr}
\EGAReferenceAlias{0.6.1.2-es}{canon-ref.0.6.1.2-fr}
\EGAReferenceAlias{0.6.1.2.fr}{canon-ref.0.6.1.2-fr}
\EGAReferenceAlias{0.6.1.2.es}{canon-ref.0.6.1.2-fr}
\EGAReferenceAlias{0.6.1.2.eq}{canon-ref.0.6.1.2-fr}
\EGAReferenceAlias{eq:0.6.1.2}{canon-ref.0.6.1.2-fr}
% BEGIN tree/worktree/es/CURRENT/EGA_I_front/frontmatter-es.tex
\begin{titlepage}
\centering
{\Large INSTITUTO\par}
\vspace{0.4em}
{\Large DE ALTOS ESTUDIOS\par}
\vspace{0.4em}
{\Large CIENTÍFICOS\par}
\vfill
{\Large\bfseries ELEMENTOS DE GEOMETRÍA ALGEBRAICA\par}
\vspace{1.5em}
{\large\itshape por A. GROTHENDIECK\par}
\vspace{0.7em}
{\itshape Redactados con la colaboración de J. DIEUDONNÉ\par}
\vspace{2em}
I
\vspace{1.5em}
{\Large\bfseries EL LENGUAJE DE LOS ESQUEMAS\par}
\vfill
{\Large 1960\par}
\vspace{1.2em}
{\Large PUBLICACIONES MATEMÁTICAS, N.\textsuperscript{o} 4\par}
\vspace{0.5em}
{\large 5, ROND-POINT BUGEAUD --- PARÍS (XVI)\par}
\end{titlepage}
\thispagestyle{empty}
\vspace*{\fill}
\begin{center}
\fbox{%
  \begin{minipage}{0.50\textwidth}
  \centering
  DEPÓSITO LEGAL
  \smallskip
  1.\textsuperscript{a} edición \dotfill 3.\textsuperscript{er} trimestre de 1960
  \smallskip
  TODOS LOS DERECHOS
  reservados para todos los países
  \smallskip
  \copyright\ 1960, \emph{Institut des Hautes Études Scientifiques}
  \end{minipage}%
}
\end{center}
\vspace*{\fill}
\clearpage

% END tree/worktree/es/CURRENT/EGA_I_front/frontmatter-es.tex
% BEGIN tree/worktree/es/CURRENT/EGA_I_front/INTRODUCTION_ES_R570.tex
% Introducción íntegra de EGA I, pp.5--9; no es el tomo completo. Fragmento editable para el proyecto del tomo.

% Unidad R565: introducción, p. 5, dedicatoria y requisitos a)--d).
% Continuación: intro-fr.tex, párrafo «Ce n'est pas le lieu...».
\oldpage[I]{5}
\section*{Introducción}

\begin{center}
\emph{A Oscar Zariski y André Weil.}
\end{center}


Este trabajo, y los numerosos trabajos que han de seguirlo, están destinados a formar un tratado sobre los fundamentos de la geometría algebraica.
En principio no presuponen ningún conocimiento particular de esta disciplina; incluso se ha comprobado que tal conocimiento, pese a sus evidentes ventajas, podía a veces resultar perjudicial para quien desea familiarizarse con el punto de vista y las técnicas expuestos aquí, por el hábito de adoptar de manera demasiado exclusiva el punto de vista birracional que dicho conocimiento lleva consigo.
En cambio, supondremos que el lector conoce bien los temas siguientes:

\begin{enumerate}
  \item[\textit{a)}] El \emph{álgebra conmutativa}, tal como se expone, por ejemplo, en los volúmenes de los \emph{Éléments} de N.~Bourbaki que están en preparación (y, hasta que aparezcan esos volúmenes, en Samuel-Zariski~\cite{I-13} y Samuel~\cite{I-11},~\cite{I-12}).
  \item[\textit{b)}] El \emph{álgebra homológica}, para la cual remitimos a Cartan-Eilenberg~\cite{I-2} (citado como~(M)) y Godement~\cite{I-4} (citado como~(G)), así como al reciente artículo de A.~Grothendieck~\cite{I-6} (citado como~(T)).
  \item[\textit{c)}] La \emph{teoría de haces}, cuyas referencias principales serán~(G) y~(T); esta teoría proporciona el lenguaje indispensable para interpretar en términos «geométricos» las nociones esenciales del álgebra conmutativa y para «globalizarlas».
  \item[\textit{d)}] Por último, será útil que el lector esté familiarizado en cierta medida con el \emph{lenguaje funtorial}, que se empleará constantemente en este tratado y para el cual podrá consultar~(M),~(G) y, sobre todo,~(T); los principios de este lenguaje y los principales resultados de la teoría general de los funtores se expondrán con mayor detalle en una obra que los autores de este tratado están preparando.
\end{enumerate}

\bigskip
\begin{center}*\quad*\quad*\end{center}
\bigskip

% Unidad R566: introducción, líneas francesas 24--35, pp. 5--6.
% Traducción: OpenAI Codex, GPT-6.1 Sol, esfuerzo Ultra.
No corresponde en esta Introducción ofrecer una descripción más o menos sumaria del punto de vista de los «esquemas» en geometría algebraica, ni la larga lista de razones que hicieron necesaria su adopción, en particular la aceptación sistemática de elementos nilpotentes en los anillos locales de las «variedades» que consideramos (lo que necesariamente relega a un segundo plano la noción de aplicación racional, en favor de la de aplicación regular o «morfismo»).
Este tratado se propone precisamente desarrollar de manera sistemática el lenguaje de los «esquemas» y demostrará, así lo esperamos, su necesidad.
Aunque sería fácil hacerlo, tampoco intentaremos ofrecer aquí

\oldpage[I]{6}
una introducción «intuitiva» a las nociones desarrolladas en el capítulo primero.
El lector que desee tener una visión preliminar de las materias de este tratado podrá consultar la conferencia de A.~Grothendieck en el Congreso Internacional de Matemáticos celebrado en Edimburgo en 1958~\cite{I-7}, así como la exposición~\cite{I-8} del mismo autor.
El trabajo~\cite{I-14} (citado como~(FAC)) de J.-P.~Serre también puede considerarse una exposición intermedia entre el punto de vista clásico y el de los esquemas en geometría algebraica; por ello, su lectura puede constituir una excelente preparación para la de nuestros \emph{Éléments}.

\bigskip
\begin{center}*\quad*\quad*\end{center}
\bigskip

% Unidad R567: introducción, líneas francesas 37--68, pp. 6--7.
% Traducción: OpenAI Codex, GPT-6.1 Sol, esfuerzo Ultra.
A título informativo, damos a continuación el plan general previsto para este tratado, sujeto, por lo demás, a modificaciones posteriores, sobre todo en lo que respecta a los últimos capítulos:

\medskip
\noindent\begin{tabular}{@{}p{0.22\linewidth}p{0.74\linewidth}@{}}
  Capítulo primero. & --- El lenguaje de los esquemas.\\
  --- II. & --- Estudio global elemental de algunas clases de morfismos.\\
  --- III. & --- Cohomología de los haces algebraicos coherentes. Aplicaciones.\\
  --- IV. & --- Estudio local de los morfismos.\\
  --- V. & --- Procedimientos elementales de construcción de esquemas.\\
  --- VI. & --- Técnica de descenso. Método general de construcción de esquemas.\\
  --- VII. & --- Esquemas de grupos, espacios fibrados principales.\\
  --- VIII. & --- Estudio diferencial de los espacios fibrados.\\
  --- IX. & --- El grupo fundamental.\\
  --- X. & --- Residuos y dualidad.\\
  --- XI. & --- Teorías de intersección, clases de Chern, teorema de Riemann-Roch.\\
  --- XII. & --- Esquemas abelianos y esquemas de Picard.\\
  --- XIII. & --- Cohomología de Weil.\\
\end{tabular}
\medskip

En principio, todos los capítulos se consideran abiertos a ampliaciones, y siempre podrán añadírseles párrafos suplementarios más adelante; estos párrafos aparecerán en fascículos separados, para reducir los inconvenientes del modo de publicación adoptado.
Cuando uno de estos párrafos esté previsto o en preparación al publicarse un capítulo, se mencionará en el sumario de ese capítulo, aunque determinadas prioridades de urgencia hagan que su publicación efectiva tenga lugar mucho más tarde.
Para comodidad del lector, ofrecemos en un «Capítulo~0» diversos complementos de álgebra conmutativa, álgebra homológica y teoría de haces, utilizados en los capítulos de este tratado, que son más o menos conocidos, pero para los cuales no ha sido posible dar referencias cómodas de consultar.
Se recomienda al lector acudir al capítulo~0 solo mientras lee el tratado propiamente dicho, y en la medida en que los resultados a los que remitimos no le sean

\oldpage[I]{7}
suficientemente familiares.
Pensamos, además, que así la lectura de este tratado podrá ser para el principiante un buen método de familiarizarse con el álgebra conmutativa y el álgebra homológica, cuyo estudio, cuando no se acompaña de aplicaciones tangibles, resulta tedioso, e incluso deprimente, para un número considerable de personas.

\bigskip
\begin{center}*\quad*\quad*\end{center}
\bigskip

% Unidad R568: introducción, líneas francesas 70--93, pp. 7--8.
% Traducción: OpenAI Codex, GPT-6.1 Sol, esfuerzo Ultra.
Queda fuera de nuestra competencia ofrecer en esta introducción una reseña histórica, siquiera sumaria, de las nociones y los resultados expuestos.
El texto solo contendrá las referencias que se consideren especialmente útiles para su comprensión, y solo indicaremos el origen de los resultados más importantes.
Al menos en su forma, los temas tratados en nuestra obra son bastante nuevos; esto explica la escasez de referencias a los padres de la geometría algebraica del siglo~\textsc{xix} y de comienzos del siglo~\textsc{xx}, cuyos trabajos solo conocemos de oídas.
Conviene, sin embargo, decir aquí unas palabras sobre las obras que han influido más directamente en los autores y han contribuido al desarrollo del punto de vista de los esquemas.
Hay que citar, en primer lugar, el trabajo fundamental (FAC) de J.-P.~Serre, que ha servido de introducción a la geometría algebraica a más de un joven interesado en ella (entre ellos, uno de los autores de este tratado), desalentado por la aridez de los clásicos \emph{Foundations} de A.~Weil~\cite{I-18}.
Allí se demuestra por primera vez que la «topología de Zariski» de una variedad algebraica «abstracta» es plenamente adecuada para aplicarle ciertas técnicas de la topología algebraica y dar lugar, en particular, a una teoría cohomológica.
Además, la definición de variedad algebraica que allí se da es la que se presta de manera más natural a la extensión de esta noción que desarrollamos aquí\footnote{Como nos ha señalado J.-P.~Serre, conviene recordar que la idea de definir la estructura de variedad mediante un haz de anillos se debe a H.~Cartan, quien tomó esta idea como punto de partida de su teoría de los espacios analíticos.
Por supuesto, al igual que en geometría algebraica, sería importante admitir, en «geometría analítica», los elementos nilpotentes en los anillos locales de los espacios analíticos.
H.~Grauert ha abordado recientemente esta extensión de la definición de H.~Cartan y J.-P.~Serre~\cite{I-5}, y cabe esperar que pronto aparezca una exposición sistemática de geometría analítica en este marco general.
Es evidente, por otra parte, que las nociones y técnicas desarrolladas en este tratado siguen teniendo sentido en geometría analítica, aunque en esta última teoría hay que prever dificultades técnicas más considerables.
Cabe prever que la geometría algebraica, gracias a la sencillez de sus métodos, podrá servir como una especie de modelo formal para futuros desarrollos de la teoría de los espacios analíticos.}.
El propio Serre había observado, además, que la teoría cohomológica de las variedades algebraicas afines podía trasladarse sin dificultad sustituyendo las álgebras afines sobre un cuerpo por anillos conmutativos cualesquiera.
Por tanto, los capítulos~I y~II de este tratado y los dos primeros párrafos del capítulo~III pueden considerarse, en lo esencial, transposiciones fáciles, a este marco ampliado, de los principales resultados de (FAC) y de un artículo posterior del mismo autor~\cite{I-15}.
También hemos obtenido gran provecho del \emph{Séminaire de Géométrie algébrique} de C.~Chevalley~\cite{I-1}; en particular, el uso sistemático de los «conjuntos construibles» introducidos por él ha resultado muy útil en la teoría de esquemas (cf. cap.~IV).
También hemos tomado de él el estudio de los morfismos desde el punto

\oldpage[I]{8}
de vista de la dimensión (cap.~IV), que se traslada al marco de los esquemas sin cambios apreciables.
Conviene señalar, además, que la noción de «esquemas de anillos locales», introducida por Chevalley, se presta de manera natural a una extensión de la geometría algebraica (que, sin embargo, no tiene toda la flexibilidad y generalidad que pretendemos darle aquí); para la relación entre esta noción y nuestra teoría, véase el capítulo primero, \S~8.
M.~Nagata ha desarrollado una extensión de este tipo en una serie de memorias~\cite{I-9} que contienen numerosos resultados especiales sobre geometría algebraica sobre anillos de Dedekind\footnote{Entre los trabajos que se aproximan a nuestro punto de vista en geometría algebraica, señalemos el importante trabajo de E.~Kähler~\cite{I-22} y una nota reciente de Chow e Igusa~\cite{I-3}, que retoman, en el marco de la teoría de Nagata-Chevalley, ciertos resultados de (FAC) y dan también una fórmula de Künneth.}.

\bigskip
\begin{center}*\quad*\quad*\end{center}
\bigskip

% Unidad R569: introducción, líneas francesas 95--120, pp. 8--9.
% Traducción: OpenAI Codex, GPT-6.1 Sol, esfuerzo Ultra.
Por último, es evidente que un libro sobre geometría algebraica, y sobre todo un libro sobre sus fundamentos, recibe necesariamente la influencia, aunque sea por intermediarios, de matemáticos como O.~Zariski y A.~Weil.
En particular, la \emph{Théorie des fonctions holomorphes} de Zariski~\cite{I-20}, convenientemente flexibilizada gracias a los métodos cohomológicos y completada con un teorema de existencia (cap.~III, \S\S~4 y~5), es, junto con la técnica de descenso expuesta en el cap.~VI, una de las principales herramientas empleadas en este tratado y nos parece una de las más potentes de que se dispone en geometría algebraica.

La técnica general en la que se inscribe puede esbozarse del siguiente modo (se ofrecerá un ejemplo típico en el cap.~IX, al estudiar el grupo fundamental).
Se tiene un morfismo propio (cap.~II) $f:X\to Y$ de una variedad algebraica a otra (más generalmente, de un esquema a otro) que se quiere estudiar en una vecindad de un punto $y\in Y$, con el fin de resolver un problema P relativo a una vecindad de $y$.
Se procede por etapas sucesivas:
\begin{enumerate}
  \item[1\textsuperscript{o}] Se puede suponer que $Y$ es afín, de modo que $X$ pasa a ser un esquema definido sobre el anillo afín $A$ de $Y$, e incluso se puede sustituir $A$ por el anillo local de $y$.
  Esta reducción siempre es fácil en la práctica (cap.~V) y nos lleva al caso en que $A$ es un anillo \emph{local}.
  \item[2\textsuperscript{o}] Se estudia el problema considerado cuando $A$ es un anillo local \emph{artiniano}.
  Para que siga teniendo efectivamente sentido cuando no se supone que $A$ sea íntegro, a veces es preciso reformular el problema P, y resulta que a menudo se obtiene así una mejor comprensión del problema, de naturaleza «infinitesimal» en esta etapa.
  \item[3\textsuperscript{o}] La teoría de esquemas formales (cap.~III, \S\S~3, 4 y~5) permite pasar del caso de un anillo artiniano al de un \emph{anillo local completo}.
  \item[4\textsuperscript{o}] Por último, si $A$ es un anillo local cualquiera, la consideración de «secciones multiformes» sobre esquemas adecuados sobre $X$, que aproximan una sección «formal» dada (cap.~IV), permitirá a menudo pasar de un resultado conocido para el esquema

  \oldpage[I]{9}
  deducido de $X$ por extensión de escalares a la completación de $A$, a un resultado análogo para una extensión finita suficientemente sencilla (por ejemplo, no ramificada) de $A$.
\end{enumerate}

Este esbozo muestra la importancia del estudio sistemático de los esquemas definidos sobre un anillo artiniano $A$.
El punto de vista de Serre en su formulación de la teoría local de cuerpos de clases, y algunos trabajos recientes de Greenberg, parecen sugerir que ese estudio podría emprenderse asociando funtorialmente a un esquema $X$ de este tipo un esquema $X'$ sobre el cuerpo residual $k$ de $A$ (supuesto perfecto), de dimensión igual, en los casos favorables, a $n\dim X$, donde $n$ es la longitud de $A$.

En cuanto a la influencia de A.~Weil, baste decir que la necesidad de desarrollar las herramientas para formular con toda la generalidad deseada la definición de la «cohomología de Weil» y para abordar la demostración\footnote{Para evitar cualquier malentendido, precisemos que esta tarea apenas acaba de emprenderse en el momento en que se escribe esta introducción y, por tanto, todavía no ha llevado a la demostración de las conjeturas de Weil.} de todas las propiedades formales necesarias para establecer sus célebres conjeturas en geometría diofántica~\cite{I-19} ha sido una de las principales motivaciones para redactar este tratado, al igual que el deseo de encontrar el marco natural de las nociones y métodos habituales de la geometría algebraica y de dar a los autores la oportunidad de comprender esas nociones y técnicas.

\bigskip
\begin{center}*\quad*\quad*\end{center}
\bigskip

% Unidad R570: introducción, líneas francesas 122--132, p. 9.
% Traducción: OpenAI Codex, GPT-6.1 Sol, esfuerzo Ultra.
Para terminar, creemos útil advertir a los lectores que, al igual que los propios autores, seguramente encontrarán alguna dificultad antes de acostumbrarse al lenguaje de los esquemas y de convencerse de que las construcciones habituales sugeridas por la intuición geométrica pueden trasladarse a este lenguaje, esencialmente de una única manera razonable.
Como en muchas partes de la matemática moderna, la intuición inicial parece alejarse cada vez más del lenguaje adecuado para expresarla con toda la precisión y generalidad deseadas.
En este caso, la dificultad psicológica se debe a la necesidad de trasladar a los objetos de una categoría ya bastante distinta de la categoría de conjuntos (a saber, la categoría de preesquemas o la de preesquemas sobre un preesquema dado) nociones familiares para los conjuntos: productos cartesianos, leyes de grupo, de anillo, de módulo, fibrados, fibrados principales homogéneos, etc.
En el futuro, seguramente será difícil para el matemático eludir este nuevo esfuerzo de abstracción, quizá bastante pequeño, después de todo, en comparación con el que hicieron nuestros predecesores al familiarizarse con la teoría de conjuntos.

\bigskip
\begin{center}*\quad*\quad*\end{center}
\bigskip

Las referencias se darán según el sistema decimal; por ejemplo, en III, 4.9.3, III indica el capítulo, 4 el párrafo y 9 la sección del párrafo.
Dentro \emph{de un mismo capítulo}, se omitirá la indicación del capítulo.


% END tree/worktree/es/CURRENT/EGA_I_front/INTRODUCTION_ES_R570.tex
\setcounter{section}{0}
% BEGIN tree/worktree/es/CURRENT/EGA_I_ch0/ega0_i_1_0_1_3.tex
% Unidad R571: EGA I, capítulo 0, §1.0, (1.0.1)--(1.0.3).
% Traducción: OpenAI Codex, GPT-6.1 Sol, esfuerzo Ultra.
\clearpage
\oldpage[0\textsubscript{I}]{11}
\begin{center}
{\large CAPÍTULO 0\par}
\vspace{1.5em}
{\Large\bfseries PRELIMINARES\par}
\end{center}

\section{Anillos de fracciones}
\label{section:0.1-fr}

\setcounter{subsection}{-1}
\subsection{Anillos y álgebras}
\label{subsection:0.1.0-fr}

\begin{env}[1.0.1]
\label{0.1.0.1-fr}
Todos los anillos considerados en este tratado tendrán un \emph{elemento unidad};
todos los módulos sobre un anillo de este tipo se supondrán \emph{unitarios}; los homomorfismos de anillos siempre se supondrán \emph{que llevan el elemento unidad al elemento unidad}; salvo indicación expresa de lo contrario, se supondrá que un subanillo de un anillo $A$ \emph{contiene el elemento unidad de $A$}.
Consideraremos principalmente anillos \emph{conmutativos}, y cuando hablemos de un anillo sin más precisiones, se entenderá que se trata de un anillo conmutativo.
Si $A$ es un anillo no necesariamente conmutativo, por $A$-módulo entenderemos siempre un módulo \emph{a izquierda}, salvo indicación expresa de lo contrario.
\end{env}

\begin{env}[1.0.2]
\label{0.1.0.2-fr}
Sean $A$, $B$ dos anillos no necesariamente conmutativos, y $\varphi:A\to B$ un homomorfismo.
Todo $B$-módulo a izquierda (resp. a derecha) $M$ puede dotarse de una estructura de $A$-módulo a izquierda (resp. a derecha) poniendo $a.m=\varphi(a).m$ (resp. $m.a=m.\varphi(a)$); cuando sea necesario distinguir en $M$ las estructuras de $A$-módulo y de $B$-módulo, denotaremos por $M_{[\varphi]}$ el $A$-módulo a izquierda (resp. a derecha) así definido.
Si $L$ es un $A$-módulo, un homomorfismo $u:L\to M_{[\varphi]}$ es, por tanto, un homomorfismo de grupos conmutativos tal que $u(a.x)=\varphi(a).u(x)$ para $a\in A$, $x\in L$; también se dirá que es un \emph{$\varphi$-homomorfismo} $L\to M$, y que el par $(\varphi,u)$ (o, por abuso de lenguaje, $u$) es un \emph{dihomomorfismo} de $(A,L)$ a $(B,M)$.
Los pares $(A,L)$ formados por un anillo $A$ y un $A$-módulo $L$ forman, por tanto, una \emph{categoría} cuyos morfismos son los dihomomorfismos.
\end{env}

\begin{env}[1.0.3]
\label{0.1.0.3-fr}
Bajo las hipótesis de (1.0.2), si $\mathfrak{J}$ es un ideal a izquierda (resp. a derecha) de $A$, denotaremos por $B\mathfrak{J}$ (resp. $\mathfrak{J}B$) el ideal a izquierda (resp. a derecha) $B\varphi(\mathfrak{J})$ (resp. $\varphi(\mathfrak{J})B$) de $B$ generado por $\varphi(\mathfrak{J})$; es también la imagen del homomorfismo canónico $B\otimes_A\mathfrak{J}\to B$ (resp. $\mathfrak{J}\otimes_A B\to B$) de $B$-módulos a izquierda (resp. a derecha).
\end{env}

% END tree/worktree/es/CURRENT/EGA_I_ch0/ega0_i_1_0_1_3.tex
% BEGIN tree/worktree/es/CURRENT/EGA_I_ch0/ega0_i_1_0_4_7.tex
% Unidad R572: EGA I, capítulo 0, (1.0.4)--(1.0.7).
% Traducción: OpenAI Codex, GPT-6.1 Sol, esfuerzo Ultra.
\begin{env}[1.0.4]
\label{0.1.0.4-fr}
Si $A$ es un anillo (conmutativo) y $B$ un anillo no necesariamente conmutativo, dar una estructura de \emph{$A$-álgebra} sobre $B$ equivale a dar un homomorfismo de anillos $\varphi:A\to B$ tal que $\varphi(A)$ esté contenido en el centro de $B$.
Para todo ideal $\mathfrak{J}$ de $A$, $\mathfrak{J}B=B\mathfrak{J}$ es entonces un ideal bilateral de $B$, y para todo $B$-módulo $M$, $\mathfrak{J}M$ es entonces un $B$-módulo igual a $(B\mathfrak{J})M$.
\end{env}

\begin{env}[1.0.5]
\label{0.1.0.5-fr}
No volveremos sobre las nociones de \emph{módulo de tipo finito} y de \emph{álgebra} (conmutativa) \emph{de tipo finito}; decir que un $A$-módulo $M$ es de tipo finito significa que existe

\oldpage[0\textsubscript{I}]{12}
una sucesión exacta $A^p\to M\to 0$.
Se dice que un $A$-módulo $M$ admite una \emph{presentación finita} si es isomorfo al conúcleo de un homomorfismo $A^p\to A^q$; en otras palabras, si existe una sucesión exacta $A^p\to A^q\to M\to 0$.
Se observará que, sobre un anillo \emph{noetheriano} $A$, todo $A$-módulo de tipo finito admite una presentación finita.

Recordemos que una $A$-álgebra $B$ se dice \emph{entera sobre $A$} si todo elemento de $B$ es raíz en $B$ de un polinomio mónico con coeficientes en $A$; equivale a decir que todo elemento de $B$ está contenido en una subálgebra de $B$ que es un \emph{$A$-módulo de tipo finito}.
Cuando esto ocurre y $B$ es conmutativa, la subálgebra de $B$ generada por una parte finita de $B$ es un $A$-módulo de tipo finito; por tanto, para que el álgebra conmutativa $B$ sea entera y de tipo finito sobre $A$, es necesario y suficiente que $B$ sea un $A$-módulo de tipo finito; se dice entonces también que $B$ es una $A$-álgebra \emph{entera finita} (o simplemente \emph{finita}, si no hay riesgo de confusión).
Se observará que, en estas definiciones, no se supone que el homomorfismo $A\to B$ que define la estructura de $A$-álgebra sea inyectivo.
\end{env}

\begin{env}[1.0.6]
\label{0.1.0.6-fr}
Un anillo \emph{íntegro} es un anillo en el que el producto de una familia finita de elementos $\neq 0$ es $\neq 0$; equivale a decir que, en un anillo de este tipo, se tiene $0\neq 1$ y que el producto de dos elementos $\neq 0$ no es nulo.
Un ideal \emph{primo} de un anillo $A$ es un ideal $\mathfrak{p}$ tal que $A/\mathfrak{p}$ sea íntegro; esto implica, por tanto, que $\mathfrak{p}\neq A$.
Para que un anillo $A$ tenga al menos un ideal primo, es necesario y suficiente que $A\neq\{0\}$.
\end{env}

\begin{env}[1.0.7]
\label{0.1.0.7-fr}
Un anillo \emph{local} es un anillo $A$ en el que existe un único ideal maximal; este ideal es entonces el complementario de los elementos invertibles y contiene todos los ideales $\neq A$.
Si $A$ y $B$ son dos anillos locales, y $\mathfrak{m}$ y $\mathfrak{n}$ sus respectivos ideales maximales, se dice que un homomorfismo $\varphi:A\to B$ es \emph{local} si $\varphi(\mathfrak{m})\subset\mathfrak{n}$ (o, equivalentemente, si $\varphi^{-1}(\mathfrak{n})=\mathfrak{m}$).
Al pasar a los cocientes, un homomorfismo de este tipo define un monomorfismo del cuerpo residual $A/\mathfrak{m}$ al cuerpo residual $B/\mathfrak{n}$.
La composición de dos homomorfismos locales es un homomorfismo local.
\end{env}

% END tree/worktree/es/CURRENT/EGA_I_ch0/ega0_i_1_0_4_7.tex
% BEGIN tree/worktree/es/CURRENT/EGA_I_ch0/ega0_i_1_1_1_2.tex
\subsection{Radical de un ideal. Nilradical y radical de un anillo}
\label{subsection:0.1.1-fr}

\begin{env}[1.1.1]
\label{0.1.1.1-fr}
Sea $\mathfrak{a}$ un ideal de un anillo $A$; el \emph{radical} de $\mathfrak{a}$, denotado por $\mathfrak{r}(\mathfrak{a})$, es el conjunto de los $x\in A$ tales que $x^n\in\mathfrak{a}$ para al menos un entero $n>0$; es un ideal que contiene $\mathfrak{a}$.
Se tiene $\mathfrak{r}(\mathfrak{r}(\mathfrak{a}))=\mathfrak{r}(\mathfrak{a})$; la relación $\mathfrak{a}\subset\mathfrak{b}$ implica $\mathfrak{r}(\mathfrak{a})\subset\mathfrak{r}(\mathfrak{b})$; el radical de una intersección finita de ideales es la intersección de sus radicales.
Si $\varphi$ es un homomorfismo de un anillo $A'$ en $A$, se tiene $\mathfrak{r}(\varphi^{-1}(\mathfrak{a}))=\varphi^{-1}(\mathfrak{r}(\mathfrak{a}))$ para todo ideal $\mathfrak{a}\subset A$.
Para que un ideal sea el radical de un ideal, es necesario y suficiente que sea una intersección de ideales primos.
El radical de un ideal $\mathfrak{a}$ es la intersección de los ideales primos \emph{minimales} entre los que contienen $\mathfrak{a}$; si $A$ es noetheriano, estos ideales primos minimales son finitos en número.

El radical del ideal $(0)$ también se llama \emph{nilradical} de $A$; es el conjunto $\mathfrak{N}$ de los elementos nilpotentes de $A$.
Se dice que el anillo $A$ es \emph{reducido} si $\mathfrak{N}=(0)$; para todo anillo $A$, el cociente $A/\mathfrak{N}$ de $A$ por su nilradical es un anillo reducido.
\end{env}

\begin{env}[1.1.2]
\label{0.1.1.2-fr}
Recordemos que el \emph{radical} $\mathfrak{R}(A)$ de un anillo $A$ (no necesariamente conmutativo) es la intersección de los ideales maximales a izquierda de $A$ (y también la intersección de los ideales maximales a derecha).
El radical de $A/\mathfrak{R}(A)$ es $(0)$.
\end{env}

% END tree/worktree/es/CURRENT/EGA_I_ch0/ega0_i_1_1_1_2.tex
% BEGIN tree/worktree/es/CURRENT/EGA_I_ch0/ega0_i_1_2_1_3.tex
\oldpage[0\textsubscript{I}]{13}
\subsection{Módulos y anillos de fracciones}
\label{subsection:0.1.2-fr}

\begin{env}[1.2.1]
\label{0.1.2.1-fr}
Se dice que un subconjunto $S$ de un anillo $A$ es \emph{multiplicativo} si $1\in S$ y si el producto de dos elementos de $S$ pertenece a $S$.
Los ejemplos más importantes para lo que sigue serán: $1^\circ$ el conjunto $S_f$ de las potencias $f^n$ ($n\geq 0$) de un elemento $f\in A$; $2^\circ$ el complementario $A\setmin\mathfrak{p}$ de un ideal \emph{primo} $\mathfrak{p}$ de $A$.
\end{env}

\begin{env}[1.2.2]
\label{0.1.2.2-fr}
Sean $S$ un subconjunto multiplicativo de un anillo $A$ y $M$ un $A$-módulo; en el conjunto $M\times S$, la relación entre pares $(m_1,s_1)$, $(m_2,s_2)$:
\[
  \text{<<~existe $s\in S$ tal que $s(s_1m_2-s_2m_1)=0$~>>}
\]
es una relación de equivalencia.
Se denota por $S^{-1}M$ el conjunto cociente de $M\times S$ por esta relación, y por $m/s$ la imagen canónica en $S^{-1}M$ del par $(m,s)$; se llama aplicación \emph{canónica} de $M$ en $S^{-1}M$ a la aplicación $i_M^S:m\to m/1$ (también denotada por $i^S$).
Esta aplicación, en general, no es ni inyectiva ni sobreyectiva; su núcleo es el conjunto de los $m\in M$ tales que existe un $s\in S$ para el cual $sm=0$.

En $S^{-1}M$ se define una ley de grupo aditivo tomando
\[
  (m_1/s_1)+(m_2/s_2)=(s_2m_1+s_1m_2)/(s_1s_2)
\]
(se comprueba que esto es independiente de las expresiones de los elementos de $S^{-1}M$ considerados).
En $S^{-1}A$ se define además una ley multiplicativa tomando $(a_1/s_1)(a_2/s_2)=(a_1a_2)/(s_1s_2)$, y finalmente una ley externa en $S^{-1}M$, con $S^{-1}A$ como conjunto de operadores, poniendo $(a/s)(m/s')=(am)/(ss')$.
Se comprueba así que $S^{-1}A$ está provisto de una estructura de anillo (llamado \emph{anillo de fracciones de $A$ con denominadores en $S$}) y $S^{-1}M$ de una estructura de $S^{-1}A$-módulo (llamado \emph{módulo de fracciones de $M$ con denominadores en $S$}); para todo $s\in S$, $s/1$ es invertible en $S^{-1}A$, y su inverso es $1/s$.
La aplicación canónica $i_A^S$ (resp. $i_M^S$) es un homomorfismo de anillos (resp. un homomorfismo de $A$-módulos, considerando $S^{-1}M$ como $A$-módulo mediante el homomorfismo $i_A^S:A\to S^{-1}A$).
\end{env}

\begin{env}[1.2.3]
\label{0.1.2.3-fr}
Si $S_f=\{f^n\}_{n\geq 0}$ para un $f\in A$, se escribe $A_f$ y $M_f$ en lugar de $S_f^{-1}A$ y $S_f^{-1}M$; cuando $A_f$ se considera como álgebra sobre $A$, puede escribirse $A_f=A[1/f]$.
$A_f$ es isomorfa al álgebra cociente $A[T]/(fT-1)A[T]$.
Cuando $f=1$, $A_f$ y $M_f$ se identifican canónicamente con $A$ y $M$; si $f$ es nilpotente, $A_f$ y $M_f$ se reducen a $0$.

Cuando $S=A\setmin\mathfrak{p}$, donde $\mathfrak{p}$ es un ideal primo de $A$, se escribe $A_{\mathfrak{p}}$ y $M_{\mathfrak{p}}$ en lugar de $S^{-1}A$ y $S^{-1}M$; $A_{\mathfrak{p}}$ es un \emph{anillo local} cuyo ideal maximal $\mathfrak{q}$ está generado por $i_A^S(\mathfrak{p})$, y se tiene $(i_A^S)^{-1}(\mathfrak{q})=\mathfrak{p}$; pasando a los cocientes, $i_A^S$ da un monomorfismo del anillo íntegro $A/\mathfrak{p}$ en el cuerpo $A_{\mathfrak{p}}/\mathfrak{q}$, que se identifica con el cuerpo de fracciones de $A/\mathfrak{p}$.
\end{env}

% END tree/worktree/es/CURRENT/EGA_I_ch0/ega0_i_1_2_1_3.tex
% BEGIN tree/worktree/es/CURRENT/EGA_I_ch0/ega0_i_1_2_4_8.tex
\begin{env}[1.2.4]
\label{0.1.2.4-fr}
El anillo de fracciones $S^{-1}A$ y el homomorfismo canónico $i_A^S$ son solución de un \emph{problema de aplicación universal}: todo homomorfismo $u$ de $A$ en un anillo $B$ tal que $u(S)$ está formado por elementos \emph{invertibles} de $B$ se factoriza de una única manera
\[
  u:A\xrightarrow[i_A^S]{}S^{-1}A\xrightarrow[u^*]{}B
\]

\oldpage[0\textsubscript{I}]{14}
donde $u^*$ es un homomorfismo de anillos.
Bajo las mismas hipótesis, sean $M$ un $A$-módulo, $N$ un $B$-módulo y $v:M\to N$ un homomorfismo de $A$-módulos (para la estructura de $B$-módulo en $N$ definida por $u:A\to B$); entonces $v$ se factoriza de una única manera
\[
  v:M\xrightarrow[i_M^S]{}S^{-1}M\xrightarrow[v^*]{}N
\]
donde $v^*$ es un homomorfismo de $S^{-1}A$-módulos (para la estructura de $S^{-1}A$-módulo en $N$ definida por $u^*$).
\end{env}

\begin{env}[1.2.5]
\label{0.1.2.5-fr}
Se define un isomorfismo canónico $S^{-1}A\otimes_A M\xrightarrow{\sim}S^{-1}M$ de $S^{-1}A$-módulos haciendo corresponder al elemento $(a/s)\otimes m$ el elemento $(am)/s$; el isomorfismo inverso envía $m/s$ a $(1/s)\otimes m$.
\end{env}

\begin{env}[1.2.6]
\label{0.1.2.6-fr}
Para todo ideal $\mathfrak{a}'$ de $S^{-1}A$, $\mathfrak{a}=(i_A^S)^{-1}(\mathfrak{a}')$ es un ideal de $A$, y $\mathfrak{a}'$ es el ideal de $S^{-1}A$ generado por $i_A^S(\mathfrak{a})$, que se identifica con $S^{-1}\mathfrak{a}$ (1.3.2).
La aplicación $\mathfrak{p}'\to(i_A^S)^{-1}(\mathfrak{p}')$ es un isomorfismo, para la estructura de orden, del conjunto de los ideales \emph{primos} de $S^{-1}A$ sobre el conjunto de los ideales primos $\mathfrak{p}$ de $A$ tales que $\mathfrak{p}\cap S=\varnothing$.
Además, los anillos locales $A_{\mathfrak{p}}$ y $(S^{-1}A)_{S^{-1}\mathfrak{p}}$ son entonces canónicamente (1.5.1) isomorfos.
\end{env}

\begin{env}[1.2.7]
\label{0.1.2.7-fr}
Cuando $A$ es un anillo \emph{íntegro}, cuyo cuerpo de fracciones se denota por $K$, la aplicación canónica $i_A^S:A\to S^{-1}A$ es inyectiva para todo subconjunto multiplicativo $S$ que no contiene $0$, y $S^{-1}A$ se identifica entonces canónicamente con un subanillo de $K$ que contiene $A$.
En particular, para todo ideal primo $\mathfrak{p}$ de $A$, $A_{\mathfrak{p}}$ es un anillo local que contiene $A$, de ideal maximal $\mathfrak{p}A_{\mathfrak{p}}$, y se tiene $\mathfrak{p}A_{\mathfrak{p}}\cap A=\mathfrak{p}$.
\end{env}

\begin{env}[1.2.8]
\label{0.1.2.8-fr}
Si $A$ es un anillo \emph{reducido} (1.1.1), lo mismo ocurre con $S^{-1}A$: en efecto, si $(x/s)^n=0$ para $x\in A$, $s\in S$, esto significa que existe $s'\in S$ tal que $s'x^n=0$, de donde $(s'x)^n=0$, lo que, por hipótesis, implica $s'x=0$, y por tanto $x/s=0$.
\end{env}

% END tree/worktree/es/CURRENT/EGA_I_ch0/ega0_i_1_2_4_8.tex
% BEGIN tree/worktree/es/CURRENT/EGA_I_ch0/ega0_i_1_3.tex
\subsection{Propiedades funtoriales}
\label{subsection:0.1.3-fr}

\begin{env}[1.3.1]
\label{0.1.3.1-fr}
Sean $M$, $N$ dos $A$-módulos, y $u$ un $A$-homomorfismo $M\to N$.
Si $S$ es un subconjunto multiplicativo de $A$, se define un $S^{-1}A$-homomorfismo $S^{-1}M\to S^{-1}N$, denotado por $S^{-1}u$, poniendo $(S^{-1}u)(m/s)=u(m)/s$; si $S^{-1}M$ y $S^{-1}N$ se identifican canónicamente con $S^{-1}A\otimes_A M$ y $S^{-1}A\otimes_A N$ (1.2.5), $S^{-1}u$ se identifica con $1\otimes u$.
Si $P$ es un tercer $A$-módulo, y $v$ un $A$-homomorfismo $N\to P$, se tiene $S^{-1}(v\circ u)=(S^{-1}v)\circ(S^{-1}u)$; en otras palabras, $S^{-1}M$ es un \emph{funtor covariante en $M$}, de la categoría de los $A$-módulos en la de los $S^{-1}A$-módulos (con $A$ y $S$ fijos).
\end{env}

\begin{env}[1.3.2]
\label{0.1.3.2-fr}
El funtor $S^{-1}M$ es \emph{exacto}; en otras palabras, si la sucesión
\[
  M\xrightarrow{u}N\xrightarrow{v}P
\]
es exacta, también lo es la sucesión
\[
  S^{-1}M\xrightarrow{S^{-1}u}S^{-1}N\xrightarrow{S^{-1}v}S^{-1}P.
\]
En particular, si $u:M\to N$ es inyectivo (resp. sobreyectivo), también lo es $S^{-1}u$;

\oldpage[0\textsubscript{I}]{15}
si $N$ y $P$ son dos submódulos de $M$, $S^{-1}N$ y $S^{-1}P$ se identifican canónicamente con submódulos de $S^{-1}M$, y se tiene
\[
  S^{-1}(N+P)=S^{-1}N+S^{-1}P
  \quad\text{y}\quad
  S^{-1}(N\cap P)=(S^{-1}N)\cap(S^{-1}P).
\]
\end{env}

\begin{env}[1.3.3]
\label{0.1.3.3-fr}
Sea $(M_\alpha,\varphi_{\beta\alpha})$ un sistema inductivo de $A$-módulos; entonces $(S^{-1}M_\alpha,S^{-1}\varphi_{\beta\alpha})$ es un sistema inductivo de $S^{-1}A$-módulos.
Expresando los $S^{-1}M_\alpha$ y $S^{-1}\varphi_{\beta\alpha}$ como productos tensoriales (1.2.5 y 1.3.1), de la conmutación de las operaciones de producto tensorial y de límite inductivo se deduce que se tiene un isomorfismo canónico
\[
  S^{-1}\varinjlim M_\alpha\xrightarrow{\sim}\varinjlim S^{-1}M_\alpha
\]
lo que también se expresa diciendo que el funtor $S^{-1}M$ (en $M$) \emph{conmuta con los límites inductivos}.
\end{env}

\begin{env}[1.3.4]
\label{0.1.3.4-fr}
Sean $M$, $N$ dos $A$-módulos; existe un isomorfismo canónico \emph{funtorial} (en $M$ y $N$)
\[
  (S^{-1}M)\otimes_{S^{-1}A}(S^{-1}N)
  \xrightarrow{\sim}
  S^{-1}(M\otimes_A N)
\]
que transforma $(m/s)\otimes(n/t)$ en $(m\otimes n)/st$.
\end{env}

\begin{env}[1.3.5]
\label{0.1.3.5-fr}
Se tiene igualmente un homomorfismo \emph{funtorial} (en $M$ y $N$)
\[
  S^{-1}\operatorname{Hom}_A(M,N)
  \longrightarrow
  \operatorname{Hom}_{S^{-1}A}(S^{-1}M,S^{-1}N)
\]
que hace corresponder a $u/s$ el homomorfismo $m/t\to u(m)/st$.
Cuando $M$ tiene una \emph{presentación finita}, el homomorfismo anterior es un \emph{isomorfismo}: esto es inmediato cuando $M$ es de la forma $A^r$, y se pasa de ahí al caso general partiendo de la sucesión exacta $A^p\to A^q\to M\to 0$, y utilizando la exactitud del funtor $S^{-1}M$ y la exactitud por la izquierda del funtor $\operatorname{Hom}_A(M,N)$ en $M$.
Se observará que este caso se presenta siempre cuando $A$ es \emph{noetheriano} y el $A$-módulo $M$ es \emph{de tipo finito}.
\end{env}

% END tree/worktree/es/CURRENT/EGA_I_ch0/ega0_i_1_3.tex
% BEGIN tree/worktree/es/CURRENT/EGA_I_ch0/ega0_i_1_4.tex
\subsection{Cambio de subconjunto multiplicativo}
\label{subsection:0.1.4-fr}

\begin{env}[1.4.1]
\label{0.1.4.1-fr}
Sean $S$, $T$ dos subconjuntos multiplicativos de un anillo $A$ tales que $S\subset T$; existe un homomorfismo canónico $\rho_A^{T,S}$ (o simplemente $\rho^{T,S}$) de $S^{-1}A$ en $T^{-1}A$, que hace corresponder al elemento denotado por $a/s$ de $S^{-1}A$ el elemento denotado por $a/s$ en $T^{-1}A$; se tiene
$i_A^T=\rho_A^{T,S}\circ i_A^S$.
Para todo $A$-módulo $M$, existe igualmente una aplicación $S^{-1}A$-lineal de $S^{-1}M$ en $T^{-1}M$ (considerando este último como $S^{-1}A$-módulo mediante el homomorfismo $\rho_A^{T,S}$), que hace corresponder al elemento $m/s$ de $S^{-1}M$ el elemento $m/s$ de $T^{-1}M$; esta aplicación se denota por $\rho_M^{T,S}$, o simplemente por $\rho^{T,S}$, y se tiene también
$i_M^T=\rho_M^{T,S}\circ i_M^S$;
en la identificación canónica (1.2.5), $\rho_M^{T,S}$ se identifica con $\rho_A^{T,S}\otimes 1$.
El homomorfismo $\rho_M^{T,S}$ es un \emph{morfismo funtorial} (o transformación natural) del funtor $S^{-1}M$ en el funtor $T^{-1}M$; en otras palabras, el diagrama
\[
\begin{tikzcd}[column sep=large,row sep=large]
S^{-1}M \arrow[r,"S^{-1}u"] \arrow[d,"\rho_M^{T,S}"'] & S^{-1}N \arrow[d,"\rho_N^{T,S}"] \\
T^{-1}M \arrow[r,"T^{-1}u"] & T^{-1}N
\end{tikzcd}
\]

\oldpage[0\textsubscript{I}]{16}
es conmutativo para todo homomorfismo $u:M\to N$; se observará además que $T^{-1}u$ está completamente determinado por $S^{-1}u$, pues para $m\in M$ y $t\in T$ se tiene
\[
  (T^{-1}u)(m/t)=(t/1)^{-1}\rho^{T,S}\bigl((S^{-1}u)(m/1)\bigr).
\]
\end{env}

\begin{env}[1.4.2]
\label{0.1.4.2-fr}
Con las mismas notaciones, para dos $A$-módulos $M$, $N$, los diagramas (cf. (1.3.4) y (1.3.5))
\[
\begin{tikzcd}[column sep=large,row sep=large]
(S^{-1}M)\otimes_{S^{-1}A}(S^{-1}N) \arrow[r,"\sim"] \arrow[d]
  & S^{-1}(M\otimes_A N) \arrow[d]
  \\
(T^{-1}M)\otimes_{T^{-1}A}(T^{-1}N) \arrow[r,"\sim"]
  & T^{-1}(M\otimes_A N)
\end{tikzcd}
\]
\[
\begin{tikzcd}[column sep=large,row sep=large]
S^{-1}\operatorname{Hom}_A(M,N) \arrow[r] \arrow[d]
  & \operatorname{Hom}_{S^{-1}A}(S^{-1}M,S^{-1}N) \arrow[d] \\
T^{-1}\operatorname{Hom}_A(M,N) \arrow[r]
  & \operatorname{Hom}_{T^{-1}A}(T^{-1}M,T^{-1}N)
\end{tikzcd}
\]
son conmutativos.
\end{env}

\begin{env}[1.4.3]
\label{0.1.4.3-fr}
Hay un caso importante en el que el homomorfismo $\rho^{T,S}$ es \emph{biyectivo}, a saber, cuando todo elemento de $T$ es divisor de un elemento de $S$; se identifican entonces mediante $\rho^{T,S}$ los módulos $S^{-1}M$ y $T^{-1}M$.
Se dice que $S$ es \emph{saturado} si todo divisor en $A$ de un elemento de $S$ pertenece a $S$; sustituyendo $S$ por el conjunto $T$ de todos los divisores de los elementos de $S$ (conjunto que es multiplicativo y saturado), se ve que siempre se puede, si se desea, limitar la consideración a módulos de fracciones $S^{-1}M$ en los que $S$ sea saturado.
\end{env}

\begin{env}[1.4.4]
\label{0.1.4.4-fr}
Si $S$, $T$, $U$ son tres subconjuntos multiplicativos de $A$ tales que $S\subset T\subset U$, se tiene
\[
  \rho^{U,S}=\rho^{U,T}\circ\rho^{T,S}.
\]
\end{env}

\begin{env}[1.4.5]
\label{0.1.4.5-fr}
Consideremos una \emph{familia filtrante creciente} $(S_\alpha)$ de subconjuntos multiplicativos de $A$ (se escribirá $\alpha\leq\beta$ para $S_\alpha\subset S_\beta$), y sea $S$ el subconjunto multiplicativo $\bigcup_\alpha S_\alpha$; pongamos
$\rho_{\beta\alpha}=\rho_A^{S_\beta,S_\alpha}$ para $\alpha\leq\beta$;
en virtud de (1.4.4), los homomorfismos $\rho_{\beta\alpha}$ definen un anillo $A'$ que es el \emph{límite inductivo} del sistema inductivo de anillos $(S_\alpha^{-1}A,\rho_{\beta\alpha})$.
Sea $\rho_\alpha$ la aplicación canónica $S_\alpha^{-1}A\to A'$, y pongamos $\varphi_\alpha=\rho_A^{S,S_\alpha}$; como $\varphi_\alpha=\varphi_\beta\circ\rho_{\beta\alpha}$ para $\alpha\leq\beta$ según (1.4.4), se puede definir de forma única un homomorfismo $\varphi:A'\to S^{-1}A$ tal que el diagrama
\[
\begin{tikzcd}[column sep=huge,row sep=large]
  & S_\alpha^{-1}A
      \arrow[ddl,"\rho_\alpha"']
      \arrow[d,"\rho_{\beta\alpha}"]
      \arrow[ddr,"\varphi_\alpha"] & \\
  & S_\beta^{-1}A
      \arrow[dl,"\rho_\beta"]
      \arrow[dr,"\varphi_\beta"'] & \\
A' \arrow[rr,"\varphi"'] & & S^{-1}A
\end{tikzcd}
\qquad(\alpha\leq\beta)
\]
sea conmutativo.
De hecho, $\varphi$ es un \emph{isomorfismo}: en efecto, es inmediato por construcción que $\varphi$ es sobreyectivo.
Por otra parte, si $\rho_\alpha(a/s_\alpha)\in A'$ es tal que $\varphi(\rho_\alpha(a/s_\alpha))=0$, esto significa que $a/s_\alpha=0$ en $S^{-1}A$, es decir, que existe $s\in S$ tal que $sa=0$; pero hay un $\beta\geq\alpha$ tal que $s\in S_\beta$, y por consiguiente, como $\rho_\alpha(a/s_\alpha)=\rho_\beta(sa/ss_\alpha)=0$, se ve que $\varphi$ es inyectivo.
Se trata de igual modo el caso de un $A$-módulo $M$, y así se han definido isomorfismos canónicos
\[
  \varinjlim S_\alpha^{-1}A\xrightarrow{\sim}
  (\varinjlim S_\alpha)^{-1}A,
  \qquad
  \varinjlim S_\alpha^{-1}M\xrightarrow{\sim}
  (\varinjlim S_\alpha)^{-1}M,
\]
el segundo de los cuales es \emph{funtorial} en $M$.
\end{env}

\oldpage[0\textsubscript{I}]{17}
\begin{env}[1.4.6]
\label{0.1.4.6-fr}
Sean $S_1$, $S_2$ dos subconjuntos multiplicativos de $A$; entonces $S_1S_2$ también es un subconjunto multiplicativo de $A$.
Denotemos por $S_2'$ la imagen canónica de $S_2$ en el anillo $S_1^{-1}A$, que es un subconjunto multiplicativo de este anillo.
Para todo $A$-módulo $M$, existe entonces un isomorfismo funtorial
\[
  S_2'^{-1}(S_1^{-1}M)\xrightarrow{\sim}(S_1S_2)^{-1}M
\]
que hace corresponder a $(m/s_1)/(s_2/1)$ el elemento $m/(s_1s_2)$.
\end{env}

% END tree/worktree/es/CURRENT/EGA_I_ch0/ega0_i_1_4.tex
% BEGIN tree/worktree/es/CURRENT/EGA_I_ch0/ega0_i_1_5.tex
\subsection{Cambio de anillo}
\label{subsection:0.1.5-fr}

\begin{env}[1.5.1]
\label{0.1.5.1-fr}
Sean $A$, $A'$ dos anillos, $\varphi$ un homomorfismo $A'\to A$, y $S$ (resp. $S'$) un subconjunto multiplicativo de $A$ (resp. $A'$), tal que $\varphi(S')\subset S$; el homomorfismo compuesto
$A'\xrightarrow{\varphi}A\longrightarrow S^{-1}A$
se factoriza en
$A'\longrightarrow S'^{-1}A'\xrightarrow{\varphi^{S'}}S^{-1}A$
en virtud de (1.2.4); se tiene $\varphi^{S'}(a'/s')=\varphi(a')/\varphi(s')$.
Si $A=\varphi(A')$ y $S=\varphi(S')$, $\varphi^{S'}$ es \emph{sobreyectivo}.
Si $A'=A$ y $\varphi$ es la identidad, $\varphi^{S'}$ no es otro que el homomorfismo $\rho_A^{S,S'}$ definido en (1.4.1).
\end{env}

\begin{env}[1.5.2]
\label{0.1.5.2-fr}
Bajo las hipótesis de (1.5.1), sea $M$ un $A$-módulo.
Existe un homomorfismo canónico funtorial
\[
  \sigma:S'^{-1}(M_{[\varphi]})\longrightarrow
  (S^{-1}M)_{[\varphi^{S'}]}
\]
de $S'^{-1}A'$-módulos, que hace corresponder a todo elemento $m/s'$ de $S'^{-1}(M_{[\varphi]})$ el elemento $m/\varphi(s')$ de $(S^{-1}M)_{[\varphi^{S'}]}$; se comprueba inmediatamente que esta definición no depende de la expresión $m/s'$ del elemento considerado.
\emph{Cuando} $S=\varphi(S')$, el homomorfismo $\sigma$ es \emph{biyectivo}.
Cuando $A'=A$ y $\varphi$ es la identidad, $\sigma$ no es otro que el homomorfismo $\rho_M^{S,S'}$ definido en (1.4.1).

Tomando en particular $M=A$, el homomorfismo $\varphi$ define en $A$ una estructura de $A'$-álgebra; $S'^{-1}(A_{[\varphi]})$ está entonces provisto de una estructura de anillo, para la cual se identifica con $(\varphi(S'))^{-1}A$, y el homomorfismo
$\sigma:S'^{-1}(A_{[\varphi]})\longrightarrow S^{-1}A$
es un homomorfismo de $S'^{-1}A'$-álgebras.
\end{env}

\begin{env}[1.5.3]
\label{0.1.5.3-fr}
Sean $M$ y $N$ dos $A$-módulos; componiendo los homomorfismos definidos en (1.3.4) y (1.5.2), se obtiene un homomorfismo
\[
  (S^{-1}M\otimes_{S^{-1}A}S^{-1}N)_{[\varphi^{S'}]}
  \longleftarrow
  S'^{-1}((M\otimes_A N)_{[\varphi]})
\]
que es un isomorfismo cuando $\varphi(S')=S$.
Igualmente, componiendo los homomorfismos (1.3.5) y (1.5.2), se obtiene un homomorfismo
\[
  S'^{-1}((\operatorname{Hom}_A(M,N))_{[\varphi]})
  \longrightarrow
  (\operatorname{Hom}_{S^{-1}A}(S^{-1}M,S^{-1}N))_{[\varphi^{S'}]}
\]
que es un isomorfismo cuando $\varphi(S')=S$ y $M$ admite una presentación finita.
\end{env}

\begin{env}[1.5.4]
\label{0.1.5.4-fr}
Consideremos ahora un $A'$-módulo $N'$ y formemos el producto tensorial $N'\otimes_{A'}A_{[\varphi]}$, que puede considerarse como un $A$-módulo poniendo
$a.(n'\otimes b)=n'\otimes(ab)$.
Existe un isomorfismo funtorial de $S^{-1}A$-módulos
\[
  \tau:(S'^{-1}N')\otimes_{S'^{-1}A'}(S^{-1}A)_{[\varphi^{S'}]}
  \xrightarrow{\sim}
  S^{-1}(N'\otimes_{A'}A_{[\varphi]})
\]
\oldpage[0\textsubscript{I}]{18}
que hace corresponder al elemento $(n'/s')\otimes(a/s)$ el elemento
$(n'\otimes a)/(\varphi(s')s)$; se comprueba por separado que, cuando se
sustituye $n'/s'$ (resp. $a/s$) por otra expresión del mismo elemento,
$(n'\otimes a)/(\varphi(s')s)$ no cambia; por otra parte, puede definirse un
homomorfismo inverso de $\tau$ haciendo corresponder a $(n'\otimes a)/s$
el elemento $(n'/1)\otimes(a/s)$: se utiliza el hecho de que
$S^{-1}(N'\otimes_{A'}A_{[\varphi]})$ es canónicamente isomorfo a
$(N'\otimes_{A'}A_{[\varphi]})\otimes_A S^{-1}A$ (1.2.5), y por tanto también a
$N'\otimes_{A'}(S^{-1}A)_{[\psi]}$, denotando por $\psi$ el homomorfismo
compuesto $a'\to\varphi(a')/1$ de $A'$ en $S^{-1}A$.
\end{env}

\begin{env}[1.5.5]
\label{0.1.5.5-fr}
Si $M'$ y $N'$ son dos $A'$-módulos, componiendo los isomorfismos (1.3.4)
y (1.5.4), se obtiene un isomorfismo
\[
  S'^{-1}M'\otimes_{S'^{-1}A'}S'^{-1}N'
  \otimes_{S'^{-1}A'}S^{-1}A
  \xrightarrow{\sim}
  S^{-1}(M'\otimes_{A'}N'\otimes_{A'}A).
\]
Igualmente, si $M'$ admite una presentación finita, se tiene, en virtud de (1.3.5) y
(1.5.4), un isomorfismo
\[
  \operatorname{Hom}_{S'^{-1}A'}(S'^{-1}M',S'^{-1}N')
  \otimes_{S'^{-1}A'}S^{-1}A
  \xrightarrow{\sim}
  S^{-1}(\operatorname{Hom}_{A'}(M',N')\otimes_{A'}A).
\]
\end{env}

\begin{env}[1.5.6]
\label{0.1.5.6-fr}
Bajo las hipótesis de (1.5.1), sea $T$ (resp. $T'$) un segundo subconjunto
multiplicativo de $A$ (resp. $A'$) tal que $S\subset T$ (resp. $S'\subset T'$)
y $\varphi(T')\subset T$. Entonces el diagrama
\[
\begin{tikzcd}[column sep=huge,row sep=large]
S'^{-1}A' \arrow[r,"\varphi^{S'}"] \arrow[d,"\rho^{T',S'}"']
  & S^{-1}A \arrow[d,"\rho^{T,S}"] \\
T'^{-1}A' \arrow[r,"\varphi^{T'}"']
  & T^{-1}A
\end{tikzcd}
\]
es conmutativo. Si $M$ es un $A$-módulo, el diagrama
\[
\begin{tikzcd}[column sep=huge,row sep=large]
S'^{-1}(M_{[\varphi]}) \arrow[r,"\sigma"] \arrow[d,"\rho^{T',S'}"']
  & (S^{-1}M)_{[\varphi^{S'}]} \arrow[d,"\rho^{T,S}"] \\
T'^{-1}(M_{[\varphi]}) \arrow[r,"\sigma"']
  & (T^{-1}M)_{[\varphi^{T'}]}
\end{tikzcd}
\]
es conmutativo. Finalmente, si $N'$ es un $A'$-módulo, el diagrama
\[
\begin{tikzcd}[column sep=huge,row sep=large]
(S'^{-1}N')\otimes_{S'^{-1}A'}(S^{-1}A)_{[\varphi^{S'}]}
  \arrow[r,"\tau","{\sim}"'] \arrow[d]
  & S^{-1}(N'\otimes_{A'}A_{[\varphi]}) \arrow[d,"\rho^{T,S}"] \\
(T'^{-1}N')\otimes_{T'^{-1}A'}(T^{-1}A)_{[\varphi^{T'}]}
  \arrow[r,"\tau"',"{\sim}"]
  & T^{-1}(N'\otimes_{A'}A_{[\varphi]})
\end{tikzcd}
\]
es conmutativo, y la flecha vertical de la izquierda se obtiene aplicando
$\rho_{N'}^{T',S'}$ a $S'^{-1}N'$ y $\rho_A^{T,S}$ a $S^{-1}A$.
\end{env}

\oldpage[0\textsubscript{I}]{19}

\begin{env}[1.5.7]
\label{0.1.5.7-fr}
Sean $A''$ un tercer anillo, $\varphi':A''\to A'$ un homomorfismo
de anillos, y $S''$ un subconjunto multiplicativo de $A''$ tal que
$\varphi'(S'')\subset S'$. Pongamos $\varphi''=\varphi\circ\varphi'$; entonces se tiene
\[
  \varphi''^{S''}=\varphi^{S'}\circ\varphi'^{S''}.
\]
Sea $M$ un $A$-módulo; se tiene evidentemente
$M_{[\varphi'']}=(M_{[\varphi]})_{[\varphi']}$; si $\sigma'$ y $\sigma''$ son
los homomorfismos definidos a partir de $\varphi'$ y $\varphi''$ como $\sigma$
se define en (1.5.2) a partir de $\varphi$, se tiene la fórmula de transitividad
\[
  \sigma''=\sigma\circ\sigma'.
\]
Finalmente, sea $N''$ un $A''$-módulo; el $A$-módulo
$N''\otimes_{A''}A_{[\varphi'']}$ se identifica canónicamente con
$(N''\otimes_{A''}A'_{[\varphi']})\otimes_{A'}A_{[\varphi]}$,
e igualmente, el $S^{-1}A$-módulo
$(S''^{-1}N'')\otimes_{S''^{-1}A''}(S^{-1}A)_{[\varphi''^{S''}]}$
se identifica canónicamente con
$((S''^{-1}N'')\otimes_{S''^{-1}A''}(S'^{-1}A')_{[\varphi'^{S''}]})\otimes_{S'^{-1}A'}(S^{-1}A)_{[\varphi^{S'}]}$.
Con estas identificaciones, si $\tau'$ y $\tau''$ son los isomorfismos definidos
a partir de $\varphi'$ y $\varphi''$ como $\tau$ se define en (1.5.4) a
partir de $\varphi$, se tiene la fórmula de transitividad
\[
  \tau''=\tau\circ(\tau'\otimes 1).
\]
\end{env}

\begin{env}[1.5.8]
\label{0.1.5.8-fr}
Sea $A$ un subanillo de un anillo $B$; para todo ideal primo
\emph{minimal} $\mathfrak p$ de $A$, existe un ideal primo minimal
$\mathfrak q$ de $B$ tal que $\mathfrak p=A\cap\mathfrak q$. En efecto,
$A_{\mathfrak p}$ es un subanillo de $B_{\mathfrak p}$ (1.3.2) y tiene
\emph{un solo} ideal primo $\mathfrak p'$ (1.2.6); como $B_{\mathfrak p}$ no se
reduce a $0$, tiene al menos un ideal primo $\mathfrak q'$ y se tiene
necesariamente $\mathfrak q'\cap A_{\mathfrak p}=\mathfrak p'$; el ideal primo
$\mathfrak q_1$ de $B$, preimagen de $\mathfrak q'$, satisface por tanto
$\mathfrak q_1\cap A=\mathfrak p$, y \emph{con mayor razón} se tiene
$\mathfrak q\cap A=\mathfrak p$ para todo ideal primo minimal $\mathfrak q$
de $B$ contenido en $\mathfrak q_1$.
\end{env}

% END tree/worktree/es/CURRENT/EGA_I_ch0/ega0_i_1_5.tex
% BEGIN tree/worktree/es/CURRENT/EGA_I_ch0/ega0_i_1_6_7.tex
\subsection{Identificación del módulo $M_f$ con un límite inductivo}
\label{subsection:0.1.6-fr}

\begin{env}[1.6.1]
\label{0.1.6.1-fr}
Sean $M$ un $A$-módulo y $f$ un elemento de $A$. Consideremos una sucesión
$(M_n)$ de $A$-módulos, todos idénticos a $M$, y para todo par de enteros
$m\leq n$, sea $\varphi_{nm}$ el homomorfismo $z\to f^{n-m}z$ de $M_m$
en $M_n$; es inmediato que $((M_n),(\varphi_{nm}))$ es un \emph{sistema
inductivo} de $A$-módulos; sea $N=\varinjlim M_n$ el límite inductivo de este
sistema. Vamos a definir un $A$-isomorfismo canónico \emph{funtorial} de
$N$ sobre $M_f$. Para ello, observemos que, para todo $n$,
$\theta_n:z\to z/f^n$ es un $A$-homomorfismo de $M=M_n$ en $M_f$, y
de las definiciones se deduce que se tiene $\theta_n\circ\varphi_{nm}=\theta_m$ para
$m\leq n$. Existe, por tanto, un $A$-homomorfismo $\theta:N\to M_f$ tal que, si
$\varphi_n$ denota el homomorfismo canónico $M_n\to N$, se tiene
$\theta_n=\theta\circ\varphi_n$ para todo $n$. Como, por hipótesis, todo elemento
de $M_f$ es de la forma $z/f^n$ para al menos un $n$, es claro que
$\theta$ es sobreyectivo. Por otra parte, si $\theta(\varphi_n(z))=0$, es decir,
$z/f^n=0$, existe un entero $k>0$ tal que $f^kz=0$, y por tanto
$\varphi_{n+k,n}(z)=0$, lo que implica $\varphi_n(z)=0$. Se puede, pues,
identificar $M_f$ y $\varinjlim M_n$ mediante $\theta$.
\end{env}

\begin{env}[1.6.2]
\label{0.1.6.2-fr}
Escribamos ahora $M_{f,n}$, $\varphi_{nm}^f$ y $\varphi_n^f$ en lugar de
$M_n$, $\varphi_{nm}$ y $\varphi_n$. Sea $g$ un segundo elemento de $A$. Como
$f^n$ divide $f^ng^n$, se tiene un homomorfismo funtorial
\[
  \rho_{fg,f}:M_f\longrightarrow M_{fg}\qquad\text{(1.4.1 y 1.4.3)};
\]

\oldpage[0\textsubscript{I}]{20}

si se identifican $M_f$ y $M_{fg}$ con $\varinjlim M_{f,n}$ y
$\varinjlim M_{fg,n}$ respectivamente, $\rho_{fg,f}$ se identifica con el
\emph{límite inductivo} de las aplicaciones
$\rho_{fg,f}^n:M_{f,n}\to M_{fg,n}$ definidas por
$\rho_{fg,f}^n(z)=g^nz$. Esto se deduce inmediatamente de la
conmutatividad del diagrama
\[
\begin{tikzcd}[column sep=huge,row sep=large]
M_{f,n} \arrow[r,"\rho_{fg,f}^n"] \arrow[d,"\varphi_n^f"']
  & M_{fg,n} \arrow[d,"\varphi_n^{fg}"] \\
M_f \arrow[r,"\rho_{fg,f}"]
  & M_{fg}
\end{tikzcd}
\]
\end{env}

\subsection{Soporte de un módulo}
\label{subsection:0.1.7-fr}

\begin{env}[1.7.1]
\label{0.1.7.1-fr}
Dado un $A$-módulo $M$, se llama \emph{soporte} de $M$ y se denota por
$\operatorname{Supp}(M)$ al conjunto de los ideales primos $\mathfrak p$ de $A$
tales que $M_{\mathfrak p}\neq0$. Para que $M=0$, es necesario y suficiente que
$\operatorname{Supp}(M)=\varnothing$, pues si $M_{\mathfrak p}=0$ para todo
$\mathfrak p$, el anulador de un elemento $x\in M$ no puede estar contenido en
ningún ideal primo de $A$, y por tanto es todo $A$.
\end{env}

\begin{env}[1.7.2]
\label{0.1.7.2-fr}
Si $0\to N\to M\to P\to0$ es una sucesión exacta de $A$-módulos, se tiene
\[
  \operatorname{Supp}(M)=\operatorname{Supp}(N)\cup\operatorname{Supp}(P)
\]
pues, para todo ideal primo $\mathfrak p$ de $A$, la sucesión
$0\to N_{\mathfrak p}\to M_{\mathfrak p}\to P_{\mathfrak p}\to0$ es exacta
(1.3.2) y, para que $M_{\mathfrak p}=0$, es necesario y suficiente que
$N_{\mathfrak p}=P_{\mathfrak p}=0$.
\end{env}

\begin{env}[1.7.3]
\label{0.1.7.3-fr}
Si $M$ es suma de una familia $(M_\lambda)$ de submódulos,
$M_{\mathfrak p}$ es suma de los $(M_\lambda)_{\mathfrak p}$ para todo ideal
primo $\mathfrak p$ de $A$ (1.3.3 y 1.3.2), y por tanto
$\operatorname{Supp}(M)=\bigcup_\lambda\operatorname{Supp}(M_\lambda)$.
\end{env}

\begin{env}[1.7.4]
\label{0.1.7.4-fr}
Si $M$ es un $A$-módulo \emph{de tipo finito}, $\operatorname{Supp}(M)$ es
el conjunto de los ideales primos \emph{que contienen el anulador de $M$}. En efecto,
si $M$ es monógeno y está generado por $x$, decir que $M_{\mathfrak p}=0$ significa
que existe $s\notin\mathfrak p$ tal que $s.x=0$, y por tanto que $\mathfrak p$ no
contiene el anulador de $x$. Si ahora $M$ admite un sistema finito
$(x_i)_{1\leq i\leq n}$ de generadores y $\mathfrak a_i$ es el anulador
de $x_i$, de (1.7.3) se deduce que $\operatorname{Supp}(M)$ es el conjunto de los
$\mathfrak p$ que contienen alguno de los $\mathfrak a_i$, o, lo que es equivalente,
el conjunto de los $\mathfrak p$ que contienen $\mathfrak a=\bigcap_i\mathfrak a_i$,
que es el anulador de $M$.
\end{env}

\begin{env}[1.7.5]
\label{0.1.7.5-fr}
Si $M$ y $N$ son dos $A$-módulos \emph{de tipo finito}, se tiene
\[
  \operatorname{Supp}(M\otimes_A N)
  =\operatorname{Supp}(M)\cap\operatorname{Supp}(N).
\]
Se trata de ver que, si $\mathfrak p$ es un ideal primo de $A$, la
condición $M_{\mathfrak p}\otimes_{A_{\mathfrak p}}N_{\mathfrak p}\neq0$ es
equivalente a «~$M_{\mathfrak p}\neq0$ y $N_{\mathfrak p}\neq0$~» (teniendo
en cuenta (1.3.4)). En otras palabras, se trata de ver que, si $P$, $Q$ son dos
módulos de tipo finito sobre un anillo \emph{local} $B$, no reducidos a $0$, entonces
$P\otimes_B Q\neq0$. Sea $\mathfrak m$ el ideal maximal de $B$. En virtud del
lema de Nakayama, los espacios vectoriales $P/\mathfrak mP$ y
$Q/\mathfrak mQ$ no se reducen a $0$, y por tanto tampoco es nulo su
producto tensorial
$(P/\mathfrak mP)\otimes_{B/\mathfrak m}(Q/\mathfrak mQ)=(P\otimes_B Q)\otimes_B(B/\mathfrak m)$,
de donde se obtiene la conclusión.

En particular, si $M$ es un $A$-módulo de tipo finito, y $\mathfrak a$ un ideal
de $A$, $\operatorname{Supp}(M/\mathfrak aM)$ es el conjunto de los ideales
primos que contienen a la vez $\mathfrak a$ y el anulador $\mathfrak n$ de
$M$ (1.7.4), es decir, el conjunto de los ideales primos que contienen
$\mathfrak a+\mathfrak n$.
\end{env}

% END tree/worktree/es/CURRENT/EGA_I_ch0/ega0_i_1_6_7.tex
% BEGIN tree/worktree/es/CURRENT/EGA_I_ch0/ega0_i_2_1_1_4.tex
% R581: EGA I, 0, 2.1.1--2.1.4.
% Traducción: OpenAI Codex, GPT-6.1 Sol, esfuerzo Ultra.
\oldpage[0\textsubscript{I}]{21}

\section{Espacios irreducibles. Espacios noetherianos}
\label{section:0.2-fr}

\subsection{Espacios irreducibles}
\label{subsection:0.2.1-fr}

\begin{env}[2.1.1]
\label{0.2.1.1-fr}
Se dice que un espacio topológico $X$ es \emph{irreducible} si no es vacío
y no es unión de dos subespacios cerrados distintos de $X$. Esto
equivale a decir que $X\neq\varnothing$ y que la intersección de dos
abiertos no vacíos de $X$ (y, por tanto, de un número finito de ellos) no es
vacía, o que todo abierto no vacío es denso en todo el espacio, o que toda parte
cerrada $\neq X$ tiene \emph{interior vacío}, o, finalmente, que todo abierto de $X$ es
\emph{conexo}.
\end{env}

\begin{env}[2.1.2]
\label{0.2.1.2-fr}
Para que un subespacio $Y$ de un espacio topológico $X$ sea irreducible, es
necesario y suficiente que su adherencia $\overline Y$ sea irreducible. En
particular, todo subespacio que es la adherencia $\overline{\{x\}}$ de un
subespacio reducido a un punto es irreducible; expresaremos la relación
$y\in\overline{\{x\}}$ (equivalente a
$\overline{\{y\}}\subset\overline{\{x\}}$) diciendo que $y$ es
\emph{una especialización de $x$}, o que $x$ es \emph{una generalización de $y$}.
Cuando en un espacio irreducible $X$ existe un punto $x$ tal que
$X=\overline{\{x\}}$, diremos que $x$ es un \emph{punto genérico} de $X$.
Todo abierto no vacío de $X$ contiene entonces a $x$, y todo subespacio que contiene
a $x$ admite a $x$ como punto genérico.
\end{env}

\begin{env}[2.1.3]
\label{0.2.1.3-fr}
Recordemos que se llama \emph{espacio de Kolmogoroff} a un espacio topológico $X$
que satisface el axioma de separación siguiente:

\smallskip
\noindent
\textup{(T$_0$)} Si $x\neq y$ son dos puntos cualesquiera de $X$, existe
un conjunto abierto que contiene a uno de los puntos $x$, $y$ y no al otro.

Si un espacio de Kolmogoroff irreducible admite un punto genérico, admite
\emph{uno solo}, pues todo abierto no vacío contiene todo punto
genérico.

Recordemos que un espacio topológico $X$ se llama \emph{cuasicompacto} si de todo
recubrimiento abierto de $X$ se puede extraer un subrecubrimiento finito de $X$ (o,
lo que es equivalente, si toda familia filtrante decreciente de conjuntos
cerrados no vacíos tiene una intersección no vacía). Si $X$ es un espacio
cuasicompacto, toda parte cerrada no vacía $A$ de $X$ contiene un conjunto
cerrado no vacío \emph{minimal} $M$, pues el conjunto de las partes cerradas no vacías
de $A$ es inductivo para la relación $\supset$; si, además, $X$ es un espacio
de Kolmogoroff, $M$ se reduce necesariamente a un solo punto (o, como se
dice por abuso de lenguaje, es un \emph{punto cerrado}).
\end{env}

\begin{env}[2.1.4]
\label{0.2.1.4-fr}
En un espacio irreducible $X$, todo subespacio abierto no vacío $U$ es
irreducible, y si $X$ admite un punto genérico $x$, $x$ es también un punto
genérico de $U$.

Sea $(U_\alpha)$ un recubrimiento (con conjunto de índices no vacío) de un
espacio topológico $X$, formado por abiertos no vacíos; para que $X$ sea
irreducible, es necesario y suficiente que $U_\alpha$ sea irreducible para todo
$\alpha$, y que $U_\alpha\cap U_\beta\neq\varnothing$ para cualesquiera
$\alpha$, $\beta$. La condición es evidentemente necesaria; para ver que
es suficiente, basta probar que, si $V$ es un abierto no vacío de $X$,
$V\cap U_\alpha$ no es vacío para todo $\alpha$, pues entonces $V\cap U_\alpha$
es denso en $U_\alpha$ para todo $\alpha$, y, por tanto, $V$ es denso en
$X$. Ahora bien, existe al menos un índice $\gamma$ tal que
$V\cap U_\gamma\neq\varnothing$; por tanto, $V\cap U_\gamma$ es denso en
$U_\gamma$, y como, para todo $\alpha$,
$U_\alpha\cap U_\gamma\neq\varnothing$, también se tiene
$V\cap U_\alpha\cap U_\gamma\neq\varnothing$.
\end{env}

% END tree/worktree/es/CURRENT/EGA_I_ch0/ega0_i_2_1_1_4.tex
% BEGIN tree/worktree/es/CURRENT/EGA_I_ch0/ega0_i_2_1_5_8.tex
% R582: EGA I, 0, 2.1.5--2.1.8.
% Traducción: OpenAI Codex, GPT-6.1 Sol, esfuerzo Ultra.
\oldpage[0\textsubscript{I}]{22}

\begin{env}[2.1.5]
\label{0.2.1.5-fr}
Sean $X$ un espacio irreducible y $f$ una aplicación continua de $X$ en un
espacio topológico $Y$. Entonces $f(X)$ es irreducible y, si $x$ es un punto
genérico de $X$, $f(x)$ es un punto genérico de $f(X)$ y, por tanto, también de
$\overline{f(X)}$. En particular, si además $Y$ es irreducible y tiene un solo
punto genérico $y$, para que $f(X)$ sea denso en todo el espacio, es necesario y suficiente
que $f(x)=y$.
\end{env}

\begin{env}[2.1.6]
\label{0.2.1.6-fr}
Todo subespacio irreducible de un espacio topológico $X$ está contenido en un
subespacio irreducible \emph{maximal}, que es necesariamente cerrado. Los
subespacios irreducibles maximales de $X$ se llaman \emph{componentes
irreducibles} de $X$. Si $Z_1$, $Z_2$ son dos componentes irreducibles
distintas del espacio $X$, $Z_1\cap Z_2$ es un conjunto cerrado \emph{de interior vacío}
en cada uno de los subespacios $Z_1$, $Z_2$; en particular, si una componente
irreducible de $X$ admite un punto genérico (2.1.2), tal punto no puede
pertenecer a ninguna otra componente irreducible. Si $X$ tiene un número
\emph{finito} de componentes irreducibles $Z_i$ ($1\leq i\leq n$) y si, para
cada $i$, se pone
\[
  U_i=Z_i\cap\complement\Bigl(\bigcup_{j\neq i}Z_j\Bigr),
\]
los $U_i$ son abiertos, irreducibles, disjuntos dos a dos y su
unión es densa en $X$.

Sea $U$ una parte abierta de un espacio topológico $X$. Si $Z$ es una parte
irreducible de $X$ que intersecta a $U$, $Z\cap U$ es abierto y denso en $Z$,
y, por tanto, irreducible; recíprocamente, para toda parte cerrada irreducible $Y$ de
$U$, la adherencia $\overline Y$ de $Y$ en $X$ es irreducible y
$\overline Y\cap U=Y$. Se concluye que hay una \emph{correspondencia
biyectiva} entre las componentes irreducibles de $U$ y las componentes
irreducibles de $X$ que intersectan a $U$.
\end{env}

\begin{env}[2.1.7]
\label{0.2.1.7-fr}
Si un espacio topológico $X$ es unión de un número \emph{finito} de
subespacios irreducibles cerrados $Y_i$, las componentes irreducibles de $X$
son los elementos maximales del conjunto de los $Y_i$, pues, si $Z$ es una parte
cerrada irreducible de $X$, $Z$ es unión de los $Z\cap Y_i$, de donde se deduce
inmediatamente que $Z$ debe estar contenido en uno de los $Y_i$. Sea $Y$ un subespacio
de un espacio topológico $X$ y supongamos que $Y$ tiene un número finito de
componentes irreducibles $Y_i$ ($1\leq i\leq n$); entonces las adherencias
$\overline{Y}_i$ en $X$ son las componentes irreducibles de $\overline Y$.
\end{env}

\begin{env}[2.1.8]
\label{0.2.1.8-fr}
Sea $Y$ un espacio irreducible que admite un solo punto genérico $y$. Sean
$X$ un espacio topológico y $f$ una aplicación continua de $X$ en $Y$.
Entonces, para toda componente irreducible $Z$ de $X$ que intersecta a $f^{-1}(y)$,
$f(Z)$ es denso en $Y$. La recíproca no es necesariamente cierta;
sin embargo, si $Z$ admite un punto genérico $z$ y si $f(Z)$ es denso en
$Y$, necesariamente se tiene $f(z)=y$ (2.1.5); además, $Z\cap f^{-1}(y)$ es
entonces la adherencia de $\{z\}$ en $f^{-1}(y)$ y, por tanto, es irreducible y,
como toda parte irreducible de $f^{-1}(y)$ que contiene a $z$ está
necesariamente contenida en $Z$ (2.1.6), $z$ es un punto genérico de
$Z\cap f^{-1}(y)$. Como toda componente irreducible de $f^{-1}(y)$ está
contenida en una componente irreducible de $X$, se ve que, si toda
componente irreducible $Z$ de $X$ que intersecta a $f^{-1}(y)$ admite un punto
genérico, hay una \emph{correspondencia biyectiva} entre el conjunto de
estas componentes y el conjunto de las componentes irreducibles
$Z\cap f^{-1}(y)$ de $f^{-1}(y)$, siendo los puntos genéricos de $Z$
idénticos a los de $Z\cap f^{-1}(y)$.
\end{env}

% END tree/worktree/es/CURRENT/EGA_I_ch0/ega0_i_2_1_5_8.tex
% BEGIN tree/worktree/es/CURRENT/EGA_I_ch0/ega0_i_2_2.tex
% R583: EGA I, 0, 2.2.1--2.2.5.
% Traducción: OpenAI Codex, GPT-6.1 Sol, esfuerzo Ultra.
\oldpage[0\textsubscript{I}]{23}

\subsection{Espacios noetherianos}
\label{subsection:0.2.2-fr}

\begin{env}[2.2.1]
\label{0.2.2.1-fr}
Se dice que un espacio topológico $X$ es \emph{noetheriano} si el conjunto de los
abiertos de $X$ satisface la condición \emph{maximal}, o, lo que equivale a lo
mismo, si el conjunto de los cerrados de $X$ satisface la condición \emph{minimal}. Se
dice que $X$ es \emph{localmente noetheriano} si todo $x\in X$ admite una
vecindad que es un subespacio noetheriano.
\end{env}

\begin{env}[2.2.2]
\label{0.2.2.2-fr}
Sea $E$ un conjunto ordenado que satisface la condición \emph{minimal} y sea
$\boldsymbol{P}$ una propiedad de los elementos de $E$ sujeta a la condición siguiente: si
$a\in E$ es tal que, para todo $x<a$, $\boldsymbol{P}(x)$ es cierta, entonces $\boldsymbol{P}(a)$ es
cierta. En estas condiciones, \emph{$\boldsymbol{P}(x)$ es cierta para todo $x\in E$}
(«principio de inducción noetheriana»). En efecto, sea $F$ el conjunto de los
$x\in E$ para los que $\boldsymbol{P}(x)$ es falsa; si $F$ no fuera vacío, tendría un
elemento minimal $a$ y, como entonces $\boldsymbol{P}(x)$ es cierta para todo $x<a$, $\boldsymbol{P}(a)$
también sería cierta, lo que es contradictorio.

Aplicaremos este principio, en particular, cuando $E$ sea un \emph{conjunto
de partes cerradas de un espacio noetheriano}.
\end{env}

\begin{env}[2.2.3]
\label{0.2.2.3-fr}
Todo subespacio de un espacio noetheriano es noetheriano. Recíprocamente, todo
espacio topológico que es unión finita de subespacios noetherianos es noetheriano.
\end{env}

\begin{env}[2.2.4]
\label{0.2.2.4-fr}
Todo espacio noetheriano es cuasicompacto; recíprocamente, todo espacio topológico
en el que todo abierto es cuasicompacto es noetheriano.
\end{env}

\begin{env}[2.2.5]
\label{0.2.2.5-fr}
Un espacio noetheriano tiene un número \emph{finito} de componentes
irreducibles, como se ve por inducción noetheriana.
\end{env}

% END tree/worktree/es/CURRENT/EGA_I_ch0/ega0_i_2_2.tex
% BEGIN tree/worktree/es/CURRENT/EGA_I_ch0/ega0_i_3_1_1_3.tex
% R585: EGA I, 0, 3.1.1--3.1.3, con ambas notas.
% Traducción: OpenAI Codex, GPT-6.1 Sol, esfuerzo Ultra.
\section{Complementos sobre los haces}
\label{section:0.3-fr}

\subsection{Haces con valores en una categoría}
\label{subsection:0.3.1-fr}

\begin{env}[3.1.1]
\label{0.3.1.1-fr}
Sean $\boldsymbol{K}$ una categoría, $(A_\alpha)_{\alpha\in I}$,
$(A_{\alpha\beta})_{(\alpha,\beta)\in I\times I}$ dos familias de objetos de
$\boldsymbol{K}$ tales que $A_{\beta\alpha}=A_{\alpha\beta}$, y
$(\rho_{\alpha\beta})_{(\alpha,\beta)\in I\times I}$ una familia de morfismos
$\rho_{\alpha\beta}:A_\alpha\to A_{\alpha\beta}$. Diremos que un par
formado por un objeto $A$ de $\boldsymbol{K}$ y una familia de morfismos
$\rho_\alpha:A\to A_\alpha$ es \emph{solución del problema universal} definido
por los datos de las familias $(A_\alpha)$, $(A_{\alpha\beta})$ y
$(\rho_{\alpha\beta})$ si, para todo objeto $B$ de $\boldsymbol{K}$, la aplicación
que a todo $f\in\operatorname{Hom}(B,A)$ hace corresponder la familia
$(\rho_\alpha\circ f)\in\prod_\alpha\operatorname{Hom}(B,A_\alpha)$ es una
\emph{biyección} de $\operatorname{Hom}(B,A)$ sobre el conjunto de las $(f_\alpha)$
tales que
$\rho_{\alpha\beta}\circ f_\alpha=\rho_{\beta\alpha}\circ f_\beta$ para todo
par de índices $(\alpha,\beta)$. Se ve inmediatamente que, si existe tal
solución, es única salvo isomorfismo.
\end{env}

\begin{env}[3.1.2]
\label{0.3.1.2-fr}
No recordaremos la definición de un \emph{prehaz}
$U\to\mathcal{F}(U)$ sobre un espacio topológico $X$, con valores en una
categoría $\boldsymbol{K}$ (G, I, 1.9); diremos que tal prehaz es un
\emph{haz con valores en $\boldsymbol{K}$} si satisface el axioma
siguiente:

\smallskip
\noindent
\textup{(F)} \emph{Para todo recubrimiento $(U_\alpha)$ de un abierto $U$ de $X$
por abiertos $U_\alpha$ contenidos en $U$, si se denota por $\rho_\alpha$
\textup{(resp. $\rho_{\alpha\beta}$)} el morfismo de restricción}
\[
  \mathcal{F}(U)\to\mathcal{F}(U_\alpha)
  \qquad
  \bigl(\text{resp. }\mathcal{F}(U_\alpha)
  \to\mathcal{F}(U_\alpha\cap U_\beta)\bigr),
\]

\oldpage[0\textsubscript{I}]{24}

\emph{el par formado por $\mathcal{F}(U)$ y la familia $(\rho_\alpha)$ es
solución del problema universal para $(\mathcal{F}(U_\alpha))$,
$(\mathcal{F}(U_\alpha\cap U_\beta))$ y $(\rho_{\alpha\beta})$ \textup{(3.1.1)}}%
\footnote{Es un caso particular de la noción general de \emph{límite
proyectivo} (no filtrante) (\emph{véase} (T, I, 1.8) y el libro en preparación
anunciado en la Introducción).}.

Equivale a decir que, para todo objeto $T$ de $\boldsymbol{K}$, la
familia $U\to\operatorname{Hom}(T,\mathcal{F}(U))$ es un \emph{haz
de conjuntos}.
\end{env}

\begin{env}[3.1.3]
\label{0.3.1.3-fr}
Supongamos que $\boldsymbol{K}$ es la categoría definida por una «especie de
estructura con morfismos» $\Sigma$, siendo los objetos de $\boldsymbol{K}$, por tanto,
los conjuntos dotados de estructuras de especie $\Sigma$ y los morfismos los de
$\Sigma$. Supongamos que la categoría $\boldsymbol{K}$ satisface además la
condición siguiente:

\smallskip
\noindent
\textup{(E)} Si $(A,(\rho_\alpha))$ es solución de un problema de aplicación
universal \emph{en la categoría $\boldsymbol{K}$} para unas familias
$(A_\alpha)$, $(A_{\alpha\beta})$, $(\rho_{\alpha\beta})$, entonces es también
solución del problema de aplicación universal para las mismas familias
\emph{en la categoría de conjuntos} (es decir, cuando se consideran $A$,
los $A_\alpha$ y $A_{\alpha\beta}$ como conjuntos, los $\rho_\alpha$ y
$\rho_{\alpha\beta}$ como aplicaciones)%
\footnote{Se puede demostrar que esto significa también que el funtor canónico
$\boldsymbol{K}\to(\mathrm{Ens})$ \emph{conmuta con los límites proyectivos} (no
necesariamente filtrantes).}.

En estas condiciones, la condición (F) implica que, considerado como
prehaz de \emph{conjuntos}, $U\to\mathcal{F}(U)$ es un \emph{haz}.
Además, para que una aplicación $u:T\to\mathcal{F}(U)$ sea un morfismo de
$\boldsymbol{K}$, es necesario y suficiente, en virtud de (F), que cada aplicación
$\rho_\alpha\circ u$ sea un morfismo $T\to\mathcal{F}(U_\alpha)$, lo que
significa que la estructura de especie $\Sigma$ sobre $\mathcal{F}(U)$ es
\emph{estructura inicial} para los morfismos $\rho_\alpha$. Recíprocamente,
supongamos que un prehaz $U\to\mathcal{F}(U)$ sobre $X$, con valores en
$\boldsymbol{K}$, es un \emph{haz de conjuntos} y satisface la condición
anterior; es claro entonces que satisface (F), por lo que es un
\emph{haz con valores en $\boldsymbol{K}$}.
\end{env}

% END tree/worktree/es/CURRENT/EGA_I_ch0/ega0_i_3_1_1_3.tex
% BEGIN tree/worktree/es/CURRENT/EGA_I_ch0/ega0_i_3_1_4_6.tex
% R586: EGA I, 0, 3.1.4--3.1.6.
% Traducción: OpenAI Codex, GPT-6.1 Sol, esfuerzo Ultra.
\begin{env}[3.1.4]
\label{0.3.1.4-fr}
Cuando $\Sigma$ es la especie de estructura de grupo o de anillo, el hecho de que
el prehaz $U\to\mathcal{F}(U)$ con valores en $\boldsymbol{K}$ sea un
haz de \emph{conjuntos} implica \emph{ipso facto} que es un
haz \emph{con valores en $\boldsymbol{K}$} (es decir, un haz de
grupos o de anillos en el sentido de (G))%
\footnote{Esto se debe a que, en la categoría $\boldsymbol{K}$, todo morfismo
que es una \emph{biyección} (como aplicación de conjuntos) es un
\emph{isomorfismo}. Esto ya no es cierto cuando $\boldsymbol{K}$ es, por ejemplo,
la categoría de los espacios topológicos.}. Pero esto deja de ser así
cuando, por ejemplo, $\boldsymbol{K}$ es la categoría de los \emph{anillos
topológicos} (con las representaciones continuas como morfismos): un
haz con valores en $\boldsymbol{K}$ es un haz de anillos
$U\to\mathcal{F}(U)$ tal que, para todo abierto $U$ y todo recubrimiento de $U$
por abiertos $U_\alpha\subset U$, la topología del anillo $\mathcal{F}(U)$
sea \emph{la menos fina} que hace continuas las representaciones
$\mathcal{F}(U)\to\mathcal{F}(U_\alpha)$. En este caso se dirá que
$U\to\mathcal{F}(U)$, considerado como haz de anillos (sin topología), es
\emph{subyacente} al haz de anillos topológicos
$U\to\mathcal{F}(U)$. Por tanto, los morfismos
$u_V:\mathcal{F}(V)\to\mathcal{G}(V)$ ($V$ abierto arbitrario de $X$) de
haces de anillos topológicos son homomorfismos de los haces de
anillos subyacentes tales que $u_V$ es \emph{continuo} para todo abierto
$V\subset X$; para distinguirlos de los homomorfismos cualesquiera de los
haces de anillos subyacentes, se los llamará homomorfismos
\emph{continuos} de haces de anillos topológicos. Se tienen definiciones y
convenciones análogas para los haces de espacios topológicos o de grupos
topológicos.
\end{env}

\oldpage[0\textsubscript{I}]{25}

\begin{env}[3.1.5]
\label{0.3.1.5-fr}
Es claro que, para toda categoría $\boldsymbol{K}$, si $\mathcal{F}$ es un
prehaz (resp. un haz) sobre $X$ con valores en $\boldsymbol{K}$ y $U$ es un
abierto de $X$, los $\mathcal{F}(V)$ para los abiertos $V\subset U$ constituyen
un prehaz (resp. un haz) con valores en $\boldsymbol{K}$, que se
llama prehaz (resp. haz) \emph{inducido} por $\mathcal{F}$ sobre $U$ y
se denota por $\mathcal{F}|U$.

Para todo morfismo $u:\mathcal{F}\to\mathcal{G}$ de prehaces sobre $X$ con
valores en $\boldsymbol{K}$, se denotará por $u|U$ el morfismo
$\mathcal{F}|U\to\mathcal{G}|U$ formado por los $u_V$ para $V\subset U$.
\end{env}

\begin{env}[3.1.6]
\label{0.3.1.6-fr}
Supongamos ahora que la categoría $\boldsymbol{K}$ admite \emph{límites
inductivos} (T, 1.8); entonces, para todo prehaz (y, en particular, todo
haz) $\mathcal{F}$ sobre $X$ con valores en $\boldsymbol{K}$ y todo $x\in X$,
se puede definir el \emph{tallo} $\mathcal{F}_x$ como el objeto de $\boldsymbol{K}$
que es el límite inductivo de los $\mathcal{F}(U)$ según el conjunto filtrante (para
$\supset$) de las vecindades abiertas $U$ de $x$ en $X$, y para los morfismos
$\rho_U^V:\mathcal{F}(V)\to\mathcal{F}(U)$. Si
$u:\mathcal{F}\to\mathcal{G}$ es un morfismo de prehaces con valores en
$\boldsymbol{K}$, se define, para todo $x\in X$, el morfismo
$u_x:\mathcal{F}_x\to\mathcal{G}_x$ como el límite inductivo de los
$u_U:\mathcal{F}(U)\to\mathcal{G}(U)$ según el conjunto de las vecindades abiertas
de $x$; así se define $\mathcal{F}_x$ como funtor covariante en
$\mathcal{F}$, con valores en $\boldsymbol{K}$, para todo $x\in X$.

Cuando $\boldsymbol{K}$ está además definida por una especie de estructura con
morfismos $\Sigma$, se siguen llamando \emph{secciones sobre $U$} de un
\emph{haz $\mathcal{F}$} con valores en $\boldsymbol{K}$ los elementos de
$\mathcal{F}(U)$, y se escribe entonces $\Gamma(U,\mathcal{F})$ en lugar de
$\mathcal{F}(U)$; para $s\in\Gamma(U,\mathcal{F})$ y $V$ abierto contenido en
$U$, se escribe $s|V$ en lugar de $\rho_V^U(s)$; para todo $x\in U$, la imagen
canónica de $s$ en $\mathcal{F}_x$ es el \emph{germen} de $s$ en el punto $x$,
denotado por $s_x$ (\emph{nunca emplearemos la notación $s(x)$ en este sentido},
pues esta notación se reserva para otra noción relativa a los haces
particulares que se considerarán en este Tratado (5.5.1)).

Si entonces $u:\mathcal{F}\to\mathcal{G}$ es un morfismo de haces con valores
en $\boldsymbol{K}$, se escribirá $u(s)$ en lugar de $u_V(s)$ para todo
$s\in\Gamma(V,\mathcal{F})$.

Si $\mathcal{F}$ es un haz de grupos conmutativos, de anillos o de
módulos, se dice que el conjunto de los $x\in X$ tales que
$\mathcal{F}_x\neq\{0\}$ es el \emph{soporte} de $\mathcal{F}$, denotado por
$\operatorname{Supp}(\mathcal{F})$; este conjunto no es necesariamente
cerrado en $X$.

Cuando $\boldsymbol{K}$ está definida por una especie de estructura con morfismos,
\emph{nos abstendremos sistemáticamente de recurrir al punto de
vista de los «espacios étalé»} en lo relativo a los haces con valores en
$\boldsymbol{K}$; es decir, nunca consideraremos un haz como
un espacio topológico (ni siquiera como el conjunto unión de sus tallos), y
tampoco consideraremos un morfismo
$u:\mathcal{F}\to\mathcal{G}$ de tales haces sobre $X$ como una aplicación
continua de espacios topológicos.
\end{env}

% END tree/worktree/es/CURRENT/EGA_I_ch0/ega0_i_3_1_4_6.tex
% BEGIN tree/worktree/es/CURRENT/EGA_I_ch0/ega0_i_3_2_1.tex
% R587: EGA I, 0, 3.2.1.
% Traducción: OpenAI Codex, GPT-6.1 Sol, esfuerzo Ultra.
\subsection{Prehaces sobre una base de abiertos}
\label{subsection:0.3.2-fr}

\begin{env}[3.2.1]
\label{0.3.2.1-fr}
En lo que sigue nos restringiremos a categorías $\boldsymbol{K}$
que admitan \emph{límites proyectivos} (generalizados, es decir,
correspondientes a conjuntos preordenados no necesariamente filtrantes,
cf. (T, 1.8)). Sean $X$ un espacio topológico y $\mathfrak{B}$ una base
de abiertos para la topología de $X$. Llamaremos \emph{prehaz sobre
$\mathfrak{B}$, con valores en $\boldsymbol{K}$}, a una familia de objetos
$\mathcal{F}(U)\in\boldsymbol{K}$, asociados a cada $U\in\mathfrak{B}$, y a una
familia de morfismos
$\rho_U^V:\mathcal{F}(V)\to\mathcal{F}(U)$ definidos para todo par $(U,V)$
de elementos de $\mathfrak{B}$ tales que $U\subset V$,

\oldpage[0\textsubscript{I}]{26}

con las condiciones $\rho_U^U=\text{identidad}$ y
$\rho_U^W=\rho_U^V\circ\rho_V^W$ si $U$, $V$, $W$ en $\mathfrak{B}$ son
tales que $U\subset V\subset W$. Se le puede asociar un \emph{prehaz con
valores en $\boldsymbol{K}$}~: $U\to\mathcal{F}'(U)$ en el sentido usual, tomando
para todo abierto $U$,
\[
  \mathcal{F}'(U)=\varprojlim\mathcal{F}(V),
\]
donde $V$ recorre el conjunto ordenado (para $\subset$, \emph{no filtrante} en
general) de los conjuntos $V\in\mathfrak{B}$ tales que $V\subset U$, pues los
$\mathcal{F}(V)$ forman un sistema proyectivo para los $\rho_V^W$
($V\subset W\subset U$, $V\in\mathfrak{B}$, $W\in\mathfrak{B}$). En efecto,
si $U$, $U'$ son dos abiertos de $X$ tales que $U\subset U'$, se define
$\rho'_U{}^{U'}$ como el límite proyectivo (para $V\subset U$) de los morfismos
canónicos $\mathcal{F}'(U')\to\mathcal{F}(V)$, es decir, como el único
morfismo $\mathcal{F}'(U')\to\mathcal{F}'(U)$ que, compuesto con los
morfismos canónicos $\mathcal{F}'(U)\to\mathcal{F}(V)$, da los morfismos
canónicos $\mathcal{F}'(U')\to\mathcal{F}(V)$; la comprobación de la
transitividad de los $\rho'_U{}^{U'}$ es entonces inmediata. Además, si
$U\in\mathfrak{B}$, el morfismo canónico
$\mathcal{F}'(U)\to\mathcal{F}(U)$ es un isomorfismo que permite identificar
estos dos objetos%
\footnote{Si $X$ es un espacio \emph{noetheriano}, se puede aún definir
$\mathcal{F}'(U)$ y demostrar que es un prehaz (en el sentido usual)
suponiendo solo que $\boldsymbol{K}$ admite límites proyectivos
para los sistemas proyectivos \emph{finitos}. En efecto, si $U$ es un abierto
cualquiera de $X$, existe un recubrimiento \emph{finito} $(V_i)$ de $U$ formado
por conjuntos de $\mathfrak{B}$; para todo par $(i,j)$ de índices, sea
$(V_{ijk})$ un recubrimiento finito de $V_i\cap V_j$ formado por conjuntos de
$\mathfrak{B}$. Sea $I$ el conjunto formado por los $i$ y las ternas $(i,j,k)$,
ordenado por las únicas relaciones $i>(i,j,k)$, $j>(i,j,k)$; se toma entonces
como $\mathcal{F}'(U)$ el límite proyectivo del sistema de los $\mathcal{F}(V_i)$
y $\mathcal{F}(V_{ijk})$; se comprueba fácilmente que no depende de los
recubrimientos $(V_i)$ y $(V_{ijk})$ y que $U\to\mathcal{F}'(U)$ es un
prehaz.}.
\end{env}

% END tree/worktree/es/CURRENT/EGA_I_ch0/ega0_i_3_2_1.tex
% BEGIN tree/worktree/es/CURRENT/EGA_I_ch0/ega0_i_3_2_2.tex
% R588: EGA I, 0, 3.2.2; source-faithful draft, editorial review pending.
% Traducción: OpenAI Codex, GPT-6.1 Sol, esfuerzo Ultra.
\begin{env}[3.2.2]
\label{0.3.2.2-fr}
Para que el prehaz $\mathcal{F}'$ así definido sea un \emph{haz},
es necesario y suficiente que el prehaz $\mathcal{F}$ sobre $\mathfrak{B}$
satisfaga la condición~:

\smallskip
\noindent
\textup{(F\textsubscript{0})} \emph{Para todo recubrimiento $(U_\alpha)$ de
$U\in\mathfrak{B}$ por conjuntos $U_\alpha\in\mathfrak{B}$ contenidos en
$U$, y para todo objeto $T\in\boldsymbol{K}$, la aplicación que a cada
$f\in\operatorname{Hom}(T,\mathcal{F}(U))$ hace corresponder la familia
\[
  (\rho_{U_\alpha}^{U}\circ f)
  \in\prod_\alpha\operatorname{Hom}(T,\mathcal{F}(U_\alpha))
\]
es una biyección de $\operatorname{Hom}(T,\mathcal{F}(U))$ sobre el conjunto de las
familias $(f_\alpha)$ tales que
$\rho_V^{U_\alpha}\circ f_\alpha=\rho_V^{U_\beta}\circ f_\beta$ para todo
par de índices $(\alpha,\beta)$ y todo $V\in\mathfrak{B}$ tal que
$V\subset U_\alpha\cap U_\beta$}%
\footnote{Esto significa también que el par formado por $\mathcal{F}(U)$ y los
$\rho_\alpha=\rho_{U_\alpha}^{U}$ es \emph{solución del problema universal}
definido en (3.1.1) por los datos
$A_\alpha=\mathcal{F}(U_\alpha)$,
$A_{\alpha\beta}=\prod\mathcal{F}(V)$ (para los $V\in\mathfrak{B}$ tales que
$V\subset U_\alpha\cap U_\beta$) y
$\rho_{\alpha\beta}=(\rho_V''):\mathcal{F}(U_\alpha)\to
\prod\mathcal{F}(V)$ definido por la condición de que, para
$V\in\mathfrak{B}$, $V'\in\mathfrak{B}$, $W\in\mathfrak{B}$,
$V\cup V'\subset U_\alpha\cap U_\beta$, $W\subset V\cap V'$,
$\rho_W^V\circ\rho_V''=\rho_W^{V'}\circ\rho_{V'}''$.}.

La condición es evidentemente necesaria. Para demostrar que es suficiente,
consideremos primero una segunda base $\mathfrak{B}'$ de la topología de $X$,
\emph{contenida en $\mathfrak{B}$}, y demostremos que, si $\mathcal{F}''$
denota el prehaz deducido de la subfamilia
$(\mathcal{F}(V))_{V\in\mathfrak{B}'}$, $\mathcal{F}''$ es
\emph{canónicamente isomorfo} a $\mathcal{F}'$. En efecto, en primer lugar el
límite proyectivo (para $V\in\mathfrak{B}'$, $V\subset U$) de los morfismos
canónicos $\mathcal{F}'(U)\to\mathcal{F}(V)$ es un morfismo
$\mathcal{F}'(U)\to\mathcal{F}''(U)$ para todo abierto $U$. Si
$U\in\mathfrak{B}$, este morfismo es un isomorfismo, pues por hipótesis los
morfismos canónicos $\mathcal{F}''(U)\to\mathcal{F}(V)$ para
$V\in\mathfrak{B}'$, $V\subset U$, se factorizan en
\[
  \mathcal{F}''(U)\to\mathcal{F}(U)\to\mathcal{F}(V),
\]
y es inmediato comprobar que las composiciones de los morfismos
$\mathcal{F}(U)\to\mathcal{F}''(U)$ y
$\mathcal{F}''(U)\to\mathcal{F}(U)$ así definidos son las identidades. Así,
para todo abierto $U$, los morfismos
$\mathcal{F}''(U)\to\mathcal{F}''(W)=\mathcal{F}(W)$ para
$W\in\mathfrak{B}$ y $W\subset U$ satisfacen las condiciones que caracterizan el
límite proyectivo de los $\mathcal{F}(W)$ ($W\in\mathfrak{B}$, $W\subset U$), lo
que demuestra nuestra afirmación, habida cuenta de la unicidad de un límite proyectivo
salvo isomorfismo.

Dicho esto, sean $U$ un abierto cualquiera de $X$, $(U_\alpha)$ un recubrimiento
de $U$ por abiertos contenidos en $U$, y sea $\mathfrak{B}'$ la
subfamilia de $\mathfrak{B}$ constituida por los conjuntos

\oldpage[0\textsubscript{I}]{27}

de $\mathfrak{B}$ contenidos en algún $U_\alpha$; es claro que
$\mathfrak{B}'$ es todavía una base de la topología de $U$, por lo que
$\mathcal{F}'(U)$ (resp. $\mathcal{F}'(U_\alpha)$) es el límite proyectivo de los
$\mathcal{F}(V)$ para $V\in\mathfrak{B}'$ y $V\subset U$ (resp.
$V\subset U_\alpha$); el axioma (F) se comprueba entonces inmediatamente por la
definición del límite proyectivo.

Cuando se satisface (F\textsubscript{0}), diremos por abuso de lenguaje que
el prehaz $\mathcal{F}$ sobre la base $\mathfrak{B}$ es \emph{un haz}.
\end{env}

% END tree/worktree/es/CURRENT/EGA_I_ch0/ega0_i_3_2_2.tex
% BEGIN tree/worktree/es/CURRENT/EGA_I_ch0/ega0_i_3_2_3_5.tex
% R589: EGA I, Chapter0, whole3.2.3-5; source-faithful translation.
% Traducción: OpenAI Codex, GPT-6.1 Sol, esfuerzo Ultra.
\begin{env}[3.2.3]
\label{0.3.2.3-fr}
Sean $\mathcal{F}$ y $\mathcal{G}$ dos prehaces sobre la base
$\mathfrak{B}$, con valores en $\boldsymbol{K}$; se define un \emph{morfismo}
$u:\mathcal{F}\to\mathcal{G}$ como una familia
$(u_V)_{V\in\mathfrak{B}}$ de morfismos
$u_V:\mathcal{F}(V)\to\mathcal{G}(V)$ que satisfacen las condiciones
usuales de compatibilidad con los morfismos de restricción $\rho_V^W$. Con
la notación de (3.2.1), se deduce un morfismo
$u':\mathcal{F}'\to\mathcal{G}'$ de los prehaces (ordinarios) correspondientes
tomando como $u_U'$ el límite proyectivo de los $u_V$ para
$V\in\mathfrak{B}$ y $V\subset U$; la verificación de las condiciones de
compatibilidad con los $\rho'_U{}^{U'}$ se deduce del carácter funtorial del
límite proyectivo.
\end{env}

\begin{env}[3.2.4]
\label{0.3.2.4-fr}
Si la categoría $\boldsymbol{K}$ admite límites inductivos y $\mathcal{F}$
es un prehaz sobre la base $\mathfrak{B}$, con valores en $\boldsymbol{K}$,
para todo $x\in X$ las vecindades de $x$ que pertenecen a $\mathfrak{B}$ forman
un conjunto cofinal (para $\supset$) en el conjunto de las vecindades de $x$;
por tanto, si $\mathcal{F}'$ es el prehaz (ordinario) correspondiente a
$\mathcal{F}$, el tallo $\mathcal{F}_x'$ es igual a
\[
  \varinjlim_{\mathfrak{B}}\mathcal{F}(V)
\]
tomado sobre el conjunto de los $V\in\mathfrak{B}$ que contienen a $x$.
Si $u:\mathcal{F}\to\mathcal{G}$ es un morfismo de prehaces sobre
$\mathfrak{B}$ con valores en $\boldsymbol{K}$ y
$u':\mathcal{F}'\to\mathcal{G}'$ es el morfismo correspondiente de prehaces
ordinarios, $u_x'$ es, del mismo modo, el límite inductivo de los morfismos
$u_V:\mathcal{F}(V)\to\mathcal{G}(V)$ para $V\in\mathfrak{B}$ y $x\in V$.
\end{env}

\begin{env}[3.2.5]
\label{0.3.2.5-fr}
Volvamos a las condiciones generales de (3.2.1). Si $\mathcal{F}$ es un
\emph{haz} ordinario con valores en $\boldsymbol{K}$ y $\mathcal{F}_1$ es el
haz \emph{sobre $\mathfrak{B}$} obtenido al restringir $\mathcal{F}$ a
$\mathfrak{B}$, el haz ordinario $\mathcal{F}_1'$ obtenido a partir de
$\mathcal{F}_1$ mediante el procedimiento de (3.2.1) es canónicamente isomorfo a
$\mathcal{F}$, por la condición (F) y las propiedades de unicidad del
límite proyectivo. Habitualmente se identificarán $\mathcal{F}$ y
$\mathcal{F}_1'$.

Si $\mathcal{G}$ es un segundo haz (ordinario) sobre $X$ con valores en
$\boldsymbol{K}$ y $u:\mathcal{F}\to\mathcal{G}$ es un morfismo, la observación
anterior muestra que los datos de los
$u_V:\mathcal{F}(V)\to\mathcal{G}(V)$ \emph{solo para los $V\in\mathfrak{B}$}
determinan completamente $u$; recíprocamente, dados los $u_V$
para $V\in\mathfrak{B}$, basta verificar la conmutatividad del diagrama
con los morfismos de restricción $\rho_V^W$ para $V\in\mathfrak{B}$,
$W\in\mathfrak{B}$ y $V\subset W$, para que exista un morfismo $u'$ y uno
solo de $\mathcal{F}$ a $\mathcal{G}$ tal que $u_V'=u_V$ para todo
$V\in\mathfrak{B}$ (3.2.3).
\end{env}

% END tree/worktree/es/CURRENT/EGA_I_ch0/ega0_i_3_2_3_5.tex
% BEGIN tree/worktree/es/CURRENT/EGA_I_ch0/ega0_i_3_2_6.tex
% R590: EGA I, Chapter0, whole3.2.6 including printed28 continuation.
% Traducción: OpenAI Codex, GPT-6.1 Sol, esfuerzo Ultra.
\begin{env}[3.2.6]
\label{0.3.2.6-fr}
Supongamos siempre que $\boldsymbol{K}$ admite límites proyectivos. Entonces la
categoría de los \emph{haces sobre $X$ con valores en $\boldsymbol{K}$} también admite
\emph{límites proyectivos}; si $(\mathcal{F}_\lambda)$ es un sistema
proyectivo de haces sobre $X$ con valores en $\boldsymbol{K}$, las fórmulas
\[
  \mathcal{F}(U)=\varprojlim_\lambda\mathcal{F}_\lambda(U)
\]
definen, en efecto, un prehaz con valores en $\boldsymbol{K}$, y la
verificación del axioma (F) se deduce de la transitividad de los límites
proyectivos; es inmediato que $\mathcal{F}$ es entonces el límite proyectivo
de los $\mathcal{F}_\lambda$.

Cuando $\boldsymbol{K}$ es la categoría de los conjuntos, para todo sistema
proyectivo $(\mathfrak{H}_\lambda)$ tal

\oldpage[0\textsubscript{I}]{28}

que $\mathfrak{H}_\lambda$ sea un \emph{subhaz} de
$\mathcal{F}_\lambda$ para todo $\lambda$,
$\varprojlim\mathfrak{H}_\lambda$ se identifica canónicamente con un
\emph{subhaz} de $\varprojlim\mathcal{F}_\lambda$. Si
$\boldsymbol{K}$ es la categoría de los grupos conmutativos, el funtor covariante
$\varprojlim\mathcal{F}_\lambda$ es \emph{aditivo} y \emph{exacto por la izquierda}.
\end{env}

% END tree/worktree/es/CURRENT/EGA_I_ch0/ega0_i_3_2_6.tex
% BEGIN tree/worktree/es/CURRENT/EGA_I_ch0/ega0_i_3_3_1.tex
% R591: EGA I, Chapter0, whole section3.3 heading and3.3.1; source-faithful translation.
% Traducción: OpenAI Codex, GPT-6.1 Sol, esfuerzo Ultra.
\subsection{Pegado de haces}
\label{subsection:0.3.3-fr}

\begin{env}[3.3.1]
\label{0.3.3.1-fr}
Supongamos, como antes, que la categoría $\boldsymbol{K}$ admite límites proyectivos
(generalizados). Sean $X$ un espacio topológico y
$\mathfrak{U}=(U_\lambda)_{\lambda\in L}$ un cubrimiento abierto de $X$; para
cada $\lambda\in L$, sea $\mathcal{F}_\lambda$ un haz sobre
$U_\lambda$, con valores en $\boldsymbol{K}$~; para todo par de índices
$(\lambda,\mu)$, supongamos dado un \emph{isomorfismo}
\[
  \theta_{\lambda\mu}:
  \mathcal{F}_\mu|(U_\lambda\cap U_\mu)
  \xrightarrow{\sim}
  \mathcal{F}_\lambda|(U_\lambda\cap U_\mu);
\]
además, supongamos que, para toda terna $(\lambda,\mu,\nu)$, al denotar por
$\theta'_{\lambda\mu}$, $\theta'_{\mu\nu}$, $\theta'_{\lambda\nu}$ las
restricciones de $\theta_{\lambda\mu}$, $\theta_{\mu\nu}$,
$\theta_{\lambda\nu}$ a $U_\lambda\cap U_\mu\cap U_\nu$, se tenga
\[
  \theta'_{\lambda\nu}
  =\theta'_{\lambda\mu}\circ\theta'_{\mu\nu}
\]
(\emph{condición de pegado} para los $\theta_{\lambda\mu}$). Entonces,
existe un haz $\mathcal{F}$ sobre $X$, con valores en $\boldsymbol{K}$, y para
cada $\lambda$ un isomorfismo
\[
  \eta_\lambda:\mathcal{F}|U_\lambda\xrightarrow{\sim}\mathcal{F}_\lambda
\]
tales que, para todo par $(\lambda,\mu)$, al denotar por
$\eta'_\lambda$ y $\eta'_\mu$ las restricciones de $\eta_\lambda$ y
$\eta_\mu$ a $U_\lambda\cap U_\mu$, se tenga
\[
  \theta_{\lambda\mu}=\eta'_\lambda\circ{\eta'_\mu}^{-1};
\]
además, $\mathcal{F}$ y los $\eta_\lambda$ quedan determinados, por estas
condiciones, hasta un único isomorfismo. La unicidad se deduce, en efecto,
inmediatamente de (3.2.5). Para probar la existencia de $\mathcal{F}$, denotemos por
$\mathfrak{B}$ la base de abiertos formada por los abiertos contenidos en al menos un
$U_\lambda$ y, para cada $U\in\mathfrak{B}$, elijamos (mediante la
función $\boldsymbol{\tau}$ de Hilbert) uno de los $\mathcal{F}_\lambda(U)$ para uno de los
$\lambda$ tales que $U\subset U_\lambda$~; si se denota este objeto por
$\mathcal{F}(U)$, los $\rho_U^V$ para $U\subset V$, $U\in\mathfrak{B}$,
$V\in\mathfrak{B}$ se definen de manera evidente (mediante los
$\theta_{\lambda\mu}$), y la condición de transitividad se sigue de la
condición de pegado~; además, la verificación de
(F\textsubscript{0}) es inmediata; por tanto, el prehaz sobre $\mathfrak{B}$
así definido es efectivamente un haz, y de él se obtiene mediante el procedimiento general
(3.2.1) un haz (ordinario), también denotado por $\mathcal{F}$, que resuelve el
problema. Se dice que $\mathcal{F}$ se obtiene por \emph{pegado de los
$\mathcal{F}_\lambda$ mediante los $\theta_{\lambda\mu}$} y habitualmente se identificarán
$\mathcal{F}_\lambda$ y $\mathcal{F}|U_\lambda$ mediante
$\eta_\lambda$.

Es claro que todo haz $\mathcal{F}$ sobre $X$ con valores en
$\boldsymbol{K}$ puede considerarse obtenido por pegado de los haces
$\mathcal{F}_\lambda=\mathcal{F}|U_\lambda$ (donde $(U_\lambda)$ es un
cubrimiento abierto arbitrario de $X$), mediante los isomorfismos
$\theta_{\lambda\mu}$ que se reducen a la identidad.
\end{env}

% END tree/worktree/es/CURRENT/EGA_I_ch0/ega0_i_3_3_1.tex
% BEGIN tree/worktree/es/CURRENT/EGA_I_ch0/ega0_i_3_3_2_3.tex
% R592: whole EGA I Chapter0 3.3.2–3; source-faithful translation.
% OpenAI Codex, GPT-6.1 Sol, Ultra; no human review claimed.
\begin{env}[3.3.2]
\label{0.3.3.2-fr}
Con las mismas notaciones, sea $\mathcal{G}_\lambda$ un segundo haz sobre
$U_\lambda$ (para todo $\lambda\in L$), con valores en $\boldsymbol{K}$, y supongamos
dado para todo par $(\lambda,\mu)$ un isomorfismo
\[
  \omega_{\lambda\mu}:
  \mathcal{G}_\mu|(U_\lambda\cap U_\mu)
  \xrightarrow{\sim}
  \mathcal{G}_\lambda|(U_\lambda\cap U_\mu),
\]
estos isomorfismos verifican la condición de pegado. Supongamos finalmente dado
para cada $\lambda$ un morfismo
$u_\lambda:\mathcal{F}_\lambda\to\mathcal{G}_\lambda$, y que los diagramas
\begin{equation}
\begin{tikzcd}[column sep=huge,row sep=large]
\mathcal{F}_\mu|(U_\lambda\cap U_\mu)
  \arrow[r,"u_\mu"] \arrow[d]
  & \mathcal{G}_\mu|(U_\lambda\cap U_\mu) \arrow[d] \\
\mathcal{F}_\lambda|(U_\lambda\cap U_\mu)
  \arrow[r,"u_\lambda"']
  & \mathcal{G}_\lambda|(U_\lambda\cap U_\mu)
\end{tikzcd}
\tag{3.3.2.1}\label{0.3.3.2.1-fr}
\end{equation}
sean conmutativos. Entonces, si $\mathcal{G}$ se obtiene por pegado de los
$\mathcal{G}_\lambda$ mediante los $\omega_{\lambda\mu}$, existe un
morfismo $u:\mathcal{F}\to\mathcal{G}$ y uno solo tal que los diagramas

\oldpage[0\textsubscript{I}]{29}

\[
\begin{tikzcd}[column sep=huge,row sep=large]
\mathcal{F}|U_\lambda
  \arrow[r,"u|U_\lambda"] \arrow[d]
  & \mathcal{G}|U_\lambda \arrow[d] \\
\mathcal{F}_\lambda
  \arrow[r,"u_\lambda"']
  & \mathcal{G}_\lambda
\end{tikzcd}
\]
sean conmutativos~; esto se deduce inmediatamente de (3.2.3). La correspondencia entre
la familia $(u_\lambda)$ y $u$ es una biyección funtorial del subconjunto de
$\prod_\lambda\operatorname{Hom}(\mathcal{F}_\lambda,\mathcal{G}_\lambda)$
que verifica las condiciones (3.3.2.1) sobre
$\operatorname{Hom}(\mathcal{F},\mathcal{G})$.
\end{env}

\begin{env}[3.3.3]
\label{0.3.3.3-fr}
Con las notaciones de (3.3.1), sea $V$ un abierto de $X$~; es inmediato que
las restricciones a $V\cap U_\lambda\cap U_\mu$ de los $\theta_{\lambda\mu}$
satisfacen la condición de pegado para los haces inducidos
$\mathcal{F}_\lambda|(V\cap U_\lambda)$ y que el haz sobre $V$ obtenido por
pegado de estos últimos se identifica canónicamente con $\mathcal{F}|V$.
\end{env}

% END tree/worktree/es/CURRENT/EGA_I_ch0/ega0_i_3_3_2_3.tex
% BEGIN tree/worktree/es/CURRENT/EGA_I_ch0/ega0_i_3_4_1_3.tex
% R593: whole EGA I Chapter0 3.4 heading and3.4.1–3.
% OpenAI Codex, GPT-6.1 Sol, Ultra; no human review claimed.
\subsection{Imágenes directas de prehaces}
\label{subsection:0.3.4-fr}

\begin{env}[3.4.1]
\label{0.3.4.1-fr}
Sean $X$, $Y$ dos espacios topológicos y $\psi:X\to Y$ una aplicación
continua. Sea $\mathcal{F}$ un prehaz sobre $X$ con valores en una
categoría $\boldsymbol{K}$~; para todo abierto $U\subset Y$, sea
\[
  \mathcal{G}(U)=\mathcal{F}(\psi^{-1}(U)),
\]
y, si $U$, $V$ son dos partes abiertas de $Y$ tales que $U\subset V$, sea
$\rho_U^V$ el morfismo
\[
  \mathcal{F}(\psi^{-1}(V))\to\mathcal{F}(\psi^{-1}(U));
\]
es inmediato que los $\mathcal{G}(U)$ y los $\rho_U^V$ definen un
\emph{prehaz} sobre $Y$ con valores en $\boldsymbol{K}$, que se llama
la \emph{imagen directa de $\mathcal{F}$ por $\psi$} y se denota por
$\psi_*(\mathcal{F})$. Si $\mathcal{F}$ es un haz, se verifica inmediatamente
el axioma (F) para el prehaz $\mathcal{G}(U)$, por lo que
$\psi_*(\mathcal{F})$ es un haz.
\end{env}

\begin{env}[3.4.2]
\label{0.3.4.2-fr}
Sean $\mathcal{F}_1$, $\mathcal{F}_2$ dos prehaces sobre $X$ con valores
en $\boldsymbol{K}$, y sea $u:\mathcal{F}_1\to\mathcal{F}_2$ un morfismo.
Cuando $U$ recorre el conjunto de las partes abiertas de $Y$, la familia de
morfismos
\[
  u_{\psi^{-1}(U)}:
  \mathcal{F}_1(\psi^{-1}(U))\to\mathcal{F}_2(\psi^{-1}(U))
\]
satisface las condiciones de compatibilidad con los morfismos de restricción
y define, por tanto, un morfismo
\[
  \psi_*(u):\psi_*(\mathcal{F}_1)\to\psi_*(\mathcal{F}_2).
\]
Si $v:\mathcal{F}_2\to\mathcal{F}_3$ es un morfismo de $\mathcal{F}_2$ en
un tercer prehaz sobre $X$ con valores en $\boldsymbol{K}$, se tiene
\[
  \psi_*(v\circ u)=\psi_*(v)\circ\psi_*(u);
\]
en otras palabras, $\psi_*(\mathcal{F})$ es un \emph{funtor covariante} en
$\mathcal{F}$, de la categoría de los prehaces (resp. haces) sobre $X$ con
valores en $\boldsymbol{K}$ a la de los prehaces (resp. haces) sobre
$Y$ con valores en $\boldsymbol{K}$.
\end{env}

\begin{env}[3.4.3]
\label{0.3.4.3-fr}
Sean $Z$ un tercer espacio topológico y $\psi':Y\to Z$ una aplicación
continua, y sea $\psi''=\psi'\circ\psi$. Es claro que se tiene
\[
  \psi_*''(\mathcal{F})=\psi_*'(\psi_*(\mathcal{F}))
\]
para todo prehaz $\mathcal{F}$ sobre $X$ con valores en $\boldsymbol{K}$~; además,
para todo morfismo $u:\mathcal{F}\to\mathcal{G}$ entre tales prehaces,
se tiene
\[
  \psi_*''(u)=\psi_*'(\psi_*(u)).
\]
En otras palabras, $\psi_*''$ es el \emph{compuesto} de los funtores $\psi_*'$
y $\psi_*$, lo que puede escribirse
\[
  (\psi'\circ\psi)_*=\psi_*'\circ\psi_*.
\]

Además, para todo abierto $U$ de $Y$, la imagen por la restricción
$\psi|\psi^{-1}(U)$ del prehaz inducido
$\mathcal{F}|\psi^{-1}(U)$ no es otro que el prehaz inducido
$\psi_*(\mathcal{F})|U$.
\end{env}

% END tree/worktree/es/CURRENT/EGA_I_ch0/ega0_i_3_4_1_3.tex
% BEGIN tree/worktree/es/CURRENT/EGA_I_ch0/ega0_i_3_4_4_5.tex
% R594: whole EGA I Chapter0 3.4.4–5.
% OpenAI Codex, GPT-6.1 Sol, Ultra; no human review claimed.
\begin{env}[3.4.4]
\label{0.3.4.4-fr}
Supongamos que la categoría $\boldsymbol{K}$ admite límites inductivos, y sea
$\mathcal{F}$ un prehaz sobre $X$ con valores en $\boldsymbol{K}$~; para todo
$x\in X$, los morfismos
\[
  \Gamma(\psi^{-1}(U),\mathcal{F})\to\mathcal{F}_x
\]
($U$ vecindad abierta de $\psi(x)$ en $Y$) forman un sistema inductivo, que
da, por paso

\oldpage[0\textsubscript{I}]{30}

al límite, un morfismo
\[
  \psi_x:(\psi_*(\mathcal{F}))_{\psi(x)}\to\mathcal{F}_x
\]
de los tallos~; en general, este morfismo no es \emph{ni inyectivo ni sobreyectivo}.
Es funtorial~; en efecto, si
$u:\mathcal{F}_1\to\mathcal{F}_2$ es un morfismo de prehaces sobre $X$ con
valores en $\boldsymbol{K}$, el diagrama
\[
\begin{tikzcd}[column sep=huge,row sep=large]
(\psi_*(\mathcal{F}_1))_{\psi(x)}
  \arrow[r,"\psi_x"] \arrow[d,"{(\psi_*(u))_{\psi(x)}}"']
  & (\mathcal{F}_1)_x \arrow[d,"u_x"] \\
(\psi_*(\mathcal{F}_2))_{\psi(x)}
  \arrow[r,"\psi_x"']
  & (\mathcal{F}_2)_x
\end{tikzcd}
\]
es conmutativo. Si $Z$ es un tercer espacio topológico,
$\psi':Y\to Z$ una aplicación continua y $\psi''=\psi'\circ\psi$, se tiene
\[
  \psi_x''=\psi_x\circ\psi'_{\psi(x)}
\]
para $x\in X$.
\end{env}

\begin{env}[3.4.5]
\label{0.3.4.5-fr}
Bajo las hipótesis de (3.4.4), supongamos además que $\psi$ es un
\emph{homeomorfismo} de $X$ sobre el subespacio $\psi(X)$ de $Y$. Entonces, para
todo $x\in X$, $\psi_x$ es un \emph{isomorfismo}. Esto se aplica en
particular a la inyección canónica $j$ de una parte $X$ de $Y$ en $Y$.
\end{env}

% END tree/worktree/es/CURRENT/EGA_I_ch0/ega0_i_3_4_4_5.tex
% BEGIN tree/worktree/es/CURRENT/EGA_I_ch0/ega0_i_3_4_6.tex
% R595: whole EGA I Chapter0 3.4.6.
% OpenAI Codex, GPT-6.1 Sol, Ultra; no human review claimed.
\begin{env}[3.4.6]
\label{0.3.4.6-fr}
Supongamos que $\boldsymbol{K}$ es la categoría de los grupos, o de los anillos, etc.
Si $\mathcal{F}$ es un haz sobre $X$ con valores en $\boldsymbol{K}$, de
soporte $S$, y si $y\notin\overline{\psi(S)}$, de la definición de
$\psi_*(\mathcal{F})$ se deduce que
\[
  (\psi_*(\mathcal{F}))_y=\{0\},
\]
es decir, el soporte de $\psi_*(\mathcal{F})$ está contenido en
$\overline{\psi(S)}$~; pero no está necesariamente contenido en
$\psi(S)$. Bajo las mismas hipótesis, si $j$ es la inyección canónica de una
parte $X$ de $Y$ en $Y$, el haz $j_*(\mathcal{F})$ induce
$\mathcal{F}$ sobre $X$~; si además $X$ es \emph{cerrada} en $Y$,
$j_*(\mathcal{F})$ es el haz sobre $Y$ que induce $\mathcal{F}$ sobre $X$
y $0$ sobre $Y\setmin X$ (G, II, 2.9.2), pero en general es distinto de este último
cuando se supone que $X$ es localmente cerrada pero no cerrada.
\end{env}

% END tree/worktree/es/CURRENT/EGA_I_ch0/ega0_i_3_4_6.tex
% BEGIN tree/worktree/es/CURRENT/EGA_I_ch0/ega0_i_3_5_1.tex
% R596: whole EGA I Chapter0 3.5 heading and3.5.1.
% OpenAI Codex, GPT-6.1 Sol, Ultra; no human review claimed.
\subsection{Imágenes inversas de prehaces}

\begin{env}[3.5.1]
\label{0.3.5.1-fr}
Bajo las hipótesis de (3.4.1), si $\mathcal{F}$ (resp. $\mathcal{G}$) es un
prehaz sobre $X$ (resp. $Y$) con valores en $\boldsymbol{K}$, todo morfismo
$u:\mathcal{G}\to\psi_*(\mathcal{F})$ de prehaces sobre $Y$ se llama también
un $\psi$-\emph{morfismo} de $\mathcal{G}$ en $\mathcal{F}$, y se denota
asimismo por $\mathcal{G}\to\mathcal{F}$. Se denota también por
$\operatorname{Hom}_\psi(\mathcal{G},\mathcal{F})$ el conjunto
$\operatorname{Hom}_Y(\mathcal{G},\psi_*(\mathcal{F}))$ de los
$\psi$-morfismos de $\mathcal{G}$ en $\mathcal{F}$.

Para todo par $(U,V)$, donde $U$ es un abierto de $X$ y $V$ un abierto de $Y$
tal que $\psi(U)\subset V$, se tiene un morfismo
$u_{U,V}:\mathcal{G}(V)\to\mathcal{F}(U)$ componiendo el morfismo de
restricción
$\mathcal{F}(\psi^{-1}(V))\to\mathcal{F}(U)$ y el morfismo
$u_V:\mathcal{G}(V)\to\psi_*(\mathcal{F})(V)=\mathcal{F}(\psi^{-1}(V))$~; es
inmediato que estos morfismos hacen conmutativos los diagramas
\begin{equation*}
\begin{tikzcd}[column sep=huge,row sep=large]
\mathcal{G}(V) \arrow[r,"u_{U,V}"] \arrow[d]
  & \mathcal{F}(U) \arrow[d] \\
\mathcal{G}(V') \arrow[r,"u_{U',V'}"']
  & \mathcal{F}(U')
\end{tikzcd}
\tag{3.5.1.1}\label{0.3.5.1.1-fr}
\end{equation*}
para $U'\subset U$, $V'\subset V$, $\psi(U')\subset V'$. Recíprocamente, el
dato de una familia $(u_{U,V})$ de morfismos que hacen conmutativos los
diagramas (3.5.1.1) define un $\psi$-morfismo $u$, pues basta tomar
\[
  u_V=u_{\psi^{-1}(V),V}.
\]

\oldpage[0\textsubscript{I}]{31}

Si la categoría $\boldsymbol{K}$ admite límites proyectivos (generalizados), y
si $\mathfrak{B}$, $\mathfrak{B}'$ son bases de las topologías de $X$ y
$Y$ respectivamente, para definir un $\psi$-morfismo $u$ de \emph{haces},
basta dar los $u_{U,V}$ para
$U\in\mathfrak{B}$, $V\in\mathfrak{B}'$ y $\psi(U)\subset V$, que verifican las
condiciones de compatibilidad (3.5.1.1) para $U$, $U'$ en $\mathfrak{B}$ y
$V$, $V'$ en $\mathfrak{B}'$~; basta, en efecto, definir $u_W$, para
todo abierto $W\subset Y$, como límite proyectivo de los $u_{U,V}$ para
$V\in\mathfrak{B}'$ y $V\subset W$, $U\in\mathfrak{B}$ y
$\psi(U)\subset V$.

Cuando la categoría $\boldsymbol{K}$ admite límites inductivos, se tiene, para
todo $x\in X$, un morfismo
\[
  \mathcal{G}(V)\to\mathcal{F}(\psi^{-1}(V))\to\mathcal{F}_x
\]
para toda vecindad abierta $V$ de $\psi(x)$ en $Y$, y estos morfismos
forman un sistema inductivo que da, por paso al límite, un morfismo
\[
  \mathcal{G}_{\psi(x)}\to\mathcal{F}_x.
\]
\end{env}

% END tree/worktree/es/CURRENT/EGA_I_ch0/ega0_i_3_5_1.tex
% BEGIN tree/worktree/es/CURRENT/EGA_I_ch0/ega0_i_3_5_2_3.tex
% R597: whole EGA I Chapter0 3.5.2–3.
% OpenAI Codex, GPT-6.1 Sol, Ultra; no human review claimed.
\begin{env}[3.5.2]
\label{0.3.5.2-fr}
Bajo las hipótesis de (3.4.3), sean $\mathcal{F}$, $\mathcal{G}$,
$\mathcal{H}$ prehaces con valores en $\boldsymbol{K}$ sobre $X$, $Y$, $Z$
respectivamente, y sean
$u:\mathcal{G}\to\psi_*(\mathcal{F})$,
$v:\mathcal{H}\to\psi'_*(\mathcal{G})$ un $\psi$-morfismo y un
$\psi'$-morfismo respectivamente. Se deduce de ellos un $\psi''$-morfismo
\[
  w:\mathcal{H}\xrightarrow{v}\psi'_*(\mathcal{G})
  \xrightarrow{\psi'_*(u)}
  \psi'_*(\psi_*(\mathcal{F}))=\psi''_*(\mathcal{F}),
\]
que se llama, por definición, el \emph{compuesto} de $u$ y de $v$. Se puede,
pues, considerar los pares $(X,\mathcal{F})$ formados por un espacio topológico
$X$ y un prehaz $\mathcal{F}$ sobre $X$ (con valores en $\boldsymbol{K}$)
como una \emph{categoría}, cuyos morfismos son los pares
$(\psi,\theta):(X,\mathcal{F})\to(Y,\mathcal{G})$ formados por una aplicación
continua $\psi:X\to Y$ y un $\psi$-morfismo
$\theta:\mathcal{G}\to\mathcal{F}$.
\end{env}

\begin{env}[3.5.3]
\label{0.3.5.3-fr}
Sean $\psi:X\to Y$ una aplicación continua y $\mathcal{G}$ un
\emph{prehaz} sobre $Y$ con valores en $\boldsymbol{K}$. Llamaremos
\emph{imagen inversa de $\mathcal{G}$ por $\psi$} a un par
$(\mathcal{G}',\rho)$, donde $\mathcal{G}'$ es un \emph{haz} sobre $X$ con
valores en $\boldsymbol{K}$, y
$\rho:\mathcal{G}\to\mathcal{G}'$ es un $\psi$-morfismo (es decir, un
homomorfismo $\mathcal{G}\to\psi_*(\mathcal{G}'))$ tales que, para todo
\emph{haz} $\mathcal{F}$ sobre $X$ con valores en $\boldsymbol{K}$,
la aplicación
\begin{equation*}
  \operatorname{Hom}_X(\mathcal{G}',\mathcal{F})
  \longrightarrow \operatorname{Hom}_\psi(\mathcal{G},\mathcal{F})
  =\operatorname{Hom}_Y(\mathcal{G},\psi_*(\mathcal{F}))
\tag{3.5.3.1}\label{0.3.5.3.1-fr}
\end{equation*}
que transforma $v$ en $\psi_*(v)\circ\rho$ sea una \emph{biyección}~; esta
aplicación, al ser funtorial en $\mathcal{F}$, define entonces un
isomorfismo de funtores en $\mathcal{F}$. Al ser el par
$(\mathcal{G}',\rho)$ solución de un problema universal, se sabe que está
\emph{determinado salvo un isomorfismo único} cuando existe. Se escribirá
entonces $\mathcal{G}'=\psi^*(\mathcal{G})$, $\rho=\rho_{\mathcal{G}}$, y por
abuso de lenguaje se dirá que $\psi^*(\mathcal{G})$ es \emph{el haz imagen
inversa} de $\mathcal{G}$ por $\psi$, entendiéndose que
$\psi^*(\mathcal{G})$ se considera provisto del $\psi$-\emph{morfismo
canónico}
$\rho_{\mathcal{G}}:\mathcal{G}\to\psi^*(\mathcal{G})$, es decir, del
\emph{homomorfismo canónico} de prehaces sobre $Y$~:
\begin{equation*}
  \rho_{\mathcal{G}}:\mathcal{G}\to\psi_*(\psi^*(\mathcal{G})).
\tag{3.5.3.2}\label{0.3.5.3.2-fr}
\end{equation*}
Para todo homomorfismo
$v:\psi^*(\mathcal{G})\to\mathcal{F}$ (donde $\mathcal{F}$ es un haz sobre
$X$ con valores en $\boldsymbol{K}$), se pondrá
$v^\flat=\psi_*(v)\circ\rho_{\mathcal{G}}:
\mathcal{G}\to\psi_*(\mathcal{F})$. Por definición, \emph{todo} morfismo de
prehaces $u:\mathcal{G}\to\psi_*(\mathcal{F})$ es de la forma $v^\flat$ para
un único $v$, que se denotará por $u^\sharp$. En otras palabras, todo
morfismo $u:\mathcal{G}\to\psi_*(\mathcal{F})$ de prehaces se factoriza
de manera única como
\begin{equation*}
  u:\mathcal{G}\xrightarrow{\rho_{\mathcal{G}}}
  \psi_*(\psi^*(\mathcal{G}))
  \xrightarrow{\psi_*(u^\sharp)}\psi_*(\mathcal{F}).
\tag{3.5.3.3}\label{0.3.5.3.3-fr}
\end{equation*}
\end{env}

% END tree/worktree/es/CURRENT/EGA_I_ch0/ega0_i_3_5_2_3.tex
% BEGIN tree/worktree/es/CURRENT/EGA_I_ch0/ega0_i_3_5_4.tex
% R598: whole EGA I Chapter0 3.5.4.
% OpenAI Codex, GPT-6.1 Sol, Ultra; no human review claimed.
\oldpage[0\textsubscript{I}]{32}

\begin{env}[3.5.4]
\label{0.3.5.4-fr}
Supongamos ahora que la categoría $\boldsymbol{K}$ sea tal\footnote{En
el libro citado en la Introducción, daremos condiciones muy
generales sobre la categoría $\boldsymbol{K}$ que aseguren la existencia de las imágenes
inversas de prehaces con valores en $\boldsymbol{K}$.} que \emph{todo}
prehaz $\mathcal{G}$ sobre $Y$ con valores en $\boldsymbol{K}$ admita una
imagen inversa por $\psi$, que denotaremos por $\psi^*(\mathcal{G})$.

Veremos que se puede definir $\psi^*(\mathcal{G})$ como \emph{funtor
covariante} en $\mathcal{G}$, de la categoría de prehaces sobre $Y$ con
valores en $\boldsymbol{K}$ a la de haces sobre $X$ con valores en
$\boldsymbol{K}$, de tal manera que el isomorfismo $v\to v^\flat$ sea un
\emph{isomorfismo de bifuntores}
\begin{equation*}
  \operatorname{Hom}_X(\psi^*(\mathcal{G}),\mathcal{F})
  \xrightarrow{\sim}
  \operatorname{Hom}_Y(\mathcal{G},\psi_*(\mathcal{F}))
\tag{3.5.4.1}\label{0.3.5.4.1-fr}
\end{equation*}
en $\mathcal{G}$ y $\mathcal{F}$.

En efecto, para todo morfismo
$w:\mathcal{G}_1\to\mathcal{G}_2$ de prehaces sobre $Y$ con valores en
$\boldsymbol{K}$, consideremos el morfismo compuesto
\[
  \mathcal{G}_1\xrightarrow{w}\mathcal{G}_2
  \xrightarrow{\rho_{\mathcal{G}_2}}
  \psi_*(\psi^*(\mathcal{G}_2))~;
\]
le corresponde un morfismo
$(\rho_{\mathcal{G}_2}\circ w)^\sharp:
\psi^*(\mathcal{G}_1)\to\psi^*(\mathcal{G}_2)$, que denotaremos por
$\psi^*(w)$. Se tiene, pues, en virtud de (3.5.3.3)
\begin{equation*}
  \psi_*(\psi^*(w))\circ\rho_{\mathcal{G}_1}
  =\rho_{\mathcal{G}_2}\circ w.
\tag{3.5.4.2}\label{0.3.5.4.2-fr}
\end{equation*}
Para todo morfismo $u:\mathcal{G}_2\to\psi_*(\mathcal{F})$, donde
$\mathcal{F}$ es un haz sobre $X$ con valores en $\boldsymbol{K}$, se tiene,
según (3.5.3.3), (3.5.4.2) y la definición de $v^\flat$
\[
  (u^\sharp\circ\psi^*(w))^\flat
  =\psi_*(u^\sharp)\circ\psi_*(\psi^*(w))\circ\rho_{\mathcal{G}_1}
  =\psi_*(u^\sharp)\circ\rho_{\mathcal{G}_2}\circ w
  =u\circ w
\]
o, equivalentemente,
\begin{equation*}
  (u\circ w)^\sharp=u^\sharp\circ\psi^*(w).
\tag{3.5.4.3}\label{0.3.5.4.3-fr}
\end{equation*}
Si se toma, en particular, como $u$ un morfismo
\[
  \mathcal{G}_2\xrightarrow{w'}\mathcal{G}_3
  \xrightarrow{\rho_{\mathcal{G}_3}}
  \psi_*(\psi^*(\mathcal{G}_3)),
\]
resulta
\[
  \psi^*(w'\circ w)
  =(\rho_{\mathcal{G}_3}\circ w'\circ w)^\sharp
  =(\rho_{\mathcal{G}_3}\circ w')^\sharp\circ\psi^*(w)
  =\psi^*(w')\circ\psi^*(w),
\]
de donde se sigue nuestra afirmación.

Finalmente, para todo haz $\mathcal{F}$ sobre $X$ con valores en
$\boldsymbol{K}$, sea $i_{\mathcal{F}}$ el morfismo identidad de
$\psi_*(\mathcal{F})$ y denotemos por
\[
  \sigma_{\mathcal{F}}:
  \psi^*(\psi_*(\mathcal{F}))\to\mathcal{F}
\]
el morfismo $(i_{\mathcal{F}})^\sharp$~; la fórmula (3.5.4.3) da, en
particular, la factorización
\begin{equation*}
  u^\sharp:\psi^*(\mathcal{G})
  \xrightarrow{\psi^*(u)}
  \psi^*(\psi_*(\mathcal{F}))
  \xrightarrow{\sigma_{\mathcal{F}}}\mathcal{F}
\tag{3.5.4.4}\label{0.3.5.4.4-fr}
\end{equation*}
para todo morfismo $u:\mathcal{G}\to\psi_*(\mathcal{F})$. Diremos que el
morfismo $\sigma_{\mathcal{F}}$ es \emph{canónico}.
\end{env}

% END tree/worktree/es/CURRENT/EGA_I_ch0/ega0_i_3_5_4.tex
% BEGIN tree/worktree/es/CURRENT/EGA_I_ch0/ega0_i_3_5_5.tex
% R599: whole EGA I Chapter0 3.5.5.
% OpenAI Codex, GPT-6.1 Sol, Ultra; no human review claimed.
\begin{env}[3.5.5]
\label{0.3.5.5-fr}
Sea $\psi':Y\to Z$ una aplicación continua, y supongamos que todo
prehaz $\mathcal{H}$ sobre $Z$ con valores en $\boldsymbol{K}$ admita una
imagen inversa $\psi'^*(\mathcal{H})$ por $\psi'$. Entonces (con las
hipótesis de (3.5.4)) todo prehaz $\mathcal{H}$ sobre $Z$ con valores en
$\boldsymbol{K}$ admite una imagen inversa por $\psi''=\psi'\circ\psi$ y se tiene
un isomorfismo canónico funtorial
\begin{equation*}
  \psi''^*(\mathcal{H})\xrightarrow{\sim}
  \psi^*(\psi'^*(\mathcal{H}))
\tag{3.5.5.1}\label{0.3.5.5.1-fr}
\end{equation*}

\oldpage[0\textsubscript{I}]{33}

Esto se deduce, en efecto, inmediatamente de las definiciones, teniendo en cuenta que
$\psi''_*=\psi'_*\circ\psi_*$. Además, si
$u:\mathcal{G}\to\psi_*(\mathcal{F})$ es un $\psi$-morfismo,
$v:\mathcal{H}\to\psi'_*(\mathcal{G})$ un $\psi'$-morfismo, y
$w=\psi'_*(u)\circ v$ su compuesto (3.5.2), se comprueba inmediatamente que
$w^\sharp$ es el morfismo compuesto
\[
  w^\sharp:\psi^*(\psi'^*(\mathcal{H}))
  \xrightarrow{\psi^*(v^\sharp)}
  \psi^*(\mathcal{G})
  \xrightarrow{u^\sharp}\mathcal{F}.
\]
\end{env}

% END tree/worktree/es/CURRENT/EGA_I_ch0/ega0_i_3_5_5.tex
% BEGIN tree/worktree/es/CURRENT/EGA_I_ch0/ega0_i_3_5_6.tex
% R600: whole EGA I Chapter0 3.5.6.
% OpenAI Codex, GPT-6.1 Sol, Ultra; no human review claimed.
\begin{env}[3.5.6]
\label{0.3.5.6-fr}
Tomemos en particular como $\psi$ la aplicación identidad
$1_X:X\to X$. Entonces, si existe la imagen inversa por $\psi$ de un prehaz
$\mathcal{F}$ sobre $X$ con valores en $\boldsymbol{K}$, se dice que esta
imagen inversa es el \emph{haz asociado al prehaz $\mathcal{F}$}.
Todo morfismo $u:\mathcal{F}\to\mathcal{F}'$ de $\mathcal{F}$ en un
\emph{haz} $\mathcal{F}'$ con valores en $\boldsymbol{K}$ se factoriza, por tanto,
de manera única como
\[
  \mathcal{F}\xrightarrow{\rho_{\mathcal{F}}}
  1_X^*(\mathcal{F})\xrightarrow{u^\sharp}\mathcal{F}'.
\]
\end{env}

% END tree/worktree/es/CURRENT/EGA_I_ch0/ega0_i_3_5_6.tex
% BEGIN tree/worktree/es/CURRENT/EGA_I_ch0/ega0_i_3_6_1.tex
% EGA I, Chapitre 0, §3.6 and complete3.6.1; French1794-1806.
% GPT-6.1 Sol / Ultra; source-linked draft, R601.
\subsection{Haces simples y haces localmente simples}

\begin{env}[3.6.1]
\label{0.3.6.1-fr}
Diremos que un \emph{prehaz} $\mathcal{F}$ sobre $X$, con valores en
$\boldsymbol{K}$, es \emph{constante} si los morfismos canónicos
$\mathcal{F}(X)\to\mathcal{F}(U)$ son \emph{isomorfismos} para todo
abierto no vacío $U\subset X$~; se observará que $\mathcal{F}$ no es
necesariamente un haz. Se dice que un \emph{haz} es \emph{simple}
si está asociado (3.5.6) a un prehaz constante. Se dice que un haz
$\mathcal{F}$ es \emph{localmente simple} si todo $x\in X$ admite una vecindad
abierto $U$ tal que $\mathcal{F}|U$ sea simple.
\end{env}

% END tree/worktree/es/CURRENT/EGA_I_ch0/ega0_i_3_6_1.tex
% BEGIN tree/worktree/es/CURRENT/EGA_I_ch0/ega0_i_3_6_2.tex
% EGA I Chapter0 complete3.6.2 with entire proof; French1808-1842.
% GPT-6.1 Sol / Ultra; source-linked draft, R602.
\begin{env}[3.6.2]
\label{0.3.6.2-fr}
Supongamos que $X$ es \emph{irreducible} (2.1.1)~; entonces las propiedades
siguientes son equivalentes~:
\begin{enumerate}
  \item[a)] \emph{$\mathcal{F}$ es un prehaz constante sobre $X$}~;
  \item[b)] \emph{$\mathcal{F}$ es un haz simple sobre $X$}~;
  \item[c)] \emph{$\mathcal{F}$ es un haz localmente simple sobre $X$}.
\end{enumerate}

En efecto, sea $\mathcal{F}$ un prehaz constante sobre $X$~; si $U$, $V$
son dos abiertos no vacíos de $X$, $U\cap V$ no es vacío, por lo que
$\mathcal{F}(X)\to\mathcal{F}(U)\to\mathcal{F}(U\cap V)$ y
$\mathcal{F}(X)\to\mathcal{F}(U)$ son isomorfismos, y lo mismo ocurre
con $\mathcal{F}(U)\to\mathcal{F}(U\cap V)$; análogamente,
$\mathcal{F}(V)\to\mathcal{F}(U\cap V)$ es un isomorfismo. De ello se deduce
inmediatamente que el axioma (F) de (3.1.2) se verifica, que $\mathcal{F}$ es
isomorfo a su haz asociado y, por tanto, que \emph{a)} implica \emph{b)}.

Sea ahora $(U_\alpha)$ un recubrimiento abierto de $X$ por abiertos
no vacíos y sea $\mathcal{F}$ un haz sobre $X$ tal que
$\mathcal{F}|U_\alpha$ sea simple para todo $\alpha$~; como $U_\alpha$ es
irreducible, $\mathcal{F}|U_\alpha$ es un prehaz constante en virtud de
lo anterior. Como $U_\alpha\cap U_\beta$ no es vacío,
$\mathcal{F}(U_\alpha)\to\mathcal{F}(U_\alpha\cap U_\beta)$ y
$\mathcal{F}(U_\beta)\to\mathcal{F}(U_\alpha\cap U_\beta)$ son
isomorfismos, de donde se deduce un isomorfismo canónico
$\theta_{\alpha\beta}:\mathcal{F}(U_\alpha)\to\mathcal{F}(U_\beta)$ para todo
par de índices. Pero entonces, al aplicar la condición (F) para $U=X$,
se ve que, para todo índice $\alpha_0$, $\mathcal{F}(U_{\alpha_0})$ y los
$\theta_{\alpha_0\alpha}$ son soluciones del problema universal, lo que (en
virtud de la unicidad) implica que
$\mathcal{F}(X)\to\mathcal{F}(U_{\alpha_0})$ es un isomorfismo y demuestra,
por tanto, que \emph{c)} implica \emph{a)}.
\end{env}

% END tree/worktree/es/CURRENT/EGA_I_ch0/ega0_i_3_6_2.tex
% BEGIN tree/worktree/es/CURRENT/EGA_I_ch0/ega0_i_3_7_1.tex
\oldpage[0\textsubscript{I}]{34}

\subsection{Imágenes inversas de prehaces de grupos o de anillos}

\begin{env}[3.7.1]
\label{0.3.7.1-fr}
Vamos a mostrar que, cuando se toma como $\boldsymbol{K}$ la categoría de los
conjuntos, la imagen inversa por $\psi$ de todo prehaz $\mathcal{G}$ con
valores en $\boldsymbol{K}$ \emph{existe siempre} (con las notaciones y las
hipótesis sobre $X$, $Y$, $\psi$ de (3.5.3)). En efecto, para todo abierto
$U\subset X$, definamos $\mathcal{G}'(U)$ del modo siguiente. Un elemento $s'$
de $\mathcal{G}'(U)$ es una familia $(s'_x)_{x\in U}$, donde
$s'_x\in\mathcal{G}_{\psi(x)}$ para todo $x\in U$ y donde, para todo $x\in U$,
se cumple la condición siguiente: existen una vecindad abierta $V$ de
$\psi(x)$ en $Y$, una vecindad $W\subset\psi^{-1}(V)\cap U$ de $x$ y un
elemento $s\in\mathcal{G}(V)$ tales que $s'_z=s_{\psi(z)}$ para todo $z\in W$.
Se verifica inmediatamente que $U\to\mathcal{G}'(U)$ satisface los axiomas
de los \emph{haces}.

Sea ahora $\mathcal{F}$ un haz de conjuntos sobre $X$, y sean
$u:\mathcal{G}\to\psi_*(\mathcal{F})$,
$v:\mathcal{G}'\to\mathcal{F}$ morfismos. Se definen $u^\sharp$ y $v^\flat$
del modo siguiente. Si $s'$ es una sección de $\mathcal{G}'$ sobre una vecindad
$U$ de $x\in X$ y si $V$ es una vecindad abierta de $\psi(x)$ y
$s\in\mathcal{G}(V)$ es tal que $s'_z=s_{\psi(z)}$ para $z$ en una vecindad de
$x$ contenido en $\psi^{-1}(V)\cap U$, se toma
$u_x^\sharp(s'_x)=u_{\psi(x)}(s_{\psi(x)})$. Del mismo modo, si
$s\in\mathcal{G}(V)$ ($V$ abierto en $Y$), $v^\flat(s)$ es la sección de
$\mathcal{F}$ sobre $\psi^{-1}(V)$, imagen por $v$ de la sección $s'$ de
$\mathcal{G}'$ tal que $s'_x=s_{\psi(x)}$ para todo $x\in\psi^{-1}(V)$.
Además, el homomorfismo canónico (3.5.3)
$\rho:\mathcal{G}\to\psi_*(\psi^*(\mathcal{G}))$ se define del modo siguiente.
Para todo abierto $V\subset Y$ y toda sección $s\in\Gamma(V,\mathcal{G})$,
$\rho(s)$ es la sección $(s_{\psi(x)})_{x\in\psi^{-1}(V)}$ de
$\psi^*(\mathcal{G})$ sobre $\psi^{-1}(V)$. La comprobación de las relaciones
$(u^\sharp)^\flat=u$, $(v^\flat)^\sharp=v$ y $v^\flat=\psi_*(v)\circ\rho$
es inmediata y demuestra nuestra afirmación.

Se verifica que, si $w:\mathcal{G}_1\to\mathcal{G}_2$ es un homomorfismo de
prehaces de conjuntos sobre $Y$, $\psi^*(w)$ se describe explícitamente del
modo siguiente. Si $s'=(s'_x)_{x\in U}$ es una sección de
$\psi^*(\mathcal{G}_1)$ sobre un abierto $U$ de $X$, $(\psi^*(w))(s')$ es
la familia $(w_{\psi(x)}(s'_x))_{x\in U}$. Por último, es inmediato que para
todo abierto $V$ de $Y$, la imagen inversa de $\mathcal{G}|V$ por la
restricción de $\psi$ a $\psi^{-1}(V)$ es idéntica al haz inducido
$\psi^*(\mathcal{G})|\psi^{-1}(V)$.

Cuando $\psi$ es la identidad $1_X$, se recupera la definición de un haz de
conjuntos asociado a un prehaz (G, II, 1.2). Las consideraciones anteriores se
aplican sin cambios cuando $\boldsymbol{K}$ es la categoría de los grupos o
de los anillos (no necesariamente conmutativos).

Cuando $X$ es una parte cualquiera de un espacio topológico $Y$, y $j$ es la
inyección canónica $X\to Y$, para todo haz $\mathcal{G}$ sobre $Y$ con valores
en una categoría $\boldsymbol{K}$ se llama haz \emph{inducido} sobre $X$ por
$\mathcal{G}$ a la imagen inversa $j^*(\mathcal{G})$ (cuando existe). Para los
haces de conjuntos (o de grupos, o de anillos) se recupera la definición
usual (G, II, 1.5).
\end{env}

% END tree/worktree/es/CURRENT/EGA_I_ch0/ega0_i_3_7_1.tex
% BEGIN tree/worktree/es/CURRENT/EGA_I_ch0/ega0_i_3_7_2.tex
\begin{env}[3.7.2]
\label{0.3.7.2-fr}
Conservando las notaciones y las hipótesis de (3.5.3), supongamos que
$\mathcal{G}$ es un \emph{haz} de grupos (resp. de anillos) sobre $Y$.
La definición de las secciones de $\psi^*(\mathcal{G})$ (3.7.1) muestra
(teniendo en cuenta (3.4.4)) que el homomorfismo de tallos
$\psi_x\circ\rho_{\psi(x)}:\mathcal{G}_{\psi(x)}\to
(\psi^*(\mathcal{G}))_x$ es un \emph{isomorfismo funtorial en
$\mathcal{G}$}, que permite identificar estos dos tallos. Con esta
identificación, $u_x^\sharp$ es idéntico al homomorfismo definido en
(3.5.1) y, en particular, se tiene
$\operatorname{Supp}(\psi^*(\mathcal{G}))=
\psi^{-1}(\operatorname{Supp}(\mathcal{G}))$.

Una consecuencia inmediata de este resultado es que \emph{el funtor
$\psi^*(\mathcal{G})$ es exacto en $\mathcal{G}$} en la categoría abeliana
de los haces de grupos conmutativos.
\end{env}

% END tree/worktree/es/CURRENT/EGA_I_ch0/ega0_i_3_7_2.tex
% BEGIN tree/worktree/es/CURRENT/EGA_I_ch0/ega0_i_3_8_1.tex
\oldpage[0\textsubscript{I}]{35}

\subsection{Haces de espacios pseudodiscretos}

\begin{env}[3.8.1]
\label{0.3.8.1-fr}
Sea $X$ un espacio topológico cuya topología admite una base
$\mathfrak{B}$ formada por conjuntos abiertos \emph{cuasicompactos}. Sea
$\mathcal{F}$ un \emph{haz de conjuntos} sobre $X$. Si se dota a cada
$\mathcal{F}(U)$ de la topología \emph{discreta},
$U\to\mathcal{F}(U)$ es un \emph{prehaz de espacios topológicos}.
Vamos a ver que existe un \emph{haz de espacios topológicos
$\mathcal{F}'$ asociado} a $\mathcal{F}$ (3.5.6) tal que
$\Gamma(U,\mathcal{F}')$ es el espacio discreto $\mathcal{F}(U)$ para todo
abierto \emph{cuasicompacto} $U$. Para ello bastará mostrar que el
prehaz $U\to\mathcal{F}(U)$ de espacios topológicos discretos
\emph{sobre $\mathfrak{B}$} verifica la condición $(\mathrm{F}_0)$ de (3.2.2)
y, más generalmente, que si $U$ es un abierto cuasicompacto y
$(U_\alpha)$ es un recubrimiento de $U$ por conjuntos de $\mathfrak{B}$,
la topología menos fina $\mathcal{T}$ sobre $\Gamma(U,\mathcal{F})$ que hace
continuas las aplicaciones
$\Gamma(U,\mathcal{F})\to\Gamma(U_\alpha,\mathcal{F})$ es la topología
\emph{discreta}. Ahora bien, hay un número finito de índices $\alpha_i$ tales
que $U=\bigcup_i U_{\alpha_i}$. Sea $s\in\Gamma(U,\mathcal{F})$ y sea $s_i$
su imagen en $\Gamma(U_{\alpha_i},\mathcal{F})$. La intersección de las
imágenes inversas de los conjuntos $\{s_i\}$ es, por definición, una vecindad
de $s$ para $\mathcal{T}$. Como $\mathcal{F}$ es un haz de conjuntos y los
$U_{\alpha_i}$ recubren $U$, esta intersección se reduce a $s$, lo cual
demuestra nuestra afirmación.

Se observará que, si $U$ es un abierto no cuasicompacto de $X$, el espacio
topológico $\Gamma(U,\mathcal{F}')$ sigue teniendo a $\Gamma(U,\mathcal{F})$
como conjunto subyacente. Su topología, que en general no es discreta, es
la menos fina que hace continuas las aplicaciones
$\Gamma(U,\mathcal{F})\to\Gamma(V,\mathcal{F})$, para $V\in\mathfrak{B}$ y
$V\subset U$ (siendo discretos los $\Gamma(V,\mathcal{F})$).

Las consideraciones precedentes se aplican sin modificación a los haces de
grupos o de anillos (no necesariamente conmutativos), y les asocian
respectivamente haces de \emph{grupos topológicos} o de \emph{anillos
topológicos}. Para abreviar, diremos que el haz $\mathcal{F}'$ es el haz de
\emph{espacios} (resp. de \emph{grupos}, de \emph{anillos})
\emph{pseudodiscretos} asociado al haz de conjuntos (resp. de grupos,
de anillos) $\mathcal{F}$.
\end{env}

% END tree/worktree/es/CURRENT/EGA_I_ch0/ega0_i_3_8_1.tex
% BEGIN tree/worktree/es/CURRENT/EGA_I_ch0/ega0_i_3_8_2.tex
\begin{env}[3.8.2]
\label{0.3.8.2-fr}
Sean $\mathcal{F}$, $\mathcal{G}$ dos haces de conjuntos (resp. de grupos,
de anillos) sobre $X$, y $u:\mathcal{F}\to\mathcal{G}$ un homomorfismo.
Entonces $u$ es también un homomorfismo \emph{continuo}
$\mathcal{F}'\to\mathcal{G}'$, donde $\mathcal{F}'$ y $\mathcal{G}'$ designan
los haces pseudodiscretos asociados a $\mathcal{F}$ y $\mathcal{G}$.
Esto se deduce, en efecto, de (3.2.5).
\end{env}

% END tree/worktree/es/CURRENT/EGA_I_ch0/ega0_i_3_8_2.tex
% BEGIN tree/worktree/es/CURRENT/EGA_I_ch0/ega0_i_3_8_3.tex
\begin{env}[3.8.3]
\label{0.3.8.3-fr}
Sean $\mathcal{F}$ un haz de conjuntos, $\mathcal{H}$ un subhaz de
$\mathcal{F}$, y $\mathcal{F}'$ y $\mathcal{H}'$ los haces pseudodiscretos
asociados a $\mathcal{F}$ y $\mathcal{H}$, respectivamente. Entonces, para
todo abierto $U\subset X$, $\Gamma(U,\mathcal{H}')$ es \emph{cerrado} en
$\Gamma(U,\mathcal{F}')$. En efecto, es la intersección de las imágenes
inversas de los $\Gamma(V,\mathcal{H})$ (para $V\in\mathfrak{B}$,
$V\subset U$) por las aplicaciones continuas
$\Gamma(U,\mathcal{F})\to\Gamma(V,\mathcal{F})$, y
$\Gamma(V,\mathcal{H})$ es cerrado en el espacio discreto
$\Gamma(V,\mathcal{F})$.
\end{env}

% END tree/worktree/es/CURRENT/EGA_I_ch0/ega0_i_3_8_3.tex
% BEGIN tree/worktree/es/CURRENT/EGA_I_ch0/ega0_i_4_1_1.tex
\section{Espacios anillados}
\label{section:0.4-fr}

\subsection{Espacios anillados, $\mathcal{A}$-Módulos, $\mathcal{A}$-Álgebras}
\label{subsection:0.4.1-fr}

\begin{env}[4.1.1]
\label{0.4.1.1-fr}
Un \emph{espacio anillado} (resp. \emph{topológicamente anillado}) es un par
$(X,\mathcal{A})$ formado por un espacio topológico $X$ y un haz de anillos
(no necesariamente conmutativos) (resp. un haz de anillos topológicos)
$\mathcal{A}$ sobre $X$. Se dice que $X$ es el espacio
\oldpage[0\textsubscript{I}]{36}
topológico \emph{subyacente} al espacio anillado $(X,\mathcal{A})$, y que
$\mathcal{A}$ es el \emph{haz estructural}. Este último se denota por
$\mathcal{O}_X$, y su tallo en un punto $x\in X$ se denota por $\mathcal{O}_{X,x}$
o simplemente por $\mathcal{O}_x$ cuando no haya lugar a confusión.

Se designará por $1$ o $e$ la \emph{sección unidad} de $\mathcal{O}_X$
sobre $X$ (elemento unidad de $\Gamma(X,\mathcal{O}_X)$).

Como en este Tratado consideraremos sobre todo haces de anillos
\emph{conmutativos}, se sobreentenderá, cuando hablemos de un espacio
anillado $(X,\mathcal{A})$ sin más precisión, que $\mathcal{A}$ es un haz
de anillos conmutativos.

Los espacios anillados cuyo haz estructural no es necesariamente conmutativo
(resp. los espacios topológicamente anillados) forman una \emph{categoría}
si se define un \emph{morfismo}
$(X,\mathcal{A})\to(Y,\mathcal{B})$ como un par
$(\psi,\theta)=\Psi$ formado por una aplicación continua $\psi:X\to Y$ y un
$\psi$-\emph{morfismo} $\theta:\mathcal{B}\to\mathcal{A}$ (3.5.1) de
haces de anillos (resp. de haces de anillos topológicos).
El \emph{compuesto} de un segundo morfismo
$\Psi'=(\psi',\theta'):(Y,\mathcal{B})\to(Z,\mathcal{C})$ y de $\Psi$, que se denota
por $\Psi''=\Psi'\circ\Psi$, es el morfismo $(\psi'',\theta'')$ en el que
$\psi''=\psi'\circ\psi$, y $\theta''$ es el compuesto de $\theta$ y $\theta'$
(igual a $\psi'_*(\theta)\circ\theta'$, cf. 3.5.2). Para los espacios
anillados, recordemos que entonces se tiene
$\theta''^\sharp=\theta^\sharp\circ\psi^*(\theta'^\sharp)$ (3.5.5). Por tanto, si
$\theta'^\sharp$ y $\theta^\sharp$ son homomorfismos \emph{inyectivos}
(resp. \emph{sobreyectivos}), también lo es $\theta''^\sharp$, pues
$\psi_x\circ\rho_{\psi(x)}$ es un isomorfismo para todo $x\in X$
(3.7.2). De lo anterior se deduce de inmediato que, cuando $\psi$ es una
aplicación continua \emph{inyectiva} y $\theta^\sharp$ es un homomorfismo
\emph{sobreyectivo} de haces de anillos, el morfismo
$(\psi,\theta)$ es un \emph{monomorfismo} (T, 1.1) en la categoría de los
espacios anillados.

Por abuso de lenguaje, se sustituirá a menudo $\psi$ por $\Psi$ en la
notación, por ejemplo escribiendo $\Psi^{-1}(U)$ en lugar de $\psi^{-1}(U)$
para una parte $U$ de $Y$, cuando no haya riesgo de confusión.
\end{env}

% END tree/worktree/es/CURRENT/EGA_I_ch0/ega0_i_4_1_1.tex
% BEGIN tree/worktree/es/CURRENT/EGA_I_ch0/ega0_i_4_1_2.tex
\begin{env}[4.1.2]
\label{0.4.1.2-fr}
Para toda parte $M$ de $X$, el par $(M,\mathcal{A}|M)$ es evidentemente un
espacio anillado, llamado \emph{inducido} sobre $M$ por el espacio anillado
$(X,\mathcal{A})$ (también llamado \emph{restricción} de
$(X,\mathcal{A})$ a $M$). Si $j$ es la inyección canónica $M\to X$ y
$\omega$ la aplicación identidad de $\mathcal{A}|M$, entonces $(j,\omega^\flat)$
es un monomorfismo $(M,\mathcal{A}|M)\to(X,\mathcal{A})$ de espacios anillados,
llamado \emph{inyección canónica}. El compuesto de un morfismo
$\Psi:(X,\mathcal{A})\to(Y,\mathcal{B})$ y de esta inyección se llama
\emph{restricción} de $\Psi$ a $M$.
\end{env}

% END tree/worktree/es/CURRENT/EGA_I_ch0/ega0_i_4_1_2.tex
% BEGIN tree/worktree/es/CURRENT/EGA_I_ch0/ega0_i_4_1_3.tex
\begin{env}[4.1.3]
\label{0.4.1.3-fr}
No volveremos sobre la definición de los $\mathcal{A}$-\emph{Módulos} o
\emph{haces algebraicos} sobre un espacio anillado $(X,\mathcal{A})$
(G, II, 2.2). Cuando $\mathcal{A}$ es un haz de anillos no necesariamente
conmutativos, por $\mathcal{A}$-Módulo se entenderá siempre
«$\mathcal{A}$-Módulo a la izquierda», salvo indicación expresa en contrario.
Los sub-$\mathcal{A}$-Módulos de $\mathcal{A}$ se llamarán
\emph{haces de ideales} (a la izquierda, a la derecha o biláteros) en
$\mathcal{A}$ o $\mathcal{A}$-\emph{Ideales}.

Cuando $\mathcal{A}$ es un haz de anillos conmutativos, si en la definición de
los $\mathcal{A}$-Módulos se sustituye en todas partes la estructura de
\emph{módulo} por la de \emph{álgebra}, se obtiene la definición de una
$\mathcal{A}$-\emph{Álgebra} sobre $X$. Equivalentemente, una
$\mathcal{A}$-Álgebra (no necesariamente conmutativa) es un
$\mathcal{A}$-Módulo $\mathcal{C}$ provisto de un homomorfismo de
$\mathcal{A}$-Módulos
$\varphi:\mathcal{C}\otimes_{\mathcal{A}}\mathcal{C}\to\mathcal{C}$ y de una
sección $e$ sobre $X$ que satisfacen las dos condiciones siguientes.
1\textsuperscript{o} El diagrama
\oldpage[0\textsubscript{I}]{37}
\[
\begin{tikzcd}[column sep=large,row sep=large]
\mathcal{C}\otimes_{\mathcal{A}}\mathcal{C}\otimes_{\mathcal{A}}\mathcal{C}
  \arrow[r,"\varphi\otimes 1"]
  \arrow[d,"1\otimes\varphi"']
&
\mathcal{C}\otimes_{\mathcal{A}}\mathcal{C}
  \arrow[d,"\varphi"]
\\
\mathcal{C}\otimes_{\mathcal{A}}\mathcal{C}
  \arrow[r,"\varphi"']
&
\mathcal{C}
\end{tikzcd}
\]
es conmutativo. 2\textsuperscript{o} Para todo abierto $U\subset X$ y toda
sección $s\in\Gamma(U,\mathcal{C})$, se tiene
$\varphi((e|U)\otimes s)=\varphi(s\otimes(e|U))=s$. Decir que
$\mathcal{C}$ es una $\mathcal{A}$-Álgebra conmutativa equivale a decir,
además, que el diagrama
\[
\begin{tikzcd}[column sep=large,row sep=large]
\mathcal{C}\otimes_{\mathcal{A}}\mathcal{C}
  \arrow[rr,"\sigma"]
  \arrow[dr,"\varphi"']
&&
\mathcal{C}\otimes_{\mathcal{A}}\mathcal{C}
  \arrow[dl,"\varphi"]
\\
& \mathcal{C} &
\end{tikzcd}
\]
es conmutativo, donde $\sigma$ designa la simetría canónica del producto
tensorial $\mathcal{C}\otimes_{\mathcal{A}}\mathcal{C}$.

Los homomorfismos de $\mathcal{A}$-Álgebras se definen del mismo modo que los
homomorfismos de $\mathcal{A}$-Módulos en (G, II, 2.2), aunque, naturalmente,
ya no forman un grupo abeliano.

Si $\mathcal{M}$ es un sub-$\mathcal{A}$-Módulo de una
$\mathcal{A}$-Álgebra $\mathcal{C}$, la \emph{sub-$\mathcal{A}$-Álgebra de
$\mathcal{C}$ engendrada por $\mathcal{M}$} es la suma de las imágenes de los
homomorfismos $\bigotimes^n\mathcal{M}\to\mathcal{C}$ (para los
$n\geqslant 0$). Es también el haz asociado al prehaz
$U\to\mathcal{B}(U)$ de álgebras, donde $\mathcal{B}(U)$ es la subálgebra de
$\Gamma(U,\mathcal{C})$ engendrada por el submódulo
$\Gamma(U,\mathcal{M})$.
\end{env}

% END tree/worktree/es/CURRENT/EGA_I_ch0/ega0_i_4_1_3.tex
% BEGIN tree/worktree/es/CURRENT/EGA_I_ch0/ega0_i_4_1_4.tex
\begin{env}[4.1.4]
\label{0.4.1.4-fr}
Se dice que un haz de anillos $\mathcal{A}$ sobre un espacio topológico $X$
es \emph{reducido} en un punto $x$ de $X$ si el tallo $\mathcal{A}_x$ es un
anillo \emph{reducido} (1.1.1). Se dice que $\mathcal{A}$ es \emph{reducido}
si lo es en todo punto de $X$. Recordemos que un anillo $A$ se llama
\emph{regular} si cada uno de los anillos locales $A_{\mathfrak p}$ (donde
$\mathfrak p$ recorre el conjunto de los ideales primos de $A$) es un anillo
local regular. Diremos que un haz de anillos $\mathcal{A}$ sobre $X$
es \emph{regular en un punto} $x$ (resp. \emph{regular}) si el tallo
$\mathcal{A}_x$ es un anillo regular (resp. si $\mathcal{A}$ es regular
en todo punto). Por último, diremos que un haz de anillos $\mathcal{A}$ sobre
$X$ es \emph{normal en un punto} $x$ (resp. \emph{normal}) si su tallo
$\mathcal{A}_x$ es un anillo \emph{íntegro e integralmente cerrado} (resp. si
$\mathcal{A}$ es normal en todo punto). Diremos que un espacio anillado
$(X,\mathcal{A})$ tiene una de las propiedades anteriores si el haz de anillos
$\mathcal{A}$ tiene esa propiedad.

Un haz de \emph{anillos graduados} $\mathcal{A}$ es, por definición, un
haz de anillos que es suma directa (G, II, 2.7) de una familia
$(\mathcal{A}_n)_{n\in\mathbf Z}$ de haces de grupos abelianos, con las
condiciones $\mathcal{A}_m\mathcal{A}_n\subset\mathcal{A}_{m+n}$. Un
$\mathcal{A}$-\emph{Módulo graduado} es un $\mathcal{A}$-Módulo
$\mathcal{F}$ que es suma directa de una familia
$(\mathcal{F}_n)_{n\in\mathbf Z}$ de haces de grupos abelianos y satisface
las condiciones $\mathcal{A}_m\mathcal{F}_n\subset\mathcal{F}_{m+n}$. Equivalentemente,
se tiene
$(\mathcal{A}_m)_x(\mathcal{A}_n)_x\subset(\mathcal{A}_{m+n})_x$
(resp.
$(\mathcal{A}_m)_x(\mathcal{F}_n)_x\subset(\mathcal{F}_{m+n})_x$)
en todo punto $x$.
\end{env}

% END tree/worktree/es/CURRENT/EGA_I_ch0/ega0_i_4_1_4.tex
% BEGIN tree/worktree/es/CURRENT/EGA_I_ch0/ega0_i_4_1_5.tex
\begin{env}[4.1.5]
\label{0.4.1.5-fr}
Dado un espacio anillado $(X,\mathcal{A})$ (no necesariamente conmutativo),
no recordaremos aquí las definiciones de los bifuntores
$\mathcal{F}\otimes_{\mathcal{A}}\mathcal{G}$,
$\mathcal{H}\!om_{\mathcal{A}}(\mathcal{F},\mathcal{G})$ y
$\operatorname{Hom}_{\mathcal{A}}(\mathcal{F},\mathcal{G})$
(G, II, 2.8 y 2.2) en las categorías de los $\mathcal{A}$-Módulos a la izquierda
o a la derecha (según el caso), con valores en la categoría de los haces de
grupos abelianos (o, en términos más
\oldpage[0\textsubscript{I}]{38}
generales, en la de los $\mathcal{C}$-Módulos, si $\mathcal{C}$ es el centro de
$\mathcal{A}$). El tallo
$(\mathcal{F}\otimes_{\mathcal{A}}\mathcal{G})_x$ en todo punto $x\in X$
se identifica canónicamente con
$\mathcal{F}_x\otimes_{\mathcal{A}_x}\mathcal{G}_x$, y se define un
homomorfismo canónico funtorial
\[
  (\mathcal{H}\!om_{\mathcal{A}}(\mathcal{F},\mathcal{G}))_x
  \longrightarrow
  \operatorname{Hom}_{\mathcal{A}_x}(\mathcal{F}_x,\mathcal{G}_x)
\]
que, en general, no es ni inyectivo ni sobreyectivo. Los bifuntores que
acabamos de considerar son aditivos y, en particular, conmutan con las sumas
directas finitas. El funtor $\mathcal{F}\otimes_{\mathcal{A}}\mathcal{G}$ es
exacto por la derecha en $\mathcal{F}$ y en $\mathcal{G}$, conmuta con los
límites inductivos, y $\mathcal{A}\otimes_{\mathcal{A}}\mathcal{G}$
(resp. $\mathcal{F}\otimes_{\mathcal{A}}\mathcal{A}$) se identifica
canónicamente con $\mathcal{G}$ (resp. $\mathcal{F}$). Los funtores
$\mathcal{H}\!om_{\mathcal{A}}(\mathcal{F},\mathcal{G})$ y
$\operatorname{Hom}_{\mathcal{A}}(\mathcal{F},\mathcal{G})$ son
\emph{exactos por la izquierda} en $\mathcal{F}$ y $\mathcal{G}$.
Más precisamente, si se tiene una sucesión exacta de la forma
$0\to\mathcal{G}'\to\mathcal{G}\to\mathcal{G}''$, la sucesión
\[
  0\longrightarrow
  \mathcal{H}\!om_{\mathcal{A}}(\mathcal{F},\mathcal{G}')
  \longrightarrow
  \mathcal{H}\!om_{\mathcal{A}}(\mathcal{F},\mathcal{G})
  \longrightarrow
  \mathcal{H}\!om_{\mathcal{A}}(\mathcal{F},\mathcal{G}'')
\]
es exacta, y si se tiene una sucesión exacta de la forma
$\mathcal{F}'\to\mathcal{F}\to\mathcal{F}''\to0$, la sucesión
\[
  0\longrightarrow
  \mathcal{H}\!om_{\mathcal{A}}(\mathcal{F}'',\mathcal{G})
  \longrightarrow
  \mathcal{H}\!om_{\mathcal{A}}(\mathcal{F},\mathcal{G})
  \longrightarrow
  \mathcal{H}\!om_{\mathcal{A}}(\mathcal{F}',\mathcal{G})
\]
es exacta. El funtor $\operatorname{Hom}$ tiene propiedades análogas.
Además, $\mathcal{H}\!om_{\mathcal{A}}(\mathcal{A},\mathcal{G})$ se identifica
canónicamente con $\mathcal{G}$. Por último, para todo abierto $U\subset X$,
se tiene
\[
  \Gamma\bigl(U,\mathcal{H}\!om_{\mathcal{A}}(\mathcal{F},\mathcal{G})\bigr)
  =
  \operatorname{Hom}_{\mathcal{A}|U}(\mathcal{F}|U,\mathcal{G}|U).
\]

Para todo $\mathcal{A}$-Módulo a la izquierda $\mathcal{F}$ (resp. a la
derecha), se llama \emph{dual} de $\mathcal{F}$ y se denota por
$\check{\mathcal{F}}$ el $\mathcal{A}$-Módulo a la derecha (resp. a la
izquierda) $\mathcal{H}\!om_{\mathcal{A}}(\mathcal{F},\mathcal{A})$.

Por último, si $\mathcal{A}$ es un haz de anillos conmutativos y
$\mathcal{F}$ es un $\mathcal{A}$-Módulo,
$U\to\bigwedge^p\Gamma(U,\mathcal{F})$ es un prehaz cuyo haz asociado es un
$\mathcal{A}$-Módulo que se denota por
$\bigwedge^p\mathcal{F}$ y se llama la \emph{potencia exterior $p$-ésima
de $\mathcal{F}$}. Se comprueba fácilmente que la aplicación canónica del
prehaz $U\to\bigwedge^p\Gamma(U,\mathcal{F})$ en el haz asociado
$\bigwedge^p\mathcal{F}$ es \emph{inyectiva} y que, para todo
$x\in X$, se tiene
$(\bigwedge^p\mathcal{F})_x=\bigwedge^p(\mathcal{F}_x)$. Es claro que
$\bigwedge^p\mathcal{F}$ es un funtor covariante en $\mathcal{F}$.
\end{env}

% END tree/worktree/es/CURRENT/EGA_I_ch0/ega0_i_4_1_5.tex
% BEGIN tree/worktree/es/CURRENT/EGA_I_ch0/ega0_i_4_1_6.tex
\begin{env}[4.1.6]
\label{0.4.1.6-fr}
Supongamos que $\mathcal{A}$ es un haz de anillos no necesariamente
conmutativos, que $\mathcal{J}$ es un haz de ideales a la izquierda de
$\mathcal{A}$ y que $\mathcal{F}$ es un $\mathcal{A}$-Módulo a la izquierda.
Se denota entonces por $\mathcal{J}\mathcal{F}$ el sub-$\mathcal{A}$-Módulo de
$\mathcal{F}$ que es la imagen de $\mathcal{J}\otimes_{\mathbf Z}\mathcal{F}$
(donde $\mathbf Z$ es el haz asociado al prehaz constante $U\to\mathbf Z$) por
la aplicación canónica $\mathcal{J}\otimes_{\mathbf Z}\mathcal{F}\to\mathcal{F}$.
Es claro que, para todo $x\in X$, se tiene
$(\mathcal{J}\mathcal{F})_x=\mathcal{J}_x\mathcal{F}_x$. Cuando
$\mathcal{A}$ es conmutativo, $\mathcal{J}\mathcal{F}$ es también la imagen
canónica de $\mathcal{J}\otimes_{\mathcal{A}}\mathcal{F}\to\mathcal{F}$.
Es inmediato que $\mathcal{J}\mathcal{F}$ es también el $\mathcal{A}$-Módulo
asociado al prehaz $U\to\Gamma(U,\mathcal{J})\Gamma(U,\mathcal{F})$. Si
$\mathcal{J}_1$, $\mathcal{J}_2$ son dos haces de ideales a la izquierda de
$\mathcal{A}$, se tiene
$\mathcal{J}_1(\mathcal{J}_2\mathcal{F})=(\mathcal{J}_1\mathcal{J}_2)\mathcal{F}$.
\end{env}

% END tree/worktree/es/CURRENT/EGA_I_ch0/ega0_i_4_1_6.tex
% BEGIN tree/worktree/es/CURRENT/EGA_I_ch0/ega0_i_4_1_7.tex
\begin{env}[4.1.7]
\label{0.4.1.7-fr}
Sea $(X_\lambda,\mathcal{A}_\lambda)_{\lambda\in L}$ una familia de espacios
anillados. Para todo par $(\lambda,\mu)$, supongamos dados una parte abierta
$V_{\lambda\mu}$ de $X_\lambda$ y un isomorfismo de espacios anillados
\[
  \varphi_{\lambda\mu}:
  (V_{\mu\lambda},\mathcal{A}_\mu|V_{\mu\lambda})
  \xrightarrow{\sim}
  (V_{\lambda\mu},\mathcal{A}_\lambda|V_{\lambda\mu}),
\]
con $V_{\lambda\lambda}=X_\lambda$ y $\varphi_{\lambda\lambda}$ igual a la
identidad. Supongamos además que, para toda terna $(\lambda,\mu,\nu)$, al
denotar por $\varphi'_{\mu\lambda}$ la restricción de
$\varphi_{\mu\lambda}$ a $V_{\lambda\mu}\cap V_{\lambda\nu}$,
$\varphi'_{\mu\lambda}$ sea un isomorfismo de
\[
  (V_{\lambda\mu}\cap V_{\lambda\nu},
   \mathcal{A}_\lambda|(V_{\lambda\mu}\cap V_{\lambda\nu}))
\]
sobre
\[
  (V_{\mu\nu}\cap V_{\mu\lambda},
   \mathcal{A}_\mu|(V_{\mu\nu}\cap V_{\mu\lambda}))
\]
y que se tenga
$\varphi'_{\lambda\nu}=\varphi'_{\lambda\mu}\circ\varphi'_{\mu\nu}$
(\emph{condición de pegado} para los $\varphi_{\lambda\mu}$). Se puede
considerar entonces, en primer lugar, el espacio topológico obtenido por
pegado (mediante los $\varphi_{\lambda\mu}$) de los $X_\lambda$
\oldpage[0\textsubscript{I}]{39}
a lo largo de los $V_{\lambda\mu}$. Si se identifica $X_\lambda$ con la parte
abierta $X'_\lambda$ correspondiente en $X$, las hipótesis implican que los
tres conjuntos $V_{\lambda\mu}\cap V_{\lambda\nu}$,
$V_{\mu\nu}\cap V_{\mu\lambda}$,
$V_{\nu\lambda}\cap V_{\nu\mu}$ se identifican con
$X'_\lambda\cap X'_\mu\cap X'_\nu$. Se puede entonces transportar a
$X'_\lambda$ la estructura de espacio anillado de $X_\lambda$. Si
$\mathcal{A}'_\lambda$ es el haz de anillos obtenido al transportar
$\mathcal{A}_\lambda$, los $\mathcal{A}'_\lambda$ satisfacen la condición de
pegado (3.3.1) y, por tanto, definen un haz de anillos $\mathcal{A}$ sobre
$X$. Se dice que $(X,\mathcal{A})$ es el espacio anillado obtenido por
\emph{pegado de los $(X_\lambda,\mathcal{A}_\lambda)$ a lo largo de los
$V_{\lambda\mu}$}, mediante los $\varphi_{\lambda\mu}$.
\end{env}

% END tree/worktree/es/CURRENT/EGA_I_ch0/ega0_i_4_1_7.tex
% BEGIN tree/worktree/es/CURRENT/EGA_I_ch0/ega0_i_4_2_1.tex
\subsection{Imagen directa de un $\mathcal{A}$-Módulo}
\label{subsection:0.4.2-fr}

\begin{env}[4.2.1]
\label{0.4.2.1-fr}
Sean $(X,\mathcal{A})$, $(Y,\mathcal{B})$ dos espacios anillados y
$\Psi=(\psi,\theta)$ un morfismo
$(X,\mathcal{A})\to(Y,\mathcal{B})$. Entonces $\psi_*(\mathcal{A})$ es un
haz de anillos sobre $Y$ y $\theta$ un homomorfismo
$\mathcal{B}\to\psi_*(\mathcal{A})$ de haces de anillos. Sea ahora
$\mathcal{F}$ un $\mathcal{A}$-Módulo. Su imagen directa
$\psi_*(\mathcal{F})$ es un haz de grupos abelianos sobre $Y$. Además,
para todo abierto $U\subset Y$,
\[
  \Gamma(U,\psi_*(\mathcal{F}))
  =\Gamma(\psi^{-1}(U),\mathcal{F})
\]
está provisto de una estructura de módulo sobre el anillo
$\Gamma(U,\psi_*(\mathcal{A}))=\Gamma(\psi^{-1}(U),\mathcal{A})$.
Como las aplicaciones bilineales que definen estas estructuras son
compatibles con las operaciones de restricción, definen sobre
$\psi_*(\mathcal{F})$ una estructura de $\psi_*(\mathcal{A})$-Módulo.
El homomorfismo $\theta:\mathcal{B}\to\psi_*(\mathcal{A})$ permite entonces
definir también sobre $\psi_*(\mathcal{F})$ una estructura de
$\mathcal{B}$-\emph{Módulo}. Diremos que este $\mathcal{B}$-Módulo es la
\emph{imagen directa de $\mathcal{F}$ por el morfismo $\Psi$}, y lo
denotaremos por $\Psi_*(\mathcal{F})$. Si $\mathcal{F}_1$, $\mathcal{F}_2$
son dos $\mathcal{A}$-Módulos sobre $X$ y $u$ es un
$\mathcal{A}$-homomorfismo $\mathcal{F}_1\to\mathcal{F}_2$, es inmediato
(al considerar las secciones sobre los abiertos de $Y$) que $\psi_*(u)$
es un $\psi_*(\mathcal{A})$-homomorfismo
$\psi_*(\mathcal{F}_1)\to\psi_*(\mathcal{F}_2)$ y, \emph{a fortiori}, un
$\mathcal{B}$-homomorfismo
$\Psi_*(\mathcal{F}_1)\to\Psi_*(\mathcal{F}_2)$. Como
$\mathcal{B}$-homomorfismo, lo denotaremos por $\Psi_*(u)$. Se ve, pues, que
$\Psi_*$ es un \emph{funtor covariante} de la categoría de los
$\mathcal{A}$-Módulos a la de los $\mathcal{B}$-Módulos. Además, es
inmediato que este funtor es \emph{exacto por la izquierda} (G, II, 2.12).

En $\psi_*(\mathcal{A})$, la estructura de $\mathcal{B}$-Módulo y la
estructura de haz de anillos definen una estructura de
$\mathcal{B}$-Álgebra. Se denotará por $\Psi_*(\mathcal{A})$ esta
$\mathcal{B}$-Álgebra.
\end{env}

% END tree/worktree/es/CURRENT/EGA_I_ch0/ega0_i_4_2_1.tex
% BEGIN tree/worktree/es/CURRENT/EGA_I_ch0/ega0_i_4_2_2.tex
\begin{env}[4.2.2]
\label{0.4.2.2-fr}
Sean $\mathcal{M}$, $\mathcal{N}$ dos $\mathcal{A}$-Módulos. Para todo
abierto $U$ de $Y$, se tiene una aplicación canónica
\[
  \Gamma(\psi^{-1}(U),\mathcal{M})
  \times\Gamma(\psi^{-1}(U),\mathcal{N})
  \longrightarrow
  \Gamma(\psi^{-1}(U),\mathcal{M}\otimes_{\mathcal{A}}\mathcal{N})
\]
que es bilineal sobre el anillo
$\Gamma(\psi^{-1}(U),\mathcal{A})=\Gamma(U,\psi_*(\mathcal{A}))$ y,
\emph{a fortiori}, sobre $\Gamma(U,\mathcal{B})$. Define, por tanto,
un homomorfismo
\[
  \Gamma(U,\Psi_*(\mathcal{M}))
  \otimes_{\Gamma(U,\mathcal{B})}
  \Gamma(U,\Psi_*(\mathcal{N}))
  \longrightarrow
  \Gamma(U,\Psi_*(\mathcal{M}\otimes_{\mathcal{A}}\mathcal{N}))
\]
y, como se comprueba de inmediato que estos homomorfismos son compatibles
con las operaciones de restricción, dan un homomorfismo canónico funtorial
de $\mathcal{B}$-Módulos
\[
  \Psi_*(\mathcal{M})\otimes_{\mathcal{B}}\Psi_*(\mathcal{N})
  \longrightarrow
  \Psi_*(\mathcal{M}\otimes_{\mathcal{A}}\mathcal{N})
  \tag{4.2.2.1}\label{0.4.2.2.1-fr}
\]
\oldpage[0\textsubscript{I}]{40}
que, en general, no será ni inyectivo ni sobreyectivo. Si $\mathcal{P}$ es
un tercer $\mathcal{A}$-Módulo, se comprueba de inmediato que el diagrama
\[
\begin{tikzcd}[column sep=large,row sep=large]
\Psi_*(\mathcal{M})\otimes_{\mathcal{B}}\Psi_*(\mathcal{N})
  \otimes_{\mathcal{B}}\Psi_*(\mathcal{P})
  \arrow[r]
  \arrow[d]
&
\Psi_*(\mathcal{M}\otimes_{\mathcal{A}}\mathcal{N})
  \otimes_{\mathcal{B}}\Psi_*(\mathcal{P})
  \arrow[d]
\\
\Psi_*(\mathcal{M})\otimes_{\mathcal{B}}
  \Psi_*(\mathcal{N}\otimes_{\mathcal{A}}\mathcal{P})
  \arrow[r]
&
\Psi_*(\mathcal{M}\otimes_{\mathcal{A}}\mathcal{N}
  \otimes_{\mathcal{A}}\mathcal{P})
\end{tikzcd}
\tag{4.2.2.2}\label{0.4.2.2.2-fr}
\]
es conmutativo.
\end{env}

% END tree/worktree/es/CURRENT/EGA_I_ch0/ega0_i_4_2_2.tex
% BEGIN tree/worktree/es/CURRENT/EGA_I_ch0/ega0_i_4_2_3.tex
\begin{env}[4.2.3]
\label{0.4.2.3-fr}
Sean $\mathcal{M}$, $\mathcal{N}$ dos $\mathcal{A}$-Módulos. Para todo
abierto $U\subset Y$, se tiene por definición
\[
  \Gamma\bigl(\psi^{-1}(U),
    \mathcal{H}\!om_{\mathcal{A}}(\mathcal{M},\mathcal{N})\bigr)
  =
  \operatorname{Hom}_{\mathcal{A}|V}
    (\mathcal{M}|V,\mathcal{N}|V),
\]
donde se ha puesto, para abreviar, $V=\psi^{-1}(U)$. La aplicación
$u\to\Psi_*(u)$ es un homomorfismo
\[
  \operatorname{Hom}_{\mathcal{A}|V}(\mathcal{M}|V,\mathcal{N}|V)
  \longrightarrow
  \operatorname{Hom}_{\mathcal{B}|U}
    (\Psi_*(\mathcal{M})|U,\Psi_*(\mathcal{N})|U)
\]
para las estructuras de $\Gamma(U,\mathcal{B})$-módulos. Como estos
homomorfismos son compatibles con las operaciones de restricción, definen
un homomorfismo canónico funtorial de $\mathcal{B}$-Módulos
\[
  \Psi_*\bigl(\mathcal{H}\!om_{\mathcal{A}}(\mathcal{M},\mathcal{N})\bigr)
  \longrightarrow
  \mathcal{H}\!om_{\mathcal{B}}
    (\Psi_*(\mathcal{M}),\Psi_*(\mathcal{N})).
  \tag{4.2.3.1}\label{0.4.2.3.1-fr}
\]
\end{env}

% END tree/worktree/es/CURRENT/EGA_I_ch0/ega0_i_4_2_3.tex
% BEGIN tree/worktree/es/CURRENT/EGA_I_ch0/ega0_i_4_2_4.tex
\begin{env}[4.2.4]
\label{0.4.2.4-fr}
Si $\mathcal{C}$ es una $\mathcal{A}$-Álgebra, el homomorfismo compuesto
\[
  \Psi_*(\mathcal{C})\otimes_{\mathcal{B}}\Psi_*(\mathcal{C})
  \longrightarrow
  \Psi_*(\mathcal{C}\otimes_{\mathcal{A}}\mathcal{C})
  \longrightarrow
  \Psi_*(\mathcal{C})
\]
define sobre $\Psi_*(\mathcal{C})$ una estructura de
$\mathcal{B}$-\emph{Álgebra}, como se deduce de (4.2.2.2). Del mismo modo,
se ve que, si $\mathcal{M}$ es un $\mathcal{C}$-Módulo, $\Psi_*(\mathcal{M})$
está provisto canónicamente de una estructura de
$\Psi_*(\mathcal{C})$-Módulo.
\end{env}

% END tree/worktree/es/CURRENT/EGA_I_ch0/ega0_i_4_2_4.tex
% BEGIN tree/worktree/es/CURRENT/EGA_I_ch0/ega0_i_4_2_5.tex
\begin{env}[4.2.5]
\label{0.4.2.5-fr}
Consideremos en particular el caso en que $X$ es un subespacio
\emph{cerrado} de $Y$ y $\psi$ es la inyección canónica $j:X\to Y$. Si
$\mathcal{B}'=\mathcal{B}|X=j^*(\mathcal{B})$ es la restricción del haz
de anillos $\mathcal{B}$ a $X$, un $\mathcal{A}$-Módulo $\mathcal{M}$ puede
considerarse como $\mathcal{B}'$-Módulo mediante el homomorfismo
$\theta^\sharp:\mathcal{B}'\to\mathcal{A}$. Entonces $\Psi_*(\mathcal{M})$
es el $\mathcal{B}$-Módulo que induce $\mathcal{M}$ sobre $X$ y $0$ fuera
de él. Si $\mathcal{N}$ es un segundo $\mathcal{A}$-Módulo,
$\Psi_*(\mathcal{M})\otimes_{\mathcal{B}}\Psi_*(\mathcal{N})$ se identifica,
por tanto, con $\Psi_*(\mathcal{M}\otimes_{\mathcal{B}'}\mathcal{N})$ y
$\mathcal{H}\!om_{\mathcal{B}}(\Psi_*(\mathcal{M}),\Psi_*(\mathcal{N}))$
con
$\Psi_*\bigl(\mathcal{H}\!om_{\mathcal{B}'}(\mathcal{M},\mathcal{N})\bigr)$.
\end{env}

% END tree/worktree/es/CURRENT/EGA_I_ch0/ega0_i_4_2_5.tex
% BEGIN tree/worktree/es/CURRENT/EGA_I_ch0/ega0_i_4_2_6.tex
\begin{env}[4.2.6]
\label{0.4.2.6-fr}
Sean $(Z,\mathcal{C})$ un tercer espacio anillado y
$\Psi'=(\psi',\theta')$ un morfismo
$(Y,\mathcal{B})\to(Z,\mathcal{C})$. Si $\Psi''$ es el morfismo compuesto
$\Psi'\circ\Psi$, es claro que se tiene
$\Psi''_*=\Psi'_*\circ\Psi_*$.
\end{env}

% END tree/worktree/es/CURRENT/EGA_I_ch0/ega0_i_4_2_6.tex
% BEGIN tree/worktree/es/CURRENT/EGA_I_ch0/ega0_i_4_3_1.tex
\subsection{Imagen inversa de un $\mathcal{B}$-Módulo}
\label{subsection:0.4.3-fr}

\begin{env}[4.3.1]
\label{0.4.3.1-fr}
Con las hipótesis y la notación de (4.2.1), sean $\mathcal{G}$ un
$\mathcal{B}$-Módulo y $\psi^*(\mathcal{G})$ su imagen inversa (3.7.1), que
es, por tanto, un haz de grupos abelianos sobre $X$. La definición de las
secciones de $\psi^*(\mathcal{G})$ y de $\psi^*(\mathcal{B})$ (3.7.1)
muestra que $\psi^*(\mathcal{G})$ está provisto canónicamente de una
estructura de $\psi^*(\mathcal{B})$-Módulo. Por otra parte, el
homomorfismo $\theta^\sharp:\psi^*(\mathcal{B})\to\mathcal{A}$ dota a
$\mathcal{A}$ de una estructura de $\psi^*(\mathcal{B})$-Módulo. Para
evitar confusiones, cuando sea preciso, lo denotaremos por
$\mathcal{A}_{[\theta]}$. El producto tensorial
$\psi^*(\mathcal{G})\otimes_{\psi^*(\mathcal{B})}\mathcal{A}_{[\theta]}$
está provisto entonces de una estructura de $\mathcal{A}$-Módulo.
Diremos que este $\mathcal{A}$-Módulo es \emph{la imagen
\oldpage[0\textsubscript{I}]{41}
inversa de $\mathcal{G}$ por el morfismo}
$\Psi$ y lo denotaremos por $\Psi^*(\mathcal{G})$. Si $\mathcal{G}_1$,
$\mathcal{G}_2$ son dos $\mathcal{B}$-Módulos sobre $Y$ y $v$ es un
$\mathcal{B}$-homomorfismo $\mathcal{G}_1\to\mathcal{G}_2$,
$\psi^*(v)$, como se comprueba de inmediato, es un
$\psi^*(\mathcal{B})$-homomorfismo de $\psi^*(\mathcal{G}_1)$ en
$\psi^*(\mathcal{G}_2)$. Por consiguiente, $\psi^*(v)\otimes1$ es un
$\mathcal{A}$-homomorfismo $\Psi^*(\mathcal{G}_1)\to\Psi^*(\mathcal{G}_2)$,
que denotaremos por $\Psi^*(v)$. Así queda definido $\Psi^*$ como un
\emph{funtor covariante} de la categoría de los $\mathcal{B}$-Módulos a la
de los $\mathcal{A}$-Módulos. Aquí, a diferencia de $\psi^*$, este funtor
solo es \emph{exacto por la derecha} en general, ya que el producto tensorial
con $\mathcal{A}$ es un funtor exacto por la derecha en la categoría de los
$\psi^*(\mathcal{B})$-Módulos.

Para todo $x\in X$, se tiene
\[
  (\Psi^*(\mathcal{G}))_x
  =\mathcal{G}_{\psi(x)}\otimes_{\mathcal{B}_{\psi(x)}}\mathcal{A}_x,
\]
en virtud de (3.7.2). Por tanto, el soporte de $\Psi^*(\mathcal{G})$ está
contenido en $\psi^{-1}(\operatorname{Supp}\mathcal{G})$.
\end{env}

% END tree/worktree/es/CURRENT/EGA_I_ch0/ega0_i_4_3_1.tex
% BEGIN tree/worktree/es/CURRENT/EGA_I_ch0/ega0_i_4_3_2.tex
\begin{env}[4.3.2]
\label{0.4.3.2-fr}
Sea $(\mathcal{G}_\lambda)$ un sistema inductivo de $\mathcal{B}$-Módulos
y sea $\mathcal{G}=\varinjlim\mathcal{G}_\lambda$ su límite inductivo. Los
homomorfismos canónicos $\mathcal{G}_\lambda\to\mathcal{G}$ definen
$\psi^*(\mathcal{B})$-homomorfismos
$\psi^*(\mathcal{G}_\lambda)\to\psi^*(\mathcal{G})$, que dan un
homomorfismo canónico
\[
  \varinjlim\psi^*(\mathcal{G}_\lambda)
  \longrightarrow
  \psi^*(\mathcal{G}).
\]
Como el tallo en un punto de un límite inductivo de haces es el límite
inductivo de los tallos en ese mismo punto (G, II, 1.11), el homomorfismo
canónico anterior es \emph{biyectivo} (3.7.2). Además, el producto tensorial
conmuta con los límites inductivos de haces, por lo que se tiene un
\emph{isomorfismo canónico funtorial}
\[
  \varinjlim\Psi^*(\mathcal{G}_\lambda)
  \xrightarrow{\sim}
  \Psi^*\bigl(\varinjlim\mathcal{G}_\lambda\bigr)
\]
de $\mathcal{A}$-Módulos.

Por otra parte, para una suma directa finita $\bigoplus_i\mathcal{G}_i$ de
$\mathcal{B}$-Módulos, es claro que
\[
  \psi^*\bigl(\bigoplus_i\mathcal{G}_i\bigr)
  =\bigoplus_i\psi^*(\mathcal{G}_i),
\]
de modo que, al tomar el producto tensorial con $\mathcal{A}_{[\theta]}$,
\[
  \Psi^*\bigl(\bigoplus_i\mathcal{G}_i\bigr)
  =\bigoplus_i\Psi^*(\mathcal{G}_i).
  \tag{4.3.2.1}\label{0.4.3.2.1-fr}
\]
Al pasar al límite inductivo, se deduce, en virtud de lo anterior, que la
igualdad anterior sigue siendo válida para una suma directa
\emph{cualquiera}.
\end{env}

% END tree/worktree/es/CURRENT/EGA_I_ch0/ega0_i_4_3_2.tex
% BEGIN tree/worktree/es/CURRENT/EGA_I_ch0/ega0_i_4_3_3.tex
\begin{env}[4.3.3]
\EGAReaderTarget{canon-ref.0.4.3.3-fr}{0.4.3.3-fr}
Sean $\mathcal{G}_1$, $\mathcal{G}_2$ dos $\mathcal{B}$-Módulos. De la
definición de las imágenes inversas de haces de grupos abelianos (3.7.1)
se deduce de inmediato un homomorfismo canónico
\[
  \psi^*(\mathcal{G}_1)\otimes_{\psi^*(\mathcal{B})}\psi^*(\mathcal{G}_2)
  \longrightarrow
  \psi^*(\mathcal{G}_1\otimes_{\mathcal{B}}\mathcal{G}_2)
\]
de $\psi^*(\mathcal{B})$-Módulos. Como el tallo en un punto de un producto
tensorial de haces es el producto tensorial de los tallos en ese punto
(G, II, 2.8), se deduce de (3.7.2) que el homomorfismo anterior es, de
hecho, un \emph{isomorfismo}. Al tomar el producto tensorial con
$\mathcal{A}$, se obtiene un \emph{isomorfismo canónico} funtorial
\[
  \Psi^*(\mathcal{G}_1)\otimes_{\mathcal{A}}\Psi^*(\mathcal{G}_2)
  \xrightarrow{\sim}
  \Psi^*(\mathcal{G}_1\otimes_{\mathcal{B}}\mathcal{G}_2).
  \tag{4.3.3.1}\label{0.4.3.3.1-fr}
\]
\end{env}

% END tree/worktree/es/CURRENT/EGA_I_ch0/ega0_i_4_3_3.tex
% BEGIN tree/worktree/es/CURRENT/EGA_I_ch0/ega0_i_4_3_4.tex
\begin{env}[4.3.4]
\label{0.4.3.4-fr}
Sea $\mathcal{C}$ una $\mathcal{B}$-Álgebra. Dar la estructura de
Álgebra sobre $\mathcal{C}$ equivale a dar un $\mathcal{B}$-homomorfismo
\[
  \mathcal{C}\otimes_{\mathcal{B}}\mathcal{C}\longrightarrow\mathcal{C}
\]
que satisfaga las condiciones de asociatividad y de conmutatividad
(condiciones que se verifican en cada tallo). El isomorfismo anterior
permite transportar este homomorfismo a un homomorfismo de
$\mathcal{A}$-Módulos
\[
  \Psi^*(\mathcal{C})\otimes_{\mathcal{A}}\Psi^*(\mathcal{C})
  \longrightarrow\Psi^*(\mathcal{C})
\]
que satisface las mismas condiciones. Así, $\Psi^*(\mathcal{C})$ queda
provisto de una estructura de $\mathcal{A}$-Álgebra. En particular, de
las definiciones se deduce de inmediato que la $\mathcal{A}$-Álgebra
$\Psi^*(\mathcal{B})$ es \emph{igual a} $\mathcal{A}$ (salvo isomorfismo
canónico).

Del mismo modo, si $\mathcal{M}$ es un $\mathcal{C}$-Módulo, dar esta
estructura de Módulo equivale
\oldpage[0\textsubscript{I}]{42}
a dar un $\mathcal{B}$-homomorfismo
$\mathcal{C}\otimes_{\mathcal{B}}\mathcal{M}\to\mathcal{M}$ que satisfaga
la condición de asociatividad. Por transporte se obtiene así una
estructura de $\Psi^*(\mathcal{C})$-Módulo sobre $\Psi^*(\mathcal{M})$.
\end{env}

% END tree/worktree/es/CURRENT/EGA_I_ch0/ega0_i_4_3_4.tex
% BEGIN tree/worktree/es/CURRENT/EGA_I_ch0/ega0_i_4_3_5.tex
\begin{env}[4.3.5]
\label{0.4.3.5-fr}
Sea $\mathcal{J}$ un haz de ideales de $\mathcal{B}$. Como el funtor
$\psi^*$ es exacto, el $\psi^*(\mathcal{B})$-Módulo $\psi^*(\mathcal{J})$
se identifica canónicamente con un haz de ideales de $\psi^*(\mathcal{B})$.
La inyección canónica $\psi^*(\mathcal{J})\to\psi^*(\mathcal{B})$ da
entonces un homomorfismo de $\mathcal{A}$-Módulos
\[
  \Psi^*(\mathcal{J})
  =\psi^*(\mathcal{J})\otimes_{\psi^*(\mathcal{B})}\mathcal{A}_{[\theta]}
  \longrightarrow\mathcal{A}~;
\]
denotaremos por $\Psi^*(\mathcal{J})\mathcal{A}$, o por
$\mathcal{J}\mathcal{A}$ si no hay riesgo de confusión, la imagen de
$\Psi^*(\mathcal{J})$ por este homomorfismo. Se tiene, pues, por definición
\[
  \mathcal{J}\mathcal{A}
  =\theta^\sharp\bigl(\psi^*(\mathcal{J})\bigr)\mathcal{A}
\]
y, en particular, para todo $x\in X$,
\[
  (\mathcal{J}\mathcal{A})_x
  =\theta_x^\sharp(\mathcal{J}_{\psi(x)})\mathcal{A}_x,
\]
teniendo en cuenta la identificación canónica de los tallos de
$\psi^*(\mathcal{J})$ con los de $\mathcal{J}$ (3.7.2). Si
$\mathcal{J}_1$, $\mathcal{J}_2$ son dos ideales de $\mathcal{B}$, se tiene
\[
  (\mathcal{J}_1\mathcal{J}_2)\mathcal{A}
  =\mathcal{J}_1(\mathcal{J}_2\mathcal{A})
  =(\mathcal{J}_1\mathcal{A})(\mathcal{J}_2\mathcal{A}).
\]
Si $\mathcal{F}$ es un $\mathcal{A}$-Módulo, se pondrá
$\mathcal{J}\mathcal{F}=(\mathcal{J}\mathcal{A})\mathcal{F}$.
\end{env}

% END tree/worktree/es/CURRENT/EGA_I_ch0/ega0_i_4_3_5.tex
% BEGIN tree/worktree/es/CURRENT/EGA_I_ch0/ega0_i_4_3_6.tex
\begin{env}[4.3.6]
\label{0.4.3.6-fr}
Sean $(Z,\mathcal{C})$ un tercer espacio anillado y
$\Psi'=(\psi',\theta')$ un morfismo $(Y,\mathcal{B})\to(Z,\mathcal{C})$.
Si $\Psi''$ es el morfismo compuesto $\Psi'\circ\Psi$, de la definición
(4.3.1) y de (4.3.3.1) se deduce que
${\Psi''}^*=\Psi^*\circ{\Psi'}^*$.
\end{env}

% END tree/worktree/es/CURRENT/EGA_I_ch0/ega0_i_4_3_6.tex
% BEGIN tree/worktree/es/CURRENT/EGA_I_ch0/ega0_i_4_4_1.tex
\subsection{Relaciones entre imágenes directas e imágenes inversas}
\label{subsection:0.4.4-fr}

\begin{env}[4.4.1]
\label{0.4.4.1-fr}
Con las hipótesis y la notación de (4.2.1), sea $\mathcal{G}$ un
$\mathcal{B}$-Módulo. Por definición, un homomorfismo
$u:\mathcal{G}\to\Psi_*(\mathcal{F})$ de $\mathcal{B}$-Módulos se llama
también \emph{$\Psi$-morfismo de $\mathcal{G}$ en $\mathcal{F}$}, o
simplemente \emph{homomorfismo de $\mathcal{G}$ en $\mathcal{F}$}, y se
escribe $u:\mathcal{G}\to\mathcal{F}$ cuando no hay lugar a confusión.
Dar tal homomorfismo equivale a dar, para todo par $(U,V)$ en el que
$U$ es un abierto de $X$ y $V$ un abierto de $Y$ con
$\psi(U)\subset V$, un \emph{homomorfismo}
\[
  u_{U,V}:\Gamma(V,\mathcal{G})\longrightarrow\Gamma(U,\mathcal{F})
\]
\emph{de $\Gamma(V,\mathcal{B})$-módulos}, donde $\Gamma(U,\mathcal{F})$
se considera como $\Gamma(V,\mathcal{B})$-módulo mediante el homomorfismo
de anillos
$\theta_{U,V}:\Gamma(V,\mathcal{B})\to\Gamma(U,\mathcal{A})$. Además,
los $u_{U,V}$ deben hacer conmutativos los diagramas (3.5.1.1). Por lo demás,
para definir $u$ basta dar los $u_{U,V}$ cuando $U$ (resp. $V$) recorre
una base $\mathfrak{B}$ (resp. $\mathfrak{B}'$) de la topología de $X$
(resp. $Y$) y verificar la conmutatividad de (3.5.1.1) para estas
restricciones.
\end{env}

% END tree/worktree/es/CURRENT/EGA_I_ch0/ega0_i_4_4_1.tex
% BEGIN tree/worktree/es/CURRENT/EGA_I_ch0/ega0_i_4_4_2.tex
\begin{env}[4.4.2]
\label{0.4.4.2-fr}
Bajo las hipótesis de (4.2.1) y (4.2.6), sean $\mathcal{H}$ un
$\mathcal{C}$-Módulo y $v:\mathcal{H}\to\Psi'_*(\mathcal{G})$ un
$\Psi'$-morfismo. Entonces
\[
  w:\mathcal{H}\xrightarrow{v}\Psi'_*(\mathcal{G})
  \xrightarrow{\Psi'_*(u)}\Psi'_*(\Psi_*(\mathcal{F}))
\]
es un $\Psi''$-morfismo que se llama el \emph{compuesto} de $u$ y $v$.
\end{env}

% END tree/worktree/es/CURRENT/EGA_I_ch0/ega0_i_4_4_2.tex
% BEGIN tree/worktree/es/CURRENT/EGA_I_ch0/ega0_i_4_4_3_8.tex
\begin{env}[4.4.3]
\label{0.4.4.3-fr}
Vamos a ver ahora que se puede definir un \emph{isomorfismo}
canónico de \emph{bifuntores} en $\mathcal{F}$ y $\mathcal{G}$
\[
  \operatorname{Hom}_{\mathcal{A}}(\Psi^*(\mathcal{G}),\mathcal{F})
  \xrightarrow{\sim}
  \operatorname{Hom}_{\mathcal{B}}(\mathcal{G},\Psi_*(\mathcal{F}))
  \tag{4.4.3.1}\label{0.4.4.3.1-fr}
\]
que designaremos por $v\to v_\theta^\flat$ (o simplemente
$v\to v^\flat$ si no es posible ninguna confusión); denotaremos por
$u\to u_\theta^\sharp$, o $u\to u^\sharp$ el isomorfismo inverso.
Esta definición es la siguiente: al componer
$v:\Psi^*(\mathcal{G})\to\mathcal{F}$ con la aplicación canónica
$\psi^*(\mathcal{G})\to\Psi^*(\mathcal{G})$, se obtiene un homomorfismo de
haces de grupos $v':\psi^*(\mathcal{G})\to\mathcal{F}$, que es también un
homomorfismo de $\psi^*(\mathcal{B})$-módulos. Se deduce de él (3.7.1) un
homomorfismo
${v'}^\flat:\mathcal{G}\to\psi_*(\mathcal{F})=\Psi_*(\mathcal{F})$, que es
también un homomorfismo de $\mathcal{B}$-Módulos como se
\oldpage[0\textsubscript{I}]{43}
verifica sin dificultad; se toma $v_\theta^\flat={v'}^\flat$. Del mismo modo, de
$u:\mathcal{G}\to\Psi_*(\mathcal{F})$, que es un homomorfismo de
$\mathcal{B}$-Módulos, se deduce (3.7.1) un homomorfismo
$u^\sharp:\psi^*(\mathcal{G})\to\mathcal{F}$ de
$\psi^*(\mathcal{B})$-Módulos, de donde, por tensorización con $\mathcal{A}$, un
homomorfismo de $\mathcal{A}$-Módulos
$\Psi^*(\mathcal{G})\to\mathcal{F}$, que designaremos por
$u_\theta^\sharp$. Es inmediato verificar que
$(u_\theta^\sharp)_\theta^\flat=u$ y
$(v_\theta^\flat)_\theta^\sharp=v$, así como el carácter funtorial en
$\mathcal{F}$ del isomorfismo $v\to v_\theta^\flat$. El carácter
funtorial en $\mathcal{G}$ de $u\to u_\theta^\sharp$ se deduce entonces
formalmente como en (3.5.4) (razonamiento que demostraría también el
carácter funtorial de $\Psi^*$ establecido directamente en (4.3.1)).

Si se toma como $v$ el homomorfismo identidad de $\Psi^*(\mathcal{G})$,
$v_\theta^\flat$ es un homomorfismo
\[
  \rho_{\mathcal{G}}:\mathcal{G}\longrightarrow
  \Psi_*(\Psi^*(\mathcal{G}))~;
  \tag{4.4.3.2}\label{0.4.4.3.2-fr}
\]
si se toma como $u$ el homomorfismo identidad de $\Psi_*(\mathcal{F})$,
$u_\theta^\sharp$ es un homomorfismo
\[
  \sigma_{\mathcal{F}}:\Psi^*(\Psi_*(\mathcal{F}))
  \longrightarrow\mathcal{F}~;
  \tag{4.4.3.3}\label{0.4.4.3.3-fr}
\]
estos homomorfismos se llamarán \emph{canónicos}. En general no son ni
inyectivos ni sobreyectivos. Se tienen factorizaciones canónicas análogas a
(3.5.3.3) y (3.5.4.4).

Se observará que si $s$ es una sección de $\mathcal{G}$ sobre un abierto
$V$ de $Y$, $\rho_{\mathcal{G}}(s)$ es la sección $s'\otimes1$ de
$\Psi^*(\mathcal{G})$ sobre $\psi^{-1}(V)$, siendo $s'$ tal que
$s'_x=s_{\psi(x)}$ para todo $x\in\psi^{-1}(V)$. Observemos también que si
$u:\mathcal{G}\to\psi_*(\mathcal{F})$ es un homomorfismo, define para
todo $x\in X$ un homomorfismo
$u_x:\mathcal{G}_{\psi(x)}\to\mathcal{F}_x$ sobre los tallos, obtenido al
componer
$(u^\sharp)_x:(\Psi^*(\mathcal{G}))_x\to\mathcal{F}_x$ y el homomorfismo
canónico $s_x\to s_x\otimes1$ de $\mathcal{G}_{\psi(x)}$ en
$(\Psi^*(\mathcal{G}))_x
=\mathcal{G}_{\psi(x)}\otimes_{\mathcal{B}_{\psi(x)}}\mathcal{A}_x$.
El homomorfismo $u_x$ se obtiene también por paso al límite inductivo a
partir de los homomorfismos
\[
  \Gamma(V,\mathcal{G})\xrightarrow{u}
  \Gamma(\psi^{-1}(V),\mathcal{F})\longrightarrow\mathcal{F}_x,
\]
donde $V$ recorre las vecindades de $\psi(x)$.
\end{env}

\begin{env}[4.4.4]
\label{0.4.4.4-fr}
Sean $\mathcal{F}_1$, $\mathcal{F}_2$ unos $\mathcal{A}$-Módulos,
$\mathcal{G}_1$, $\mathcal{G}_2$ unos $\mathcal{B}$-Módulos, $u_i$ ($i=1,2$)
un homomorfismo de $\mathcal{G}_i$ en $\mathcal{F}_i$. Designaremos por
$u_1\otimes u_2$ el homomorfismo
\[
  u:\mathcal{G}_1\otimes_{\mathcal{B}}\mathcal{G}_2
  \longrightarrow
  \mathcal{F}_1\otimes_{\mathcal{A}}\mathcal{F}_2
\]
tal que $u^\sharp=(u_1)^\sharp\otimes(u_2)^\sharp$ (teniendo en cuenta
(4.3.3.1)); se verifica que $u$ es también el compuesto
\[
  \mathcal{G}_1\otimes_{\mathcal{B}}\mathcal{G}_2
  \longrightarrow
  \Psi_*(\mathcal{F}_1)\otimes_{\mathcal{B}}\Psi_*(\mathcal{F}_2)
  \longrightarrow
  \Psi_*(\mathcal{F}_1\otimes_{\mathcal{A}}\mathcal{F}_2),
\]
donde la primera flecha es el producto tensorial ordinario
$u_1\otimes_{\mathcal{B}}u_2$ y la segunda el homomorfismo canónico
(4.2.2.1).
\end{env}

\begin{env}[4.4.5]
\label{0.4.4.5-fr}
Sean $(\mathcal{G}_\lambda)_{\lambda\in L}$ un sistema inductivo de
$\mathcal{B}$-Módulos, y, para todo $\lambda\in L$, sea $u_\lambda$ un
homomorfismo $\mathcal{G}_\lambda\to\Psi_*(\mathcal{F})$, que forman un sistema
inductivo; pongamos
$\mathcal{G}=\varinjlim\mathcal{G}_\lambda$ y
$u=\varinjlim u_\lambda$; entonces los $(u_\lambda)^\sharp$ forman un sistema
inductivo de homomorfismos $\Psi^*(\mathcal{G}_\lambda)\to\mathcal{F}$, y el
límite inductivo de este sistema no es otro que $u^\sharp$.
\end{env}

\begin{env}[4.4.6]
\label{0.4.4.6-fr}
Sean $\mathcal{M}$, $\mathcal{N}$ dos $\mathcal{B}$-Módulos, $V$ un abierto
de $Y$, $U=\psi^{-1}(V)$; la aplicación $v\to\Psi^*(v)$ es un
homomorfismo
\[
  \operatorname{Hom}_{\mathcal{B}|V}(\mathcal{M}|V,\mathcal{N}|V)
  \longrightarrow
  \operatorname{Hom}_{\mathcal{A}|U}
  (\Psi^*(\mathcal{M})|U,\Psi^*(\mathcal{N})|U)
\]
\begin{sloppypar}
para las estructuras de $\Gamma(V,\mathcal{B})$-módulo
($\operatorname{Hom}_{\mathcal{A}|U}
(\Psi^*(\mathcal{M})|U,\allowbreak\Psi^*(\mathcal{N})|U)$ está normalmente provisto de una
estructura de $\Gamma(U,\psi^*(\mathcal{B}))$-módulo, y gracias al
homomorfismo
\oldpage[0\textsubscript{I}]{44}
canónico (3.7.2)
$\Gamma(V,\mathcal{B})\to\Gamma(U,\psi^*(\mathcal{B}))$, es, por tanto, también un
$\Gamma(V,\mathcal{B})$-módulo).
\end{sloppypar}
Se ve inmediatamente que estos homomorfismos son
compatibles con las restricciones, y por consiguiente definen un homomorfismo
canónico funtorial
\[
  \gamma:\mathcal{H}\!om_{\mathcal{B}}(\mathcal{M},\mathcal{N})
  \longrightarrow
  \Psi_*\bigl(\mathcal{H}\!om_{\mathcal{A}}
  (\Psi^*(\mathcal{M}),\Psi^*(\mathcal{N}))\bigr)~;
\]
a este homomorfismo le corresponde, por tanto, también el homomorfismo
\[
  \gamma^\sharp:
  \Psi^*\bigl(\mathcal{H}\!om_{\mathcal{B}}(\mathcal{M},\mathcal{N})\bigr)
  \longrightarrow
  \mathcal{H}\!om_{\mathcal{A}}
  (\Psi^*(\mathcal{M}),\Psi^*(\mathcal{N}))
\]
y estos homomorfismos canónicos son funtoriales en $\mathcal{M}$ y
$\mathcal{N}$.
\end{env}

\begin{env}[4.4.7]
\label{0.4.4.7-fr}
Supongamos que $\mathcal{F}$ (resp. $\mathcal{G}$) es una
$\mathcal{A}$-Álgebra (resp. una $\mathcal{B}$-Álgebra). Si
$u:\mathcal{G}\to\Psi_*(\mathcal{F})$ es un homomorfismo de
$\mathcal{B}$-Álgebras, $u^\sharp$ es un homomorfismo
$\Psi^*(\mathcal{G})\to\mathcal{F}$ de $\mathcal{A}$-Álgebras; esto resulta
de la conmutatividad del diagrama
\[
\begin{tikzcd}[column sep=huge,row sep=large]
\mathcal{G}\otimes_{\mathcal{B}}\mathcal{G}
  \arrow[r] \arrow[d]
  & \mathcal{G} \arrow[d,"u"] \\
\Psi_*(\mathcal{F}\otimes_{\mathcal{A}}\mathcal{F})
  \arrow[r]
  & \Psi_*(\mathcal{F})
\end{tikzcd}
\]
y de (4.4.4). Del mismo modo, si $v:\Psi^*(\mathcal{G})\to\mathcal{F}$ es un
homomorfismo de $\mathcal{A}$-Álgebras,
$v^\flat:\mathcal{G}\to\Psi_*(\mathcal{F})$ es un homomorfismo de
$\mathcal{B}$-Álgebras.
\end{env}

\begin{env}[4.4.8]
\label{0.4.4.8-fr}
Sean $(Z,\mathcal{C})$ un tercer espacio anillado,
$\Psi'=(\psi',\theta')$ un morfismo $(Y,\mathcal{B})\to(Z,\mathcal{C})$, y
$\Psi'':(X,\mathcal{A})\to(Z,\mathcal{C})$ el morfismo compuesto
$\Psi'\circ\Psi$. Sea $\mathcal{H}$ un $\mathcal{C}$-Módulo, $v'$ un
homomorfismo de $\mathcal{H}$ en $\mathcal{G}$; el compuesto
$v''=v\circ v'$ es por definición el homomorfismo de $\mathcal{H}$ en
$\mathcal{F}$ definido por
\[
  \mathcal{H}\xrightarrow{v'}\Psi'_*(\mathcal{G})
  \xrightarrow{\Psi'_*(v)}\Psi'_*(\Psi_*(\mathcal{F}))~;
\]
se verifica que ${v''}^\sharp$ es el homomorfismo
\[
  \Psi^*({\Psi'}^*(\mathcal{H}))
  \xrightarrow{\Psi^*({v'}^\sharp)}
  \Psi^*(\mathcal{G})
  \xrightarrow{v^\sharp}\mathcal{F}.
\]
\end{env}


% END tree/worktree/es/CURRENT/EGA_I_ch0/ega0_i_4_4_3_8.tex
% BEGIN tree/worktree/es/CURRENT/EGA_I_ch0/ega0_i_5_1.tex
\section{Haces cuasicoherentes y haces coherentes}
\label{section:0.5-fr}

\subsection{Haces cuasicoherentes}
\label{subsection:0.5.1-fr}

\begin{env}[5.1.1]
\label{0.5.1.1-fr}
Sean $(X,\mathcal{O}_X)$ un espacio anillado, $\mathcal{F}$ un
$\mathcal{O}_X$-Módulo. El dato de un homomorfismo
$u:\mathcal{O}_X\to\mathcal{F}$ de $\mathcal{O}_X$-Módulos equivale al
de la sección $s=u(1)\in\Gamma(X,\mathcal{F})$. En efecto, cuando $s$ está
dada, para toda sección $t\in\Gamma(U,\mathcal{O}_X)$, se tiene necesariamente
$u(t)=t\cdot(s|U)$; se dice que $u$ está \emph{definido por la sección $s$}. Si
ahora $I$ es un conjunto cualquiera de índices, consideremos el haz
suma directa $\mathcal{O}_X^{(I)}$, y para todo $i\in I$, sea $h_i$
la inyección canónica del $i$-ésimo factor en $\mathcal{O}_X^{(I)}$; se sabe
que $u\to(u\circ h_i)$ es un isomorfismo de
$\operatorname{Hom}_{\mathcal{O}_X}(\mathcal{O}_X^{(I)},\mathcal{F})$ sobre el
producto
$\bigl(\operatorname{Hom}_{\mathcal{O}_X}(\mathcal{O}_X,\mathcal{F})\bigr)^I$.
Hay, pues, una correspondencia biunívoca canónica entre los homomorfismos
$u:\mathcal{O}_X^{(I)}\to\mathcal{F}$ y las \emph{familias de secciones
$(s_i)_{i\in I}$ de $\mathcal{F}$ sobre $X$}. El homomorfismo $u$
correspondiente a $(s_i)$ aplica un elemento
$(a_i)\in(\Gamma(U,\mathcal{O}_X))^{(I)}$ en
$\sum_{i\in I}a_i\cdot(s_i|U)$.

Se dice que $\mathcal{F}$ está \emph{generado por la familia $(s_i)$} si
el homomorfismo $\mathcal{O}_X^{(I)}\to\mathcal{F}$ definido
\oldpage[0\textsubscript{I}]{45}
por esta familia es \emph{sobreyectivo} (es decir, si, para todo
$x\in X$, $\mathcal{F}_x$ es un $\mathcal{O}_x$-módulo generado por los
$(s_i)_x$). Se dice que $\mathcal{F}$ está \emph{generado por sus secciones
sobre $X$} si está generado por la familia de todas estas secciones (o
por una subfamilia), es decir, si existe un homomorfismo sobreyectivo
$\mathcal{O}_X^{(I)}\to\mathcal{F}$ para un $I$ adecuado.

Se observará que un $\mathcal{O}_X$-Módulo $\mathcal{F}$ puede ser tal que
exista un punto $x_0\in X$ para el cual $\mathcal{F}|U$ no esté generado por
sus secciones sobre $U$, \emph{cualquiera que sea la vecindad abierta $U$ de
$x_0$}: basta tomar $X=\mathbb{R}$, como $\mathcal{O}_X$ el haz
simple $\mathbb{Z}$, como $\mathcal{F}$ el subhaz algebraico de
$\mathcal{O}_X$ tal que $\mathcal{F}_0=\{0\}$, $\mathcal{F}_x=\mathbb{Z}$ para
$x\ne0$, y, finalmente, $x_0=0$: la única sección de $\mathcal{F}|U$ sobre
$U$ es $0$ para una vecindad $U$ de $0$.
\end{env}

\begin{env}[5.1.2]
\label{0.5.1.2-fr}
Sea $f:X\to Y$ un morfismo de espacios anillados. Si $\mathcal{F}$ es un
$\mathcal{O}_X$-Módulo generado por sus secciones sobre $X$, entonces
el homomorfismo canónico
$f^*(f_*(\mathcal{F}))\to\mathcal{F}$ (4.4.3.3) es \emph{sobreyectivo}: en
efecto, con las notaciones de (5.1.1), $s_i\otimes1$ es una sección de
$f^*(f_*(\mathcal{F}))$ sobre $X$, y su imagen en $\mathcal{F}$ es
$s_i$. El ejemplo de (5.1.1), en el que $f$ es la identidad, muestra que el recíproco de
esta proposición es inexacto en general.
\end{env}

\begin{env}[5.1.3]
\label{0.5.1.3-fr}
Se dice que un $\mathcal{O}_X$-Módulo $\mathcal{F}$ es \emph{cuasicoherente}
si, para todo $x\in X$, hay una vecindad abierta $U$ de $x$ tal que
$\mathcal{F}|U$ sea isomorfo al \emph{conúcleo} de un homomorfismo de la
forma
$\mathcal{O}_X^{(I)}|U\to\mathcal{O}_X^{(J)}|U$, donde $I$ y $J$ son
conjuntos arbitrarios de índices. Es claro que el propio $\mathcal{O}_X$ es
un $\mathcal{O}_X$-Módulo cuasicoherente, y que toda suma directa finita de
$\mathcal{O}_X$-Módulos cuasicoherentes es un $\mathcal{O}_X$-Módulo
cuasicoherente. Se dice que una $\mathcal{O}_X$-\emph{Álgebra} $\mathcal{A}$ es
\emph{cuasicoherente} si lo es como $\mathcal{O}_X$-Módulo.
\end{env}

\begin{env}[5.1.4]
\label{0.5.1.4-fr}
Sea $f:X\to Y$ un morfismo de espacios anillados. Si $\mathcal{G}$ es un
$\mathcal{O}_Y$-Módulo cuasicoherente, entonces $f^*(\mathcal{G})$ es un
$\mathcal{O}_X$-Módulo cuasicoherente. En efecto, para todo $x\in X$, hay
una vecindad abierta $V$ de $f(x)$ en $Y$ tal que $\mathcal{G}|V$ es el
conúcleo de un homomorfismo
$\mathcal{O}_Y^{(I)}|V\to\mathcal{O}_Y^{(J)}|V$. Si $U=f^{-1}(V)$, y si
$f_U$ es la restricción de $f$ a $U$, se tiene
$f^*(\mathcal{G})|U=f_U^*(\mathcal{G}|V)$; como $f_U^*$ es exacto por la derecha
y conmuta con las sumas directas, $f_U^*(\mathcal{G}|V)$ es conúcleo de un
homomorfismo
$\mathcal{O}_X^{(I)}|U\to\mathcal{O}_X^{(J)}|U$.
\end{env}


% END tree/worktree/es/CURRENT/EGA_I_ch0/ega0_i_5_1.tex
% BEGIN tree/worktree/es/CURRENT/EGA_I_ch0/ega0_i_5_2.tex
\subsection{Haces de tipo finito}
\label{subsection:0.5.2-fr}

\begin{env}[5.2.1]
\label{0.5.2.1-fr}
Se dice que un $\mathcal{O}_X$-Módulo $\mathcal{F}$ es \emph{de tipo finito} si,
para todo $x\in X$, existe una vecindad abierta $U$ de $x$ tal que
$\mathcal{F}|U$ está generado por una familia \emph{finita} de secciones
sobre $U$, o, equivalentemente, es isomorfo a un haz cociente de un
haz de la forma $(\mathcal{O}_X|U)^p$, donde $p$ es finito. Todo haz
cociente de un haz de tipo finito es de tipo finito, así como toda suma
directa finita y todo producto tensorial finito de haces de tipo finito. Un
$\mathcal{O}_X$-Módulo de tipo finito no es necesariamente cuasicoherente,
como se puede ver para el $\mathcal{O}_X$-Módulo
$\mathcal{O}_X/\mathcal{F}$, donde $\mathcal{F}$ es el ejemplo de (5.1.1). Si
$\mathcal{F}$ es de tipo finito, $\mathcal{F}_x$ es un
$\mathcal{O}_x$-módulo de tipo finito para todo $x\in X$, pero el ejemplo
(5.1.1) muestra que esta condición necesaria no es en general
suficiente.
\end{env}

\begin{env}[5.2.2]
\label{0.5.2.2-fr}
Sea $\mathcal{F}$ un $\mathcal{O}_X$-Módulo \emph{de tipo finito}. Si $s_i$
$(1\leq i\leq n)$ son secciones de $\mathcal{F}$ sobre una vecindad
abierta $U$ de un punto $x\in X$ y si los $(s_i)_x$ generan
$\mathcal{F}_x$, existe una vecindad abierta $V\subset U$ de $x$ tal que
los $(s_i)_y$ generan $\mathcal{F}_y$ para todo $y\in V$ (FAC, I, 2, 12,
prop.~1). Se concluye, en particular, que el soporte de $\mathcal{F}$ es
\emph{cerrado}.

\oldpage[0\textsubscript{I}]{46}

Del mismo modo, si $u:\mathcal{F}\to\mathcal{G}$ es un homomorfismo tal que
$u_x=0$, entonces existe una vecindad $U$ de $x$ tal que $u_y=0$ para todo
$y\in U$.
\end{env}

\begin{env}[5.2.3]
\label{0.5.2.3-fr}
Supongamos que $X$ es \emph{cuasicompacto}, y sean $\mathcal{F}$, $\mathcal{G}$ dos
$\mathcal{O}_X$-Módulos tales que $\mathcal{G}$ es de \emph{tipo finito},
$u:\mathcal{F}\to\mathcal{G}$ un homomorfismo \emph{sobreyectivo}. Supongamos
además que $\mathcal{F}$ es límite inductivo de un sistema inductivo
$(\mathcal{F}_\lambda)$ de $\mathcal{O}_X$-Módulos. Entonces existe un índice
$\mu$ tal que el homomorfismo $\mathcal{F}_\mu\to\mathcal{G}$ es
\emph{sobreyectivo}. En efecto, para todo $x\in X$, existe un sistema finito de
secciones $s_i$ de $\mathcal{G}$ sobre una vecindad abierta $U(x)$ de $x$
tal que los $(s_i)_y$ generan $\mathcal{G}_y$ para todo $y\in U(x)$; hay,
pues, una vecindad abierta $V(x)\subset U(x)$ de $x$ y $n$ secciones $t_i$ de
$\mathcal{F}$ sobre $V(x)$ tales que
$s_i|V(x)=u(t_i)$ para todo $i$; se puede suponer además que los $t_i$ son
las imágenes canónicas de secciones de un mismo haz
$\mathcal{F}_{\lambda(x)}$ sobre $V(x)$. Cubramos entonces $X$ con un
número finito de vecindades $V(x_k)$, y sea $\mu$ un índice superior a todos
los $\lambda(x_k)$; es claro que este índice responde a la cuestión.

Suponiendo todavía que $X$ es cuasicompacto, sea $\mathcal{F}$ un
$\mathcal{O}_X$-Módulo de tipo finito generado por sus secciones sobre $X$
(5.1.1); entonces $\mathcal{F}$ está generado por una subfamilia \emph{finita}
de estas secciones: basta, en efecto, cubrir $X$ con un número finito de
vecindades abiertas $U_k$ tales que, para cada $k$, hay un número finito de
secciones $s_{ik}$ de $\mathcal{F}$ sobre $X$ cuyas restricciones a
$U_k$ generan $\mathcal{F}|U_k$; es claro que los $s_{ik}$ generan
entonces $\mathcal{F}$.
\end{env}

\begin{env}[5.2.4]
\label{0.5.2.4-fr}
Sea $f:X\to Y$ un morfismo de espacios anillados. Si $\mathcal{G}$ es un
$\mathcal{O}_Y$-Módulo de tipo finito, entonces $f^*(\mathcal{G})$ es un
$\mathcal{O}_X$-módulo de tipo finito. En efecto, para todo $x\in X$, hay una
vecindad abierta $V$ de $f(x)$ en $Y$ y un homomorfismo sobreyectivo
$v:\mathcal{O}_Y^p|V\to\mathcal{G}|V$. Si $U=f^{-1}(V)$ y si $f_U$ es la
restricción de $f$ a $U$, se tiene
$f^*(\mathcal{G})|U=f_U^*(\mathcal{G}|V)$; como $f_U^*$ es exacto por la derecha
(4.3.1) y conmuta con las sumas directas (4.3.2), $f_U^*(v)$ es un
homomorfismo sobreyectivo
$\mathcal{O}_X^p|U\to f^*(\mathcal{G})|U$.
\end{env}

\begin{env}[5.2.5]
\label{0.5.2.5-fr}
Se dice que un $\mathcal{O}_X$-Módulo $\mathcal{F}$ \emph{admite una presentación
finita} si, para todo $x\in X$, existe una vecindad abierta $U$ de $x$ tal
que $\mathcal{F}|U$ es isomorfo a un \emph{conúcleo de un
$(\mathcal{O}_X|U)$-homomorfismo}
$\mathcal{O}_X^p|U\to\mathcal{O}_X^q|U$, siendo $p$ y $q$ dos enteros $>0$.
Tal $\mathcal{O}_X$-Módulo es, por tanto, de tipo finito y cuasicoherente. Si
$f:X\to Y$ es un morfismo de espacios anillados, y si $\mathcal{G}$ es un
$\mathcal{O}_Y$-Módulo que admite una presentación finita,
$f^*(\mathcal{G})$ admite una presentación finita, como muestra el
razonamiento de (5.1.4).
\end{env}

\begin{env}[5.2.6]
\label{0.5.2.6-fr}
Sea $\mathcal{F}$ un $\mathcal{O}_X$-Módulo que admite una presentación finita
(5.2.5); entonces, para todo $\mathcal{O}_X$-Módulo $\mathcal{H}$,
el homomorfismo canónico funtorial
\[
  \bigl(\mathcal{H}\!om_{\mathcal{O}_X}(\mathcal{F},\mathcal{H})\bigr)_x
  \longrightarrow
  \operatorname{Hom}_{\mathcal{O}_x}(\mathcal{F}_x,\mathcal{H}_x)
\]
es \emph{biyectivo} (T, 4.1.1).
\end{env}

\begin{env}[5.2.7]
\label{0.5.2.7-fr}
Sean $\mathcal{F}$, $\mathcal{G}$ dos $\mathcal{O}_X$-Módulos que tienen una
presentación finita. Si, para un $x\in X$, $\mathcal{F}_x$ y $\mathcal{G}_x$
son dos $\mathcal{O}_x$-módulos \emph{isomorfos}, entonces existe una
vecindad abierta $U$ de $x$ tal que $\mathcal{F}|U$ y $\mathcal{G}|U$ son
\emph{isomorfos}. En efecto, si
$\varphi:\mathcal{F}_x\to\mathcal{G}_x$ y
$\psi:\mathcal{G}_x\to\mathcal{F}_x$ son dos isomorfismos inversos, existe,
según (5.2.6), una vecindad abierta $V$ de $x$ y una sección $u$
(resp. $v$) de
$\mathcal{H}\!om_{\mathcal{O}_X}(\mathcal{F},\mathcal{G})$ (resp.
$\mathcal{H}\!om_{\mathcal{O}_X}(\mathcal{G},\mathcal{F})$) sobre $V$
tal

\oldpage[0\textsubscript{I}]{47}

que $u_x=\varphi$ (resp. $v_x=\psi$). Como $(u\circ v)_x$ y
$(v\circ u)_x$ son los automorfismos identidad, existe una vecindad
abierta $U\subset V$ de $x$ tal que $(u\circ v)|U$ y $(v\circ u)|U$ son los
automorfismos identidad, de donde la proposición.
\end{env}


% END tree/worktree/es/CURRENT/EGA_I_ch0/ega0_i_5_2.tex
% BEGIN tree/worktree/es/CURRENT/EGA_I_ch0/ega0_i_5_3.tex
\subsection{Haces coherentes}
\label{subsection:0.5.3-fr}

\begin{env}[5.3.1]
\label{0.5.3.1-fr}
Se dice que un $\mathcal{O}_X$-Módulo $\mathcal{F}$ es \emph{coherente} si
verifica las dos condiciones siguientes:
\begin{enumerate}
  \item[$a)$] $\mathcal{F}$ es de tipo finito.
  \item[$b)$] Para todo abierto $U\subset X$, todo entero $n>0$, y todo
    homomorfismo
    $u:\mathcal{O}_X^n|U\to\mathcal{F}|U$, el núcleo de $u$ es de tipo finito.
\end{enumerate}
Se observará que estas dos condiciones son de carácter \emph{local}.

Para la mayoría de las demostraciones de las propiedades de los haces coherentes
que se recuerdan a continuación, cf. (FAC), I, 2.
\end{env}

\begin{env}[5.3.2]
\label{0.5.3.2-fr}
Todo $\mathcal{O}_X$-Módulo coherente admite una presentación finita (5.2.5); el
recíproco no es necesariamente exacto, pues el propio $\mathcal{O}_X$ no es
necesariamente un $\mathcal{O}_X$-Módulo coherente.

Todo sub-$\mathcal{O}_X$-Módulo \emph{de tipo finito} de un
$\mathcal{O}_X$-Módulo coherente es coherente; toda suma directa \emph{finita}
de $\mathcal{O}_X$-Módulos coherentes es un $\mathcal{O}_X$-Módulo coherente.
\end{env}

\begin{env}[5.3.3]
\label{0.5.3.3-fr}
Si $0\to\mathcal{F}\to\mathcal{G}\to\mathcal{H}\to0$ es una sucesión exacta de
$\mathcal{O}_X$-Módulos y si dos de estos $\mathcal{O}_X$-Módulos son
coherentes, lo mismo ocurre con el tercero.
\end{env}

\begin{env}[5.3.4]
\label{0.5.3.4-fr}
Si $\mathcal{F}$ y $\mathcal{G}$ son dos $\mathcal{O}_X$-Módulos coherentes,
$u:\mathcal{F}\to\mathcal{G}$ un homomorfismo, $\operatorname{Im}(u)$,
$\operatorname{Ker}(u)$ y $\operatorname{Coker}(u)$ son
$\mathcal{O}_X$-Módulos coherentes. En particular, si $\mathcal{F}$ y
$\mathcal{G}$ son sub-$\mathcal{O}_X$-Módulos coherentes de un
$\mathcal{O}_X$-Módulo coherente, $\mathcal{F}+\mathcal{G}$ y
$\mathcal{F}\cap\mathcal{G}$ son coherentes.
\end{env}

\begin{env}[5.3.5]
\label{0.5.3.5-fr}
Si $\mathcal{F}$ y $\mathcal{G}$ son dos $\mathcal{O}_X$-Módulos coherentes,
lo mismo ocurre con $\mathcal{F}\otimes_{\mathcal{O}_X}\mathcal{G}$ y con
$\mathcal{H}\!om_{\mathcal{O}_X}(\mathcal{F},\mathcal{G})$.
\end{env}

\begin{env}[5.3.6]
\label{0.5.3.6-fr}
Sean $\mathcal{F}$ un $\mathcal{O}_X$-Módulo coherente, $\mathcal{J}$ un
haz coherente de ideales de $\mathcal{O}_X$. Entonces el
$\mathcal{O}_X$-Módulo $\mathcal{J}\mathcal{F}$ es coherente, como imagen
de $\mathcal{J}\otimes_{\mathcal{O}_X}\mathcal{F}$ por el homomorfismo
canónico
$\mathcal{J}\otimes_{\mathcal{O}_X}\mathcal{F}\to\mathcal{F}$
(5.3.4 y 5.3.5).
\end{env}

\begin{env}[5.3.7]
\label{0.5.3.7-fr}
Se dice que una $\mathcal{O}_X$-Álgebra $\mathcal{A}$ es \emph{coherente} si
$\mathcal{A}$ es un $\mathcal{O}_X$-Módulo coherente. En particular,
$\mathcal{O}_X$ es un \emph{haz coherente de anillos} si y solo si,
para todo abierto $U\subset X$ y todo homomorfismo de la forma
$u:\mathcal{O}_X^p|U\to\mathcal{O}_X|U$, el núcleo de $u$ es un
$(\mathcal{O}_X|U)$-Módulo de tipo finito.

Si $\mathcal{O}_X$ es un haz coherente de anillos, todo
$\mathcal{O}_X$-Módulo $\mathcal{F}$ que admite una presentación finita (5.2.5)
es coherente, en virtud de (5.3.4).

El \emph{anulador} de un $\mathcal{O}_X$-Módulo $\mathcal{F}$ es el núcleo
$\mathcal{J}$ del homomorfismo canónico
$\mathcal{O}_X\to\mathcal{H}\!om_{\mathcal{O}_X}(\mathcal{F},\mathcal{F})$
que a toda sección $s\in\Gamma(U,\mathcal{O}_X)$ hace corresponder la
multiplicación por $s$ en
$\operatorname{Hom}(\mathcal{F}|U,\mathcal{F}|U)$; si $\mathcal{O}_X$ es
coherente y si $\mathcal{F}$ es un $\mathcal{O}_X$-Módulo coherente,
$\mathcal{J}$ es coherente (5.3.4 y 5.3.5) y para todo $x\in X$,
$\mathcal{J}_x$ es el anulador de $\mathcal{F}_x$ (5.2.6).
\end{env}

\oldpage[0\textsubscript{I}]{48}

\begin{env}[5.3.8]
\label{0.5.3.8-fr}
Supongamos que $\mathcal{O}_X$ es coherente; sean $\mathcal{F}$ un
$\mathcal{O}_X$-Módulo coherente, $x$ un punto de $X$, $M$ un submódulo de
tipo finito de $\mathcal{F}_x$; existe entonces una vecindad abierta $U$ de $x$
y un sub-$(\mathcal{O}_X|U)$-Módulo coherente $\mathcal{G}$ de
$\mathcal{F}|U$ tal que $\mathcal{G}_x=M$ (T, 4.1, lema 1).

Este resultado, unido a las propiedades de los sub-$\mathcal{O}_X$-Módulos de un
$\mathcal{O}_X$-Módulo coherente, impone condiciones necesarias a los anillos
$\mathcal{O}_x$ para que $\mathcal{O}_X$ sea coherente. Por ejemplo (5.3.4),
la intersección de dos ideales de tipo finito de $\mathcal{O}_x$ debe ser también
un ideal de tipo finito.
\end{env}

\begin{env}[5.3.9]
\label{0.5.3.9-fr}
Supongamos que $\mathcal{O}_X$ es coherente, y sea $M$ un $\mathcal{O}_x$-módulo
que admite una presentación finita, por tanto isomorfo a un conúcleo de un
homomorfismo $\varphi:\mathcal{O}_x^p\to\mathcal{O}_x^q$; entonces existe una
vecindad abierta $U$ de $x$ y un $(\mathcal{O}_X|U)$-Módulo coherente
$\mathcal{F}$ tal que $\mathcal{F}_x$ es isomorfo a $M$. En efecto, en virtud
de (5.2.6), existe una sección $u$ de
$\mathcal{H}\!om_{\mathcal{O}_X}(\mathcal{O}_X^p,\mathcal{O}_X^q)$ tal que $u_x=\varphi$; el conúcleo
$\mathcal{F}$ del homomorfismo
$u:\mathcal{O}_X^p|U\to\mathcal{O}_X^q|U$ responde a la cuestión (5.3.4).
\end{env}

\begin{env}[5.3.10]
\label{0.5.3.10-fr}
Supongamos que $\mathcal{O}_X$ es coherente, y sea $\mathcal{J}$ un haz coherente
de ideales de $\mathcal{O}_X$. Para que un
$(\mathcal{O}_X/\mathcal{J})$-Módulo $\mathcal{F}$ sea coherente, es necesario y
suficiente que sea coherente como $\mathcal{O}_X$-Módulo. En particular,
$\mathcal{O}_X/\mathcal{J}$ es un haz coherente de anillos.
\end{env}

\begin{env}[5.3.11]
\label{0.5.3.11-fr}
Sea $f:X\to Y$ un morfismo de espacios anillados, y supongamos que
$\mathcal{O}_X$ es coherente; entonces, para todo $\mathcal{O}_Y$-Módulo coherente
$\mathcal{G}$, $f^*(\mathcal{G})$ es un $\mathcal{O}_X$-Módulo coherente. En
efecto, con las notaciones de (5.2.4), se puede suponer que $\mathcal{G}|V$ es
conúcleo de un homomorfismo
$v:\mathcal{O}_Y^q|V\to\mathcal{O}_Y^p|V$; como $f_U^*$ es exacto por la derecha,
$f^*(\mathcal{G})|U=f_U^*(\mathcal{G}|V)$ es conúcleo del homomorfismo
$f_U^*(v):\mathcal{O}_X^q|U\to\mathcal{O}_X^p|U$, de donde nuestra afirmación.
\end{env}

\begin{env}[5.3.12]
\label{0.5.3.12-fr}
Sean $Y$ una parte cerrada de $X$, $j:Y\to X$ la inyección canónica,
$\mathcal{O}_Y$ un haz de anillos sobre $Y$, y pongamos
$\mathcal{O}_X=j_*(\mathcal{O}_Y)$. Para que un $\mathcal{O}_Y$-Módulo
$\mathcal{G}$ sea de tipo finito (resp. cuasicoherente, coherente), es necesario y
suficiente que $j_*(\mathcal{G})$ sea un $\mathcal{O}_X$-Módulo de tipo finito
(resp. cuasicoherente, coherente).
\end{env}


% END tree/worktree/es/CURRENT/EGA_I_ch0/ega0_i_5_3.tex
% BEGIN tree/worktree/es/CURRENT/EGA_I_ch0/ega0_i_5_4.tex
\subsection{Haces localmente libres}
\label{subsection:0.5.4-fr}

\begin{env}[5.4.1]
\label{0.5.4.1-fr}
Sea $X$ un espacio anillado. Se dice que un $\mathcal{O}_X$-Módulo $\mathcal{F}$
es \emph{localmente libre} si, para todo $x\in X$, existe una vecindad
abierta $U$ de $x$ tal que $\mathcal{F}|U$ es isomorfo a un
$(\mathcal{O}_X|U)$-Módulo de la forma $\mathcal{O}_X^{(I)}|U$, donde $I$ puede
depender de $U$. Si para todo $U$, $I$ es finito, se dice que $\mathcal{F}$ es
\emph{de rango finito}; si para todo $U$, $I$ tiene el mismo número finito de elementos
$n$, se dice que $\mathcal{F}$ es \emph{de rango $n$}. Un
$\mathcal{O}_X$-Módulo localmente libre de rango 1 se llama también
\emph{invertible} (cf. (5.4.3)). Si $\mathcal{F}$ es un
$\mathcal{O}_X$-Módulo localmente libre de rango finito, para todo $x\in X$,
$\mathcal{F}_x$ es un $\mathcal{O}_x$-módulo libre de rango finito $n(x)$, y
existe una vecindad $U$ de $x$ tal que $\mathcal{F}|U$ es de rango $n(x)$;
si $X$ es conexo, $n(x)$ es, por tanto, \emph{constante}.

Es claro que todo haz localmente libre es cuasicoherente, y si
$\mathcal{O}_X$ es un haz coherente de anillos, todo
$\mathcal{O}_X$-Módulo localmente libre de rango finito es coherente.

Si $\mathcal{L}$ es localmente libre,
$\mathcal{L}\otimes_{\mathcal{O}_X}\mathcal{F}$ es un funtor \emph{exacto}
en $\mathcal{F}$ en la categoría de los $\mathcal{O}_X$-Módulos.

Tendremos que considerar sobre todo $\mathcal{O}_X$-Módulos localmente libres
de rango finito,

\oldpage[0\textsubscript{I}]{49}

y cuando hablemos de haces localmente libres sin precisar, se
sobreentenderá que son de \emph{rango finito}.
\end{env}

\begin{env}[5.4.2]
\label{0.5.4.2-fr}
Si $\mathcal{L}$, $\mathcal{F}$ son dos $\mathcal{O}_X$-Módulos, se tiene un
homomorfismo canónico funtorial
\[
  \check{\mathcal{L}}\otimes_{\mathcal{O}_X}\mathcal{F}
  =\mathcal{H}\!om_{\mathcal{O}_X}(\mathcal{L},\mathcal{O}_X)
   \otimes_{\mathcal{O}_X}\mathcal{F}
  \longrightarrow
  \mathcal{H}\!om_{\mathcal{O}_X}(\mathcal{L},\mathcal{F})
  \tag{5.4.2.1}\label{0.5.4.2.1-fr}
\]
definido de la manera siguiente: para todo abierto $U$, a todo par $(u,t)$, donde
\[
  u\in\Gamma\bigl(U,
    \mathcal{H}\!om_{\mathcal{O}_X}(\mathcal{L},\mathcal{O}_X)\bigr)
  =\operatorname{Hom}(\mathcal{L}|U,\mathcal{O}_X|U)
  \quad\text{y}\quad
  t\in\Gamma(U,\mathcal{F}),
\]
se hace corresponder el elemento de
$\operatorname{Hom}(\mathcal{L}|U,\mathcal{F}|U)$ que, para todo $x\in U$,
hace corresponder a $s_x\in\mathcal{L}_x$ el elemento $u_x(s_x)t_x$ de
$\mathcal{F}_x$. Si $\mathcal{L}$ es \emph{localmente libre de rango finito},
este homomorfismo es \emph{biyectivo}; como la propiedad es local, se puede,
en efecto, limitar al caso en que $\mathcal{L}=\mathcal{O}_X^n$; como para todo
$\mathcal{O}_X$-Módulo $\mathcal{G}$,
$\mathcal{H}\!om_{\mathcal{O}_X}(\mathcal{O}_X^n,\mathcal{G})$ es
canónicamente isomorfo a $\mathcal{G}^n$, se reduce al caso
$\mathcal{L}=\mathcal{O}_X$, que es inmediato.
\end{env}

\begin{env}[5.4.3]
\EGAReaderTarget{canon-ref.0.5.4.3-fr}{0.5.4.3-fr}
Si $\mathcal{L}$ es invertible, lo mismo ocurre con su dual
$\check{\mathcal{L}}=
\mathcal{H}\!om_{\mathcal{O}_X}(\mathcal{L},\mathcal{O}_X)$, pues se reduce
inmediatamente (siendo la cuestión local) al caso $\mathcal{L}=\mathcal{O}_X$. Además,
se tiene un isomorfismo canónico
\[
  \mathcal{H}\!om_{\mathcal{O}_X}(\mathcal{L},\mathcal{O}_X)
  \otimes_{\mathcal{O}_X}\mathcal{L}
  \xrightarrow{\sim}
  \mathcal{O}_X
  \tag{5.4.3.1}\label{0.5.4.3.1-fr}
\]
en efecto, según (5.4.2), basta definir un isomorfismo canónico
$\mathcal{H}\!om_{\mathcal{O}_X}(\mathcal{L},\mathcal{L})
\xrightarrow{\sim}\mathcal{O}_X$. Ahora bien, para \emph{todo}
$\mathcal{O}_X$-Módulo $\mathcal{F}$, se tiene un homomorfismo canónico
$\mathcal{O}_X\longrightarrow
\mathcal{H}\!om_{\mathcal{O}_X}(\mathcal{F},\mathcal{F})$ (5.3.7). Queda
por probar que si $\mathcal{F}=\mathcal{L}$ es invertible, este homomorfismo
es biyectivo, y como la cuestión es local, se reduce al caso
$\mathcal{L}=\mathcal{O}_X$, que es inmediato.

Por lo que precede, se pone
$\mathcal{L}^{-1}=
\mathcal{H}\!om_{\mathcal{O}_X}(\mathcal{L},\mathcal{O}_X)$, y se dice que
$\mathcal{L}^{-1}$ es el \emph{inverso} de $\mathcal{L}$. La terminología de
\og haz invertible\fg{} puede justificarse de la manera siguiente cuando $X$ se
reduce a un punto y $\mathcal{O}_X$ es un anillo \emph{local} $A$ de ideal
maximal $\mathfrak m$; si $M$ y $M'$ son dos $A$-módulos (siendo $M$ de
tipo finito) tales que $M\otimes_A M'$ es isomorfo a $A$, como
$(A/\mathfrak m)\otimes_A(M\otimes_A M')$ se identifica con
$(M/\mathfrak mM)\otimes_{A/\mathfrak m}(M'/\mathfrak mM')$, este último
producto tensorial de espacios vectoriales sobre el cuerpo $A/\mathfrak m$ es
isomorfo a $A/\mathfrak m$, lo que exige que $M/\mathfrak mM$ y
$M'/\mathfrak mM'$ sean de dimensión 1. Para todo elemento $z\in M$
que no pertenece a $\mathfrak mM$, se tiene, pues, $M=Az+\mathfrak mM$, lo que
implica $M=Az$ según el lema de Nakayama, siendo $M$ de tipo finito.
Por lo demás, como el anulador de $z$ anula $M\otimes_A M'$, isomorfo a
$A$, este anulador es $\{0\}$, y $M$ es, por consiguiente, \emph{isomorfo a $A$}.
En el caso general, esto muestra que si $\mathcal{L}$ es un
$\mathcal{O}_X$-Módulo de tipo finito, tal que existe un
$\mathcal{O}_X$-Módulo $\mathcal{F}$ para el cual
$\mathcal{L}\otimes_{\mathcal{O}_X}\mathcal{F}$ es isomorfo a
$\mathcal{O}_X$, y si además los anillos $\mathcal{O}_x$ son anillos
locales, entonces $\mathcal{L}_x$ es un $\mathcal{O}_x$-módulo isomorfo a
$\mathcal{O}_x$ para todo $x\in X$. Si se supone que $\mathcal{O}_X$ y $\mathcal{L}$ son
\emph{coherentes}, se concluye, pues, que $\mathcal{L}$ es invertible
en virtud de (5.2.7).
\end{env}

\begin{env}[5.4.4]
\label{0.5.4.4-fr}
Si $\mathcal{L}$ y $\mathcal{L}'$ son dos $\mathcal{O}_X$-Módulos
invertibles, lo mismo ocurre con
$\mathcal{L}\otimes_{\mathcal{O}_X}\mathcal{L}'$, pues, siendo la cuestión
local, se puede suponer que $\mathcal{L}=\mathcal{O}_X$, y el resultado es
entonces trivial. Para todo entero $n\geq1$, se designa por
$\mathcal{L}^{\otimes n}$ el producto tensorial de $n$ haces

\oldpage[0\textsubscript{I}]{50}

idénticos a $\mathcal{L}$; se pone por convenio
$\mathcal{L}^{\otimes0}=\mathcal{O}_X$ y para $n\geq1$,
$\mathcal{L}^{\otimes(-n)}=(\mathcal{L}^{-1})^{\otimes n}$. Con estas
notaciones, existe entonces un \emph{isomorfismo canónico funtorial}
\[
  \mathcal{L}^{\otimes m}\otimes_{\mathcal{O}_X}\mathcal{L}^{\otimes n}
  \xrightarrow{\sim}
  \mathcal{L}^{\otimes(n+m)}
  \tag{5.4.4.1}\label{0.5.4.4.1-fr}
\]
cualesquiera que sean los enteros racionales $m,n$: en efecto, en virtud de las
definiciones, se reduce inmediatamente al caso en que $m=-1$, $n=1$, y
el isomorfismo en cuestión ha sido entonces definido en (5.4.3).
\end{env}

\begin{env}[5.4.5]
\EGAReaderTarget{canon-ref.0.5.4.5-fr}{0.5.4.5-fr}
Sea $f:Y\to X$ un morfismo de espacios anillados. Si $\mathcal{L}$ es un
$\mathcal{O}_X$-Módulo localmente libre (resp. invertible),
$f^*(\mathcal{L})$ es un $\mathcal{O}_Y$-Módulo localmente libre (resp.
invertible): esto resulta inmediatamente de que las imágenes inversas de dos
$\mathcal{O}_X$-Módulos localmente isomorfos son localmente isomorfas, de
que $f^*$ conmuta con las sumas directas finitas y de que
$f^*(\mathcal{O}_X)=\mathcal{O}_Y$ (4.3.4). Además, se sabe que se tiene un
homomorfismo canónico funtorial
\[
  f^*(\check{\mathcal{L}})
  \longrightarrow
  \bigl(f^*(\mathcal{L})\bigr)\check{}
\]
(4.4.6), y cuando $\mathcal{L}$ es localmente libre, este homomorfismo es
\emph{biyectivo}: en efecto, se reduce de nuevo al caso en que
$\mathcal{L}=\mathcal{O}_X$, que es trivial. Se concluye que si
$\mathcal{L}$ es invertible, $f^*(\mathcal{L}^{\otimes n})$ se identifica
canónicamente con $(f^*(\mathcal{L}))^{\otimes n}$ para todo entero racional
$n$.
\end{env}

\begin{env}[5.4.6]
\label{0.5.4.6-fr}
Sea $\mathcal{L}$ un $\mathcal{O}_X$-Módulo invertible; se designa por
$\Gamma_*(X,\mathcal{L})$ o simplemente $\Gamma_*(\mathcal{L})$ el grupo
abeliano suma directa
$\bigoplus_{n\in\mathbb{Z}}\Gamma(X,\mathcal{L}^{\otimes n})$; se lo dota
de una estructura de \emph{anillo graduado}, haciendo corresponder al par
$(s_n,s_m)$, donde $s_n\in\Gamma(X,\mathcal{L}^{\otimes n})$,
$s_m\in\Gamma(X,\mathcal{L}^{\otimes m})$, la sección de
$\mathcal{L}^{\otimes(n+m)}$ sobre $X$ que corresponde canónicamente
(5.4.4.1) a la sección $s_n\otimes s_m$ de
$\mathcal{L}^{\otimes n}\otimes_{\mathcal{O}_X}\mathcal{L}^{\otimes m}$;
la asociatividad de esta multiplicación se verifica de manera inmediata. Es
claro que $\Gamma_*(X,\mathcal{L})$ es un funtor covariante en
$\mathcal{L}$, con valores en la categoría de los anillos graduados.

Si ahora $\mathcal{F}$ es un $\mathcal{O}_X$-Módulo cualquiera, se pone
\[
  \Gamma_*(\mathcal{L},\mathcal{F})
  =\bigoplus_{n\in\mathbb{Z}}
   \Gamma\bigl(X,
     \mathcal{F}\otimes_{\mathcal{O}_X}\mathcal{L}^{\otimes n}\bigr).
\]
Se dota a este grupo abeliano de una estructura de \emph{módulo graduado} sobre
el anillo graduado $\Gamma_*(\mathcal{L})$ de la manera siguiente: al par
$(s_n,u_m)$, donde $s_n\in\Gamma(X,\mathcal{L}^{\otimes n})$ y
$u_m\in\Gamma(X,
\mathcal{F}\otimes\mathcal{L}^{\otimes m})$, se hace
corresponder la sección de
$\mathcal{F}\otimes\mathcal{L}^{\otimes(m+n)}$ que corresponde
canónicamente (5.4.4.1) a $s_n\otimes u_m$; la verificación de los axiomas de
los módulos es inmediata. Para $X$ y $\mathcal{L}$ fijos,
$\Gamma_*(\mathcal{L},\mathcal{F})$ es un funtor covariante en
$\mathcal{F}$ con valores en la categoría de los
$\Gamma_*(\mathcal{L})$-módulos graduados; para $X$ y $\mathcal{F}$ fijos,
es un funtor covariante en $\mathcal{L}$ con valores en la categoría de los
grupos abelianos.

Si $f:Y\to X$ es un morfismo de espacios anillados, el homomorfismo canónico
(4.4.3.2)
\[
  \rho:\mathcal{L}^{\otimes n}
  \longrightarrow
  f_*\bigl(f^*(\mathcal{L}^{\otimes n})\bigr)
\]
define un homomorfismo de grupos abelianos
\[
  \Gamma(X,\mathcal{L}^{\otimes n})
  \longrightarrow
  \Gamma\bigl(Y,f^*(\mathcal{L}^{\otimes n})\bigr),
\]
y como
$f^*(\mathcal{L}^{\otimes n})=(f^*(\mathcal{L}))^{\otimes n}$, resulta de las
definiciones de los homomorfismos canónicos (4.4.3.2) y (5.4.4.1) que los
homomorfismos anteriores definen un \emph{homomorfismo funtorial
de anillos graduados}
$\Gamma_*(\mathcal{L})\to\Gamma_*(f^*(\mathcal{L}))$. El mismo homomorfismo
canónico (4.4.3) define asimismo un homomorfismo de grupos abelianos
\[
  \Gamma\bigl(X,
    \mathcal{F}\otimes\mathcal{L}^{\otimes n}\bigr)
  \longrightarrow
  \Gamma\bigl(Y,
    f^*(\mathcal{F}\otimes\mathcal{L}^{\otimes n})\bigr),
\]
y como
\[
  f^*(\mathcal{F}\otimes\mathcal{L}^{\otimes n})
  =f^*(\mathcal{F})\otimes
   (f^*(\mathcal{L}))^{\otimes n}
  \quad\text{(4.3.3.1),}
\]

\oldpage[0\textsubscript{I}]{51}

estos homomorfismos (para $n$ variable) definen un
\emph{dihomomorfismo de módulos graduados}
\[
  \Gamma_*(\mathcal{L},\mathcal{F})
  \longrightarrow
  \Gamma_*(f^*(\mathcal{L}),f^*(\mathcal{F})).
\]
\end{env}

\begin{env}[5.4.7]
\EGAReaderTarget{canon-ref.0.5.4.7-fr}{0.5.4.7-fr}
Se puede mostrar que existe un \emph{conjunto} $\mathfrak{M}$ (denotado también
por $\mathfrak{M}(X)$) de $\mathcal{O}_X$-Módulos invertibles tal que todo
$\mathcal{O}_X$-Módulo invertible es isomorfo a un elemento de
$\mathfrak{M}$ y a uno solo\footnote{Véase el libro en preparación citado en
la Introducción.}; se define en $\mathfrak{M}$ una ley de composición
haciendo corresponder a dos elementos $\mathcal{L}$, $\mathcal{L}'$ de
$\mathfrak{M}$ el único elemento de $\mathfrak{M}$ isomorfo a
$\mathcal{L}\otimes\mathcal{L}'$. Con esta ley de
composición, $\mathfrak{M}$ es un \emph{grupo isomorfo al grupo de
cohomología} $H^1(X,\mathcal{O}_X^*)$, donde $\mathcal{O}_X^*$ es el
subhaz de $\mathcal{O}_X$ tal que $\Gamma(U,\mathcal{O}_X^*)$ es el
grupo de los elementos invertibles del anillo $\Gamma(U,\mathcal{O}_X)$ para
todo abierto $U\subset X$ ($\mathcal{O}_X^*$ es, por tanto, un haz de grupos
abelianos \emph{multiplicativos}).

Para ello se observará que, para todo abierto $U\subset X$, el grupo de las secciones
$\Gamma(U,\mathcal{O}_X^*)$ se identifica canónicamente con el \emph{grupo de los
automorfismos} del $(\mathcal{O}_X|U)$-Módulo $\mathcal{O}_X|U$,
haciendo corresponder la identificación a una sección $\varepsilon$ de
$\mathcal{O}_X^*$ sobre $U$ el automorfismo $u$ de $\mathcal{O}_X|U$
tal que $u_x(s_x)=\varepsilon_xs_x$ para todo $x\in U$ y todo
$s_x\in\mathcal{O}_x$. Sea entonces $\mathfrak{U}=(U_\lambda)$ un recubrimiento
abierto de $X$; el dato, para todo par de índices $(\lambda,\mu)$, de un
automorfismo $\theta_{\lambda\mu}$ de
$\mathcal{O}_X|(U_\lambda\cap U_\mu)$ equivale a darse una
$1$-\emph{cocadena} del recubrimiento $\mathfrak{U}$, con valores en
$\mathcal{O}_X^*$, y decir que los $\theta_{\lambda\mu}$ verifican la
condición de pegado (3.3.1) significa que la cocadena correspondiente es
un \emph{cociclo}. Del mismo modo, el dato, para todo $\lambda$, de un automorfismo
$\omega_\lambda$ de $\mathcal{O}_X|U_\lambda$ equivale al dato de una
$0$-cocadena del recubrimiento $\mathfrak{U}$, con valores en
$\mathcal{O}_X^*$, y su \emph{coborde} corresponde a la familia de los
automorfismos
$(\omega_\lambda|U_\lambda\cap U_\mu)\circ
 (\omega_\mu|U_\lambda\cap U_\mu)^{-1}$.
Se puede hacer corresponder a todo $1$-cociclo de $\mathfrak{U}$ con valores
en $\mathcal{O}_X^*$ el elemento de $\mathfrak{M}$ isomorfo al
$\mathcal{O}_X$-Módulo invertible obtenido por pegado a partir de la
familia de automorfismos $(\theta_{\lambda\mu})$ correspondiente a este cociclo,
y a dos cociclos cohomólogos corresponden dos elementos iguales de
$\mathfrak{M}$ (3.3.2); en otras palabras, se ha definido así una aplicación
$\varphi_{\mathfrak{U}}:H^1(\mathfrak{U},\mathcal{O}_X^*)\to\mathfrak{M}$.
Además, si $\mathfrak{B}$ es un segundo recubrimiento abierto de $X$, más fino
que $\mathfrak{U}$, el diagrama
\[
\begin{tikzcd}[column sep=huge,row sep=large]
H^1(\mathfrak{U},\mathcal{O}_X^*)
  \arrow[dd]
  \arrow[rrd,"\varphi_{\mathfrak{U}}"] & & \\
& & \mathfrak{M} \\
H^1(\mathfrak{B},\mathcal{O}_X^*)
  \arrow[rru,swap,"\varphi_{\mathfrak{B}}"] & &
\end{tikzcd}
\]
donde la flecha vertical es el homomorfismo canónico (G, II, 5.7), es
conmutativo, como resulta de (3.3.3). Por paso al límite inductivo, se
obtiene, pues, una aplicación
$H^1(X,\mathcal{O}_X^*)\to\mathfrak{M}$, identificándose el grupo de cohomología de \v{C}ech
$\check{H}^1(X,\mathcal{O}_X^*)$, como se sabe, con el primer grupo
de cohomología $H^1(X,\mathcal{O}_X^*)$ (G, II, 5.9, cor. del teorema~5.9.1).
Esta aplicación es \emph{sobreyectiva}: en efecto, por definición, para todo
$\mathcal{O}_X$-Módulo invertible $\mathcal{L}$, hay un recubrimiento abierto
$(U_\lambda)$ de $X$ tal que $\mathcal{L}$ se obtiene por pegado de los
haces $\mathcal{O}_X|U_\lambda$ (3.3.1). Es también
\emph{inyectiva}, pues basta probarlo para las aplicaciones
$H^1(\mathfrak{U},\mathcal{O}_X^*)\to\mathfrak{M}$, y esto resulta entonces de
(3.3.2). Queda por mostrar que

\oldpage[0\textsubscript{I}]{52}

la biyección así definida es un homomorfismo de grupos. Ahora bien, dados
dos $\mathcal{O}_X$-Módulos invertibles $\mathcal{L}$, $\mathcal{L}'$, hay
un recubrimiento abierto $(U_\lambda)$ tal que $\mathcal{L}|U_\lambda$ y
$\mathcal{L}'|U_\lambda$ son isomorfos a $\mathcal{O}_X|U_\lambda$ para
todo $\lambda$; hay, pues, para cada índice $\lambda$ un elemento
$a_\lambda$ (resp. $a'_\lambda$) de $\Gamma(U_\lambda,\mathcal{L})$ (resp.
$\Gamma(U_\lambda,\mathcal{L}'))$ tal que los elementos de
$\Gamma(U_\lambda,\mathcal{L})$ (resp. $\Gamma(U_\lambda,\mathcal{L}'))$
son los $s_\lambda\cdot a_\lambda$ (resp.
$s_\lambda\cdot a'_\lambda$), donde $s_\lambda$ recorre
$\Gamma(U_\lambda,\mathcal{O}_X)$. Los cociclos correspondientes
$(\varepsilon_{\lambda\mu})$, $(\varepsilon'_{\lambda\mu})$ son tales que
$s_\lambda\cdot a_\lambda=s_\mu\cdot a_\mu$ (resp.
$s_\lambda\cdot a'_\lambda=s_\mu\cdot a'_\mu$) sobre
$U_\lambda\cap U_\mu$ es equivalente a
$s_\lambda=\varepsilon_{\lambda\mu}s_\mu$ (resp.
$s_\lambda=\varepsilon'_{\lambda\mu}s_\mu$) sobre
$U_\lambda\cap U_\mu$. Como las secciones de
$\mathcal{L}\otimes_{\mathcal{O}_X}\mathcal{L}'$ sobre $U_\lambda$
son las sumas finitas de los
$s_\lambda s'_\lambda\cdot(a_\lambda\otimes a'_\lambda)$ donde $s_\lambda$ y
$s'_\lambda$ recorren $\Gamma(U_\lambda,\mathcal{O}_X)$, es claro que
el cociclo $(\varepsilon_{\lambda\mu}\varepsilon'_{\lambda\mu})$ corresponde
a $\mathcal{L}\otimes_{\mathcal{O}_X}\mathcal{L}'$, lo que termina la
demostración\footnote{Para una forma general de este resultado, véase el libro
citado en la nota de la p.~51.}.
\end{env}

\begin{env}[5.4.8]
\label{0.5.4.8-fr}
Sea $f=(\psi,\omega)$ un morfismo $Y\to X$ de espacios anillados. El funtor
$f^*$, restringido a la categoría de los $\mathcal{O}_X$-Módulos invertibles,
define una aplicación (denotada también por $f^*$ por abuso de lenguaje) del conjunto
$\mathfrak{M}(X)$ en el conjunto $\mathfrak{M}(Y)$. Por otra parte, se tiene un
homomorfismo canónico (T, 3.2.2)
\[
  H^1(X,\mathcal{O}_X^*)\longrightarrow H^1(Y,\mathcal{O}_Y^*).
  \tag{5.4.8.1}\label{0.5.4.8.1-fr}
\]
Cuando se identifican canónicamente (5.4.7) $\mathfrak{M}(X)$ y
$H^1(X,\mathcal{O}_X^*)$ (resp. $\mathfrak{M}(Y)$ y
$H^1(Y,\mathcal{O}_Y^*)$), el homomorfismo (5.4.8.1) \emph{se identifica con
la aplicación $f^*$}. En efecto, si $\mathcal{L}$ procede de un cociclo
$(\varepsilon_{\lambda\mu})$ correspondiente a un recubrimiento abierto
$(U_\lambda)$ de $X$, basta mostrar que $f^*(\mathcal{L})$ procede
de un cociclo cuya clase de cohomología es imagen por (5.4.8.1) de la de
$(\varepsilon_{\lambda\mu})$. Ahora bien, si $\theta_{\lambda\mu}$ es
el automorfismo de $\mathcal{O}_X|(U_\lambda\cap U_\mu)$ que corresponde a
$\varepsilon_{\lambda\mu}$, es claro que $f^*(\mathcal{L})$ se obtiene por
pegado de los $\mathcal{O}_Y|\psi^{-1}(U_\lambda)$ mediante los
automorfismos $f^*(\theta_{\lambda\mu})$, y basta, pues, verificar que
estos últimos corresponden al cociclo
$(\omega^\sharp(\varepsilon_{\lambda\mu}))$, lo que resulta inmediatamente de las
definiciones (se ha identificado $\varepsilon_{\lambda\mu}$ con su imagen canónica
por $\rho$ (3.7.2), sección de $\psi^*(\mathcal{O}_X^*)$ sobre
$\psi^{-1}(U_\lambda\cap U_\mu)$).
\end{env}

\begin{env}[5.4.9]
\label{0.5.4.9-fr}
Sean $\mathcal{E}$, $\mathcal{F}$ dos $\mathcal{O}_X$-Módulos,
suponiéndose $\mathcal{F}$ \emph{localmente libre}, y sea $\mathcal{G}$ un
$\mathcal{O}_X$-Módulo \emph{extensión de $\mathcal{F}$ por $\mathcal{E}$},
es decir, tal que existe una sucesión exacta
\[
  0\longrightarrow\mathcal{E}\xrightarrow{j}\mathcal{G}
  \xrightarrow{p}\mathcal{F}\longrightarrow 0.
\]
Entonces, para todo $x\in X$, existe una vecindad abierta $U$ de $x$ tal que
$\mathcal{G}|U$ es isomorfo a la \emph{suma directa}
$\mathcal{E}|U\oplus\mathcal{F}|U$. Se puede, en efecto, limitar al caso en que
$\mathcal{F}=\mathcal{O}_X^n$; sean $e_i$ ($1\leq i\leq n$) las secciones
canónicas (5.5.5) de $\mathcal{O}_X^n$; existe entonces una vecindad abierta
$U$ de $x$ y $n$ secciones $s_i$ de $\mathcal{G}$ sobre $U$ tales que
$p(s_i|U)=e_i|U$ para $1\leq i\leq n$. Dicho esto, sea $f$
el homomorfismo $\mathcal{F}|U\to\mathcal{G}|U$ definido por las secciones
$s_i|U$ (5.1.1). Es inmediato que para todo abierto $V\subset U$, y toda
sección $s\in\Gamma(V,\mathcal{G})$ se tiene
$s-f(p(s))\in\Gamma(V,\mathcal{E})$, de donde nuestra afirmación.
\end{env}

\begin{env}[5.4.10]
\label{0.5.4.10-fr}
Sean $f:X\to Y$ un morfismo de espacios anillados, $\mathcal{F}$ un
$\mathcal{O}_X$-Módulo, $\mathcal{L}$ un $\mathcal{O}_Y$-Módulo localmente
libre de rango finito. Existe entonces un isomorfismo canónico
\[
  f_*(\mathcal{F})\otimes_{\mathcal{O}_Y}\mathcal{L}
  \xrightarrow{\sim}
  f_*\bigl(\mathcal{F}\otimes_{\mathcal{O}_X}f^*(\mathcal{L})\bigr)
  \tag{5.4.10.1}\label{0.5.4.10.1-fr}
\]

\oldpage[0\textsubscript{I}]{53}

En efecto, para todo $\mathcal{O}_Y$-Módulo $\mathcal{L}$, se tiene un
homomorfismo canónico
\[
  f_*(\mathcal{F})\otimes_{\mathcal{O}_Y}\mathcal{L}
  \xrightarrow{1\otimes\rho}
  f_*(\mathcal{F})\otimes_{\mathcal{O}_Y}f_*(f^*(\mathcal{L}))
  \xrightarrow{\alpha}
  f_*\bigl(\mathcal{F}\otimes_{\mathcal{O}_X}f^*(\mathcal{L})\bigr)
\]
siendo $\rho$ el homomorfismo (4.4.3.2) y $\alpha$ el homomorfismo (4.2.2.1).
Para mostrar que cuando $\mathcal{L}$ es localmente libre este homomorfismo
es biyectivo, basta, siendo la cuestión local sobre $Y$, considerar el
caso en que $\mathcal{L}=\mathcal{O}_Y^n$; además, como $f_*$ y $f^*$ conmutan con las
sumas directas finitas, se puede suponer $n=1$, y en este caso la proposición
resulta inmediatamente de las definiciones y de la relación
$f^*(\mathcal{O}_Y)=\mathcal{O}_X$.
\end{env}


% END tree/worktree/es/CURRENT/EGA_I_ch0/ega0_i_5_4.tex
% BEGIN tree/worktree/es/CURRENT/EGA_I_ch0/ega0_i_5_5.tex
\subsection{Haces sobre un espacio localmente anillado}
\label{subsection:0.5.5-fr}

\begin{env}[5.5.1]
\label{0.5.5.1-fr}
Diremos que un espacio anillado $(X,\mathcal{O}_X)$ es un \emph{espacio
anillado en anillos locales} si, para todo $x\in X$, $\mathcal{O}_x$ es un
anillo local; estos espacios anillados serán con mucho los más frecuentes de los
que se considerarán en este trabajo. Designaremos entonces por
$\mathfrak{m}_x$ el \emph{ideal maximal} de $\mathcal{O}_x$, por $\boldsymbol{k}(x)$ el
\emph{cuerpo residual} $\mathcal{O}_x/\mathfrak{m}_x$; para todo
$\mathcal{O}_X$-Módulo $\mathcal{F}$, todo abierto $U$ de $X$, todo punto
$x\in U$ y toda sección $f\in\Gamma(U,\mathcal{F})$, designaremos por
$f(x)$ la \emph{clase} del germen
$f_x\in\mathcal{F}_x$ mod.~$\mathfrak{m}_x\mathcal{F}_x$, y diremos que
es el \emph{valor} de $f$ en el punto $x$. La relación $f(x)=0$ significa,
pues, que $f_x\in\mathfrak{m}_x\mathcal{F}_x$; cuando se cumple, se dirá
(por abuso de lenguaje) que $f$ \emph{se anula en $x$}. Se tendrá cuidado de no
confundir esta relación con $f_x=0$.
\end{env}

\begin{env}[5.5.2]
\label{0.5.5.2-fr}
Sean $X$ un espacio localmente anillado, $\mathcal{L}$ un
$\mathcal{O}_X$-Módulo invertible, $f$ una sección de $\mathcal{L}$ sobre
$X$. Hay entonces \emph{equivalencia} entre las tres propiedades siguientes
en un punto $x\in X$:
\begin{enumerate}
  \item[a)] \emph{$f_x$ es un generador de $\mathcal{L}_x$};
  \item[b)] \emph{$f_x\notin\mathfrak{m}_x\mathcal{L}_x$} (es decir,
    $f(x)\neq0$);
  \item[c)] \emph{Existe una sección $g$ de $\mathcal{L}^{-1}$ sobre
    una vecindad abierta $V$ de $x$ tal que la imagen canónica de $f\otimes
    g$ en $\Gamma(V,\mathcal{O}_X)$ \textup{(5.4.3)} es la sección unidad}.
\end{enumerate}

En efecto, siendo la cuestión local, se puede limitar al caso en que
$\mathcal{L}=\mathcal{O}_X$; la equivalencia de $a)$ y $b)$ es entonces
evidente, y es claro que $c)$ implica $b)$. Por otra parte, si
$f_x\notin\mathfrak{m}_x$, $f_x$ es invertible en $\mathcal{O}_x$, sea
$f_xg_x=1_x$. Por definición de los gérmenes de secciones, esto quiere decir que
existe una vecindad $V$ de $x$ y una sección $g$ de $\mathcal{O}_X$ sobre
$V$ tal que $fg=1$ en $V$, de donde $c)$.

Resulta inmediatamente de la condición $c)$ que el conjunto $X_f$ de los $x$
que verifican las condiciones equivalentes $a)$, $b)$, $c)$ es \emph{abierto} en
$X$; según la terminología introducida en (5.5.1), es el conjunto de los
$x$ donde $f$ \emph{no se anula}.
\end{env}

\begin{env}[5.5.3]
\label{0.5.5.3-fr}
Bajo las hipótesis de (5.5.2), sea $\mathcal{L}'$ un segundo
$\mathcal{O}_X$-Módulo invertible; entonces, si
$f\in\Gamma(X,\mathcal{L})$, $g\in\Gamma(X,\mathcal{L}')$, se tiene
\[
  X_f\cap X_g=X_{f\otimes g}
\]

En efecto, se reduce inmediatamente al caso en que
$\mathcal{L}=\mathcal{L}'=\mathcal{O}_X$ (siendo la cuestión local); como
$f\otimes g$ se identifica entonces canónicamente con el producto $fg$, la proposición
es evidente.
\end{env}

\oldpage[0\textsubscript{I}]{54}

\begin{env}[5.5.4]
\label{0.5.5.4-fr}
Sea $\mathcal{F}$ un $\mathcal{O}_X$-Módulo localmente libre de rango $n$;
es inmediato que $\bigwedge^p\mathcal{F}$ es un $\mathcal{O}_X$-Módulo
localmente libre de rango $\binom{n}{p}$ si $p\leq n$, reducido a $0$ si $p>n$,
pues la cuestión es local y se reduce al caso en que
$\mathcal{F}=\mathcal{O}_X^n$; además, para todo $x\in X$,
\[
  (\bigwedge^p\mathcal{F})_x/
  \mathfrak{m}_x(\bigwedge^p\mathcal{F})_x
\]
es un espacio vectorial de dimensión $\binom{n}{p}$ sobre $\boldsymbol{k}(x)$, que
se identifica canónicamente con
\[
  \bigwedge^p(\mathcal{F}_x/\mathfrak{m}_x\mathcal{F}_x).
\]
Sean
$s_1,\ldots,s_p$ unas secciones de $\mathcal{F}$ sobre un abierto $U$ de $X$, y
sea $s=s_1\wedge\ldots\wedge s_p$, que es una sección de
$\bigwedge^p\mathcal{F}$ sobre $U$ (4.1.5); se tiene
$s(x)=s_1(x)\wedge\ldots\wedge s_p(x)$, y, por consiguiente, decir que
$s_1(x),\ldots,s_p(x)$ son \emph{linealmente dependientes} significa que
$s(x)=0$. Se concluye que el \emph{conjunto de los $x\in U$ tales que
$s_1(x),\ldots,s_p(x)$ son linealmente independientes es abierto en
$X$}: basta, en efecto, reduciéndose al caso en que
$\mathcal{F}=\mathcal{O}_X^n$, aplicar (5.5.2) a la sección imagen de $s$
por una de las proyecciones de
$\bigwedge^p\mathcal{F}=\mathcal{O}_X^{\binom{n}{p}}$ sobre los
$\binom{n}{p}$ factores.

En particular, si $s_1,\ldots,s_n$ son $n$ secciones de $\mathcal{F}$
sobre $U$ tales que $s_1(x),\ldots,s_n(x)$ son linealmente
independientes en todo punto $x\in U$, el homomorfismo
$u:\mathcal{O}_X^n|U\to\mathcal{F}|U$ definido por los $s_i$ (5.1.1) es un
\emph{isomorfismo}: en efecto, se puede restringir al caso en que
$\mathcal{F}=\mathcal{O}_X^n$ y en que se identifican canónicamente
$\bigwedge^n\mathcal{F}$ y $\mathcal{O}_X$;
$s=s_1\wedge\ldots\wedge s_n$ es entonces una sección de $\mathcal{O}_X$
\emph{invertible} sobre $U$, y se define un homomorfismo inverso
de $u$ mediante las fórmulas de Cramer.
\end{env}

\begin{env}[5.5.5]
\label{0.5.5.5-fr}
Sean $\mathcal{E}$, $\mathcal{F}$ dos $\mathcal{O}_X$-Módulos localmente
libres (de rango finito), y sea $u:\mathcal{E}\to\mathcal{F}$ un
homomorfismo. Para que exista una vecindad $U$ de $x\in X$ tal que $u|U$
sea \emph{inyectivo} y $\mathcal{F}|U$ sea \emph{suma directa de
$u(\mathcal{E})|U$ y de un sub-$(\mathcal{O}_X|U)$-Módulo localmente libre
$\mathcal{G}$}, es necesario y suficiente que
$u_x:\mathcal{E}_x\to\mathcal{F}_x$ dé, por paso a los cocientes, un
homomorfismo \emph{inyectivo} de espacios vectoriales
\[
  \mathcal{E}_x/\mathfrak{m}_x\mathcal{E}_x
  \longrightarrow
  \mathcal{F}_x/\mathfrak{m}_x\mathcal{F}_x.
\]
La condición es, en efecto, \emph{necesaria}, pues $\mathcal{F}_x$ es entonces
suma directa de los $\mathcal{O}_x$-módulos libres $u_x(\mathcal{E}_x)$ y
$\mathcal{G}_x$, por tanto $\mathcal{F}_x/\mathfrak{m}_x\mathcal{F}_x$ es suma
directa de
$u_x(\mathcal{E}_x)/\mathfrak{m}_x u_x(\mathcal{E}_x)$ y de
$\mathcal{G}_x/\mathfrak{m}_x\mathcal{G}_x$. La condición es
\emph{suficiente}, pues se puede limitar al caso en que
$\mathcal{E}=\mathcal{O}_X^m$; sean $s_1,\ldots,s_m$ las imágenes por $u$ de las
secciones $e_i$ de $\mathcal{O}_X^m$ tales que $(e_i)_y$ es igual al $i$-ésimo
elemento de la base canónica de $\mathcal{O}_y^m$ para todo $y\in X$
(\emph{secciones canónicas} de $\mathcal{O}_X^m$); por hipótesis,
$s_1(x),\ldots,s_m(x)$ son linealmente independientes, por lo que, si $\mathcal{F}$
es de rango $n$, existen $n-m$ secciones $s_{m+1},\ldots,s_n$ de
$\mathcal{F}$ sobre una vecindad $V$ de $x$ tales que los $s_i(x)$
$(1\leq i\leq n)$ forman una base de
$\mathcal{F}_x/\mathfrak{m}_x\mathcal{F}_x$. Existe entonces (5.5.4) una
vecindad $U\subset V$ de $x$ tal que los $s_i(y)$ $(1\leq i\leq n)$ forman
una base de $\mathcal{F}_y/\mathfrak{m}_y\mathcal{F}_y$ para todo $y\in U$,
y se concluye (5.5.4) que hay un isomorfismo de $\mathcal{F}|U$ sobre
$\mathcal{O}_X^n|U$, que aplica los $s_i|U$ $(1\leq i\leq m)$ en los
$e_i|U$, lo que termina la demostración.
\end{env}


% END tree/worktree/es/CURRENT/EGA_I_ch0/ega0_i_5_5.tex
% BEGIN tree/worktree/es/CURRENT/EGA_I_ch0/ega0_i_6_0.tex
\section{Planitud}
\label{section:0.6-fr}

\begin{env}[6.0]
\label{0.6.0.0-fr}
La noción de planitud se debe a J.-P.~Serre [16]; en lo que sigue, se omiten
las demostraciones de los resultados que se exponen en el \emph{Álgebra
conmutativa} de N.~Bourbaki, a la que remitimos al lector. Suponemos que
todos los anillos son conmutativos\footnote{Véase la exposición citada de
N.~Bourbaki para la generalización de la mayoría de los resultados al caso no
conmutativo.}.

\oldpage[0\textsubscript{I}]{55}

Si $M$, $N$ son dos $A$-módulos, $M'$ (resp. $N'$) un submódulo de $M$
(resp. $N$), denotaremos por $\operatorname{Im}(M'\otimes_A N')$ el submódulo
de $M\otimes_A N$, imagen de la aplicación canónica
$M'\otimes_A N'\to M\otimes_A N$.
\end{env}


% END tree/worktree/es/CURRENT/EGA_I_ch0/ega0_i_6_0.tex
% BEGIN tree/worktree/es/CURRENT/EGA_I_ch0/ega0_i_6_1.tex
\subsection{Módulos planos}
\label{subsection:0.6.1-fr}

\begin{env}[6.1.1]
\label{0.6.1.1-fr}
Sea $M$ un $A$-módulo. Las propiedades siguientes son equivalentes:
\begin{enumerate}
  \item[$a)$] El funtor $M\otimes_A N$ en $N$ es exacto en la categoría
    de los $A$-módulos;
  \item[$b)$] $\operatorname{Tor}_i^A(M,N)=0$ para todo $i>0$ y todo
    $A$-módulo $N$;
  \item[$c)$] $\operatorname{Tor}_1^A(M,N)=0$ para todo $A$-módulo $N$.
\end{enumerate}

Cuando $M$ verifica estas condiciones, se dice que $M$ es un
$A$-\emph{módulo plano}. Es claro que todo $A$-módulo libre es plano.

Para que $M$ sea un $A$-módulo plano, basta que para todo ideal
$\mathfrak{J}$ de $A$, de \emph{tipo finito}, la aplicación canónica
$M\otimes_A\mathfrak{J}\to M\otimes_A A=M$ sea \emph{inyectiva}.
\end{env}

\begin{env}[6.1.2]
\EGAReaderTarget{canon-ref.0.6.1.2-fr}{0.6.1.2-fr}
Todo límite inductivo de $A$-módulos planos es un $A$-módulo plano. Para
que una suma directa $\bigoplus_{\lambda\in L}M_\lambda$ de $A$-módulos sea
un $A$-módulo plano, es necesario y suficiente que cada uno de los $A$-módulos
$M_\lambda$ sea plano. En particular, todo $A$-módulo proyectivo es plano.

Sea
\[
  0\longrightarrow M'\longrightarrow M\longrightarrow M''\longrightarrow 0
\]
una sucesión exacta de $A$-módulos, tal que $M''$ es \emph{plano}. Entonces, para
todo $A$-módulo $N$, la sucesión
\[
  0\longrightarrow M'\otimes N\longrightarrow M\otimes N
  \longrightarrow M''\otimes N\longrightarrow 0
\]
es exacta. Además, para que $M$ sea plano, es entonces necesario y suficiente que
$M'$ lo sea (pero puede ocurrir que $M$ y $M'$ sean planos sin que
$M''=M/M'$ lo sea).
\end{env}

\begin{env}[6.1.3]
\label{0.6.1.3-fr}
Sean $M$ un $A$-módulo plano, $N$ un $A$-módulo cualquiera; para dos
submódulos $N'$, $N''$ de $N$, se tiene entonces
\[
  \operatorname{Im}(M\otimes(N'+N''))
  =\operatorname{Im}(M\otimes N')+\operatorname{Im}(M\otimes N'')
\]
\[
  \operatorname{Im}(M\otimes(N'\cap N''))
  =\operatorname{Im}(M\otimes N')\cap\operatorname{Im}(M\otimes N'')
\]
(imágenes tomadas en $M\otimes N$).
\end{env}

\begin{env}[6.1.4]
\label{0.6.1.4-fr}
Sean $M$, $N$ dos $A$-módulos, $M'$ (resp. $N'$) un submódulo de $M$
(resp. $N$), y supongamos que uno de los módulos $M/M'$, $N/N'$ es plano.
Entonces se tiene
\[
  \operatorname{Im}(M'\otimes N')
  =\operatorname{Im}(M'\otimes N)\cap\operatorname{Im}(M\otimes N')
\]
(imágenes en $M\otimes N$). En particular, si $\mathfrak{J}$ es un ideal
de $A$ y si $M/M'$ es plano, se tiene
$\mathfrak{J}M'=M'\cap\mathfrak{J}M$.
\end{env}


% END tree/worktree/es/CURRENT/EGA_I_ch0/ega0_i_6_1.tex
% BEGIN tree/worktree/es/CURRENT/EGA_I_ch0/ega0_i_6_2.tex
\subsection{Cambio de anillos}
\label{subsection:0.6.2-fr}

Cuando un grupo aditivo $M$ está provisto de varias estructuras de módulo con
respecto a anillos $A$, $B$, \ldots, en lugar de decir que $M$ es plano
como $A$-módulo, $B$-módulo, \ldots, diremos a veces también que $M$
es $A$-\emph{plano}, $B$-\emph{plano}, \ldots

\begin{env}[6.2.1]
\label{0.6.2.1-fr}
Sean $A$ y $B$ dos anillos, $M$ un $A$-módulo, $N$ un
$(A,B)$-bimódulo. Si $M$ es plano y si $N$ es $B$-plano, entonces
$M\otimes_A N$ es $B$-plano. En particular, si $M$ y $N$ son dos
$A$-módulos planos, $M\otimes_A N$ es un $A$-módulo plano. Si $B$ es una
$A$-álgebra y si $M$ es
\oldpage[0\textsubscript{I}]{56}
un $A$-módulo plano, el $B$-módulo $M_{(B)}=M\otimes_A B$ es plano. Finalmente, si
$B$ es una $A$-álgebra que es plana como $A$-módulo, y si $N$ es
un $B$-módulo plano, entonces $N$ es también $A$-plano.
\end{env}

\begin{env}[6.2.2]
\label{0.6.2.2-fr}
Sean $A$ un anillo, $B$ una $A$-álgebra plana como $A$-módulo.
Sean $M$, $N$ dos $A$-módulos, tales que $M$ admite una presentación
finita; entonces el homomorfismo canónico
\[
  \operatorname{Hom}_A(M,N)\otimes_A B
  \longrightarrow
  \operatorname{Hom}_B(M\otimes_A B,N\otimes_A B)
  \tag{6.2.2.1}
\]
(que transforma $u\otimes b$ en el homomorfismo
$m\otimes b'\to u(m)\otimes b'b$) es un isomorfismo.
\end{env}

\begin{env}[6.2.3]
\label{0.6.2.3-fr}
Sea $(A_\lambda,\varphi_{\mu\lambda})$ un sistema inductivo filtrante
de anillos; sea $A=\varinjlim A_\lambda$. Sea, por otra parte, para cada
$\lambda$, $M_\lambda$ un $A_\lambda$-módulo y para $\lambda\leq\mu$ sea
$\theta_{\mu\lambda}:M_\lambda\to M_\mu$ un
$\varphi_{\mu\lambda}$-homomorfismo, tal que
$(M_\lambda,\theta_{\mu\lambda})$ es un sistema inductivo;
$M=\varinjlim M_\lambda$ es entonces un $A$-módulo. Dicho esto, si para todo
$\lambda$, $M_\lambda$ es un $A_\lambda$-módulo \emph{plano}, entonces $M$ es
un $A$-módulo \emph{plano}. En efecto, sea $\mathfrak{J}$ un ideal \emph{de
tipo finito} de $A$; por definición del límite inductivo, existe un
índice $\lambda$ y un ideal $\mathfrak{J}_\lambda$ de $A_\lambda$ tales que
$\mathfrak{J}=\mathfrak{J}_\lambda A$. Si se pone
$\mathfrak{J}'_\mu=\mathfrak{J}_\lambda A_\mu$ para $\mu\geq\lambda$, se tiene
también $\mathfrak{J}=\varinjlim\mathfrak{J}'_\mu$ (donde $\mu$ recorre los
índices $\geq\lambda$), de donde (siendo el funtor $\varinjlim$ exacto y
conmutando con el producto tensorial)
\[
  M\otimes_A\mathfrak{J}
  =\varinjlim(M_\mu\otimes_{A_\mu}\mathfrak{J}'_\mu)
  =\varinjlim\mathfrak{J}'_\mu M_\mu
  =\mathfrak{J}M.
\]
\end{env}


% END tree/worktree/es/CURRENT/EGA_I_ch0/ega0_i_6_2.tex
% BEGIN tree/worktree/es/CURRENT/EGA_I_ch0/ega0_i_6_3.tex
\subsection{Localización de la planitud}
\label{subsection:0.6.3-fr}

\begin{env}[6.3.1]
\label{0.6.3.1-fr}
Si $A$ es un anillo, $S$ una parte multiplicativa de $A$, $S^{-1}A$ es
un $A$-\emph{módulo plano}. En efecto, para todo $A$-módulo $N$,
$N\otimes_A S^{-1}A$ se identifica con $S^{-1}N$ (1.2.5) y se sabe (1.3.2)
que $S^{-1}N$ es un funtor exacto en $N$.

Si ahora $M$ es un $A$-módulo plano,
$S^{-1}M=M\otimes_A S^{-1}A$ es un $S^{-1}A$-módulo plano (6.2.1), por tanto
es también $A$-plano en virtud de lo que precede y de (6.2.1). En particular,
si $P$ es un $S^{-1}A$-módulo, se lo puede considerar como un $A$-módulo
isomorfo a $S^{-1}P$; para que $P$ sea $A$-plano, es necesario y suficiente
que sea $S^{-1}A$-plano.
\end{env}

\begin{env}[6.3.2]
\label{0.6.3.2-fr}
Sean $A$ un anillo, $B$ una $A$-álgebra, $T$ una parte multiplicativa de
$B$. Si $P$ es un $B$-módulo que es $A$-\emph{plano}, $T^{-1}P$ es $A$-\emph{plano}. En
efecto, para todo $A$-módulo $N$, se tiene
\[
  (T^{-1}P)\otimes_A N
  =(T^{-1}B\otimes_B P)\otimes_A N
  =T^{-1}B\otimes_B(P\otimes_A N)
  =T^{-1}(P\otimes_A N)~;
\]
ahora bien, $T^{-1}(P\otimes_A N)$ es un funtor exacto en $N$, por ser compuesto de los
dos funtores exactos $P\otimes_A N$ (en $N$) y $T^{-1}Q$ (en $Q$). Si $S$
es una parte multiplicativa de $A$ cuya imagen en $B$ está
\emph{contenida en $T$}, $T^{-1}P$ es igual a $S^{-1}(T^{-1}P)$, por tanto es
también $S^{-1}A$-plano en virtud de (6.3.1).
\end{env}

\begin{env}[6.3.3]
\label{0.6.3.3-fr}
Sean $\varphi:A\to B$ un homomorfismo de anillos, $M$ un $B$-módulo. Las
propiedades siguientes son equivalentes:
\begin{enumerate}
  \item[$a)$] $M$ es un $A$-módulo plano.
  \item[$b)$] Para todo ideal maximal $\mathfrak{n}$ de $B$,
    $M_\mathfrak{n}$ es un $A$-módulo plano.
  \item[$c)$] Para todo ideal maximal $\mathfrak{n}$ de $B$, poniendo
    $\mathfrak{m}=\varphi^{-1}(\mathfrak{n})$, $M_\mathfrak{n}$ es un
    $A_\mathfrak{m}$-módulo plano.
\end{enumerate}

En efecto, como $M_\mathfrak{n}=(M_\mathfrak{n})_\mathfrak{m}$,
la equivalencia de $b)$ y $c)$ resulta de (6.3.1), y el hecho de que $a)$
implica $b)$ es un caso particular de (6.3.2). Queda por ver que $b)$
implica $a)$,
\oldpage[0\textsubscript{I}]{57}
es decir, que para todo homomorfismo inyectivo $u:N'\to N$ de
$A$-módulos, el homomorfismo
$v=1\otimes u:M\otimes_A N'\to M\otimes_A N$ es inyectivo. Ahora bien, $v$ es también
un homomorfismo de $B$-módulos, y se sabe que para que sea inyectivo,
basta que para todo ideal maximal $\mathfrak{n}$ de $B$,
\[
  v_\mathfrak{n}:(M\otimes_A N')_\mathfrak{n}
  \longrightarrow(M\otimes_A N)_\mathfrak{n}
\]
sea inyectivo. Pero como
\[
  (M\otimes_A N)_\mathfrak{n}
  =B_\mathfrak{n}\otimes_B(M\otimes_A N)
  =M_\mathfrak{n}\otimes_A N,
\]
$v_\mathfrak{n}$ no es otro que el homomorfismo
\[
  1\otimes u:M_\mathfrak{n}\otimes_A N'
  \longrightarrow M_\mathfrak{n}\otimes_A N,
\]
que es inyectivo puesto que $M_\mathfrak{n}$ es $A$-plano.

En particular (haciendo $B=A$), para que un $A$-módulo $M$ sea plano, es
necesario y suficiente que $M_\mathfrak{m}$ sea $A_\mathfrak{m}$-plano para todo
ideal maximal $\mathfrak{m}$ de $A$.
\end{env}

\begin{env}[6.3.4]
\label{0.6.3.4-fr}
Sea $M$ un $A$-módulo; si $M$ es plano, y si $f\in A$ no es divisor
de $0$ en $A$, $f$ no anula ningún elemento $\neq0$ de $M$, pues
el homomorfismo $m\to f.m$ se escribe $1\otimes u$, donde $u$ es la
multiplicación $a\to f.a$ en $A$ y $M$ está identificado con
$M\otimes_A A$; si $u$ es inyectiva, lo mismo ocurre, pues, con
$1\otimes u$ puesto que $M$ es plano. En particular, si $A$ es
\emph{íntegro}, $M$ es \emph{sin torsión}.

Recíprocamente, supongamos que $A$ es íntegro, $M$ sin torsión, y supongamos que para
todo ideal maximal $\mathfrak{m}$ de $A$, $A_\mathfrak{m}$ es un
\emph{anillo de valoración discreta}; entonces $M$ es $A$-\emph{plano}. Basta, en
efecto (6.3.3), probar que $M_\mathfrak{m}$ es $A_\mathfrak{m}$-plano, y
se puede, pues, suponer que $A$ es ya un anillo de valoración discreta. Pero
como $M$ es límite inductivo de sus submódulos de tipo finito, y estos
últimos son sin torsión, se puede además limitar al caso en que $M$ es de
tipo finito (6.1.2). La proposición resulta en este caso de que $M$ es un
$A$-módulo libre.

En particular, si $A$ es un anillo \emph{íntegro},
$\varphi:A\to B$ un homomorfismo de anillos que hace de $B$ un $A$-módulo
\emph{plano} y $\neq\{0\}$, $\varphi$ es necesariamente \emph{inyectivo}.
Recíprocamente, si $B$ es íntegro, $A$ un subanillo de $B$ y si para todo
ideal maximal $\mathfrak{m}$ de $A$, $A_\mathfrak{m}$ es un anillo de
valoración discreta, $B$ es $A$-\emph{plano}.
\end{env}


% END tree/worktree/es/CURRENT/EGA_I_ch0/ega0_i_6_3.tex
% BEGIN tree/worktree/es/CURRENT/EGA_I_ch0/ega0_i_6_4.tex
\subsection{Módulos fielmente planos}
\label{subsection:0.6.4-fr}

\begin{env}[6.4.1]
\label{0.6.4.1-fr}
Para un $A$-módulo $M$, las cuatro propiedades siguientes son equivalentes:
\begin{enumerate}
  \item[$a)$] Para que una sucesión $N'\to N\to N''$ de $A$-módulos sea exacta,
    es necesario y suficiente que la sucesión
    $M\otimes N'\to M\otimes N\to M\otimes N''$ sea exacta;
  \item[$b)$] $M$ es plano y para todo $A$-módulo $N$, la relación
    $M\otimes N=0$ implica $N=0$;
  \item[$c)$] $M$ es plano y para todo homomorfismo $v:N\to N'$ de
    $A$-módulos, la relación $1_M\otimes v=0$ implica $v=0$, siendo $1_M$
    el automorfismo identidad de $M$;
  \item[$d)$] $M$ es plano y para todo ideal maximal $\mathfrak{m}$ de
    $A$, $\mathfrak{m}M\neq M$.
\end{enumerate}

Cuando $M$ verifica estas condiciones, se dice que $M$ es un $A$-módulo
\emph{fielmente plano}; $M$ es entonces necesariamente un módulo
\emph{fiel}. Además, si $u:N\to N'$ es un homomorfismo de $A$-módulos,
para que $u$ sea inyectivo (resp. sobreyectivo, biyectivo), es necesario y suficiente
que $1\otimes u:M\otimes N\to M\otimes N'$ lo sea.
\end{env}

\oldpage[0\textsubscript{I}]{58}

\begin{env}[6.4.2]
\label{0.6.4.2-fr}
Un módulo libre $\neq\{0\}$ es fielmente plano; lo mismo ocurre con la
suma directa de un módulo plano y de un módulo fielmente plano. Si $S$ es
una parte multiplicativa de $A$, $S^{-1}A$ es un $A$-módulo fielmente
plano solo si $S$ está formado por elementos invertibles (por tanto $S^{-1}A=A$).
\end{env}

\begin{env}[6.4.3]
\label{0.6.4.3-fr}
Sea
\[
  0\longrightarrow M'\longrightarrow M\longrightarrow M''\longrightarrow 0
\]
una sucesión exacta de $A$-módulos; si $M'$ y $M''$ son planos, y si uno
de ellos es fielmente plano, entonces $M$ es fielmente plano.
\end{env}

\begin{env}[6.4.4]
\label{0.6.4.4-fr}
Sean $A$ y $B$ dos anillos, $M$ un $A$-módulo, $N$ un
$(A,B)$-bimódulo. Si $M$ es fielmente plano y si $N$ es un $B$-módulo
fielmente plano, entonces $M\otimes_A N$ es un $B$-módulo fielmente plano. En
particular, si $M$ y $N$ son dos $A$-módulos fielmente planos, lo mismo
ocurre con $M\otimes_A N$. Si $B$ es una $A$-álgebra y si $M$ es un
$A$-módulo fielmente plano, el $B$-módulo $M_{(B)}$ es fielmente plano.
\end{env}

\begin{env}[6.4.5]
\label{0.6.4.5-fr}
Si $M$ es un $A$-módulo fielmente plano y si $S$ es una parte
multiplicativa de $A$, $S^{-1}M$ es un $S^{-1}A$-módulo fielmente plano,
puesto que $S^{-1}M=M\otimes_A(S^{-1}A)$ (6.4.4). Recíprocamente, si para todo
ideal maximal $\mathfrak{m}$ de $A$, $M_\mathfrak{m}$ es un
$A_\mathfrak{m}$-módulo fielmente plano, $M$ es un $A$-módulo fielmente
plano, pues $M$ es $A$-plano (6.3.3), y se tiene
\[
  M_\mathfrak{m}/\mathfrak{m}M_\mathfrak{m}
  =(M\otimes_A A_\mathfrak{m})
  \otimes_{A_\mathfrak{m}}
  (A_\mathfrak{m}/\mathfrak{m}A_\mathfrak{m})
  =M\otimes_A(A/\mathfrak{m})
  =M/\mathfrak{m}M
\]
por tanto, la hipótesis implica que $M/\mathfrak{m}M\neq0$ para todo ideal
maximal $\mathfrak{m}$ de $A$, lo que prueba nuestra afirmación (6.4.1).
\end{env}


% END tree/worktree/es/CURRENT/EGA_I_ch0/ega0_i_6_4.tex
% BEGIN tree/worktree/es/CURRENT/EGA_I_ch0/ega0_i_6_5.tex
\subsection{Restricción de escalares}
\label{subsection:0.6.5-fr}

\begin{env}[6.5.1]
\label{0.6.5.1-fr}
Sean $A$ un anillo, $\varphi:A\to B$ un homomorfismo de anillos que hace
de $B$ una $A$-álgebra. Supongamos que existe un $B$-módulo $N$ que es un
$A$-módulo \emph{fielmente plano}. Entonces, para todo $A$-módulo $M$,
el homomorfismo $x\to 1\otimes x$ de $M$ en
$B\otimes_A M=M_{(B)}$ es \emph{inyectivo}. En particular, $\varphi$ es
inyectivo; para todo ideal $\mathfrak{a}$ de $A$, se tiene
$\varphi^{-1}(\mathfrak{a}B)=\mathfrak{a}$; para todo ideal maximal (resp.
primo) $\mathfrak{m}$ de $A$, existe un ideal maximal (resp. primo)
$\mathfrak{n}$ de $B$ tal que $\varphi^{-1}(\mathfrak{n})=\mathfrak{m}$.
\end{env}

\begin{env}[6.5.2]
\label{0.6.5.2-fr}
Cuando se cumple la condición de (6.5.1), se identifica $A$ con un
subanillo de $B$ por $\varphi$ y, más generalmente, para todo $A$-módulo
$M$, se identifica $M$ con un sub-$A$-módulo de $M_{(B)}$. Se observará que si
$B$ es entonces \emph{noetheriano}, lo mismo ocurre con $A$, pues la aplicación
$\mathfrak{a}\to\mathfrak{a}B$ es una inyección creciente del
conjunto de los ideales de $A$ en el conjunto de los ideales de $B$;
la existencia de una sucesión infinita estrictamente creciente de ideales de $A$
implicaría, por tanto, la existencia de una sucesión análoga de ideales de $B$.
\end{env}


% END tree/worktree/es/CURRENT/EGA_I_ch0/ega0_i_6_5.tex
% BEGIN tree/worktree/es/CURRENT/EGA_I_ch0/ega0_i_6_6.tex
\subsection{Anillos fielmente planos}
\label{subsection:0.6.6-fr}

\begin{env}[6.6.1]
\label{0.6.6.1-fr}
Sea $\varphi:A\to B$ un homomorfismo de anillos que hace de $B$ una
$A$-álgebra. Las cinco propiedades siguientes son equivalentes:
\begin{enumerate}
  \item[$a)$] $B$ es un $A$-módulo fielmente plano (es decir,
    $M_{(B)}$ es un funtor \emph{exacto} y \emph{fiel} en $M$).
  \item[$b)$] El homomorfismo $\varphi$ es inyectivo y el $A$-módulo
    $B/\varphi(A)$ es plano.

\oldpage[0\textsubscript{I}]{59}

  \item[$c)$] El $A$-módulo $B$ es plano (es decir, el funtor
    $M_{(B)}$ es \emph{exacto}), y para todo $A$-módulo $M$, el homomorfismo
    $x\to1\otimes x$ de $M$ en $M_{(B)}$ es inyectivo.
  \item[$d)$] El $A$-módulo $B$ es plano y para todo ideal $\mathfrak a$ de
    $A$, se tiene $\varphi^{-1}(\mathfrak aB)=\mathfrak a$.
  \item[$e)$] El $A$-módulo $B$ es plano y para todo ideal maximal
    $\mathfrak m$ de $A$, existe un ideal maximal $\mathfrak n$ de $B$ tal
    que $\varphi^{-1}(\mathfrak n)=\mathfrak m$.
\end{enumerate}

Cuando se cumplen estas condiciones, se identifica $A$ con un subanillo de $B$.
\end{env}

\begin{env}[6.6.2]
\label{0.6.6.2-fr}
Sean $A$ un anillo \emph{local}, $\mathfrak m$ su ideal maximal, $B$ una
$A$-álgebra tal que $\mathfrak mB\neq B$ (lo que ocurre, por ejemplo, cuando
$B$ es un anillo local y $A\to B$ un homomorfismo \emph{local}). Si $B$ es un
$A$-módulo \emph{plano}, $B$ es un $A$-módulo \emph{fielmente plano}. En efecto, como
$\mathfrak mB\neq B$, hay un ideal maximal $\mathfrak n$ de $B$ que contiene
$\mathfrak mB$; como $\varphi^{-1}(\mathfrak n)\cap A$ contiene
$\mathfrak m$ y no contiene $1$, se tiene
$\varphi^{-1}(\mathfrak n)=\mathfrak m$ y se aplica el criterio $e)$ de (6.6.1).
Bajo las condiciones indicadas, se ve, pues, que si $B$ es
noetheriano, lo mismo ocurre con $A$ (6.5.2).
\end{env}

\begin{env}[6.6.3]
\label{0.6.6.3-fr}
Sea $B$ una $A$-álgebra que es un $A$-módulo fielmente plano. Para todo
$A$-módulo $M$ y todo sub-$A$-módulo $M'$ de $M$, se tiene (identificando $M$
con un sub-$A$-módulo de $M_{(B)}$) $M'=M\cap M'_{(B)}$. Para que $M$ sea un
$A$-módulo plano (resp. fielmente plano), es necesario y suficiente que $M_{(B)}$
sea un $B$-módulo plano (resp. fielmente plano).
\end{env}

\begin{env}[6.6.4]
\label{0.6.6.4-fr}
Sean $B$ una $A$-álgebra, $N$ un $B$-módulo fielmente plano. Para que $B$
sea un $A$-módulo plano (resp. fielmente plano), es necesario y suficiente que $N$
lo sea.

En particular, sea $C$ una $B$-álgebra; si el anillo $C$ es fielmente
plano sobre $B$ y $B$ fielmente plano sobre $A$, entonces $C$ es fielmente plano
sobre $A$; si $C$ es fielmente plano sobre $B$ y sobre $A$, entonces $B$ es
fielmente plano sobre $A$.
\end{env}


% END tree/worktree/es/CURRENT/EGA_I_ch0/ega0_i_6_6.tex
% BEGIN tree/worktree/es/CURRENT/EGA_I_ch0/ega0_i_6_7.tex
\subsection{Morfismos planos de espacios anillados}
\label{subsection:0.6.7-fr}

\begin{env}[6.7.1]
\label{0.6.7.1-fr}
Sea $f:X\to Y$ un morfismo de espacios anillados, y sea $\mathcal F$ un
$\mathcal O_X$-Módulo. Se dice que $\mathcal F$ es \emph{$f$-plano} (o $Y$-\emph{plano}
cuando no hay que temer ninguna confusión acerca de $f$) \emph{en un punto} $x\in X$ si
$\mathcal F_x$ es un $\mathcal O_{f(x)}$-módulo plano; se dice que
$\mathcal F$ es \emph{$f$-plano sobre} $y\in Y$ si $\mathcal F$ es $f$-plano
en todos los puntos $x\in f^{-1}(y)$; se dice que $\mathcal F$ es \emph{$f$-plano}
si $\mathcal F$ es $f$-plano en todos los puntos de $X$. Se dice que el
morfismo $f$ es \emph{plano en} $x\in X$ (resp. \emph{plano sobre} $y\in Y$, resp.
\emph{plano}) si $\mathcal O_X$ es $f$-plano en $x$ (resp. $f$-plano sobre $y$,
resp. $f$-plano).
\end{env}

\begin{env}[6.7.2]
\label{0.6.7.2-fr}
Con las notaciones de (6.7.1), si $\mathcal F$ es $f$-plano en $x$, para toda
vecindad abierta $U$ de $y=f(x)$, el funtor
$(f^*(\mathcal G)\otimes_{\mathcal O_X}\mathcal F)_x$ en $\mathcal G$ es
\emph{exacto} en la categoría de los $(\mathcal O_Y|U)$-Módulos; en efecto, este
tallo se identifica canónicamente con
$\mathcal G_y\otimes_{\mathcal O_y}\mathcal F_x$, y nuestra afirmación resulta
de la definición. En particular, si $f$ es un morfismo \emph{plano}, el funtor
$f^*$ es \emph{exacto} en la categoría de los $\mathcal O_Y$-Módulos.
\end{env}

\begin{env}[6.7.3]
\label{0.6.7.3-fr}
Recíprocamente, supongamos que el haz de anillos $\mathcal O_Y$ es \emph{coherente}, y
supongamos que para \emph{toda} vecindad abierta $U$ de $y$, el funtor
$(f^*(\mathcal G)\otimes_{\mathcal O_X}\mathcal F)_x$ es exacto en
$\mathcal G$ en la categoría de los $(\mathcal O_Y|U)$-Módulos \emph{coherentes}.
Entonces $\mathcal F$ es \emph{$f$-plano en $x$}. En efecto, basta probar que
para todo ideal de tipo finito $\mathfrak J$ de $\mathcal O_y$,
el homomorfismo canónico
$\mathfrak J\otimes_{\mathcal O_y}\mathcal F_x\to\mathcal F_x$ es inyectivo
(6.1.1). Ahora bien, se sabe (5.3.8) que existe entonces una vecindad

\oldpage[0\textsubscript{I}]{60}

abierta $U$ de $y$ y un haz coherente de ideales $\mathcal J$ de
$\mathcal O_Y|U$ tal que $\mathcal J_y=\mathfrak J$, de donde la conclusión.
\end{env}

\begin{env}[6.7.4]
\label{0.6.7.4-fr}
Los resultados de (6.1) sobre los módulos planos se transcriben inmediatamente en
proposiciones sobre los haces $f$-planos en un punto:

Si $0\to\mathcal F'\to\mathcal F\to\mathcal F''\to0$ es una sucesión exacta
de $\mathcal O_X$-Módulos y si $\mathcal F''$ es $f$-plano en el punto
$x\in X$, entonces, para toda vecindad abierta $U$ de $y=f(x)$ y todo
$(\mathcal O_Y|U)$-Módulo $\mathcal G$, la sucesión
\[
  0\longrightarrow
  (f^*(\mathcal G)\otimes_{\mathcal O_X}\mathcal F')_x
  \longrightarrow
  (f^*(\mathcal G)\otimes_{\mathcal O_X}\mathcal F)_x
  \longrightarrow
  (f^*(\mathcal G)\otimes_{\mathcal O_X}\mathcal F'')_x
  \longrightarrow0
\]
es exacta. Para que $\mathcal F$ sea $f$-plano en $x$, es entonces necesario y suficiente
que $\mathcal F'$ lo sea. Se tienen enunciados análogos para las nociones
correspondientes de $\mathcal O_X$-Módulo $f$-plano sobre $y\in Y$, o
de $\mathcal O_X$-Módulo $f$-plano.
\end{env}

\begin{env}[6.7.5]
\label{0.6.7.5-fr}
Sean $f:X\to Y$, $g:Y\to Z$ dos morfismos de espacios anillados; sean
$x\in X$, $y=f(x)$, $\mathcal F$ un $\mathcal O_X$-Módulo. Si $\mathcal F$
es $f$-plano en el punto $x$ y si el morfismo $g$ es plano en el punto $y$, entonces
$\mathcal F$ es $(g\circ f)$-plano en $x$ (6.2.1). En particular, si $f$ y
$g$ son morfismos planos, $g\circ f$ es plano.
\end{env}

\begin{env}[6.7.6]
\label{0.6.7.6-fr}
Sean $X$, $Y$ dos espacios anillados, $f:X\to Y$ un morfismo \emph{plano}. Entonces
el homomorfismo canónico de bifuntores (4.4.6)
\[
  f^*(\mathcal{H}om_{\mathcal O_Y}(\mathcal F,\mathcal G))
  \longrightarrow
  \mathcal{H}om_{\mathcal O_X}(f^*(\mathcal F),f^*(\mathcal G))
  \tag{6.7.6.1}
  \label{0.6.7.6.1-fr}
\]
es un \emph{isomorfismo} cuando $\mathcal F$ admite una \emph{presentación finita}
(5.2.5).

En efecto, siendo la cuestión local, se puede suponer que existe una sucesión
exacta $\mathcal O_Y^m\to\mathcal O_Y^n\to\mathcal F\to0$. Ahora bien, los dos
miembros de (6.7.6.1) son funtores exactos por la izquierda en $\mathcal F$ en
virtud de la hipótesis sobre $f$; se reduce entonces a demostrar la proposición
cuando $\mathcal F=\mathcal O_Y$, caso en que es trivial.
\end{env}

\begin{env}[6.7.8]
\label{0.6.7.8-fr}
Se dice que un morfismo $f:X\to Y$ de espacios anillados es \emph{fielmente plano} si
$f$ es \emph{sobreyectivo} y si, para todo $x\in X$, $\mathcal O_x$ es un
$\mathcal O_{f(x)}$-módulo \emph{fielmente plano}. Cuando $X$ y $Y$ son
espacios localmente anillados (5.5.1), equivale a decir que el
morfismo $f$ es \emph{sobreyectivo} y \emph{plano} (6.6.2). Cuando $f$ es fielmente plano,
$f^*$ es un funtor \emph{exacto} y \emph{fiel} en la categoría de los
$\mathcal O_Y$-Módulos (6.6.1, $a)$) y para que un $\mathcal O_Y$-Módulo
$\mathcal G$ sea $Y$-plano, es necesario y suficiente que $f^*(\mathcal G)$ lo
sea (6.6.3).
\end{env}


% END tree/worktree/es/CURRENT/EGA_I_ch0/ega0_i_6_7.tex
% BEGIN tree/worktree/es/CURRENT/EGA_I_ch0/ega0_i_7_1.tex
\section{Anillos ádicos}
\label{section:0.7-fr}

\subsection{Anillos admisibles}
\label{subsection:0.7.1-fr}

\begin{env}[7.1.1]
\label{0.7.1.1-fr}
Recordemos que en un anillo topológico $A$ (no necesariamente separado), se
dice que un elemento $x$ es \emph{topológicamente nilpotente} si $0$ es \emph{un} límite de
la sucesión $(x^n)_{n\geq0}$. Se dice que un anillo topológico $A$ está
\emph{linealmente topologizado} si existe un sistema fundamental de vecindades de
$0$ en $A$ formado por \emph{ideales} (necesariamente \emph{abiertos}).
\end{env}

\begin{definition}[7.1.2]
\label{0.7.1.2-fr}
En un anillo linealmente topologizado $A$, se dice que un ideal $\mathfrak J$
es un \textup{ideal de definición} si $\mathfrak J$ es abierto y si, para toda
vecindad $V$ de $0$, existe un entero $n>0$ tal

\oldpage[0\textsubscript{I}]{61}

que $\mathfrak J^n\subset V$ \textup{(lo que se expresa, por abuso de lenguaje, diciendo
que la sucesión $(\mathfrak J^n)$ \emph{tiende a} $0$)}. Se dice que un anillo linealmente
topologizado $A$ es \emph{preadmisible} si existe en $A$ un ideal de definición;
se dice que $A$ es \emph{admisible} si es preadmisible y si además es
separado y completo.
\end{definition}

Es claro que si $\mathfrak J$ es un ideal de definición, $\mathfrak L$ un
ideal abierto de $A$, $\mathfrak J\cap\mathfrak L$ es también un ideal de
definición; los ideales de definición de un anillo preadmisible $A$ forman,
pues, un \emph{sistema fundamental de vecindades de} $0$.

\begin{lemma}[7.1.3]
\label{0.7.1.3-fr}
Sea $A$ un anillo linealmente topologizado.
\begin{enumerate}
  \item[{\upshape(i)}] Para que $x\in A$ sea topológicamente nilpotente, es
    necesario y suficiente que para todo ideal abierto $\mathfrak J$ de $A$, la imagen
    canónica de $x$ en $A/\mathfrak J$ sea nilpotente. El conjunto
    $\mathfrak T$ de los elementos topológicamente nilpotentes de $A$ es un ideal.
  \item[{\upshape(ii)}] Supongamos además que $A$ es preadmisible, y sea
    $\mathfrak J$ un ideal de definición de $A$. Para que $x\in A$ sea
    topológicamente nilpotente, es necesario y suficiente que su imagen canónica en
    $A/\mathfrak J$ sea nilpotente; el ideal $\mathfrak T$ es la imagen
    inversa del nilradical de $A/\mathfrak J$ y es, por tanto, abierto.
\end{enumerate}
\end{lemma}

(i) se deduce inmediatamente de las definiciones. Para probar (ii), basta
observar que para toda vecindad $V$ de $0$ en $A$, existe $n>0$ tal que
$\mathfrak J^n\subset V$; si $x\in A$ es tal que $x^m\in\mathfrak J$, se tiene
$x^{mq}\in V$ para $q\geq n$, por tanto $x$ es topológicamente nilpotente.

\begin{proposition}[7.1.4]
\label{0.7.1.4-fr}
Sean $A$ un anillo preadmisible, $\mathfrak J$ un ideal de definición de
$A$.
\begin{enumerate}
  \item[{\upshape(i)}] Para que un ideal $\mathfrak J'$ de $A$ esté contenido
    en un ideal de definición, es necesario y suficiente que exista un entero
    $n>0$ tal que ${\mathfrak J'}^n\subset\mathfrak J$.
  \item[{\upshape(ii)}] Para que un $x\in A$ esté contenido en un ideal de
    definición, es necesario y suficiente que sea topológicamente nilpotente.
\end{enumerate}
\end{proposition}

(i) Si ${\mathfrak J'}^n\subset\mathfrak J$, para toda vecindad abierta $V$ de
$0$ en $A$, existe $m$ tal que $\mathfrak J^m\subset V$, por tanto
${\mathfrak J'}^{mn}\subset V$.

(ii) La condición es evidentemente necesaria; es suficiente, pues si se
cumple, existe $n$ tal que $x^n\in\mathfrak J$, por tanto
$\mathfrak J'=\mathfrak J+Ax$ es un ideal de definición, puesto que es abierto,
y que ${\mathfrak J'}^n\subset\mathfrak J$.

\begin{corollary}[7.1.5]
\label{0.7.1.5-fr}
En un anillo preadmisible $A$, un ideal primo abierto contiene todos los
ideales de definición.
\end{corollary}

\begin{corollary}[7.1.6]
\label{0.7.1.6-fr}
Siendo las notaciones y las hipótesis las de \textup{(7.1.4)}, las propiedades siguientes
de un ideal $\mathfrak J_0$ de $A$ son equivalentes:
\begin{enumerate}
  \item[$a)$] $\mathfrak J_0$ es el mayor ideal de definición de $A$;
  \item[$b)$] $\mathfrak J_0$ es un ideal de definición maximal;
  \item[$c)$] $\mathfrak J_0$ es un ideal de definición tal que el anillo
    $A/\mathfrak J_0$ es reducido.
\end{enumerate}
Para que exista un ideal $\mathfrak J_0$ con estas propiedades, es necesario y
suficiente que el nilradical de $A/\mathfrak J$ sea nilpotente; $\mathfrak J_0$
es entonces igual al ideal $\mathfrak T$ de los elementos topológicamente nilpotentes
de $A$.
\end{corollary}

Es claro que $a)$ implica $b)$, y $b)$ implica $c)$ en virtud de
(7.1.4, (ii)) y (7.1.3, (ii)); por la misma razón, $c)$ implica $a)$. La
última afirmación resulta de (7.1.4, (i)) y (7.1.3, (ii)).

Cuando $\mathfrak T/\mathfrak J$, nilradical de $A/\mathfrak J$, es
nilpotente, se denota por $A_{\mathrm{red}}$ el anillo cociente (reducido)
$A/\mathfrak T$.

\oldpage[0\textsubscript{I}]{62}

\begin{corollary}[7.1.7]
\label{0.7.1.7-fr}
Un anillo preadmisible noetheriano admite un mayor ideal de definición.
\end{corollary}

\begin{corollary}[7.1.8]
\label{0.7.1.8-fr}
Si un anillo preadmisible $A$ es tal que, para un ideal de definición
$\mathfrak J$, las potencias $\mathfrak J^n$ ($n>0$) forman un sistema
fundamental de vecindades de $0$, lo mismo ocurre con las potencias
${\mathfrak J'}^n$ de todo ideal de definición $\mathfrak J'$ de $A$.
\end{corollary}

\begin{definition}[7.1.9]
\label{0.7.1.9-fr}
Se dice que un anillo preadmisible $A$ es \emph{preádico} si existe un ideal de
definición $\mathfrak J$ de $A$ tal que los $\mathfrak J^n$ forman un sistema
fundamental de vecindades de $0$ en $A$ \textup{(o, lo que equivale a lo mismo, tal que
los $\mathfrak J^n$ son \emph{abiertos})}. Se llama anillo \emph{ádico} a un anillo
preádico separado y completo.
\end{definition}

Si $\mathfrak J$ es un ideal de definición de un anillo preádico (resp.
ádico) $A$, se dice también que $A$ es un anillo
\emph{$\mathfrak J$-preádico} (resp. \emph{$\mathfrak J$-ádico}), y que su
topología es la topología \emph{$\mathfrak J$-preádica} (resp.
\emph{$\mathfrak J$-ádica}). Más generalmente, si $M$ es un $A$-módulo, la
topología sobre $M$ que tiene como sistema fundamental de vecindades de $0$ los
submódulos $\mathfrak J^nM$ se llama topología
\emph{$\mathfrak J$-preádica} (resp. \emph{$\mathfrak J$-ádica}). En virtud
de (7.1.8), estas topologías son independientes de la elección del ideal de definición
$\mathfrak J$.

\begin{proposition}[7.1.10]
\label{0.7.1.10-fr}
Sean $A$ un anillo admisible, $\mathfrak J$ un ideal de definición de $A$.
Entonces $\mathfrak J$ está contenido en el radical de $A$.
\end{proposition}

Este enunciado equivale a cualquiera de los corolarios siguientes:

\begin{corollary}[7.1.11]
\label{0.7.1.11-fr}
Para todo $x\in\mathfrak J$, $1+x$ es invertible en $A$.
\end{corollary}

\begin{corollary}[7.1.12]
\label{0.7.1.12-fr}
Para que $f\in A$ sea invertible en $A$, es necesario y suficiente que su imagen
canónica en $A/\mathfrak J$ sea invertible en $A/\mathfrak J$.
\end{corollary}

\begin{corollary}[7.1.13]
\label{0.7.1.13-fr}
Para todo $A$-módulo $M$ de tipo finito, la relación $M=\mathfrak JM$
(equivalente a $M\otimes_A(A/\mathfrak J)=0$) implica $M=0$.
\end{corollary}

\begin{corollary}[7.1.14]
\label{0.7.1.14-fr}
Sea $u:M\to N$ un homomorfismo de $A$-módulos, siendo $N$ de tipo finito;
para que $u$ sea sobreyectivo, es necesario y suficiente que
$u\otimes1:M\otimes_A(A/\mathfrak J)\to N\otimes_A(A/\mathfrak J)$ lo sea.
\end{corollary}

En efecto, la equivalencia de (7.1.10) y (7.1.11) resulta de Bourbaki,
\emph{Alg.}, cap.~VIII, \textsection6, n\textsuperscript{o}~3, teorema~1, y
la de (7.1.10) y (7.1.13) de \emph{loc. cit.}, teorema~2; el hecho de que
(7.1.10) implica (7.1.14) resulta de \emph{loc. cit.}, cor.~4 de la prop.~6;
por otra parte, (7.1.14) implica (7.1.13) aplicándolo al homomorfismo nulo.
Finalmente, (7.1.10) implica que si $f$ es invertible en $A/\mathfrak J$, $f$
no está contenido en ningún ideal maximal de $A$, por tanto $f$ es invertible en
$A$, es decir, (7.1.10) implica (7.1.12); y, recíprocamente, (7.1.12)
implica (7.1.11).

Todo se reduce, pues, a demostrar (7.1.11). Ahora bien, como $A$ es separado y completo y
la sucesión $(\mathfrak J^n)$ tiende a $0$, es inmediato que la serie
$\sum_{n=0}^{\infty}(-1)^n x^n$ es convergente en $A$ y que, si $y$ es
su suma, se tiene $y(1+x)=1$.


% END tree/worktree/es/CURRENT/EGA_I_ch0/ega0_i_7_1.tex
% BEGIN tree/worktree/es/CURRENT/EGA_I_ch0/ega0_i_7_2.tex
\subsection{Anillos ádicos y límites proyectivos}
\label{subsection:0.7.2-fr}

\begin{env}[7.2.1]
\label{0.7.2.1-fr}
Todo límite proyectivo de anillos \emph{discretos} es evidentemente un anillo
linealmente topologizado, separado y completo. Recíprocamente, sea $A$ un anillo
linealmente topologizado, y sea $(\mathfrak J_\lambda)$ un sistema fundamental
de vecindades abiertas de $0$ en $A$ formado

\oldpage[0\textsubscript{I}]{63}

por ideales. Las aplicaciones canónicas
$\varphi_\lambda:A\to A/\mathfrak J_\lambda$ forman un sistema proyectivo de
representaciones continuas y definen, por tanto, una representación continua
$\varphi:A\to\varprojlim A/\mathfrak J_\lambda$; si $A$ es \emph{separado},
$\varphi$ es un isomorfismo topológico de $A$ sobre un
subanillo denso en $\varprojlim A/\mathfrak J_\lambda$; si, además,
$A$ es \emph{completo}, $\varphi$ es un isomorfismo topológico de $A$ \emph{sobre}
$\varprojlim A/\mathfrak J_\lambda$.
\end{env}

\begin{lemma}[7.2.2]
\label{0.7.2.2-fr}
Para que un anillo linealmente topologizado sea admisible, es necesario y suficiente
que sea isomorfo a un límite proyectivo
$A=\varprojlim A_\lambda$, donde $(A_\lambda,u_{\lambda\mu})$ es un sistema
proyectivo de anillos discretos que tiene como conjunto de índices un conjunto ordenado
filtrante $L$ (para $\leq$) que admite un menor elemento denotado por $0$ y satisface
las condiciones siguientes: 1\textsuperscript{o} los
$u_\lambda:A\to A_\lambda$ son sobreyectivos; 2\textsuperscript{o} el núcleo
$\mathfrak J_\lambda$ de $u_{0\lambda}:A_\lambda\to A_0$ es nilpotente.
Cuando esto ocurre, el núcleo $\mathfrak J$ de $u_0:A\to A_0$ es igual a
$\varprojlim\mathfrak J_\lambda$.
\end{lemma}

La necesidad de la condición resulta de (7.2.1), tomando como
$(\mathfrak J_\lambda)$ un sistema fundamental de vecindades de $0$ formado
por ideales de definición contenidos en uno de ellos $\mathfrak J_0$ y aplicando
(7.1.4, (i)). El recíproco resulta de la definición de un límite proyectivo y
de (7.1.2), y la última afirmación es inmediata.

\begin{env}[7.2.3]
\label{0.7.2.3-fr}
Sean $A$ un anillo topológico \emph{admisible}, $\mathfrak J$ un ideal de $A$
contenido en un ideal de definición (es decir (7.1.4), tal que
$(\mathfrak J^n)$ tiende a $0$); se puede considerar sobre $A$ la topología
de anillo que tiene como sistema fundamental de vecindades de $0$ las potencias
$\mathfrak J^n$ ($n>0$); la llamaremos también la topología
\emph{$\mathfrak J$-preádica}. La hipótesis de que $A$ es admisible implica que
$\bigcap_n\mathfrak J^n=0$, por tanto la topología $\mathfrak J$-preádica sobre $A$
es \emph{separada}; sea
$\widehat A=\varprojlim A/\mathfrak J^n$ el completado de $A$ para esta topología
(donde los $A/\mathfrak J^n$ están provistos de la topología discreta), y designemos
por $u$ el homomorfismo de anillos $A\to\widehat A$ (no necesariamente
continuo), límite proyectivo de la sucesión de homomorfismos
$u_n:A\to A/\mathfrak J^n$. Por otra parte, la topología $\mathfrak J$-preádica
sobre $A$ es más fina que la topología dada $\mathcal T$ sobre $A$; como $A$
es separado y completo para $\mathcal T$, se puede prolongar por continuidad
la aplicación identidad de $A$ (provisto de la topología $\mathfrak J$-preádica)
en $A$ provisto de $\mathcal T$; esto da una representación continua
$v:\widehat A\to A$.
\end{env}

\begin{proposition}[7.2.4]
\label{0.7.2.4-fr}
Si $A$ es un anillo admisible y $\mathfrak J$ está contenido en un ideal de
definición de $A$, $A$ es separado y completo para la topología
$\mathfrak J$-preádica.
\end{proposition}

En efecto, con las notaciones de (7.2.3), es inmediato que $v\circ u$ es
la aplicación identidad de $A$. Por otra parte,
$u_n\circ v:\widehat A\to A/\mathfrak J^n$ es la prolongación por continuidad
(para la topología $\mathfrak J$-preádica sobre $A$ y la topología discreta sobre
$A/\mathfrak J^n$) de la aplicación canónica $u_n$; es decir, es
la aplicación canónica de
$\widehat A=\varprojlim_k A/\mathfrak J^k$ sobre $A/\mathfrak J^n$;
$u\circ v$ es, por tanto, límite proyectivo de esta sucesión de aplicaciones,
es decir, por definición la aplicación identidad de $\widehat A$; esto
demuestra la proposición.

\begin{corollary}[7.2.5]
\label{0.7.2.5-fr}
Bajo las hipótesis de \textup{(7.2.3)}, las condiciones siguientes son equivalentes:
\begin{enumerate}
  \item[{\upshape a)}] El homomorfismo $u$ es continuo;

\oldpage[0\textsubscript{I}]{64}

  \item[{\upshape b)}] El homomorfismo $v$ es bicontinuo;
  \item[{\upshape c)}] $A$ es un anillo $\mathfrak J$-ádico.
\end{enumerate}
\end{corollary}

\begin{corollary}[7.2.6]
\label{0.7.2.6-fr}
Sean $A$ un anillo admisible, $\mathfrak J$ un ideal de definición de $A$.
Para que $A$ sea noetheriano, es necesario y suficiente que $A/\mathfrak J$ sea
noetheriano y que $\mathfrak J/\mathfrak J^2$ sea un
$(A/\mathfrak J)$-módulo de tipo finito.
\end{corollary}

Estas condiciones son evidentemente necesarias. Recíprocamente, supongamos que se
verifican; como en virtud de (7.2.4) $A$ es completo para la topología
$\mathfrak J$-preádica, para que sea noetheriano, es necesario y suficiente que
el anillo graduado asociado $\operatorname{grad}(A)$ (para la filtración de los
$\mathfrak J^n$) lo sea ([1], p.~18-07, teorema~4). Ahora bien, sean
$a_1,\ldots,a_n$ unos elementos de $\mathfrak J$ cuyas clases mod.
$\mathfrak J^2$ son generadores de $\mathfrak J/\mathfrak J^2$ como
$A/\mathfrak J$-módulo. Es inmediato por inducción que las clases mod.
$\mathfrak J^{m+1}$ de los monomios de grado total $m$ en los $a_i$
$(1\leq i\leq n)$ forman un sistema de generadores del
$(A/\mathfrak J)$-módulo $\mathfrak J^m/\mathfrak J^{m+1}$. Se concluye que
$\operatorname{grad}(A)$ es un anillo isomorfo a un cociente de
$(A/\mathfrak J)[T_1,\ldots,T_n]$ ($T_i$ indeterminadas), lo que termina la
demostración.

\begin{proposition}[7.2.7]
\label{0.7.2.7-fr}
Sea $(A_i,u_{ij})$ un sistema proyectivo $(i\in\mathbf N)$ de anillos discretos,
y para todo entero $i$, sea $\mathfrak J_i$ el núcleo en $A_i$ del
homomorfismo $u_{0i}:A_i\to A_0$. Se supone que:
\begin{enumerate}
  \item[{\upshape a)}] Para $i\leq j$, $u_{ij}$ es sobreyectivo y su núcleo es
    $\mathfrak J_j^{i+1}$ \textup{(por tanto $A_i$ es isomorfo a
    $A_j/\mathfrak J_j^{i+1}$)}.
  \item[{\upshape b)}] $\mathfrak J_1/\mathfrak J_1^2$ ($=\mathfrak J_1$) es un
    módulo de tipo finito sobre $A_0=A_1/\mathfrak J_1$.
\end{enumerate}
Sea $A=\varprojlim_i A_i$, y para todo entero $n\geq0$, sean $u_n$
el homomorfismo canónico $A\to A_n$, $\mathfrak J^{(n+1)}\subset A$ su
núcleo. En estas condiciones:
\begin{enumerate}
  \item[{\upshape(i)}] $A$ es un anillo ádico, que tiene como ideal de definición
    $\mathfrak J=\mathfrak J^{(1)}$.
  \item[{\upshape(ii)}] Se tiene $\mathfrak J^{(n)}=\mathfrak J^n$ para todo $n\geq1$.
  \item[{\upshape(iii)}] $\mathfrak J/\mathfrak J^2$ es isomorfo a
    $\mathfrak J_1=\mathfrak J_1/\mathfrak J_1^2$, y es, por consiguiente, un módulo
    de tipo finito sobre $A_0=A/\mathfrak J$.
\end{enumerate}
\end{proposition}

La hipótesis de sobreyectividad de los $u_{ij}$ implica que $u_n$ es sobreyectivo;
además, la hipótesis a) implica que $\mathfrak J_j^{j+1}=0$, por tanto $A$ es
un anillo admisible (7.2.2); por definición, los $\mathfrak J^{(n)}$ forman
un sistema fundamental de vecindades de $0$ en $A$, por tanto (ii) implica (i).
Además, se tiene $\mathfrak J=\varprojlim_i\mathfrak J_i$ y las aplicaciones
$\mathfrak J\to\mathfrak J_i$ son sobreyectivas, por tanto (ii) implica (iii), y
se reduce a probar (ii). Por definición, $\mathfrak J^{(n)}$ está formado
por los elementos $(x_k)_{k\geq0}$ de $A$ tales que $x_k=0$ para $k<n$, por tanto
$\mathfrak J^{(n)}\mathfrak J^{(m)}\subset\mathfrak J^{(n+m)}$, es decir,
los $\mathfrak J^{(n)}$ constituyen una \emph{filtración} de $A$. Por otra
parte, $\mathfrak J^{(n)}/\mathfrak J^{(n+1)}$ es isomorfo a la proyección de
$\mathfrak J^{(n)}$ sobre $A_n$; como
$\mathfrak J^{(n)}=\varprojlim_{i>n}\mathfrak J_i^n$, esta proyección no es
otra que $\mathfrak J_n^n$, que es un módulo sobre
$A_0=A_n/\mathfrak J_n$. Sean entonces $a_j=(a_{jk})_{k\geq0}$ unos $r$ elementos de
$\mathfrak J=\mathfrak J^{(1)}$ tales que $a_{11},\ldots,a_{r1}$ forman un
sistema de generadores de $\mathfrak J_1$ sobre $A_0$; vamos a ver que
el conjunto $S_n$ de los monomios de grado total $n$ en los $a_j$ genera el ideal
$\mathfrak J^{(n)}$ de $A$. Como $\mathfrak J_i^{i+1}=0$, es claro, en primer
lugar, que $S_n\subset\mathfrak J^{(n)}$; puesto que $A$ es completo para la
filtración $(\mathfrak J^{(m)})$, basta probar que el conjunto
$\overline S_n$ de las clases mod. $\mathfrak J^{(n+1)}$ de los elementos de $S_n$
genera el módulo graduado $\operatorname{grad}(\mathfrak J^{(n)})$ sobre
el anillo graduado $\operatorname{grad}(A)$ para la filtración anterior
([1], p.~18-06, lema); en virtud de la definición de la multiplicación en
$\operatorname{grad}(A)$,

\oldpage[0\textsubscript{I}]{65}

bastará probar que para todo $m$, $\overline S_m$ es un sistema de
generadores del $A_0$-módulo
$\mathfrak J^{(m)}/\mathfrak J^{(m+1)}$, o, equivalentemente, que $\mathfrak J_m^m$ está
generado por los monomios de grado $m$ en los $a_{jm}$ ($1\leq j\leq r$).
Para ello, queda por mostrar que $\mathfrak J_m$ está (como
$A_m$-módulo) generado por los monomios de grado $\leq m$ con respecto a los
$a_{jm}$; siendo la proposición evidente por definición para $m=1$,
razonemos por inducción sobre $m$, y sea $\mathfrak J'_m$ el sub-$A_m$-módulo
de $\mathfrak J_m$ generado por estos monomios. La relación
$\mathfrak J_{m-1}=\mathfrak J_m/\mathfrak J_m^m$ y la hipótesis de inducción
prueban que $\mathfrak J_m=\mathfrak J'_m+\mathfrak J_m^m$, de donde, puesto que
$\mathfrak J_m^{m+1}=0$, se obtiene
$\mathfrak J_m^m={\mathfrak J'_m}^{\,m}$, y finalmente
$\mathfrak J_m=\mathfrak J'_m$.

\begin{corollary}[7.2.8]
\label{0.7.2.8-fr}
Bajo las condiciones de \textup{(7.2.7)}, para que $A$ sea noetheriano, es necesario y
suficiente que $A_0$ lo sea.
\end{corollary}

Esto resulta inmediatamente de (7.2.6).

\begin{proposition}[7.2.9]
\label{0.7.2.9-fr}
Supongamos verificadas las hipótesis de \textup{(7.2.7)}: para todo entero $i$, sea
$M_i$ un $A_i$-módulo, y, para $i\leq j$, sea
$v_{ij}:M_j\to M_i$ un $u_{ij}$-homomorfismo, tales que $(M_i,v_{ij})$ es un
sistema proyectivo. Supongamos además que $M_0$ es un $A_0$-módulo de tipo
finito, que $v_{ij}$ es sobreyectivo y que su núcleo es
$\mathfrak J_j^{i+1}M_j$. Entonces $M=\varprojlim M_i$ es un $A$-módulo de tipo
finito, y el núcleo del $u_n$-homomorfismo sobreyectivo $v_n:M\to M_n$ es
$\mathfrak J^{n+1}M$ \textup{(de modo que $M_n$ se identifica con
$M/\mathfrak J^{n+1}M=M\otimes_A(A/\mathfrak J^{n+1})$)}.
\end{proposition}

Sean $z_h=(z_{hk})_{k\geq0}$ un sistema de $s$ elementos de $M$ tales que los
$z_{h0}$ ($1\leq h\leq s$) forman un sistema de generadores de $M_0$; se va a
mostrar que los $z_h$ generan el $A$-módulo $M$. El $A$-módulo $M$ es
separado y completo para la filtración formada por los $M^{(n)}$, donde $M^{(n)}$
es el conjunto de los $y=(y_k)_{k\geq0}$ de $M$ tales que $y_k=0$ para $k<n$;
es claro que se tiene $\mathfrak J^{(n)}M\subset M^{(n)}$ y que
$M^{(n)}/M^{(n+1)}=\mathfrak J_n^nM_n$. Se reduce, pues, a mostrar que las
clases de los $z_h$ mod. $M^{(1)}$ generan el módulo graduado
$\operatorname{grad}(M)$ (para la filtración anterior) sobre el anillo graduado
$\operatorname{grad}(A)$ ([1], p.~18-06, lema); para ello, se comprueba
fácilmente que basta de nuevo probar que los $z_{hn}$ ($1\leq h\leq s$)
generan el $A_n$-módulo $M_n$. Se razona de nuevo por inducción sobre
$n$, siendo la proposición evidente por definición para $n=0$; la relación
$M_{n-1}=M_n/\mathfrak J_n^nM_n$ y la hipótesis de inducción muestran que si
$M'_n$ es el submódulo de $M_n$ generado por los $z_{hn}$, se tiene
$M_n=M'_n+\mathfrak J_n^nM_n$, y como $\mathfrak J_n$ es nilpotente, esto
implica $M_n=M'_n$. El mismo razonamiento de paso a los módulos graduados
asociados muestra que la aplicación canónica de $\mathfrak J^{(n)}M$ en
$M^{(n)}$ es sobreyectiva (por tanto biyectiva), es decir, que
$\mathfrak J^{(n)}M=\mathfrak J^nM$ es el núcleo de $M\to M_{n-1}$.

\begin{corollary}[7.2.10]
\label{0.7.2.10-fr}
Sea $(N_i,w_{ij})$ un segundo sistema proyectivo de $A_i$-módulos que verifica las
condiciones de \textup{(7.2.9)}, y sea $N=\varprojlim N_i$. Hay una correspondencia
biunívoca entre los sistemas proyectivos $(h_i)$ de $A_i$-homomorfismos
$h_i:M_i\to N_i$ y los homomorfismos de $A$-módulos $h:M\to N$ (que son
necesariamente continuos para las topologías $\mathfrak J$-ádicas).
\end{corollary}

Es claro que si $h:M\to N$ es un $A$-homomorfismo, se tiene
$h(\mathfrak J^nM)\subset\mathfrak J^nN$, de donde la continuidad de $h$; por
paso a los cocientes, corresponde, pues, a $h$ un sistema proyectivo de
$A_i$-homomorfismos $h_i:M_i\to N_i$, del que $h$ es el límite proyectivo,
de donde el corolario.

\begin{remark}[7.2.11]
\label{0.7.2.11-fr}
Sea $A$ un anillo ádico que tiene un ideal de definición $\mathfrak J$ tal que
$\mathfrak J/\mathfrak J^2$ es un $(A/\mathfrak J)$-módulo de tipo finito;
es claro que los $A_i=A/\mathfrak J^{i+1}$ verifican
\oldpage[0\textsubscript{I}]{66}
las condiciones de (7.2.7); como $A$ es el límite proyectivo de los $A_i$, se
ve que la prop. (7.2.7) da la descripción de \emph{todos} los anillos
ádicos del tipo considerado (y en particular de todos los anillos
\emph{ádicos noetherianos}).
\end{remark}

\begin{example}[7.2.12]
\label{0.7.2.12-fr}
Sean $B$ un anillo, $\mathfrak J$ un ideal de $B$ tal que
$\mathfrak J/\mathfrak J^2$ es un módulo de tipo finito sobre
$B/\mathfrak J$ (o sobre $B$, lo que equivale a lo mismo); pongamos
$A=\varprojlim_n B/\mathfrak J^{n+1}$; $A$ es el separado completado de $B$
provisto de la topología $\mathfrak J$-preádica. Si
$A_n=B/\mathfrak J^{n+1}$, es inmediato que los $A_n$ verifican las
condiciones de (7.2.7); por tanto $A$ es un anillo ádico y si
$\overline{\mathfrak J}$ es la adherencia en $A$ de la imagen canónica de
$\mathfrak J$, $\overline{\mathfrak J}$ es un ideal de definición de $A$,
$\overline{\mathfrak J}^n$ es la adherencia de la imagen canónica de
$\mathfrak J^n$, $A/\overline{\mathfrak J}^n$ se identifica con
$B/\mathfrak J^n$ y
$\overline{\mathfrak J}/\overline{\mathfrak J}^2$ es isomorfo a
$\mathfrak J/\mathfrak J^2$ como
$(A/\overline{\mathfrak J})$-módulo. Del mismo modo, si $N$ es tal que
$N/\mathfrak JN$ es un $B$-módulo de tipo finito, y si se pone
$M_i=N/\mathfrak J^{i+1}N$, $M=\varprojlim M_i$ es un $A$-módulo de tipo
finito, isomorfo al separado completado de $N$ para la topología
$\mathfrak J$-preádica, $\overline{\mathfrak J}^nM$ se identifica con
la adherencia de la imagen canónica de $\mathfrak J^nN$, y
$M/\overline{\mathfrak J}^nM$ con $N/\mathfrak J^nN$.
\end{example}


% END tree/worktree/es/CURRENT/EGA_I_ch0/ega0_i_7_2.tex
% BEGIN tree/worktree/es/CURRENT/EGA_I_ch0/ega0_i_7_3.tex
\subsection{Anillos preádicos noetherianos}
\label{subsection:0.7.3-fr}

\begin{env}[7.3.1]
\label{0.7.3.1-fr}
Sean $A$ un anillo, $\mathfrak J$ un ideal de $A$, $M$ un $A$-módulo;
designaremos por
$\widehat A=\varprojlim_n A/\mathfrak J^n$ (resp.
$\widehat M=\varprojlim_n M/\mathfrak J^nM$) el separado completado de $A$
(resp. $M$) para la topología $\mathfrak J$-preádica. Sea
$M'\xrightarrow{u}M\xrightarrow{v}M''\to0$
una sucesión exacta de $A$-módulos; como
$M/\mathfrak J^nM=M\otimes_A(A/\mathfrak J^n)$, la sucesión
\[
  M'/\mathfrak J^nM'\xrightarrow{u_n}
  M/\mathfrak J^nM\xrightarrow{v_n}
  M''/\mathfrak J^nM''\to0
\]
es exacta para todo $n$. Además, como
$v(\mathfrak J^nM)=\mathfrak J^nv(M)=\mathfrak J^nM''$,
$\widehat v=\varprojlim v_n$ es sobreyectiva (Bourbaki,
\emph{Top. gén.}, cap.~IX, 2.\textsuperscript{a} ed., p.~60, cor.~2).
Por otra parte, si $z=(z_k)$ es un elemento del núcleo de $\widehat v$, para todo
entero $k$, existe un $z'_k\in M'/\mathfrak J^kM'$ tal que
$u_k(z'_k)=z_k$; se concluye que existe
$z'=(z'_n)\in\widehat{M}'$ tal que las $k$ primeras componentes de
$\widehat u(z')$ coinciden con las de $z$; es decir, la imagen por
$\widehat u$ de $\widehat{M}'$ es \emph{densa} en el núcleo de
$\widehat v$.

Si se supone que $A$ es \emph{noetheriano}, lo mismo ocurre con $\widehat A$, en
virtud de (7.2.12), siendo entonces $\mathfrak J/\mathfrak J^2$ un $A$-módulo de
tipo finito. Además:
\end{env}

\begin{namedtheorem}[7.3.2]{Teorema de Krull}
\label{0.7.3.2-fr}
Sean $A$ un anillo \textup{noetheriano}, $\mathfrak J$ un ideal de $A$, $M$ un
$A$-módulo \textup{de tipo finito}, $M'$ un submódulo de $M$; entonces la topología
inducida sobre $M'$ por la topología $\mathfrak J$-preádica de $M$ es idéntica
a la topología $\mathfrak J$-preádica de $M'$.
\end{namedtheorem}

Esto resulta inmediatamente del

\begin{namedlemma}[7.3.2.1]{Lema de Artin-Rees}
\label{0.7.3.2.1-fr}
Bajo las hipótesis de \textup{(7.3.2)}, existe un entero $p$ tal que, para
$n\geq p$, se tiene
\[
  M'\cap\mathfrak J^nM=\mathfrak J^{n-p}
  (M'\cap\mathfrak J^pM)
\]
\end{namedlemma}

Para la demostración, véase ([1], p.~2-04).

\oldpage[0\textsubscript{I}]{67}

\begin{corollary}[7.3.3]
\label{0.7.3.3-fr}
Bajo las hipótesis de \textup{(7.3.2)}, la aplicación canónica
$M\otimes_A\widehat A\to\widehat M$ es biyectiva, y el funtor
$M\otimes_A\widehat A$ es exacto en $M$ en la categoría de los $A$-módulos de
tipo finito; por consiguiente, el separado completado $\mathfrak J$-ádico $\widehat A$
es un $A$-módulo plano \textup{(6.1.1)}.
\end{corollary}

Observemos primero que en la categoría de los $A$-módulos de tipo finito,
$\widehat M$ es un funtor \emph{exacto} en $M$. Sea, en efecto,
$0\to M'\xrightarrow{u}M\xrightarrow{v}M''\to0$
una sucesión exacta; se sabe ya que
$\widehat v:\widehat M\to\widehat{M}''$ es sobreyectivo (7.3.1); por otra parte,
si $i$ es el homomorfismo canónico $M\to\widehat M$, resulta del teorema de
Krull que la adherencia en $\widehat M$ de $i(u(M'))$ se identifica con el separado
completado de $M'$ para la topología $\mathfrak J$-preádica; por tanto
$\widehat u$ es inyectivo, y en virtud de (7.3.1) la imagen de $\widehat u$ es
igual al núcleo de $\widehat v$.

Dicho esto, la aplicación canónica $M\otimes_A\widehat A\to\widehat M$
se obtiene pasando al límite proyectivo en las aplicaciones
$M\otimes_A\widehat A\to M\otimes_A(A/\mathfrak J^n)=M/\mathfrak J^nM$.
Es claro que esta aplicación es biyectiva cuando $M$ es de la forma
$A^p$. Si $M$ es un $A$-módulo de tipo finito, se tiene una sucesión exacta
$A^p\to A^q\to M\to0$, de donde, en virtud de la exactitud \emph{por la derecha} de los
funtores $M\otimes_A\widehat A$ y $\widehat M$ (en $M$) en la categoría
de los $A$-módulos de tipo finito, el diagrama conmutativo
\[
\begin{tikzcd}[column sep=large,row sep=large]
  A^p\otimes\widehat A \arrow[r] \arrow[d] &
  A^q\otimes\widehat A \arrow[r] \arrow[d] &
  M\otimes\widehat A \arrow[r] \arrow[d] &
  0 \\
  \widehat A^{\,p} \arrow[r] &
  \widehat A^{\,q} \arrow[r] &
  \widehat M \arrow[r] &
  0
\end{tikzcd}
\]
donde las dos filas son exactas y las dos primeras flechas verticales son
isomorfismos; se obtiene inmediatamente la conclusión.

\begin{corollary}[7.3.4]
\label{0.7.3.4-fr}
Sean $A$ un anillo noetheriano, $\mathfrak J$ un ideal de $A$, $M$, $N$ dos
$A$-módulos de tipo finito; se tienen isomorfismos canónicos funtoriales
\[
  (M\otimes_A N)^\wedge
  \xrightarrow{\sim}
  \widehat M\otimes_{\widehat A}\widehat N,
  \qquad
  (\operatorname{Hom}_A(M,N))^\wedge
  \xrightarrow{\sim}
  \operatorname{Hom}_{\widehat A}(\widehat M,\widehat N)
\]
\end{corollary}

Esto resulta de (7.3.3), (6.2.1) y (6.2.2).

\begin{corollary}[7.3.5]
\label{0.7.3.5-fr}
Sean $A$ un anillo noetheriano, $\mathfrak J$ un ideal de $A$. Las
condiciones siguientes son equivalentes:
\begin{enumerate}
  \item[$a)$] $\mathfrak J$ está contenido en el radical de $A$.
  \item[$b)$] $\widehat A$ es un $A$-módulo fielmente plano \textup{(6.4.1)}.
  \item[$c)$] Todo $A$-módulo de tipo finito es separado para la topología
  $\mathfrak J$-preádica.
  \item[$d)$] Todo submódulo de un $A$-módulo de tipo finito es cerrado para la
  topología $\mathfrak J$-preádica.
\end{enumerate}
\end{corollary}

Como $\widehat A$ es un $A$-módulo plano, las condiciones $b)$ y $c)$ son
equivalentes, pues $b)$ equivale a decir que si $M$ es un $A$-módulo de tipo
finito, la aplicación canónica
$M\to\widehat M=M\otimes_A\widehat A$ es inyectiva (6.6.1, $c)$). Es
inmediato que $c)$ implica $d)$, pues si $N$ es un submódulo de un $A$-módulo
$M$ de tipo finito, $M/N$ es separado para la topología $\mathfrak J$-preádica,
por tanto $N$ es cerrado en $M$. Mostremos que $d)$ implica $a)$: si
$\mathfrak m$ es un ideal maximal de $A$, $\mathfrak m$ es cerrado en $A$
para la topología $\mathfrak J$-preádica, por tanto $\mathfrak m=\bigcap_{p\geq0}(\mathfrak m+\mathfrak J^p)$,
y como $\mathfrak m+\mathfrak J^p$ es necesariamente igual a $A$ o a
$\mathfrak m$, se tiene $\mathfrak m+\mathfrak J^p=\mathfrak m$ para $p$ suficientemente
\oldpage[0\textsubscript{I}]{68}
grande, de donde $\mathfrak J^p\subset\mathfrak m$ y
$\mathfrak J\subset\mathfrak m$ puesto que $\mathfrak m$ es primo. Finalmente,
$a)$ implica $b)$: sea, en efecto, $P$ la adherencia de $\{0\}$ en un
$A$-módulo $M$ de tipo finito, para la topología $\mathfrak J$-preádica; en
virtud del teorema de Krull (7.3.2), la topología inducida sobre $P$ por la topología
$\mathfrak J$-preádica de $M$ es la topología $\mathfrak J$-preádica de
$P$, por tanto $\mathfrak JP=P$; como $P$ es de tipo finito, resulta del lema
de Nakayama que $P=0$ (estando $\mathfrak J$ contenido en el radical de $A$).

Se observará que las condiciones de (7.3.5) se cumplen cuando $A$ es un
\emph{anillo local noetheriano} y $\mathfrak J\neq A$ un ideal cualquiera de
$A$.

\begin{corollary}[7.3.6]
\label{0.7.3.6-fr}
Si $A$ es un anillo $\mathfrak J$-ádico noetheriano, todo $A$-módulo de tipo
finito es separado y completo para la topología $\mathfrak J$-preádica.
\end{corollary}

Como se tiene entonces $\widehat A=A$, esto resulta inmediatamente de (7.3.3).

Se concluye que la prop. (7.2.9) da la descripción de \emph{todos} los
módulos de tipo finito sobre un anillo ádico noetheriano.

\begin{corollary}[7.3.7]
\label{0.7.3.7-fr}
Bajo las hipótesis de \textup{(7.3.2)}, el núcleo de la aplicación canónica
$M\to\widehat M=M\otimes_A\widehat A$ es el conjunto de los $x\in M$ anulados por
un elemento de $1+\mathfrak J$.
\end{corollary}

En efecto, para que $x\in M$ pertenezca a este núcleo, es necesario y suficiente que
el separado completado del submódulo $Ax$ se reduzca a 0 (por (7.3.2)), es decir,
que $x\in\mathfrak Jx$.


% END tree/worktree/es/CURRENT/EGA_I_ch0/ega0_i_7_3.tex
% BEGIN tree/worktree/es/CURRENT/EGA_I_ch0/ega0_i_7_4.tex
\subsection{Módulos cuasifinitos sobre los anillos locales}
\label{subsection:0.7.4-fr}

\begin{definition}[7.4.1]
\label{0.7.4.1-fr}
Dado un anillo local $A$, de ideal maximal $\mathfrak m$, se dice que un
$A$-módulo $M$ es cuasifinito (sobre $A$) si $M/\mathfrak mM$ es de rango finito
sobre el cuerpo residual $k=A/\mathfrak m$.
\end{definition}

Cuando $A$ es \emph{noetheriano}, el separado completado $\widehat M$ de $M$ para la
topología $\mathfrak m$-preádica es entonces un
$\widehat A$-\emph{módulo de tipo finito}; en efecto, como
$\mathfrak m/\mathfrak m^2$ es entonces un $A$-módulo de tipo finito, esto resulta
de (7.2.12) y de la hipótesis sobre $M/\mathfrak mM$.

En particular, si se supone además que $A$ es \emph{completo} y $M$
\emph{separado} para la topología $\mathfrak m$-preádica (es decir,
$\bigcap_n\mathfrak m^nM=0$), el propio $M$ es un
$A$-\emph{módulo de tipo finito}: en efecto, $\widehat M$ es entonces un
$A$-módulo de tipo finito, y como $M$ se identifica con un submódulo de
$\widehat M$, $M$ es también de tipo finito (y, por lo demás, idéntico a su
completado en virtud de (7.3.6)).

\begin{proposition}[7.4.2]
\label{0.7.4.2-fr}
Sean $A$, $B$ dos anillos locales, $\mathfrak m$, $\mathfrak n$ sus ideales
maximales, y supongamos que $B$ es noetheriano. Sean $\varphi:A\to B$ un
homomorfismo local, $M$ un $B$-módulo de tipo finito. Si $M$ es un $A$-módulo
cuasifinito, las topologías $\mathfrak m$-preádica y $\mathfrak n$-preádica
sobre $M$ son idénticas, por tanto separadas.
\end{proposition}

Observemos que por hipótesis $M/\mathfrak mM$ es de \emph{longitud finita}
como $A$-módulo, por tanto también \emph{a fortiori} como $B$-módulo. Se
concluye que $\mathfrak n$ es el \emph{único ideal primo} de $B$ que contiene
el anulador de $M/\mathfrak mM$: en efecto, se reduce inmediatamente, en virtud
de (1.7.4) y (1.7.2), al caso en que $M/\mathfrak mM$ es \emph{simple}, por tanto
necesariamente isomorfo a $B/\mathfrak n$, y nuestra afirmación es evidente
en este último caso. Por otra parte, como $M$ es un $B$-módulo de tipo finito,
los ideales primos que contienen el anulador de $M/\mathfrak mM$ son los
que contienen $\mathfrak mB+\mathfrak b$, designando por $\mathfrak b$
el anulador del $B$-módulo $M$ (1.7.5). Como $B$ es noetheriano, se concluye
([11], p.~127, cor.~4) que $\mathfrak mB+\mathfrak b$ es un ideal
\oldpage[0\textsubscript{I}]{69}
de definición de $B$, es decir, que existe $k>0$ tal que
$\mathfrak n^k\subset\mathfrak mB+\mathfrak b\subset\mathfrak n$; por consiguiente,
para todo $h>0$
\[
  \mathfrak n^{hk}M
  \subset(\mathfrak mB+\mathfrak b)^hM
  =\mathfrak m^hM
  \subset\mathfrak n^hM
\]
lo que prueba que las topologías $\mathfrak m$-preádica y
$\mathfrak n$-preádica sobre $M$ son las mismas; la segunda es, por lo demás,
separada en virtud de (7.3.5).

\begin{corollary}[7.4.3]
\label{0.7.4.3-fr}
Bajo las hipótesis de \textup{(7.4.2)}, si además $A$ es noetheriano y completo para
la topología $\mathfrak m$-preádica, $M$ es un $A$-módulo de tipo finito.
\end{corollary}

En efecto, $M$ es entonces separado para la topología $\mathfrak m$-preádica, y
nuestra afirmación resulta de la observación hecha en (7.4.1).

\begin{env}[7.4.4]
\label{0.7.4.4-fr}
El caso de aplicación de (7.4.2) que es más importante es aquel en que el propio $B$
es un $A$-\emph{módulo cuasifinito}, lo que equivale a decir que
$B/\mathfrak mB$ es una \emph{álgebra de rango finito sobre}
$k=A/\mathfrak m$; esta condición puede, por lo demás, descomponerse en la
conjunción de las dos siguientes, en virtud de lo que precede:
\begin{enumerate}
  \item[(i)] $\mathfrak mB$ \emph{es un ideal de definición de} $B$;
  \item[(ii)] $B/\mathfrak n$ \emph{es una extensión de rango finito del cuerpo}
  $A/\mathfrak m$.
\end{enumerate}

Cuando esto ocurre, todo $B$-módulo de tipo finito es evidentemente un
$A$-módulo cuasifinito.
\end{env}

\begin{corollary}[7.4.5]
\label{0.7.4.5-fr}
Bajo las hipótesis de \textup{(7.4.2)}, si $\mathfrak b$ es el anulador del
$B$-módulo $M$, $B/\mathfrak b$ es un $A$-módulo cuasifinito.
\end{corollary}

Supongamos $M\neq0$ (en caso contrario el corolario es evidente). Se puede considerar $M$
como un módulo sobre el anillo local noetheriano $B/\mathfrak b$; estando entonces su anulador
reducido a $0$, la demostración de (7.4.2) muestra que
$\mathfrak m(B/\mathfrak b)$ es un ideal de definición de
$B/\mathfrak b$. Por otra parte, $M/\mathfrak nM$ es un espacio vectorial de
rango finito sobre $A/\mathfrak m$, por ser un cociente de $M/\mathfrak mM$, que es
por hipótesis de rango finito sobre $A/\mathfrak m$; como $M\neq0$, se tiene
$M\neq\mathfrak nM$ en virtud del lema de Nakayama; como
$M/\mathfrak nM$ es un espacio vectorial $\neq0$ sobre $B/\mathfrak n$, el hecho
de que sea de rango finito sobre $A/\mathfrak m$ implica que $B/\mathfrak n$ es
también de rango finito sobre $A/\mathfrak m$; la conclusión resulta, pues, de (7.4.4)
aplicado al anillo $B/\mathfrak b$.


% END tree/worktree/es/CURRENT/EGA_I_ch0/ega0_i_7_4.tex
% BEGIN tree/worktree/es/CURRENT/EGA_I_ch0/ega0_i_7_5.tex
\subsection{Anillos de series formales restringidas}
\label{subsection:0.7.5-fr}

\begin{env}[7.5.1]
\label{0.7.5.1-fr}
Sea $A$ un anillo topológico, linealmente topologizado, separado y completo;
sea $(\mathfrak J_\lambda)$ un sistema fundamental de vecindades de $0$ en
$A$ formado por ideales (abiertos), de modo que $A$ se identifica canónicamente con
$\varprojlim A/\mathfrak J_\lambda$ (7.2.1). Para todo $\lambda$, sea
$B_\lambda=(A/\mathfrak J_\lambda)[T_1,\ldots,T_r]$, donde los $T_i$ son
indeterminadas; es claro que los $B_\lambda$ forman un sistema proyectivo
de anillos discretos. Pondremos
$\varprojlim B_\lambda=A\{T_1,\ldots,T_r\}$, y vamos a ver que este anillo
topológico es independiente del sistema fundamental de ideales
$(\mathfrak J_\lambda)$ considerado. De manera precisa, sea $A'$ el subanillo
del anillo de series formales $A[\![T_1,\ldots,T_r]\!]$ formado por las series
formales $\sum_\alpha c_\alpha T^\alpha$ (con
$\alpha=(\alpha_1,\ldots,\alpha_r)\in\mathbf{N}^r$) tales que
$\lim c_\alpha=0$ (según el filtro de los complementarios de las partes
finitas de $\mathbf{N}^r$); diremos que estas series son las series
formales \emph{restringidas} en los $T_i$, con coeficientes en $A$.
\oldpage[0\textsubscript{I}]{70}
Para toda vecindad $V$ de $0$ en $A$, sea $V'$ el conjunto de los
$x=\sum_\alpha c_\alpha T^\alpha\in A'$ tales que $c_\alpha\in V$ para todo
$\alpha$. Se verifica inmediatamente que los $V'$ forman un sistema fundamental de
vecindades de $0$ que define sobre $A'$ una topología de anillo \emph{separada};
vamos a definir canónicamente un \emph{isomorfismo topológico} del
anillo $A\{T_1,\ldots,T_r\}$ sobre $A'$. Para todo
$\alpha\in\mathbf{N}^r$ y todo $\lambda$, sea
$\varphi_{\lambda,\alpha}$ la aplicación de
$(A/\mathfrak J_\lambda)[T_1,\ldots,T_r]$ en
$A/\mathfrak J_\lambda$ que a todo polinomio del primer anillo hace
corresponder el coeficiente de $T^\alpha$ en ese polinomio. Es claro que
los $\varphi_{\lambda,\alpha}$ forman un sistema proyectivo de homomorfismos
de $(A/\mathfrak J_\lambda)$-módulos, cuyo límite proyectivo es un
homomorfismo continuo
$\varphi_\alpha:A\{T_1,\ldots,T_r\}\to A$; vamos a ver que, para todo
$y\in A\{T_1,\ldots,T_r\}$, la serie formal
$\sum_\alpha\varphi_\alpha(y)T^\alpha$ es \emph{restringida}. En efecto, si
$y_\lambda$ es la componente de $y$ en $B_\lambda$, y si se designa por
$H_\lambda$ el conjunto finito de los $\alpha\in\mathbf{N}^r$ para los cuales los
coeficientes del polinomio $y_\lambda$ no son nulos, se tiene
$\varphi_{\mu,\alpha}(y_\mu)\in\mathfrak J_\lambda$ para
$\mathfrak J_\mu\subset\mathfrak J_\lambda$ y $\alpha\notin H_\lambda$ y,
por paso al límite, $\varphi_\alpha(y)\in\mathfrak J_\lambda$ para
$\alpha\notin H_\lambda$. Se define, pues, un homomorfismo de anillos
$\varphi:A\{T_1,\ldots,T_r\}\to A'$ poniendo
$\varphi(y)=\sum_\alpha\varphi_\alpha(y)T^\alpha$, y es inmediato que
$\varphi$ es continuo. Recíprocamente, si $\theta_\lambda$ es el homomorfismo
canónico $A\to A/\mathfrak J_\lambda$, para todo elemento
$z=\sum_\alpha c_\alpha T^\alpha\in A'$ y todo $\lambda$, solo hay un
número finito de índices $\alpha$ tales que $\theta_\lambda(c_\alpha)\neq0$, y
por consiguiente
$\psi_\lambda(z)=\sum_\alpha\theta_\lambda(c_\alpha)T^\alpha$ pertenece a
$B_\lambda$; los $\psi_\lambda$ son continuos y forman un sistema
proyectivo de homomorfismos cuyo límite proyectivo es un homomorfismo
continuo $\psi:A'\to A\{T_1,\ldots,T_r\}$; queda finalmente por verificar que
$\varphi\circ\psi$ y $\psi\circ\varphi$ son los automorfismos identidad,
lo que es inmediato.
\end{env}

\begin{env}[7.5.2]
\label{0.7.5.2-fr}
Identificaremos $A\{T_1,\ldots,T_r\}$ con el anillo $A'$ de las series formales
restringidas mediante los isomorfismos definidos en (7.5.1). Los isomorfismos
canónicos
\[
  ((A/\mathfrak J_\lambda)[T_1,\ldots,T_r])
  [T_{r+1},\ldots,T_s]
  \xrightarrow{\sim}
  (A/\mathfrak J_\lambda)[T_1,\ldots,T_s]
\]
definen, por paso al límite proyectivo, un isomorfismo canónico
\[
  (A\{T_1,\ldots,T_r\})\{T_{r+1},\ldots,T_s\}
  \xrightarrow{\sim}
  A\{T_1,\ldots,T_s\}
\]
\end{env}

\begin{env}[7.5.3]
\label{0.7.5.3-fr}
Para todo homomorfismo continuo $u:A\to B$ de $A$ en un anillo linealmente
topologizado $B$, separado y completo, y todo sistema $(b_1,\ldots,b_r)$ de $r$
elementos de $B$, existe \emph{un homomorfismo continuo y uno solo}
$\overline u:A\{T_1,\ldots,T_r\}\to B$, tal que $\overline u(a)=u(a)$ para
todo $a\in A$ y $\overline u(T_j)=b_j$ para $1\leq j\leq r$. Basta, en
efecto, tomar
\[
  \overline u\Bigl(\sum_\alpha c_\alpha T^\alpha\Bigr)
  =\sum_\alpha u(c_\alpha)b_1^{\alpha_1}\ldots b_r^{\alpha_r}~;
\]
las verificaciones del hecho de que la familia
$(u(c_\alpha)b_1^{\alpha_1}\ldots b_r^{\alpha_r})$ es sumable en $B$ y
de la continuidad de $\overline u$ son inmediatas y se dejan al lector. Se
observará que esta propiedad (para $B$ y los $b_j$ arbitrarios)
\emph{caracteriza} el anillo topológico $A\{T_1,\ldots,T_r\}$ salvo un
isomorfismo único.
\end{env}

\begin{proposition}[7.5.4]
\label{0.7.5.4-fr}
\begin{enumerate}
  \item[\textup{(i)}] Si $A$ es un anillo admisible, lo mismo ocurre con
  $A'=A\{T_1,\ldots,T_r\}$.
  \item[\textup{(ii)}] Sean $A$ un anillo ádico, $\mathfrak J$ un ideal de definición
  de $A$ tal que $\mathfrak J/\mathfrak J^2$ es de tipo finito
\oldpage[0\textsubscript{I}]{71}
  sobre $A/\mathfrak J$. Si se pone $\mathfrak J'=\mathfrak JA'$, $A'$ es entonces
  un anillo $\mathfrak J'$-ádico, y
  $\mathfrak J'/\mathfrak J'^{\,2}$ es de tipo finito sobre
  $A'/\mathfrak J'$. Si, además, $A$ es noetheriano, lo mismo ocurre con
  $A'$.
\end{enumerate}
\end{proposition}

\begin{enumerate}
  \item[(i)] Si $\mathfrak J$ es un ideal de $A$, $\mathfrak J'$ el ideal de
  $A'$ formado por los $\sum_\alpha c_\alpha T^\alpha$ tales que
  $c_\alpha\in\mathfrak J$ para todo $\alpha$, entonces
  $(\mathfrak J')^n\subset(\mathfrak J^n)'$; si $\mathfrak J$ es un ideal de
  definición de $A$, $\mathfrak J'$ es, por tanto, un ideal de definición de $A'$.

  \item[(ii)] Pongamos $A_i=A/\mathfrak J^{i+1}$, y para $i\leq j$, sea
  $u_{ij}$ el homomorfismo canónico
  $A/\mathfrak J^{j+1}\to A/\mathfrak J^{i+1}$; pongamos
  $A'_i=A_i[T_1,\ldots,T_r]$, y sea $u'_{ij}$ el homomorfismo
  $A'_j\to A'_i$ ($i\leq j$) obtenido aplicando $u_{ij}$ a los coeficientes
  de los polinomios de $A'_j$. Mostremos que el sistema proyectivo
  $(A'_i,u'_{ij})$ verifica las condiciones de (7.2.7); como $\mathfrak J'$ es
  el núcleo de $A'\to A'_0$, esto demostrará la primera afirmación de (ii). Ahora bien,
  es claro que los $u'_{ij}$ son sobreyectivos; el núcleo $\mathfrak J'_i$
  de $u'_{0i}$ es el conjunto de los polinomios de
  $A_i[T_1,\ldots,T_r]$ cuyos coeficientes están en
  $\mathfrak J/\mathfrak J^{i+1}$; en particular, $\mathfrak J'_1$ es
  el conjunto de los polinomios de $A_1[T_1,\ldots,T_r]$ cuyos coeficientes
  están en $\mathfrak J/\mathfrak J^2$. Como $\mathfrak J/\mathfrak J^2$
  es de tipo finito sobre $A_1=A/\mathfrak J^2$, se ve ya que
  $\mathfrak J'_1/\mathfrak J_1'^{\,2}$ es un módulo de tipo finito sobre $A'_1$
  (o, lo que equivale a lo mismo, sobre
  $A'_0=A'_1/\mathfrak J'_1$). Mostremos que el núcleo de $u'_{ij}$ es
  $\mathfrak J_j'^{\,i+1}$. Es evidente que
  $\mathfrak J_j'^{\,i+1}$ está contenido en este núcleo. Por otra parte, sean
  $a_1,\ldots,a_m$ unos elementos de $\mathfrak J$ cuyas clases
  mod.~$\mathfrak J^2$ generan $\mathfrak J/\mathfrak J^2$; se verifica
  inmediatamente que las clases mod.~$\mathfrak J^{j+1}$ de los monomios de grado
  $\leq j$ en los $a_k$ ($1\leq k\leq m$) generan
  $\mathfrak J/\mathfrak J^{j+1}$, y las clases de los monomios de grado $>i$
  y $\leq j$ generan, por tanto,
  $\mathfrak J^{i+1}/\mathfrak J^{j+1}$; un monomio en los $T_k$ que tiene tal
  elemento como coeficiente es, por tanto, producto de $i+1$ elementos de
  $\mathfrak J'_j$, lo que establece nuestra afirmación. Finalmente, si $A$ es
  noetheriano, lo mismo ocurre con
  $A'/\mathfrak J'=(A/\mathfrak J)[T_1,\ldots,T_r]$, por tanto $A'$ es noetheriano
  (7.2.8).
\end{enumerate}

\begin{proposition}[7.5.5]
\label{0.7.5.5-fr}
Sean $A$ un anillo noetheriano $\mathfrak J$-ádico, $B$ un anillo topológico
admisible, $\varphi:A\to B$ un homomorfismo continuo, que hace de $B$ una
$A$-álgebra. Las condiciones siguientes son equivalentes:
\begin{enumerate}
  \item[a)] $B$ es noetheriano y $\mathfrak JB$-ádico, y
  $B/\mathfrak JB$ es un álgebra de tipo finito sobre $A/\mathfrak J$.

  \item[b)] $B$ es topológicamente $A$-isomorfo a $\varprojlim B_n$, donde
  $B_n=B_m/\mathfrak J^{n+1}B_m$ para $m\geq n$, y $B_1$ es un álgebra de
  tipo finito sobre $A_1=A/\mathfrak J^2$.

  \item[c)] $B$ es topológicamente $A$-isomorfo a un cociente de un álgebra
  de la forma\linebreak $A\{T_1,\ldots,T_r\}$ por un ideal (necesariamente cerrado en virtud
  de \textup{(7.3.6)} y \textup{(7.5.4 (ii))}).
\end{enumerate}
\end{proposition}

Como $A$ es noetheriano, lo mismo ocurre con
$A'=A\{T_1,\ldots,T_r\}$ (7.5.4), por tanto \emph{c)} implica que $B$ es noetheriano;
como $\mathfrak J'=\mathfrak JA'$ es una vecindad abierta de 0 en $A'$ tal
que los $\mathfrak J'^{n}$ forman un sistema fundamental de vecindades de 0,
las imágenes $\mathfrak J^nB$ de los $\mathfrak J'^{n}$ forman un sistema
fundamental de vecindades de 0 en $B$, y como se sabe que $B$ es separado y
completo, $B$ es un anillo $\mathfrak JB$-ádico. Finalmente, $B/\mathfrak JB$ es
un álgebra (sobre $A/\mathfrak J$) cociente de
$A'/\mathfrak JA'=(A/\mathfrak J)[T_1,\ldots,T_r]$, por tanto es de tipo finito, lo
que termina de probar que \emph{c)} implica \emph{a)}.

Si $B$ es $\mathfrak JB$-ádico y noetheriano, $B$ es isomorfo a
$\varprojlim B_n$, donde $B_n=B/\mathfrak J^{n+1}B$ (7.2.11), y
$\mathfrak JB/\mathfrak J^2B$ es un módulo de tipo finito sobre
$B/\mathfrak JB$. Sea $(a_j)_{1\leq j\leq s}$ un sistema de generadores del
$(B/\mathfrak JB)$-módulo $\mathfrak JB/\mathfrak J^2B$, y
$(c_i)_{1\leq i\leq r}$ un sistema de elementos de $B/\mathfrak J^2B$ cuyas
clases
\oldpage[0\textsubscript{I}]{72}
mod.~$\mathfrak JB/\mathfrak J^2B$ forman un sistema de generadores de la
$(A/\mathfrak J)$-álgebra $B/\mathfrak JB$; se ve inmediatamente que los $c_i,a_j$
forman un sistema de generadores de la $(A/\mathfrak J^2)$-álgebra
$B/\mathfrak J^2B$, por tanto \emph{a)} implica \emph{b)}.

Queda por probar que \emph{b)} implica \emph{c)}. La hipótesis implica que $B_1$ es un
anillo noetheriano, y como $B_1=B_2/\mathfrak J^2B_2$, se tiene
$\mathfrak J^2B_1=0$, por tanto
$\mathfrak JB_1=\mathfrak JB_1/\mathfrak J^2B_1$ es un $B_0$-módulo de tipo
finito. Las condiciones de (7.2.7) son, pues, verificadas por el sistema proyectivo
$(B_n)$ y $B$ es un anillo $\mathfrak JB$-ádico. Sea
$(c_i)_{1\leq i\leq r}$ un sistema finito de elementos de $B$ cuyas clases
mod.~$\mathfrak JB$ generan la $(A/\mathfrak J)$-álgebra
$B/\mathfrak JB$, y cuyas combinaciones lineales con coeficientes en
$\mathfrak J$ son tales que sus clases mod.~$\mathfrak J^2B$ generan
el $B_0$-módulo $\mathfrak JB/\mathfrak J^2B$. Existe un
$A$-homomorfismo continuo $u$ de $A'=A\{T_1,\ldots,T_r\}$ en $B$ que se
reduce a $\varphi$ en $A$ y es tal que $u(T_i)=c_i$ para $1\leq i\leq r$
(7.5.3); si probamos que $u$ es \emph{sobreyectivo}, \emph{c)} quedará establecido, pues de
$u(A')=B$ se deducirá $u(\mathfrak J^nA')=\mathfrak J^nB$, es decir, $u$
será un morfismo estricto de anillos topológicos y $B$ será, por tanto, isomorfo a un
cociente de $A'$ por un ideal cerrado. Ahora bien, como $B$ es completo para la topología
$\mathfrak JB$-ádica, basta ([1], p.~18-07) mostrar que
el homomorfismo $\operatorname{grad}(A')\to\operatorname{grad}(B)$ deducido
canónicamente de $u$ para las filtraciones $\mathfrak J$-ádicas sobre $A'$ y
$B$ es sobreyectivo. Pero, por definición, los homomorfismos
$A'/\mathfrak JA'\to B/\mathfrak JB$ y
$\mathfrak JA'/\mathfrak J^2A'\to\mathfrak JB/\mathfrak J^2B$ deducidos de $u$
son sobreyectivos; por inducción sobre $n$, se deduce inmediatamente que lo mismo
ocurre con $\mathfrak JA'/\mathfrak J^nA'\to\mathfrak JB/\mathfrak J^nB$, y
\emph{a fortiori} con
$\mathfrak J^nA'/\mathfrak J^{n+1}A'\to
\mathfrak J^nB/\mathfrak J^{n+1}B$, lo que termina la demostración.


% END tree/worktree/es/CURRENT/EGA_I_ch0/ega0_i_7_5.tex
% BEGIN tree/worktree/es/CURRENT/EGA_I_ch0/ega0_i_7_6.tex
\subsection{Anillos completos de fracciones}
\label{subsection:0.7.6-fr}

\begin{env}[7.6.1]
\label{0.7.6.1-fr}
Sean $A$ un anillo linealmente topologizado, $(\mathfrak J_\lambda)$ un
sistema fundamental de vecindades de 0 en $A$ formado por ideales, $S$ una parte
multiplicativa de $A$. Sea $u_\lambda$ el homomorfismo canónico
$A\to A_\lambda=A/\mathfrak J_\lambda$, y para
$\mathfrak J_\mu\subset\mathfrak J_\lambda$, sea $u_{\lambda\mu}$
el homomorfismo canónico $A_\mu\to A_\lambda$. Pongamos
$S_\lambda=u_\lambda(S)$, de modo que
$u_{\lambda\mu}(S_\mu)=S_\lambda$. Los $u_{\lambda\mu}$ dan
canónicamente homomorfismos sobreyectivos
$S_\mu^{-1}A_\mu\to S_\lambda^{-1}A_\lambda$, para los cuales estos anillos
forman un sistema proyectivo; designemos por $A\{S^{-1}\}$ el límite
proyectivo de este sistema. Esta definición no depende del sistema fundamental
de vecindades $(\mathfrak J_\lambda)$ elegido; en efecto:
\end{env}

\begin{proposition}[7.6.2]
\label{0.7.6.2-fr}
El anillo $A\{S^{-1}\}$ es topológicamente isomorfo al separado completado del
anillo $S^{-1}A$ para la topología en la que un sistema fundamental de vecindades
de 0 está formado por los $S^{-1}\mathfrak J_\lambda$.
\end{proposition}

En efecto, si $v_\lambda$ es el homomorfismo canónico
$S^{-1}A\to S_\lambda^{-1}A_\lambda$ deducido de $u_\lambda$, el núcleo de
$v_\lambda$ es $S^{-1}\mathfrak J_\lambda$ y $v_\lambda$ es sobreyectivo,
de donde la proposición (7.2.1).

\begin{corollary}[7.6.3]
\label{0.7.6.3-fr}
Si $S'$ es la imagen canónica de $S$ en el separado completado $\widehat A$ de
$A$, $A\{S^{-1}\}$ se identifica canónicamente con
$\widehat A\{{S'}^{-1}\}$.
\end{corollary}

Se observará que aun si $A$ es separado y completo, no ocurre lo mismo con
$S^{-1}A$ para la topología definida por los
$S^{-1}\mathfrak J_\lambda$, como se ve, por ejemplo, tomando como $S$
el conjunto de los $f^n$ ($n\geq0$), donde $f$ es topológicamente nilpotente y no
nilpotente: en efecto, $S^{-1}A$ no se reduce a 0 y, por otra parte, para todo
$\lambda$ existe $n$ tal que $f^n\in\mathfrak J_\lambda$, por tanto
$1=f^n/f^n\in S^{-1}\mathfrak J_\lambda$ y
$S^{-1}\mathfrak J_\lambda=S^{-1}A$.

\oldpage[0\textsubscript{I}]{73}

\begin{corollary}[7.6.4]
\label{0.7.6.4-fr}
Si, en $A$, $0$ no es adherente a $S$, el anillo $A\{S^{-1}\}$ no se
reduce a $0$.
\end{corollary}

En efecto, $0$ no es adherente a $\{1\}$ en el anillo $S^{-1}A$; en caso contrario,
se tendría $1\in S^{-1}\mathfrak J_\lambda$ para todo ideal abierto
$\mathfrak J_\lambda$ de $A$, y resultaría que
$\mathfrak J_\lambda\cap S\neq\varnothing$ para todo $\lambda$, contrariamente
a la hipótesis.

\begin{env}[7.6.5]
\label{0.7.6.5-fr}
Diremos que $A\{S^{-1}\}$ es el \emph{anillo completo de las fracciones} de $A$
que tienen sus denominadores en $S$. Con las notaciones anteriores, es claro
que la imagen inversa de $S^{-1}\mathfrak J_\lambda$ en $A$ contiene
$\mathfrak J_\lambda$, por tanto la aplicación canónica $A\to S^{-1}A$ es
continua, y si se la compone con la aplicación canónica
$S^{-1}A\to A\{S^{-1}\}$, se obtiene un homomorfismo canónico continuo
$A\to A\{S^{-1}\}$, límite proyectivo de los homomorfismos
$A\to S_\lambda^{-1}A_\lambda$.
\end{env}

\begin{env}[7.6.6]
\label{0.7.6.6-fr}
El par formado por $A\{S^{-1}\}$ y por la aplicación canónica
$A\to A\{S^{-1}\}$ está caracterizado por la \emph{propiedad universal}
siguiente: todo homomorfismo continuo $u$ de $A$ en un anillo linealmente
topologizado $B$, separado y completo, tal que $u(S)$ está formado por elementos
invertibles en $B$, se factoriza de una sola manera en
$A\to A\{S^{-1}\}\xrightarrow{u'}B$, donde $u'$ es continuo. En efecto, $u$ se
factoriza de una sola manera en $A\to S^{-1}A\xrightarrow{v'}B$; como para
todo ideal abierto $\mathfrak K$ de $B$, $u^{-1}(\mathfrak K)$ contiene un
$\mathfrak J_\lambda$, ${v'}^{-1}(\mathfrak K)$ contiene necesariamente
$S^{-1}\mathfrak J_\lambda$, por tanto $v'$ es continuo; puesto que $B$ es separado y
completo, $v'$ se factoriza de una sola manera en
$S^{-1}A\to A\{S^{-1}\}\xrightarrow{u'}B$, donde $u'$ es continuo; de donde nuestra
afirmación.
\end{env}

\begin{env}[7.6.7]
\label{0.7.6.7-fr}
Sean $B$ un segundo anillo linealmente topologizado, $T$ una parte
multiplicativa de $B$, $\varphi:A\to B$ un homomorfismo continuo tal que
$\varphi(S)\subset T$. Según lo que precede, el homomorfismo continuo
$A\xrightarrow{\varphi}B\to B\{T^{-1}\}$ se factoriza de manera única en
$A\to A\{S^{-1}\}\xrightarrow{\varphi'}B\{T^{-1}\}$, donde $\varphi'$ es
continuo. En particular, si $B=A$ y si $\varphi$ es la identidad, se ve que
para $S\subset T$ se tiene un homomorfismo continuo
$\rho^{T,S}:A\{S^{-1}\}\to A\{T^{-1}\}$ obtenido por paso al separado completado
a partir de $S^{-1}A\to T^{-1}A$; si $U$ es una tercera parte
multiplicativa de $A$ tal que $S\subset T\subset U$, se tiene
$\rho^{U,S}=\rho^{U,T}\circ\rho^{T,S}$.
\end{env}

\begin{env}[7.6.8]
\label{0.7.6.8-fr}
Sean $S_1$, $S_2$ dos partes multiplicativas de $A$, y sea $S_2'$ la imagen
canónica de $S_2$ en $A\{S_1^{-1}\}$; se tiene entonces un isomorfismo
topológico canónico
$A\{(S_1S_2)^{-1}\}\xrightarrow{\sim}A\{S_1^{-1}\}\{{S_2'}^{-1}\}$,
como se ve partiendo del isomorfismo canónico
$(S_1S_2)^{-1}A\xrightarrow{\sim}{S_2''}^{-1}(S_1^{-1}A)$ (donde $S_2''$ es
la imagen canónica de $S_2$ en $S_1^{-1}A$), que es bicontinuo.
\end{env}

\begin{env}[7.6.9]
\label{0.7.6.9-fr}
Sea $\mathfrak a$ un ideal \emph{abierto} de $A$; se puede suponer que
$\mathfrak J_\lambda\subset\mathfrak a$ para todo $\lambda$, y, por consiguiente,
$S^{-1}\mathfrak J_\lambda\subset S^{-1}\mathfrak a$ en el anillo
$S^{-1}A$, es decir, $S^{-1}\mathfrak a$ es un ideal \emph{abierto} de
$S^{-1}A$; designaremos por $\mathfrak a\{S^{-1}\}$ su separado completado,
igual a $\varprojlim(S^{-1}\mathfrak a/S^{-1}\mathfrak J_\lambda)$, que es un
ideal \emph{abierto} de $A\{S^{-1}\}$, isomorfo a la adherencia de la imagen
canónica de $S^{-1}\mathfrak a$. Además, \emph{el anillo discreto
$A\{S^{-1}\}/\mathfrak a\{S^{-1}\}$ es canónicamente isomorfo a
$S^{-1}A/S^{-1}\mathfrak a=S^{-1}(A/\mathfrak a)$}. Recíprocamente, si
$\mathfrak a'$ es un ideal abierto de $A\{S^{-1}\}$, $\mathfrak a'$ contiene
un ideal de la forma $\mathfrak J_\lambda\{S^{-1}\}$, por tanto es la imagen
inversa de un ideal de $S^{-1}A/S^{-1}\mathfrak J_\lambda$, que es
necesariamente (1.2.6) de la forma $S^{-1}\mathfrak a$, donde
$\mathfrak a\supset\mathfrak J_\lambda$. Se concluye que se tiene
$\mathfrak a'=\mathfrak a\{S^{-1}\}$. En particular (1.2.6):
\end{env}

\begin{proposition}[7.6.10]
\label{0.7.6.10-fr}
La aplicación $\mathfrak p\to\mathfrak p\{S^{-1}\}$ es una biyección
creciente del conjunto de los ideales \textup{primos abiertos} $\mathfrak p$ de $A$ tales
que $\mathfrak p\cap S=\varnothing$ sobre el conjunto de los ideales \textup{primos abiertos}
\oldpage[0\textsubscript{I}]{74}
de $A\{S^{-1}\}$; además, el cuerpo de fracciones de
$A\{S^{-1}\}/\mathfrak p\{S^{-1}\}$ es canónicamente isomorfo al de
$A/\mathfrak p$.
\end{proposition}

\begin{proposition}[7.6.11]
\label{0.7.6.11-fr}
\begin{enumerate}
  \item[\rm(i)] Si $A$ es un anillo admisible, lo mismo ocurre con
    $A'=A\{S^{-1}\}$ y para todo ideal de definición $\mathfrak J$ de $A$,
    $\mathfrak J'=\mathfrak J\{S^{-1}\}$ es un ideal de definición de $A'$.
  \item[\rm(ii)] Sean $A$ un anillo ádico, $\mathfrak J$ un ideal de
    definición de $A$ tal que $\mathfrak J/\mathfrak J^2$ es de tipo finito
    sobre $A/\mathfrak J$; entonces $A'$ es un anillo $\mathfrak J'$-ádico y
    $\mathfrak J'/\mathfrak J'^2$ es de tipo finito sobre
    $A'/\mathfrak J'$. Si además $A$ es noetheriano, lo mismo ocurre con
    $A'$.
\end{enumerate}
\end{proposition}

\begin{enumerate}
  \item[\rm(i)] Si $\mathfrak J$ es un ideal de definición en $A$, es
    claro que $S^{-1}\mathfrak J$ es un ideal de definición en el anillo
    topológico $S^{-1}A$, pues se tiene
    $(S^{-1}\mathfrak J)^n=S^{-1}\mathfrak J^n$. Sea $A''$ el anillo separado
    asociado a $S^{-1}A$, $\mathfrak J''$ la imagen de $S^{-1}\mathfrak J$ en
    $A''$; la imagen de $S^{-1}\mathfrak J^n$ es $\mathfrak J''^n$, por tanto
    $\mathfrak J''^n$ tiende a 0 en $A''$; como $\mathfrak J'$ es
    la adherencia de $\mathfrak J''$ en $A'$, $\mathfrak J'^n$ está contenido
    en la adherencia de $\mathfrak J''^n$, por tanto tiende a 0 en $A'$.
  \item[\rm(ii)] Pongamos $A_i=A/\mathfrak J^{i+1}$, y para $i\leq j$, sea
    $u_{ij}$ el homomorfismo canónico
    $A/\mathfrak J^{j+1}\to A/\mathfrak J^{i+1}$; sea $S_i$ la imagen
    canónica de $S$ en $A_i$, y pongamos $A_i'=S_i^{-1}A_i$; sea finalmente
    $u_{ij}':A_j'\to A_i'$ el homomorfismo deducido canónicamente de $u_{ij}$.
    Mostremos que el sistema proyectivo $(A_i',u_{ij}')$ verifica las condiciones
    de la prop.~(7.2.7): es claro que los $u_{ij}'$ son sobreyectivos;
    por otra parte, el núcleo de $u_{ij}'$ es
    $S_j^{-1}(\mathfrak J^{i+1}/\mathfrak J^{j+1})$ (1.3.2), igual a
    $\mathfrak J_j'^{i+1}$, donde
    $\mathfrak J_j'=S_j^{-1}(\mathfrak J/\mathfrak J^{j+1})$; finalmente,
    $\mathfrak J_1'/\mathfrak J_1'^2=S_1^{-1}(\mathfrak J/\mathfrak J^2)$,
    y como $\mathfrak J/\mathfrak J^2$ es de tipo finito sobre
    $A/\mathfrak J^2$, $\mathfrak J_1'/\mathfrak J_1'^2$ es de tipo finito
    sobre $A_1'$. Finalmente, si $A$ es noetheriano, lo mismo ocurre con
    $A_0'=S_0^{-1}(A/\mathfrak J)$, lo que, por el corolario (7.2.8), termina de demostrar la proposición
    (7.6.11).
\end{enumerate}

\begin{corollary}[7.6.12]
\label{0.7.6.12-fr}
Bajo las hipótesis de \textup{(7.6.11, (ii))}, se tiene
$(\mathfrak J\{S^{-1}\})^n=\mathfrak J^n\{S^{-1}\}$.
\end{corollary}

Esto resulta, en efecto, de (7.2.7) y de la demostración de (7.6.11).

\begin{proposition}[7.6.13]
\label{0.7.6.13-fr}
Sean $A$ un anillo ádico noetheriano, $S$ una parte multiplicativa de $A$;
entonces $A\{S^{-1}\}$ es un $A$-módulo plano.
\end{proposition}

En efecto, si $\mathfrak J$ es un ideal de definición de $A$,
$A\{S^{-1}\}$ es el separado completado del anillo noetheriano $S^{-1}A$ provisto
de la topología $S^{-1}\mathfrak J$-preádica; por consiguiente (7.3.3),
$A\{S^{-1}\}$ es un $S^{-1}A$-módulo plano; como $S^{-1}A$ es un
$A$-módulo plano (6.3.1), la proposición resulta de la transitividad de la
planitud (6.2.1).

\begin{corollary}[7.6.14]
\label{0.7.6.14-fr}
Bajo las hipótesis de \textup{(7.6.13)}, sea $S'\subset S$ una segunda parte
multiplicativa de $A$; entonces $A\{S^{-1}\}$ es un
$A\{{S'}^{-1}\}$-módulo plano.
\end{corollary}

En efecto (7.6.8), $A\{S^{-1}\}$ se identifica canónicamente con
$A\{{S'}^{-1}\}\{S_0^{-1}\}$, donde $S_0$ es la imagen canónica de $S$ en
$A\{{S'}^{-1}\}$, y $A\{{S'}^{-1}\}$ es noetheriano (7.6.11).

\begin{env}[7.6.15]
\label{0.7.6.15-fr}
Para todo elemento $f$ de un anillo linealmente topologizado $A$,
designaremos por $A_{\{f\}}$ el anillo completo de fracciones $A\{S_f^{-1}\}$,
donde $S_f$ es el conjunto multiplicativo de los $f^n$ ($n\geq 0$); para todo
ideal abierto $\mathfrak a$ de $A$, escribiremos $\mathfrak a_{\{f\}}$ en
lugar de $\mathfrak a\{S_f^{-1}\}$. Si $g$ es un segundo elemento de $A$, se tiene
un homomorfismo canónico continuo $A_{\{f\}}\to A_{\{fg\}}$ (7.6.7).
Cuando $f$ recorre una parte multiplicativa $S$ de $A$, los $A_{\{f\}}$
forman, pues, un sistema inductivo filtrante de anillos para los homomorfismos
anteriores; pondremos
$A_{\{S\}}=\varinjlim_{f\in S}A_{\{f\}}$. Para todo $f\in S$, se tiene un
homomorfismo $A_{\{f\}}\to A\{S^{-1}\}$ (7.6.7), y
\oldpage[0\textsubscript{I}]{75}
estos homomorfismos forman un sistema inductivo; por paso al límite
inductivo, definen, por tanto, un homomorfismo canónico
$A_{\{S\}}\to A\{S^{-1}\}$.
\end{env}

\begin{proposition}[7.6.16]
\label{0.7.6.16-fr}
Si $A$ es un anillo noetheriano, $A\{S^{-1}\}$ es un módulo plano sobre
$A_{\{S\}}$.
\end{proposition}

En efecto (7.6.14), $A\{S^{-1}\}$ es plano sobre cada uno de los anillos
$A_{\{f\}}$ para $f\in S$, y la conclusión resulta de (6.2.3).

\begin{proposition}[7.6.17]
\label{0.7.6.17-fr}
Sea $\mathfrak p$ un ideal primo abierto en un anillo admisible $A$, y
sea $S=A\setmin\mathfrak p$. Entonces los anillos $A\{S^{-1}\}$ y $A_{\{S\}}$ son
anillos locales, el homomorfismo canónico
$A_{\{S\}}\to A\{S^{-1}\}$ es local y los cuerpos residuales de $A_{\{S\}}$
y $A\{S^{-1}\}$ son canónicamente isomorfos al cuerpo de fracciones de
$A/\mathfrak p$.
\end{proposition}

En efecto, sea $\mathfrak J\subset\mathfrak p$ un ideal de definición de
$A$; se tiene $S^{-1}\mathfrak J\subset S^{-1}\mathfrak p=\mathfrak p
A_{\mathfrak p}$, por tanto $A_{\mathfrak p}/S^{-1}\mathfrak J$ es un anillo
local; se concluye de (7.1.12), (7.6.9) y (7.6.11, (i)) que
$A\{S^{-1}\}$ es un anillo local. Pongamos
$\mathfrak m=\varinjlim_{f\in S}\mathfrak p_{\{f\}}$, que es un ideal de
$A_{\{S\}}$; vamos a ver que todo elemento de $A_{\{S\}}$ que no pertenece
a $\mathfrak m$ es invertible. En efecto, tal elemento es la imagen en
$A_{\{S\}}$ de un elemento $z\in A_{\{f\}}$ que no pertenece a
$\mathfrak p_{\{f\}}$, para un $f\in S$; su imagen canónica $z_0$ en
$A_{\{f\}}/\mathfrak J_{\{f\}}=S_f^{-1}(A/\mathfrak J)$ no pertenece, pues,
a $S_f^{-1}(\mathfrak p/\mathfrak J)$ (7.6.9), lo que significa que
$z_0=\overline{x}/\overline{f}^k$, donde $x\notin\mathfrak p$ y
$\overline{x},\overline{f}$ son las clases de $x,f$ mod.~$\mathfrak J$.
Como $x\in S$, se tiene $g=xf\in S$, y en $S_g^{-1}A$,
$y_0=x^{k+1}/g^k$, imagen canónica de $x/f^k\in S_f^{-1}A$, admite un
inverso $x^{k-1}f^{2k}/g^k$. Esto implica \emph{a fortiori} que la imagen de
$y_0$ en $S_g^{-1}A/S_g^{-1}\mathfrak J$ es invertible, por tanto (7.6.9 y
7.1.12) la imagen canónica $y$ de $z$ en $A_{\{g\}}$ es invertible;
la imagen de $z$ en $A_{\{S\}}$ (igual a la de $y$) es, por consiguiente,
invertible. Se ve, pues, que $A_{\{S\}}$ es un anillo local de ideal maximal
$\mathfrak m$; además, la imagen de $\mathfrak p_{\{f\}}$ en
$A\{S^{-1}\}$ está contenida en el ideal maximal $\mathfrak p\{S^{-1}\}$ de
este anillo; \emph{a fortiori}, la imagen de $\mathfrak m$ en
$A\{S^{-1}\}$ está contenida en $\mathfrak p\{S^{-1}\}$, por tanto
el homomorfismo canónico $A_{\{S\}}\to A\{S^{-1}\}$ es local. Finalmente,
como todo elemento de $A\{S^{-1}\}/\mathfrak p\{S^{-1}\}$ es imagen de un
elemento de un anillo $S_f^{-1}A$ para un $f\in S$ adecuado,
el homomorfismo $A_{\{S\}}\to A\{S^{-1}\}/\mathfrak p\{S^{-1}\}$ es
sobreyectivo, y da, por tanto, por paso a los cocientes un isomorfismo de los cuerpos
residuales.

\begin{corollary}[7.6.18]
\label{0.7.6.18-fr}
Bajo las hipótesis de \textup{(7.6.17)}, si se supone además que $A$ es un anillo
ádico noetheriano, los anillos locales $A\{S^{-1}\}$ y $A_{\{S\}}$ son
noetherianos, y $A\{S^{-1}\}$ es un $A_{\{S\}}$-módulo fielmente plano.
\end{corollary}

Se sabe ya (7.6.11, (ii)) que $A\{S^{-1}\}$ es noetheriano y
$A_{\{S\}}$-plano (7.6.16); como el homomorfismo
$A_{\{S\}}\to A\{S^{-1}\}$ es local, se concluye que $A\{S^{-1}\}$ es
un $A_{\{S\}}$-módulo fielmente plano (6.6.2), y, por consiguiente, que $A_{\{S\}}$
es noetheriano (6.5.2).


% END tree/worktree/es/CURRENT/EGA_I_ch0/ega0_i_7_6.tex
% BEGIN tree/worktree/es/CURRENT/EGA_I_ch0/ega0_i_7_7.tex
\subsection{Productos tensoriales completados}
\label{subsection:0.7.7-fr}

\begin{env}[7.7.1]
\label{0.7.7.1-fr}
Sean $A$ un anillo linealmente topologizado, $M$, $N$, dos $A$-módulos
linealmente topologizados. Sean $\mathfrak J$, $V$, $W$ unas vecindades
abiertas de $0$ en $A$, $M$, $N$ respectivamente, que sean
$A$-módulos, y tales que $\mathfrak J.M\subset V$, $\mathfrak J.N\subset W$,
de modo que $M/V$ y $N/W$ pueden considerarse como
$(A/\mathfrak J)$-módulos. Cuando $\mathfrak J$, $V$, $W$ recorren los
sistemas de vecindades abiertas que verifican las condiciones anteriores, es
inmediato que los módulos
$(M/V)\otimes_{A/\mathfrak J}(N/W)$ forman un sistema proyectivo de módulos
\oldpage[0\textsubscript{I}]{76}
sobre el sistema proyectivo de anillos $A/\mathfrak J$; por paso al límite
proyectivo, se deduce, pues, un módulo sobre el separado completado $\widehat A$
de $A$, que se llama el \emph{producto tensorial completado} de $M$ y de
$N$ y que se denota por $(M\otimes_A N)^\wedge$. Si se observa que $M/V$
es canónicamente isomorfo a $\widehat M/\overline V$, donde $\widehat M$ es
el separado completado de $M$ y $\overline V$ la adherencia en $\widehat M$ de
la imagen de $V$, se ve que el producto tensorial completado
$(M\otimes_A N)^\wedge$ se identifica canónicamente con
$(\widehat M\otimes_{\widehat A}\widehat N)^\wedge$, que se denota también por
$\widehat M\widehat\otimes_{\widehat A}\widehat N$.
\end{env}

\begin{env}[7.7.2]
\label{0.7.7.2-fr}
\begin{sloppypar}
Con las notaciones anteriores, los productos tensoriales
$(M/V)\otimes_A(N/W)$ y
$(M/V)\otimes_{A/\mathfrak J}(N/W)$ se identifican canónicamente;
se identifican, por tanto, también con
$(M\otimes_A N)/(\Im(V\otimes_A N)+\Im(M\otimes_A W))$.
Se concluye que $(M\otimes_A N)^\wedge$ es el \emph{separado completado del
$A$-módulo $M\otimes_A N$, provisto de la topología para la cual los
submódulos}
\[
  \Im(V\otimes_A N)+\Im(M\otimes_A W)
\]
\emph{forman un sistema fundamental de vecindades de $0$} (recorriendo $V$ y $W$
el conjunto de los submódulos abiertos de $M$ y $N$
respectivamente); diremos, para abreviar, que esta topología es el
\emph{producto tensorial} de las topologías dadas sobre $M$ y $N$.
\end{sloppypar}
\end{env}

\begin{env}[7.7.3]
\label{0.7.7.3-fr}
Sean $M'$, $N'$ dos $A$-módulos linealmente topologizados,
$u:M\to M'$, $v:N\to N'$ dos homomorfismos continuos; es inmediato
que $u\otimes v$ es continuo para las topologías producto tensorial sobre
$M\otimes N$ y $M'\otimes N'$ respectivamente; por paso al separado
completado, se deduce un homomorfismo continuo
$(M\otimes N)^\wedge\to(M'\otimes N')^\wedge$, que designaremos
por $u\widehat\otimes v$; $(M\otimes_A N)^\wedge$ es, por tanto, un
\emph{bifuntor} en $M$ y $N$ en la categoría de los $A$-módulos
linealmente topologizados.
\end{env}

\begin{env}[7.7.4]
\label{0.7.7.4-fr}
Se define del mismo modo el producto tensorial completado de un número finito cualquiera
de $A$-módulos linealmente topologizados; es inmediato que este producto
posee las propiedades usuales de asociatividad y de conmutatividad.
\end{env}

\begin{env}[7.7.5]
\label{0.7.7.5-fr}
Si $B$, $C$ son dos $A$-álgebras topológicas linealmente topologizadas,
la topología producto tensorial sobre $B\otimes_A C$ tiene como sistema fundamental
de vecindades de $0$ los \emph{ideales}
$\Im(\mathfrak K\otimes_A C)+\Im(B\otimes_A\mathfrak L)$ del álgebra
$B\otimes_A C$, recorriendo $\mathfrak K$ (resp. $\mathfrak L$) el conjunto
de los ideales abiertos de $B$ (resp. $C$). Por consiguiente, $(B\otimes_A C)^\wedge$
está provisto de una estructura de \emph{$\widehat A$-álgebra topológica}, límite
proyectivo del sistema proyectivo de $(A/\mathfrak J)$-álgebras
$(B/\mathfrak K)\otimes_{A/\mathfrak J}(C/\mathfrak L)$ ($\mathfrak J$
ideal abierto de $A$ tal que $\mathfrak J.B\subset\mathfrak K$,
$\mathfrak J.C\subset\mathfrak L$; siempre existe uno). Se dice que
esta álgebra es el \emph{producto tensorial completado} de las álgebras $B$ y
$C$.
\end{env}

\begin{env}[7.7.6]
\label{0.7.7.6-fr}
Los homomorfismos de $A$-álgebras $b\to b\otimes 1$,
$c\to 1\otimes c$ de $B$ y $C$ en $B\otimes_A C$ son continuos
cuando se dota a esta última álgebra de la topología producto tensorial;
por composición con el homomorfismo canónico de $B\otimes_A C$ en su
separado completado, dan, pues, \emph{homomorfismos canónicos}
$\rho:B\to(B\otimes_A C)^\wedge$, $\sigma:C\to(B\otimes_A C)^\wedge$.
El álgebra $(B\otimes_A C)^\wedge$ y los homomorfismos $\rho$ y $\sigma$
poseen además la \emph{propiedad universal} siguiente: para toda
$A$-álgebra linealmente topologizada, separada y completa $D$, y todo
par de $A$-homomorfismos continuos $u:B\to D$, $v:C\to D$, existe un
$A$-homomorfismo continuo y uno solo
$w:(B\otimes_A C)^\wedge\to D$ tal que $u=w\circ\rho$ y
$v=w\circ\sigma$. En efecto, existe ya un $A$-homomorfismo único
$w_0:B\otimes_A C\to D$ tal que $u(b)=w_0(b\otimes 1)$ y
$v(c)=w_0(1\otimes c)$, y todo se reduce a probar que $w_0$ es \emph{continuo},
pues dará entonces un

\oldpage[0\textsubscript{I}]{77}

homomorfismo continuo $w$ por paso al separado completado. Ahora bien, si
$\mathfrak M$ es un ideal abierto de $D$, existen por hipótesis ideales
abiertos $\mathfrak K\subset B$, $\mathfrak L\subset C$ tales que
$u(\mathfrak K)\subset\mathfrak M$, $v(\mathfrak L)\subset\mathfrak M$;
la imagen por $w_0$ de
$\Im(\mathfrak K\otimes C)+\Im(B\otimes\mathfrak L)$ está también
contenida en $\mathfrak M$, de donde nuestra afirmación.
\end{env}

\begin{proposition}[7.7.7]
\label{0.7.7.7-fr}
Si $B$ y $C$ son dos $A$-álgebras preadmisibles,
$(B\otimes_A C)^\wedge$ es admisible, y si $\mathfrak K$ \textup{(resp.
$\mathfrak L$)} es un ideal de definición de $B$ \textup{(resp. $C$)}, la adherencia
en $(B\otimes_A C)^\wedge$ de la imagen canónica de
$\mathfrak H=\Im(\mathfrak K\otimes C)+\Im(B\otimes\mathfrak L)$ es
un ideal de definición.
\end{proposition}

Basta mostrar que $\mathfrak H^n$ tiende a $0$ para la topología
producto tensorial, lo que resulta inmediatamente de la inclusión
\[
  \mathfrak H^{2n}\subset
  \Im(\mathfrak K^n\otimes C)+\Im(B\otimes\mathfrak L^n)
\]

\begin{proposition}[7.7.8]
\label{0.7.7.8-fr}
Sean $A$ un anillo preádico, $\mathfrak J$ un ideal de definición de
$A$, $M$ un $A$-módulo \emph{de tipo finito}, provisto de la topología
$\mathfrak J$-preádica. Para toda $A$-álgebra topológica ádica y
noetheriana $B$, $B\otimes_A M$ se identifica con el producto tensorial completado
$(B\otimes_A M)^\wedge$.
\end{proposition}

Si $\mathfrak K$ es un ideal de definición de $B$, existe por hipótesis
un entero $m$ tal que $\mathfrak J^mB\subset\mathfrak K$, por tanto
$\Im(B\otimes_A\mathfrak J^{nm}M)=
\Im(\mathfrak J^{nm}B\otimes_A M)\subset
\Im(\mathfrak K^nB\otimes_A M)=\mathfrak K^n(B\otimes_A M)$; se
concluye que sobre $B\otimes_A M$, el producto tensorial de las topologías de $B$
y de $M$ es la topología $\mathfrak K$-preádica. Como $B\otimes_A M$
es un $B$-módulo de tipo finito, la proposición resulta inmediatamente de (7.3.6).


% END tree/worktree/es/CURRENT/EGA_I_ch0/ega0_i_7_7.tex
% BEGIN tree/worktree/es/CURRENT/EGA_I_ch0/ega0_i_7_8.tex
\subsection{Topologías sobre los módulos de homomorfismos}
\label{subsection:0.7.8-fr}

\begin{env}[7.8.1]
\label{0.7.8.1-fr}
Sean $A$ un anillo \emph{$\mathfrak J$-ádico noetheriano}, $M$ y $N$ dos
$A$-módulos de tipo finito, provistos de la topología $\mathfrak J$-preádica;
se sabe (7.3.6) que son separados y completos; además, todo
$A$-homomorfismo $M\to N$ es automáticamente continuo, y el $A$-módulo
$\operatorname{Hom}_A(M,N)$ es de tipo finito. Para todo entero $i\geq 0$, pongamos
$A_i=A/\mathfrak J^{i+1}$, $M_i=M/\mathfrak J^{i+1}M$,
$N_i=N/\mathfrak J^{i+1}N$; para $i\leq j$, todo homomorfismo
$u_j:M_j\to N_j$ aplica $\mathfrak J^{i+1}M_j$ en
$\mathfrak J^{i+1}N_j$, por tanto da por paso a los cocientes un
homomorfismo $u_i:M_i\to N_i$, lo que define un homomorfismo canónico
$\operatorname{Hom}_{A_j}(M_j,N_j)\to\operatorname{Hom}_{A_i}(M_i,N_i)$; además, los
$\operatorname{Hom}_{A_i}(M_i,N_i)$ constituyen un \emph{sistema proyectivo} para estos
homomorfismos, y resulta de (7.2.10) que hay un isomorfismo
canónico
\[
  \varphi:\operatorname{Hom}_A(M,N)\longrightarrow
  \varprojlim_i\operatorname{Hom}_{A_i}(M_i,N_i).
\]
Además:
\end{env}

\begin{proposition}[7.8.2]
\label{0.7.8.2-fr}
Si $M$ y $N$ son módulos de tipo finito sobre un anillo $\mathfrak J$-ádico
noetheriano $A$, los submódulos $\operatorname{Hom}_A(M,\mathfrak J^{i+1}N)$ forman un
sistema fundamental de vecindades de $0$ en $\operatorname{Hom}_A(M,N)$ para la
topología $\mathfrak J$-ádica, y el homomorfismo canónico
\[
  \varphi:\operatorname{Hom}_A(M,N)\longrightarrow
  \varprojlim_i\operatorname{Hom}_{A_i}(M_i,N_i)
\]
es un isomorfismo topológico.
\end{proposition}

Se puede, en efecto, considerar $M$ como cociente de un $A$-módulo libre $L$ de
tipo finito, y, por consiguiente, identificar $\operatorname{Hom}_A(M,N)$ con un submódulo de
$\operatorname{Hom}_A(L,N)$; en esta identificación,
$\operatorname{Hom}_A(M,\mathfrak J^{i+1}N)$ es la intersección de $\operatorname{Hom}_A(M,N)$ y de
$\operatorname{Hom}_A(L,\mathfrak J^{i+1}N)$; como la topología inducida sobre
$\operatorname{Hom}_A(M,N)$ por la topología $\mathfrak J$-ádica de $\operatorname{Hom}_A(L,N)$ es
la topología $\mathfrak J$-ádica (7.3.2), se reduce a demostrar la
primera afirmación

\oldpage[0\textsubscript{I}]{78}

para $M=L=A^m$; pero entonces $\operatorname{Hom}_A(L,N)=N^m$,
$\operatorname{Hom}_A(L,\mathfrak J^{i+1}N)=(\mathfrak J^{i+1}N)^m=
\mathfrak J^{i+1}N^m$ y el resultado es evidente. Para establecer la segunda
afirmación, observemos que la imagen de $\operatorname{Hom}_A(M,\mathfrak J^{i+1}N)$ en
$\operatorname{Hom}_{A_j}(M_j,N_j)$ es nula para $j\leq i$, por tanto $\varphi$ es
continuo; recíprocamente, la imagen inversa en $\operatorname{Hom}_A(M,N)$ del $0$ de
$\operatorname{Hom}_{A_i}(M_i,N_i)$ es $\operatorname{Hom}_A(M,\mathfrak J^{i+1}N)$, por tanto $\varphi$
es bicontinuo.

Si se supone solamente que $A$ es un anillo \emph{$\mathfrak J$-preádico
noetheriano}, $M$ y $N$ dos $A$-módulos de tipo finito, separados para la
topología $\mathfrak J$-preádica, la demostración anterior muestra que
la primera afirmación de (7.8.2) sigue siendo válida, y que $\varphi$ es un
isomorfismo topológico de $\operatorname{Hom}_A(M,N)$ sobre un submódulo de
$\varprojlim_i\operatorname{Hom}_{A_i}(M_i,N_i)$.

\begin{proposition}[7.8.3]
\label{0.7.8.3-fr}
Bajo las hipótesis de \textup{(7.8.2)}, el conjunto de los homomorfismos inyectivos
\textup{(resp.} sobreyectivos, biyectivos) de $M$ en $N$ es una parte abierta de
$\operatorname{Hom}_A(M,N)$.
\end{proposition}

En efecto, en virtud de (7.3.5) y (7.1.14), para que $u$ sea sobreyectivo, es
necesario y suficiente que el homomorfismo correspondiente
$u_0:M/\mathfrak JM\to N/\mathfrak JN$ lo sea, y el conjunto de los
homomorfismos sobreyectivos de $M$ en $N$ es, por tanto, imagen inversa por
la aplicación continua $\operatorname{Hom}_A(M,N)\to\operatorname{Hom}_{A_0}(M_0,N_0)$ de una parte
de un espacio discreto. Mostremos ahora que el conjunto de los homomorfismos
inyectivos es abierto; sea $v$ tal homomorfismo y pongamos $M'=v(M)$;
por el lema de Artin-Rees (7.3.2.1), existe un entero $k\geq 0$ tal que
$M'\cap\mathfrak J^{m+k}N\subset\mathfrak J^mM'$ para todo $m>0$; vamos
a ver que para $w\in\mathfrak J^{k+1}\operatorname{Hom}_A(M,N)$,
$u=v+w$ es inyectivo, lo que terminará la demostración. En efecto, sea
$x\in M$ tal que $u(x)=0$; probemos que para todo $i\geq 0$ la relación
$x\in\mathfrak J^iM$ implica $x\in\mathfrak J^{i+1}M$; resultará de ello
que
\[
  x\in\bigcap_{i\geq 0}\mathfrak J^iM=(0).
\]
En efecto, se tiene entonces $w(x)\in\mathfrak J^{i+k+1}N$, y, por otra parte,
$w(x)=-v(x)\in M'$, por tanto
$v(x)\in M'\cap\mathfrak J^{i+k+1}N\subset\mathfrak J^{i+1}M'$, y como
$v$ es un isomorfismo de $M$ sobre $M'$, $x\in\mathfrak J^{i+1}M$;
c.q.d.

\begin{flushright}
\emph{(Continuará.)}
\end{flushright}

% END tree/worktree/es/CURRENT/EGA_I_ch0/ega0_i_7_8.tex
% BEGIN tree/worktree/es/CURRENT/EGA_I_ch0/R586_SOURCE_NOTE_es.tex
\paragraph{Nota sobre la lectura de 3.1.6.}
La edición corriente aplica el erratum publicado en EGA II, \emph{Errata et Addenda}, Lista 1, página 217: en la frase que abrevia $u_V(s)$ como $u(s)$, se sustituye $\Gamma(U,\mathcal{F})$ por $\Gamma(V,\mathcal{F})$. Así el dominio de la sección coincide con el componente aplicado. La lectura diplomática conserva $\Gamma(U,\mathcal{F})$. Las apariciones anteriores de $U$ y el componente $u_V$ se mantienen. Esta es una corrección publicada, ya admitida por el control canónico D186; no se presenta como un hallazgo nuevo.

% END tree/worktree/es/CURRENT/EGA_I_ch0/R586_SOURCE_NOTE_es.tex
% BEGIN tree/worktree/es/CURRENT/EGA_I_ch0/R587_SOURCE_NOTE_es.tex
\paragraph{Nota sobre el recubrimiento finito de 3.2.1.}
La nota impresa en la página 26 afirma que la construcción mediante recubrimientos finitos es independiente del recubrimiento para cualquier prehaz sobre la base. Esto falla sin una condición de pegado. En el espacio discreto $X=\{a,b\}$, con todos sus abiertos como base, tomemos $\mathcal{F}(X)=\{0,1\}$ y $\mathcal{F}(W)=\{0\}$ para cada abierto propio $W$, con las restricciones únicas. El recubrimiento formado por $X$ da un límite de dos elementos; el recubrimiento por $\{a\}$ y $\{b\}$ da uno de un elemento. La edición corriente añade a esa nota la condición $(F_0)$ de (3.2.2), que asegura el pegado necesario para la independencia. La lectura diplomática conserva la afirmación impresa completa. Se trata de una enmienda editorial propuesta aquí, sin atribución a un erratum publicado.

% END tree/worktree/es/CURRENT/EGA_I_ch0/R587_SOURCE_NOTE_es.tex
% BEGIN tree/worktree/es/CURRENT/EGA_I_ch0/R588_SOURCE_NOTE_es.tex
\paragraph{Notas de fuente sobre 3.2.2.}
En la nota de la página 26, la compatibilidad impresa no determina por sí sola las componentes $\rho_V''$: incluso para el haz de dos elementos sobre un espacio de un punto, sus cuatro endomorfismos satisfacen esa compatibilidad. La edición corriente especifica las componentes canónicas $\rho_V''=\rho_V^{U_\alpha}$ y conserva la identidad de compatibilidad. Esta aclaración es una propuesta editorial local; no modifica el criterio de haz.

En la página 27, la familia subordinada al recubrimiento de $U$ es una base de $U$. Se aplica la corrección ya aceptada EGA1-CANON-ERR-0001, propagada en D193, que sustituye $X$ por $U$ solo en esa afirmación. La prueba se entiende restringiendo primero el prehaz a los abiertos básicos contenidos en $U$ y aplicando allí el primer paso. El primer paso sigue teniendo $X$ como espacio ambiente. La lectura diplomática conserva ambos textos impresos; no se atribuye esta corrección a un erratum publicado.

% END tree/worktree/es/CURRENT/EGA_I_ch0/R588_SOURCE_NOTE_es.tex
% BEGIN tree/worktree/es/CURRENT/EGA_I_ch0/R590_SOURCE_NOTE_es.tex
% R590 local editorial clarification; no accepted Canon disposition.
La edición actual precisa «subsistema proyectivo»: las inclusiones de los subhaces son compatibles con todos los morfismos de transición. Es una aclaración editorial local de la hipótesis implícita; la edición diplomática conserva la redacción impresa. No se afirma aceptación canónica ni revisión humana.

% END tree/worktree/es/CURRENT/EGA_I_ch0/R590_SOURCE_NOTE_es.tex
\chapter*{CAPÍTULO I\\EL LENGUAJE DE LOS ESQUEMAS}\addcontentsline{toc}{chapter}{Capítulo I: El lenguaje de los esquemas}
\section*{Sumario}

\begin{longtable}{ll}
  \hyperref[section:I.1]{\textsection1}.   & Esquemas afines.\\
  \hyperref[section:I.2]{\textsection2}.   & Preesquemas y morfismos de preesquemas.\\
  \hyperref[section:I.3]{\textsection3}.   & Productos de preesquemas.\\
  \hyperref[section:I.4]{\textsection4}.   & Subpreesquemas y morfismos de inmersión.\\
  \hyperref[section:I.5]{\textsection5}.   & Preesquemas reducidos; condición de separación.\\
  \hyperref[section:I.6]{\textsection6}.   & Condiciones de finitud.\\
  \hyperref[section:I.7]{\textsection7}.   & Aplicaciones racionales.\\
  \hyperref[section:I.8]{\textsection8}.   & Los esquemas de Chevalley.\\
  \hyperref[section:I.9]{\textsection9}.   & Complementos sobre los haces cuasicoherentes.\\
  \hyperref[section:I.10]{\textsection10}. & Esquemas formales.
\end{longtable}
\bigskip

\oldpage[I]{79}
\phantomsection\label{chapter:I-fr}
En los \textsection\textsection\hyperref[section:I.1]{1}--\hyperref[section:I.8]{8} apenas hacemos más que desarrollar un lenguaje, el que se utilizará en todo lo que sigue.
Observemos, no obstante, que, conforme al espíritu general de este tratado, los \textsection\textsection\hyperref[section:I.7]{7}--\hyperref[section:I.8]{8} se utilizarán menos que los demás y de manera menos esencial; solo se habla de los esquemas de Chevalley para enlazar con el lenguaje de Chevalley \cite{I-1} y Nagata \cite{I-9}.
En el \textsection\hyperref[section:I.9]{9} damos definiciones y resultados sobre los haces cuasicoherentes, algunos de los cuales ya no se limitan a traducir al lenguaje ``geométrico'' nociones conocidas de álgebra conmutativa, sino que son ya de naturaleza global; serán indispensables, desde los capítulos siguientes, para el estudio global de los morfismos.
Finalmente, en el \textsection\hyperref[section:I.10]{10} introducimos una generalización de la noción de esquema, que nos servirá de intermediaria en el Capítulo~III para formular y demostrar cómodamente los resultados fundamentales del estudio cohomológico de los morfismos propios;
por otra parte, señalemos que la noción de esquema formal parece indispensable para expresar ciertos hechos de la ``teoría de módulos'' (problemas de clasificación de variedades algebraicas).
Los resultados del \textsection\hyperref[section:I.10]{10} no se utilizarán antes del \textsection3 del Capítulo~III, y se recomienda omitir su lectura hasta entonces.
\bigskip

% BEGIN source_es/ega1/ega1-1.tex
\setcounter{section}{0}
\section{Esquemas afines}
\label{section:I.1}

\setcounter{subsection}{0}
\subsection{El espectro primo de un anillo}
\label{subsection:I.1.1}


\oldpage[I]{80}
\begin{env}[1.1.1]
\label{I.1.1.1}

\emph{Notación}. Sean $A$ un anillo (conmutativo) y $M$ un $A$-módulo.
En este capítulo y en los siguientes utilizaremos constantemente las notaciones siguientes:
\begin{itemize}
  \item $\operatorname{Spec}(A)={}$ \emph{conjunto de los ideales primos} de $A$, llamado también \emph{espectro primo} de $A$; para $x\in X=\operatorname{Spec}(A)$, será a menudo conveniente escribir $\mathfrak{j}_x$ en lugar de $x$.
    Para que $\operatorname{Spec}(A)$ sea \emph{vacío}, es necesario y suficiente que $A$ sea el anillo $0$.
  \item $A_x=A_{\mathfrak{j}_x}={}$ \emph{anillo (local) de fracciones $S^{-1}A$}, donde $S=A-\mathfrak{j}_x$.
  \item $\mathfrak{m}_x=\mathfrak{j}_xA_{\mathfrak{j}_x}={}$ \emph{ideal maximal de $A_x$}.
  \item $\boldsymbol{k}(x)=A_x/\mathfrak{m}_x={}$ \emph{campo residual de $A_x$}, canónicamente isomorfo al campo de fracciones del dominio entero $A/\mathfrak{j}_x$, con el cual se lo identifica.
  \item $f(x)={}$ \emph{clase de $f$ mód.~$\mathfrak{j}_x$} en $A/\mathfrak{j}_x\subset\boldsymbol{k}(x)$, para $f\in A$ y $x\in X$.
    Se dice también que $f(x)$ es el \emph{valor} de $f$ en un punto $x\in\operatorname{Spec}(A)$; las relaciones $f(x)=0$ y $f\in\mathfrak{j}_x$ son \emph{equivalentes}.
  \item $M_x=M\otimes_AA_x={}$ \emph{módulo de fracciones con denominadores en $A-\mathfrak{j}_x$}.
  \item $\mathfrak{r}(E)={}$ \emph{radical del ideal de $A$ engendrado por un subconjunto $E$ de $A$}.
  \item $V(E)={}$ \emph{conjunto de los $x\in X$ tales que $E\subset\mathfrak{j}_x$} (o conjunto de los $x\in X$ tales que $f(x)=0$ para todo $f\in E$), para $E\subset A$.
    Se tiene entonces
    \begin{equation}
  \mathfrak{r}(E)=\bigcap_{x\in V(E)}\mathfrak{j}_x.
  \tag{1.1.1.1}\label{I.1.1.1.1-fr}
\end{equation}
  \item $V(f)=V(\{f\})$ para $f\in A$.
  \item $D(f)=X-V(f)={}$ \emph{conjunto de los $x\in X$ en los que $f(x)\neq0$}.
\end{itemize}
\end{env}

\begin{proposition}[1.1.2]
\label{I.1.1.2}
Se tienen las propiedades siguientes:
\begin{enumerate}
  \item[{\rm(i)}] $V(0)=X$, $V(1)=\varnothing$.
  \item[{\rm(ii)}] La relación $E\subset E'$ implica $V(E)\supset V(E')$.
  \item[{\rm(iii)}] Para toda familia $(E_\lambda)$ de subconjuntos de $A$, \[
    V\left(\bigcup_\lambda E_\lambda\right)
      =V\left(\sum_\lambda E_\lambda\right)
      =\bigcap_\lambda V(E_\lambda).
  \].
  \item[{\rm(iv)}] $V(EE')=V(E)\cup V(E')$.
  \item[{\rm(v)}] $V(E)=V(\mathfrak{r}(E))$.
\end{enumerate}
\end{proposition}

\begin{proof}
Las propiedades (i), (ii) y (iii) son triviales, y (v) resulta de (ii) y de la fórmula \sref{I.1.1.1.1}.
Es evidente que $V(EE')\supset V(E)\cup V(E')$;
recíprocamente, si $x\notin V(E)$ y $x\notin V(E')$, existen $f\in E$ y $f'\in E'$ tales que $f(x)\neq0$ y $f'(x)\neq0$ en $\boldsymbol{k}(x)$, de donde $f(x)f'(x)\neq0$, es decir, $x\notin V(EE')$, lo que demuestra (iv).
\end{proof}

La Proposición \sref{I.1.1.2} muestra, entre otras cosas, que los conjuntos de la forma $V(E)$ (donde $E$ recorre los subconjuntos de $A$) son los \emph{conjuntos cerrados} de una topología en $X$, que llamaremos \emph{topología espectral}
\footnote{La introducción de esta topología en geometría algebraica se debe a Zariski.
Por ello esta topología se denomina habitualmente ``topología de Zariski'' de $X$.};
salvo mención expresa en contrario, se supondrá siempre que $X=\operatorname{Spec}(A)$ está provisto de la topología espectral.


\oldpage[I]{81}
\begin{env}[1.1.3]
\label{I.1.1.3}

Para cada subconjunto $Y$ de $X$, denotaremos por $\mathfrak{j}(Y)$ el conjunto de los $f\in A$ tales que $f(y)=0$ para todo $y\in Y$;
equivalentemente, $\mathfrak{j}(Y)$ es la intersección de los ideales primos $\mathfrak{j}_y$ para $y\in Y$.
Es claro que la relación $Y\subset Y'$ implica $\mathfrak{j}(Y)\supset\mathfrak{j}(Y')$ y que se tiene
\begin{equation}
  \mathfrak{j}\left(\bigcup_\lambda Y_\lambda\right)
    =\bigcap_\lambda\mathfrak{j}(Y_\lambda)
  \tag{1.1.3.1}\label{I.1.1.3.1-fr}
\end{equation}
para toda familia $(Y_\lambda)$ de subconjuntos de $X$.
Finalmente se tiene
\begin{equation}
  \mathfrak{j}(\{x\})=\mathfrak{j}_x.
  \tag{1.1.3.2}\label{I.1.1.3.2-fr}
\end{equation}
\end{env}

\begin{proposition}[1.1.4]
\label{I.1.1.4}
\medskip\noindent
\begin{enumerate}
  \item[{\rm(i)}] Para todo subconjunto $E$ de $A$, se tiene $\mathfrak{j}(V(E))=\mathfrak{r}(E)$.
  \item[{\rm(ii)}] Para todo subconjunto $Y$ de $X$, $V(\mathfrak{j}(Y))=\overline{Y}$, la clausura de $Y$ en $X$.
\end{enumerate}
\end{proposition}

\begin{proof}
(i) es una consecuencia inmediata de las definiciones y de \sref{I.1.1.1.1};
por otra parte, $V(\mathfrak{j}(Y))$ es cerrado y contiene a $Y$;
recíprocamente, si $Y\subset V(E)$, se tiene $f(y)=0$ para $f\in E$ y todo $y\in Y$, luego $E\subset\mathfrak{j}(Y)$, $V(E)\supset V(\mathfrak{j}(Y))$, lo que demuestra (ii).
\end{proof}

\begin{corollary}[1.1.5]
\label{I.1.1.5}
Los subconjuntos cerrados de $X=\operatorname{Spec}(A)$ y los ideales de $A$ iguales a sus radicales (dicho de otro modo, los que son intersecciones de ideales primos) se corresponden biyectivamente mediante las aplicaciones que invierten las inclusiones $Y\to\mathfrak{j}(Y)$, $\mathfrak{a}\to V(\mathfrak{a})$;
a la unión $Y_1\cup Y_2$ de dos subconjuntos cerrados corresponde $\mathfrak{j}(Y_1)\cap\mathfrak{j}(Y_2)$, y a la intersección de una familia cualquiera $(Y_\lambda)$ de subconjuntos cerrados corresponde el radical de la suma de los $\mathfrak{j}(Y_\lambda)$.
\end{corollary}

\begin{corollary}[1.1.6]
\label{I.1.1.6}
Si $A$ es un anillo noetheriano, $X=\operatorname{Spec}(A)$ es un espacio noetheriano.
\end{corollary}

Se observará que el recíproco de este corolario es inexacto, como muestra el ejemplo de un anillo íntegro no noetheriano que tiene un solo ideal primo $\neq\{0\}$, por ejemplo un anillo de valoración no discreta de rango 1.

Como ejemplo de un anillo $A$ cuyo espectro no es un espacio noetheriano, puede considerarse el anillo $\mathcal{C}(Y)$ de las funciones reales continuas sobre un espacio compacto infinito $Y$;
se sabe que, como conjunto, $Y$ se identifica con el conjunto de los ideales maximales de $A$, y es fácil ver que la topología inducida sobre $Y$ por la de $X=\operatorname{Spec}(A)$ es la topología original de $Y$.
Como $Y$ no es un espacio noetheriano, lo mismo ocurre con $X$.

\begin{corollary}[1.1.7]
\label{I.1.1.7}
Para todo $x\in X$, la clausura de $\{x\}$ es el conjunto de los $y\in X$ tales que $\mathfrak{j}_x\subset\mathfrak{j}_y$.
Para que $\{x\}$ sea cerrado, es necesario y suficiente que $\mathfrak{j}_x$ sea maximal.
\end{corollary}

\begin{corollary}[1.1.8]
\label{I.1.1.8}
El espacio $X=\operatorname{Spec}(A)$ es un espacio de Kolmogórov.
\end{corollary}

\begin{proof}
Si $x$ e $y$ son dos puntos distintos de $X$, se tiene o bien $\mathfrak{j}_x\not\subset\mathfrak{j}_y$ o bien $\mathfrak{j}_y\not\subset\mathfrak{j}_x$, por lo que uno de los puntos $x$, $y$ no pertenece a la clausura del otro.
\end{proof}

\begin{env}[1.1.9]
\label{I.1.1.9}
Según la Proposición \sref{I.1.1.2}[iv], para dos elementos $f$, $g$ de $A$, se tiene
\begin{equation}
  D(fg)=D(f)\cap D(g)
  \tag{1.1.9.1}\label{I.1.1.9.1-fr}
\end{equation}
Notemos asimismo que la igualdad $D(f)=D(g)$ significa, según la Proposición~\sref{I.1.1.4}[i] y la Proposición~\sref{I.1.1.2}[v], que $\mathfrak{r}(f)=\mathfrak{r}(g)$, o que los ideales primos minimales que contienen $(f)$ y $(g)$ son los mismos;
en particular, esto también ocurre cuando $f=ug$, donde $u$ es invertible.
\end{env}


\oldpage[I]{82}
\begin{proposition}[1.1.10]
\label{I.1.1.10}
\medskip\noindent

\begin{enumerate}
  \item[{\rm(i)}] Cuando $f$ recorre $A$, los conjuntos $D(f)$ forman una base de la topología de $X$.
  \item[{\rm(ii)}] Para todo $f\in A$, $D(f)$ es cuasicompacto.
    En particular, $X=D(1)$ es cuasicompacto.
\end{enumerate}
\end{proposition}

\begin{proof}
\medskip\noindent
\begin{enumerate}
  \item[(i)] Sea $U$ un abierto de $X$;
    por definición, se tiene $U=X-V(E)$, donde $E$ es un subconjunto de $A$, y $V(E)=\bigcap_{f\in E}V(f)$, de donde $U=\bigcup_{f\in E}D(f)$.
  \item[(ii)] Por (i), basta demostrar que, si $(f_\lambda)_{\lambda\in L}$ es una familia de elementos de $A$ tal que $D(f)\subset\bigcup_{\lambda\in L}D(f_\lambda)$, entonces existe un subconjunto finito $J$ de $L$ tal que $D(f)\subset\bigcup_{\lambda\in J}D(f_\lambda)$.
    Sea $\mathfrak{a}$ el ideal de $A$ engendrado por los $f_\lambda$;
    por hipótesis se tiene $V(f)\supset V(\mathfrak{a})$, luego $\mathfrak{r}(f)\subset\mathfrak{r}(\mathfrak{a})$;
    como $f\in\mathfrak{r}(f)$, existe un entero $n\geq0$ tal que $f^n\in\mathfrak{a}$.
    Pero entonces $f^n$ pertenece al ideal $\mathfrak{b}$ engendrado por la subfamilia finita $(f_\lambda)_{\lambda\in J}$, y se tiene $V(f)=V(f^n)\supset V(\mathfrak{b})=\bigcap_{\lambda\in J}V(f_\lambda)$, es decir, $D(f)\subset\bigcup_{\lambda\in J}D(f_\lambda)$.
\end{enumerate}
\end{proof}

\begin{proposition}[1.1.11]
\label{I.1.1.11}
Para todo ideal $\mathfrak{a}$ de $A$, $\operatorname{Spec}(A/\mathfrak{a})$ se identifica canónicamente con el subespacio cerrado $V(\mathfrak{a})$ de $\operatorname{Spec}(A)$.
\end{proposition}

\begin{proof}
Se sabe que existe una correspondencia biyectiva canónica (que respeta la estructura de orden dada por la inclusión) entre los ideales (resp. los ideales primos) de $A/\mathfrak{a}$ y los ideales (resp. los ideales primos) de $A$ que contienen a $\mathfrak{a}$.
\end{proof}

Recordemos que el conjunto $\mathfrak{N}$ de los elementos nilpotentes de $A$ (el \emph{nilradical} de $A$) es un ideal igual a $\mathfrak{r}(0)$, la intersección de todos los ideales primos de $A$ \sref[0]{0.1.1.1}.

\begin{corollary}[1.1.12]
\label{I.1.1.12}
Los espacios topológicos $\operatorname{Spec}(A)$ y $\operatorname{Spec}(A/\mathfrak{N})$ son canónicamente homeomorfos.
\end{corollary}

\begin{proposition}[1.1.13]
\label{I.1.1.13}
Para que $X=\operatorname{Spec}(A)$ sea irreducible \sref[0]{0.2.1.1}, es necesario y suficiente que el anillo $A/\mathfrak{N}$ sea un dominio entero \textup{(o, equivalentemente, que el ideal $\mathfrak{N}$ sea \emph{primo})}.
\end{proposition}

\begin{proof}
En virtud del Corolario \sref{I.1.1.12}, se puede limitar al caso en que $\mathfrak{N}=0$.
Si $X$ es reducible, existen dos subconjuntos cerrados propios $Y_1$ e $Y_2$ de $X$ tales que $X=Y_1\cup Y_2$, de donde $\mathfrak{j}(X)=\mathfrak{j}(Y_1)\cap\mathfrak{j}(Y_2)=0$, pues los ideales $\mathfrak{j}(Y_1)$ y $\mathfrak{j}(Y_2)$ son distintos de $(0)$ \sref{I.1.1.5};
por tanto, $A$ no es un dominio entero.
Recíprocamente, si existen elementos $f\neq0$, $g\neq0$ de $A$ tales que $fg=0$, se tiene $V(f)\neq X$, $V(g)\neq X$ (pues la intersección de todos los ideales primos de $A$ es $(0)$), y $X=V(fg)=V(f)\cup V(g)$.
\end{proof}

\begin{corollary}[1.1.14]
\label{I.1.1.14}
\medskip\noindent
\begin{enumerate}
  \item[{\rm(i)}] En la correspondencia biyectiva entre los subconjuntos cerrados de $X=\operatorname{Spec}(A)$ y los ideales de $A$ iguales a sus radicales, los subconjuntos cerrados irreducibles de $X$ corresponden a los ideales primos de $A$.
    En particular, las componentes irreducibles de $X$ corresponden a los ideales primos minimales de $A$.
  \item[{\rm(ii)}] La aplicación $x\to\overline{\{x\}}$ establece una correspondencia biyectiva entre $X$ y el conjunto de los subconjuntos cerrados irreducibles de $X$ \textup{(es decir, \emph{todo subconjunto cerrado irreducible de $X$ admite exactamente un punto genérico})}.
\end{enumerate}
\end{corollary}

\begin{proof}
(i) resulta inmediatamente de \sref{I.1.1.13} y \sref{I.1.1.11};
y, para demostrar (ii), se puede, en virtud de \sref{I.1.1.11}, limitar al caso en que $X$ es irreducible;
entonces, según la Proposición~\sref{I.1.1.13}, existe un ideal primo $\mathfrak{N}$ de $A$ contenido en todos los demás, que corresponde al punto genérico
\oldpage[I]{83}
de $X$;
además, $X$ admite a lo sumo un punto genérico, puesto que es un espacio de Kolmogórov (\sref{I.1.1.8} y \sref[0]{0.2.1.3}).
\end{proof}

\begin{proposition}[1.1.15]
\label{I.1.1.15}
Si $\mathfrak{J}$ es un ideal de $A$ contenido en el radical de Jacobson $\mathfrak{R}(A)$, la única
vecindad de $V(\mathfrak{J})$ en $X=\operatorname{Spec}(A)$ es todo el espacio $X$.
\end{proposition}

\begin{proof}
Por definición, todo ideal maximal de $A$ pertenece a $V(\mathfrak{J})$.
Como todo ideal $\mathfrak{a}\neq A$ de $A$ está contenido en un ideal maximal, se tiene $V(\mathfrak{a})\cap V(\mathfrak{J})\neq\varnothing$, de donde resulta la proposición.
\end{proof}

\subsection{Propiedades funtoriales de los espectros primos de los anillos}
\label{subsection:I.1.2}

\begin{env}[1.2.1]
\label{I.1.2.1}
Sean $A$, $A'$ dos anillos, y
\[
  \varphi:A'\longrightarrow A
\]
un homomorfismo de anillos.
Para todo ideal primo $x=\mathfrak{j}_x\in\operatorname{Spec}(A)=X$, el anillo $A'/\varphi^{-1}(\mathfrak{j}_x)$ es canónicamente isomorfo a un subanillo de $A/\mathfrak{j}_x$, y por tanto es un dominio entero, o, dicho de otro modo, $\varphi^{-1}(\mathfrak{j}_x)$ es un ideal primo de $A'$;
lo denotaremos por ${}^a\varphi(x)$, y hemos definido así una aplicación
\[
  {}^a\varphi:X=\operatorname{Spec}(A)\longrightarrow
  X'=\operatorname{Spec}(A')
\]
(denotada también por $\operatorname{Spec}(\varphi)$), que llamaremos la aplicación \emph{asociada} al homomorfismo $\varphi$.
Denotaremos por $\varphi^x$ el homomorfismo inyectivo de $A'/\varphi^{-1}(\mathfrak{j}_x)$ en $A/\mathfrak{j}_x$ inducido por $\varphi$ al pasar a los cocientes, así como su prolongación canónica a un monomorfismo de campos
\[
  \varphi^x:\boldsymbol{k}({}^a\varphi(x))\longrightarrow\boldsymbol{k}(x)~;
\]
para todo $f'\in A'$, se tiene, por definición,
\begin{equation}
  \varphi^x\bigl(f'({}^a\varphi(x))\bigr)=\bigl(\varphi(f')\bigr)(x)
  \qquad(x\in X).
  \tag{1.2.1.1}\label{I.1.2.1.1-fr}
\end{equation}
\end{env}

\begin{proposition}[1.2.2]
\label{I.1.2.2}
\medskip\noindent
\begin{enumerate}
  \item[{\rm(i)}] Para todo subconjunto $E'$ de $A'$, se tiene
    \begin{equation}
    {}^a\varphi^{-1}(V(E'))=V(\varphi(E'))
    \tag{1.2.2.1}\label{I.1.2.2.1-fr}
  \end{equation}
    y, en particular, para todo $f'\in A'$
    \begin{equation}
    {}^a\varphi^{-1}(D(f'))=D(\varphi(f')).
    \tag{1.2.2.2}\label{I.1.2.2.2-fr}
  \end{equation}
  \item[{\rm(ii)}] Para todo ideal $\mathfrak{a}$ de $A$, se tiene
    \begin{equation}
    \overline{{}^a\varphi(V(\mathfrak{a}))}
    =V(\varphi^{-1}(\mathfrak{a})).
    \tag{1.2.2.3}\label{I.1.2.2.3-fr}
  \end{equation}
\end{enumerate}
\end{proposition}

\begin{proof}
La relación ${}^a\varphi(x)\in V(E')$ equivale, por definición, a $E'\subset\varphi^{-1}(\mathfrak{j}_x)$, por tanto a $\varphi(E')\subset\mathfrak{j}_x$, y finalmente a $x\in V(\varphi(E'))$, de donde (i).
Para demostrar (ii), se puede suponer que $\mathfrak{a}$ es igual a su radical, puesto que $V(\mathfrak{r}(\mathfrak{a}))=V(\mathfrak{a})$ \sref{I.1.1.2}[(v)] y $\varphi^{-1}(\mathfrak{r}(\mathfrak{a}))
=\mathfrak{r}(\varphi^{-1}(\mathfrak{a}))$;
si se pone $Y=V(\mathfrak{a})$ y $\mathfrak{a}'=\mathfrak{j}({}^a\varphi(Y))$, se tiene $\overline{{}^a\varphi(Y)}=V(\mathfrak{a}')$ (\sref{I.1.1.4}[(ii)]).
La relación $f'\in\mathfrak{a}'$ equivale, por definición, a $f'(x')=0$ para todo $x'\in{}^a\varphi(Y)$, luego, en virtud de la fórmula~\sref{I.1.2.1.1}, equivale asimismo a $\varphi(f')(x)=0$ para todo $x\in Y$, o a $\varphi(f')\in\mathfrak{j}(Y)=\mathfrak{a}$, puesto que $\mathfrak{a}$ es igual a su radical;
de donde (ii).
\end{proof}

\begin{corollary}[1.2.3]
\label{I.1.2.3}
La aplicación ${}^a\varphi$ es continua.
\end{corollary}

Observemos que, si $A''$ es un tercer anillo y $\varphi'$ un homomorfismo $A''\to A'$, entonces se tiene ${}^a(\varphi\circ\varphi')={}^a\varphi'\circ{}^a\varphi$;
este resultado, junto con el Corolario~\sref{I.1.2.3}, significa que $\operatorname{Spec}(A)$ es un \emph{funtor contravariante} en $A$, de la categoría de los anillos a la de los espacios topológicos.

\oldpage[I]{84}
\begin{corollary}[1.2.4]
\label{I.1.2.4}
Supongamos que $\varphi$ es tal que todo $f\in A$ puede escribirse como $f=h\varphi(f')$, donde $h$ es invertible en $A$ \textup{(lo que sucede, en particular, cuando $\varphi$ es \emph{sobreyectivo})}.
Entonces ${}^a\varphi$ es un homeomorfismo de $X$ sobre ${}^a\varphi(X)$.
\end{corollary}

\begin{proof}
Demostremos que, para todo subconjunto $E\subset A$, existe un subconjunto $E'$ de $A'$ tal que $V(E)=V(\varphi(E'))$;
en virtud del axioma ($(T_0)$) \sref{I.1.1.8} y de la fórmula \sref{I.1.2.2.1}, esto implicará, en primer lugar, que ${}^a\varphi$ es inyectiva y, después, de nuevo por \sref{I.1.2.2.1}, que ${}^a\varphi$ es un homeomorfismo.
Ahora bien, para cada $f\in E$, basta elegir $f'\in A'$ tal que $h\varphi(f')=f$, con $h$ invertible en $A$;
el conjunto $E'$ de estos elementos $f'$ es precisamente el que buscábamos.
\end{proof}

\begin{env}[1.2.5]
\label{I.1.2.5}
En particular, cuando $\varphi$ es el homomorfismo canónico de $A$ sobre un anillo cociente $A/\mathfrak{a}$, se recupera \sref{I.1.1.12-fr}, y ${}^a\varphi$ es la \emph{inyección canónica} de $V(\mathfrak{a})$, identificado con $\operatorname{Spec}(A/\mathfrak{a})$, en $X=\operatorname{Spec}(A)$.
\end{env}

Otro caso particular de \sref{I.1.2.4}:
\begin{corollary}[1.2.6]
\label{I.1.2.6}
Si $S$ es un conjunto multiplicativo de $A$, el espectro $\operatorname{Spec}(S^{-1}A)$ se identifica canónicamente (con su topología) con el subespacio de $X=\operatorname{Spec}(A)$ formado por los $x$ tales que $\mathfrak{j}_x\cap S=\varnothing$.
\end{corollary}

\begin{proof}
Se sabe por \sref[0]{0.1.2.6} que los ideales primos de $S^{-1}A$ son los ideales $S^{-1}\mathfrak{j}_x$ tales que $\mathfrak{j}_x\cap S=\varnothing$, y que se tiene $\mathfrak{j}_x=(i_A^S)^{-1}(S^{-1}\mathfrak{j}_x)$.
Basta entonces aplicar el Corolario~\sref{I.1.2.4} a $i_A^S$.
\end{proof}

\begin{corollary}[1.2.7]
\label{I.1.2.7}
Para que ${}^a\varphi(X)$ sea densa en $X'$, es necesario y suficiente que todo elemento del núcleo $\operatorname{Ker}\varphi$ sea nilpotente.
\end{corollary}

\begin{proof}
Aplicando la fórmula~\sref{I.1.2.2.3} al ideal $\mathfrak{a}=(0)$, se tiene $\overline{{}^a\varphi(X)}=V(\operatorname{Ker}\varphi)$, y para que $V(\operatorname{Ker}\varphi)=X'$ se cumpla, es necesario y suficiente que $\operatorname{Ker}\varphi$ esté contenido en todos los ideales primos de $A'$, o, equivalentemente, en el nilradical $\mathfrak{N}'$ de $A'$.
\end{proof}

\subsection{Haz asociado a un módulo}
\label{subsection:I.1.3}
\label{I.1.3}

\begin{env}[1.3.1]
\label{I.1.3.1}
Sean $A$ un anillo conmutativo, $M$ un $A$-módulo, $f$ un elemento de $A$ y $S_f$ el conjunto multiplicativo formado por los $f^n$, donde $n\geq0$.
Recordemos que ponemos $A_f=S_f^{-1}A$, $M_f=S_f^{-1}M$.
Si $S'_f$ es el subconjunto multiplicativo saturado de $A$ formado por los $g\in A$ que dividen a un elemento de $S_f$, se sabe que $A_f$ y $M_f$ se identifican canónicamente con ${S'_f}^{-1}A$ y ${S'_f}^{-1}M$ \sref[0]{0.1.4.3}.
\end{env}

\begin{lemma}[1.3.2]
\label{I.1.3.2}
Las condiciones siguientes son equivalentes:
\begin{enumerate}
  \item[{\textit a)}] $g\in S'_f$;
  \item[{\textit b)}] $S'_g\subset S'_f$;
  \item[{\textit c)}] $f\in\mathfrak{r}(g)$;
  \item[{\textit d)}] $\mathfrak{r}(f)\subset\mathfrak{r}(g)$;
  \item[{\textit e)}] $V(g)\subset V(f)$;
  \item[{\textit f)}] $D(f)\subset D(g)$.
\end{enumerate}
\end{lemma}

\begin{proof}
Esto resulta inmediatamente de las definiciones y de \sref{I.1.1.5}.
\end{proof}

\begin{env}[1.3.3]
\label{I.1.3.3}
Si $D(f)=D(g)$, el Lema~\sref{I.1.3.2}[\textit{b})] muestra que $M_f=M_g$.
Más generalmente, si $D(f)\supset D(g)$, entonces $S'_f\subset S'_g$, y se sabe \sref[0]{0.1.4.1} que existe un homomorfismo canónico y funtorial
\[
  \rho_{g,f}:M_f\longrightarrow M_g,
\]
y, si $D(f)\supset D(g)\supset D(h)$, se tiene \sref[0]{0.1.4.4}
\begin{equation}
  \rho_{h,g}\circ\rho_{g,f}=\rho_{h,f}.
  \tag{1.3.3.1}\label{I.1.3.3.1-fr}
\end{equation}
\end{env}

\oldpage[I]{85}
Cuando $f$ recorre los elementos de $A-\mathfrak{j}_x$ (para un $x$ dado de $X=\operatorname{Spec}(A)$), los conjuntos $S'_f$ constituyen una familia filtrante creciente de subconjuntos de $A-\mathfrak{j}_x$, pues, para dos elementos $f$ y $g$ de $A-\mathfrak{j}_x$, $S'_f$ y $S'_g$ están contenidos en $S'_{fg}$;
como la unión de los $S'_f$ para $f\in A-\mathfrak{j}_x$ es $A-\mathfrak{j}_x$, se concluye \sref[0]{0.1.4.5} que el $A_x$-módulo $M_x$ se identifica canónicamente con el \emph{límite inductivo} $\varinjlim M_f$, respecto de la familia de homomorfismos $(\rho_{g,f})$.
Denotaremos por
\[
  \rho_x^f:M_f\longrightarrow M_x
\]
el homomorfismo canónico para $f\in A-\mathfrak{j}_x$ (o, lo que es equivalente, $x\in D(f)$).

\begin{definition}[1.3.4]
\label{I.1.3.4}
{\itshape Se llama haz estructural del espectro primo $X=\operatorname{Spec}(A)$ \textup{(resp. \emph{haz asociado al $A$-módulo} $M$)}, y se denota por $\widetilde{A}$ o $\mathscr{O}_X$ \textup{(resp. $\widetilde{M}$)}, al haz de anillos \textup{(resp.} al $\widetilde{A}$-módulo) asociado al prehaz $D(f)\to A_f$ \textup{(resp. $D(f)\to M_f$)}, definido sobre la base $\mathfrak{B}$ de $X$ formada por los $D(f)$ para $f\in A$ \textup{(\sref{I.1.1.10}, (\hyperref[0.3.2.1]{\textbf{0},~3.2.1} y \hyperref[0.3.5.6]{3.5.6}))}.}
\end{definition}

Se ha visto \sref[0]{0.3.2.4} que el tallo $\widetilde{A}_x$ (resp. $\widetilde{M}_x$) \emph{se identifica con el anillo $A_x$} (resp. \emph{con el $A_x$-módulo $M_x$});
denotaremos por
\[
  \theta_f:A_f\longrightarrow\Gamma(D(f),\widetilde{A})
\]
\[
  \text{(resp. }\theta_f:M_f\longrightarrow
  \Gamma(D(f),\widetilde{M})\text{),}
\]
la aplicación canónica, de modo que, para todo $x\in D(f)$ y todo $\xi\in M_f$, se tiene
\begin{equation}
  (\theta_f(\xi))_x=\rho_x^f(\xi).
  \tag{1.3.4.1}\label{I.1.3.4.1-fr}
\end{equation}

\begin{proposition}[1.3.5]
\label{I.1.3.5}
$\widetilde{M}$ es un funtor covariante exacto en $M$, de la categoría de $A$-módulos a la categoría de $\widetilde{A}$-módulos.
\end{proposition}

\begin{proof}
En efecto, sean $M$, $N$ dos $A$-módulos y $u$ un homomorfismo $M\to N$;
para todo $f\in A$, a $u$ le corresponde canónicamente un homomorfismo $u_f$ del $A_f$-módulo $M_f$ en el $A_f$-módulo $N_f$, y el diagrama (para $D(g)\subset D(f)$)
\[
  \xymatrix{
    M_f\ar[r]^{u_f}\ar[d]_{\rho_{g,f}} & N_f\ar[d]^{\rho_{g,f}}\\
    M_g\ar[r]_{u_g} & N_g
  }
\]
es conmutativo \sref[0]{0.1.4.1};
estos homomorfismos definen, por tanto, un homomorfismo de $\widetilde{A}$-módulos $\widetilde{u}:\widetilde{M}\to\widetilde{N}$ \sref[0]{0.3.2.3}.
Además, para todo $x\in X$, $\widetilde{u}_x$ es el límite inductivo de los $u_f$ para $x\in D(f)$ ($f\in A$), y, por consiguiente \sref[0]{0.1.4.5}, si se identifican canónicamente $\widetilde{M}_x$ y $\widetilde{N}_x$ con $M_x$ y $N_x$ respectivamente, entonces $\widetilde{u}_x$ se identifica con el homomorfismo $u_x$ inducido canónicamente por $u$.
Si $P$ es un tercer $A$-módulo, $v$ un homomorfismo $N\to P$ y $w=v\circ u$, es inmediato que $w_x=v_x\circ u_x$, luego $\widetilde{w}=\widetilde{v}\circ\widetilde{u}$.
Queda así bien definido un \emph{funtor covariante (en $M$)} $\widetilde{M}$, de la categoría de $A$-módulos a la de $\widetilde{A}$-módulos.
\emph{Este funtor es exacto}, pues, para todo $x\in X$, $M_x$ es un funtor exacto en $M$ \sref[0]{0.1.3.2};
además, se tiene $\operatorname{Supp}(M)=\operatorname{Supp}(\widetilde{M})$, en virtud de las definiciones (\hyperref[0.1.7.1]{\textbf{0},~1.7.1} y \hyperref[0.3.1.6]{3.1.6}).
\end{proof}

\oldpage[I]{86}
\begin{proposition}[1.3.6]
\label{I.1.3.6}
Para todo $f\in A$, el subconjunto abierto $D(f)\subset X$ se identifica canónicamente con el espectro primo $\operatorname{Spec}(A_f)$, y el haz $\widetilde{M_f}$ asociado al $A_f$-módulo $M_f$ se identifica canónicamente con la restricción $\widetilde{M}|D(f)$.
\end{proposition}

\begin{proof}
La primera afirmación es un caso particular de \sref{I.1.2.6}.
Además, si $g\in A$ es tal que $D(g)\subset D(f)$, entonces $M_g$ se identifica canónicamente con el módulo de fracciones de $M_f$ cuyos denominadores son las potencias de la imagen canónica de $g$ en $A_f$ \sref[0]{0.1.4.6}.
La identificación canónica de $\widetilde{M_f}$ con $\widetilde{M}|D(f)$ resulta entonces de las definiciones.
\end{proof}

\begin{theorem}[1.3.7]
\label{I.1.3.7}
Para todo $A$-módulo $M$ y todo $f\in A$, el homomorfismo
\[
  \theta_f:M_f\longrightarrow\Gamma(D(f),\widetilde{M})
\]
es biyectivo \textup{(en otras palabras, el prehaz $D(f)\to M_f$ es un \emph{haz})}.
En particular, $M$ se identifica con $\Gamma(X,\widetilde{M})$ mediante $\theta_1$.
\end{theorem}

\begin{proof}
Se observará que, si $M=A$, entonces $\theta_f$ es un homomorfismo de anillos;
el Teorema~\sref{I.1.3.7} implicará, por tanto, que, si se identifican los anillos $A_f$ y $\Gamma(D(f),\widetilde{A})$ mediante $\theta_f$, el homomorfismo $\theta_f:M_f\to\Gamma(D(f),\widetilde{M})$ será un isomorfismo de \emph{módulos}.

Mostremos primero que $\theta_f$ es \emph{inyectivo}.
En efecto, si $\xi\in M_f$ es tal que $\theta_f(\xi)=0$, esto significa que, para todo ideal primo $\mathfrak{p}$ de $A_f$, existe $h\notin\mathfrak{p}$ tal que $h\xi=0$;
como el anulador de $\xi$ no está contenido en ningún ideal primo de $A_f$, es todo el anillo $A_f$ y, por tanto, $\xi=0$.

Queda por demostrar que $\theta_f$ es \emph{sobreyectivo};
se puede reducir al caso en que $f=1$ (deduciéndose entonces el caso general al ``localizar'', mediante \sref{I.1.3.6}).
Sea, pues, $s$ una sección de $\widetilde{M}$ sobre $X$;
en virtud de \sref{I.1.3.4} y \sref{I.1.1.10}[ii], existe un recubrimiento \emph{finito} $(D(f_i))_{i\in I}$ de $X$ ($f_i\in A$) tal que, para todo $i\in I$, la restricción $s_i=s|D(f_i)$ es de la forma $\theta_{f_i}(\xi_i)$, donde $\xi_i\in M_{f_i}$.
Si $i$, $j$ son dos índices de $I$, al expresar que las restricciones de $s_i$ y $s_j$ a $D(f_i)\cap D(f_j)=D(f_i f_j)$ son iguales, resulta, por definición de $\widetilde{M}$, que
\begin{equation}
  \rho_{f_i f_j,f_i}(\xi_i)=\rho_{f_i f_j,f_j}(\xi_j).
  \tag{1.3.7.1}\label{I.1.3.7.1-fr}
\end{equation}
Por definición, se puede escribir, para todo $i\in I$, $\xi_i=z_i/f_i^{n_i}$, donde $z_i\in M$, y, como $I$ es finito, multiplicando cada $z_i$ por una potencia de $f_i$, se puede suponer que todos los $n_i$ son iguales a un mismo $n$.
Entonces, por definición, \sref{I.1.3.7.1} significa que existe un entero $m_{ij}\geq0$ tal que $(f_i f_j)^{m_{ij}}(f_j^n z_i-f_i^n z_j)=0$, y se puede suponer además que todos los $m_{ij}$ son iguales a un mismo entero $m$;
reemplazando entonces los $z_i$ por $f_i^m z_i$, se reduce al caso en que $m=0$, es decir, al caso en que se tiene
\begin{equation}
  f_j^n z_i=f_i^n z_j
  \tag{1.3.7.2}\label{I.1.3.7.2-fr}
\end{equation}
para cualesquiera $i$, $j$.
Se tiene $D(f_i^n)=D(f_i)$ y, como los $D(f_i)$ forman un recubrimiento de $X$, el ideal engendrado por los $f_i^n$ es $A$;
en otras palabras, existen elementos $g_i\in A$ tales que $\sum_i g_i f_i^n=1$.
Consideremos entonces el elemento $z=\sum_i g_i z_i$ de $M$;
por \sref{I.1.3.7.2}, se tiene $f_i^n z=\sum_j g_j f_i^n z_j=(\sum_j g_j f_j^n)z_i=z_i$, de donde, por definición, $\xi_i=z/1$ en $M_{f_i}$.
Se concluye
\oldpage[I]{87}
que los $s_i$ son las restricciones a $D(f_i)$ de $\theta_1(z)$, lo que demuestra que $s=\theta_1(z)$ y termina la demostración.
\end{proof}

\begin{corollary}[1.3.8]
\label{I.1.3.8}
Sean $M$ y $N$ dos $A$-módulos;
el homomorfismo canónico $u\to\widetilde{u}$ de $\operatorname{Hom}_A(M,N)$ en $\operatorname{Hom}_{\widetilde{A}}(\widetilde{M},\widetilde{N})$ es biyectivo.
En particular, las relaciones $M=0$ y $\widetilde{M}=0$ son equivalentes.
\end{corollary}

\begin{proof}
Consideremos, en efecto, el homomorfismo canónico $v\to\Gamma(v)$ de $\operatorname{Hom}_{\widetilde{A}}(\widetilde{M},\widetilde{N})$ en $\operatorname{Hom}_{\Gamma(\widetilde{A})}
(\Gamma(\widetilde{M}),\Gamma(\widetilde{N}))$;
este último módulo se identifica canónicamente con $\operatorname{Hom}_A(M,N)$, por el Teorema~\sref{I.1.3.7}.
Queda por comprobar que $u\to\widetilde{u}$ y $v\to\Gamma(v)$ son recíprocos;
es evidente que $\Gamma(\widetilde{u})=u$ por definición de $\widetilde{u}$;
por otra parte, si se pone $u=\Gamma(v)$ para $v\in\operatorname{Hom}_{\widetilde{A}}(\widetilde{M},\widetilde{N})$, entonces la aplicación $w:\Gamma(D(f),\widetilde{M})\to\Gamma(D(f),\widetilde{N})$ inducida canónicamente por $v$ es tal que el diagrama
\[
  \xymatrix{
    M\ar[r]^u\ar[d]_{\rho_{f,1}} & N\ar[d]^{\rho_{f,1}}\\
    M_f\ar[r]_w & N_f
  }
\]
es conmutativo;
por tanto, necesariamente $w=u_f$ para todo $f\in A$ \sref[0]{0.1.2.4}, lo que demuestra que $(\Gamma(v))\supertilde=v$.
\end{proof}

\begin{corollary}[1.3.9]
\label{I.1.3.9}
\medskip\noindent
\begin{enumerate}
  \item[{\rm(i)}] Sea $u$ un homomorfismo de un $A$-módulo $M$ en un $A$-módulo $N$;
    entonces los haces asociados a $\operatorname{Ker}u$, $\operatorname{Im}u$ y $\operatorname{Coker}u$ son, respectivamente, $\operatorname{Ker}\widetilde{u}$, $\operatorname{Im}\widetilde{u}$ y $\operatorname{Coker}\widetilde{u}$.
    En particular, para que $\widetilde{u}$ sea inyectivo \textup{(resp. \emph{sobreyectivo, biyectivo})}, es necesario y suficiente que $u$ lo sea.
  \item[{\rm(ii)}] Si $M$ es un límite inductivo \textup{(resp. una \emph{suma directa})} de una familia de $A$-módulos $(M_\lambda)$, entonces $\widetilde{M}$ es el límite inductivo \textup{(resp. la \emph{suma directa})} de la familia $(\widetilde{M_\lambda})$, mediante un isomorfismo canónico.
\end{enumerate}
\end{corollary}

\begin{proof}
\medskip\noindent
\begin{enumerate}
  \item[(i)] Basta aplicar el hecho de que $\widetilde{M}$ es un funtor exacto en $M$ \sref{I.1.3.5} a las dos sucesiones exactas de $A$-módulos
    \[
  0\to\operatorname{Ker}u\to M\to\operatorname{Im}u\to0
\]
    \[
  0\to\operatorname{Im}u\to N\to\operatorname{Coker}u\to0
\]
    La segunda afirmación resulta entonces del Teorema~\sref{I.1.3.7}.
  \item[(ii)] Sea $(M_\lambda,g_{\mu\lambda})$ un sistema inductivo de $A$-módulos, de límite inductivo $M$, y sea $g_\lambda$ el homomorfismo canónico $M_\lambda\to M$.
    Como se tiene $\widetilde{g_{\nu\mu}}\circ\widetilde{g_{\mu\lambda}}
=\widetilde{g_{\nu\lambda}}$ y $\widetilde{g_\lambda}=\widetilde{g_\mu}\circ
\widetilde{g_{\mu\lambda}}$ para $\lambda\leq\mu\leq\nu$, resulta que $(\widetilde{M_\lambda},\widetilde{g_{\mu\lambda}})$ es un sistema inductivo de haces sobre $X$, y, si se denota por $h_\lambda$ el homomorfismo canónico $\widetilde{M_\lambda}\to\varinjlim\widetilde{M_\lambda}$, existe un único homomorfismo $v:\varinjlim\widetilde{M_\lambda}\to\widetilde{M}$ tal que $v\circ h_\lambda=\widetilde{g_\lambda}$.
    Para ver que $v$ es biyectivo, basta comprobar, para todo $x\in X$, que $v_x$ es una biyección de $(\varinjlim\widetilde{M_\lambda})_x$ en $\widetilde{M}_x$;
    pero $\widetilde{M}_x=M_x$, y
    \[
  (\varinjlim\widetilde{M_\lambda})_x
  =\varinjlim(\widetilde{M_\lambda})_x
  =\varinjlim(M_\lambda)_x=M_x
  \quad(\hyperref[0.1.3.3-fr]{\textbf{0}, 1.3.3}).
\]
    Por otra parte, de las definiciones resulta que $(\widetilde{g_\lambda})_x$ y
    $(h_\lambda)_x$ son ambos iguales a la aplicación canónica de $(M_\lambda)_x$ en $M_x$;
    como $(\widetilde{g_\lambda})_x=v_x\circ(h_\lambda)_x$, $v_x$ es la identidad.
\oldpage[I]{88}
    Finalmente, si $M$ es la suma directa de dos $A$-módulos $N$ y $P$, es inmediato que $\widetilde{M}=\widetilde{N}\oplus\widetilde{P}$;
    siendo toda suma directa el límite inductivo de sumas directas finitas, quedan demostradas las afirmaciones de (ii).
\end{enumerate}
\end{proof}

Se observará que el Corolario~\sref{I.1.3.8} demuestra que los haces isomorfos a haces asociados a $A$-módulos forman una \emph{categoría abeliana} (T, I, 1.4).

Se observará también que del Corolario~\sref{I.1.3.9} resulta que, si $M$ es un $A$-módulo \emph{de tipo finito} (es decir, si existe un homomorfismo sobreyectivo $A^n\to M$), entonces existe un homomorfismo sobreyectivo $\widetilde{A}^n\to\widetilde{M}$ o, en otras palabras, el $\widetilde{A}$-módulo $\widetilde{M}$ está \emph{engendrado por una familia finita de secciones sobre $X$} \sref[0]{0.5.1.1}, y recíprocamente.

\begin{env}[1.3.10]
\label{I.1.3.10}
Si $N$ es un submódulo de un $A$-módulo $M$, la inyección canónica $j:N\to M$ da, por \sref{I.1.3.9}, un homomorfismo inyectivo $\widetilde{N}\to\widetilde{M}$, que permite identificar canónicamente $\widetilde{N}$ con un \emph{$\widetilde{A}$-submódulo} de $\widetilde{M}$;
supondremos siempre efectuada esta identificación.
Si $N$ y $P$ son dos submódulos de $M$, entonces se tiene
\begin{equation}
  (N+P)\supertilde=\widetilde{N}+\widetilde{P}
  \tag{1.3.10.1}\label{I.1.3.10.1-fr}
\end{equation}
\begin{equation}
  (N\cap P)\supertilde=\widetilde{N}\cap\widetilde{P}
  \tag{1.3.10.2}\label{I.1.3.10.2-fr}
\end{equation}
pues $N+P$ y $N\cap P$ son, respectivamente, la imagen del homomorfismo canónico
$N\oplus P\to M$ y el núcleo del homomorfismo canónico $M\to(M/N)\oplus(M/P)$, y basta aplicar \sref{I.1.3.9}.

De \sref{I.1.3.10.1} y \sref{I.1.3.10.2} se concluye que, si $\widetilde{N}=\widetilde{P}$, entonces $N=P$.
\end{env}

\begin{corollary}[1.3.11]
\label{I.1.3.11}
En la categoría de los haces isomorfos a haces asociados a $A$-módulos, el funtor $\Gamma$ es exacto.
\end{corollary}

\begin{proof}
Sea $\widetilde{M}\xrightarrow{\widetilde{u}}\widetilde{N}
\xrightarrow{\widetilde{v}}\widetilde{P}$ una sucesión exacta correspondiente a dos homomorfismos $u:M\to N$ y $v:N\to P$ de $A$-módulos.
Si $Q=\operatorname{Im}u$ y $R=\operatorname{Ker}v$, se tiene $\widetilde{Q}=\operatorname{Im}\widetilde{u}
=\operatorname{Ker}\widetilde{v}=\widetilde{R}$ (Corolario~\sref{I.1.3.9}), luego $Q=R$.
\end{proof}

\begin{corollary}[1.3.12]
\label{I.1.3.12}
Sean $M$ y $N$ dos $A$-módulos.
\begin{enumerate}
  \item[{\rm(i)}] El haz asociado a $M\otimes_A N$ se identifica canónicamente con $\widetilde{M}\otimes_{\widetilde{A}}\widetilde{N}$.
  \item[{\rm(ii)}] Si, además, $M$ admite una presentación finita, entonces el haz asociado a $\operatorname{Hom}_A(M,N)$ se identifica canónicamente con $\mathscr{H}\!om_{\widetilde{A}}(\widetilde{M},\widetilde{N})$.
\end{enumerate}
\end{corollary}

\begin{proof}
\medskip\noindent
\begin{enumerate}
  \item[(i)] El haz $\mathscr{F}=\widetilde{M}\otimes_{\widetilde{A}}\widetilde{N}$ está asociado al prehaz
    \[
  U\to\mathscr{F}(U)=\Gamma(U,\widetilde{M})
  \otimes_{\Gamma(U,\widetilde{A})}\Gamma(U,\widetilde{N})
\]
    donde $U$ recorre la base \sref{I.1.1.10}[i] de $X$ formada por los $D(f)$, con $f\in A$.
    Ahora bien, $\mathscr{F}(D(f))$ se identifica canónicamente con $M_f\otimes_{A_f}N_f$, por \sref{I.1.3.7} y \sref{I.1.3.6}.
    Por otra parte, se sabe que el $A_f$-módulo $M_f\otimes_{A_f}N_f$ es canónicamente
    isomorfo a $(M\otimes_A N)_f$ \sref[0]{0.1.3.4}, que a su vez es canónicamente isomorfo a $\Gamma(D(f),(M\otimes_A N)\supertilde)$ (Teorema~\sref{I.1.3.7} y Proposición~\sref{I.1.3.6}).
    Además, se comprueba inmediatamente que los isomorfismos canónicos
    \[
  \mathscr{F}(D(f))\xrightarrow{\sim}\Gamma(D(f),(M\otimes_A N)\supertilde)
\]
\oldpage[I]{89}
    así obtenidos satisfacen las condiciones de compatibilidad con respecto a los operadores de restricción \sref[0]{0.1.4.2}, por lo que definen un isomorfismo canónico funtorial
    \[
  \widetilde{M}\otimes_{\widetilde{A}}\widetilde{N}
  \xrightarrow{\sim}(M\otimes_A N)\supertilde
\]
  \item[(ii)] El haz $\mathscr{G}=\mathscr{H}\!om_{\widetilde{A}}(\widetilde{M},\widetilde{N})$ está asociado al prehaz
    \[
  U\longrightarrow\mathscr{G}(U)
  =\operatorname{Hom}_{\widetilde{A}|U}
  (\widetilde{M}|U,\widetilde{N}|U)
\]
    donde $U$ recorre la base de $X$ formada por los $D(f)$.
    Ahora bien, $\mathscr{G}(D(f))$ se identifica canónicamente con $\operatorname{Hom}_{A_f}(M_f,N_f)$ (Proposición~\sref{I.1.3.6} y
    Corolario~\sref{I.1.3.8}), que, en virtud de las hipótesis sobre $M$, se identifica canónicamente con $(\operatorname{Hom}_A(M,N))_f$ \sref[0]{0.1.3.5}.
    Finalmente, $(\operatorname{Hom}_A(M,N))_f$ se identifica canónicamente con $\Gamma(D(f),(\operatorname{Hom}_A(M,N))\supertilde)$ (Proposición~\sref{I.1.3.6} y Teorema~\sref{I.1.3.7}), y los isomorfismos canónicos $\mathscr{G}(D(f))\to
\Gamma(D(f),(\operatorname{Hom}_A(M,N))\supertilde)$ así obtenidos son compatibles con los operadores de restricción \sref[0]{0.1.4.2};
    definen, por tanto, un isomorfismo canónico \[
  \mathscr{H}\!om_{\widetilde{A}}(\widetilde{M},\widetilde{N})
  \xrightarrow{\sim}(\operatorname{Hom}_A(M,N))\supertilde.
\].
\end{enumerate}
\end{proof}

\begin{env}[1.3.13]
\label{I.1.3.13}
Sea ahora $B$ una $A$-álgebra (conmutativa); esto puede interpretarse diciendo que $B$ es un $A$-módulo y que se han dado un elemento $e\in B$ y un $A$-homomorfismo $\varphi:B\otimes_A B\to B$, de modo que: $1^\circ$ los diagramas
\[
  \xymatrix{
    B\otimes_A B\otimes_A B\ar[r]^{\varphi\otimes 1}
      \ar[d]_{1\otimes\varphi} &
    B\otimes_A B\ar[d]^\varphi & &
    B\otimes_A B\ar[rr]^\sigma\ar[rd]_\varphi & &
    B\otimes_A B\ar[dl]^\varphi\\
    B\otimes_A B\ar[r]_\varphi &
    B & & & B
  }
\]
(donde $\sigma$ es la simetría canónica) son conmutativos; $2^\circ$ $\varphi(e\otimes x)=\varphi(x\otimes e)=x$.
Por el Corolario~\sref{I.1.3.12}, el homomorfismo $\widetilde{\varphi}:\widetilde{B}\otimes_{\widetilde{A}}\widetilde{B}
\to\widetilde{B}$ de $\widetilde{A}$-módulos satisface las condiciones análogas y define, por tanto, una estructura de $\widetilde{A}$-\emph{álgebra} sobre $\widetilde{B}$.
Análogamente, dar un $B$-módulo $N$ equivale a dar un $A$-módulo $N$ y un $A$-homomorfismo $\psi:B\otimes_A N\to N$ tal que el diagrama
\[
  \xymatrix{
    B\otimes_A B\otimes_A N\ar[r]^{\varphi\otimes 1}
      \ar[d]_{1\otimes\psi} &
    B\otimes_A N\ar[d]^\psi\\
    B\otimes_A N\ar[r]_\psi & N
  }
\]
sea conmutativo y $\psi(e\otimes n)=n$;
el homomorfismo $\widetilde{\psi}:\widetilde{B}\otimes_{\widetilde{A}}\widetilde{N}
\to\widetilde{N}$ satisface la condición análoga y define, por tanto, una estructura de $\widetilde{B}$-\emph{módulo} sobre $\widetilde{N}$.

De la misma manera, se ve que, si $u:B\to B'$ (resp. $v:N\to N'$) es un homomorfismo de $A$-álgebras (resp. de $B$-módulos), entonces $\widetilde{u}$ (resp. $\widetilde{v}$) es un homomorfismo de $\widetilde{A}$-álgebras (resp. de $\widetilde{B}$-módulos), y $\operatorname{Ker}\widetilde{u}$ es un $\widetilde{B}$-ideal (resp. $\operatorname{Ker}\widetilde{v}$, $\operatorname{Coker}\widetilde{v}$ e $\operatorname{Im}\widetilde{v}$ son $\widetilde{B}$-módulos).
Si $N$ es un $B$-módulo, entonces $\widetilde{N}$ es un $\widetilde{B}$-módulo de tipo finito si y solo si $N$ es un $B$-módulo de tipo finito \sref[0]{0.5.2.3}.

\oldpage[I]{90}
Si $M$, $N$ son dos $B$-módulos, entonces el $\widetilde{B}$-módulo $\widetilde{M}\otimes_{\widetilde{B}}\widetilde{N}$ se identifica canónicamente con $(M\otimes_B N)\supertilde$;
asimismo, $\mathscr{H}\!om_{\widetilde{B}}(\widetilde{M},\widetilde{N})$ se identifica canónicamente con $(\operatorname{Hom}_B(M,N))\supertilde$ cuando $M$ admite una presentación finita;
las demostraciones son las mismas que las del Corolario~\sref{I.1.3.12}.

Si $\mathfrak{J}$ es un ideal de $B$ y $N$ es un $B$-módulo, entonces se tiene $(\mathfrak{J}N)\supertilde=\widetilde{\mathfrak{J}}\cdot\widetilde{N}$.

Finalmente, si $B$ es una $A$-álgebra \emph{graduada} por los $A$-submódulos $B_n$ ($n\in\mathbb{Z}$), entonces la $\widetilde{A}$-álgebra $\widetilde{B}$, suma directa de los $\widetilde{A}$-módulos $\widetilde{B}_n$ \sref{I.1.3.9}, está graduada por estos $\widetilde{A}$-submódulos, pues el axioma de la graduación expresa que la imagen del homomorfismo $B_m\otimes B_n\to B$ está contenida en $B_{m+n}$.
Análogamente, si $M$ es un $B$-módulo graduado por los submódulos $M_n$, entonces $\widetilde{M}$ es un $\widetilde{B}$-módulo graduado por los $\widetilde{M}_n$.
\end{env}

\begin{env}[1.3.14]
\label{I.1.3.14}
Si $B$ es una $A$-álgebra y $M$ un submódulo de $B$, entonces la $\widetilde{A}$-subálgebra de $\widetilde{B}$ engendrada por $\widetilde{M}$ \sref[0]{0.4.1.3} es la $\widetilde{A}$-subálgebra $\widetilde{C}$, donde se denota por $C$ la subálgebra de $B$ engendrada por $M$.
En efecto, $C$ es la suma de los submódulos de $B$ que son las imágenes de los homomorfismos $\bigotimes^n M\to B$ ($n\geq0$), y basta aplicar \sref{I.1.3.9} y \sref{I.1.3.12}.
\end{env}

\subsection{Haces cuasicoherentes sobre un espectro primo}
\label{subsection:I.1.4}

\begin{theorem}[1.4.1]
\label{I.1.4.1}
Sean $X$ el espectro primo de un anillo $A$, $V$ un subconjunto abierto cuasicompacto de $X$ y $\mathscr{F}$ un $(\mathscr{O}_X|V)$-módulo.
Las cuatro condiciones siguientes son equivalentes.
\begin{enumerate}
  \item[\textup{a)}] Existe un $A$-módulo $M$ tal que $\mathscr{F}$ sea isomorfo a $\widetilde{M}|V$.
  \item[\textup{b)}] Existe un cubrimiento abierto finito $(V_i)$ de $V$ por conjuntos de la forma $D(f_i)$ ($f_i\in A$) contenidos en $V$, tal que, para cada $i$, $\mathscr{F}|V_i$ sea isomorfo a un haz de la forma $\widetilde{M}_i$, donde $M_i$ es un $A_{f_i}$-módulo.
  \item[\textup{c)}] El haz $\mathscr{F}$ es cuasicoherente \sref[0]{0.5.1.3}.
  \item[\textup{d)}] Se satisfacen las dos propiedades siguientes:
    \begin{enumerate}
      \item[\textup{d\,1)}] Para todo $f\in A$ tal que $D(f)\subset V$, y para toda sección $s\in\Gamma(D(f),\mathscr{F})$, existe un entero $n\geq0$ tal que $f^n s$ se prolonga a una sección de $\mathscr{F}$ sobre $V$.
      \item[\textup{d\,2)}] Para todo $f\in A$ tal que $D(f)\subset V$ y para toda sección $t\in\Gamma(V,\mathscr{F})$ tal que la restricción de $t$ a $D(f)$ es $0$, existe un entero $n\geq0$ tal que $f^n t=0$.
    \end{enumerate}
\end{enumerate}
\end{theorem}
(En el enunciado de las condiciones (d1) y (d2), se han identificado tácitamente $A$ y $\Gamma(\widetilde{A})$ mediante el Teorema~\sref{I.1.3.7}).

\begin{proof}
El hecho de que (a) implica (b) resulta inmediatamente de la Proposición~\sref{I.1.3.6} y del hecho de que los $D(f_i)$ forman una base de la topología de $X$ \sref{I.1.1.10}.
Como todo $A$-módulo es isomorfo al conúcleo de un homomorfismo de la forma $A^{(I)}\to A^{(J)}$, \sref{I.1.3.9} demuestra que todo haz asociado a un $A$-módulo es cuasicoherente;
por tanto, (b) implica (c).
Recíprocamente, si $\mathscr{F}$ es cuasicoherente, todo $x\in V$ posee una vecindad de la forma $D(f)\subset V$ tal que $\mathscr{F}|D(f)$ es isomorfo al conúcleo de un homomorfismo $\widetilde{A}_f^{(I)}\to\widetilde{A}_f^{(J)}$, y por tanto también al haz asociado al módulo $N$, conúcleo del homomorfismo correspondiente
$A_f^{(I)}\to A_f^{(J)}$ (Corolarios~\sref{I.1.3.8} y \sref{I.1.3.9});
como $V$ es cuasicompacto, es claro que (c) implica (b).

\oldpage[I]{91}
Para demostrar que (b) implica (d1) y (d2), supongamos primero que $V=D(g)$ para cierto $g\in A$, y que $\mathscr{F}$ es isomorfo al haz $\widetilde{N}$ asociado a un $A_g$-módulo $N$;
reemplazando $X$ por $V$ y $A$ por $A_g$ \sref{I.1.3.6}, se puede reducir al caso en que $g=1$.
Entonces $\Gamma(D(f),\widetilde{N})$ y $N_f$ se identifican canónicamente entre sí (Proposición~\sref{I.1.3.6} y Teorema~\sref{I.1.3.7}), de modo que una sección $s\in\Gamma(D(f),\widetilde{N})$ se identifica con un elemento de la forma $z/f^n$, donde $z\in N$;
la sección $f^n s$ se identifica con el elemento $z/1$ de $N_f$ y, por consiguiente, es la restricción a $D(f)$ de la sección de $\widetilde{N}$ sobre $X$ que se identifica con el elemento $z\in N$;
de donde resulta (d1) en este caso.
Análogamente, $t\in\Gamma(X,\widetilde{N})$ se identifica con un elemento $z'\in N$, la restricción de $t$ a $D(f)$ se identifica con la imagen $z'/1$ de $z'$ en $N_f$, y decir que esta imagen es cero significa que existe un $n\geq0$ tal que $f^n z'=0$ en $N$, o, equivalentemente, $f^n t=0$.

Para terminar de demostrar que (b) implica (d1) y (d2), basta establecer el lema siguiente.
\begin{lemma}[1.4.1.1]
\label{I.1.4.1.1}
Supongamos que $V$ es la unión finita de conjuntos de la forma $D(g_i)$, y que todos los haces $\mathscr{F}|D(g_i)$ y $\mathscr{F}|(D(g_i)\cap D(g_j))=\mathscr{F}|D(g_i g_j)$ satisfacen \textup{d\,1)} y \textup{d\,2)};
entonces $\mathscr{F}$ posee las dos propiedades siguientes:
\begin{enumerate}
  \item[\textup{d'\,1)}] Para todo $f\in A$ y para toda sección $s\in\Gamma(D(f)\cap V,\mathscr{F})$, existe un entero $n\geq0$ tal que $f^n s$ se prolonga a una sección de $\mathscr{F}$ sobre $V$.
  \item[\textup{d'\,2)}] Para todo $f\in A$ y para toda sección $t\in\Gamma(V,\mathscr{F})$ tal que la restricción de $t$ a $D(f)\cap V$ es $0$, existe un entero $n\geq0$ tal que $f^n t=0$.
\end{enumerate}
\end{lemma}

Demostremos primero (d'2): como $D(f)\cap D(g_i)=D(fg_i)$, existe, para cada $i$, un entero $n_i$ tal que la restricción de $(fg_i)^{n_i}t$ a $D(g_i)$ es cero:
como la imagen de $g_i$ en $A_{g_i}$ es invertible, la restricción de $f^{n_i}t$ a $D(g_i)$ es también cero;
tomando como $n$ el mayor de los $n_i$, queda demostrado (d'2).

Para demostrar (d'1), apliquemos (d1) al haz $\mathscr{F}|D(g_i)$: existen un entero $n_i\geq0$ y una sección $s_i'$ de $\mathscr{F}$ sobre $D(g_i)$ que prolonga la restricción de $(fg_i)^{n_i}s$ a $D(fg_i)$;
como la imagen de $g_i$ en $A_{g_i}$ es invertible, existe una sección $s_i$ de $\mathscr{F}$ sobre $D(g_i)$ tal que $s_i'=g_i^{n_i}s_i$, y $s_i$ prolonga la
restricción de $f^{n_i}s$ a $D(fg_i)$;
además, se puede suponer que todos los $n_i$ son iguales a un mismo entero $n$.
Por construcción, la restricción de $s_i-s_j$ a $D(f)\cap D(g_i)\cap D(g_j)=D(fg_i g_j)$ es cero;
aplicando (d2) al haz $\mathscr{F}|D(g_i g_j)$, existe un entero $m_{ij}\geq0$ tal que la restricción a $D(g_i g_j)$ de $(fg_i g_j)^{m_{ij}}(s_i-s_j)$ es cero;
como la imagen de $g_i g_j$ en $A_{g_i g_j}$ es invertible, la restricción de $f^{m_{ij}}(s_i-s_j)$ a $D(g_i g_j)$ es cero.
Se puede suponer entonces que todos los $m_{ij}$ son iguales a un mismo entero $m$, y por tanto existe una sección $s'\in\Gamma(V,\mathscr{F})$ que prolonga todas las secciones $f^m s_i$;
por consiguiente, esta sección prolonga $f^{n+m}s$, y así queda demostrado (d'1).

Queda por demostrar que (d1) y (d2) implican (a).
Demostremos primero que (d1) y (d2) implican que estas condiciones se satisfacen para todo haz $\mathscr{F}|D(g)$, donde $g\in A$ es tal que $D(g)\subset V$.
Esto es evidente para (d1);
por otra parte, si $t\in\Gamma(D(g),\mathscr{F})$ es tal que su restricción a $D(f)\subset D(g)$ es cero, existe, por (d1), un entero $m\geq0$ tal que $g^m t$
\oldpage[I]{92}
se prolonga a una sección $s$ de $\mathscr{F}$ sobre $V$;
aplicando (d2), se ve que existe un entero $n\geq0$ tal que $f^n g^m t=0$, y, como la imagen de $g$ en $A_g$ es invertible, $f^n t=0$.

Así pues, como $V$ es cuasicompacto, el Lema~\sref{I.1.4.1.1} demuestra que
se satisfacen las condiciones (d'1) y (d'2).
Consideremos entonces el $A$-módulo $M=\Gamma(V,\mathscr{F})$, y definamos un homomorfismo de $\widetilde{A}$-módulos $u:\widetilde{M}\to j_*(\mathscr{F})$, donde $j$ es la inclusión canónica $V\to X$.
Como los $D(f)$ forman una base de la topología de $X$, basta, para cada $f\in A$, definir un homomorfismo $u_f:M_f\to\Gamma(D(f),j_*(\mathscr{F}))=\Gamma(D(f)\cap V,\mathscr{F})$, con las condiciones de compatibilidad usuales \sref[0]{0.3.2.5}.
Como la imagen canónica de $f$ en $A_f$ es invertible, el homomorfismo de restricción $M=\Gamma(V,\mathscr{F})\to\Gamma(D(f)\cap V,\mathscr{F})$ se factoriza como $M\longrightarrow M_f\xrightarrow{u_f}\Gamma(D(f)\cap V,\mathscr{F})$ \sref[0]{0.1.2.4}, y la verificación de estas condiciones de compatibilidad para $D(g)\subset D(f)$ es inmediata.
Dicho esto, demostremos que la condición (d'1) (resp. (d'2)) implica que cada $u_f$ es sobreyectivo (resp. inyectivo), lo que demuestra que $u$ es \emph{biyectivo}, y por consiguiente que $\mathscr{F}$ es la restricción a $V$ de un $\widetilde{A}$-módulo isomorfo a $\widetilde{M}$.
Si $s\in\Gamma(D(f)\cap V,\mathscr{F})$, existe, por (d'1), un entero $n\geq0$ tal que $f^n s$ se prolonga a una sección $z\in M$;
se tiene entonces $u_f(z/f^n)=s$, de modo que $u_f$ es sobreyectivo.
Análogamente, si $z\in M$ es tal que $u_f(z/1)=0$, esto significa que la restricción a $D(f)\cap V$ de la sección $z$ es cero;
según (d'2), existe un entero $n\geq0$ tal que $f^n z=0$, por tanto $z/1=0$ en $M_f$, y así $u_f$ es inyectivo.
\end{proof}

\begin{corollary}[1.4.2]
\label{I.1.4.2}
Todo haz cuasicoherente sobre un subconjunto abierto cuasicompacto de $X$ es inducido por un haz cuasicoherente sobre $X$.
\end{corollary}

\begin{corollary}[1.4.3]
\label{I.1.4.3}
Toda $\mathscr{O}_X$-álgebra cuasicoherente sobre $X=\operatorname{Spec}(A)$ es isomorfa a una $\mathscr{O}_X$-álgebra de la forma $\widetilde{B}$, donde $B$ es un álgebra sobre $A$;
todo $\widetilde{B}$-módulo cuasicoherente es isomorfo a un $\widetilde{B}$-módulo de la forma $\widetilde{N}$, donde $N$ es un $B$-módulo.
\end{corollary}

\begin{proof}
En efecto, una $\mathscr{O}_X$-álgebra cuasicoherente es un $\mathscr{O}_X$-módulo cuasicoherente, y por tanto es de la forma $\widetilde{B}$, donde $B$ es un $A$-módulo;
el hecho de que $B$ sea un $A$-álgebra resulta de la caracterización de la estructura de $\mathscr{O}_X$-álgebra mediante el homomorfismo $\widetilde{B}\otimes_{\widetilde{A}}\widetilde{B}\to\widetilde{B}$ de $\widetilde{A}$-módulos, así como del Corolario~\sref{I.1.3.12}.
Si $\mathscr{G}$ es un $\widetilde{B}$-módulo cuasicoherente, basta demostrar que $\mathscr{G}$ también es un $\widetilde{A}$-módulo cuasicoherente y concluir luego del mismo modo;
como la cuestión es local, se puede, restringiéndose a un subconjunto abierto de $X$ de la forma $D(f)$, suponer que $\mathscr{G}$ es el conúcleo de un homomorfismo $\widetilde{B}^{(I)}\to\widetilde{B}^{(J)}$ de $\widetilde{B}$-módulos (y \emph{a fortiori} de $\widetilde{A}$-módulos);
la proposición resulta entonces de los Corolarios~\sref{I.1.3.8} y \sref{I.1.3.9}.
\end{proof}

\subsection{Haces coherentes sobre un espectro primo}
\label{subsection:I.1.5}

\begin{theorem}[1.5.1]
\label{I.1.5.1}
Sean $A$ un anillo \emph{noetheriano}, $X=\operatorname{Spec}(A)$ su espectro primo, $V$ un abierto de $X$ y $\mathscr{F}$ un $(\mathscr{O}_X|V)$-módulo.
Las condiciones siguientes son equivalentes.
\begin{enumerate}
  \item[\textup{a)}] $\mathscr{F}$ es coherente.
  \item[\textup{b)}] $\mathscr{F}$ es de tipo finito y cuasicoherente.
  \item[\textup{c)}] Existe un $A$-módulo $M$ de tipo finito tal que $\mathscr{F}$ sea isomorfo al haz $\widetilde{M}|V$.
\end{enumerate}
\end{theorem}

\begin{proof}
\oldpage[I]{93}
(a) implica trivialmente (b).
Para ver que (b) implica (c), observemos primero que, como $V$ es cuasicompacto \sref[0]{0.2.2.3}, $\mathscr{F}$ es isomorfo a un haz $\widetilde{N}|V$, donde $N$ es un $A$-módulo \sref{I.1.4.1}.
Se tiene $N=\varinjlim M_\lambda$, donde $M_\lambda$ recorre el conjunto de los $A$-submódulos de $N$ de tipo finito; por tanto \sref{I.1.3.9} $\mathscr{F}=\widetilde{N}|V=\varinjlim\widetilde{M}_\lambda|V$;
pero, como $\mathscr{F}$ es de tipo finito y $V$ es cuasicompacto, existe un índice $\lambda$ tal que $\mathscr{F}=\widetilde{M}_\lambda|V$ \sref[0]{0.5.2.3}.

Por último, demostremos que (c) implica (a).
Está claro que $\mathscr{F}$ es entonces de tipo finito (\sref{I.1.3.6} y \sref{I.1.3.9});
además, como la cuestión es local, se puede limitar al caso en que $V=D(f)$, $f\in A$.
Como $A_f$ es noetheriano, se ve que basta demostrar que el núcleo de un homomorfismo $\widetilde{A}^n\to\widetilde{M}$, donde $M$ es un $A$-módulo, es de tipo finito.
Pero tal homomorfismo es de la forma $\widetilde{u}$, donde $u$ es un homomorfismo $A^n\to M$ \sref{I.1.3.8}, y, si $P=\operatorname{Ker}u$, entonces $\widetilde{P}=\operatorname{Ker}\widetilde{u}$ \sref{I.1.3.9}.
Como $A$ es noetheriano, $P$ es de tipo finito, lo que termina la demostración.
\end{proof}

\begin{corollary}[1.5.2]
\label{I.1.5.2}
Bajo las hipótesis de \sref{I.1.5.1}, el haz $\mathscr{O}_X$ es un haz coherente de anillos.
\end{corollary}

\begin{corollary}[1.5.3]
\label{I.1.5.3}
Bajo las hipótesis de \sref{I.1.5.1}, todo haz coherente sobre un abierto de $X$ es inducido por un haz coherente sobre $X$.
\end{corollary}

\begin{corollary}[1.5.4]
\label{I.1.5.4}
Bajo las hipótesis de \sref{I.1.5.1}, todo $\mathscr{O}_X$-módulo cuasicoherente $\mathscr{F}$ es el límite inductivo de los $\mathscr{O}_X$-submódulos coherentes de $\mathscr{F}$.
\end{corollary}

\begin{proof}
En efecto, $\mathscr{F}=\widetilde{M}$, donde $M$ es un $A$-módulo, y $M$ es el límite inductivo de sus submódulos de tipo finito;
la demostración concluye aplicando \sref{I.1.3.9} y \sref{I.1.5.1}.
\end{proof}

\subsection{Propiedades funtoriales de los haces cuasicoherentes sobre un espectro primo}
\label{subsection:I.1.6}

\begin{env}[1.6.1]
\label{I.1.6.1}
Sean $A$, $A'$ dos anillos,
\[
  \varphi:A'\to A
\]
un homomorfismo, y
\[
  {}^a\varphi:X=\operatorname{Spec}(A)\to
  X'=\operatorname{Spec}(A')
\]
la aplicación continua asociada a $\varphi$ \sref{I.1.2.1}.
Definiremos un \emph{homomorfismo canónico}
\[
  \widetilde{\varphi}:\mathscr{O}_{X'}\to
  {}^a\varphi_*(\mathscr{O}_X)
\]
de haces de anillos.
Para todo $f'\in A'$, pongamos $f=\varphi(f')$;
se tiene ${}^a\varphi^{-1}(D(f'))=D(f)$ \sref{I.1.2.2.2}.
Los anillos $\Gamma(D(f'),\widetilde{A}')$ y $\Gamma(D(f),\widetilde{A})$ se identifican respectivamente con $A'_{f'}$ y $A_f$ (\sref{I.1.3.6} y \sref{I.1.3.7}). El homomorfismo $\varphi$ define canónicamente un homomorfismo $\varphi_{f'}:A'_{f'}\to A_f$ \sref[0]{0.1.5.1}; en otras palabras, se tiene un homomorfismo de anillos
\[
  \Gamma(D(f'),\widetilde{A}')\to
  \Gamma({}^a\varphi^{-1}(D(f')),\widetilde{A})
  =\Gamma(D(f'),{}^a\varphi_*(\widetilde{A}))
\]
\oldpage[I]{94}
Además, estos homomorfismos satisfacen las condiciones usuales de compatibilidad: para $D(f')\supset D(g')$, el diagrama
\[
  \xymatrix{
    \Gamma(D(f'),\widetilde{A}')\ar[r]\ar[d] &
    \Gamma(D(f'),{}^a\varphi_*(\widetilde{A}))\ar[d]\\
    \Gamma(D(g'),\widetilde{A}')\ar[r] &
    \Gamma(D(g'),{}^a\varphi_*(\widetilde{A}))
  }
\]
es conmutativo \sref[0]{0.1.5.1};
se ha definido así un homomorfismo de $\mathscr{O}_{X'}$-álgebras, pues los $D(f')$ forman una base de la topología de $X'$ \sref[0]{0.3.2.3}.
El par $\Phi=({}^a\varphi,\widetilde{\varphi})$ es, por tanto, un \emph{morfismo} de espacios anillados
\[
  \Phi:(X,\mathscr{O}_X)\to(X',\mathscr{O}_{X'})
\]
\sref[0]{0.4.1.1}.

Notemos además que, si se pone $x'={}^a\varphi(x)$, entonces el homomorfismo $\widetilde{\varphi}_x^{\sharp}$ \sref[0]{0.3.7.1} no es otro que el homomorfismo
\[
  \varphi_x:A'_{x'}\to A_x
\]
inducido canónicamente por $\varphi:A'\to A$ \sref[0]{0.1.5.1}.
En efecto, todo $z'\in A'_{x'}$ puede escribirse como $g'/f'$, donde $f'$, $g'$ pertenecen a $A'$ y $f'\notin\mathfrak{j}_{x'}$;
$D(f')$ es entonces una vecindad de $x'$ en $X'$, y el homomorfismo $\Gamma(D(f'),\widetilde{A}')\to
\Gamma({}^a\varphi^{-1}(D(f')),\widetilde{A})$ inducido por $\widetilde{\varphi}$ no es otro que $\varphi_{f'}$;
al considerar la sección $s'\in\Gamma(D(f'),\widetilde{A}')$ correspondiente a $g'/f'\in A'_{f'}$, se obtiene $\widetilde{\varphi}_x^{\sharp}(z')=\varphi(g')/\varphi(f')$ en $A_x$.
\end{env}

\begin{example}[1.6.2]
\label{I.1.6.2}
Sean $S$ un conjunto multiplicativo de $A$ y $\varphi$ el homomorfismo canónico $A\to S^{-1}A$;
ya se ha visto \sref{I.1.2.6} que ${}^a\varphi$ es un \emph{homeomorfismo} de $Y=\operatorname{Spec}(S^{-1}A)$ sobre el subespacio de $X=\operatorname{Spec}(A)$ formado por los $x$ tales que $\mathfrak{j}_x\cap S=\varnothing$.
Además, para todo $x$ de este subespacio, que es por tanto de la forma ${}^a\varphi(y)$ con $y\in Y$, el homomorfismo $\widetilde{\varphi}_y^{\sharp}:\mathscr{O}_x\to\mathscr{O}_y$ es
\emph{biyectivo} \sref[0]{0.1.2.6};
en otras palabras, $\mathscr{O}_Y$ se identifica con el haz inducido sobre $Y$ por $\mathscr{O}_X$.
\end{example}

\begin{proposition}[1.6.3]
\label{I.1.6.3}
Para todo $A$-módulo $M$, existe un isomorfismo canónico funtorial del $\mathscr{O}_{X'}$-módulo $(M_{[\varphi]})^{\sim}$ con la imagen directa $\Phi_*(\widetilde{M})$.
\end{proposition}

\begin{proof}
Para abreviar, pongamos $M'=M_{[\varphi]}$ y, para todo $f'\in A'$, pongamos asimismo $f=\varphi(f')$.
Los módulos de secciones $\Gamma(D(f'),\widetilde{M}')$ y $\Gamma(D(f),\widetilde{M})$ se identifican respectivamente con los módulos $M'_{f'}$ y $M_f$ (sobre $A'_{f'}$ y $A_f$, respectivamente);
además, el $A'_{f'}$-módulo $(M_f)_{[\varphi_{f'}]}$ es canónicamente isomorfo a $M'_{f'}$ \sref[0]{0.1.5.2}.
Se tiene así un isomorfismo funtorial de $\Gamma(D(f'),\widetilde{A}')$-módulos: $\Gamma(D(f'),\widetilde{M}')\xrightarrow{\sim}
\Gamma({}^a\varphi^{-1}(D(f')),\widetilde{M})_{[\varphi_{f'}]}$, y estos isomorfismos satisfacen las condiciones usuales de compatibilidad con las restricciones \sref[0]{0.1.5.6}, por lo que definen el isomorfismo funtorial buscado.
Más precisamente, si $u:M_1\to M_2$ es un homomorfismo de $A$-módulos, puede considerarse como un homomorfismo $(M_1)_{[\varphi]}\to(M_2)_{[\varphi]}$ de $A'$-módulos;
si se denota este homomorfismo por $u_{[\varphi]}$, entonces $\Phi_*(\widetilde{u})$ se identifica con $(u_{[\varphi]})^{\sim}$.
\end{proof}

Esta demostración prueba asimismo que, para toda $A$-\emph{álgebra} $B$, el isomorfismo canónico funtorial
\oldpage[I]{95}
$(B_{[\varphi]})^{\sim}\to\Phi_*(\widetilde{B})$ es un isomorfismo de $\mathscr{O}_{X'}$-\emph{álgebras};
si $M$ es un $B$-módulo, el isomorfismo canónico funtorial $(M_{[\varphi]})^{\sim}\xrightarrow{\sim}\Phi_*(\widetilde{M})$ es un isomorfismo de $\Phi_*(\widetilde{B})$-módulos.

\begin{corollary}[1.6.4]
\label{I.1.6.4}
El funtor imagen directa $\Phi_*$ es exacto en la categoría de los $\mathscr{O}_X$-módulos cuasicoherentes.
\end{corollary}

\begin{proof}
En efecto, es claro que $M_{[\varphi]}$ es un funtor exacto en $M$ y que $\widetilde{M}'$ es un
funtor exacto en $M'$ \sref{I.1.3.5}.
\end{proof}

\begin{proposition}[1.6.5]
\label{I.1.6.5}
Sea $N'$ un $A'$-módulo y sea $N$ el $A$-módulo $N'\otimes_{A'}A_{[\varphi]}$;
entonces existe un isomorfismo canónico funtorial del $\mathscr{O}_X$-módulo $\Phi^*(\widetilde{N}')$ con $\widetilde{N}$.
\end{proposition}

\begin{proof}
Observemos primero que $j:z'\to z'\otimes 1$ es un $A'$-homomorfismo de $N'$ en $N_{[\varphi]}$: en efecto, por definición, para $f'\in A'$, se tiene $(f'z')\otimes 1=z'\otimes\varphi(f')=\varphi(f')(z'\otimes 1)$.
Se deduce de \sref{I.1.3.8} que existe un homomorfismo $\widetilde{j}:\widetilde{N}'\to(N_{[\varphi]})^{\sim}$ de $\mathscr{O}_{X'}$-módulos, y por \sref{I.1.6.3} se puede considerar que $\widetilde{j}$ aplica $\widetilde{N}'$ en $\Phi_*(\widetilde{N})$.
A este homomorfismo $\widetilde{j}$ le corresponde canónicamente un homomorfismo $h=\widetilde{j}^{\sharp}$ de $\Phi^*(\widetilde{N}')$ en $\widetilde{N}$ \sref[0]{0.4.4.3};
se verá que, en cada tallo, $h_x$ es \emph{biyectivo}.
Pongamos $x'={}^a\varphi(x)$ y sea $f'\in A'$ tal que $x'\in D(f')$;
sea $f=\varphi(f')$.
El anillo $\Gamma(D(f),\widetilde{A})$ se identifica con $A_f$, y los módulos
$\Gamma(D(f),\widetilde{N})$ y $\Gamma(D(f'),\widetilde{N}')$ se identifican con $N_f$ y $N'_{f'}$, respectivamente;
sea una sección $s'\in\Gamma(D(f'),\widetilde{N}')$, identificada con $n'/{f'}^p$ ($n'\in N'$), y sea $s$ su imagen por $\widetilde{j}$ en $\Gamma(D(f),\widetilde{N})$;
$s$ se identifica con $(n'\otimes 1)/f^p$.
Por otra parte, sea una sección $t\in\Gamma(D(f),\widetilde{A})$, identificada con $g/f^q$ ($g\in A$);
entonces, por definición, se tiene $h_x(s'_{x'}\otimes t_x)=t_x\cdot s_x$ \sref[0]{0.4.4.3}.
Pero se puede identificar canónicamente $N_f$ con $N'_{f'}\otimes_{A'_{f'}}(A_f)_{[\varphi_{f'}]}$ \sref[0]{0.1.5.4};
$s$ corresponde entonces al elemento $(n'/{f'}^p)\otimes 1$, y la sección $y\to t_y\cdot s_y$ corresponde a $(n'/{f'}^p)\otimes(g/f^q)$.
Los diagramas de compatibilidad de \sref[0]{0.1.5.6} muestran que $h_x$ no es otro que el isomorfismo canónico
\[
  \label{I.1.6.5.1-fr}
  N'_{x'}\otimes_{A'_{x'}}(A_x)_{[\varphi_x]}
  \xrightarrow{\sim}N_x=(N'\otimes_{A'}A_{[\varphi]})_x.
  \tag{1.6.5.1}
\]

Además, sea $v:N'_1\to N'_2$ un homomorfismo de $A'$-módulos;
como $\widetilde{v}_{x'}=v_{x'}$ para todo $x'\in X'$, se deduce inmediatamente de lo anterior que $\Phi^*(\widetilde{v})$ se identifica canónicamente con $(v\otimes 1)^{\sim}$, lo que termina la demostración de \sref{I.1.6.5}.
\end{proof}

Si $B'$ es una $A'$-álgebra, el isomorfismo canónico de $\Phi^*(\widetilde{B}')$ con $(B'\otimes_{A'}A_{[\varphi]})^{\sim}$ es un isomorfismo de $\mathscr{O}_X$-álgebras;
si, además, $N'$ es un $B'$-módulo, entonces el isomorfismo canónico de $\Phi^*(\widetilde{N}')$ con $(N'\otimes_{A'}A_{[\varphi]})^{\sim}$ es un isomorfismo de $\Phi^*(\widetilde{B}')$-módulos.

\begin{corollary}[1.6.6]
\label{I.1.6.6}
Las secciones de $\Phi^*(\widetilde{N}')$ que son imágenes canónicas de las secciones $s'$, donde $s'$ recorre el $A'$-módulo $\Gamma(\widetilde{N}')$, generan el $A$-módulo $\Gamma(\Phi^*(\widetilde{N'}))$.
\end{corollary}

\begin{proof}
En efecto, estas imágenes se identifican con los elementos $z'\otimes 1$ de $N$, cuando se identifican $N'$ y $N$ con $\Gamma(\widetilde{N}')$ y $\Gamma(\widetilde{N})$, respectivamente \sref{I.1.3.7}, y $z'$ recorre $N'$.
\end{proof}

\begin{env}[1.6.7]
\label{I.1.6.7}
En la demostración de \sref{I.1.6.5}, se ha probado de paso que la aplicación canónica (\textbf{0},~4.4.3.2) $\rho:\widetilde{N}'\to\Phi_*(\Phi^*(\widetilde{N}'))$ no es otra que el homomorfismo $\widetilde{j}$,
\oldpage[I]{96}
donde $j:N'\to N'\otimes_{A'}A_{[\varphi]}$ es el homomorfismo $z'\to z'\otimes 1$.
Del mismo modo, la aplicación canónica (\textbf{0},~4.4.3.3) $\sigma:\Phi^*(\Phi_*(\widetilde{M}))\to\widetilde{M}$ no es otra que $\widetilde{p}$, donde $p:M_{[\varphi]}\otimes_{A'}A_{[\varphi]}\to M$ es el homomorfismo canónico que envía cada producto tensorial $z\otimes a$ ($z\in M$, $a\in A$) a $a\cdot z$;
esto se deduce inmediatamente de las definiciones (\EGAUnavailableHyperref{0.3.7.1}{\textbf{0},~3.7.1}, \EGAUnavailableHyperref{0.4.4.3}{\textbf{0},~4.4.3} y \hyperref[I.1.3.7]{\textbf{I},~1.3.7}).

Se concluye (\textbf{0},~\EGAUnavailableHyperref{0.4.4.3}{4.4.3} y \EGAUnavailableHyperref{0.3.5.4.4}{3.5.4.4}) que, si $v:N'\to M_{[\varphi]}$ es un $A'$-homomorfismo, entonces $\widetilde{v}^{\sharp}=(v\otimes 1)^{\sim}$.
\end{env}

\begin{env}[1.6.8]
\label{I.1.6.8}
Sean $N'_1$ y $N'_2$ dos $A'$-módulos, y supongamos que $N'_1$ admite una \emph{presentación finita};
se deduce entonces de \sref{I.1.6.7} y \sref{I.1.3.12}[ii] que el
homomorfismo canónico \sref[0]{0.4.4.6}
\[
  \Phi^*\bigl(\mathscr{H}\!om_{\widetilde{A}'}
    (\widetilde{N}'_1,\widetilde{N}'_2)\bigr)
  \longrightarrow
  \mathscr{H}\!om_{\widetilde A}
    \bigl(\Phi^*(\widetilde{N}'_1),\Phi^*(\widetilde{N}'_2)\bigr)
\]
no es otro que $\widetilde{\gamma}$, donde $\gamma$ denota el homomorfismo canónico de $A$-módulos \[
  \operatorname{Hom}_{A'}(N'_1,N'_2)\otimes_{A'}A
  \longrightarrow
  \operatorname{Hom}_A(N'_1\otimes_{A'}A,N'_2\otimes_{A'}A).
\].
\end{env}

\begin{env}[1.6.9]
\label{I.1.6.9}
Sean $\mathfrak{J}'$ un ideal de $A'$ y $M$ un $A$-módulo;
como, por definición, $\widetilde{\mathfrak{J}}'\widetilde M$ es la imagen del homomorfismo canónico $\Phi^*(\widetilde{\mathfrak{J}}')\otimes_{\widetilde A}\widetilde M
\to\widetilde M$, se
deduce de la Proposición~\sref{I.1.6.5} y del Corolario~\sref{I.1.3.12}[i] que
$\widetilde{\mathfrak{J}}'\widetilde M$ se identifica canónicamente con $(\mathfrak{J}'M)^{\sim}$;
en particular, $\Phi^*(\widetilde{\mathfrak{J}}')\widetilde A$ se identifica con $(\mathfrak{J}'A)^{\sim}$ y, teniendo en cuenta la exactitud por la derecha del funtor $\Phi^*$,
la $\widetilde A$-álgebra $\Phi^*((A'/\mathfrak{J}')^{\sim})$ se identifica con $(A/\mathfrak{J}'A)^{\sim}$.
\end{env}

\begin{env}[1.6.10]
\label{I.1.6.10}
Sea $A''$ un tercer anillo, sea $\varphi'$ un homomorfismo $A''\to A'$ y pongamos $\varphi''=\varphi\circ\varphi'$.
Se deduce inmediatamente de las definiciones que ${}^a\varphi''=({}^a\varphi')\circ({}^a\varphi)$ y $\widetilde{\varphi}''=\widetilde\varphi\circ\widetilde{\varphi}'$ \sref[0]{0.1.5.7}. Se concluye que $\Phi''=\Phi'\circ\Phi$;
en otras palabras, $(\operatorname{Spec} A,\widetilde A)$ es un \emph{funtor} en $A$ de la categoría de los anillos a la de los espacios anillados.
\end{env}
	
\subsection{Caracterización de los morfismos de esquemas afines}
\label{subsection:I.1.7}

\begin{definition}[1.7.1]
\label{I.1.7.1}
Se dice que un espacio anillado $(X,\mathscr{O}_X)$ es un esquema afín si es isomorfo a un espacio anillado de la forma $(\operatorname{Spec}(A),\widetilde A)$, donde $A$ es un anillo;
se dice entonces que $\Gamma(X,\mathscr{O}_X)$, que se identifica canónicamente con el anillo $A$ \sref{I.1.3.7}, es el anillo del esquema afín $(X,\mathscr{O}_X)$, y se lo denota por $A(X)$ cuando no haya confusión.
\end{definition}

Por abuso de lenguaje, cuando hablemos del \emph{esquema afín $\operatorname{Spec}(A)$}, nos referiremos siempre al espacio anillado $(\operatorname{Spec}(A),\widetilde A)$.
\begin{env}[1.7.2]
\label{I.1.7.2}
Sean $A$ y $B$ dos anillos, y $(X,\mathscr{O}_X)$ y $(Y,\mathscr{O}_Y)$ los esquemas afines correspondientes a los espectros primos $X=\operatorname{Spec}(A)$, $Y=\operatorname{Spec}(B)$.
Se ha visto \sref{I.1.6.1} que a todo homomorfismo de anillos $\varphi:B\to A$ corresponde un morfismo $\Phi=({}^a\varphi,\widetilde\varphi)=\operatorname{Spec}(\varphi):
(X,\mathscr{O}_X)\to(Y,\mathscr{O}_Y)$.
Se observará que $\varphi$ está completamente determinado por $\Phi$, pues se tiene, por definición, \[
  \varphi=\Gamma(\widetilde\varphi):
  \Gamma(\widetilde B)\longrightarrow
  \Gamma({}^a\varphi_*(\widetilde A))=\Gamma(\widetilde A).
\].
\end{env}

\begin{theorem}[1.7.3]
\label{I.1.7.3}
\footnote{\emph{[Trad.] Véase \hyperref[subsection:I.1.8]{(1.8)} y la nota a pie de página correspondiente.}}
Sean $(X,\mathscr{O}_X)$ y $(Y,\mathscr{O}_Y)$ dos esquemas afines.
Para que un morfismo de espacios anillados $(\psi,\theta):(X,\mathscr{O}_X)\to(Y,\mathscr{O}_Y)$ sea de la forma $({}^a\varphi,\widetilde\varphi)$, donde $\varphi$ es un homomorfismo de anillos: $A(Y)\to A(X)$, es necesario y suficiente que, para todo $x\in X$, $\theta_x^{\sharp}$ sea un homomorfismo \textup{local}: $\mathscr{O}_{\psi(x)}\to\mathscr{O}_x$.
\end{theorem}

\begin{proof}
\oldpage[I]{97}
Pongamos $A=A(X)$, $B=A(Y)$.
La condición es necesaria, pues se ha visto \sref{I.1.6.1} que $\widetilde\varphi_x^{\sharp}$ es el homomorfismo de $B_{{}^a\varphi(x)}$ en $A_x$ inducido canónicamente por $\varphi$, y, por definición, ${}^a\varphi(x)=\varphi^{-1}(\mathfrak{j}_x)$, por lo que este homomorfismo es local.

Demostremos ahora que la condición es suficiente.
Por definición, $\theta$ es un homomorfismo $\mathscr{O}_Y\to\psi_*(\mathscr{O}_X)$, y se obtiene canónicamente un homomorfismo de anillos
\[
  \varphi=\Gamma(\theta):
  B=\Gamma(Y,\mathscr{O}_Y)
  \longrightarrow
  \Gamma(Y,\psi_*(\mathscr{O}_X))=\Gamma(X,\mathscr{O}_X)=A.
\]

La hipótesis sobre $\theta_x^{\sharp}$ significa que este homomorfismo induce, por paso a los cocientes, un monomorfismo $\theta^x$ del campo residual $\boldsymbol{k}(\psi(x))$ en el campo residual $\boldsymbol{k}(x)$, tal que, para toda sección $f\in\Gamma(Y,\mathscr{O}_Y)=B$, se tiene $\theta^x(f(\psi(x)))=\varphi(f)(x)$.
La relación $f(\psi(x))=0$ es, por tanto, equivalente a $\varphi(f)(x)=0$, lo que significa que $\mathfrak{j}_{\psi(x)}=\mathfrak{j}_{{}^a\varphi(x)}$; equivalentemente, $\psi(x)={}^a\varphi(x)$ para todo $x\in X$, o bien $\psi={}^a\varphi$.
Se sabe también que el diagrama
\phantomsection\label{I.1.7.3.diagram-fr}
\[
  \xymatrix{
    B=\Gamma(Y,\mathscr{O}_Y)\ar[r]^\varphi\ar[d] &
    \Gamma(X,\mathscr{O}_X)=A\ar[d]\\
    B_{\psi(x)}\ar[r]_{\theta_x^{\sharp}} & A_x
  }
\]
es conmutativo \sref[0]{0.3.7.2}, lo que significa que $\theta_x^{\sharp}$ es igual al homomorfismo $\varphi_x:B_{\psi(x)}\to A_x$ inducido canónicamente por $\varphi$ \sref[0]{0.1.5.1}.
Como los homomorfismos $\theta_x^{\sharp}$ determinan por completo $\theta^{\sharp}$, y por tanto $\theta$ \sref[0]{0.3.7.1}, se concluye que $\theta=\widetilde\varphi$, por la definición de $\widetilde\varphi$ \sref{I.1.6.1}.
\end{proof}

Decimos que un morfismo $(\psi,\theta)$ de espacios anillados que satisface la condición de \sref{I.1.7.3} es un \emph{morfismo de esquemas afines}.

\begin{corollary}[1.7.4]
\label{I.1.7.4}
Si $(X,\mathscr{O}_X)$ y $(Y,\mathscr{O}_Y)$ son esquemas afines, existe un isomorfismo canónico del conjunto de morfismos de esquemas afines $\operatorname{Hom}((X,\mathscr{O}_X),(Y,\mathscr{O}_Y))$ al conjunto de homomorfismos de anillos de $B$ en $A$, donde $A=\Gamma(\mathscr{O}_X)$ y $B=\Gamma(\mathscr{O}_Y)$.
\end{corollary}

También se puede decir que los funtores $(\operatorname{Spec}(A),\widetilde A)$ en $A$ y $\Gamma(X,\mathscr{O}_X)$ en $(X,\mathscr{O}_X)$ definen una \emph{equivalencia} entre la categoría de los anillos conmutativos y la categoría opuesta de los esquemas afines (T, I, 1.2).

\begin{corollary}[1.7.5]
\label{I.1.7.5}
Si $\varphi:B\to A$ es sobreyectivo, el morfismo correspondiente $({}^a\varphi,\widetilde\varphi)$ es un monomorfismo de espacios anillados \emph{(cf.~\sref{I.4.1.7})}.
\end{corollary}

\begin{proof}
En efecto, se sabe que ${}^a\varphi$ es inyectivo \sref{I.1.2.5}, y, puesto que $\varphi$ es
sobreyectivo, para todo $x\in X$, $\varphi_x^{\sharp}:B_{{}^a\varphi(x)}\to A_x$, que se obtiene de $\varphi$ por paso a los anillos de fracciones, es también sobreyectivo \sref[0]{0.1.5.1};
de ahí la conclusión \sref[0]{0.4.1.1}.
\end{proof}

\subsection{Morfismos de espacios localmente anillados a esquemas afines.}
\label{subsection:I.1.8}

\oldpage[II]{217}
Debido a una observación de J.~Tate, los enunciados del Teorema~\sref{I.1.7.3} y de la Proposición~\sref{I.2.2.4} pueden generalizarse como sigue:\footnote{\emph{[Trad.] La sección siguiente (I.1.8) fue añadida en las erratas de EGA~II, de ahí el cambio temporal en los números de página, que remiten a EGA~II.}}

\begin{proposition}[1.8.1]
\label{I.1.8.1}
Sean $(S,\sh{O}_S)$ un esquema afín y $(X,\sh{O}_X)$ un espacio localmente anillado.
Entonces existe una biyección canónica del conjunto de homomorfismos de anillos
\oldpage[II]{218}
$\Gamma(S,\sh{O}_S)\to\Gamma(X,\sh{O}_X)$ sobre el conjunto de morfismos de espacios anillados $(\psi,\theta):(X,\sh{O}_X)\to(S,\sh{O}_S)$ tales que, para cada $x\in X$, $\theta_x^\sharp$ sea un homomorfismo local $\sh{O}_{\psi(x)}\to\sh{O}_x$.
\end{proposition}

\begin{proof}
Observemos ante todo que, si $(X,\sh{O}_X)$ y $(S,\sh{O}_S)$ son dos espacios anillados cualesquiera, un morfismo $(\psi,\theta)$ de $(X,\sh{O}_X)$ a $(S,\sh{O}_S)$ define canónicamente un homomorfismo de anillos $\Gamma(\theta):\Gamma(S,\sh{O}_S)\to\Gamma(X,\sh{O}_X)$, y por tanto una primera aplicación
\[
  \label{I.1.8.1.1}
  \rho:\Hom((X,\sh{O}_X),(S,\sh{O}_S))\to\Hom(\Gamma(S,\sh{O}_S),\Gamma(X,\sh{O}_X)).
  \tag{1.8.1.1}
\]
Recíprocamente, bajo las hipótesis del enunciado, pongamos $A=\Gamma(S,\sh{O}_S)$ y consideremos un homomorfismo de anillos $\vphi:A\to\Gamma(X,\sh{O}_X)$.
Para cada $x\in X$, está claro que el conjunto de los $f\in A$ tales que $\vphi(f)(x)=0$ es un \emph{ideal primo} de $A$, pues $\sh{O}_x/\mathfrak{m}_x=\kres(x)$ es un campo;
es por tanto un elemento de $S=\Spec(A)$, que denotaremos por ${}^a\vphi(x)$.
  Además, para cada $f\in A$, se tiene, por definición \sref[0]{0.5.5.2}, que ${}^a\vphi^{-1}(D(f))=X_f$, lo que demuestra que ${}^a\vphi$ es una \emph{aplicación continua} $X\to S$.
Definimos entonces un homomorfismo
\[
  \widetilde{\vphi}:\sh{O}_S\to{}^a\vphi_*(\sh{O}_X)
\]
de $\sh{O}_S$-módulos;
para cada $f\in A$, se tiene $\Gamma(D(f),\sh{O}_S)=A_f$ \sref{I.1.3.6};
para cada $s\in A$, hacemos corresponder al elemento $s/f\in A_f$ el elemento $(\vphi(s)|X_f)(\vphi(f)|X_f)^{-1}$ de $\Gamma(X_f,\sh{O}_X)=\Gamma(D(f),{}^a\vphi_*(\sh{O}_X))$, y se comprueba inmediatamente (pasando de $D(f)$ a $D(fg)$) que esto define efectivamente un homomorfismo de $\sh{O}_S$-módulos, y por tanto un morfismo $({}^a\vphi,\widetilde{\vphi})$ de espacios anillados.
Además, con las mismas notaciones, y poniendo para abreviar $y={}^a\vphi(x)$, se ve inmediatamente \sref[0]{0.3.7.1} que $\widetilde{\vphi}_x^\sharp(s_y/f_y)=(\vphi(s)_x)(\vphi(f)_x)^{-1}$;
como la relación $s_y\in\mathfrak{m}_y$ equivale, por definición, a $\vphi(s)_x\in\mathfrak{m}_x$, vemos que $\widetilde{\vphi}_x^\sharp$ es un homomorfismo \emph{local} $\sh{O}_y\to\sh{O}_x$, y hemos definido así una segunda aplicación $\sigma:\Hom(\Gamma(S,\sh{O}_S),\Gamma(X,\sh{O}_X))\to\mathfrak{L}$, donde $\mathfrak{L}$ es el conjunto de los morfismos $(\psi,\theta):(X,\sh{O}_X)\to(S,\sh{O}_S)$ tales que $\theta_x^\sharp$ es local para cada $x\in X$.
Queda por demostrar que $\sigma$ y $\rho$ (restringida a $\mathfrak{L}$) son inversas una de otra;
la definición de $\widetilde{\vphi}$ muestra inmediatamente que $\Gamma(\widetilde{\vphi})=\vphi$ y, por consiguiente, que $\rho\circ\sigma$ es la identidad.
Para ver que $\sigma\circ\rho$ es la identidad, partamos de un morfismo $(\psi,\theta)\in\mathfrak{L}$ y pongamos $\vphi=\Gamma(\theta)$;
las hipótesis sobre $\theta_x^\sharp$ significan que este morfismo induce, al pasar a los cocientes, un monomorfismo $\theta^x:\kres(\psi(x))\to\kres(x)$ tal que, para toda sección $f\in A=\Gamma(S,\sh{O}_S)$, se tiene $\theta^x(f(\psi(x)))=\vphi(f)(x)$;
la igualdad $f(\psi(x))=0$ equivale por tanto a $\vphi(f)(x)=0$, lo que demuestra que ${}^a\vphi=\psi$.
Por otra parte, de las definiciones se deduce que el diagrama
\[
  \xymatrix{
    A\ar[r]^\vphi\ar[d] &
    \Gamma(X,\sh{O}_X)\ar[d]\\
    A_{\psi(x)}\ar[r]_{\theta_x^\sharp} &
    \sh{O}_x
  }
\]
es conmutativo, y lo mismo sucede con el diagrama análogo en el que $\theta_x^\sharp$ se sustituye por $\widetilde{\vphi}_x^\sharp$, de donde $\widetilde{\vphi}_x^\sharp=\theta_x^\sharp$ \sref[0]{0.1.2.4} y, por consiguiente, $\widetilde{\vphi}=\theta$.
\end{proof}

\begin{env}[1.8.2]
\label{I.1.8.2}
Cuando $(X,\sh{O}_X)$ e $(Y,\sh{O}_Y)$ sean espacios \emph{localmente} anillados, consideraremos los morfismos $(\psi,\theta):(X,\sh{O}_X)\to(Y,\sh{O}_Y)$ tales que, para cada $x\in X$, $\theta_x^\sharp$ sea un homomorfismo \emph{local} $\sh{O}_{\psi(x)}\to\sh{O}_x$.
En adelante, cuando hablemos
\oldpage[II]{219}
de un \emph{morfismo de espacios localmente anillados}, se tratará siempre de un morfismo como el anterior;
con esta definición de los morfismos, está claro que los espacios localmente anillados forman una \emph{categoría};
para cualesquiera dos objetos $X$ e $Y$ de esta categoría, $\Hom(X,Y)$ denota así el conjunto de morfismos de espacios localmente anillados de $X$ a $Y$ (el conjunto denotado por $\mathfrak{L}$ en \sref{I.1.8.1});
cuando consideremos el conjunto de \emph{morfismos de espacios anillados} de $X$ a $Y$, lo denotaremos por $\Hom_\text{rs}(X,Y)$ para evitar toda confusión.
La aplicación \sref{I.1.8.1.1} se escribe entonces
\[
  \label{I.1.8.2.1}
  \rho:\Hom_\text{rs}(X,Y)\to\Hom(\Gamma(Y,\sh{O}_Y),\Gamma(X,\sh{O}_X))
  \tag{1.8.2.1}
\]
y su restricción
\[
  \label{I.1.8.2.2}
  \rho':\Hom(X,Y)\to\Hom(\Gamma(Y,\sh{O}_Y),\Gamma(X,\sh{O}_X))
  \tag{1.8.2.2}
\]
es una aplicación \emph{funtorial} en $X$ e $Y$ en la categoría de los espacios localmente anillados.
\end{env}

\begin{corollary}[1.8.3]
\label{I.1.8.3}
Sea $(Y,\sh{O}_Y)$ un espacio localmente anillado.
Para que $Y$ sea un esquema afín, es necesario y suficiente que, para cada espacio localmente anillado $(X,\sh{O}_X)$, la aplicación \sref{I.1.8.2.2} sea biyectiva.
\end{corollary}

\begin{proof}
La Proposición~\sref{I.1.8.1} demuestra que la condición es necesaria.
Recíprocamente, si suponemos que la condición se satisface y ponemos $A=\Gamma(Y,\sh{O}_Y)$, resulta de las hipótesis y de \sref{I.1.8.1} que los funtores $X\to\Hom(X,Y)$ y $X\to\Hom(X,\Spec(A))$, de la categoría de los espacios localmente anillados a la de los conjuntos, son \emph{isomorfos}.
Sabemos que esto implica la existencia de un isomorfismo canónico $Y\to\Spec(A)$ (cf.~\textbf{0},~8).\footnote{\emph{[Trad.] El texto francés imprime aquí $X\to\Spec(A)$; el funtor representado y la conclusión exigen $Y\to\Spec(A)$.}}
\end{proof}

\begin{env}[1.8.4]
\label{I.1.8.4}
Sea $S=\Spec(A)$ un esquema afín;
denotemos por $(S',A')$ el espacio anillado cuyo espacio subyacente está \emph{reducido a un punto} y cuyo haz estructural $A'$ es el haz (necesariamente simple) sobre $S'$ definido por el anillo $A$.
Sea $\pi:S\to S'$ la única aplicación de $S$ a $S'$;
por otra parte, observemos que, para cada abierto $U$ de $S$, existe una aplicación canónica $\Gamma(S',A')=\Gamma(S,\sh{O}_S)\to\Gamma(U,\sh{O}_S)$ que define así un \emph{$\pi$-morfismo} $\iota:A'\to\sh{O}_S$ de haces de anillos.
Hemos definido de este modo canónicamente un \emph{morfismo de espacios anillados $i=(\pi,\iota):(S,\sh{O}_S)\to(S',A')$}.
Para cada $A$-módulo $M$, denotamos por $M'$ el haz simple sobre $S'$ definido por $M$, que es evidentemente un $A'$-módulo.
Está claro que $i_*(\widetilde{M})=M'$ \sref{I.1.3.7}.
\end{env}

\begin{lemma}[1.8.5]
\label{I.1.8.5}
Con las notaciones de \sref{I.1.8.4}, para cada $A$-módulo $M$, el homomorfismo canónico funtorial de $\sh{O}_S$-módulos \sref[0]{0.4.4.3.3}
\[
\label{I.1.8.5.1}
  i^*(i_*(\widetilde{M}))\to\widetilde{M}
  \tag{1.8.5.1}
\]
es un isomorfismo.
\end{lemma}

\begin{proof}
En efecto, los dos miembros de \sref{I.1.8.5.1} son exactos por la derecha (pues el funtor $M\mapsto i_*(\widetilde{M})$ es evidentemente exacto) y conmutan con las sumas directas;
considerando $M$ como el conúcleo de un homomorfismo $A^{(I)}\to A^{(J)}$, podemos reducirnos a demostrar el lema en el caso $M=A$, y en este caso es evidente.
\end{proof}

\begin{corollary}[1.8.6]
\label{I.1.8.6}
Sean $(X,\sh{O}_X)$ un espacio anillado y $u:X\to S$ un morfismo de espacios anillados.
\oldpage[II]{220}
Para cada $A$-módulo $M$, se tiene (con las notaciones de \sref{I.1.8.4}) un isomorfismo canónico funtorial de $\sh{O}_X$-módulos
\[
\label{I.1.8.6.1}
  u^*(\widetilde{M})\isoto u^*(i^*(M'))
  \tag{1.8.6.1}
\]
\end{corollary}

\begin{corollary}[1.8.7]
\label{I.1.8.7}
Bajo las hipótesis de \sref{I.1.8.6}, se tiene, para cada $A$-módulo $M$ y cada $\sh{O}_X$-módulo $\sh{F}$, un isomorfismo canónico, funtorial en $M$ y $\sh{F}$,
\[
\label{I.1.8.7.1}
  \Hom_{\sh{O}_S}(\widetilde{M},u_*(\sh{F}))\isoto\Hom_A(M,\Gamma(X,\sh{F})).
  \tag{1.8.7.1}
\]
\end{corollary}

\begin{proof}
Se tiene, en virtud de \sref[0]{0.4.4.3} y del Lema~\sref{I.1.8.5}, un isomorfismo canónico
de bifuntores
\[
  \Hom_{\sh{O}_S}(\widetilde{M},u_*(\sh{F}))\isoto\Hom_{A'}(M',i_*(u_*(\sh{F})))
\]
y está claro que el segundo miembro es exactamente $\Hom_A(M,\Gamma(X,\sh{F}))$.
Observemos que el homomorfismo canónico \sref{I.1.8.7.1} hace corresponder a cada homomorfismo de $\sh{O}_S$-módulos $h:\widetilde{M}\to u_*(\sh{F})$ (es decir, a cada $u$-morfismo $\widetilde{M}\to\sh{F}$) el homomorfismo de $A$-módulos $\Gamma(h):M\to\Gamma(S,u_*(\sh{F}))=\Gamma(X,\sh{F})$.
\end{proof}

\begin{env}[1.8.8]
\label{I.1.8.8}
Con las notaciones de \sref{I.1.8.4}, está claro \sref[0]{0.4.1.1} que cada morfismo de espacios anillados $(\psi,\theta):X\to S'$ equivale al dato de un homomorfismo de anillos $A\to\Gamma(X,\sh{O}_X)$.
Podemos interpretar así la Proposición~\sref{I.1.8.1} como la definición de una biyección canónica $\Hom(X,S)\isoto\Hom(X,S')$ (donde se entiende que el segundo miembro es el conjunto de morfismos de espacios anillados, pues en general $A$ no es un anillo local).
Más generalmente, si $X$ e $Y$ son espacios localmente anillados y si $(Y',A')$ es el espacio anillado cuyo espacio subyacente está reducido a un punto y cuyo haz de anillos $A'$ es el haz simple definido por el anillo $\Gamma(Y,\sh{O}_Y)$, podemos interpretar \sref{I.1.8.2.1} como una aplicación
\[
  \label{I.1.8.8.1}
  \rho:\Hom_\text{rs}(X,Y)\to\Hom(X,Y')
  \tag{1.8.8.1}
\]
El resultado del Corolario~\sref{I.1.8.3} se interpreta diciendo que los esquemas afines se caracterizan, entre los espacios localmente anillados, como aquellos para los cuales la restricción de $\rho$ a
$\Hom(X,Y)$:
\[
  \label{I.1.8.8.2}
  \rho':\Hom(X,Y)\to\Hom(X,Y')
  \tag{1.8.8.2}
\]
es \emph{biyectiva} para \emph{todo} espacio localmente anillado $X$.
En un capítulo posterior generalizaremos esta definición, lo que permitirá asociar a \emph{todo} espacio anillado $Z$ (y no solamente a un espacio anillado cuyo espacio subyacente esté reducido a un punto) un espacio localmente anillado que llamaremos $\Spec(Z)$;
este será el punto de partida de una teoría «relativa» de preesquemas sobre espacios anillados cualesquiera, que extenderá los resultados del Capítulo~I.
\end{env}

\begin{env}[1.8.9]
\label{I.1.8.9}
Podemos considerar que los pares $(X,\sh{F})$ formados por un espacio localmente anillado $X$ y un $\sh{O}_X$-módulo $\sh{F}$ constituyen una categoría, siendo un \emph{morfismo} de esta categoría un par $(u,h)$ formado por un morfismo de espacios localmente anillados
\oldpage[II]{221}
$u:X\to Y$ y un $u$-morfismo $h:\sh{G}\to\sh{F}$ de módulos;
estos morfismos (para $(X,\sh{F})$ e $(Y,\sh{G})$ fijos) forman un conjunto que denotamos por $\Hom((X,\sh{F}),(Y,\sh{G}))$;
la aplicación $(u,h)\to(\rho'(u),\Gamma(h))$ es una aplicación canónica
\[
  \label{I.1.8.9.1}
  \Hom((X,\sh{F}),(Y,\sh{G}))\to\Hom((\Gamma(Y,\sh{O}_Y),\Gamma(Y,\sh{G})),(\Gamma(X,\sh{O}_X),\Gamma(X,\sh{F})))
  \tag{1.8.9.1}
\]
\emph{funtorial} en $(X,\sh{F})$ e $(Y,\sh{G})$, siendo el segundo miembro el conjunto de dihomomorfismos correspondientes a los anillos y módulos considerados \sref[0]{0.1.0.2}.
\end{env}

\begin{corollary}[1.8.10]
\label{I.1.8.10}
Sean $Y$ un espacio localmente anillado y $\sh{G}$ un $\sh{O}_Y$-módulo.
Para que $Y$ sea un esquema afín y $\sh{G}$ sea un $\sh{O}_Y$-módulo cuasicoherente, es necesario y suficiente que, para cada par $(X,\sh{F})$ formado por un espacio localmente anillado $X$ y un $\sh{O}_X$-módulo $\sh{F}$, la aplicación canónica \sref{I.1.8.9.1} sea biyectiva.
\end{corollary}

Dejamos al lector el desarrollo de la demostración, inspirada en la de \sref{I.1.8.3} y que utiliza \sref{I.1.8.1} y \sref{I.1.8.7}.

\begin{remark}[1.8.11]
\label{I.1.8.11}
Los enunciados \sref{I.1.7.3}, \sref{I.1.7.4} y \sref{I.2.2.4} son casos particulares de \sref{I.1.8.1}, al igual que la definición de \sref{I.1.6.1};
análogamente, \sref{I.2.2.5} se deduce de \sref{I.1.8.7}.
El Corolario~\sref{I.1.8.7} implica también \sref{I.1.6.3} (y, por consiguiente, \sref{I.1.6.4}) como caso particular, pues, si $X$ es un esquema afín y $\Gamma(X,\sh{F})=N$, los funtores $M\mapsto\Hom_{\sh{O}_S}(\widetilde{M},u_*(\widetilde{N}))$ y $M\mapsto\Hom_{\sh{O}_S}(\widetilde{M},(N_{[\vphi]})\supertilde)$ (donde $\vphi:A\to\Gamma(X,\sh{O}_X)$ corresponde a $u$) son isomorfos, por los Corolarios~\sref{I.1.8.7} y \sref{I.1.3.8}.
Por último, \sref{I.1.6.5} (y, por consiguiente, \sref{I.1.6.6}) se deducen de \sref{I.1.8.6} y del hecho de que, para cada $f\in A$, los $A_f$-módulos $N'\otimes_{A'}A_f$ y $(N'\otimes_{A'}A)_f$ (con las notaciones de \sref{I.1.6.5}) son canónicamente isomorfos.
\end{remark}
% END source_es/ega1/ega1-1.tex

% BEGIN source_es/ega1/ega1-2.tex
\section{Preesquemas y morfismos de preesquemas}
\label{section:I.2}

\subsection{Definición de los preesquemas}
\label{subsection:I.2.1}

\begin{env}[2.1.1]
\label{I.2.1.1}
Dado un espacio anillado $(X,\mathscr{O}_X)$, se dice que una parte abierta $V$ de $X$ es un \emph{abierto afín} si el espacio anillado $(V,\mathscr{O}_X|V)$ es un esquema afín \sref{I.1.7.1}.
\end{env}

\begin{definition}[2.1.2]
\label{I.2.1.2}
Se llama preesquema a un espacio anillado $(X,\mathscr{O}_X)$ tal que todo punto de $X$ admite una vecindad abierta afín.
\end{definition}


\oldpage[I]{98}
\begin{proposition}[2.1.3]
\label{I.2.1.3}

Si $(X,\mathscr{O}_X)$ es un preesquema, sus abiertos afines forman una base de la topología de $X$.
\end{proposition}

\begin{proof}
Si $V$ es una vecindad abierta cualquiera de $x\in X$, existe por hipótesis una vecindad abierta $W$ de $x$ tal que $(W,\mathscr{O}_X|W)$ es un esquema afín;
denotemos por $A$ su anillo.
En el espacio $W$, $V\cap W$ es una vecindad abierta de $x$;
por tanto, existe $f\in A$ tal que $D(f)$ es una vecindad abierta de $x$ contenida en $V\cap W$ \sref{I.1.1.10}[i].
El espacio anillado $(D(f),\mathscr{O}_X|D(f))$ es entonces un esquema afín cuyo anillo es isomorfo a $A_f$ \sref{I.1.3.6}, de donde resulta la proposición.
\end{proof}

\begin{proposition}[2.1.4]
\label{I.2.1.4}
El espacio subyacente de un preesquema es un espacio de Kolmogórov.
\end{proposition}

\begin{proof}
Si $x,y$ son dos puntos distintos de un preesquema $X$, está claro que existe una vecindad abierta de uno de estos puntos que no contiene al otro cuando $x,y$ no pertenecen a un mismo abierto afín; y, cuando pertenecen a un mismo abierto afín, esto resulta de \sref{I.1.1.8}.
\end{proof}

\begin{proposition}[2.1.5]
\label{I.2.1.5}
Si $(X,\mathscr{O}_X)$ es un preesquema, toda parte cerrada irreducible de $X$ admite un único punto genérico, y la aplicación $x\to\overline{\{x\}}$ es, por tanto, una biyección de $X$ sobre el conjunto de sus partes cerradas irreducibles.
\end{proposition}

\begin{proof}
Si $Y$ es una parte cerrada irreducible de $X$, $y\in Y$, y $U$ es una vecindad abierta afín de $y$ en $X$, entonces $U\cap Y$ es denso en $Y$ y también irreducible (\hyperref[0.2.1.1]{\textbf{0},~2.1.1} y \hyperref[0.2.1.4]{2.1.4});
por tanto, según el Corolario~\sref{I.1.1.14}, $U\cap Y$ es la clausura en $U$ de un punto $x$, y, por consiguiente, $Y\subset\overline{U}$ es la clausura de $x$ en $X$.
La unicidad del punto genérico de $Y$ resulta de la Proposición~\sref{I.2.1.4} y de \sref[0]{0.2.1.3}.\footnote{\emph{[Trad.] El texto francés imprime aquí «el punto genérico de $X$»; la proposición y la frase precedente conciernen al punto genérico de $Y$.}}
\end{proof}

\begin{env}[2.1.6]
\label{I.2.1.6}
Si $Y$ es una parte cerrada irreducible de $X$ e $y$ es su punto genérico, el anillo local $\mathscr{O}_y$ (denotado también por $\mathscr{O}_{X/Y}$) se llama \emph{anillo local de $X$ a lo largo de $Y$}, o \emph{anillo local de $Y$ en $X$}.

Si $X$ mismo es irreducible y $x$ es su punto genérico, se dice que $\mathscr{O}_x$ es el \emph{anillo de funciones racionales sobre $X$} (cf.~\textsection\hyperref[section:I.7]{7}).
\end{env}

\begin{proposition}[2.1.7]
\label{I.2.1.7}
Si $(X,\mathscr{O}_X)$ es un preesquema, para toda parte abierta $U$ de $X$, el espacio anillado $(U,\mathscr{O}_X|U)$ es un preesquema.
\end{proposition}

\begin{proof}
Esto resulta directamente de la Definición~\sref{I.2.1.2} y de la Proposición~\sref{I.2.1.3}.
\end{proof}

Se dice que $(U,\mathscr{O}_X|U)$ es el preesquema \emph{inducido} sobre $U$ por $(X,\mathscr{O}_X)$, o la \emph{restricción} de $(X,\mathscr{O}_X)$ a $U$.

\begin{env}[2.1.8]
\label{I.2.1.8}
Se dice que un preesquema $(X,\mathscr{O}_X)$ es \emph{irreducible} (resp. \emph{conexo}) si el espacio subyacente $X$ es irreducible (resp. conexo).
Se dice que un preesquema es \emph{íntegro} si es \emph{irreducible y reducido} (cf. \sref{I.5.1.4}).
Se dice que un preesquema $(X,\mathscr{O}_X)$ es \emph{localmente íntegro} si todo $x\in X$ admite una vecindad abierta $U$ tal que el preesquema inducido sobre $U$ por $(X,\mathscr{O}_X)$ es íntegro.
\end{env}

\subsection{Morfismos de preesquemas}
\label{subsection:I.2.2}
\label{I.2.2}

\begin{definition}[2.2.1]
\label{I.2.2.1}
Dados dos preesquemas $(X,\mathscr{O}_X)$ y $(Y,\mathscr{O}_Y)$, se llama morfismo (de preesquemas) de $(X,\mathscr{O}_X)$ en $(Y,\mathscr{O}_Y)$ a todo morfismo de espacios anillados $(\psi,\theta)$ tal que, para todo $x\in X$, $\theta_x^\sharp$ es un homomorfismo local $\mathscr{O}_{\psi(x)}\longrightarrow\mathscr{O}_x$.
\end{definition}

\oldpage[I]{99}
Por paso a los cocientes, la aplicación $\theta_x^\sharp:\mathscr{O}_{\psi(x)}\to\mathscr{O}_x$ proporciona, por tanto, un monomorfismo $\theta^x:\boldsymbol{k}(\psi(x))\to\boldsymbol{k}(x)$, lo que permite considerar $\boldsymbol{k}(x)$ como una \emph{extensión} del campo $\boldsymbol{k}(\psi(x))$.

\begin{env}[2.2.2]
\label{I.2.2.2}
El compuesto $(\psi'',\theta'')$ de dos morfismos de preesquemas $(\psi,\theta)$ y $(\psi',\theta')$ es también un morfismo de preesquemas, como resulta de la fórmula ${\theta''}^\sharp=\theta^\sharp\circ\psi^*({\theta'}^\sharp)$ \sref[0]{0.3.5.5}.
De aquí se concluye que los preesquemas forman una \emph{categoría};
siguiendo la notación habitual, se escribirá $\operatorname{Hom}(X,Y)$ para denotar el conjunto de los morfismos de un preesquema $X$ en un preesquema $Y$.
\end{env}

\begin{example}[2.2.3]
\label{I.2.2.3}
Si $U$ es una parte abierta de $X$, la inyección canónica \sref[0]{0.4.1.2} del preesquema inducido $(U,\mathscr{O}_X|U)$ en $(X,\mathscr{O}_X)$ es un morfismo de preesquemas;
es, además, un \emph{monomorfismo} de espacios anillados (y \emph{a fortiori} un monomorfismo de preesquemas), como resulta directamente de \sref[0]{0.4.1.1}.
\end{example}

\begin{proposition}[2.2.4]
\label{I.2.2.4}
\footnote{\emph{[Trad.] Véase \hyperref[subsection:I.1.8]{(1.8)} y la nota a pie de página correspondiente.}}
Sea $(X,\mathscr{O}_X)$ un preesquema y sea $(S,\mathscr{O}_S)$ un esquema afín asociado a un anillo $A$.
Existe entonces una correspondencia biyectiva canónica entre los morfismos de preesquemas de $(X,\mathscr{O}_X)$ en $(S,\mathscr{O}_S)$ y los homomorfismos de anillos de $A$ en $\Gamma(X,\mathscr{O}_X)$.
\end{proposition}

\begin{proof}
Observemos primero que, si $(X,\mathscr{O}_X)$ y $(Y,\mathscr{O}_Y)$ son dos espacios anillados cualesquiera,
un morfismo $(\psi,\theta)$ de $(X,\mathscr{O}_X)$ en $(Y,\mathscr{O}_Y)$ define canónicamente un homomorfismo de anillos $\Gamma(\theta):\Gamma(Y,\mathscr{O}_Y)\to
\Gamma(Y,\psi_*(\mathscr{O}_X))=\Gamma(X,\mathscr{O}_X)$.
En el caso que nos ocupa, todo se reduce a demostrar que todo homomorfismo $\varphi:A\to\Gamma(X,\mathscr{O}_X)$ es de la forma $\Gamma(\theta)$ para un único $\theta$.
Ahora bien, por hipótesis, existe un recubrimiento $(V_\alpha)$ de $X$ por abiertos afines;
al componer $\varphi$ con el homomorfismo de restricción $\Gamma(X,\mathscr{O}_X)\to
\Gamma(V_\alpha,\mathscr{O}_X|V_\alpha)$, se obtiene un homomorfismo $\varphi_\alpha:A\to\Gamma(V_\alpha,\mathscr{O}_X|V_\alpha)$ que corresponde a un único morfismo $(\psi_\alpha,\theta_\alpha)$ del preesquema $(V_\alpha,\mathscr{O}_X|V_\alpha)$ en $(S,\mathscr{O}_S)$, en virtud del Teorema~\sref{I.1.7.3}.
Además, para cada par de índices $(\alpha,\beta)$, todo punto de $V_\alpha\cap V_\beta$ admite una vecindad abierta afín $W$ contenida en $V_\alpha\cap V_\beta$ \sref{I.2.1.3};
está claro que, al componer $\varphi_\alpha$ y $\varphi_\beta$ con los homomorfismos de restricción a $W$, se obtiene el mismo homomorfismo $\Gamma(S,\mathscr{O}_S)\to\Gamma(W,\mathscr{O}_X|W)$; por tanto, en virtud de las relaciones $(\theta_\alpha^\sharp)_x=(\varphi_\alpha)_x$ para todo $x\in V_\alpha$ y todo $\alpha$ \sref{I.1.6.1}, las restricciones a $W$ de los morfismos $(\psi_\alpha,\theta_\alpha)$ y $(\psi_\beta,\theta_\beta)$ coinciden.
Se concluye que existe un único morfismo de espacios anillados $(\psi,\theta):(X,\mathscr{O}_X)\to(S,\mathscr{O}_S)$ cuya restricción a cada $V_\alpha$ es $(\psi_\alpha,\theta_\alpha)$; está claro que se trata de un morfismo de preesquemas y que $\Gamma(\theta)=\varphi$.

Sea $u:A\to\Gamma(X,\mathscr{O}_X)$ un homomorfismo de anillos, y sea $v=(\psi,\theta)$ el morfismo correspondiente $(X,\mathscr{O}_X)\to(S,\mathscr{O}_S)$.
Para todo $f\in A$, se tiene
\begin{equation}
  \psi^{-1}(D(f))=X_{u(f)}
  \tag{2.2.4.1}\label{I.2.2.4.1-fr}
\end{equation}
con la notación de \sref[0]{0.5.5.2} relativa al haz localmente libre $\mathscr{O}_X$.
En efecto, basta verificar esta fórmula cuando $X$ mismo es afín, y entonces no es otra que \sref{I.1.2.2.2}.
\end{proof}

\begin{proposition}[2.2.5]
\label{I.2.2.5}
Bajo las hipótesis de la Proposición~\sref{I.2.2.4}, sea $\varphi:A\to\Gamma(X,\mathscr{O}_X)$ un homomorfismo de anillos, sea $f:(X,\mathscr{O}_X)\to(S,\mathscr{O}_S)$ el morfismo de preesquemas correspondiente, sea $\mathscr{G}$ \textup{(resp. $\mathscr{F}$)} un $\mathscr{O}_X$-módulo cuasicoherente \textup{(resp.} un $\mathscr{O}_S$-módulo cuasicoherente), y sea $M=\Gamma(S,\mathscr{F})$.
Existe entonces una correspondencia biyectiva canónica
\oldpage[I]{100}
entre los $f$-morfismos $\mathscr{F}\to\mathscr{G}$ \sref[0]{0.4.4.1} y los $A$-homomorfismos $M\to(\Gamma(X,\mathscr{G}))_{[\varphi]}$.
\end{proposition}

\begin{proof}
Razonando como en la Proposición~\sref{I.2.2.4}, se reduce al caso en que $X$ es afín, y la proposición resulta entonces de la Proposición~\sref{I.1.6.3} y del Corolario~\sref{I.1.3.8}.
\end{proof}

\begin{env}[2.2.6]
\label{I.2.2.6}
Se dice que un morfismo de preesquemas $(\psi,\theta):(X,\mathscr{O}_X)\to(Y,\mathscr{O}_Y)$ es \emph{abierto} (resp. \emph{cerrado}) si, para toda parte abierta $U$ de $X$ (resp. toda parte cerrada $F$ de $X$), $\psi(U)$ es abierta en $Y$ (resp. $\psi(F)$ cerrada en $Y$).
Se dice que $(\psi,\theta)$ es \emph{dominante} si $\psi(X)$ es denso en $Y$, y \emph{sobreyectivo} si $\psi$ es sobreyectiva.
Se observará que estas condiciones solo dependen de la aplicación continua $\psi$.
\end{env}

\begin{proposition}[2.2.7]
\label{I.2.2.7}
Sean
\[
  f=(\psi,\theta):(X,\mathscr{O}_X)\to(Y,\mathscr{O}_Y),
  \qquad
  g=(\psi',\theta'):(Y,\mathscr{O}_Y)\to(Z,\mathscr{O}_Z)
\]
dos morfismos de preesquemas.
\begin{enumerate}
  \item[{\rm(i)}] Si $f$ y $g$ son ambos abiertos \textup{(resp.} cerrados, dominantes, sobreyectivos),
    también lo es $g\circ f$.
  \item[{\rm(ii)}] Si $f$ es sobreyectivo y $g\circ f$ es cerrado, entonces $g$ es cerrado.
  \item[{\rm(iii)}] Si $g\circ f$ es sobreyectivo, entonces $g$ es sobreyectivo.
\end{enumerate}
\end{proposition}

\begin{proof}
Las afirmaciones (i) y (iii) son evidentes.
Pongamos $g\circ f=(\psi'',\theta'')$.
Si $F$ es una parte cerrada de $Y$, entonces $\psi^{-1}(F)$ es una parte cerrada de $X$, por lo que $\psi''(\psi^{-1}(F))$ es cerrada en $Z$;
pero, como $\psi$ es sobreyectiva, $\psi(\psi^{-1}(F))=F$, de modo que $\psi''(\psi^{-1}(F))=\psi'(F)$, lo que demuestra (ii).
\end{proof}

\begin{proposition}[2.2.8]
\label{I.2.2.8}
Sea $f=(\psi,\theta)$ un morfismo $(X,\mathscr{O}_X)\to(Y,\mathscr{O}_Y)$, y sea $(U_\alpha)$ un recubrimiento abierto de $Y$.
Para que $f$ sea abierto \textup{(resp.} cerrado, sobreyectivo, dominante), es necesario y suficiente que su restricción a cada preesquema inducido $(\psi^{-1}(U_\alpha),\mathscr{O}_X|\psi^{-1}(U_\alpha))$, considerada como un morfismo de preesquemas de este preesquema inducido en el preesquema inducido $(U_\alpha,\mathscr{O}_Y|U_\alpha)$, sea un morfismo abierto \textup{(resp.} cerrado, sobreyectivo, dominante).
\end{proposition}

\begin{proof}
La proposición resulta directamente de las definiciones, teniendo en cuenta que una parte $F$ de $Y$ es cerrada (resp. abierta, densa) en $Y$ si y solo si cada $F\cap U_\alpha$ es cerrada (resp. abierta, densa) en $U_\alpha$.
\end{proof}

\begin{env}[2.2.9]
\label{I.2.2.9}
Sean $(X,\mathscr{O}_X)$ y $(Y,\mathscr{O}_Y)$ dos preesquemas;
se supone que $X$ e $Y$ tienen un mismo número \emph{finito} de componentes irreducibles $X_i$ (resp. $Y_i$) $(1\leq i\leq n)$;
sean $\xi_i$ (resp. $\eta_i$) el punto genérico de $X_i$ (resp. $Y_i$) \sref{I.2.1.5}.
Se dice que un morfismo
\[
  f=(\psi,\theta):(X,\mathscr{O}_X)\to(Y,\mathscr{O}_Y)
\]
es \emph{birracional} si, para todo $i$, $\psi^{-1}(\eta_i)=\{\xi_i\}$ y $\theta_{\xi_i}^{\sharp}:\mathscr{O}_{\eta_i}\to\mathscr{O}_{\xi_i}$ es un \emph{isomorfismo}.
Está claro que un morfismo birracional es dominante \sref[0]{0.2.1.8}, y por tanto es sobreyectivo si además es cerrado.
\end{env}

\begin{notation}[2.2.10]
\label{I.2.2.10}
En todo lo que sigue, cuando no haya riesgo de confusión, \emph{suprimiremos} el haz estructural (resp. el morfismo de haces estructurales) de la notación de un preesquema (resp. de un morfismo de preesquemas).
Si $U$ es una parte abierta del espacio subyacente $X$ de un preesquema, siempre que se hable de $U$ como preesquema entenderemos el preesquema inducido sobre $U$.
\end{notation}


\oldpage[I]{101}
\subsection{Pegado de preesquemas}
\label{subsection:I.2.3}

\begin{env}[2.3.1]
\label{I.2.3.1}

De la Definición~\sref{I.2.1.2} se deduce que todo espacio anillado obtenido por \emph{pegado} de preesquemas \sref[0]{0.4.1.7} es también un preesquema.
En particular, puesto que todo preesquema admite, por definición, un recubrimiento por abiertos afines, se ve que todo preesquema puede obtenerse por \emph{pegado de esquemas afines}.
\end{env}

\begin{example}[2.3.2]
\label{I.2.3.2}
Sean $K$ un cuerpo, $B=K[s]$ y $C=K[t]$ dos anillos de polinomios en una indeterminada sobre $K$, y sean $X_1=\operatorname{Spec}(B)$ y $X_2=\operatorname{Spec}(C)$, que son dos esquemas afines isomorfos.
En $X_1$ (resp. $X_2$), sea $U_{12}$ (resp. $U_{21}$) el abierto afín $D(s)$ (resp. $D(t)$) cuyo anillo $B_s$ (resp. $C_t$) está formado por las fracciones racionales de la forma $f(s)/s^m$ (resp. $g(t)/t^n$), con $f\in B$ (resp. $g\in C$).
Sea $u_{12}$ el isomorfismo de preesquemas $U_{21}\to U_{12}$ correspondiente \sref{I.2.2.4} al isomorfismo de $B_s$ en $C_t$ que a $f(s)/s^m$ hace corresponder la fracción racional $f(1/t)/(1/t^m)$.
Se pueden pegar $X_1$ y $X_2$ a lo largo de $U_{12}$ y $U_{21}$ mediante $u_{12}$, pues evidentemente no hay condición de pegado.
Más adelante volveremos a encontrar el preesquema $X$ así obtenido como caso particular de un método general de construcción \sref[II]{II.2.4.3}.
Mostremos aquí solamente que $X$ \emph{no es un esquema afín};
esto se deducirá de que el anillo $\Gamma(X,\mathscr{O}_X)$ es \emph{isomorfo a} $K$, por lo que su espectro se reduce a un punto.
En efecto, una sección de $\mathscr{O}_X$ sobre $X$ tiene una restricción sobre $X_1$ (resp. $X_2$), identificado con un abierto afín de $X$, que es un polinomio $f(s)$ (resp. $g(t)$), y de las definiciones se deduce que se debe tener $g(t)=f(1/t)$, lo cual solo es posible si $f=g\in K$.
\end{example}

\subsection{Esquemas locales}
\label{subsection:I.2.4}

\begin{env}[2.4.1]
\label{I.2.4.1}
Se dice que un esquema afín es un \emph{esquema local} si su anillo $A$ es local;
existe entonces, en $X=\operatorname{Spec}(A)$, un único \emph{punto cerrado $a$}, y para cualquier otro $b\in X$ se tiene $a\in\overline{\{b\}}$ \sref{I.1.1.7}.
\end{env}

Para todo preesquema $Y$ y todo punto $y\in Y$, el esquema local $\operatorname{Spec}(\mathscr{O}_y)$
se denomina \emph{esquema local de $Y$ en el punto $y$}.
Sea $V$ un abierto afín de $Y$ que contiene a $y$, y sea $B$ el anillo del esquema afín $V$;
entonces $\mathscr{O}_y$ se identifica canónicamente con $B_y$ \sref{I.1.3.4}, y el homomorfismo canónico $B\to B_y$ corresponde por tanto \sref{I.1.6.1} a un morfismo de preesquemas $\operatorname{Spec}(\mathscr{O}_y)\to V$.
Si se compone este morfismo con la inyección canónica $V\to Y$, se obtiene un morfismo $\operatorname{Spec}(\mathscr{O}_y)\to Y$ que es \emph{independiente} del abierto afín $V$ (que contiene a $y$) elegido: en efecto, si $V'$ es otro abierto afín que contiene a $y$, existe un tercer abierto afín $W$ que contiene a $y$ y tal que $W\subset V\cap V'$ \sref{I.2.1.3};
se puede, pues, suponer que $V\subset V'$, y entonces, si $B'$ es el anillo de $V'$, basta observar que el diagrama
\phantomsection\label{I.2.4.1.diagram-fr}
\[
  \xymatrix{
    B'\ar[rr]\ar[dr] & & B\ar[dl]\\
    & \mathscr{O}_y
  }
\]
es conmutativo \sref[0]{0.1.5.1}.
El morfismo
\[
  \operatorname{Spec}(\mathscr{O}_y)\to Y
\]
así definido se denomina {\it canónico}.


\oldpage[I]{102}
\begin{proposition}[2.4.2]
\label{I.2.4.2}

Sea $(Y,\mathscr{O}_Y)$ un preesquema;
para todo $y\in Y$, sea $(\psi,\theta)$ el morfismo canónico $(\operatorname{Spec}(\mathscr{O}_y),\widetilde{\mathscr{O}}_y)
\to(Y,\mathscr{O}_Y)$.
Entonces $\psi$ es un homeomorfismo de $\operatorname{Spec}(\mathscr{O}_y)$ sobre el subespacio $S_y$ de $Y$ formado por los $z$ tales que $y\in\overline{\{z\}}$ \textup{(dicho de otro modo, las} generalizaciones \textup{de} $y$ \sref[0]{0.2.1.2}\textup{)};
además, si $z=\psi(\mathfrak{p})$, entonces $\theta_z^{\sharp}:\mathscr{O}_z\to(\mathscr{O}_y)_{\mathfrak{p}}$ es un isomorfismo;
$(\psi,\theta)$ es, por tanto, un monomorfismo de espacios anillados.
\end{proposition}

\begin{proof}
Como el único punto cerrado $a$ de $\operatorname{Spec}(\mathscr{O}_y)$ pertenece a la clausura de todo punto de este espacio y $\psi(a)=y$, la imagen de $\operatorname{Spec}(\mathscr{O}_y)$ por la aplicación continua $\psi$ está contenida en $S_y$.
Como $S_y$ está contenido en todo abierto afín que contiene a $y$, se puede reducir al caso en que $Y$ es un esquema afín;
pero, en este caso, la proposición resulta de \sref{I.1.6.2}.
\end{proof}

\emph{Se ve, pues, \sref{I.2.1.5} que existe una correspondencia biunívoca entre $\operatorname{Spec}(\mathscr{O}_y)$ y el conjunto de los subconjuntos cerrados irreducibles de $Y$ que contienen a $y$.}

\begin{corollary}[2.4.3]
\label{I.2.4.3}
Para que $y\in Y$ sea el punto genérico de una componente irreducible de $Y$, es necesario y suficiente que el único ideal primo del anillo local $\mathscr{O}_y$ sea su ideal maximal \textup{(dicho de otro modo, que} $\mathscr{O}_y$ \textup{sea} de dimensión cero).
\end{corollary}

\begin{proposition}[2.4.4]
\label{I.2.4.4}
Sean $(X,\mathscr{O}_X)$ un esquema local de anillo $A$, $a$ su único punto cerrado, y $(Y,\mathscr{O}_Y)$ un preesquema.
Todo morfismo $u=(\psi,\theta):(X,\mathscr{O}_X)\to(Y,\mathscr{O}_Y)$ se factoriza entonces de manera única como $X\to\operatorname{Spec}(\mathscr{O}_{\psi(a)})\to Y$, donde la segunda flecha designa el morfismo canónico, y la primera corresponde a un homomorfismo local $\mathscr{O}_{\psi(a)}\to A$.
Esto establece una correspondencia biunívoca canónica entre el conjunto de los morfismos $(X,\mathscr{O}_X)\to(Y,\mathscr{O}_Y)$ y el conjunto de los homomorfismos locales $\mathscr{O}_y\to A$ $(y\in Y)$.
\end{proposition}

En efecto, para todo $x\in X$ se tiene $a\in\overline{\{x\}}$, luego $\psi(a)\in\overline{\{\psi(x)\}}$, lo que muestra que $\psi(X)$ está contenido en todo abierto afín que contiene a $\psi(a)$.
Basta, pues, considerar el caso en que $(Y,\mathscr{O}_Y)$ es un esquema afín de anillo $B$, y entonces se tiene $u=({}^a\varphi,\widetilde{\varphi})$, donde $\varphi\in\operatorname{Hom}(B,A)$ \sref{I.1.7.3}.
Además, se tiene $\varphi^{-1}(\mathfrak{j}_a)=\mathfrak{j}_{\psi(a)}$, y, por consiguiente, la imagen por $\varphi$ de todo elemento de $B-\mathfrak{j}_{\psi(a)}$ es invertible en el anillo local $A$;
la factorización afirmada en el enunciado resulta entonces de la propiedad universal del anillo de fracciones \sref[0]{0.1.2.4}.
Recíprocamente, a todo homomorfismo local $\mathscr{O}_y\to A$ le corresponde un único morfismo $(\psi,\theta):X\to\operatorname{Spec}(\mathscr{O}_y)$ tal que $\psi(a)=y$ \sref{I.1.7.3}, y, componiéndolo con el morfismo canónico $\operatorname{Spec}(\mathscr{O}_y)\to Y$, se obtiene un morfismo $X\to Y$, lo que concluye la demostración de la proposición.

\begin{env}[2.4.5]
\label{I.2.4.5}
Los esquemas afines cuyo anillo es un cuerpo $K$ tienen un espacio subyacente reducido a un solo punto.
Si $A$ es un anillo local de ideal maximal $\mathfrak{m}$, todo homomorfismo local $A\to K$ tiene núcleo igual a $\mathfrak{m}$ y se factoriza, por tanto, como $A\to A/\mathfrak{m}\to K$, donde la segunda flecha es un monomorfismo.
Los morfismos $\operatorname{Spec}(K)\to\operatorname{Spec}(A)$ corresponden, pues, biunívocamente a los monomorfismos de cuerpos $A/\mathfrak{m}\to K$.
\end{env}

Sea $(Y,\mathscr{O}_Y)$ un preesquema;
para todo $y\in Y$ y todo ideal $\mathfrak{a}_y$ de $\mathscr{O}_y$, el homomorfismo canónico $\mathscr{O}_y\to\mathscr{O}_y/\mathfrak{a}_y$ define un morfismo $\operatorname{Spec}(\mathscr{O}_y/\mathfrak{a}_y)
\to\operatorname{Spec}(\mathscr{O}_y)$;
si se compone con el morfismo canónico $\operatorname{Spec}(\mathscr{O}_y)\to Y$, se obtiene un morfismo $\operatorname{Spec}(\mathscr{O}_y/\mathfrak{a}_y)\to Y$, que también se denomina \textit{canónico}.
Para $\mathfrak{a}_y=\mathfrak{m}_y$, es decir, $\mathscr{O}_y/\mathfrak{a}_y=\boldsymbol{k}(y)$, la Proposición~\sref{I.2.4.4} da por tanto:


\oldpage[I]{103}
\begin{corollary}[2.4.6]
\label{I.2.4.6}

Sean $(X,\mathscr{O}_X)$ un esquema local cuyo anillo $K$ es un cuerpo, $\xi$ el único punto de $X$, y $(Y,\mathscr{O}_Y)$ un preesquema.
Entonces todo morfismo $u:(X,\mathscr{O}_X)\to(Y,\mathscr{O}_Y)$ se factoriza de manera única como $X\to\operatorname{Spec}(\boldsymbol{k}(\psi(\xi)))\to Y$, donde la segunda flecha designa el morfismo canónico, y la primera corresponde a un monomorfismo $\boldsymbol{k}(\psi(\xi))\to K$.
Esto establece una correspondencia biunívoca canónica entre el conjunto de los morfismos $(X,\mathscr{O}_X)\to(Y,\mathscr{O}_Y)$ y el conjunto de los monomorfismos $\boldsymbol{k}(y)\to K$ $(y\in Y)$.
\end{corollary}

\begin{corollary}[2.4.7]
\label{I.2.4.7}
Para todo $y\in Y$, todo morfismo canónico $\operatorname{Spec}(\mathscr{O}_y/\mathfrak{a}_y)\to Y$ es un monomorfismo de espacios anillados.
\end{corollary}

\begin{proof}
Ya se ha visto esto cuando $\mathfrak{a}_y=0$ \sref{I.2.4.2}, y basta aplicar el Corolario~\sref{I.1.7.5}.
\end{proof}

\begin{remark}[2.4.8]
\label{I.2.4.8}
Sea $X$ un esquema local, y sea $a$ su único punto cerrado.
Como todo abierto afín que contiene a $a$ es necesariamente todo $X$, todo $\mathscr{O}_X$-módulo \emph{invertible} \sref[0]{0.5.4.1} es necesariamente \emph{isomorfo a $\mathscr{O}_X$} (o, como también se dice, es \emph{trivial}).
Esta propiedad no se cumple en general para un esquema afín arbitrario $\operatorname{Spec}(A)$;
se verá en el Capítulo~V que, si $A$ es un anillo normal, se cumple cuando $A$ es un dominio de factorización única.
\end{remark}

\subsection{Preesquemas sobre un preesquema}
\label{subsection:I.2.5}

\begin{definition}[2.5.1]
\label{I.2.5.1}
Dado un preesquema $S$, se dice que el dato de un preesquema $X$ y de un morfismo de preesquemas $\varphi:X\to S$ define un preesquema $X$ \emph{sobre el preesquema $S$}, o un $S$-preesquema;
se dice que $S$ es el preesquema de base del $S$-preesquema $X$.
El morfismo $\varphi$ se denomina morfismo estructural del $S$-preesquema $X$.
Cuando $S$ es un esquema afín de anillo $A$, se dice también que $X$, provisto de $\varphi$, es un preesquema \textup{sobre el anillo} $A$ (o un \emph{$A$-preesquema}).
\end{definition}

De \sref{I.2.2.4} resulta que el dato de un preesquema sobre un anillo $A$ equivale al de un preesquema $(X,\mathscr{O}_X)$ cuyo haz estructural $\mathscr{O}_X$ es un haz de \emph{$A$-álgebras}.
\emph{Por tanto, todo preesquema puede considerarse siempre, de manera única, como un $\mathbf{Z}$-preesquema}.

Si $\varphi:X\to S$ es el morfismo estructural de un $S$-preesquema $X$, se
dice que un punto $x\in X$ está \emph{sobre un punto $s\in S$} si $\varphi(x)=s$. Se
dice que $X$ \emph{domina} a $S$ si $\varphi$ es un morfismo dominante \sref{I.2.2.6}.

\begin{env}[2.5.2]
\label{I.2.5.2}
Sean $X$ e $Y$ dos $S$-preesquemas;
se dice que un morfismo de preesquemas $u:X\to Y$ es un \emph{morfismo de preesquemas sobre $S$} (o un \emph{$S$-morfismo}) si el diagrama
\phantomsection\label{I.2.5.2.diagram-fr}
\[
  \xymatrix{
    X\ar[rr]^u\ar[dr] & & Y\ar[dl]\\
    & S &
  }
\]
(donde las flechas diagonales son los morfismos estructurales) es conmutativo: esto implica que, para todo $s\in S$ y todo $x\in X$ sobre $s$, $u(x)$ debe estar también sobre $s$.
\end{env}

De esta definición resulta inmediatamente que la composición de dos $S$-morfismos cualesquiera es un $S$-morfismo;
los $S$-preesquemas forman, por tanto, una \emph{categoría}.

Se denota por $\operatorname{Hom}_S(X,Y)$ el conjunto de los $S$-morfismos de un $S$-preesquema $X$ a un $S$-preesquema $Y$;
el morfismo identidad de un $S$-preesquema $X$ se denota por $1_X$.

Cuando $S$ es un esquema afín de anillo $A$, se dice también \emph{$A$-morfismo} en lugar de $S$-morfismo.


\oldpage[I]{104}
\begin{env}[2.5.3]
\label{I.2.5.3}

Si $X$ es un $S$-preesquema y $v:X'\to X$ un morfismo de preesquemas, la composición $X'\xrightarrow{v}X\to S$ dota a $X'$ de una estructura de $S$-preesquema;
en particular, todo preesquema inducido sobre un abierto $U$ de $X$ puede considerarse como un $S$-preesquema mediante la inyección canónica.

Si $u:X\to Y$ es un $S$-morfismo de $S$-preesquemas, la restricción de $u$ a todo preesquema inducido sobre un abierto $U$ de $X$ es, por tanto, un $S$-morfismo $U\to Y$.
Recíprocamente, sea $(U_\alpha)$ un recubrimiento abierto de $X$, y, para cada $\alpha$, sea $u_\alpha:U_\alpha\to Y$ un $S$-morfismo;
si, para todo par de índices $(\alpha,\beta)$, las restricciones de $u_\alpha$ y $u_\beta$ a $U_\alpha\cap U_\beta$ coinciden, existe un único $S$-morfismo $X\to Y$ cuya restricción a cada $U_\alpha$ es $u_\alpha$.

Si $u:X\to Y$ es un $S$-morfismo tal que $u(X)\subset V$, donde $V$ es un abierto de $Y$, entonces $u$, considerado como morfismo de $X$ a $V$, es también un $S$-morfismo.
\end{env}

\begin{env}[2.5.4]
\label{I.2.5.4}
Sea $S'\to S$ un morfismo de preesquemas;
para todo $S'$-preesquema $X$, la composición $X\to S'\to S$ dota a $X$ de una estructura de $S$-preesquema.
Recíprocamente, supongamos que $S'$ es el preesquema inducido sobre un abierto de $S$;
sea $X$ un $S$-preesquema y supongamos que el morfismo estructural $f:X\to S$ es tal que $f(X)\subset S'$;
entonces puede considerarse $X$ como un $S'$-preesquema.
En este último caso, si $Y$ es otro $S$-preesquema cuyo morfismo estructural envía también el espacio subyacente a $S'$, todo $S$-morfismo de $X$ a $Y$ es también un $S'$-morfismo.
\end{env}

\begin{env}[2.5.5]
\label{I.2.5.5}
Si $X$ es un $S$-preesquema,\footnote{La fuente francesa imprime aquí «$S$-morfismo»; el morfismo estructural $X\to S$ y la definición de una sección requieren un $S$-preesquema.} de morfismo estructural $\varphi:X\to S$, se denomina \emph{$S$-sección de $X$} a todo $S$-morfismo de $S$ a $X$, es decir, a todo morfismo de preesquemas $\psi:S\to X$ tal que $\varphi\circ\psi$ sea la identidad de $S$.
Se denota por $\Gamma(X/S)$ el conjunto de las $S$-secciones de $X$.
\end{env}
% END source_es/ega1/ega1-2.tex

% BEGIN source_es/ega1/ega1-3.tex
\section{Productos de preesquemas}
\label{section:I.3}

\subsection{Suma de preesquemas}
\label{subsection:I.3.1}
\label{I.3.1}

\phantomsection\label{I.3.1.1}
Sea $(X_\alpha)$ una familia cualquiera de preesquemas; sea $X$ un espacio topológico que es la
\emph{suma} de los espacios subyacentes $X_\alpha$; $X$ es entonces la unión de subespacios
abiertos $X'_\alpha$ disjuntos dos a dos, y para cada $\alpha$ existe un homeomorfismo
$\varphi_\alpha$ de $X_\alpha$ sobre $X'_\alpha$. Si se dota cada uno de los $X'_\alpha$ del haz
$(\varphi_\alpha)_*(\mathscr{O}_{X_\alpha})$, es claro que $X$ se convierte en un preesquema, al que
se llama la \emph{suma} de la familia de preesquemas $(X_\alpha)$ y que se denota por
$\coprod_\alpha X_\alpha$. Si $Y$ es un preesquema, la aplicación $f\to(f\circ\varphi_\alpha)$ es una
\emph{biyección} del conjunto $\operatorname{Hom}(X,Y)$ sobre el conjunto producto $\prod_\alpha\operatorname{Hom}(X_\alpha,Y)$.
En particular, si los $X_\alpha$ son $S$-preesquemas, con morfismos estructurales $\psi_\alpha$, entonces $X$ es un $S$-preesquema mediante el único morfismo $\psi:X\to S$ tal que
$\psi\circ\varphi_\alpha=\psi_\alpha$ para cada $\alpha$. La suma de dos preesquemas $X$ e $Y$ se denota por $X\amalg Y$.
Es inmediato que, si $X=\operatorname{Spec}(A)$ e $Y=\operatorname{Spec}(B)$, entonces $X\amalg Y$ se identifica canónicamente con $\operatorname{Spec}(A\times B)$.

\subsection{Producto de preesquemas}
\label{subsection:I.3.2}

\begin{definition}[3.2.1]
\label{I.3.2.1}
Dados dos $S$-preesquemas $X$ e $Y$, se dice que un triple $(Z,p_1,p_2)$, formado por un
$S$-preesquema $Z$ y por dos $S$-morfismos $p_1:Z\to X$ y $p_2:Z\to Y$, es un producto de los
\oldpage[I]{105}
$S$-preesquemas $X$ e $Y$ si, para todo $S$-preesquema $T$, la aplicación
$f\to(p_1\circ f,p_2\circ f)$ es una biyección del conjunto de los $S$-morfismos de $T$ en
$Z$ sobre el conjunto de los pares formados por un $S$-morfismo $T\to X$ y un $S$-morfismo $T\to Y$
\textup{(en otras palabras, una biyección}
\[
  \operatorname{Hom}_S(T,Z)\xrightarrow{\sim}
  \operatorname{Hom}_S(T,X)\times\operatorname{Hom}_S(T,Y))
\]
\end{definition}

Esta es la noción general de \emph{producto} de dos objetos en una categoría, aplicada
a la categoría de los $S$-preesquemas (T, I, 1.1); en particular, un producto de dos $S$-preesquemas
es \emph{único} salvo un $S$-isomorfismo único. En virtud de esta unicidad, la mayoría de las veces
se denotará un producto de dos $S$-preesquemas $X$ e $Y$ por $X\times_S Y$ (o
simplemente por $X\times Y$, cuando no haya posibilidad de confusión), omitiendo en la notación los morfismos $p_1$ y $p_2$
(las \emph{proyecciones canónicas} de $X\times_S Y$ sobre $X$ e $Y$, respectivamente).
Si $g:T\to X$ y $h:T\to Y$ son $S$-morfismos, se denota por $(g,h)_S$, o
simplemente por $(g,h)$, el $S$-morfismo $f:T\to X\times_S Y$ tal que $p_1\circ f=g$,
$p_2\circ f=h$. Si $X'$ e $Y'$ son dos $S$-preesquemas, y $p'_1$ y $p'_2$ son las proyecciones canónicas
de $X'\times_S Y'$ (que se supone que existe), y $u:X'\to X$ y
$v:Y'\to Y$ son $S$-morfismos, se escribirá $u\times_S v$ (o simplemente $u\times v$)
para designar el $S$-morfismo $(u\circ p'_1,v\circ p'_2)_S$ de $X'\times_S Y'$ en
$X\times_S Y$.

Cuando $S$ es un esquema afín de anillo $A$, en las notaciones anteriores se sustituye a menudo $S$
por $A$.

\begin{proposition}[3.2.2]
\label{I.3.2.2}
Sean $X$, $Y$ y $S$ tres esquemas afines, cuyos anillos respectivos son $B$, $C$ y $A$. Sea
$Z=\operatorname{Spec}(B\otimes_A C)$, y sean $p_1$ y $p_2$ los $S$-morfismos correspondientes \sref{I.2.2.4} a
los $A$-homomorfismos canónicos $u:b\to b\otimes1$ y $v:c\to1\otimes c$, de $B$
y $C$, respectivamente, en $B\otimes_A C$; entonces $(Z,p_1,p_2)$ es un producto de $X$ e $Y$.
\end{proposition}

\begin{proof}
En virtud de \sref{I.2.2.4}, basta comprobar que, si a cada $A$-homomorfismo
$f:B\otimes_A C\to L$ (donde $L$ es una $A$-álgebra) se le asocia el par
$(f\circ u,f\circ v)$, se define así una biyección
\[
  \operatorname{Hom}_A(B\otimes_A C,L)\xrightarrow{\sim}
  \operatorname{Hom}_A(B,L)\times\operatorname{Hom}_A(C,L)\footnotemark,
\]
\footnotetext{La notación $\operatorname{Hom}_A$
designa aquí el conjunto de homomorfismos de $A$-\emph{álgebras}.}, lo que resulta inmediatamente
de las definiciones y de la relación $b\otimes c=(b\otimes1)(1\otimes c)$.
\end{proof}

\begin{corollary}[3.2.3]
\label{I.3.2.3}
Sea $T$ un esquema afín de anillo $D$, y sea $\alpha=({}^a\rho,\widetilde{\rho})$
\textup{(resp. $\beta=({}^a\sigma,\widetilde{\sigma})$)} un $S$-morfismo $T\to X$ \textup{(resp. $T\to Y$)},
donde $\rho$ \textup{(resp. $\sigma$)} es un $A$-homomorfismo de $B$ \textup{(resp. $C$)} en $D$; entonces
$(\alpha,\beta)_S=({}^a\tau,\widetilde{\tau})$, donde $\tau$ es el homomorfismo
$B\otimes_A C\to D$ tal que $\tau(b\otimes c)=\rho(b)\sigma(c)$.
\end{corollary}

\begin{proposition}[3.2.4]
\label{I.3.2.4}
Sea $f:S'\to S$ un \textup{monomorfismo} de preesquemas \textup{(T, I, 1.1)}, y sean $X$ e $Y$
$S'$-preesquemas, considerados también como $S$-preesquemas mediante $f$. Todo producto de
los $S$-preesquemas $X$ e $Y$ es entonces un producto de los $S'$-preesquemas $X$ e $Y$, y recíprocamente.
\end{proposition}

\begin{proof}
Sean $\varphi:X\to S'$ y $\psi:Y\to S'$ los morfismos estructurales. Si $T$ es un $S$-preesquema, y
$u:T\to X$ y $v:T\to Y$ son $S$-morfismos, se tiene por definición
$f\circ\varphi\circ u=f\circ\psi\circ v=\theta$, morfismo estructural de $T$; la hipótesis
sobre $f$ implica que $\varphi\circ u=\psi\circ v=\theta'$, de modo que se puede considerar $T$ como un
$S'$-preesquema cuyo morfismo estructural es $\theta'$, y $u$ y $v$ como $S'$-morfismos. La
conclusión de la proposición resulta inmediatamente, teniendo en cuenta \sref{I.3.2.1}.
\end{proof}

\begin{corollary}[3.2.5]
\label{I.3.2.5}
Sean $X$ e $Y$ dos $S$-preesquemas, con morfismos estructurales $\varphi:X\to S$ y $\psi:Y\to S$, y sea
$S'$ una parte abierta de $S$ tal que $\varphi(X)\subset S'$ y $\psi(Y)\subset S'$. Todo producto
de los $S$-preesquemas $X$ e $Y$ es entonces también un producto de los $S'$-preesquemas $X$ e $Y$, y
recíprocamente.
\end{corollary}

\oldpage[I]{106}
Basta aplicar \sref{I.3.2.4} a la inyección canónica $S'\to S$.

\begin{theorem}[3.2.6]
\label{I.3.2.6}
Dados dos $S$-preesquemas $X$ e $Y$, existe un producto $X\times_S Y$.
\end{theorem}

Procederemos en varias etapas.
\begin{lemma}[3.2.6.1]
\label{I.3.2.6.1}
Sean $(Z,p,q)$ un producto de $X$ e $Y$, y $U$ y $V$ partes abiertas de $X$ e $Y$,
respectivamente. Si se pone $W=p^{-1}(U)\cap q^{-1}(V)$, la terna constituida por $W$ y
las restricciones de $p$ y $q$ a $W$ \emph{(consideradas como los morfismos $W\to U$ y $W\to V$,
respectivamente)} es un producto de $U$ y $V$.
\end{lemma}

En efecto, si $T$ es un $S$-preesquema, se puede identificar los $S$-morfismos $T\to W$ con los
$S$-morfismos $T\to Z$ que aplican $T$ en $W$. Entonces, si $g:T\to U$ y $h:T\to V$ son dos
$S$-morfismos cualesquiera, se puede considerarlos como $S$-morfismos de $T$ en $X$ y en $Y$,
respectivamente, y, por hipótesis, existe un único $S$-morfismo $f:T\to Z$ tal que $g=p\circ f$ y
$h=q\circ f$. Como $p(f(T))\subset U$ y $q(f(T))\subset V$, se tiene
\[
  f(T)\subset p^{-1}(U)\cap q^{-1}(V)=W
\]
de donde resulta nuestra afirmación.

\begin{lemma}[3.2.6.2]
\label{I.3.2.6.2}
Sean $Z$ un $S$-preesquema, $p:Z\to X$ y $q:Z\to Y$ dos $S$-morfismos, $(U_\alpha)$ un recubrimiento
abierto de $X$ y $(V_\lambda)$ un recubrimiento abierto de $Y$. Se supone que, para cada par
$(\alpha,\lambda)$, el $S$-preesquema
$W_{\alpha\lambda}=p^{-1}(U_\alpha)\cap q^{-1}(V_\lambda)$ y las restricciones de $p$ y
$q$ a $W_{\alpha\lambda}$ constituyen un producto de $U_\alpha$ y $V_\lambda$. Entonces $(Z,p,q)$ es un
producto de $X$ e $Y$.
\end{lemma}

Mostremos primero que, si $f_1$ y $f_2$ son $S$-morfismos $T\to Z$, las relaciones
$p\circ f_1=p\circ f_2$ y $q\circ f_1=q\circ f_2$ implican que $f_1=f_2$. En efecto, $Z$ es la unión
de los $W_{\alpha\lambda}$, por lo que los $f_1^{-1}(W_{\alpha\lambda})$ forman un recubrimiento abierto de $T$,
y análogamente para $f_2^{-1}(W_{\alpha\lambda})$. Además, se tiene
\[
\begin{aligned}
  f_1^{-1}(W_{\alpha\lambda})
  &=f_1^{-1}(p^{-1}(U_\alpha))\cap
    f_1^{-1}(q^{-1}(V_\lambda)) \\
  &=f_2^{-1}(p^{-1}(U_\alpha))\cap
    f_2^{-1}(q^{-1}(V_\lambda))
   =f_2^{-1}(W_{\alpha\lambda})
\end{aligned}
\]
por hipótesis, y todo se reduce, pues, a observar que las restricciones de $f_1$ y $f_2$ a
$f_1^{-1}(W_{\alpha\lambda})=f_2^{-1}(W_{\alpha\lambda})$ son idénticas para cada par de
índices. Pero, como estas restricciones pueden considerarse $S$-morfismos de
$f_1^{-1}(W_{\alpha\lambda})$ en $W_{\alpha\lambda}$, nuestra afirmación se deduce de las
hipótesis y de la Definición~\sref{I.3.2.1}.

Supongamos ahora dados dos $S$-morfismos $g:T\to X$ y $h:T\to Y$. Pongamos
$T_{\alpha\lambda}=g^{-1}(U_\alpha)\cap h^{-1}(V_\lambda)$; los $T_{\alpha\lambda}$ forman entonces un
recubrimiento abierto de $T$. Por hipótesis, existe un $S$-morfismo $f_{\alpha\lambda}$ tal que
$p\circ f_{\alpha\lambda}$ y $q\circ f_{\alpha\lambda}$ son, respectivamente, las restricciones de
$g$ y $h$ a $T_{\alpha\lambda}$. Mostremos además que las restricciones de
$f_{\alpha\lambda}$ y de $f_{\beta\mu}$ a $T_{\alpha\lambda}\cap T_{\beta\mu}$ coinciden,
lo que terminará de probar el Lema~\sref{I.3.2.6.2}. Las imágenes de
$T_{\alpha\lambda}\cap T_{\beta\mu}$ por $f_{\alpha\lambda}$ y $f_{\beta\mu}$ están
contenidas en $W_{\alpha\lambda}\cap W_{\beta\mu}$ por definición. Como
\[
  W_{\alpha\lambda}\cap W_{\beta\mu}
  =p^{-1}(U_\alpha\cap U_\beta)\cap
   q^{-1}(V_\lambda\cap V_\mu)
\]
se sigue del Lema~\sref{I.3.2.6.1} que $W_{\alpha\lambda}\cap W_{\beta\mu}$ y las
restricciones de $p$ y $q$ a este preesquema constituyen un \emph{producto} de $U_\alpha\cap U_\beta$
y $V_\lambda\cap V_\mu$. Como $p\circ f_{\alpha\lambda}$ y $p\circ f_{\beta\mu}$ coinciden
en $T_{\alpha\lambda}\cap T_{\beta\mu}$, y lo mismo sucede con $q\circ f_{\alpha\lambda}$ y
$q\circ f_{\beta\mu}$, se ve que $f_{\alpha\lambda}$ y $f_{\beta\mu}$ coinciden en
$T_{\alpha\lambda}\cap T_{\beta\mu}$.\qed


\oldpage[I]{107}
\begin{lemma}[3.2.6.3]
\label{I.3.2.6.3}

Sean $(U_\alpha)$ un recubrimiento abierto de $X$ y $(V_\lambda)$ un recubrimiento abierto de $Y$, y supongamos
que, para cada par $(\alpha,\lambda)$, existe un producto de $U_\alpha$ y $V_\lambda$;
entonces existe un producto de $X$ e $Y$.
\end{lemma}

Aplicando el Lema~\sref{I.3.2.6.1} a los abiertos $U_\alpha\cap U_\beta$ y
$V_\lambda\cap V_\mu$, se ve que existe un producto de los $S$-preesquemas inducidos,
respectivamente, por $X$ e $Y$ sobre estos abiertos; además, la unicidad del producto
muestra que, si se pone $i=(\alpha,\lambda)$ y $j=(\beta,\mu)$, existe un isomorfismo
canónico $h_{ij}$ (resp. $h_{ji}$) de este producto sobre un $S$-preesquema $W_{ij}$
(resp. $W_{ji}$) inducido por $U_\alpha\times_S V_\lambda$ (resp. $U_\beta\times_S V_\mu$) sobre
un abierto; entonces $f_{ij}=h_{ij}\circ h_{ji}^{-1}$ es un isomorfismo de $W_{ji}$ sobre
$W_{ij}$. Además, para un tercer par $k=(\gamma,\nu)$, se tiene
$f_{ik}=f_{ij}\circ f_{jk}$ en
$W_{ki}\cap W_{kj}$, aplicando el Lema \sref{I.3.2.6.1} a
los abiertos $U_\alpha\cap U_\beta\cap U_\gamma$ y $V_\lambda\cap V_\mu\cap V_\nu$ de
$U_\beta$ y $V_\mu$, respectivamente. Existe, por consiguiente, un preesquema $Z$, un recubrimiento
abierto $(Z_i)$ del espacio subyacente de $Z$ y, para cada $i$, un isomorfismo $g_i$ del
preesquema inducido
$Z_i$ sobre el preesquema $U_\alpha\times_S V_\lambda$, tal que, para cada par $(i,j)$, se tiene
$f_{ij}=g_i\circ g_j^{-1}$ \sref{I.2.3.1}; además, se tiene $g_i(Z_i\cap Z_j)=W_{ij}$.
Sean $p_i$ y $q_i$ las proyecciones, y $\theta_i$ el morfismo estructural, del
$S$-preesquema $U_\alpha\times_S V_\lambda$. Se ve inmediatamente que
$p_i\circ g_i=p_j\circ g_j$ en $Z_i\cap Z_j$, y análogamente para los otros dos morfismos. Se puede
definir así los morfismos de preesquemas $p:Z\to X$ (resp. $q:Z\to Y$, $\theta:Z\to S$) mediante
la condición de que $p$ (resp. $q$, $\theta$) coincida con $p_i\circ g_i$
(resp. $q_i\circ g_i$, $\theta_i\circ g_i$) en cada uno de los $Z_i$; $Z$, provisto de
$\theta$, es entonces un $S$-preesquema. Mostremos ahora que
$Z'_i=p^{-1}(U_\alpha)\cap q^{-1}(V_\lambda)$ es igual a $Z_i$. Para cada índice
$j=(\beta,\mu)$, se tiene $Z_j\cap Z'_i=g_j^{-1}(p_j^{-1}(U_\alpha)\cap q_j^{-1}(V_\lambda))$.
Se tiene
\[
  p_j^{-1}(U_\alpha)\cap q_j^{-1}(V_\lambda)
  =p_j^{-1}(U_\alpha\cap U_\beta)\cap
   q_j^{-1}(V_\lambda\cap V_\mu)~;
\]
Por el Lema~\sref{I.3.2.6.1}, las restricciones de $p_j$ y $q_j$ a
$p_j^{-1}(U_\alpha)\cap q_j^{-1}(V_\lambda)$ definen sobre este $S$-preesquema la estructura de un
producto de $U_\alpha\cap U_\beta$ y $V_\lambda\cap V_\mu$; pero la unicidad del
producto implica entonces que $p_j^{-1}(U_\alpha)\cap q_j^{-1}(V_\lambda)=W_{ji}$. Por consiguiente, se
tiene $Z_j\cap Z'_i=Z_j\cap Z_i$ para cada $j$, y, por tanto, $Z'_i=Z_i$. Se deduce entonces del
Lema~\sref{I.3.2.6.2} que $(Z,p,q)$ es un producto de $X$ e $Y$.

\begin{lemma}[3.2.6.4]
\label{I.3.2.6.4}
Sean $\varphi:X\to S$ y $\psi:Y\to S$ los morfismos estructurales de $X$ e $Y$, $(S_i)$ un
recubrimiento abierto de $S$, y pongamos $X_i=\varphi^{-1}(S_i)$ e $Y_i=\psi^{-1}(S_i)$. Si existe cada uno
de los productos $X_i\times_S Y_i$, entonces existe $X\times_S Y$.
\end{lemma}

Por el lema~\sref{I.3.2.6.3}, basta probar que existen los productos
$X_i\times_S Y_j$ para todos $i$ y $j$. Pongamos
$X_{ij}=X_i\cap X_j=\varphi^{-1}(S_i\cap S_j)$, $Y_{ij}=Y_i\cap Y_j=\psi^{-1}(S_i\cap S_j)$;
por el lema~\sref{I.3.2.6.1}, existe el producto $Z_{ij}=X_{ij}\times_S Y_{ij}$. Observemos ahora
que, si $T$ es un $S$-preesquema y $g:T\to X_i$ y $h:T\to Y_j$ son $S$-morfismos,
entonces necesariamente $\varphi(g(T))=\psi(h(T))\subset S_i\cap S_j$ por la
definición de $S$-morfismo, y por tanto $g(T)\subset X_{ij}$ y $h(T)\subset Y_{ij}$; es
entonces inmediato que $Z_{ij}$ es el producto de $X_i$ e $Y_j$.

\begin{env}[3.2.6.5]
\label{I.3.2.6.5}
Podemos ya terminar la demostración del teorema~\sref{I.3.2.6}. Si $S$ es un
\emph{esquema afín}, existen recubrimientos $(U_\alpha)$ y $(V_\lambda)$ de $X$ e $Y$, respectivamente,
formados por partes abiertas afines; como $U_\alpha\times_S V_\lambda$ existe en virtud de
\sref{I.3.2.2}, el lema~\sref{I.3.2.6.3} muestra que existe $X\times_S Y$. Si $S$
es un preesquema cualquiera, existe un recubrimiento $(S_i)$ de $S$ formado por partes abiertas afines. Si
$\varphi:X\to S$ y $\psi:Y\to S$ son los morfismos estructurales, y si se pone
$X_i=\varphi^{-1}(S_i)$ e $Y_i=\psi^{-1}(S_i)$, entonces existen los productos $X_i\times_{S_i}Y_i$, por lo
\oldpage[I]{108}
anterior; pero entonces también existen los productos $X_i\times_S Y_i$ \sref{I.3.2.5}, y
el lema~\sref{I.3.2.6.4} muestra por tanto que existe $X\times_S Y$.\qed
\end{env}

\begin{corollary}[3.2.7]
\label{I.3.2.7}
Sea $Z=X\times_S Y$ el producto de dos $S$-preesquemas, sean $p$ y $q$ las proyecciones de $Z$
sobre $X$ y sobre $Y$, respectivamente, y sea $\varphi$ \textup{(resp. $\psi$)} el morfismo estructural de
$X$ (resp. de $Y$). Sea $S'$ una
parte abierta de $S$, y sea $U$ \textup{(resp. $V$)} una parte abierta de $X$ (resp. de $Y$) contenida en
$\varphi^{-1}(S')$ \textup{(resp.} en $\psi^{-1}(S')$\textup{)}. Entonces el producto $U\times_{S'}V$ se identifica
canónicamente con el preesquema inducido por $Z$ sobre $p^{-1}(U)\cap q^{-1}(V)$ (considerado como un
$S'$-preesquema). Además, si $f:T\to X$ y $g:T\to Y$ son $S$-morfismos tales que
$f(T)\subset U$ y $g(T)\subset V$, entonces el $S'$-morfismo $(f,g)_{S'}$ se identifica con
$(f,g)_S$ considerado como un morfismo con valores en $p^{-1}(U)\cap q^{-1}(V)$.
\end{corollary}

\begin{proof}
Esto resulta del corolario~\sref{I.3.2.5} y del lema~\sref{I.3.2.6.1}.
\end{proof}

\begin{env}[3.2.8]
\label{I.3.2.8}
Sean $(X_\alpha)$ e $(Y_\lambda)$ familias de $S$-preesquemas, y sean $X$ \textup{(resp. $Y$)} la suma
de la familia $(X_\alpha)$ (resp. de $(Y_\lambda)$) \hyperref[subsection:I.3.1]{(3.1)}.
Entonces $X\times_S Y$ se identifica con la \emph{suma} de la familia $(X_\alpha\times_S Y_\lambda)$; esto resulta inmediatamente del lema~\sref{I.3.2.6.3}.
\end{env}

\begin{env}[3.2.9]
\label{I.3.2.9}
\footnote{\emph{[Trad.] \sref{I.3.2.9} procede de la fe de erratas de EGA~II, en la página 221, de ahí el cambio en la numeración de las páginas.}}
\oldpage[II]{221}
Se deduce de \sref{I.1.8.1} que podemos precisar \sref{I.3.2.2} de la manera
siguiente: $Z=\Spec(B\otimes_A C)$ no solo es un producto de $X=\Spec(B)$ e $Y=\Spec(C)$ en la
categoría de los \emph{$S$-preesquemas}, sino también en la categoría de los \emph{espacios localmente anillados
sobre $S$} (con una definición de los $S$-morfismos calcada sobre la de \sref{I.2.5.2}). La
demostración de \sref{I.3.2.6} prueba también que, para dos $S$-preesquemas cualesquiera $X$ e $Y$, el
preesquema $X\times_S Y$ no solo es el producto de $X$ e $Y$ en la categoría de los
$S$-preesquemas, sino también en la categoría de los espacios localmente anillados sobre el preesquema $S$.
\end{env}

\subsection{Propiedades formales del producto; cambio del preesquema base}
\label{subsection:I.3.3}

\begin{env}[3.3.1]
\label{I.3.3.1}
El lector observará que todas las propiedades enunciadas en esta sección, salvo
\sref{I.3.3.13}, \sref{I.3.3.15} y la afirmación sobre las sumas de la última frase
de \sref{I.3.3.10}, son válidas sin modificación en \emph{toda categoría}, siempre que
los productos que intervienen en los enunciados \emph{existan} (pues está claro
que las nociones de $S$-objeto y de $S$-morfismo pueden definirse exactamente
como en \hyperref[subsection:I.2.5]{(2.5)} para todo objeto $S$ de la categoría).
\end{env}

\begin{env}[3.3.2]
\label{I.3.3.2}
En primer lugar, $X\times_S Y$ es un \emph{bifuntor covariante} en $X$ e $Y$ en la
categoría de los $S$-preesquemas: basta, en efecto, observar que el diagrama
\[
  \xymatrix{
    X\times Y\ar[r]^{f\times 1}\ar[d] &
    X'\times Y\ar[r]^{f'\times 1}\ar[d] &
    X''\times Y\ar[d]\\
    X\ar[r]_f & X'\ar[r]_{f'} & X''
  }
\]
es conmutativo.
\end{env}

\begin{proposition}[3.3.3]
\label{I.3.3.3}
Para todo $S$-preesquema $X$, la primera \textup{(resp.} segunda) proyección de
$X\times_S S$ \textup{(resp. $S\times_S X$)} es un isomorfismo funtorial de
$X\times_S S$ \textup{(resp. $S\times_S X$)} sobre $X$, cuyo isomorfismo inverso es
$(1_X,\varphi)_S$ \textup{(resp. $(\varphi,1_X)_S$)}, donde se denota por $\varphi$ el
morfismo estructural $X\to S$; se puede escribir, por tanto, salvo isomorfismo
canónico,
\[
  X\times_S S=S\times_S X=X.
\]
\end{proposition}

\begin{proof}
Basta probar que la terna $(X,1_X,\varphi)$ es un producto de $X$ y $S$. Si $T$ es
un $S$-preesquema, el único $S$-morfismo de $T$ en $S$ es necesariamente el
morfismo estructural $\psi:T\to S$. Si $f$ es un $S$-morfismo de $T$ en $X$,
se tiene necesariamente $\psi=\varphi\circ f$, de donde resulta la afirmación.
\end{proof}

\begin{corollary}[3.3.4]
\label{I.3.3.4}
Sean $X$ e $Y$ dos $S$-preesquemas, y sean $\varphi:X\to S$ y $\psi:Y\to S$ sus
morfismos estructurales. Si se identifican canónicamente $X$ con $X\times_S S$
e $Y$ con $S\times_S Y$, las proyecciones $X\times_S Y\to X$ y
$X\times_S Y\to Y$ se identifican respectivamente con $1_X\times\psi$ y
$\varphi\times 1_Y$.
\end{corollary}

La verificación es inmediata y se deja al lector.

\begin{env}[3.3.5]
\label{I.3.3.5}
Se puede definir, del mismo modo que en \hyperref[subsection:I.3.2]{(3.2)}, el producto de un
\oldpage[I]{109}
número finito cualquiera $n$ de $S$-preesquemas, y la existencia de estos productos se deduce
de \sref{I.3.2.6} por inducción sobre $n$, observando que
$(X_1\times_S X_2\times\cdots\times_S X_{n-1})\times_S X_n$ satisface la
definición del producto. La unicidad del producto implica, como en toda
categoría, sus propiedades de \emph{conmutatividad} y de \emph{asociatividad}. Si, por
ejemplo, $p_1$, $p_2$ y $p_3$ denotan las proyecciones de
$X_1\times X_2\times X_3$, y se identifica este preesquema con
$(X_1\times X_2)\times X_3$, la proyección sobre $X_1\times X_2$ se
identifica con $(p_1,p_2)_S$.
\end{env}

\begin{env}[3.3.6]
\label{I.3.3.6}
Sean $S$ y $S'$ dos preesquemas, y sea $\varphi:S'\to S$ un morfismo, que hace de $S'$ un
$S$-preesquema. Para todo $S$-preesquema $X$, consideremos el producto $X\times_S S'$,
y sean $p$ y $\pi'$ sus proyecciones sobre $X$ y $S'$ respectivamente. Provisto
de $\pi'$, este producto es un $S'$-preesquema; cuando se lo considera como tal, se
lo denota por $X_{(S')}$ o $X_{(\varphi)}$, y se dice que es el preesquema
obtenido por \emph{cambio de base} de $S$ a $S'$ mediante el morfismo
$\varphi$, o la \emph{imagen inversa} de $X$ por $\varphi$. Se observará que, si $\pi$ es
el morfismo estructural de $X$ y $\theta$ el morfismo estructural de
$X\times_S S'$, considerado como $S$-preesquema, el diagrama
\phantomsection\label{I.3.3.6.base-change-diagram-fr}
\[
  \xymatrix{
    X\ar[d]_{\pi} &
    X_{(S')}\ar[l]_{p}\ar[ld]_{\theta}\ar[d]^{\pi'}\\
    S & S'\ar[l]^{\varphi}
  }
\]
es conmutativo.
\end{env}

\begin{env}[3.3.7]
\label{I.3.3.7}
Con las notaciones de \sref{I.3.3.6}, para todo $S$-morfismo $f:X\to Y$, se
denota también por $f_{(S')}$ el $S'$-morfismo $f\times_S 1:X_{(S')}\to Y_{(S')}$, y
se dice que $f_{(S')}$ es el \emph{cambio de base} (o la \emph{imagen inversa}) de
$f$ por $\varphi$. Por tanto, $X_{(S')}$ es un \emph{funtor covariante} en $X$, de
la categoría de los $S$-preesquemas en la de los $S'$-preesquemas.
\end{env}

\begin{env}[3.3.8]
\label{I.3.3.8}
El preesquema $X_{(S')}$ puede considerarse también como solución de un
\emph{problema de aplicación universal}: todo $S'$-preesquema $T$ es asimismo un
$S$-preesquema por medio de $\varphi$; todo $S$-morfismo $g:T\to X$ se escribe entonces
de manera única como $g=p\circ f$, donde $f$ es un $S'$-morfismo
$T\to X_{(S')}$, como resulta de la definición del producto aplicada a los
$S$-morfismos $g$\footnote{\emph{[Trad.] El texto francés imprime aquí $f$;
la pareja que interviene en la definición es $(g,\psi)$.}} y
$\psi:T\to S'$ (el morfismo estructural de $T$).
\end{env}

\begin{proposition}[3.3.9]
\label{I.3.3.9}
\emph{(``Transitividad del cambio de base'')}. Sea $S''$ un preesquema, y
$\varphi':S''\to S'$ un morfismo. Para todo $S$-preesquema $X$, existe un
isomorfismo canónico funtorial del $S''$-preesquema
$(X_{(\varphi)})_{(\varphi')}$ sobre el $S''$-preesquema $X_{(\varphi\circ\varphi')}$.
\end{proposition}

\begin{proof}
Sea $T$ un $S''$-preesquema, $\psi$ su morfismo estructural, y $g$ un
$S$-morfismo de $T$ en $X$ (considerando $T$ como $S$-preesquema con
morfismo estructural $\varphi\circ\varphi'\circ\psi$). Puesto que $T$ es también un $S'$-preesquema
con morfismo estructural $\varphi'\circ\psi$, se puede escribir $g=p\circ g'$, donde
$g'$ es un $S'$-morfismo $T\to X_{(\varphi)}$, y luego $g'=p'\circ g''$, donde
$g''$ es un $S''$-morfismo $T\to(X_{(\varphi)})_{(\varphi')}$:
\phantomsection\label{I.3.3.9.transitivity-diagram-fr}
\[
  \xymatrix{
    X\ar[d]_{\pi} &
    X_{(\varphi)}\ar[l]_{p}\ar[d]_{\pi'} &
    (X_{(\varphi)})_{(\varphi')}\ar[l]_{p'}\ar[d]_{\pi''}\\
    S & S'\ar[l]^{\varphi} & S''\ar[l]^{\varphi'}
  }
\]
\oldpage[I]{110}
Por tanto, el resultado se sigue de la unicidad de la solución de un problema de
aplicación universal.
\end{proof}

Este resultado puede escribirse como la igualdad (salvo un isomorfismo canónico)
$(X_{(S')})_{(S'')}=X_{(S'')}$ (si no hay riesgo de confusión), o también como
\begin{equation}
  (X\times_S S')\times_{S'}S''=X\times_S S''~;
  \tag{3.3.9.1}\label{I.3.3.9.1-fr}
\end{equation}
la naturaleza funtorial del isomorfismo definido en \sref{I.3.3.9} puede
también expresarse mediante la fórmula de \emph{transitividad} de los morfismos de
cambio de base
\begin{equation}
  (f_{(S')})_{(S'')}=f_{(S'')}
  \tag{3.3.9.2}\label{I.3.3.9.2-fr}
\end{equation}
para cada $S$-morfismo $f:X\to Y$.

\begin{corollary}[3.3.10]
\label{I.3.3.10}
Si $X$ e $Y$ son $S$-preesquemas, existe un isomorfismo canónico funtorial
del $S'$-preesquema $X_{(S')}\times_{S'}Y_{(S')}$ sobre el $S'$-preesquema
$(X\times_S Y)_{(S')}$. 
\end{corollary}

\begin{proof}
Se tiene, salvo isomorfismos canónicos,
\[
  (X\times_S S')\times_{S'}(Y\times_S S')
  =X\times_S(Y\times_S S')=(X\times_S Y)\times_S S'
\]
en virtud de \sref{I.3.3.9.1} y de la asociatividad de los productos de
$S$-preesquemas.
\end{proof}

La naturaleza funtorial del isomorfismo definido en
el corolario~\sref{I.3.3.10} puede expresarse mediante la fórmula
\begin{equation}
  (u_{(S')},v_{(S')})_{S'}=((u,v)_S)_{(S')}
  \tag{3.3.10.1}\label{I.3.3.10.1-fr}
\end{equation}
para cada par de $S$-morfismos $u:T\to X$, $v:T\to Y$.

En otros términos, el funtor de cambio de base $X_{(S')}$ \emph{conmuta con
los productos}; también conmuta con las sumas \sref{I.3.2.8}.

\begin{corollary}[3.3.11]
\label{I.3.3.11}
Sea $Y$ un $S$-preesquema, y sea $f:X\to Y$ un morfismo que hace de $X$ un
$Y$-preesquema (y, por consiguiente, también un $S$-preesquema). El preesquema $X_{(S')}$
se identifica entonces con el producto $X\times_Y Y_{(S')}$, identificándose la proyección
$X\times_Y Y_{(S')}\to Y_{(S')}$ con $f_{(S')}$. 
\end{corollary}

\begin{proof}
Sea $\psi:Y\to S$ el morfismo estructural de $Y$; se tiene el diagrama conmutativo
\phantomsection\label{I.3.3.11.diagram-fr}
\[
  \xymatrix{
    S'\ar[d] &
    Y_{(S')}\ar[l]_{\psi_{(S')}}\ar[d] &
    X_{(S')}\ar[l]_{f_{(S')}}\ar[d]\\
    S & Y\ar[l]^{\psi} & X\ar[l]^{f}
  }
\]
Ahora bien, $Y_{(S')}$ se identifica con $S'_{(\psi)}$, y $X_{(S')}$ con
$S'_{(\psi\circ f)}$; teniendo en cuenta \sref{I.3.3.9} y
\sref{I.3.3.4}, se deduce así el corolario.
\end{proof}
\begin{env}[3.3.12]
\label{I.3.3.12}
Sean $f:X\to X'$ y $g:Y\to Y'$ $S$-morfismos que son \emph{monomorfismos}
de preesquemas (T, I, 1.1); entonces $f\times_S g$ es un \emph{monomorfismo}. En efecto,
si $p$ y $q$ son las proyecciones de $X\times_S Y$, $p'$ y $q'$ las proyecciones de
$X'\times_S Y'$, y $u$ y $v$ son $S$-morfismos $T\to X\times_S Y$, entonces la
ecuación $(f\times_S g)\circ u=(f\times_S g)\circ v$ implica que
$p'\circ(f\times_S g)\circ u=p'\circ(f\times_S g)\circ v$, o, dicho de otro modo, que
$f\circ p\circ u=f\circ p\circ v$, y como $f$ es un monomorfismo,
$p\circ u=p\circ v$; usando el hecho de que $g$ es un monomorfismo, se obtiene de manera análoga
$q\circ u=q\circ v$, de donde $u=v$.

\oldpage[I]{111}
Se sigue que, para cada cambio de base $S'\to S$,
\phantomsection\label{I.3.3.12.base-change-monomorphism-fr}
\[
  f_{(S')}:X_{(S')}\to X'_{(S')}
\]
es un monomorfismo.\footnote{\emph{[Trad.] La fórmula francesa imprime
$Y_{(S')}$; puesto que 3.3.12 supone $f:X\to X'$, el codominio tras el cambio de base es
$X'_{(S')}$.}}
\end{env}

\begin{env}[3.3.13]
\label{I.3.3.13}
Sean $S$ y $S'$ esquemas afines de anillos $A$ y $A'$ (respectivamente); un morfismo
$S'\to S$ corresponde entonces a un homomorfismo de anillos $A\to A'$. Si $X$ es un
$S$-preesquema, se denota por $X_{(A')}$ o $X\otimes_A A'$ el $S'$-preesquema
$X_{(S')}$; cuando $X$ es también afín de anillo $B$, $X_{(A')}$ es afín de anillo
$B_{(A')}=B\otimes_A A'$ obtenido extendiendo los escalares del $A$-álgebra
$B$ a $A'$.
\end{env}

\begin{env}[3.3.14]
\label{I.3.3.14}
Con las notaciones de \sref{I.3.3.6}, para cada \emph{$S$-morfismo}
$f:S'\to X$, $f'=(f,1_{S'})_S$ es un $S'$-morfismo $S'\to X'=X_{(S')}$ tal que
$p\circ f'=f$, $\pi'\circ f'=1_{S'}$, o, en otras palabras, una
\emph{$S'$-sección de $X'$}; recíprocamente, si $f'$ es tal $S'$-sección, entonces
$f=p\circ f'$ es un
$S$-morfismo $S'\to X$. Se define así una \emph{correspondencia biyectiva}
canónica
\phantomsection\label{I.3.3.14.correspondence-fr}
\[
  \operatorname{Hom}_S(S',X)\overset{\sim}{\longrightarrow}
  \operatorname{Hom}_{S'}(S',X')
\]
Se dice que $f'$ es el \emph{morfismo grafo} de $f$, y se lo denota por
$\Gamma_f$.
\end{env}

\begin{env}[3.3.15]
\label{I.3.3.15}
Dado un preesquema $X$, que siempre se puede considerar como un $\mathbf{Z}$-preesquema,
se sigue, en particular, de \sref{I.3.3.14} que las $X$-\emph{secciones} de
$X\otimes_{\mathbf{Z}}\mathbf{Z}[T]$ (donde $T$ es una indeterminada)
corresponden biyectivamente a \emph{morfismos} $X\to\mathbf{Z}[T]$\footnote{\emph{[Trad.]
El texto francés imprime $\bb{Z}[T]\to X$; la Proposición~\sref{I.3.3.14} y los
homomorfismos de anillos locales subsiguientes $A_\alpha[T]\to A_\alpha$ requieren el morfismo de
esquemas $X\to\bb{Z}[T]$.}}. Demostremos que estas
$X$-secciones también corresponden biyectivamente a \emph{secciones del haz estructural
$\mathcal{O}_X$ sobre $X$}. En efecto, sea $(U_\alpha)$ un recubrimiento de $X$ por abiertos
afines; sea $u:X\to X\otimes_{\mathbf{Z}}\mathbf{Z}[T]$ un $X$-morfismo, y sea
$u_\alpha$ su restricción a $U_\alpha$; si $A_\alpha$ es el anillo del esquema
afín $U_\alpha$, entonces $U_\alpha\otimes_{\mathbf{Z}}\mathbf{Z}[T]$ es un esquema afín
de anillo $A_\alpha[T]$ \sref{I.3.2.2}, y $u_\alpha$ corresponde canónicamente a un
$A_\alpha$-homomorfismo $A_\alpha[T]\to A_\alpha$
\sref{I.1.7.3}. Ahora bien, puesto que tal homomorfismo está completamente determinado por
los datos de la imagen de $T$ en $A_\alpha$, sea
$s_\alpha\in A_\alpha=\Gamma(U_\alpha,\mathcal{O}_X)$, y si se supone que las
restricciones de $u_\alpha$ y $u_\beta$ a un abierto afín
$V\subset U_\alpha\cap U_\beta$ coinciden, se ve inmediatamente que
$s_\alpha$ y $s_\beta$ coinciden en $V$; por tanto, la familia $(s_\alpha)$ está formada
por las restricciones a $U_\alpha$ de una sección $s$ de $\mathcal{O}_X$ sobre $X$;
recíprocamente, es claro que tal sección define una familia $(u_\alpha)$ de
morfismos que son las restricciones a $U_\alpha$ de un $X$-morfismo
$X\to X\otimes_{\mathbf{Z}}\mathbf{Z}[T]$. Este resultado se generaliza en
\sref[II]{II.1.7.12}.
\end{env}

\subsection{Puntos de un preesquema con valores en un preesquema; puntos geométricos}
\label{subsection:I.3.4}

\begin{env}[3.4.1]
\label{I.3.4.1}
Sea $X$ un preesquema; para cada preesquema $T$, se denota entonces por $X(T)$ el conjunto
$\operatorname{Hom}(T,X)$ de morfismos $T\to X$, y los elementos de este conjunto se llaman
\emph{los puntos de $X$ con valores en $T$}. Si se asocia a cada morfismo
$f:T\to T'$ la aplicación $u'\to u'\circ f$ de $X(T')$ en $X(T)$, se ve,
para $X$ fijo, que $X(T)$ es un \emph{funtor contravariante en $T$}, de la
categoría de los preesquemas en la de los conjuntos. Además, cada morfismo de preesquemas
$g:X\to Y$ define un homomorfismo funtorial $X(T)\to Y(T)$, que envía
$v\in X(T)$ a $g\circ v$.
\end{env}

\begin{env}[3.4.2]
\label{I.3.4.2}
Dados los conjuntos $P$, $Q$ y $R$, y las aplicaciones $\varphi:P\to R$ y $\psi:Q\to R$, se define el
\emph{producto fibrado de $P$ y $Q$ sobre $R$} (con respecto a $\varphi$ y $\psi$) como el subconjunto
del
\oldpage[I]{112}
conjunto producto $P\times Q$ formado por los pares $(p,q)$ tales que $\varphi(p)=\psi(q)$; se lo denota por $P\times_R Q$. La Definición~\sref{I.3.2.1} del producto
de $S$-preesquemas puede interpretarse, con la notación de \sref{I.3.4.1},
mediante la fórmula
\phantomsection\label{I.3.4.2.1-fr}
\[
  (X\times_S Y)(T)=X(T)\times_{S(T)}Y(T)
  \tag{3.4.2.1}
\]
donde las aplicaciones $X(T)\to S(T)$ y $Y(T)\to S(T)$ corresponden a los morfismos
estructurales $X\to S$ y $Y\to S$.
\end{env}

\begin{env}[3.4.3]
\label{I.3.4.3}
Si se nos da un preesquema $S$ y se consideran solamente los $S$-preesquemas y los $S$-morfismos,
se denotará por $X(T)_S$ el conjunto $\operatorname{Hom}_S(T,X)$ de $S$-morfismos $T\to X$, y suprimiremos
el subíndice $S$ cuando no haya posibilidad de confusión; se dice que los elementos de $X(T)_S$
son los \emph{puntos} (o \emph{$S$-puntos}, cuando exista la posibilidad de confusión)
\emph{del $S$-preesquema $X$ con valores en el $S$-preesquema $T$}. En particular, una
\emph{$S$-sección} de $X$ no es otra cosa que un \emph{punto de $X$ con valores en $S$}. La
fórmula \sref{I.3.4.2.1} se puede escribir entonces como
\phantomsection\label{I.3.4.3.1-fr}
\[
  (X\times_S Y)(T)_S=X(T)_S\times Y(T)_S~;
  \tag{3.4.3.1}
\]
más generalmente, si $Z$ es un $S$-preesquema, y $X$, $Y$ y $T$ son $Z$-preesquemas (por tanto,
\emph{ipso facto}, $S$-preesquemas), se tiene
\phantomsection\label{I.3.4.3.2-fr}
\[
  (X\times_Z Y)(T)_S=X(T)_S\times_{Z(T)_S}Y(T)_S.
  \tag{3.4.3.2}
\]

Se observará que, para demostrar que un triple $(W,r,s)$ formado por un $S$-preesquema $W$ y
los $S$-morfismos $r:W\to X$ y $s:W\to Y$ es un producto de $X$ e $Y$ (sobre $Z$), basta, por
definición, comprobar que, para \emph{cada $S$-preesquema $T$}, el diagrama
\phantomsection\label{I.3.4.3.product-diagram-fr}
\[
  \xymatrix{
    W(T)_S\ar[r]^{r'}\ar[d]_{s'} & X(T)_S\ar[d]^{\varphi'}\\
    Y(T)_S\ar[r]_{\psi'} & Z(T)_S
  }
\]
hace de $W(T)_S$ el producto fibrado de $X(T)_S$ y $Y(T)_S$ sobre $Z(T)_S$, donde $r'$ y $s'$
corresponden a $r$ y $s$, y $\varphi'$ y $\psi'$ a los morfismos estructurales $\varphi:X\to Z$ y
$\psi:Y\to Z$.
\end{env}

\begin{env}[3.4.4]
\label{I.3.4.4}
Cuando $T$ (resp. $S$) de lo anterior es un esquema afín de anillo $B$ (resp. $A$), se sustituye
$T$ (resp. $S$) por $B$ (resp. $A$) en la notación anterior, y se llama entonces a los elementos de
$X(B)$ los \emph{puntos de $X$ con valores en el anillo $B$}, y a los elementos de $X(B)_A$ los
\emph{puntos del $A$-preesquema $X$ con valores en el $A$-álgebra $B$}. Se observará que
$X(B)$ y $X(B)_A$ son funtores \emph{covariantes} en $B$. Se escribirá asimismo $X(T)_A$ para el
conjunto de puntos del $A$-preesquema $X$ con valores en el $A$-preesquema $T$.
\end{env}

\begin{env}[3.4.5]
\label{I.3.4.5}
Consideremos, en particular, el caso en que $T$ es de la forma $\operatorname{Spec}(A)$, donde $A$ es un
\emph{anillo local}; los elementos de $X(A)$ corresponden biyectivamente a \emph{homomorfismos
locales} $\mathcal{O}_x\to A$ para $x\in X$ \sref{I.2.4.4}; se dice que el punto $x$ del espacio
subyacente de $X$ es la \emph{localidad}\footnote{\emph{[Trad.] También decimos que
\emph{el punto geométrico está sobre este $x$}.}} del punto de $X$ con valores en $A$ al que
corresponde.

Más precisamente, se definen los \emph{puntos geométricos} de un preesquema $X$ como los \emph{puntos de
$X$ con valores en un cuerpo $K$}: los datos de tal punto equivalen a los datos de su
\oldpage[I]{113}
localidad $x$ en el espacio subyacente de $X$, y de una \emph{extensión} $K$ de $\boldsymbol{k}(x)$;
$K$ se llamará el \emph{cuerpo de valores} del punto geométrico correspondiente, y se dice que este
punto geométrico está \emph{localizado en $x$}. Se define así una aplicación $X(K)\to X$,
que envía un punto geométrico con valores en $K$ a su localidad.

Si $S'=\operatorname{Spec}(K)$ es un $S$-preesquema (en otras palabras, si $K$ se considera como una extensión
del cuerpo residual $\boldsymbol{k}(s)$, donde $s\in S$), y si $X$ es un $S$-preesquema, entonces un
elemento de $X(K)_S$, o, como se dice, un \emph{punto geométrico de $X$ situado sobre $s$ con valores
en $K$}, consiste en los datos de un $\boldsymbol{k}(s)$-monomorfismo del cuerpo residual $\boldsymbol{k}(x)$
en $K$, donde $x$ es un punto de $X$ \emph{situado sobre $s$} (por tanto $\boldsymbol{k}(x)$ es una
extensión de $\boldsymbol{k}(s)$).

En particular, si $S=\operatorname{Spec}(K)=\{\xi\}$, \emph{los puntos geométricos de $X$ con valores en
$K$ se pueden identificar con los puntos $x\in X$ tales que $\boldsymbol{k}(x)=K$}; se dice que estos últimos
son los \emph{puntos racionales sobre $K$ del $K$-preesquema $X$}; si $K'$ es una extensión de
$K$, entonces los puntos geométricos de $X$ con valores en $K'$ corresponden biyectivamente a los
puntos $K'$-racionales de $X'=X_{(K')}$ \sref{I.3.3.14}.
\end{env}

\begin{lemma}[3.4.6]
\label{I.3.4.6}
Sean $X_i$ ($1\leq i\leq n$) $S$-preesquemas, $s$ un punto de $S$, y $x_i$
($1\leq i\leq n$) puntos de $X_i$ situados sobre $s$. Entonces existe una extensión
$K$ de $\boldsymbol{k}(s)$ y un punto geométrico del producto
$Y=X_1\times_S X_2\times\ldots\times_S X_n$, con valores en $K$, cuyas proyecciones estén
localizadas en los $x_i$.
\end{lemma}

\begin{proof}
Existen $\boldsymbol{k}(s)$-monomorfismos $\boldsymbol{k}(x_i)\to K$, todos en una misma extensión $K$ de $\boldsymbol{k}(s)$ (Bourbaki, \emph{Alg.}, chap.~V, \textsection4, prop.~2).
Las composiciones $\boldsymbol{k}(s)\to \boldsymbol{k}(x_i)\to K$ son todas idénticas y, por tanto, los morfismos $\operatorname{Spec}(K)\to X_i$ correspondientes a los morfismos $\boldsymbol{k}(x_i)\to K$ son todos $S$-morfismos; se concluye así que definen un único morfismo $\operatorname{Spec}(K)\to Y$.
Si $y$ es el punto correspondiente de $Y$, está claro que su proyección sobre cada uno de los $X_i$ es $x_i$.
\end{proof}

\begin{proposition}[3.4.7]
\label{I.3.4.7}
Sean $X_i$ ($1\leq i\leq n$) $S$-preesquemas y, para cada índice $i$, sea $x_i$ un punto de $X_i$.
Para que exista un punto $y$ de $Y=X_1\times_S X_2\times\ldots\times_S X_n$ cuya imagen sea $x_i$ por la $i$-ésima proyección para cada $1\leq i\leq n$, es necesario y suficiente que todos los $x_i$ estén situados sobre un mismo punto $s$ de $S$.
\end{proposition}

\begin{proof}
La condición es evidentemente necesaria; el Lema~\sref{I.3.4.6} demuestra que es suficiente.
\end{proof}

En otras palabras, si se denota por \phantomsection\label{I.3.4.7.underlying-map-fr}
$(X)$ el conjunto subyacente de $X$, se ve que se tiene una aplicación \emph{sobreyectiva} canónica $(X\times_S Y)\to (X)\times_{(S)}(Y)$; es necesario señalar que esta aplicación \emph{no es inyectiva} en general; dicho de otro modo, \emph{puede haber varios puntos distintos $z$ en $X\times_S Y$ que tengan las mismas proyecciones $x\in X$ e $y\in Y$}; ya se ha visto esto cuando $S$, $X$ e $Y$ son espectros primos de cuerpos $k$, $K$ y $K'$ (respectivamente), pues el producto tensorial $K\otimes_k K'$ tiene, en general, varios ideales primos distintos (cf.~\sref{I.3.4.9}).

\begin{corollary}[3.4.8]
\label{I.3.4.8}
Sea $f:X\to Y$ un $S$-morfismo, y sea $f_{(S')}:X_{(S')}\to Y_{(S')}$ el $S'$-morfismo obtenido de $f$ mediante la extensión de base $S'\to S$.
Sea $p$ (\textup{resp.} $q$) la proyección $X_{(S')}\to X$ (\textup{resp.} $Y_{(S')}\to Y$); para todo subconjunto $M$ de $X$, se tiene
\phantomsection\label{I.3.4.8.base-change-image-fr}
\[
  q^{-1}(f(M))=f_{(S')}(p^{-1}(M))
\]
\end{corollary}

\oldpage[I]{114}
\begin{proof}
En efecto, por \sref{I.3.3.11}, $X_{(S')}$ puede identificarse con el producto $X\times_Y Y_{(S')}$ gracias al diagrama conmutativo
\phantomsection\label{I.3.4.8.proof-diagram-fr}
\[
  \xymatrix{
    X\ar[d]_f & X_{(S')}\ar[l]_p\ar[d]^{f_{(S')}}\\
    Y & Y_{(S')}\ar[l]^q
  }
\]
En virtud de \sref{I.3.4.7}, la ecuación $q(y')=f(x)$ para $x\in M$ y $y'\in Y_{(S')}$ es equivalente a la existencia de algún $x'\in X_{(S')}$ tal que $p(x')=x$ y $f_{(S')}(x')=y'$, de donde se sigue el corolario.
\end{proof}

El Lema~\sref{I.3.4.6} se puede precisar de la siguiente manera:
\begin{proposition}[3.4.9]
\label{I.3.4.9}
Sean $X$ e $Y$ $S$-preesquemas, $x$ un punto de $X$ e $y$ un punto de $Y$, sobre el mismo punto $s\in S$.
El conjunto de puntos de $X\times_S Y$ cuyas proyecciones son $x$ e $y$ está en correspondencia biyectiva canónica con el conjunto de tipos de extensiones compuestas de $\boldsymbol{k}(x)$ y $\boldsymbol{k}(y)$ consideradas como extensiones de $\boldsymbol{k}(s)$ \textup{(Bourbaki, \emph{Alg.}, chap.~VIII, \textsection~8, prop.~2)}.
\end{proposition}

\begin{proof}
Sea $p$ (resp. $q$) la proyección de $X\times_S Y$ sobre $X$ (resp. $Y$), y sea $E$ el subespacio $p^{-1}(x)\cap q^{-1}(y)$ del espacio subyacente de $X\times_S Y$.
En primer lugar, observemos que los morfismos $\operatorname{Spec}(\boldsymbol{k}(x))\to S$ y $\operatorname{Spec}(\boldsymbol{k}(y))\to S$ se factorizan como $\operatorname{Spec}(\boldsymbol{k}(x))\to\operatorname{Spec}(\boldsymbol{k}(s))\to S$ y $\operatorname{Spec}(\boldsymbol{k}(y))\to\operatorname{Spec}(\boldsymbol{k}(s))\to S$; puesto que $\operatorname{Spec}(\boldsymbol{k}(s))\to S$ es un monomorfismo \sref{I.2.4.7}, se deduce de \sref{I.3.2.4} que se tiene
\phantomsection\label{I.3.4.9.tensor-product-fibre-fr}
\[
  P=\operatorname{Spec}(\boldsymbol{k}(x))\times_S\operatorname{Spec}(\boldsymbol{k}(y))
  =\operatorname{Spec}(\boldsymbol{k}(x))\times_{\operatorname{Spec}(\boldsymbol{k}(s))}
   \operatorname{Spec}(\boldsymbol{k}(y))
  =\operatorname{Spec}\bigl(\boldsymbol{k}(x)\otimes_{\boldsymbol{k}(s)}\boldsymbol{k}(y)\bigr)
\]
Definiremos dos aplicaciones, $\alpha:P_0\to E$ y $\beta:E\to P_0$, inversas entre sí (donde $P_0$ denota el conjunto subyacente del preesquema $P$).
Si $i:\operatorname{Spec}(\boldsymbol{k}(x))\to X$ y $j:\operatorname{Spec}(\boldsymbol{k}(y))\to Y$ son los morfismos canónicos \sref{I.2.4.5}, se tomará $\alpha$ como la aplicación de los espacios subyacentes correspondiente al morfismo $i\times_S j$.
Por otro lado, todo $z\in E$ define, por hipótesis, dos $\boldsymbol{k}(s)$-monomorfismos, $\boldsymbol{k}(x)\to \boldsymbol{k}(z)$ y $\boldsymbol{k}(y)\to \boldsymbol{k}(z)$, y, por tanto, un $\boldsymbol{k}(s)$-homomorfismo\footnote{\emph{[Trad.] El texto francés imprime «monomorfismo» aquí. El mapa inducido desde el producto tensorial no tiene por qué ser inyectivo; solo se necesita un homomorfismo para definir el morfismo de espectros que sigue.}} $\boldsymbol{k}(x)\otimes_{\boldsymbol{k}(s)}\boldsymbol{k}(y)\to \boldsymbol{k}(z)$ y, por consiguiente, un morfismo $\operatorname{Spec}(\boldsymbol{k}(z))\to P$; $\beta(z)$ será la imagen de $z$ en $P_0$ mediante este morfismo.
La verificación de que $\alpha\circ\beta$ y $\beta\circ\alpha$ son las aplicaciones identidad se deduce de \sref{I.2.4.5} y de la definición del producto \sref{I.3.2.1}.
Finalmente, se sabe que $P_0$ está en correspondencia biyectiva con el conjunto de tipos de extensiones compuestas de $\boldsymbol{k}(x)$ y $\boldsymbol{k}(y)$ (Bourbaki, \emph{Alg.}, chap.~VIII, \textsection8, prop.~1).
\end{proof}

\subsection{Sobreyectividad e inyectividad}
\label{subsection:I.3.5}

\begin{env}[3.5.1]
\label{I.3.5.1}
De manera general, consideremos una propiedad $\boldsymbol{P}$ de los morfismos de preesquemas, y las dos proposiciones siguientes:
\begin{enumerate}
  \item[\textup{(i)}] \itshape Si $f:X\to X'$ y $g:Y\to Y'$ son $S$-morfismos que poseen la propiedad $\boldsymbol{P}$, entonces $f\times_S g$ también posee la propiedad $\boldsymbol{P}$.
  \upshape\item[\textup{(ii)}] \itshape Si $f:X\to Y$ es un $S$-morfismo que posee la propiedad $\boldsymbol{P}$, entonces todo $S'$-morfismo $f_{(S')}:X_{(S')}\to Y_{(S')}$ obtenido de $f$ por extensión del preesquema de base también posee la propiedad $\boldsymbol{P}$.
\end{enumerate}

\oldpage[I]{115}Puesto que $f_{(S')}=f\times_S 1_{S'}$, se ve que, si para todo preesquema $X$ la \emph{identidad} $1_X$ posee la propiedad $\boldsymbol{P}$, entonces (i) implica (ii); por otra parte, puesto que $f\times_S g$ es el morfismo compuesto
\phantomsection\label{I.3.5.1.product-composition-fr}
\[
  X\times_S Y\xrightarrow{f\times 1_Y}X'\times_S Y
  \xrightarrow{1_{X'}\times g}X'\times_S Y'
\]
se ve que, si el \emph{compuesto} de dos morfismos que poseen la propiedad $\boldsymbol{P}$ también posee esta propiedad, entonces (ii) implica (i).
\end{env}

Una primera aplicación de esta observación es
\begin{proposition}[3.5.2]
\label{I.3.5.2}
\medskip\noindent
\begin{enumerate}
  \item[{\rm(i)}] Si $f:X\to X'$ y $g:Y\to Y'$ son $S$-morfismos sobreyectivos, entonces $f\times_S g$ es sobreyectivo.
  \item[{\rm(ii)}] Si $f:X\to Y$ es un $S$-morfismo sobreyectivo, entonces $f_{(S')}$ es sobreyectivo para toda extensión de base $S'$.
\end{enumerate}
\end{proposition}

\begin{proof}
Como el compuesto de dos morfismos sobreyectivos es sobreyectivo, basta demostrar (ii); pero esta proposición se sigue inmediatamente de \sref{I.3.4.8} aplicado a $M=X$.
\end{proof}

\begin{proposition}[3.5.3]
\label{I.3.5.3}
Para que un morfismo $f:X\to Y$ sea sobreyectivo, es necesario y suficiente que, para todo cuerpo $K$ y todo morfismo $\operatorname{Spec}(K)\to Y$, exista una extensión $K'$ de $K$ y un morfismo $\operatorname{Spec}(K')\to X$ que hagan conmutar el diagrama siguiente:
\phantomsection\label{I.3.5.3.surjectivity-diagram-fr}
\[
  \xymatrix{
    X\ar[d]_f & \operatorname{Spec}(K')\ar[l]\ar[d]\\
    Y & \operatorname{Spec}(K)\ar[l]
  }
\]
\end{proposition}

\begin{proof}
La condición es suficiente porque, para todo $y\in Y$, basta aplicarla a un morfismo $\operatorname{Spec}(K)\to Y$ correspondiente a un monomorfismo $\boldsymbol{k}(y)\to K$, siendo $K$ una extensión de $\boldsymbol{k}(y)$ \sref{I.2.4.6}.
Recíprocamente, supongamos que $f$ es sobreyectivo, y sea $y\in Y$ la imagen del único punto de $\operatorname{Spec}(K)$; existe algún $x\in X$ tal que $f(x)=y$; consideremos el monomorfismo correspondiente $\boldsymbol{k}(y)\to\boldsymbol{k}(x)$ \sref{I.2.2.1}; basta entonces tomar $K'$ una extensión de $\boldsymbol{k}(y)$ tal que existan monomorfismos sobre $\boldsymbol{k}(y)$ de $\boldsymbol{k}(x)$ y de $K$ en $K'$ (Bourbaki, \emph{Alg.}, cap.~V, §4, prop.~2); el morfismo $\operatorname{Spec}(K')\to X$ correspondiente a $\boldsymbol{k}(x)\to K'$ es exactamente el que se busca.
\end{proof}

Con el lenguaje introducido en \sref{I.3.4.5}, también se puede decir que \emph{todo punto geométrico de $Y$ con valores en $K$ proviene de un punto geométrico de $X$ con valores en una extensión de $K$}.

\begin{definition}[3.5.4]
\label{I.3.5.4}
Se dice que un morfismo $f:X\to Y$ de preesquemas es \emph{universalmente inyectivo}, o un \emph{morfismo radicial}, si, para todo cuerpo $K$, la aplicación correspondiente $X(K)\to Y(K)$ es inyectiva.
\end{definition}

Se sigue inmediatamente de las definiciones que todo \emph{monomorfismo de preesquemas} (T,~1.1) es radicial.

\begin{env}[3.5.5]
\label{I.3.5.5}
Para que un morfismo $f:X\to Y$ sea radicial, basta que la condición de la Definición~\sref{I.3.5.4} se cumpla para todo cuerpo \emph{algebraicamente cerrado}.
En efecto, si $K$ es un cuerpo arbitrario, y $K'$ una extensión algebraicamente cerrada de $K$, entonces el diagrama
\phantomsection\label{I.3.5.5.algebraic-closure-diagram-fr}
\[
  \xymatrix{
    X(K)\ar[r]^\alpha\ar[d]_\varphi & Y(K)\ar[d]^{\varphi'}\\
    X(K')\ar[r]_{\alpha'} & Y(K')
  }
\]
\oldpage[I]{116}conmuta, donde $\varphi$ y $\varphi'$ proceden del morfismo $\operatorname{Spec}(K')\to\operatorname{Spec}(K)$, y $\alpha$ y $\alpha'$ corresponden a $f$.
Sin embargo, $\varphi$ es inyectiva, y también $\alpha'$ lo es, por hipótesis; por lo tanto $\alpha$ es necesariamente inyectiva.
\end{env}

\par\medskip
\begin{proposition}[3.5.6]
\label{I.3.5.6}
Sean $f:X\to Y$ y $g:Y\to Z$ dos morfismos de preesquemas.
\begin{enumerate}
  \item[{\rm(i)}] Si $f$ y $g$ son radiciales, entonces $g\circ f$ también lo es.
  \item[{\rm(ii)}] Recíprocamente, si $g\circ f$ es radicial, entonces $f$ también lo es.
\end{enumerate}
\end{proposition}

\begin{proof}
Teniendo en cuenta la Definición~\sref{I.3.5.4}, la proposición se reduce a las afirmaciones correspondientes para las aplicaciones $X(K)\to Y(K)\to Z(K)$, y estas afirmaciones son evidentes.
\end{proof}

\begin{proposition}[3.5.7]
\label{I.3.5.7}
\medskip\noindent
\begin{enumerate}
  \item[{\rm(i)}] Si los $S$-morfismos $f:X\to X'$ y $g:Y\to Y'$ son radiciales, entonces $f\times_S g$ también lo es.
  \item[{\rm(ii)}] Si el $S$-morfismo $f:X\to Y$ es radicial, entonces también lo es $f_{(S')}:X_{(S')}\to Y_{(S')}$ para toda extensión $S'\to S$ del preesquema base.
\end{enumerate}
\end{proposition}

\begin{proof}
Dado \sref{I.3.5.1}, basta demostrar (i).
Se ha visto en \sref{I.3.4.2.1} que
\phantomsection\label{I.3.5.7.field-valued-product-fr}
\[
  (X\times_S Y)(K)=X(K)\times_{S(K)}Y(K),\qquad
  (X'\times_S Y')(K)=X'(K)\times_{S(K)}Y'(K)
\]
con la aplicación $(X\times_S Y)(K)\to(X'\times_S Y')(K)$ correspondiente a $f\times_S g$ identificándose así con $(u,v)\to(f\circ u,g\circ v)$, y la proposición se sigue inmediatamente.
\end{proof}

\begin{proposition}[3.5.8]
\label{I.3.5.8}
Para que un morfismo $f=(\psi,\theta):X\to Y$ sea radicial, es necesario y suficiente que $\psi$ sea inyectiva y que el monomorfismo $\theta^x:\boldsymbol{k}(\psi(x))\to\boldsymbol{k}(x)$ haga de $\boldsymbol{k}(x)$ una extensión radicial de $\boldsymbol{k}(\psi(x))$ para todo $x\in X$.
\end{proposition}

\begin{proof}
Supongamos que $f$ es radicial y mostremos primero que la ecuación $\psi(x_1)=\psi(x_2)=y$ implica necesariamente que $x_1=x_2$.
En efecto, existe un cuerpo $K$ que extiende $\boldsymbol{k}(y)$, junto con $\boldsymbol{k}(y)$-monomorfismos $\boldsymbol{k}(x_1)\to K$ y $\boldsymbol{k}(x_2)\to K$ (Bourbaki, \emph{Alg.}, cap.~V, §4, prop.~2); los morfismos correspondientes $u_1:\operatorname{Spec}(K)\to X$ y $u_2:\operatorname{Spec}(K)\to X$ son entonces tales que $f\circ u_1=f\circ u_2$, y por hipótesis $u_1=u_2$, lo cual implica, en particular, que $x_1=x_2$.
Ahora consideremos $\boldsymbol{k}(x)$ como una extensión de $\boldsymbol{k}(\psi(x))$ mediante $\theta^x$: si $\boldsymbol{k}(x)$ no es una extensión radicial de $\boldsymbol{k}(\psi(x))$, existen dos $\boldsymbol{k}(\psi(x))$-monomorfismos distintos de $\boldsymbol{k}(x)$ en una extensión algebraicamente cerrada $K$ de $\boldsymbol{k}(\psi(x))$, y los dos morfismos correspondientes $\operatorname{Spec}(K)\to X$ contradirían la hipótesis.
Recíprocamente, teniendo en cuenta \sref{I.2.4.6}, es inmediato que las condiciones enunciadas son suficientes para que $f$ sea radicial.
\end{proof}

\begin{corollary}[3.5.9]
\label{I.3.5.9}
Si $A$ es un anillo y $S$ es una parte multiplicativa de $A$, entonces el morfismo canónico $\operatorname{Spec}(S^{-1}A)\to\operatorname{Spec}(A)$ es radicial.
\end{corollary}

\begin{proof}
En efecto, este morfismo es un monomorfismo \sref{I.1.6.2}.
\end{proof}

\begin{corollary}[3.5.10]
\label{I.3.5.10}
Sea $f:X\to Y$ un morfismo radicial, $g:Y'\to Y$ un morfismo y $X'=X_{(Y')}=X\times_Y Y'$.
Entonces el morfismo radicial $f_{(Y')}$ \sref{I.3.5.7}[ii] es una biyección del espacio subyacente de $X'$ sobre $g^{-1}(f(X))$; además, para todo cuerpo $K$, el conjunto $X'(K)$ se identifica con el subconjunto de $Y'(K)$ imagen inversa por la aplicación $Y'(K)\to Y(K)$ (correspondiente a $g$) del subconjunto $X(K)$ de $Y(K)$.
\end{corollary}

\begin{proof}
La primera afirmación se sigue inmediatamente de \sref{I.3.5.8} y \sref{I.3.4.8}; la segunda, de la conmutatividad del diagrama siguiente:\oldpage[I]{117}
\phantomsection\label{I.3.5.10.base-change-points-diagram-fr}
\[
  \xymatrix{
    X'(K)\ar[r]\ar[d] & Y'(K)\ar[d]\\
    X(K)\ar[r] & Y(K)
  }
\]
\end{proof}

\begin{remark}[3.5.11]
\label{I.3.5.11}
Decimos que un morfismo $f=(\psi,\theta)$ de preesquemas es \emph{inyectivo} si la aplicación $\psi$ es inyectiva.
Para que un morfismo $f=(\psi,\theta):X\to Y$ sea radicial, es necesario y suficiente que, para todo morfismo $Y'\to Y$, el morfismo $f_{(Y')}:X_{(Y')}\to Y'$ sea inyectivo (lo que justifica la terminología de morfismo \emph{universalmente inyectivo}).
De hecho, la condición es necesaria en virtud de \sref{I.3.5.7}[ii] y \sref{I.3.5.8}.
Recíprocamente, la condición implica primero que $\psi$ es inyectiva; si, para algún $x\in X$, el monomorfismo $\theta^x:\boldsymbol{k}(\psi(x))\to\boldsymbol{k}(x)$ no fuera radicial, existiría una extensión $K$ de $\boldsymbol{k}(\psi(x))$ y dos morfismos distintos $\operatorname{Spec}(K)\to X$ correspondientes al mismo morfismo $\operatorname{Spec}(K)\to Y$ \sref{I.3.5.8}.
Pero entonces, tomando $Y'=\operatorname{Spec}(K)$, habría dos $Y'$-secciones distintas de $X_{(Y')}$ \sref{I.3.3.14}, lo que contradice la hipótesis de que $f_{(Y')}$ es inyectivo.
\end{remark}

\subsection{Fibras}
\label{subsection:I.3.6}

\begin{proposition}[3.6.1]
\label{I.3.6.1}
Sea $f:X\to Y$ un morfismo, $y$ un punto de $Y$, y $\mathfrak{a}_y$ un ideal de definición de $\mathcal{O}_y$ para la topología $\mathfrak{m}_y$-preádica.
Entonces la proyección $p:X\times_Y\operatorname{Spec}(\mathcal{O}_y/\mathfrak{a}_y)\to X$ es un homeomorfismo del espacio subyacente al preesquema $X\times_Y\operatorname{Spec}(\mathcal{O}_y/\mathfrak{a}_y)$ sobre la fibra $f^{-1}(y)$ provista de la topología inducida por la del espacio subyacente a $X$.
\end{proposition}

\begin{proof}
Como $\operatorname{Spec}(\mathcal{O}_y/\mathfrak{a}_y)\to Y$ es radicial (\sref{I.3.5.4} y \sref{I.2.4.7}), como $\operatorname{Spec}(\mathcal{O}_y/\mathfrak{a}_y)$ se reduce a un solo punto, y como el ideal $\mathfrak{m}_y/\mathfrak{a}_y$ es nilpotente por hipótesis \sref{I.1.1.12}, ya se sabe (\sref{I.3.5.10} y \sref{I.3.3.4}) que $p$ identifica, en cuanto \emph{conjunto}, el espacio subyacente a $X\times_Y\operatorname{Spec}(\mathcal{O}_y/\mathfrak{a}_y)$ con $f^{-1}(y)$; todo se reduce a demostrar que $p$ es un homeomorfismo.
En virtud de \sref{I.3.2.7}, la cuestión es local en $X$ e $Y$, y se puede por tanto suponer que $X=\operatorname{Spec}(B)$ e $Y=\operatorname{Spec}(A)$, siendo $B$ una $A$-álgebra.
El morfismo $p$ corresponde entonces al homomorfismo $1\otimes\varphi:B\to B\otimes_A A'$, donde $A'=A_y/\mathfrak{a}_y$ y $\varphi$ es la aplicación canónica de $A$ en $A'$.
Todo elemento de $B\otimes_A A'$ se escribe entonces en la forma
\phantomsection\label{I.3.6.1.localization-fraction-fr}
\[
  \sum_i b_i\otimes\varphi(a_i)/\varphi(s)
  =\left(\sum_i(a_i b_i\otimes 1)\right)(1\otimes\varphi(s))^{-1},
\]
donde $s\notin\mathfrak{j}_y$, y se aplica la Proposición~\sref{I.1.2.4}.
\end{proof}

\begin{env}[3.6.2]
\label{I.3.6.2}
En todo el resto de este Tratado, cuando consideremos una \emph{fibra} $f^{-1}(y)$ de un morfismo como dotada de estructura de $\boldsymbol{k}(y)$-preesquema, \emph{se tratará siempre del preesquema obtenido transportando la estructura de $X\times_Y\operatorname{Spec}(\boldsymbol{k}(y))$ mediante la proyección en $X$}.
Escribiremos también este último producto como $X\otimes_Y \boldsymbol{k}(y)$, o $X\otimes_{\mathcal{O}_Y}\boldsymbol{k}(y)$; más generalmente, si $B$ es una $\mathcal{O}_y$-álgebra, denotaremos por $X\otimes_Y B$ o $X\otimes_{\mathcal{O}_Y}B$ el producto $X\times_Y\operatorname{Spec}(B)$.
\end{env}

Con la convención precedente, se sigue de \sref{I.3.5.10} que los puntos de $X$ con valores en una extensión $K$ de $\boldsymbol{k}(y)$ se identifican con los \emph{puntos de $f^{-1}(y)$ con valores en $K$}.

\begin{env}[3.6.3]
\label{I.3.6.3}
Sean $f:X\to Y$ y $g:Y\to Z$ dos morfismos, y $h=g\circ f$ su compuesto; para todo $z\in Z$, la fibra $h^{-1}(z)$ es un preesquema isomorfo a
\phantomsection\label{I.3.6.3.fibre-composition-fr}
\[
  X\times_Z\operatorname{Spec}(\boldsymbol{k}(z))
  =(X\times_Y Y)\times_Z\operatorname{Spec}(\boldsymbol{k}(z))
  =X\times_Y g^{-1}(z)
\]
\oldpage[I]{118}En particular, si $U$ es una parte abierta de $X$, el preesquema inducido sobre $U\cap f^{-1}(y)$ por el preesquema $f^{-1}(y)$ es isomorfo a $f_U^{-1}(y)$ (siendo $f_U$ la restricción de $f$ a $U$).
\end{env}

\begin{proposition}[3.6.4) \mdseries\upshape («~transitividad de las fibras~»]
\label{I.3.6.4}
%
Sean $f:X\to Y$ y $g:Y'\to Y$ dos morfismos; pongamos $X'=X\times_Y Y'=X_{(Y')}$ y $f'=f_{(Y')}:X'\to Y'$.
Para todo $y'\in Y'$, si se pone $y=g(y')$, el preesquema $f'^{-1}(y')$ es isomorfo a $f^{-1}(y)\otimes_{\boldsymbol{k}(y)}\boldsymbol{k}(y')$.
\end{proposition}

\begin{proof}
En efecto, basta observar que los dos preesquemas $(X\otimes_Y\boldsymbol{k}(y))\otimes_{\boldsymbol{k}(y)}\boldsymbol{k}(y')$ y $(X\times_Y Y')\otimes_{Y'}\boldsymbol{k}(y')$ son ambos canónicamente isomorfos a $X\times_Y\operatorname{Spec}(\boldsymbol{k}(y'))$ por \sref{I.3.3.9.1}.
\end{proof}

En particular, si $V$ es una vecindad abierta de $y$ en $Y$, y se denota por $f_V$ la restricción de $f$ al preesquema inducido sobre $f^{-1}(V)$, los preesquemas $f^{-1}(y)$ y $f_V^{-1}(y)$ se identifican canónicamente.

\par\medskip
\begin{proposition}[3.6.5]
\label{I.3.6.5}
Sea $f:X\to Y$ un morfismo, $y$ un punto de $Y$, $Z$ el preesquema local $\operatorname{Spec}(\mathcal{O}_y)$, y $p=(\psi,\theta)$ la proyección $X\times_Y Z\to X$; entonces $p$ es un homeomorfismo del espacio subyacente a $X\times_Y Z$ sobre el subespacio $f^{-1}(Z)$ de $X$ \textup{(cuando el espacio subyacente a $Z$ se identifica con un subespacio de $Y$, cf.~\sref{I.2.4.2})}, y, para todo $t\in X\times_Y Z$, poniendo $z=\psi(t)$, $\theta_t^\sharp$ es un isomorfismo de $\mathcal{O}_z$ sobre $\mathcal{O}_t$.
\end{proposition}

\begin{proof}
Como $Z$ (identificado con un subespacio de $Y$) está contenido en todo abierto afín que contiene a $y$ \sref{I.2.4.2}, se puede, como en \sref{I.3.6.1}, reducir al caso en que $X=\operatorname{Spec}(A)$ e $Y=\operatorname{Spec}(B)$ son esquemas afines, siendo $A$ una $B$-álgebra.
Entonces $X\times_Y Z$ es el espectro primo de $A\otimes_B B_y$, y este anillo se identifica canónicamente con $S^{-1}A$, donde $S$ es la imagen de $B\setmin\mathfrak{j}_y$ en $A$ \sref[0]{0.1.5.2}; como $p$ corresponde entonces al homomorfismo canónico $A\to S^{-1}A$, la proposición se sigue de \sref{I.1.6.2}.
\end{proof}

\subsection[Aplicación: reducción de un preesquema mod.~\texorpdfstring{$\mathfrak{J}$}{J}.]{Aplicación: reducción de un preesquema mod.~\texorpdfstring{$\mathfrak{J}$\protect\footnotemark}{J}.}
\label{subsection:I.3.7}
\label{I.3.7}

\footnotetext{Esta sección, que hace uso de nociones y resultados posteriores del Capítulo~I y del Capítulo~II, no se utilizará en lo que sigue de este Tratado, y está destinada únicamente a los lectores familiarizados con la geometría algebraica clásica.}

\begin{env}[3.7.1]
\label{I.3.7.1}
Sean $A$ un anillo, $X$ un $A$-preesquema, y $\mathfrak{J}$ un ideal de $A$; entonces $X_0=X\otimes_A(A/\mathfrak{J})$ es un $(A/\mathfrak{J})$-preesquema, del que se dice a veces que se deduce de $X$ por \emph{reducción} mod.~$\mathfrak{J}$.
\end{env}

\begin{env}[3.7.2]
\label{I.3.7.2}
Esta terminología se utiliza sobre todo cuando $A$ es un \emph{anillo local} y $\mathfrak{J}$ su ideal maximal, de modo que $X_0$ es un preesquema sobre el cuerpo residual $k=A/\mathfrak{J}$ de $A$.

Cuando además $A$ es íntegro, con cuerpo de fracciones $K$, se puede considerar el $K$-preesquema $X'=X\otimes_A K$.
Por un abuso de lenguaje del que no haremos uso, se acostumbraba hasta ahora decir que $X_0$ se \emph{deduce de $X'$} por reducción mod.~$\mathfrak{J}$.
En los casos en que se utilizaba este lenguaje, $A$ era un anillo local de dimensión 1 (la mayoría de las veces un anillo de valuación discreta), y se sobreentendía (de manera más o menos explícita) que el $K$-preesquema dado $X'$ era un subpreesquema cerrado de un $K$-preesquema $P'$ (de hecho, un espacio proyectivo del tipo $\mathbf{P}_K^r$, cf. \sref[II]{II.4.1.1}), él mismo de la forma $P'=P\otimes_A K$, donde $P$ es un $A$-preesquema dado (de hecho, el $A$-esquema $\mathbf{P}_A^r$, con la notación de \sref[II]{II.4.1.1}).
En nuestro lenguaje, la definición de $X_0$ a partir de $X'$ se formula como sigue:

Se considera el esquema afín $Y=\operatorname{Spec}(A)$, formado por dos puntos, el único punto cerrado
\oldpage[I]{119}
$y=\mathfrak{J}$ y el punto genérico $(0)$, siendo así el conjunto unitario $U$ del punto genérico un abierto $U=\operatorname{Spec}(K)$ de $Y$.
Si $X$ es un $A$-preesquema (o, dicho de otro modo, un $Y$-preesquema), entonces $X\otimes_A K=X'$ no es sino el preesquema inducido por $X$ sobre $\psi^{-1}(U)$, denotando por $\psi$ el morfismo estructural $X\to Y$.
En particular, si $\varphi$ es el morfismo estructural $P\to Y$, un subpreesquema cerrado $X'$ de $P'=\varphi^{-1}(U)$ es un subpreesquema (localmente cerrado) de $P$.
Si $P$ es noetheriano (por ejemplo, si $A$ es noetheriano y $P$ es de tipo finito sobre $A$), existe un menor subpreesquema cerrado $X=\overline{X}'$ de $P$ que contiene a $X'$ \sref{I.9.5.10}, y $X'$ es el preesquema inducido por $X$ sobre el abierto $\varphi^{-1}(U)\cap X$, y es por tanto isomorfo a $X\otimes_A K$ \sref{I.9.5.10}.
\emph{La inmersión de $X'$ en $P'=P\otimes_A K$ permite así considerar canónicamente $X'$ como de la forma $X'=X\otimes_A K$, donde $X$ es un $A$-preesquema.}
Se puede entonces considerar el preesquema reducido mod.~$\mathfrak{J}$, $X_0=X\otimes_A k$, que no es sino la fibra $\psi^{-1}(y)$ del punto cerrado $y$.
Hasta ahora, a falta de una terminología adecuada, se había evitado introducir explícitamente el $A$-preesquema $X$.
Conviene señalar, no obstante, que todos los enunciados que se hacen habitualmente sobre el preesquema «reducido mod.~$\mathfrak{J}$» $X_0$ deben considerarse como consecuencias de enunciados más completos relativos al propio $X$, y no pueden formularse ni comprenderse satisfactoriamente sino interpretándolos de este modo.
Parece además que las hipótesis hechas se reducen siempre a hipótesis sobre el propio $X$ (independientemente del dato previo de una inmersión de $X'$ en $\mathbf{P}_K^r$), lo que permite dar enunciados más intrínsecos.
\end{env}

\begin{env}[3.7.3]
\label{I.3.7.3}
Señalemos por último un hecho muy particular, que sin duda ha contribuido a retrasar la clarificación conceptual de la situación considerada aquí: si $A$ es un anillo de valuación discreta, y si $X$ es \emph{propio} sobre $A$ (lo que ocurre en efecto si $X$ es un subpreesquema cerrado de algún $\mathbf{P}_A^r$, cf. \sref[II]{II.5.5.4}), entonces los puntos de $X$ con valores en $A$ y los puntos de $X'$ con valores en $K$ están en correspondencia biyectiva \sref[II]{II.7.3.8}.
Por esta razón se ha creído a menudo demostrar resultados relativos a $X'$, cuando en realidad se demostraban enunciados relativos a $X$; estos permanecen válidos (en esta forma) cuando ya no se supone que el anillo local de base tenga dimensión 1.
\end{env}
% END source_es/ega1/ega1-3.tex

% BEGIN source_es/ega1/ega1-4.tex
\section{Subpreesquemas y morfismos de inmersión}
\label{section:I.4}

\subsection{Subpreesquemas}
\label{subsection:I.4.1}

\begin{env}[4.1.1]
\label{I.4.1.1}
Como la noción de haz cuasicoherente \sref[0]{0.5.1.3} es local, un $\mathscr{O}_X$-módulo cuasicoherente $\mathscr{F}$ sobre un preesquema $X$ puede definirse mediante la condición siguiente: para todo abierto afín $V$ de $X$, $\mathscr{F}|V$ es isomorfo al haz asociado a un $\Gamma(V,\mathscr{O}_X)$-módulo \sref{I.1.4.1}.
Es claro que, sobre un preesquema $X$, el haz estructural $\mathscr{O}_X$ es cuasicoherente, y que los núcleos, conúcleos e imágenes de homomorfismos de $\mathscr{O}_X$-módulos cuasicoherentes, así como los límites inductivos y las sumas directas de $\mathscr{O}_X$-módulos cuasicoherentes, son también cuasicoherentes (Teorema~\sref{I.1.3.7} y Corolario~\sref{I.1.3.9}).
\end{env}

\begin{proposition}[4.1.2]
\label{I.4.1.2}
Sean $X$ un preesquema e $\mathscr{I}$ un haz cuasicoherente de ideales de $\mathscr{O}_X$.
Entonces el soporte $Y$ del haz $\mathscr{O}_X/\mathscr{I}$ es cerrado, y, si se denota por $\mathscr{O}_Y$ la restricción de $\mathscr{O}_X/\mathscr{I}$ a $Y$, entonces $(Y,\mathscr{O}_Y)$ es un preesquema.
\end{proposition}

\begin{proof}
\oldpage[I]{120}
Evidentemente basta \sref{I.2.1.3} considerar el caso en que $X$ es un esquema afín y demostrar que, en este caso, $Y$ es cerrado en $X$ y es un \emph{esquema afín}.
En efecto, si $X=\operatorname{Spec}(A)$, entonces $\mathscr{O}_X=\widetilde{A}$ y $\mathscr{I}=\widetilde{\mathfrak{J}}$, donde $\mathfrak{J}$ es un ideal de $A$ \sref{I.1.4.1};
$Y$ es entonces igual a la parte cerrada $V(\mathfrak{J})$ de $X$ y se identifica con el espectro primo del anillo $B=A/\mathfrak{J}$ \sref{I.1.1.11};
además, si $\varphi$ es el homomorfismo canónico $A\to B=A/\mathfrak{J}$, entonces la imagen directa ${}^a\varphi_*(\widetilde{B})$ se identifica canónicamente con el haz $\widetilde{A}/\widetilde{\mathfrak{J}}=\mathscr{O}_X/\mathscr{I}$ (Proposición~\sref{I.1.6.3} y Corolario~\sref{I.1.3.9}), lo que termina la demostración.
\end{proof}

Decimos que $(Y,\mathscr{O}_Y)$ es el \emph{subpreesquema} de $(X,\mathscr{O}_X)$ \emph{definido por el haz de ideales $\mathscr{I}$};
este es un caso particular de la noción más general de \emph{subpreesquema}:

\begin{definition}[4.1.3]
\label{I.4.1.3}
{\itshape Se dice que un espacio anillado $(Y,\mathscr{O}_Y)$ es un subpreesquema de un preesquema $(X,\mathscr{O}_X)$ si:
\begin{enumerate}
  \item[$1^\circ$] $Y$ es un subespacio localmente cerrado de $X$;
  \item[$2^\circ$] si $U$ denota el mayor abierto de $X$ que contiene a $Y$ y tal que $Y$ sea cerrado en $U$ (\textup{equivalentemente,} el complemento en $X$ de la frontera de $Y$ respecto de $\overline{Y}$), entonces $(Y,\mathscr{O}_Y)$ es un subpreesquema de $(U,\mathscr{O}_X|U)$ definido por un haz cuasicoherente de ideales de $\mathscr{O}_X|U$.
\end{enumerate}
Se dice que el subpreesquema $(Y,\mathscr{O}_Y)$ de $(X,\mathscr{O}_X)$ es cerrado si $Y$ es cerrado en $X$ (en cuyo caso $U=X$).}
\end{definition}

De esta definición y de la Proposición~\sref{I.4.1.2} se sigue inmediatamente que los subpreesquemas cerrados de $X$ están en \emph{correspondencia biyectiva} canónica con los haces cuasicoherentes de ideales $\mathscr{J}$ de $\mathscr{O}_X$, pues si dos de estos haces $\mathscr{J}$ y $\mathscr{J}'$ tienen el mismo soporte (cerrado) $Y$, y si las restricciones de $\mathscr{O}_X/\mathscr{J}$ y $\mathscr{O}_X/\mathscr{J}'$ a $Y$ son idénticas, entonces $\mathscr{J}'=\mathscr{J}$.

\begin{env}[4.1.4]
\label{I.4.1.4}
Sean $(Y,\mathscr{O}_Y)$ un subpreesquema de $X$, $U$ el mayor abierto de $X$ tal que $Y$ sea cerrado (y, por tanto, esté contenido) en $Y$, y $V$ un abierto de $X$ contenido en $U$;
entonces $V\cap Y$ es cerrado en $V$.
Además, si $Y$ está definido por el haz cuasicoherente de ideales $\mathscr{J}$ de $\mathscr{O}_X|U$, entonces $\mathscr{J}|V$ es un haz cuasicoherente de ideales de $\mathscr{O}_X|V$, y es inmediato que el preesquema inducido por $Y$ sobre $Y\cap V$ es el subpreesquema cerrado de $V$ definido por el haz de ideales $\mathscr{J}|V$.
Recíprocamente:
\end{env}

\begin{proposition}[4.1.5]
\label{I.4.1.5}
Sea $(Y,\mathscr{O}_Y)$ un espacio anillado tal que $Y$ sea un subespacio de $X$ y que exista un recubrimiento $(V_\alpha)$ de $Y$ por abiertos de $X$ tales que, para todo $\alpha$, $Y\cap V_\alpha$ sea cerrado en $V_\alpha$, y que el espacio anillado $(Y\cap V_\alpha,\mathscr{O}_Y|(Y\cap V_\alpha))$ sea un subpreesquema cerrado del preesquema inducido sobre $V_\alpha$ por $X$.
Entonces $(Y,\mathscr{O}_Y)$ es un subpreesquema de $X$.
\end{proposition}

\begin{proof}
Las hipótesis implican que $Y$ es localmente cerrado en $X$ y que el mayor abierto $U$ que contiene a $Y$ y en el cual $Y$ es cerrado contiene todos los $V_\alpha$;
se puede, pues, reducir al caso en que $U=X$ e $Y$ es cerrado en $X$.
Se define entonces un haz cuasicoherente de ideales $\mathscr{J}$ de $\mathscr{O}_X$ tomando como $\mathscr{J}|V_\alpha$ el haz de ideales de $\mathscr{O}_X|V_\alpha$ que define el subpreesquema cerrado $(Y\cap V_\alpha,\mathscr{O}_Y|(Y\cap V_\alpha))$, y, para todo abierto $W$ de $X$ disjunto de $Y$, $\mathscr{J}|W=\mathscr{O}_X|W$.
Se verifica inmediatamente, por la Definición~\sref{I.4.1.3} y \sref{I.4.1.4}, que existe un único haz de ideales $\mathscr{J}$ que satisface estas condiciones y que este define el subpreesquema cerrado $(Y,\mathscr{O}_Y)$.
\end{proof}

En particular, el preesquema \emph{inducido} (por $X$) sobre una \emph{parte abierta} de $X$ es un \emph{subpreesquema} de $X$.

\begin{proposition}[4.1.6]
\label{I.4.1.6}
Un subpreesquema (\textup{resp.} un subpreesquema cerrado) de un subpreesquema
\oldpage[I]{121}
(\textup{resp.} un subpreesquema cerrado) de $X$ se identifica canónicamente con un subpreesquema (\textup{resp.} un subpreesquema cerrado) de $X$.
\end{proposition}

\begin{proof}
Como una parte localmente cerrada de un subespacio localmente cerrado de $X$ es un subespacio localmente cerrado de $X$, está claro por \sref{I.4.1.5} que la cuestión es local y que se puede suponer, por tanto, que $X$ es afín;
la proposición resulta entonces de la identificación canónica de $A/\mathfrak{J}'$ con $(A/\mathfrak{J})/(\mathfrak{J}'/\mathfrak{J})$, cuando $\mathfrak{J}$, $\mathfrak{J}'$ son dos ideales de un anillo $A$ tales que $\mathfrak{J}\subset\mathfrak{J}'$.
\end{proof}

En lo sucesivo haremos siempre la identificación precedente.

\begin{env}[4.1.7]
\label{I.4.1.7}
Sea $Y$ un subpreesquema de un preesquema $X$, y designemos por $\psi$ la inyección canónica $Y\to X$ de los \emph{espacios subyacentes};
se sabe que la imagen inversa $\psi^*(\mathscr{O}_X)$ es la restricción $\mathscr{O}_X|Y$ \sref[0]{0.3.7.1}.
Si, para todo $y\in Y$, se designa por $\omega_y$ el homomorfismo canónico $(\mathscr{O}_X)_y\to(\mathscr{O}_Y)_y$, estos homomorfismos son los inducidos en los tallos por un homomorfismo \emph{sobreyectivo} $\omega$ de haces de anillos $\mathscr{O}_X|Y\to\mathscr{O}_Y$: en efecto, basta comprobarlo localmente sobre $Y$, es decir, se puede suponer que $X$ es afín y que el subpreesquema $Y$ es cerrado; si, en este caso, $\mathscr{I}$ es el haz de ideales de $\mathscr{O}_X$ que define $Y$, los $\omega_y$ no son otra cosa que los homomorfismos inducidos en los tallos por $\mathscr{O}_X|Y\to(\mathscr{O}_X/\mathscr{I})|Y$.
Hemos definido así un \emph{monomorfismo de espacios anillados} \sref[0]{0.4.1.1} $j=(\psi,\omega^\flat)$, que es evidentemente un morfismo $Y\to X$ de preesquemas \sref{I.2.2.1}, y que llamaremos \emph{morfismo de inyección canónico}.

Si $f:X\to Z$ es un morfismo, diremos asimismo que el morfismo compuesto $Y\xrightarrow{j}X\xrightarrow{f}Z$ es la \emph{restricción} de $f$ al subpreesquema $Y$.
\end{env}

\begin{env}[4.1.8]
\label{I.4.1.8}
Conforme a las definiciones generales (T, I, 1.1), diremos que un morfismo de preesquemas $f:Z\to X$ está \emph{dominado}\footnote{\emph{[Trad.] Traducimos aquí el francés «majoré» por «dominado». Se trata de la relación de preorden definida por la condición de que $f$ se factorice por $j$, no de una cota numérica.}} por el morfismo de inyección $j:Y\to X$ de un subpreesquema $Y$ de $X$ si $f$ se factoriza en $Z\xrightarrow{g}Y\xrightarrow{j}X$, donde $g$ es un morfismo de preesquemas;
además, $g$ es necesariamente \emph{único}, puesto que $j$ es un monomorfismo.
\end{env}

\begin{proposition}[4.1.9]
\label{I.4.1.9}
Para que un morfismo $f:Z\to X$ esté dominado por un morfismo de inyección $j:Y\to X$, es necesario y suficiente que $f(Z)\subset Y$ y que, para todo $z\in Z$, poniendo $y=f(z)$, el homomorfismo $(\mathscr{O}_X)_y\to\mathscr{O}_z$ correspondiente a $f$ se factorice en $(\mathscr{O}_X)_y\to(\mathscr{O}_Y)_y\to\mathscr{O}_z$ (\textup{o, equivalentemente, que} el núcleo de $(\mathscr{O}_X)_y\to\mathscr{O}_z$ contenga al de $(\mathscr{O}_X)_y\to(\mathscr{O}_Y)_y$).
\end{proposition}

\begin{proof}
Las condiciones son evidentemente necesarias.
Para ver que son suficientes, se puede reducir al caso en que $Y$ es un subpreesquema \emph{cerrado} de $X$, reemplazando, si es preciso, $X$ por un abierto $U$ tal que $Y$ sea cerrado en $U$ \sref{I.4.1.3};
$Y$ queda entonces definido por un haz cuasicoherente de ideales $\mathscr{I}$ de $\mathscr{O}_X$.
Pongamos $f=(\psi,\theta)$, y sea $\mathscr{J}$ el haz de ideales de $\psi^*(\mathscr{O}_X)$ que es el núcleo de $\theta^\sharp:\psi^*(\mathscr{O}_X)\to\mathscr{O}_Z$;
teniendo en cuenta las propiedades del funtor $\psi^*$ \sref[0]{0.3.7.2}, las hipótesis implican que, para todo $z\in Z$, se tiene $(\psi^*(\mathscr{I}))_z\subset\mathscr{J}_z$ y, por consiguiente, que $\psi^*(\mathscr{I})\subset\mathscr{J}$.
Por tanto, $\theta^\sharp$ se factoriza en
\[
  \psi^*(\mathscr{O}_X)\longrightarrow
  \psi^*(\mathscr{O}_X)/\psi^*(\mathscr{I})
  =\psi^*(\mathscr{O}_X/\mathscr{I})
  \xrightarrow{\omega}\mathscr{O}_Z
\]
siendo la primera flecha el homomorfismo canónico.
Sea $\psi'$ la aplicación continua $Z\to Y$ que coincide con $\psi$;
está claro que se tiene ${\psi'}^*(\mathscr{O}_Y)=\psi^*(\mathscr{O}_X/\mathscr{I})$;
por otra parte, $\omega$ es evidentemente un homomorfismo local, por lo que $g=(\psi',\omega^\flat)$ es un morfismo
\oldpage[I]{122}
de preesquemas $Z\to Y$ \sref{I.2.2.1} que, en virtud de lo anterior, satisface $f=j\circ g$, lo que demuestra la proposición.
\end{proof}

\begin{corollary}[4.1.10]
\label{I.4.1.10}
Para que un morfismo de inyección $Z\to X$ esté dominado por el morfismo de inyección $Y\to X$, es necesario y suficiente que $Z$ sea un subpreesquema de $Y$.
\end{corollary}

Se escribe entonces $Z\leq Y$, y esta relación es evidentemente una \emph{relación de orden} en el conjunto de subpreesquemas de $X$.

\subsection{Morfismos de inmersión.}
\label{subsection:I.4.2}

\begin{definition}[4.2.1]
\label{I.4.2.1}
{\itshape Se dice que un morfismo $f:Y\to X$ es una inmersión (\textup{resp.} una inmersión cerrada, una inmersión abierta) si se factoriza en $Y\xrightarrow{g}Z\xrightarrow{j}X$, donde $g$ es un isomorfismo, $Z$ es un subpreesquema de $X$ (\textup{resp.} un subpreesquema cerrado, un subpreesquema inducido sobre un abierto), y $j$ es el morfismo de inyección.}
\end{definition}

El subpreesquema $Z$ y el isomorfismo $g$ quedan entonces determinados de manera \emph{única}, pues si $Z'$ es un segundo subpreesquema de $X$, $j'$ la inyección $Z'\to X$, y $g'$ un isomorfismo $Y\to Z'$ tal que $j\circ g=j'\circ g'$, entonces se tiene $j'=j\circ g\circ {g'}^{-1}$, de donde $Z'\leq Z$ \sref{I.4.1.10}, y se demuestra del mismo modo que $Z\leq Z'$, por tanto $Z'=Z$, y, puesto que $j$ es un monomorfismo de preesquemas, $g'=g$.

Se dice que $f=j\circ g$ es la \emph{factorización canónica} de la inmersión $f$, y que el subpreesquema $Z$ y el isomorfismo $g$ están \emph{asociados a $f$}.

Es claro que una inmersión es un \emph{monomorfismo} de preesquemas \sref{I.4.1.7} y, \emph{a fortiori}, un \emph{morfismo radicial} \sref{I.3.5.4}.

\begin{proposition}[4.2.2]
\label{I.4.2.2}
\medskip\noindent
\begin{enumerate}
  \item[\textup{a)}] Para que un morfismo $f=(\psi,\theta):Y\to X$ sea una inmersión abierta, es necesario y suficiente que $\psi$ sea un homeomorfismo de $Y$ sobre una parte abierta de $X$, y que, para todo $y\in Y$, el homomorfismo $\theta_y^\sharp:\mathscr{O}_{\psi(y)}\to\mathscr{O}_y$ sea biyectivo.
  \item[\textup{b)}] Para que un morfismo $f=(\psi,\theta):Y\to X$ sea una inmersión (\textup{resp.} una inmersión cerrada), es necesario y suficiente que $\psi$ sea un homeomorfismo de $Y$ sobre una parte localmente cerrada (\textup{resp.} cerrada) de $X$, y que, para todo $y\in Y$, el homomorfismo $\theta_y^\sharp:\mathscr{O}_{\psi(y)}\to\mathscr{O}_y$ sea sobreyectivo.
\end{enumerate}
\end{proposition}

\begin{proof}
\medskip\noindent
\begin{enumerate}
  \item[\emph{a)}] Las condiciones son evidentemente necesarias.
    Recíprocamente, si se satisfacen, está claro que $\theta^\sharp$ es un isomorfismo de $\psi^*(\mathscr{O}_X)$ sobre $\mathscr{O}_Y$\footnote{\emph{[Trad.] La fuente francesa invierte aquí el origen y el destino. Por la definición de un morfismo de preesquemas y por los homomorfismos de tallos mostrados en la Proposición~\hyperref[I.4.2.2]{4.2.2}(a), $\theta^\sharp$ tiene por origen $\psi^*(\sh{O}_X)$ y por destino $\sh{O}_Y$.}}, y $\psi^*(\mathscr{O}_X)$ es el haz obtenido por ``transporte de estructura'' mediante $\psi^{-1}$ a partir de $\mathscr{O}_X|\psi(Y)$;
    de donde la conclusión.
  \item[\emph{b)}] Las condiciones son evidentemente necesarias; demostremos que son suficientes.
    Consideremos primero el caso particular en que se supone que $X$ es un esquema afín y que $Z=\psi(Y)$ está \emph{cerrado} en $X$.
    Se sabe entonces \sref[0]{0.3.4.6} que $\psi_*(\mathscr{O}_Y)$ tiene soporte $Z$, y que su restricción a $Z$, denotada por $\mathscr{O}'_Z$, hace que el espacio anillado $(Z,\mathscr{O}'_Z)$ se obtenga de $(Y,\mathscr{O}_Y)$ por transporte de estructura mediante el homeomorfismo $\psi$, considerado como aplicación de $Y$ sobre $Z$.
    Demostremos ahora que $f_*(\mathscr{O}_Y)=\psi_*(\mathscr{O}_Y)$ es un módulo \emph{cuasicoherente} sobre $\mathscr{O}_X$.
    En efecto, para todo $x\notin Z$, $\psi_*(\mathscr{O}_Y)$ restringido a una vecindad adecuada de $x$ es nulo.
    Si, por el contrario, $x\in Z$, entonces $x=\psi(y)$ para un único $y\in Y$;
    sea $V$ una vecindad abierta afín de $y$ en $Y$;
    $\psi(V)$ es entonces abierto en $Z$ y, por tanto, igual a la intersección de $Z$ con un abierto $U$ de $X$, y la restricción a $U$ de $\psi_*(\mathscr{O}_Y)$ se identifica con la restricción a $U$ de la imagen directa
\oldpage[I]{123}
    $(\psi_V)_*(\mathscr{O}_Y|V)$, donde $\psi_V$ es la restricción de $\psi$ a $V$.
    La restricción del morfismo $(\psi,\theta)$ a $(V,\mathscr{O}_Y|V)$ es un morfismo de este preesquema a $(X,\mathscr{O}_X)$ y, por consiguiente, es de la forma $({}^a\varphi,\widetilde{\varphi})$, donde $\varphi$ es el homomorfismo del anillo $A=\Gamma(X,\mathscr{O}_X)$ al anillo $\Gamma(V,\mathscr{O}_Y)$ \sref{I.1.7.3};
    se deduce que $(\psi_V)_*(\mathscr{O}_Y|V)$ es un $\mathscr{O}_X$-módulo cuasicoherente \sref{I.1.6.3}, lo que demuestra nuestra afirmación, por el carácter local de los haces cuasicoherentes.
    Además, la hipótesis de que $\psi$ es un homeomorfismo implica \sref[0]{0.3.4.5} que, para todo $y\in Y$, $\psi_y$ es un isomorfismo $(\psi_*(\mathscr{O}_Y))_{\psi(y)}\to\mathscr{O}_y$;
    como el diagrama
    \phantomsection\label{I.4.2.2.theta-diagram-fr}
\[
    \xymatrix{
      \mathscr{O}_{\psi(y)}\ar[r]^{\theta_{\psi(y)}}
        \ar[d]_{\psi_y\circ\rho_{\psi(y)}} &
      (\psi_*(\mathscr{O}_Y))_{\psi(y)}\ar[d]^{\psi_y}\\
      (\psi^*(\mathscr{O}_X))_y\ar[r]^{\theta_y^\sharp} &
      \mathscr{O}_y
    }
  \]
    es conmutativo y las flechas verticales son los isomorfismos \sref[0]{0.3.7.2}, la hipótesis de que $\theta_y^\sharp$ es sobreyectivo implica que $\theta_{\psi(y)}$ también es sobreyectivo.
    Como el soporte de $\psi_*(\mathscr{O}_Y)$ es $Z=\psi(Y)$, $\theta$ es un homomorfismo \emph{sobreyectivo} de $\mathscr{O}_X=\widetilde{A}$ sobre el $\mathscr{O}_X$-módulo cuasicoherente $f_*(\mathscr{O}_Y)$.
    Por consiguiente, existe un único isomorfismo $\omega$ de un haz cociente $\widetilde{A}/\widetilde{\mathfrak{J}}$ (donde $\mathfrak{J}$ es un ideal de $A$) sobre $f_*(\mathscr{O}_Y)$ que, compuesto con el homomorfismo canónico $\widetilde{A}\to\widetilde{A}/\widetilde{\mathfrak{J}}$, da $\theta$ \sref{I.1.3.8};
    si $\mathscr{O}_Z$ denota la restricción de $\widetilde{A}/\widetilde{\mathfrak{J}}$ a $Z$, entonces $(Z,\mathscr{O}_Z)$ es un subpreesquema de $(X,\mathscr{O}_X)$, y $f$ se factoriza a través de la inyección canónica de este subpreesquema en $X$ y del isomorfismo $(\psi_0,\omega_0)$, donde $\psi_0$ es $\psi$ considerado como aplicación de $Y$ sobre $Z$, y $\omega_0$ la restricción de $\omega$ a $\mathscr{O}_Z$.

    Pasemos ahora al caso general.
    Sea $U$ un abierto afín de $X$ tal que $U\cap\psi(Y)$ sea cerrado en $U$ y no vacío.
    Restringiendo $f$ al preesquema inducido por $Y$ sobre el abierto $\psi^{-1}(U)$ y considerándolo como un morfismo de este preesquema al preesquema inducido por $X$ sobre $U$, se reduce al primer caso;
    la restricción de $f$ a $\psi^{-1}(U)$ es, por tanto, una inmersión cerrada $\psi^{-1}(U)\to U$, que se factoriza canónicamente como $j_U\circ g_U$, donde $g_U$ es un isomorfismo del preesquema $\psi^{-1}(U)$ sobre un subpreesquema $Z_U$ de $U$, y $j_U$ es la inyección canónica $Z_U\to U$.
    Sea $V$ otro abierto afín de $X$ tal que $V\subset U$;
    puesto que la restricción $Z'_V$ de $Z_U$ a $V$ es un subpreesquema del preesquema $V$, la restricción de $f$ a $\psi^{-1}(V)$ se factoriza como $j'_V\circ g'_V$, donde $j'_V$ es la inyección canónica $Z'_V\to V$ y $g'_V$ es un isomorfismo de $\psi^{-1}(V)$ sobre $Z'_V$.
    Por la unicidad de la factorización canónica de una inmersión \sref{I.4.2.1}, se tiene necesariamente $Z'_V=Z_V$ y $g'_V=g_V$.
    Se deduce \sref{I.4.1.5} que existe un subpreesquema $Z$ de $X$ cuyo espacio subyacente es $\psi(Y)$ y cuya restricción a cada $U\cap\psi(Y)$ es $Z_U$;
    los $g_U$ son entonces las restricciones a $\psi^{-1}(U)$ de un isomorfismo $g:Y\to Z$ tal que $f=j\circ g$, donde $j$ es la inyección canónica $Z\to X$.
\end{enumerate}
\end{proof}

\begin{corollary}[4.2.3]
\label{I.4.2.3}
Sea $X$ un esquema afín.
Para que un morfismo $f=(\psi,\theta):Y\to X$ sea una inmersión cerrada, es necesario y suficiente que $Y$ sea un esquema afín y que el homomorfismo $\Gamma(\theta):\Gamma(\mathscr{O}_X)\to\Gamma(\mathscr{O}_Y)$\footnote{\emph{[Trad.] La fuente francesa imprime $\Gamma(\psi)$. Como $\psi$ es la aplicación continua y $\theta$ el homomorfismo de haces, el homomorfismo inducido en las secciones globales es $\Gamma(\theta)$.}} sea sobreyectivo.
\end{corollary}

\begin{corollary}[4.2.4]
\label{I.4.2.4}
\medskip\noindent
\begin{enumerate}
  \item[\textup{a)}] Sean $f$ un morfismo $Y\to X$ y $(V_\lambda)$ un recubrimiento de $f(Y)$ por abiertos de $X$.
    Para que $f$ sea una inmersión (\textup{resp.} una inmersión abierta), es necesario y suficiente
\oldpage[I]{124}
    que su restricción a cada uno de los preesquemas inducidos $f^{-1}(V_\lambda)$ sea una inmersión (\textup{resp.} una inmersión abierta) en $V_\lambda$.
  \item[\textup{b)}] Sean $f$ un morfismo $Y\to X$ y $(V_\lambda)$ un recubrimiento abierto de $X$.
    Para que $f$ sea una inmersión cerrada, es necesario y suficiente que su restricción a cada uno de los preesquemas inducidos $f^{-1}(V_\lambda)$ sea una inmersión cerrada en $V_\lambda$.
\end{enumerate}
\end{corollary}

\begin{proof}
Sea $f=(\psi,\theta)$;
en el caso \emph{a)}, $\theta_y^\sharp$ es sobreyectivo (resp. biyectivo) para todo $y\in Y$, y en el caso \emph{b)} $\theta_y^\sharp$ es sobreyectivo para todo $y\in Y$;
basta, pues, comprobar, en el caso \emph{a)}, que $\psi$ es un homeomorfismo de $Y$ sobre una parte localmente cerrada (resp. abierta) de $X$, y, en el caso \emph{b)}, un homeomorfismo de $Y$ sobre una parte cerrada de $X$.
Ahora bien, $\psi$ es evidentemente inyectiva y transforma toda vecindad de $y$ en $Y$ en una vecindad de $\psi(y)$ en $\psi(Y)$ para todo $y\in Y$, en virtud de la hipótesis;
en el caso \emph{a)}, $\psi(Y)\cap V_\lambda$ es localmente cerrado (resp. abierto) en $V_\lambda$, por lo que $\psi(Y)$ es localmente cerrado (resp. abierto) en la unión de los $V_\lambda$, y \emph{a fortiori} en $X$;
en el caso \emph{b)}, $\psi(Y)\cap V_\lambda$ es cerrado en $V_\lambda$, por lo que $\psi(Y)$ es cerrado en $X$, ya que $X=\bigcup_\lambda V_\lambda$.
\end{proof}

\begin{proposition}[4.2.5]
\label{I.4.2.5}
El compuesto de dos inmersiones cualesquiera (\textup{resp.} de dos inmersiones abiertas, de dos inmersiones cerradas) es una inmersión (\textup{resp.} una inmersión abierta, una inmersión cerrada).
\end{proposition}

\begin{proof}
Esto resulta trivialmente de \sref{I.4.1.6}.
\end{proof}

\subsection{Producto de inmersiones}
\label{subsection:I.4.3}
\label{I.4.3}

\begin{proposition}[4.3.1]
\label{I.4.3.1}
Sean $\alpha:X'\to X$, $\beta:Y'\to Y$ dos $S$-morfismos; si $\alpha$ y $\beta$ son inmersiones (\textup{resp.} inmersiones abiertas, inmersiones cerradas), entonces $\alpha\times_S\beta$ es una inmersión (\textup{resp.} una inmersión abierta, una inmersión cerrada).
Además, si $\alpha$ (\textup{resp.}~$\beta$) identifica $X'$ (\textup{resp.}~$Y'$) con un subpreesquema $X''$ (\textup{resp.}~$Y''$) de $X$ (\textup{resp.}~$Y$), entonces $\alpha\times_S\beta$ identifica el espacio subyacente de $X'\times_S Y'$ con el subespacio $p^{-1}(X'')\cap q^{-1}(Y'')$ del espacio subyacente de $X\times_S Y$, donde $p$ y $q$ denotan las proyecciones de $X\times_S Y$ sobre $X$ e $Y$, respectivamente.
\end{proposition}

\begin{proof}
Habida cuenta de la Definición~\sref{I.4.2.1}, se puede restringir al caso en que $X'$ e $Y'$ son subpreesquemas, y $\alpha$ y $\beta$ son los morfismos de inyección.
La proposición ya ha sido establecida para los subpreesquemas inducidos sobre los abiertos \sref{I.3.2.7};
como todo subpreesquema es un subpreesquema cerrado de un preesquema inducido sobre un abierto \sref{I.4.1.3}, se reduce al caso en que $X'$ e $Y'$ son subpreesquemas \emph{cerrados}.

Demostremos primero que se puede suponer que $S$ es \emph{afín}.
Sea $(S_\lambda)$ un recubrimiento de $S$ formado por abiertos afines;
si $\varphi$ y $\psi$ son los morfismos estructurales de $X$ e $Y$, sean $X_\lambda=\varphi^{-1}(S_\lambda)$ e $Y_\lambda=\psi^{-1}(S_\lambda)$.
La restricción $X'_\lambda$ (resp. $Y'_\lambda$) de $X'$ (resp. $Y'$) a $X_\lambda\cap X'$ (resp. $Y_\lambda\cap Y'$) es un subpreesquema cerrado de $X_\lambda$ (resp. $Y_\lambda$), los preesquemas $X_\lambda$, $Y_\lambda$, $X'_\lambda$ e $Y'_\lambda$ pueden considerarse como $S_\lambda$-preesquemas, y los productos $X_\lambda\times_S Y_\lambda$ y $X_\lambda\times_{S_\lambda}Y_\lambda$ (resp. $X'_\lambda\times_S Y'_\lambda$ y $X'_\lambda\times_{S_\lambda}Y'_\lambda$) son idénticos \sref{I.3.2.5}.
Si la proposición es verdadera cuando $S$ es afín, la restricción de $\alpha\times_S\beta$ a cada uno de los $X'_\lambda\times_S Y'_\lambda$ es, por tanto, una inmersión \sref{I.3.2.7}.
Como el producto $X'_\lambda\times_S Y'_\mu$ (resp. $X_\lambda\times_S Y_\mu$) se identifica con $(X'_\lambda\cap X'_\mu)\times_S(Y'_\lambda\cap Y'_\mu)$ (resp. $(X_\lambda\cap X_\mu)\times_S(Y_\lambda\cap Y_\mu)$) \sref{I.3.2.6.4}, la restricción de $\alpha\times_S\beta$
\oldpage[I]{125}
a cada uno de los $X'_\lambda\times_S Y'_\mu$ es también una inmersión;
lo mismo vale entonces para $\alpha\times_S\beta$ en virtud de \sref{I.4.2.4}.

Demostremos ahora que se puede suponer que $X$ e $Y$ son \emph{afines}.
En efecto, sea $(U_i)$ (resp. $(V_j)$) un recubrimiento de $X$ (resp. $Y$) por abiertos afines, y sea $X'_i$ (resp. $Y'_j$) la restricción de $X'$ (resp. $Y'$) a $X'\cap U_i$ (resp. $Y'\cap V_j$), que es un subpreesquema cerrado de $U_i$ (resp. $V_j$);
$U_i\times_S V_j$ se identifica con la restricción de $X\times_S Y$ a $p^{-1}(U_i)\cap q^{-1}(V_j)$ \sref{I.3.2.7};
de igual modo, si $p'$ y $q'$ son las proyecciones de $X'\times_S Y'$, entonces $X'_i\times_S Y'_j$ se identifica con la restricción de $X'\times_S Y'$ a ${p'}^{-1}(X'_i)\cap{q'}^{-1}(Y'_j)$.
Pongamos $\gamma=\alpha\times_S\beta$;
se tiene, por definición, $p\circ\gamma=\alpha\circ p'$ y $q\circ\gamma=\beta\circ q'$;
como $X'_i=\alpha^{-1}(U_i)$ e $Y'_j=\beta^{-1}(V_j)$, se tiene también ${p'}^{-1}(X'_i)=\gamma^{-1}(p^{-1}(U_i))$ y ${q'}^{-1}(Y'_j)=\gamma^{-1}(q^{-1}(V_j))$, de donde
\[
  {p'}^{-1}(X'_i)\cap{q'}^{-1}(Y'_j)
  =\gamma^{-1}\bigl(p^{-1}(U_i)\cap q^{-1}(V_j)\bigr)
  =\gamma^{-1}(U_i\times_S V_j)
\]
y se concluye entonces como en la primera parte del razonamiento.

Supongamos, pues, que $X$, $Y$ y $S$ son afines, y sean $B$, $C$ y $A$ sus anillos respectivos.
Entonces $B$ y $C$ son $A$-álgebras, y $X'$ e $Y'$ son esquemas afines cuyos anillos son las álgebras cociente $B'$ y $C'$ de $B$ y $C$, respectivamente.
Además, se tiene $\alpha=({}^a\rho,\widetilde{\rho})$ y $\beta=({}^a\sigma,\widetilde{\sigma})$, donde $\rho$ y $\sigma$ son, respectivamente, los homomorfismos canónicos $B\to B'$ y $C\to C'$ \sref{I.1.7.3}.
En estas condiciones, se sabe que $X\times_S Y$ (resp. $X'\times_S Y'$) es un esquema afín de anillo $B\otimes_A C$ (resp. $B'\otimes_A C'$), y $\alpha\times_S\beta=({}^a\tau,\widetilde{\tau})$, donde $\tau$ es el homomorfismo $\rho\otimes\sigma$ de $B\otimes_A C$ en $B'\otimes_A C'$ (Proposición~\sref{I.3.2.2} y Corolario~\sref{I.3.2.3});
como este homomorfismo es sobreyectivo, $\alpha\times_S\beta$ es una inmersión.
Además, si $\mathfrak{b}$ (resp. $\mathfrak{c}$) es el núcleo de $\rho$ (resp. $\sigma$), entonces el núcleo de $\tau$ es
la suma de los ideales engendrados por $u(\mathfrak{b})$ y $v(\mathfrak{c})$, donde $u$ (resp. $v$) es el homomorfismo $b\to b\otimes1$ (resp. $c\to1\otimes c$).
Como $p=({}^a u,\widetilde{u})$ y $q=({}^a v,\widetilde{v})$, este núcleo corresponde, en el espectro primo de $B\otimes_A C$, al conjunto cerrado $p^{-1}(X')\cap q^{-1}(Y')$, en virtud de \sref{I.1.2.2.1} y de la Proposición~\sref{I.1.1.2}[iii], lo que concluye la demostración.
\end{proof}

\begin{corollary}[4.3.2]
\label{I.4.3.2}
Si $f:X\to Y$ es una inmersión (\textup{resp.} una inmersión abierta, una inmersión cerrada) y un $S$-morfismo, entonces $f_{(S')}$ es una inmersión (\textup{resp.} una inmersión abierta, una inmersión cerrada) para toda extensión $S'\to S$ del preesquema base.
\end{corollary}

\subsection{Imagen inversa de un subpreesquema}
\label{subsection:I.4.4}

\begin{proposition}[4.4.1]
\label{I.4.4.1}
Sean $f:X\to Y$ un morfismo, $Y'$ un subpreesquema (\textup{resp.} un subpreesquema cerrado, un preesquema inducido sobre un abierto) de $Y$, y $j:Y'\to Y$ el morfismo de inyección.
Entonces la proyección $p:X\times_Y Y'\to X$ es una inmersión (\textup{resp.} una inmersión cerrada, una inmersión abierta);
el espacio subyacente del subpreesquema de $X$ asociado a $p$ es $f^{-1}(Y')$;
además, si $j'$ es el morfismo de inyección de este subpreesquema en $X$, entonces, para todo morfismo $h:Z\to X$, el compuesto $f\circ h:Z\to Y$ está dominado por $j$ si y solo si $h$ está dominado por $j'$.
\end{proposition}

\begin{proof}
Como $p=1_X\times_Y j$ \sref{I.3.3.4}, la primera aserción resulta de la Proposición~\sref{I.4.3.1};
la segunda es un caso particular del Corolario~\sref{I.3.5.10} (intercambiando los papeles de $X$ y de $Y'$).
Finalmente, si se tiene $f\circ h=j\circ h'$, donde $h'$ es un morfismo $Z\to Y'$, de la definición del producto resulta que se tiene $h=p\circ u$, donde $u$ es un morfismo $Z\to X\times_Y Y'$, de donde se sigue la última aserción.
\end{proof}

Diremos que el subpreesquema de $X$ así definido es la \emph{imagen inversa} del subpreesquema $Y'$ de $Y$ por el morfismo $f$, terminología que concuerda con la introducida
\oldpage[I]{126}
de manera más general en \sref{I.3.3.6}.
Cuando hablemos de $f^{-1}(Y')$ como de un subpreesquema de $X$, se tratará siempre de este subpreesquema.

Cuando los preesquemas $f^{-1}(Y')$ y $X$ son idénticos, $j'$ es la identidad y todo morfismo $h:Z\to X$ está, por tanto, dominado por $j'$; por tanto, \emph{el morfismo $f:X\to Y$ se factoriza como $X\xrightarrow{g}Y'\xrightarrow{j}Y$}.

Cuando $y$ es un punto \emph{cerrado} de $Y$ y $Y'=\operatorname{Spec}(\boldsymbol{k}(y))$ es el menor subpreesquema cerrado de $Y$ que tiene $\{y\}$ como espacio subyacente \sref{I.4.1.9}, el subpreesquema cerrado $f^{-1}(Y')$ es canónicamente isomorfo a la \emph{fibra} $f^{-1}(y)$ definida en \sref{I.3.6.2}, con la cual se lo identificará.

\begin{corollary}[4.4.2]
\label{I.4.4.2}
Sean $f:X\to Y$ y $g:Y\to Z$ dos morfismos, y sea $h=g\circ f$ su compuesto.
Para todo subpreesquema $Z'$ de $Z$, los subpreesquemas $f^{-1}(g^{-1}(Z'))$ y $h^{-1}(Z')$ de $X$ son idénticos.
\end{corollary}

\begin{proof}
Esto resulta de la existencia del isomorfismo canónico $X\times_Y(Y\times_Z Z')\xrightarrow{\sim}X\times_Z Z'$ \sref{I.3.3.9.1}.
\end{proof}

\begin{corollary}[4.4.3]
\label{I.4.4.3}
Sean $X'$ y $X''$ dos subpreesquemas de $X$, y sean $j':X'\to X$ y $j'':X''\to X$ sus morfismos de inyección;
entonces ${j'}^{-1}(X'')$ y ${j''}^{-1}(X')$ son ambos iguales al ínfimo $\inf(X',X'')$ de $X'$ y $X''$ para la relación de orden entre subpreesquemas, y son canónicamente isomorfos a $X'\times_X X''$.
\end{corollary}

\begin{proof}
Esto se deduce inmediatamente de la Proposición~\sref{I.4.4.1} y del Corolario~\sref{I.4.1.10}.
\end{proof}

\begin{corollary}[4.4.4]
\label{I.4.4.4}
Sea $f:X\to Y$ un morfismo, y sean $Y'$ e $Y''$ dos subpreesquemas de $Y$;
entonces se tiene \[
  f^{-1}(\inf(Y',Y''))=\inf(f^{-1}(Y'),f^{-1}(Y'')).
\].
\end{corollary}

\begin{proof}
Esto se deduce de la existencia del isomorfismo canónico entre $(X\times_Y Y')\times_X(X\times_Y Y'')$ y $X\times_Y(Y'\times_Y Y'')$ \sref{I.3.3.9.1}.
\end{proof}

\begin{proposition}[4.4.5]
\label{I.4.4.5}
Sea $f:X\to Y$ un morfismo, y sea $Y'$ un subpreesquema cerrado de $Y$ definido por un haz cuasicoherente de ideales $\mathscr{K}$ de $\mathscr{O}_Y$ \sref{I.4.1.3};
el subpreesquema cerrado $f^{-1}(Y')$ de $X$ está entonces definido por el haz cuasicoherente de ideales $f^*(\mathscr{K})\mathscr{O}_X$ de $\mathscr{O}_X$.
\end{proposition}

\begin{proof}
La cuestión es evidentemente local en $X$ e $Y$;
basta, pues, observar que si \erratum[II]{$A$ es una $B$-álgebra} y $\mathfrak{K}$ un ideal de $B$, entonces se tiene $A\otimes_B(B/\mathfrak{K})=A/\mathfrak{K}A$, y aplicar \sref{I.1.6.9}.
\end{proof}

\begin{corollary}[4.4.6]
\label{I.4.4.6}
Sean $X'$ un subpreesquema cerrado de $X$ definido por un haz cuasicoherente de ideales $\mathscr{J}$ de $\mathscr{O}_X$, e $i$ la inyección $X'\to X$;
el morfismo $f\circ i$, restricción de $f$ a $X'$, está dominado por la inyección $j:Y'\to Y$ \textup{(en otras palabras, se factoriza como $j\circ g$, con $g$ un morfismo $X'\to Y'$)} si y solo si $f^*(\mathscr{K})\mathscr{O}_X\subset\mathscr{J}$.
\end{corollary}

\begin{proof}
Basta aplicar la Proposición~\sref{I.4.4.1} a $i$, teniendo en cuenta la Proposición~\sref{I.4.4.5}.
\end{proof}

\subsection{Inmersiones locales e isomorfismos locales}
\label{subsection:I.4.5}

\begin{definition}[4.5.1]
\label{I.4.5.1}
{\itshape Sea $f:X\to Y$ un morfismo de preesquemas.
Se dice que $f$ es una inmersión local en un punto $x\in X$ si existen una vecindad abierta $U$ de $x$ en $X$ y una vecindad abierta $V$ de $f(x)$ en $Y$ tales que la restricción de $f$ al preesquema inducido $U$ sea una inmersión cerrada de $U$ en el preesquema inducido $V$.
Se dice que $f$ es una inmersión local si $f$ es una inmersión local en todo punto de $X$.}
\end{definition}

\begin{definition}[4.5.2]
\label{I.4.5.2}
Se dice que un morfismo $f:X\to Y$ es un isomorfismo local en
\oldpage[I]{127}
un punto $x\in X$ si existe una vecindad abierta $U$ de $x$ en $X$ tal que la restricción de $f$ al preesquema inducido $U$ sea una inmersión abierta de $U$ en $Y$.
Se dice que $f$ es un isomorfismo local si $f$ es un isomorfismo local en todo punto de $X$.
\end{definition}

\begin{env}[4.5.3]
\label{I.4.5.3}
Una inmersión (resp. una inmersión cerrada) $f:X\to Y$ puede caracterizarse como una inmersión local tal que $f$ sea un homeomorfismo del espacio subyacente a $X$ sobre un subconjunto (resp. un subconjunto cerrado) de $Y$.
Una inmersión abierta $f$ puede caracterizarse como un isomorfismo local \emph{inyectivo}.
\end{env}

\begin{proposition}[4.5.4]
\label{I.4.5.4}
Sea $X$ un preesquema irreducible y $f:X\to Y$ un morfismo inyectivo dominante.
Si $f$ es una inmersión local, entonces $f$ es una inmersión y $f(X)$ es abierto en $Y$.
\end{proposition}

\begin{proof}
Sea $x\in X$, y sean $U$ una vecindad abierta de $x$ y $V$ una vecindad abierta de $f(x)$ en $Y$ tales que la restricción de $f$ a $U$ sea una inmersión cerrada en $V$;
como $U$ es denso en $X$, $f(U)$ es denso en $Y$ por hipótesis, por lo que $f(U)=V$, y $f$ es un homeomorfismo de $U$ sobre $V$;
la hipótesis de que $f$ es inyectivo implica que $f^{-1}(V)=U$, de donde resulta la proposición.
\end{proof}

\begin{proposition}[4.5.5]
\label{I.4.5.5}
\medskip\noindent
\begin{enumerate}
  \item[{\rm(i)}] El compuesto de dos inmersiones locales (\textup{resp.} de dos isomorfismos locales) es una inmersión local (\textup{resp.} un isomorfismo local).
  \item[{\rm(ii)}] Sean $f:X\to X'$ y $g:Y\to Y'$ dos $S$-morfismos.
    Si $f$ y $g$ son inmersiones locales (\textup{resp.} isomorfismos locales), entonces $f\times_S g$ es asimismo una inmersión local (resp. un isomorfismo local).
  \item[{\rm(iii)}] Si un $S$-morfismo $f$ es una inmersión local (\textup{resp.} un isomorfismo local), entonces $f_{(S')}$ también lo es para toda extensión $S'\to S$ del preesquema de base.
\end{enumerate}
\end{proposition}

\begin{proof}
Según \sref{I.3.5.1}, basta demostrar (i) y (ii).

(i) resulta inmediatamente de la transitividad de las inmersiones cerradas (resp. abiertas) \sref{I.4.2.5}\footnote{\emph{[Trad.] La fuente francesa cita aquí (4.2.4); la transitividad es la Proposición~\hyperref[I.4.2.5]{4.2.5}.}} y del hecho de que, si $f$ es un homeomorfismo de $X$ sobre un subconjunto cerrado de $Y$, para todo abierto $U\subset X$, $f(U)$ es abierto en $f(X)$; por tanto, existe un subconjunto abierto $V$ de $Y$ tal que $f(U)=V\cap f(X)$ y, en consecuencia, $f(U)$ es cerrado en $V$.

Para demostrar (ii), sea $z\in X\times_S Y$, pongamos $z'=(f\times_S g)(z)$\footnote{\emph{[Trad.] La fuente francesa utiliza $z$ y $z'$ más abajo sin introducirlos; esta inserción proporciona el punto necesario y su imagen.}}, sean $p$ y $q$ las proyecciones de $X\times_S Y$, y sean $p'$ y $q'$ las proyecciones de $X'\times_S Y'$.
Existen, por hipótesis, vecindades abiertas $U$, $U'$, $V$ y $V'$ de $x=p(z)$, $x'=p'(z')$, $y=q(z)$ e $y'=q'(z')$, respectivamente, tales que las restricciones de $f$ y $g$ a $U$ y $V$, respectivamente, sean inmersiones cerradas (resp. abiertas) en $U'$ y $V'$, respectivamente.
Como los espacios subyacentes de $U\times_S V$ y $U'\times_S V'$ pueden identificarse con las vecindades abiertas $p^{-1}(U)\cap q^{-1}(V)$ y ${p'}^{-1}(U')\cap{q'}^{-1}(V')$ de $z$ y $z'$, respectivamente, \sref{I.3.2.7}, la proposición se deduce de la Proposición~\sref{I.4.3.1}.
\end{proof}
% END source_es/ega1/ega1-4.tex

% BEGIN source_es/ega1/ega1-5.tex
\section{Preesquemas reducidos; condición de separación}
\label{section:I.5}

\subsection{Preesquemas reducidos}
\label{subsection:I.5.1}

\begin{proposition}[5.1.1]
\label{I.5.1.1}
Sean $(X,\mathscr{O}_X)$ un preesquema y $\mathscr{B}$ una $\mathscr{O}_X$-álgebra cuasicoherente.
Entonces existe un único $\mathscr{O}_X$-módulo cuasicoherente $\mathscr{N}$ cuyo tallo $\mathscr{N}_x$ en todo punto $x\in X$ es el nilradical del anillo $\mathscr{B}_x$.
Cuando $X$ es afín y, por consiguiente, $\mathscr{B}=\widetilde{B}$, donde $B$ es un álgebra sobre $A(X)$, se tiene $\mathscr{N}=\widetilde{\mathfrak{N}}$, donde $\mathfrak{N}$ es el nilradical de $B$.
\end{proposition}

\begin{proof}
\oldpage[I]{128}
Como la cuestión es local, basta demostrar la última aserción.
Se sabe que $\widetilde{\mathfrak{N}}$ es un $\mathscr{O}_X$-módulo cuasicoherente \sref{I.1.4.1} y que su tallo en un punto $x\in X$ es el ideal $\mathfrak{N}_x$ del anillo de fracciones $B_x$;
queda por demostrar que el nilradical de $B_x$ está contenido en $\mathfrak{N}_x$, pues la inclusión opuesta es evidente.
Sea $z/s$ un elemento del nilradical de $B_x$, con $z\in B$ y $s\notin\mathfrak{j}_x$;
por hipótesis, existe un entero $k$ tal que $(z/s)^k=0$, lo cual implica que existe un $t\notin\mathfrak{j}_x$ tal que $tz^k=0$.
Se concluye que $(tz)^k=0$ y, por consiguiente, que $z/s=(tz)/(ts)$ pertenece a $\mathfrak{N}_x$.
\end{proof}

Diremos que el $\mathscr{O}_X$-módulo cuasicoherente $\mathscr{N}$ así definido es el \emph{nilradical} de la $\mathscr{O}_X$-álgebra $\mathscr{B}$; en particular, denotaremos por $\mathscr{N}_X$ el nilradical de $\mathscr{O}_X$.

\begin{corollary}[5.1.2]
\label{I.5.1.2}
Sea $X$ un preesquema;
el subpreesquema cerrado de $X$ definido por el haz de ideales $\mathscr{N}_X$ es el único subpreesquema reducido \sref[0]{0.4.1.4} de $X$ que tiene $X$ como espacio subyacente;
es también el menor subpreesquema de $X$ que tiene $X$ como espacio subyacente.
\end{corollary}

\begin{proof}
Como el haz estructural del subpreesquema cerrado $Y$ definido por $\mathscr{N}_X$ es $\mathscr{O}_X/\mathscr{N}_X$, es inmediato que $Y$ es reducido y tiene $X$ como espacio subyacente, pues $\mathscr{N}_x\ne\mathscr{O}_x$ para todo $x\in X$.
Para demostrar las demás aserciones, observemos que un subpreesquema $Z$ de $X$ que tenga $X$ como espacio subyacente está definido por un haz de ideales $\mathscr{J}$ \sref{I.4.1.3} tal que $\mathscr{J}_x\ne\mathscr{O}_x$ para todo $x\in X$.
Se puede restringir al caso en que $X$ es afín, digamos $X=\operatorname{Spec}(A)$ y $\mathscr{J}=\widetilde{\mathfrak{J}}$, donde $\mathfrak{J}$ es un ideal de $A$;
entonces, para todo $x\in X$, se tiene $\mathfrak{J}\subset\mathfrak{j}_x$\footnote{\emph{[Trad.] La fuente francesa imprime aquí $\mathfrak{I}_x\subset\mathfrak{j}_x$, comparación mal tipada entre un ideal localizado de $A_x$ y un ideal de $A$; el argumento exige $\mathfrak{I}\subset\mathfrak{j}_x$.}}, de modo que $\mathfrak{J}$ está contenido en todo ideal primo de $A$ y, por tanto, también en su intersección $\mathfrak{N}$, el nilradical de $A$.
Esto demuestra que $Y$ es el menor subpreesquema de $X$ que tiene $X$ como espacio subyacente \sref{I.4.1.9};
además, si $Z$ es distinto de $Y$, se tiene necesariamente $\mathscr{J}_x\ne\mathscr{N}_x$ para al menos un $x\in X$ y, por \sref{I.5.1.1}, $Z$ no es reducido.
\end{proof}

\begin{definition}[5.1.3]
\label{I.5.1.3}
Se llama preesquema reducido asociado a un preesquema $X$, y se denota por $X_{\mathrm{red}}$, al único subpreesquema reducido de $X$ que tiene $X$ como espacio subyacente.
\end{definition}

Decir que un preesquema $X$ es reducido significa, pues, que $X=X_{\mathrm{red}}$.

\begin{proposition}[5.1.4]
\label{I.5.1.4}
Para que el espectro primo de un anillo $A$ sea un preesquema reducido (\textup{resp.} íntegro) \sref{I.2.1.8}\footnote{\emph{[Trad.] La fuente francesa remite aquí a (2.1.7), que trata de la restricción a abiertos; la definición pertinente de preesquema íntegro es (2.1.8).}}, es necesario y suficiente que $A$ sea un anillo reducido (\textup{resp.} íntegro).
\end{proposition}

\begin{proof}
En efecto, de \sref{I.5.1.1} resulta inmediatamente que la condición $\mathscr{N}=(0)$ es necesaria y suficiente para que $X=\operatorname{Spec}(A)$ sea reducido;
la aserción relativa a los anillos íntegros es entonces consecuencia de \sref{I.1.1.13}.
\end{proof}

Como todo anillo de fracciones $\ne\{0\}$ de un anillo íntegro es íntegro, de \sref{I.5.1.4} resulta que, para todo preesquema \emph{localmente íntegro} $X$, $\mathscr{O}_x$ es un anillo \emph{íntegro} para todo $x\in X$.
La recíproca es verdadera siempre que el espacio subyacente a $X$ sea \emph{localmente noetheriano}:
en efecto, $X$ es entonces reducido y, si $U$ es un abierto afín de $X$, que es un espacio noetheriano, $U$ solo tiene un número finito de componentes irreducibles, por lo que su anillo $A$ solo tiene un número finito de ideales primos minimales \sref{I.1.1.14}.
Si dos de las componentes irreducibles de $U_i$ tuviesen un punto común $x$, entonces $\mathscr{O}_x$ tendría al menos dos ideales primos minimales distintos y, por tanto, no sería íntegro;
las componentes de $U_i$ son, pues, abiertos disjuntos dos a dos, y cada una de ellas es íntegra.

\begin{env}[5.1.5]
\label{I.5.1.5}
Sea $f=(\psi,\theta):X\to Y$ un morfismo de preesquemas;
\oldpage[I]{129}
el homomorfismo $\theta_x^\sharp:\mathscr{O}_{\psi(x)}\to\mathscr{O}_x$ envía todo elemento nilpotente de $\mathscr{O}_{\psi(x)}$ a un elemento nilpotente de $\mathscr{O}_x$;
al pasar a los cocientes, $\theta^\sharp$ induce un homomorfismo
\[
  \omega:\psi^*(\mathscr{O}_Y/\mathscr{N}_Y)
  \longrightarrow \mathscr{O}_X/\mathscr{N}_X
\]
Es claro que, para todo $x\in X$, $\omega_x:\mathscr{O}_{\psi(x)}/\mathscr{N}_{\psi(x)}
\to\mathscr{O}_x/\mathscr{N}_x$ es un homomorfismo local, de modo que $(\psi,\omega^\flat)$ es un morfismo de preesquemas $X_{\mathrm{red}}\to Y_{\mathrm{red}}$, que denotaremos por $f_{\mathrm{red}}$ y llamaremos el morfismo \emph{reducido} asociado a $f$.
Es inmediato que, para dos morfismos $f:X\to Y$ y $g:Y\to Z$, se tiene $(g\circ f)_{\mathrm{red}}=g_{\mathrm{red}}\circ f_{\mathrm{red}}$, de modo que se ha definido $X_{\mathrm{red}}$ como un \emph{funtor}, \emph{covariante} en $X$.

La definición precedente muestra que el diagrama
\[
  \xymatrix{
    X_{\mathrm{red}}\ar[r]^{f_{\mathrm{red}}}\ar[d] &
    Y_{\mathrm{red}}\ar[d]\\
    X\ar[r]_f & Y
  }
\]
es conmutativo, siendo las flechas verticales los morfismos de inyección;
dicho de otro modo, $X_{\mathrm{red}}\to X$ es un morfismo \emph{funtorial}.
Se observará en particular que, si $X$ es reducido, todo morfismo $f:X\to Y$ se factoriza como $X\xrightarrow{f_{\mathrm{red}}}Y_{\mathrm{red}}\to Y$;
dicho de otro modo, $f$ está dominado por el morfismo de inyección $Y_{\mathrm{red}}\to Y$.
\end{env}

\begin{proposition}[5.1.6]
\label{I.5.1.6}
Sea $f:X\to Y$ un morfismo;
si $f$ es sobreyectivo (\textup{resp.} radicial, una inmersión, una inmersión cerrada, una inmersión abierta, una inmersión local, un isomorfismo local), entonces $f_{\mathrm{red}}$ también lo es.
Recíprocamente, si $f_{\mathrm{red}}$ es sobreyectivo (\textup{resp.} radicial), entonces $f$ también lo es.
\end{proposition}

\begin{proof}
La proposición es trivial si $f$ es sobreyectivo;
si $f$ es radicial, resulta del hecho de que, para todo $x\in X$, el cuerpo $\boldsymbol{k}(x)$ es el mismo para los preesquemas $X$ y $X_{\mathrm{red}}$ \sref{I.3.5.8}.
Por último, si $f=(\psi,\theta)$ es una inmersión, una inmersión cerrada o una inmersión local (resp. una inmersión abierta o un isomorfismo local), la proposición resulta del hecho de que, si $\theta_x^\sharp$ es sobreyectivo (resp. biyectivo), también lo es el homomorfismo obtenido al pasar a los cocientes por los nilradicales de $\mathscr{O}_{\psi(x)}$ y $\mathscr{O}_x$ (\sref{I.5.1.2} y \sref{I.4.2.2}) (cf. \sref{I.5.5.12}).
\end{proof}

\begin{proposition}[5.1.7]
\label{I.5.1.7}
Si $X$ e $Y$ son dos $S$-preesquemas, entonces los preesquemas $X_{\mathrm{red}}\times_{S_{\mathrm{red}}}Y_{\mathrm{red}}$ y $X_{\mathrm{red}}\times_S Y_{\mathrm{red}}$ son idénticos, y se identifican canónicamente con un subpreesquema de $X\times_S Y$ que tiene el mismo espacio subyacente que dicho producto.
\end{proposition}

\begin{proof}
La identificación canónica de $X_{\mathrm{red}}\times_S Y_{\mathrm{red}}$ con un subpreesquema de $X\times_S Y$ que tiene el mismo espacio subyacente resulta de \sref{I.4.3.1}.
Además, si $\varphi$ y $\psi$ son, respectivamente, los morfismos estructurales $X_{\mathrm{red}}\to S$ e $Y_{\mathrm{red}}\to S$, estos se factorizan a través de $S_{\mathrm{red}}$ \sref{I.5.1.5}, y, como $S_{\mathrm{red}}\to S$ es un monomorfismo, la primera afirmación de la proposición resulta de \sref{I.3.2.4}.
\end{proof}

\begin{corollary}[5.1.8]
\label{I.5.1.8}
Los preesquemas $(X\times_S Y)_{\mathrm{red}}$ y $(X_{\mathrm{red}}\times_{S_{\mathrm{red}}}Y_{\mathrm{red}})_{\mathrm{red}}$ se identifican canónicamente entre sí.
\end{corollary}

\begin{proof}
Esto resulta de \sref{I.5.1.2} y \sref{I.5.1.7}.
\end{proof}

Se observará que, aunque $X$ e $Y$ sean preesquemas reducidos, $X\times_S Y$ no es necesariamente reducido, pues el producto tensorial de dos álgebras reducidas puede tener elementos nilpotentes.


\oldpage[I]{130}
\begin{proposition}[5.1.9]
\label{I.5.1.9}

Sea $X$ un preesquema, y sea $\mathscr{J}$ un haz cuasicoherente de ideales de $\mathscr{O}_X$ tal que $\mathscr{J}^n=0$ para algún entero $n>0$.
Sea $X_0$ el subpreesquema cerrado $(X,\mathscr{O}_X/\mathscr{J})$ de $X$;
para que $X$ sea un esquema afín, es necesario y suficiente que $X_0$ sea un esquema afín.
\end{proposition}

La condición es evidentemente necesaria; probemos que es suficiente.
Si se pone $X_k=(X,\mathscr{O}_X/\mathscr{J}^{k+1})$, basta demostrar por inducción sobre $k$ que $X_k$ es afín, con lo cual se reduce al caso en que $\mathscr{J}^2=0$.
Pongamos
\[
  A=\Gamma(X,\mathscr{O}_X),\qquad
  A_0=\Gamma(X_0,\mathscr{O}_{X_0})
      =\Gamma(X,\mathscr{O}_X/\mathscr{J}).
\]
El homomorfismo canónico $\mathscr{O}_X\to\mathscr{O}_X/\mathscr{J}$ induce un homomorfismo de anillos $\varphi:A\to A_0$.
Veremos a continuación que $\varphi$ es \emph{sobreyectivo}, de modo que
\[
  0\to\Gamma(X,\mathscr{J})\to\Gamma(X,\mathscr{O}_X)
  \to\Gamma(X,\mathscr{O}_X/\mathscr{J})\to 0
  \tag{5.1.9.1}\label{I.5.1.9.1-fr}
\]
es una sucesión \emph{exacta}.
Supongamos demostrado este punto y veamos que de él se sigue la proposición.
Observemos que $\mathfrak{K}=\Gamma(X,\mathscr{J})$ es un ideal de cuadrado cero en $A$, y por tanto un módulo sobre $A_0=A/\mathfrak{K}$.
Por hipótesis, se tiene $X_0=\operatorname{Spec}(A_0)$ y, como los espacios subyacentes de $X_0$ y $X$ son idénticos, $\mathfrak{K}=\Gamma(X_0,\mathscr{J})$;
además, como $\mathscr{J}^2=0$, $\mathscr{J}$ es un $(\mathscr{O}_X/\mathscr{J})$-módulo cuasicoherente, de modo que $\mathscr{J}\cong\widetilde{\mathfrak{K}}$ y $\mathfrak{K}_x=\mathscr{J}_x$ para todo $x\in X_0$ \sref{I.1.4.1}.
En estas condiciones, sea $X'=\operatorname{Spec}(A)$, y consideremos el morfismo $f=(\psi,\theta):X\to X'$ de preesquemas correspondiente a la aplicación identidad $A\to\Gamma(X,\mathscr{O}_X)$ \sref{I.2.2.4}.
Para todo subconjunto abierto afín $V$ de $X$, el diagrama
\[
  \xymatrix{
    A\ar[r]\ar[d] & \Gamma(V,\mathscr{O}_X|V)\ar[d]\\
    A_0=A/\mathfrak{K}\ar[r] & \Gamma(V,\mathscr{O}_{X_0}|V)
  }
\]
es conmutativo, de donde se deduce que el diagrama
\[
  \xymatrix{
    X' & X\ar[l]_f\\
    X'_0\ar[u]^{j'} & X_0\ar[l]^{f_0}\ar[u]_j
  }
\]
también es conmutativo (siendo $X'_0$ el subpreesquema cerrado de $X'$ definido por el haz cuasicoherente de ideales $\widetilde{\mathfrak{K}}$, y siendo $j$ y $j'$ los morfismos de inyección canónicos).
Pero, como $X_0$ es afín, $f_0$ es un isomorfismo y, como las aplicaciones continuas subyacentes a $j$ y $j'$ son las aplicaciones identidad, se ve inmediatamente que $\psi:X\to X'$ es un homeomorfismo.
Además, la igualdad $\mathfrak{K}_x=\mathscr{J}_x$ muestra que la restricción de $\theta^\sharp:\psi^*(\mathscr{O}_{X'})\to\mathscr{O}_X$ es un \emph{isomorfismo} de $\psi^*(\widetilde{\mathfrak{K}})$ sobre $\mathscr{J}$;
por otra parte, al pasar a los cocientes, $\theta^\sharp$ da un \emph{isomorfismo} $\psi^*(\mathscr{O}_{X'}/\widetilde{\mathfrak{K}})
\to\mathscr{O}_X/\mathscr{J}$, pues $f_0$ es un isomorfismo;
se concluye entonces de inmediato, por el lema de los cinco (M,~I,~1.1), que $\theta^\sharp$ es a su vez un isomorfismo y, por tanto, que $f$ es un \emph{isomorfismo}, y, en consecuencia, que $X$ es afín.
Todo se reduce, pues, a demostrar la exactitud de \sref{I.5.1.9.1}, que se seguirá de demostrar que 
\oldpage[I]{131}
$H^1(X,\mathscr{J})=0$.

Pero $H^1(X,\mathscr{J})=H^1(X_0,\mathscr{J})$, y hemos visto que $\mathscr{J}$ es un $\mathscr{O}_{X_0}$-módulo cuasicoherente.
Nuestra demostración se seguirá, por tanto, del
\begin{lemma}[5.1.9.2]
\label{I.5.1.9.2}
Si $Y$ es un esquema afín y $\mathscr{F}$ un $\mathscr{O}_Y$-módulo cuasicoherente, entonces $H^1(Y,\mathscr{F})=0$.
\end{lemma}

\begin{proof}
Este lema se demostrará en el Capítulo~III, §1, como consecuencia del teorema más general según el cual $H^i(Y,\mathscr{F})=0$ para todo $i>0$.
Para dar una demostración independiente, observemos que $H^1(Y,\mathscr{F})$ puede identificarse con el módulo $\operatorname{Ext}^1_{\mathscr{O}_Y}
(Y;\mathscr{O}_Y,\mathscr{F})$ de las clases de extensiones del $\mathscr{O}_Y$-módulo $\mathscr{O}_Y$ por el $\mathscr{O}_Y$-módulo $\mathscr{F}$ (T,~4.2.3);
todo se reduce, pues, a demostrar que tal extensión $\mathscr{G}$ es trivial.
Ahora bien, para todo $y\in Y$, existe una vecindad $V$ de $y$ en $Y$ tal que $\mathscr{G}|V$ es isomorfo a $\mathscr{F}|V\oplus\mathscr{O}_Y|V$\footnote{\emph{[Trad.] La fuente francesa imprime aquí $\sh{F}|Y$; la restricción necesaria sobre la vecindad $V$ es $\sh{F}|V$.}} \sref[0]{0.5.4.9};
de ello se deduce que $\mathscr{G}$ es un módulo \emph{cuasicoherente} sobre $\mathscr{O}_Y$.
Si $A$ es el anillo de $Y$, entonces se tiene $\mathscr{F}=\widetilde{M}$ y $\mathscr{G}=\widetilde{N}$, donde $M$ y $N$ son $A$-módulos, y, por hipótesis, $N$ es una extensión del $A$-módulo $A$ por el $A$-módulo $M$ \sref{I.1.3.11}.
Como esta extensión es necesariamente trivial, queda demostrado el lema, y con ello también \sref{I.5.1.9}.
\end{proof}

\begin{corollary}[5.1.10]
\label{I.5.1.10}
Sea $X$ un preesquema tal que $\mathscr{N}_X$ sea nilpotente.
Para que $X$ sea un esquema afín, es necesario y suficiente que $X_{\mathrm{red}}$ sea un esquema afín.
\end{corollary}

\subsection{Existencia de un subpreesquema con un espacio subyacente dado}
\label{subsection:I.5.2}

\begin{proposition}[5.2.1]
\label{I.5.2.1}
Para todo subespacio localmente cerrado $Y$ del espacio subyacente a un preesquema $X$, existe un único subpreesquema reducido de $X$ que tiene $Y$ como espacio subyacente.
\end{proposition}

\begin{proof}
La unicidad resulta de \sref{I.5.1.2}, por lo que solo queda demostrar la existencia del preesquema en cuestión.

Si $X$ es afín, de anillo $A$, e $Y$ es cerrado en $X$, la proposición es inmediata:
$\mathfrak{j}(Y)$ es el mayor ideal $\mathfrak{a}\subset A$ tal que $V(\mathfrak{a})=Y$, y es igual a su radical \sref{I.1.1.4}[i], por lo que $A/\mathfrak{j}(Y)$ es un anillo reducido.

En el caso general, para todo abierto afín $U\subset X$ tal que $U\cap Y$ sea cerrado en $U$, consideremos el subpreesquema cerrado $Y_U$ de $U$ definido por el haz de ideales asociado al ideal $\mathfrak{j}(U\cap Y)$ de $A(U)$.
Este subpreesquema es reducido.
Afirmamos que, si $V$ es un abierto afín de $X$ contenido en $U$, entonces $Y_V$ es \emph{inducido} por $Y_U$ sobre $V\cap Y$.
En efecto, este preesquema inducido es un subpreesquema cerrado reducido de $V$ que tiene $V\cap Y$ como espacio subyacente;
la unicidad de $Y_V$ demuestra, por tanto, la afirmación.
\end{proof}

\begin{proposition}[5.2.2]
\label{I.5.2.2}
Sea $X$ un preesquema reducido, sea $f:X\to Y$ un morfismo, y sea $Z$ un subpreesquema cerrado de $Y$ tal que $f(X)\subset Z$.
Entonces $f$ se factoriza como $X\xrightarrow{g}Z\xrightarrow{j}Y$, donde $j$ es el morfismo de inyección.
\end{proposition}

\begin{proof}
De las hipótesis resulta que el subpreesquema cerrado $f^{-1}(Z)$ de $X$ tiene todo $X$ como espacio subyacente \sref{I.4.4.1};
como $X$ es reducido, este subpreesquema cerrado coincide con $X$ \sref{I.5.1.2}, y la proposición resulta entonces de \sref{I.4.4.1}.
\end{proof}

\begin{corollary}[5.2.3]
\label{I.5.2.3}
Sea $X$ un subpreesquema reducido de un preesquema $Y$; si $Z$ es el subpreesquema cerrado reducido de $Y$ que tiene $\overline{X}$ como espacio subyacente, entonces $X$ es un subpreesquema inducido sobre un abierto de $Z$.
\end{corollary}

\begin{proof}
\oldpage[I]{132}
Existe, en efecto, un abierto $U$ de $Y$ tal que $X=U\cap\overline{X}$;
como, por \sref{I.5.2.2}, $X$ es un subpreesquema reducido de $Z$, el subpreesquema $X$ es inducido por $Z$ sobre el subespacio abierto $X$ en virtud de la unicidad \sref{I.5.2.1}.
\end{proof}

\begin{corollary}[5.2.4]
\label{I.5.2.4}
Sea $f:X\to Y$ un morfismo, y sea $X'$ (\textup{resp.} $Y'$) un subpreesquema cerrado de $X$ (\textup{resp.} $Y$) definido por un haz cuasicoherente de ideales $\mathscr{J}$ (\textup{resp.} $\mathscr{K}$) de $\mathscr{O}_X$ (\textup{resp.} $\mathscr{O}_Y$).
Supongamos que $X'$ es reducido y que $f(X')\subset Y'$.
Entonces $f^*(\mathscr{K})\mathscr{O}_X\subset\mathscr{J}$.
\end{corollary}

\begin{proof}
Como, por \sref{I.5.2.2}, la restricción de $f$ a $X'$ se factoriza como $X'\to Y'\to Y$, basta aplicar \sref{I.4.4.6}.
\end{proof}

\subsection{Diagonal; grafo de un morfismo}
\label{subsection:I.5.3}
\label{I.5.3}

\begin{env}[5.3.1]
\label{I.5.3.1}
Sea $X$ un $S$-preesquema;
se llama \emph{morfismo diagonal} de $X$ en $X\times_S X$, y se denota por $\Delta_{X|S}$, o $\Delta_X$, o incluso $\Delta$ si no puede haber confusión, al $S$-morfismo $(1_X,1_X)_S$, es decir, al único $S$-morfismo $\Delta_X$ tal que
\[
  p_1\circ\Delta_X=p_2\circ\Delta_X=1_X
  \tag{5.3.1.1}\label{I.5.3.1.1-fr}
\]
donde $p_1$ y $p_2$ son las proyecciones de $X\times_S X$ (Definición~\sref{I.3.2.1}).
Si $f:T\to X$ y $g:T\to Y$ son $S$-morfismos, se verifica inmediatamente que
\[
  (f,g)_S=(f\times_S g)\circ\Delta_{T|S}
  \tag{5.3.1.2}\label{I.5.3.1.2-fr}
\]

El lector observará que la definición precedente y los resultados enunciados desde \sref{I.5.3.1} hasta \sref{I.5.3.8} son válidos en toda categoría, \emph{siempre que los productos que allí figuran existan} en esta categoría.
\end{env}

\begin{proposition}[5.3.2]
\label{I.5.3.2}
Sean $X$ e $Y$ dos $S$-preesquemas;
si se identifica canónicamente $(X\times Y)\times(X\times Y)$ con $(X\times X)\times(Y\times Y)$, entonces el morfismo $\Delta_{X\times Y}$ se identifica con $\Delta_X\times\Delta_Y$.
\end{proposition}

\begin{proof}
En efecto, si $p_1$, $q_1$ son las primeras proyecciones $X\times X\to X$, $Y\times Y\to Y$, entonces la proyección sobre el primer factor $(X\times Y)\times(X\times Y)\to X\times Y$ se identifica con $p_1\times q_1$, y se tiene
\[
  (p_1\times q_1)\circ(\Delta_X\times\Delta_Y)
  =(p_1\circ\Delta_X)\times(q_1\circ\Delta_Y)=1_{X\times Y}
\]
y se razona del mismo modo para las proyecciones sobre el segundo factor.
\end{proof}

\phantomsection
\label{I.5.3.3}

\begin{corollary}[5.3.4]
\label{I.5.3.4}
Para toda extensión de base $S'\to S$, $\Delta_{X_{(S')}}$ se identifica canónicamente con $(\Delta_X)_{(S')}$.
\end{corollary}

\begin{proof}
Basta observar que $(X\times_S X)_{(S')}$ se identifica canónicamente con $X_{(S')}\times_{S'}X_{(S')}$ \sref{I.3.3.10}.
\end{proof}

\begin{proposition}[5.3.5]
\label{I.5.3.5}
Sean $X$ e $Y$ dos $S$-preesquemas, y sea $\varphi:S\to T$ un morfismo de preesquemas que permite considerar todo $S$-preesquema como un $T$-preesquema.
Sean $f:X\to S$ y $g:Y\to S$\footnote{\emph{[Trad.] La fuente francesa imprime solamente $Y\to S$ aquí, omitiendo el nombre $g$; la identidad siguiente $\pi=f\circ p=g\circ q$ identifica este segundo morfismo estructural como $g$.}} los morfismos estructurales, $p$ y $q$ las proyecciones de $X\times_S Y$, y $\pi=f\circ p=g\circ q$ el morfismo estructural $X\times_S Y\to S$.
Entonces el diagrama
\[
  \xymatrix{
    X\times_S Y\ar[r]^{(p,q)_T}\ar[d]_\pi &
    X\times_T Y\ar[d]^{f\times_T g}\\
    S\ar[r]_{\Delta_{S|T}} &
    S\times_T S
  }
  \tag{5.3.5.1}\label{I.5.3.5.1-fr}
\]
\oldpage[I]{133}
es conmutativo e identifica $X\times_S Y$ con el producto de los $(S\times_T S)$-preesquemas $S$ y $X\times_T Y$, identificándose las proyecciones con $\pi$ y $(p,q)_T$.
\end{proposition}

\begin{proof}
En virtud de \sref{I.3.4.3}, se reduce a demostrar la proposición correspondiente \emph{en la categoría de conjuntos}, sustituyendo $X$, $Y$ y $S$ por $X(Z)_T$, $Y(Z)_T$ y $S(Z)_T$, respectivamente, siendo $Z$ un $T$-preesquema arbitrario.
Pero, en la categoría de conjuntos, la verificación es inmediata y se deja al lector.
\end{proof}

\begin{corollary}[5.3.6]
\label{I.5.3.6}
El morfismo $(p,q)_T$ se identifica (poniendo $P=S\times_T S$) con $1_{X\times_T Y}\times_P\Delta_S$.
\end{corollary}

\begin{proof}
Esto resulta de \sref{I.5.3.5} y \sref{I.3.3.4}.
\end{proof}

\begin{corollary}[5.3.7]
\label{I.5.3.7}
Si $f:X\to Y$ es un $S$-morfismo, entonces el diagrama
\[
  \xymatrix{
    X\ar[r]^{(1_X,f)_S}\ar[d]_f &
    X\times_S Y\ar[d]^{f\times_S 1_Y}\\
    Y\ar[r]_{\Delta_Y} &
    Y\times_S Y
  }
\]
es conmutativo e identifica $X$ con el producto de los $(Y\times_S Y)$-preesquemas $Y$ y $X\times_S Y$.
\end{corollary}

\begin{proof}
Basta aplicar \sref{I.5.3.5}, sustituyendo $S$ por $Y$ y $T$ por $S$, y observando que $X\times_Y Y=X$ \sref{I.3.3.3}.
\end{proof}

\begin{proposition}[5.3.8]
\label{I.5.3.8}
Para que $f:X\to Y$ sea un monomorfismo de preesquemas, es necesario y suficiente que $\Delta_{X|Y}$ sea un isomorfismo de $X$ sobre $X\times_Y X$.
\end{proposition}

\begin{proof}
En efecto, $f$ es un monomorfismo si y solo si, para todo $Y$-preesquema $Z$, la aplicación correspondiente $f':X(Z)_Y\to Y(Z)_Y$ es inyectiva.
Como $Y(Z)_Y$ consta de un solo elemento, esto equivale a decir que $X(Z)_Y$ tiene a lo sumo un elemento\footnote{\emph{[Trad.] La fuente francesa afirma que $X(Z)_Y$ se reduce igualmente a un solo elemento; la inyectividad en el conjunto unitario $Y(Z)_Y$ solo implica que tiene a lo sumo un elemento, admitiendo el caso vacío.}}, o, equivalentemente, que la aplicación diagonal de $X(Z)_Y$ en $X(Z)_Y\times X(Z)_Y$ es una biyección.
Este último producto es precisamente el conjunto $(X\times_Y X)(Z)_Y$ \hyperref[I.3.4.3.1-fr]{(3.4.3.1)}, de modo que esto equivale a que $\Delta_{X|Y}$ sea un isomorfismo.
\end{proof}

\begin{proposition}[5.3.9]
\label{I.5.3.9}
El morfismo diagonal $\Delta_X$ es una inmersión de $X$ en $X\times_S X$.
\end{proposition}

\begin{proof}
Para tener una demostración correcta, basta utilizar \sref{I.4.2.4}[$a$]: para todo $x\in X$, escojamos una vecindad abierta afín $V\subset S$ de su imagen y una vecindad abierta afín $U$ de $x$ en $X$ cuya imagen esté contenida en $V$. Entonces $U\times_SU=U\times_VU$ es una vecindad abierta afín de $\Delta_X(x)$;\footnote{\emph{[Trad.] La formulación publicada en EGA~III.2, Erratas y adiciones (lista~2), elemento $\mathrm{Err}_{\mathrm{III}},10$, afirma esto para toda vecindad afín $U$ sin exigir que su morfismo estructural se factorice por un abierto afín de $S$; la hipótesis aquí explicitada es necesaria para este paso.}}
teniendo en cuenta \sref{I.5.3.16} (cuya demostración solo utiliza la definición \sref{I.5.3.1.1}), se reduce a demostrar \sref{I.5.3.9} cuando $S=\operatorname{Spec}(B)$ y $X=\operatorname{Spec}(A)$ son esquemas afines; es claro entonces, en virtud de \sref{I.5.3.1.1}, que $\Delta_X$ corresponde al homomorfismo canónico $A\otimes_BA\to A$ que transforma $x\otimes y$ en $xy$; siendo este homomorfismo sobreyectivo, $\Delta_X$ es en este caso una inmersión cerrada \sref{I.4.2.3}, lo que termina la demostración.\footnote{La demostración impresa en la fuente francesa es insuficiente, pues no
demuestra que $\Delta_X(X)$ sea localmente cerrado en $X\times_S X$.}\footnote{\emph{[Trad.] EGA~III.2, Erratas y adiciones (lista~2), elemento $\mathrm{Err}_{\mathrm{III}},10$, identifica explícitamente esta laguna y proporciona la demostración afín-local traducida aquí.}}
\end{proof}

Se dice que el subpreesquema de $X\times_S X$ asociado a la inmersión $\Delta_X$ \sref{I.4.2.1} es \emph{la diagonal} de $X\times_S X$.

\begin{corollary}[5.3.10]
\label{I.5.3.10}
Bajo las hipótesis de \sref{I.5.3.5}, $(p,q)_T$ es una inmersión.
\end{corollary}

\begin{proof}
Esto resulta de \sref{I.5.3.6} y \sref{I.4.3.1}.
\end{proof}

Se dice (bajo las hipótesis de \sref{I.5.3.5}) que $(p,q)_T$ es la \emph{inmersión canónica} de $X\times_S Y$ en $X\times_T Y$.

\begin{corollary}[5.3.11]
\label{I.5.3.11}
Sean $X$ e $Y$ dos $S$-preesquemas, y sea $f:X\to Y$ un $S$-morfismo;
entonces el morfismo grafo $\Gamma_f=(1_X,f)_S$ de $f$ \sref{I.3.3.14} es una inmersión de $X$ en $X\times_S Y$.
\end{corollary}

\begin{proof}
Este es un caso particular del Corolario \sref{I.5.3.10}, donde se sustituye $S$ por $Y$ y $T$ por $S$ (cf. \sref{I.5.3.7}).
\end{proof}

\oldpage[I]{134}
El subpreesquema de $X\times_S Y$ asociado a la inmersión $\Gamma_f$ \sref{I.4.2.1} se llama \emph{el grafo} del morfismo $f$;
los subpreesquemas de $X\times_S Y$ que son grafos de morfismos $X\to Y$ se caracterizan por la propiedad de que la restricción a tal subpreesquema $G$ de la proyección $p_1:X\times_S Y\to X$ es un \emph{isomorfismo} $g$ de $G$ sobre $X$:
$G$ es el grafo del morfismo $p_2\circ g^{-1}$, donde $p_2$ es la proyección $X\times_S Y\to Y$.

Tomando, en particular, $X=S$, los $S$-morfismos $S\to Y$ (que no son sino las $S$-\emph{secciones} de $Y$ \sref{I.2.5.5}) son iguales a sus morfismos grafo;
los subpreesquemas de $Y$ que son grafos de $S$-secciones (en otras palabras, los que son isomorfos a $S$ por la restricción del morfismo estructural $Y\to S$) se llaman asimismo \emph{imágenes} de estas secciones o, por abuso de lenguaje, las \emph{$S$-secciones} de $Y$.

\begin{corollary}[5.3.12]
\label{I.5.3.12}
Con las hipótesis y la notación de \sref{I.5.3.11}, para todo morfismo $g:S'\to S$, sea $f'$ la imagen inversa de $f$ por $g$ \sref{I.3.3.7};
entonces $\Gamma_{f'}$ es la imagen inversa de $\Gamma_f$ por $g$.
\end{corollary}

\begin{proof}
Este es un caso particular de \sref{I.3.3.10.1}.
\end{proof}

\begin{corollary}[5.3.13]
\label{I.5.3.13}
Sean $f:X\to Y$ y $g:Y\to Z$ dos morfismos;
si $g\circ f$ es una inmersión (\textup{resp.} una inmersión local), entonces $f$ también lo es.
\end{corollary}

\begin{proof}
En efecto, $f$ se factoriza como $X\xrightarrow{\Gamma_f}X\times_Z Y\xrightarrow{p_2}Y$.
Además, $p_2$ puede identificarse con $(g\circ f)\times_Z 1_Y$ \sref{I.3.3.4};
si $g\circ f$ es una inmersión (resp. una inmersión local), entonces $p_2$ también lo es (\sref{I.4.3.1} y \sref{I.4.5.5}), y, como $\Gamma_f$ es una inmersión \sref{I.5.3.11}, se concluye, por \sref{I.4.2.5}\footnote{\emph{[Trad.] La fuente francesa imprime aquí \sref{I.4.2.4}; EGA~III.2, Erratas y adiciones (lista~2), prescribe explícitamente su sustitución por \sref{I.4.2.5}.}} (resp. \sref{I.4.5.5}).
\end{proof}

\begin{corollary}[5.3.14]
\label{I.5.3.14}
Sean $j:X\to Y$ y $g:X\to Z$ dos $S$-morfismos.
Si $j$ es una inmersión (\textup{resp.} una inmersión local), entonces $(j,g)_S$ también lo es.
\end{corollary}

\begin{proof}
En efecto, si $p:Y\times_S Z\to Y$ es la proyección sobre el primer factor, entonces se tiene $j=p\circ(j,g)_S$, y basta aplicar \sref{I.5.3.13}.
\end{proof}

\begin{proposition}[5.3.15]
\label{I.5.3.15}
Si $f:X\to Y$ es un $S$-morfismo, entonces el diagrama
\[
  \xymatrix{
    X\ar[r]^{\Delta_X}\ar[d]_f &
    X\times_S X\ar[d]^{f\times_S f}\\
    Y\ar[r]_{\Delta_Y} &
    Y\times_S Y
  }
  \tag{5.3.15.1}\label{I.5.3.15.1-fr}
\]
es conmutativo \textup{(dicho de otro modo, $\Delta_X$ es un morfismo funtorial en la categoría de los preesquemas)}.
\end{proposition}

\begin{proof}
La verificación es inmediata y se deja al lector.
\end{proof}

\begin{corollary}[5.3.16]
\label{I.5.3.16}
Si $X$ es un subpreesquema de $Y$, entonces la diagonal $\Delta_X(X)$ se identifica con un subpreesquema de $\Delta_Y(Y)$, y su espacio subyacente se identifica con
\[
  \Delta_Y(Y)\cap p_1^{-1}(X)=\Delta_Y(Y)\cap p_2^{-1}(X)
\]
($p_1$ y $p_2$ son las proyecciones de $Y\times_S Y$).
\end{corollary}

\begin{proof}
Apliquemos \sref{I.5.3.15} al morfismo de inyección $f:X\to Y$; se sabe entonces que $f\times_S f$ es una inmersión que identifica el espacio subyacente de $X\times_S X$ con el subespacio $p_1^{-1}(X)\cap p_2^{-1}(X)$ de $Y\times_S Y$ \sref{I.4.3.1};
además, si $z\in\Delta_Y(Y)\cap p_1^{-1}(X)$, entonces se tiene $z=\Delta_Y(y)$
\oldpage[I]{135}
y $y=p_1(z)\in X$, por lo que $y=f(y)$, y $z=\Delta_Y(f(y))$ pertenece a $\Delta_X(X)$ en virtud de la conmutatividad de \sref{I.5.3.15.1}.
\end{proof}

\begin{corollary}[5.3.17]
\label{I.5.3.17}
Sean $f_1:Y\to X$ y $f_2:Y\to X$ dos $S$-morfismos, y sea $y$ un punto de $Y$ tal que $f_1(y)=f_2(y)=x$ y que los homomorfismos $\boldsymbol{k}(x)\to\boldsymbol{k}(y)$ correspondientes a $f_1$ y $f_2$ sean idénticos.
Entonces, si $f=(f_1,f_2)_S$, el punto $f(y)$ pertenece a la diagonal $\Delta_{X|S}(X)$.
\end{corollary}

\begin{proof}
Los dos homomorfismos $\boldsymbol{k}(x)\to\boldsymbol{k}(y)$ correspondientes a $f_i$ ($i=1,2$) definen dos $S$-morfismos $g_i:\operatorname{Spec}(\boldsymbol{k}(y))\to\operatorname{Spec}(\boldsymbol{k}(x))$ tales que los diagramas
\[
  \xymatrix{
    \operatorname{Spec}(\boldsymbol{k}(y))\ar[r]^{g_i}\ar[d] &
    \operatorname{Spec}(\boldsymbol{k}(x))\ar[d]\\
    Y\ar[r]_{f_i} & X
  }
\]
son conmutativos.
El diagrama
\[
  \xymatrix{
    \operatorname{Spec}(\boldsymbol{k}(y))\ar[rr]^{(g_1,g_2)_S}\ar[d] & &
    \operatorname{Spec}(\boldsymbol{k}(x))\times_S\operatorname{Spec}(\boldsymbol{k}(x))\ar[d]\\
    Y\ar[rr]_{(f_1,f_2)_S} & & X\times_S X
  }
\]
es, por tanto, también conmutativo.
Pero de la igualdad $g_1=g_2$ resulta que la imagen por $(g_1,g_2)_S$ del único punto de $\operatorname{Spec}(\boldsymbol{k}(y))$ pertenece a la diagonal de $\operatorname{Spec}(\boldsymbol{k}(x))\times_S\operatorname{Spec}(\boldsymbol{k}(x))$;
la conclusión resulta entonces de \sref{I.5.3.15}.
\end{proof}

\subsection{Morfismos separados y preesquemas separados}
\label{subsection:I.5.4}
\label{I.5.4}

\begin{definition}[5.4.1]
\label{I.5.4.1}
{\itshape Se dice que un morfismo de preesquemas $f:X\to Y$ es separado si el morfismo diagonal $X\to X\times_Y X$ es una inmersión \emph{cerrada}; se dice también entonces que $X$ es un preesquema separado sobre $Y$, o un $Y$-esquema.
Se dice que un preesquema $X$ es separado si es separado sobre $\operatorname{Spec}(\mathbf{Z})$;
se dice también entonces que $X$ es un esquema\footnote{\emph{[Trad.] Repetimos aquí la advertencia formulada al comienzo mismo de esta traducción: las primeras versiones de EGA emplean \emph{preesquema} para lo que hoy suele llamarse esquema, y \emph{esquema} para lo que hoy suele llamarse esquema separado.} El propio Grothendieck dijo más tarde que era preferible la terminología moderna, pero hemos decidido mantener esta traducción «históricamente fiel» usando la nomenclatura antigua.} \textup{(cf. \sref{I.5.5.7})}.}
\end{definition}

En virtud de \sref{I.5.3.9}, para que $X$ sea separado sobre $Y$, es necesario y suficiente que $\Delta_X(X)$ sea un \emph{subespacio cerrado} del espacio subyacente a $X\times_Y X$.

\begin{proposition}[5.4.2]
\label{I.5.4.2}
Sea $S\to T$ un morfismo separado.
Si $X$ e $Y$ son dos $S$-preesquemas, entonces la inmersión canónica $X\times_S Y\to X\times_T Y$ \sref{I.5.3.10} es cerrada.
\end{proposition}

\begin{proof}
En efecto, remitiéndonos al diagrama de \sref{I.5.3.5.1}, se ve que $(p,q)_T$ puede considerarse obtenido de $\Delta_{S|T}$ por la extensión $f\times_T g:X\times_T Y\to S\times_T S$ del preesquema de base $S\times_T S$;
la proposición resulta entonces de \sref{I.4.3.2}.
\end{proof}

\begin{corollary}[5.4.3]
\label{I.5.4.3}
Sea $Y$ un $S$-esquema, y sea $f:X\to Y$ un $S$-morfismo.
Entonces el morfismo grafo $\Gamma_f:X\to X\times_S Y$ \sref{I.5.3.11} es una inmersión cerrada.
\end{corollary}

\begin{proof}
Este es el caso particular de \sref{I.5.4.2} en el que se sustituye $S$ por $Y$ y $T$ por $S$.
\end{proof}

\begin{corollary}[5.4.4]
\label{I.5.4.4}
Sean $f:X\to Y$ y $g:Y\to Z$ dos morfismos, siendo $g$ separado.
Si $g\circ f$ es una inmersión cerrada, entonces $f$ también lo es.
\end{corollary}

\begin{proof}
La demostración a partir de \sref{I.5.4.3} es la misma que la de \sref{I.5.3.13} a partir de \sref{I.5.3.11}.
\end{proof}

\oldpage[I]{136}
\begin{corollary}[5.4.5]
\label{I.5.4.5}
Sea $Z$ un $S$-esquema, y sean $j:X\to Y$ y $g:X\to Z$ dos $S$-morfismos.
Si $j$ es una inmersión cerrada, entonces $(j,g)_S:X\to Y\times_S Z$ también lo es.
\end{corollary}

\begin{proof}
La demostración a partir de \sref{I.5.4.4} es la misma que la de \sref{I.5.3.14} a partir de \sref{I.5.3.13}.
\end{proof}

\begin{corollary}[5.4.6]
\label{I.5.4.6}
Si $X$ es un $S$-esquema, entonces toda $S$-sección de $X$ \sref{I.2.5.5} es una inmersión cerrada.
\end{corollary}

\begin{proof}
Si $\varphi:X\to S$ es el morfismo estructural, y $\psi:S\to X$ una $S$-sección de $X$, basta aplicar \sref{I.5.4.5} a $\varphi\circ\psi=1_S$.
\end{proof}

\begin{corollary}[5.4.7]
\label{I.5.4.7}
Sea $S$ un preesquema íntegro de punto genérico $s$, y sea $X$ un $S$-esquema.
Si dos $S$-secciones $f$, $g$ de $X$ son tales que $f(s)=g(s)$, entonces $f=g$.
\end{corollary}

\begin{proof}
En efecto, si $x=f(s)=g(s)$, entonces los homomorfismos $\boldsymbol{k}(x)\to\boldsymbol{k}(s)$ correspondientes a $f$ y $g$ son necesariamente idénticos.
Si $h=(f,g)_S$, se deduce, pues, de \sref{I.5.3.17} que $h(s)$ pertenece a la diagonal $Z=\Delta_X(X)$;
pero como $S=\overline{\{s\}}$, y como $Z$ es cerrado por hipótesis, se tiene $h(S)\subset Z$.
Resulta entonces de \sref{I.5.2.2} que $h$ se factoriza como $S\to Z\to X\times_S X$, y se concluye así que $f=g$, por definición de la diagonal.
\end{proof}

\begin{remark}[5.4.8]
\label{I.5.4.8}
Si se supone, recíprocamente, que la conclusión de \sref{I.5.4.3} es válida cuando $f=1_Y$, entonces se concluye que $Y$ es separado sobre $S$;
del mismo modo, si se supone que la conclusión de \sref{I.5.4.4}\footnote{\emph{[Trad.] La fuente francesa remite aquí a (5.4.5), que concierne a un morfismo par; el argumento con $p_1\circ\Delta_Y=1_Y$ utiliza el criterio de composición del Corolario~\hyperref[I.5.4.4]{5.4.4}.}} se aplica a los dos morfismos \[
  Y\xrightarrow{\Delta_Y}Y\times_Z Y\xrightarrow{p_1}Y,
\], entonces se concluye que $\Delta_Y$ es una inmersión cerrada y, por tanto, que $Y$ es separado sobre $Z$;
por último, si se supone que la conclusión de \sref{I.5.4.6} es válida para la $Y$-sección $\Delta_Y$ del $Y$-preesquema $Y\times_S Y\to Y$, entonces esto implica que $Y$ es separado sobre $S$.
\end{remark}

\subsection{Criterios de separación}
\label{subsection:I.5.5}

\begin{proposition}[5.5.1]
\label{I.5.5.1}
\medskip\noindent
\begin{enumerate}
  \item[{\rm(i)}] Todo monomorfismo de preesquemas (y, en particular, toda inmersión) es un morfismo separado.
  \item[{\rm(ii)}] El compuesto de dos morfismos separados es separado.
  \item[{\rm(iii)}] Si $f:X\to X'$ y $g:Y\to Y'$ son $S$-morfismos separados, entonces $f\times_S g$ es separado.
  \item[{\rm(iv)}] Si $f:X\to Y$ es un $S$-morfismo separado, entonces el $S'$-morfismo $f_{(S')}$ es separado para toda extensión $S'\to S$ del preesquema de base.
  \item[{\rm(v)}] Si el compuesto $g\circ f$ es separado, entonces $f$ es separado.
  \item[{\rm(vi)}] Para que un morfismo $f$ sea separado, es necesario y suficiente que $f_{\mathrm{red}}$ \sref{I.5.1.5} lo sea.
\end{enumerate}
\end{proposition}

\begin{proof}
La afirmación~(i) resulta inmediatamente de \sref{I.5.3.8}.
Si $f:X\to Y$ y $g:Y\to Z$ son morfismos, entonces el diagrama
\[
  \xymatrix{
    X\ar[rr]^{\Delta_{X|Z}}\ar[dr]_{\Delta_{X|Y}} & &
    X\times_Z X\\
    & X\times_Y X\ar[ur]_j
  }
  \tag{5.5.1.1}\label{I.5.5.1.1-fr}
\]
es conmutativo (donde $j$ denota la inmersión canónica \sref{I.5.3.10}), como se comprueba inmediatamente.
Si $f$ y $g$ son separados, entonces $\Delta_{X|Y}$ es una inmersión cerrada por definición, y $j$ es una inmersión cerrada por \sref{I.5.4.2}, de donde $\Delta_{X|Z}$ es una inmersión cerrada por \sref{I.4.2.5}\footnote{\emph{[Trad.] La fuente francesa remite aquí a (4.2.4), que es el criterio local; la transitividad de las inmersiones cerradas utilizada en este paso es la Proposición~\hyperref[I.4.2.5]{4.2.5}.}}, lo que
\oldpage[I]{137}
demuestra~(ii).
En virtud de~(i) y~(ii), (iii) y~(iv) son equivalentes \sref{I.3.5.1}, de modo que basta demostrar (iv).
Ahora bien, $X_{(S')}\times_{Y_{(S')}}X_{(S')}$ se identifica canónicamente con $(X\times_Y X)\times_Y Y_{(S')}$ por \sref{I.3.3.11} y \hyperref[I.3.3.9.1]{(3.3.9.1)}, y se ve inmediatamente que el morfismo diagonal $\Delta_{X_{(S')}}$ puede identificarse entonces con $\Delta_X\times_Y 1_{Y_{(S')}}$;
la proposición se sigue entonces de \sref{I.4.3.1}.

Para demostrar (v), consideremos, como en \sref{I.5.3.13}, la factorización $X\xrightarrow{\Gamma_f}X\times_Z Y\xrightarrow{p_2}Y$ de $f$, observando que $p_2=(g\circ f)\times_Z 1_Y$;
la hipótesis (que $g\circ f$ es separado) implica que $p_2$ es separado, por (iii) y (i), y, como $\Gamma_f$ es una inmersión, $\Gamma_f$ es separado, por (i), de donde $f$ es separado, por (ii).
Finalmente, para demostrar (vi), recordemos que los preesquemas $X_{\mathrm{red}}\times_{Y_{\mathrm{red}}}X_{\mathrm{red}}$ y $X_{\mathrm{red}}\times_Y X_{\mathrm{red}}$ se identifican canónicamente entre sí \sref{I.5.1.7};
si se denota por $j$ la inyección $X_{\mathrm{red}}\to X$, entonces el diagrama
\[
  \xymatrix{
    X_{\mathrm{red}}\ar[r]^{\Delta_{X_{\mathrm{red}}}}\ar[d]_j &
    X_{\mathrm{red}}\times_Y X_{\mathrm{red}}\ar[d]^{j\times_Y j}\\
    X\ar[r]_{\Delta_X} & X\times_Y X
  }
\]
es conmutativo \sref{I.5.3.15}, y la proposición se sigue del hecho de que las flechas verticales son homeomorfismos de los espacios subyacentes \sref{I.4.3.1}.
\end{proof}

\begin{corollary}[5.5.2]
\label{I.5.5.2}
Si $f:X\to Y$ es separado, entonces la restricción de $f$ a todo subpreesquema de $X$ es separada.
\end{corollary}

\begin{proof}
Esto se deduce de \sref{I.5.5.1}[i y ii].
\end{proof}

\begin{corollary}[5.5.3]
\label{I.5.5.3}
Si $X$ e $Y$ son $S$-preesquemas tales que $Y$ sea separado sobre $S$, entonces $X\times_S Y$ es separado sobre $X$.
\end{corollary}

\begin{proof}
Este es un caso particular de \sref{I.5.5.1}[iv].
\end{proof}

\begin{proposition}[5.5.4]
\label{I.5.5.4}
Sea $X$ un preesquema, y supongamos que su espacio subyacente es la unión de una familia \emph{finita} de partes cerradas $X_k$ ($(1\leq k\leq n)$);
para cada $k$, consideremos el subpreesquema reducido de $X$ que tiene a $X_k$ como espacio subyacente \sref{I.5.2.1}, y denotémoslo también por $X_k$.
Sea $f:X\to Y$ un morfismo, y, para cada $k$, sea $Y_k$ una parte cerrada de $Y$ tal que $f(X_k)\subset Y_k$;
se denota asimismo por $Y_k$ el subpreesquema reducido de $Y$ que tiene a $Y_k$ como espacio subyacente, de modo que la restricción $X_k\to Y$ de $f$ a $X_k$ se factoriza como $X_k\xrightarrow{f_k}Y_k\to Y$ \sref{I.5.2.2}.
Para que $f$ sea separado, es necesario y suficiente que todos los $f_k$ sean separados.
\end{proposition}

\begin{proof}
La necesidad se deduce de \sref{I.5.5.1}[i, ii y v].
Recíprocamente, si se satisface la condición del enunciado, entonces cada una de las restricciones $X_k\to Y$ de $f$ es separada \hyperref[I.5.5.1]{(5.5.1, (i) y (ii))};
si $p_1$ y $p_2$ son las proyecciones de $X\times_Y X$, entonces el subespacio $\Delta_{X_k}(X_k)$ se identifica con el subespacio $\Delta_X(X)\cap p_1^{-1}(X_k)$ del espacio subyacente de $X\times_Y X$ \sref{I.5.3.16};
estos subespacios son cerrados en $X\times_Y X$, y por tanto también lo es su unión $\Delta_X(X)$.
\end{proof}

Supongamos, en particular, que los $X_k$ son las \emph{componentes irreducibles} de $X$;
entonces se puede suponer que los $Y_k$ son componentes irreducibles de $Y$ \sref[0]{0.2.1.5};
por tanto, la Proposición~\sref{I.5.5.4} reduce, en este caso, la noción de separación al caso de los preesquemas \emph{íntegros} \sref{I.2.1.8}\footnote{\emph{[Trad.] La fuente francesa remite aquí a (2.1.7), que trata de la restricción a abiertos; la definición de preesquema íntegro es (2.1.8).}}.

\oldpage[I]{138}
\begin{proposition}[5.5.5]
\label{I.5.5.5}
Sea $(Y_\lambda)$ un recubrimiento abierto de un preesquema $Y$;
para que un morfismo $f:X\to Y$ sea separado, es necesario y suficiente que cada una de sus restricciones $f^{-1}(Y_\lambda)\to Y_\lambda$ sea separada.
\end{proposition}

\begin{proof}
Si se pone $X_\lambda=f^{-1}(Y_\lambda)$, todo se reduce, teniendo en cuenta \sref{I.4.2.4}[\textit{b)}] y la identificación de los productos $X_\lambda\times_Y X_\lambda$ y $X_\lambda\times_{Y_\lambda}X_\lambda$ \sref{I.3.2.5}, a demostrar que los $X_\lambda\times_Y X_\lambda$ forman un recubrimiento de $X\times_Y X$.
Si ahora se pone $Y_{\lambda\mu}=Y_\lambda\cap Y_\mu$ y $X_{\lambda\mu}=X_\lambda\cap X_\mu=f^{-1}(Y_{\lambda\mu})$, entonces $X_\lambda\times_Y X_\mu$ se identifica con el producto $X_{\lambda\mu}\times_{Y_{\lambda\mu}}X_{\lambda\mu}$ \sref{I.3.2.6.4}, por tanto también con $X_{\lambda\mu}\times_Y X_{\lambda\mu}$ \sref{I.3.2.5}, y finalmente con un abierto de $X_\lambda\times_Y X_\lambda$, lo que establece nuestra afirmación \sref{I.3.2.7}.
\end{proof}

La Proposición~\sref{I.5.5.5}\footnote{\emph{[Trad.] La fuente francesa cita aquí (5.5.4); la reducción mediante un recubrimiento abierto de $Y$ utiliza la Proposición~\hyperref[I.5.5.5]{5.5.5}.}} permite, tomando un recubrimiento de $Y$ por abiertos afines, reducir el estudio de los morfismos separados al de los morfismos separados con valores en esquemas afines.

\begin{proposition}[5.5.6]
\label{I.5.5.6}
Sean $Y$ un esquema \emph{afín}, $X$ un preesquema y $(U_\alpha)$ un recubrimiento de $X$ por abiertos afines.
Para que un morfismo $f:X\to Y$ sea separado, es necesario y suficiente que $U_\alpha\cap U_\beta$ sea, para todo par de índices $(\alpha,\beta)$, un abierto afín, \emph{y} que el anillo $\Gamma(U_\alpha\cap U_\beta,\mathscr{O}_X)$ esté engendrado por la reunión de las imágenes canónicas de los anillos $\Gamma(U_\alpha,\mathscr{O}_X)$ y $\Gamma(U_\beta,\mathscr{O}_X)$.
\end{proposition}

\begin{proof}
Los $U_\alpha\times_Y U_\beta$ forman un recubrimiento abierto de $X\times_Y X$ \sref{I.3.2.7};
denotando las proyecciones de $X\times_Y X$ por $p$ y $q$, se tiene
\[
  \begin{aligned}
    \Delta_X^{-1}(U_\alpha\times_Y U_\beta)
      &=\Delta_X^{-1}\bigl(p^{-1}(U_\alpha)\cap q^{-1}(U_\beta)\bigr)\\
      &=\Delta_X^{-1}\bigl(p^{-1}(U_\alpha)\bigr)
        \cap\Delta_X^{-1}\bigl(q^{-1}(U_\beta)\bigr)
        =U_\alpha\cap U_\beta~;
  \end{aligned}
\]
de modo que todo se reduce a demostrar que la restricción de $\Delta_X$ a $U_\alpha\cap U_\beta$ es una inmersión cerrada en $U_\alpha\times_Y U_\beta$.
Ahora bien, esta restricción no es otra que $(j_\alpha,j_\beta)_Y$, donde $j_\alpha$ (resp. $j_\beta$) denota el morfismo de inyección de $U_\alpha\cap U_\beta$ en $U_\alpha$ (resp. $U_\beta$), como se sigue de las definiciones.
Como $U_\alpha\times_Y U_\beta$ es un esquema afín cuyo anillo es canónicamente isomorfo a $\Gamma(U_\alpha,\mathscr{O}_X)
 \otimes_{\Gamma(Y,\mathscr{O}_Y)}
 \Gamma(U_\beta,\mathscr{O}_X)$ \sref{I.3.2.2}, se ve que $U_\alpha\cap U_\beta$ debe ser un esquema afín y que la aplicación $h_\alpha\otimes h_\beta\to h_\alpha h_\beta$ del anillo $A(U_\alpha\times_Y U_\beta)$ en $\Gamma(U_\alpha\cap U_\beta,\mathscr{O}_X)$ debe ser sobreyectiva \sref{I.4.2.3}, lo que concluye la demostración.
\end{proof}

\begin{corollary}[5.5.7]
\label{I.5.5.7}
Un esquema afín es separado \textup{(y, por tanto, es un \emph{esquema}, lo que justifica la terminología de \sref{I.5.4.1})}.
\end{corollary}

\begin{corollary}[5.5.8]
\label{I.5.5.8}
Sea $Y$ un esquema afín;
para que $f:X\to Y$ sea un morfismo separado, es necesario y suficiente que $X$ sea separado \textup{(dicho de otro modo, que $X$ sea un \emph{esquema})}.
\end{corollary}

\begin{proof}
En efecto, se observa que el criterio de \sref{I.5.5.6} no depende de $f$.
\end{proof}

\begin{corollary}[5.5.9]
\label{I.5.5.9}
Para que un morfismo $f:X\to Y$ sea separado, es necesario que el preesquema inducido $f^{-1}(U)$ sea separado para todo abierto $U$ sobre el que $Y$ induce un preesquema separado, y basta que esto se cumpla para todo abierto afín $U\subset Y$.
\end{corollary}

\begin{proof}
La necesidad de la condición se deduce de \sref{I.5.5.5}\footnote{\emph{[Trad.] En las dos referencias a (5.5.5) de esta demostración, la fuente francesa imprime (5.5.4); el argumento aplica el carácter local sobre la base establecido en la Proposición~\hyperref[I.5.5.5]{5.5.5}.}} y \sref{I.5.5.1}[ii];
la suficiencia se deduce de \sref{I.5.5.5} y \sref{I.5.5.8}, teniendo en cuenta la existencia de recubrimientos de $Y$ por abiertos afines.
\end{proof}

En particular, si $X$ e $Y$ son esquemas afines, entonces \emph{todo} morfismo $X\to Y$ es separado.

\begin{proposition}[5.5.10]
\label{I.5.5.10}
Sea $Y$ un \emph{esquema}, y sea $f:X\to Y$ un morfismo.
Para todo abierto afín $U$ de $X$ y todo abierto afín $V$ de $Y$, $U\cap f^{-1}(V)$ es afín.
\end{proposition}

\begin{proof}
Sean $p_1$ y $p_2$ las proyecciones de $X\times_{\mathbf{Z}}Y$;
el subespacio $U\cap f^{-1}(V)$ es la imagen de $\Gamma_f(X)\cap p_1^{-1}(U)\cap p_2^{-1}(V)$ por $p_1$.
Pero $p_1^{-1}(U)\cap p_2^{-1}(V)$ se identifica con el espacio subyacente al
\oldpage[I]{139}
preesquema $U\times_{\mathbf{Z}}V$ \sref{I.3.2.7}, y es, por tanto, un esquema afín \sref{I.3.2.2};
como $\Gamma_f(X)$ es cerrado en $X\times_{\mathbf{Z}}Y$ \sref{I.5.4.3}, $\Gamma_f(X)\cap p_1^{-1}(U)\cap p_2^{-1}(V)$ es cerrado en $U\times_{\mathbf{Z}}V$, de modo que el preesquema inducido por el subpreesquema de $X\times_{\mathbf{Z}}Y$ asociado a $\Gamma_f$ \sref{I.4.2.1} sobre el abierto $\Gamma_f(X)\cap p_1^{-1}(U)\cap p_2^{-1}(V)$ de su espacio subyacente es un subpreesquema cerrado de un esquema afín y, por tanto, un esquema afín \sref{I.4.2.3}.
La proposición resulta entonces del hecho de que $\Gamma_f$ es una inmersión.
\end{proof}

\begin{examples}[5.5.11]
\label{I.5.5.11}
El preesquema del Ejemplo~\sref{I.2.3.2} (``la recta proyectiva sobre un cuerpo $K$'') es \emph{separado}, pues, para el recubrimiento $(X_1,X_2)$ de $X$ por abiertos afines, $X_1\cap X_2=U_{12}$ es afín y $\Gamma(U_{12},\mathscr{O}_X)$, el anillo de las fracciones racionales de la forma $f(s)/s^m$, con $f\in K[s]$, está generado por $K[s]$ y $1/s$, por lo que se satisfacen las condiciones de \sref{I.5.5.6}.

Con la misma elección de $X_1$, $X_2$, $U_{12}$ y $U_{21}$ que en el Ejemplo~\sref{I.2.3.2}, tomemos ahora como $u_{12}$ el isomorfismo que envía $f(s)$ a $f(t)$;
se obtiene entonces, por pegado, un preesquema \emph{íntegro no separado} $X$, pues se satisface la primera condición de \sref{I.5.5.6}, pero no la segunda.
Es inmediato aquí que $\Gamma(X,\mathscr{O}_X)\to\Gamma(X_1,\mathscr{O}_X)=K[s]$ es un isomorfismo;
el isomorfismo inverso define un morfismo sobreyectivo $f:X\to\operatorname{Spec}(K[s])$, y, para todo $y\in\operatorname{Spec}(K[s])$ tal que $\mathfrak{j}_y\neq(0)$, $f^{-1}(y)$ consta de un solo punto, mientras que, para $\mathfrak{j}_y=(0)$, $f^{-1}(y)$ consta de \emph{dos} puntos distintos\footnote{\emph{[Trad.] La fuente francesa imprime $(0)$ en ambas condiciones sobre el ideal. Para el origen duplicado de $\Spec(K[s])$, el ideal primo pertinente es $(s)$; $(0)$ es el punto genérico.}} (se dice que $X$ es la ``recta afín sobre $K$ con el punto $0$ duplicado'').

También se puede dar un ejemplo en el que no se satisface la primera de las dos condiciones de \sref{I.5.5.6}.
Observemos primero que, en el espectro primo $Y$ del anillo $A=K[s,t]$ de polinomios en dos indeterminadas sobre un cuerpo $K$, el abierto $U$ dado por la unión de $D(s)$ y $D(t)$ \emph{no es un abierto afín}.
En efecto, si $z$ es una sección de $\mathscr{O}_Y$ sobre $U$, existen dos enteros $m\geq0$, $n\geq0$ tales que $s^m z$ y $t^n z$ son las restricciones de polinomios en $s$ y $t$ a $U$ \sref{I.1.4.1}, lo cual solo es posible, evidentemente, si la sección $z$ se prolonga a una sección sobre todo $Y$, identificada con un polinomio en $s$ y $t$.
Si $U$ fuese un abierto afín, el morfismo de inyección $U\to Y$ sería un isomorfismo \sref{I.1.7.3}, lo cual es contradictorio, pues $U\neq Y$.

Teniendo en cuenta lo anterior, tomemos dos esquemas afines $Y_1$ e $Y_2$, espectros primos de los anillos $A_1=K[s_1,t_1]$ y $A_2=K[s_2,t_2]$, respectivamente;
tomemos $U_{12}=D(s_1)\cup D(t_1)$ y $U_{21}=D(s_2)\cup D(t_2)$, y tomemos como $u_{12}$ la restricción de un isomorfismo $Y_2\to Y_1$ a $U_{21}$ correspondiente al isomorfismo de anillos que envía $f(s_1,t_1)$ a $f(s_2,t_2)$;
se obtiene entonces un ejemplo en el que no se satisface la primera condición de \sref{I.5.5.6}, mientras que la segunda sí se satisface\footnote{\emph{[Trad.] La fuente francesa afirma que no se satisface ninguna de las dos condiciones de \sref{I.5.5.6}. Sin embargo, para $U=D(s)\cup D(t)$ se tiene $\Gamma(U,\sh{O}_Y)=K[s,t]$, de modo que la segunda condición sí se satisface; solo falla la afinidad de $U$.}} (el preesquema íntegro así obtenido se llama ``plano afín sobre $K$ con el punto $0$ duplicado'').
\end{examples}

\begin{remark}[5.5.12]
\label{I.5.5.12}
Dada una propiedad \textbf{\textit P} de los morfismos de preesquemas, consideremos las proposiciones siguientes.
\begin{enumerate}
  \item[(i)] \emph{Toda inmersión cerrada tiene la propiedad~\textbf{\textup{P}}.}
  \item[(ii)] \emph{La composición de dos morfismos cualesquiera que tengan ambos la propiedad~\textbf{\textup{P}} tiene también la propiedad~\textbf{\textup{P}}.}
  \item[(iii)] \emph{Si $f:X\to X'$ y $g:Y\to Y'$ son $S$-morfismos que tienen la propiedad~\textbf{\textup{P}}, entonces $f\times_S g$ tiene la propiedad~\textbf{\textup{P}}.}
\oldpage[I]{140}
  \item[(iv)] \emph{Si $f:X\to Y$ es un $S$-morfismo que tiene la propiedad~\textbf{P}, todo $S'$-morfismo $f_{(S')}$ obtenido mediante una extensión $S'\to S$ del preesquema base tiene también la propiedad~\textbf{P}.}
  \item[(v)] \emph{Si la composición $g\circ f$ de dos morfismos $f:X\to Y$ y $g:Y\to Z$ tiene la propiedad~\textbf{P}, y $g$ es separado, entonces $f$ tiene la propiedad~\textbf{P}.}
  \item[(vi)] \emph{Si un morfismo $f:X\to Y$ tiene la propiedad~\textbf{P}, entonces $f_{\mathrm{red}}$ también la tiene \sref{I.5.1.5}.}
\end{enumerate}
Si se supone que (i) y (ii) son verdaderas, entonces (iii) y (iv) son \emph{equivalentes}, y (v) y (vi) son \emph{consecuencias} de (i), (ii) y (iii).

La primera aserción ya ha sido demostrada \sref{I.3.5.1}.
Consideremos la factorización \sref{I.5.3.13} $X\xrightarrow{\Gamma_f}X\times_Z Y\xrightarrow{p_2}Y$ de $f$;
la igualdad $p_2=(g\circ f)\times_Z 1_Y$ muestra que, si $g\circ f$ tiene la propiedad~\textbf{\textit{P}}, entonces $p_2$ también la tiene, en virtud de (iii);
si $g$ es separado, entonces $\Gamma_f$ es una inmersión cerrada \sref{I.5.4.3}, y tiene también, por tanto, la propiedad~\textbf{\textit{P}}, en virtud de (i);
finalmente, por (ii), $f$ tiene la propiedad~\textbf{\textit{P}}.

Consideremos, por último, el diagrama conmutativo
\[
  \xymatrix{
    X_{\mathrm{red}}\ar[r]^{f_{\mathrm{red}}}\ar[d] &
    Y_{\mathrm{red}}\ar[d]\\
    X\ar[r]_f &
    Y
  }
\]
donde las flechas verticales son inmersiones cerradas \sref{I.5.1.5} y tienen, por tanto, la propiedad~\textbf{\textit{P}}, en virtud de (i).
La hipótesis de que $f$ tiene la propiedad~\textbf{\textit{P}} implica, por (ii), que $X_{\mathrm{red}}\xrightarrow{f_{\mathrm{red}}}Y_{\mathrm{red}}\to Y$ tiene la propiedad~\textbf{\textit{P}};
finalmente, como una inmersión cerrada es separada \sref{I.5.5.1}[i], $f_{\mathrm{red}}$ tiene la propiedad~\textbf{\textit{P}}, en virtud de (v).

Se observará que, si se consideran las proposiciones
\begin{enumerate}
  \item[(i')] \emph{Toda inmersión tiene la propiedad~\textbf{P}};
  \item[(v')] \emph{Si $g\circ f$ tiene la propiedad~\textbf{P}, entonces $f$ también la tiene};
\end{enumerate}
los argumentos anteriores muestran que (v') es consecuencia de (i'), (ii) y (iii).
\end{remark}

\begin{env}[5.5.13]
\label{I.5.5.13}
Se observará que (v) y (vi) son también consecuencias de (i), (iii) y
\begin{enumerate}
  \item[(ii')] \emph{Si $j:X\to Y$ es una inmersión cerrada y $g:Y\to Z$ es un morfismo que tiene la propiedad~\textbf{P}, entonces $g\circ j$ tiene la propiedad~\textbf{P}}.
\end{enumerate}
Análogamente, (v') es consecuencia de (i'), (iii) y
\begin{enumerate}
  \item[(ii'')] \emph{Si $j:X\to Y$ es una inmersión y $g:Y\to Z$ es un morfismo que tiene la propiedad~\textbf{P}, entonces $g\circ j$ tiene la propiedad \textbf{P}}.
\end{enumerate}
Esto resulta inmediatamente de los argumentos de \sref{I.5.5.12}.
\end{env}
% END source_es/ega1/ega1-5.tex

% BEGIN source_es/ega1/ega1-6.tex
\section{Condiciones de finitud}
\label{section:I.6}

\subsection{Preesquemas noetherianos y localmente noetherianos}
\label{subsection:I.6.1}

\begin{definition}[6.1.1]
\label{I.6.1.1}
{\itshape Se dice que un preesquema $X$ es noetheriano \textup{(resp.} localmente noetheriano) si es una unión finita \textup{(resp.} una unión) de abiertos afines $V_\alpha$ tales que el anillo del esquema inducido sobre cada $V_\alpha$ es noetheriano.}
\end{definition}

Se sigue inmediatamente de \sref{I.1.5.2} que, si $X$ es localmente noetheriano, el haz estructural $\mathscr{O}_X$ es un \emph{haz coherente de anillos}, pues la cuestión es local.
Todo \emph{$\mathscr{O}_X$-submódulo}
\oldpage[I]{141}
(resp. \emph{todo $\mathscr{O}_X$-módulo cociente}) \emph{cuasicoherente} de un módulo \emph{coherente} sobre $\mathscr{O}_X$, denotado por $\mathscr{F}$, es \emph{coherente}, pues la cuestión es de nuevo local y basta aplicar \sref{I.1.5.1}, \sref{I.1.4.1} y \sref{I.1.3.10}, junto con el hecho de que un submódulo (resp. módulo cociente) de un módulo de tipo finito sobre un anillo noetheriano es de tipo finito.
En particular, todo \emph{haz cuasicoherente de ideales} de $\mathscr{O}_X$ es \emph{coherente}.

Si un preesquema $X$ es una unión finita (resp. una unión) de abiertos $W_\lambda$ de modo que los preesquemas inducidos sobre los $W_\lambda$ sean noetherianos (resp. localmente noetherianos), es claro que $X$ es noetheriano (resp. localmente noetheriano).

\begin{proposition}[6.1.2]
\label{I.6.1.2}
Para que un preesquema $X$ sea noetheriano, es necesario y suficiente que sea localmente noetheriano y que su espacio subyacente sea cuasicompacto.
El espacio subyacente de $X$ es entonces noetheriano.
\end{proposition}

\begin{proof}
La primera aserción se sigue inmediatamente de las definiciones y de \sref{I.1.1.10}[ii].
La segunda se sigue de \sref{I.1.1.6} y del hecho de que todo espacio que es unión finita de subespacios noetherianos es a su vez noetheriano \sref[0]{0.2.2.3}.
\end{proof}

\begin{proposition}[6.1.3]
\label{I.6.1.3}
Sea $X$ un esquema afín de anillo $A$.
Las condiciones siguientes son equivalentes:
\textup{a)} $X$ es noetheriano;
\textup{b)} $X$ es localmente noetheriano;
\textup{c)} $A$ es noetheriano.
\end{proposition}

\begin{proof}
La equivalencia entre \emph{a)} y \emph{b)} se sigue de \sref{I.6.1.2} y del hecho de que el espacio subyacente de todo esquema afín es cuasicompacto \sref{I.1.1.10}; es claro, además, que \emph{c)} implica \emph{a)}.
Para ver que \emph{a)} implica \emph{c)}, observemos que existe un recubrimiento finito $(V_i)$ de $X$ por abiertos afines tal que el anillo $A_i$ del preesquema inducido sobre $V_i$ es noetheriano.
Sea, pues, $(\mathfrak{a}_n)$ una sucesión creciente de ideales de $A$; mediante la correspondencia biyectiva canónica \sref{I.1.3.7}, le corresponde una sucesión creciente $(\widetilde{\mathfrak{a}}_n)$ de haces de ideales en $\widetilde{A}=\mathscr{O}_X$;
para ver que la sucesión $(\mathfrak{a}_n)$ es estacionaria, basta demostrar que la sucesión $(\widetilde{\mathfrak{a}}_n)$ lo es.
Ahora bien, la restricción $\widetilde{\mathfrak{a}}_n|V_i$ es un haz cuasicoherente de ideales de $\mathscr{O}_X|V_i$, por ser la imagen inversa de $\widetilde{\mathfrak{a}}_n$ por la inyección canónica $V_i\to X$ \sref[0]{0.5.1.4};
$\widetilde{\mathfrak{a}}_n|V_i$ es, pues, de la forma $\widetilde{\mathfrak{a}}_{ni}$, donde $\mathfrak{a}_{ni}$ es un ideal de $A_i$ \sref{I.1.3.7}.
Como $A_i$ es noetheriano, la sucesión $(\mathfrak{a}_{ni})$ es estacionaria para todo $i$, de donde resulta la proposición.
\end{proof}

Se observará que el argumento anterior demuestra también que \emph{si $X$ es un preesquema noetheriano, toda sucesión creciente de haces coherentes de ideales de $\mathscr{O}_X$ es estacionaria}.

\begin{proposition}[6.1.4]
\label{I.6.1.4}
Todo subpreesquema de un preesquema noetheriano \textup{(resp.} localmente noetheriano) es noetheriano \textup{(resp.} localmente noetheriano).
\end{proposition}

\begin{proof}
Basta hacer la demostración para un preesquema noetheriano $X$;
además, por la definición~\sref{I.6.1.1}, se puede también reducir al caso en que $X$ es un esquema afín.
Como todo subpreesquema de $X$ es un subpreesquema cerrado de un preesquema inducido sobre un abierto \sref{I.4.1.3}, se puede reducir al caso de un subpreesquema $Y$ cerrado o inducido sobre un abierto de $X$.
La demostración en el caso en que $Y$ es cerrado es inmediata, pues, si $A$ es el anillo de $X$, se sabe que $Y$ es un esquema afín de anillo $A/\mathfrak{J}$, donde $\mathfrak{J}$ es un ideal de $A$ \sref{I.4.2.3};
como $A$ es noetheriano \sref{I.6.1.3}, también lo es $A/\mathfrak{J}$.

Supongamos ahora que $Y$ es abierto en $X$;
el espacio subyacente de $Y$ es noetheriano \sref{I.6.1.2}, por tanto cuasicompacto y, en consecuencia, una unión finita de abiertos $D(f_i)$ ($f_i\in A$);
\oldpage[I]{142}
todo se reduce a demostrar la proposición en el caso en que $Y=D(f)$ con $f\in A$.
Pero entonces $Y$ es un esquema afín cuyo anillo es isomorfo a $A_f$ \sref{I.1.3.6};
como $A$ es noetheriano \sref{I.6.1.3}, también lo es $A_f$.
\end{proof}

\begin{env}[6.1.5]
\label{I.6.1.5}
Se observará que el \emph{producto} de dos $S$-preesquemas noetherianos no es necesariamente noetheriano, incluso si los preesquemas son afines, pues el producto tensorial de dos álgebras noetherianas no es necesariamente un anillo noetheriano (cf. \sref{I.6.3.8}).
\end{env}

\begin{proposition}[6.1.6]
\label{I.6.1.6}
Si $X$ es un preesquema noetheriano, el nilradical $\mathscr{N}_X$ de $\mathscr{O}_X$ es nilpotente.
\end{proposition}

\begin{proof}
Se puede, en efecto, recubrir $X$ por un número finito de abiertos afines $U_i$, y basta demostrar que existen enteros $n_i$ tales que $(\mathscr{N}_X|U_i)^{n_i}=0$;
si $n$ es el mayor de los $n_i$, se tendrá entonces $\mathscr{N}_X^n=0$.
Se puede, pues, reducir al caso en que $X=\operatorname{Spec}(A)$ es afín, siendo $A$ un anillo noetheriano;
por \sref{I.5.1.1} y \sref{I.1.3.13}, basta observar que el nilradical de $A$ es nilpotente (\cite[p.~127, cor.~4]{I-11}).
\end{proof}

\begin{corollary}[6.1.7]
\label{I.6.1.7}
Sea $X$ un preesquema noetheriano;
para que $X$ sea un esquema afín, es necesario y suficiente que $X_{\mathrm{red}}$ sea afín.
\end{corollary}

\begin{proof}
Esto se sigue de \sref{I.6.1.6} y \sref{I.5.1.10}.
\end{proof}

\begin{lemma}[6.1.8]
\label{I.6.1.8}
Sean $X$ un espacio topológico, $x$ un punto de $X$ y $U$ una vecindad abierta de $x$ que solo tiene un número finito de componentes irreducibles.
Entonces existe una vecindad $V$ de $x$ tal que toda vecindad abierta de $x$ contenida en $V$ sea conexa.
\end{lemma}

\begin{proof}
Sean $U_i$ ($1\leq i\leq m$) las componentes irreducibles de $U$ que no contienen a $x$;
el complemento (en $U$) de la unión de los $U_i$ es una vecindad abierta $V$ de $x$ en $U$, y por tanto también en $X$;
es además la intersección con $U$ del complemento (en $X$) de la unión de las componentes irreducibles de $X$ que no contienen a $x$ \sref[0]{0.2.1.6}\footnote{\emph{[Trad.] La fuente francesa identifica aquí $V$ directamente con el complemento en $X$ de la unión de las componentes irreducibles que no contienen a $x$, omitiendo la intersección con $U$. La igualdad impresa es falsa, por ejemplo, si $X$ es irreducible y $U\subsetneq X$: entonces aquel complemento es $X$, mientras que $V=U$.}}.
Sea entonces $W$ una vecindad abierta de $x$ contenida en $V$.
Las componentes irreducibles de $W$ son las intersecciones con $W$ de las componentes irreducibles de $U$ que intersecan $W$ \sref[0]{0.2.1.6}, por lo que estas componentes contienen a $x$;
como son conexas, $W$ también lo es.
\end{proof}

\begin{corollary}[6.1.9]
\label{I.6.1.9}
Un espacio topológico localmente noetheriano es localmente conexo (lo que implica, entre otras cosas, que sus componentes conexas son abiertas).
\end{corollary}

\begin{proposition}[6.1.10]
\label{I.6.1.10}
Sea $X$ un espacio topológico localmente noetheriano.
Las condiciones siguientes son equivalentes.
\begin{enumerate}
  \item[\textup{a)}] Las componentes irreducibles de $X$ son abiertas.
  \item[\textup{b)}] Las componentes irreducibles de $X$ son exactamente sus componentes conexas.
  \item[\textup{c)}] Las componentes conexas de $X$ son irreducibles.
  \item[\textup{d)}] Dos componentes irreducibles distintas de $X$ tienen intersección vacía.
\end{enumerate}
Por último, si $X$ es un preesquema, estas condiciones son también equivalentes a
\begin{enumerate}
  \item[\textup{e)}] Para todo $x\in X$, $\operatorname{Spec}(\mathscr{O}_x)$ es irreducible (o, dicho de otro modo, el nilradical de $\mathscr{O}_x$ es primo).
\end{enumerate}
\end{proposition}

\begin{proof}
Es inmediato que \emph{a)} implica \emph{b)}, pues todo espacio irreducible es conexo, y \emph{a)} implica que las componentes irreducibles de $X$ son los conjuntos a la vez abiertos y cerrados.
Es trivial que \emph{b)} implica \emph{c)}; recíprocamente, un conjunto cerrado $F$ que contiene
\oldpage[I]{143}
una componente conexa $C$ de $X$, y distinta de $C$, no puede ser irreducible, pues, al no ser conexo, este conjunto es la unión de dos conjuntos no vacíos disjuntos, a la vez abiertos y cerrados en $F$ y, por tanto, cerrados en $X$; en consecuencia, \emph{c)} implica \emph{b)}.
De aquí se deduce inmediatamente que \emph{c)} implica \emph{d)}, puesto que dos componentes conexas distintas no tienen ningún punto en común.

Hasta ahora no hemos utilizado el hecho de que $X$ sea localmente noetheriano.
Supongamos ahora que así es, y demostremos que \emph{d)} implica \emph{a)}: por \sref[0]{0.2.1.6}, se puede limitar al caso en que el espacio $X$ sea noetheriano y, por tanto, solo tenga un número finito de componentes irreducibles.
Como son cerradas y disjuntas dos a dos, son abiertas.

Por último, la equivalencia entre \emph{d)} y \emph{e)} es válida incluso sin suponer que el espacio subyacente del preesquema $X$ sea localmente noetheriano.
En efecto, se puede limitar al caso en que $X=\operatorname{Spec}(A)$ sea afín, por \sref[0]{0.2.1.6};
decir que $x$ está contenido en una única componente irreducible de $X$ equivale a decir que $\mathfrak{j}_x$ contiene un único ideal primo minimal de $A$ \sref{I.1.1.14}, lo que equivale a decir que $\mathfrak{j}_x\mathscr{O}_x$ contiene un único ideal primo minimal de $\mathscr{O}_x$, de donde se sigue la conclusión.
\end{proof}

\begin{corollary}[6.1.11]
\label{I.6.1.11}
Sea $X$ un espacio localmente noetheriano.
Para que $X$ sea irreducible, es necesario y suficiente que $X$ sea conexo y no vacío, y que cualesquiera dos componentes irreducibles distintas de $X$ tengan intersección vacía.
Si $X$ es un preesquema, esta última condición equivale a que $\operatorname{Spec}(\mathscr{O}_x)$ sea irreducible para todo $x\in X$.
\end{corollary}

\begin{proof}
La segunda afirmación ya se ha demostrado en \sref{I.6.1.10};
solo queda, por tanto, demostrar que las condiciones de la primera afirmación son suficientes.
Pero, por \sref{I.6.1.10}, estas condiciones implican que las componentes irreducibles de $X$ son exactamente sus componentes conexas y, puesto que $X$ es conexo y no vacío, es irreducible.
\end{proof}

\begin{corollary}[6.1.12]
\label{I.6.1.12}
Sea $X$ un preesquema localmente noetheriano.
Para que $X$ sea íntegro, es necesario y suficiente que $X$ sea conexo y no vacío\footnote{\emph{[Trad.] La fuente francesa omite «y no vacío». Sin esta condición, el caso vacío satisfaría vacíamente la condición sobre todos los anillos locales, aunque no sería irreducible; además, el Corolario~\sref{I.6.1.11} exige explícitamente que $X$ sea no vacío.}}, y que $\mathscr{O}_x$ sea íntegro para todo $x\in X$.
\end{corollary}

\begin{proposition}[6.1.13]
\label{I.6.1.13}
Sea $X$ un preesquema localmente noetheriano, y sea $x\in X$ un punto tal que el nilradical $\mathscr{N}_x$ de $\mathscr{O}_x$ sea primo \textup{(resp.} tal que $\mathscr{O}_x$ sea reducido, \textup{resp.} íntegro);
entonces existe una vecindad abierta $U$ de $x$ que sea irreducible \textup{(resp.} reducida, \textup{resp.} íntegra).
\end{proposition}

\begin{proof}
Basta considerar dos casos: aquel en que $\mathscr{N}_x$ es primo y aquel en que $\mathscr{N}_x=0$; la tercera hipótesis es la conjunción de las dos primeras.
Si $\mathscr{N}_x$ es primo, entonces $x$ pertenece a una única componente irreducible $Y$ de $X$ \sref{I.6.1.10};
la unión de las componentes irreducibles de $X$ que no contienen a $x$ es cerrada (pues el conjunto de estas componentes es localmente finito), y el complemento $U$ de esta unión es por tanto abierto y está contenido en $Y$, y por consiguiente es irreducible \sref[0]{0.2.1.6}.
Si $\mathscr{N}_x=0$, también se tiene $\mathscr{N}_y=0$ para todo $y$ en una vecindad de $x$, porque $\mathscr{N}$ es cuasicoherente \sref{I.5.1.1}, y por tanto coherente, ya que $X$ es localmente noetheriano, y la conclusión se sigue entonces de \sref[0]{0.5.2.2}.
\end{proof}

\subsection{Preesquemas artinianos}
\label{subsection:I.6.2}
\label{I.6.2}

\begin{definition}[6.2.1]
\label{I.6.2.1}
{\itshape Se dice que un preesquema es artiniano si es afín y su anillo es artiniano.}
\end{definition}

\oldpage[I]{144}
\begin{proposition}[6.2.2]
\label{I.6.2.2}
Dado
un preesquema $X$, las condiciones siguientes son equivalentes:
\begin{enumerate}
  \item[\textit{a)}] $X$ es un esquema artiniano;
  \item[\textit{b)}] $X$ es noetheriano y su espacio subyacente es discreto;
  \item[\textit{c)}] $X$ es noetheriano y los puntos de su espacio subyacente son cerrados \emph{(condición $\mathrm{T}_1$)}.
\end{enumerate}
Cuando se cumple cualquiera de las condiciones anteriores, el espacio subyacente a $X$ es finito, y el anillo $A$ de $X$ es el producto directo de los anillos locales (artinianos) en los puntos de $X$.
\end{proposition}

\begin{proof}
Se sabe que \emph{a)} implica la última afirmación (\cite[p.~205, th.~3]{I-13}); por tanto, todo ideal primo de $A$ es maximal y es la imagen inversa de un ideal maximal de una de las componentes locales de $A$, de modo que el espacio $X$ es finito y discreto;
\emph{a)} implica, pues, \emph{b)}, y \emph{b)} implica evidentemente \emph{c)}.
Para ver que \emph{c)} implica \emph{a)}, demostremos primero que $X$ es entonces finito;
se puede reducir, en efecto, al caso en que $X$ es afín, y se sabe que un anillo noetheriano cuyos ideales primos son todos maximales es artiniano (\cite[p.~203]{I-13}), de donde resulta nuestra afirmación.
El espacio subyacente $X$ es entonces discreto, suma topológica de un número finito de puntos $x_i$, y los anillos locales $\mathscr{O}_{x_i}=A_i$ son artinianos;
está claro que $X$ es isomorfo al esquema afín espectro primo del anillo $A$ producto directo de los $A_i$ \sref{I.1.7.3}.
\end{proof}

\subsection{Morfismos de tipo finito}
\label{subsection:I.6.3}

\begin{definition}[6.3.1]
\label{I.6.3.1}
{\itshape Se dice que un morfismo $f:X\to Y$ es de tipo finito si $Y$ es la unión de una familia $(V_\alpha)$ de abiertos afines que tienen la propiedad siguiente:
\begin{enumerate}
  \item[\textup{(P)}] $f^{-1}(V_\alpha)$ es una unión finita de abiertos afines $U_{\alpha i}$ tales que cada anillo $A(U_{\alpha i})$ es un álgebra de tipo finito sobre $A(V_\alpha)$.
\end{enumerate}
Se dice entonces que $X$ es un preesquema de tipo finito sobre $Y$, o un $Y$-preesquema de tipo finito.}
\end{definition}

\begin{proposition}[6.3.2]
\label{I.6.3.2}
Si $f:X\to Y$ es un morfismo de tipo finito, todo abierto afín $W$ de $Y$ satisface la propiedad \emph{(P)} de la definición~\sref{I.6.3.1}.
\end{proposition}

Demostremos primero
\begin{lemma}[6.3.2.1]
\label{I.6.3.2.1}
Si $T\subset Y$ es un abierto afín que satisface la propiedad \emph{(P)}, entonces, para todo $g\in A(T)$, $D(g)$ satisface también la propiedad \emph{(P)}.
\end{lemma}

\begin{proof}
Por hipótesis, $f^{-1}(T)$ es una unión finita de abiertos afines $Z_j$ tales que $A(Z_j)$ es un álgebra de tipo finito sobre $A(T)$;
sea $\varphi_j:A(T)\to A(Z_j)$ el homomorfismo de anillos correspondiente a la restricción de $f$ a $Z_j$ \sref{I.2.2.4}, y pongamos $g_j=\varphi_j(g)$;
se tiene entonces $f^{-1}(D(g))\cap Z_j=D(g_j)$ \sref{I.1.2.2.2}.
Ahora bien, $A(D(g_j))=A(Z_j)_{g_j}=A(Z_j)[1/g_j]$ es de tipo finito sobre $A(Z_j)$ y, \emph{a fortiori}, sobre $A(T)$ en virtud de la hipótesis, y por tanto también sobre $A(D(g))=A(T)[1/g]$, lo que demuestra el lema.
\end{proof}

\begin{proof}
En virtud del lema anterior, como $W$ es cuasicompacto \sref{I.1.1.10}, existe un recubrimiento finito de $W$ por conjuntos de la forma \erratum[II]{$D(g_i)$}, donde cada $g_i$ pertenece a un anillo $A(V_{\alpha(i)})$.
Cada $D(g_i)$, al ser cuasicompacto, es una unión finita de conjuntos $D(h_{ik})$, donde $h_{ik}\in A(W)$;
si $\varphi_i:A(W)\to A(D(g_i))$ es la aplicación canónica, se tiene $D(h_{ik})=D(\varphi_i(h_{ik}))$ por \sref{I.1.2.2.2}.
Por \sref{I.6.3.2.1}, cada uno de los $f^{-1}(D(h_{ik}))$ admite un recubrimiento finito por abiertos afines $U_{ijk}$ tales que los $A(U_{ijk})$ son álgebras de tipo finito sobre $A(D(h_{ik}))=A(W)[1/h_{ik}]$, de donde resulta la proposición.
\end{proof}

\oldpage[I]{145}
Se puede, por tanto, decir que la noción de preesquema de tipo finito sobre $Y$ es \emph{local sobre $Y$}.

\begin{proposition}[6.3.3]
\label{I.6.3.3}
Sean $X$ e $Y$ dos esquemas afines;
para que $X$ sea de tipo finito sobre $Y$, es necesario y suficiente que $A(X)$ sea un álgebra de tipo finito sobre $A(Y)$.
\end{proposition}

\begin{proof}
Como la condición es evidentemente suficiente, demostremos que es necesaria.
Pongamos $A=A(Y)$ y $B=A(X)$;
en virtud de \sref{I.6.3.2}, existe un recubrimiento abierto afín finito $(V_i)$ de $X$ tal que cada uno de los anillos $A(V_i)$ es un $A$-álgebra de tipo finito.
Además, como los $V_i$ son cuasicompactos, se puede recubrir cada uno de ellos por un número finito de abiertos de la forma $D(g_{ij})\subset V_i$, donde $g_{ij}\in B$;
si $\varphi_i$ es el homomorfismo $B\to A(V_i)$ correspondiente a la inyección canónica $V_i\to X$, entonces se tiene \[
  B_{g_{ij}}=(A(V_i))_{\varphi_i(g_{ij})}
  =A(V_i)[1/\varphi_i(g_{ij})],
\], por lo que $B_{g_{ij}}$ es un $A$-álgebra de tipo finito.
Se puede, pues, limitar al caso en que $V_i=D(g_i)$ con $g_i\in B$.
Por hipótesis, existe un subconjunto finito $F_i$ de $B$ y un entero $n_i\geq 0$ tales que $B_{g_i}$ es el álgebra engendrada sobre $A$ por los elementos $b_i/g_i^{n_i}$, donde $b_i$ recorre $F_i$.
Como los $g_i$ son finitos en número, se puede suponer que todos los $n_i$ son iguales a un mismo entero $n$.
Además, como los $D(g_i)$ forman un recubrimiento de $X$, el ideal engendrado en $B$ por los $g_i$ es igual a $B$ o, en otros términos, existen $h_i\in B$ tales que $\sum_i h_i g_i=1$.
Sea entonces $F$ el subconjunto finito de $B$ formado por la unión de los $F_i$, el conjunto de los $g_i$ y el conjunto de los $h_i$; demostremos que el subanillo $B'=A[F]$ de $B$ es igual a $B$.
Por hipótesis, para todo $b\in B$ y todo $i$, la imagen canónica de $b$ en $B_{g_i}$ es de la forma $b'_i/g_i^{m_i}$, donde $b'_i\in B'$;
multiplicando los $b'_i$ por potencias adecuadas de los $g_i$, se puede suponer de nuevo que todos los $m_i$ son iguales a un mismo entero $m$.
Por definición de los anillos de fracciones, existe, pues, un entero $N$ (dependiente de $b$) tal que $N\geq m$ y $g_i^N b\in B'$ para todo $i$;
pero, \emph{en el anillo $B'$}, los $g_i^N$ engendran el ideal $B'$, pues ya lo hacen los $g_i$ (y los $h_i$ pertenecen a $B'$);
existen, por tanto, $c_i\in B'$ tales que $\sum_i c_i g_i^N=1$, de donde $b=\sum_i c_i g_i^N b\in B'$, Q.E.D.
\end{proof}

\begin{proposition}[6.3.4]
\label{I.6.3.4}
\medskip\noindent
\begin{enumerate}
  \item[{\rm(i)}] Toda inmersión cerrada es de tipo finito.
  \item[{\rm(ii)}] El compuesto de dos morfismos de tipo finito es de tipo finito.
  \item[{\rm(iii)}] Si $f:X\to X'$ y $g:Y\to Y'$ son $S$-morfismos de tipo finito, entonces $f\times_S g$ es de tipo finito.
  \item[{\rm(iv)}] Si $f:X\to Y$ es un $S$-morfismo de tipo finito, entonces $f_{(S')}$ es de tipo finito para toda extensión $g:S'\to S$ del preesquema de base.
  \item[{\rm(v)}] Si el compuesto $g\circ f$ de dos morfismos es de tipo finito, siendo $g$ separado, entonces $f$ es de tipo finito.
  \item[{\rm(vi)}] Si un morfismo $f$ es de tipo finito, entonces $f_{\mathrm{red}}$ es de tipo finito.
\end{enumerate}
\end{proposition}

\begin{proof}
En virtud de \sref{I.5.5.12}, basta demostrar (i), (ii) y (iv).

Para demostrar (i), se puede limitar al caso de una inyección canónica $X\to Y$, siendo $X$ un subpreesquema cerrado de $Y$;
además, por \sref{I.6.3.2}, se puede suponer que $Y$ es afín, en cuyo caso $X$ también es afín \sref{I.4.2.3} y su anillo es isomorfo a un anillo cociente $A/\mathfrak{J}$, donde $A$ es el anillo de $Y$ y $\mathfrak{J}$ es un ideal de $A$;
como $A/\mathfrak{J}$ es de tipo finito sobre $A$, se sigue la conclusión.

Demostremos ahora (ii).
Sean $f:X\to Y$ y $g:Y\to Z$ dos morfismos de tipo finito, y sea $U$ un abierto afín de $Z$;
$g^{-1}(U)$ admite un recubrimiento finito por abiertos afines $V_i$ tales que cada $A(V_i)$ es un álgebra de tipo finito sobre $A(U)$ \sref{I.6.3.2};
del mismo modo,
\oldpage[I]{146}
cada uno de los $f^{-1}(V_i)$ admite un recubrimiento finito por abiertos afines $W_{ij}$ tales que cada $A(W_{ij})$ es un álgebra de tipo finito sobre $A(V_i)$ y, por tanto, también un álgebra de tipo finito sobre $A(U)$, de donde la conclusión.

Finalmente, para demostrar (iv), se puede limitar al caso en que $S=Y$;
entonces $f_{(S')}$ es también igual a $f_{(Y_{(S')})}$, donde se considera $f$ como un $Y$-morfismo y la extensión de base es $Y_{(S')}\to Y$ \sref{I.3.3.9}.
Sean entonces $p$ y $q$ las proyecciones $X_{(S')}\to X$ y $X_{(S')}\to S'$.
Sea $V$ un abierto afín de $S$;
$f^{-1}(V)$ es una unión finita de abiertos afines $W_i$ tales que $A(W_i)$ es un álgebra de tipo finito sobre $A(V)$ \sref{I.6.3.2}.
Sea $V'$ un abierto afín de $S'$ contenido en $g^{-1}(V)$;
como $f\circ p=g\circ q$, $q^{-1}(V')$ está contenido en la unión de los $p^{-1}(W_i)$;
por otra parte, la intersección $p^{-1}(W_i)\cap q^{-1}(V')$ se identifica con el producto $W_i\times_V V'$ \sref{I.3.2.7}, que es un esquema afín cuyo anillo es isomorfo a $A(W_i)\otimes_{A(V)}A(V')$ \sref{I.3.2.2};
este anillo es, por hipótesis, un álgebra de tipo finito sobre $A(V')$, lo que demuestra la proposición.
\end{proof}

\begin{corollary}[6.3.5]
\label{I.6.3.5}
Sea $f:X\to Y$ un morfismo de inmersión.
Si el espacio subyacente a $Y$ \textup{(resp. $X$)} es localmente noetheriano \textup{(resp.} noetheriano), entonces $f$ es de tipo finito.
\end{corollary}

\begin{proof}
Se puede suponer siempre que $Y$ es afín \sref{I.6.3.2};
si el espacio subyacente a $Y$ es localmente noetheriano, se puede suponer además que es noetheriano, y entonces el espacio subyacente a $X$, que es un subespacio suyo, también es noetheriano.
En otros términos, se puede suponer que $Y$ es afín y que el espacio subyacente a $X$ es noetheriano;
existe entonces un recubrimiento de $X$ por un número finito de abiertos afines $D(g_i)\subset Y$, donde $g_i\in A(Y)$, tales que las intersecciones $X\cap D(g_i)$ son cerradas en $D(g_i)$ (y, por tanto, esquemas afines \sref{I.4.2.3}), pues $X$ es localmente cerrado en $Y$ \sref{I.4.1.3}.
Entonces $A(X\cap D(g_i))$ es un álgebra de tipo finito sobre $A(D(g_i))$, por \sref{I.6.3.4}[i] y \sref{I.6.3.3}, y $A(D(g_i))=A(Y)_{g_i}=A(Y)[1/g_i]$ es de tipo finito sobre $A(Y)$, lo que termina la demostración.
\end{proof}

\begin{corollary}[6.3.6]
\label{I.6.3.6}
Sean $f:X\to Y$ y $g:Y\to Z$ dos morfismos.
Si $g\circ f$ es de tipo finito y se cumple que $X$ es noetheriano o $X\times_Z Y$ es localmente noetheriano, entonces $f$ es de tipo finito.
\end{corollary}

\begin{proof}
Esto resulta inmediatamente de la demostración de \sref{I.5.5.12} y de \sref{I.6.3.5} aplicada al morfismo de inmersión $\Gamma_f$.
\end{proof}

\begin{proposition}[6.3.7]
\label{I.6.3.7}
Sea $f:X\to Y$ un morfismo de tipo finito;
si $Y$ es noetheriano \textup{(resp.} localmente noetheriano), entonces $X$ es noetheriano \textup{(resp.} localmente noetheriano).
\end{proposition}

\begin{proof}
Se puede limitar a demostrar la proposición cuando $Y$ es noetheriano.
Entonces $Y$ es unión finita de abiertos afines $V_i$ tales que los $A(V_i)$ son anillos noetherianos.
En virtud de \sref{I.6.3.2}, cada uno de los $f^{-1}(V_i)$ es unión finita de abiertos afines $W_{ij}$ tales que los $A(W_{ij})$ son álgebras de tipo finito sobre $A(V_i)$ y, por tanto, anillos noetherianos;
esto demuestra que $X$ es noetheriano.
\end{proof}

\begin{corollary}[6.3.8]
\label{I.6.3.8}
Sea $X$ un preesquema de tipo finito sobre $S$.
Para toda extensión de la base $S'\to S$ tal que $S'$ sea noetheriano \textup{(resp.} localmente noetheriano), $X_{(S')}$ es noetheriano \textup{(resp.} localmente noetheriano).
\end{corollary}

\begin{proof}
Esto resulta de \sref{I.6.3.7}, pues $X_{(S')}$ es de tipo finito sobre $S'$ en virtud de \sref{I.6.3.4}[iv].
\end{proof}

También se puede decir que, en un producto $X\times_S Y$ de $S$-preesquemas, si \emph{uno} de los
\oldpage[I]{147}
factores $X$ o $Y$ es \emph{de tipo finito} sobre $S$ y \emph{el otro es noetheriano} (resp. \emph{localmente noetheriano}), entonces $X\times_S Y$ es \emph{noetheriano} (resp. \emph{localmente noetheriano}).

\begin{corollary}[6.3.9]
\label{I.6.3.9}
Sea $X$ un preesquema de tipo finito sobre un preesquema localmente noetheriano $S$.
Entonces todo $S$-morfismo $f:X\to Y$ es de tipo finito.
\end{corollary}

\begin{proof}
En efecto, se puede suponer que $S$ es noetheriano;
si $\varphi:X\to S$ y $\psi:Y\to S$ son los morfismos estructurales, se tiene $\varphi=\psi\circ f$, y $X$ es noetheriano en virtud de \sref{I.6.3.7};
por tanto, $f$ es de tipo finito en virtud de \sref{I.6.3.6}.
\end{proof}

\begin{proposition}[6.3.10]
\label{I.6.3.10}
Sea $f:X\to Y$ un morfismo de tipo finito.
Para que $f$ sea sobreyectivo, es necesario y suficiente que, para todo cuerpo \emph{algebraicamente cerrado} $\Omega$, la aplicación $X(\Omega)\to Y(\Omega)$ correspondiente a $f$ \sref{I.3.4.1} sea sobreyectiva.
\end{proposition}

\begin{proof}
La condición es suficiente, como se ve considerando, para todo $y\in Y$, una extensión algebraicamente cerrada $\Omega$ de $\boldsymbol{k}(y)$ y el diagrama conmutativo
\[
  \xymatrix{
    X\ar[dd]_f\\
    & \operatorname{Spec}(\Omega)\ar[ul]\ar[dl]\\
    Y
  }
\]
(cf. \sref{I.3.5.3}).
Recíprocamente, supongamos que $f$ es sobreyectivo y sea $g:\{\xi\}=\operatorname{Spec}(\Omega)\to Y$ un morfismo, donde $\Omega$ es un cuerpo algebraicamente cerrado.
Si se considera el diagrama
\[
  \xymatrix{
    X\ar[d]_f &
    X_{(\Omega)}\ar[l]\ar[d]^{f_{(\Omega)}}\\
    Y &
    \operatorname{Spec}(\Omega)\ar[l]
  }
\]
basta demostrar entonces que existe un \emph{punto racional sobre $\Omega$} en $X_{(\Omega)}$ (\sref{I.3.3.14}, \sref{I.3.4.3} y \sref{I.3.4.4}).
Como $f$ es sobreyectivo, $X_{(\Omega)}$ no es vacío \sref{I.3.5.10}, y como $f$ es de tipo finito, también lo es $f_{(\Omega)}$ \sref{I.6.3.4}[iv];
por consiguiente, $X_{(\Omega)}$ contiene un abierto afín no vacío $Z$ tal que $A(Z)$ es un álgebra no nula de tipo finito sobre $\Omega$.
En virtud del teorema de los ceros de Hilbert~\cite{I-21}, existe un $\Omega$-homomorfismo $A(Z)\to\Omega$ y, por tanto, una sección de $X_{(\Omega)}$ sobre $\operatorname{Spec}(\Omega)$, lo que demuestra la proposición.
\end{proof}

\subsection{Preesquemas algebraicos}
\label{subsection:I.6.4}
\label{I.6.4}

\begin{definition}[6.4.1]
\label{I.6.4.1}
{\itshape Dado un cuerpo $K$, se llama $K$-preesquema algebraico a un preesquema $X$ de tipo finito sobre $K$; $K$ se llama el cuerpo de base de $X$.
Si además $X$ es un esquema \textup{(o, lo que es equivalente \sref{I.5.5.8}, si $X$ es un $K$-\emph{esquema})}, se dice también que $X$ es un $K$-esquema algebraico.}
\end{definition}

Todo $K$-preesquema algebraico es \emph{noetheriano} \sref{I.6.3.7}.

\begin{proposition}[6.4.2]
\label{I.6.4.2}
Sea $X$ un $K$-preesquema algebraico.
Para que un punto $x\in X$ sea cerrado, es necesario y suficiente que $\boldsymbol{k}(x)$ sea una extensión algebraica de $K$ de grado finito.
\end{proposition}

\begin{proof}
Se puede suponer que $X$ es afín, siendo el anillo $A$ de $X$ una $K$-álgebra de tipo finito.
En efecto, los abiertos afines $U$ de $X$ tales que $A(U)$ es un $K$-álgebra de tipo finito forman un recubrimiento de $X$ \sref{I.6.3.1}.
Los puntos cerrados de $X$ son entonces los puntos tales que $\mathfrak{j}_x$ es
\oldpage[I]{148}
un ideal maximal de $A$, es decir, tales que $A/\mathfrak{j}_x$ es un cuerpo (necesariamente igual a $\boldsymbol{k}(x)$).
Como $A/\mathfrak{j}_x$ es un $K$-álgebra de tipo finito, se ve que, si $x$ es cerrado, $\boldsymbol{k}(x)$ es un cuerpo que es un álgebra de tipo finito sobre $K$ y, por tanto, necesariamente un $K$-álgebra de \emph{rango finito}~\cite{I-21}.
Recíprocamente, si $\boldsymbol{k}(x)$ es de rango finito sobre $K$, también lo es $A/\mathfrak{j}_x\subset \boldsymbol{k}(x)$, y como todo anillo íntegro que es además un $K$-álgebra de rango finito es un cuerpo, se tiene $A/\mathfrak{j}_x=\boldsymbol{k}(x)$ y, por consiguiente, $x$ es cerrado.
\end{proof}

\begin{corollary}[6.4.3]
\label{I.6.4.3}
Sean $K$ un cuerpo algebraicamente cerrado y $X$ un $K$-preesquema algebraico; los puntos cerrados de $X$ son entonces los puntos racionales sobre $K$ \sref{I.3.4.4} y se identifican canónicamente con los puntos de $X$ con valores en $K$.
\end{corollary}

\begin{proposition}[6.4.4]
\label{I.6.4.4}
Sea $X$ un preesquema algebraico sobre un cuerpo $K$.
Las propiedades siguientes son equivalentes.
\begin{enumerate}
  \item[{\textup{a)}}] $X$ es artiniano.
  \item[{\textup{b)}}] El espacio subyacente a $X$ es discreto.
  \item[{\textup{c)}}] El espacio subyacente a $X$ solo tiene un número finito de puntos cerrados.
  \item[{\textup{c')}}] El espacio subyacente a $X$ es finito.
  \item[{\textup{d)}}] Los puntos de $X$ son cerrados.
  \item[{\textup{e)}}] $X$ es isomorfo a $\operatorname{Spec}(A)$, donde $A$ es un $K$-álgebra de rango finito.
\end{enumerate}
\end{proposition}

\begin{proof}
Como $X$ es noetheriano, resulta de \sref{I.6.2.2} que las condiciones \emph{a)}, \emph{b)} y \emph{d)} son equivalentes e implican \emph{c)} y \emph{c')};
por otra parte, es claro que \emph{e)} implica \emph{a)}.
Queda por ver que \emph{c)} implica \emph{d)} y \emph{e)};
se puede limitar al caso en que $X$ es afín.
Entonces $A(X)$ es un $K$-álgebra de tipo finito \sref{I.6.3.3} y, por tanto, un anillo de Jacobson (\cite{I-1}, p.~3-11 y 3-12), en el que solo hay, por hipótesis, un número finito de ideales maximales.
Como una intersección finita de ideales primos solo puede ser un ideal primo si es igual a uno de los ideales primos que se intersectan, todo ideal primo de $A(X)$ es, pues, maximal, de donde \emph{d)}.
Además, se sabe entonces por \sref{I.6.2.2} que $A(X)$ es un $K$-álgebra artiniana de tipo finito y, por tanto, necesariamente de \emph{rango finito}~\cite{I-21}.
\end{proof}

\begin{env}[6.4.5]
\label{I.6.4.5}
Cuando se satisfacen las condiciones de \sref{I.6.4.4}, se dice que $X$ es un esquema \emph{finito sobre $K$} (cf. \sref[II]{II.6.1.1}), o un \emph{$K$-esquema finito}, de \emph{rango} $[A:K]$, que se denota también por $\mathit{rg}_K(X)$;
si $X$ e $Y$ son $K$-esquemas finitos, se tiene
\[
  \mathit{rg}_K(X\amalg Y)
  =\mathit{rg}_K(X)+\mathit{rg}_K(Y)
  \tag{6.4.5.1}\label{I.6.4.5.1-fr}
\]
\[
  \mathit{rg}_K(X\times_K Y)
  =\mathit{rg}_K(X)\mathit{rg}_K(Y)
  \tag{6.4.5.2}\label{I.6.4.5.2-fr}
\]
como resulta de \sref{I.3.2.2}.
\end{env}

\begin{corollary}[6.4.6]
\label{I.6.4.6}
Sea $X$ un $K$-esquema finito.
Para toda extensión $K'$ de $K$, $X\otimes_K K'$ es un $K'$-esquema finito, y su rango sobre $K'$ es igual al rango de $X$ sobre $K$.
\end{corollary}

\begin{proof}
Si $A=A(X)$, se tiene $[A\otimes_K K':K']=[A:K]$.
\end{proof}

\begin{corollary}[6.4.7]
\label{I.6.4.7}
Sea $X$ un esquema finito sobre un cuerpo $K$;
se pone $n=\sum_{x\in X}[\boldsymbol{k}(x):K]_s$ \textup{(se recuerda que, si $K'$ es una extensión de $K$, entonces $[K':K]_s$ es el \emph{rango separable} de $K'$ sobre $K$, rango de la mayor extensión algebraica separable de $K$ contenida en $K'$)};
\oldpage[I]{149}
entonces,
para toda extensión algebraicamente cerrada $\Omega$ de $K$, el espacio subyacente a $X\otimes_K\Omega$ tiene exactamente $n$ puntos, que se identifican con los puntos de $X$ con valores en $\Omega$.
\end{corollary}

\begin{proof}
Se puede limitar evidentemente al caso en que el anillo $A=A(X)$ es \emph{local} \sref{I.6.2.2};
sean $\mathfrak m$ su ideal maximal y $L=A/\mathfrak m$ su cuerpo residual, extensión algebraica de $K$.
Los puntos de $X$ con valores en $\Omega$ corresponden entonces biyectivamente a las $\Omega$-secciones de $X\otimes_K\Omega$ (\sref{I.3.4.1} y \sref{I.3.3.14}), y también a los $K$-homomorfismos de $L$ en $\Omega$ (\sref{I.1.7.3}), de donde la proposición (Bourbaki, \emph{Alg.}, cap.~V, §7, n\textsuperscript{o}~5, prop.~8), teniendo en cuenta \sref{I.6.4.3}.
\end{proof}

\begin{env}[6.4.8]
\label{I.6.4.8}
El número $n$ definido en \sref{I.6.4.7} se llama el \emph{rango separable} de $A$ (o de $X$) sobre $K$, o también el \emph{número geométrico de puntos} de $X$;
es igual al número de elementos de $X(\Omega)_K$.
Resulta inmediatamente de esta definición que, para toda extensión $K'$ de $K$, $X\otimes_K K'$ tiene el mismo número geométrico de puntos que $X$.
Si se denota este número por $n(X)$, es claro que, si $X$ e $Y$ son dos esquemas finitos sobre $K$, se tiene
\[
  n(X\amalg Y)=n(X)+n(Y).
  \tag{6.4.8.1}\label{I.6.4.8.1-fr}
\]
Bajo las mismas hipótesis, se tiene también
\[
  n(X\times_K Y)=n(X)n(Y)
  \tag{6.4.8.2}\label{I.6.4.8.2-fr}
\]
como resulta inmediatamente de la interpretación de $n(X)$ como número de elementos de $X(\Omega)_K$ y de la fórmula~\hyperref[I.3.4.3.1]{(3.4.3.1)}.
\end{env}

\begin{proposition}[6.4.9]
\label{I.6.4.9}
Sean $K$ un cuerpo, $X$ e $Y$ dos $K$-preesquemas algebraicos, $f:X\to Y$ un $K$-morfismo y $\Omega$ una extensión algebraicamente cerrada de $K$, de grado de trascendencia infinito sobre $K$.
Para que $f$ sea sobreyectivo, es necesario y suficiente que la aplicación $X(\Omega)_K\to Y(\Omega)_K$ correspondiente a $f$ \sref{I.3.4.1} sea sobreyectiva.
\end{proposition}

\begin{proof}
La necesidad resulta de \sref{I.6.3.10}, observando que $f$ es necesariamente de tipo finito \sref{I.6.3.9}.
Para ver que la condición es suficiente, se razona como en \sref{I.6.3.10}, observando que, para todo $y\in Y$, $\boldsymbol{k}(y)$ es una extensión de $K$ de tipo finito y, por tanto, es $K$-isomorfa a un subcuerpo de $\Omega$.
\end{proof}

\begin{remark}[6.4.10]
\label{I.6.4.10}
Veremos en el capítulo~IV que la conclusión de \sref{I.6.4.9} sigue siendo válida sin hipótesis relativa al grado de trascendencia de $\Omega$ sobre $K$.
\end{remark}

\begin{proposition}[6.4.11]
\label{I.6.4.11}
Si $f:X\to Y$ es un morfismo de tipo finito, entonces, para todo $y\in Y$, la fibra $f^{-1}(y)$ es un preesquema algebraico sobre el cuerpo residual $\boldsymbol{k}(y)$ y, para todo $x\in f^{-1}(y)$, $\boldsymbol{k}(x)$ es una extensión de $\boldsymbol{k}(y)$ de tipo finito.
\end{proposition}

\begin{proof}
Como $f^{-1}(y)=X\otimes_Y \boldsymbol{k}(y)$ \sref{I.3.6.3}, la proposición resulta de \sref{I.6.3.4}[iv] y de \sref{I.6.3.3}.
\end{proof}

\begin{proposition}[6.4.12]
\label{I.6.4.12}
Sean $f:X\to Y$ y $g:Y'\to Y$ dos morfismos;
pongamos $X'=X\times_Y Y'$ y sea $f'=f_{(Y')}:X'\to Y'$.
Sea $y'\in Y'$ y pongamos $y=g(y')$;
si la fibra $f^{-1}(y)$ es un esquema algebraico finito sobre $\boldsymbol{k}(y)$, entonces la fibra $f'^{-1}(y')$ es un esquema algebraico finito sobre $\boldsymbol{k}(y')$, con el mismo rango y el mismo número geométrico de puntos que $f^{-1}(y)$.
\end{proposition}

\begin{proof}
Teniendo en cuenta la transitividad de las fibras \sref{I.3.6.5}, esto resulta inmediatamente de \sref{I.6.4.6} y de \sref{I.6.4.8}.
\end{proof}

\begin{env}[6.4.13]
\label{I.6.4.13}
La Proposición~\sref{I.6.4.11} muestra que los morfismos de tipo finito corresponden intuitivamente a ``familias algebraicas de variedades algebraicas'', desempeñando los puntos de $Y$
\oldpage[I]{150}
el papel de ``parámetros'', lo que confiere a estos morfismos una significación ``geométrica''.
Los morfismos que no son de tipo finito intervendrán más adelante principalmente en cuestiones de ``cambio del preesquema de base'', por ejemplo por localización o completación.
\end{env}

\subsection{Determinación local de un morfismo}
\label{subsection:I.6.5}
\label{I.6.5}

\begin{proposition}[6.5.1]
\label{I.6.5.1}
Sean $X$ e $Y$ dos $S$-preesquemas, siendo $Y$ de tipo finito sobre $S$;
sean $x\in X$ e $y\in Y$ situados sobre un mismo punto $s\in S$.
\begin{enumerate}
  \item[{\rm(i)}] Si dos $S$-morfismos $f=(\psi,\theta)$ y $f'=(\psi',\theta')$ de $X$ en $Y$ son tales que $\psi(x)=\psi'(x)=y$, y los $\mathscr{O}_s$-homomorfismos (locales) $\theta_x^\sharp$ y ${\theta'}_x^\sharp$ de $\mathscr{O}_y$ en $\mathscr{O}_x$ son idénticos, entonces $f$ y $f'$ coinciden en una vecindad abierta de $x$.
  \item[{\rm(ii)}] Supongamos además que $S$ es localmente noetheriano.
    Para todo $\mathscr{O}_s$-homomorfismo local $\varphi:\mathscr{O}_y\to\mathscr{O}_x$, existe una vecindad abierta $U$ de $x$ en $X$ y un $S$-morfismo $f=(\psi,\theta)$ de $U$ en $Y$ tales que $\psi(x)=y$ y $\theta_x^\sharp=\varphi$.
\end{enumerate}
\end{proposition}

\begin{proof}
\medskip\noindent
\begin{enumerate}
  \item[(i)] Como la cuestión es local sobre $S$, $X$ e $Y$, se puede suponer que $S$, $X$ e $Y$ son afines, de anillos respectivos $A$, $B$ y $C$, y que $f$ y $f'$ son respectivamente de la forma $({}^a\varphi,\widetilde{\varphi})$ y $({}^a\varphi',\widetilde{\varphi}')$, donde $\varphi$ y $\varphi'$ son dos $A$-homomorfismos de $C$ en $B$ tales que $\varphi^{-1}(\mathfrak{j}_x)={\varphi'}^{-1}(\mathfrak{j}_x)
    =\mathfrak{j}_y$, y los homomorfismos $\varphi_x$ y $\varphi'_x$ de $C_y$ en $B_x$, deducidos de $\varphi$ y $\varphi'$, son idénticos;
    se puede suponer además que $C$ es un $A$-álgebra \emph{de tipo finito}.
    Sean $c_i$ ($1\leq i\leq n$) generadores del $A$-álgebra $C$, y pongamos $b_i=\varphi(c_i)$ y $b'_i=\varphi'(c_i)$;
    por hipótesis, se tiene $b_i/1=b'_i/1$ en el anillo de fracciones $B_x$ ($1\leq i\leq n$).
    Esto significa que existen elementos $s_i\in B\setmin\mathfrak{j}_x$ tales que $s_i(b_i-b'_i)=0$ para $1\leq i\leq n$, y se puede evidentemente suponer que todos los $s_i$ son iguales a un mismo elemento $g\in B\setmin\mathfrak{j}_x$.
    Se concluye de ello que se tiene $b_i/1=b'_i/1$ para $1\leq i\leq n$ en el anillo de fracciones $B_g$;
    si $i_g$ es el homomorfismo canónico $B\to B_g$, se tiene entonces $i_g\circ\varphi=i_g\circ\varphi'$;
    por tanto, las restricciones de $f$ y $f'$ a $D(g)$ son idénticas.
  \item[(ii)] Se puede reducir a la misma situación que en (i), y suponer además que el anillo $A$ es \emph{noetheriano}.
    Sean $c_i$ ($1\leq i\leq n$) generadores del $A$-álgebra $C$, y sea $\alpha:A[X_1,\ldots,X_n]\to C$ el homomorfismo del álgebra de polinomios $A[X_1,\ldots,X_n]$ sobre $C$ que transforma $X_i$ en $c_i$ para $1\leq i\leq n$.
    Sea además $i_y$ el homomorfismo canónico $C\to C_y$, y consideremos el homomorfismo compuesto
    \[
      \beta:A[X_1,\ldots,X_n]
      \xrightarrow{\alpha} C
      \xrightarrow{i_y} C_y
      \xrightarrow{\varphi} B_x.
    \]
    Denotemos por $\mathfrak{a}$ el núcleo de $\beta$;
    como $A$ es noetheriano, también lo es $A[X_1,\ldots,X_n]$, y, por tanto, $\mathfrak{a}$ admite un sistema finito de generadores $Q_j(X_1,\ldots,X_n)$ ($1\leq j\leq m$).
    Por otra parte, cada uno de los elementos $\varphi(i_y(c_i))$ puede escribirse en la forma $b_i/s_i$, donde $b_i\in B$ y $s_i\notin\mathfrak{j}_x$;
    se puede suponer además que todos los $s_i$ son iguales a un mismo elemento $g\in B\setmin\mathfrak{j}_x$.
    En estas condiciones, se tiene por hipótesis $Q_j(b_1/g,\ldots,b_n/g)=0$ en $B_x$; pongamos
    \[
      Q_j(X_1/T,\ldots,X_n/T)
      =P_j(X_1,\ldots,X_n,T)/T^{k_j}
    \]
    donde $P_j$ es homogéneo de grado $k_j$.
    Sea entonces $d_j=P_j(b_1,\ldots,b_n,g)\in B$.
    Por hipótesis, se tiene $t_jd_j=0$ para algún $t_j\in B\setmin\mathfrak{j}_x$ ($1\leq j\leq m$), y se puede evidentemente suponer que todos los $t_j$
\oldpage[I]{151}
    son iguales a un mismo elemento $h\in B\setmin\mathfrak{j}_x$;
    se concluye de ello que $P_j(hb_1,\ldots,hb_n,hg)=0$ para $1\leq j\leq m$.
    En estas condiciones, consideremos el homomorfismo $\rho$ de $A[X_1,\ldots,X_n]$ en el anillo de fracciones $B_{hg}$ que transforma $X_i$ en $hb_i/hg$ ($1\leq i\leq n$);
    la imagen de $\mathfrak{a}$ por este homomorfismo es 0, y \emph{a fortiori} también lo es la imagen del núcleo $\alpha^{-1}(0)$ por $\rho$.
    Por tanto, $\rho$ se factoriza como $A[X_1,\ldots,X_n]\xrightarrow{\alpha}C\xrightarrow{\gamma}B_{hg}$, con $\gamma(c_i)=hb_i/hg$, y es claro que, si $i_x$ es el homomorfismo canónico $B_{hg}\to B_x$, entonces el diagrama
    \[
      \xymatrix{
        C\ar[r]^\gamma\ar[d]_{i_y} & B_{hg}\ar[d]^{i_x}\\
        C_y\ar[r]_\varphi & B_x
      }
      \tag{6.5.1.1}\label{I.6.5.1.1-fr}
    \]
    es conmutativo; se tiene, pues, $\varphi=\gamma_x$, y como $\varphi$ es un homomorfismo local, ${}^a\gamma(x)=y$;
    $f=({}^a\gamma,\widetilde{\gamma})$ es, por tanto, un $S$-morfismo de la vecindad $D(hg)$ de $x$ a $Y$, como se afirmaba en la proposición.
\end{enumerate}
\end{proof}

\begin{corollary}[6.5.2]
\label{I.6.5.2}
Bajo las hipótesis de \sref{I.6.5.1}[(ii)], si además $X$ es de tipo finito sobre $S$, se puede suponer que el morfismo $f$ es de tipo finito.
\end{corollary}

\begin{proof}
Esto resulta de \sref{I.6.3.6}.
\end{proof}

\begin{corollary}[6.5.3]
\label{I.6.5.3}
Supongamos satisfechas las hipótesis de \sref{I.6.5.1}[(ii)], y supongamos además que $Y$ es íntegro y que $\varphi$ es un homomorfismo inyectivo.
Entonces se puede suponer que $f=({}^a\gamma,\widetilde{\gamma})$ donde $\gamma$ es inyectivo.
\end{corollary}

\begin{proof}
En efecto, se puede suponer $C$ íntegro \sref{I.5.1.4}, y por tanto $i_y$ inyectivo; se deduce entonces del diagrama~\sref{I.6.5.1.1} que $\gamma$ es inyectivo.
\end{proof}

\begin{proposition}[6.5.4]
\label{I.6.5.4}
Sea $f=(\psi,\theta):X\to Y$ un morfismo de tipo finito, $x$ un punto de $X$ e $y=\psi(x)$.
\begin{enumerate}
  \item[{\rm(i)}] Para que $f$ sea una inmersión local en el punto $x$ \sref{I.4.5.1}, es necesario y suficiente que $\theta_x^\sharp:\mathscr{O}_y\to\mathscr{O}_x$ sea sobreyectivo.
  \item[{\rm(ii)}] Se supone además que $Y$ es localmente noetheriano.
    Para que $f$ sea un isomorfismo local en el punto $x$ \sref{I.4.5.2}, es necesario y suficiente que $\theta_x^\sharp$ sea un isomorfismo.
\end{enumerate}
\end{proposition}

\begin{proof}
\medskip\noindent
\begin{enumerate}
  \item[(ii)] En virtud de \sref{I.6.5.1}, existe una vecindad abierta $V$ de $y$ y un morfismo $g:V\to X$ tales que $g\circ f$ (resp. $f\circ g$) está definido y coincide con la identidad en una vecindad de $x$ (resp. $y$), de donde se deduce fácilmente que $f$ es un isomorfismo local.
  \item[(i)] Como la cuestión es local sobre $X$ e $Y$, se puede suponer que $X$ e $Y$ son afines, de anillos respectivos $A$ y $B$;
    se tiene $f=({}^a\varphi,\widetilde{\varphi})$, donde $\varphi$ es un homomorfismo de anillos $B\to A$ que hace de $A$ una $B$-álgebra de tipo finito;
    se tiene $\varphi^{-1}(\mathfrak{j}_x)=\mathfrak{j}_y$, y el homomorfismo $\varphi_x:B_y\to A_x$ inducido por $\varphi$ es \emph{sobreyectivo}.
    Sea $(t_i)$ ($1\leq i\leq n$) un sistema de generadores de la $B$-álgebra $A$;
    la hipótesis sobre $\varphi_x$ implica que existen $b_i\in B$ y un $c\in B\setmin\mathfrak{j}_y$ tales que, en el anillo de fracciones $A_x$, se tiene $t_i/1=\varphi(b_i)/\varphi(c)$ para $1\leq i\leq n$.
    Por consiguiente \sref{I.1.3.3}, existe un $a\in A\setmin\mathfrak{j}_x$ tal que, si se pone $g=a\varphi(c)$, se tiene también $t_i/1=a\varphi(b_i)/g$ \emph{en el anillo de fracciones $A_g$}.
    Dicho esto, existe por hipótesis un polinomio $Q(X_1,\ldots,X_n)$, con coeficientes en el anillo $\varphi(B)$,
\oldpage[I]{152}
    tal que $a=Q(t_1,\ldots,t_n)$;
    pongamos \[
      Q(X_1/T,\ldots,X_n/T)
      =P(X_1,\ldots,X_n,T)/T^m,
    \], donde $P$ es homogéneo de grado $m$.
    En el anillo $A_g$, se tiene
    \[
      a/1
      =a^mP(\varphi(b_1),\ldots,\varphi(b_n),\varphi(c))/g^m
      =a^m\varphi(d)/g^m
    \]
    donde $d\in B$.
    Como, en $A_g$, $g/1=(a/1)(\varphi(c)/1)$ es invertible por definición, también lo son $a/1$ y $\varphi(c)/1$, y se puede por tanto escribir \[
      a/1=(\varphi(d)/1)(\varphi(c)/1)^{-m}.
    \].
    De aquí se concluye que $\varphi(d)/1$ también es invertible en $A_g$.
    Pongamos entonces $h=cd$;
    como $\varphi(h)/1$ es invertible en $A_g$, el homomorfismo compuesto \[
      B\xrightarrow{\varphi}A\longrightarrow A_g
    \] se factoriza como \[
      B\longrightarrow B_h\xrightarrow{\gamma}A_g
    \] \sref[0]{0.1.2.4}.
    Mostremos que $\gamma$ es \emph{sobreyectivo};
    basta comprobar que la imagen de $B_h$ en $A_g$ contiene los $t_i/1$ y $(g/1)^{-1}$.
    Ahora bien, se tiene \[
      (g/1)^{-1}
      =(\varphi(c)/1)^{m-1}(\varphi(d)/1)^{-1}
      =\gamma(c^m/h),
    \], y \[
      a/1=\gamma(d^{m+1}/h^m),
    \], por lo que \[
      (a\varphi(b_i))/1=\gamma(b_i d^{m+1}/h^m),
    \], y como $t_i/1=(a\varphi(b_i)/1)(g/1)^{-1}$, queda demostrada nuestra afirmación.
    La elección de $h$ implica que $\psi(D(g))\subset D(h)$, y se sabe también que la restricción de $f$ a $D(g)$ es igual a $({}^a\gamma,\widetilde{\gamma})$;
    como $\gamma$ es sobreyectivo, esta restricción es una inmersión cerrada de $D(g)$ en $D(h)$ \sref{I.4.2.3}.
\end{enumerate}
\end{proof}

\begin{corollary}[6.5.5]
\label{I.6.5.5}
Sea $f=(\psi,\theta):X\to Y$ un morfismo de tipo finito.
Se supone que $X$ es irreducible, se denota por $x$ su punto genérico y se pone $y=\psi(x)$.
\begin{enumerate}
  \item[{\rm(i)}] Para que $f$ sea una inmersión local en algún punto de $X$, es necesario y suficiente que $\theta_x^\sharp:\mathscr{O}_y\to\mathscr{O}_x$ sea sobreyectivo.
  \item[{\rm(ii)}] Se supone además que $Y$ es irreducible y localmente noetheriano.
    Para que $f$ sea un isomorfismo local en algún punto de $X$, es necesario y suficiente que $y$ sea el punto genérico de $Y$ \textup{(o, equivalentemente \sref[0]{0.2.1.4}, que $f$ sea un morfismo \emph{dominante})} y que $\theta_x^\sharp$ sea un isomorfismo \textup{(en otras palabras, que $f$ sea \emph{birracional} \sref{I.2.2.9})}.
\end{enumerate}
\end{corollary}

\begin{proof}
Es claro que (i) resulta de \sref{I.6.5.4}[i], teniendo en cuenta que todo abierto no vacío de $X$ contiene a $x$;
análogamente, (ii) resulta de \sref{I.6.5.4}[ii].
\end{proof}

\subsection{Morfismos cuasicompactos y morfismos localmente de tipo finito}
\label{subsection:I.6.6}

\begin{definition}[6.6.1]
\label{I.6.6.1}
{\itshape Se dice que un morfismo $f:X\to Y$ es cuasicompacto si la imagen inversa de todo abierto cuasicompacto de $Y$ por $f$ es cuasicompacta.}
\end{definition}

Sea $\mathfrak{B}$ una base de la topología de $Y$ formada por abiertos cuasicompactos (por ejemplo, abiertos afines);
para que $f$ sea cuasicompacto, es necesario y suficiente que la imagen inversa de todo elemento de $\mathfrak{B}$ por $f$ sea cuasicompacta (o, equivalentemente, una unión \emph{finita} de abiertos afines), pues todo abierto cuasicompacto de $Y$ es unión finita de elementos de $\mathfrak{B}$.
Por ejemplo, si $X$ es \emph{cuasicompacto} e $Y$ \emph{afín}, entonces \emph{todo} morfismo $f:X\to Y$ es cuasicompacto:
en efecto, $X$ es una unión finita de abiertos afines $U_i$, y para todo abierto afín $V$ de $Y$, $U_i\cap f^{-1}(V)$ es afín \sref{I.5.5.10}, y por tanto cuasicompacto.

Si $f:X\to Y$ es un morfismo cuasicompacto, es claro que, para todo abierto $V$ de $Y$, la restricción de $f$ a $f^{-1}(V)$ es un morfismo cuasicompacto $f^{-1}(V)\to V$.
Recíprocamente, si $(U_\alpha)$ es un recubrimiento abierto de $Y$, y $f:X\to Y$ un morfismo tal que las restricciones $f^{-1}(U_\alpha)\to U_\alpha$ sean cuasicompactas, entonces $f$ es cuasicompacto.

\begin{definition}[6.6.2]
\label{I.6.6.2}
{\itshape Se dice que un morfismo $f:X\to Y$ es localmente de tipo finito si, para todo $x\in X$, existen una vecindad abierta $U$ de $x$ y una vecindad abierta $V\supset f(U)$ de $f(x)$\footnote{\emph{[Trad.] La fuente francesa imprime aquí «de $y$» sin definir $y$; la lectura corregida explicita $y=f(x)$.}} tales que la
\oldpage[I]{153}
restricción de $f$ a $U$ sea un morfismo de tipo finito de $U$ en $V$.
Se dice también entonces que $X$ es un preesquema localmente de tipo finito sobre $Y$, o un $Y$-preesquema localmente de tipo finito.}
\end{definition}

De \sref{I.6.3.2} resulta inmediatamente que, si $f$ es localmente de tipo finito, entonces, para todo abierto $W$ de $Y$, la restricción de $f$ a $f^{-1}(W)$ es un morfismo $f^{-1}(W)\to W$ localmente de tipo finito.

Si $Y$ es localmente noetheriano y $X$ localmente de tipo finito sobre $Y$, entonces $X$ es localmente noetheriano en virtud de \sref{I.6.3.7}.

\begin{proposition}[6.6.3]
\label{I.6.6.3}
Para que un morfismo $f:X\to Y$ sea de tipo finito, es necesario y suficiente que sea cuasicompacto y localmente de tipo finito.
\end{proposition}

\begin{proof}
La necesidad de las condiciones es inmediata, teniendo en cuenta \sref{I.6.3.1} y la observación que sigue a \sref{I.6.6.1}.
Recíprocamente, supongamos satisfechas estas condiciones y sea $U$ un abierto afín de $Y$, de anillo $A$;
para todo $x\in f^{-1}(U)$, existe por hipótesis una vecindad $V(x)\subset f^{-1}(U)$ de $x$ y una vecindad $W(x)\subset U$ de $y=f(x)$ que contiene $f(V(x))$ y tal que la restricción de $f$ a $V(x)$ sea un morfismo $V(x)\to W(x)$ de tipo finito.
Reemplazando $W(x)$ por una vecindad $W_1(x)\subset W(x)$ de $y$\footnote{\emph{[Trad.] La fuente francesa imprime aquí «de $x$». Como $W_1(x)\subset W(x)\subset U\subset Y$, debe tratarse de una vecindad de $y=f(x)$.}} de la forma $D(g)$ (con $g\in A$), y $V(x)$ por $V(x)\cap f^{-1}(W_1(x))$, se puede suponer que $W(x)$ es de la forma $D(g)$, y por tanto de tipo finito sobre $U$ (pues su anillo se escribe $A[1/g]$);
por consiguiente, $V(x)$ es de tipo finito sobre $U$.
Además, $f^{-1}(U)$ es cuasicompacto por hipótesis y, por tanto, es una unión de un número finito de abiertos $V(x_i)$, lo que termina la demostración.
\end{proof}

\begin{proposition}[6.6.4]
\label{I.6.6.4}
\medskip\noindent
\begin{enumerate}
  \item[{\rm(i)}] Una inmersión $X\to Y$ es cuasicompacta si es cerrada, o si el espacio subyacente a $Y$ es localmente noetheriano, o si el espacio subyacente a $X$ es noetheriano.
  \item[{\rm(ii)}] El compuesto de dos morfismos cuasicompactos es cuasicompacto.
  \item[{\rm(iii)}] Si $f:X\to Y$ es un $S$-morfismo cuasicompacto, entonces $f_{(S')}:X_{(S')}\to Y_{(S')}$ también lo es para toda extensión $g:S'\to S$ del preesquema de base.
  \item[{\rm(iv)}] Si $f:X\to X'$ y $g:Y\to Y'$ son dos $S$-morfismos cuasicompactos, entonces $f\times_S g$ es cuasicompacto.
  \item[{\rm(v)}] Si el compuesto de dos morfismos $f:X\to Y$ y $g:Y\to Z$ es cuasicompacto, y si $g$ es separado o el espacio subyacente a $X$ es localmente noetheriano, entonces $f$ es cuasicompacto.
  \item[{\rm(vi)}] Para que un morfismo $f$ sea cuasicompacto, es necesario y suficiente que $f_{\mathrm{red}}$ lo sea.
\end{enumerate}
\end{proposition}

\begin{proof}
Observemos que (vi) es evidente, pues la propiedad de ser cuasicompacto para un morfismo solo depende de la aplicación continua correspondiente entre los espacios subyacentes.
Demostremos asimismo la parte de (v) correspondiente al caso en que el espacio subyacente a $X$ es localmente noetheriano.
Pongamos $h=g\circ f$, y sea $U$ un abierto cuasicompacto de $Y$;
$g(U)$ es cuasicompacto (pero no necesariamente abierto) en $Z$, y por tanto está contenido en una unión finita de abiertos cuasicompactos $V_j$ \sref{I.2.1.3}, y $f^{-1}(U)$ está entonces contenido en la unión de los $h^{-1}(V_j)$, que son subespacios cuasicompactos de $X$ y, por tanto, subespacios noetherianos.
Se concluye pues \sref[0]{0.2.2.3} que $f^{-1}(U)$ es un espacio noetheriano y, \emph{a fortiori}, cuasicompacto.

Para demostrar las demás aserciones, basta demostrar (i), (ii) y (iii) \sref{I.5.5.12}.
Ahora bien, (ii) es evidente, y (i) resulta de \sref{I.6.3.5} cuando el espacio $Y$ es localmente noetheriano o el espacio $X$ es noetheriano, y es evidente para una inmersión cerrada.
Para establecer (iii),
\oldpage[I]{154}
se puede limitar al caso en que $S=Y$ \sref{I.3.3.11};
pongamos $f'=f_{(S')}$, y sea $U'$ un abierto cuasicompacto de $S'$.
Para todo $s'\in U'$, sea $T$ una vecindad abierta afín de $g(s')$ en $S$, y sea $W$ una vecindad abierta afín de $s'$ contenida en $U'\cap g^{-1}(T)$;
bastará demostrar que $f'^{-1}(W)$ es cuasicompacto;
dicho de otro modo, se puede reducir a demostrar que, cuando $S$ y $S'$ son \emph{afines}, el espacio subyacente a $X\times_S S'$ es cuasicompacto.
Pero como $X$ es entonces, por hipótesis, una unión finita de abiertos afines $V_j$, $X\times_S S'$ es unión de los espacios subyacentes a los esquemas afines $V_j\times_S S'$ (\sref{I.3.2.2} y \sref{I.3.2.7}), lo que demuestra la proposición.
\end{proof}

Se observará también que, si $X=X'\amalg X''$ es la suma de dos preesquemas, un morfismo $f:X\to Y$ es cuasicompacto si y solo si sus restricciones a $X'$ y $X''$ son cuasicompactas.

\begin{proposition}[6.6.5]
\label{I.6.6.5}
Sea $f:X\to Y$ un morfismo cuasicompacto.
Para que $f$ sea dominante, es necesario y suficiente que, para todo punto genérico $y$ de una componente irreducible de $Y$, $f^{-1}(y)$ contenga un punto genérico de una componente irreducible de $X$.
\end{proposition}

\begin{proof}
Es inmediato que la condición es suficiente (incluso sin suponer que $f$ sea cuasicompacto).
Para ver que es necesaria, consideremos una vecindad abierta afín $U$ de $y$;
$f^{-1}(U)$ es cuasicompacto y, por tanto, una unión \emph{finita} de abiertos afines $V_i$, y la hipótesis de que $f$ sea dominante implica que $y$ pertenece a la adherencia \emph{en} $U$ de uno de los $f(V_i)$.
Se puede suponer evidentemente que $X$ e $Y$ son reducidos;
como la adherencia en $X$ de una componente irreducible de $V_i$ es una componente irreducible de $X$ \sref[0]{0.2.1.6}, se puede reemplazar $X$ por $V_i$, e $Y$ por el subpreesquema cerrado reducido de $U$ cuyo espacio subyacente es $\overline{f(V_i)}\cap U$ \sref{I.5.2.1}, y se reduce así a demostrar la proposición cuando $X=\operatorname{Spec}(A)$ e $Y=\operatorname{Spec}(B)$ son afines y reducidos.
Como $f$ es dominante, $B$ es un subanillo de $A$ \sref{I.1.2.7}, y la proposición resulta entonces del hecho de que todo ideal primo minimal de $B$ es la intersección con $B$ de un ideal primo minimal de $A$ \sref[0]{0.1.5.8}.
\end{proof}

\begin{proposition}[6.6.6]
\label{I.6.6.6}
\medskip\noindent
\begin{enumerate}
  \item[{\rm(i)}] Toda inmersión local es localmente de tipo finito.
  \item[{\rm(ii)}] Si dos morfismos $f:X\to Y$ y $g:Y\to Z$ son localmente de tipo finito, entonces $g\circ f$ también lo es.
  \item[{\rm(iii)}] Si $f:X\to Y$ es un $S$-morfismo localmente de tipo finito, entonces $f_{(S')}:X_{(S')}\to Y_{(S')}$ es localmente de tipo finito para toda extensión $S'\to S$ del preesquema de base.
  \item[{\rm(iv)}] Si $f:X\to X'$ y $g:Y\to Y'$ son $S$-morfismos localmente de tipo finito, entonces $f\times_S g$ es localmente de tipo finito.
  \item[{\rm(v)}] Si el compuesto $g\circ f$ de dos morfismos es localmente de tipo finito, entonces $f$ es localmente de tipo finito.
  \item[{\rm(vi)}] Si un morfismo $f$ es localmente de tipo finito, entonces $f_{\mathrm{red}}$ también lo es.
\end{enumerate}
\end{proposition}

\begin{proof}
En virtud de \sref{I.5.5.12}, basta demostrar (i), (ii) y (iii).
Si $j:X\to Y$ es una inmersión local, para todo $x\in X$ existe una vecindad abierta $V$ de $j(x)$ en $Y$ y una vecindad abierta $U$ de $x$ en $X$ tales que la restricción de $j$ a $U$ sea una inmersión cerrada $U\to V$ \sref{I.4.5.1}, y por tanto esta restricción es de tipo finito.
Para demostrar (ii), consideremos un punto $x\in X$;
por hipótesis, existe una vecindad abierta $W$ de $g(f(x))$ y una vecindad abierta $V$ de $f(x)$ tales que $g(V)\subset W$ y que $V$ es de tipo finito sobre $W$;
además, $f^{-1}(V)$ es localmente de tipo finito sobre $V$ \sref{I.6.6.2}, por lo que existe una vecindad abierta $U$ de $x$ que
\oldpage[I]{155}
está contenida en $f^{-1}(V)$ y es de tipo finito sobre $V$;
por consiguiente, se tiene $g(f(U))\subset W$, y $U$ es de tipo finito sobre $W$ \sref{I.6.3.4}[ii].
Finalmente, para demostrar (iii), se puede limitar al caso en que $Y=S$ \sref{I.3.3.11};
para todo $x'\in X'=X_{(S')}$, sean $x$ la imagen de $x'$ en $X$, $s$ la imagen de $x$ en $S$, $T$ una vecindad abierta de $s$, $T'$ su imagen inversa en $S'$ y $U$ una vecindad abierta de $x$ cuya imagen está contenida en $T$ y que es de tipo finito sobre $T$;
entonces $U\times_S T'=U\times_T T'$ es una vecindad abierta de $x'$ \sref{I.3.2.7} que es de tipo finito sobre $T'$ \sref{I.6.3.4}[iv].
\end{proof}

\begin{corollary}[6.6.7]
\label{I.6.6.7}
Sean $X$ e $Y$ dos $S$-preesquemas localmente de tipo finito sobre $S$.
Si $S$ es localmente noetheriano, entonces $X\times_S Y$ es localmente noetheriano.
\end{corollary}

\begin{proof}
En efecto, siendo $X$ localmente de tipo finito sobre $S$, es localmente noetheriano; además, $X\times_S Y$ es localmente de tipo finito sobre $X$ y, por tanto, también es localmente noetheriano.
\end{proof}

\begin{remark}[6.6.8]
\label{I.6.6.8}
La Proposición~\sref{I.6.3.10} y su demostración se extienden inmediatamente al caso en que se supone solamente que el morfismo $f$ es localmente de tipo finito.
Del mismo modo, las proposiciones~\sref{I.6.4.2} y \sref{I.6.4.9} siguen siendo válidas cuando se supone únicamente que los preesquemas $X$ e $Y$ que figuran en sus enunciados son localmente de tipo finito sobre el cuerpo $K$.
\end{remark}
% END source_es/ega1/ega1-6.tex

% BEGIN source_es/ega1/ega1-7.tex
\section{Aplicaciones racionales}
\label{section:I.7}

\subsection{Aplicaciones racionales y funciones racionales}
\label{subsection:I.7.1}

\begin{env}[7.1.1]
\label{I.7.1.1}
Sean $X$ e $Y$ dos preesquemas, $U$ y $V$ dos abiertos densos de $X$, y $f$ (resp. $g$) un morfismo de $U$ (resp. $V$) en $Y$; diremos que $f$ y $g$ son \emph{equivalentes} si coinciden en un abierto denso de $U\cap V$.
Como una intersección finita de abiertos densos de $X$ es un abierto denso de $X$, es claro que esta relación es una \emph{relación de equivalencia}.
\end{env}

\begin{definition}[7.1.2]
\label{I.7.1.2}
{\itshape Dados dos preesquemas $X$ e $Y$, se llama aplicación racional de $X$ en $Y$ a una clase de equivalencia de morfismos de abiertos densos de $X$ en $Y$, para la relación de equivalencia definida en \sref{I.7.1.1}.
Si $X$ e $Y$ son $S$-preesquemas, se dice que tal clase es una $S$-aplicación racional si existe un representante de la clase que sea un $S$-morfismo.
Se llama $S$-sección racional de $X$ a toda $S$-aplicación racional de $S$ en $X$.
Se llama función racional sobre un preesquema $X$ a toda $X$-sección racional del $X$-preesquema $X\otimes_{\mathbf{Z}}\mathbf{Z}[T]$ (donde $T$ es una indeterminada).}
\end{definition}
Por abuso de lenguaje, cuando se trate únicamente de $S$-preesquemas, se dirá ``aplicación racional'' en lugar de ``$S$-aplicación racional'' si no puede resultar ninguna confusión.

Sea $f$ una aplicación racional de $X$ en $Y$, y sea $U$ un abierto de $X$; si $f_1$ y $f_2$ son dos morfismos pertenecientes a la clase de $f$, definidos respectivamente en abiertos densos $V$ y $W$ de $X$, las restricciones $f_1|(U\cap V)$ y $f_2|(U\cap W)$ coinciden en un abierto de $U\cap V\cap W$, que es denso en $U$; la clase de morfismos $f$ define, pues, una aplicación racional de $U$ en $Y$, llamada la \emph{restricción de $f$ a $U$}, y denotada por $f|U$.

Si a todo $S$-morfismo $f:X\to Y$ se le hace corresponder la $S$-aplicación racional
\oldpage[I]{156}
a la que pertenece $f$, se obtiene una aplicación canónica de $\operatorname{Hom}_S(X,Y)$ en el conjunto de las $S$-aplicaciones racionales de $X$ en $Y$.
Se denota por $\Gamma_{\mathrm{rat}}(X/Y)$ el conjunto de las $Y$-secciones racionales de $X$, y se tiene, por tanto, una aplicación canónica $\Gamma(X/Y)\to\Gamma_{\mathrm{rat}}(X/Y)$.
Es claro además que, si $X$ e $Y$ son dos $S$-preesquemas, el conjunto de las $S$-aplicaciones racionales de $X$ en $Y$ se identifica canónicamente con $\Gamma_{\mathrm{rat}}((X\times_S Y)/X)$ \sref{I.3.3.14}.


\begin{env}[7.1.3]
\label{I.7.1.3}
Resulta también de \sref{I.7.1.2} y \sref{I.3.3.14} que las funciones racionales sobre $X$ se identifican canónicamente con las \emph{clases de equivalencia de secciones del haz estructural $\mathscr{O}_X$} sobre abiertos densos de $X$, siendo equivalentes dos de tales secciones si coinciden en algún abierto denso contenido en la intersección de sus conjuntos de definición.
En particular, resulta que las funciones racionales sobre $X$ forman un \emph{anillo} $R(X)$.
\end{env}

\begin{env}[7.1.4]
\label{I.7.1.4}
Cuando $X$ es un preesquema \emph{irreducible}, todo abierto no vacío es denso en $X$; se puede decir también que los abiertos no vacíos de $X$ son las \emph{vecindades abiertas del punto genérico} $x$ de $X$.
Decir que dos morfismos de abiertos no vacíos de $X$ en $Y$ son equivalentes significa entonces, en este caso, que tienen el \emph{mismo germen} en el punto $x$.
En otros términos, las aplicaciones racionales (resp. $S$-aplicaciones racionales) $X\to Y$ se identifican con los \emph{gérmenes de morfismos} (resp. \emph{$S$-morfismos}) de abiertos no vacíos de $X$ en $Y$ en el punto genérico $x$ de $X$.
En particular:
\end{env}

\begin{proposition}[7.1.5]
\label{I.7.1.5}
Si $X$ es un preesquema irreducible, el anillo $R(X)$ de las funciones racionales sobre $X$ se identifica canónicamente con el anillo local $\mathscr{O}_x$ del punto genérico $x$ de $X$.
Es un anillo local de dimensión~\textup{0} y, por tanto, un anillo local artiniano cuando $X$ es noetheriano; es un cuerpo cuando $X$ es íntegro y, cuando además $X$ es un esquema afín, se identifica con el cuerpo de fracciones de $A(X)$.
\end{proposition}

\begin{proof}
Por lo anterior y por la identificación de las funciones racionales con las secciones de $\mathscr{O}_X$ sobre un abierto denso, la primera afirmación no es sino la definición del tallo de un haz en un punto.
Para las demás afirmaciones, se puede limitar al caso en que $X$ es afín, de anillo $A$; entonces $\mathfrak{j}_x$ es el nilradical de $A$, y $\mathscr{O}_x$ es, por tanto, de dimensión~0; si $A$ es íntegro, entonces $\mathfrak{j}_x=(0)$, y $\mathscr{O}_x$ es, pues, el cuerpo de fracciones de $A$.
Finalmente, si $A$ es noetheriano, se sabe (\cite[p.~127, cor.~4]{I-11}) que $\mathfrak{j}_x$ es nilpotente y que $\mathscr{O}_x=A_x$ es artiniano.
\end{proof}

Si $X$ es \emph{íntegro}, el anillo $\mathscr{O}_z$ es íntegro para todo $z\in X$; todo abierto afín $U$ que contiene a $z$ contiene también a $x$, y $R(U)$, al ser igual al cuerpo de fracciones de $A(U)$, se identifica con $R(X)$; se concluye así que $R(X)$ puede identificarse también con el \emph{cuerpo de fracciones de $\mathscr{O}_z$}: la identificación canónica de $\mathscr{O}_z$ con un subanillo de $R(X)$ consiste en asociar a todo germen de una sección $s\in\mathscr{O}_z$ la única función racional sobre $X$, clase de una sección de $\mathscr{O}_X$ (necesariamente definida sobre un abierto denso) que tiene a $s$ como germen en el punto $z$.

\begin{env}[7.1.6]
\label{I.7.1.6}
Supongamos ahora que $X$ tiene un número \emph{finito} de componentes irreducibles $X_i$ ($1\leq i\leq n$) (lo que ocurre siempre que el espacio subyacente a $X$ sea \emph{noetheriano}); sea $X'_i$ el abierto de $X$ obtenido suprimiendo de $X_i$ la unión de los $X_j\cap X_i$ para $j\neq i$; $X'_i$ es irreducible, su punto genérico $x_i$ es el punto genérico de $X_i$, y los $X'_i$ son disjuntos dos a dos, siendo su unión densa en $X$ \sref[0]{0.2.1.6}.
Para
\oldpage[I]{157}
todo abierto denso $U$ de $X$, $U_i=U\cap X'_i$ es un abierto no vacío denso de $X'_i$, siendo los $U_i$ disjuntos dos a dos, y por tanto $U'=\bigcup_i U_i$ es denso en $X$.
Darse un morfismo de $U'$ en $Y$ equivale a darse, arbitrariamente, un morfismo de cada uno de los $U_i$ en $Y$.
\end{env}

Por tanto:
\begin{proposition}[7.1.7]
\label{I.7.1.7}
Sean $X$ e $Y$ dos preesquemas (\textup{resp.} $S$-preesquemas) tales que $X$ tenga un número finito de componentes irreducibles $X_i$, con puntos genéricos $x_i$ ($1\leq i\leq n$).
Si $R_i$ es el conjunto de los gérmenes de morfismos (\textup{resp.} $S$-morfismos) de abiertos de $X$ en $Y$ en el punto $x_i$, entonces el conjunto de las aplicaciones racionales (\textup{resp.} $S$-aplicaciones racionales) de $X$ en $Y$ se identifica con el producto de los $R_i$ ($1\leq i\leq n$).
\end{proposition}

\begin{corollary}[7.1.8]
\label{I.7.1.8}
Sea $X$ un preesquema noetheriano.
El anillo de las funciones racionales sobre $X$ es un anillo artiniano, cuyas componentes locales son los anillos $\mathscr{O}_{x_i}$ de los puntos genéricos $x_i$ de las componentes irreducibles de $X$.
\end{corollary}

\begin{corollary}[7.1.9]
\label{I.7.1.9}
Sean $A$ un anillo noetheriano y $X=\operatorname{Spec}(A)$.
Si $Q$ es el complemento de la unión de los ideales primos minimales de $A$, entonces el anillo de las funciones racionales sobre $X$ se identifica canónicamente con el anillo de fracciones $Q^{-1}A$.
\end{corollary}

Esto resultará del lema siguiente:
\begin{lemma}[7.1.9.1]
\label{I.7.1.9.1}
Para que un elemento $f\in A$ sea tal que $D(f)$ sea denso en $X$, es necesario y suficiente que $f\in Q$; todo abierto denso de $X$ contiene un abierto de la forma $D(f)$, donde $f\in Q$.
\end{lemma}

En efecto, supongamos demostrado el lema. Puesto que el anillo de secciones $\Gamma(D(f),\mathscr{O}_X)$ se identifica con $A_f$ \sref{I.1.3.6} y \sref{I.1.3.7}, el hecho de que los $D(f)$ con $f\in Q$ formen una familia cofinal en el conjunto ordenado (para $\supset$) de los abiertos densos de $X$, junto con la definición de las funciones racionales dada en \sref{I.7.1.1}, muestra que el anillo de las funciones racionales sobre $X$ se identifica con el límite inductivo de los $A_f$ para $f\in Q$ (para la relación de preorden ``$g$ es múltiplo de $f$''), es decir, con $Q^{-1}A$ \sref[0]{0.1.4.5}.

\begin{proof}
Para demostrar \hyperref[I.7.1.9.1]{(7.1.9.1)}, denotemos de nuevo por $X_i$ ($1\leq i\leq n$) las componentes irreducibles de $X$; si $D(f)$ es denso en $X$, entonces $D(f)\cap X_i\neq\varnothing$ para $1\leq i\leq n$, y recíprocamente; ahora bien, esto significa que $f\notin\mathfrak{p}_i$ para $1\leq i\leq n$, donde se pone $\mathfrak{p}_i=\mathfrak{j}(X_i)$, y como los $\mathfrak{p}_i$ son los ideales primos minimales de $A$ \hyperref[I.1.1.14]{(1.1.14)}, las condiciones $f\notin\mathfrak{p}_i$ ($1\leq i\leq n$) equivalen a $f\in Q$, de donde resulta la primera afirmación del lema.
Para la otra afirmación, si $U$ es un abierto denso de $X$, el complemento de $U$ es un conjunto de la forma $V(\mathfrak{a})$, donde $\mathfrak{a}$ es un ideal que no está contenido en ninguno de los $\mathfrak{p}_i$; por tanto, tampoco está contenido en su unión (\cite[p.~13]{I-10}), y existe así algún $f\in\mathfrak{a}$ que pertenece a $Q$; de donde $D(f)\subset U$, lo que concluye la demostración.
\end{proof}

\begin{env}[7.1.10]
\label{I.7.1.10}
Supongamos de nuevo que $X$ es irreducible, con punto genérico $x$.
Como todo abierto no vacío $U$ de $X$ contiene a $x$, y contiene por tanto también a todo $z\in X$ tal que $x\in\overline{\{z\}}$, todo morfismo $U\to Y$ puede componerse con el morfismo canónico $\operatorname{Spec}(\mathscr{O}_x)\to X$ \sref{I.2.4.1}; y dos morfismos en $Y$ definidos en sendos abiertos no vacíos de $X$ que coinciden en un abierto no vacío de $X$ dan, por composición, el mismo morfismo $\operatorname{Spec}(\mathscr{O}_x)\to Y$.
En otros términos, a toda aplicación racional de $X$ en $Y$ le corresponde un morfismo bien determinado $\operatorname{Spec}(\mathscr{O}_x)\to Y$.
\end{env}

\oldpage[I]{158}
\begin{proposition}[7.1.11]
\label{I.7.1.11}
Sean $X$ e $Y$ dos $S$-preesquemas; se supone que $X$ es irreducible, con punto genérico $x$, y que $Y$ es de tipo finito sobre $S$.
Dos $S$-aplicaciones racionales cualesquiera de $X$ en $Y$ a las que corresponde el mismo $S$-morfismo $\operatorname{Spec}(\mathscr{O}_x)\to Y$ son idénticas.
Si se supone además que $S$ es localmente noetheriano, entonces a todo $S$-morfismo de $\operatorname{Spec}(\mathscr{O}_x)$ en $Y$ le corresponde una única $S$-aplicación racional de $X$ en $Y$.
\end{proposition}

\begin{proof}
Teniendo en cuenta que todo abierto no vacío de $X$ es denso, esto resulta de \sref{I.6.5.1}.
\end{proof}

\begin{corollary}[7.1.12]
\label{I.7.1.12}
Se supone que $S$ es localmente noetheriano y que se satisfacen las demás hipótesis de \sref{I.7.1.11}.
Las $S$-aplicaciones racionales de $X$ en $Y$ se identifican entonces con los puntos del $S$-preesquema $Y$ con valores en el $S$-preesquema $\operatorname{Spec}(\mathscr{O}_x)$.
\end{corollary}

\begin{proof}
Esto no es sino \sref{I.7.1.11}, con la terminología introducida en \sref{I.3.4.1}.
\end{proof}

\begin{corollary}[7.1.13]
\label{I.7.1.13}
Se supone que se satisfacen las condiciones de \sref{I.7.1.12}.
Sea $s$ la imagen de $x$ en $S$.
Darse una $S$-aplicación racional de $X$ en $Y$ equivale a darse un punto $y$ de $Y$ situado sobre $s$, junto con un $\mathscr{O}_s$-homomorfismo local $\mathscr{O}_y\to\mathscr{O}_x=R(X)$.
\end{corollary}

\begin{proof}
Esto resulta de \sref{I.7.1.11} y \sref{I.2.4.4}.
\end{proof}

En particular:
\begin{corollary}[7.1.14]
\label{I.7.1.14}
Bajo las condiciones de \sref{I.7.1.12}, las $S$-aplicaciones racionales de $X$ en $Y$ dependen únicamente (para $Y$ dado) del $S$-preesquema $\operatorname{Spec}(\mathscr{O}_x)$ y, en particular, siguen siendo las mismas cuando se reemplaza $X$ por $\operatorname{Spec}(\mathscr{O}_z)$, para cualquier $z\in X$.
\end{corollary}

\begin{proof}
Como $z\in\overline{\{x\}}$, $x$ es el punto genérico de $Z=\operatorname{Spec}(\mathscr{O}_z)$, y $\mathscr{O}_{X,x}=\mathscr{O}_{Z,x}$.
\end{proof}

Cuando $X$ es íntegro, $R(X)=\mathscr{O}_x=\boldsymbol{k}(x)$ es un cuerpo \sref{I.7.1.5}; los corolarios anteriores se especializan entonces del modo siguiente:
\begin{corollary}[7.1.15]
\label{I.7.1.15}
Se supone que se satisfacen las condiciones de \sref{I.7.1.12} y, además, que $X$ es íntegro.
Sea $s$ la imagen de $x$ en $S$.
Entonces las $S$-aplicaciones racionales de $X$ en $Y$ se identifican con los puntos de $Y\otimes_S \boldsymbol{k}(s)$ con valores en la extensión $R(X)$ de $\boldsymbol{k}(s)$, o, en otros términos, cada una de ellas equivale a los datos de un punto $y\in Y$ situado sobre $s$ y de un $\boldsymbol{k}(s)$-monomorfismo de $\boldsymbol{k}(y)$ en $\boldsymbol{k}(x)=R(X)$.
\end{corollary}

\begin{proof}
En efecto, los puntos de $Y$ situados sobre $s$ se identifican con los puntos de $Y\otimes_S \boldsymbol{k}(s)$ \sref{I.3.6.3}, y los $\mathscr{O}_s$-homomorfismos locales $\mathscr{O}_y\to R(X)$ con los $\boldsymbol{k}(s)$-monomorfismos $\boldsymbol{k}(y)\to R(X)$.
\end{proof}

Más precisamente:
\begin{corollary}[7.1.16]
\label{I.7.1.16}
Sean $k$ un cuerpo, y $X$ e $Y$ dos preesquemas algebraicos sobre $k$ \sref{I.6.4.1}; se supone además que $X$ es íntegro.
Entonces las $k$-aplicaciones racionales de $X$ en $Y$ se identifican con los puntos de $Y$ con valores en la extensión $R(X)$ de $k$ \sref{I.3.4.4}.
\end{corollary}

\subsection{Dominio de definición de una aplicación racional}
\label{subsection:I.7.2}
\label{I.7.2}

\begin{env}[7.2.1]
\label{I.7.2.1}
Sean $X$ e $Y$ dos preesquemas, y $f$ una aplicación racional de $X$ en $Y$.
Se dice que $f$ está \emph{definida en un punto $x\in X$} si existe un abierto denso $U$ que contiene a $x$, y un morfismo $U\to Y$ perteneciente a la clase de equivalencia de $f$.
El conjunto de puntos $x\in X$ en los que $f$ está definida se llama \emph{dominio de definición} de $f$; es claro que es un abierto denso de $X$.
\end{env}

\oldpage[I]{159}
\begin{proposition}[7.2.2]
\label{I.7.2.2}
Sean $X$ e $Y$ dos $S$-preesquemas tales que $X$ sea reducido e $Y$ separado sobre $S$.
Sea $f$ una $S$-aplicación racional de $X$ en $Y$, de dominio de definición $U_0$.
Existe entonces un único $S$-morfismo $U_0\to Y$ perteneciente a la clase de $f$.
\end{proposition}

Como, para todo morfismo $U\to Y$ perteneciente a la clase de $f$, se tiene necesariamente $U\subset U_0$, es claro que la proposición será consecuencia de lo siguiente:

\begin{lemma}[7.2.2.1]
\label{I.7.2.2.1}
Bajo las hipótesis de \sref{I.7.2.2}, sean $U_1$ y $U_2$ dos abiertos densos de $X$, y sean $f_i:U_i\to Y$ ($i=1,2$) dos $S$-morfismos tales que existe un abierto $V\subset U_1\cap U_2$ denso en $X$, en el cual $f_1$ y $f_2$ coinciden.
Entonces $f_1$ y $f_2$ coinciden en $U_1\cap U_2$.
\end{lemma}

\begin{proof}
Se puede evidentemente reducir al caso en que $X=U_1=U_2$.
Como $X$ (y por tanto $V$) es reducido, $X$ es el menor subpreesquema cerrado de $X$ que contiene a $V$ \sref{I.5.2.2}.
Sea $g=(f_1,f_2)_S:X\to Y\times_S Y$; como, por hipótesis, la diagonal $T=\Delta_Y(Y)$ es un subpreesquema cerrado de $Y\times_S Y$, $Z=g^{-1}(T)$ es un subpreesquema cerrado de $X$ \sref{I.4.4.1}.
Si $h:V\to Y$ es la restricción común de $f_1$ y $f_2$ a $V$, la restricción de $g$ a $V$ es $g'=(h,h)_S$, que se factoriza como $g'=\Delta_Y\circ h$; como $\Delta_Y^{-1}(T)=Y$, se tiene $g'^{-1}(T)=V$, y por tanto $Z$ es un subpreesquema cerrado de $X$ que induce $V$, luego contiene a $V$, lo que implica que $Z=X$.
De la relación $g^{-1}(T)=X$, se deduce \sref{I.4.4.1} que $g$ se factoriza como $\Delta_Y\circ f$, donde $f$ es un morfismo $X\to Y$, lo que implica, por la definición del morfismo diagonal, que $f_1=f_2=f$.
\end{proof}

Es claro que el morfismo $U_0\to Y$ definido en \sref{I.7.2.2} es el único morfismo de la clase $f$ que \emph{no puede prolongarse} a un morfismo de un abierto de $X$ que contenga estrictamente a $U_0$.
\emph{Bajo las hipótesis de} \sref{I.7.2.2}, se pueden, pues, \emph{identificar} las aplicaciones racionales de $X$ en $Y$ con los morfismos \emph{no prolongables} (a abiertos estrictamente mayores) de abiertos densos de $X$ en $Y$.
Con esta identificación, la Proposición~\sref{I.7.2.2} implica:
\begin{corollary}[7.2.3]
\label{I.7.2.3}
Con las hipótesis de \sref{I.7.2.2} sobre $X$ e $Y$, sea $U$ un abierto denso de $X$.
Existe entonces una correspondencia biunívoca canónica entre los $S$-morfismos de $U$ en $Y$ y las $S$-aplicaciones racionales de $X$ en $Y$ definidas en todos los puntos de $U$.
\end{corollary}

\begin{proof}
En virtud de \sref{I.7.2.2}, para todo $S$-morfismo $f$ de $U$ en $Y$, existe una única $S$-aplicación racional $\overline f$ de $X$ en $Y$ que prolonga a $f$.
\end{proof}

\begin{corollary}[7.2.4]
\label{I.7.2.4}
Sean $S$ un esquema, $X$ un $S$-preesquema reducido, $Y$ un $S$-esquema, y $f:U\to Y$ un $S$-morfismo de un abierto denso $U$ de $X$ en $Y$.
Si $\overline f$ es la $\mathbf{Z}$-aplicación racional de $X$ en $Y$ que prolonga a $f$, entonces $\overline f$ es un $S$-morfismo (y es, por tanto, la $S$-aplicación racional de $X$ en $Y$ que prolonga a $f$).
\end{corollary}

\begin{proof}
En efecto, si $\varphi:X\to S$ y $\psi:Y\to S$ son los morfismos estructurales, $U_0$ es el dominio de definición de $\overline f$, y $j$ es la inyección $U_0\to X$, basta demostrar que $\psi\circ\overline f=\varphi\circ j$, pero esto resulta de \sref{I.7.2.2.1}, puesto que $f$ es un $S$-morfismo.
\end{proof}

\begin{corollary}[7.2.5]
\label{I.7.2.5}
Sean $X$ e $Y$ dos $S$-preesquemas; se supone que $X$ es reducido y que $X$ e $Y$ son separados sobre $S$.
Sea $p:Y\to X$ un $S$-morfismo (que hace de $Y$ un $X$-preesquema), $U$ un abierto denso de $X$, y $f$ una $U$-sección de $Y$; entonces la aplicación racional $\overline f$ de $X$ en $Y$ que prolonga a $f$ es una $X$-sección racional de $Y$.
\end{corollary}

\oldpage[I]{160}
\begin{proof}
Hay que demostrar que $p\circ\overline f$ es la identidad en el dominio de definición de $\overline f$; como $X$ es separado sobre $S$, esto resulta de nuevo de \sref{I.7.2.2.1}.
\end{proof}

\begin{corollary}[7.2.6]
\label{I.7.2.6}
Sean $X$ un preesquema reducido y $U$ un abierto denso de $X$.
Existe entonces una correspondencia biunívoca canónica entre las secciones de $\mathscr{O}_X$ sobre $U$ y las funciones racionales $f$ sobre $X$ definidas en todos los puntos de $U$.
\end{corollary}

\begin{proof}
Teniendo en cuenta \sref{I.7.2.3}, \sref{I.7.1.2} y \sref{I.7.1.3}, basta observar que el $X$-preesquema $X\otimes_{\mathbf Z}\mathbf Z[T]$ es separado sobre $X$ \sref{I.5.5.1}[iv].
\end{proof}

\begin{corollary}[7.2.7]
\label{I.7.2.7}
Sea $Y$ un preesquema reducido, $f:X\to Y$ un morfismo separado, $U$ un abierto denso de $Y$, $g:U\to f^{-1}(U)$ una $U$-sección de $f^{-1}(U)$, y $Z$ el subpreesquema reducido de $X$ que tiene $\overline{g(U)}$ como espacio subyacente \sref{I.5.2.1}.
Para que $g$ sea la restricción de una $Y$-sección de $X$ \textup{(en otras palabras \sref{I.7.2.5}, para que la aplicación racional de $Y$ en $X$ que prolonga a $g$ esté definida en todas partes)}, es necesario y suficiente que la restricción de $f$ a $Z$ sea un isomorfismo de $Z$ sobre $Y$.
\end{corollary}

\begin{proof}
La restricción de $f$ a $f^{-1}(U)$ es un morfismo separado \sref{I.5.5.1}[i], por lo que $g$ es una inmersión cerrada \sref{I.5.4.6}, y, por consiguiente, $g(U)=Z\cap f^{-1}(U)$, y el subpreesquema inducido por $Z$ sobre el abierto $g(U)$ de $Z$ es idéntico al subpreesquema cerrado de $f^{-1}(U)$ asociado a $g$ \sref{I.5.2.1}.
Es claro entonces que la condición enunciada es suficiente, pues, si se cumple, y si $f_Z:Z\to Y$ es la restricción de $f$ a $Z$ y $\overline g:Y\to Z$ es el isomorfismo inverso, $\overline g$ prolonga a $g$.
Recíprocamente, si $g$ es la restricción a $U$ de una $Y$-sección $h$ de $X$, entonces $h$ es una inmersión cerrada \sref{I.5.4.6}, de modo que $h(Y)$ es cerrado y, como está contenido en $Z$, es igual a $Z$; de \sref{I.5.2.1} resulta que $h$ es necesariamente un isomorfismo de $Y$ sobre el subpreesquema cerrado $Z$ de $X$.
\end{proof}

\begin{env}[7.2.8]
\label{I.7.2.8}
Sean $X$ e $Y$ dos $S$-preesquemas, siendo $X$ reducido e $Y$ separado sobre $S$.
Sea $f$ una $S$-aplicación racional de $X$ en $Y$, y sea $x$ un punto de $X$; se puede componer $f$ con el $S$-morfismo canónico $\operatorname{Spec}(\mathscr{O}_x)\to X$ \sref{I.2.4.1} siempre que la intersección de $\operatorname{Spec}(\mathscr{O}_x)$ con el dominio de definición de $f$ sea densa en $\operatorname{Spec}(\mathscr{O}_x)$ (identificado con el conjunto de los $z\in X$ tales que $x\in\overline{\{z\}}$ \sref{I.2.4.2}).
Esto ocurrirá en los casos siguientes:
\begin{enumerate}
  \item[1.º] $X$ es \emph{irreducible} (y, por tanto, \emph{íntegro}), pues entonces el punto genérico $\xi$ de $X$ es el punto genérico de $\operatorname{Spec}(\mathscr{O}_x)$; como el dominio de definición $U$ de $f$ contiene a $\xi$, $U\cap\operatorname{Spec}(\mathscr{O}_x)$ contiene a $\xi$ y, por tanto, es denso en $\operatorname{Spec}(\mathscr{O}_x)$.
  \item[2.º] $X$ es \emph{localmente noetheriano}; nuestra afirmación resulta entonces de lo siguiente:
\end{enumerate}
\end{env}

\begin{lemma}[7.2.8.1]
\label{I.7.2.8.1}
Sea $X$ un preesquema cuyo espacio subyacente es localmente noetheriano, y sea $x$ un punto de $X$.
Las componentes irreducibles de $\operatorname{Spec}(\mathscr{O}_x)$ son las intersecciones con $\operatorname{Spec}(\mathscr{O}_x)$ de las componentes irreducibles de $X$ que contienen a $x$.
Para que un abierto $U\subset X$ sea tal que $U\cap\operatorname{Spec}(\mathscr{O}_x)$ sea denso en $\operatorname{Spec}(\mathscr{O}_x)$, es necesario y suficiente que tenga intersección no vacía con cada una de las componentes irreducibles de $X$ que contienen a $x$ \textup{(lo que ocurre, en particular, si $U$ es \emph{denso} en $X$)}.
\end{lemma}

\begin{proof}
Basta demostrar la primera afirmación, pues la segunda resulta entonces de ella.
Como $\operatorname{Spec}(\mathscr{O}_x)$ está contenido en todo abierto afín $U$ que contiene a $x$, y como las componentes irreducibles de $U$ que contienen a $x$ son las intersecciones con $U$ de las componentes irreducibles de $X$ que contienen a $x$ \sref[0]{0.2.1.6}, se puede suponer que $X$ es afín, de anillo $A$.
Como los ideales primos de $A_x$ corresponden biunívocamente a los ideales primos
\oldpage[I]{161}
de $A$ contenidos en $\mathfrak{j}_x$ \sref[0]{0.1.2.6}, los ideales primos minimales de $A_x$ corresponden a los ideales primos minimales de $A$ contenidos en $\mathfrak{j}_x$, de donde resulta el lema \sref{I.1.1.14}.
\end{proof}

Dicho esto, supongamos que estamos en uno de los dos casos mencionados en \sref{I.7.2.8}.
Si $U$ es el dominio de definición de la $S$-aplicación racional $f$, denotemos por $f'$ la aplicación racional de $\operatorname{Spec}(\mathscr{O}_x)$ en $Y$ que coincide (teniendo en cuenta \sref{I.2.4.2}) con $f$ en $U\cap\operatorname{Spec}(\mathscr{O}_x)$; diremos que esta aplicación racional es \emph{inducida por $f$}.

\begin{proposition}[7.2.9]
\label{I.7.2.9}
Sean $S$ un preesquema localmente noetheriano, $X$ un $S$-preesquema reducido e $Y$ un $S$-esquema de tipo finito.
Se supone además que $X$ es irreducible o localmente noetheriano.
Sean entonces $f$ una $S$-aplicación racional de $X$ en $Y$, y $x$ un punto de $X$.
Para que $f$ esté definida en el punto $x$, es necesario y suficiente que la aplicación racional $f'$ de $\operatorname{Spec}(\mathscr{O}_x)$ en $Y$, inducida por $f$ \sref{I.7.2.8}, sea un morfismo.
\end{proposition}

\begin{proof}
Siendo la condición evidentemente necesaria (pues $\operatorname{Spec}(\mathscr{O}_x)$ está contenido en todo abierto que contiene a $x$), demostremos que es suficiente.
En virtud de \sref{I.6.5.1}, existe una vecindad abierta $U$ de $x$ en $X$ y un $S$-morfismo $g$ de $U$ en $Y$ que induce $f'$ sobre $\operatorname{Spec}(\mathscr{O}_x)$.
Si $X$ es irreducible, entonces $U$ es densa en $X$ y, en virtud de \sref{I.7.2.3}, se puede considerar $g$ como una $S$-aplicación racional.
Además, el punto genérico de $X$ pertenece a $\operatorname{Spec}(\mathscr{O}_x)$ y al dominio de definición de $f$, de modo que $f$ y $g$ coinciden en este punto y, por consiguiente, en un abierto no vacío de $X$ \sref{I.6.5.1}.
Pero, como $f$ y $g$ son $S$-aplicaciones racionales, son idénticas \sref{I.7.2.3}, y por tanto $f$ está definida en $x$.

Se supone ahora que $X$ es localmente noetheriano; se puede suponer que $U$ es noetheriano.
Solo hay entonces un número finito de componentes irreducibles $X_i$ de $X$ que contienen a $x$ \sref{I.7.2.8.1}, y se puede suponer que son las únicas que tienen intersección no vacía con $U$, reemplazando, si es necesario, $U$ por un abierto menor (pues solo hay un número finito de componentes irreducibles de $X$ que tienen intersección no vacía con $U$, ya que $U$ es noetheriano).
Como antes, se ve entonces que $f$ y $g$ coinciden en un abierto no vacío de cada uno de los $X_i$.
Teniendo en cuenta que cada uno de los $X_i$ está contenido en $\overline U$, consideremos el morfismo $f_1$, definido en un abierto denso de $U\cup(X\setmin\overline U)$, igual a $g$ en $U$ y a $f$ en la intersección de $X\setmin\overline U$ con el dominio de definición de $f$.
Como $U\cup(X\setmin\overline U)$ es denso en $X$, $f_1$ y $f$ coinciden en un abierto denso de $X$; y, como $f$ es una aplicación racional, $f$ es una prolongación de $f_1$ \sref{I.7.2.3}, y está por tanto definida en el punto $x$.
\end{proof}

\subsection{Haz de funciones racionales}
\label{subsection:I.7.3}

\begin{env}[7.3.1]
\label{I.7.3.1}
Sea $X$ un preesquema.
Para todo abierto $U\subset X$, denotemos por $R(U)$ el anillo de las funciones racionales sobre $U$ \sref{I.7.1.3}; es una $\Gamma(U,\mathscr{O}_X)$-álgebra.
Además, si $V\subset U$ es otro abierto de $X$, toda sección de $\mathscr{O}_X$ sobre una parte abierta densa de $U$ da, por restricción a $V$, una sección sobre una parte abierta densa de $V$; y, si dos secciones coinciden sobre una parte abierta densa de $U$, sus restricciones a $V$ coinciden sobre una parte abierta densa de $V$.
Se define así un dihomomorfismo de álgebras $\rho_{V,U}:R(U)\to R(V)$, y está claro que, si $U\supset V\supset W$ son abiertos de $X$, se tiene $\rho_{W,U}=\rho_{W,V}\circ\rho_{V,U}$; los $R(U)$ definen por tanto un \emph{prehaz} de álgebras sobre $X$.
\end{env}

\oldpage[I]{162}
\begin{definition}[7.3.2]
\label{I.7.3.2}
{\itshape Se llama haz de funciones racionales sobre un preesquema $X$, y se denota por $\mathscr{R}(X)$, a la $\mathscr{O}_X$-álgebra asociada al prehaz formado por los $R(U)$.}
\end{definition}

Para todo preesquema $X$ y todo abierto $U\subset X$, está claro que el haz inducido $\mathscr{R}(X)|U$ no es otro que $\mathscr{R}(U)$.

\begin{proposition}[7.3.3]
\label{I.7.3.3}
Sea $X$ un preesquema tal que la familia $(X_\lambda)$ de sus componentes irreducibles sea localmente finita (lo que ocurre, en particular, cuando el espacio subyacente a $X$ es localmente noetheriano).
Entonces el $\mathscr{O}_X$-módulo $\mathscr{R}(X)$ es cuasicoherente y, para todo abierto $U$ de $X$ que solo interseca un número finito de componentes $X_\lambda$, $R(U)$ es igual a $\Gamma(U,\mathscr{R}(X))$ y se identifica canónicamente con el producto directo de los anillos locales de los puntos genéricos de los $X_\lambda$ tales que $U\cap X_\lambda\neq\varnothing$.
\end{proposition}

\begin{proof}
Se puede limitar evidentemente al caso en que $X$ solo tiene un número finito de componentes irreducibles $X_i$, con puntos genéricos $x_i$ ($1\leq i\leq n$).
El hecho de que $R(U)$ se identifique canónicamente con el producto directo de los $\mathscr{O}_{x_i}=R(X_i)$ tales que $U\cap X_i\neq\varnothing$ resulta entonces de \sref{I.7.1.7}.
Mostremos además que el prehaz $U\to R(U)$ satisface los axiomas de los haces, lo que probará que $R(U)=\Gamma(U,\mathscr{R}(X))$.
En efecto, satisface (F1) por lo que precede.
Para ver que satisface (F2), consideremos un recubrimiento de un abierto $U$ de $X$ por abiertos $V_\alpha\subset U$; si los $s_\alpha\in R(V_\alpha)$ son tales que las restricciones de $s_\alpha$ y $s_\beta$ a $V_\alpha\cap V_\beta$ coinciden para todo par de índices, se concluye que, para todo índice $i$ tal que $U\cap X_i\neq\varnothing$, las componentes en $R(X_i)$ de todos los $s_\alpha$ tales que $V_\alpha\cap X_i\neq\varnothing$ son las mismas; denotando por $t_i$ esta componente, está claro que el elemento de $R(U)$ cuyas componentes son los $t_i$ tiene por restricción $s_\alpha$ en cada $V_\alpha$.
Por último, para ver que $\mathscr{R}(X)$ es cuasicoherente, se puede limitar al caso en que $X=\operatorname{Spec}(A)$ es afín; tomando como $U$ los abiertos afines de la forma $D(f)$, donde $f\in A$, resulta de lo anterior y de la Definición~\sref{I.1.3.4} que se tiene $\mathscr{R}(X)=\widetilde M$, donde $M$ es la suma directa de los $A$-módulos $A_{x_i}$.
\end{proof}

\begin{corollary}[7.3.4]
\label{I.7.3.4}
Sea $X$ un preesquema reducido que solo tiene un número finito de componentes irreducibles, y sean $X_i$ ($1\leq i\leq n$) los subpreesquemas cerrados reducidos de $X$ cuyos espacios subyacentes son las componentes irreducibles de $X$ \sref{I.5.2.1}.
Si $h_i$ es la inyección canónica $X_i\to X$, entonces $\mathscr{R}(X)$ es el producto directo de las $\mathscr{O}_X$-álgebras $(h_i)_*(\mathscr{R}(X_i))$.
\end{corollary}

\begin{corollary}[7.3.5]
\label{I.7.3.5}
Si $X$ es irreducible, todo $\mathscr{R}(X)$-módulo cuasicoherente $\mathscr{F}$ es un haz simple.
\end{corollary}

\begin{proof}
Basta demostrar que todo $x\in X$ admite una vecindad $U$ tal que $\mathscr{F}|U$ sea un haz simple \sref[0]{0.3.6.2}; en otros términos, se reduce al caso en que $X$ es afín; se puede suponer además que $\mathscr{F}$ es el conúcleo de un homomorfismo $(\mathscr{R}(X))^{(I)}\to(\mathscr{R}(X))^{(J)}$ \sref[0]{0.5.1.3}, y todo se reduce a ver que $\mathscr{R}(X)$ es un haz simple; pero esto es evidente, ya que $\Gamma(U,\mathscr{R}(X))=R(X)$ para todo abierto no vacío $U$, pues $U$ contiene el punto genérico de $X$.
\end{proof}

\begin{corollary}[7.3.6]
\label{I.7.3.6}
Si $X$ es irreducible, para todo $\mathscr{O}_X$-módulo cuasicoherente $\mathscr{F}$, $\mathscr{F}\otimes_{\mathscr{O}_X}\mathscr{R}(X)$ es un haz simple; si además $X$ es reducido (y por tanto íntegro), $\mathscr{F}\otimes_{\mathscr{O}_X}\mathscr{R}(X)$ es isomorfo a un haz de la forma $(\mathscr{R}(X))^{(I)}$.
\end{corollary}

\begin{proof}
La segunda afirmación resulta de que $R(X)$ es entonces un cuerpo.
\end{proof}

\begin{proposition}[7.3.7]
\label{I.7.3.7}
Supongamos que el preesquema $X$ es localmente íntegro o localmente
\oldpage[I]{163}
noetheriano.
Entonces $\mathscr{R}(X)$ es una $\mathscr{O}_X$-álgebra cuasicoherente; si además $X$ es reducido (lo que ocurre cuando $X$ es localmente íntegro), el homomorfismo canónico $\mathscr{O}_X\to\mathscr{R}(X)$ es inyectivo.
\end{proposition}

\begin{proof}
Como la cuestión es local, la primera afirmación resulta de \sref{I.7.3.3}; la segunda resulta inmediatamente de \sref{I.7.2.3}.
\end{proof}

\begin{env}[7.3.8]
\label{I.7.3.8}
\footnote{\emph{[Trad.] El texto original de este apartado, declarado erróneo, fue sustituido íntegramente en la fe de erratas de EGA~II.}}
Sean $X,Y$ dos preesquemas \emph{íntegros}, lo que implica que $\mathscr{R}(X)$ (resp. $\mathscr{R}(Y)$) es un $\mathscr{O}_X$-módulo cuasicoherente (resp. un $\mathscr{O}_Y$-módulo cuasicoherente) \sref{I.7.3.3}.
Sea $f:X\to Y$ un morfismo \emph{dominante}; entonces existe un homomorfismo canónico de $\mathscr{O}_X$-módulos
\[
  \tau:f^*(\mathscr{R}(Y))\to\mathscr{R}(X).
\tag{7.3.8.1}
\label{I.7.3.8.1-fr}
\]
\end{env}

\begin{proof}
Supongamos primero que $X=\operatorname{Spec}(A)$ e $Y=\operatorname{Spec}(B)$ son afines, de anillos íntegros $A$ y $B$, de modo que $f$ corresponde a un homomorfismo inyectivo $B\to A$ que se prolonga a un monomorfismo $L\to K$ del cuerpo de fracciones $L$ de $B$ en el cuerpo de fracciones $K$ de $A$.
El homomorfismo \sref{I.7.3.8.1} corresponde entonces al homomorfismo canónico \[
  L\otimes_B A\to K
\] \sref{I.1.6.5}.
En el caso general, para todo par de abiertos afines no vacíos $U\subset X$ y $V\subset Y$ tales que $f(U)\subset V$, se define, como antes, un homomorfismo $\tau_{U,V}$ y se comprueba inmediatamente que, si $U'\subset U$, $V'\subset V$, $f(U')\subset V'$, entonces $\tau_{U,V}$ prolonga $\tau_{U',V'}$, de donde resulta nuestra afirmación.
Si $x,y$ son los puntos genéricos de $X$ e $Y$, respectivamente, se tiene $f(x)=y$,
\[
  \bigl(f^*(\mathscr{R}(Y))\bigr)_x
  =\mathscr{O}_y\otimes_{\mathscr{O}_y}\mathscr{O}_x
  =\mathscr{O}_x
\]
\sref[0]{0.4.3.1} y $\tau_x$ es, por tanto, un \emph{isomorfismo}.
\end{proof}

\subsection{Haces de torsión y haces libres de torsión}
\label{subsection:I.7.4}
\label{I.7.4}

\begin{env}[7.4.1]
\label{I.7.4.1}
Sea $X$ un preesquema \emph{íntegro}.
Para todo $\mathscr{O}_X$-módulo $\mathscr{F}$, el homomorfismo canónico $\mathscr{O}_X\to\mathscr{R}(X)$ define, por tensorización, un homomorfismo (también llamado \emph{canónico}) \[
  \mathscr{F}\to
  \mathscr{F}\otimes_{\mathscr{O}_X}\mathscr{R}(X)
\] que, en cada tallo, no es otro que el homomorfismo $z\to z\otimes 1$ de $\mathscr{F}_x$ en $\mathscr{F}_x\otimes_{\mathscr{O}_x}R(X)$.
El \emph{núcleo $\mathscr{T}$} de este homomorfismo es un $\mathscr{O}_X$-submódulo de $\mathscr{F}$, llamado \emph{haz de torsión} de $\mathscr{F}$; es cuasicoherente si $\mathscr{F}$ es cuasicoherente (\sref{I.4.1.1} y \sref{I.7.3.6}).
Se dice que $\mathscr{F}$ es \emph{libre de torsión} si $\mathscr{T}=0$, y que $\mathscr{F}$ es un \emph{haz de torsión} si $\mathscr{T}=\mathscr{F}$.
Para todo $\mathscr{O}_X$-módulo $\mathscr{F}$, $\mathscr{F}/\mathscr{T}$ es libre de torsión.
Se deduce de \sref{I.7.3.5} que:
\end{env}

\begin{proposition}[7.4.2]
\label{I.7.4.2}
Si $X$ es un preesquema íntegro, todo $\mathscr{O}_X$-módulo cuasicoherente libre de torsión $\mathscr{F}$ es isomorfo a un subhaz $\mathscr{G}$ de un haz simple de la forma $(\mathscr{R}(X))^{(I)}$, generado (como $\mathscr{R}(X)$-módulo) por $\mathscr{G}$.
\end{proposition}

La cardinalidad de $I$ se llama el \emph{rango} de $\mathscr{F}$; para todo abierto afín no vacío $U$ de $X$, el rango de $\mathscr{F}$ es igual al rango de $\Gamma(U,\mathscr{F})$ como $\Gamma(U,\mathscr{O}_X)$-módulo, como se ve inmediatamente al considerar el punto genérico de $X$, contenido en $U$.
En particular:
\begin{corollary}[7.4.3]
\label{I.7.4.3}
En un preesquema íntegro $X$, todo $\mathscr{O}_X$-módulo cuasicoherente libre de torsión y de rango $1$ (en particular, todo $\mathscr{O}_X$-módulo invertible) es isomorfo a un $\mathscr{O}_X$-submódulo cuasicoherente no nulo de $\mathscr{R}(X)$, y recíprocamente.
\end{corollary}

\begin{corollary}[7.4.4]
\label{I.7.4.4}
Sean $X$ un preesquema íntegro, $\mathscr{L}$ y $\mathscr{L}'$ dos $\mathscr{O}_X$-módulos libres de torsión, y $f$ \textup{(resp. $f'$)} una sección de $\mathscr{L}$ \textup{(resp. $\mathscr{L}'$)} sobre $X$.
Para que $f\otimes f'=0$, es necesario y suficiente que una de las secciones $f$ y $f'$ sea nula.
\end{corollary}

\begin{proof}
Sea $x$ el punto genérico de $X$; por hipótesis, se tiene $(f\otimes f')_x=f_x\otimes f'_x=0$.
Como $\mathscr{L}_x$ y $\mathscr{L}'_x$ son espacios vectoriales sobre el cuerpo $\mathscr{O}_x$, la relación precedente implica $f_x=0$ o $f'_x=0$, y por consiguiente $f=0$ o $f'=0$, puesto que $\mathscr{L}$ y $\mathscr{L}'$ son libres de torsión \sref{I.7.3.5}.
\end{proof}

\begin{proposition}[7.4.5]
\label{I.7.4.5}
Sean $X$ e $Y$ dos preesquemas íntegros, y $f:X\to Y$ un morfismo dominante.
Para todo $\mathscr{O}_X$-módulo cuasicoherente libre de torsión $\mathscr{F}$, $f_*(\mathscr{F})$ es un $\mathscr{O}_Y$-módulo libre de torsión.
\end{proposition}

\oldpage[I]{164}
\begin{proof}
Como $f_*$ es exacto por la izquierda \sref[0]{0.4.2.1}, basta, en virtud de \sref{I.7.4.2}, probar la proposición en el caso en que $\mathscr{F}=(\mathscr{R}(X))^{(I)}$.
Ahora bien, todo abierto no vacío $U$ de $Y$ contiene el punto genérico de $Y$, por lo que $f^{-1}(U)$ contiene el punto genérico de $X$ \sref[0]{0.2.1.5}; por lo tanto, se tiene \[
  \Gamma(U,f_*(\mathscr{F}))
  =\Gamma(f^{-1}(U),\mathscr{F})=(R(X))^{(I)}~;
\]; en otras palabras, $f_*(\mathscr{F})$ es el haz simple cuyo tallo es $(R(X))^{(I)}$, considerado como $\mathscr{R}(Y)$-módulo, y es evidentemente libre de torsión.
\end{proof}

\begin{proposition}[7.4.6]
\label{I.7.4.6}
Sea $X$ un preesquema íntegro y $x$ su punto genérico.
Para todo $\mathscr{O}_X$-módulo cuasicoherente de tipo finito $\mathscr{F}$, las condiciones siguientes son equivalentes: \textup{a)} $\mathscr{F}$ es un haz de torsión; \textup{b)} $\mathscr{F}_x=0$; \textup{c)} $\operatorname{Supp}(\mathscr{F})\neq X$.
\end{proposition}

\begin{proof}
En virtud de \sref{I.7.3.5} y \sref{I.7.4.1}, las relaciones $\mathscr{F}_x=0$ y $\mathscr{F}\otimes_{\mathscr{O}_X}\mathscr{R}(X)=0$ son equivalentes, por lo que \emph{a)} y \emph{b)} son equivalentes; por otra parte, $\operatorname{Supp}(\mathscr{F})$ es cerrado en $X$ \sref[0]{0.5.2.2}, y como todo abierto no vacío de $X$ contiene a $x$, \emph{b)} y \emph{c)} son equivalentes.
\end{proof}

\begin{env}[7.4.7]
\label{I.7.4.7}
Se extienden (por abuso de lenguaje) las definiciones de \sref{I.7.4.1} al caso en que $X$ es un preesquema \emph{reducido} que solo tiene un número \emph{finito} de componentes irreducibles; resulta entonces de \sref{I.7.3.4} que la equivalencia entre a) y c) en \sref{I.7.4.6} sigue siendo válida para tal preesquema, a condición de sustituir c) por la condición de que $\operatorname{Supp}(\mathscr{F})$ no contenga ninguna componente irreducible de $X$.
\end{env}
% END source_es/ega1/ega1-7.tex

% BEGIN source_es/ega1/ega1-8.tex
\section{Los esquemas de Chevalley}
\label{section:I.8}

\subsection{Anillos locales emparentados}
\label{subsection:I.8.1}

Para todo anillo local $A$, denotaremos por $\mathfrak m(A)$ el ideal maximal de $A$.
\begin{lemma}[8.1.1]
\label{I.8.1.1}
Sean $A$ y $B$ dos anillos locales tales que $A\subset B$.
Las condiciones siguientes son equivalentes.
\begin{enumerate}
  \item[{\rm(i)}] $\mathfrak m(B)\cap A=\mathfrak m(A)$.
  \item[{\rm(ii)}] $\mathfrak m(A)\subset\mathfrak m(B)$.
  \item[{\rm(iii)}] $1$ no pertenece al ideal de $B$ engendrado por $\mathfrak m(A)$.
\end{enumerate}
\end{lemma}

\begin{proof}
Es evidente que (i) implica (ii) y que (ii) implica (iii);
por último, si se verifica (iii), entonces $\mathfrak m(B)\cap A$ contiene $\mathfrak m(A)$ y no contiene a $1$, por lo que es igual a $\mathfrak m(A)$.

Cuando se cumplen las condiciones equivalentes de \sref{I.8.1.1}, se dice que $B$ \emph{domina} a $A$;
esto equivale a decir que la inyección $A\to B$ es un homomorfismo \emph{local}.
Es claro que, en el conjunto de los subanillos locales de un anillo $R$, la relación de dominación es una relación de orden.
\end{proof}

\begin{env}[8.1.2]
\label{I.8.1.2}
Consideremos ahora un \emph{cuerpo} $R$.
Para todo subanillo $A$ de $R$, denotaremos por $L(A)$ el conjunto de los anillos locales $A_\mathfrak p$, donde $\mathfrak p$ recorre el espectro primo de $A$;
tales anillos locales se identifican con los subanillos de $R$ que contienen a $A$.
Puesto que $\mathfrak p=(\mathfrak p A_\mathfrak p)\cap A$, la aplicación $\mathfrak p\to A_\mathfrak p$ de $\operatorname{Spec}(A)$ en $L(A)$ es biyectiva.
\end{env}

\begin{lemma}[8.1.3]
\label{I.8.1.3}
Sean $R$ un cuerpo y $A$ un subanillo de $R$.
Para que un subanillo local $M$ de $R$ domine un anillo $A_\mathfrak p\in L(A)$, es necesario y suficiente que $A\subset M$;
el anillo local $A_\mathfrak p$ dominado por $M$ es entonces único y corresponde a $\mathfrak p=\mathfrak m(M)\cap A$.
\end{lemma}

\begin{proof}
En efecto, si $M$ domina $A_\mathfrak p$, se tiene $\mathfrak m(M)\cap A_\mathfrak p=\mathfrak p A_\mathfrak p$ según \sref{I.8.1.1}, de donde resulta la unicidad de $\mathfrak p$;
por otra parte, si $A\subset M$, entonces $\mathfrak m(M)\cap A=\mathfrak p$ es primo en $A$, y como $A\setmin\mathfrak p\subset M$, se tiene $A_\mathfrak p\subset M$ y $\mathfrak p A_\mathfrak p\subset\mathfrak m(M)$, por lo que $M$ domina $A_\mathfrak p$.
\end{proof}

\oldpage[I]{165}
\begin{lemma}[8.1.4]
\label{I.8.1.4}
Sean $R$ un cuerpo, $M$ y $N$ dos subanillos locales de $R$, y $P$ el subanillo de $R$ engendrado por $M\cup N$.
Las condiciones siguientes son equivalentes.
\begin{enumerate}
  \item[{\rm(i)}] Existe un ideal primo $\mathfrak p$ de $P$ tal que $\mathfrak m(M)=\mathfrak p\cap M$ y $\mathfrak m(N)=\mathfrak p\cap N$.
  \item[{\rm(ii)}] El ideal $\mathfrak a$ engendrado en $P$ por $\mathfrak m(M)\cup\mathfrak m(N)$ es distinto de $P$.
  \item[{\rm(iii)}] Existe un subanillo local $Q$ de $R$ que domina a la vez a $M$ y a $N$.
\end{enumerate}
\end{lemma}

\begin{proof}
Es claro que (i) implica (ii);
recíprocamente, si $\mathfrak a\ne P$, entonces $\mathfrak a$ está contenido en un ideal maximal $\mathfrak n$ de $P$, y como $1\notin\mathfrak n$, $\mathfrak n\cap M$ contiene a $\mathfrak m(M)$ y es distinto de $M$, luego $\mathfrak n\cap M=\mathfrak m(M)$, y análogamente $\mathfrak n\cap N=\mathfrak m(N)$.
Es claro que, si $Q$ domina a la vez a $M$ y a $N$, entonces $P\subset Q$ y $\mathfrak m(M)=\mathfrak m(Q)\cap M=(\mathfrak m(Q)\cap P)\cap M$, y $\mathfrak m(N)=(\mathfrak m(Q)\cap P)\cap N$, por lo que (iii) implica (i);
la recíproca es evidente tomando $Q=P_\mathfrak p$.
\end{proof}

Cuando se cumplen las condiciones de \sref{I.8.1.4}, se dice, con C.~Chevalley, que los anillos locales $M$ y $N$ están \emph{emparentados}.

\begin{proposition}[8.1.5]
\label{I.8.1.5}
Sean $A$ y $B$ dos subanillos de un cuerpo $R$, y $C$ el subanillo de $R$ engendrado por $A\cup B$.
Las condiciones siguientes son equivalentes.
\begin{enumerate}
  \item[{\rm(i)}] Para todo anillo local $Q$ que contenga a $A$ y a $B$, se tiene $A_\mathfrak p=B_\mathfrak q$, poniendo $\mathfrak p=\mathfrak m(Q)\cap A$ y $\mathfrak q=\mathfrak m(Q)\cap B$.
  \item[{\rm(ii)}] Para todo ideal primo $\mathfrak r$ de $C$, se tiene $A_\mathfrak p=B_\mathfrak q$, poniendo $\mathfrak p=\mathfrak r\cap A$ y $\mathfrak q=\mathfrak r\cap B$.
  \item[{\rm(iii)}] Si $M\in L(A)$ y $N\in L(B)$ están emparentados, son idénticos.
  \item[{\rm(iv)}] $L(A)\cap L(B)=L(C)$.
\end{enumerate}
\end{proposition}

\begin{proof}
Los lemas~\sref{I.8.1.3} y \sref{I.8.1.4} prueban que (i) y (iii) son equivalentes;
es claro que (i) implica (ii) aplicándolo a $Q=C_\mathfrak r$;
recíprocamente, (ii) implica (i), pues si $Q$ contiene a $A\cup B$, contiene a $C$, y si $\mathfrak r=\mathfrak m(Q)\cap C$, se tiene $\mathfrak p=\mathfrak r\cap A$ y $\mathfrak q=\mathfrak r\cap B$, en virtud de \sref{I.8.1.3}.
Es inmediato que (iv) implica (i), pues si $Q$ contiene a $A\cup B$, domina un anillo local $C_\mathfrak r\in L(C)$ por \sref{I.8.1.3};
por hipótesis se tiene $C_\mathfrak r\in L(A)\cap L(B)$, y \sref{I.8.1.1} y \sref{I.8.1.3} prueban que $C_\mathfrak r=A_\mathfrak p=B_\mathfrak q$.
Demostremos por último que (iii) implica (iv).
Sea $Q\in L(C)$;
$Q$ domina un $M\in L(A)$ y un $N\in L(B)$ \sref{I.8.1.3}, y por tanto $M$ y $N$, al estar emparentados, son idénticos por hipótesis.
Como se tiene entonces $C\subset M$, se sabe que $M$ domina un $Q'\in L(C)$ por \sref{I.8.1.3}; por tanto, $Q$ domina a $Q'$, de donde necesariamente, por \sref{I.8.1.3}, $Q=Q'=M$, y por tanto $Q\in L(A)\cap L(B)$.
Recíprocamente, si $Q\in L(A)\cap L(B)$, entonces $C\subset Q$, por lo que \sref{I.8.1.3} $Q$ domina un $Q''\in L(C)\subset L(A)\cap L(B)$;
$Q$ y $Q''$, al estar emparentados, son idénticos, luego $Q''=Q\in L(C)$, lo que concluye la demostración.
\end{proof}

\subsection{Anillos locales de un esquema íntegro}
\label{subsection:I.8.2}

\begin{env}[8.2.1]
\label{I.8.2.1}
Sea $X$ un preesquema \emph{íntegro}, y sea $R$ su cuerpo de funciones racionales, idéntico al anillo local del punto genérico $a$ de $X$;
para todo $x\in X$, se sabe que $\mathscr{O}_x$ se identifica canónicamente con un subanillo de $R$ \sref{I.7.1.5}, y, para toda función racional $f\in R$, el dominio de definición $\delta(f)$ de $f$ es el abierto formado por los $x\in X$ tales que
$f\in\mathscr{O}_x$.
Se deduce entonces de \sref{I.7.2.6} que, para todo abierto $U\subset X$, se tiene
\begin{equation}
  \Gamma(U,\mathscr{O}_X)=\bigcap_{x\in U}\mathscr{O}_x
  \tag{8.2.1.1}\label{I.8.2.1.1-fr}
\end{equation}
\end{env}

\oldpage[I]{166}
\begin{proposition}[8.2.2]
\label{I.8.2.2}
Sea $X$ un preesquema íntegro, y sea $R$ su cuerpo de funciones racionales. Para que $X$ sea un esquema, es necesario y suficiente que la relación «$\mathscr{O}_x$ y $\mathscr{O}_y$ están emparentados» \sref{I.8.1.4} entre puntos $x$ e $y$ de $X$, implique que $x=y$.
\end{proposition}

\begin{proof}
Supongamos que se satisface esta condición y demostremos que $X$ es separado.
Sean $U$ y $V$ dos abiertos afines distintos de $X$, de anillos $A$ y $B$, respectivamente, identificados con subanillos de $R$;
$U$ (resp. $V$) se identifica entonces \sref{I.8.1.2} con $L(A)$ (resp. $L(B)$), y la hipótesis implica \sref{I.8.1.5} que, si $C$ es el subanillo de $R$ engendrado por $A\cup B$, entonces $W=U\cap V$ se identifica con $L(A)\cap L(B)=L(C)$.
Además, se sabe (\cite{I-1}, p.~5-03, prop.~4~\emph{bis}) que todo subanillo $E$ de $R$ es igual a la intersección de los anillos locales pertenecientes a $L(E)$;
por tanto, $C$ se identifica con la intersección de los anillos $\mathscr{O}_z$ para $z\in W$, es decir \sref{I.8.2.1.1}, con $\Gamma(W,\mathscr{O}_X)$.
Consideremos entonces el subpreesquema inducido por $X$ sobre $W$;
al homomorfismo identidad $\varphi:C\to\Gamma(W,\mathscr{O}_X)$ corresponde \sref{I.2.2.4} un morfismo $\Phi=(\psi,\theta):W\to\operatorname{Spec}(C)$;
veremos que $\Phi$ es un \emph{isomorfismo} de preesquemas, de donde resultará que $W$ es un abierto \emph{afín}.
La identificación de $W$ con $L(C)=\operatorname{Spec}(C)$ muestra que $\psi$ es \emph{biyectiva}.
Por otra parte, para todo $x\in W$, $\theta_x^\sharp$ es la inyección $C_\mathfrak r\to\mathscr{O}_x$, donde $\mathfrak r=\mathfrak m_x\cap C$, y, por definición, $C_\mathfrak r$ se identifica con $\mathscr{O}_x$, por lo que $\theta_x^\sharp$ es biyectiva.
Queda, pues, por demostrar que $\psi$ es un \emph{homeomorfismo}, es decir, que para todo subconjunto cerrado $F\subset W$, $\psi(F)$ es cerrado en
$\operatorname{Spec}(C)$.
Ahora bien, $F$ es la traza sobre $W$ de un subconjunto cerrado de $U$ de la forma $V(\mathfrak a)$, donde $\mathfrak a$ es un ideal de $A$;
demostremos que $\psi(F)=V(\mathfrak a C)$, lo que probará nuestra afirmación.
En efecto, los ideales primos de $C$ que contienen $\mathfrak a C$ son los ideales primos de $C$ que contienen $\mathfrak a$, y son, por tanto, los ideales de la forma $\psi(x)=\mathfrak m_x\cap C$, donde $\mathfrak a\subset\mathfrak m_x$ y $x\in W$;
como $\mathfrak a\subset\mathfrak m_x$ equivale a $x\in V(\mathfrak a)=W\cap F$ para $x\in U$, se tiene efectivamente $\psi(F)=V(\mathfrak a C)$\footnote{\emph{[Trad.] El texto francés imprime $x\in V(\mathfrak a)=W\cap F$ para $x\in U$. Puesto que $F$ se definió inmediatamente antes como la traza de $V(\mathfrak a)$ sobre $W$, la igualdad requerida es $F=W\cap V(\mathfrak a)$ y este paso se refiere a $x\in W$.}}.

Se sigue que $X$ es separado, pues $U\cap V$ es afín y su anillo $C$ está engendrado por la unión $A\cup B$ de los anillos de $U$ y $V$ \sref{I.5.5.6}.

Recíprocamente, supongamos que $X$ es separado, y sean $x$ e $y$ puntos de $X$ tales que $\mathscr{O}_x$ y $\mathscr{O}_y$ estén emparentados.
Sea $U$ (resp. $V$) un abierto afín que contiene a $x$ (resp. $y$), de anillo $A$ (resp. $B$);
se sabe entonces que $U\cap V$ es afín y que su anillo $C$ está engendrado por $A\cup B$ \sref{I.5.5.6}.
Si $\mathfrak p=\mathfrak m_x\cap A$ y $\mathfrak q=\mathfrak m_y\cap B$, entonces $A_\mathfrak p=\mathscr{O}_x$ y $B_\mathfrak q=\mathscr{O}_y$, y, como $A_\mathfrak p$ y $B_\mathfrak q$ están emparentados, existe un ideal primo $\mathfrak r$ de $C$ tal que $\mathfrak p=\mathfrak r\cap A$ y $\mathfrak q=\mathfrak r\cap B$ \sref{I.8.1.4}.
Pero existe entonces un punto $z\in U\cap V$ tal que $\mathfrak r=\mathfrak m_z\cap C$, puesto que $U\cap V$ es afín, y se tiene evidentemente $x=z$ e $y=z$, de donde $x=y$.
\end{proof}

\begin{corollary}[8.2.3]
\label{I.8.2.3}
Sea $X$ un esquema \emph{íntegro}, y sean $x$ e $y$ puntos de $X$.
Para que $x\in\overline{\{y\}}$, es necesario y suficiente que
$\mathscr{O}_x\subset\mathscr{O}_y$, o, equivalentemente, que toda función racional
definida en $x$ también esté definida en $y$.
\end{corollary}

\begin{proof}
La condición es evidentemente necesaria, pues el dominio de definición
$\delta(f)$ de una función racional $f\in R$ es abierto; demostremos ahora
que es suficiente.
Si $\mathscr{O}_x\subset\mathscr{O}_y$, entonces existe un ideal primo $\mathfrak p$
de $\mathscr{O}_x$ tal que $\mathscr{O}_y$ domina a
$(\mathscr{O}_x)_\mathfrak p$ \sref{I.8.1.3}; pero, por \sref{I.2.4.2}, existe
algún $z\in X$ tal que $x\in\overline{\{z\}}$ y
$\mathscr{O}_z=(\mathscr{O}_x)_\mathfrak p$; como $\mathscr{O}_z$ y $\mathscr{O}_y$ están
emparentados, se tiene $z=y$ por \sref{I.8.2.2}, de donde se sigue el
corolario.
\end{proof}

\begin{corollary}[8.2.4]
\label{I.8.2.4}
Si $X$ es un esquema \emph{íntegro}, entonces la aplicación $x\to\mathscr{O}_x$ es inyectiva; en otras palabras, si $x$ e $y$ son dos puntos distintos de $X$, existe una función racional definida en uno de esos puntos y no en el otro.
\end{corollary}

\oldpage[I]{167}
\begin{proof}
Esto se sigue de \sref{I.8.2.3} y del axioma $(T_0)$ \sref{I.2.1.4}.
\end{proof}

\begin{corollary}[8.2.5]
\label{I.8.2.5}
Sea $X$ un esquema íntegro cuyo espacio subyacente es noetheriano;
cuando $f$ recorre el cuerpo $R$ de las funciones racionales sobre $X$, los conjuntos $\delta(f)$ engendran la topología de $X$.
\end{corollary}

En efecto, todo subconjunto cerrado de $X$ es entonces una unión finita de subconjuntos cerrados irreducibles, es decir, de subconjuntos de la forma $\overline{\{y\}}$ \sref{I.2.1.5}.
Ahora bien, si $x\notin\overline{\{y\}}$, existe una función racional $f$ definida en $x$ y no en $y$ \sref{I.8.2.3}; dicho de otro modo, se tiene $x\in\delta(f)$ y $\delta(f)$ no corta a $\overline{\{y\}}$.
El complemento de $\overline{\{y\}}$ es, por tanto, una unión de conjuntos de la forma $\delta(f)$ y, en virtud de la primera observación, todo abierto de $X$ es una unión de intersecciones finitas de abiertos de la forma $\delta(f)$.

\begin{env}[8.2.6]
\label{I.8.2.6}
El Corolario~\sref{I.8.2.5} muestra que la topología de $X$ queda enteramente caracterizada por la familia de anillos locales $(\mathscr{O}_x)_{x\in X}$ que tienen $R$ como cuerpo de fracciones.
Equivale asimismo a decir que los subconjuntos cerrados de $X$ se definen del modo siguiente: dado un subconjunto finito $\{x_1,\ldots,x_n\}$ de $X$, se considera el conjunto de los $y\in X$ tales que $\mathscr{O}_y\subset\mathscr{O}_{x_i}$ para al menos un índice $i$; estos conjuntos (para todas las elecciones de $\{x_1,\ldots,x_n\}$) son los subconjuntos cerrados de $X$.
Además, una vez conocida la topología de $X$, el haz estructural $\mathscr{O}_X$ queda asimismo determinado por la familia de los $\mathscr{O}_x$, puesto que $\Gamma(U,\mathscr{O}_X)=\bigcap_{x\in U}\mathscr{O}_x$, por \sref{I.8.2.1.1}.
La familia $(\mathscr{O}_x)_{x\in X}$ determina, por tanto, completamente el preesquema $X$ cuando $X$ es un esquema íntegro cuyo espacio subyacente es noetheriano.
\end{env}

\begin{proposition}[8.2.7]
\label{I.8.2.7}
Sean $X$ e $Y$ dos esquemas íntegros, $f:X\to Y$ un morfismo dominante \sref{I.2.2.6}, y $K$ \textup{(resp. $L$)} el cuerpo de funciones racionales sobre $X$ \textup{(resp. $Y$)}.
Entonces $L$ se identifica con un subcuerpo de $K$ y, para todo $x\in X$, $\mathscr{O}_{f(x)}$ es el único anillo local de $Y$ dominado por $\mathscr{O}_x$.
\end{proposition}

\begin{proof}
En efecto, si $f=(\psi,\theta)$ y $a$ es el punto genérico de $X$, $\psi(a)$ es el punto genérico de $Y$ \sref[0]{0.2.1.5};
$\theta_a^\sharp$ es, por tanto, un monomorfismo del cuerpo $L=\mathscr{O}_{\psi(a)}$ en el cuerpo $K=\mathscr{O}_a$.
Como todo abierto afín no vacío $U$ de $Y$ contiene a $\psi(a)$, resulta de \sref{I.2.2.4} que el homomorfismo $\Gamma(U,\mathscr{O}_Y)\to\Gamma(\psi^{-1}(U),\mathscr{O}_X)$ correspondiente a $f$ es la restricción de $\theta_a^\sharp$ a $\Gamma(U,\mathscr{O}_Y)$.
Por tanto, para todo $x\in X$, $\theta_x^\sharp$ es la restricción a $\mathscr{O}_{\psi(x)}$ de $\theta_a^\sharp$ y es, por consiguiente, un monomorfismo.
Se sabe además que $\theta_x^\sharp$ es un homomorfismo local; así pues, si se identifica $L$ con un subcuerpo de $K$ mediante $\theta_a^\sharp$, $\mathscr{O}_{\psi(x)}$ está dominado por $\mathscr{O}_x$ \sref{I.8.1.1};
es, además, el único anillo local de $Y$ dominado por $\mathscr{O}_x$, puesto que dos anillos locales de $Y$ que están emparentados son idénticos \sref{I.8.2.2}.
\end{proof}

\begin{proposition}[8.2.8]
\label{I.8.2.8}
Sea $X$ un preesquema \emph{irreducible}, sea $f:X\to Y$ una inmersión local \textup{(resp.} un isomorfismo local), y se supone además que $f$ es separado. Entonces $f$ es una inmersión \textup{(resp.} una inmersión abierta).
\end{proposition}

\begin{proof}
Sea $f=(\psi,\theta)$;
basta, en ambos casos, demostrar que $\psi$ es un \emph{homeomorfismo} de $X$ sobre $\psi(X)$ \sref{I.4.5.3}.
Sustituyendo $f$ por $f_{\mathrm{red}}$ (\sref{I.5.1.6} y \sref{I.5.5.1}[vi]), se puede suponer que $X$ e $Y$ son \emph{reducidos}.
Si $Y'$ es el subpreesquema cerrado reducido de $Y$ cuyo espacio subyacente es $\overline{\psi(X)}$, entonces $f$ se factoriza como $X\xrightarrow{f'}Y'\xrightarrow{j}Y$, donde $j$ es la inyección canónica \sref{I.5.2.2}.
Resulta de \sref{I.5.5.1}[v] que $f'$ sigue siendo un morfismo separado; además, $f'$ sigue siendo
\oldpage[I]{168}
una inmersión local (resp. un isomorfismo local), pues, siendo la cuestión local en $X$ e $Y$, para demostrarlo se puede limitar al caso en que $f$ es una inmersión cerrada (resp. una inmersión abierta), y nuestra afirmación se deduce entonces inmediatamente de \sref{I.4.2.2}.

Se puede, por tanto, suponer que $f$ es un morfismo \emph{dominante}, lo que implica que $Y$ también es irreducible \sref[0]{0.2.1.5}, y, por consiguiente, que $X$ e $Y$ son ambos \emph{íntegros}.
Además, siendo la cuestión local en $Y$, se puede suponer que $Y$ es un esquema afín;
como $f$ es separado, $X$ es un esquema \sref{I.5.5.1}[ii], y finalmente se está en las hipótesis de la Proposición~\sref{I.8.2.7}.
Entonces, para todo $x\in X$, $\theta_x^\sharp$ es inyectivo;
pero la hipótesis de que $f$ es una inmersión local implica que $\theta_x^\sharp$ es sobreyectivo \sref{I.4.2.2}, luego $\theta_x^\sharp$ es biyectivo, es decir, con la identificación de la Proposición~\sref{I.8.2.7}, se tiene $\mathscr{O}_{\psi(x)}=\mathscr{O}_x$.
Esto implica, por el Corolario~\sref{I.8.2.4}, que $\psi$ es una aplicación \emph{inyectiva}, lo que ya demuestra la proposición cuando $f$ es un isomorfismo local \sref{I.4.5.3}.
Si se supone solamente que $f$ es una inmersión local, para todo $x\in X$ existe una vecindad abierta $U$ de $x$ en $X$ y una vecindad abierta $V$ de $\psi(x)$ en $Y$ tales que la restricción de $\psi$ a $U$ es un homeomorfismo de $U$ sobre un subconjunto \emph{cerrado} de $V$.
Ahora bien, $U$ es densa en $X$, luego $\psi(U)$ es denso en $Y$ y \emph{a fortiori} en $V$, lo que prueba que $\psi(U)=V$;
como $\psi$ es inyectiva, $\psi^{-1}(V)=U$, y esto termina de demostrar que $\psi$ es un homeomorfismo de $X$ sobre $\psi(X)$.
\end{proof}

\subsection{Esquemas de Chevalley}
\label{subsection:I.8.3}

\begin{env}[8.3.1]
\label{I.8.3.1}
Sea $X$ un esquema íntegro \emph{noetheriano}, y sea $R$ su cuerpo de funciones racionales;
denotemos por $X'$ el conjunto de los subanillos locales $\mathscr{O}_x\subset R$, donde $x$ recorre todos los puntos de $X$.
El conjunto $X'$ satisface las tres condiciones siguientes.
\begin{enumerate}
  \item[(Sch.~1)] \emph{Para todo $M\in X'$, $R$ es el cuerpo de fracciones de $M$.}
  \item[(Sch.~2)] \emph{Existe un conjunto finito de subanillos noetherianos $A_i$ de $R$ tales que $X'=\bigcup_i L(A_i)$ y, para todo par de índices $i$, $j$, el subanillo $A_{ij}$ de $R$ generado por $A_i\cup A_j$ es un álgebra de tipo finito sobre $A_i$.}
  \item[(Sch.~3)] \emph{Dos elementos cualesquiera $M$ y $N$ de $X'$ que están emparentados son idénticos.}
\end{enumerate}
\end{env}

Se ha visto en \sref{I.8.2.1} que se satisface (Sch.~1), y (Sch.~3) se deduce de \sref{I.8.2.2}.
Para demostrar (Sch.~2), basta cubrir $X$ por un número finito de abiertos afines $U_i$ cuyos anillos sean noetherianos, y tomar $A_i=\Gamma(U_i,\mathscr{O}_X)$;
la hipótesis de que $X$ es un esquema implica que $U_i\cap U_j$ es afín, y también que $\Gamma(U_i\cap U_j,\mathscr{O}_X)=A_{ij}$ \sref{I.5.5.6};
además, como el espacio $U_i$ es noetheriano, la inmersión $U_i\cap U_j\to U_i$ es de tipo finito \sref{I.6.3.5}, luego $A_{ij}$ es un álgebra de tipo finito sobre $A_i$ \sref{I.6.3.3}.

\begin{env}[8.3.2]
\label{I.8.3.2}
Las estructuras cuyos axiomas son (Sch.~1), (Sch.~2) y (Sch.~3) generalizan los ``esquemas'' en el sentido de C.~Chevalley, quien supone además que $R$ es una extensión de tipo finito de un cuerpo $K$, y que los $A_i$ son álgebras de tipo finito sobre $K$ (lo que hace superflua una parte de (Sch.~2)) \cite{I-1}.
Recíprocamente, si se tiene tal estructura sobre un conjunto $X'$, se le puede asociar un esquema íntegro $X$ utilizando las observaciones de \sref{I.8.2.6}: el espacio subyacente de $X$ es igual a $X'$ provisto de la topología definida en \sref{I.8.2.6}, y del haz $\mathscr{O}_X$
\oldpage[I]{169}
tal que $\Gamma(U,\mathscr{O}_X)=\bigcap_{x\in U}\mathscr{O}_x$ para todo abierto $U\subset X$, con la definición evidente de los homomorfismos de restricción.
Dejamos al lector la tarea de verificar que así se obtiene un esquema íntegro cuyos anillos locales son los elementos de $X'$;
no utilizaremos este resultado en lo que sigue.
\end{env}
% END source_es/ega1/ega1-8.tex

% BEGIN source_es/ega1/ega1-9.tex
\section{Complementos sobre los haces cuasicoherentes}
\label{section:I.9}

\subsection{Producto tensorial de haces cuasicoherentes}
\label{subsection:I.9.1}

\begin{proposition}[9.1.1]
\label{I.9.1.1}
Sea $X$ un preesquema \textup{(resp.} un preesquema localmente noetheriano). Sean $\mathscr{F}$ y
$\mathscr{G}$ dos $\mathscr{O}_X$-módulos cuasicoherentes \textup{(resp.} coherentes); entonces
$\mathscr{F}\otimes_{\mathscr{O}_X}\mathscr{G}$ es cuasicoherente \textup{(resp.} coherente); además, es
de tipo finito si $\mathscr{F}$ y $\mathscr{G}$ son de tipo finito. Si
$\mathscr{F}$ admite una presentación finita y $\mathscr{G}$ es cuasicoherente
\textup{(resp.} coherente), entonces $\mathscr{H}\!om_{\mathscr{O}_X}(\mathscr{F},\mathscr{G})$ es cuasicoherente
\textup{(resp.} coherente).
\end{proposition}

\begin{proof}
Siendo la cuestión local, se puede suponer que $X$ es afín (resp. afín
noetheriano); además, si $\mathscr{F}$ es coherente, se puede suponer que es el
conúcleo de un homomorfismo $\mathscr{O}_X^m\to\mathscr{O}_X^n$. Las afirmaciones relativas a
los haces cuasicoherentes resultan de \sref{I.1.3.12} y \sref{I.1.3.9}; las
relativas a los haces coherentes resultan de \sref{I.1.5.1} y del hecho de
que, si $M$ y $N$ son módulos de tipo finito sobre un anillo noetheriano $A$,
entonces $M\otimes_A N$ y $\operatorname{Hom}_A(M,N)$ son ambos $A$-módulos de tipo finito.
\end{proof}

\begin{definition}[9.1.2]
\label{I.9.1.2}
{\itshape Sean $X$ e $Y$ dos $S$-preesquemas, $p$ y
$q$ las proyecciones de $X\times_S Y$, y $\mathscr{F}$ \textup{(resp. $\mathscr{G}$)} un
$\mathscr{O}_X$-módulo cuasicoherente \textup{(resp.} un $\mathscr{O}_Y$-módulo cuasicoherente). Se llama
producto tensorial de $\mathscr{F}$ y $\mathscr{G}$ sobre $\mathscr{O}_S$ \textup{(o} sobre $S$),
y se denota por $\mathscr{F}\otimes_{\mathscr{O}_S}\mathscr{G}$
\textup{(o} $\mathscr{F}\otimes_S\mathscr{G}$), al producto tensorial
$p^*(\mathscr{F})\otimes_{\mathscr{O}_{X\times_S Y}}q^*(\mathscr{G})$ sobre el preesquema $X\times_S Y$.}
\end{definition}

Si los $X_i$ ($1\leq i\leq n$) son $S$-preesquemas, y los $\mathscr{F}_i$ ($1\leq i\leq n$) son
$\mathscr{O}_{X_i}$-módulos cuasicoherentes, se define de la misma manera el producto tensorial
$\mathscr{F}_1\otimes_S\mathscr{F}_2\otimes_S\ldots\otimes_S\mathscr{F}_n$ sobre el
preesquema $Z=X_1\times_S X_2\times_S\ldots\times_S X_n$; es un
$\mathscr{O}_Z$-módulo \emph{cuasicoherente} en virtud de \sref{I.9.1.1} y
\sref[0]{0.5.1.4}; es \emph{coherente} si los $\mathscr{F}_i$ son coherentes \emph{y}
$Z$ es \emph{localmente noetheriano}, en virtud de \sref{I.9.1.1},
\sref[0]{0.5.3.11} y \sref{I.6.1.1}.

Se observará que, si se toma $X=Y=S$, la Definición~\sref{I.9.1.2} da de nuevo el producto
tensorial de $\mathscr{O}_S$-módulos. Además, como $q^*(\mathscr{O}_Y)=\mathscr{O}_{X\times_S Y}$
\sref[0]{0.4.3.4}, el producto $\mathscr{F}\otimes_S\mathscr{O}_Y$ se identifica
canónicamente con $p^*(\mathscr{F})$, y, de la misma manera,
$\mathscr{O}_X\otimes_S\mathscr{G}$ se identifica canónicamente con $q^*(\mathscr{G})$. En
particular, si se toma $Y=S$ y se denota por $f$ el morfismo estructural $X\to Y$,
se tiene $\mathscr{O}_X\otimes_Y\mathscr{G}=f^*(\mathscr{G})$: el producto tensorial ordinario
y la imagen inversa aparecen, por tanto, como casos particulares del producto
tensorial general.

La Definición~\sref{I.9.1.2} implica inmediatamente que, para $X$ e
$Y$ fijos, $\mathscr{F}\otimes_S\mathscr{G}$ es un \emph{bifuntor covariante aditivo y exacto por la derecha} en $\mathscr{F}$ y $\mathscr{G}$.

\begin{proposition}[9.1.3]
\label{I.9.1.3}
Sean $S$, $X$ e $Y$ tres esquemas afines de anillos,
respectivamente, $A$, $B$ y $C$, siendo $B$ y $C$ álgebras sobre $A$. Sea $M$
\textup{(resp. $N$)} un $B$-módulo \textup{(resp.} un $C$-módulo), y sea
$\mathscr{F}=\widetilde{M}$ \textup{(resp.} $\mathscr{G}=\widetilde{N}$) el
haz cuasicoherente asociado; entonces $\mathscr{F}\otimes_S\mathscr{G}$ es
canónicamente isomorfo al haz asociado al $(B\otimes_A C)$-módulo
$M\otimes_A N$.
\end{proposition}

\begin{proof}
\oldpage[I]{170}
En efecto, en virtud de \sref{I.1.6.5}, $\mathscr{F}\otimes_S\mathscr{G}$
es canónicamente isomorfo al haz asociado al $(B\otimes_A C)$-módulo
\[
  \bigl(M\otimes_B(B\otimes_A C)\bigr)
  \otimes_{B\otimes_A C}
  \bigl((B\otimes_A C)\otimes_C N\bigr)
\]
y, por los isomorfismos canónicos entre productos tensoriales, este último módulo es isomorfo a
\[
  M\otimes_B(B\otimes_A C)\otimes_C N
  =(M\otimes_B B)\otimes_A(C\otimes_C N)=M\otimes_A N.
\]
\end{proof}

\begin{proposition}[9.1.4]
\label{I.9.1.4}
Sean $f:T\to X$ y $g:T\to Y$ dos $S$-morfismos, y $\mathscr{F}$ \textup{(resp. $\mathscr{G}$)} un $\mathscr{O}_X$-módulo cuasicoherente \textup{(resp.} un $\mathscr{O}_Y$-módulo cuasicoherente).
Entonces
\[
  (f,g)_S^*(\mathscr{F}\otimes_S\mathscr{G})
  =f^*(\mathscr{F})\otimes_{\mathscr{O}_T}g^*(\mathscr{G}).
\]
\end{proposition}

\begin{proof}
Si $p$, $q$ son las proyecciones de $X\times_S Y$, la fórmula se deduce
de las igualdades $(f,g)_S^*\circ p^*=f^*$ y
$(f,g)_S^*\circ q^*=g^*$ \sref[0]{0.3.5.5}, y del hecho de que la imagen
inversa de un producto tensorial de haces algebraicos es el producto tensorial de sus imágenes
inversas \sref[0]{0.4.3.3}.
\end{proof}

\begin{corollary}[9.1.5]
\label{I.9.1.5}
Sean $f:X\to X'$ y $g:Y\to Y'$
dos $S$-morfismos, y $\mathscr{F}'$ \textup{(resp. $\mathscr{G}'$)} un
$\mathscr{O}_{X'}$-módulo cuasicoherente \textup{(resp.} un $\mathscr{O}_{Y'}$-módulo cuasicoherente). Entonces
\[
  (f\times_S g)^*(\mathscr{F}'\otimes_S\mathscr{G}')
  =f^*(\mathscr{F}')\otimes_S g^*(\mathscr{G}').
\]
\end{corollary}

\begin{proof}
Esto se deduce de \sref{I.9.1.4} y del hecho de que $f\times_S g=(f\circ p,g\circ q)_S$,
donde $p$ y $q$ son las proyecciones de $X\times_S Y$.
\end{proof}

\begin{corollary}[9.1.6]
\label{I.9.1.6}
Sean $X$, $Y$ y $Z$ tres $S$-preesquemas, y $\mathscr{F}$ \textup{(resp. $\mathscr{G}$, $\mathscr{H}$)} un
$\mathscr{O}_X$-módulo cuasicoherente \textup{(resp.} un $\mathscr{O}_Y$-módulo cuasicoherente, un
$\mathscr{O}_Z$-módulo cuasicoherente); entonces el haz $\mathscr{F}\otimes_S\mathscr{G}\otimes_S\mathscr{H}$ es la imagen inversa
de $(\mathscr{F}\otimes_S\mathscr{G})\otimes_S\mathscr{H}$ por el isomorfismo canónico de
$X\times_S Y\times_S Z$ sobre $(X\times_S Y)\times_S Z$.
\end{corollary}

\begin{proof}
Este isomorfismo se escribe $(p_1,p_2)_S\times_S p_3$, donde $p_1$, $p_2$ y $p_3$
son las proyecciones de $X\times_S Y\times_S Z$.

Análogamente, la imagen inversa de $\mathscr{G}\otimes_S\mathscr{F}$ por el isomorfismo canónico de
$X\times_S Y$ sobre $Y\times_S X$ es $\mathscr{F}\otimes_S\mathscr{G}$.
\end{proof}

\begin{corollary}[9.1.7]
\label{I.9.1.7}
Si $X$ es un $S$-preesquema, todo $\mathscr{O}_X$-módulo cuasicoherente $\mathscr{F}$ es la imagen
inversa de $\mathscr{F}\otimes_S\mathscr{O}_S$ por el isomorfismo canónico de $X$ sobre $X\times_S S$
\sref{I.3.3.3}.
\end{corollary}

\begin{proof}
Este isomorfismo es $(1_X,\varphi)_S$, donde $\varphi$ es el morfismo estructural
$X\to S$, y el corolario se deduce de \sref{I.9.1.4} y del hecho de que
$\varphi^*(\mathscr{O}_S)=\mathscr{O}_X$.
\end{proof}

\begin{env}[9.1.8]
\label{I.9.1.8}
Sea $X$ un $S$-preesquema, $\mathscr{F}$ un $\mathscr{O}_X$-módulo
cuasicoherente y $\varphi:S'\to S$ un morfismo; se denota por
$\mathscr{F}_{(\varphi)}$ o $\mathscr{F}_{(S')}$ el haz cuasicoherente
$\mathscr{F}\otimes_S\mathscr{O}_{S'}$ sobre $X\times_S S'=X_{(\varphi)}=X_{(S')}$; por tanto,
$\mathscr{F}_{(S')}=p^*(\mathscr{F})$, donde $p$ es la proyección $X_{(S')}\to X$.
\end{env}

\begin{proposition}[9.1.9]
\label{I.9.1.9}
Sea $\varphi':S''\to S'$ un morfismo.
Para todo $\mathscr{O}_X$-módulo cuasicoherente $\mathscr{F}$ sobre el $S$-preesquema $X$,
$(\mathscr{F}_{(\varphi)})_{(\varphi')}$ es la imagen inversa de
$\mathscr{F}_{(\varphi\circ\varphi')}$ por el isomorfismo canónico
$(X_{(\varphi)})_{(\varphi')}\xrightarrow{\sim}
X_{(\varphi\circ\varphi')}$
\sref{I.3.3.9}.
\end{proposition}

\begin{proof}
Esto se deduce inmediatamente de las definiciones y de \sref{I.3.3.9}, y se
escribe también
\[
  (\mathscr{F}\otimes_S\mathscr{O}_{S'})
  \otimes_{S'}\mathscr{O}_{S''}
  =\mathscr{F}\otimes_S\mathscr{O}_{S''}.
  \tag{9.1.9.1}\label{I.9.1.9.1-fr}
\]
\end{proof}

\begin{proposition}[9.1.10]
\label{I.9.1.10}
Sea $Y$ un $S$-preesquema, y sea $f:X\to Y$ un $S$-morfismo.
Para todo $\mathscr{O}_Y$-módulo cuasicoherente $\mathscr{G}$ y todo morfismo
$S'\to S$, se tiene
\[
  (f_{(S')})^*(\mathscr{G}_{(S')})
  =(f^*(\mathscr{G}))_{(S')}.
\].
\end{proposition}

\begin{proof}
Esto se deduce inmediatamente de la conmutatividad del diagrama
\oldpage[I]{171}
\[
  \xymatrix{
    X_{(S')}\ar[r]^{f_{(S')}}\ar[d] &
    Y_{(S')}\ar[d]\\
    X\ar[r]_f &
    Y
  }
\]
\end{proof}

\begin{corollary}[9.1.11]
\label{I.9.1.11}
Sean $X$ e $Y$ dos $S$-preesquemas, y sea
$\mathscr{F}$ \textup{(resp. $\mathscr{G}$)} un $\mathscr{O}_X$-módulo cuasicoherente
\textup{(resp.} un $\mathscr{O}_Y$-módulo cuasicoherente). La imagen inversa del haz
$(\mathscr{F}_{(S')})\otimes_{S'}(\mathscr{G}_{(S')})$ por el isomorfismo canónico
$(X\times_S Y)_{(S')}\xrightarrow{\sim}
(X_{(S')})\times_{S'}(Y_{(S')})$
\sref{I.3.3.10} es igual a $(\mathscr{F}\otimes_S\mathscr{G})_{(S')}$.
\end{corollary}

\begin{proof}
Si $p$ y $q$ son las proyecciones de $X\times_S Y$, el isomorfismo considerado
no es otro que $(p_{(S')},q_{(S')})_{S'}$; el corolario se deduce entonces de las
Proposiciones~\sref{I.9.1.4} y \sref{I.9.1.10}.
\end{proof}

\begin{proposition}[9.1.12]
\label{I.9.1.12}
Con las notaciones de la Definición~\sref{I.9.1.2}, sea $z$
un punto de $X\times_S Y$, y sean $x=p(z)$ e $y=q(z)$; el tallo
$(\mathscr{F}\otimes_S\mathscr{G})_z$ es isomorfo a
$(\mathscr{F}_x\otimes_{\mathscr{O}_x}\mathscr{O}_z)
\otimes_{\mathscr{O}_z}
(\mathscr{G}_y\otimes_{\mathscr{O}_y}\mathscr{O}_z)
=\mathscr{F}_x\otimes_{\mathscr{O}_x}\mathscr{O}_z
\otimes_{\mathscr{O}_y}\mathscr{G}_y$.
\end{proposition}

\begin{proof}
Como se puede reducir al caso afín, la proposición se deduce de la
fórmula~\sref{I.1.6.5.1}.
\end{proof}

\begin{corollary}[9.1.13]
\label{I.9.1.13}
Si $\mathscr{F}$ y $\mathscr{G}$ son de tipo finito, se tiene
\[
  \operatorname{Supp}(\mathscr{F}\otimes_S\mathscr{G})
  =p^{-1}(\operatorname{Supp}(\mathscr{F}))
  \cap q^{-1}(\operatorname{Supp}(\mathscr{G})).
\]
\end{corollary}

\begin{proof}
Como $p^*(\mathscr{F})$ y $q^*(\mathscr{G})$ son de tipo finito sobre
$\mathscr{O}_{X\times_S Y}$, se reduce, por la Proposición~\sref{I.9.1.12} y
por \sref[0]{0.1.7.5}, al caso particular en que $\mathscr{G}=\mathscr{O}_Y$, es
decir, a demostrar la fórmula siguiente:\footnote{\emph{[Trad.] El
original francés imprime $\Supp(p^{-1}(\sh{F}))$ en (9.1.13.1). La reducción
que acabamos de hacer identifica $\sh{F}\otimes_S\sh{O}_Y$ con $p^*(\sh{F})$, y la
comprobación siguiente $\sh{O}_z/\mathfrak{m}_x\sh{O}_z\neq0$ es el cálculo
de tallos correspondiente a esa imagen inversa; por ello escribimos $p^*(\sh{F})$.
La fórmula impresa es válida para la imagen inversa de haces de grupos abelianos; aquí formulamos la afirmación sobre la imagen inversa de módulos.}}
\[
  \operatorname{Supp}(p^*(\mathscr{F}))
  =p^{-1}(\operatorname{Supp}(\mathscr{F})).
  \tag{9.1.13.1}\label{I.9.1.13.1-fr}
\]

El mismo razonamiento que en \sref[0]{0.1.7.5} nos reduce a comprobar que, para todo
$z\in X\times_S Y$, se tiene $\mathscr{O}_z/\mathfrak{m}_x\mathscr{O}_z\neq0$ (con $x=p(z)$),
lo cual se deduce del hecho de que el homomorfismo $\mathscr{O}_x\to\mathscr{O}_z$ es \emph{local},
por hipótesis.
\end{proof}

Dejamos al lector la tarea de extender a un producto de un número cualquiera
de factores los resultados demostrados en este apartado para dos factores.

\subsection{Imagen directa de un haz cuasicoherente}
\label{subsection:I.9.2}

\begin{proposition}[9.2.1]
\label{I.9.2.1}
Sea $f:X\to Y$ un morfismo de preesquemas. Se supone que existe un
recubrimiento $(Y_\alpha)$ de $Y$ por abiertos afines con la siguiente
propiedad: cada $f^{-1}(Y_\alpha)$ admite un recubrimiento
\emph{finito} $(X_{\alpha i})$ por abiertos afines contenidos en
$f^{-1}(Y_\alpha)$, tal que cada una de las intersecciones
$X_{\alpha i}\cap X_{\alpha j}$ sea a su vez una unión \emph{finita} de
abiertos afines. En estas condiciones, para todo $\mathscr{O}_X$-módulo
cuasicoherente $\mathscr{F}$, $f_*(\mathscr{F})$ es un $\mathscr{O}_Y$-módulo
cuasicoherente.
\end{proposition}

\begin{proof}
Como la cuestión es local sobre $Y$, se puede suponer que $Y$ es igual a
uno de los $Y_\alpha$ y, por tanto, omitir los índices $\alpha$.
\begin{enumerate}
  \item[\emph{a)}] Supongamos primero que los $X_i\cap X_j$ son ellos mismos
    abiertos \emph{afines}. Pongamos $\mathscr{F}_i=\mathscr{F}|X_i$ y
    $\mathscr{F}_{ij}=\mathscr{F}|(X_i\cap X_j)$, y sean $\mathscr{F}'_i$ y
    $\mathscr{F}'_{ij}$ las imágenes directas de $\mathscr{F}_i$ y
    $\mathscr{F}_{ij}$ mediante las restricciones de $f$ a $X_i$ y a
    $X_i\cap X_j$, respectivamente; se sabe que $\mathscr{F}'_i$ y
    $\mathscr{F}'_{ij}$ son cuasicoherentes \sref{I.1.6.3}. Pongamos
    $\mathscr{G}=\bigoplus_i\mathscr{F}'_i$ y
    $\mathscr{H}=\bigoplus_{i,j}\mathscr{F}'_{ij}$; $\mathscr{G}$ y $\mathscr{H}$ son
    $\mathscr{O}_Y$-módulos cuasicoherentes. Definiremos un homomorfismo
    $u:\mathscr{G}\to\mathscr{H}$ tal que $f_*(\mathscr{F})$ sea el \emph{núcleo} de
    $u$; de ello se deducirá que $f_*(\mathscr{F})$ es cuasicoherente
    \sref{I.1.3.9}. Basta definir $u$
\oldpage[I]{172}
    como un homomorfismo de prehaces; teniendo en cuenta las definiciones
    de $\mathscr{G}$ y $\mathscr{H}$, basta, para cada abierto $W\subset Y$, definir
    un homomorfismo
    \[
  u_W:\bigoplus_i\Gamma(f^{-1}(W)\cap X_i,\mathscr{F})
  \longrightarrow
  \bigoplus_{i,j}\Gamma(f^{-1}(W)\cap X_i\cap X_j,\mathscr{F})
\]
    que satisfaga las condiciones habituales de compatibilidad al variar
    $W$. Si, para cada sección
    $s_i\in\Gamma(f^{-1}(W)\cap X_i,\mathscr{F})$, se denota por $s_{i|j}$ su
    restricción a $f^{-1}(W)\cap X_i\cap X_j$, se pone
    \[
  u_W((s_i))=(s_{i|j}-s_{j|i})
\]
    y las condiciones de compatibilidad se satisfacen de manera evidente.
    Para demostrar que el núcleo $\mathscr{R}$ de $u$ es $f_*(\mathscr{F})$,
    definamos un homomorfismo de $f_*(\mathscr{F})$ en $\mathscr{R}$ haciendo
    corresponder a cada sección
    $s\in\Gamma(f^{-1}(W),\mathscr{F})$ la familia $(s_i)$, donde $s_i$ es la
    restricción de $s$ a $f^{-1}(W)\cap X_i$; los axiomas (F1) y (F2) de
    los haces (G,~II,~1.1) implican que este homomorfismo es
    \emph{biyectivo}, lo que termina la demostración en este caso.
  \item[\emph{b)}] En el caso general, se aplica el mismo razonamiento una vez
    establecido que los $\mathscr{F}'_{ij}$ son cuasicoherentes. Ahora bien,
    por hipótesis, $X_i\cap X_j$ es una unión finita de abiertos afines
    $X_{ijk}$; y como los $X_{ijk}$ son abiertos afines \emph{en un
    esquema}, la intersección de dos cualesquiera de ellos vuelve a ser un
    abierto afín \sref{I.5.5.6}. Se reduce así al primer caso, y
    queda demostrada la Proposición~\sref{I.9.2.1}.
\end{enumerate}
\end{proof}

\begin{corollary}[9.2.2]
\label{I.9.2.2}
La conclusión de la Proposición~\sref{I.9.2.1} es válida en cada uno de los casos siguientes:
\begin{enumerate}
  \item[\textup{a)}] $f$ es separado y cuasicompacto;
  \item[\textup{b)}] $f$ es separado y de tipo finito;
  \item[\textup{c)}] $f$ es cuasicompacto, y el espacio subyacente a $X$ es localmente noetheriano.
\end{enumerate}
\end{corollary}

\begin{proof}
En el caso~\emph{a)}, los $X_{\alpha i}\cap X_{\alpha j}$ son afines \sref{I.5.5.6}.
El caso~\emph{b)} es un caso particular del caso~\emph{a)} \sref{I.6.6.3}. Por último, en el caso~\emph{c)}, se puede reducir al caso en que $Y$ es afín y el espacio subyacente a $X$ es noetheriano; entonces $X$ admite un recubrimiento finito de abiertos afines $(X_i)$, y los $X_i\cap X_j$, al ser cuasicompactos, son uniones finitas de abiertos afines \sref{I.2.1.3}.
\end{proof}

\subsection{Prolongación de secciones de haces cuasicoherentes}
\label{subsection:I.9.3}

\begin{theorem}[9.3.1]
\label{I.9.3.1}
Sea $X$ un preesquema cuyo espacio subyacente sea noetheriano, o un esquema cuyo espacio
subyacente sea cuasicompacto. Sea $\mathscr{L}$ un $\mathscr{O}_X$-módulo invertible \sref[0]{0.5.4.1}, $f$ una
sección de $\mathscr{L}$ sobre $X$, $X_f$ el abierto formado por los $x\in X$ tales que $f(x)\neq0$
\sref[0]{0.5.5.1}, y $\mathscr{F}$ un $\mathscr{O}_X$-módulo cuasicoherente.
\begin{enumerate}
  \item[{\rm(i)}] Si $s\in\Gamma(X,\mathscr{F})$ es tal que $s|X_f=0$, existe un entero $n>0$ tal que $s\otimes f^{\otimes n}=0$.
  \item[{\rm(ii)}] Para toda sección $s\in\Gamma(X_f,\mathscr{F})$, existe un entero $n>0$ tal que $s\otimes f^{\otimes n}$ se prolonga a una sección de $\mathscr{F}\otimes\mathscr{L}^{\otimes n}$ sobre $X$.
\end{enumerate}
\end{theorem}

\begin{proof}
\medskip\noindent
\begin{enumerate}
  \item[(i)] Como el espacio subyacente a $X$ es cuasicompacto y, por tanto, unión de
    un número finito de abiertos afines $U_i$ tales que $\mathscr{L}|U_i$ es isomorfo a
    $\mathscr{O}_X|U_i$, se puede reducir al caso en que $X$ es afín y $\mathscr{L}=\mathscr{O}_X$.
    En este caso, $f$ se identifica con un elemento de $A(X)$ y se tiene
    $X_f=D(f)$; $s$ se identifica con un elemento de un $A(X)$-módulo $M$, y
    $s|X_f$ con el elemento correspondiente de $M_f$; el resultado es entonces trivial,
    recordando la definición de módulo de fracciones.
\oldpage[I]{173}
  \item[(ii)] De nuevo, $X$ es unión finita de abiertos afines $U_i$ ($1\leq i\leq r$)
    tales que $\mathscr{L}|U_i\cong\mathscr{O}_X|U_i$, y, para todo $i$,
    $(s\otimes f^{\otimes n})|(U_i\cap X_f)$ se identifica (mediante el isomorfismo
    precedente) con $(f|(U_i\cap X_f))^n(s|(U_i\cap X_f))$. Se sabe entonces
    \sref{I.1.4.1} que existe un entero $n>0$ tal que, para todo
    $i$, $(s\otimes f^{\otimes n})|(U_i\cap X_f)$ se prolonga a una sección $s_i$ de
    $\mathscr{F}\otimes\mathscr{L}^{\otimes n}$ sobre $U_i$. Sea $s_{i|j}$ la restricción
    de $s_i$ a $U_i\cap U_j$; por definición, se tiene $s_{i|j}-s_{j|i}=0$ en
    $X_f\cap U_i\cap U_j$. Ahora bien, si el espacio subyacente a $X$ es noetheriano, $U_i\cap U_j$ es
    cuasicompacto; si $X$ es un esquema, $U_i\cap U_j$ es un abierto afín
    \sref{I.5.5.6} y, por tanto, también cuasicompacto. En virtud de (i), existe, pues,
    un entero $m$ (independiente de $i$ y $j$) tal que
    $(s_{i|j}-s_{j|i})\otimes f^{\otimes m}=0$. Se deduce inmediatamente que
    existe una sección $s'$ de $\mathscr{F}\otimes\mathscr{L}^{\otimes(n+m)}$ sobre $X$ que
    induce $s_i\otimes f^{\otimes m}$ sobre cada $U_i$, e induce
    $s\otimes f^{\otimes(n+m)}$ sobre $X_f$.
\end{enumerate}
\end{proof}

Los corolarios que siguen dan una interpretación del Teorema~\sref{I.9.3.1} en un lenguaje
más algebraico:
\begin{corollary}[9.3.2]
\label{I.9.3.2}
Bajo las hipótesis de \sref{I.9.3.1}, consideremos el anillo graduado $A_* = \Gamma_*(\mathscr{L})$
y el $A_*$-módulo graduado $M_* = \Gamma_*(\mathscr{L},\mathscr{F})$ \sref[0]{0.5.4.6}. Si $f\in A_n$,
donde $n\in\mathbb{Z}$, se tiene un isomorfismo canónico
\[
  \Gamma(X_f,\mathscr{F})\xrightarrow{\sim}((M_*)_f)_0
\] \textup{(el subgrupo del módulo de fracciones
$(M_*)_f$ formado por los elementos de grado $0$)}.
\end{corollary}

\begin{corollary}[9.3.3]
\label{I.9.3.3}
Se suponen verificadas las hipótesis de \sref{I.9.3.1}, y Se suponen además que
$\mathscr{L}=\mathscr{O}_X$. Entonces, poniendo $A=\Gamma(X,\mathscr{O}_X)$ y $M=\Gamma(X,\mathscr{F})$, el $A_f$-módulo
$\Gamma(X_f,\mathscr{F})$ es canónicamente isomorfo a $M_f$.
\end{corollary}

\begin{proposition}[9.3.4]
\label{I.9.3.4}
Sean $X$ un preesquema noetheriano, $\mathscr{F}$ un $\mathscr{O}_X$-módulo coherente y $\mathscr{J}$ un
haz coherente de ideales de $\mathscr{O}_X$, tales que el soporte de $\mathscr{F}$ esté contenido en el
de $\mathscr{O}_X/\mathscr{J}$. Entonces existe un entero $n>0$ tal que $\mathscr{J}^n\mathscr{F}=0$.
\end{proposition}

\begin{proof}
Como $X$ es unión finita de abiertos afines cuyos anillos son noetherianos, se puede suponer
que $X$ es afín, de anillo $A$ noetheriano; entonces $\mathscr{F}=\widetilde M$, donde
$M=\Gamma(X,\mathscr{F})$ es un $A$-módulo de tipo finito, y $\mathscr{J}=\widetilde{\mathfrak{J}}$,
donde $\mathfrak{J}=\Gamma(X,\mathscr{J})$ es un ideal de $A$ (\sref{I.1.4.1} y
\sref{I.1.5.1}). Como $A$ es noetheriano, $\mathfrak{J}$ admite un sistema finito de
generadores $f_i$ ($1\leq i\leq m$). Por hipótesis, toda sección de $\mathscr{F}$ sobre
$X$ es nula en cada uno de los $D(f_i)$; si $s_j$ ($1\leq j\leq q$) son secciones de
$\mathscr{F}$ que engendran $M$, existe entonces un entero $h$, independiente de $i$ y $j$,
tal que $f_i^h s_j=0$ \sref{I.1.4.1}, de donde $f_i^h s=0$ para todo $s\in M$. Se concluye
que, si $n=mh$, se tiene $\mathfrak{J}^nM=0$, y por consiguiente el $\mathscr{O}_X$-módulo correspondiente
$\mathscr{J}^n\mathscr{F}=(\mathfrak{J}^nM)\supertilde$ \sref{I.1.3.13} es nulo.
\end{proof}

\begin{corollary}[9.3.5]
\label{I.9.3.5}
Con las hipótesis de \sref{I.9.3.4}, existe un subpreesquema cerrado $Y$ de $X$ cuyo
espacio subyacente es el soporte de $\mathscr{O}_X/\mathscr{J}$ y tal que, si $j:Y\to X$ es la
inyección canónica, se tiene $\mathscr{F}=j_*(j^*(\mathscr{F}))$.
\end{corollary}

\begin{proof}
Observemos primero que los soportes de $\mathscr{O}_X/\mathscr{J}$ y de $\mathscr{O}_X/\mathscr{J}^n$ son
los mismos, pues, si $\mathscr{J}_x=\mathscr{O}_x$, se tiene también $\mathscr{J}_x^n=\mathscr{O}_x$, y,
por otra parte, $\mathscr{J}_x^n\subset\mathscr{J}_x$ para todo $x\in X$. Se puede, por tanto, en
virtud de \sref{I.9.3.4}, suponer que $\mathscr{J}\mathscr{F}=0$; se toma entonces por $Y$ el
subpreesquema cerrado de $X$ definido por $\mathscr{J}$, y, como $\mathscr{F}$ es entonces un
$(\mathscr{O}_X/\mathscr{J})$-módulo, la conclusión es inmediata.
\end{proof}

\oldpage[I]{174}
\subsection{Prolongación de los haces cuasicoherentes}
\label{subsection:I.9.4}

\begin{env}[9.4.1]
\label{I.9.4.1}
Sean $X$ un espacio topológico, $\mathscr{F}$ un haz de conjuntos (resp. de grupos, de anillos)
sobre $X$, $U$ un abierto de $X$, $\psi:U\to X$ la inyección canónica y $\mathscr{G}$ un subhaz
de $\mathscr{F}|U=\psi^*(\mathscr{F})$. Como $\psi_*$ es exacto por la izquierda, $\psi_*(\mathscr{G})$
es un subhaz de $\psi_*(\psi^*(\mathscr{F}))$; si se considera el homomorfismo canónico $\rho:\mathscr{F}\to\psi_*(\psi^*(\mathscr{F}))$ \sref[0]{0.3.5.3}, denotaremos por $\overline{\mathscr{G}}$ el subhaz
$\rho^{-1}(\psi_*(\mathscr{G}))$ de $\mathscr{F}$. Resulta
inmediatamente de las definiciones que, para todo abierto $V$ de $X$,
$\Gamma(V,\overline{\mathscr{G}})$ está formado por las secciones $s\in\Gamma(V,\mathscr{F})$ cuya
restricción a $V\cap U$ es una sección de $\mathscr{G}$ sobre $V\cap U$. Se tiene, por tanto,
$\overline{\mathscr{G}}|U=\psi^*(\overline{\mathscr{G}})=\mathscr{G}$, y $\overline{\mathscr{G}}$ es el
\emph{mayor} subhaz de $\mathscr{F}$ que induce $\mathscr{G}$ sobre $U$; diremos que
$\overline{\mathscr{G}}$ es la \emph{prolongación canónica} del subhaz $\mathscr{G}$ de $\mathscr{F}|U$
a un subhaz de $\mathscr{F}$.
\end{env}

\begin{proposition}[9.4.2]
\label{I.9.4.2}
Sean $X$ un preesquema y $U$ una parte abierta de $X$ tal que la inyección canónica
$j:U\to X$ sea un morfismo cuasicompacto \emph{(lo que se verificará para \emph{todo} $U$ si
el espacio subyacente a $X$ es \emph{localmente noetheriano}
\sref{I.6.6.4}[i])}. Entonces:
\begin{enumerate}
  \item[{\rm(i)}] para todo $(\mathscr{O}_X|U)$-módulo cuasicoherente $\mathscr{G}$, $j_*(\mathscr{G})$
    es un $\mathscr{O}_X$-módulo cuasicoherente, y
    $j_*(\mathscr{G})|U=j^*(j_*(\mathscr{G}))=\mathscr{G}$;
  \item[{\rm(ii)}] para todo $\mathscr{O}_X$-módulo cuasicoherente $\mathscr{F}$ y todo
    $(\mathscr{O}_X|U)$-submódulo cuasicoherente $\mathscr{G}$, la prolongación canónica
    $\overline{\mathscr{G}}$ de $\mathscr{G}$ \sref{I.9.4.1} es un
    $\mathscr{O}_X$-submódulo cuasicoherente de $\mathscr{F}$.
\end{enumerate}
\end{proposition}

\begin{proof}
Si $j=(\psi,\theta)$ (siendo $\psi$ la inyección $U\to X$ de los espacios subyacentes),
por definición se tiene $j_*(\mathscr{G})=\psi_*(\mathscr{G})$ para todo
$(\mathscr{O}_X|U)$-módulo $\mathscr{G}$, y además
$j^*(\mathscr{H})=\psi^*(\mathscr{H})=\mathscr{H}|U$ para todo $\mathscr{O}_X$-módulo $\mathscr{H}$, en virtud
de la definición del preesquema inducido sobre un abierto. Por tanto, (i) es un caso
particular de \hyperref[I.9.2.2]{(9.2.2, \emph{a)})}; por la misma razón, $j_*(j^*(\mathscr{F}))$ es
cuasicoherente y, como $\overline{\mathscr{G}}$ es la imagen inversa de $j_*(\mathscr{G})$ por el
homomorfismo $\rho:\mathscr{F}\to j_*(j^*(\mathscr{F}))$, (ii) resulta de \sref{I.4.1.1}.
\end{proof}

Se observará que la hipótesis de que el morfismo $j:U\to X$ sea cuasicompacto también se
verifica cuando el abierto $U$ es \emph{cuasicompacto} y $X$ es un \emph{esquema}: en
efecto, $U$ es entonces una unión finita de abiertos afines $U_i$, y, para todo abierto
afín $V$ de $X$, $V\cap U_i$ es un abierto afín \sref{I.5.5.6}, y por tanto
cuasicompacto.

\begin{corollary}[9.4.3]
\label{I.9.4.3}
Sean $X$ un preesquema y $U$ un abierto cuasicompacto de $X$ tal que el morfismo de
inyección $j:U\to X$ sea cuasicompacto. Supongamos además que todo $\mathscr{O}_X$-módulo
cuasicoherente sea límite inductivo de sus $\mathscr{O}_X$-submódulos cuasicoherentes de tipo
finito \emph{(lo que sucede cuando $X$ es un \emph{esquema afín})}. Sean entonces
$\mathscr{F}$ un $\mathscr{O}_X$-módulo cuasicoherente y $\mathscr{G}$ un
$(\mathscr{O}_X|U)$-submódulo cuasicoherente \textup{de tipo finito} de $\mathscr{F}|U$. Existe entonces
un $\mathscr{O}_X$-submódulo cuasicoherente \textup{de tipo finito} $\mathscr{G}'$ de $\mathscr{F}$ tal que
$\mathscr{G}'|U=\mathscr{G}$.
\end{corollary}

\begin{proof}
En efecto, se tiene $\mathscr{G}=\overline{\mathscr{G}}|U$, y $\overline{\mathscr{G}}$ es cuasicoherente
por \sref{I.9.4.2}, luego es límite inductivo de sus $\mathscr{O}_X$-submódulos cuasicoherentes
de tipo finito $\mathscr{H}_\lambda$. Por consiguiente, $\mathscr{G}$ es límite inductivo de los
$\mathscr{H}_\lambda|U$ y, por tanto, igual a uno de los $\mathscr{H}_\lambda|U$, puesto que es de
tipo finito \sref[0]{0.5.2.3}.
\end{proof}

\begin{remark}[9.4.4]
\label{I.9.4.4}
Supongamos que, para \emph{todo} abierto afín $U\subset X$, el morfismo de inyección
$U\to X$ sea cuasicompacto. Entonces, si la conclusión de \sref{I.9.4.3} se verifica
para todo abierto afín $U$ y todo $(\mathscr{O}_X|U)$-submódulo cuasicoherente de tipo finito
$\mathscr{G}$ de $\mathscr{F}|U$,
\oldpage[I]{175}
se sigue que $\mathscr{F}$ es límite inductivo de sus $\mathscr{O}_X$-submódulos cuasicoherentes
de tipo finito. En efecto, para todo abierto afín $U\subset X$, se tiene
$\mathscr{F}|U=\widetilde{M}$, donde $M$ es un $A(U)$-módulo, y, como este último es límite
inductivo de sus submódulos de tipo finito, $\mathscr{F}|U$ es límite inductivo de sus
$(\mathscr{O}_X|U)$-submódulos cuasicoherentes de tipo finito \sref{I.1.3.9}. Ahora bien,
por hipótesis, cada uno de estos submódulos está inducido sobre $U$ por un
$\mathscr{O}_X$-submódulo cuasicoherente de tipo finito $\mathscr{G}_{\lambda,U}$ de $\mathscr{F}$.
Las sumas finitas de los $\mathscr{G}_{\lambda,U}$ son asimismo $\mathscr{O}_X$-módulos
cuasicoherentes de tipo finito, pues la cuestión es local y el caso en que $X$ es afín
fue tratado en \sref{I.1.3.10}; entonces es claro que $\mathscr{F}$ es límite inductivo de
estas sumas finitas, de donde la afirmación.
\end{remark}

\begin{corollary}[9.4.5]
\label{I.9.4.5}
Bajo las hipótesis de \sref{I.9.4.3}, para todo $(\mathscr{O}_X|U)$-módulo cuasicoherente
de tipo finito $\mathscr{G}$, existe un $\mathscr{O}_X$-módulo cuasicoherente de tipo finito
$\mathscr{G}'$ tal que $\mathscr{G}'|U=\mathscr{G}$.
\end{corollary}

\begin{proof}
Como $\mathscr{F}=j_*(\mathscr{G})$ es cuasicoherente \sref{I.9.4.2} y
$\mathscr{F}|U=\mathscr{G}$, basta aplicar \sref{I.9.4.3} a $\mathscr{F}$.
\end{proof}

\begin{lemma}[9.4.6]
\label{I.9.4.6}
Sean $X$ un preesquema, $L$ un conjunto bien ordenado,
$(V_\lambda)_{\lambda\in L}$ un recubrimiento de $X$ por abiertos afines y $U$ un
abierto de $X$; para todo $\lambda\in L$, se pone
$W_\lambda=\bigcup_{\mu<\lambda}V_\mu$. Se supone que:
\begin{enumerate}
  \item[\textup{1\textsuperscript{o}}] Para todo $\lambda\in L$,
    $V_\lambda\cap W_\lambda$ es cuasicompacto;
  \item[\textup{2\textsuperscript{o}}] El morfismo de inmersión $U\to X$ es cuasicompacto.
\end{enumerate}
Entonces, para todo $\mathscr{O}_X$-módulo cuasicoherente $\mathscr{F}$ y todo
$(\mathscr{O}_X|U)$-submódulo cuasicoherente \textup{de tipo finito} $\mathscr{G}$ de $\mathscr{F}|U$, existe
un $\mathscr{O}_X$-submódulo cuasicoherente \textup{de tipo finito} $\mathscr{G}'$ de $\mathscr{F}$ tal que
$\mathscr{G}'|U=\mathscr{G}$.
\end{lemma}

\begin{proof}
Sea $\kappa$ el tipo de orden de $L$ e indexemos el recubrimiento dado por $\lambda<\kappa$.\footnote{\emph{[Ed.] La inducción impresa omite el caso inicial y, si $L$ tiene un elemento máximo, el paso final. Los pasos $0\leq\lambda\leq\kappa$ completan la demostración sin modificar las hipótesis.}}
Para $0\leq\lambda\leq\kappa$, pongamos $W_\lambda=\bigcup_{\mu<\lambda}V_\mu$ y $U_\lambda=U\cup W_\lambda$.
Se construye por inducción transfinita un $(\mathscr{O}_X|U_\lambda)$-submódulo cuasicoherente de tipo finito
$\mathscr{G}'_\lambda$ de $\mathscr{F}|U_\lambda$, tal que $\mathscr{G}'_\lambda|U_\mu=\mathscr{G}'_\mu$ para $\mu<\lambda$ y
$\mathscr{G}'_\lambda|U=\mathscr{G}$.
Se comienza con $\mathscr{G}'_0=\mathscr{G}$ sobre $U_0=U$; puesto que $U_\kappa=X$,
el módulo $\mathscr{G}'=\mathscr{G}'_\kappa$ proporcionará lo buscado.
Supongamos, pues, que $0<\lambda\leq\kappa$ y que los $\mathscr{G}'_\mu$ se han construido con estas propiedades
para $\mu<\lambda$.
Si $\lambda$ es un ordinal límite, se toma como $\mathscr{G}'_\lambda$ el único $(\mathscr{O}_X|U_\lambda)$-submódulo de
$\mathscr{F}|U_\lambda$ tal que $\mathscr{G}'_\lambda|U_\mu=\mathscr{G}'_\mu$ para todo $\mu<\lambda$.
Los $U_\mu$ con $\mu<\lambda$ recubren entonces $U_\lambda$, y la cuasicoherencia y el tipo finito son propiedades locales.
Si, por el contrario, $\lambda=\mu+1$, se tiene $U_\lambda=U_\mu\cup V_\mu$, y basta
definir un $(\mathscr{O}_X|V_\mu)$-submódulo cuasicoherente de tipo finito
$\mathscr{G}''_\mu$ de $\mathscr{F}|V_\mu$ tal que
\[
  \mathscr{G}''_\mu|(U_\mu\cap V_\mu)
  =\mathscr{G}'_\mu|(U_\mu\cap V_\mu)~;
\]
se tomará después como $\mathscr{G}'_\lambda$ el $(\mathscr{O}_X|U_\lambda)$-submódulo de
$\mathscr{F}|U_\lambda$ tal que $\mathscr{G}'_\lambda|U_\mu=\mathscr{G}'_\mu$ y
$\mathscr{G}'_\lambda|V_\mu=\mathscr{G}''_\mu$ \sref[0]{0.3.3.1}. Ahora bien, como $V_\mu$ es
afín, la existencia de $\mathscr{G}''_\mu$ queda asegurada por \sref{I.9.4.3} en cuanto se
haya demostrado que $U_\mu\cap V_\mu$ es cuasicompacto; pero $U_\mu\cap V_\mu$ es la
unión de $U\cap V_\mu$ y $W_\mu\cap V_\mu$, ambos cuasicompactos en virtud de las
hipótesis.
\end{proof}

\begin{theorem}[9.4.7]
\label{I.9.4.7}
Sean $X$ un preesquema y $U$ un abierto de $X$. Se supone verificada una de las
condiciones siguientes:
\begin{enumerate}
  \item[\textup{a)}] El espacio subyacente a $X$ es localmente noetheriano.
  \item[\textup{b)}] $X$ es un esquema cuasicompacto y $U$ un abierto cuasicompacto.
\end{enumerate}
Entonces, para todo $\mathscr{O}_X$-módulo cuasicoherente $\mathscr{F}$ y todo
$(\mathscr{O}_X|U)$-submódulo cuasicoherente \textup{de tipo finito} $\mathscr{G}$ de $\mathscr{F}|U$, existe
un $\mathscr{O}_X$-submódulo cuasicoherente \textup{de tipo finito} $\mathscr{G}'$ de $\mathscr{F}$ tal que
$\mathscr{G}'|U=\mathscr{G}$.
\end{theorem}

\oldpage[I]{176}
\begin{proof}
Sea $(V_\lambda)_{\lambda\in L}$ un recubrimiento de $X$ por abiertos afines,
suponiéndose $L$ finito en el caso \emph{b)}; dotando a $L$ de una estructura de conjunto
bien ordenado, basta comprobar que se satisfacen las condiciones de \sref{I.9.4.6}.
Esto es evidente bajo la hipótesis \emph{a)}, pues los espacios $V_\lambda$ son noetherianos.
Bajo la hipótesis \emph{b)}, los $V_\lambda\cap V_\mu$ son afines \sref{I.5.5.6}, y por
tanto cuasicompactos; como $L$ es finito, $V_\lambda\cap W_\lambda$ es cuasicompacto.
De aquí el teorema.
\end{proof}

\begin{corollary}[9.4.8]
\label{I.9.4.8}
Bajo las hipótesis de \sref{I.9.4.7}, para todo $(\mathscr{O}_X|U)$-módulo cuasicoherente
de tipo finito $\mathscr{G}$, existe un $\mathscr{O}_X$-módulo cuasicoherente de tipo finito
$\mathscr{G}'$ tal que $\mathscr{G}'|U=\mathscr{G}$.
\end{corollary}

\begin{proof}
Basta aplicar \sref{I.9.4.7} a $\mathscr{F}=j_*(\mathscr{G})$, que es cuasicoherente
\sref{I.9.4.2} y tal que $\mathscr{F}|U=\mathscr{G}$.
\end{proof}

\begin{corollary}[9.4.9]
\label{I.9.4.9}
Sea $X$ un preesquema cuyo espacio subyacente es localmente noetheriano, o un esquema
cuasicompacto. Entonces todo $\mathscr{O}_X$-módulo cuasicoherente es límite inductivo de
sus $\mathscr{O}_X$-submódulos cuasicoherentes de tipo finito.
\end{corollary}

\begin{proof}
Esto resulta de \sref{I.9.4.7} y de la observación \sref{I.9.4.4}.
\end{proof}

\begin{corollary}[9.4.10]
\label{I.9.4.10}
Bajo las hipótesis de \sref{I.9.4.9}, si un $\mathscr{O}_X$-módulo cuasicoherente
$\mathscr{F}$ es tal que todo $\mathscr{O}_X$-submódulo cuasicoherente de tipo finito de $\mathscr{F}$
está generado por sus secciones sobre $X$, entonces $\mathscr{F}$ está generado por sus
secciones sobre $X$.
\end{corollary}

\begin{proof}
Sea $U$ una vecindad abierta afín de un punto $x\in X$, y sea $s$ una sección de
$\mathscr{F}$ sobre $U$; el $\mathscr{O}_X$-submódulo $\mathscr{G}$ de $\mathscr{F}|U$ generado por $s$ es
cuasicoherente y de tipo finito, por lo que existe un $\mathscr{O}_X$-submódulo cuasicoherente
de tipo finito $\mathscr{G}'$ de $\mathscr{F}$ tal que $\mathscr{G}'|U=\mathscr{G}$ \sref{I.9.4.7}. Por
hipótesis, existen entonces un número finito de secciones $t_i$ de $\mathscr{G}'$ sobre $X$ y
secciones $a_i$ de $\mathscr{O}_X$ sobre una vecindad $V\subset U$ de $x$ tales que
$s|V=\sum_i a_i.(t_i|V)$, lo que demuestra el corolario.
\end{proof}

\subsection{Imagen cerrada de un preesquema; adherencia de un subpreesquema}
\label{subsection:I.9.5}

\begin{proposition}[9.5.1]
\label{I.9.5.1}
Sea $f:X\to Y$ un morfismo de preesquemas tal que $f_*(\mathscr{O}_X)$ sea un
$\mathscr{O}_Y$-módulo cuasicoherente (lo que ocurre si $f$ es cuasicompacto y, además,
o bien $f$ es separado o bien $X$ es localmente noetheriano \sref{I.9.2.2}). Entonces
existe un subpreesquema cerrado más pequeño $Y'$ de $Y$ tal que $f$ factoriza por la inyección
canónica $j:Y'\to Y$ \textup{(o, equivalentemente \sref{I.4.4.1}, tal que el
subpreesquema $f^{-1}(Y')$ de $X$ sea \emph{idéntico a} $X$)}.
\end{proposition}

Más precisamente:
\begin{corollary}[9.5.2]
\label{I.9.5.2}
Bajo las condiciones de \sref{I.9.5.1}, sean $f=(\psi,\theta)$ y $\mathscr{J}$ el
núcleo (cuasicoherente) del homomorfismo $\theta:\mathscr{O}_Y\to f_*(\mathscr{O}_X)$. Entonces el
subpreesquema cerrado $Y'$ de $Y$ definido por $\mathscr{J}$ satisface las condiciones de
\sref{I.9.5.1}.
\end{corollary}

\begin{proof}
Como el funtor $\psi^*$ es exacto, la factorización canónica
$\theta:\mathscr{O}_Y\to\mathscr{O}_Y/\mathscr{J}
\xrightarrow{\theta'}\psi_*(\mathscr{O}_X)$ proporciona (\textbf{0},~3.5.4.3)
una factorización
$\theta^\sharp:\psi^*(\mathscr{O}_Y)
\to\psi^*(\mathscr{O}_Y)/\psi^*(\mathscr{J})
\xrightarrow{{\theta'}^\sharp}\mathscr{O}_X$; como $\theta_x^\sharp$ es un homomorfismo local para
todo $x\in X$, lo mismo ocurre con ${\theta'_x}^\sharp$; si se denota por $\psi_0$ la
aplicación continua $\psi$ considerada como aplicación de $X$ a \erratum[II]{$Y'$}, y por $\theta_0$ la restricción
$\theta'|Y':(\mathscr{O}_Y/\mathscr{J})|Y'
\to\psi_*(\mathscr{O}_X)|Y'=(\psi_0)_*(\mathscr{O}_X)$, se ve que
$f_0=(\psi_0,\theta_0)$ es un morfismo de preesquemas $X\to Y'$ \sref{I.2.2.1} tal que
$f=j\circ f_0$. Ahora bien, si $X''$
\oldpage[I]{177}
es un segundo subpreesquema cerrado de $Y$, definido por un haz de ideales cuasicoherente $\mathscr{J}'$ de
$\mathscr{O}_Y$, y tal que $f$ factoriza por la inyección $j':X''\to Y$, se debe tener en primer lugar $X''\supset\psi(X)$ y, por tanto, $Y'\subset X''$, puesto que $X''$ es cerrado.
Además, para todo $y\in X''$, $\theta$ debe factorizar como
$\mathscr{O}_y\to\mathscr{O}_y/\mathscr{J}'_y
\to(\psi_*(\mathscr{O}_X))_y$, lo que, por definición, implica
$\mathscr{J}'_y\subset\mathscr{J}_y$, y, por consiguiente, $Y'$ es un subpreesquema cerrado de $X''$
\sref{I.4.1.10}.
\end{proof}

\begin{definition}[9.5.3]
\label{I.9.5.3}
{\itshape Cuando existe un subpreesquema cerrado más pequeño $Y'$ de $Y$ tal que $f$ factoriza
por la inyección canónica $j:Y'\to Y$, se dice que $Y'$ es el
preesquema \textup{imagen cerrada} de $X$ por el morfismo $f$.}
\end{definition}

\begin{proposition}[9.5.4]
\label{I.9.5.4}
Si $f_*(\mathscr{O}_X)$ es un $\mathscr{O}_Y$-módulo cuasicoherente, entonces el espacio subyacente
de la imagen cerrada de $X$ por $f$ es la adherencia $\overline{f(X)}$ en $Y$.
\end{proposition}

\begin{proof}
Como el soporte de $f_*(\mathscr{O}_X)$ está contenido en $\overline{f(X)}$, se tiene (con
la notación de \sref{I.9.5.2}) $\mathscr{J}_y=\mathscr{O}_y$ para
$y\notin\overline{f(X)}$, por lo que el soporte de $\mathscr{O}_Y/\mathscr{J}$ está contenido en
$\overline{f(X)}$. Además, este soporte es cerrado y contiene $f(X)$:
en efecto, si $y\in f(X)$, el elemento unidad del anillo $(\psi_*(\mathscr{O}_X))_y$ no es
cero, pues es el germen en $y$ de la sección
\[
  1\in\Gamma(X,\mathscr{O}_X)=\Gamma(Y,\psi_*(\mathscr{O}_X))~;
\]
como es la imagen por $\theta$ del elemento unidad de $\mathscr{O}_y$, este último
no pertenece a $\mathscr{J}_y$, por tanto $\mathscr{O}_y/\mathscr{J}_y\neq0$; esto concluye la
demostración.
\end{proof}

\begin{proposition}[9.5.5]
\label{I.9.5.5}
\emph{(Transitividad de las imágenes cerradas).} Sean $f:X\to Y$ y $g:Y\to Z$ dos
morfismos de preesquemas; se supone que la imagen cerrada $Y'$ de $X$ por $f$ existe y que,
si $g'$ es la restricción de $g$ a $Y'$, la imagen cerrada $Z'$ de $Y'$ por $g'$ existe.
Entonces la imagen cerrada de $X$ por $g\circ f$ existe y es igual a $Z'$.
\end{proposition}

\begin{proof}
Basta \sref{I.9.5.1} demostrar que $Z'$ es el subpreesquema cerrado más pequeño
$Z_1$ de $Z$ tal que el subpreesquema cerrado $(g\circ f)^{-1}(Z_1)$ de $X$
(igual a $f^{-1}(g^{-1}(Z_1))$ por el Corolario~\sref{I.4.4.2}) sea igual a
$X$; equivale a decir que $Z'$ es el subpreesquema cerrado más pequeño de $Z$ tal que
$f$ factoriza por la inyección $g^{-1}(Z_1)\to Y$ \sref{I.4.4.1}. En virtud de
la existencia de la imagen cerrada $Y'$, todo $Z_1$ con esta propiedad es tal que
$g^{-1}(Z_1)$ factoriza por $Y'$, lo que equivale a decir que
$j^{-1}(g^{-1}(Z_1))=g'^{-1}(Z_1)=Y'$ denotando por $j$ la inyección
$Y'\to Y$. Por definición de $Z'$, se concluye que $Z'$ es el subpreesquema cerrado
más pequeño de $Z$ que satisface la condición anterior.
\end{proof}

\begin{corollary}[9.5.6]
\label{I.9.5.6}
Sea $f:X\to Y$ un morfismo sobre $S$ tal que $Y$ sea la imagen cerrada de $X$ por $f$.
Sea $Z$ un esquema sobre $S$; si dos morfismos sobre $S$ $g_1$, $g_2$ de $Y$ en $Z$
satisfacen $g_1\circ f=g_2\circ f$, entonces $g_1=g_2$.
\end{corollary}

\begin{proof}
Sea $h=(g_1,g_2)_S:Y\to Z\times_S Z$; como la diagonal $T=\Delta_Z(Z)$ es un
subpreesquema cerrado de $Z\times_S Z$, $Y'=h^{-1}(T)$ es un subpreesquema cerrado de $Y$
\sref{I.4.4.1}. Pongamos $u=g_1\circ f=g_2\circ f$; entonces, por definición del producto,
$h'=h\circ f=(u,u)_S$, por lo que $h\circ f=\Delta_Z\circ u$; como
$\Delta_Z^{-1}(T)=Z$, se tiene $h'^{-1}(T)=u^{-1}(Z)=X$, y, por tanto, $f^{-1}(Y')=X$.
De aquí se concluye \sref{I.4.4.1} que $f$ factoriza por la inyección canónica $Y'\to Y$,
así que $Y'=Y$ por hipótesis; luego, por \sref{I.4.4.1}, $h$ se factoriza como
$\Delta_Z\circ v$, donde $v$ es un morfismo $Y\to Z$, lo que implica que $g_1=g_2=v$.
\end{proof}

\begin{remark}[9.5.7]
\label{I.9.5.7}
Si $X$ e $Y$ son $S$-esquemas, la proposición~\sref{I.9.5.6} implica que,
\oldpage[I]{178}
cuando $Y$ es la imagen cerrada de $X$ por $f$, $f$ es un \emph{epimorfismo} en la categoría de $S$-\emph{esquemas} (T,~1.1).
Mostraremos en el Capítulo~V que, recíprocamente, si existe la imagen cerrada $Y'$ de $X$
por $f$ y si $f$ es un epimorfismo de $S$-esquemas, entonces necesariamente $Y'=Y$.
\end{remark}

\begin{proposition}[9.5.8]
\label{I.9.5.8}
Supongamos verificadas las hipótesis de \sref{I.9.5.1}, y sea $Y'$ la imagen cerrada de $X$ por $f$. Para todo abierto $V$ de $Y$, sea
$f_V:f^{-1}(V)\to V$ la restricción de $f$; entonces la imagen cerrada de
$f^{-1}(V)$ por $f_V$ en $V$ existe y es igual al preesquema inducido por
$Y'$ en el abierto $V\cap Y'$ de $Y'$ \textup{(en otras palabras, al subpreesquema
$\inf(V,Y')$ de $Y$ \sref{I.4.4.3})}.
\end{proposition}

\begin{proof}
Pongamos $X'=f^{-1}(V)$; como la imagen directa de $\mathscr{O}_{X'}$ por
$f_V$ no es otra que la restricción de $f_*(\mathscr{O}_X)$ a $V$, es claro
que el núcleo $\mathscr{J}'$ del homomorfismo
$\mathscr{O}_V\to(f_V)_*(\mathscr{O}_{X'})$ es la restricción de
$\mathscr{J}$ a $V$, de donde se sigue inmediatamente la proposición.
\end{proof}

Se verá que este resultado puede entenderse diciendo que la formación de la imagen cerrada conmuta con una extensión $Y_1\to Y$ del preesquema base, que es una \emph{inmersión abierta}.
Veremos en el Capítulo~IV que lo mismo ocurre para una extensión $Y_1\to Y$ que es un morfismo \emph{plano}, siempre que $f$ sea separado y cuasicompacto.

\begin{proposition}[9.5.9]
\label{I.9.5.9}
Sea $f:X\to Y$ un morfismo tal que exista la imagen cerrada $Y'$ de $X$ por $f$.
\begin{enumerate}
  \item[{\rm(i)}] Si $X$ es reducido, también lo es $Y'$.
  \item[{\rm(ii)}] Si se satisfacen las hipótesis de la Proposición~\sref{I.9.5.1} y $X$ es irreducible \textup{(resp.} íntegro), también lo es $Y'$.
\end{enumerate}
\end{proposition}

\begin{proof}
Por hipótesis, el morfismo $f$ factoriza como $X\xrightarrow{g}Y'\xrightarrow{j}Y$, donde $j$ es la inyección canónica.
Como $X$ es reducido, $g$ factoriza como $X\xrightarrow{h}Y'_{\mathrm{red}}\xrightarrow{j'}Y'$ donde $j'$ es la inyección canónica \sref{I.5.2.2}, y se sigue entonces de la definición de $Y'$ que $Y'_{\mathrm{red}}=Y'$.
Si además se satisfacen las condiciones de la Proposición~\sref{I.9.5.1}, se sigue de \sref{I.9.5.4} que $f(X)$ es denso en $Y'$; si $X$ es irreducible, también lo es $Y'$ \sref[0]{0.2.1.5}.
La afirmación relativa a los preesquemas íntegros resulta de la conjunción de las otras dos.
\end{proof}

\begin{proposition}[9.5.10]
\label{I.9.5.10}
Sea $Y$ un subpreesquema de un preesquema $X$, tal que la inyección canónica $i:Y\to X$ sea un morfismo cuasicompacto.
Existe entonces un subpreesquema cerrado mínimo $\overline{Y}$ de $X$ que contiene a $Y$; su espacio subyacente es la adherencia del espacio subyacente de $Y$; este último es abierto en su adherencia, y el preesquema $Y$ está inducido sobre este abierto por $\overline{Y}$.
\end{proposition}

\begin{proof}
Basta aplicar la Proposición~\sref{I.9.5.1} a la inyección $i$\footnote{\emph{[Trad.] La fuente francesa imprime $j$ aquí, pero la proposición define la inyección canónica como $i:Y\to X$.}}, que es separada \sref{I.5.5.1} y cuasicompacta por hipótesis; \sref{I.9.5.1} prueba entonces la existencia de $\overline{Y}$, y \sref{I.9.5.4} muestra que su espacio subyacente es la adherencia de $Y$ en $X$; como $Y$ es localmente cerrado en $X$, es abierto en $\overline{Y}$, y la última afirmación proviene de aplicar \sref{I.9.5.8} a un abierto $V$ de $X$ tal que $Y$ sea cerrado en $V$.
\end{proof}

Con estas notaciones, si la inyección $V\to X$ es cuasicompacta, y si $\mathscr{J}$ es el haz cuasicoherente de ideales de $\mathscr{O}_X|V$ que define el subpreesquema cerrado $Y$ de $V$, se sigue de la Proposición~\sref{I.9.5.1} que el haz cuasicoherente de ideales de $\mathscr{O}_X$ que define $\overline{Y}$ es la prolongación canónica \sref{I.9.4.1} $\overline{\mathscr{J}}$ de $\mathscr{J}$, porque es claramente el subhaz cuasicoherente de ideales más grande de $\mathscr{O}_X$ que induce $\mathscr{J}$ sobre $V$.

\oldpage[I]{179}
\begin{corollary}[9.5.11]
\label{I.9.5.11}
Bajo las hipótesis de la Proposición~\sref{I.9.5.10}, toda sección de
$\mathscr{O}_{\overline{Y}}$ sobre un abierto $V$ de $\overline{Y}$ que sea nula en
$V\cap Y$ es nula.
\end{corollary}

\begin{proof}
En virtud de \sref{I.9.5.8}, se puede reducir al caso $V=\overline{Y}$.
Si se tiene en cuenta que las secciones de $\mathscr{O}_{\overline{Y}}$ sobre
$\overline{Y}$ corresponden canónicamente a las secciones sobre $\overline{Y}$ de
$\overline{Y}\otimes_{\mathbb{Z}}\mathbb{Z}[T]$ \sref{I.3.3.15} y que este último
es separado sobre $\overline{Y}$, el corolario aparece como un caso particular
de \sref{I.9.5.6}.
\end{proof}

Cuando existe un subpreesquema cerrado mínimo $Y'$ de $X$ que contiene un
subpreesquema $Y$ de $X$, se dice que $Y'$ es la \emph{adherencia} de $Y$ en
$X$, cuando no haya confusión.

\subsection{Haces cuasicoherentes de álgebras; cambio del haz estructural}
\label{subsection:I.9.6}
\label{I.9.6}

\begin{proposition}[9.6.1]
\label{I.9.6.1}
Sea $X$ un preesquema y $\mathscr{B}$ un álgebra cuasicoherente sobre $\mathscr{O}_X$ \sref[0]{0.5.1.3}.
Para que un $\mathscr{B}$-módulo $\mathscr{F}$ sea cuasicoherente (sobre el espacio anillado $(X,\mathscr{B})$), es necesario y suficiente que $\mathscr{F}$ sea un $\mathscr{O}_X$-módulo cuasicoherente.
\end{proposition}

\begin{proof}
Como la cuestión es local, se puede suponer que $X$ es afín, dado por un anillo $A$, y entonces $\mathscr{B}=\widetilde{B}$, donde $B$ es un álgebra sobre $A$ \sref{I.1.4.3}.
Si $\mathscr{F}$ es cuasicoherente sobre el espacio anillado $(X,\mathscr{B})$, se puede también suponer que $\mathscr{F}$ es el conúcleo de un $\mathscr{B}$-homomorfismo $\mathscr{B}^{(I)}\to\mathscr{B}^{(J)}$; como este homomorfismo es también un $\mathscr{O}_X$-homomorfismo de $\mathscr{O}_X$-módulos, y $\mathscr{B}^{(I)}$ y $\mathscr{B}^{(J)}$ son $\mathscr{O}_X$-módulos cuasicoherentes \sref{I.1.3.9}[ii], $\mathscr{F}$ es también un $\mathscr{O}_X$-módulo cuasicoherente \sref{I.1.3.9}[i].

Por el contrario, si $\mathscr{F}$ es un $\mathscr{O}_X$-módulo cuasicoherente, se tiene $\mathscr{F}=\widetilde{M}$, donde $M$ es un $B$-módulo \sref{I.1.4.3}; $M$ es isomorfo al conúcleo de un $B$-homomorfismo $B^{(I)}\to B^{(J)}$, por lo que $\mathscr{F}$ es un $\mathscr{B}$-módulo isomorfo al conúcleo del homomorfismo correspondiente $\mathscr{B}^{(I)}\to\mathscr{B}^{(J)}$ \sref{I.1.3.13}, lo que termina la demostración.
\end{proof}

En particular, si $\mathscr{F}$ y $\mathscr{G}$ son dos $\mathscr{B}$-módulos cuasicoherentes, entonces $\mathscr{F}\otimes_{\mathscr{B}}\mathscr{G}$ es un $\mathscr{B}$-módulo cuasicoherente; lo mismo ocurre con $\mathscr{H}\!om_{\mathscr{B}}(\mathscr{F},\mathscr{G})$ cuando además se supone que $\mathscr{F}$ admite una presentación finita \sref{I.1.3.13}.

\begin{env}[9.6.2]
\label{I.9.6.2}
Dado un preesquema $X$, diremos que una $\mathscr{O}_X$-álgebra cuasicoherente $\mathscr{B}$ es \emph{de tipo finito} si, para todo $x\in X$, existe una vecindad abierta \emph{afín} $U$ de $x$ tal que $\Gamma(U,\mathscr{B})=B$ es un álgebra de tipo finito sobre $\Gamma(U,\mathscr{O}_X)=A$.
Entonces $\mathscr{B}|U=\widetilde{B}$ y, para todo $f\in A$, la $(\mathscr{O}_X|D(f))$-álgebra inducida $\mathscr{B}|D(f)$ es de tipo finito, pues es isomorfa a $(B_f)\supertilde$, y $B_f=B\otimes_A A_f$ es claramente un álgebra de tipo finito sobre $A_f$.
Como los $D(f)$ forman una base de la topología de $U$, se concluye que si $\mathscr{B}$ es un álgebra cuasicoherente de $\mathscr{O}_X$ de tipo finito, entonces, para todo abierto $V$ de $X$, $\mathscr{B}|V$ es un álgebra cuasicoherente de $(\mathscr{O}_X|V)$ de tipo finito.
\end{env}

\begin{proposition}[9.6.3]
\label{I.9.6.3}
Sea $X$ un preesquema localmente noetheriano.
Entonces toda $\mathscr{O}_X$-álgebra cuasicoherente $\mathscr{B}$ de tipo finito es un haz coherente de anillos \sref[0]{0.5.3.7}.
\end{proposition}

\begin{proof}
Se puede limitar de nuevo al caso en que $X$ es un esquema afín dado por un anillo noetheriano $A$, y $\mathscr{B}=\widetilde{B}$, siendo $B$ un álgebra de tipo finito sobre $A$; por tanto, $B$ es un anillo noetheriano.
Con esto, queda por demostrar que el núcleo $\mathscr{N}$ de un $\mathscr{B}$-homomorfismo $\mathscr{B}^m\to\mathscr{B}$ es un $\mathscr{B}$-
\oldpage[I]{180}
módulo de tipo finito; pero es isomorfo (como $\mathscr{B}$-módulo) a $\widetilde{N}$, donde $N$ es el núcleo del homomorfismo correspondiente de $B$-módulos $B^m\to B$ \sref{I.1.3.13}.
Como $B$ es noetheriano, el submódulo $N$ de $B^m$ es un $B$-módulo de tipo finito, por lo que existe un homomorfismo $B^p\to B^m$ cuya imagen es $N$; como la sucesión $B^p\to B^m\to B$ es exacta, también lo es la sucesión correspondiente $\mathscr{B}^p\to\mathscr{B}^m\to\mathscr{B}$ \sref{I.1.3.5}, y como $\mathscr{N}$ es la imagen de $\mathscr{B}^p\to\mathscr{B}^m$ \sref{I.1.3.9}[i], queda demostrada la proposición.
\end{proof}

\begin{corollary}[9.6.4]
\label{I.9.6.4}
Bajo las hipótesis de \sref{I.9.6.3}, para que un $\mathscr{B}$-módulo $\mathscr{F}$ sea coherente, es necesario y suficiente que $\mathscr{G}$ sea un $\mathscr{O}_X$-módulo cuasicoherente y un $\mathscr{B}$-módulo de tipo finito.
Si este es el caso, y si $\mathscr{G}$ es un $\mathscr{B}$-submódulo o un $\mathscr{B}$-módulo cociente de $\mathscr{F}$, entonces, para que $\mathscr{G}$ sea un $\mathscr{B}$-módulo coherente, es necesario y suficiente que sea un $\mathscr{O}_X$-módulo cuasicoherente.
\end{corollary}

\begin{proof}
Teniendo en cuenta \sref{I.9.6.1}, las condiciones sobre $\mathscr{F}$ son claramente necesarias; mostremos que son suficientes.
Se puede limitar al caso en que $X$ es afín, dado por un anillo noetheriano $A$, donde $\mathscr{B}=\widetilde{B}$, con $B$ un álgebra de tipo finito sobre $A$, donde $\mathscr{F}=\widetilde{M}$, con $M$ un $B$-módulo, y donde existe un $\mathscr{B}$-homomorfismo sobreyectivo $\mathscr{B}^m\to\mathscr{F}\to 0$.
Se tiene entonces la sucesión exacta correspondiente $B^m\to M\to 0$, por lo que $M$ es un $B$-módulo de tipo finito; además, el núcleo $P$ del homomorfismo $B^m\to M$ es entonces un $B$-módulo de tipo finito, ya que $B$ es noetheriano.
Se concluye de \sref{I.1.3.13} que $\mathscr{F}$ es el conúcleo de un $\mathscr{B}$-homomorfismo $\mathscr{B}^n\to\mathscr{B}^m$, y por tanto es coherente, ya que $\mathscr{B}$ es un haz coherente de anillos \sref[0]{0.5.3.4}.
El mismo razonamiento muestra que un $\mathscr{B}$-submódulo cuasicoherente (resp. un $\mathscr{B}$-módulo cociente) de $\mathscr{F}$ es de tipo finito, de donde se sigue la segunda parte del corolario.
\end{proof}

\begin{proposition}[9.6.5]
\label{I.9.6.5}
Sea $X$ un esquema cuasicompacto, o un preesquema cuyo espacio subyacente es noetheriano.
Para toda $\mathscr{O}_X$-álgebra cuasicoherente $\mathscr{B}$ de tipo finito, existe un $\mathscr{O}_X$-submódulo cuasicoherente $\mathscr{E}$ de $\mathscr{B}$, de tipo finito, tal que $\mathscr{E}$ genera \sref[0]{0.4.1.3} la $\mathscr{O}_X$-álgebra $\mathscr{B}$.
\end{proposition}

\begin{proof}
En efecto, por hipótesis, existe un recubrimiento finito $(U_i)$ de $X$ formado por abiertos afines tales que $\Gamma(U_i,\mathscr{B})=B_i$ sea un álgebra de tipo finito sobre $\Gamma(U_i,\mathscr{O}_X)=A_i$; sea $E_i$ un $A_i$-submódulo de tipo finito de $B_i$ que genere el $A_i$-álgebra $B_i$; en virtud de \sref{I.9.4.7}, existe un $\mathscr{O}_X$-submódulo $\mathscr{E}_i$ de $\mathscr{B}$, cuasicoherente y de tipo finito, tal que $\mathscr{E}_i|U_i=\widetilde{E}_i$.
Es claro que la suma $\mathscr{E}$ de los $\mathscr{E}_i$ es el objeto buscado.
\end{proof}

\begin{proposition}[9.6.6]
\label{I.9.6.6}
Sea $X$ un preesquema cuyo espacio subyacente es localmente noetheriano, o un esquema cuasicompacto.
Entonces toda $\mathscr{O}_X$-álgebra cuasicoherente $\mathscr{B}$ es el límite inductivo de sus $\mathscr{O}_X$-subálgebras cuasicoherentes de tipo finito.
\end{proposition}

\begin{proof}
En efecto, se deduce de \sref{I.9.4.9} que $\mathscr{B}$ es el límite inductivo (como $\mathscr{O}_X$-módulo) de sus $\mathscr{O}_X$-sub\emph{módulos} cuasicoherentes de tipo finito; estos últimos generan $\mathscr{O}_X$-sub\emph{álgebras} cuasicoherentes de $\mathscr{B}$ de tipo finito \sref{I.1.3.14}, y por tanto $\mathscr{B}$ es \emph{a fortiori} su límite inductivo.
\end{proof}
% END source_es/ega1/ega1-9.tex

% BEGIN source_es/ega1/ega1-10.tex
\section{Esquemas formales}
\label{section:I.10}

\subsection{Esquemas formales afines}
\label{subsection:I.10.1}

\begin{env}[10.1.1]
\label{I.10.1.1}
Sea $A$ un anillo topológico \emph{admisible} \sref[0]{0.7.1.2}; para todo ideal de definición $\mathfrak{J}$ de $A$, $\operatorname{Spec}(A/\mathfrak{J})$ se identifica con el subespacio cerrado $V(\mathfrak{J})$ de $\operatorname{Spec}(A)$ \sref{I.1.1.11}, conjunto de los ideales primos \emph{abiertos} de $A$; este espacio topológico no depende
\oldpage[I]{181}
del ideal de definición $\mathfrak{J}$ considerado; denotémoslo $\mathfrak{X}$. Sea $(\mathfrak{J}_\lambda)$ un sistema fundamental de vecindades de 0 en $A$, formado por ideales de definición, y para cada $\lambda$, sea $\mathscr{O}_\lambda$ el haz estructural de $\operatorname{Spec}(A/\mathfrak{J}_\lambda)$; este haz es inducido sobre $\mathfrak{X}$ por $\widetilde{A}/\widetilde{\mathfrak{J}}_\lambda$ (que es nulo fuera de $\mathfrak{X}$).
Para $\mathfrak{J}_\mu\subset\mathfrak{J}_\lambda$, el homomorfismo canónico $A/\mathfrak{J}_\mu\to A/\mathfrak{J}_\lambda$ define por tanto un homomorfismo $u_{\lambda\mu}:\mathscr{O}_\mu\to\mathscr{O}_\lambda$ de haces de anillos \sref{I.1.6.1}, y $(\mathscr{O}_\lambda)$ es un \emph{sistema proyectivo de haces de anillos} para estos homomorfismos.
Como la topología de $\mathfrak{X}$ admite una base formada por abiertos cuasicompactos, se puede asociar a cada $\mathscr{O}_\lambda$ un \emph{haz pseudodiscreto de anillos topológicos} \sref[0]{0.3.8.1} que tiene $\mathscr{O}_\lambda$ como haz de anillos subyacente (sin topología), y que denotamos también por $\mathscr{O}_\lambda$; y los $\mathscr{O}_\lambda$ forman de nuevo un \emph{sistema proyectivo de haces de anillos topológicos} \sref[0]{0.3.8.2}.
Denotamos por $\mathscr{O}_{\mathfrak{X}}$ el \emph{haz de anillos topológicos} sobre $\mathfrak{X}$, límite proyectivo del sistema $(\mathscr{O}_\lambda)$; para todo abierto \emph{cuasicompacto} $U$ de $\mathfrak{X}$, $\Gamma(U,\mathscr{O}_{\mathfrak{X}})$ es por tanto el anillo topológico límite proyectivo del sistema de anillos \emph{discretos} $\Gamma(U,\mathscr{O}_\lambda)$ \sref[0]{0.3.2.6}.
\end{env}

\begin{definition}[10.1.2]
\label{I.10.1.2}
Dado un anillo topológico \emph{admisible} $A$, se llama espectro formal de $A$, y se lo denota por $\operatorname{Spf}(A)$, al subespacio cerrado $\mathfrak{X}$ de $\operatorname{Spec}(A)$ formado por los ideales primos \emph{abiertos} de $A$.
Se dice que un espacio topológicamente anillado es un \emph{esquema formal afín} si es isomorfo a un espectro formal $\operatorname{Spf}(A)=\mathfrak{X}$ provisto del haz de anillos topológicos $\mathscr{O}_{\mathfrak{X}}$, que es el límite proyectivo de haces de anillos topológicos pseudodiscretos $(\widetilde{A}/\widetilde{\mathfrak{J}}_\lambda)|\mathfrak{X}$, donde $\mathfrak{J}_\lambda$ recorre el conjunto filtrante de los ideales de definición de $A$.
\end{definition}

Cuando hablemos de un \emph{espectro formal $\mathfrak{X}=\operatorname{Spf}(A)$} como de un esquema formal afín, se tratará siempre del espacio topológicamente anillado $(\mathfrak{X},\mathscr{O}_{\mathfrak{X}})$, donde $\mathscr{O}_{\mathfrak{X}}$ está definido como arriba.

Se observará que todo \emph{esquema afín} $X=\operatorname{Spec}(A)$ puede considerarse como un esquema formal afín de una única manera, considerando $A$ como un anillo topológico discreto: los anillos topológicos $\Gamma(U,\mathscr{O}_X)$ son entonces discretos siempre que $U$ sea cuasicompacto (pero no, en general, cuando $U$ sea un abierto cualquiera de $X$).

\begin{proposition}[10.1.3]
\label{I.10.1.3}
Si $\mathfrak{X}=\operatorname{Spf}(A)$, donde $A$ es un anillo admisible, entonces $\Gamma(\mathfrak{X},\mathscr{O}_{\mathfrak{X}})$ es topológicamente isomorfo a $A$.
\end{proposition}

\begin{proof}
En efecto, como $\mathfrak{X}$ es cerrado en $\operatorname{Spec}(A)$, es cuasicompacto y, por tanto, $\Gamma(\mathfrak{X},\mathscr{O}_{\mathfrak{X}})$ es topológicamente isomorfo al límite proyectivo de los anillos discretos $\Gamma(\mathfrak{X},\mathscr{O}_\lambda)$; pero $\Gamma(\mathfrak{X},\mathscr{O}_\lambda)$ es isomorfo a $A/\mathfrak{J}_\lambda$ \sref{I.1.3.7}; como $A$ es separado y completo, es topológicamente isomorfo a $\varprojlim A/\mathfrak{J}_\lambda$ \sref[0]{0.7.2.1}, de donde se sigue la proposición.
\end{proof}

\begin{proposition}[10.1.4]
\label{I.10.1.4}
Sea $A$ un anillo admisible, $\mathfrak{X}=\operatorname{Spf}(A)$ y, para todo $f\in A$, sea $\mathfrak{D}(f)=D(f)\cap\mathfrak{X}$; entonces el espacio topológicamente anillado $(\mathfrak{D}(f),\mathscr{O}_{\mathfrak{X}}|\mathfrak{D}(f))$ es isomorfo al espectro formal afín $\operatorname{Spf}(A_{\{f\}})$ \sref[0]{0.7.6.15}.
\end{proposition}

\begin{proof}
Para todo ideal de definición $\mathfrak{J}$ de $A$, el anillo discreto $S_f^{-1}A/S_f^{-1}\mathfrak{J}$ se identifica canónicamente con $A_{\{f\}}/\mathfrak{J}_{\{f\}}$ \sref[0]{0.7.6.9}, de modo que, por \sref{I.1.2.5} y \sref{I.1.2.6}, el espacio topológico $\operatorname{Spf}(A_{\{f\}})$ se identifica canónicamente con $\mathfrak{D}(f)$.
Además, para todo abierto cuasicompacto $U$ de $\mathfrak{X}$ contenido en $\mathfrak{D}(f)$, $\Gamma(U,\mathscr{O}_\lambda)$ se puede identificar con el módulo de secciones del haz estructural de $\operatorname{Spec}(S_f^{-1}A/S_f^{-1}\mathfrak{J}_\lambda)$ sobre $U$ \sref{I.1.3.6}; por tanto, si se pone $\mathfrak{Y}=\operatorname{Spf}(A_{\{f\}})$, $\Gamma(U,\mathscr{O}_{\mathfrak{X}})$ se puede identificar con el módulo de secciones $\Gamma(U,\mathscr{O}_{\mathfrak{Y}})$, lo que demuestra la proposición.
\end{proof}

\oldpage[I]{182}
\begin{env}[10.1.5]
\label{I.10.1.5}
Como haz de anillos \emph{sin topología}, el haz estructural $\mathscr{O}_{\mathfrak{X}}$ de $\operatorname{Spf}(A)$ admite, para todo $x\in\mathfrak{X}$, un tallo que, por \sref{I.10.1.4}, se puede identificar con el límite inductivo $\varinjlim A_{\{f\}}$ para los $f\notin\mathfrak{j}_x$. Por tanto, según \sref[0]{0.7.6.17} y \sref[0]{0.7.6.18}:
\end{env}

\begin{proposition}[10.1.6]
\label{I.10.1.6}
Para todo $x\in\mathfrak{X}=\operatorname{Spf}(A)$, el tallo $\mathscr{O}_x$ es un anillo local cuyo campo residual es isomorfo a $\boldsymbol{k}(x)=A_x/\mathfrak{j}_xA_x$.
Si además $A$ es ádico y noetheriano, $\mathscr{O}_x$ es un anillo noetheriano.
\end{proposition}

Como $\boldsymbol{k}(x)$ no se reduce a 0, se concluye de este resultado que el \emph{soporte} del haz de anillos $\mathscr{O}_{\mathfrak{X}}$ es \emph{igual a $\mathfrak{X}$}.

\subsection{Morfismos de esquemas formales afines}
\label{subsection:I.10.2}

\begin{env}[10.2.1]
\label{I.10.2.1}
Sean $A$, $B$ dos anillos admisibles, y sea $\varphi:B\to A$ un homomorfismo \emph{continuo}.
La aplicación continua $ {}^a\varphi:\operatorname{Spec}(A)\to\operatorname{Spec}(B)$ \sref{I.1.2.1} lleva entonces $\mathfrak{X}=\operatorname{Spf}(A)$ a $\mathfrak{Y}=\operatorname{Spf}(B)$, pues la imagen inversa por $\varphi$ de un ideal primo abierto de $A$ es un ideal primo abierto de $B$.
Por otra parte, para todo $g\in B$,
$\varphi$ define un homomorfismo continuo
$\Gamma(\mathfrak{D}(g),\mathscr{O}_{\mathfrak{Y}})\to
\Gamma(\mathfrak{D}(\varphi(g)),\mathscr{O}_{\mathfrak{X}})$
en virtud de \sref{I.10.1.4},
\sref{I.10.1.3}
y \sref[0]{0.7.6.7}; como estos homomorfismos satisfacen las condiciones de compatibilidad para las restricciones correspondientes al paso de $g$ a un múltiplo de $g$, y como $\mathfrak{D}(\varphi(g))={}^a\varphi^{-1}(\mathfrak{D}(g))$, definen un homomorfismo \emph{continuo} de haces de anillos topológicos $\mathscr{O}_{\mathfrak{Y}}\to
{}^a\varphi_*(\mathscr{O}_{\mathfrak{X}})$ \sref[0]{0.3.2.5} que denotamos por $\widetilde{\varphi}$; así se ha obtenido un morfismo $\Phi=({}^a\varphi,\widetilde{\varphi})$ de espacios topológicamente anillados $\mathfrak{X}\to\mathfrak{Y}$.
Se observará que, como homomorfismo de haces de anillos sin topología, $\widetilde{\varphi}$ define un homomorfismo $\widetilde{\varphi}_x^\sharp:
\mathscr{O}_{{}^a\varphi(x)}\to\mathscr{O}_x$ sobre los tallos, para todo $x\in\mathfrak{X}$.
\end{env}

\begin{proposition}[10.2.2]
\label{I.10.2.2}
Sean $A$, $B$ dos anillos topológicos admisibles, y sean $\mathfrak{X}=\operatorname{Spf}(A)$, $\mathfrak{Y}=\operatorname{Spf}(B)$. Para que un morfismo $u=(\psi,\theta):\mathfrak{X}\to\mathfrak{Y}$ de espacios topológicamente anillados sea de la forma $({}^a\varphi,\widetilde{\varphi})$, donde $\varphi$ es un homomorfismo continuo de anillos $B\to A$, es necesario y suficiente que, para todo $x\in\mathfrak{X}$, $\theta_x^\sharp$ sea un homomorfismo local $\mathscr{O}_{\psi(x)}\to\mathscr{O}_x$.
\end{proposition}

\begin{proof}
La condición es necesaria: sean $\mathfrak{p}=\mathfrak{j}_x\in\operatorname{Spf}(A)$ y $\mathfrak{q}=\varphi^{-1}(\mathfrak{j}_x)$; si $g\notin\mathfrak{q}$, entonces $\varphi(g)\notin\mathfrak{p}$, y es inmediato que el homomorfismo $B_{\{g\}}\to A_{\{\varphi(g)\}}$ inducido por $\varphi$ \sref[0]{0.7.6.7} envía $\mathfrak{q}_{\{g\}}$ a un subconjunto de $\mathfrak{p}_{\{\varphi(g)\}}$; pasando al límite inductivo, se ve (teniendo en cuenta \sref{I.10.1.5} y \sref[0]{0.7.6.17}) que $\widetilde{\varphi}_x^\sharp$ es un homomorfismo local.

Recíprocamente, sea $(\psi,\theta)$ un morfismo que satisface la condición de la proposición; en virtud de \sref{I.10.1.3}, $\theta$ define un homomorfismo continuo de anillos
\[
  \varphi=\Gamma(\theta):
  B=\Gamma(\mathfrak{Y},\mathscr{O}_{\mathfrak{Y}})
  \longrightarrow
  \Gamma(\mathfrak{X},\mathscr{O}_{\mathfrak{X}})=A.
\]
En virtud de la hipótesis sobre $\theta$, para que la sección $\varphi(g)$ de $\mathscr{O}_{\mathfrak{X}}$ sobre $\mathfrak{X}$ tenga un germen invertible en el punto $x$, es necesario y suficiente que $g$ tenga un germen invertible en el punto $\psi(x)$. Pero, por \sref[0]{0.7.6.17}, las secciones de $\mathscr{O}_{\mathfrak{X}}$ (resp. $\mathscr{O}_{\mathfrak{Y}}$) sobre $\mathfrak{X}$ (resp. $\mathfrak{Y}$) cuyo germen no es invertible en el punto $x$ (resp. $\psi(x)$) son exactamente los elementos de $\mathfrak{j}_x$
\oldpage[I]{183}
(resp. $\mathfrak{j}_{\psi(x)}$); la observación anterior muestra, por tanto, que ${}^a\varphi=\psi$. Finalmente, para todo $g\in B$, el diagrama
\[
  \xymatrix{
    B=\Gamma(\mathfrak{Y},\mathscr{O}_{\mathfrak{Y}})
      \ar[r]^\varphi\ar[d] &
    \Gamma(\mathfrak{X},\mathscr{O}_{\mathfrak{X}})=A\ar[d]\\
    B_{\{g\}}=\Gamma(\mathfrak{D}(g),\mathscr{O}_{\mathfrak{Y}})
      \ar[r]_{\Gamma(\theta_{\mathfrak{D}(g)})} &
    \Gamma(\mathfrak{D}(\varphi(g)),\mathscr{O}_{\mathfrak{X}})
      =A_{\{\varphi(g)\}}
  }
\]
es conmutativo; según la propiedad universal de los anillos completos de fracciones \sref[0]{0.7.6.6}, $\theta_{\mathfrak{D}(g)}$ es igual a $\widetilde{\varphi}_{\mathfrak{D}(g)}$ para todo $g\in B$, y por \sref[0]{0.3.2.5} se tiene $\theta=\widetilde{\varphi}$.
\end{proof}

Diremos que un morfismo $(\psi,\theta)$ de espacios topológicamente anillados que satisface la condición de la Proposición~\sref{I.10.2.2} es un \emph{morfismo de esquemas formales afines}. Se puede decir que los funtores $\operatorname{Spf}(A)$ en $A$ y $\Gamma(\mathfrak{X},\mathscr{O}_{\mathfrak{X}})$ en $\mathfrak{X}$ definen una \emph{equivalencia} entre la categoría de anillos admisibles y la categoría dual de la categoría de esquemas formales afines (T, I, 1.2).

\begin{env}[10.2.3]
\label{I.10.2.3}
Como caso particular de \sref{I.10.2.2}, observemos que, para $f\in A$, la inyección canónica del esquema formal afín inducido por $\mathfrak{X}$ sobre $\mathfrak{D}(f)$ corresponde al homomorfismo canónico continuo $A\to A_{\{f\}}$.
Bajo las hipótesis de la Proposición~\sref{I.10.2.2}, sean $h$ un elemento de $B$ y $g$ un elemento de $A$ múltiplo de $\varphi(h)$;
se tiene entonces $\psi(\mathfrak{D}(g))\subset\mathfrak{D}(h)$; la restricción de $u$ a $\mathfrak{D}(g)$, considerada como morfismo de $\mathfrak{D}(g)$ en $\mathfrak{D}(h)$, es el único morfismo $v$ que hace conmutativo el diagrama
\[
  \xymatrix{
    \mathfrak{D}(g)\ar[r]^v\ar[d] &
    \mathfrak{D}(h)\ar[d]\\
    \mathfrak{X}\ar[r]_u &
    \mathfrak{Y}
  }
\]
Este morfismo corresponde al único homomorfismo continuo $\varphi':B_{\{h\}}\to A_{\{g\}}$ \sref[0]{0.7.6.7} que hace conmutativo el diagrama
\[
  \xymatrix{
    A\ar[d] &
    B\ar[l]_\varphi\ar[d]\\
    A_{\{g\}} &
    B_{\{h\}}\ar[l]^{\varphi'}
  }
\]
\end{env}

\subsection{Ideales de definición de un esquema formal afín}
\label{subsection:I.10.3}

\begin{env}[10.3.1]
\label{I.10.3.1}
Sean $A$ un anillo admisible, $\mathfrak{J}$ un ideal abierto de $A$ y $\mathfrak{X}$ el esquema formal afín $\operatorname{Spf}(A)$.
Sea $(\mathfrak{J}_\lambda)$ el conjunto de los ideales de definición de $A$ contenidos en $\mathfrak{J}$; entonces $\widetilde{\mathfrak{J}}/\widetilde{\mathfrak{J}}_\lambda$ es un haz de ideales de $\widetilde{A}/\widetilde{\mathfrak{J}}_\lambda$.
Designemos por $\mathfrak{J}^\Delta$ el límite proyectivo de los haces inducidos sobre $\mathfrak{X}$ por $\widetilde{\mathfrak{J}}/\widetilde{\mathfrak{J}}_\lambda$, que se identifica con un \emph{haz de ideales} de $\mathscr{O}_{\mathfrak{X}}$ \sref[0]{0.3.2.6}.
Para todo $f\in A$, $\Gamma(\mathfrak{D}(f),\mathfrak{J}^\Delta)$ es el límite proyectivo de los $S_f^{-1}\mathfrak{J}/S_f^{-1}\mathfrak{J}_\lambda$, es decir, se identifica con el ideal abierto $\mathfrak{J}_{\{f\}}$ del anillo $A_{\{f\}}$ \sref[0]{0.7.6.9}, y, en particular, $\Gamma(\mathfrak{X},\mathfrak{J}^\Delta)=\mathfrak{J}$; se deduce (pues los $\mathfrak{D}(f)$ forman una base de la topología de $\mathfrak{X}$) que
\[
  \mathfrak{J}^\Delta|\mathfrak{D}(f)
  =(\mathfrak{J}_{\{f\}})^\Delta
  \tag{10.3.1.1}\label{I.10.3.1.1-fr}
\]
\end{env}

\oldpage[I]{184}
\begin{env}[10.3.2]
\label{I.10.3.2}
Con las notaciones de \sref{I.10.3.1}, para todo $f\in A$, la aplicación canónica de $A_{\{f\}}=\Gamma(\mathfrak{D}(f),\mathscr{O}_{\mathfrak{X}})$ en $\Gamma\bigl(\mathfrak{D}(f),
(\widetilde{A}/\widetilde{\mathfrak{J}})|\mathfrak{X}\bigr)
=S_f^{-1}A/S_f^{-1}\mathfrak{J}$ es \emph{sobreyectiva} y tiene por núcleo $\Gamma(\mathfrak{D}(f),\mathfrak{J}^\Delta)=\mathfrak{J}_{\{f\}}$ \sref[0]{0.7.6.9}; estas aplicaciones definen, por tanto, un homomorfismo continuo \emph{sobreyectivo}, llamado \emph{canónico}, del haz de anillos topológicos $\mathscr{O}_{\mathfrak{X}}$ sobre el haz de anillos discretos $(\widetilde{A}/\widetilde{\mathfrak{J}})|\mathfrak{X}$, cuyo núcleo es $\mathfrak{J}^\Delta$; este homomorfismo no es otro que $\widetilde{\varphi}$ \sref{I.10.2.1}, donde $\varphi$ es el homomorfismo continuo $A\to A/\mathfrak{J}$; el morfismo $({}^a\varphi,\widetilde{\varphi}):
\operatorname{Spec}(A/\mathfrak{J})\to\mathfrak{X}$ de esquemas formales afines (donde ${}^a\varphi$ es un homeomorfismo sobre el subespacio cerrado $V(\mathfrak{J})$ de $\mathfrak{X}$)\footnote{\emph{[Ed.] Para un ideal abierto cualquiera, la imagen no es necesariamente todo $\mathfrak{X}$. Por ejemplo, si $A=k[t]$ es discreto y $\mathfrak{J}=(t)$, la imagen es el cerrado propio $V(t)$. La afirmación impresa es válida cuando $\mathfrak{J}$ es un ideal de definición.}} también se llama \emph{canónico}.
Se tiene, pues, según lo anterior, un \emph{isomorfismo canónico}
\[
  \mathscr{O}_{\mathfrak{X}}/\mathfrak{J}^\Delta
  \xrightarrow{\sim}
  (\widetilde{A}/\widetilde{\mathfrak{J}})|\mathfrak{X}
  \tag{10.3.2.1}\label{I.10.3.2.1-fr}
\]

Está claro (en virtud de $\Gamma(\mathfrak{X},\mathfrak{J}^\Delta)=\mathfrak{J}$) que la aplicación $\mathfrak{J}\to\mathfrak{J}^\Delta$ es \emph{estrictamente creciente}; según lo anterior, para $\mathfrak{J}\subset\mathfrak{J}'$, el haz ${\mathfrak{J}'}^\Delta/\mathfrak{J}^\Delta$ es canónicamente isomorfo a $\widetilde{\mathfrak{J}}'/\widetilde{\mathfrak{J}}
=(\mathfrak{J}'/\mathfrak{J})\supertilde$.
\end{env}

\begin{env}[10.3.3]
\label{I.10.3.3}
Siendo todavía las hipótesis y notaciones las de \sref{I.10.3.1}, diremos que un haz de ideales $\mathscr{J}$ de $\mathscr{O}_{\mathfrak{X}}$ es un \emph{haz de ideales de definición} de $\mathfrak{X}$ (o un \emph{ideal de definición} de $\mathfrak{X}$) si, para todo $x\in\mathfrak{X}$, existe una vecindad abierta de $x$ de la forma $\mathfrak{D}(f)$, donde $f\in A$, tal que $\mathscr{J}|\mathfrak{D}(f)$ sea de la forma $\mathfrak{H}^\Delta$, donde $\mathfrak{H}$ es un ideal de definición de $A_{\{f\}}$.
\end{env}

\begin{proposition}[10.3.4]
\label{I.10.3.4}
Para todo $f\in A$, todo ideal de definición de $\mathfrak{X}$ induce un ideal de definición de $\mathfrak{D}(f)$.
\end{proposition}

\begin{proof}
Esto resulta de \sref{I.10.3.1.1}.
\end{proof}

\begin{proposition}[10.3.5]
\label{I.10.3.5}
Si $A$ es un anillo admisible, todo ideal de definición de $\mathfrak{X}=\operatorname{Spf}(A)$ es de la forma $\mathfrak{J}^\Delta$, donde $\mathfrak{J}$ es un ideal de definición de $A$ unívocamente determinado.
\end{proposition}

\begin{proof}
En efecto, sea $\mathscr{J}$ un ideal de definición de $\mathfrak{X}$; por hipótesis, y puesto que $\mathfrak{X}$ es cuasicompacto, existen un número finito de elementos $f_i\in A$ tales que los $\mathfrak{D}(f_i)$ recubren $\mathfrak{X}$ y que $\mathscr{J}|\mathfrak{D}(f_i)=\mathfrak{H}_i^\Delta$, donde $\mathfrak{H}_i$ es un ideal de definición de $A_{\{f_i\}}$.
Para todo $i$, existe, pues, un ideal abierto $\mathfrak{K}_i$ de $A$ tal que $(\mathfrak{K}_i)_{\{f_i\}}=\mathfrak{H}_i$ \sref[0]{0.7.6.9}; sea $\mathfrak{K}$ un ideal de definición de $A$ contenido en todos los $\mathfrak{K}_i$.
La imagen canónica de $\mathscr{J}/\mathfrak{K}^\Delta$ en el haz estructural $(A/\mathfrak{K})\supertilde$ de $\operatorname{Spec}(A/\mathfrak{K})$ \sref{I.10.3.2} es, por tanto, tal que su restricción a $\mathfrak{D}(f_i)$ es igual a la de $(\mathfrak{K}_i/\mathfrak{K})\supertilde$;
se concluye que esta imagen canónica es un haz \emph{cuasicoherente} sobre $\operatorname{Spec}(A/\mathfrak{K})$, y por tanto de la forma $(\mathfrak{J}/\mathfrak{K})\supertilde$, donde $\mathfrak{J}$ es un ideal de $A$ que contiene a $\mathfrak{K}$ \sref{I.1.4.1}, de donde $\mathscr{J}=\mathfrak{J}^\Delta$ \sref{I.10.3.2};
además, como para todo $i$ existe un entero $n_i$ tal que $\mathfrak{H}_i^{n_i}\subset\mathfrak{K}_{\{f_i\}}$, designando por $n$ el mayor de los $n_i$, se tendrá $(\mathscr{J}/\mathfrak{K}^\Delta)^n=0$, y, por consiguiente \sref{I.10.3.2}, $((\mathfrak{J}/\mathfrak{K})\supertilde)^n=0$, de donde finalmente $(\mathfrak{J}/\mathfrak{K})^n=0$ \sref{I.1.3.13}, lo que prueba que $\mathfrak{J}$ es un ideal de definición de $A$ \sref[0]{0.7.1.4}.
\end{proof}

\begin{proposition}[10.3.6]
\label{I.10.3.6}
Sean $A$ un anillo ádico y $\mathfrak{J}$ un ideal de definición de $A$ tal que $\mathfrak{J}/\mathfrak{J}^2$ sea un $A/\mathfrak{J}$-módulo de tipo finito. Para todo entero $n>0$, se tiene entonces $(\mathfrak{J}^\Delta)^n=(\mathfrak{J}^n)^\Delta$.
\end{proposition}

\begin{proof}
En efecto, para todo $f\in A$, se tiene (puesto que $\mathfrak{J}^n$ es un ideal abierto)
\[
  \bigl(\Gamma(\mathfrak{D}(f),\mathfrak{J}^\Delta)\bigr)^n
  =(\mathfrak{J}_{\{f\}})^n
  =(\mathfrak{J}^n)_{\{f\}}
  =\Gamma\bigl(\mathfrak{D}(f^n),(\mathfrak{J}^n)^\Delta\bigr)
\]
\oldpage[I]{185}
en virtud de \sref{I.10.3.1.1} y \sref[0]{0.7.6.12}.
Como $(\mathfrak{J}^\Delta)^n$ está asociado al prehaz $U\to(\Gamma(U,\mathfrak{J}^\Delta))^n$ \sref[0]{0.4.1.6}, la conclusión se sigue\footnote{\emph{[Trad.] La fuente francesa impresa dice aquí «le corollaire», aunque 10.3.6 está enunciado como proposición.}} de que los $\mathfrak{D}(f)$ forman una base de la topología de $\mathfrak{X}$.
\end{proof}

\begin{env}[10.3.7]
\label{I.10.3.7}
Se dice que una familia $(\mathscr{J}_\lambda)$ de ideales de definición de $\mathfrak{X}$ es un \emph{sistema fundamental de ideales de definición} si todo ideal de definición de $\mathfrak{X}$ contiene uno de los $\mathscr{J}_\lambda$; como $\mathscr{J}_\lambda=\mathfrak{J}_\lambda^\Delta$, equivale a decir que los $\mathfrak{J}_\lambda$ forman un \emph{sistema fundamental de vecindades de} 0 en $A$.
Sea $(f_\alpha)$ una familia de elementos de $A$ tal que los $\mathfrak{D}(f_\alpha)$ recubren $\mathfrak{X}$.
Si $(\mathscr{J}_\lambda)$ es una familia filtrante decreciente de haces de ideales de $\mathscr{O}_{\mathfrak{X}}$ tal que, para todo $\alpha$, la familia $(\mathscr{J}_\lambda|\mathfrak{D}(f_\alpha))$ es un sistema fundamental de ideales de definición de $\mathfrak{D}(f_\alpha)$, entonces $(\mathscr{J}_\lambda)$ es un sistema fundamental de ideales de definición de $\mathfrak{X}$.
En efecto, para todo ideal de definición $\mathscr{J}$ de $\mathfrak{X}$, existe un recubrimiento finito de $\mathfrak{X}$ por $\mathfrak{D}(f_i)$ tal que, para todo $i$, $\mathscr{J}_{\lambda_i}|\mathfrak{D}(f_i)$ es un ideal de definición de $\mathfrak{D}(f_i)$ contenido en $\mathscr{J}|\mathfrak{D}(f_i)$.
Si $\mu$ es un índice tal que $\mathscr{J}_\mu\subset\mathscr{J}_{\lambda_i}$ para todo $i$, se sigue de \sref{I.10.3.3} que $\mathscr{J}_\mu$ es un ideal de definición de $\mathfrak{X}$, evidentemente contenido en $\mathscr{J}$, de donde la afirmación.
\end{env}

\subsection{Preesquemas formales y morfismos de preesquemas formales}
\label{subsection:I.10.4}

\begin{env}[10.4.1]
\label{I.10.4.1}
Dado un espacio topológicamente anillado $\mathfrak{X}$, se dice que un abierto $U\subset\mathfrak{X}$ es un \emph{abierto formal afín} (resp.~un \emph{abierto formal afín ádico}, resp.~un \emph{abierto formal afín noetheriano}) si el espacio topológicamente anillado inducido por $\mathfrak{X}$ sobre $U$ es un esquema formal afín (resp.~un esquema de ese tipo cuyo anillo es ádico, resp.~ádico y noetheriano).
\end{env}

\begin{definition}[10.4.2]
\label{I.10.4.2}
{\itshape Se llama preesquema formal a un espacio topológicamente anillado $\mathfrak{X}$ en el que todo punto admite una vecindad abierta formal afín.
Se dice que el preesquema formal $\mathfrak{X}$ es ádico \textup{(resp.} localmente noetheriano) si todo punto de $\mathfrak{X}$ admite una vecindad abierta formal afín ádica \textup{(resp.} noetheriana).
Se dice que $\mathfrak{X}$ es noetheriano si es localmente noetheriano y si su espacio subyacente es cuasicompacto (y, por tanto, noetheriano).}
\end{definition}

\begin{proposition}[10.4.3]
\label{I.10.4.3}
Si $\mathfrak{X}$ es un preesquema formal \textup{(resp.} un preesquema formal localmente noetheriano), los abiertos formales afines \textup{(resp.} afines noetherianos) forman una base de la topología de $\mathfrak{X}$.
\end{proposition}

\begin{proof}
Esto se deduce de \sref{I.10.4.2} y \sref{I.10.1.4}, teniendo en cuenta que, si $A$ es un anillo ádico noetheriano, lo mismo ocurre con $A_{\{f\}}$ para todo $f\in A$ \sref[0]{0.7.6.11}.
\end{proof}

\begin{corollary}[10.4.4]
\label{I.10.4.4}
Si $\mathfrak{X}$ es un preesquema formal \textup{(resp.} un preesquema formal localmente noetheriano, \textup{resp.} noetheriano), el espacio topológicamente anillado inducido sobre todo abierto de $\mathfrak{X}$ sigue siendo un preesquema formal \textup{(resp.} un preesquema formal localmente noetheriano, \textup{resp.} noetheriano).
\end{corollary}

\begin{definition}[10.4.5]
\label{I.10.4.5}
Dados dos preesquemas formales $\mathfrak{X}$ y $\mathfrak{Y}$, se llama morfismo (de preesquemas formales) de $\mathfrak{X}$ en $\mathfrak{Y}$ a todo morfismo $(\psi,\theta)$ de espacios topológicamente anillados tal que, para todo $x\in\mathfrak{X}$, $\theta_x^\sharp$ sea un homomorfismo local $\mathscr{O}_{\psi(x)}\to\mathscr{O}_x$.
\end{definition}

Es inmediato que el compuesto de dos morfismos de preesquemas formales sigue siendo un morfismo de este tipo; por tanto, los preesquemas formales forman una \emph{categoría}, y se denota por $\operatorname{Hom}(\mathfrak{X},\mathfrak{Y})$ el conjunto de morfismos de un preesquema formal $\mathfrak{X}$ en un preesquema formal $\mathfrak{Y}$.

\oldpage[I]{186}
Si $U$ es una parte abierta de $\mathfrak{X}$, la inyección canónica en $\mathfrak{X}$ del preesquema formal inducido por $\mathfrak{X}$ sobre $U$ es un morfismo de preesquemas formales (e incluso un \emph{monomorfismo} de espacios topológicamente anillados \sref[0]{0.4.1.1}).

\begin{proposition}[10.4.6]
\label{I.10.4.6}
Sean $\mathfrak{X}$ un preesquema formal y $\mathfrak{S}=\operatorname{Spf}(A)$ un esquema formal afín.
Existe una correspondencia biunívoca canónica entre los morfismos del preesquema formal $\mathfrak{X}$ en el preesquema formal $\mathfrak{S}$ y los homomorfismos continuos del anillo $A$ en el anillo topológico $\Gamma(\mathfrak{X},\mathscr{O}_{\mathfrak{X}})$.
\end{proposition}

\begin{proof}
La demostración es la misma que la de \sref{I.2.2.4}, sustituyendo ``homomorfismo'' por ``homomorfismo continuo'', ``abierto afín'' por ``abierto formal afín'', y utilizando \sref{I.10.2.2} en lugar de \sref{I.1.7.3}; dejamos los detalles al lector.
\end{proof}

\begin{env}[10.4.7]
\label{I.10.4.7}
Dados un preesquema formal $\mathfrak{S}$, un preesquema formal $\mathfrak{X}$ y un morfismo $\varphi:\mathfrak{X}\to\mathfrak{S}$, se dice que estos datos definen un preesquema formal \emph{$\mathfrak{X}$ sobre $\mathfrak{S}$}, o un \emph{$\mathfrak{S}$-preesquema formal}; $\varphi$ se llama el \emph{morfismo estructural} del $\mathfrak{S}$-preesquema $\mathfrak{X}$.
Si $\mathfrak{S}=\operatorname{Spf}(A)$, donde $A$ es un anillo admisible, se dice también que el $\mathfrak{S}$-preesquema formal $\mathfrak{X}$ es un \emph{$A$-preesquema formal}, o un preesquema formal \emph{sobre $A$}.
Todo preesquema formal puede considerarse como un preesquema formal sobre $\mathbf{Z}$ (dotado de la topología discreta).

Si $\mathfrak{X}$ y $\mathfrak{Y}$ son dos $\mathfrak{S}$-preesquemas formales, se dice que un morfismo $u:\mathfrak{X}\to\mathfrak{Y}$ es un \emph{$\mathfrak{S}$-morfismo} si el diagrama
\[
  \xymatrix{
    \mathfrak{X}\ar[rr]^u\ar[rd] & &
    \mathfrak{Y}\ar[ld]\\
    & \mathfrak{S}
  }
\]
(donde las flechas oblicuas son los morfismos estructurales) es conmutativo.
Con esta definición, los $\mathfrak{S}$-preesquemas formales (para $\mathfrak{S}$ fijo) forman una \emph{categoría}.
Se denota por $\operatorname{Hom}_{\mathfrak{S}}(\mathfrak{X},\mathfrak{Y})$ el conjunto de los $\mathfrak{S}$-morfismos del $\mathfrak{S}$-preesquema formal $\mathfrak{X}$ en el $\mathfrak{S}$-preesquema formal $\mathfrak{Y}$.
Cuando $\mathfrak{S}=\operatorname{Spf}(A)$, se dice también \emph{$A$-morfismo} en lugar de \emph{$\mathfrak{S}$-morfismo}.
\end{env}

\begin{env}[10.4.8]
\label{I.10.4.8}
Como todo esquema afín puede considerarse como un esquema formal afín \sref{I.10.1.2}, todo preesquema (usual) puede considerarse como un preesquema formal.
Además, de \sref{I.10.4.5} se deduce que, para los preesquemas \emph{usuales}, los morfismos (resp. $S$-morfismos) de preesquemas \emph{formales} coinciden con los morfismos (resp. $S$-morfismos) definidos en \textsection\hyperref[section:I.2]{2}.
\end{env}

\subsection{Ideales de definición de los preesquemas formales}
\label{subsection:I.10.5}

\begin{env}[10.5.1]
\label{I.10.5.1}
Sea $\mathfrak{X}$ un preesquema formal; se dice que un $\mathscr{O}_{\mathfrak{X}}$-ideal $\mathscr{J}$ es un \emph{haz de ideales de definición} (o un \emph{ideal de definición}) de $\mathfrak{X}$ si todo $x\in\mathfrak{X}$ posee una vecindad abierta formal afín $U$ tal que $\mathscr{J}|U$ sea un ideal de definición del esquema formal afín inducido por $\mathfrak{X}$ sobre $U$ \sref{I.10.3.3}; en virtud de \sref{I.10.3.1.1} y \sref{I.10.4.3}, para todo abierto $V\subset\mathfrak{X}$, $\mathscr{J}|V$ es entonces un ideal de definición del preesquema formal inducido por $\mathfrak{X}$ sobre $V$.

Se dice que una familia $(\mathscr{J}_\lambda)$ de ideales de definición de $\mathfrak{X}$ es un \emph{sistema fundamental}
\oldpage[I]{187}
\emph{de ideales de definición} si existe un recubrimiento $(U_\alpha)$ de $\mathfrak{X}$ por abiertos formales afines tal que, para todo $\alpha$, la familia de los $\mathscr{J}_\lambda|U_\alpha$ sea un sistema fundamental de ideales de definición \sref{I.10.3.7}\footnote{\emph{[Trad.] La fuente francesa impresa remite aquí a 10.3.6; la definición citada es 10.3.7.}} del esquema formal afín inducido por $\mathfrak{X}$ sobre $U_\alpha$.
De la observación final de \sref{I.10.3.7} se deduce que, cuando $\mathfrak{X}$ es un esquema formal afín, esta definición coincide con la dada en \sref{I.10.3.7}.
Para todo abierto $V$ de $\mathfrak{X}$, las restricciones $\mathscr{J}_\lambda|V$ forman entonces un sistema fundamental de ideales de definición del preesquema formal inducido sobre $V$, en virtud de \sref{I.10.3.1.1}.
Si $\mathfrak{X}$ es un preesquema formal \emph{localmente noetheriano} y $\mathscr{J}$ es un ideal de definición de $\mathfrak{X}$, de \sref{I.10.3.6} se deduce que las potencias $\mathscr{J}^n$ forman un sistema fundamental de ideales de definición de $\mathfrak{X}$.
\end{env}

\begin{env}[10.5.2]
\label{I.10.5.2}
Sean $\mathfrak{X}$ un preesquema formal y $\mathscr{J}$ un ideal de definición de $\mathfrak{X}$.
Entonces el espacio anillado $(\mathfrak{X},\mathscr{O}_{\mathfrak{X}}/\mathscr{J})$ es un \emph{preesquema} (usual), que es afín (resp. localmente noetheriano, resp. noetheriano) cuando $\mathfrak{X}$ es un esquema formal afín (resp. un preesquema formal localmente noetheriano, resp. noetheriano); en efecto, se reduce inmediatamente al caso afín, y entonces la proposición ya fue demostrada en \sref{I.10.3.2}.
Además, si $\theta:\mathscr{O}_{\mathfrak{X}}\to
\mathscr{O}_{\mathfrak{X}}/\mathscr{J}$ es el homomorfismo canónico, entonces $u=(1_{\mathfrak{X}},\theta)$ es un \emph{morfismo} (llamado \emph{canónico}) de preesquemas formales $(\mathfrak{X},\mathscr{O}_{\mathfrak{X}}/\mathscr{J})\to
(\mathfrak{X},\mathscr{O}_{\mathfrak{X}})$, pues, también aquí, esto se vio en el caso afín \sref{I.10.3.2}, al cual se reduce inmediatamente.
\end{env}

\begin{proposition}[10.5.3]
\label{I.10.5.3}
Sean $\mathfrak{X}$ un preesquema formal y $(\mathscr{J}_\lambda)$ un sistema fundamental de ideales de definición de $\mathfrak{X}$.
Entonces el haz de anillos topológicos $\mathscr{O}_{\mathfrak{X}}$ es el límite proyectivo de los haces pseudodiscretos de anillos \sref[0]{0.3.8.1} $\mathscr{O}_{\mathfrak{X}}/\mathscr{J}_\lambda$.
\end{proposition}

\begin{proof}
Como la topología de $\mathfrak{X}$ admite una base de abiertos formales afines cuasicompactos \sref{I.10.4.3}, se reduce al caso afín, donde la proposición es consecuencia de la Proposición \sref{I.10.3.5}, de \sref{I.10.3.2} y de la Definición \sref{I.10.1.1}.
\end{proof}

No es seguro que todo preesquema formal admita ideales de definición.
Sin embargo:
\begin{proposition}[10.5.4]
\label{I.10.5.4}
Sea $\mathfrak{X}$ un preesquema formal localmente noetheriano.
Existe un mayor ideal de definición $\mathscr{T}$ de $\mathfrak{X}$; es el único ideal de definición $\mathscr{J}$ tal que el preesquema $(\mathfrak{X},\mathscr{O}_{\mathfrak{X}}/\mathscr{J})$ sea reducido.
Si $\mathscr{J}$ es un ideal de definición de $\mathfrak{X}$, entonces $\mathscr{T}$ es la imagen inversa, por $\mathscr{O}_{\mathfrak{X}}\to
\mathscr{O}_{\mathfrak{X}}/\mathscr{J}$, del nilradical de $\mathscr{O}_{\mathfrak{X}}/\mathscr{J}$.
\end{proposition}

\begin{proof}
Supongamos primero que $\mathfrak{X}=\operatorname{Spf}(A)$, donde $A$ es un anillo ádico noetheriano.
La existencia y las propiedades de $\mathscr{T}$ se deducen inmediatamente de \sref{I.10.3.4} y \sref{I.5.1.1}, teniendo en cuenta la existencia y las propiedades del mayor ideal de definición de $A$ (\sref[0]{0.7.1.6} y \sref[0]{0.7.1.7}).

Para probar la existencia y las propiedades de $\mathscr{T}$ en el caso general, basta mostrar que, si $U\supset V$ son dos abiertos formales afines noetherianos de $\mathfrak{X}$, el mayor ideal de definición $\mathscr{T}_U$ de $U$ induce el mayor ideal de definición $\mathscr{T}_V$ de $V$; pero, como $(V,(\mathscr{O}_{\mathfrak{X}}|V)/(\mathscr{T}_U|V))$ es reducido, esto se deduce de lo anterior.
\end{proof}

Se denota por $\mathfrak{X}_{\mathrm{red}}$ el preesquema (usual) reducido $(\mathfrak{X},\mathscr{O}_{\mathfrak{X}}/\mathscr{T})$).

\begin{corollary}[10.5.5]
\label{I.10.5.5}
Sean $\mathfrak{X}$ un preesquema formal localmente noetheriano y $\mathscr{T}$ el mayor ideal de definición de $\mathfrak{X}$; para todo abierto $V$ de $\mathfrak{X}$, $\mathscr{T}|V$ es el mayor ideal de definición del preesquema formal inducido por $\mathfrak{X}$ sobre $V$.
\end{corollary}

\oldpage[I]{188}

\begin{proposition}[10.5.6]
\label{I.10.5.6}
Sean $\mathfrak{X}$ y $\mathfrak{Y}$ dos preesquemas formales, $\mathscr{J}$ \textup{(resp. $\mathscr{K}$)} un ideal de definición de $\mathfrak{X}$ \textup{(resp. $\mathfrak{Y}$)}, y $f:\mathfrak{X}\to\mathfrak{Y}$ un morfismo de preesquemas formales.
\begin{enumerate}
  \item[\rm{(i)}] Si $f^*(\mathscr{K})\mathscr{O}_{\mathfrak{X}}\subset\mathscr{J}$, existe un único morfismo $f':(\mathfrak{X},\mathscr{O}_{\mathfrak{X}}/\mathscr{J})\to
    (\mathfrak{Y},\mathscr{O}_{\mathfrak{Y}}/\mathscr{K})$ de preesquemas usuales que hace conmutativo el diagrama
    \[
      \label{I.10.5.6.1-fr}
      \xymatrix{
        (\mathfrak{X},\mathscr{O}_{\mathfrak{X}})\ar[r]^f &
        (\mathfrak{Y},\mathscr{O}_{\mathfrak{Y}})\\
        (\mathfrak{X},\mathscr{O}_{\mathfrak{X}}/\mathscr{J})
          \ar[r]_{f'}\ar[u] &
        (\mathfrak{Y},\mathscr{O}_{\mathfrak{Y}}/\mathscr{K})\ar[u]
      }
      \tag{10.5.6.1}
    \]
    donde las flechas verticales son los morfismos canónicos.
  \item[\rm{(ii)}] Supongamos que $\mathfrak{X}=\operatorname{Spf}(A)$ y $\mathfrak{Y}=\operatorname{Spf}(B)$ son esquemas formales afines, $\mathscr{J}=\mathfrak{J}^{\Delta}$ y $\mathscr{K}=\mathfrak{K}^{\Delta}$, donde $\mathfrak{J}$ \textup{(resp. $\mathfrak{K}$)} es un ideal de definición de $A$ \textup{(resp. $B$)}, y $f=({}^a\varphi,\widetilde{\varphi})$, donde $\varphi:B\to A$ es un homomorfismo continuo;
    para que se tenga $f^*(\mathscr{K})\mathscr{O}_{\mathfrak{X}}\subset\mathscr{J}$, es necesario y suficiente que $\varphi(\mathfrak{K})\subset\mathfrak{J}$, y entonces $f'$ es el morfismo $({}^a\varphi',\widetilde{\varphi}')$, donde $\varphi':B/\mathfrak{K}\to A/\mathfrak{J}$ es el homomorfismo deducido de $\varphi$ por paso a los cocientes.
\end{enumerate}
\end{proposition}

\begin{proof}
\medskip\noindent
\begin{enumerate}
  \item[\rm{(i)}] Si $f=(\psi,\theta)$, la hipótesis implica que la imagen por $\theta^\sharp:\psi^*(\mathscr{O}_{\mathfrak{Y}})\to
    \mathscr{O}_{\mathfrak{X}}$ del haz de ideales $\psi^*(\mathscr{K})$ de $\psi^*(\mathscr{O}_{\mathfrak{Y}})$ está contenida en $\mathscr{J}$ \sref[0]{0.4.3.5}.
    Pasando a los cocientes, se deduce por tanto de $\theta^\sharp$ un homomorfismo de haces de anillos
    \[
      \omega:\psi^*(\mathscr{O}_{\mathfrak{Y}}/\mathscr{K})
      =\psi^*(\mathscr{O}_{\mathfrak{Y}})/\psi^*(\mathscr{K})
      \to\mathscr{O}_{\mathfrak{X}}/\mathscr{J}
    \]
    además, como para todo $x\in\mathfrak{X}$, $\theta_x^\sharp$ es un homomorfismo \emph{local}, lo mismo ocurre con $\omega_x$.
    El morfismo de espacios anillados $(\psi,\omega^\flat)$ es pues \sref{I.2.2.1} el único morfismo $f'$ de espacios anillados que responde a la cuestión.
  \item[\rm{(ii)}] La correspondencia canónica y funtorial entre los morfismos de preesquemas formales afines y los homomorfismos continuos de anillos \sref{I.10.2.2} muestra que, en el caso considerado, la relación $f^*(\mathscr{K})\mathscr{O}_{\mathfrak{X}}\subset\mathscr{J}$ implica que $f'=({}^a\varphi',\widetilde{\varphi}')$, donde $\varphi':B/\mathfrak{K}\to A/\mathfrak{J}$ es el único homomorfismo que hace conmutativo el diagrama
    \[
      \label{I.10.5.6.2-fr}
      \xymatrix{
        B\ar[r]^\varphi\ar[d] & A\ar[d]\\
        B/\mathfrak{K}\ar[r]_{\varphi'} & A/\mathfrak{J}
      }
      \tag{10.5.6.2}
    \]
    La existencia de $\varphi'$ implica por tanto que $\varphi(\mathfrak{K})\subset\mathfrak{J}$.
    Recíprocamente, si se cumple esta condición, denotando por $\varphi'$ el único homomorfismo que hace conmutativo el diagrama \sref{I.10.5.6.2} y poniendo $f'=({}^a\varphi',\widetilde{\varphi}')$, es claro que el diagrama \sref{I.10.5.6.1} es conmutativo; la consideración de los homomorfismos ${}^a\varphi^*(\mathscr{O}_{\mathfrak{Y}})\to
    \mathscr{O}_{\mathfrak{X}}$ y ${}^a{\varphi'}^*(\mathscr{O}_{\mathfrak{Y}}/\mathscr{K})\to
    \mathscr{O}_{\mathfrak{X}}/\mathscr{J}$ correspondientes respectivamente a $f$ y $f'$ muestra entonces que esto implica la relación $f^*(\mathscr{K})\mathscr{O}_{\mathfrak{X}}\subset\mathscr{J}$.
\end{enumerate}
\end{proof}

Es claro que la correspondencia $f\to f'$ definida anteriormente es \emph{funtorial}.

\subsection{Preesquemas formales como límites inductivos de preesquemas}
\label{subsection:I.10.6}

\begin{env}[10.6.1]
\label{I.10.6.1}
Sean $\mathfrak{X}$ un preesquema formal y $(\mathscr{J}_\lambda)$ un sistema fundamental de ideales de definición de $\mathfrak{X}$; para todo $\lambda$, sea $f_\lambda$ el morfismo canónico $(\mathfrak{X},\mathscr{O}_{\mathfrak{X}}/\mathscr{J}_\lambda)
\to\mathfrak{X}$ \sref{I.10.5.2}; para $\mathscr{J}_\mu\subset\mathscr{J}_\lambda$, el homomorfismo canónico $\mathscr{O}_{\mathfrak{X}}/\mathscr{J}_\mu\to
\mathscr{O}_{\mathfrak{X}}/\mathscr{J}_\lambda$ define un morfismo canónico
\oldpage[I]{189}
$f_{\mu\lambda}:(\mathfrak{X},
\mathscr{O}_{\mathfrak{X}}/\mathscr{J}_\lambda)\to
(\mathfrak{X},\mathscr{O}_{\mathfrak{X}}/\mathscr{J}_\mu)$ de preesquemas (usuales) tal que $f_\lambda=f_\mu\circ f_{\mu\lambda}$.
Los preesquemas $X_\lambda=(\mathfrak{X},
\mathscr{O}_{\mathfrak{X}}/\mathscr{J}_\lambda)$ y los morfismos $f_{\mu\lambda}$ constituyen por tanto (en virtud de \sref{I.10.4.8}) un \emph{sistema inductivo} en la categoría de los preesquemas formales.
\end{env}

\begin{proposition}[10.6.2]
\label{I.10.6.2}
Con las notaciones de \sref{I.10.6.1}, el preesquema formal $\mathfrak{X}$ y los morfismos $f_\lambda$ constituyen un límite inductivo \textup{(T,~I,~1.8)} del sistema $(X_\lambda,f_{\mu\lambda})$ en la categoría de los preesquemas formales.
\end{proposition}

\begin{proof}
Sea $\mathfrak{Y}$ un preesquema formal y, para cada índice $\lambda$, sea
\[
  g_\lambda=(\psi_\lambda,\theta_\lambda):X_\lambda\to\mathfrak{Y}
\]
un morfismo tal que $g_\lambda=g_\mu\circ f_{\mu\lambda}$ para $\mathscr{J}_\mu\subset\mathscr{J}_\lambda$.
Esta última condición y la definición de los $X_\lambda$ implican primero que todos los $\psi_\lambda$ son idénticos a una misma aplicación continua $\psi:\mathfrak{X}\to\mathfrak{Y}$ de los espacios subyacentes; además, los homomorfismos $\theta_\lambda^\sharp:\psi^*(\mathscr{O}_{\mathfrak{Y}})\to
\mathscr{O}_{X_\lambda}=
\mathscr{O}_{\mathfrak{X}}/\mathscr{J}_\lambda$ forman un \emph{sistema proyectivo} de homomorfismos de haces de anillos.
Pasando al límite proyectivo, se deduce por tanto un homomorfismo $\omega:\psi^*(\mathscr{O}_{\mathfrak{Y}})\to
  \varprojlim\mathscr{O}_{\mathfrak{X}}/\mathscr{J}_\lambda
  =\mathscr{O}_{\mathfrak{X}}$, y es claro que el morfismo $g=(\psi,\omega^\flat)$ de \emph{espacios anillados} es el \emph{único} morfismo que hace conmutativos los diagramas
\[
  \label{I.10.6.2.1-fr}
  \xymatrix{
    X_\lambda\ar[rr]^{g_\lambda}\ar[rd]_{f_\lambda} &&
    \mathfrak{Y}\\
    & \mathfrak{X}\ar[ru]_g
  }
  \tag{10.6.2.1}
\]
Queda por probar que $g$ es un morfismo de \emph{preesquemas formales}; como la cuestión es local sobre $\mathfrak{X}$ y $\mathfrak{Y}$, se puede suponer que $\mathfrak{X}=\operatorname{Spf}(A)$ y $\mathfrak{Y}=\operatorname{Spf}(B)$, siendo $A$ y $B$ anillos admisibles, y que $\mathscr{J}_\lambda=\mathfrak{J}_\lambda^{\Delta}$, donde $(\mathfrak{J}_\lambda)$ es un sistema fundamental de ideales de definición de $A$ \sref{I.10.3.5}; como $A=\varprojlim A/\mathfrak{J}_\lambda$, la existencia de un morfismo $g$ de esquemas formales afines que hace conmutativos los diagramas \sref{I.10.6.2.1} se deduce entonces de la correspondencia biyectiva \sref{I.10.2.2} entre morfismos de esquemas formales afines y homomorfismos continuos de anillos, y de la definición del límite proyectivo.
Pero la unicidad de $g$ como morfismo de espacios anillados muestra que coincide con el morfismo denotado del mismo modo al comienzo de la demostración.
\end{proof}

La proposición siguiente establece, bajo ciertas condiciones suplementarias, la existencia del límite inductivo de un sistema inductivo dado de preesquemas (usuales) en la categoría de los preesquemas formales:
\begin{proposition}[10.6.3]
\label{I.10.6.3}
Sean $\mathfrak{X}$ un espacio topológico y $(\mathscr{O}_i,u_{ji})$ un sistema proyectivo de haces de anillos sobre $\mathfrak{X}$, cuyo conjunto de índices es $\mathbf{N}$.
Sea $\mathscr{J}_i$ el núcleo de $u_{0i}:\mathscr{O}_i\to\mathscr{O}_0$.
Se supone que:
\begin{enumerate}
  \item[\textup{a)}] el espacio anillado $(\mathfrak{X},\mathscr{O}_i)$ es un preesquema $X_i$;
  \item[\textup{b)}] para todo $x\in\mathfrak{X}$\footnote{\emph{[Trad.] La fuente francesa impresa dice $x\in X$; la corrección publicada en EGA~II sustituye $X$ por $\mathfrak{X}$.}} y todo $i$, existe una vecindad abierta $U_i$ de $x$ en $\mathfrak{X}$ tal que la restricción $\mathscr{J}_i|U_i$ es nilpotente; y
  \item[\textup{c)}] los homomorfismos $u_{ji}$ son sobreyectivos.
\end{enumerate}

\oldpage[I]{190}
Sea $\mathscr{O}_{\mathfrak{X}}$ el haz de anillos topológicos dado por el límite proyectivo de los haces de anillos pseudodiscretos $\mathscr{O}_i$, y sea $u_i:\mathscr{O}_{\mathfrak{X}}\to\mathscr{O}_i$ el homomorfismo canónico.
Entonces el espacio topológicamente anillado $(\mathfrak{X},\mathscr{O}_{\mathfrak{X}})$ es un preesquema formal; los homomorfismos $u_i$ son sobreyectivos; sus núcleos $\mathscr{J}^{(i)}$ forman un sistema fundamental de haces de ideales de definición de $\mathfrak{X}$, y $\mathscr{J}^{(0)}$ es el límite proyectivo de los haces de ideales $\mathscr{J}_i$.
\end{proposition}

\begin{proof}
Observemos primero que, sobre cada tallo, $u_{ji}$ es un homomorfismo sobreyectivo y \emph{a fortiori} un homomorfismo local; por tanto, $v_{ij}=(1_{\mathfrak{X}},u_{ji})$ es un morfismo de preesquemas $X_j\to X_i$ ($i\geq j$) \sref{I.2.2.1}.
Supongamos primero que cada $X_i$ es un esquema afín de anillo $A_i$.
Existe un homomorfismo de \emph{anillos} $\varphi_{ji}:A_i\to A_j$ tal que $u_{ji}=\widetilde{\varphi}_{ji}$ \sref{I.1.7.3}; por consiguiente \sref{I.1.6.3}, el haz $\mathscr{O}_j$ es un $\mathscr{O}_i$-módulo cuasicoherente sobre $X_i$ (para la ley externa definida por $u_{ji}$), asociado a $A_j$ considerado como $A_i$-módulo mediante $\varphi_{ji}$.
Para todo $f\in A_i$, sea $f'=\varphi_{ji}(f)$; por hipótesis, los abiertos $D(f)$ y $D(f')$ son idénticos en $\mathfrak{X}$, y el homomorfismo de $\Gamma(D(f),\mathscr{O}_i)=(A_i)_f$ en $\Gamma(D(f),\mathscr{O}_j)=(A_j)_{f'}$ correspondiente a $u_{ji}$ no es otro que $(\varphi_{ji})_f$ \sref{I.1.6.1}.
Pero, cuando se considera $A_j$ como $A_i$-módulo, $(A_j)_{f'}$ es el $(A_i)_f$-módulo $(A_j)_f$, por lo que también se tiene $u_{ji}=\widetilde{\varphi}_{ji}$, considerando esta vez $\varphi_{ji}$ como homomorfismo de $A_i$-\emph{módulos}.
Entonces, como $u_{ji}$ es sobreyectivo, se concluye que $\varphi_{ji}$ también lo es \sref{I.1.3.9}; y si $\mathfrak{J}_{ji}$ es el núcleo de $\varphi_{ji}$, el núcleo de $u_{ji}$ es un $\mathscr{O}_i$-módulo cuasicoherente igual a $\widetilde{\mathfrak{J}}_{ji}$.
En particular, se tiene $\mathscr{J}_i=\widetilde{\mathfrak{J}}_i$, donde $\mathfrak{J}_i$ es el núcleo de $\varphi_{0i}:A_i\to A_0$.
La hipótesis~\emph{b)} implica que $\mathscr{J}_i$ es \emph{nilpotente}: en efecto, como $\mathfrak{X}$ es cuasicompacto, se puede recubrir $\mathfrak{X}$ por un número finito de abiertos $U_k$ tales que $(\mathscr{J}_i|U_k)^{n_k}=0$ y, tomando como $n$ el mayor de los $n_k$, se tiene $\mathscr{J}_i^n=0$.
Se concluye así que $\mathfrak{J}_i$ es nilpotente \sref{I.1.3.13}.
Entonces el anillo $A=\varprojlim A_i$ es admisible \sref[0]{0.7.2.2}, el homomorfismo canónico $\varphi_i:A\to A_i$ es sobreyectivo y su núcleo $\mathfrak{J}^{(i)}$ es igual al límite proyectivo de los $\mathfrak{J}_{ik}$ para $k\geq i$; los $\mathfrak{J}^{(i)}$ forman un sistema fundamental de vecindades de 0 en $A$.
Las afirmaciones de la Proposición~\sref{I.10.6.3} se deducen en este caso de \sref{I.10.1.1} y \sref{I.10.3.2}, pues $(\mathfrak{X},\mathscr{O}_{\mathfrak{X}})$ no es otro que $\operatorname{Spf}(A)$.

Siempre en este mismo caso particular, observemos que, si $f=(f_i)$ es un elemento del límite proyectivo $A=\varprojlim A_i$, todos los abiertos $D(f_i)$ (abiertos afines en $X_i$) se identifican con el abierto $\mathfrak{D}(f)$ de $\mathfrak{X}$, y el preesquema inducido por $X_i$ sobre $\mathfrak{D}(f)$ se identifica, por tanto, con el esquema afín $\operatorname{Spec}((A_i)_{f_i})$.

En el caso general, observemos primero que, para todo abierto cuasicompacto $U$ de $\mathfrak{X}$, cada uno de los $\mathscr{J}_i|U$ es nilpotente, como muestra el razonamiento anterior.
Vamos a ver que, para todo $x\in\mathfrak{X}$, existe una vecindad abierta $U$ de $x$ en $\mathfrak{X}$ que es un \emph{abierto afín} para \emph{todos} los $X_i$.
En efecto, tomemos $U$ abierto afín para $X_0$ y observemos que $\mathscr{O}_{X_0}=\mathscr{O}_{X_i}/\mathscr{J}_i$.
Como, por lo anterior, $\mathscr{J}_i|U$ es nilpotente, $U$ es también abierto afín para cada $X_i$, en virtud de la Proposición~\sref{I.5.1.9}.
En estas condiciones, para todo $U$ que satisfaga las condiciones precedentes, el estudio anterior del caso afín muestra que $(U,\mathscr{O}_{\mathfrak{X}}|U)$ es un preesquema formal cuyos $\mathscr{J}^{(i)}|U$ forman un sistema fundamental de haces de ideales de definición, y que $\mathscr{J}^{(0)}|U$ es el límite proyectivo de los $\mathscr{J}_i|U$; de donde se obtiene la conclusión.
\end{proof}

\begin{corollary}[10.6.4]
\label{I.10.6.4}
Supongamos que, para $i\geq j$, el núcleo de $u_{ji}$ es $\mathscr{J}_i^{j+1}$ y que $\mathscr{J}_1/\mathscr{J}_1^2$
\oldpage[I]{191}
es de tipo finito sobre $\mathscr{O}_0=\mathscr{O}_1/\mathscr{J}_1$.
Entonces $\mathfrak{X}$ es un preesquema formal ádico y, si $\mathscr{J}^{(n)}$ es el núcleo de $\mathscr{O}_{\mathfrak{X}}\to\mathscr{O}_n$, se tiene $\mathscr{J}^{(n)}=\mathscr{J}^{n+1}$ y $\mathscr{J}/\mathscr{J}^2$ es isomorfo a $\mathscr{J}_1$.
Si además $X_0$ es localmente noetheriano \textup{(resp.} noetheriano), entonces $\mathfrak{X}$ es localmente noetheriano \textup{(resp.} noetheriano).
\end{corollary}

\begin{proof}
Como los espacios subyacentes a $\mathfrak{X}$ y $X_0$ son los mismos, la cuestión es local y se puede suponer que todos los $X_i$ son afines; teniendo en cuenta las relaciones $\mathscr{J}_{ij}=\widetilde{\mathfrak{J}}_{ji}$ (con las notaciones de la Proposición~\sref{I.10.6.3}), el problema se reduce inmediatamente a las afirmaciones correspondientes de la Proposición~\sref[0]{0.7.2.7} y del Corolario~\sref[0]{0.7.2.8}, observando que $\mathfrak{J}_1/\mathfrak{J}_1^2$ es entonces un $A_0$-módulo de tipo finito \sref{I.1.3.9}.
\end{proof}

En particular, \emph{todo preesquema formal localmente noetheriano $\mathfrak{X}$} es límite inductivo de una sucesión $(X_n)$ de preesquemas (usuales) localmente noetherianos que satisfacen las condiciones de la Proposición~\sref{I.10.6.3} y del Corolario~\sref{I.10.6.4}: basta considerar un haz de ideales de definición $\mathscr{J}$ de $\mathfrak{X}$ \sref{I.10.5.4} y tomar $X_n=(\mathfrak{X},
\mathscr{O}_{\mathfrak{X}}/\mathscr{J}^{n+1})$ (\sref{I.10.5.1} y Proposición~\sref{I.10.6.2}).

\begin{corollary}[10.6.5]
\label{I.10.6.5}
Sea $A$ un anillo admisible.
Para que el esquema formal afín $\mathfrak{X}=\operatorname{Spf}(A)$ sea noetheriano, es necesario y suficiente que $A$ sea ádico y noetheriano.
\end{corollary}

\begin{proof}
La condición es evidentemente suficiente.
Recíprocamente, supongamos que $\mathfrak{X}$ es noetheriano, y sean $\mathfrak{J}$ un ideal de definición de $A$ y $\mathscr{J}=\mathfrak{J}^\Delta$ el haz de ideales de definición correspondiente de $\mathfrak{X}$.
Los preesquemas (usuales) $X_n=(\mathfrak{X},
\mathscr{O}_{\mathfrak{X}}/\mathscr{J}^{n+1})$ son entonces afines y noetherianos, luego los anillos $A_n=A/\mathfrak{J}^{n+1}$ son noetherianos \sref{I.6.1.3}; de donde se concluye que $\mathfrak{J}/\mathfrak{J}^2$ es un $A/\mathfrak{J}$-módulo de tipo finito.
Como los $\mathscr{J}^n$ forman un sistema fundamental de haces de ideales de definición de $\mathfrak{X}$ \sref{I.10.5.1}, se tiene $\mathscr{O}_{\mathfrak{X}}
=\varprojlim(\mathscr{O}_{\mathfrak{X}}/\mathscr{J}^n)$ \sref{I.10.5.3}; se concluye así \sref{I.10.1.3} que $A$ es topológicamente isomorfo a $\varprojlim A/\mathfrak{J}^n$ y, por tanto, es ádico y noetheriano \sref[0]{0.7.2.8}.
\end{proof}

\begin{remark}[10.6.6]
\label{I.10.6.6}
Con las notaciones de la Proposición~\sref{I.10.6.3}, sea $\mathscr{F}_i$ un $\mathscr{O}_i$-módulo, y supongamos dado, para $i\geq j$, un $v_{ij}$-morfismo $\theta_{ji}:\mathscr{F}_i\to\mathscr{F}_j$, de modo que $\theta_{kj}\circ\theta_{ji}=\theta_{ki}$ para $k\leq j\leq i$.
Como la aplicación continua subyacente a $v_{ij}$ es la identidad, $\theta_{ji}$ es un homomorfismo de haces de grupos abelianos sobre el espacio $\mathfrak{X}$; además, si $\mathscr{F}$ es el límite proyectivo del sistema proyectivo $(\mathscr{F}_i)$ de haces de grupos abelianos, el hecho de que los $\theta_{ji}$ sean $v_{ij}$-morfismos permite definir sobre $\mathscr{F}$ una estructura de $\mathscr{O}_{\mathfrak{X}}$-módulo por paso al límite proyectivo; provisto de esta estructura, diremos que $\mathscr{F}$ es el \emph{límite proyectivo} (para los $\theta_{ji}$) del sistema de $\mathscr{O}_i$-módulos $(\mathscr{F}_i)$.
En el caso particular en que $v_{ij}^*(\mathscr{F}_i)=\mathscr{F}_j$ y $\theta_{ji}$ es la \emph{identidad}, diremos, para abreviar, que $\mathscr{F}$ es el límite proyectivo de un sistema $(\mathscr{F}_i)$ tal que $v_{ij}^*(\mathscr{F}_i)=\mathscr{F}_j$ para $j\leq i$ (sin mencionar los $\theta_{ji}$).
\end{remark}

\begin{env}[10.6.7]
\label{I.10.6.7}
Sean $\mathfrak{X}$ y $\mathfrak{Y}$ dos preesquemas formales, $\mathscr{J}$ (resp. $\mathscr{K}$) un ideal de definición de $\mathfrak{X}$ (resp. $\mathfrak{Y}$), y $f:\mathfrak{X}\to\mathfrak{Y}$ un morfismo tal que $f^*(\mathscr{K})\mathscr{O}_{\mathfrak{X}}\subset\mathscr{J}$.
Se tiene entonces, para todo entero $n>0$, $f^*(\mathscr{K}^n)\mathscr{O}_{\mathfrak{X}}
=(f^*(\mathscr{K})\mathscr{O}_{\mathfrak{X}})^n
\subset\mathscr{J}^n$; por tanto, $f$ induce \sref{I.10.5.6} un morfismo de preesquemas (usuales) $f_n:X_n\to Y_n$, poniendo $X_n=(\mathfrak{X},
\mathscr{O}_{\mathfrak{X}}/\mathscr{J}^{n+1})$ y $Y_n=(\mathfrak{Y},
\mathscr{O}_{\mathfrak{Y}}/\mathscr{K}^{n+1})$, y se sigue inmediatamente de las definiciones que los diagramas
\[
  \label{I.10.6.7.1-fr}
  \xymatrix{
    X_m\ar[r]^{f_m}\ar[d] &
    Y_m\ar[d]\\
    X_n\ar[r]_{f_n} &
    Y_n
  }
  \tag{10.6.7.1}
\]
\oldpage[I]{192}
son conmutativos para $m\leq n$; dicho de otro modo, $(f_n)$ es un \emph{sistema inductivo} de morfismos.
\end{env}

\begin{env}[10.6.8]
\label{I.10.6.8}
Recíprocamente, sea $(X_n)$ (resp. $(Y_n)$) un sistema inductivo de preesquemas (usuales) que satisfaga las condiciones~\emph{b)} y \emph{c)} de la Proposición~\sref{I.10.6.3}, y sea $\mathfrak{X}$ (resp. $\mathfrak{Y}$) su límite inductivo.
Por definición de los límites inductivos, toda sucesión $(f_n)$ de morfismos $X_n\to Y_n$ que forme un sistema inductivo admite un \emph{límite inductivo $f:\mathfrak{X}\to\mathfrak{Y}$}, que es el único morfismo de preesquemas formales para el cual los diagramas
\[
  \xymatrix{
    X_n\ar[r]^{f_n}\ar[d] &
    Y_n\ar[d]\\
    \mathfrak{X}\ar[r]_f &
    \mathfrak{Y}
  }
\]
son conmutativos.
\end{env}

\begin{proposition}[10.6.9]
\label{I.10.6.9}
Sean $\mathfrak{X}$ y $\mathfrak{Y}$ dos preesquemas formales localmente noetherianos, y $\mathscr{J}$ \textup{(resp. $\mathscr{K}$)} un ideal de definición de $\mathfrak{X}$ \textup{(resp. $\mathfrak{Y}$)}; la aplicación $f\to(f_n)$ definida en \sref{I.10.6.7} es una biyección del conjunto de los morfismos $f:\mathfrak{X}\to\mathfrak{Y}$ tales que $f^*(\mathscr{K})\mathscr{O}_{\mathfrak{X}}\subset\mathscr{J}$ sobre el conjunto de las sucesiones $(f_n)$ de morfismos que hacen conmutativos los diagramas \sref{I.10.6.7.1}.
\end{proposition}

\begin{proof}
Si $f$ es el límite inductivo de una sucesión tal, queda por mostrar que $f^*(\mathscr{K})\mathscr{O}_{\mathfrak{X}}\subset\mathscr{J}$.
Como la cuestión es local sobre $\mathfrak{X}$ y $\mathfrak{Y}$, se puede reducir al caso en que $\mathfrak{X}=\operatorname{Spf}(A)$ y $\mathfrak{Y}=\operatorname{Spf}(B)$ son afines, siendo $A$ y $B$ anillos ádicos noetherianos, y $\mathscr{J}=\mathfrak{J}^\Delta$, $\mathscr{K}=\mathfrak{K}^\Delta$, donde $\mathfrak{J}$ (resp. $\mathfrak{K}$) es un ideal de definición de $A$ (resp. $B$).
Se tiene entonces $X_n=\operatorname{Spec}(A_n)$, $Y_n=\operatorname{Spec}(B_n)$, con $A_n=A/\mathfrak{J}^{n+1}$ y $B_n=B/\mathfrak{K}^{n+1}$, en virtud de \sref{I.10.3.6} y \sref{I.10.3.2}; además, $f_n=({}^a\varphi_n,\widetilde{\varphi}_n)$, donde los homomorfismos $\varphi_n:B_n\to A_n$ forman un sistema proyectivo, y por tanto $f=({}^a\varphi,\widetilde{\varphi})$, donde $\varphi=\varprojlim\varphi_n$.
La conmutatividad del diagrama \sref{I.10.6.7.1} para $m=0$ da entonces la condición $\varphi_n(\mathfrak{K}/\mathfrak{K}^{n+1})
\subset\mathfrak{J}/\mathfrak{J}^{n+1}$ para todo $n$; pasando al límite proyectivo, se obtiene $\varphi(\mathfrak{K})\subset
\mathfrak{J}$, lo cual implica $f^*(\mathscr{K})\mathscr{O}_{\mathfrak{X}}\subset\mathscr{J}$ \sref{I.10.5.6}[ii].
\end{proof}

\begin{corollary}[10.6.10]
\label{I.10.6.10}
Sean $\mathfrak{X}$ y $\mathfrak{Y}$ dos preesquemas formales localmente noetherianos, y sea $\mathscr{T}$ el mayor ideal de definición de $\mathfrak{X}$ \sref{I.10.5.4}.
\begin{enumerate}
  \item[\rm{(i)}] Para todo ideal de definición $\mathscr{K}$ de $\mathfrak{Y}$ y todo morfismo $f:\mathfrak{X}\to\mathfrak{Y}$, se tiene $f^*(\mathscr{K})\mathscr{O}_{\mathfrak{X}}\subset\mathscr{T}$.
  \item[\rm{(ii)}] Hay una correspondencia biyectiva canónica entre $\operatorname{Hom}(\mathfrak{X},\mathfrak{Y})$ y el conjunto de las sucesiones $(f_n)$ de morfismos que hacen conmutativos los diagramas \sref{I.10.6.7.1}, donde $X_n=(\mathfrak{X},
  \mathscr{O}_{\mathfrak{X}}/\mathscr{T}^{n+1})$ e $Y_n=(\mathfrak{Y},
  \mathscr{O}_{\mathfrak{Y}}/\mathscr{K}^{n+1})$.
\end{enumerate}
\end{corollary}

\begin{proof}
(ii) se deduce inmediatamente de (i) y de la Proposición~\sref{I.10.6.9}.
Para demostrar (i), se puede reducir al caso en que $\mathfrak{X}=\operatorname{Spf}(A)$ y $\mathfrak{Y}=\operatorname{Spf}(B)$, siendo $A$ y $B$ noetherianos, y $\mathscr{T}=\mathfrak{T}^\Delta$, $\mathscr{K}=\mathfrak{K}^\Delta$, donde $\mathfrak{T}$ es el mayor ideal de definición de $A$ y $\mathfrak{K}$ un ideal de definición de $B$.
Sea $f=({}^a\varphi,\widetilde{\varphi})$, donde $\varphi:B\to A$ es un homomorfismo continuo; como los elementos de $\mathfrak{K}$ son topológicamente nilpotentes \sref[0]{0.7.1.4}[(ii)], también lo son los de $\varphi(\mathfrak{K})$, por lo que $\varphi(\mathfrak{K})\subset\mathfrak{T}$, puesto que $\mathfrak{T}$ es el conjunto de los elementos topológicamente nilpotentes de $A$ \sref[0]{0.7.1.6}; de donde se obtiene la conclusión en virtud de \sref{I.10.5.6}[ii].
\end{proof}

\begin{corollary}[10.6.11]
\label{I.10.6.11}
Sean $\mathfrak{S}$, $\mathfrak{X}$ y $\mathfrak{Y}$ tres preesquemas formales localmente noetherianos, y sean $f:\mathfrak{X}\to\mathfrak{S}$ y $g:\mathfrak{Y}\to\mathfrak{S}$ los morfismos que hacen de $\mathfrak{X}$ y $\mathfrak{Y}$ preesquemas formales sobre $\mathfrak{S}$.
Sea $\mathscr{J}$ \textup{(resp. $\mathscr{K}$, $\mathscr{L}$)} un ideal de definición de $\mathfrak{S}$ \textup{(resp. $\mathfrak{X}$, $\mathfrak{Y}$)}, y supongamos que $f^*(\mathscr{J})\mathscr{O}_{\mathfrak{X}}\subset\mathscr{K}$ y $g^*(\mathscr{J})\mathscr{O}_{\mathfrak{Y}}=\mathscr{L}$; pongamos $S_n=(\mathfrak{S},
\mathscr{O}_{\mathfrak{S}}/\mathscr{J}^{n+1})$, $X_n=(\mathfrak{X},
\mathscr{O}_{\mathfrak{X}}/\mathscr{K}^{n+1})$ y $Y_n=(\mathfrak{Y},
\mathscr{O}_{\mathfrak{Y}}/\mathscr{L}^{n+1})$.
Hay entonces una correspondencia biyectiva canónica
\oldpage[I]{193}
entre $\operatorname{Hom}_{\mathfrak{S}}(\mathfrak{X},\mathfrak{Y})$ y el conjunto de las sucesiones $(u_n)$ de $S_n$-morfismos $u_n:X_n\to Y_n$ que hacen conmutativos los diagramas \sref{I.10.6.7.1}.
\end{corollary}

\begin{proof}
Para todo $\mathfrak{S}$-morfismo $u:\mathfrak{X}\to\mathfrak{Y}$, se tiene por definición $f=g\circ u$, de donde
\[
  u^*(\mathscr{L})\mathscr{O}_{\mathfrak{X}}
  =u^*(g^*(\mathscr{J})\mathscr{O}_{\mathfrak{Y}})
    \mathscr{O}_{\mathfrak{X}}
  =f^*(\mathscr{J})\mathscr{O}_{\mathfrak{X}}
  \subset\mathscr{K}
\]
y el corolario se deduce entonces de la Proposición~\sref{I.10.6.9}.
\end{proof}

Se observará que, para $m\leq n$, el dato de un morfismo $f_n:X_n\to Y_n$ determina un único morfismo $f_m:X_m\to Y_m$ que hace conmutativo el diagrama \sref{I.10.6.7.1}, como se ve inmediatamente reduciendo al caso afín; se ha definido así una aplicación $\varphi_{mn}:\operatorname{Hom}_{S_n}(X_n,Y_n)
\to\operatorname{Hom}_{S_m}(X_m,Y_m)$, y los $\operatorname{Hom}_{S_n}(X_n,Y_n)$ forman, para los $\varphi_{mn}$, un \emph{sistema proyectivo de conjuntos}; el Corolario~\sref{I.10.6.11} también se enuncia diciendo que existe una biyección canónica
\[
  \operatorname{Hom}_{\mathfrak{S}}(\mathfrak{X},\mathfrak{Y})
  \xrightarrow{\sim}
  \varprojlim_n\operatorname{Hom}_{S_n}(X_n,Y_n).
\]

\subsection{Producto de preesquemas formales}
\label{subsection:I.10.7}

\begin{env}[10.7.1]
\label{I.10.7.1}
Sea $\mathfrak{S}$ un preesquema formal; como los preesquemas formales sobre $\mathfrak{S}$ forman una categoría, se puede definir la noción de \emph{producto} de preesquemas formales sobre $\mathfrak{S}$.
\end{env}

\begin{proposition}[10.7.2]
\label{I.10.7.2}
Sean $\mathfrak{X}=\operatorname{Spf}(B)$ y $\mathfrak{Y}=\operatorname{Spf}(C)$ dos esquemas formales afines sobre un esquema formal afín $\mathfrak{S}=\operatorname{Spf}(A)$.
Sean $\mathfrak{Z}=\operatorname{Spf}(B\widehat{\otimes}_A C)$, $p_1$ y $p_2$ los $\mathfrak{S}$-morfismos que corresponden, según \sref{I.10.2.2}, a los $A$-homomorfismos canónicos (continuos) $\rho$ y $\sigma$ de $B$ y $C$ en $B\widehat{\otimes}_A C$; entonces $(\mathfrak{Z},p_1,p_2)$ es un producto de los esquemas formales afines sobre $\mathfrak{S}$, $\mathfrak{X}$ y $\mathfrak{Y}$.
\end{proposition}

\begin{proof}
Por \sref{I.10.4.6}, basta comprobar que la aplicación que asocia a cada $A$-homomorfismo continuo $\varphi:B\widehat{\otimes}_A C\to D$ (donde $D$ es un anillo admisible que es una $A$-álgebra topológica) el par $(\varphi\circ\rho,\varphi\circ\sigma)$ es una biyección
\[
  \operatorname{Hom}_A(B\widehat{\otimes}_A C,D)
  \xrightarrow{\sim}
  \operatorname{Hom}_A(B,D)\times\operatorname{Hom}_A(C,D)
\]
que no es sino la propiedad universal del producto tensorial completado \sref[0]{0.7.7.6}.
\end{proof}

\begin{proposition}[10.7.3]
\label{I.10.7.3}
Dados dos $\mathfrak{S}$-preesquemas formales $\mathfrak{X}$ y $\mathfrak{Y}$, existe su producto $\mathfrak{X}\times_{\mathfrak{S}}\mathfrak{Y}$.
\end{proposition}

\begin{proof}
La demostración es idéntica a la de \sref{I.3.2.6}, sustituyendo los esquemas afines (resp. los abiertos afines) por los esquemas formales afines (resp. los abiertos formales afines), y la Proposición \sref{I.3.2.2} por la Proposición \sref{I.10.7.2}.
\end{proof}

Todas las propiedades formales del producto de preesquemas (\sref{I.3.2.7} y \sref{I.3.2.8}, de \sref{I.3.3.1} a \sref{I.3.3.12}) son válidas sin modificación alguna para el producto de preesquemas formales.

\begin{env}[10.7.4]
\label{I.10.7.4}
Sean $\mathfrak{S}$, $\mathfrak{X}$ y $\mathfrak{Y}$ preesquemas formales, y sean $f:\mathfrak{X}\to\mathfrak{S}$ y $g:\mathfrak{Y}\to\mathfrak{S}$ dos morfismos.
Supongamos que existen, respectivamente en $\mathfrak{S}$, $\mathfrak{X}$ y $\mathfrak{Y}$, tres sistemas fundamentales de haces de ideales de definición $(\mathscr{J}_\lambda)$, $(\mathscr{K}_\lambda)$ y $(\mathscr{L}_\lambda)$, todos con el mismo conjunto de índices $I$, tales que $f^*(\mathscr{J}_\lambda)\mathscr{O}_{\mathfrak{X}}
\subset\mathscr{K}_\lambda$ y $g^*(\mathscr{J}_\lambda)\mathscr{O}_{\mathfrak{Y}}
\subset\mathscr{L}_\lambda$ para todo $\lambda$.
Pongamos $S_\lambda=(\mathfrak{S},
\mathscr{O}_{\mathfrak{S}}/\mathscr{J}_\lambda)$, $X_\lambda=(\mathfrak{X},
\mathscr{O}_{\mathfrak{X}}/\mathscr{K}_\lambda)$ e $Y_\lambda=(\mathfrak{Y},
\mathscr{O}_{\mathfrak{Y}}/\mathscr{L}_\lambda)$; si $\mathscr{J}_\mu\subset\mathscr{J}_\lambda$, $\mathscr{K}_\mu\subset\mathscr{K}_\lambda$ y $\mathscr{L}_\mu\subset\mathscr{L}_\lambda$, observemos que $S_\lambda$ (resp. $X_\lambda$, $Y_\lambda$) es un subpreesquema cerrado de $S_\mu$ (resp. $X_\mu$, $Y_\mu$) con el \emph{mismo}
\oldpage[I]{194}
espacio subyacente \sref{I.10.6.1}.
Como $S_\lambda\to S_\mu$ es un monomorfismo de preesquemas, se ve primero que los productos $X_\lambda\times_{S_\lambda}Y_\lambda$ y $X_\lambda\times_{S_\mu}Y_\lambda$ son idénticos \sref{I.3.2.4}, y después que $X_\lambda\times_{S_\mu}Y_\lambda$ se identifica con un subpreesquema cerrado de $X_\mu\times_{S_\mu}Y_\mu$ con el \emph{mismo} espacio subyacente \sref{I.4.3.1}.
Así pues, el producto $\mathfrak{X}\times_{\mathfrak{S}}\mathfrak{Y}$ es el \emph{límite inductivo} de los preesquemas usuales $X_\lambda\times_{S_\lambda}Y_\lambda$: en efecto, como en \sref{I.10.6.2}, se puede reducir al caso en que $\mathfrak{S}$, $\mathfrak{X}$ y $\mathfrak{Y}$ son esquemas formales afines.
Teniendo en cuenta \sref{I.10.5.6}[ii] y las hipótesis sobre los sistemas fundamentales de haces de ideales de definición de $\mathfrak{S}$, $\mathfrak{X}$ y $\mathfrak{Y}$, se ve inmediatamente que nuestra afirmación resulta de la definición del producto tensorial completado de dos álgebras \sref[0]{0.7.7.1}.

Además, sea $\mathfrak{Z}$ un $\mathfrak{S}$-preesquema formal, sea $(\mathscr{M}_\lambda)$ un sistema fundamental de ideales de definición de $\mathfrak{Z}$ cuyo conjunto de índices es $I$, y sean $u:\mathfrak{Z}\to\mathfrak{X}$ y $v:\mathfrak{Z}\to\mathfrak{Y}$ dos $\mathfrak{S}$-morfismos tales que $u^*(\mathscr{K}_\lambda)\mathscr{O}_{\mathfrak{Z}}
\subset\mathscr{M}_\lambda$ y $v^*(\mathscr{L}_\lambda)\mathscr{O}_{\mathfrak{Z}}
\subset\mathscr{M}_\lambda$.
Si se pone $Z_\lambda=(\mathfrak{Z},
\mathscr{O}_{\mathfrak{Z}}/\mathscr{M}_\lambda)$, y si $u_\lambda:Z_\lambda\to X_\lambda$ y $v_\lambda:Z_\lambda\to Y_\lambda$ son los $S_\lambda$-morfismos correspondientes a $u$ y $v$ \sref{I.10.5.6}, se comprueba inmediatamente que $(u,v)_{\mathfrak{S}}$ es el límite inductivo de los $S_\lambda$-morfismos $(u_\lambda,v_\lambda)_{S_\lambda}$.

Las consideraciones de este apartado se aplican, en particular, cuando $\mathfrak{S}$, $\mathfrak{X}$ y $\mathfrak{Y}$ son localmente noetherianos, tomando como sistemas fundamentales de haces de ideales de definición los sistemas formados por las potencias de un haz de ideales de definición \sref{I.10.5.1}.
Sin embargo, se observará que $\mathfrak{X}\times_{\mathfrak{S}}\mathfrak{Y}$ no es necesariamente localmente noetheriano (véase, no obstante, \sref{I.10.13.5}).
\end{env}

\subsection{Completado formal de un preesquema a lo largo de una parte cerrada}
\label{subsection:I.10.8}

\begin{env}[10.8.1]
\label{I.10.8.1}
Sea $X$ un preesquema (usual) \emph{localmente noetheriano}, y sea $X'$ una parte cerrada del espacio subyacente a $X$; designemos por $\Phi$ el conjunto de haces \emph{coherentes} de ideales $\mathscr{J}$ de $\mathscr{O}_X$ tales que el soporte de $\mathscr{O}_X/\mathscr{J}$ sea $X'$.
El conjunto $\Phi$ no es vacío (\sref{I.5.2.1}, \sref{I.4.1.4}, \sref{I.6.1.1}); lo ordenaremos mediante la relación $\supset$.
\end{env}

\begin{lemma}[10.8.2]
\label{I.10.8.2}
El conjunto ordenado $\Phi$ es filtrante; si $X$ es noetheriano, para todo $\mathscr{J}_0\in\Phi$, el conjunto de las potencias $\mathscr{J}_0^n$ ($n>0$) es cofinal en $\Phi$.
\end{lemma}

\begin{proof}
En efecto, si $\mathscr{J}_1$ y $\mathscr{J}_2$ pertenecen a $\Phi$, y si se pone $\mathscr{J}=\mathscr{J}_1\cap\mathscr{J}_2$, entonces $\mathscr{J}$ es coherente porque $\mathscr{O}_X$ es coherente (\sref{I.6.1.1} y \sref[0]{0.5.3.4}), y se tiene $\mathscr{J}_x=(\mathscr{J}_1)_x\cap(\mathscr{J}_2)_x$, para todo $x\in X$; por tanto, $\mathscr{J}_x=\mathscr{O}_x$ para $x\notin X'$, y $\mathscr{J}_x\neq\mathscr{O}_x$ para $x\in X'$, lo que demuestra que $\mathscr{J}\in\Phi$.
Por otra parte, si $X$ es noetheriano y $\mathscr{J}_0$ y $\mathscr{J}$ pertenecen a $\Phi$, existe un entero $n>0$ tal que $\mathscr{J}_0^n(\mathscr{O}_X/\mathscr{J})=0$ \sref{I.9.3.4}, lo que implica que $\mathscr{J}_0^n\subset\mathscr{J}$.
\end{proof}

\begin{env}[10.8.3]
\label{I.10.8.3}
Sea ahora $\mathscr{F}$ un $\mathscr{O}_X$-módulo \emph{coherente}; para todo $\mathscr{J}\in\Phi$, $\mathscr{F}\otimes_{\mathscr{O}_X}(\mathscr{O}_X/\mathscr{J})$ es un $\mathscr{O}_X$-módulo coherente \sref{I.9.1.1}, cuyo soporte está contenido en $X'$, y que identificaremos casi siempre con su restricción a $X'$.
Cuando $\mathscr{J}$ recorre $\Phi$, estos haces forman un \emph{sistema proyectivo} de haces de grupos abelianos.
\end{env}

\begin{definition}[10.8.4]
\label{I.10.8.4}
{\itshape Dada una parte cerrada $X'$ de un preesquema localmente noetheriano $X$ y un $\mathscr{O}_X$-módulo coherente $\mathscr{F}$, se llama completado de $\mathscr{F}$ a lo largo de $X'$, y se denota por $\mathscr{F}_{/X'}$ (o por $\widehat{\mathscr{F}}$ cuando no haya riesgo de confusión), a la restricción a $X'$ del haz
\oldpage[I]{195}
$\varprojlim_{\Phi}
(\mathscr{F}\otimes_{\mathscr{O}_X}
(\mathscr{O}_X/\mathscr{J}))$; se dice que sus secciones sobre $X'$ son las secciones formales de $\mathscr{F}$ a lo largo de $X'$.}
\end{definition}

Es inmediato que, para todo abierto $U\subset X$, se tiene $(\mathscr{F}|U)_{/(U\cap X')}=(\mathscr{F}_{/X'})|(U\cap X')$.

Por paso al límite proyectivo, está claro que $(\mathscr{O}_X)_{/X'}$ es un haz de anillos, y que $\mathscr{F}_{/X'}$ puede considerarse como un $(\mathscr{O}_X)_{/X'}$-módulo.
Además, como existe una base de la topología de $X'$ formada por abiertos cuasicompactos, se puede considerar $(\mathscr{O}_X)_{/X'}$ (resp. $\mathscr{F}_{/X'}$) como un \emph{haz de anillos topológicos} (resp. de \emph{grupos topológicos}), límite proyectivo de los haces de anillos (resp. grupos) \emph{pseudodiscretos} $\mathscr{O}_X/\mathscr{J}$ (resp. $\mathscr{F}\otimes_{\mathscr{O}_X}
(\mathscr{O}_X/\mathscr{J})=\mathscr{F}/\mathscr{J}\mathscr{F}$), y, por paso al límite proyectivo, $\mathscr{F}_{/X'}$ se convierte entonces en un \emph{$(\mathscr{O}_X)_{/X'}$-módulo topológico} (\sref[0]{0.3.8.1} y \sref[0]{0.3.8.2}); recordemos que, para todo abierto \emph{cuasicompacto} $U\subset X$, $\Gamma(U\cap X',(\mathscr{O}_X)_{/X'})$ (resp. $\Gamma(U\cap X',\mathscr{F}_{/X'})$) es entonces el límite proyectivo de los anillos (resp. grupos) discretos $\Gamma(U,\mathscr{O}_X/\mathscr{J})$ (resp. $\Gamma(U,\mathscr{F}/\mathscr{J}\mathscr{F})$).

Si ahora $u:\mathscr{F}\to\mathscr{G}$ es un homomorfismo de $\mathscr{O}_X$-módulos, se deducen canónicamente homomorfismos $u_{\mathscr{J}}:
\mathscr{F}\otimes_{\mathscr{O}_X}(\mathscr{O}_X/\mathscr{J})
\to
\mathscr{G}\otimes_{\mathscr{O}_X}(\mathscr{O}_X/\mathscr{J})$ para todo $\mathscr{J}\in\Phi$, y estos homomorfismos forman un sistema proyectivo.
Por paso al límite proyectivo y restricción a $X'$, dan un $(\mathscr{O}_X)_{/X'}$-homomorfismo continuo $\mathscr{F}_{/X'}\to\mathscr{G}_{/X'}$, denotado por $u_{/X'}$ o $\widehat u$, y llamado el \emph{completado} del homomorfismo $u$ a lo largo de $X'$.
Está claro que, si $v:\mathscr{G}\to\mathscr{H}$ es un segundo homomorfismo de $\mathscr{O}_X$-módulos, se tiene $(v\circ u)_{/X'}=(v_{/X'})\circ(u_{/X'})$; por tanto, $\mathscr{F}_{/X'}$ es un \emph{funtor aditivo covariante} en $\mathscr{F}$, de la categoría de los $\mathscr{O}_X$-módulos coherentes a la categoría de los $(\mathscr{O}_X)_{/X'}$-módulos topológicos.

\begin{proposition}[10.8.5]
\label{I.10.8.5}
El soporte de $(\mathscr{O}_X)_{/X'}$ es $X'$; el espacio topológicamente anillado $(X',(\mathscr{O}_X)_{/X'})$ es un preesquema formal localmente noetheriano y, si $\mathscr{J}\in\Phi$, entonces $\mathscr{J}_{/X'}$ es un haz de ideales de definición de este preesquema formal.
Si $X=\operatorname{Spec}(A)$ es un esquema afín de anillo noetheriano, $\mathscr{J}=\widetilde{\mathfrak{J}}$, donde $\mathfrak{J}$ es un ideal de $A$, y $X'=V(\mathfrak{J})$, entonces $(X',(\mathscr{O}_X)_{/X'})$ se identifica canónicamente con $\operatorname{Spf}(\widehat A)$, donde $\widehat A$ es el completado separado de $A$ para la topología $\mathfrak{J}$-preádica.
\end{proposition}

\begin{proof}
Evidentemente, basta demostrar la última afirmación.
Se sabe por \sref[0]{0.7.3.3} que el completado separado $\widehat{\mathfrak{J}}$ de $\mathfrak{J}$ para la topología $\mathfrak{J}$-preádica se identifica con el ideal $\mathfrak{J}\widehat A$ de $\widehat A$, y que $\widehat A$ es un anillo noetheriano $\widehat{\mathfrak{J}}$-ádico tal que $\widehat A/\widehat{\mathfrak{J}}^n=A/\mathfrak{J}^n$ \sref[0]{0.7.2.6}.
Esta última igualdad muestra que los ideales primos abiertos de $\widehat A$ son los ideales $\widehat{\mathfrak p}=\mathfrak p\widehat A$, donde $\mathfrak p$ es un ideal primo de $A$ que contiene a $\mathfrak{J}$, y que $\widehat{\mathfrak p}\cap A=\mathfrak p$; por tanto, $\operatorname{Spf}(\widehat A)=X'$.
Como $\mathscr{O}_X/\mathscr{J}^n=(A/\mathfrak{J}^n)^{\sim}$, la proposición se deduce inmediatamente de las definiciones.
\end{proof}

Se dice que el preesquema formal así definido es el \emph{completado de $X$ a lo largo de $X'$}, y se denota por $X_{/X'}$ o por $\widehat X$ cuando no haya riesgo de confusión.
Cuando se toma $X'=X$, se puede tomar $\mathscr{J}=0$, y se tiene por tanto $X_{/X}=X$.

Está claro que, si $U$ es un subpreesquema inducido sobre un abierto de $X$, entonces $U_{/(U\cap X')}$ se identifica canónicamente con el subpreesquema formal inducido por $X_{/X'}$ sobre el abierto $U\cap X'$ de $X'$.

\begin{corollary}[10.8.6]
\label{I.10.8.6}
El preesquema (usual) $\widehat X_{\mathrm{red}}$ es el único subpreesquema reducido de $X$ cuyo espacio subyacente es $X'$ \sref{I.5.2.1}.
Para que $\widehat X$ sea noetheriano, es necesario y suficiente que $\widehat X_{\mathrm{red}}$ lo sea, y basta que $X$ lo sea.
\end{corollary}

\oldpage[I]{196}

\begin{proof}
Como la determinación de $\widehat X_{\mathrm{red}}$ es local \sref{I.10.5.4}, se puede suponer que $X$ es un esquema afín de anillo noetheriano; con las notaciones de \sref{I.10.8.5}, el ideal $\mathfrak{T}$ de los elementos topológicamente nilpotentes de $\widehat A$ es la imagen inversa, por la aplicación canónica $\widehat A\to\widehat A/\widehat{\mathfrak{J}}=A/\mathfrak{J}$, del nilradical de $A/\mathfrak{J}$ \sref[0]{0.7.1.3}; por tanto, $\widehat A/\mathfrak{T}$ es isomorfo al cociente de $A/\mathfrak{J}$ por su nilradical.
La primera afirmación resulta entonces de \sref{I.10.5.4} y \sref{I.5.1.1}.
Si $\widehat X_{\mathrm{red}}$ es noetheriano, su espacio subyacente $X'$ también lo es; por tanto, los $X'_n=\operatorname{Spec}(\mathscr{O}_X/\mathscr{J}^n)$ son noetherianos \sref{I.6.1.2}, y lo mismo ocurre con $\widehat X$ \sref{I.10.6.4}; la recíproca es inmediata por \sref{I.6.1.2}.
\end{proof}

\begin{env}[10.8.7]
\label{I.10.8.7}
Los homomorfismos canónicos $\mathscr{O}_X\to\mathscr{O}_X/\mathscr{J}$ (para $\mathscr{J}\in\Phi$) forman un sistema proyectivo y dan, al pasar al límite proyectivo, un homomorfismo de haces de anillos $\theta:\mathscr{O}_X\to\psi_*((\mathscr{O}_X)_{/X'})
=\varprojlim_\Phi(\mathscr{O}_X/\mathscr{J})$, donde $\psi$ designa la inyección canónica $X'\to X$ de los espacios subyacentes.
Denotaremos por $i$ (o $i_X$) el morfismo (llamado \emph{canónico})
\[
  (\psi,\theta):X_{/X'}\to X
\]
de espacios anillados.

Por tensorización, para todo $\mathscr{O}_X$-módulo coherente $\mathscr{F}$, los homomorfismos canónicos $\mathscr{O}_X\to\mathscr{O}_X/\mathscr{J}$ dan homomorfismos $\mathscr{F}\to
\mathscr{F}\otimes_{\mathscr{O}_X}(\mathscr{O}_X/\mathscr{J})$ de $\mathscr{O}_X$-módulos que forman también un sistema proyectivo y dan, por tanto, al pasar al límite proyectivo, un homomorfismo canónico funtorial $\gamma:\mathscr{F}\to\psi_*(\mathscr{F}_{/X'})$ de $\mathscr{O}_X$-módulos.
\end{env}

\begin{proposition}[10.8.8]
\label{I.10.8.8}
\medskip\noindent
\begin{enumerate}
  \item[\rm{(i)}] El funtor $\mathscr{F}_{/X'}$ (en $\mathscr{F}$) es exacto.
  \item[\rm{(ii)}] El homomorfismo funtorial $\gamma^\sharp:i^*(\mathscr{F})\to\mathscr{F}_{/X'}$ de $(\mathscr{O}_X)_{/X'}$-módulos es un isomorfismo.
\end{enumerate}
\end{proposition}

\begin{proof}
\medskip\noindent
\begin{enumerate}
  \item[\rm{(i)}] Basta demostrar que, si $0\to\mathscr{F}'\to\mathscr{F}\to\mathscr{F}''\to0$ es una sucesión exacta de $\mathscr{O}_X$-módulos coherentes y si $U$ es un abierto afín de $X$ de anillo noetheriano $A$, entonces la sucesión
\[
      0\to\Gamma(U\cap X',\mathscr{F}'_{/X'})
      \to\Gamma(U\cap X',\mathscr{F}_{/X'})
      \to\Gamma(U\cap X',\mathscr{F}''_{/X'})\to0
    \]
es exacta.
Se tiene $\mathscr{F}|U=\widetilde M$, $\mathscr{F}'|U=\widetilde{M}'$ y $\mathscr{F}''|U=\widetilde{M}''$, donde $M$, $M'$ y $M''$ son tres $A$-módulos de tipo finito tales que la sucesión $0\to M'\to M\to M''\to0$ es exacta (\sref{I.1.5.1} y \sref{I.1.3.11}); sea $\mathscr{J}\in\Phi$ y sea $\mathfrak{J}$ un ideal de $A$ tal que $\mathscr{J}|U=\widetilde{\mathfrak{J}}$.
Se tiene entonces
\[
      \Gamma(U\cap X',
      \mathscr{F}\otimes_{\mathscr{O}_X}\mathscr{O}_X/\mathscr{J}^n)
      =M\otimes_A(A/\mathfrak{J}^n)
    \]
\sref{I.1.3.12}; por tanto, por definición del límite proyectivo, se tiene
\[
      \Gamma(U\cap X',\mathscr{F}_{/X'})
      =\varprojlim_n(M\otimes_A(A/\mathfrak{J}^n))=\widehat M
    \]
el completado separado de $M$ para la topología $\mathfrak{J}$-preádica, y análogamente
\[
      \Gamma(U\cap X',\mathscr{F}'_{/X'})=\widehat{M}',\qquad
      \Gamma(U\cap X',\mathscr{F}''_{/X'})=\widehat{M}''
    \]
nuestra afirmación resulta entonces de que, siendo $A$ noetheriano, el funtor $\widehat M$ en $M$ es exacto sobre la categoría de los $A$-módulos de tipo finito \sref[0]{0.7.3.3}.
\oldpage[I]{197}
  \item[\rm{(ii)}] La cuestión es local, por lo que se puede suponer que se tiene una sucesión exacta $\mathscr{O}_X^m\to\mathscr{O}_X^n\to\mathscr{F}\to0$ \sref[0]{0.5.3.2}; como $\gamma^\sharp$ es funtorial y los funtores $i^*(\mathscr{F})$ y $\mathscr{F}_{/X'}$ son exactos por la derecha (por (i) y \sref[0]{0.4.3.1}), se tiene el diagrama conmutativo
\[
      \xymatrix{
        i^*(\mathscr{O}_X^m)\ar[r]\ar[d]_{\gamma^\sharp} &
        i^*(\mathscr{O}_X^n)\ar[r]\ar[d]_{\gamma^\sharp} &
        i^*(\mathscr{F})\ar[r]\ar[d]_{\gamma^\sharp} &
        0\\
        (\mathscr{O}_X^m)_{/X'}\ar[r] &
        (\mathscr{O}_X^n)_{/X'}\ar[r] &
        \mathscr{F}_{/X'}\ar[r] &
        0
      }
      \tag{10.8.8.1}\label{I.10.8.8.1-fr}
    \]
cuyas filas son exactas.
Además, los funtores $i^*(\mathscr{F})$ y $\mathscr{F}_{/X'}$ conmutan con las sumas directas finitas (\sref[0]{0.3.2.6} y \sref[0]{0.4.3.2}), por lo que se reduce a demostrar nuestra afirmación para $\mathscr{F}=\mathscr{O}_X$.
Se tiene $i^*(\mathscr{O}_X)=(\mathscr{O}_X)_{/X'}=
    \mathscr{O}_{\widehat X}$ \sref[0]{0.4.3.4}, y $\gamma^\sharp$ es un homomorfismo de $\mathscr{O}_{\widehat X}$-\emph{módulos}; basta, pues, comprobar que $\gamma^\sharp$ transforma en sí misma la sección unidad de $\mathscr{O}_{\widehat X}$ sobre un abierto de $X'$, lo cual es inmediato y muestra que, en este caso, $\gamma^\sharp$ es la identidad.
\end{enumerate}
\end{proof}

\begin{corollary}[10.8.9]
\label{I.10.8.9}
El morfismo $i:X_{/X'}\to X$ de espacios anillados es plano.
\end{corollary}

\begin{proof}
Esto resulta de \sref[0]{0.6.7.3} y de \sref{I.10.8.8}[i].
\end{proof}

\begin{corollary}[10.8.10]
\label{I.10.8.10}
Si $\mathscr{F}$ y $\mathscr{G}$ son $\mathscr{O}_X$-módulos coherentes, existen isomorfismos canónicos funtoriales (en $\mathscr{F}$ y $\mathscr{G}$)
\[
  (\mathscr{F}_{/X'})\otimes_{(\mathscr{O}_X)_{/X'}}
  (\mathscr{G}_{/X'})
  \xrightarrow{\sim}
  (\mathscr{F}\otimes_{\mathscr{O}_X}\mathscr{G})_{/X'}
  \tag{10.8.10.1}\label{I.10.8.10.1-fr}
\]
\[
  (\mathscr{H}\!om_{\mathscr{O}_X}(\mathscr{F},\mathscr{G}))_{/X'}
  \xrightarrow{\sim}
  \mathscr{H}\!om_{(\mathscr{O}_X)_{/X'}}
  (\mathscr{F}_{/X'},\mathscr{G}_{/X'})
  \tag{10.8.10.2}\label{I.10.8.10.2-fr}
\]
\end{corollary}

\begin{proof}
Esto resulta de la identificación canónica de $i^*(\mathscr{F})$ con $\mathscr{F}_{/X'}$; la existencia del primer isomorfismo es entonces un resultado válido para todos los morfismos de espacios anillados \sref[0]{0.4.3.3.1}, y la del segundo es un resultado válido para todos los morfismos planos \sref[0]{0.6.7.6}, por \sref{I.10.8.9}.
\end{proof}

\begin{proposition}[10.8.11]
\label{I.10.8.11}
Para todo $\mathscr{O}_X$-módulo coherente $\mathscr{F}$, el núcleo del homomorfismo canónico $\Gamma(X,\mathscr{F})\to\Gamma(X',\mathscr{F}_{/X'})$ inducido por $\mathscr{F}\to\mathscr{F}_{/X'}$ está formado por las secciones nulas en alguna vecindad de $X'$.
\end{proposition}

\begin{proof}
De la definición de $\mathscr{F}_{/X'}$ resulta que la imagen canónica de una sección tal es nula.
Recíprocamente, si $s\in\Gamma(X,\mathscr{F})$ tiene imagen nula en $\Gamma(X',\mathscr{F}_{/X'})$, basta ver que todo $x\in X'$ admite una vecindad en $X$ en la que $s$ es nula, y se puede reducir así al caso en que $X=\operatorname{Spec}(A)$ es afín, $A$ es noetheriano, $X'=V(\mathfrak{J})$ para algún ideal $\mathfrak{J}$ de $A$, y $\mathscr{F}=\widetilde M$ para algún $A$-módulo $M$ de tipo finito.
Entonces $\Gamma(X',\mathscr{F}_{/X'})$ es el completado separado $\widehat M$ de $M$ para la topología $\mathfrak{J}$-preádica, y el homomorfismo $\Gamma(X,\mathscr{F})\to\Gamma(X',\mathscr{F}_{/X'})$ es el homomorfismo canónico $M\to\widehat M$.
Se sabe por \sref[0]{0.7.3.7} que el núcleo de este homomorfismo es el conjunto de los $z\in M$ anulados por un elemento de $1+\mathfrak{J}$.
Por tanto, se tiene $(1+f)s=0$ para algún $f\in\mathfrak{J}$; para todo $x\in X'$ se deduce $(1_x+f_x)s_x=0$, y como $1_x+f_x$ es invertible en $\mathscr{O}_x$ (pues $\mathfrak{J}_x\mathscr{O}_x$ está contenido en el ideal maximal de $\mathscr{O}_x$), se tiene $s_x=0$, lo que demuestra la proposición.
\end{proof}

\begin{corollary}[10.8.12]
\label{I.10.8.12}
El soporte de $\mathscr{F}_{/X'}$ es igual a $\operatorname{Supp}(\mathscr{F})\cap X'$.
\end{corollary}

\begin{proof}
Está claro que $\mathscr{F}_{/X'}$ es un $(\mathscr{O}_X)_{/X'}$-módulo de tipo finito (\sref{I.10.8.8}[ii] y \sref[0]{0.5.2.4}),
\oldpage[I]{198}
por lo que su soporte es cerrado \sref[0]{0.5.2.2} y está evidentemente contenido en $\operatorname{Supp}(\mathscr{F})\cap X'$.
Para demostrar que es igual a este último conjunto, se reduce de inmediato a probar que la igualdad $\Gamma(X',\mathscr{F}_{/X'})=0$ implica $\operatorname{Supp}(\mathscr{F})\cap X'=\varnothing$; esto resulta de \sref{I.10.8.11} y de \sref{I.1.4.1}.
\end{proof}

\begin{corollary}[10.8.13]
\label{I.10.8.13}
Sea $u:\mathscr{F}\to\mathscr{G}$ un homomorfismo de $\mathscr{O}_X$-módulos coherentes.
Para que $u_{/X'}:\mathscr{F}_{/X'}\to\mathscr{G}_{/X'}$ sea nulo, es necesario y suficiente que $u$ sea nulo en una vecindad de $X'$.
\end{corollary}

\begin{proof}
En efecto, por \sref{I.10.8.8}[ii], $u_{/X'}$ se identifica con $i^*(u)$; por tanto, si se considera $u$ como una sección sobre $X$ del haz $\mathscr{H}=\mathscr{H}\!om_{\mathscr{O}_X}(\mathscr{F},\mathscr{G})$, $u_{/X'}$ es la sección sobre $X'$ de $i^*(\mathscr{H})=\mathscr{H}_{/X'}$ que le corresponde canónicamente (\sref{I.10.8.10.2} y \sref[0]{0.4.4.6}).
Basta, pues, aplicar \sref{I.10.8.11} al $\mathscr{O}_X$-módulo coherente $\mathscr{H}$.
\end{proof}

\begin{corollary}[10.8.14]
\label{I.10.8.14}
Sea $u:\mathscr{F}\to\mathscr{G}$ un homomorfismo de $\mathscr{O}_X$-módulos coherentes.
Para que $u_{/X'}$ sea un monomorfismo \textup{(resp.} un epimorfismo), es necesario y suficiente que $u$ sea un monomorfismo \textup{(resp.} un epimorfismo) en una vecindad de $X'$.
\end{corollary}

\begin{proof}
Sean $\mathscr{P}$ y $\mathscr{N}$ el conúcleo y el núcleo, respectivamente, de $u$, de modo que se tiene la sucesión exacta \[
  0\to\mathscr{N}\xrightarrow{v}\mathscr{F}\xrightarrow{u}
  \mathscr{G}\xrightarrow{w}\mathscr{P}\to0,
\], y por tanto, en virtud de \sref{I.10.8.8}[i], la sucesión exacta
\[
  0\to\mathscr{N}_{/X'}\xrightarrow{v_{/X'}}
  \mathscr{F}_{/X'}\xrightarrow{u_{/X'}}
  \mathscr{G}_{/X'}\xrightarrow{w_{/X'}}
  \mathscr{P}_{/X'}\to0.
\]
Si $u_{/X'}$ es un monomorfismo (resp. un epimorfismo), se tiene $v_{/X'}=0$ (resp. $w_{/X'}=0$); por tanto, existe una vecindad de $X'$ en la que $v=0$ (resp. $w=0$), en virtud de \sref{I.10.8.13}.
\end{proof}

\subsection{Prolongación de un morfismo a los completados}
\label{subsection:I.10.9}

\begin{env}[10.9.1]
\label{I.10.9.1}
Sean $X$ e $Y$ dos preesquemas (usuales) localmente noetherianos, $f:X\to Y$ un morfismo, y $X'$ (resp. $Y'$) una parte cerrada del espacio subyacente $X$ (resp. $Y$), tales que $f(X')\subset Y'$.
Sea $\mathscr{J}$ (resp. $\mathscr{K}$) un haz de ideales de $\mathscr{O}_X$ (resp. $\mathscr{O}_Y$) tal que el soporte de $\mathscr{O}_X/\mathscr{J}$ (resp. $\mathscr{O}_Y/\mathscr{K}$) sea $X'$ (resp. $Y'$) y que $f^*(\mathscr{K})\mathscr{O}_X\subset\mathscr{J}$; se observará que siempre existen tales haces de ideales, pues, por ejemplo, puede tomarse como $\mathscr{J}$ el mayor haz de ideales de $\mathscr{O}_X$ que define un subpreesquema de $X$ cuyo espacio subyacente es $X'$ \sref{I.5.2.1}, y la hipótesis $f(X')\subset Y'$ implica entonces $f^*(\mathscr{K})\mathscr{O}_X\subset\mathscr{J}$ \sref{I.5.2.4}.
Para todo entero $n>0$ se tiene $f^*(\mathscr{K}^n)\mathscr{O}_X\subset\mathscr{J}^n$ \sref[0]{0.4.3.5}; por consiguiente \sref{I.4.4.6}, si se pone \[
  X'_n=(X',\mathscr{O}_X/\mathscr{J}^{n+1}),\qquad
  Y'_n=(Y',\mathscr{O}_Y/\mathscr{K}^{n+1}),
\], $f$ induce un morfismo $f_n:X'_n\to Y'_n$, y es inmediato que los $f_n$ forman un sistema inductivo.
Denotaremos su límite inductivo \sref{I.10.6.8} por $\widehat f:X_{/X'}\to Y_{/Y'}$, y diremos (por abuso de lenguaje) que $\widehat f$ es la \emph{prolongación de $f$ a los completados de $X$ e $Y$ a lo largo de $X'$ e $Y'$}.
Se comprueba inmediatamente que este morfismo no depende de la elección de los haces de ideales $\mathscr{J}$ y $\mathscr{K}$ que satisfacen las condiciones anteriores.
Basta, en efecto, verlo cuando $X$ e $Y$ son esquemas afines noetherianos de anillos $A$ y $B$, respectivamente; entonces $\mathscr{J}=\widetilde{\mathfrak{J}}$ y $\mathscr{K}=\widetilde{\mathfrak{K}}$, donde $\mathfrak{J}$ (resp. $\mathfrak{K}$) es un ideal de $A$ (resp. $B$), y $f$ corresponde a un homomorfismo de anillos $\varphi:B\to A$ tal que $\varphi(\mathfrak{K})\subset\mathfrak{J}$ (\sref{I.4.4.6} y \sref{I.1.7.4}); \emph{$\widehat f$ es entonces el morfismo que corresponde \sref{I.10.2.2} al homomorfismo continuo $\widehat\varphi:\widehat B\to\widehat A$}, donde $\widehat A$ (resp. $\widehat B$) es el completado separado de $A$ (resp. $B$) para la topología $\mathfrak{J}$-preádica (resp. $\mathfrak{K}$-preádica) \sref{I.10.6.8}; se sabe que, si se reemplaza $\mathscr{J}$ por otro
\oldpage[I]{199}
haz de ideales $\mathscr{J}'=\widetilde{\mathfrak{J}}'$ tal que el soporte de $\mathscr{O}_X/\mathscr{J}'$ sea $X'$, entonces las topologías $\mathfrak{J}$-preádica y $\mathfrak{J}'$-preádica sobre $A$ son las mismas \sref{I.10.8.2}.

Se observará que, por definición, la aplicación continua $X'\to Y'$ de los espacios subyacentes a $X_{/X'}$ e $Y_{/Y'}$ que corresponde a $\widehat f$ no es sino la restricción de $f$ a $X'$.
\end{env}

\begin{env}[10.9.2]
\label{I.10.9.2}
De la definición anterior resulta inmediatamente que el diagrama de morfismos de espacios anillados
\[
  \label{I.10.9.2.1-fr}
  \xymatrix{
    \widehat X\ar[r]^{\widehat f}\ar[d]_{i_X} &
    \widehat Y\ar[d]^{i_Y}\\
    X\ar[r]_f &
    Y
  }
\]
es conmutativo, siendo las flechas verticales los morfismos canónicos \sref{I.10.8.7}.
\end{env}

\begin{env}[10.9.3]
\label{I.10.9.3}
Sean $Z$ un tercer preesquema, $g:Y\to Z$ un morfismo y $Z'$ una parte cerrada de $Z$ tal que $g(Y')\subset Z'$.
Si $\widehat g$ designa el completado del morfismo $g$ a lo largo de $Y'$ y $Z'$, resulta inmediatamente de \sref{I.10.9.1} que
$(g\circ f){\,\widehat{\,}\,}=\widehat g\circ\widehat f$.
\end{env}

\begin{proposition}[10.9.4]
\label{I.10.9.4}
Sean $X$ e $Y$ dos $S$-preesquemas localmente noetherianos, siendo $Y$ de tipo finito sobre $S$.
Sean $f$ y $g$ dos $S$-morfismos de $X$ en $Y$ tales que $f(X')\subset Y'$ y $g(X')\subset Y'$.
Para que $\widehat f=\widehat g$, es necesario y suficiente que $f$ y $g$ coincidan en una vecindad de $X'$.
\end{proposition}

\begin{proof}
La condición es evidentemente suficiente (sin hipótesis de finitud sobre $Y$).
Para ver que es necesaria, observemos primero que la hipótesis $\widehat f=\widehat g$ implica $f(x)=g(x)$ para todo $x\in X'$.
Por otra parte, al ser la cuestión local, se puede suponer que $X$ e $Y$ son, respectivamente, vecindades abiertas afines de $x$ y de $y=f(x)=g(x)$, con anillos noetherianos, que $S$ es afín y que $\Gamma(Y,\mathscr{O}_Y)$ es una $\Gamma(S,\mathscr{O}_S)$-álgebra de tipo finito \sref{I.6.3.3}.
Entonces $f$ y $g$ corresponden a dos $\Gamma(S,\mathscr{O}_S)$-homomorfismos $\rho$ y $\sigma$ de $\Gamma(Y,\mathscr{O}_Y)$ en $\Gamma(X,\mathscr{O}_X)$ \sref{I.1.7.3}, y, por hipótesis, las prolongaciones por continuidad de estos homomorfismos al completado separado de $\Gamma(Y,\mathscr{O}_Y)$ son iguales.
Se deduce de \sref{I.10.8.11} que, para toda sección $s\in\Gamma(Y,\mathscr{O}_Y)$, las secciones $\rho(s)$ y $\sigma(s)$ coinciden en una vecindad de $X'$ (dependiente de $s$); como $\Gamma(Y,\mathscr{O}_Y)$ es un álgebra de tipo finito sobre $\Gamma(S,\mathscr{O}_S)$, se deduce inmediatamente que existe una vecindad $V$ de $X'$ tal que $\rho(s)$ y $\sigma(s)$ coinciden en $V$ para \emph{toda} sección $s\in\Gamma(Y,\mathscr{O}_Y)$.
Si $h\in\Gamma(X,\mathscr{O}_X)$\footnote{\emph{[Trad.] La fuente francesa impresa tiene aquí $\Gamma(Y,\sh{O}_X)$; el dominio debe ser $X$ para que $D(h)\subset X$ tenga sentido.}} es tal que $D(h)$ es una vecindad de $x$ contenida en $V$, se deduce de lo anterior y de \sref{I.1.4.1}[\emph{d)}] que $f$ y $g$ coinciden en $D(h)$.
\end{proof}

\begin{proposition}[10.9.5]
\label{I.10.9.5}
Bajo las hipótesis de \sref{I.10.9.1}, para todo $\mathscr{O}_Y$-módulo coherente $\mathscr{G}$, existe un isomorfismo canónico funtorial de $(\mathscr{O}_X)_{/X'}$-módulos
\[
  (f^*(\mathscr{G}))_{/X'}\xrightarrow{\sim}
  \widehat f^*(\mathscr{G}_{/Y'})
\]
\end{proposition}

\begin{proof}
Si se identifica canónicamente $(f^*(\mathscr{G}))_{/X'}$ con $i_X^*(f^*(\mathscr{G}))$, y $\widehat f^*(\mathscr{G}_{/Y'})$ con $\widehat f^*(i_Y^*(\mathscr{G}))$ \sref{I.10.8.8}, la proposición resulta inmediatamente de la conmutatividad del diagrama de \sref{I.10.9.2}.
\end{proof}

\begin{env}[10.9.6]
\label{I.10.9.6}
Sean ahora $\mathscr{F}$ un $\mathscr{O}_X$-módulo coherente y $\mathscr{G}$ un $\mathscr{O}_Y$-módulo coherente.
Si $u:\mathscr{G}\to\mathscr{F}$ es un $f$-morfismo de $\mathscr{G}$ en $\mathscr{F}$, le corresponde un $\mathscr{O}_X$-homomorfismo $u^\sharp:f^*(\mathscr{G})\to\mathscr{F}$, y por completación un $(\mathscr{O}_X)_{/X'}$-homomorfismo continuo $(u^\sharp)_{/X'}:(f^*(\mathscr{G}))_{/X'}\to\mathscr{F}_{/X'}$; en virtud de \sref{I.10.9.5}, existe un $\widehat f$-morfismo $v:\mathscr{G}_{/Y'}\to\mathscr{F}_{/X'}$,
\oldpage[I]{200}
y uno solo, tal que $v^\sharp=(u^\sharp)_{/X'}$.
Si se consideran las ternas $(\mathscr{F},X,X')$ (siendo $\mathscr{F}$ un $\mathscr{O}_X$-módulo coherente y $X'$ una parte cerrada de $X$) como una \emph{categoría}, cuyos morfismos $(\mathscr{F},X,X')\to(\mathscr{G},Y,Y')$ consisten en un morfismo de preesquemas $f:X\to Y$ tal que $f(X')\subset Y'$ y un $f$-morfismo $u:\mathscr{G}\to\mathscr{F}$, puede decirse que $(X_{/X'},\mathscr{F}_{/X'})$ es un \emph{funtor} en $(\mathscr{F},X,X')$ con valores en la categoría de los pares $(\mathfrak{Z},\mathscr{H})$ formados por un preesquema formal localmente noetheriano $\mathfrak{Z}$ y un $\mathscr{O}_{\mathfrak{Z}}$-módulo $\mathscr{H}$, siendo los morfismos de esta última categoría los pares formados por un morfismo $g$ de preesquemas formales y un $g$-morfismo.
\end{env}

\begin{proposition}[10.9.7]
\label{I.10.9.7}
Sean $S$, $X$ e $Y$ tres preesquemas localmente noetherianos, $g:X\to S$ y $h:Y\to S$ dos morfismos, $S'$ una parte cerrada de $S$, y $X'$ \textup{(resp. $Y'$)} una parte cerrada de $X$ \textup{(resp. $Y$)} tal que $g(X')\subset S'$ \textup{(resp. $h(Y')\subset S'$)}; sea $Z=X\times_S Y$; supongamos que $Z$ es localmente noetheriano, y sea $Z'=p^{-1}(X')\cap q^{-1}(Y')$, donde $p$ y $q$ son las proyecciones de $X\times_S Y$.
En estas condiciones, el completado $Z_{/Z'}$ se identifica con el producto $(X_{/X'})\times_{S_{/S'}}(Y_{/Y'})$ de los $S_{/S'}$-preesquemas formales, identificándose los morfismos estructurales con $\widehat g$ y $\widehat h$, y las proyecciones con $\widehat p$ y $\widehat q$.
\end{proposition}

\begin{proof}
Es inmediato que la cuestión es local para $S$, $X$ e $Y$, y se reduce así al caso en que $S=\operatorname{Spec}(A)$, $X=\operatorname{Spec}(B)$, $Y=\operatorname{Spec}(C)$, $S'=V(\mathfrak{J})$, $X'=V(\mathfrak{K})$ e $Y'=V(\mathfrak{L})$, donde $\mathfrak{J}$, $\mathfrak{K}$ y $\mathfrak{L}$ son ideales tales que $\varphi(\mathfrak{J})\subset\mathfrak{K}$ y $\psi(\mathfrak{J})\subset\mathfrak{L}$, denotando por $\varphi$ y $\psi$ los homomorfismos $A\to B$ y $A\to C$ que corresponden, respectivamente, a $g$ y $h$.
Se sabe entonces que $Z=\operatorname{Spec}(B\otimes_A C)$ y que $Z'=V(\mathfrak{M})$, donde $\mathfrak{M}$ es el ideal $\operatorname{Im}(\mathfrak{K}\otimes_A C)
+\operatorname{Im}(B\otimes_A\mathfrak{L})$.
La conclusión se deduce de \sref{I.10.7.2} y del hecho de que el producto tensorial completado $(\widehat B\otimes_{\widehat A}\widehat C)^\wedge$ (donde $\widehat A$, $\widehat B$ y $\widehat C$ son, respectivamente, los completados separados de $A$, $B$ y $C$ para las topologías $\mathfrak{J}$-, $\mathfrak{K}$- y $\mathfrak{L}$-preádicas) es el completado separado del producto tensorial $B\otimes_A C$ para la topología $\mathfrak{M}$-preádica \sref[0]{0.7.7.2}.
\end{proof}

Se observará además que, si $T$ es un $S$-preesquema localmente noetheriano, $u:T\to X$ y $v:T\to Y$ son dos $S$-morfismos, y $T'$ es una parte cerrada de $T$ tal que $u(T')\subset X'$ y $v(T')\subset Y'$, entonces la prolongación a los completados $((u,v)_S)^\wedge$ se identifica con $(\widehat u,\widehat v)_{S_{/S'}}$.

\begin{corollary}[10.9.8]
\label{I.10.9.8}
Sean $X$ e $Y$ dos $S$-preesquemas localmente noetherianos tales que $X\times_S Y$ sea localmente noetheriano; sean $S'$ una parte cerrada de $S$, y $X'$ \textup{(resp. $Y'$)} una parte cerrada de $X$ \textup{(resp. $Y$)} cuya imagen en $S$ está contenida en $S'$.
Para todo $S$-morfismo $f:X\to Y$ tal que $f(X')\subset Y'$, el morfismo grafo $\Gamma_{\widehat f}$ se identifica con la prolongación $(\Gamma_f)^\wedge$ del morfismo grafo de $f$.
\end{corollary}

\begin{corollary}[10.9.9]
\label{I.10.9.9}
Sean $X$ e $Y$ dos preesquemas localmente noetherianos, $f:X\to Y$ un morfismo, $Y'$ una parte cerrada de $Y$, y $X'=f^{-1}(Y')$.
Entonces el preesquema formal $X_{/X'}$ se identifica, mediante el diagrama conmutativo
\[
  \label{I.10.9.9.1-fr}
  \xymatrix{
    X\ar[d]_f &
    X_{/X'}\ar[l]\ar[d]^{\widehat f}\\
    Y &
    Y_{/Y'}\ar[l]
  }
\]
con el producto $X\times_Y(Y_{/Y'})$ de preesquemas formales.
\end{corollary}

\begin{proof}
Basta aplicar \sref{I.10.9.7}, sustituyendo $S$ y $S'$ por $Y$, y $X$ y $X'$ por $X$.
\end{proof}


\oldpage[I]{201}
\begin{remark}[10.9.10]
\label{I.10.9.10}

Si $X$ es la suma disjunta $X_1\amalg X_2$ \hyperref[subsection:I.3.1]{(3.1)}, y $X'$ es la unión $X'_1\cup X'_2$, donde $X'_i$ es una parte cerrada de $X_i$ ($i=1,2$), se ve inmediatamente que
$X_{/X'}={X_1}_{/X'_1}\amalg {X_2}_{/X'_2}$
\end{remark}

\subsection{Aplicación a los haces coherentes sobre los esquemas formales afines}
\label{subsection:I.10.10}

\begin{env}[10.10.1]
\label{I.10.10.1}
En todo este parágrafo, $A$ designará un \emph{anillo ádico noetheriano}, y $\mathfrak{J}$ un ideal de definición de $A$.
Sean $X=\operatorname{Spec}(A)$ y $\mathfrak{X}=\operatorname{Spf}(A)$, identificándose este último con la parte cerrada $V(\mathfrak{J})$ de $X$ \sref{I.10.1.2}.
Además, las Definiciones \sref{I.10.1.2} y \sref{I.10.8.4} muestran que el \emph{esquema formal afín $\mathfrak{X}$} es idéntico al \emph{completado $X_{/\mathfrak{X}}$} del esquema afín $X$ a lo largo de la parte cerrada $\mathfrak{X}$ de su espacio subyacente.
A todo $\mathscr{O}_X$-módulo coherente $\mathscr{F}$ corresponde, por tanto, un $\mathscr{O}_{\mathfrak{X}}$-módulo de tipo finito $\mathscr{F}_{/\mathfrak{X}}$, que es además un haz de módulos topológicos sobre el haz de anillos topológicos $\mathscr{O}_{\mathfrak{X}}$.
Pero todo $\mathscr{O}_X$-módulo coherente $\mathscr{F}$ es de la forma $\widetilde M$, donde $M$ es un $A$-módulo de tipo finito \sref{I.1.5.1}; pondremos $(\widetilde M)_{/\mathfrak{X}}=M^\Delta$.
Además, si $u:M\to N$ es un $A$-homomorfismo de $A$-módulos de tipo finito, le corresponde un homomorfismo $\widetilde u:\widetilde M\to\widetilde N$ y, por consiguiente, también un homomorfismo continuo $\widetilde u_{/\mathfrak{X}}:(\widetilde M)_{/\mathfrak{X}}\to
(\widetilde N)_{/\mathfrak{X}}$, que denotaremos por $u^\Delta$.
Es inmediato que $(v\circ u)^\Delta=v^\Delta\circ u^\Delta$; se ha definido así un \emph{funtor aditivo covariante $M^\Delta$} de la categoría de los $A$-módulos de tipo finito en la de los $\mathscr{O}_{\mathfrak{X}}$-módulos de tipo finito.
Cuando $A$ es un anillo \emph{discreto}, se tiene $M^\Delta=\widetilde M$.
\end{env}

\begin{proposition}[10.10.2]
\label{I.10.10.2}
\medskip\noindent
\begin{enumerate}
  \item[\rm{(i)}] $M^\Delta$ es un funtor exacto en $M$, y existe un isomorfismo canónico funtorial de $A$-módulos $\Gamma(\mathfrak{X},M^\Delta)\xrightarrow{\sim}M$.
  \item[\rm{(ii)}] Si $M$ y $N$ son dos $A$-módulos de tipo finito, existen isomorfismos canónicos funtoriales
    \[
      (M\otimes_A N)^\Delta
      \xrightarrow{\sim}
      M^\Delta\otimes_{\mathscr{O}_{\mathfrak{X}}}N^\Delta
      \tag{10.10.2.1}\label{I.10.10.2.1-fr}
    \]
    \[
      (\operatorname{Hom}_A(M,N))^\Delta
      \xrightarrow{\sim}
      \mathscr{H}\!om_{\mathscr{O}_{\mathfrak{X}}}
      (M^\Delta,N^\Delta)
      \tag{10.10.2.2}\label{I.10.10.2.2-fr}
    \]
  \item[\rm{(iii)}] La aplicación $u\to u^\Delta$ es un isomorfismo funtorial
    \[
      \operatorname{Hom}_A(M,N)
      \xrightarrow{\sim}
      \operatorname{Hom}_{\mathscr{O}_{\mathfrak{X}}}
      (M^\Delta,N^\Delta)
      \tag{10.10.2.3}\label{I.10.10.2.3-fr}
    \]
\end{enumerate}
\end{proposition}

\begin{proof}
La exactitud de $M^\Delta$ resulta de la exactitud de los funtores $\widetilde M$ \sref{I.1.3.5} y $\mathscr{F}_{/X'}$ \sref{I.10.8.8}.
Por definición, $\Gamma(\mathfrak{X},M^\Delta)$ es el completado separado del $A$-módulo $\Gamma(X,\widetilde M)=M$ para la topología $\mathfrak{J}$-preádica; pero, como $A$ es completo y $M$ es de tipo finito, se sabe \sref[0]{0.7.3.6} que $M$ es separado y completo, lo que termina de demostrar (i).
El isomorfismo \sref{I.10.10.2.1} (resp. \sref{I.10.10.2.2}) proviene de la composición de los isomorfismos \sref{I.1.3.12}[i] y \sref{I.10.8.10.1} (resp. \sref{I.1.3.12}[ii] y \sref{I.10.8.10.2}).
Finalmente, como $\operatorname{Hom}_A(M,N)$ es un $A$-módulo de tipo finito, se le puede aplicar (i), que identifica $\Gamma(\mathfrak{X},(\operatorname{Hom}_A(M,N))^\Delta)$ con $\operatorname{Hom}_A(M,N)$, y utilizar \sref{I.10.10.2.2}, lo que demuestra que el homomorfismo \sref{I.10.10.2.3} es un isomorfismo.
\end{proof}

De la Proposición~\sref{I.10.10.2} se deduce toda una serie de consecuencias análogas a las deducidas del Teorema~\sref{I.1.3.7} y del Corolario~\sref{I.1.3.12}, cuya formulación dejamos al lector.

\oldpage[I]{202}
Se observará que la propiedad de exactitud de $M^\Delta$, aplicada a la sucesión exacta $0\to\mathfrak{J}\to A\to A/\mathfrak{J}\to0$, muestra que el haz de ideales de $\mathscr{O}_{\mathfrak{X}}$ designado aquí por $\mathfrak{J}^\Delta$ coincide con el que se había denotado del mismo modo en \sref{I.10.3.1}, en virtud de \sref{I.10.3.2}.

\begin{proposition}[10.10.3]
\label{I.10.10.3}
Bajo las hipótesis de \sref{I.10.10.1}, $\mathscr{O}_{\mathfrak{X}}$ es un haz coherente de anillos.
\end{proposition}

\begin{proof}
Si $f\in A$, se sabe que $A_{\{f\}}$ es un anillo ádico noetheriano \sref[0]{0.7.6.11} y, como la cuestión es local, se reduce \sref{I.10.1.4} a demostrar que el núcleo de un homomorfismo $v:\mathscr{O}_{\mathfrak{X}}^n\to\mathscr{O}_{\mathfrak{X}}$ es un $\mathscr{O}_{\mathfrak{X}}$-módulo de tipo finito.
Se tiene entonces $v=u^\Delta$, donde $u$ es un $A$-homomorfismo $A^n\to A$ \sref{I.10.10.2}; como $A$ es noetheriano, el núcleo de $u$ es de tipo finito o, lo que es lo mismo, existe un homomorfismo $A^m\xrightarrow{w}A^n$ tal que la sucesión $A^m\xrightarrow{w}A^n\xrightarrow{u}A$ sea exacta.
Se concluye \sref{I.10.10.2} que la sucesión $\mathscr{O}_{\mathfrak{X}}^m\xrightarrow{w^\Delta}
\mathscr{O}_{\mathfrak{X}}^n\xrightarrow{v}
\mathscr{O}_{\mathfrak{X}}$ es exacta, lo que demuestra que el núcleo de $v$ es de tipo finito.
\end{proof}

\begin{env}[10.10.4]
\label{I.10.10.4}
Con las notaciones anteriores, pongamos $A_n=A/\mathfrak{J}^{n+1}$, y sea $X_n$ el esquema afín $\operatorname{Spec}(A_n)=(\mathfrak{X},
\mathscr{O}_{\mathfrak{X}}/\mathscr{J}^{n+1})$, siendo $\mathscr{J}=\mathfrak{J}^\Delta$ el haz de ideales de definición de $\mathscr{O}_{\mathfrak{X}}$ correspondiente al ideal $\mathfrak{J}$.
Sea $u_{mn}$ el morfismo de preesquemas $X_m\to X_n$ correspondiente al homomorfismo canónico $A_n\to A_m$ para $m\leq n$; el esquema formal $\mathfrak{X}$ es el límite inductivo de los $X_n$ para los $u_{mn}$ \sref{I.10.6.3}.
\end{env}

\begin{proposition}[10.10.5]
\label{I.10.10.5}
Bajo las hipótesis de \sref{I.10.10.1}, sea $\mathscr{F}$ un $\mathscr{O}_{\mathfrak{X}}$-módulo.
Las condiciones siguientes son equivalentes:
\begin{enumerate}
  \item[\textup{a)}] $\mathscr{F}$ es un $\mathscr{O}_{\mathfrak{X}}$-módulo coherente;
  \item[\textup{b)}] $\mathscr{F}$ es isomorfo al límite proyectivo \sref{I.10.6.6} de una sucesión $(\mathscr{F}_n)$ de $\mathscr{O}_{X_n}$-módulos coherentes tales que $u_{mn}^*(\mathscr{F}_n)=\mathscr{F}_m$; y
  \item[\textup{c)}] existe un $A$-módulo de tipo finito $M$ (determinado salvo isomorfismo canónico por la Proposición~\sref{I.10.10.2}[(i)]) tal que $\mathscr{F}$ es isomorfo a $M^\Delta$.
\end{enumerate}
\end{proposition}

\begin{proof}
Demostremos primero que \emph{b)} implica \emph{c)}.
Se tiene $\mathscr{F}_n=\widetilde{M}_n$, donde $M_n$ es un $A_n$-módulo de tipo finito, y la hipótesis implica que $M_m=M_n\otimes_{A_n}A_m$ para $m\leq n$ \sref{I.1.6.5}; los $M_n$ forman, pues, un sistema proyectivo para los dihomomorfismos canónicos $M_n\to M_m$ ($m\leq n$), y de la definición de los $A_n$ resulta inmediatamente que este sistema proyectivo satisface las condiciones de \sref[0]{0.7.2.9}; por consiguiente, su límite proyectivo $M$ es un $A$-módulo de tipo finito tal que $M_n=M\otimes_A A_n$ para todo $n$.
Se deduce que $\mathscr{F}_n$ está inducido sobre $X_n$ por $\widetilde M\otimes_{\mathscr{O}_X}
(\mathscr{O}_X/\widetilde{\mathfrak{J}}^{n+1})$ y, por tanto, $\mathscr{F}=M^\Delta$ por definición \sref{I.10.8.4}.

Recíprocamente, \emph{c)} implica \emph{b)}; en efecto, si $u_n$ es el morfismo de inmersión $X_n\to X$, entonces $u_n^*(\widetilde M)=(M\otimes_A A_n)^{\sim}$ está inducido sobre $X_n$ por $\widetilde M\otimes_{\mathscr{O}_X}
(\mathscr{O}_X/\widetilde{\mathfrak{J}}^{n+1})$, y $M^\Delta=\varprojlim u_n^*(\widetilde M)$ por definición \sref{I.10.8.4}; como $u_m=u_n\circ u_{mn}$ para $m\leq n$, los $\mathscr{F}_n=u_n^*(\widetilde M)$ satisfacen las condiciones de \emph{b)}, de donde resulta la afirmación.

Demostremos ahora que \emph{c)} implica \emph{a)}: en efecto, por definición se tiene $\mathscr{O}_{\mathfrak{X}}=A^\Delta$; como $M$ es el conúcleo de un homomorfismo $A^m\to A^n$, se deduce de la Proposición~\sref{I.10.10.2} que $M^\Delta$ es el conúcleo de un homomorfismo $\mathscr{O}_{\mathfrak{X}}^m\to\mathscr{O}_{\mathfrak{X}}^n$ y, como el haz de anillos $\mathscr{O}_{\mathfrak{X}}$ es coherente \sref{I.10.10.3}, lo mismo ocurre con $M^\Delta$ \sref[0]{0.5.3.4}.
\oldpage[I]{203}
Finalmente, \emph{a)} implica \emph{b)}.
Considerado como $\mathscr{O}_{\mathfrak{X}}$-módulo, se tiene $\mathscr{O}_{X_n}=\mathscr{O}_{\mathfrak{X}}/
\mathscr{J}^{n+1}=A_n^\Delta$; ahora bien, $\mathscr{F}_n=\mathscr{F}
\otimes_{\mathscr{O}_{\mathfrak{X}}}\mathscr{O}_{X_n}$ es un $\mathscr{O}_{\mathfrak{X}}$-módulo coherente \sref[0]{0.5.3.5} y, como también es un $\mathscr{O}_{X_n}$-módulo y $\mathscr{J}^{n+1}$ es coherente, se concluye que $\mathscr{F}_n$ es un $\mathscr{O}_{X_n}$-módulo coherente \sref[0]{0.5.3.10}; además, es inmediato que $u_{mn}^*(\mathscr{F}_n)=\mathscr{F}_m$ para $m\leq n$ (recordando que la aplicación continua $X_m\to X_n$ de los espacios subyacentes es la identidad de $\mathfrak{X}$).
El haz $\mathscr{G}=\varprojlim\mathscr{F}_n$ es, por tanto, un $\mathscr{O}_{\mathfrak{X}}$-módulo coherente, puesto que se ha visto que \emph{b)} implica \emph{a)}. Los homomorfismos canónicos $\mathscr{F}\to\mathscr{F}_n$ forman un sistema proyectivo que, al pasar al límite, da un homomorfismo canónico $w:\mathscr{F}\to\mathscr{G}$, y todo se reduce a demostrar que $w$ es biyectivo.
Como la cuestión es ahora \emph{local}, basta limitarse al caso en que $\mathscr{F}$ es el conúcleo de un homomorfismo $\mathscr{O}_{\mathfrak{X}}^p\to\mathscr{O}_{\mathfrak{X}}^q$; dado que este homomorfismo es de la forma $v^\Delta$, donde $v$ es un homomorfismo $A^p\to A^q$ \sref{I.10.10.2-fr}\footnote{\emph{[Trad.] La fuente francesa impresa tiene aquí $A^m\to A^n$; los rangos del homomorfismo precedente exigen $A^p\to A^q$.}}, se tiene que $\mathscr{F}$ es isomorfo a $M^\Delta$, donde $M=\operatorname{Coker}v$ \sref{I.10.10.2}.
Se tiene entonces, en virtud de la Proposición~\sref{I.10.10.2}, $\mathscr{F}_n=M^\Delta\otimes_{\mathscr{O}_{\mathfrak{X}}}A_n^\Delta
=(M\otimes_A A_n)^\Delta$ y, como la topología $\mathfrak{J}$-ádica sobre $M\otimes_A A_n$ es discreta, se tiene $(M\otimes_A A_n)^\Delta=(M\otimes_A A_n)^{\sim}$ (como $\mathscr{O}_{X_n}$-módulo); se ha visto antes que $M^\Delta=\varprojlim\mathscr{F}_n$, y $w$ es, pues, la identidad en este caso.
\end{proof}

\begin{corollary}[10.10.6]
\label{I.10.10.6}
Si $\mathscr{F}$ satisface la condición \textup{b)} de la Proposición~\sref{I.10.10.5}, el sistema proyectivo $(\mathscr{F}_n)$ es isomorfo al sistema de los $\mathscr{F}\otimes_{\mathscr{O}_{\mathfrak{X}}}\mathscr{O}_{X_n}$.
\end{corollary}

\begin{env}[10.10.7]
\label{I.10.10.7}
Sean ahora $A$ y $B$ dos anillos ádicos noetherianos, y $\varphi:B\to A$ un homomorfismo continuo; se designará por $\mathfrak{J}$ (resp. $\mathfrak{K}$) un ideal de definición de $A$ (resp. de $B$) tal que $\varphi(\mathfrak{K})\subset\mathfrak{J}$, y se pondrá $X=\operatorname{Spec}(A)$, $Y=\operatorname{Spec}(B)$, $\mathfrak{X}=\operatorname{Spf}(A)$ y $\mathfrak{Y}=\operatorname{Spf}(B)$.
Sea $f:X\to Y$ el morfismo de preesquemas correspondiente a $\varphi$ \sref{I.1.6.1}, y $\widehat f:\mathfrak{X}\to
\mathfrak{Y}$ su prolongación a los completados \sref{I.10.9.1}, que es también el morfismo de preesquemas formales correspondiente a $\varphi$ \sref{I.10.2.2}.
\end{env}

\begin{proposition}[10.10.8]
\label{I.10.10.8}
Para todo $B$-módulo $N$ de tipo finito, existe un isomorfismo canónico funtorial de $\mathscr{O}_{\mathfrak{X}}$-módulos
\[
  \widehat f^*(N^\Delta)
  \xrightarrow{\sim}
  (N\otimes_B A)^\Delta
\]
\end{proposition}

\begin{proof}
Designando por $i_X:\mathfrak{X}\to X$ e $i_Y:\mathfrak{Y}\to Y$ los morfismos canónicos, se tiene \sref{I.10.8.8}, salvo isomorfismos canónicos funtoriales, $N^\Delta=i_Y^*(\widetilde N)$ y
\[
  (N\otimes_B A)^\Delta
  =i_X^*((N\otimes_B A)^{\sim})
  =i_X^*(f^*(\widetilde N))
\]
\sref{I.1.6.5}; por tanto, la proposición se deduce de la conmutatividad del diagrama de \sref{I.10.9.2}.
\end{proof}

\begin{corollary}[10.10.9]
\label{I.10.10.9}
Para todo ideal $\mathfrak b$ de $B$, se tiene $\widehat f^*(\mathfrak b^\Delta)\mathscr{O}_{\mathfrak{X}}
=(\mathfrak bA)^\Delta$.
\end{corollary}

\begin{proof}
Sea $j$ la inyección canónica $\mathfrak b\to B$, a la cual corresponde la inyección canónica $j^\Delta:\mathfrak b^\Delta\to\mathscr{O}_{\mathfrak{Y}}$ de haces de $\mathscr{O}_{\mathfrak{Y}}$-módulos; por definición, $\widehat f^*(\mathfrak b^\Delta)\mathscr{O}_{\mathfrak{X}}$ es la imagen del homomorfismo $\widehat f^*(j^\Delta):\widehat f^*(\mathfrak b^\Delta)\to
\mathscr{O}_{\mathfrak{X}}=\widehat f^*(\mathscr{O}_{\mathfrak{Y}})$; pero este homomorfismo se identifica con $(j\otimes1)^\Delta:(\mathfrak b\otimes_B A)^\Delta\to
\mathscr{O}_{\mathfrak{X}}=(B\otimes_B A)^\Delta$ por la Proposición~\sref{I.10.10.8}.
Como la imagen de $j\otimes1$ es el ideal $\mathfrak bA$ de $A$, la imagen de $(j\otimes1)^\Delta$ es, por tanto, $(\mathfrak bA)^\Delta$, en virtud de la Proposición~\sref{I.10.10.2}, de donde la conclusión.
\end{proof}

\oldpage[I]{204}
\subsection{Haces coherentes sobre los preesquemas formales}
\label{subsection:I.10.11}

\begin{proposition}[10.11.1]
\label{I.10.11.1}
Si $\mathfrak{X}$ es un preesquema formal localmente noetheriano, el haz de anillos $\mathscr{O}_{\mathfrak{X}}$ es coherente, y todo haz de ideales de definición de $\mathfrak{X}$ es coherente.
\end{proposition}

\begin{proof}
La cuestión es local, por lo que se reduce al caso de un esquema formal afín noetheriano, y la proposición se deduce entonces de las Proposiciones~\sref{I.10.10.3} y~\sref{I.10.10.5}.
\end{proof}

\begin{env}[10.11.2]
\label{I.10.11.2}
Sea $\mathfrak{X}$ un preesquema formal localmente noetheriano, $\mathscr{J}$ un haz de ideales de definición de $\mathfrak{X}$, y $X_n$ el preesquema (usual) localmente noetheriano $(\mathfrak{X},\mathscr{O}_{\mathfrak{X}}/\mathscr{J}^{n+1})$, de modo que $\mathfrak{X}$ es el \emph{límite inductivo} de la sucesión $(X_n)$ para los morfismos canónicos $u_{mn}:X_m\to X_n$ \sref{I.10.6.3}.
Con estas notaciones:
\end{env}

\begin{theorem}[10.11.3]
\label{I.10.11.3}
Para que un $\mathscr{O}_{\mathfrak{X}}$-módulo $\mathscr{F}$ sea coherente, es necesario y suficiente que sea isomorfo al límite proyectivo de una sucesión $(\mathscr{F}_n)$, donde los $\mathscr{F}_n$ son $\mathscr{O}_{X_n}$-módulos coherentes tales que $u_{mn}^*(\mathscr{F}_n)=\mathscr{F}_m$ para $m\leq n$ \sref{I.10.6.6}.
El sistema proyectivo $(\mathscr{F}_n)$ es entonces isomorfo al sistema de los $u_n^*(\mathscr{F})=\mathscr{F}
\otimes_{\mathscr{O}_{\mathfrak{X}}}\mathscr{O}_{X_n}$, donde $u_n$ es el morfismo canónico $X_n\to\mathfrak{X}$.
\end{theorem}

\begin{proof}
La cuestión es local, por lo que se reduce al caso en que $\mathfrak{X}$ es un esquema formal afín noetheriano, y el teorema es entonces consecuencia de la Proposición~\sref{I.10.10.5} y del Corolario~\sref{I.10.10.6}.
\end{proof}

Se puede decir, por tanto, que \emph{el dato de un $\mathscr{O}_{\mathfrak{X}}$-módulo coherente equivale al de un sistema proyectivo $(\mathscr{F}_n)$ de $\mathscr{O}_{X_n}$-módulos coherentes tales que $u_{mn}^*(\mathscr{F}_n)=\mathscr{F}_m$ para $m\leq n$}.

\begin{corollary}[10.11.4]
\label{I.10.11.4}
Si $\mathscr{F}$ y $\mathscr{G}$ son dos $\mathscr{O}_{\mathfrak{X}}$-módulos coherentes, se puede definir \textup{(con las notaciones del Teorema~\sref{I.10.11.3})} un isomorfismo canónico funtorial
\[
  \label{I.10.11.4.1-fr}
  \operatorname{Hom}_{\mathscr{O}_{\mathfrak{X}}}
  (\mathscr{F},\mathscr{G})
  \xrightarrow{\sim}
  \varprojlim_n
  \operatorname{Hom}_{\mathscr{O}_{X_n}}
  (\mathscr{F}_n,\mathscr{G}_n)
  \tag{10.11.4.1}
\]
\end{corollary}

\begin{proof}
El límite proyectivo del segundo miembro se entiende con respecto a las aplicaciones $\theta_n\to u_{mn}^*(\theta_n)$ ($m\leq n$) de $\operatorname{Hom}_{\mathscr{O}_{X_n}}
(\mathscr{F}_n,\mathscr{G}_n)$ en $\operatorname{Hom}_{\mathscr{O}_{X_m}}
(\mathscr{F}_m,\mathscr{G}_m)$.
El homomorfismo \sref{I.10.11.4.1} hace corresponder a un elemento $\theta\in\operatorname{Hom}_{\mathscr{O}_{\mathfrak{X}}}
(\mathscr{F},\mathscr{G})$ la sucesión $(u_n^*(\theta))$; se define un homomorfismo inverso al precedente haciendo corresponder al sistema proyectivo $(\theta_n)\in\varprojlim_n
\operatorname{Hom}_{\mathscr{O}_{X_n}}
(\mathscr{F}_n,\mathscr{G}_n)$ su límite proyectivo en $\operatorname{Hom}_{\mathscr{O}_{\mathfrak{X}}}
(\mathscr{F},\mathscr{G})$, teniendo en cuenta el Teorema~\sref{I.10.11.3}.
\end{proof}

\begin{corollary}[10.11.5]
\label{I.10.11.5}
Para que un homomorfismo $\theta:\mathscr{F}\to\mathscr{G}$ sea sobreyectivo, es necesario y suficiente que lo sea el homomorfismo correspondiente $\theta_0=u_0^*(\theta):\mathscr{F}_0\to\mathscr{G}_0$.
\end{corollary}

\begin{proof}
La cuestión es local, por lo que se reduce al caso en que $\mathfrak{X}=\operatorname{Spf}(A)$, siendo $A$ un anillo ádico noetheriano, $\mathscr{F}=M^\Delta$, $\mathscr{G}=N^\Delta$ y $\theta=u^\Delta$, donde $M$ y $N$ son $A$-módulos de tipo finito y $u$ un homomorfismo $M\to N$; se tiene además $\theta_0=\widetilde{u}_0$, donde $u_0$ es el homomorfismo $u\otimes1:M\otimes_A A/\mathfrak{J}\to
N\otimes_A A/\mathfrak{J}$; la conclusión resulta de que $\theta$ y $u$ (resp. $\theta_0$ y $u_0$) son simultáneamente sobreyectivos (\sref{I.1.3.9} y \sref{I.10.10.2}) y de que $u$ y $u_0$ son simultáneamente sobreyectivos \sref[0]{0.7.1.14}.
\end{proof}

\begin{env}[10.11.6]
\label{I.10.11.6}
El Teorema~\sref{I.10.11.3} muestra que se puede considerar todo $\mathscr{O}_{\mathfrak{X}}$-módulo coherente $\mathscr{F}$ como un $\mathscr{O}_{\mathfrak{X}}$-\emph{módulo topológico}, considerándolo como límite proyectivo de los haces de grupos \emph{pseudodiscretos} $\mathscr{F}_n$ \sref[0]{0.3.8.1}.
Se deduce entonces del Corolario~\sref{I.10.11.4} que todo homomorfismo $u:\mathscr{F}\to\mathscr{G}$ de $\mathscr{O}_{\mathfrak{X}}$-módulos coherentes es automáticamente \emph{continuo}
\oldpage[I]{205}
\sref[0]{0.3.8.2}.
Además, si $\mathscr{H}$ es un $\mathscr{O}_{\mathfrak{X}}$-submódulo coherente de un $\mathscr{O}_{\mathfrak{X}}$-módulo coherente $\mathscr{F}$, entonces, para todo abierto $U\subset\mathfrak{X}$, $\Gamma(U,\mathscr{H})$ es un subgrupo \emph{cerrado} del grupo topológico $\Gamma(U,\mathscr{F})$, pues, como el funtor $\Gamma$ es exacto por la izquierda, $\Gamma(U,\mathscr{H})$ es el núcleo del homomorfismo $\Gamma(U,\mathscr{F})\to\Gamma(U,\mathscr{F}/\mathscr{H})$, que es \emph{continuo} por lo anterior, ya que $\mathscr{F}/\mathscr{H}$ es coherente \sref[0]{0.5.3.4}; la afirmación resulta de que $\Gamma(U,\mathscr{F}/\mathscr{H})$ es un grupo topológico separado.
\end{env}

\begin{proposition}[10.11.7]
\label{I.10.11.7}
Sean $\mathscr{F}$ y $\mathscr{G}$ dos $\mathscr{O}_{\mathfrak{X}}$-módulos coherentes.
Se pueden definir \textup{(con las notaciones del Teorema~\sref{I.10.11.3})} isomorfismos canónicos funtoriales de $\mathscr{O}_{\mathfrak{X}}$-módulos topológicos \sref{I.10.11.6}
\[
  \label{I.10.11.7.1-fr}
  \mathscr{F}\otimes_{\mathscr{O}_{\mathfrak{X}}}\mathscr{G}
  \xrightarrow{\sim}
  \varprojlim_n
  (\mathscr{F}_n\otimes_{\mathscr{O}_{X_n}}\mathscr{G}_n)
  \tag{10.11.7.1}
\]
\[
  \label{I.10.11.7.2-fr}
  \mathscr{H}\!om_{\mathscr{O}_{\mathfrak{X}}}
  (\mathscr{F},\mathscr{G})
  \xrightarrow{\sim}
  \varprojlim_n
  \mathscr{H}\!om_{\mathscr{O}_{X_n}}(\mathscr{F}_n,\mathscr{G}_n)
  \tag{10.11.7.2}
\]
\end{proposition}

\begin{proof}
La existencia del isomorfismo \sref{I.10.11.7.1} resulta de la fórmula
\[
  \mathscr{F}_n\otimes_{\mathscr{O}_{X_n}}\mathscr{G}_n
  =
  (\mathscr{F}\otimes_{\mathscr{O}_{\mathfrak{X}}}
  \mathscr{O}_{X_n})
  \otimes_{\mathscr{O}_{X_n}}
  (\mathscr{G}\otimes_{\mathscr{O}_{\mathfrak{X}}}
  \mathscr{O}_{X_n})
  =
  (\mathscr{F}\otimes_{\mathscr{O}_{\mathfrak{X}}}\mathscr{G})
  \otimes_{\mathscr{O}_{\mathfrak{X}}}\mathscr{O}_{X_n}
\]
y del Teorema~\sref{I.10.11.3}.
El isomorfismo \sref{I.10.11.7.2}, donde ambos miembros se consideran como haces de módulos sin topología, resulta de la definición de las secciones de $\mathscr{H}\!om_{\mathscr{O}_{\mathfrak{X}}}
(\mathscr{F},\mathscr{G})$ y $\mathscr{H}\!om_{\mathscr{O}_{X_n}}(\mathscr{F}_n,\mathscr{G}_n)$, y de la existencia del isomorfismo \sref{I.10.11.4.1}, aplicado al preesquema inducido sobre un abierto formal afín noetheriano arbitrario de $\mathfrak{X}$.
Queda por probar que el isomorfismo \sref{I.10.11.7.2} es bicontinuo sobre un conjunto cuasicompacto, y se reduce así al caso en que $\mathfrak{X}=\operatorname{Spf}(A)$, siendo $A$ un anillo ádico noetheriano, y por tanto, en virtud de \sref{I.10.10.5}, a aquel en que $\mathscr{F}=M^\Delta$ y $\mathscr{G}=N^\Delta$, siendo $M$ y $N$ dos $A$-módulos de tipo finito; teniendo en cuenta \sref{I.10.10.2.1}, \sref{I.10.10.2.3} y el Corolario~\sref{I.1.3.12}[ii], se reduce a mostrar que el isomorfismo canónico $\operatorname{Hom}_A(M,N)\xrightarrow{\sim}
\varprojlim_n\operatorname{Hom}_{A_n}(M_n,N_n)$ (con $M_n=M\otimes_A A_n$ y $N_n=N\otimes_A A_n$) es continuo, lo cual ya se probó en \sref[0]{0.7.8.2}.
\end{proof}

\begin{env}[10.11.8]
\label{I.10.11.8}
Como $\operatorname{Hom}_{\mathscr{O}_{\mathfrak{X}}}
(\mathscr{F},\mathscr{G})$ es el grupo de secciones del haz de grupos topológicos $\mathscr{H}\!om_{\mathscr{O}_{\mathfrak{X}}}
(\mathscr{F},\mathscr{G})$, está provisto de una topología de grupo.
Si $\mathfrak{X}$ es \emph{noetheriano}, se deduce de \sref{I.10.11.7.2} que los subgrupos $\operatorname{Hom}_{\mathscr{O}_{\mathfrak{X}}}
(\mathscr{F},\mathscr{J}^n\mathscr{G})$ (para $n$ arbitrario) forman un sistema fundamental de vecindades de $0$ en este grupo.
\end{env}

\begin{proposition}[10.11.9]
\label{I.10.11.9}
Sean $\mathfrak{X}$ un preesquema formal noetheriano, y $\mathscr{F}$ y $\mathscr{G}$ dos $\mathscr{O}_{\mathfrak{X}}$-módulos coherentes.
En el grupo topológico $\operatorname{Hom}_{\mathscr{O}_{\mathfrak{X}}}
(\mathscr{F},\mathscr{G})$, los homomorfismos sobreyectivos \textup{(resp.} inyectivos, biyectivos) forman una parte abierta.
\end{proposition}

\begin{proof}
En virtud del Corolario~\sref{I.10.11.5}, el conjunto de los homomorfismos sobreyectivos de $\operatorname{Hom}_{\mathscr{O}_{\mathfrak{X}}}
(\mathscr{F},\mathscr{G})$ es la imagen inversa, por la aplicación continua $\operatorname{Hom}_{\mathscr{O}_{\mathfrak{X}}}
(\mathscr{F},\mathscr{G})\to
\operatorname{Hom}_{\mathscr{O}_{X_0}}(\mathscr{F}_0,\mathscr{G}_0)$, de una parte del grupo discreto $\operatorname{Hom}_{\mathscr{O}_{X_0}}(\mathscr{F}_0,\mathscr{G}_0)$, de donde resulta la primera afirmación.
Para demostrar la segunda, recubramos $\mathfrak{X}$ por un número finito de abiertos formales afines noetherianos $U_i$.
Para que $\theta\in\operatorname{Hom}_{\mathscr{O}_{\mathfrak{X}}}
(\mathscr{F},\mathscr{G})$ sea inyectivo, es necesario y suficiente que todas sus imágenes por las aplicaciones (continuas) de restricción $\operatorname{Hom}_{\mathscr{O}_{\mathfrak{X}}}
(\mathscr{F},\mathscr{G})\to
\operatorname{Hom}_{\mathscr{O}_{\mathfrak{X}}|U_i}
(\mathscr{F}|U_i,\mathscr{G}|U_i)$ sean inyectivas; se reduce así al caso afín, y esto ya se probó en \sref[0]{0.7.8.3}.
\end{proof}

\oldpage[I]{206}
\subsection{Morfismos ádicos de preesquemas formales}
\label{subsection:I.10.12}

\begin{env}[10.12.1]
\label{I.10.12.1}
Sean $\mathfrak{X}$ y $\mathfrak{S}$ dos preesquemas formales \emph{localmente noetherianos};
diremos que un morfismo $f:\mathfrak{X}\to\mathfrak{S}$ es \emph{ádico} si existe un ideal de definición $\mathscr{J}$ de $\mathfrak{S}$ tal que $\mathscr{K}=f^*(\mathscr{J})\mathscr{O}_{\mathfrak{X}}$ sea un ideal de definición de $\mathfrak{X}$;
se dice también entonces que $\mathfrak{X}$ es un \emph{$\mathfrak{S}$-preesquema ádico} (para $f$).
Cuando así ocurre, para \emph{todo} ideal de definición $\mathscr{J}_1$ de $\mathfrak{S}$, $\mathscr{K}_1=f^*(\mathscr{J}_1)\mathscr{O}_{\mathfrak{X}}$ es un ideal de definición de $\mathfrak{X}$.
En efecto, como la cuestión es local, se puede suponer que $\mathfrak{X}$ y $\mathfrak{S}$ son afines y noetherianos;
existe entonces un entero $n$ tal que $\mathscr{J}^n\subset\mathscr{J}_1$ y $\mathscr{J}_1^n\subset\mathscr{J}$ (\sref{I.10.3.6} y \sref[0]{0.7.1.4}), de donde $\mathscr{K}^n\subset\mathscr{K}_1$ y $\mathscr{K}_1^n\subset\mathscr{K}$.
La primera de estas relaciones prueba que $\mathscr{K}_1=\mathfrak{K}_1^\Delta$, donde $\mathfrak{K}_1$ es un ideal abierto de $A=\Gamma(\mathfrak{X},\mathscr{O}_{\mathfrak{X}})$, y la segunda prueba que $\mathfrak{K}_1$ es un ideal de definición de $A$ \sref[0]{0.7.1.4}, de donde resulta nuestra afirmación.
\end{env}

Se deduce inmediatamente de lo anterior que, si $\mathfrak{X}$ y $\mathfrak{Y}$ son dos $\mathfrak{S}$-preesquemas ádicos, \emph{todo $\mathfrak{S}$-morfismo $u:\mathfrak{X}\to\mathfrak{Y}$ es ádico}:
en efecto, si $f:\mathfrak{X}\to\mathfrak{S}$ y $g:\mathfrak{Y}\to\mathfrak{S}$ son los morfismos estructurales, y $\mathscr{J}$ es un ideal de definición de $\mathfrak{S}$, se tiene $f=g\circ u$, y por tanto $u^*(g^*(\mathscr{J})\mathscr{O}_{\mathfrak{Y}})
\mathscr{O}_{\mathfrak{X}}=f^*(\mathscr{J})
\mathscr{O}_{\mathfrak{X}}$ es un ideal de definición de $\mathfrak{X}$ y, por hipótesis, $g^*(\mathscr{J})\mathscr{O}_{\mathfrak{Y}}$ es un ideal de definición de $\mathfrak{Y}$.

\begin{env}[10.12.2]
\label{I.10.12.2}
En lo que sigue, supondremos fijados un preesquema formal localmente noetheriano $\mathfrak{S}$ y un ideal de definición $\mathscr{J}$ de $\mathfrak{S}$;
pondremos $S_n=(\mathfrak{S},\mathscr{O}_{\mathfrak{S}}/\mathscr{J}^{n+1})$.
Los $\mathfrak{S}$-preesquemas ádicos (localmente noetherianos) forman evidentemente una \emph{categoría}.
Diremos que un sistema inductivo $(X_n)$ de $S_n$-preesquemas (usuales) localmente noetherianos es un \emph{$(S_n)$-sistema inductivo ádico} si los morfismos estructurales $f_n:X_n\to S_n$ son tales que, para $m\leq n$, los diagramas
\[
  \label{I.10.12.2.1-fr}
  \xymatrix{
    X_n \ar[d]_{f_n}
    & X_m \ar[l] \ar[d]^{f_m}\\
    S_n
    & S_m \ar[l]
  }
  \tag{10.12.2.1}
\]
sean conmutativos e \emph{identifiquen $X_m$ con el producto $X_n\times_{S_n}S_m=(X_n)_{(S_m)}$}.
Los sistemas inductivos ádicos forman una \emph{categoría}:
basta, en efecto, definir un morfismo $(X_n)\to(Y_n)$ de tales sistemas como un \emph{sistema inductivo de $S_n$-morfismos $u_n:X_n\to Y_n$} tal que $u_m$ se identifique con $(u_n)_{(S_m)}$ para $m\leq n$.
Sentado esto:
\end{env}

\begin{theorem}[10.12.3]
\label{I.10.12.3}
Hay una equivalencia canónica entre la categoría de los $\mathfrak{S}$-preesquemas ádicos y la categoría de los $(S_n)$-sistemas inductivos ádicos.
\end{theorem}

La equivalencia en cuestión se obtiene del modo siguiente:
si $\mathfrak{X}$ es un $\mathfrak{S}$-preesquema ádico y $f:\mathfrak{X}\to\mathfrak{S}$ es el morfismo estructural, entonces $\mathscr{K}=f^*(\mathscr{J})\mathscr{O}_{\mathfrak{X}}$ es un ideal de definición de $\mathfrak{X}$, y se hace corresponder a $\mathfrak{X}$ el sistema inductivo de los $X_n=(\mathfrak{X},\mathscr{O}_{\mathfrak{X}}/\mathscr{K}^{n+1})$, siendo $f_n:X_n\to S_n$ el morfismo estructural correspondiente a $f$ \sref{I.10.5.6}.
Mostremos primero que $(X_n)$ es un \emph{sistema inductivo ádico}:
si $f=(\psi,\theta)$, se tiene $\psi^*(\mathscr{J})\mathscr{O}_{\mathfrak{X}}=\mathscr{K}$, luego $\psi^*(\mathscr{J}^n)\mathscr{O}_{\mathfrak{X}}=\mathscr{K}^n$ para todo $n$, y (por la exactitud del funtor $\psi^*$) $\mathscr{K}^{m+1}/\mathscr{K}^{n+1}
=\psi^*(\mathscr{J}^{m+1}/\mathscr{J}^{n+1})
(\mathscr{O}_{\mathfrak{X}}/\mathscr{K}^{n+1})$ para $m\leq n$;
nuestra conclusión resulta, pues, de \sref{I.4.4.5}.
Además, se verifica inmediatamente que a un $\mathfrak{S}$-morfismo $u:\mathfrak{X}\to\mathfrak{Y}$ de $\mathfrak{S}$-preesquemas ádicos corresponde (con notaciones evidentes)
\oldpage[I]{207}
un sistema inductivo de $S_n$-morfismos $u_n:X_n\to Y_n$ tal que $u_m$ se identifica con $(u_n)_{(S_m)}$ para $m\leq n$.

El hecho de que así se haya definido efectivamente una equivalencia resultará de la proposición más precisa siguiente:

\begin{proposition}[10.12.3.1]
\label{I.10.12.3.1}
Sea $(X_n)$ un sistema inductivo de $S_n$-preesquemas;
se supone que los morfismos estructurales $f_n:X_n\to S_n$ son tales que los diagramas de \sref{I.10.12.2.1} sean conmutativos e identifiquen $X_m$ con $X_n\times_{S_n}S_m$ para $m\leq n$.
Entonces el sistema inductivo $(X_n)$ satisface las condiciones \textup{b)} y \textup{c)} de \sref{I.10.6.3};
sea $\mathfrak{X}$ su límite inductivo y sea $f:\mathfrak{X}\to\mathfrak{S}$ el morfismo límite inductivo del sistema inductivo $(f_n)$.
Entonces, si $X_0$ es localmente noetheriano, $\mathfrak{X}$ es localmente noetheriano y $f$ es un morfismo ádico.
\end{proposition}

\begin{proof}
Como el haz de ideales de $\mathscr{O}_{S_n}$ que define el subpreesquema $S_m$ de $S_n$ es nilpotente, lo mismo ocurre, en virtud de \sref{I.4.4.5}, con el haz de ideales de $\mathscr{O}_{X_n}$ que define el subpreesquema $X_m$ de $X_n$; por tanto, se satisfacen las condiciones de \sref{I.10.6.3}.
La cuestión es, pues, local sobre $\mathfrak{X}$ y $\mathfrak{S}$, y se puede suponer que $\mathfrak{S}=\operatorname{Spf}(A)$, $\mathscr{J}=\mathfrak{J}^\Delta$ (siendo $A$ un anillo $\mathfrak{J}$-ádico noetheriano) y $X_n=\operatorname{Spec}(B_n)$;
si $A_n=A/\mathfrak{J}^{n+1}$, la hipótesis implica que $B_0$ es noetheriano y, si se pone $\mathfrak{J}_n=\mathfrak{J}/\mathfrak{J}^{n+1}$, entonces $B_m=B_n/\mathfrak{J}_n^{m+1}B_n$.
El núcleo de $B_n\to B_0$ es, pues, $\mathfrak{K}_n=\mathfrak{J}_nB_n$, y el núcleo de $B_n\to B_m$ es $\mathfrak{K}_n^{m+1}$ para $m\leq n$;
además, como $A_1$ es noetheriano, $\mathfrak{J}_1$ es de tipo finito sobre $A_1$, luego $\mathfrak{K}_1=\mathfrak{K}_1/\mathfrak{K}_1^2$ es de tipo finito sobre $B_1$ y, \emph{a fortiori}, de tipo finito sobre $B_0=B_1/\mathfrak{K}_1$;
el hecho de que $\mathfrak{X}$ sea noetheriano resulta entonces de \sref{I.10.6.4};
si $B=\varprojlim B_n$, se tiene $\mathfrak{X}=\operatorname{Spf}(B)$ y, si $\mathfrak{K}$ es el núcleo de $B\to B_0$, entonces $B_n=B/\mathfrak{K}^{n+1}$.
Si $\rho_n:A/\mathfrak{J}^{n+1}\to B/\mathfrak{K}^{n+1}$ es el homomorfismo correspondiente a $f_n$, se tiene
\[
  \mathfrak{K}/\mathfrak{K}^{n+1}
  =(B/\mathfrak{K}^{n+1})\rho_n(\mathfrak{J}/\mathfrak{J}^{n+1})
\]
Como el homomorfismo $\rho:A\to B$ correspondiente a $f$ es igual a $\varprojlim\rho_n$, el ideal $\mathfrak{J}B$ de $B$ es denso en $\mathfrak{K}$; como todo ideal de $B$ es cerrado \sref[0]{0.7.3.5}, se tiene también $\mathfrak{K}=\mathfrak{J}B$.
Si $\mathscr{K}=\mathfrak{K}^\Delta$, la igualdad $f^*(\mathscr{J})\mathscr{O}_{\mathfrak{X}}=\mathscr{K}$ resulta entonces de \sref{I.10.10.9} y termina la demostración.
\end{proof}

\begin{env}[10.12.3.2]
\label{I.10.12.3.2}
La equivalencia precedente proporciona, para dos $\mathfrak{S}$-preesquemas ádicos $\mathfrak{X}$ y $\mathfrak{Y}$, una \emph{biyección canónica}
\[
  \operatorname{Hom}_{\mathfrak{S}}(\mathfrak{X},\mathfrak{Y})
  \xrightarrow{\sim}
  \varprojlim_n\operatorname{Hom}_{S_n}(X_n,Y_n)
\]
donde el límite proyectivo se toma con respecto a las aplicaciones $u_n\to(u_n)_{(S_m)}$ para $m\leq n$.
\end{env}

\subsection{Morfismos de tipo finito}
\label{subsection:I.10.13}

\begin{proposition}[10.13.1]
\label{I.10.13.1}
Sea $\mathfrak{Y}$ un preesquema formal localmente noetheriano, $\mathscr{K}$ un ideal de definición de $\mathfrak{Y}$, y sea $f:\mathfrak{X}\to\mathfrak{Y}$ un morfismo de preesquemas formales.
Las condiciones siguientes son equivalentes.
\begin{enumerate}
  \item[\textup{a)}] $\mathfrak{X}$ es localmente noetheriano, $f$ es un morfismo ádico \sref{I.10.12.1} y, si se pone $\mathscr{J}=f^*(\mathscr{K})\mathscr{O}_{\mathfrak{X}}$, el morfismo $f_0:(\mathfrak{X},\mathscr{O}_{\mathfrak{X}}/\mathscr{J})
    \to(\mathfrak{Y},\mathscr{O}_{\mathfrak{Y}}/\mathscr{K})$ deducido de $f$ es de tipo finito.
  \item[\textup{b)}] $\mathfrak{X}$ es localmente noetheriano y es el límite inductivo de un $(Y_n)$-sistema inductivo ádico $(X_n)$ tal que el morfismo $X_0\to Y_0$ sea de tipo finito.
\oldpage[I]{208}
  \item[\textup{c)}] Todo punto de $\mathfrak{Y}$ posee una vecindad abierta formal afín noetheriana $V$ que satisface la propiedad siguiente:
\begin{enumerate}
  \item[\textup{(Q)}] $f^{-1}(V)$ es la unión de una familia finita de abiertos formales afines noetherianos $U_i$ tales que el anillo ádico noetheriano $\Gamma(U_i,\mathscr{O}_{\mathfrak{X}})$ sea topológicamente isomorfo al cociente de un álgebra de series formales restringidas \sref[0]{0.7.5.1} sobre $\Gamma(V,\mathscr{O}_{\mathfrak{Y}})$ por un ideal (necesariamente cerrado).
\end{enumerate}
\end{enumerate}
\end{proposition}

\begin{proof}
Es inmediato que \emph{a)} implica \emph{b)}, en virtud de \sref{I.10.12.3}.
Para demostrar que \emph{b)} implica \emph{c)}, se puede, puesto que la cuestión es local sobre $\mathfrak{Y}$, suponer que $\mathfrak{Y}=\operatorname{Spf}(B)$, donde $B$ es noetheriano y ádico;
sea $\mathscr{K}=\mathfrak{K}^\Delta$, siendo $\mathfrak{K}$ un ideal de definición de $B$.
Como, por hipótesis, $X_0$ es de tipo finito sobre $Y_0$, $X_0$ es la unión de un número finito de abiertos afines $U_i$ tales que el anillo $A_{i0}$ del esquema afín inducido por $X_0$ sobre $U_i$ sea un álgebra de tipo finito sobre el anillo $B/\mathfrak{K}$ de $Y_0$ \sref{I.6.3.2}.
En virtud de \sref{I.5.1.9}, $U_i$ es también un abierto afín de cada uno de los preesquemas noetherianos $X_n$ y, si $A_{in}$ es el anillo del esquema afín inducido por $X_n$ sobre $U_i$, la hipótesis \emph{b)} implica, para $m\leq n$, que $A_{im}$ es isomorfo a $A_{in}/\mathfrak{K}^{m+1}A_{in}$.
Por consiguiente, el preesquema formal inducido por $\mathfrak{X}$ sobre $U_i$ es isomorfo a $\operatorname{Spf}(A_i)$, donde $A_i=\varprojlim_n A_{in}$ \sref{I.10.6.4};
$A_i$ es un anillo $\mathfrak{K}A_i$-ádico, y $A_i/\mathfrak{K}A_i$, que es isomorfo a $A_{i0}$, es un álgebra de tipo finito sobre $B/\mathfrak{K}$.
Se concluye, pues, de \sref[0]{0.7.5.5} que $A_i$ es topológicamente isomorfo a un cociente de un álgebra de series formales restringidas sobre $B$ (por un ideal necesariamente cerrado, puesto que tal álgebra es noetheriana \sref[0]{0.7.5.4}).

Para demostrar que \emph{c)} implica \emph{a)}, se puede limitar al caso en que $\mathfrak{X}=\operatorname{Spf}(A)$ es también afín, siendo $A$ un anillo ádico noetheriano isomorfo a un cociente de un álgebra de series formales restringidas sobre $B$ por un ideal cerrado.
Entonces, por \sref[0]{0.7.5.5}, $A/\mathfrak{K}A$ es un álgebra de tipo finito sobre $B/\mathfrak{K}$ y $\mathfrak{K}A=\mathfrak{J}$ es un ideal de definición de $A$; por tanto, en virtud de \sref{I.10.10.9}, se satisfacen las condiciones de \emph{a)}.
\end{proof}

Se observará que, si se satisfacen las condiciones de la Proposición~\sref{I.10.13.1}, la propiedad~\emph{a)} es válida para \emph{todo} ideal de definición $\mathscr{K}$ de $\mathfrak{Y}$ (en virtud de \emph{c)}) y, por consiguiente, en la propiedad~\emph{b)}, \emph{todos} los $f_n$ son morfismos de tipo finito.

\begin{corollary}[10.13.2]
\label{I.10.13.2}
Si se satisfacen las condiciones de \sref{I.10.13.1}, todo abierto formal afín noetheriano $V$ de $\mathfrak{Y}$ posee la propiedad~\textup{(Q)} y, si $\mathfrak{Y}$ es noetheriano, también lo es $\mathfrak{X}$.
\end{corollary}

\begin{proof}
Esto resulta inmediatamente de \sref{I.10.13.1} y \sref{I.6.3.2}.
\end{proof}

\begin{definition}[10.13.3]
\label{I.10.13.3}
{\itshape Cuando se satisfacen las propiedades equivalentes \textup{a), b) y c)} de \sref{I.10.13.1}, se dice que el morfismo $f$ es de tipo finito, o que $\mathfrak{X}$ es un $\mathfrak{Y}$-preesquema formal de tipo finito, o un preesquema formal de tipo finito sobre $\mathfrak{Y}$.}
\end{definition}

\begin{corollary}[10.13.4]
\label{I.10.13.4}
Sean $\mathfrak{X}=\operatorname{Spf}(A)$ y $\mathfrak{Y}=\operatorname{Spf}(B)$ dos esquemas formales afines noetherianos;
para que $\mathfrak{X}$ sea de tipo finito sobre $\mathfrak{Y}$, es necesario y suficiente que el anillo ádico noetheriano $A$ sea isomorfo al cociente de un álgebra de series formales restringidas sobre $B$ por un ideal cerrado.
\end{corollary}

\begin{proof}
Con las notaciones de \sref{I.10.13.1}, si $\mathfrak{X}$ es de tipo finito sobre $\mathfrak{Y}$, entonces $A/\mathfrak{K}A$ es un $(B/\mathfrak{K})$-álgebra de tipo finito en virtud de \sref{I.6.3.3}, y $\mathfrak{K}A$ es un ideal de definición de $A$ \sref{I.10.10.9}.
Se concluye entonces por \sref[0]{0.7.5.5}.
\end{proof}

\begin{proposition}[10.13.5]
\label{I.10.13.5}
\medskip\noindent
\begin{enumerate}
  \item[\rm{(i)}] El compuesto de dos morfismos de preesquemas formales de tipo finito es de tipo finito.
\oldpage[I]{209}
  \item[\rm{(ii)}] Sean $\mathfrak{X}$, $\mathfrak{S}$ y $\mathfrak{S}'$ tres preesquemas formales localmente noetherianos \textup{(resp.} noetherianos\textup{)}, y sean $f:\mathfrak{X}\to\mathfrak{S}$ y $g:\mathfrak{S}'\to\mathfrak{S}$ dos morfismos.
    Si $f$ es de tipo finito, $\mathfrak{X}\times_{\mathfrak{S}}\mathfrak{S}'$ es localmente noetheriano \textup{(resp.} noetheriano\textup{)} y es de tipo finito sobre $\mathfrak{S}'$.
  \item[\rm{(iii)}] Sea $\mathfrak{S}$ un preesquema formal localmente noetheriano y sean $\mathfrak{X}'$ y $\mathfrak{Y}'$ dos $\mathfrak{S}$-preesquemas formales localmente noetherianos tales que $\mathfrak{X}'\times_{\mathfrak{S}}\mathfrak{Y}'$ sea localmente noetheriano.
    Si $\mathfrak{X}$ y $\mathfrak{Y}$ son $\mathfrak{S}$-preesquemas formales localmente noetherianos, y $f:\mathfrak{X}\to\mathfrak{X}'$ y $g:\mathfrak{Y}\to\mathfrak{Y}'$ son dos $\mathfrak{S}$-morfismos de tipo finito, $\mathfrak{X}\times_{\mathfrak{S}}\mathfrak{Y}$ es localmente noetheriano y $f\times_{\mathfrak{S}}g$ es un $\mathfrak{S}$-morfismo de tipo finito.
\end{enumerate}
\end{proposition}

\begin{proof}
Por el razonamiento formal de \sref{I.3.5.1}, (iii) se deduce de (i) y (ii); basta, pues, probar (i) y (ii).

Sean $\mathfrak{X}$, $\mathfrak{Y}$ y $\mathfrak{Z}$ tres preesquemas formales localmente noetherianos, y sean $f:\mathfrak{X}\to\mathfrak{Y}$ y $g:\mathfrak{Y}\to\mathfrak{Z}$ dos morfismos de tipo finito.
Si $\mathscr{L}$ es un ideal de definición de $\mathfrak{Z}$, $\mathscr{K}=g^*(\mathscr{L})\mathscr{O}_{\mathfrak{Y}}$ es un ideal de definición de $\mathfrak{Y}$ y $\mathscr{J}=f^*(g^*(\mathscr{L}))\mathscr{O}_{\mathfrak{X}}$ es un ideal de definición de $\mathfrak{X}$.
Pongamos $X_0=(\mathfrak{X},\mathscr{O}_{\mathfrak{X}}/\mathscr{J})$, $Y_0=(\mathfrak{Y},\mathscr{O}_{\mathfrak{Y}}/\mathscr{K})$ y $Z_0=(\mathfrak{Z},\mathscr{O}_{\mathfrak{Z}}/\mathscr{L})$, y sean $f_0:X_0\to Y_0$ y $g_0:Y_0\to Z_0$ los morfismos correspondientes a $f$ y $g$, respectivamente.
Como, por hipótesis, $f_0$ y $g_0$ son de tipo finito, también lo es $g_0\circ f_0$ \sref{I.6.3.4}, que corresponde a $g\circ f$;
por tanto, $g\circ f$ es de tipo finito en virtud de \sref{I.10.13.1}.

En las condiciones de (ii), $\mathfrak{S}$ (resp. $\mathfrak{X}$, $\mathfrak{S}'$) es el límite inductivo de una sucesión $(S_n)$ (resp. $(X_n)$, $(S'_n)$) de preesquemas localmente noetherianos, y se puede suponer \sref{I.10.13.1} que $X_m=X_n\times_{S_n}S_m$ para $m\leq n$.
El preesquema formal $\mathfrak{X}\times_{\mathfrak{S}}\mathfrak{S}'$ es entonces el límite inductivo de los preesquemas $X_n\times_{S_n}S'_n$ \sref{I.10.7.4}, y se tiene
\[
  X_m\times_{S_m}S'_m
  =(X_n\times_{S_n}S_m)\times_{S_m}S'_m
  =(X_n\times_{S_n}S'_n)\times_{S'_n}S'_m.
\]
Además, $X_0\times_{S_0}S'_0$ es localmente noetheriano, puesto que $X_0$ es de tipo finito sobre $S_0$ \sref{I.6.3.8}.
Se concluye primero, por \sref{I.10.12.3.1}, que $\mathfrak{X}\times_{\mathfrak{S}}\mathfrak{S}'$ es localmente noetheriano;
además, como $X_0\times_{S_0}S'_0$ es de tipo finito sobre $S'_0$ \sref{I.6.3.8}, resulta de \sref{I.10.12.3.1} y \sref{I.10.13.1} que $\mathfrak{X}\times_{\mathfrak{S}}\mathfrak{S}'$ es de tipo finito sobre $\mathfrak{S}'$, lo que termina de probar (ii) (la afirmación relativa a los preesquemas noetherianos es consecuencia inmediata de \sref{I.6.3.8}).
\end{proof}

\begin{corollary}[10.13.6]
\label{I.10.13.6}
Bajo las hipótesis de \sref{I.10.9.9}, si $f$ es un morfismo de tipo finito, también lo es su prolongación $\widehat f$ a los completados.
\end{corollary}

\subsection{Subpreesquemas cerrados de los preesquemas formales}
\label{subsection:I.10.14}

\begin{proposition}[10.14.1]
\label{I.10.14.1}
Sea $\mathfrak{X}$ un preesquema formal localmente noetheriano y sea $\mathscr{A}$ un haz coherente de ideales de $\mathscr{O}_{\mathfrak{X}}$.
Si $\mathfrak{Y}$ es el soporte (cerrado) de $\mathscr{O}_{\mathfrak{X}}/\mathscr{A}$, el espacio topológicamente anillado $(\mathfrak{Y},(\mathscr{O}_{\mathfrak{X}}/\mathscr{A})|\mathfrak{Y})$ es un preesquema formal localmente noetheriano, que es noetheriano si $\mathfrak{X}$ lo es.
\end{proposition}

\begin{proof}
Notemos que $\mathscr{O}_{\mathfrak{X}}/\mathscr{A}$ es coherente en virtud de \sref{I.10.10.3} y \sref[0]{0.5.3.4}; por tanto, su soporte $\mathfrak{Y}$ es cerrado \sref[0]{0.5.2.2}.
Sea $\mathscr{J}$ un ideal de definición de $\mathfrak{X}$ y sea $X_n=(\mathfrak{X},\mathscr{O}_{\mathfrak{X}}/\mathscr{J}^{n+1})$;
el haz de anillos $\mathscr{O}_{\mathfrak{X}}/\mathscr{A}$ es el límite proyectivo de los haces $\mathscr{O}_{\mathfrak{X}}/(\mathscr{A}+\mathscr{J}^{n+1})
=(\mathscr{O}_{\mathfrak{X}}/\mathscr{A})
\otimes_{\mathscr{O}_{\mathfrak{X}}}
(\mathscr{O}_{\mathfrak{X}}/\mathscr{J}^{n+1})$ \sref{I.10.11.3}, todos los cuales tienen por soporte $\mathfrak{Y}$.
El haz $(\mathscr{A}+\mathscr{J}^{n+1})/\mathscr{J}^{n+1}$ es un $\mathscr{O}_{\mathfrak{X}}$-módulo coherente, puesto que $\mathscr{J}^{n+1}$ es coherente; por tanto, $(\mathscr{A}+\mathscr{J}^{n+1})/\mathscr{J}^{n+1}$ es también un $(\mathscr{O}_{\mathfrak{X}}/\mathscr{J}^{n+1})$-módulo coherente \sref[0]{0.5.3.10};
si $Y_n$ es el subpreesquema cerrado de $X_n$ definido por este haz de ideales, es inmediato que $(\mathfrak{Y},(\mathscr{O}_{\mathfrak{X}}/\mathscr{A})|\mathfrak{Y})$
\oldpage[I]{210}
es el preesquema formal límite inductivo de los $Y_n$ y, como se satisfacen las condiciones de \sref{I.10.6.4}, esto prueba que dicho preesquema formal es localmente noetheriano, y noetheriano si $\mathfrak{X}$ lo es (pues entonces $Y_0$ lo es en virtud de \sref{I.6.1.4}).
\end{proof}

\begin{definition}[10.14.2]
\label{I.10.14.2}
{\itshape Se llama subpreesquema cerrado de un preesquema formal $\mathfrak{X}$ a todo preesquema formal de la forma $(\mathfrak{Y},(\mathscr{O}_{\mathfrak{X}}/\mathscr{A})|\mathfrak{Y})$, donde $\mathscr{A}$ es un \erratum[II]{ideal coherente de $\mathscr{O}_{\mathfrak{X}}$};
se dice que este preesquema es el subpreesquema cerrado definido por $\mathscr{A}$.}
\end{definition}

Es claro que la correspondencia así definida entre los \erratum[II]{ideales coherentes de $\mathscr{O}_{\mathfrak{X}}$} y los subpreesquemas cerrados de $\mathfrak{X}$ es biyectiva.

El morfismo de espacios topológicamente anillados $j=(\psi,\theta):\mathfrak{Y}\to\mathfrak{X}$, donde $\psi$ es la inyección $\mathfrak{Y}\to\mathfrak{X}$ y $\theta^\sharp$ el homomorfismo canónico $\mathscr{O}_{\mathfrak{X}}\to\mathscr{O}_{\mathfrak{X}}/\mathscr{A}$, es evidentemente \sref{I.10.4.5} un morfismo de preesquemas formales, al que se llama la \emph{inyección canónica} de $\mathfrak{Y}$ en $\mathfrak{X}$.
Se observará que, si $\mathfrak{X}=\operatorname{Spf}(A)$, donde $A$ es noetheriano y ádico, entonces $\mathscr{A}=\mathfrak{a}^\Delta$, donde $\mathfrak{a}$ es un ideal de $A$ \sref{I.10.10.5}; de lo anterior se sigue inmediatamente que $\mathfrak{Y}=\operatorname{Spf}(A/\mathfrak{a})$, salvo isomorfismo, y que $j$ corresponde \sref{I.10.2.2} al homomorfismo canónico $A\to A/\mathfrak{a}$.

Se dice que un morfismo $f:\mathfrak{Z}\to\mathfrak{X}$ de preesquemas formales localmente noetherianos es una \emph{inmersión cerrada} si se factoriza como $\mathfrak{Z}\xrightarrow{g}\mathfrak{Y}\xrightarrow{j}\mathfrak{X}$, donde $g$ es un isomorfismo de $\mathfrak{Z}$ sobre un subpreesquema cerrado $\mathfrak{Y}$ de $\mathfrak{X}$ y $j$ es la inyección canónica.
Como $j$ es un monomorfismo de espacios anillados, $g$ y $\mathfrak{Y}$ son necesariamente \emph{únicos}.

\begin{proposition}[10.14.3]
\label{I.10.14.3}
Una inmersión cerrada es un morfismo de tipo finito.
\end{proposition}

\begin{proof}
Se reduce inmediatamente al caso en que $\mathfrak{X}$ es un esquema formal afín $\operatorname{Spf}(A)$ y $\mathfrak{Y}=\operatorname{Spf}(A/\mathfrak{a})$;
la proposición resulta entonces de la Proposición~\sref{I.10.13.1}[\emph{c)}].
\end{proof}

\begin{lemma}[10.14.4]
\label{I.10.14.4}
Sea $f:\mathfrak{Y}\to\mathfrak{X}$ un morfismo de preesquemas formales localmente noetherianos, y sea $(U_\alpha)$ un recubrimiento de $f(\mathfrak{Y})$ por abiertos formales afines noetherianos de $\mathfrak{X}$ tales que los $f^{-1}(U_\alpha)$ sean abiertos formales afines noetherianos de $\mathfrak{Y}$.
Para que $f$ sea una inmersión cerrada, es necesario y suficiente que $f(\mathfrak{Y})$ sea un subconjunto cerrado de $\mathfrak{X}$ y que, para todo $\alpha$, la restricción de $f$ a $f^{-1}(U_\alpha)$ corresponda \sref{I.10.4.6} a un homomorfismo sobreyectivo $\Gamma(U_\alpha,\mathscr{O}_{\mathfrak{X}})\to
\Gamma(f^{-1}(U_\alpha),\mathscr{O}_{\mathfrak{Y}})$.
\end{lemma}

\begin{proof}
Las condiciones son evidentemente necesarias.
Recíprocamente, si se satisfacen, y si se designa por $\mathfrak{a}_\alpha$ el núcleo de $\Gamma(U_\alpha,\mathscr{O}_{\mathfrak{X}})\to
\Gamma(f^{-1}(U_\alpha),\mathscr{O}_{\mathfrak{Y}})$, se define un haz coherente de ideales $\mathscr{A}$ de $\mathscr{O}_{\mathfrak{X}}$ tomando $\mathscr{A}|U_\alpha=\mathfrak{a}_\alpha^\Delta$ y tomando $\mathscr{A}$ igual a $\mathscr{O}_{\mathfrak{X}}$ en el complemento de la unión de los $U_\alpha$.
\footnote{\emph{[Trad.] El texto francés imprime que $\sh{A}$ es nulo en este complemento y, en la frase siguiente, que el propio $\mathfrak{a}_\alpha^\Delta$ tiene soporte $U_\alpha\cap f(\mathfrak{Y})$. El ideal definidor debe ser el ideal unidad fuera de $f(\mathfrak{Y})$, mientras que su cociente tiene ese soporte; ambas correcciones son necesarias para que los haces se peguen en el subpreesquema cerrado cuyo espacio subyacente es $f(\mathfrak{Y})$.}}
En efecto, como $f(\mathfrak{Y})$ es cerrado y el soporte del cociente $(\mathscr{O}_{\mathfrak{X}}|U_\alpha)/\mathfrak{a}_\alpha^\Delta$ es $U_\alpha\cap f(\mathfrak{Y})$, todo se reduce a probar que $\mathfrak{a}_\alpha^\Delta$ y $\mathfrak{a}_\beta^\Delta$ inducen el mismo haz sobre todo abierto formal afín noetheriano $V\subset U_\alpha\cap U_\beta$.
Pero la restricción a $f^{-1}(U_\alpha)$ de $f$ es una inmersión cerrada de este preesquema formal en $U_\alpha$, $f^{-1}(V)$ es un abierto formal afín noetheriano de $f^{-1}(U_\alpha)$, y la restricción de $f$ a $f^{-1}(V)$ es una inmersión cerrada;
si $\mathfrak{b}$ es el núcleo del homomorfismo sobreyectivo $\Gamma(V,\mathscr{O}_{\mathfrak{X}})\to
\Gamma(f^{-1}(V),\mathscr{O}_{\mathfrak{Y}})$ correspondiente a esta restricción, es inmediato \sref{I.10.10.2} que $\mathfrak{a}_\alpha^\Delta$ induce $\mathfrak{b}^\Delta$ sobre $V$.
Definido así el haz de ideales $\mathscr{A}$, es claro que $f=j\circ g$\footnote{\emph{[Trad.] La fuente francesa imprime $f=g\circ j$; con $g:\mathfrak{Y}\to\mathfrak{Z}$ y $j:\mathfrak{Z}\to\mathfrak{X}$, el compuesto bien tipado es $j\circ g$.}}, donde $j:\mathfrak{Z}\to\mathfrak{X}$ es la inyección canónica del subpreesquema cerrado $\mathfrak{Z}$ de $\mathfrak{X}$ definido por $\mathscr{A}$, y que $g$ es un isomorfismo de $\mathfrak{Y}$ sobre $\mathfrak{Z}$.
\end{proof}

\begin{proposition}[10.14.5]
\label{I.10.14.5}
\medskip\noindent
\begin{enumerate}
  \item[\rm{(i)}] Si $f:\mathfrak{Z}\to\mathfrak{Y}$ y $g:\mathfrak{Y}\to\mathfrak{X}$ son inmersiones cerradas de preesquemas formales localmente noetherianos, $g\circ f$ es una inmersión cerrada.
\oldpage[I]{211}
  \item[\rm{(ii)}] Sean $\mathfrak{X}$, $\mathfrak{Y}$ y $\mathfrak{S}$ tres preesquemas formales localmente noetherianos, sea $f:\mathfrak{X}\to\mathfrak{S}$ una inmersión cerrada y sea $g:\mathfrak{Y}\to\mathfrak{S}$ un morfismo.
    Entonces el morfismo $\mathfrak{X}\times_{\mathfrak{S}}\mathfrak{Y}\to\mathfrak{Y}$ es una inmersión cerrada.
  \item[\rm{(iii)}] Sea $\mathfrak{S}$ un preesquema formal localmente noetheriano, y sean $\mathfrak{X}'$ y $\mathfrak{Y}'$ dos $\mathfrak{S}$-preesquemas formales localmente noetherianos tales que $\mathfrak{X}'\times_{\mathfrak{S}}\mathfrak{Y}'$ sea localmente noetheriano.
    Si $\mathfrak{X}$ y $\mathfrak{Y}$ son $\mathfrak{S}$-preesquemas formales localmente noetherianos, y $f:\mathfrak{X}\to\mathfrak{X}'$ y $g:\mathfrak{Y}\to\mathfrak{Y}'$ son $\mathfrak{S}$-morfismos que son inmersiones cerradas, entonces $f\times_{\mathfrak{S}}g$ es una inmersión cerrada.
\end{enumerate}
\end{proposition}

\begin{proof}
En virtud de \sref{I.3.5.1}, basta de nuevo probar (i) y (ii).

Para probar (i), se puede suponer que $\mathfrak{Y}$ (resp. $\mathfrak{Z}$) es un subpreesquema cerrado de $\mathfrak{X}$ (resp. $\mathfrak{Y}$), definido por un haz coherente $\mathscr{J}$ (resp. $\mathscr{K}$) de ideales de $\mathscr{O}_{\mathfrak{X}}$ (resp. $\mathscr{O}_{\mathfrak{Y}}$);
si $\psi$ es la inyección $\mathfrak{Y}\to\mathfrak{X}$ de los espacios subyacentes, entonces $\psi_*(\mathscr{K})$ es un haz coherente de ideales de $\psi_*(\mathscr{O}_{\mathfrak{Y}})
=\mathscr{O}_{\mathfrak{X}}/\mathscr{J}$ \sref[0]{0.5.3.12}, y por tanto también un $\mathscr{O}_{\mathfrak{X}}$-módulo coherente \sref[0]{0.5.3.10};
el núcleo $\mathscr{K}_1$ de $\mathscr{O}_{\mathfrak{X}}\to
(\mathscr{O}_{\mathfrak{X}}/\mathscr{J})/\psi_*(\mathscr{K})$ es, pues, un haz coherente de ideales de $\mathscr{O}_{\mathfrak{X}}$ \sref[0]{0.5.3.4}, y $\mathscr{O}_{\mathfrak{X}}/\mathscr{K}_1$ es isomorfo a $\psi_*(\mathscr{O}_{\mathfrak{Y}}/\mathscr{K})$, lo que prueba que $\mathfrak{Z}$ es isomorfo a un subpreesquema cerrado de $\mathfrak{X}$.

Para probar (ii), se puede reducir inmediatamente al caso en que $\mathfrak{S}=\operatorname{Spf}(A)$, $\mathfrak{X}=\operatorname{Spf}(B)$ y $\mathfrak{Y}=\operatorname{Spf}(C)$, donde $A$ es un anillo noetheriano $\mathfrak{J}$-ádico, $B=A/\mathfrak{a}$ (siendo $\mathfrak{a}$ un ideal de $A$) y $C$ es una $A$-álgebra topológica ádica y noetheriana.
Todo se reduce entonces a probar que el homomorfismo $C\to C\widehat{\otimes}_A(A/\mathfrak{a})$ es \emph{sobreyectivo}:
pero $A/\mathfrak{a}$ es un $A$-módulo de tipo finito, y su topología es la topología $\mathfrak{J}$-ádica;
resulta entonces de \sref[0]{0.7.7.8} que $C\widehat{\otimes}_A(A/\mathfrak{a})$ se identifica con $C\otimes_A(A/\mathfrak{a})=C/\mathfrak{a}C$, de donde se sigue la afirmación.
\end{proof}

\begin{corollary}[10.14.6]
\label{I.10.14.6}
Bajo las hipótesis de \sref{I.10.14.5}[(ii)], sean $p:\mathfrak{X}\times_{\mathfrak{S}}\mathfrak{Y}\to\mathfrak{X}$ y $q:\mathfrak{X}\times_{\mathfrak{S}}\mathfrak{Y}\to\mathfrak{Y}$ las proyecciones, de modo que el diagrama
\[
  \xymatrix{
    \mathfrak{X} \ar[d]_{f}
    & \mathfrak{X}\times_{\mathfrak{S}}\mathfrak{Y}
      \ar[l]_{p} \ar[d]^{q}\\
    \mathfrak{S}
    & \mathfrak{Y} \ar[l]^{g}
  }
\]
es conmutativo.
Para todo $\mathscr{O}_{\mathfrak{X}}$-módulo coherente $\mathscr{F}$, se tiene entonces un isomorfismo canónico de $\mathscr{O}_{\mathfrak{Y}}$-módulos
\[
  \label{I.10.14.6.1-fr}
  u:g^*(f_*(\mathscr{F}))
  \xrightarrow{\sim}
  q_*(p^*(\mathscr{F})).
  \tag{10.14.6.1}
\]
\end{corollary}

\begin{proof}
Se sabe que definir un homomorfismo $g^*(f_*(\mathscr{F}))\to q_*(p^*(\mathscr{F}))$ equivale a definir un homomorfismo $f_*(\mathscr{F})\to g_*(q_*(p^*(\mathscr{F})))
=f_*(p_*(p^*(\mathscr{F})))$ \sref[0]{0.4.4.3}:
tomaremos como $u$ el homomorfismo que corresponde por esta adjunción a $f_*(\rho)$,\footnote{\emph{[Trad.] La fuente francesa escribe $u=f_*(\rho)$; estos morfismos viven sobre espacios distintos y tienen haces fuente y meta distintos, siendo $f_*(\rho)$ el adjunto de $u$.}} donde $\rho$ es el homomorfismo canónico $\mathscr{F}\to p_*(p^*(\mathscr{F}))$ \sref[0]{0.4.4.3}.
Para ver que $u$ es un isomorfismo, se puede reducir inmediatamente al caso en que $\mathfrak{S}$, $\mathfrak{X}$ y $\mathfrak{Y}$ son los espectros formales de anillos ádicos noetherianos $A$, $B$ y $C$, respectivamente, que satisfacen las condiciones de \sref{I.10.14.5}[(ii)];
se tiene entonces $\mathscr{F}=M^\Delta$, donde $M$ es un $(A/\mathfrak{a})$-módulo de tipo finito \sref{I.10.10.5}, y los dos miembros de \sref{I.10.14.6.1} se identifican, respectivamente, en virtud de \sref{I.10.10.8}, con $(C\otimes_A M)^\Delta$ y $((C/\mathfrak{a}C)\otimes_{A/\mathfrak{a}}M)^\Delta$; de ahí el corolario, pues $(C/\mathfrak{a}C)\otimes_{A/\mathfrak{a}}M
=(C\otimes_A(A/\mathfrak{a}))\otimes_{A/\mathfrak{a}}M$ se identifica canónicamente con $C\otimes_A M$.
\end{proof}

\oldpage[I]{212}
\begin{corollary}[10.14.7]
\label{I.10.14.7}
Sean $X$ un preesquema usual localmente noetheriano, $Y$ un subpreesquema cerrado de $X$, $j$ la inyección canónica $Y\to X$, $X'$ un subconjunto cerrado de $X$ y $Y'=Y\cap X'$;
entonces $\widehat{j}:Y_{/Y'}\to X_{/X'}$ es una inmersión cerrada y, para todo $\mathscr{O}_Y$-módulo coherente $\mathscr{F}$, se tiene
\[
  \widehat{j}_*(\mathscr{F}_{/Y'})=(j_*(\mathscr{F}))_{/X'}.
\]
\end{corollary}

\begin{proof}
Como $Y'=j^{-1}(X')$, basta utilizar \sref{I.10.9.9} y aplicar \sref{I.10.14.5} y \sref{I.10.14.6}.
\end{proof}

\subsection{Preesquemas formales separados}
\label{subsection:I.10.15}

\begin{definition}[10.15.1]
\label{I.10.15.1}
Sean $\mathfrak{S}$ un preesquema formal, $\mathfrak{X}$ un $\mathfrak{S}$-preesquema formal y $f:\mathfrak{X}\to\mathfrak{S}$ el morfismo estructural.
Se llama morfismo diagonal $\Delta_{\mathfrak{X}|\mathfrak{S}}:\mathfrak{X}\to
\mathfrak{X}\times_{\mathfrak{S}}\mathfrak{X}$ (denotado también por $\Delta_{\mathfrak{X}}$) al morfismo $(1_{\mathfrak{X}},1_{\mathfrak{X}})_{\mathfrak{S}}$.
Se dice que $\mathfrak{X}$ es separado sobre $\mathfrak{S}$, o que es un $\mathfrak{S}$-esquema formal, o que $f$ es un morfismo separado, si la imagen del espacio subyacente a $\mathfrak{X}$ por $\Delta_{\mathfrak{X}}$ es un subconjunto cerrado del espacio subyacente a $\mathfrak{X}\times_{\mathfrak{S}}\mathfrak{X}$.
Se dice que un preesquema formal $\mathfrak{X}$ es separado, o que es un esquema formal, si es separado sobre $\mathbf{Z}$.
\end{definition}

\begin{proposition}[10.15.2]
\label{I.10.15.2}
Supongamos que los preesquemas formales $\mathfrak{S}$ y $\mathfrak{X}$ son, respectivamente, límites inductivos de sucesiones $(S_n)$ y $(X_n)$ de preesquemas usuales, y que el morfismo $f:\mathfrak{X}\to\mathfrak{S}$ es el límite inductivo de una sucesión de morfismos $f_n:X_n\to S_n$.
Para que $f$ sea separado, es necesario y suficiente que el morfismo $f_0:X_0\to S_0$ lo sea.
\end{proposition}

\begin{proof}
En efecto, $\Delta_{\mathfrak{X}|\mathfrak{S}}$ es entonces el límite inductivo de la sucesión de morfismos $\Delta_{X_n|S_n}$ \sref{I.10.7.4}; la imagen por $\Delta_{\mathfrak{X}|\mathfrak{S}}$ del espacio subyacente $\mathfrak{X}$ (resp. el espacio subyacente $\mathfrak{X}\times_{\mathfrak{S}}\mathfrak{X}$) es idéntica a la imagen por $\Delta_{X_0|S_0}$ del espacio subyacente $X_0$ (resp. al espacio subyacente $X_0\times_{S_0}X_0$);
de donde se sigue la conclusión.
\end{proof}

\begin{proposition}[10.15.3]
\label{I.10.15.3}
Se supone en lo que sigue que todos los preesquemas formales \textup{(resp.} morfismos de preesquemas formales) considerados son límites inductivos de sucesiones de preesquemas usuales \textup{(resp.} de morfismos de preesquemas usuales).
\begin{enumerate}
    \item[\rm{(i)}] El compuesto de dos morfismos separados es separado.
    \item[\rm{(ii)}] Si $f:\mathfrak{X}\to\mathfrak{X}'$ y $g:\mathfrak{Y}\to\mathfrak{Y}'$ son dos $\mathfrak{S}$-morfismos separados, entonces $f\times_{\mathfrak{S}}g$ es separado.
    \item[\rm{(iii)}] Si $f:\mathfrak{X}\to\mathfrak{Y}$ es un $\mathfrak{S}$-morfismo separado, entonces el $\mathfrak{S}'$-morfismo $f_{(\mathfrak{S}')}$ es separado para toda extensión $\mathfrak{S}'\to\mathfrak{S}$ del preesquema formal base.
    \item[\rm{(iv)}] Si el compuesto $g\circ f$ de dos morfismos es separado, entonces $f$ es separado.
\end{enumerate}
\emph{(Se sobreentiende en este enunciado que, si un mismo preesquema formal $\mathfrak{Z}$ interviene varias veces en una misma proposición, se lo considera como límite inductivo de la \emph{misma} sucesión $(Z_n)$ de preesquemas usuales en todos los lugares donde figura, y los morfismos de $\mathfrak{Z}$ en otro preesquema formal (resp. de un preesquema formal en $\mathfrak{Z}$) como límites inductivos de morfismos de los $Z_n$ en preesquemas usuales (resp. de preesquemas usuales en los $Z_n$)).}
\end{proposition}

\begin{proof}
Con las notaciones de \sref{I.10.15.2}, se tiene, en efecto, $(g\circ f)_0=g_0\circ f_0$ y $(f\times_{\mathfrak{S}}g)_0=f_0\times_{S_0}g_0$;
las afirmaciones de \sref{I.10.15.3} son entonces consecuencias inmediatas de \sref{I.10.15.2} y de las afirmaciones correspondientes de \sref{I.5.5.1} para los preesquemas usuales.
\end{proof}

Dejamos al lector el cuidado de enunciar, para la misma clase de preesquemas formales y de morfismos que en \sref{I.10.15.3}, las proposiciones correspondientes a \sref{I.5.5.5},
\oldpage[I]{213}
\sref{I.5.5.9} y \sref{I.5.5.10} (sustituyendo «abierto afín» por «abierto formal afín que satisface la condición \emph{b)} de \sref{I.10.6.3}»).

Un razonamiento análogo muestra también que todo esquema formal afín \emph{noetheriano} es separado, lo que justifica la terminología.

\begin{proposition}[10.15.4]
\label{I.10.15.4}
Sean $\mathfrak{S}$ un preesquema formal localmente noetheriano, y $\mathfrak{X}$ y $\mathfrak{Y}$ dos $\mathfrak{S}$-preesquemas formales localmente noetherianos tales que $\mathfrak{X}$ o $\mathfrak{Y}$ sea de tipo finito sobre $\mathfrak{S}$ (de modo que $\mathfrak{X}\times_{\mathfrak{S}}\mathfrak{Y}$ sea localmente noetheriano) y que $\mathfrak{Y}$ sea separado sobre $\mathfrak{S}$.
Sea $f:\mathfrak{X}\to\mathfrak{Y}$ un $\mathfrak{S}$-morfismo; entonces el morfismo grafo $\Gamma_f=(1_{\mathfrak{X}},f)_{\mathfrak{S}}:\mathfrak{X}\to
\mathfrak{X}\times_{\mathfrak{S}}\mathfrak{Y}$ es una inmersión cerrada.
\end{proposition}

\begin{proof}
Se puede suponer que $\mathfrak{S}$ es el límite inductivo de una sucesión $(S_n)$ de preesquemas localmente noetherianos, que $\mathfrak{X}$ (resp. $\mathfrak{Y}$) es el límite inductivo de una sucesión $(X_n)$ (resp. $(Y_n)$) de $S_n$-preesquemas, y que $f$ es el límite inductivo de una sucesión de $S_n$-morfismos $f_n:X_n\to Y_n$;
entonces $\mathfrak{X}\times_{\mathfrak{S}}\mathfrak{Y}$ es el límite inductivo de la sucesión $(X_n\times_{S_n}Y_n)$, y $\Gamma_f$ el límite inductivo de la sucesión $(\Gamma_{f_n})$ \sref{I.10.7.4};
por hipótesis, $Y_0$ es separado sobre $S_0$ \sref{I.10.15.2}, por lo que el espacio $\Gamma_{f_0}(X_0)$ es un subespacio cerrado de $X_0\times_{S_0}Y_0$;
como los espacios subyacentes a $\mathfrak{X}\times_{\mathfrak{S}}\mathfrak{Y}$ (resp. $\Gamma_f(\mathfrak{X})$) y a $X_0\times_{S_0}Y_0$ (resp. $\Gamma_{f_0}(X_0)$) son idénticos, se ve ya que $\Gamma_f(\mathfrak{X})$ es un subespacio \emph{cerrado} de $\mathfrak{X}\times_{\mathfrak{S}}\mathfrak{Y}$.
Observemos ahora que, cuando $(U,V)$ recorre el conjunto de pares formados por un abierto formal afín noetheriano $U$ (resp. $V$) de $\mathfrak{X}$ (resp. $\mathfrak{Y}$) tales que $f(U)\subset V$, los abiertos $U\times_{\mathfrak{S}}V$ forman un recubrimiento de $\Gamma_f(\mathfrak{X})$ en $\mathfrak{X}\times_{\mathfrak{S}}\mathfrak{Y}$ y, si $f_U:U\to V$ es la restricción de $f$ a $U$, entonces $\Gamma_{f_U}:U\to U\times_{\mathfrak{S}}V$ es la restricción de $\Gamma_f$ a $U$.
Si demostramos que $\Gamma_{f_U}$ es una inmersión cerrada, entonces $\Gamma_f$ será una inmersión cerrada \sref{I.10.14.4}; en otras palabras, se queda reducido al caso en que $\mathfrak{S}=\operatorname{Spf}(A)$, $\mathfrak{X}=\operatorname{Spf}(B)$ y $\mathfrak{Y}=\operatorname{Spf}(C)$ son afines (siendo $A$, $B$ y $C$ anillos ádicos noetherianos), y $f$ corresponde a un $A$-homomorfismo continuo $\varphi:C\to B$;
entonces $\Gamma_f$ corresponde al único homomorfismo continuo $\omega:B\widehat{\otimes}_A C\to B$ que, compuesto con los homomorfismos canónicos $B\to B\widehat{\otimes}_A C$ y $C\to B\widehat{\otimes}_A C$, da respectivamente la identidad y $\varphi$.
Pero es claro que $\omega$ es \emph{sobreyectivo}, de donde se sigue la afirmación.
\end{proof}

\begin{corollary}[10.15.5]
\label{I.10.15.5}
Sean $\mathfrak{S}$ un preesquema formal localmente noetheriano y $\mathfrak{X}$ un $\mathfrak{S}$-preesquema de tipo finito;
para que $\mathfrak{X}$ sea separado sobre $\mathfrak{S}$, es necesario y suficiente que el morfismo diagonal $\mathfrak{X}\to\mathfrak{X}\times_{\mathfrak{S}}\mathfrak{X}$ sea una inmersión cerrada.
\end{corollary}

\begin{proposition}[10.15.6]
\label{I.10.15.6}
Una inmersión cerrada $j:\mathfrak{Y}\to\mathfrak{X}$ de preesquemas formales localmente noetherianos es un morfismo separado.
\end{proposition}

\begin{proof}
Con las notaciones de \sref{I.10.14.2}, $j_0:Y_0\to X_0$ es una inmersión cerrada, por tanto un morfismo separado, y basta aplicar \sref{I.10.15.2}.
\end{proof}

\begin{proposition}[10.15.7]
\label{I.10.15.7}
Sean $X$ un preesquema (usual) localmente noetheriano, $X'$ un subconjunto cerrado de $X$ y $\widehat{X}=X_{/X'}$.
Para que $\widehat{X}$ sea separado, es necesario y suficiente que $\widehat{X}_{\mathrm{red}}$ lo sea, y basta que $X$ sea separado.
\end{proposition}

\begin{proof}
Con las notaciones de \sref{I.10.8.5}, para que $\widehat{X}$ sea separado, es necesario y suficiente que $X'_0$ lo sea \sref{I.10.15.2}; como $\widehat{X}_{\mathrm{red}}=(X'_0)_{\mathrm{red}}$, esto equivale a que $\widehat{X}_{\mathrm{red}}$ sea separado \sref{I.5.5.1}[(vi)].
\end{proof}
% END source_es/ega1/ega1-10.tex


% BEGIN source_es/ega1/ega1-backmatter-index-notation.tex
% Proyección TeX, basada ante todo en la fuente, del índice de notaciones de EGA I.
% Linaje de borrador: Markdown comparativo de Jacob C. Reinhold; rige el escaneo francés de NUMDAM.
\clearpage
\section*{Índice de notaciones}
\phantomsection
\label{section:index-notation}
\begingroup
\small
\setlength{\columnsep}{2em}
\begin{multicols}{2}
\raggedcolumns
\providecommand{\EGAIndexLine}[2]{\par\noindent\hangindent=1em\hangafter=1 #1\ \dotfill\ \textup{#2}\par}

\phantomsection\label{section:index-notation-chapter0}
\medskip\noindent\textbf{Capítulo 0}\par\nopagebreak
\EGAIndexLine{$M_{[\phi]}$}{(0.1.0.2)}
\EGAIndexLine{$B\mathfrak{J}$, $\mathfrak{J}B$}{(0.1.0.3)}
\EGAIndexLine{$r(\mathfrak{a})$}{(0.1.1.1)}
\EGAIndexLine{$\mathfrak{R}(A)$}{(0.1.1.2)}
\EGAIndexLine{$S_{f}$}{(0.1.2.1)}
\EGAIndexLine{$S^{-1}A$, $S^{-1}M$, $m/s$, $i_A^S$, $i_M^S$, $i^S$}{(0.1.2.2)}
\EGAIndexLine{$A_{f}$, $M_{f}$, $A_{\mathfrak{p}}$, $M_{\mathfrak{p}}$}{(0.1.2.3)}
\EGAIndexLine{$S^{-1}u$}{(0.1.3.1)}
\EGAIndexLine{$\rho^{T,S}_{A}$, $\rho^{T,S}_{M}$, $\rho^{T,S}$}{(0.1.4.1)}
\EGAIndexLine{$\operatorname{Supp}(M)$}{(0.1.7.1)}
\EGAIndexLine{$\mathcal{F}|U$, $u|U$}{(0.3.1.5)}
\EGAIndexLine{$\mathcal{F}_{x}$, $s_{x}$, $\Gamma(U, \mathcal{F})$, $u(s)$, $\operatorname{Supp}(\mathcal{F})$}{(0.3.1.6)}
\EGAIndexLine{$\psi_{*}(\mathcal{F})$}{(0.3.4.1)}
\EGAIndexLine{$\psi_{*}(u)$}{(0.3.4.2)}
\EGAIndexLine{$\psi_{x}$}{(0.3.4.4)}
\EGAIndexLine{$\mathcal{G} \to \mathcal{F}$}{(0.3.5.1)}
\EGAIndexLine{$u^{\sharp}$, $v^{\flat}$, $\rho_{\mathcal{G}}$}{(0.3.5.3)}
\EGAIndexLine{$\psi^*(\mathcal{G})$, $\psi^*(v)$, $\sigma_{\mathcal{F}}$}{(0.3.5.5)}
\EGAIndexLine{$\mathcal{O}_{X}$, $\mathcal{O}_{X,x}$, $\mathcal{O}_{x}$, $1$, $e$}{(0.4.1.1)}
\EGAIndexLine{$\check{\mathcal{F}}$, $\bigwedge^{p} \mathcal{F}$}{(0.4.1.4)}
\EGAIndexLine{$\mathcal{J}\mathcal{F}$}{(0.4.1.5)}
\EGAIndexLine{$\Psi_{*}(\mathcal{F})$, $\Psi_{*}(u)$}{(0.4.2.1)}
\EGAIndexLine{$\Psi_{*}(\mathcal{C})$}{(0.4.2.4)}
\EGAIndexLine{$\Psi^*(\mathcal{G})$, $\Psi^*(v)$}{(0.4.3.1)}
\EGAIndexLine{$\Psi^*(\mathcal{C})$}{(0.4.3.4)}
\EGAIndexLine{$\Psi^*(\mathcal{J})\mathcal{O}_{X}$, $\mathcal{J}^{c}_{X}$}{(0.4.3.5)}
\EGAIndexLine{$\mathcal{G} \to \mathcal{F}$}{(0.4.4.1)}
\EGAIndexLine{$u_{\theta}^{\sharp}$, $u^{\sharp}$, $v_{\theta}^{\flat}$, $v^{\flat}$, $\rho_{\mathcal{G}}$, $\sigma_{\mathcal{F}}$}{(0.4.4.3)}
\EGAIndexLine{$u_{1} \otimes u_{2}$}{(0.4.4.4)}
\EGAIndexLine{$\mathcal{L}^{-1}$}{(0.5.4.3)}
\EGAIndexLine{$\mathcal{L}^{\otimes n}$}{(0.5.4.4)}
\EGAIndexLine{$\Gamma_{*}(\mathcal{L})$, $\Gamma_{*}(\mathcal{L}, \mathcal{F})$}{(0.5.4.6)}
\EGAIndexLine{$\mathcal{O}_{X}^{*}$}{(0.5.4.7)}
\EGAIndexLine{$\mathfrak{m}_{x}$, $\kappa(x)$, $f(x)$}{(0.5.5.1)}
\EGAIndexLine{$X_{f}$}{(0.5.5.2)}
\EGAIndexLine{$\operatorname{Im}(M' \otimes_{A} N')$}{(0.6.0)}
\EGAIndexLine{$\widehat A$, $\widehat M$}{(0.7.2.3) y (0.7.3.1)}
\EGAIndexLine{$A\{T_{1}, \cdots, T_{r}\}$}{(0.7.5.1)}
\EGAIndexLine{$A\{S^{-1}\}$}{(0.7.6.1)}
\EGAIndexLine{$\mathfrak{a}\{S^{-1}\}$}{(0.7.6.9)}
\EGAIndexLine{$A_{\{f\}}$, $\mathfrak{a}_{\{f\}}$}{(0.7.6.14)}
\EGAIndexLine{$A_{\{S\}}$}{(0.7.6.14)}
\EGAIndexLine{$(M \otimes_{A} N)^{\wedge}$, $M \widehat{\otimes}_{A} N$}{(0.7.7.1)}
\EGAIndexLine{$u \widehat{\otimes} v$}{(0.7.7.3)}

\phantomsection\label{section:index-notation-chapterI}
\medskip\noindent\textbf{Capítulo I}\par\nopagebreak
\EGAIndexLine{$\operatorname{Spec}(A)$, $\mathfrak{j}_{x}$, $\mathfrak{m}_{x}$, $\kappa(x)$, $f(x)$, $M_{x}$, $r(E)$, $V(E)$, $V(f)$, $D(f)$}{(I, \hyperref[I.1.1.1]{1.1.1})}
\EGAIndexLine{$\mathfrak{j}(Y)$}{(I, \hyperref[I.1.1.3]{1.1.3})}
\EGAIndexLine{${}^{a}\phi$}{(I, \hyperref[I.1.2.1]{1.2.1})}
\EGAIndexLine{$S_{f}'$}{(I, \hyperref[I.1.3.1]{1.3.1})}
\EGAIndexLine{$\rho_{g,f}$}{(I, \hyperref[I.1.3.3]{1.3.3})}
\EGAIndexLine{$\widetilde A$, $\widetilde{M}$, $\theta_{f}$}{(I, \hyperref[I.1.3.4]{1.3.4})}
\EGAIndexLine{$\widetilde u$}{(I, \hyperref[I.1.3.5]{1.3.5})}
\EGAIndexLine{$\widetilde{\phi}$}{(I, \hyperref[I.1.6.1]{1.6.1})}
\EGAIndexLine{$A(X)$}{(I, \hyperref[I.1.7.1]{1.7.1})}
\EGAIndexLine{$\mathcal{O}_{X/Y}$}{(I, \hyperref[I.2.1.6]{2.1.6})}
\EGAIndexLine{$\operatorname{Hom}(X, Y)$}{(I, \hyperref[I.2.2.1]{2.2.1})}
\EGAIndexLine{$\operatorname{Hom}_{S}(X, Y)$, $1_X$}{(I, \hyperref[I.2.5.2]{2.5.2})}
\EGAIndexLine{$\Gamma(X/S)$}{(I, \hyperref[I.2.5.5]{2.5.5})}
\EGAIndexLine{$X \amalg Y$}{(I, \hyperref[I.3.1.1]{3.1.1})}
\EGAIndexLine{$X \times_{S} Y$, $X \times Y$, $(g, h)_{S}$, $u \times_{S} v$, $u \times v$}{(I, \hyperref[I.3.2.1]{3.2.1})}
\EGAIndexLine{$X \times_{A} Y$, $X \otimes_{A} B$, $(g, h)_{A}$, $u \times_{A} v$}{(I, \hyperref[I.3.2.1]{3.2.1})}
\EGAIndexLine{$X_{(S')}$}{(I, \hyperref[I.3.3.6]{3.3.6})}
\EGAIndexLine{$f_{(S')}$}{(I, \hyperref[I.3.3.7]{3.3.7})}
\EGAIndexLine{$\Gamma_{f}$}{(I, \hyperref[I.3.3.14]{3.3.14})}
\EGAIndexLine{$X(T)$}{(I, \hyperref[I.3.4.1]{3.4.1})}
\EGAIndexLine{$P \times_{R} Q$}{(I, \hyperref[I.3.4.2]{3.4.2})}
\EGAIndexLine{$X(T)_{S}$}{(I, \hyperref[I.3.4.3]{3.4.3})}
\EGAIndexLine{$X(B)$, $X(B)_{A}$}{(I, \hyperref[I.3.4.4]{3.4.4})}
\EGAIndexLine{$X \otimes_{y} B$, $X \otimes_{\mathcal{O}_y} B$}{(I, \hyperref[I.3.6.3]{3.6.3})}
\EGAIndexLine{$Z \leq Y$}{(I, \hyperref[I.4.1.10]{4.1.10})}
\EGAIndexLine{$f^{-1}(Y)$}{(I, \hyperref[I.4.4.1]{4.4.1})}
\EGAIndexLine{$\mathcal{N}_{X}$}{(I, \hyperref[I.5.1.1]{5.1.1})}
\EGAIndexLine{$X_{\mathrm{red}}$}{(I, \hyperref[I.5.1.3]{5.1.3})}
\EGAIndexLine{$f_{\mathrm{red}}$}{(I, \hyperref[I.5.1.5]{5.1.5})}
\EGAIndexLine{$\Delta_{X/S}$, $\Delta_{X}$, $\Delta$}{(I, \hyperref[I.5.3.1]{5.3.1})}
\EGAIndexLine{$\operatorname{rg}_{K}(X)$}{(I, \hyperref[I.6.4.5]{6.4.5})}
\EGAIndexLine{$n(X)$}{(I, \hyperref[I.6.4.8]{6.4.8})}
\EGAIndexLine{$\Gamma_{\mathrm{rat}}(X/Y)$}{(I, \hyperref[I.7.1.2]{7.1.2})}
\EGAIndexLine{$R(X)$}{(I, \hyperref[I.7.1.3]{7.1.3})}
\EGAIndexLine{$\mathcal{R}(X)$}{(I, \hyperref[I.7.3.1]{7.3.1})}
\EGAIndexLine{$L(A)$}{(I, \hyperref[I.8.1.2]{8.1.2})}
\EGAIndexLine{$\delta(f)$}{(I, \hyperref[I.8.2.1]{8.2.1})}
\EGAIndexLine{$\mathcal{F} \otimes_{\mathcal{O}_{S}} \mathcal{G}$, $\mathcal{F} \otimes_{S} \mathcal{G}$}{(I, \hyperref[I.9.1.2]{9.1.2})}
\EGAIndexLine{$\overline{\mathcal{G}}$}{(I, \hyperref[I.9.4.1]{9.4.1})}
\EGAIndexLine{$\overline Y$}{(I, \hyperref[I.9.5.11]{9.5.11})}
\EGAIndexLine{$\operatorname{Spf}(A)$, $\mathcal{O}_{\mathfrak{X}}$ ($\mathfrak{X}=\operatorname{Spf}(A)$)}{(I, \hyperref[I.10.1.2]{10.1.2})}
\EGAIndexLine{$\mathfrak{D}(f)$}{(I, \hyperref[I.10.1.4]{10.1.4})}
\EGAIndexLine{${}^{a}\phi$, $\widetilde{\phi}$}{(I, \hyperref[I.10.2.1]{10.2.1})}
\EGAIndexLine{$\mathfrak{J}^{\Delta}$}{(I, \hyperref[I.10.3.1]{10.3.1})}
\EGAIndexLine{$\mathfrak{X} \times_{\mathfrak{S}} \mathfrak{Y}$}{(I, \hyperref[I.10.7.3]{10.7.3})}
\EGAIndexLine{$\mathcal{F}_{/X'}$, $\widehat{\mathcal{F}}$, $u_{/X'}$, $\widehat{u}$}{(I, \hyperref[I.10.8.4]{10.8.4})}
\EGAIndexLine{$X_{/X'}$, $\widehat{X}$}{(I, \hyperref[I.10.8.5]{10.8.5})}
\EGAIndexLine{$\widehat{f}$}{(I, \hyperref[I.10.9.1]{10.9.1})}
\EGAIndexLine{$M^{\Delta}$, $u^{\Delta}$}{(I, \hyperref[I.10.10.1]{10.10.1})}
\EGAIndexLine{$\Delta_{\mathfrak{S}|\mathfrak{X}}$, $\Delta_{\mathfrak{X}}$}{(I, \hyperref[I.10.15.1]{10.15.1})}
\end{multicols}
\endgroup
% END source_es/ega1/ega1-backmatter-index-notation.tex

% BEGIN source_es/ega1/ega1-backmatter-index-terminology.tex
% Proyección española, basada ante todo en la fuente, del índice terminológico de EGA I.
% Rige el escaneo francés de NUMDAM. Las entradas conservan el orden del índice
% de la fuente para que cada fila pueda cotejarse directamente con las pp. 219--224 impresas.
\clearpage
\section*{Índice terminológico}
\phantomsection
\label{section:index-terminology}
\begingroup
\footnotesize
\setlength{\columnsep}{2em}
\begin{multicols}{2}
\raggedcolumns
\providecommand{\EGATermIndexLine}[2]{\par\noindent\hangindent=1em\hangafter=1 #1\ \dotfill\ \textup{#2}\par}

\noindent\emph{El orden de las entradas sigue el índice de la fuente francesa;
por ello, las traducciones españolas no siempre aparecen en orden alfabético español.}\par\medskip

\EGATermIndexLine{Adherencia de un subpreesquema}{(I, \hyperref[I.9.5.11]{9.5.11})}
\EGATermIndexLine{$\mathcal O_X$-álgebra}{(0.4.1.3)}
\EGATermIndexLine{Anillo ádico; anillo $\mathfrak J$-ádico}{(0.7.1.9)}
\EGATermIndexLine{Anillo admisible}{(0.7.1.2)}
\EGATermIndexLine{Anillo completo de fracciones}{(0.7.6.5)}
\EGATermIndexLine{Anillo de fracciones}{(0.1.2.2)}
\EGATermIndexLine{Anillo de funciones racionales}{(I, \hyperref[I.2.1.6]{2.1.6} y \hyperref[I.7.1.3]{7.1.3})}
\EGATermIndexLine{Anillo de un esquema afín}{(I, \hyperref[I.1.7.1]{1.7.1})}
\EGATermIndexLine{Anillo íntegro}{(0.1.0.6)}
\EGATermIndexLine{Anillo topologizado linealmente}{(0.7.1.1)}
\EGATermIndexLine{Anillo local}{(0.1.0.7)}
\EGATermIndexLine{Anillo local dominante}{(I, \hyperref[I.8.1.1]{8.1.1})}
\EGATermIndexLine{Anillo local de $X$ a lo largo de $Y$; anillo local de $Y$ en $X$}{(I, \hyperref[I.2.1.6]{2.1.6})}
\EGATermIndexLine{Anillo preádico; anillo $\mathfrak J$-preádico}{(0.7.1.9)}
\EGATermIndexLine{Anillo preadmisible}{(0.7.1.2)}
\EGATermIndexLine{Anillo reducido}{(0.1.1.1)}
\EGATermIndexLine{Anillo regular}{(0.4.1.3)}
\EGATermIndexLine{Anillos locales emparentados}{(I, \hyperref[I.8.1.4]{8.1.4})}
\EGATermIndexLine{Anulador de un $\mathcal O_X$-módulo}{(0.5.3.7)}
\EGATermIndexLine{Aplicación de espectros asociada a un homomorfismo de anillos}{(I, \hyperref[I.1.2.1]{1.2.1})}
\EGATermIndexLine{Aplicación racional; $S$-aplicación racional}{(I, \hyperref[I.7.1.2]{7.1.2})}
\EGATermIndexLine{Aplicación racional definida en un punto}{(I, \hyperref[I.7.2.1]{7.2.1})}
\EGATermIndexLine{Aplicación racional inducida sobre un subconjunto abierto}{(I, \hyperref[I.7.1.2]{7.1.2})}
\EGATermIndexLine{Aplicación racional inducida sobre $\operatorname{Spec}(\mathcal O_x)$}{(I, \hyperref[I.7.2.8]{7.2.8})}

\EGATermIndexLine{$\mathcal O_X$-módulo coherente}{(0.5.3.1)}
\EGATermIndexLine{$\mathcal O_X$-álgebra coherente}{(0.5.3.6)}
\EGATermIndexLine{Completado de un $\mathcal O_X$-módulo y de un homomorfismo de $\mathcal O_X$-módulos a lo largo de un subconjunto cerrado}{(I, \hyperref[I.10.8.4]{10.8.4})}
\EGATermIndexLine{Completado de un preesquema a lo largo de un subconjunto cerrado}{(I, \hyperref[I.10.8.5]{10.8.5})}
\EGATermIndexLine{Componente irreducible}{(0.2.1.5)}
\EGATermIndexLine{Compuesto de un $\psi$-morfismo y un $\psi'$-morfismo}{(0.3.5.2)}
\EGATermIndexLine{Compuesto de un $\Psi$-morfismo y un $\Psi'$-morfismo}{(0.4.4.2)}
\EGATermIndexLine{Condición de pegado}{(0.3.3.1 y 0.4.1.6)}
\EGATermIndexLine{Cuerpo de base de un preesquema algebraico}{(I, \hyperref[I.6.4.1]{6.4.1})}
\EGATermIndexLine{Cuerpo residual de un punto geométrico}{(I, \hyperref[I.3.4.5]{3.4.5})}

\EGATermIndexLine{Diagonal de $X\times_S X$}{(I, \hyperref[I.5.3.9]{5.3.9})}
\EGATermIndexLine{Dihomomorfismo}{(0.1.0.2)}
\EGATermIndexLine{Dominio de definición de una aplicación racional}{(I, \hyperref[I.7.2.1]{7.2.1})}
\EGATermIndexLine{Dual de un $\mathcal O_X$-módulo}{(0.4.1.4)}

\EGATermIndexLine{Elemento topológicamente nilpotente}{(0.7.1.1)}
\EGATermIndexLine{$\mathcal O_X$-módulo generado por una familia de secciones}{(0.5.1.2)}
\EGATermIndexLine{Lugar de ceros de una sección}{(0.5.5.1)}
\EGATermIndexLine{Álgebra íntegra; álgebra íntegra finita}{(0.1.0.5)}
\EGATermIndexLine{Espacio anillado}{(0.4.1.1)}
\EGATermIndexLine{Espacio localmente anillado}{(0.5.5.1)}
\EGATermIndexLine{Espacio anillado inducido sobre un subconjunto abierto}{(0.4.1.2)}
\EGATermIndexLine{Espacio anillado normal, reducido o regular}{(0.4.1.3)}
\EGATermIndexLine{Espacio anillado obtenido por pegado}{(0.4.1.6)}
\EGATermIndexLine{Espacio de Kolmogórov}{(0.2.1.2)}
\EGATermIndexLine{Espacio irreducible}{(0.2.1.1)}
\EGATermIndexLine{Espacio noetheriano}{(0.2.2.1)}
\EGATermIndexLine{Espacio cuasicompacto}{(0.2.2.4)}
\EGATermIndexLine{Espacio subyacente de un espacio anillado}{(0.4.1.1)}
\EGATermIndexLine{Espacio topológicamente anillado}{(0.4.1.1)}

\EGATermIndexLine{Haz algebraico sobre un espacio anillado}{(0.4.1.3)}
\EGATermIndexLine{Haz asociado a un prehaz}{(0.3.5.7)}
\EGATermIndexLine{Haz asociado a un $A$-módulo sobre $\operatorname{Spec}(A)$}{(I, \hyperref[I.1.3.4]{1.3.4})}
\EGATermIndexLine{Haz con valores en una categoría}{(0.3.1.1)}
\EGATermIndexLine{Haz coherente de anillos}{(0.5.3.5)}
\EGATermIndexLine{Haz de anillos normal en un punto, normal, reducido en un punto, reducido, regular en un punto o regular}{(0.4.1.3)}
\EGATermIndexLine{Haz graduado de anillos}{(0.4.1.3)}
\EGATermIndexLine{Haz de funciones racionales}{(I, \hyperref[I.7.3.1]{7.3.1})}
\EGATermIndexLine{Haz de torsión}{(I, \hyperref[I.7.4.1]{7.4.1})}
\EGATermIndexLine{Haz de ideales}{(0.4.1.3)}
\EGATermIndexLine{Haz de ideales de definición}{(I, \hyperref[I.10.3.3]{10.3.3} y \hyperref[I.10.5.1]{10.5.1})}
\EGATermIndexLine{Haz inducido}{(0.3.7.1)}
\EGATermIndexLine{Haz localmente simple}{(0.3.6.1)}
\EGATermIndexLine{Haz obtenido por pegado}{(0.3.3.1)}
\EGATermIndexLine{Haz pseudodiscreto}{(0.3.8.1)}
\EGATermIndexLine{Haz simple}{(0.3.6.1)}
\EGATermIndexLine{Haz estructural de un espacio anillado}{(0.4.1.1)}
\EGATermIndexLine{Haz estructural de un esquema afín}{(I, \hyperref[I.1.3.4]{1.3.4})}
\EGATermIndexLine{Haz sobre una base de subconjuntos abiertos}{(0.3.2.2)}
\EGATermIndexLine{Tallo de un haz}{(0.3.1.6)}
\EGATermIndexLine{Fibra de un morfismo de preesquemas}{(I, \hyperref[I.3.6.3]{3.6.3})}
\EGATermIndexLine{Álgebra finita}{(0.1.0.5)}
\EGATermIndexLine{Función racional}{(I, \hyperref[I.7.1.2]{7.1.2})}

\EGATermIndexLine{Generalización de un punto}{(0.2.1.2)}
\EGATermIndexLine{$\mathcal O_X$-módulo graduado}{(0.4.1.3)}
\EGATermIndexLine{Grafo de un morfismo}{(I, \hyperref[I.5.3.11]{5.3.11})}

\EGATermIndexLine{Homomorfismo continuo de haces de anillos topológicos}{(0.3.1.4)}
\EGATermIndexLine{Homomorfismo definido por una sección}{(0.5.1.1)}
\EGATermIndexLine{Homomorfismo local de anillos locales}{(0.1.0.7)}
\EGATermIndexLine{$\phi$-homomorfismo de módulos}{(0.1.0.2)}

\EGATermIndexLine{Ideal de definición de un anillo admisible}{(0.7.1.2)}
\EGATermIndexLine{$\mathcal O_X$-ideal}{(0.4.1.3)}
\EGATermIndexLine{$\mathcal O_{\mathfrak X}$-ideal de definición de un preesquema formal $\mathfrak X$}{(I, \hyperref[I.10.3.3]{10.3.3} y \hyperref[I.10.5.1]{10.5.1})}
\EGATermIndexLine{Ideal primo}{(0.1.0.6)}
\EGATermIndexLine{Imagen directa de un $\mathcal O_X$-módulo}{(0.4.2.1)}
\EGATermIndexLine{Imagen directa de un prehaz}{(0.3.4.1)}
\EGATermIndexLine{Imagen de una $S$-sección}{(I, \hyperref[I.5.3.11]{5.3.11})}
\EGATermIndexLine{Imagen cerrada de un preesquema por un morfismo}{(I, \hyperref[I.9.5.3]{9.5.3})}
\EGATermIndexLine{Imagen inversa de un $\mathcal O_Y$-módulo}{(0.4.3.1)}
\EGATermIndexLine{Imagen inversa de un $S$-morfismo}{(I, \hyperref[I.3.3.7]{3.3.7})}
\EGATermIndexLine{Imagen inversa de un prehaz}{(0.3.5.3)}
\EGATermIndexLine{Imagen inversa de un $S$-preesquema}{(I, \hyperref[I.3.3.6]{3.3.6})}
\EGATermIndexLine{Imagen inversa de un subpreesquema}{(I, \hyperref[I.4.4.1]{4.4.1})}
\EGATermIndexLine{Inmersión, inmersión cerrada o inmersión abierta de preesquemas}{(I, \hyperref[I.4.2.1]{4.2.1})}
\EGATermIndexLine{Inmersión cerrada de preesquemas formales}{(I, \hyperref[I.10.14.2]{10.14.2})}
\EGATermIndexLine{Inmersión local de preesquemas}{(I, \hyperref[I.4.5.1]{4.5.1})}
\EGATermIndexLine{Inyección canónica de un espacio anillado inducido sobre un subconjunto abierto}{(0.4.1.2)}
\EGATermIndexLine{Inyección canónica de un subpreesquema}{(I, \hyperref[I.4.1.7]{4.1.7})}
\EGATermIndexLine{Inyección canónica de un subpreesquema formal}{(I, \hyperref[I.10.14.2]{10.14.2})}
\EGATermIndexLine{Inyección geométrica}{(I, \hyperref[I.3.5.4]{3.5.4})}
\EGATermIndexLine{Inverso de un $\mathcal O_X$-módulo invertible}{(0.5.4.3)}
\EGATermIndexLine{$\mathcal O_X$-módulo invertible}{(0.5.4.1)}
\EGATermIndexLine{Isomorfismo asociado a una inmersión}{(I, \hyperref[I.4.2.1]{4.2.1})}
\EGATermIndexLine{Isomorfismo local de preesquemas}{(I, \hyperref[I.4.5.1]{4.5.1})}

\EGATermIndexLine{Límite proyectivo de $\mathcal O_{X_n}$-módulos}{(I, \hyperref[I.10.6.6]{10.6.6})}
\EGATermIndexLine{$\mathcal O_X$-módulo localmente libre}{(0.5.4.1)}
\EGATermIndexLine{Carácter local de un punto con valores en un anillo local}{(I, \hyperref[I.3.4.5]{3.4.5})}

\EGATermIndexLine{Módulo de fracciones}{(0.1.2.2)}
\EGATermIndexLine{Módulo fielmente plano}{(0.6.4.1)}
\EGATermIndexLine{Módulo plano}{(0.6.1.1)}
\EGATermIndexLine{Módulo $A$-plano}{(0.6.2)}
\EGATermIndexLine{Módulo cuasifinito}{(0.7.4.1)}
\EGATermIndexLine{$\mathcal O_X$-módulo}{(0.4.1.3)}
\EGATermIndexLine{Morfismo de espacios anillados o topológicamente anillados}{(0.4.1.1)}
\EGATermIndexLine{Morfismo fielmente plano}{(0.6.7.8)}
\EGATermIndexLine{Morfismo plano}{(0.6.7.1)}
\EGATermIndexLine{Morfismo de prehaces definidos sobre una base de subconjuntos abiertos}{(0.3.2.3)}
\EGATermIndexLine{Morfismo de preesquemas}{(I, \hyperref[I.2.2.1]{2.2.1})}
\EGATermIndexLine{Morfismo de $S$-preesquemas o de $A$-preesquemas}{(I, \hyperref[I.2.5.2]{2.5.2})}
\EGATermIndexLine{Morfismo birracional}{(I, 6.5.6)}
\EGATermIndexLine{Morfismo de tipo finito}{(I, \hyperref[I.6.3.1]{6.3.1})}
\EGATermIndexLine{Morfismo diagonal}{(I, \hyperref[I.5.3.1]{5.3.1})}
\EGATermIndexLine{Morfismo dominante}{(I, \hyperref[I.2.2.6]{2.2.6})}
\EGATermIndexLine{Morfismo cerrado}{(I, \hyperref[I.2.2.6]{2.2.6})}
\EGATermIndexLine{Morfismo grafo de un morfismo}{(I, \hyperref[I.3.3.14]{3.3.14})}
\EGATermIndexLine{Morfismo localmente de tipo finito}{(I, \hyperref[I.6.6.2]{6.6.2})}
\EGATermIndexLine{Morfismo dominado por otro}{(I, \hyperref[I.4.1.8]{4.1.8})}
\EGATermIndexLine{Morfismo abierto}{(I, \hyperref[I.2.2.6]{2.2.6})}
\EGATermIndexLine{Morfismo cuasicompacto}{(I, \hyperref[I.6.6.1]{6.6.1})}
\EGATermIndexLine{Morfismo radicial}{(I, \hyperref[I.3.5.4]{3.5.4})}
\EGATermIndexLine{Morfismo reducido}{(I, \hyperref[I.5.1.5]{5.1.5})}
\EGATermIndexLine{Morfismo separado}{(I, \hyperref[I.5.4.1]{5.4.1})}
\EGATermIndexLine{Morfismo estructural de un $S$-preesquema}{(I, \hyperref[I.2.5.1]{2.5.1})}
\EGATermIndexLine{Morfismo sobreyectivo}{(I, \hyperref[I.2.2.6]{2.2.6})}
\EGATermIndexLine{Morfismo universalmente inyectivo}{(I, \hyperref[I.3.5.4]{3.5.4})}
\EGATermIndexLine{Morfismos equivalentes}{(I, \hyperref[I.7.1.1]{7.1.1})}
\EGATermIndexLine{$A$-morfismo o $S$-morfismo de preesquemas}{(I, \hyperref[I.2.5.2]{2.5.2})}
\EGATermIndexLine{Morfismo de preesquemas formales}{(I, \hyperref[I.10.4.5]{10.4.5})}
\EGATermIndexLine{Morfismo ádico de preesquemas formales}{(I, \hyperref[I.10.12.1]{10.12.1})}
\EGATermIndexLine{Morfismo diagonal de preesquemas formales}{(I, \hyperref[I.10.15.1]{10.15.1})}
\EGATermIndexLine{Morfismo de tipo finito de preesquemas formales}{(I, \hyperref[I.10.13.3]{10.13.3})}
\EGATermIndexLine{Morfismo separado de preesquemas formales}{(I, \hyperref[I.10.15.1]{10.15.1})}
\EGATermIndexLine{Morfismo estructural de un $\mathfrak S$-preesquema formal}{(I, \hyperref[I.10.4.7]{10.4.7})}
\EGATermIndexLine{$A$-morfismo o $\mathfrak S$-morfismo de preesquemas formales}{(I, \hyperref[I.10.4.7]{10.4.7})}
\EGATermIndexLine{$\psi$-morfismo de prehaces}{(0.3.5.1)}
\EGATermIndexLine{$\Psi$-morfismo de un $\mathcal O_Y$-módulo a un $\mathcal O_X$-módulo}{(0.4.4.1)}

\EGATermIndexLine{Nilradical de un anillo}{(0.1.1.1)}
\EGATermIndexLine{Nilradical de una $\mathcal O_X$-álgebra}{(I, \hyperref[I.5.1.1]{5.1.1})}
\EGATermIndexLine{Número geométrico de puntos de un preesquema}{(I, \hyperref[I.6.4.8]{6.4.8})}

\EGATermIndexLine{Subconjunto abierto afín}{(I, \hyperref[I.2.1.1]{2.1.1})}
\EGATermIndexLine{Subconjunto abierto formal afín; subconjunto abierto formal afín ádico; subconjunto abierto formal afín noetheriano}{(I, \hyperref[I.10.4.1]{10.4.1})}

\EGATermIndexLine{Subconjunto multiplicativo de un anillo}{(0.1.2.1)}
\EGATermIndexLine{Subconjunto multiplicativo saturado}{(0.1.4.3)}
\EGATermIndexLine{Módulo $f$-plano sobre $\mathcal O_X$}{(0.6.7.1)}
\EGATermIndexLine{Punto de un preesquema con valores en un anillo}{(I, \hyperref[I.3.4.4]{3.4.4})}
\EGATermIndexLine{Punto de un preesquema con valores en un preesquema}{(I, \hyperref[I.3.4.1]{3.4.1})}
\EGATermIndexLine{Punto de un $A$-preesquema con valores en una $A$-álgebra}{(I, \hyperref[I.3.4.4]{3.4.4})}
\EGATermIndexLine{Punto de un $S$-preesquema situado sobre un punto $s\in S$}{(I, \hyperref[I.2.5.1]{2.5.1})}
\EGATermIndexLine{Punto de un $S$-preesquema con valores en un $S$-preesquema}{(I, \hyperref[I.3.4.3]{3.4.3})}
\EGATermIndexLine{Punto cerrado}{(0.2.2.6)}
\EGATermIndexLine{Punto genérico}{(0.2.1.2)}
\EGATermIndexLine{Punto geométrico de un preesquema}{(I, \hyperref[I.3.4.5]{3.4.5})}
\EGATermIndexLine{Punto geométrico situado sobre $s$}{(I, \hyperref[I.3.4.5]{3.4.5})}
\EGATermIndexLine{Punto geométrico localizado en $x$}{(I, \hyperref[I.3.4.5]{3.4.5})}
\EGATermIndexLine{Punto $K$-racional}{(I, \hyperref[I.3.4.5]{3.4.5})}
\EGATermIndexLine{Prehaz constante}{(0.3.6.1)}
\EGATermIndexLine{Prehaz sobre una base de subconjuntos abiertos}{(0.3.2.1)}
\EGATermIndexLine{Preesquema}{(I, \hyperref[I.2.1.2]{2.1.2})}
\EGATermIndexLine{Preesquema artiniano}{(I, \hyperref[I.6.2.1]{6.2.1})}
\EGATermIndexLine{Preesquema conexo}{(I, \hyperref[I.2.1.7]{2.1.7})}
\EGATermIndexLine{Preesquema de base}{(I, \hyperref[I.2.5.1]{2.5.1})}
\EGATermIndexLine{Preesquema obtenido por reducción módulo $\mathfrak J$}{(I, \hyperref[I.3.7.1]{3.7.1})}
\EGATermIndexLine{Preesquema inducido sobre un subconjunto abierto}{(I, \hyperref[I.2.1.8]{2.1.8})}
\EGATermIndexLine{Preesquema íntegro}{(I, \hyperref[I.2.1.7]{2.1.7})}
\EGATermIndexLine{Preesquema irreducible}{(I, \hyperref[I.2.1.7]{2.1.7})}
\EGATermIndexLine{Preesquema localmente íntegro}{(I, \hyperref[I.2.1.7]{2.1.7})}
\EGATermIndexLine{Preesquema localmente noetheriano}{(I, \hyperref[I.6.1.1]{6.1.1})}
\EGATermIndexLine{Preesquema reducido asociado a un preesquema}{(I, \hyperref[I.5.1.3]{5.1.3})}
\EGATermIndexLine{$A$-preesquema; preesquema sobre $A$}{(I, \hyperref[I.2.5.1]{2.5.1})}
\EGATermIndexLine{$K$-preesquema algebraico}{(I, \hyperref[I.6.4.5]{6.4.5})}
\EGATermIndexLine{$S$-preesquema; preesquema sobre $S$}{(I, \hyperref[I.2.5.1]{2.5.1})}
\EGATermIndexLine{$S$-preesquema dominante}{(I, \hyperref[I.2.5.1]{2.5.1})}
\EGATermIndexLine{Preesquema de tipo finito sobre $S$; $S$-preesquema de tipo finito}{(I, \hyperref[I.6.3.1]{6.3.1})}
\EGATermIndexLine{Preesquema obtenido por extensión del preesquema de base}{(I, \hyperref[I.3.3.6]{3.3.6})}
\EGATermIndexLine{Preesquema separado sobre $S$}{(I, \hyperref[I.5.4.1]{5.4.1})}
\EGATermIndexLine{Preesquema formal}{(I, \hyperref[I.10.4.2]{10.4.2})}
\EGATermIndexLine{Preesquema formal ádico}{(I, \hyperref[I.10.4.2]{10.4.2})}
\EGATermIndexLine{Preesquema formal localmente noetheriano}{(I, \hyperref[I.10.4.2]{10.4.2})}
\EGATermIndexLine{Preesquema formal noetheriano}{(I, \hyperref[I.10.4.2]{10.4.2})}
\EGATermIndexLine{Preesquema formal sobre $A$; $A$-preesquema formal}{(I, \hyperref[I.10.4.7]{10.4.7})}
\EGATermIndexLine{Preesquema formal sobre $\mathfrak S$; $\mathfrak S$-preesquema formal}{(I, \hyperref[I.10.4.7]{10.4.7})}
\EGATermIndexLine{Preesquema formal de tipo finito sobre $\mathfrak S$; $\mathfrak S$-preesquema formal de tipo finito}{(I, \hyperref[I.10.13.3]{10.13.3})}
\EGATermIndexLine{Preesquema formal separado sobre $S$}{(I, \hyperref[I.10.15.1]{10.15.1})}
\EGATermIndexLine{Presentación finita (módulo que admite una)}{(0.1.0.5)}
\EGATermIndexLine{Presentación finita ($\mathcal O_X$-módulo que admite una)}{(0.5.2.5)}
\EGATermIndexLine{Principio de inducción noetheriana}{(0.2.2.2)}
\EGATermIndexLine{Producto de $S$-preesquemas}{(I, \hyperref[I.3.2.1]{3.2.1})}
\EGATermIndexLine{Producto de $\mathfrak S$-preesquemas formales}{(I, \hyperref[I.10.7.1]{10.7.1})}
\EGATermIndexLine{Producto fibrado de conjuntos}{(I, \hyperref[I.3.4.2]{3.4.2})}
\EGATermIndexLine{Producto tensorial de haces sobre preesquemas distintos}{(I, \hyperref[I.9.1.2]{9.1.2})}
\EGATermIndexLine{Producto tensorial completado de álgebras}{(0.7.7.5)}
\EGATermIndexLine{Producto tensorial completado de módulos}{(0.7.7.2)}
\EGATermIndexLine{Proyecciones canónicas de un producto}{(I, \hyperref[I.3.2.1]{3.2.1})}
\EGATermIndexLine{Extensión canónica de un submódulo}{(I, \hyperref[I.9.4.1]{9.4.1})}
\EGATermIndexLine{Extensión de un morfismo a los completados}{(I, \hyperref[I.10.9.1]{10.9.1})}
\EGATermIndexLine{$p$-ésima potencia exterior de un $\mathcal O_X$-módulo}{(0.4.1.4)}

\EGATermIndexLine{$\mathcal O_X$-módulo cuasicoherente}{(0.5.1.3)}
\EGATermIndexLine{$\mathcal O_X$-álgebra cuasicoherente}{(0.5.1.3)}

\EGATermIndexLine{Radical de un ideal}{(0.1.1.1)}
\EGATermIndexLine{Radical de un anillo}{(0.1.1.2)}
\EGATermIndexLine{Rango de un $\mathcal O_X$-módulo localmente libre}{(0.5.4.1)}
\EGATermIndexLine{Rango de un $\mathcal O_X$-módulo sin torsión}{(I, \hyperref[I.7.4.2]{7.4.2})}
\EGATermIndexLine{Rango de un $K$-esquema finito}{(I, \hyperref[I.6.4.5]{6.4.5})}
\EGATermIndexLine{Rango separable de un $K$-esquema finito}{(I, \hyperref[I.6.4.8]{6.4.8})}
\EGATermIndexLine{Restricción de un espacio anillado a un subconjunto abierto}{(0.4.1.2)}
\EGATermIndexLine{Restricción de un morfismo a un subconjunto abierto}{(0.4.1.2)}
\EGATermIndexLine{Restricción de un morfismo de preesquemas a un subpreesquema}{(I, \hyperref[I.4.1.7]{4.1.7})}
\EGATermIndexLine{Restricción de un preesquema a un subconjunto abierto}{(I, \hyperref[I.2.1.8]{2.1.8})}
\EGATermIndexLine{Restricción de una aplicación racional a un subconjunto abierto}{(I, \hyperref[I.7.1.2]{7.1.2})}

\EGATermIndexLine{Esquema}{(I, \hyperref[I.5.4.1]{5.4.1})}
\EGATermIndexLine{Esquema afín}{(I, \hyperref[I.1.7.1]{1.7.1})}
\EGATermIndexLine{Esquema local; esquema local en un punto de un preesquema}{(I, \hyperref[I.2.4.1]{2.4.1})}
\EGATermIndexLine{$K$-esquema algebraico}{(I, \hyperref[I.6.4.1]{6.4.1})}
\EGATermIndexLine{Esquema finito sobre $K$; $K$-esquema finito}{(I, \hyperref[I.6.4.5]{6.4.5})}
\EGATermIndexLine{$S$-esquema}{(I, \hyperref[I.5.4.1]{5.4.1})}
\EGATermIndexLine{Esquema formal afín}{(I, \hyperref[I.10.1.2]{10.1.2})}
\EGATermIndexLine{$\mathfrak S$-esquema formal}{(I, \hyperref[I.10.15.1]{10.15.1})}
\EGATermIndexLine{Secciones de un haz}{(0.3.1.6)}
\EGATermIndexLine{$S$-sección de un $S$-preesquema}{(I, \hyperref[I.2.5.5]{2.5.5} y \hyperref[I.5.3.11]{5.3.11})}
\EGATermIndexLine{$S$-sección racional de un $S$-preesquema}{(I, \hyperref[I.7.1.2]{7.1.2})}
\EGATermIndexLine{Sección unidad de $\mathcal O_X$}{(0.4.1.1)}
\EGATermIndexLine{Series formales restringidas}{(0.7.5.1)}
\EGATermIndexLine{Suma de preesquemas}{(I, \hyperref[I.3.1.1]{3.1.1}; el índice de la fuente imprime 3.6.1)}
\EGATermIndexLine{$\mathcal O_X$-subálgebra generada por un submódulo}{(0.4.1.3)}
\EGATermIndexLine{Subpreesquema}{(I, \hyperref[I.4.1.3]{4.1.3})}
\EGATermIndexLine{Subpreesquema asociado a una inmersión}{(I, \hyperref[I.4.2.1]{4.2.1})}
\EGATermIndexLine{Subpreesquema cerrado}{(I, \hyperref[I.4.1.3]{4.1.3})}
\EGATermIndexLine{Subpreesquema cerrado definido por un haz de ideales}{(I, \hyperref[I.4.1.2]{4.1.2})}
\EGATermIndexLine{Subpreesquema formal cerrado}{(I, \hyperref[I.10.14.2]{10.14.2})}
\EGATermIndexLine{Especialización de un punto}{(0.2.1.2)}
\EGATermIndexLine{Espectro de un anillo}{(I, \hyperref[I.1.1.1]{1.1.1})}
\EGATermIndexLine{Espectro formal de un anillo admisible}{(I, \hyperref[I.10.1.2]{10.1.2})}
\EGATermIndexLine{Soporte de un módulo}{(0.1.7.1)}
\EGATermIndexLine{Soporte de un haz de grupos}{(0.3.1.6)}
\EGATermIndexLine{Sistema fundamental de ideales de definición}{(I, \hyperref[I.10.3.6]{10.3.6} y \hyperref[I.10.5.1]{10.5.1})}
\EGATermIndexLine{Sistema inductivo ádico de preesquemas}{(I, \hyperref[I.10.12.2]{10.12.2})}

\EGATermIndexLine{Topología $\mathfrak J$-ádica; topología $\mathfrak J$-preádica}{(0.7.1.9 y 0.7.2.3)}
\EGATermIndexLine{Topología espectral}{(I, \hyperref[I.1.1.2]{1.1.2})}
\EGATermIndexLine{$\mathcal O_X$-módulo invertible trivial}{(I, \hyperref[I.2.4.8]{2.4.8})}
\EGATermIndexLine{$\mathcal O_X$-módulo de tipo finito}{(0.5.2.1)}

\EGATermIndexLine{Valor de una sección en un punto}{(0.5.5.1)}

\medskip\noindent\textit{Nota terminológica.} La fuente emplea «ádico» únicamente
para anillos o módulos separados y completos; este uso difiere algo del habitual.
\end{multicols}
\endgroup
% END source_es/ega1/ega1-backmatter-index-terminology.tex



\clearpage

\oldpage[I]{214}
\begin{thebibliography}{22}

\bibitem{I-1}
H.~\textsc{Cartan} y C.~\textsc{Chevalley},
Séminaire de l'École Normale Supérieure,
8\textsuperscript{o} año (1955-56), \emph{Géométrie algébrique}.

\bibitem{I-2}
H.~\textsc{Cartan} y S.~\textsc{Eilenberg},
\emph{Homological Algebra}, Princeton Math. Series (Princeton University
Press), 1956.

\bibitem{I-3}
W.~L.~\textsc{Chow} y J.~\textsc{Igusa},
Cohomology theory of varieties over rings,
\emph{Proc. Nat. Acad. Sci. U.S.A.}, t.~XLIV (1958), p.~1244-1248.

\bibitem{I-4}
R.~\textsc{Godement}, \emph{Théorie des faisceaux}, Actual. Scient. y Ind.,
n\textsuperscript{o}~1252, Paris (Hermann), 1958.

\bibitem{I-5}
H.~\textsc{Grauert},
Ein Theorem der analytischen Garbentheorie und die Modulräume komplexer
Strukturen, \emph{Publ. Math. Inst. Hautes Études Scient.},
n\textsuperscript{o}~5, 1960.

\bibitem{I-6}
A.~\textsc{Grothendieck}, Sur quelques points d'algèbre homologique,
\emph{Tôhoku Math. Journ.}, t.~IX (1957), p.~119-221.

\bibitem{I-7}
A.~\textsc{Grothendieck},
Cohomology theory of abstract algebraic varieties,
\emph{Proc. Intern. Congress of Math.}, p.~103-118, Edinburgh (1958).

\bibitem{I-8}
A.~\textsc{Grothendieck},
Géométrie formelle y géométrie algébrique,
\emph{Séminaire Bourbaki}, 11\textsuperscript{o} año (1958-59), exposición 182.

\bibitem{I-9}
M.~\textsc{Nagata},
A general theory of algebraic geometry over Dedekind domains,
\emph{Amer. Math. Journ.}~: I, t.~LXXVIII, p.~78-116 (1956)~; II, t.~LXXX,
p.~382-420 (1958).

\bibitem{I-10}
D.~G.~\textsc{Northcott}, \emph{Ideal theory}, Cambridge Univ. Press, 1953.

\bibitem{I-11}
P.~\textsc{Samuel}, \emph{Commutative algebra} (Notas de D.~Herzig), Cornell
Univ., 1953.

\bibitem{I-12}
P.~\textsc{Samuel}, \emph{Algèbre locale}, Mém. Sci. Math.,
n\textsuperscript{o}~123, Paris, 1953.

\bibitem{I-13}
P.~\textsc{Samuel} y O.~\textsc{Zariski},
\emph{Commutative algebra}, 2 vol., New York (Van Nostrand), 1958-60.

\bibitem{I-14}
J.-P.~\textsc{Serre}, Faisceaux algébriques cohérents,
\emph{Ann. of Math.}, t.~LXI (1955), p.~197-278.

\bibitem{I-15}
J.-P.~\textsc{Serre},
Sur la cohomologie des variétés algébriques,
\emph{Journ. de Math.} (9), t.~XXXVI (1957), p.~1-16.

\bibitem{I-16}
J.-P.~\textsc{Serre},
Géométrie algébrique y géométrie analytique,
\emph{Ann. Inst. Fourier}, t.~VI (1955-56), p.~1-42.

\bibitem{I-17}
J.-P.~\textsc{Serre},
Sur la dimension homologique des anneaux y des modules noethériens,
\emph{Proc. Intern. Symp. on Alg. Number theory}, p.~176-189, Tokyo-Nikko,
1955.

\bibitem{I-18}
A.~\textsc{Weil}, \emph{Foundations of algebraic geometry}, Amer. Math. Soc.
Coll. Publ., n\textsuperscript{o}~29, 1946.

\bibitem{I-19}
A.~\textsc{Weil}, Numbers of solutions of equations in finite fields,
\emph{Bull. Amer. Math. Soc.}, t.~LV (1949), p.~497-508.

\bibitem{I-20}
O.~\textsc{Zariski},
\emph{Theory y applications of holomorphic functions on algebraic varieties
over arbitrary ground fields}, Mem. Amer. Math. Soc., n\textsuperscript{o}~5
(1951).

\bibitem{I-21}
O.~\textsc{Zariski}, A new proof of Hilbert's Nullstellensatz,
\emph{Bull. Amer. Math. Soc.}, t.~LIII (1947), p.~362-368.

\bibitem{I-22}
E.~\textsc{Kähler}, Geometria Arithmetica,
\emph{Ann. di Mat.} (4), t.~XLV (1958), p.~1-368.

\end{thebibliography}

\oldpage[II]{217}
\phantomsection
\addcontentsline{toc}{section}{Erratas y adiciones (Lista 1)}
\section*{Erratas y adiciones}
\begin{center}
\large (Lista 1)
\end{center}

\subsection*{Capítulo 0}

\paragraph{(0, 2.1.3)}
En la línea 15 contando desde abajo de la p.~21, reemplazar
«inter-sección» por «intersección».

\paragraph{(0, 2.1.6)}
En la línea 13 de la p.~22, reemplazar
$Z_i\cap\complement(\bigcup\limits_{j\ne i}Z_j)$ par
$\bigcup\limits_{j\ne i}Z_j$.

\paragraph{(0, 3.1.6)}
En la línea 17 contando desde abajo de la p.~25, reemplazar
$\Gamma(U,\sh{F})$ por $\Gamma(V,\sh{F})$.

\paragraph{(0, 4.1.1)}
En la línea 12 de la p.~36, reemplazar $\sh{G}\to\sh{F}$ por
$\sh{B}\to\sh{A}$.

\paragraph{(0, 4.2.1)}
En la línea 11 contando desde abajo de la p.~39, reemplazar «haces» por
«haz».

\paragraph{(0, 5.1.2)}
Después de la línea 14 de la p.~45, añadir:

Si $\sh{G}$ es un $\sh{O}_Y$-módulo generado por sus secciones sobre
$Y$, entonces $f^*(\sh{G})$ está generado por sus secciones sobre $X$,
pues $f^*$ es un funtor exacto por la derecha.

\paragraph{(0, 5.2.4)}
En la línea 18 contando desde abajo de la p.~46, reemplazar $f_U$ por
$f_U^*$.

\paragraph{(0, 5.3.4)}
Antes de la línea 16 contando desde abajo de la p.~47, añadir:

Si
$\sh{A}\to\sh{B}\to\sh{C}\to\sh{D}\to\sh{E}$
es una sucesión exacta de $\sh{O}_X$-módulos, y si
$\sh{A},\sh{B},\sh{D},\sh{E}$ son coherentes, entonces $\sh{C}$ es
coherente.

\paragraph{(0, 5.4.1)}
Después de la línea 2 de la p.~49, añadir:

Supongamos que $\sh{O}_X$ es \emph{coherente}, y sea $\sh{F}$ un
$\sh{O}_X$-módulo \emph{coherente}. Si, en un punto $x\in X$, $\sh{F}_x$ es
un $\sh{O}_x$-módulo \emph{libre de rango $n$}, existe una vecindad $U$ de $x$
tal que $\sh{F}|U$ es \emph{localmente libre de rango $n$}; en efecto,
$\sh{F}_x$ es entonces isomorfo a $\sh{O}_x^n$, y la afirmación se
deduce de (5.2.7).

\paragraph{(0, 6.6.2)}
Reemplazar las líneas 10 y 11 de la p.~59 por:

En efecto, esto se deduce de (6.4.1, \emph{d)}).

\paragraph{(0, 6.7.1)}
Antes de la línea 10 contando desde abajo de la p.~59, añadir:

Si $f$ es un morfismo plano, también se dice que $X$ es
\emph{plano sobre $Y$}, o \emph{$Y$-plano}.

\paragraph{(0, 7.2.9)}
En la línea 11 contando desde abajo de la p.~65, reemplazar $M_n$ por
$M_{n-1}$.

\subsection*{Capítulo I}

\paragraph{(I, 1.2.1)}
En la línea 18 de la p.~83, reemplazar $\boldsymbol{K}$ por $\boldsymbol{k}$ (dos veces).

\paragraph{(I, 1.2.7)}
En la línea 16 contando desde abajo de la p.~84, reemplazar $A$ por
$A'$ (dos veces) y $\mathfrak N$ por $\mathfrak N'$; en la línea 17 contando desde
abajo, reemplazar $X$ por $X'$.


\paragraph{(I, 1.7.3) y (I, 2.2.4)}
Los enunciados dados bajo estos números pueden generalizarse, de
acuerdo con una observación de J.~Tate, como sigue.

\subsection*{1.8. Morfismos de espacios localmente anillados a esquemas afines}
\label{ega2:addendum:subsection:I.1.8-fr}

\begin{proposition}[1.8.1]
\label{ega2:addendum:I.1.8.1-fr}
Sean $(S,\sh{O}_S)$ un esquema afín y $(X,\sh{O}_X)$ un espacio
localmente anillado. Existe una biyección canónica del conjunto de
homomorfismos de anillos
\oldpage[II]{218}
$\Gamma(S,\sh{O}_S)\to\Gamma(X,\sh{O}_X)$
al conjunto de morfismos de espacios anillados
$(\psi,\theta):(X,\sh{O}_X)\to(S,\sh{O}_S)$ tales que, para todo
$x\in X$, $\theta_x^\sharp$ sea un homomorfismo local: $\sh{O}_{\psi(x)}\to\sh{O}_x$.
\end{proposition}

\begin{proof}
Observemos primero que, si $(X,\sh{O}_X)$ y $(S,\sh{O}_S)$ son dos
espacios anillados cualesquiera, un morfismo $(\psi,\theta)$ de
$(X,\sh{O}_X)$ a $(S,\sh{O}_S)$ define canónicamente un homomorfismo
de anillos
$\Gamma(\theta):\Gamma(S,\sh{O}_S)\to\Gamma(X,\sh{O}_X)$;
y, por tanto, una primera aplicación
\begin{equation}
\label{ega2:addendum:I.1.8.1.1-fr}
  \rho:\Hom((X,\sh{O}_X),(S,\sh{O}_S))
  \to\Hom(\Gamma(S,\sh{O}_S),\Gamma(X,\sh{O}_X)).
\tag{1.8.1.1}
\end{equation}
Recíprocamente, bajo las hipótesis de la proposición, pongamos
$A=\Gamma(S,\sh{O}_S)$ y consideremos un homomorfismo de anillos
$\vphi:A\to\Gamma(X,\sh{O}_X)$. Para todo $x\in X$, el conjunto de
los $f\in A$ tales que $\vphi(f)(x)=0$ es un \emph{ideal primo} de $A$, pues
$\sh{O}_x/\mathfrak m_x=\boldsymbol{k}(x)$ es un cuerpo. Es, por tanto, un
elemento de $S=\Spec(A)$, que denotaremos asimismo por
${}^a\vphi(x)$. Además, para todo $f\in A$, por definición (\textbf 0, 5.5.2),
${}^a\vphi^{-1}(D(f))=X_f$,
lo que prueba que ${}^a\vphi$ es una aplicación \emph{continua} $X\to S$.

Definamos después un homomorfismo
\[
  \widetilde{\vphi}:\sh{O}_S\to{}^a\vphi_*(\sh{O}_X)
\] de $\sh{O}_S$-módulos; para todo $f\in A$, se tiene
$\Gamma(D(f),\sh{O}_S)=A_f$ (1.3.6); para todo $s\in A$, se hace corresponder a $s/f\in A_f$ el elemento
$(\vphi(s)|X_f)(\vphi(f)|X_f)^{-1}$ de $\Gamma(X_f,\sh{O}_X)=\Gamma(D(f),{}^a\vphi_*(\sh{O}_X))$,
y se comprueba inmediatamente (pasando de $D(f)$ a $D(fg)$) que esto define un homomorfismo de $\sh{O}_S$-módulos, de donde, un morfismo de espacios
anillados $({}^a\vphi,\widetilde{\vphi})$. Con la misma notación, y
escribiendo $y={}^a\vphi(x)$, se tiene (\textbf 0, 3.7.1)
$\widetilde{\vphi}_x^\sharp(s_y/f_y)=(\vphi(s)_x)(\vphi(f)_x)^{-1}$;
Como $s_y\in\mathfrak m_y$ equivale, por definición, a
$\vphi(s)_x\in\mathfrak m_x$, se ve que $\widetilde{\vphi}_x^\sharp$ es un homomorfismo \emph{local} $\sh{O}_y\to\sh{O}_x$. Se ha
definido así una segunda aplicación
$\sigma:\Hom(\Gamma(S,\sh{O}_S),\Gamma(X,\sh{O}_X))\to\mathfrak L$,
donde $\mathfrak L$ denota el conjunto de morfismos
$(\psi,\theta):(X,\sh{O}_X)\to(S,\sh{O}_S)$ para los cuales
$\theta_x^\sharp$ es local para todo $x\in X$.

Queda probar que $\sigma$ y la restricción de $\rho$ a $\mathfrak L$
son inversas. La definición de $\widetilde{\vphi}$ da inmediatamente
$\Gamma(\widetilde{\vphi})=\vphi$, y por tanto $\rho\circ\sigma$ es la
identidad. Para ver que $\sigma\circ\rho$ es la identidad, partamos de
$(\psi,\theta)\in\mathfrak L$ y sea $\vphi=\Gamma(\theta)$. La
hipótesis sobre $\theta_x^\sharp$ proporciona, al pasar a los cuerpos
residuales, un monomorfismo
$\theta^x:\boldsymbol{k}(\psi(x))\to\boldsymbol{k}(x)$
tal que, para todo $f\in A=\Gamma(S,\sh{O}_S)$,
$\theta^x(f(\psi(x)))=\vphi(f)(x)$;
Así, $f(\psi(x))=0$ equivale a $\vphi(f)(x)=0$, lo que prueba que
${}^a\vphi=\psi$. Por otra parte, las definiciones hacen conmutativo
el diagrama
\[
  \xymatrix{
    A\ar[r]^\vphi\ar[d] &
    \Gamma(X,\sh{O}_X)\ar[d]\\
    A_{\psi(x)}\ar[r]_{\theta_x^\sharp} &
    \sh{O}_x
  }
\]
y también es conmutativo el diagrama análogo en el que
$\theta_x^\sharp$ se reemplaza por $\widetilde{\vphi}_x^\sharp$. De
(\textbf 0, 1.2.4) se sigue que
$\widetilde{\vphi}_x^\sharp=\theta_x^\sharp$, y, por consiguiente,
$\widetilde{\vphi}=\theta$.
\end{proof}

\begin{env}[1.8.2]
\label{ega2:addendum:I.1.8.2-fr}
Cuando $(X,\sh{O}_X)$ e $(Y,\sh{O}_Y)$ son espacios \emph{localmente}
anillados, consideraremos los morfismos
$(\psi,\theta):(X,\sh{O}_X)\to(Y,\sh{O}_Y)$ tales que, para todo
$x\in X$, $\theta_x^\sharp$ sea un homomorfismo \emph{local}: $\sh{O}_{\psi(x)}\to\sh{O}_x$.

Siempre que en lo sucesivo hablemos de un
\oldpage[II]{219}
\emph{morfismo de espacios localmente anillados}, entenderemos siempre
un morfismo de este tipo. Con esta definición de los morfismos, los
espacios localmente anillados forman una \emph{categoría}. Para dos objetos
$X,Y$ de esta categoría, $\Hom(X,Y)$ denotará, por tanto, el conjunto
de morfismos de espacios localmente anillados de $X$ a $Y$ (el
conjunto denotado por $\mathfrak L$ en \sref{ega2:addendum:I.1.8.1-fr}). Cuando queramos
considerar el conjunto de \emph{morfismos de espacios anillados} de $X$ a
$Y$, escribiremos $\Hom_{\mathit{an}}(X,Y)$ para evitar toda
ambigüedad. La aplicación \sref{ega2:addendum:I.1.8.1.1-fr} se escribe entonces
\begin{equation}
\label{ega2:addendum:I.1.8.2.1-fr}
  \rho:\Hom_{\mathit{an}}(X,Y)
  \to\Hom(\Gamma(Y,\sh{O}_Y),\Gamma(X,\sh{O}_X))
\tag{1.8.2.1}
\end{equation}
y su restricción
\begin{equation}
\label{ega2:addendum:I.1.8.2.2-fr}
  \rho':\Hom(X,Y)
  \to\Hom(\Gamma(Y,\sh{O}_Y),\Gamma(X,\sh{O}_X))
\tag{1.8.2.2}
\end{equation}
es \emph{funtorial} en $X$ e $Y$ en la categoría de los espacios localmente
anillados.
\end{env}

\begin{corollary}[1.8.3]
\label{ega2:addendum:I.1.8.3-fr}
Sea $(Y,\sh{O}_Y)$ un espacio localmente anillado. Para que $Y$ sea un
esquema afín, es necesario y suficiente que, para todo espacio
localmente anillado $(X,\sh{O}_X)$, la aplicación \sref{ega2:addendum:I.1.8.2.2-fr} sea
biyectiva.
\end{corollary}

\begin{proof}
La Proposición~\sref{ega2:addendum:I.1.8.1-fr} muestra que la condición es necesaria.
Recíprocamente, se supone que se cumple y se pone
$A=\Gamma(Y,\sh{O}_Y)$. La hipótesis y \sref{ega2:addendum:I.1.8.1-fr} muestran que los
funtores
$X\to\Hom(X,Y)$ y $X\to\Hom(X,\Spec(A))$
de los espacios localmente anillados a los conjuntos son \emph{isomorfos}.
Esto implica la existencia de un isomorfismo canónico
$Y\to\Spec(A)$ (cf.~\textbf 0, 8).
\end{proof}

\begin{env}[1.8.4]
\label{ega2:addendum:I.1.8.4-fr}
Sea $S=\Spec(A)$ un esquema afín. Denotemos por $(S',A')$ el espacio
anillado cuyo espacio subyacente \emph{consta de un punto} y cuyo haz
estructural $A'$ es el haz (necesariamente simple) sobre $S'$ definido
por el anillo $A$. Sea $\pi:S\to S'$ la única aplicación de $S$ en $S'$; observemos, por otra parte, que para todo
abierto $U$ de $S$ se tiene una aplicación canónica
$\Gamma(S',A')=\Gamma(S,\sh{O}_S)\to\Gamma(U,\sh{O}_S)$
que define un \emph{$\pi$-morfismo} $\iota:A'\to\sh{O}_S$ de haces de anillos.
Se obtiene así un \emph{morfismo} canónico \emph{de espacios anillados}
$i=(\pi,\iota):(S,\sh{O}_S)\to(S',A')$.
Para todo $A$-módulo $M$, denotemos por $M'$ el haz simple sobre $S'$
definido por $M$, que es evidentemente un $A'$-módulo. Está claro que
$i_*(\widetilde M)=M'$ (1.3.7).
\end{env}

\begin{lemma}[1.8.5]
\label{ega2:addendum:I.1.8.5-fr}
Con la notación de \sref{ega2:addendum:I.1.8.4-fr}, para todo $A$-módulo $M$, el
$\sh{O}_S$-homomorfismo canónico y funtorial \textup{(\textbf{0}, 4.4.3.3)}
\begin{equation}
\label{ega2:addendum:I.1.8.5.1-fr}
  i^*(i_*(\widetilde M))\to\widetilde M
\tag{1.8.5.1}
\end{equation}
es un isomorfismo.
\end{lemma}

\begin{proof}
Ambos miembros de \sref{ega2:addendum:I.1.8.5.1-fr} son exactos por la derecha
(siendo evidentemente exacto el funtor
$M\to i_*(\widetilde M)$) y conmutan con las sumas directas.
Considerando $M$ como el conúcleo de un homomorfismo
$A^{(I)}\to A^{(J)}$, se reduce uno al caso $M=A$, en el cual la
afirmación es inmediata.
\end{proof}

\begin{corollary}[1.8.6]
\label{ega2:addendum:I.1.8.6-fr}
Sea $(X,\sh{O}_X)$ un espacio anillado y sea $u:X\to S$ un morfismo de
\oldpage[II]{220}
espacios anillados. Para todo $A$-módulo $M$, existe, con la notación
de \sref{ega2:addendum:I.1.8.4-fr}, un isomorfismo canónico y funtorial de
$\sh{O}_X$-módulos
\begin{equation}
\label{ega2:addendum:I.1.8.6.1-fr}
  u^*(\widetilde M)\xrightarrow{\sim} u^*(i^*(M'))
\tag{1.8.6.1}
\end{equation}
\end{corollary}


\begin{corollary}[1.8.7]
\label{ega2:addendum:I.1.8.7-fr}
Bajo las hipótesis de \sref{ega2:addendum:I.1.8.6-fr}, para todo $A$-módulo $M$ y todo
$\sh{O}_X$-módulo $\sh{F}$ existe un isomorfismo canónico, funtorial
en $M$ y $\sh{F}$,
\begin{equation}
\label{ega2:addendum:I.1.8.7.1-fr}
  \Hom_{\sh{O}_S}(\widetilde M,u_*(\sh{F}))
  \xrightarrow{\sim}\Hom_A(M,\Gamma(X,\sh{F})).
\tag{1.8.7.1}
\end{equation}
\end{corollary}

\begin{proof}
Por (\textbf 0, 4.4.3) y el Lema~\sref{ega2:addendum:I.1.8.5-fr}, existe un isomorfismo
canónico de bifuntores
\[
  \Hom_{\sh{O}_S}(\widetilde M,u_*(\sh{F}))
  \xrightarrow{\sim}\Hom_{A'}(M',i_*(u_*(\sh{F})))
\]
El segundo miembro es claramente $\Hom_A(M,\Gamma(X,\sh{F}))$. Bajo
el homomorfismo canónico \sref{ega2:addendum:I.1.8.7.1-fr}, a un
$\sh{O}_S$-homomorfismo
$h:\widetilde M\to u_*(\sh{F})$ (equivalentemente, a un $u$-morfismo
$\widetilde M\to\sh{F}$) le corresponde el $A$-homomorfismo $\Gamma(h):M\to\Gamma(S,u_*(\sh{F}))=\Gamma(X,\sh{F})$.
\end{proof}

\begin{env}[1.8.8]
\label{ega2:addendum:I.1.8.8-fr}
Con la notación de \sref{ega2:addendum:I.1.8.4-fr}, se deduce claramente de
(\textbf 0, 4.1.1) que todo morfismo de espacios anillados
$(\psi,\theta):X\to S'$ equivale al dato de un homomorfismo de
anillos $A\to\Gamma(X,\sh{O}_X)$. Por tanto, la
Proposición~\sref{ega2:addendum:I.1.8.1-fr} puede interpretarse como una biyección
canónica
$\Hom(X,S)\xrightarrow{\sim}\Hom(X,S')$
donde en el segundo miembro se entienden morfismos de espacios
anillados, pues $A$ no es en general un anillo local.

Más generalmente, sean $X,Y$ espacios localmente anillados, y sea
$(Y',A')$ el espacio anillado cuyo espacio subyacente consta de un
punto y cuyo haz de anillos $A'$ es el haz simple definido por
$\Gamma(Y,\sh{O}_Y)$. Entonces \sref{ega2:addendum:I.1.8.2.1-fr} puede interpretarse
como una aplicación
\begin{equation}
\label{ega2:addendum:I.1.8.8.1-fr}
  \rho:\Hom_{\mathit{an}}(X,Y)\to\Hom(X,Y')
\tag{1.8.8.1}
\end{equation}
El Corolario~\sref{ega2:addendum:I.1.8.3-fr} dice que los esquemas afines se
caracterizan, entre los espacios localmente anillados, como aquellos
para los cuales la restricción de $\rho$ a $\Hom(X,Y)$,
\begin{equation}
\label{ega2:addendum:I.1.8.8.2-fr}
  \rho':\Hom(X,Y)\to\Hom(X,Y')
\tag{1.8.8.2}
\end{equation}
es \emph{biyectiva} para \emph{todo} espacio localmente anillado $X$. En un capítulo
posterior generalizaremos esta definición, asociando así a un espacio
anillado \emph{arbitrario} $Z$ (y no solamente a uno cuyo espacio subyacente
consta de un punto) un espacio localmente anillado que volveremos a
denotar por $\Spec(Z)$. Este será el punto de partida de una teoría
«relativa» de los preesquemas sobre espacios anillados arbitrarios,
que extiende los resultados del Capítulo~I.
\end{env}

\begin{env}[1.8.9]
\label{ega2:addendum:I.1.8.9-fr}
Se pueden considerar los pares $(X,\sh{F})$ formados por un espacio localmente anillado $X$ y un $\sh{O}_X$-módulo $\sh{F}$ como formando una categoría, siendo un \emph{morfismo} $(X,\sh{F})\to(Y,\sh{G})$ de esta categoría un par $(u,h)$ formado por un morfismo de
\oldpage[II]{221}
espacios localmente anillados $u:X\to Y$ y de un $u$-morfismo
$h:\sh{G}\to\sh{F}$ de módulos; estos
morfismos (para $(X,\sh{F})$ e $(Y,\sh{G})$ fijados) forman un conjunto
que denotaremos por $\Hom((X,\sh{F}),(Y,\sh{G}))$; la aplicación
$(u,h)\to(\rho'(u),\Gamma(h))$ es una aplicación canónica
\begin{equation}
\label{ega2:addendum:I.1.8.9.1-fr}
\begin{split}
  \Hom((X,\sh{F}),(Y,\sh{G}))\to
  \Hom\bigl(
    &(\Gamma(Y,\sh{O}_Y),\Gamma(Y,\sh{G})),\\
    &(\Gamma(X,\sh{O}_X),\Gamma(X,\sh{F}))
  \bigr)
\end{split}
\tag{1.8.9.1}
\end{equation}
\emph{funtorial} en $(X,\sh{F})$ y $(Y,\sh{G})$, siendo el segundo miembro
el conjunto de los dihomomorfismos correspondientes a los anillos y a
los módulos considerados (\textbf 0, 1.0.2).
\end{env}

\begin{corollary}[1.8.10]
\label{ega2:addendum:I.1.8.10-fr}
Sean $Y$ un espacio localmente anillado y $\sh{G}$ un
$\sh{O}_Y$-módulo. Para que $Y$ sea un esquema afín y $\sh{G}$ un
$\sh{O}_Y$-módulo cuasicoherente, es necesario y suficiente que, para
toda pareja $(X,\sh{F})$ formada por un espacio localmente anillado $X$
y un $\sh{O}_X$-módulo $\sh{F}$, la aplicación canónica
\sref{ega2:addendum:I.1.8.9.1-fr} sea biyectiva.
\end{corollary}

Dejamos al lector la tarea de desarrollar el razonamiento, modelado
sobre el de \sref{ega2:addendum:I.1.8.3-fr}, en el cual se utilizan \sref{ega2:addendum:I.1.8.1-fr} y
\sref{ega2:addendum:I.1.8.7-fr}.

\begin{remark}[1.8.11]
\label{ega2:addendum:I.1.8.11-fr}
Los enunciados (1.7.3), (1.7.4) y (2.2.4), así como la
definición (1.6.1), son casos particulares de \sref{ega2:addendum:I.1.8.1-fr}; del
mismo modo, (2.2.5) se deduce de \sref{ega2:addendum:I.1.8.7-fr}. El
Corolario~\sref{ega2:addendum:I.1.8.7-fr} proporciona también (1.6.3) (y, por tanto,
(1.6.4)) como caso particular. En efecto, si $X$ es afín y
$\Gamma(X,\sh{F})=N$, los funtores
$M\to
  \Hom_{\sh{O}_S}(\widetilde M,u_*(\widetilde N))$ y $M\to
  \Hom_{\sh{O}_S}(\widetilde M,(N_{[\vphi]})\supertilde)$
(donde $\vphi:A\to\Gamma(X,\sh{O}_X)$ corresponde a $u$) son
isomorfos por \sref{ega2:addendum:I.1.8.7-fr} y (1.3.8). Finalmente, (1.6.5) (y,
por tanto, (1.6.6)) se deduce de \sref{ega2:addendum:I.1.8.6-fr} y del hecho de que,
para todo $f\in A$, los $A_f$-módulos
$N'\otimes_{A'}A_f$ y $(N'\otimes_{A'}A)_f$
(con las notaciones de (1.6.5)) son canónicamente isomorfos.
\end{remark}

\paragraph{Insértese después de (I, 3.2.8):}

\begin{env}[3.2.9]
\label{ega2:addendum:I.3.2.9-fr}
De \sref{ega2:addendum:I.1.8.1-fr} se deduce que (3.2.2) puede precisarse del modo
siguiente. El preesquema
$Z=\Spec(B\otimes_A C)$
no solo es un producto de $X=\Spec(B)$ e $Y=\Spec(C)$ en la categoría
de los $S$-\emph{preesquemas}, sino también en la categoría de los \emph{espacios
localmente anillados} \emph{sobre} $S$ (con la definición de los
$S$-morfismos modelada sobre (2.5.2)). La demostración de (3.2.6)
muestra asimismo que, para cualesquiera $S$-preesquemas $X,Y$, el
preesquema $X\times_S Y$ no solo es el producto de $X$ y $Y$ en la categoría de los
$S$-preesquemas, sino también en la categoría de los espacios
localmente anillados sobre el preesquema $S$.
\end{env}

\paragraph{(I, 4.4.5)}
En la línea 14 contando desde abajo de la p.~126, sustitúyase $B$ por
$A$; en la línea 13 contando desde abajo, sustitúyase «álgebra sobre
$A$» por «álgebra sobre $B$».

\paragraph{(I, 5.3.5)}
En la línea 2 contando desde abajo de la p.~132, sustitúyase $Y\to S$
por $g:Y\to S$.

\paragraph{(I, 6.3.2.1)}
En la línea 6 contando desde abajo de la p.~144, léase:
$D(g_i)\subset W$.

\paragraph{(I, 7.3.8)}
Sustitúyase el texto actual, que es erróneo, por el siguiente:

Sean $X,Y$ dos preesquemas \emph{íntegros}; entonces $\sh{R}(X)$ y
$\sh{R}(Y)$ son, respectivamente, un $\sh{O}_X$-módulo
cuasicoherente y un $\sh{O}_Y$-módulo cuasicoherente (7.3.3). Si
$f:X\to Y$ es un morfismo \emph{dominante}, existe un homomorfismo canónico
de $\sh{O}_X$-módulos
\oldpage[II]{222}
\begin{equation}
\label{ega2:addendum:I.7.3.8.1-fr}
  \tau:f^*(\sh{R}(Y))\to\sh{R}(X).
\tag{7.3.8.1}
\end{equation}

Supongamos primero que $X=\Spec(A)$ e $Y=\Spec(B)$ son afines, con
$A$ y $B$ anillos íntegros. Entonces $f$ corresponde a un homomorfismo
inyectivo $B\to A$, que se prolonga en un monomorfismo $L\to K$ del
cuerpo de fracciones $L$ de $B$ en el cuerpo de fracciones $K$ de $A$.
El homomorfismo \sref{ega2:addendum:I.7.3.8.1-fr} corresponde entonces al homomorfismo
canónico
$L\otimes_B A\to K$
de (1.6.5).

En el caso general, para toda pareja de abiertos afines no vacíos
$U\subset X$, $V\subset Y$ tales que $f(U)\subset V$, la construcción
anterior define un homomorfismo $\tau_{U,V}$. Si $U'\subset U$,
$V'\subset V$ y $f(U')\subset V'$, entonces $\tau_{U,V}$ prolonga a
$\tau_{U',V'}$; la afirmación se obtiene por pegado. Si $x,y$ son los
puntos genéricos de $X$ y $Y$, respectivamente, entonces $f(x)=y$ y
\[
  \bigl(f^*(\sh{R}(Y))\bigr)_x
  =\sh{O}_y\otimes_{\sh{O}_y}\sh{O}_x
  =\sh{O}_x
\]
por (\textbf 0, 4.3.1) de modo que $\tau_x$ es un \emph{isomorfismo}.

\paragraph{(I, 9.5.2)}
En las tres últimas líneas de la p.~176 y las cinco primeras de la
p.~177, sustitúyase en todas partes $X'$ por $Y'$ y $X''$ por $Y''$.

\paragraph{(I, 10.5.4)}
En la línea 11 contando desde abajo de la p.~187, sustitúyase (10.3.4)
por (10.3.5).

\paragraph{(I, 10.6.3)}
En la línea 3 contando desde abajo de la p.~189, sustitúyase $X$ por
$\mathfrak X$.

\paragraph{(I, 10.14.2)}
En la línea 5 de la p.~210, sustitúyase «$\sh{O}_{\mathfrak X}$-\emph{módulo
coherente}» por «\emph{ideal coherente de} $\sh{O}_{\mathfrak X}$»; en la
línea 7, sustitúyase «$\sh{O}_{\mathfrak X}$-módulos coherentes» por
«ideales coherentes de $\sh{O}_{\mathfrak X}$».

\end{document}
